\documentclass[11pt]{article}
\usepackage[T1]{fontenc}
\usepackage{lmodern}
\usepackage[margin=1in]{geometry}
\usepackage{amsmath,amssymb,amsthm,mathtools,mathrsfs,bm}
\usepackage{enumitem,booktabs,tabularx,array,microtype}
\usepackage[hidelinks]{hyperref}
\usepackage{xurl}
\usepackage[nameinlink,capitalise]{cleveref}
\setlength{\emergencystretch}{2em}
\raggedbottom
\numberwithin{equation}{section}
\newtheorem{theorem}{Theorem}[section]
\newtheorem{proposition}[theorem]{Proposition}
\newtheorem{lemma}[theorem]{Lemma}
\newtheorem{corollary}[theorem]{Corollary}
\theoremstyle{definition}
\newtheorem{definition}[theorem]{Definition}
\newtheorem{inputhyp}[theorem]{Inherited input}
\newtheorem{example}[theorem]{Example}
\crefname{inputhyp}{inherited input}{inherited inputs}
\Crefname{inputhyp}{Inherited input}{Inherited inputs}
\theoremstyle{remark}
\newtheorem{remark}[theorem]{Remark}
\newcommand{\R}{\mathbb R}
\newcommand{\C}{\mathbb C}
\newcommand{\N}{\mathbb N}
\newcommand{\Z}{\mathbb Z}
\newcommand{\T}{\mathbb T}
\newcommand{\E}{\mathbb E}
\newcommand{\PP}{\mathbb P}
\newcommand{\one}{\mathbf1}
\newcommand{\dd}{\mathrm d}
\newcommand{\ii}{\mathrm i}
\newcommand{\tr}{\operatorname{tr}}
\newcommand{\Var}{\operatorname{Var}}
\newcommand{\Cov}{\operatorname{Cov}}
\newcommand{\Ran}{\operatorname{Ran}}
\newcommand{\Ker}{\operatorname{Ker}}
\newcommand{\Riem}{\operatorname{Riem}}
\newcommand{\norm}[1]{\lVert#1\rVert}
\newcommand{\abs}[1]{\lvert#1\rvert}
\newcommand{\ip}[2]{\langle#1,#2\rangle}
\newcommand{\Sh}{S_h^{\mathrm{sh}}}
\newcommand{\Ucal}{\mathcal U}
\newcolumntype{L}[1]{>{\raggedright\arraybackslash}p{#1}}
\title{\textbf{Same-Source Regenerative Control and\ Classical Einstein--Standard-Model Closure}\\[0.4em]
\large Joint record selection, quantitative reader tails, and an innovation obstruction\\[0.3em]
\normalsize Research notes, Version 1}
\author{}
\date{September 2026}
\begin{document}
\maketitle

% BEGIN IMMUTABLE VERSION 1 BODY
\section*{Context and purpose of this proof lineage}
This file starts a new append-only research lineage.  The problem is to connect
three already developed components on \emph{one literal accepted source family}:
the regenerative positive-energy machinery of the Grand Tensor monograph,
the same-action finite-prefix stationarity estimate, and the full-record
classical Einstein--Standard-Model closure theorem.  The question is not to
construct another Fourier--midpoint regulator or to reprove the general
compactness and coupled PDE results.  It is to determine what additional
same-source estimates actually permit these components to be composed.

The two monographs and the cumulative classical closure file are treated as
specified sources, not as a license to strengthen their quantifiers.
The regenerative results control states and declared source energies;
the prefix result controls a mean squared Euler row along actual accepted
comparisons; the classical closure result tests a \emph{single complete
spacetime record}.  These are different objects.  In particular, a record
selected for its small Euler row need not be a record with a controlled
high-frequency tail.

The principal new conclusions proved here are the following.  First, a joint
occupation-score selection rule combines energy and stationarity estimates
without requiring every candidate to have a high-order regularity bound.
Second, a finite source-to-reader norm identifies exactly which energy controls
the weighted spacetime Fourier tails used by the existing closure theorem;
a Sobolev specialization gives explicit cutoff powers and a sharp threshold
for this particular transfer estimate.  Third, a conditional innovation
inequality gives a necessary noise scale for the prefix-balance route.
Together these results yield a same-source quantitative closure criterion with
an explicit failure probability.  They do not verify the criterion for an
unspecified primitive law.

\paragraph{Status of proofs.}
Every result called a theorem, proposition, lemma, or corollary in this file is
proved below, except the two statements explicitly labelled \emph{Inherited
input}.  The concluding Einstein--SM statement is conditional on those
inherited inputs and on the explicitly displayed finite-source certificates.
No independent verification of the entire coupled PDE argument in the supplied
files is asserted.  The numerical diagnostics described at the end are checks,
not substitutes for proofs or a machine-checked formalization.

\paragraph{Continuation convention.}
Subsequent iterations are to preserve this mathematical body, append new work
at the marked insertion point before \verb|\end{document}|, and increment the
filename suffix from \verb|_V1.tex| to \verb|_V2.tex| and so on.  Corrections to
an inherited or earlier assertion must be stated explicitly in a later
section rather than silently rewriting this body.

\tableofcontents

\section{Source audit and inherited mathematical inputs}
\label{v1:sec:audit}

\subsection{What is already established in the supplied sources}
The relevant source results are identified by their literal LaTeX labels.
This avoids dependence on page numbers that may change after compilation.
\begin{center}\small
\begin{tabularx}{\textwidth}{@{}L{0.20\textwidth}>{\raggedright\arraybackslash}X@{}}
\toprule
Source & Results used, not reproved here\\
\midrule
Grand Tensor V076 \cite{GT} &
\texttt{thm:renewal-Lyapunov-tightness};
\texttt{thm:renewal-shorted-Hodge};
\texttt{def:renewal-regenerative-packet};
\texttt{thm:renewal-regenerative-continuum};
\texttt{thm:renewal-finite-type-closure}.\\[0.3em]
Classical closure V27 \cite{CC} &
\texttt{v14:thm:master}, the populated comparison branch;
\texttt{v20:thm:balance}, the same-action transient estimate;
\texttt{v26:prop:potential}, the distinction between a Fisher metric and a
Bregman Hessian;
\texttt{v27:thm:source} and \texttt{v27:thm:closure}, the full-record transfer.\\[0.3em]
SM--Einstein monograph \cite{SM} &
\texttt{ass:closure} and \texttt{thm:main}, the separate conditional quantum
continuum chain; \texttt{thm:coexact}, the small-chart gauge reconstruction.
These are not replacements for the classical same-source estimates below.\\
\bottomrule
\end{tabularx}
\end{center}

The Grand Tensor regenerative theorem explicitly leaves domain-specific
variational residuals and semantic identification separate.  Its positive
spatial screen has dimension-three spectral counting; it is not a bound for
arbitrarily high \emph{spacetime} derivatives of the physical record.
The classical closure file already proves finite initial-data generation on a
comparison branch, and finite-prefix stationarity from an observed balance.
Neither needs to be reconstructed in this lineage.

The comparison branch, the explicitly designed operational realization, and
an independently specified native source are also distinct.  In particular,
\texttt{v17:thm:master} in \cite{CC} proves an existence statement for a specified
computationally marked law.  Its accompanying quantifier audit does not identify
that law with every independently given Renewal source.  The present notes
retain this distinction.

\subsection{Normalization and the two inputs}
All source coordinates are realified when needed.  At fixed spacing $h$, use
the physical mass Hilbert norm on the full free nodal variation space:
\begin{equation}
 \norm v_{0,h}^2=h^4\sum_n\abs{v(n)}^2.
 \label{v1:eq:mass}
\end{equation}
An action gradient is its Riesz gradient in this norm.  The full variation bank
includes the independent metric and temporal-gauge rows; it is not a projected
Gauss or harmonic tangent bank.  A single calibrated action is kept fixed
throughout an accepted comparison prefix.

\begin{inputhyp}[Same-action prefix estimate from the supplied closure file]
\label[inputhyp]{v1:in:prefix}
Let $\mathcal A_h$ be the actual action on a convex coordinate chart and let
$p_h=\nabla_h\Psi_h$, where
\begin{equation}
 g_*I\preceq Dp_h\preceq g^*I,\qquad
 \norm{D E_h}_{0,h\to0,h}\le K_0h^{-2},\qquad
 E_h=\nabla_h\mathcal A_h.
 \label{v1:eq:prefix-geometry}
\end{equation}
Here $g_*,g^*,K_0$ are independent of $h$.  Put
$\upsilon_h=\upsilon_0h^2$, with $\upsilon_0K_0<g_*/2$.  For an actual accepted
sequence $\theta_0,\ldots,\theta_J$ whose segments lie in the chart, define
\begin{equation}
 b_j=E_h(\theta_{j+1})+
 \upsilon_h^{-1}\bigl(p_h(\theta_{j+1})-p_h(\theta_j)\bigr).
 \label{v1:eq:balance}
\end{equation}
The source result gives constants $C_P,C_B$ independent of $h,J$ such that
\begin{equation}
 \frac1J\sum_{j<J}\norm{E_h(\theta_j)}_{0,h}^2
 \le \frac{C_Ph^{-2}D_h}{J}
       +\frac{C_B}{J}\sum_{j<J}\norm{b_j}_{0,h}^2,
 \label{v1:eq:prefix-input}
\end{equation}
whenever $\mathcal A_h(\theta_0)-\mathcal A_h(\theta_J)\le D_h$.
The allowed span $D_h$ is not assumed cutoff independent here.
\end{inputhyp}

This is exactly the role of \texttt{v20:thm:balance}; the descent calculation is
not reproved.  A changing action would require the explicit work term already
identified in \texttt{v20:cor:work}.  Uniform stationary covariance does not by
itself supply the scalar potential $\Psi_h$ in
\eqref{v1:eq:prefix-geometry}.  The obstruction in
\texttt{v26:prop:potential} is retained.

For the second input, use the auxiliary four-dimensional periodic readout box
of \cite{CC}, containing a buffered physical cylinder
$Q=[0,T]\times\T^3$.  Periodicity in the auxiliary time coordinate is a readout
convention outside that cylinder, not periodic physical evolution.  Let
$\Lambda_h\subset\Z^4$ be the odd-grid Fourier set with Nyquist scale
$N_h\asymp h^{-1}$.  The primitive label interpolant $u$ contains every recorded
mode, and the physical field is $z=\mathscr D(u)$ for the fixed decoder.
For $1\le K<N_h$, set
\begin{equation}
 \tau_{h,j}(K;u)=
 \sum_{\substack{\ell\in\Lambda_h\\\abs\ell_\infty>K}}
 \left(1+\frac{\abs\ell_1}{K}\right)^j\abs{\widehat u(\ell)}.
 \label{v1:eq:tau}
\end{equation}
For $k\ge4$ put $m=k+2$.  Denote the actual initial and harmonic errors by
\begin{align}
 i_{h,k}&=\norm{\Ucal_h(0)-\Ucal_*(0)}_{H_x^k},\notag\\
 \gamma_{h,k}&=\norm{C(g_h)}_{L_t^2H_x^{k+1}}+
                  \norm{\partial_t C(g_h)}_{L_t^2H_x^k}.
 \label{v1:eq:alignment}
\end{align}
The reference constrained initial datum is fixed; only its local physical
solution is supplied by the inherited existence theorem.  The candidate is
in temporal internal gauge on the evaluated cylinder.  Norms and coefficient
margins are the positive comparison norms and fixed calibrated margins of the
source theorem.

\begin{inputhyp}[Full-record classical closure from the supplied closure file]
\label[inputhyp]{v1:in:closure}
For the specified, unchanged raw common action and decoder, suppose
$\norm{E_h^{\rm raw}(y_h)}_{0,h}\le C\sigma_h$, the compact nodal chart margins
hold, $hK\to0$, and $\tau_{h,m}\le\tau_*$.  Define
\begin{align}
 \epsilon_0&=\sigma_h+hK^3,\qquad
 \epsilon_c=\epsilon_0+K^2\tau_{h,2},\notag\\
 L_{h,k}&=1+\log\left(2+\frac{C_kK^{k+5}}{\epsilon_c}\right),\notag\\
 F_{h,k}&=\epsilon_cK^{k+1}L_{h,k}^{k+1}+K^{k+2}\tau_{h,m},\qquad
 D_{h,k}=i_{h,k}+\gamma_{h,k}+F_{h,k}.
 \label{v1:eq:closure-budget}
\end{align}
For sufficiently small $D_{h,k}$ and $\epsilon_c$, the full, unfiltered record
satisfies
\begin{align}
 \mathfrak E_k(z_h,z_*)&:=
 \norm{\Ucal_h-\Ucal_*}_{C_tH_x^k}
 +\norm{\Riem(g_h)-\Riem(g_*)}_{L_t^2H_x^{k-1}}\notag\\
 &\hspace{2em}+\norm{T^{\rm SM}(z_h)-T^{\rm SM}(z_*)}_{L_t^2H_x^k}
 \le C D_{h,k}.
 \label{v1:eq:inherited-closure}
\end{align}
Its zero-order physical equation residual is at most $C\epsilon_0$, and the
literal Cartan reconstruction discrepancy is at most $ChK^3$ in the source's
$L^2$ norm.  Existence of the complete finite record and the common chart and
calibration remain inputs.
\end{inputhyp}

We use this input as stated, including its dependence on the earlier
actual-jet identity and local physical existence theorem.  The new mathematics
below supplies a way to choose one actual record and to bound its
$\sigma_h,\tau_{h,j},i_{h,k},\gamma_{h,k}$ from same-source information.
It is not a new proof of \eqref{v1:eq:inherited-closure}.

\section{A joint selector on the actual accepted record law}
\label{v1:sec:joint}

\subsection{Why separate selection is insufficient}
An energy estimate for a probability law and a small minimum action slope
along its sample path do not imply that the minimizer has controlled energy.
The distinction persists for independent draws and does not require subtle
long-range dependence.

\begin{example}[Rare rough records can dominate minimum-slope selection]
\label[example]{v1:ex:bias}
For $0<h<1/2$, let independent records take values $r,s$ with
$\PP(r)=h^2$, $\PP(s)=1-h^2$.  Assign
\[
 V_h(r)=h^{-2},\quad V_h(s)=0,\qquad
 g_h(r)=0,\quad g_h(s)=h.
\]
Then $\E V_h=1$ and $\E g_h^2\le h^2$.  Among
$J_h=\lceil h^{-4}\rceil$ records, a minimum-slope selector chooses $r$ whenever
$r$ occurs.  That probability tends to one, and its selected energy is
$h^{-2}$.
\end{example}
\begin{proof}
The two expectation statements follow by summing the two atoms.  The
probability that $r$ is absent is
$(1-h^2)^{J_h}\le\exp(-h^2J_h)\to0$.  The strict inequality
$g_h(r)<g_h(s)$ proves the selection assertion.
\end{proof}

This example is a logical selection obstruction, not a model of the
Einstein--SM action.  A separate minimum for each estimate is not enough.
The next result selects all required properties at one common index.

\begin{lemma}[Deterministic joint occupation selection]
\label[lemma]{v1:lem:selector}
Let $g_j,V_j,W_j\ge0$ for $0\le j<J$.  For thresholds $s,R,a>0$ define
\begin{equation}
 S_j=\frac{g_j^2}{s^2}+\frac{V_j}{R^2}+\frac{W_j}{a^2}.
 \label{v1:eq:score}
\end{equation}
If $J^{-1}\sum_j S_j\le1$, a minimizing index $j_*$ satisfies
\begin{equation}
 g_{j_*}\le s,\qquad V_{j_*}\le R^2,\qquad W_{j_*}\le a^2.
 \label{v1:eq:joint-output}
\end{equation}
The index is one of the observed records; no averaging or modification of its
field coordinates is involved.
\end{lemma}
\begin{proof}
$S_{j_*}\le J^{-1}\sum_jS_j\le1$, and all three terms in
\eqref{v1:eq:score} are nonnegative.
\end{proof}

\subsection{Actual finite tables and retained failures}
At spacing $h$, let $X_0,\ldots,X_J$ be records of an actual finite process.
A state is allowed to retain the entire finite prefix, so the formulation
covers a finite non-Markovian source after this literal enlargement.
Let $K_{h,j}(x,y)$ be its actual transition probabilities.  Each table state
$x$ has a declared physical field-coordinate readout $\theta_h(x)$ in the
raw-action chart.  In the formulas below, $E_h(x)$ and $p_h(x)$ abbreviate
$E_h(\theta_h(x))$ and $p_h(\theta_h(x))$; $b_h(x,y)$ uses these same
readouts.  Extra history registers are not added to the physical variation
bank.  Legal segments mean the segments between these field-coordinate
readouts.  Choose a
predeclared legal-edge mask $\ell_{h,j}(x,y)\in\{0,1\}$ for the raw-action
chart, margins, and action-identification requirements.  No regularity-tail
bound is silently included in this mask.  Define
\begin{align}
 K_{h,j}^{\rm leg}(x,y)&=K_{h,j}(x,y)\ell_{h,j}(x,y),\notag\\
 \mu_{j+1}(y)&=\sum_x\mu_j(x)K_{h,j}^{\rm leg}(x,y),\qquad
 p_h^{\rm exit}=1-\sum_y\mu_J(y).
 \label{v1:eq:killed}
\end{align}
The initial subprobability $\mu_0$ includes any illegal initial mass as
failure.  Let $\Omega_j$ be survival through stage $j$ and $\Omega=\Omega_J$.
These are the original probabilities, not probabilities conditioned on
success.

Take
\begin{equation}
 g_h(x)=\norm{E_h(x)}_{0,h},\qquad
 W_h(x)=\bigl(i_{h,k}(x)+\gamma_{h,k}(x)\bigr)^2,
 \label{v1:eq:gW}
\end{equation}
and a nonnegative native energy readout $V_h(x)$.  Define the finite occupation
and transition budgets
\begin{align}
 M_{h,J}&=\frac1J\sum_{j<J}\sum_x\mu_j(x)V_h(x),\notag\\
 A_{h,J}&=\frac1J\sum_{j<J}\sum_x\mu_j(x)W_h(x),\notag\\
 B_{h,J}&=\frac1J\sum_{j<J}\sum_{x,y}
   \mu_j(x)K_{h,j}^{\rm leg}(x,y)\norm{b_h(x,y)}_{0,h}^2.
 \label{v1:eq:occupations}
\end{align}
Here $b_h(x,y)$ is the full same-action discrepancy
\eqref{v1:eq:balance}.  Every sum is on the given source's actual finite table.
In particular the symbols do not prescribe a new sampler.

\begin{theorem}[Same-record probabilistic selection without uniform prefix regularity]
\label[theorem]{v1:thm:joint}
Assume \cref{v1:in:prefix} applies on surviving prefixes with endpoint span at
most $D_h$.  Fix positive thresholds $s_h,R_h,a_h$.  On $\Omega$, choose the
observed record minimizing \eqref{v1:eq:score}, with
$g_j=g_h(X_j)$, $V_j=V_h(X_j)$, $W_j=W_h(X_j)$; retain failure on
$\Omega^c$.  Then the probability that this rule fails to output a record
satisfying all of \eqref{v1:eq:joint-output} is at most
\begin{equation}
 p_h^{\rm exit}
 +\frac{C_Ph^{-2}D_h}{J s_h^2}
 +\frac{C_B B_{h,J}}{s_h^2}
 +\frac{M_{h,J}}{R_h^2}
 +\frac{A_{h,J}}{a_h^2}.
 \label{v1:eq:joint-failure}
\end{equation}
No stationary distribution, mixing assumption, or independence between
records is required.  In particular there is no factor $J$ multiplying the
energy-occupation term.
\end{theorem}
\begin{proof}
Let $\overline S=J^{-1}\sum_{j<J}S_j$ on surviving prefixes.  By
\cref{v1:lem:selector}, failure on $\Omega$ implies
$\overline S>1$.  Therefore
\[
 \PP(\Omega\text{ and selection fails})
 \le \E[\one_\Omega\overline S].
\]
Multiply \eqref{v1:eq:prefix-input} by $\one_\Omega$ and take expectations.
The action term is at most $C_Ph^{-2}D_h/J$.  Since
$\Omega\subseteq\Omega_{j+1}$, the nonnegative balance contribution is bounded
by $C_B B_{h,J}$.  Similarly,
\[
 \frac1J\sum_{j<J}\E[\one_\Omega V_h(X_j)]\le M_{h,J},\qquad
 \frac1J\sum_{j<J}\E[\one_\Omega W_h(X_j)]\le A_{h,J}.
\]
Divide by the respective squared thresholds and add the probability of
$\Omega^c$ to obtain \eqref{v1:eq:joint-failure}.
\end{proof}

\begin{remark}[Exactly what has changed]
The raw coordinate chart and same-action segment estimates must still hold
on the surviving prefix.  The theorem does \emph{not} require every candidate
to satisfy the high-order physical energy, spectral-tail, or initial/gauge
threshold.  Only the selected candidate must do so.  Replacing this argument
by a union bound for the event that every candidate has small energy would
unnecessarily cost a factor $J$ and would miss the purpose of the selection.
\end{remark}

\subsection{Transient drift supplies the occupation energy}
The stationary Lyapunov estimate in \cite{GT} is already available.
For the joint selector the required object is instead a transient occupation
average under the actual, possibly killed table.

\begin{lemma}[Transient occupation bound from the same finite drift]
\label[lemma]{v1:lem:drift}
Suppose $0<d_h\le d_{\max}$, $r_h=e^{-\kappa d_h}$, and for every legal state
and every $j$,
\begin{equation}
 \sum_yK_{h,j}^{\rm leg}(x,y)V_h(y)
 \le r_hV_h(x)+B d_h,
 \label{v1:eq:drift}
\end{equation}
with $\kappa,B>0$ independent of $h$.  Put
$M_{0,h}=\sum_x\mu_0(x)V_h(x)$.  Then
\begin{equation}
 \sum_x\mu_j(x)V_h(x)
 \le r_h^jM_{0,h}+\frac{B d_h(1-r_h^j)}{1-r_h},
 \label{v1:eq:transient}
\end{equation}
and
\begin{equation}
 M_{h,J}\le
 \frac{M_{0,h}(1-r_h^J)}{J(1-r_h)}+
 \frac{B d_h}{1-r_h}
 \le M_{0,h}+\frac{Be^{\kappa d_{\max}}}{\kappa}.
 \label{v1:eq:drift-occupation}
\end{equation}
If the same inequality holds for $W_h$, with initial expectation
$O(\rho_h^2)$ and injection $B_W\rho_h^2d_h$, then
$A_{h,J}=O(\rho_h^2)$ uniformly in $J$.
\end{lemma}
\begin{proof}
Multiplication of \eqref{v1:eq:drift} by $\mu_j(x)$ and summation gives
$M_{j+1}\le r_hM_j+B d_h\sum_x\mu_j(x)\le r_hM_j+B d_h$.
Iterating the scalar recursion proves \eqref{v1:eq:transient}.
Summing its first term geometrically and bounding its second term by
$B d_h/(1-r_h)$ proves the first inequality in
\eqref{v1:eq:drift-occupation}.  Use
$1-e^{-x}\ge xe^{-x}$ and
$(1-r_h^J)/(J(1-r_h))\le1$ for the second.
The identical argument applied to $W_h$ proves the last assertion.
\end{proof}

\begin{remark}[The clocks are different]
The parameter $d_h$ is the duration associated with one native renewal cycle
in a drift estimate.  The parameter $\upsilon_h$ is an action-comparison
scale.  The parameter $h$ is a spacetime record resolution.  The index $j$
labels candidate spacetime records, not physical time slices within one such
record.  None of these quantities is identified with another unless the
actual source supplies that identification.
\end{remark}

\begin{remark}[A state drift is not automatically a record drift]
A positive map on a finite noncommutative algebra can control
$\omega_X(V_X)$ without providing \eqref{v1:eq:drift} for a proposed classical
record process.  Application requires an actual record instrument or table
whose expectation of the declared readout agrees with the corresponding
source expectation, and whose conditional drift is the displayed one.
Changing the observation instrument or inserting a preparation/gate changes
the source law and must be recorded.  This compatibility is not inferred from
the word ``regenerative.''
\end{remark}

\section{From a native source energy to the required spacetime tails}
\label{v1:sec:reader}

\subsection{An exact finite reader test}
Let $\mathsf H_h$ be the finite Hilbert space of actual source coordinates and
let $S_h\succeq0$ be a supplied protected-soft energy.  The generic shorting
identity producing such an $S_h$ is already in \cite{GT}.  What matters here is
its relation to the \emph{actual Fourier reader}, not an unrelated operator
with a compact resolvent.

For a fixed set of tail modes, let $T_{h,\ell}:\mathsf H_h\to\C^d$ be the
linear Fourier coefficient readers, with known offsets $u^{\rm soft}$ handled
separately.  Set
\[
 \mathcal T_{h,K,m}x=(T_{h,\ell}x)_{\abs\ell_\infty>K},\qquad
 w_\ell=(1+\abs\ell_1/K)^m.
\]
All statements below apply to the realification of the complex coefficient
space.  Write $S_h^{\dagger/2}=(S_h^\dagger)^{1/2}$.

\begin{proposition}[Exact energy-to-tail constant for a specified reader]
\label[proposition]{v1:prop:reader}
There is a finite constant $c$ such that
\begin{equation}
 \sum_{\abs\ell_\infty>K}w_\ell\norm{T_{h,\ell}x}
 \le c\,\ip{x}{S_hx}^{1/2}\qquad(x\in\mathsf H_h)
 \label{v1:eq:reader-bound}
\end{equation}
if and only if $\mathcal T_{h,K,m}\Ker S_h=0$.  When this holds, its optimal
constant is
\begin{align}
 c_{h,K,m}
 &=\norm{\mathcal T_{h,K,m}S_h^{\dagger/2}}_{2\to\ell^1(w)},\notag\\
 c_{h,K,m}^2
 &=\max_{\norm{\eta_\ell}\le w_\ell}
 \ip{\mathcal T_{h,K,m}^*\eta}
      {S_h^\dagger\mathcal T_{h,K,m}^*\eta}.
 \label{v1:eq:reader-exact}
\end{align}
Consequently, for a record with $\ip{x}{S_hx}\le V_h(x)$ and primitive reader
$u=u^{\rm soft}+\mathcal R_hx$,
\begin{equation}
 \tau_{h,m}(K;u)
 \le \tau_{h,m}(K;u^{\rm soft})+c_{h,K,m}\sqrt{V_h(x)}.
 \label{v1:eq:reader-tail}
\end{equation}
\end{proposition}
\begin{proof}
If $x\in\Ker S_h$, the right side of \eqref{v1:eq:reader-bound} is zero, so
necessity follows.  Conversely restrict to $(\Ker S_h)^\perp$ and write
$x=S_h^{\dagger/2}y$, with
$\norm y^2=\ip{x}{S_hx}$.  Finite dimensionality gives the first formula in
\eqref{v1:eq:reader-exact}.  Weighted block $\ell^1$ duality gives
\[
 \sum_\ell w_\ell\norm{v_\ell}
 =\max_{\norm{\eta_\ell}\le w_\ell}\Re\ip{\eta}{v}.
\]
Interchanging the two suprema over $y$ and $\eta$ and evaluating the Hilbert
space supremum proves the second formula.  The offset conclusion is the
triangle inequality applied to the literal coefficient reader.
\end{proof}

This criterion permits direct finite-matrix verification without postulating
a Sobolev identification.  Its constant may diverge with $h$; a fixed-cutoff
finite norm is not a uniform continuum estimate.  If an uncontrolled protected
mode contributes to the physical tail, the kernel condition fails.  It must
then be retained with its own bound rather than silently deleted.

\subsection{A quantitative Sobolev specialization}
Use normalized Haar measure on $\T^4$ and
\begin{equation}
 \norm u_{H^q}^2=
 \sum_{\ell\in\Lambda_h}(1+\abs\ell_2^2)^q\abs{\widehat u(\ell)}^2.
 \label{v1:eq:Hq}
\end{equation}
The next calculation gives the exact derivative demand of this route into
\cref{v1:in:closure}.  It is stronger than a qualitative compact screen and
concerns the full time--space reader.

\begin{theorem}[Sharp weighted Fourier-tail transfer]
\label[theorem]{v1:thm:Fourier}
For $m\ge0$ and $q>m+2$, there is a constant $C_{q,m}$ independent of
$h,K,N_h$ such that every represented trigonometric record satisfies
\begin{equation}
 \tau_{h,m}(K;u)\le C_{q,m}K^{2-q}\norm u_{H^q},\qquad K\ge1.
 \label{v1:eq:Fourier-tail}
\end{equation}
The exponent $2-q$ is optimal for a uniform estimate on the $H^q$ unit ball.
At fixed $K$, no cutoff-independent $H^q$-to-$\tau_{h,m}$ bound exists when
$q\le m+2$.
\end{theorem}
\begin{proof}
Cauchy--Schwarz gives the exact bound
\begin{equation}
 \tau_{h,m}(K;u)\le A_{N_h,K,m,q}\norm u_{H^q},\qquad
 A_{N,K,m,q}^2=
 \sum_{K<\abs\ell_\infty\le N}
 \frac{(1+\abs\ell_1/K)^{2m}}{(1+\abs\ell_2^2)^q}.
 \label{v1:eq:exact-tail-sum}
\end{equation}
On the tail, $1+\abs\ell_1/K\le C\abs\ell_2/K$.
Partition the tail into dyadic shells with radius comparable to $2^aK$.
A shell contains at most $C(2^aK)^4$ lattice points.  Its contribution is at
most
\[
 C K^{-2m}(2^aK)^{2m-2q+4}.
\]
The series over $a\ge0$ converges exactly when $q>m+2$, and its sum is at most
$C K^{4-2q}$.  Taking square roots proves
\eqref{v1:eq:Fourier-tail}.

For sharpness, consider the symmetric shell
$\{\ell:K<\abs\ell_\infty\le2K\}$, with $N\ge2K$, and take a scalar reader
whose Fourier coefficients there are a common positive multiple of
$K^{-q-2}$.  The shell has comparable to $K^4$ modes; its $H^q$ norm is bounded
above and below by positive constants, whereas its weighted tail is comparable
to $K^{2-q}$.  The coefficients can be chosen conjugate symmetric, so the
example is real.

Finally, the Cauchy--Schwarz constant in \eqref{v1:eq:exact-tail-sum} is
attained after normalization by scalar coefficients proportional to
$(1+\abs\ell_1/K)^m(1+\abs\ell_2^2)^{-q}$ on the finite tail.
Lower shell counts comparable to the fourth power of the radius show that
$A_{N,K,m,q}$ diverges as $N\to\infty$ for fixed $K$ when $q\le m+2$.
At equality, each sufficiently large dyadic shell contributes a comparable
positive amount to its square, giving logarithmic divergence.
\end{proof}

\begin{corollary}[A finite positive-form certificate for the Sobolev reader]
\label[corollary]{v1:cor:reader-Sobolev}
Let $\mathsf L_h$ denote the diagonal Fourier multiplier
$(1+\abs\ell_2^2)^{1/2}$ on the reader space.  If, on the actual retained
source coordinates,
\begin{equation}
 \mathcal R_h^*\mathsf L_h^{2q}\mathcal R_h\preceq C_R S_h,
 \qquad \norm{u_h^{\rm soft}}_{H^q}\le C_R,
 \label{v1:eq:reader-LMI}
\end{equation}
then
\begin{equation}
 \norm{u_h}_{H^q}\le C\bigl(1+\sqrt{V_h}\bigr),\qquad
 \tau_{h,j}(K;u_h)\le C K^{2-q}\bigl(1+\sqrt{V_h}\bigr)
 \quad(0\le j\le m<q-2).
 \label{v1:eq:reader-Hq}
\end{equation}
All operator adjoints and positive-form comparisons use the declared
source and Fourier Hilbert structures.
\end{corollary}
\begin{proof}
The quadratic form inequality gives
$\norm{\mathsf L_h^q\mathcal R_hx}^2\le C_R\ip{x}{S_hx}\le C_RV_h$.
Add the offset and apply \cref{v1:thm:Fourier} to each indicated order.
\end{proof}

\begin{remark}[No hidden derivative gain]
The positive-form test \eqref{v1:eq:reader-LMI} is a statement about a specific
source energy and its \emph{physical reader}.  It does not follow just from a
Poincar\'e gap or a dimension-three spatial Weyl estimate.  In particular it
includes all auxiliary time frequencies needed by \eqref{v1:eq:tau}.
No spectral projection is applied to the record.  If only a weaker or
cutoff-dependent reader estimate is available, its actual constant must be
inserted in \eqref{v1:eq:reader-tail}; it cannot be replaced by $C_R$.
\end{remark}

\section{The same-source closure theorem and an explicit admissible regime}
\label{v1:sec:composition}

\subsection{General finite-tail composition}
\begin{theorem}[Joint native selection composed with full-record closure]
\label[theorem]{v1:thm:composition}
Consider one family of actual finite accepted tables as in
\cref{v1:sec:joint}, with the same raw action, full physical normalization,
decoder, reference constrained initial datum, and chart realization as
\cref{v1:in:closure}.  Suppose the occupation budgets
\eqref{v1:eq:occupations} and the prefix estimate hold.  Suppose also that the
native energy reader satisfies, on legal records,
\begin{equation}
 \tau_{h,j}(K_h;u_h(x))
 \le t_{h,j}+c_{h,j}\sqrt{V_h(x)},\qquad j=2,m,
 \label{v1:eq:general-reader}
\end{equation}
where $m=k+2$, and $t_{h,j},c_{h,j}$ are known finite reader bounds.
Set $T_{h,j}=t_{h,j}+c_{h,j}R_h$, and form the deterministic upper budget
\begin{align}
 \bar\epsilon_c&=s_h+hK_h^3+K_h^2T_{h,2},\notag\\
 \bar L_h&=1+\log\left(2+\frac{C_kK_h^{k+5}}{\bar\epsilon_c}\right),\notag\\
 \bar F_h&=\bar\epsilon_cK_h^{k+1}\bar L_h^{k+1}
                        +K_h^{k+2}T_{h,m},\qquad
 \bar D_h=a_h+\bar F_h.
 \label{v1:eq:composed-budget}
\end{align}
Assume $K_h\to\infty$, $hK_h\to0$, $T_{h,m}\to0$,
$\bar\epsilon_c\to0$, and $\bar D_h\to0$.
Then, except on an event bounded by \eqref{v1:eq:joint-failure}, the joint
selector outputs an actually observed record such that
\begin{equation}
 \mathfrak E_k(z_{h,j_*},z_*)\le C\bar D_h.
 \label{v1:eq:composed-result}
\end{equation}
Its physical zero-order residual is at most $C(s_h+hK_h^3)$ and its literal
Cartan discrepancy is at most $ChK_h^3$.
If both the failure bounds and $\bar D_{h_n}$ are summable on a cofinal chain,
then under any common coupling the selected successful records eventually have
summable adjacent state, curvature and stress differences almost surely.
\end{theorem}
\begin{proof}
Outside the failure event of \cref{v1:thm:joint}, the selected record has
$g_h\le s_h$, $V_h\le R_h^2$, and $\sqrt{W_h}\le a_h$.
Hence \eqref{v1:eq:general-reader} yields $\tau_{h,j}\le T_{h,j}$ and
\eqref{v1:eq:gW} gives $i_{h,k}+\gamma_{h,k}\le a_h$.
Use the \emph{upper} bounds $s_h,T_{h,j}$ to define the admissible positive
budgets in \cref{v1:in:closure}.  Its statement permits such certified upper
budgets; it is unnecessary to compare a logarithmic expression evaluated at
two different actual residuals.  These are precisely
\eqref{v1:eq:composed-budget}.  For small $h$, all smallness conditions in that
input hold, proving \eqref{v1:eq:composed-result} and the residual statements.

The first Borel--Cantelli lemma makes only finitely many cutoff failures occur
when their upper bounds are summable.  On the successful tail, all records are
compared to the same fixed physical solution.  In each displayed norm an
adjacent difference is bounded by $C(\bar D_{h_n}+\bar D_{h_{n+1}})$.
Summing proves the assertion.  No independence between cutoffs is used.
\end{proof}

\begin{remark}[What is and is not automatic]
The new theorem combines three properties on one \emph{selected} native
record.  It does not assert that every accepted record converges, that the last
record is good, or that the energy drift makes a raw gravitational Cauchy
update stable.  Basic chart exits, action identity, and the full-record
reader realization still have to be supplied by the actual source.  Neither
construction of a comparison law nor reweighting a law after observing its
outcomes verifies these hypotheses for a different primitive source.
\end{remark}

\subsection{A polynomial regime that discharges the tail hypotheses}
The next corollary expresses the composition in quantities with explicit
cutoff powers.  It is intentionally not optimized in derivative order.

\begin{corollary}[Occupation energy gives a quantitative selected Einstein--SM branch]
\label[corollary]{v1:cor:powers}
Let $k\ge4$, $m=k+2$, and assume the reader estimate
\eqref{v1:eq:reader-Hq} holds for a fixed $q>k+5$.  Suppose
$M_{h,J}\le M$ and, for some positive numbers
$a,b,c,\epsilon$ and $d_A\ge0$,
\begin{align}
 D_h&\le C h^{-d_A},&
 J_h&\ge c_Jh^{-(2+d_A+2b+2\epsilon)},\notag\\
 B_{h,J_h}&\le C h^{2b+2\epsilon},&
 A_{h,J_h}&\le C h^{2c+2\epsilon},&
 p_h^{\rm exit}&\le C h^{2\epsilon}.
 \label{v1:eq:power-source}
\end{align}
Choose thresholds
\begin{equation}
 s_h=h^b,\qquad R_h=h^{-a},\qquad a_h=h^c,\qquad
 K_h\asymp h^{-\beta},
 \label{v1:eq:power-thresholds}
\end{equation}
where
\begin{equation}
 0<\beta<\min\left\{\frac{b}{k+1},\frac1{k+4}\right\},\qquad
 0<a<\beta(q-k-5).
 \label{v1:eq:power-window}
\end{equation}
Then the failure probability is at most
$C(h^{2a}+h^{2\epsilon})$.  On success,
\begin{equation}
 \begin{split}
 \mathfrak E_k(z_{h,j_*},z_*)&\le C h^\rho(1+\abs{\log h})^{k+1},\\
 \rho&=\min\{c,\ b-\beta(k+1),\ 1-\beta(k+4),\
                    \beta(q-k-5)-a\}>0.
 \end{split}
 \label{v1:eq:power-rate}
\end{equation}
Along $h_n\asymp2^{-n}$, this yields eventual almost-sure selected closure and
summable adjacent differences in the inherited state, curvature and stress
norms.
\end{corollary}
\begin{proof}
Substitute \eqref{v1:eq:power-source} and \eqref{v1:eq:power-thresholds} into
\eqref{v1:eq:joint-failure}.  The prefix term is bounded by
\[C h^{-2-d_A}h^{2+d_A+2b+2\epsilon}h^{-2b}=C h^{2\epsilon}.\]
The balance and alignment terms have the same bound, and the energy term is
$M h^{2a}$.

On success, \eqref{v1:eq:reader-Hq} gives, for $j=2,m$,
\begin{equation}
 T_{h,j}\le C h^{-a}K_h^{2-q}.
 \label{v1:eq:power-tail}
\end{equation}
Since $q>k+5$ and \eqref{v1:eq:power-window} holds, this tends to zero and
\[
 \bar\epsilon_c\le C\bigl(h^b+hK_h^3+h^{-a}K_h^{4-q}\bigr)\longrightarrow0.
\]
By definition $\bar\epsilon_c\ge hK_h^3$, after fixing its upper constants
at least one.  Thus $\bar L_h\le C(1+\abs{\log h})$.
Equation \eqref{v1:eq:composed-budget} is bounded by the sum of
\begin{align*}
 &h^bK_h^{k+1}(1+\abs{\log h})^{k+1},\\
 &hK_h^{k+4}(1+\abs{\log h})^{k+1},\\
 &h^{-a}K_h^{k+5-q}(1+\abs{\log h})^{k+1},\\
 &h^{-a}K_h^{k+4-q},\\
 &h^c.
\end{align*}
The fourth exponent is larger by $\beta$ than the third.  The remaining
exponents are exactly those in \eqref{v1:eq:power-rate}.
All required quantities therefore tend to zero, and
\cref{v1:thm:composition} applies.  Every positive power of a dyadic spacing
times a fixed logarithmic power is summable, as are the failure bounds.
\end{proof}

\begin{proposition}[Sharp derivative threshold for this budget composition]
\label[proposition]{v1:prop:threshold}
For the sole information $\norm{u_h}_{H^q}\le R_h$ and the transfer
\eqref{v1:eq:Fourier-tail}, the worst-case term
$K_h^{k+3}\tau_{h,2}$ in the inherited forcing budget is of order
$R_hK_h^{k+5-q}$.  If $R_h$ is bounded below by a positive constant, this
information alone cannot force that budget term to vanish when $q\le k+5$.
\end{proposition}
\begin{proof}
The upper estimate follows from \cref{v1:thm:Fourier} with $m=2$ whenever
$q>4$.  The shell construction in its proof, scaled to $H^q$ norm $R_h$,
gives a matching lower estimate.  If $q\le4$, the cutoff-independent
weighted-tail estimate is already unavailable.  Thus $q=k+5$ is the limiting
exponent for a vanishing bound by this route when $R_h$ does not tend to zero.
\end{proof}

This is a threshold for the stated \emph{budget argument}, not a theorem that
physical Einstein--SM solutions require $H^{k+5+}$ data.  Cancellations,
better action-specific reader bounds, or a different residual transfer can
improve the demand.  No such improvement is inferred here.

\begin{example}[A fully specified set of compatible cutoff powers]
\label[example]{v1:ex:powers}
Take
\begin{equation}
 k=4,\quad m=6,\quad q=14,\quad
 \beta=a=\epsilon=\frac1{16},\quad b=1,\quad c=\frac14.
 \label{v1:eq:explicit-powers}
\end{equation}
Then it suffices that the actual source satisfy
\begin{align}
 M_{h,J}&=O(1),& B_{h,J}&=O(h^{17/8}),&
 A_{h,J}&=O(h^{5/8}),\notag\\
 p_h^{\rm exit}&=O(h^{1/8}),&
 D_h&=O(h^{-d_A}),&
 J_h&\gtrsim h^{-(d_A+33/8)}.
 \label{v1:eq:explicit-source-powers}
\end{align}
On the selected successful record,
\[
 \sigma_h\le h,\qquad
 \tau_{h,2},\tau_{h,6}\le C h^{11/16},\qquad
 i_{h,4}+\gamma_{h,4}\le h^{1/4}.
\]
The closure error is
$O(h^{1/4}(1+\abs{\log h})^5)$, and in particular $O(h^{1/8})$.
The failure probability is $O(h^{1/8})$; the physical zero-order residual
and raw Cartan discrepancy are $O(h^{13/16})$.
\end{example}
\begin{proof}
The source exponents are direct substitutions in
\eqref{v1:eq:power-source}.  The reader exponent is
$\beta(q-2)-a=12/16-1/16=11/16$.
The four candidate exponents in \eqref{v1:eq:power-rate} are
$1/4,11/16,1/2,1/4$.  The minimum is $1/4$.
Also $hK_h^3\asymp h^{13/16}$ dominates $h$ as $h\downarrow0$.
A spare positive power absorbs the fifth logarithmic power.
\end{proof}

\paragraph{Interpretation of the example.}
These are compatible \emph{certificate powers}, not a newly verified primitive
law.  In particular, an $H^{14}$-controlling reader energy and its actual
regenerative drift have not been proved in this file for the physical source
in the monographs.  The example shows precisely what would suffice, including
the accepted-comparison count and the probability cost, without claiming
those powers have already been observed or derived from that law.

\section{Conditional innovations and the physical balance floor}
\label{v1:sec:innovation}

The previous composition requires a small mean \emph{squared realized}
balance.  The regenerative packet distinguishes predictable mean control from
fresh fluctuations and permits retained stochastic sources.  The following
calculation quantifies why these cannot automatically be substituted into a
pathwise stationarity estimate.

\begin{theorem}[Exact conditional balance decomposition and innovation bounds]
\label[theorem]{v1:thm:innovation}
Let $\mathsf H$ be a finite real Hilbert space and $\Theta\subset\mathsf H$ a
convex set.  Suppose $p=\nabla\Psi$ is $g_*$-strongly monotone and
$g^*$-Lipschitz on $\Theta$, while $E$ is $L$-Lipschitz.  Let $\upsilon>0$
satisfy $m_\upsilon:=g_*/\upsilon-L>0$, and define
\[
 H_\upsilon(y)=E(y)+\upsilon^{-1}p(y),\qquad
 L_\upsilon=L+g^*/\upsilon.
\]
Let $X,Y$ be the physical field-coordinate readouts of an actual consecutive
pair of records in $\Theta$, and let
$\mathcal F$ contain $X$ and the observed past.  Assume the following moments
are finite, and put
\begin{align}
 b&=H_\upsilon(Y)-\upsilon^{-1}p(X),&
 \mu&=\E[b\mid\mathcal F],\notag\\
 \nu&=\E\bigl[\norm{Y-\E[Y\mid\mathcal F]}^2\mid\mathcal F\bigr].
 \label{v1:eq:conditional-moments}
\end{align}
Then
\begin{equation}
 \norm\mu^2+m_\upsilon^2\nu
 \le \E[\norm b^2\mid\mathcal F]
 \le \norm\mu^2+L_\upsilon^2\nu.
 \label{v1:eq:innovation-floor}
\end{equation}
In particular, a small predictable balance $\mu$ does not remove the
conditional innovation term.
\end{theorem}
\begin{proof}
For $y,y'\in\Theta$, strong monotonicity of $p$ and Lipschitz continuity of
$E$ give
\[
 \ip{H_\upsilon(y)-H_\upsilon(y')}{y-y'}
 \ge (g_*/\upsilon-L)\norm{y-y'}^2.
\]
Cauchy--Schwarz implies
$\norm{H_\upsilon(y)-H_\upsilon(y')}\ge m_\upsilon\norm{y-y'}$.
The upper Lipschitz bound is immediate.

Condition on $\mathcal F$ and take conditionally independent copies $Y,Y'$.
This is only a product-measure device for computing the conditional moments;
it does not require the source to generate independent consecutive records.
The Hilbert variance identity gives
\begin{align*}
 \E[\norm{H_\upsilon(Y)-\E[H_\upsilon(Y)\mid\mathcal F]}^2\mid\mathcal F]
 &=\frac12\E[\norm{H_\upsilon(Y)-H_\upsilon(Y')}^2\mid\mathcal F],\\
 \nu&=\frac12\E[\norm{Y-Y'}^2\mid\mathcal F].
\end{align*}
Apply the lower and upper Lipschitz estimates and use
\[
 \E[\norm b^2\mid\mathcal F]=\norm\mu^2+
 \E[\norm{b-\mu}^2\mid\mathcal F].
\]
Since $p(X)$ is measurable with respect to $\mathcal F$, subtracting it does
not change the conditional variance.  This proves
\eqref{v1:eq:innovation-floor}.
\end{proof}

\begin{corollary}[Necessary and sufficient noise scales for the prefix route]
\label[corollary]{v1:cor:noise-scale}
Under \eqref{v1:eq:prefix-geometry} and $\upsilon_h=\upsilon_0h^2$,
there are fixed $c_-,c_+>0$ such that
\begin{equation}
 \norm{\mu_j}^2+c_-h^{-4}\nu_j
 \le \E[\norm{b_j}_{0,h}^2\mid\mathcal F_j]
 \le \norm{\mu_j}^2+c_+h^{-4}\nu_j.
 \label{v1:eq:h-noise-floor}
\end{equation}
Thus an occupation-mean balance bound $O(h^{2b+2\epsilon})$ necessarily
requires an occupation-mean physical innovation trace
$O(h^{4+2b+2\epsilon})$ on a legal process.  Conversely, a predictable-balance
mean square $O(h^{2b+2\epsilon})$ and an innovation trace
$O(h^{4+2b+2\epsilon})$ suffice for that balance bound.
\end{corollary}
\begin{proof}
The lower and upper constants in \cref{v1:thm:innovation} are respectively
$(g_*/\upsilon_0-K_0)h^{-2}$ and
$(g^*/\upsilon_0+K_0)h^{-2}$.  They have strictly positive fixed lower
coefficient by the assumed choice of $\upsilon_0$.  Squaring and averaging
\eqref{v1:eq:innovation-floor} proves both assertions.
\end{proof}

For a killed table, the same statement is applied to the normalized legal
conditional law at a row with positive legal mass and then multiplied by that
mass.  Exits remain counted in $p_h^{\rm exit}$; they are not removed by
renormalizing the full experiment.

\begin{corollary}[A conditional obstruction to a naive regenerative composition]
\label[corollary]{v1:cor:retained-noise}
Suppose the actual \emph{physical record coordinates} have a fresh conditional
innovation trace at least $c d_h$ per legal accepted comparison, in occupation
mean, with $c>0$ independent of $h$.  If the hypotheses of
\cref{v1:cor:noise-scale} hold, then
\begin{equation}
 B_{h,J}\ge c' h^{-4}d_h
 \label{v1:eq:retained-noise}
\end{equation}
for a fully legal process.  Consequently the sufficient balance scale
$B_{h,J}=O(h^{2b+2\epsilon})$ is impossible for this route unless
$d_h=O(h^{4+2b+2\epsilon})$.
\end{corollary}
\begin{proof}
Average the lower bound in \eqref{v1:eq:h-noise-floor} and discard its
nonnegative predictable-mean term.
\end{proof}

The innovation trace is taken in an orthonormal basis for the physical mass
inner product, so its quadrature weights are already included.

This is \emph{not} a nonexistence theorem for classical limits of noisy sources.
It rules out an unjustified use of this particular squared-balance certificate.
Noise may be smaller, may lie in genuinely unread protected coordinates,
may be removed by a source-defined observation protocol whose errors are
retained, or may lead to a stochastic continuum rather than the deterministic
classical branch.  None of these alternatives permits deletion of a
fluctuation that remains in the observed full physical record.

\begin{example}[Zero mean balance with a nonzero realized floor]
\label[example]{v1:ex:linear-noise}
In one dimension let $p(y)=y$ and $E(y)=\lambda y$, $\lambda\ge0$.
For a fixed previous state $X$, choose an actual conditional law
\[
 Y=\frac{X}{1+\upsilon\lambda}+\sqrt{d}\,\xi,
 \qquad \E[\xi\mid X]=0,\quad \E[\xi^2\mid X]=1.
\]
Then
\[
 \E[b\mid X]=0,\qquad
 \E[b^2\mid X]=(\lambda+\upsilon^{-1})^2d.
\]
With $\upsilon\asymp h^2$, this is of order $h^{-4}d$.
\end{example}
\begin{proof}
Substitute the expression for $Y$ into
$b=(\lambda+\upsilon^{-1})Y-\upsilon^{-1}X$.
The deterministic terms cancel and
$b=(\lambda+\upsilon^{-1})\sqrt d\,\xi$.
\end{proof}

\paragraph{Relation to the source covariance theorem.}
The fluctuation--dissipation result in \cite{GT} controls a declared source
covariance.  To use \eqref{v1:eq:retained-noise}, that covariance must first be
identified with the conditional innovation in the \emph{same physical record
coordinates}.  Stationary variance and one-step innovation variance are not
the same quantity.  Likewise, a shorted predictable source norm is not the
unshorted squared norm in \eqref{v1:eq:balance}.  The present floor identifies
the exact place where an illicit substitution would occur.

\section{What the first iteration resolves, and the next source-specific test}
\label{v1:sec:ledger}

\subsection{A precise nonduplicating conclusion}
The main mathematical result of this iteration is the composition
\[
 \begin{gathered}
 \text{one actual finite accepted-record law}\\
 +\text{same-action prefix balance}\\
 +\text{native occupation-energy control}\\
 +\text{a verified physical reader-tail map}\\
 \Downarrow\\[-0.3em]
 \text{one jointly selected record with small full Euler row,}\\
 \text{controlled spacetime leakage, and aligned initial/gauge data}\\
 \Downarrow\\[-0.3em]
 \text{quantitative classical closure by the inherited PDE theorem}.
 \end{gathered}
\]
The first downward implication, with its probability estimate and rate
thresholds, is proved here.  The second uses \cref{v1:in:closure}.
The record is never replaced by a projected field, a mean of candidates, or a
trajectory of a newly inserted numerical equation.

This removes a genuine logical mismatch: neither the smallest slope nor a
stationary mean energy separately identifies a usable classical record.
A same-record occupation selector does, without requiring high-order bounds
on every member of a long prefix.  It also exposes two quantitative costs that
are not discharged by citing the regenerative packet: the physical
space--time reader norm, and the realized innovation contribution to the
full action balance.

\subsection{The remaining finite certificates must be checked, not renamed}
For an independently specified native source, the following tests remain.
\begin{enumerate}[label=\textup{(N\arabic*)},leftmargin=3em]
\item \textbf{Actual law and action identity.}
The recorded transition tables or instrument, its clock, decoder, collars,
and the full common action must be the same objects used in the source
construction.  A covariance or a binary accepted bit alone does not determine
these tables or the full affine action covector.
\item \textbf{Physical reader and transient drift.}
Evaluate \eqref{v1:eq:reader-exact} or verify a suitable positive-form estimate
such as \eqref{v1:eq:reader-LMI}.  Verify its drift under the actual table,
including time frequencies.  For initial and harmonic alignment verify the
corresponding $W_h$ occupation bound, or its own drift.  Neither positivity
nor a spatial Poincar\'e gap supplies these identifications automatically.
\item \textbf{Realized balance and noise.}
Evaluate the actual conditional mean of \eqref{v1:eq:balance} and its
innovation trace in physical mass coordinates.  Theorem
\ref{v1:thm:innovation} gives both a sufficient bound and a necessary floor.
The required scalar potential and its Hessian bounds must be checked in the
actual coordinates; a source Fisher matrix cannot simply be declared to be
that Hessian.
\item \textbf{Legal-domain and clock control.}
Bound the exit probability and the accepted-prefix span $D_h$ and length
$J_h$.  The resulting operation count is not automatically a physical-time
bound.  The common initial datum and retained calibration must be unchanged
across cutoffs if one claims a single cofinal limit.
\end{enumerate}

These are a sharply specified set of source-level obligations.  This iteration
does not assert that the uploaded monographs already verify them
simultaneously, nor that an arbitrary rule must satisfy them.  In particular,
there is no claim of universal microscopic attraction or of a quantum
Einstein--Standard-Model construction.

\subsection{Priority for the next iteration}
The next calculation should use the actual source's current/outgoing mixed
panel and physical reader, not a redesigned probability law.  Compute the
conditional innovation trace and the reader constant first.  If the trace
violates \eqref{v1:eq:h-noise-floor} at the proposed cutoff powers, the
squared-balance route cannot be closed with those clocks and observations;
additional smoothing or a different theorem would have to be justified as
part of the source.  If it passes, the next task is the positive-form or
exact reader norm and its transient drift, including the harmonic alignment
readout.  Repeating the midpoint population, the generic Hodge short, the
prefix descent proof, or the coupled first-exit PDE theorem would not answer
these questions.

\subsection{Verification record}
The algebraic exponent example is reproducible with rational arithmetic.
A separate diagnostic script evaluates those fractions, verifies the
conditional-variance bounds on positive finite-dimensional linear examples,
and checks the exact weighted-reader dual formula on a small real matrix.
Those checks are only diagnostics.  The proofs are the Hilbert variance
identity, finite occupation inequalities, the exact reader duality, dyadic
Fourier counting, and the explicit substitutions above.  There is no claim
of machine-checked verification or of a simulation of the complete coupled
native source.

\begin{thebibliography}{9}
\raggedright
\bibitem{GT}
A.~P\'elissier,
\emph{Renewal Grand Tensor Monograph}, V076, supplied file
\nolinkurl{Renewal_Grand_Tensor_Monograph_V076(7).tex}.
The regenerative drift, shorted Hodge, spatial screen, and finite-type results
are used with their stated physical identifications and uniformity hypotheses.

\bibitem{CC}
A.~P\'elissier,
\emph{Renewal Classical Einstein--SM Closure Research Notes}, cumulative V27,
supplied file
\nolinkurl{Renewal_Classical_Einstein_SM_Closure_Research_Notes_V27(3).tex}.
The exact labels used as inputs are listed in \cref{v1:sec:audit}.

\bibitem{SM}
A.~P\'elissier,
\emph{SM--Einstein Continuum Monograph}, supplied file
\nolinkurl{SM_Einstein_Continuum_Monograph_Updated(3)(6).tex}.
Its joint quantum-continuum requirements are distinguished from the classical
record-level results of this file.

\bibitem{HM}
M.~Hairer and J.~C.~Mattingly,
\emph{Yet another look at Harris' ergodic theorem for Markov chains},
\href{https://arxiv.org/abs/0810.2777}{arXiv:0810.2777}.
This is background for the distinction between Lyapunov drift and ergodic
conclusions.  No Harris theorem or mixing conclusion is invoked in the new
proofs here.
\end{thebibliography}

% END IMMUTABLE VERSION 1 BODY
% APPEND VERSION 2 HERE, BEFORE THE FINAL END-DOCUMENT COMMAND.

% BEGIN IMMUTABLE VERSION 2 BODY
\clearpage
\begin{center}
{\Large\bfseries Version 2: Action increments before force-balance norms}\\[0.5em]
{\large An alternative same-source closure route with retained action-neutral noise}
\end{center}
\addcontentsline{toc}{part}{Version 2: Action increments and retained neutral noise}

\section{Upstream audit and the reason for changing the stationarity entrance}
\label{v2:sec:scope}

Version 1 separates two obligations: the physical reader must transfer native
energy to spacetime tails, and the accepted transitions must supply enough
stationarity on the same observed record. Its innovation inequality shows that
small \emph{realized Bregman force balance} is a stringent second obligation.
With comparison scale $\upsilon_h\asymp h^2$, an innovation trace $\nu_h$
contributes order $h^{-4}\nu_h$ to that particular certificate. This is not a
necessary condition for small action gradients under every possible source law.

There is an upstream alternative. A transition may move the full physical
record, and even have a substantial innovation, while preserving its action.
Charging the norm of that motion divided by $\upsilon_h$ can then be irrelevant
to whether the action gradient is small. The object examined in this iteration
is instead the \emph{conditional increment of the independently fixed signed
common action}. The law, its physical readout, and the action are not replaced.

\paragraph{Nonduplication audit.}
The positive-Fisher saddle obstruction is already present in \cite{CC}, at
\texttt{v23:thm:index}, \texttt{v23:thm:response}, and
\texttt{v23:thm:synthesis}. In particular, a positive-proximal full-history
response is not a general attracting mechanism for the Lorentzian action.
Those statements are not proved again here. The outgoing/noise Schur identity
and the regenerative Lyapunov estimates are likewise already in \cite{GT}, at
\texttt{thm:renewal-noise-shorted-Ward} and
\texttt{thm:renewal-Lyapunov-tightness}. The explicit phase-count generator used
in \cref{v2:sec:primitive} is the one supplied by
\texttt{thm:renewal-interchange-actual-audit}; its construction and OU limit are
not repeated. Neither a new midpoint regulator nor another general
Einstein--matter energy estimate is introduced.

The additions are a stopped-law action-drift entrance to the joint selector,
its derivation from two checkable transition mechanisms, an exact stationary
identity explaining the extra cost in the Bregman certificate, and a finite
pilot separating the two entrances. A final calculation tests the alternative
entrance on an explicit primitive already supplied in the monograph. All
results in this appended body have proofs below. The Einstein--SM conclusion
still uses \cref{v1:in:closure} with its original hypotheses.

\paragraph{Attribution and scope.}
Conditional-drift telescoping, Taylor expansion, and energy estimates are
standard tools. Entropy dissipation of a Markov law is also distinct from
stationarity of its individual field records; see, for example,
\cite{v2:KMS}. The familiar proximal strict-saddle phenomenon has independent
precedents such as \cite{v2:DD}. No general priority claim is made for these
principles. The purpose of the present proofs is the exact same-source
composition, physical normalization, and distinction between action-neutral
noise and unobserved or discarded noise. The concrete pilot is explicitly
identified as a diagnostic model, not as a verified primitive Einstein--SM law.

\section{Conditional action drift selects one complete physical record}
\label{v2:sec:drift-selection}

Retain the actual finite table, legal-edge mask, surviving subprobabilities
$\mu_j$, and failure probability $p_h^{\rm exit}$ of
\eqref{v1:eq:killed}. Write $\theta_h(x)$ for the full physical field-coordinate
reader and
\[
 a_h(x)=\mathcal A_h(\theta_h(x)),\qquad
 g_h(x)=\norm{\nabla_h\mathcal A_h(\theta_h(x))}_{0,h}.
\]
In this section the lower-case function $a_h(x)$ denotes the action value;
thresholds for initial/gauge alignment will be denoted $\alpha_h$. The action
is the same independently calibrated common action as in the closure input.
Suppose on the legal chart
\begin{equation}
 A_h^-\le a_h(x)\le A_h^+,\qquad D_h=A_h^+-A_h^-<\infty.
 \label{v2:eq:action-span}
\end{equation}
Neither $D_h$ nor the action's sign is assumed independent of the cutoff.

\begin{definition}[Stopped action increment on the given table]
\label[definition]{v2:def:action-increment}
For an active legal state define
\begin{equation}
 \mathfrak d_{h,j}(x)
 =\sum_y K_{h,j}^{\rm leg}(x,y)\,[a_h(y)-a_h(x)].
 \label{v2:eq:action-increment}
\end{equation}
On the first illegal transition the diagnostic stopped process retains the
\emph{last legal action value} and thereafter freezes it. The original exit
outcome remains a failure; no law is conditioned on survival. Consequently
\eqref{v2:eq:action-increment} is its actual conditional expected action
increment while it is active.
\end{definition}

If $q_j(x)=\sum_yK_{h,j}^{\rm leg}(x,y)$, then
$\mathfrak d_{h,j}=K_{h,j}^{\rm leg}a_h-q_j a_h$, not
$K_{h,j}^{\rm leg}a_h-a_h$. The latter would falsely credit
$(1-q_j)a_h$ as an action decrease on killing. This distinction matters for a
signed action and is not removed by shifting it: \eqref{v2:eq:action-increment}
is invariant under any constant shift of $a_h$.

\begin{theorem}[Stopped occupation slope bound from actual action drift]
\label[theorem]{v2:thm:occupation}
Suppose $\delta_h>0$, a constant $\kappa_A>0$ independent of the cutoff,
and a nonnegative finite readout
$r_{h,j}(x)$ satisfy, on every active legal row,
\begin{equation}
 \mathfrak d_{h,j}(x)
 \le-\kappa_A\delta_h g_h(x)^2+\delta_h r_{h,j}(x).
 \label{v2:eq:action-drift}
\end{equation}
Put
\begin{equation}
 \mathcal R_{h,J}=\frac1J\sum_{j<J}\sum_x\mu_j(x)r_{h,j}(x).
 \label{v2:eq:work-occupation}
\end{equation}
Then
\begin{equation}
 \frac1J\sum_{j<J}\sum_x\mu_j(x)g_h(x)^2
 \le\frac{D_h}{\kappa_A\delta_hJ}
                  +\frac{\mathcal R_{h,J}}{\kappa_A}.
 \label{v2:eq:slope-occupation}
\end{equation}
This conclusion uses no Bregman potential, no full force-balance vector,
no invariant distribution, and no conditional-innovation bound.
\end{theorem}
\begin{proof}
For each initially legal realization retain the action until an exit is
attempted, and then retain its last legal value. Denote the resulting bounded
random variable at stage $j$ by $Z_j$. Illegal initial mass may be assigned
zero after first replacing $a_h$ by $a_h-A_h^-$; this shift does not change
\eqref{v2:eq:action-increment}. Thus $0\le Z_j\le D_h$ on every realization.
Conditioning at an active row and summing gives
\[
 \E[Z_{j+1}-Z_j]=\sum_x\mu_j(x)\mathfrak d_{h,j}(x).
\]
Sum \eqref{v2:eq:action-drift} over the prefix. The left side telescopes,
while $\E[Z_0-Z_J]\le D_h$. Hence
\[
 \kappa_A\delta_h\sum_{j<J,x}\mu_j(x)g_h(x)^2
 \le D_h+\delta_h\sum_{j<J,x}\mu_j(x)r_{h,j}(x).
\]
Division by $\kappa_A\delta_hJ$ proves the estimate. The frozen diagnostic
retains, rather than suppresses, every exit in $p_h^{\rm exit}$.
\end{proof}

\begin{theorem}[Same-record selection from action drift and native energy]
\label[theorem]{v2:thm:joint}
Assume \cref{v2:thm:occupation}, and retain the nonnegative energy $V_h$, the
alignment readout $W_h=(i_{h,k}+\gamma_{h,k})^2$, and their actual occupation
budgets $M_{h,J},A_{h,J}$ from \eqref{v1:eq:occupations}. On a surviving prefix
select an observed index minimizing
\begin{equation}
 \mathcal S_j=\frac{g_h(X_j)^2}{s_h^2}
                   +\frac{V_h(X_j)}{R_h^2}
                   +\frac{W_h(X_j)}{\alpha_h^2},
 \qquad s_h,R_h,\alpha_h>0.
 \label{v2:eq:score}
\end{equation}
The probability of failing to output one record with
\begin{equation}
 g_h\le s_h,\qquad V_h\le R_h^2,\qquad
 i_{h,k}+\gamma_{h,k}\le\alpha_h
 \label{v2:eq:selected}
\end{equation}
is at most
\begin{equation}
 \boxed{\quad
 p_h^{\rm exit}
 +\frac{D_h}{\kappa_A\delta_hJ s_h^2}
 +\frac{\mathcal R_{h,J}}{\kappa_A s_h^2}
 +\frac{M_{h,J}}{R_h^2}
 +\frac{A_{h,J}}{\alpha_h^2}.\quad}
 \label{v2:eq:failure}
\end{equation}
The record and all of its noise are retained. In particular, the failure bound
contains no term proportional to $h^{-4}$ times its innovation trace.
\end{theorem}
\begin{proof}
Use the deterministic selector of \cref{v1:lem:selector}; it is not reproved.
Failure on $\Omega_J$ implies $J^{-1}\sum_{j<J}\mathcal S_j>1$. The expected
score on this event is bounded by its nonnegative active-prefix occupation
sum because $\Omega_J\subseteq\Omega_j$ for every $j<J$. Apply
\eqref{v2:eq:slope-occupation} to the first term and the definitions of
$M_{h,J},A_{h,J}$ to the other two. Markov's inequality and the retained exit
probability give \eqref{v2:eq:failure}.
\end{proof}

\begin{corollary}[Action-drift entrance to the inherited full-record limit]
\label[corollary]{v2:cor:closure}
Use the same raw action, decoder, calibration, physical variation bank, and
constrained reference initial datum as \cref{v1:in:closure}. Assume the reader
bounds \eqref{v1:eq:general-reader}. Form $T_{h,j}$ and the deterministic
forcing budget of \eqref{v1:eq:composed-budget}, replacing the alignment
threshold there by $\alpha_h$. If all of that budget's smallness conditions
hold, then outside an event bounded by \eqref{v2:eq:failure},
\begin{equation}
 \mathfrak E_k(z_{h,j_*},z_*)\le C(\alpha_h+\bar F_h).
 \label{v2:eq:closure}
\end{equation}
The zero-order equation and Cartan conclusions are the ones in
\cref{v1:in:closure}. Summable failure and error bounds on a common cofinal
chain give eventual almost-sure selected closure and summable adjacent state,
curvature, and stress differences. The prefix-balance input
\cref{v1:in:prefix} is not needed for this corollary.
\end{corollary}
\begin{proof}
Equation \eqref{v2:eq:selected} bounds the full raw row, native energy, and
initial/gauge error of the same record. The reader bound gives
$\tau_{h,j}\le t_{h,j}+c_{h,j}R_h$ at that record. These are exactly the upper
budgets required by \cref{v1:in:closure}; invoke that input, not a new PDE
argument. Borel--Cantelli and the triangle inequality through the same limit
prove the last assertion as in \cref{v1:thm:composition}.
\end{proof}

\begin{corollary}[Explicit cutoff powers for the alternative entrance]
\label[corollary]{v2:cor:powers}
Retain the Sobolev-reader hypothesis and threshold choices in
\cref{v1:cor:powers}, with $\alpha_h=h^c$. Let
$\delta_h\ge c_\delta h^d$, $D_h\le C h^{-d_A}$, and suppose
\begin{equation}
 \begin{aligned}
 J_h&\ge c_J h^{-(d+d_A+2b+2\epsilon)},&
 \mathcal R_{h,J_h}&\le C h^{2b+2\epsilon},\\
 M_{h,J_h}&\le C,& A_{h,J_h}&\le C h^{2c+2\epsilon},\\
 p_h^{\rm exit}&\le C h^{2\epsilon}.&&
 \end{aligned}
 \label{v2:eq:powers}
\end{equation}
With the window \eqref{v1:eq:power-window}, the same error exponent $\rho$ in
\eqref{v1:eq:power-rate} holds, and the failure probability is
$O(h^{2a}+h^{2\epsilon})$. For $d=2$ the required comparison-count exponent
is the one in Version 1, but $\mathcal R_{h,J}$ replaces $B_{h,J}$.
\end{corollary}
\begin{proof}
The action-span term in \eqref{v2:eq:failure} is bounded by
$C h^{-d-d_A}h^{d+d_A+2b+2\epsilon}h^{-2b}=C h^{2\epsilon}$.
The work, energy, and alignment terms give the other stated powers.
The reader and forcing calculation is unchanged from
\cref{v1:cor:powers}, since the selected thresholds are identical. Its proof
therefore gives $\rho>0$ and \eqref{v2:eq:closure}.
\end{proof}

\begin{remark}[An alternative source condition, not a concealed verification]
Equation \eqref{v2:eq:action-drift} is not inferred from the existence of the
common action, a stationary state, a Hodge gap, or small entropy production.
It must be derived from the given transitions and the independently fixed
action. Merely defining $r_{h,j}$ as the positive part of a failed inequality
is not evidence that its occupation is small. The next sections give two
specific, testable ways to establish it and an actual-source test where it
fails. The cycle clock, comparison scale $\delta_h$, and spacetime resolution
$h$ remain distinct unless the native protocol identifies them.
\end{remark}

\section{What the actual transition must do to the action}
\label{v2:sec:mechanisms}

All norms and adjoints in this section are in the physical mass Hilbert
structure. Statements in Euclidean coordinates are understood after its fixed
isometric identification. No Hessian interpretation of the source Fisher
matrix is needed.

\subsection{Exact signed action work, before a worst-case noise estimate}

\begin{proposition}[Conditional drift from mean motion and signed curvature work]
\label[proposition]{v2:prop:work}
Let $\mathcal A\in C^2(\Theta)$ on a convex chart and $E=\nabla\mathcal A$.
At an actual active row $x$, add the illegal mass to a diagnostic self-transition
at $x$, so its mass is one and its action increment is exactly
\eqref{v2:eq:action-increment}. Write $\Delta=Y-x$ and suppose the conditional
mean has the independently tested incidence
\begin{equation}
 m(x):=\E[\Delta\mid x]
       =-\delta M(x)E(x)+\delta e(x),\qquad
 \frac{M+M^*}{2}\succeq\lambda I,
 \label{v2:eq:mean-incidence}
\end{equation}
where $\lambda>0$. Define the exact signed work
\begin{equation}
 \mathcal W(x)
 =\E\left[\int_0^1(1-t)D^2\mathcal A(x+t\Delta)
                              [\Delta,\Delta]dt\,\middle|\,x\right].
 \label{v2:eq:work}
\end{equation}
Then \eqref{v2:eq:action-drift} holds with
\begin{equation}
 \kappa_A=\lambda/2,\qquad
 r(x)=\frac{\norm{e(x)}^2}{2\lambda}+\frac{[\mathcal W(x)]_+}{\delta},
 \qquad [t]_+=\max\{t,0\}.
 \label{v2:eq:work-remainder}
\end{equation}
Only the action work of the conditional innovation is charged, not its entire
variance by definition.
\end{proposition}
\begin{proof}
Taylor's integral identity gives exactly
\[
 \E[\mathcal A(Y)-\mathcal A(x)\mid x]
       =\ip{E(x)}{m(x)}+\mathcal W(x).
\]
Insert \eqref{v2:eq:mean-incidence}. The symmetric part of $M$ gives
$\ip{E}{ME}\ge\lambda\norm E^2$, and
$\ip E e\le\lambda\norm E^2/2+\norm e^2/(2\lambda)$.
The stated inequality follows. Adding illegal mass at $\Delta=0$ changes
neither the legal increment sum nor this identity.
\end{proof}

\begin{proposition}[Moment-resolved action work and the generic noise cost]
\label[proposition]{v2:prop:work-moments}
Let $Q(x)=\Cov(\Delta\mid x)$ and $H(x)=D^2\mathcal A(x)$. If the Hessian is
$L_3$-Lipschitz on the connecting segments, then
\begin{equation}
 \mathcal W(x)=\tfrac12\ip{m(x)}{H(x)m(x)}
                    +\tfrac12\tr(H(x)Q(x))+\mathcal E_3(x),
 \quad
 |\mathcal E_3(x)|\le\frac{L_3}{6}\E[\norm\Delta^3\mid x].
 \label{v2:eq:work-moments}
\end{equation}
In particular, for a quadratic action this formula is exact with
$\mathcal E_3=0$ and retains the sign of the covariance--Hessian contraction.

Alternatively, if only $\norm{D^2\mathcal A}\le L$ and $\norm M\le\Lambda$
are known, and $L\delta\Lambda^2\le\lambda/4$, then
\eqref{v2:eq:action-drift} holds with
\begin{equation}
 \kappa_A=\lambda/4,\qquad
 r(x)=\left(\frac1{2\lambda}+L\delta\right)\norm{e(x)}^2
                            +\frac{L}{2\delta}\tr Q(x).
 \label{v2:eq:generic-work-cost}
\end{equation}
Thus for $L\asymp h^{-2}$ and $\delta\asymp h^2$, a worst-case estimate still
costs $h^{-4}\tr Q$. The alternative route improves that cost only when the
actual action-work structure is verified more sharply.
\end{proposition}
\begin{proof}
Subtract the Hessian at $x$ inside \eqref{v2:eq:work}. Its Lipschitz estimate
gives a remainder at most
$L_3\E\norm\Delta^3\int_0^1t(1-t)dt=L_3\E\norm\Delta^3/6$.
The identity
$\E[\Delta\otimes\Delta\mid x]=m\otimes m+Q$ gives the first formula.
For the second assertion use
\[
 \mathcal W\le\frac L2(\norm m^2+\tr Q)
 \le L\delta^2\Lambda^2\norm E^2+L\delta^2\norm e^2
                                           +\frac L2\tr Q.
\]
Combine this with the proof of \cref{v2:prop:work} and absorb
$L\delta^2\Lambda^2\norm E^2$ into the negative drift. The step restriction
leaves $-\lambda\delta\norm E^2/4$.
\end{proof}

\paragraph{Why a covariance panel alone is insufficient.}
The quadratic formula can be evaluated from a complete conditional first and
second moment panel once the \emph{same action Hessian} is identified. For a
nonquadratic action the displayed third-moment error, or the exact action
increment itself, must also be retained. An unrelated positive Hodge energy
cannot replace the signed $H(x)$ in \eqref{v2:eq:work-moments}. Negative action
work may assist stationarity selection but does not by itself control the
positive physical reader energy or prevent chart exit.

\subsection{An action-neutral part of the actual transition}

\begin{definition}[Action-neutral transition component]
\label[definition]{v2:def:neutral}
A normalized finite kernel $N_x(z,dy)$ is action-neutral for the independently
fixed action if
\begin{equation}
 \int \mathcal A(y)N_x(z,dy)=\mathcal A(z).
 \label{v2:eq:neutral}
\end{equation}
Pathwise invariance, $\mathcal A(y)=\mathcal A(z)$ on its support, is a
sufficient stronger property. Neither condition says that $y=z$, that the
physical reader is constant, or that the conditional variance vanishes.
\end{definition}

\begin{theorem}[Retained action-neutral noise does not enter the drift budget]
\label[theorem]{v2:thm:neutral}
Suppose the \emph{actual} legal transition is a verified composition
\begin{equation}
 K(x,dy)=\int D(x,dz)N_x(z,dy),
 \label{v2:eq:factorization}
\end{equation}
where both kernels are normalized on the same chart. Suppose $c>0$ and $\epsilon_D\ge0$ satisfy
\begin{align}
 \int[\mathcal A(z)-\mathcal A(x)]D(x,dz)
       &\le-c\delta\norm{E(x)}^2+\delta\epsilon_D(x),\\
 \int\mathcal A(y)N_x(z,dy)&\le\mathcal A(z)+\zeta(x,z),
 \qquad \zeta\ge0.
 \label{v2:eq:component-drifts}
\end{align}
Then \eqref{v2:eq:action-drift} holds for the full transition with
\begin{equation}
 \kappa_A=c,\qquad
 r(x)=\epsilon_D(x)+\delta^{-1}\int\zeta(x,z)D(x,dz).
 \label{v2:eq:neutral-remainder}
\end{equation}
For an exactly action-neutral second component, its contribution is zero
regardless of its innovation variance. Every outcome $y$, not $z$ or a
conditional mean, is the record used by \cref{v2:thm:joint}.
\end{theorem}
\begin{proof}
Use \eqref{v2:eq:factorization} to write the full expected action increment as
\[
 \int D(x,dz)\left\{\mathcal A(z)-\mathcal A(x)
             +\int[\mathcal A(y)-\mathcal A(z)]N_x(z,dy)\right\}.
\]
The two inequalities in \eqref{v2:eq:component-drifts} give the claimed drift.
There is no estimate of $\norm{y-z}$ in this computation. The selector still
checks $V_h(y)$ and $W_h(y)$, so action-neutral noise that ruins regularity or
alignment is not declared harmless to the continuum theorem.
\end{proof}

\begin{corollary}[A directly verifiable deterministic first component]
\label[corollary]{v2:cor:descent-component}
If $D(x,dz)$ is the point mass at $z=x-\delta E(x)$, the connecting segment
lies in the chart, $\norm{D^2\mathcal A}\le L$, and $\delta L\le1$, then
\cref{v2:thm:neutral} applies with $c=1/2$ and $\epsilon_D=0$.
\end{corollary}
\begin{proof}
Taylor's integral identity gives
$\mathcal A(x-\delta E)-\mathcal A(x)
\le-\delta\norm E^2+(L\delta^2/2)\norm E^2
\le-\delta\norm E^2/2$.
\end{proof}

\begin{remark}[No algorithm is silently added to a native source]
The factorization in \eqref{v2:eq:factorization} must be the source's observed
composition or an independently justified identity for its kernel. It is not
obtained by declaring an arbitrary accepted step to be gradient descent.
Likewise a finite exact symmetry of the common action gives an action-neutral
kernel only when the given source actually applies that symmetry. The full
metric, temporal-gauge, Higgs, and fermionic derivatives remain in $E_h$.
There is no projection onto symmetry tangents in any of these statements.
Exact lattice covariance may verify \eqref{v2:eq:neutral} for an identified
symmetry layer; it does not establish \eqref{v2:eq:factorization}, energy
preservation, or initial/gauge alignment for an unspecified source.
\end{remark}

\section{An exact distinction between small action gradients and small Bregman balance}
\label{v2:sec:stationary-identity}

\begin{theorem}[Stationary coupling identity for a constant source metric]
\label[theorem]{v2:thm:stationary-identity}
Let $\mathsf H$ be a finite real Hilbert space, $\mathcal A\in C^2(\Theta)$
on a convex chart, and $\norm{D^2\mathcal A}\le L$. Let $X,Y$ be any coupling
with the same marginal distribution in $\Theta$, with the following
expectations finite. Take $p(x)=x$, $0<\upsilon L<1$, and put
\[
 b=E(Y)+\upsilon^{-1}(Y-X),\qquad
 f(x)=\tfrac12\norm x^2+\upsilon\mathcal A(x).
\]
For the Bregman divergence
$D_f(x,y)=f(x)-f(y)-\ip{\nabla f(y)}{x-y}$,
\begin{equation}
 \boxed{\quad
 \E\norm b^2=\E\norm{E(Y)}^2+
                       \frac{2}{\upsilon^2}\E D_f(X,Y).
 \quad}
 \label{v2:eq:stationary-identity}
\end{equation}
Consequently
\begin{equation}
 \E\norm b^2\ge\E\norm{E(Y)}^2+
          \frac{1-\upsilon L}{\upsilon^2}\E\norm{Y-X}^2.
 \label{v2:eq:stationary-lower}
\end{equation}
Reversibility, independence, and a Markov assumption are unnecessary.
\end{theorem}
\begin{proof}
Since $\upsilon b=\nabla f(Y)-X$ and
$\upsilon E(Y)=\nabla f(Y)-Y$, subtracting the two squared norms gives
\[
 \upsilon^2\bigl(\norm b^2-\norm{E(Y)}^2\bigr)
 =\norm X^2-\norm Y^2-2\ip{\nabla f(Y)}{X-Y}.
\]
Equal marginal laws cancel the expectation of the first two terms, and also
$\E[f(X)-f(Y)]$. The remaining expectation is $2\E D_f(X,Y)$, proving
\eqref{v2:eq:stationary-identity}. The Hessian of $f$ is bounded below by
$(1-\upsilon L)I$, hence its integral Taylor formula gives
$D_f(X,Y)\ge(1-\upsilon L)\norm{X-Y}^2/2$.
\end{proof}

\begin{remark}[Interpretation and coordinate boundary]
The squared balance certificate controls not only stationarity of the action
but also motion of the records. A process supported on a nontrivial critical
orbit can have $E=0$ and a strictly positive second term in
\eqref{v2:eq:stationary-identity}. This is a precise reason not to treat the
innovation floor in Version 1 as a universal requirement for classical
stationarity. The result is stated for a constant physical source metric;
a fixed linear change of coordinates gives the corresponding constant-metric
formula with all gradients and norms transformed consistently. A variable
Fisher matrix is not silently replaced by such a constant metric.
\end{remark}

\section{A fully finite pilot separating the two stationarity entrances}
\label{v2:sec:pilot}

This section is an explicit diagnostic. It is not the native Einstein--SM
action, does not construct a spacetime metric, and is not offered as a solution
to the microscopic identification problem. Its role is to prove that the new
entrance is genuinely less restrictive with respect to some retained physical
noise; the distinction is not just a renamed hypothesis.

Let $e_1=(1,0)$, $0<a_0\le1/4$, and $0<h\le1/8$. On the convex chart
$\Theta=\overline B(e_1,1/2)\subset\R^2$ define
\begin{equation}
 \mathcal A(x)=\tfrac12(\norm x-1)^2,
 \qquad E(x)=(\norm x-1)\frac{x}{\norm x},
 \qquad \delta_h=h^2,\quad\omega_h=h^2.
 \label{v2:eq:pilot-action}
\end{equation}
The chart avoids the origin. At every active step set
\begin{equation}
 z=x-\delta_h E(x),\qquad
 Y=\norm z\,(\cos\omega_h,\xi\sin\omega_h),
 \qquad \PP(\xi=1\mid\mathcal F_j)=\PP(\xi=-1\mid\mathcal F_j)=1/2,
 \label{v2:eq:pilot-rule}
\end{equation}
starting at $X_0=(1+a_0)e_1$. For a fixed finite horizon $J_h$ the reachable
states, including the stage register, form a finite table with at most
$1+2J_h$ states. Thus the rule defines a finite process even though its action
has a smooth off-table extension. That extension is fixed before execution.

\begin{theorem}[A retained-noise finite population with small gradients but non-small balance]
\label[theorem]{v2:thm:pilot}
The source \eqref{v2:eq:pilot-rule} has the following properties.
\begin{enumerate}[label=\textup{(\roman*)},leftmargin=2.8em]
\item Every reachable state stays in $\Theta$, and
\begin{equation}
 r_j:=\norm{X_j}=1+a_0(1-h^2)^j,\qquad
 g_j:=\norm{E(X_j)}=a_0(1-h^2)^j.
 \label{v2:eq:pilot-radius}
\end{equation}
The actual action drift obeys \eqref{v2:eq:action-drift} with
$\kappa_A=1/2$, $\delta_h=h^2$, and $r_{h,j}=0$.
\item Its conditional physical innovation satisfies, uniformly in $j$,
\begin{equation}
 \nu_j=\E[\norm{X_{j+1}-\E[X_{j+1}\mid\mathcal F_j]}^2
                    \mid\mathcal F_j]
       =r_{j+1}^2\sin^2(h^2)\asymp h^4.
 \label{v2:eq:pilot-noise}
\end{equation}
With $p(x)=x$ and $\upsilon_h=h^2$, its actual squared Bregman balance has
\begin{equation}
 c\le\E[\norm{b_j}^2\mid\mathcal F_j]\le C,
 \qquad
 \frac1J\sum_{j<J}\E\norm{b_j}^2\asymp1.
 \label{v2:eq:pilot-balance}
\end{equation}
\item For any fixed $\epsilon>0$ and
$J_h=\lceil h^{-(4+2\epsilon)}\rceil$, every realization satisfies
\begin{align}
 \frac1{J_h}\sum_{j<J_h}g_j^2&\le C h^{2+2\epsilon},\\
 \frac1{J_h}\sum_{j<J_h}\norm{X_j-e_1}^2
                           &\le C(h^{2+2\epsilon}+h^4).
 \label{v2:eq:pilot-occupation}
\end{align}
Its final record has
$\norm{X_{J_h}-e_1}\le a_0\exp(-h^2J_h)+C h^2=O(h^2)$.
\item The reader $u_j(t,x)=X_j$, constant on an auxiliary spacetime torus,
has bounded $H^q$ norm for every fixed $q$ and zero tails for $K\ge1$.
With $V_j=1+\norm{X_j}^2$, $W_j=\norm{X_j-e_1}^2$, thresholds
$s_h=\alpha_h=h$ and $R_h=h^{-a}$, $a>0$, the joint selector succeeds for
every realization once $h$ is sufficiently small. The alignment quantity in
this statement is the pilot distance $W_j$, not an Einstein gauge residual.
\end{enumerate}
\end{theorem}
\begin{proof}
On $\Theta$ the radial Hessian eigenvalue of $\mathcal A$ is $1$ and the
tangential eigenvalue is $1-1/\norm x$, lying in $[-1,1]$. Thus the gradient
is $1$-Lipschitz. The deterministic step changes radius $r$ to
$1+(1-h^2)(r-1)$. The random step preserves that radius exactly, proving
\eqref{v2:eq:pilot-radius}. Reachable points satisfy
\[
 \norm{X_j-e_1}\le(r_j-1)+r_j\norm{(\cos\omega_h,\pm\sin\omega_h)-e_1}
 \le a_0+(1+a_0)h^2<1/2.
\]
The deterministic connecting segment is in $\Theta$, and the action change is
\[
 \mathcal A(X_{j+1})-\mathcal A(X_j)
 =-h^2(1-h^2/2)g_j^2\le-\tfrac12h^2g_j^2.
\]
This proves (i), including drift for the \emph{actual} two-outcome rule.

The two outputs have common first coordinate and opposite second coordinates,
so their conditional variance is \eqref{v2:eq:pilot-noise}. Since
$1\le r_j\le5/4$ and $\sin(h^2)\asymp h^2$, this is uniformly of order $h^4$.
The lower bound in \eqref{v2:eq:pilot-balance} follows from
\cref{v1:thm:innovation} with $L=g_*=g^*=1$:
$(h^{-2}-1)^2\nu_j\ge c$. For the upper bound,
$\norm{E(X_{j+1})}\le1/4$ and
\[
 \norm{X_{j+1}-X_j}\le h^2 a_0+2(1+a_0)\sin(h^2)\le C h^2
\]
for $j\ge1$; the same bound holds at $j=0$ with one rather than two angular
increments. Hence $\norm{b_j}\le C$ on every outcome. Averaging proves (ii).

The geometric sum gives
\[
 \sum_{j<J}g_j^2\le\frac{a_0^2}{1-(1-h^2)^2}\le C h^{-2}.
\]
Moreover
$\norm{X_j-e_1}^2=(r_j-1)^2+2r_j(1-\cos h^2)$ for $j\ge1$,
with the second term absent at $j=0$. This is at most $g_j^2+C h^4$.
Division by $J_h$ proves \eqref{v2:eq:pilot-occupation}.
The elementary bound $(1-h^2)^J\le e^{-h^2J}$ proves the final-record rate.
For (iv), only the zero Fourier coefficient is nonzero. The occupation score
is bounded on every realization by
\[
 C h^{2\epsilon}+C h^2+C h^{2a},
\]
which is less than one for small $h$. Apply \cref{v1:lem:selector}.
\end{proof}

\begin{remark}[What the pilot does and does not establish]
The innovation scale $h^4$ is larger than the $h^6$ scale required by the
Version 1 $O(h^2)$ squared-balance entrance; indeed
\eqref{v2:eq:pilot-balance} proves failure of that entrance. Nevertheless
small gradients and aligned, smooth records are generated by an explicit
finite law. No record has been replaced by its mean or by a projected field.
The improvement comes from exact action neutrality, not from denying the
noise. The action is a radial diagnostic potential, not the full Lorentzian
Einstein--SM functional. Therefore the pilot proves strict separation of
sufficient mechanisms, not completion of the native Einstein--SM program.
Its long stage horizon is an operation count, not an identified physical time.
\end{remark}

\section{Testing the alternative entrance on an explicit upstream primitive}
\label{v2:sec:primitive}

The Grand Tensor monograph supplies a definite source rather than an abstract
regenerative condition in \texttt{thm:renewal-interchange-actual-audit}. With
$n=N^3$ phase sites, the phase-count quotient $k\in\{0,\ldots,n\}$ has the
stated generator
\begin{equation}
 (\mathcal L_n f)(k)
  =a(n-k)[f(k+1)-f(k)]+bk[f(k-1)-f(k)],
 \quad a=\frac{4\lambda}{5},\quad b=\frac{2\lambda}{3}.
 \label{v2:eq:count-generator}
\end{equation}
Its binomial invariant law has parameter
$p=a/(a+b)=6/11$. These are inherited source data, not a newly designed law.
Swaps, including the stated modulation, leave the count unchanged.
We calculate the action-drift and innovation consequences for this precise
quotient. This calculation does not identify a count with a Lorentzian metric.

\begin{proposition}[The source's bare phase action cannot supply a vanishing slope-drift error]
\label[proposition]{v2:prop:count-action}
In the explicitly declared density coordinate $r=k/n$, with Euclidean scalar
metric, the action per site inherited from the product phase law is
\begin{equation}
 \bar{\mathcal A}(r)=-(1-r)\log(5/11)-r\log(6/11),\qquad
 \bar E(r)=-\ell,\quad \ell=\log(6/5)>0.
 \label{v2:eq:count-action}
\end{equation}
Let $P_d=e^{d\mathcal L_n}$, $d>0$, and $\lambda_0=a+b$. Then
\begin{equation}
 (P_d\bar{\mathcal A}-\bar{\mathcal A})(r)
    =\ell(1-e^{-\lambda_0d})(r-p).
 \label{v2:eq:count-drift}
\end{equation}
If the actual sampled quotient obeys an action-drift inequality of the form
\eqref{v2:eq:action-drift}, with a fixed $\kappa_A>0$, then its invariant
occupation necessarily satisfies
\begin{equation}
 \E_{\pi_n} r_{h,j}\ge\kappa_A\ell^2.
 \label{v2:eq:count-obstruction}
\end{equation}
Thus its own bare phase action does not furnish a vanishing stationarity-drift
error in this fixed physical scalar normalization, regardless of $n$.
\end{proposition}
\begin{proof}
From \eqref{v2:eq:count-generator},
$\mathcal L_n r=a-(a+b)r=-\lambda_0(r-p)$.
The finite matrix exponential therefore gives
$P_dr=p+e^{-\lambda_0d}(r-p)$. Apply this to the affine function
\eqref{v2:eq:count-action} to obtain \eqref{v2:eq:count-drift}.
Under its invariant law, the expectation of the action increment is zero.
Averaging \eqref{v2:eq:action-drift}, with $g^2=\ell^2$, gives
$0\le-\kappa_A\delta_h\ell^2+\delta_h\E_{\pi_n}r_{h,j}$.
Division by $\delta_h>0$ proves \eqref{v2:eq:count-obstruction}.
\end{proof}

\begin{proposition}[Exact innovation trace of the sampled count reader]
\label[proposition]{v2:prop:count-noise}
At stationarity the same finite source satisfies
\begin{equation}
 \E_{\pi_n}\Var(r_{j+1}\mid r_j)
   =\frac{p(1-p)}{n}(1-e^{-2\lambda_0d}).
 \label{v2:eq:count-noise}
\end{equation}
In the illustrative identification $n=h^{-3}$ and sampling interval $d=h^2$,
this is asymptotic to $2\lambda_0p(1-p)h^5$. For the test potential $\Psi(r)=r^2/2$ and $\upsilon_h\asymp h^2$,
the same scalar prefix-balance test has an innovation contribution of order
$h$, not $O(h^2)$. This is a diagnostic choice of metric and potential, not
an identification with the coupled physical source potential.
\end{proposition}
\begin{proof}
The binomial law gives $\Var_{\pi_n}(r)=p(1-p)/n$.
The conditional mean calculated in the preceding proof has variance
$e^{-2\lambda_0d}p(1-p)/n$. The total-variance identity subtracts that quantity
from the stationary variance, giving \eqref{v2:eq:count-noise}. Expand the
exponential for the stated small interval. Here $\bar E$ is constant, so the conditional variance of
$\bar E(r_{j+1})+\upsilon_h^{-1}(r_{j+1}-r_j)$ is exactly
$\upsilon_h^{-2}\Var(r_{j+1}\mid r_j)$. Its stationary mean is therefore
of the asserted order. The normalization and the observable have not been
changed during this calculation.
\end{proof}

\paragraph{Consequences of this actual-source calculation.}
The monograph already distinguishes this source's free-energy concentration
and stochastic reaction--diffusion limit from an independently fixed Einstein
action. Equations \eqref{v2:eq:count-obstruction} and
\eqref{v2:eq:count-noise} quantify that distinction for the two stationarity
entrances of this lineage. They do not prove failure of an unobserved coupled
Einstein--SM source, and the per-site scalar gradient is not its full physical
Euler covector. They show why the explicit interchange primitive cannot be
silently used to discharge either entrance for a different action. Changing
the reference entropy, action normalization, reader, or transition law would
be a separate source modification, not a proof about this one.

\section{A finite action-drift readout and the remaining native calculation}
\label{v2:sec:ledger}

\begin{proposition}[Robust finite-row certification]
\label[proposition]{v2:prop:row-audit}
Suppose on an actual row there are certified enclosures
\begin{equation}
 |\widehat{\mathfrak d}(x)-\mathfrak d(x)|\le e_d(x),\qquad
 |\widehat g(x)^2-g(x)^2|\le e_g(x),
 \label{v2:eq:audit-errors}
\end{equation}
and a nonnegative certified $\widehat r(x)$ with a verified inequality
$\widehat{\mathfrak d}(x)\le-\kappa_A\delta_h\widehat g(x)^2+
\delta_h\widehat r(x)$.
Then the true drift obeys \eqref{v2:eq:action-drift} with
\begin{equation}
 r(x)=\widehat r(x)+e_d(x)/\delta_h+\kappa_A e_g(x).
 \label{v2:eq:audit-r}
\end{equation}
For a finite known successor set, if the independently evaluated action
increments satisfy $|a(y)-a(x)|\le M_A(x)$ and the table approximation has
$\sum_y|\widehat K^{\rm leg}(x,y)-K^{\rm leg}(x,y)|\le e_K(x)$, the table
contribution to $e_d(x)$ is at most $M_A(x)e_K(x)$.
\end{proposition}
\begin{proof}
Insert the two enclosures in the verified inequality, using
$\widehat g^2\ge g^2-e_g$. The table bound follows by multiplying each
probability error by its actual action increment and summing absolute values.
No unrecorded successor or postselected renormalization is covered by this
last bound.
\end{proof}

\paragraph{What is proved in this iteration.}
The same-source composition can now be entered through either of two distinct
stationarity mechanisms:
\[
 \begin{array}{c}
 \text{small realized full Bregman balance}\\
 \text{or}\\
 \text{a verified conditional drift of the same signed action}.
 \end{array}
\]
The second mechanism does not require a Bregman Hessian or a small full
innovation trace. Its new source obligation is the action-increment estimate
\eqref{v2:eq:action-drift}, derived rather than merely renamed on the
factorized action-neutral class. The finite pilot proves that this class can
pass when the Version 1 balance entrance fails. The explicit inherited phase
primitive, on the other hand, fails the vanishing-error action-drift test for
its own bare per-site action. Thus the argument has not manufactured an
Einstein action from a generic stationary law.

\paragraph{What remains independent.}
The physical reader-to-tail bound, its occupation energy, initial and harmonic
alignment, legal-record existence, and exact common-action identification are
unchanged obligations. A noise layer neutral for the action may still change
all of those quantities. In particular, the positive-Fisher saddle obstruction
already proved in \cite{CC} remains in force: the alternative entrance is not
a proof that descent on the unrestricted Lorentzian history action has a
regular attracting neighborhood. No general primitive-source attraction or
quantum continuum statement has been established.

\paragraph{The next nonduplicating source target.}
For the independently supplied coupled source, the relevant observation is
now the joint record of the initial full field, the resolved transition,
the terminal full field, and the \emph{same-action increment}. First compute
\eqref{v2:eq:action-increment}, or verify a literal action-neutral composition
and its complementary action drift. If only an outgoing covariance is
available, \eqref{v2:eq:work-moments} states exactly which action-Hessian and
higher-moment information is additionally needed. The noise-short Ward panel
alone does not supply that incidence. Only after this action-level test passes
should its occupation control be combined with the physical reader-energy
and alignment estimates of Version 1. In the absence of that source row,
another generic compactness theorem or a newly designed finite sampler would
not identify the intended native law.

\paragraph{Verification.}
The accompanying diagnostic script checks the stopped-action telescope with
retained exits, the stationary coupling identity for an indefinite quadratic
action and a nonreversible equal-marginal coupling, the exact signed work
formula, the count-reader variance, and the radial pilot's drift, variance,
and incompatible balance scaling. These are finite checks of the displayed
proofs, not a numerical verification of the coupled Einstein--SM system.
The complete Version 1 source text before its final end-document command is
preserved verbatim. All added labels have the prefix \texttt{v2:}.

\begingroup
\renewcommand{\refname}{Additional references for Version 2}
\begin{thebibliography}{9}
\raggedright
\bibitem{v2:KMS}
I.~Karatzas, J.~Maas, and W.~Schachermayer,
\emph{Trajectorial dissipation and gradient flow for the relative entropy in
Markov chains}, 2020, revised 2022.
\href{https://arxiv.org/abs/2005.14177}{arXiv:2005.14177}.
Background for distinguishing Markov-law entropy dissipation from a
fieldwise variational equation; no theorem from it is an unstated hypothesis.
\bibitem{v2:DD}
D.~Davis and D.~Drusvyatskiy,
\emph{Proximal methods avoid active strict saddles of weakly convex functions},
2019, revised 2021.
\href{https://arxiv.org/abs/1912.07146}{arXiv:1912.07146}.
Background already recognized in the supplied closure notes; the saddle
obstruction is not rederived as a new result here.
\end{thebibliography}
\endgroup
% END IMMUTABLE VERSION 2 BODY
% APPEND VERSION 3 HERE, BEFORE THE FINAL END-DOCUMENT COMMAND.


% BEGIN IMMUTABLE VERSION 3 BODY
\clearpage
\begin{center}
{\Large\bfseries Version 3: Action-resolved moments and the physical reader}\\[0.5em]
{\large Exact Higgs work, transient all-mode leakage, and a selection-cost obstruction}
\end{center}
\addcontentsline{toc}{part}{Version 3: Action moments and native reader obstruction}

\section{Upstream audit: evaluate the source before adding a closure criterion}
\label{v3:sec:scope}

The alternative stationarity entrance in Version 2 is useful only after its
signed action drift has been checked on the actual transition.  In addition,
action-neutral noise need not be neutral for the physical reader energy.
This iteration calculates these two issues rather than adding another generic
PDE or compactness implication.

\paragraph{Results already present, and not claimed again.}
The supplied Grand Tensor monograph \cite{GT}, at
\nolinkurl{thm:GT-target-sandwich} and
\nolinkurl{thm:GT-one-history-projection}, already gives sharp positive-extension
intervals for a scalar writer and the obstruction to assembling an unmeasured
same-history coupling from its marginals.  We therefore do not reprove a
general moment-problem, linear-programming, or source-identifiability theorem.
Its \nolinkurl{thm:renewal-interchange-actual-audit} already computes the phase
count, the product action, and the reversible law of the explicit
renewal--interchange regulator.  Its
\nolinkurl{thm:renewal-interchange-first-chaos},
\nolinkurl{thm:renewal-interchange-OU}, and
\nolinkurl{thm:renewal-fixed-coupling-OU} already address reaction--diffusion and
fluctuation limits.  Those limits, and Version 2's bare-action gradient and
count-innovation calculations, are not rederived here.

The same monograph explicitly leaves the accepted initial determining field,
resolved update, terminal field, and same-action increment of the \emph{locked
spacetime source} as an unexported joint row.  That source cannot be silently
identified with the separately specified interchange regulator.  We retain
this distinction.  No numerical values for the missing coupled row are
invented in this iteration.

\paragraph{Specific additions.}
For the actual minimal finite Higgs action with the other fields fixed, we
identify an exact compressed conditional-moment bank for its action increment.
We then give two positive finite successor laws with identical moments through
order three but opposite drift of that same Higgs action.  This is an
obstruction in an Einstein--SM action sector, not another radial diagnostic
action.  For the explicit upstream phase generator, we prove its exact
nonstationary product entrance, quantify the all-mode spatial leakage, and
show that selecting a regular record after a fixed positive physical time
requires exponentially many observations in a precise sampling class.
An actual-rate entrance bound strengthens this to exponentially costly
physical waiting even with adaptive trajectory observations.  Finally, an exact reader-spectrum formula and a full-variation-bank conorm
specify what must actually be measured to use the Version 2 entrance.

\paragraph{Attribution and logical strength.}
Polynomial expansion, product-measure forward equations, Fourier Parseval,
and bounded-variable concentration are standard tools.  The concentration
bound used below is proved in its required finite form; see
\cite{v3:Hoeffding} for its classical origin.  General interacting-particle
hydrodynamics also distinguishes density and fluctuation readers, for example
\cite{v3:Funaki}; no theorem from that literature is used as a substitute for
the direct calculations below.  The contribution here is the source-specific
moment requirement, the full-record selection obstruction, and the exact
physical normalization of their interface with Versions 1--2.  The terminal
Einstein--SM implication remains the inherited input
\cref{v1:in:closure}, not a newly established native microscopic law.

\section{An exact action-moment bank for the minimal Higgs sector}
\label{v3:sec:Higgs}

Fix a finite coframe, gauge links, spinor and dual-spinor values, and the
calibrated coefficients of the minimal local action.  Realify the Higgs
variables, and write their space as
\[
 \mathcal H_h=\bigoplus_{x\in\Lambda_h}\R^d,
 \qquad (u,v)_h=h^4\sum_x u_x\cdot v_x.
\]
All gradients and covariance operators below use this physical mass inner
product.  Let $w_x>0$ be the fixed cell-volume density including $h^4$.
The restriction of the full action to the Higgs variables has the exact form
\begin{equation}
 \mathcal A_h(H)=C_h+\tfrac12(H,B_hH)_h+(J_h,H)_h
       -\lambda\sum_x w_x( |H_x|^2-v^2)^2,
 \qquad B_h=B_h^*.
 \label{v3:eq:Higgs-action}
\end{equation}
The signed covariant kinetic form gives $B_h$, and the fixed fermionic
Yukawa term is affine in the real Higgs variables and gives $J_h$.  All other
sectors are constant during this restricted update.  In particular $B_h$ is
not assumed positive.  The representation follows because fixed gauge-link
covariant differences are linear in $H$ and the minimal Higgs potential is
quartic.  It is valid for this fixed-field restriction of the supplied local
action, not for arbitrary simultaneous metric, gauge, and Higgs changes.

At a current physical field $H$, let $D$ be its actual decoded Higgs increment
on an accepted row of this type.  Illegal successors are represented by
$D=0$ solely for the stopped-increment calculation of
\cref{v2:def:action-increment}; their probability is still charged to the
exit event.  A primitive coordinate increment may be substituted for $D$
only if its physical Higgs decoder is the stated affine map.  Every
conditional expectation in this section is taken on the complete observed
current source record and its resolved Higgs-only label; the notation
$\E[\,\cdot\mid H]$ suppresses the other fixed coordinates and that label.
It does not assert that the Higgs field alone is a Markov-sufficient record.
Define
\begin{align}
 E_H&=\nabla_h\mathcal A_h(H),& L_H&=D E_H(H),\\
 m&=\E[D\mid H],&
 C&=\E[(D-m)\otimes_h(D-m)\mid H],\\
 c_3&=\E\left[\sum_xw_x(H_x\cdot D_x)|D_x|^2\,\middle|\,H\right],&
 c_4&=\E\left[\sum_xw_x|D_x|^4\,\middle|\,H\right].
 \label{v3:eq:moment-bank}
\end{align}
Here $(u\otimes_h u)v=(u,v)_h u$, so $C$ is the covariance operator rather
than an unweighted covariance matrix in unnormalized nodal coordinates.

\begin{theorem}[Exact conditional work for a Higgs-only transition]
\label[theorem]{v3:thm:Higgs-work}
For every finite conditional successor law on the stated chart,
\begin{equation}
 \boxed{
 \begin{aligned}
 \mathfrak d_H(H)
 &:=\E[\mathcal A_h(H+D)-\mathcal A_h(H)\mid H]\\
 &= (E_H,m)_h+\tfrac12(m,L_Hm)_h
       +\tfrac12\tr(L_HC)-4\lambda c_3-\lambda c_4.
 \end{aligned}}
 \label{v3:eq:Higgs-work}
\end{equation}
There is no Taylor remainder.  The action drift is therefore determined by
the decoded mean and second moment together with the two displayed
third- and fourth-order scalar contractions.  The latter are local
same-event Higgs contractions; a full fourth-order covariance tensor is
unnecessary.
\end{theorem}
\begin{proof}
For real vectors $H_x,D_x$, direct multiplication gives
\[
 \begin{split}
 (|H_x+D_x|^2-v^2)^2={}&(|H_x|^2-v^2)^2
 +4(|H_x|^2-v^2)H_x\cdot D_x\\
 &+2(|H_x|^2-v^2)|D_x|^2+4(H_x\cdot D_x)^2\\
 &+4(H_x\cdot D_x)|D_x|^2+|D_x|^4.
 \end{split}
\]
The linear and quadratic terms in this identity, together with the kinetic
and Yukawa variations, are exactly $(E_H,D)_h+\tfrac12(D,L_HD)_h$.
The two remaining terms are the cubic and quartic terms in
\eqref{v3:eq:Higgs-work}.  Finally,
\[
 \E[(D,L_HD)_h\mid H]=(m,L_Hm)_h+\tr(L_HC),
\]
by expansion of $D=m+(D-m)$ and the defining covariance identity in an
orthonormal physical-mass basis.  Taking conditional expectations proves the
formula, including any atom at $D=0$.
\end{proof}

\begin{corollary}[A finite scalar audit with explicit uncertainty]
\label[corollary]{v3:cor:Higgs-audit}
Put $q_1=(E_H,m)_h$ and
$q_2=(m,L_Hm)_h+\tr(L_HC)$.  Suppose the same-event source supplies estimates
$\widehat q_1,\widehat q_2,\widehat c_3,\widehat c_4$ with respective
certified absolute errors $e_1,e_2,e_3,e_4$.  Then
\begin{equation}
 \left|\mathfrak d_H-
 (\widehat q_1+\tfrac12\widehat q_2
                  -4\lambda\widehat c_3-\lambda\widehat c_4)\right|
 \le e_1+\tfrac12e_2+4\lambda e_3+\lambda e_4.
 \label{v3:eq:moment-uncertainty}
\end{equation}
Thus a Higgs-sector drift inequality can be certified with four scalar
conditional action contractions.  This does not certify the full physical
gradient unless the other variation directions are also controlled.
\end{corollary}
\begin{proof}
Subtract the two exact linear combinations and use the triangle inequality.
The last sentence follows because $E_H$ is only the Higgs component of the
common-action gradient; it is not the independent metric or gauge row.
\end{proof}

\begin{remark}[What has and has not been reduced]
The Hessian contraction in $q_2$ uses the \emph{same signed action} and its
covariant kinetic incidence.  An unrelated positive outgoing Gram cannot
replace it.  A directly read action increment can replace all four
contractions, consistently with the already proved target-sandwich theorem
in \cite{GT}.  The present result supplies the exact action-specific
contractions when the source exports moments rather than action increments.
Neither choice reconstructs an unprovided successor law.
\end{remark}

\section{Identical source moments can give opposite Higgs-action drift}
\label{v3:sec:moment-obstruction}

Consider the constant real Higgs direction on a flat, unit-volume finite
periodic evaluation box, with identity gauge transport and zero spinors.
Keep the coframe and every other field fixed.  Covariant differences vanish
on this direction.  With $\lambda=v=1$ its \emph{finite} common-action
restriction, up to a fixed additive gravitational term, is
\begin{equation}
 a(u)=-(u^2-1)^2.
 \label{v3:eq:literal-quartic}
\end{equation}
This is an off-shell restriction of the declared minimal action; no Einstein
solution or special scalar field equation is used.

\begin{proposition}[Moment-matched rows with opposite same-action drift]
\label[proposition]{v3:prop:quartic-witness}
At $u=1/2$ there are two finite successor laws with identical moments through
order three for the complete scalar increment, but with expected action
increments of opposite signs.  They can be chosen to have strictly positive
probability on the same five-point successor alphabet.  Consequently the
incoming/outgoing first- and second-moment Gram, even supplemented with all
third moments of this scalar row, does not determine the sign of the action
drift required by \cref{v2:thm:occupation}.
\end{proposition}
\begin{proof}
At $u=1/2$, exact expansion yields
\begin{equation}
 a(1/2+D)-a(1/2)=\tfrac32D+\tfrac12D^2-2D^3-D^4.
 \label{v3:eq:quartic-expansion}
\end{equation}
First take the increment law $P$ giving $D=\pm1/2$ probability $1/2$ each,
and the law $Q$ giving $D=\pm1$ probability $1/8$ each and $D=0$
probability $3/4$.  Both have
\[
 \E D=\E D^3=0,\qquad \E D^2=\tfrac14.
\]
Their fourth moments are $1/16$ and $1/4$.  Equation
\eqref{v3:eq:quartic-expansion} gives
\begin{equation}
 \mathfrak d_P=\tfrac1{16}>0,
 \qquad \mathfrak d_Q=-\tfrac18<0.
 \label{v3:eq:opposite-drifts}
\end{equation}
For a common strictly positive alphabet, let $U$ be uniform on
$\{-1,-1/2,0,1/2,1\}$ and set
$P^+=(7/8)P+(1/8)U$, $Q^+=(7/8)Q+(1/8)U$.
Both laws are positive at all five points, have the same first three
moments, and give
\[
 \mathfrak d_{P^+}=\frac{21}{640}>0,
 \qquad \mathfrak d_{Q^+}=-\frac{21}{160}<0,
\]
since $\mathfrak d_U=-7/40$.  The incoming field is the same deterministic
$1/2$, so every joint incoming/outgoing moment of total degree at most three
also agrees.  Appending unchanged field coordinates does not change this
conclusion.  The current scalar action gradient is $a'(1/2)=3/2$ for both
rows; neither row is being made stationary by a different normalization.
\end{proof}

\begin{theorem}[Sharp scalar work interval from a variance and a jump cap]
\label[theorem]{v3:thm:sharp-quartic-interval}
Let an action restriction at a fixed current point be
\[
 a(D)-a(0)=a_1D+a_2D^2+a_3D^3-\lambda D^4,
 \qquad \lambda>0.
\]
Among symmetric conditional increment laws with $|D|\le R$ and
$\E D^2=s_2\in[0,R^2]$, the exact possible action-drift interval is
\begin{equation}
 \boxed{
 a_2s_2-\lambda R^2s_2
 \ \le\ \mathfrak d\ \le\
 a_2s_2-\lambda s_2^2.}
 \label{v3:eq:sharp-work-interval}
\end{equation}
Both endpoints and every intermediate value are attained by finite laws.
For the data of \cref{v3:prop:quartic-witness}, with $R=1$ and $s_2=1/4$,
this interval is exactly $[-1/8,1/16]$.
\end{theorem}
\begin{proof}
Symmetry removes the odd terms.  Jensen's inequality gives
$\E D^4\ge (\E D^2)^2=s_2^2$, and the pointwise inequality
$D^4\le R^2D^2$ gives the other bound.  The symmetric two-point law
$D=\pm\sqrt{s_2}$ realizes the lower fourth moment.  The law with masses
$s_2/(2R^2)$ at each of $\pm R$ and the remaining mass at zero realizes the
upper fourth moment when $R>0$.  Mixtures of these laws preserve the variance
and interpolate every fourth moment between the endpoints.  The zero-radius
case is immediate.  Multiplication by $-\lambda$ reverses the fourth-moment
order and proves the stated work interval.
\end{proof}

\paragraph{Interpretation.}
This calculation identifies an actual missing row, rather than inferring
failure of an unobserved coupled source.  A covariance-only source audit
cannot even decide the sign of the minimal Higgs work on the displayed
regular finite chart.  A measured fourth contraction or action increment
removes this particular ambiguity.  Exact action neutrality of a different
noise mechanism and small noise variance are neither assumed nor required
for this conclusion.

\section{Transient spatial leakage of the explicit upstream phase source}
\label{v3:sec:phase}

We now keep the literal phase rates from \cite{GT}, rather than design an
Einstein--SM generator.  Write $\eta_x=\one_{\{P_x\}}$ on
$V_N=(\Z/N\Z)^3$, take odd $N\ge3$, and put $n=N^3$, $h=N^{-1}$.
For a swap edge $e=\{x,y\}$ write $\eta^e$ for the exchanged configuration.
The actual source, including its stated one-direction modulation, has the
form
\begin{equation}
 \begin{split}
 (\mathcal L_Nf)(\eta)={}&\sum_x[a(1-\eta_x)+b\eta_x]
                       [f(\eta^x)-f(\eta)]\\
 &+\sum_e c_{N,e}(\eta)[f(\eta^e)-f(\eta)],
 \quad a=\frac{4\lambda_0}{5},\quad b=\frac{2\lambda_0}{3},
 \end{split}
 \label{v3:eq:phase-generator}
\end{equation}
where $\eta^x$ flips the indicated bit and $\lambda_0>0$ is the primitive
rate parameter (not the Higgs coupling).  The source's swap rates satisfy
\begin{equation}
 c_{N,e}(\eta)\ge0,\qquad c_{N,e}(\eta^e)=c_{N,e}(\eta).
 \label{v3:eq:swap-symmetry}
\end{equation}
For example, $c_{N,e}=\kappa N^2$ on the unmodulated edges, and on the
positive $e_1$ edge from $x$ the factor
$1+\theta\varphi_{x+2e_1}$ is unchanged by the swap when the environment
site is distinct.  These are the source hypotheses already displayed in
\texttt{def:modulated-renewal-interchange}.  No restriction on the size of
the finite swap rates is needed in the following equal-time calculation.

\begin{theorem}[Exact product-law entrance from constant-density initial data]
\label[theorem]{v3:thm:transient-product}
Start the actual finite process \eqref{v3:eq:phase-generator} from independent
Bernoulli variables of common parameter $r_0\in[0,1]$.  At every deterministic
time $t\ge0$, its full configuration law is exactly
\begin{equation}
 \mu_{N,t}=\operatorname{Bern}(r(t))^{\otimes n},\qquad
 r(t)=p+(r_0-p)e^{-(a+b)t},\qquad p=\frac6{11}.
 \label{v3:eq:transient-product}
\end{equation}
In particular, starting from the deterministic all-$H$ configuration gives
$r_0=0$, but at every fixed $t>0$ the variance per point is
$v(t)=r(t)(1-r(t))>0$.  This conclusion includes the source's allowed
nonlinear swap modulation.
\end{theorem}
\begin{proof}
Let $\nu_r=\operatorname{Bern}(r)^{\otimes n}$.  For $0<r<1$ it assigns
the same probability to $\eta$ and $\eta^e$, since a swap preserves the
number of occupied sites.  Equation \eqref{v3:eq:swap-symmetry} therefore
makes each swap generator reversible for \emph{every} $\nu_r$, so
$\nu_r\mathcal L_N^{\rm swap}=0$.  The flip generator is a sum of
independent two-state generators.  Differentiating the finite product
probabilities shows
\[
 \frac{\dd}{\dd t}\nu_{r(t)}
   =\nu_{r(t)}\mathcal L_N^{\rm flip}
 \quad\text{when}\quad r'=a-(a+b)r.
\]
Thus the displayed product curve solves the full finite forward equation,
with its specified initial law.  Uniqueness for this finite linear ODE proves
\eqref{v3:eq:transient-product}.  Continuity at $r_0=0,1$ includes the
constant initial configurations.  The argument proves only equal-time
product laws; the process at different times is not asserted independent.
\end{proof}

For the \emph{full point reader} put $\xi_x(t)=\eta_x(t)-r(t)$.  Define
\begin{equation}
 \widehat\xi_N(k,t)=\frac1n\sum_{x\in V_N}\xi_x(t)e^{-2\pi\ii k\cdot x/N},
 \quad k\in\mathcal B_N:=\{-\tfrac{N-1}{2},\ldots,\tfrac{N-1}{2}\}^3,
 \label{v3:eq:DFT}
\end{equation}
and let $u_N(t,x)=\sum_{k\in\mathcal B_N}\widehat\xi_N(k,t)
 e^{2\pi\ii k\cdot x}$ be its all-mode trigonometric interpolation on the
unit three-torus.  The normalization gives
$\|u_N(t)\|_2^2=n^{-1}\sum_x\xi_x(t)^2$.
Let $P_K$ retain the Fourier cube $|k|_\infty\le K$, and put
$d_{N,K}=\#\{k\in\mathcal B_N:|k|_\infty\le K\}$.

\begin{theorem}[Full-reader leakage from an initially constant configuration]
\label[theorem]{v3:thm:phase-leakage}
For every deterministic $t>0$ and $K<(N-1)/2$,
\begin{equation}
 \E\|(I-P_K)u_N(t)\|_2^2
       =v(t)\left(1-\frac{d_{N,K}}n\right),
 \qquad
 \E\|P_Ku_N(t)\|_2^2=v(t)\frac{d_{N,K}}n.
 \label{v3:eq:phase-tail-mean}
\end{equation}
If $K_N/N\to0$, then
\begin{equation}
 \|(I-P_{K_N})u_N(t)\|_2^2\longrightarrow v(t)>0
 \quad\text{in }L^1(\PP).
 \label{v3:eq:phase-tail-limit}
\end{equation}
The same convergence holds after integrating over a fixed interval
$[t_0,T]$ with $t_0>0$, with limit $\int_{t_0}^T v(t)\dd t$.
No independence of different observation times is required.  All fixed
smooth spatial tests of $u_N(t)$ nevertheless converge to zero in mean
square.  Hence the full point-reader family has no strong $L^2$ limit in
probability on this nontrivial positive-time entrance.
\end{theorem}
\begin{proof}
By \cref{v3:thm:transient-product}, equal-time centered coordinates are
independent with variance $v(t)$.  Orthogonality of the discrete characters
therefore gives
\begin{equation}
 \E[\widehat\xi_N(k,t)\overline{\widehat\xi_N(l,t)}]
           =\frac{v(t)}n\one_{\{k=l\}}.
 \label{v3:eq:Fourier-covariance}
\end{equation}
Summing over either Fourier set proves \eqref{v3:eq:phase-tail-mean}.
The full norm is an average of $n$ independent bounded variables, with mean
$v(t)$ and variance at most $1/n$.  It converges to $v(t)$ in $L^1$.
The low-mode norm is nonnegative with expectation tending to zero, since
$d_{N,K_N}/n\to0$.  Subtract it from the full norm to obtain
\eqref{v3:eq:phase-tail-limit}.  The estimates are uniform in $t$, so
integration and the triangle inequality in $L^1$ prove the time-integrated
statement without a temporal independence assertion.
For a smooth test $f$, Parseval and \eqref{v3:eq:Fourier-covariance} give
$\E|\langle u_N,f\rangle|^2\le v(t)\|f\|_2^2/n$.
Any strong $L^2$ limit in probability would therefore be zero, by testing on
a countable determining set, whereas its norm would have to tend to
$\sqrt{v(t)}>0$.  This is a contradiction.
\end{proof}

\begin{remark}[Meaning of the centering and of the physical clock]
Subtracting the deterministic spatial constant $r(t)$ only removes the zero
Fourier mode; it does not alter any high-frequency assertion.  The conclusions
therefore also hold for the uncentered phase reader in the nonzero modes.
Theorem \ref{v3:thm:transient-product} concerns deterministic physical-time
sampling of the inherited continuous-time generator.  Sampling at jump times
or conditioning on an accepted event can reweight the law and requires its
own table.  Neither reweighting nor temporal smoothing is hidden here.
\end{remark}

\section{An exponentially costly search for a smooth observed phase record}
\label{v3:sec:rare-selection}

The preceding expectation obstruction does not, by itself, exclude finding
a rare regular candidate in a long observed prefix.  We now quantify that
possibility.  The proof uses only finite product probabilities and a covering
of a predetermined low-mode space, not an unproved mixing assumption.

\begin{lemma}[Finite bounded-variable and low-rank projection bounds]
\label[lemma]{v3:lem:projection-concentration}
Let $\xi_1,\ldots,\xi_n$ be independent centered Bernoulli$(r)$ variables
and let $P$ be a real orthogonal projection of rank $d$ on $\R^n$.  For $u>0$,
\begin{equation}
 \PP\{\|P\xi\|_2\ge u\}\le 2\,5^d e^{-u^2/2}.
 \label{v3:eq:projection-concentration}
\end{equation}
If $v=r(1-r)\ge v_*>0$, then
\begin{equation}
 \PP\left\{\frac{\|(I-P)\xi\|_2^2}{n}\le\frac{v_*}{4}\right\}
 \le e^{-nv_*^2/2}+2\exp(d\log5-nv_*/8).
 \label{v3:eq:small-tail-probability}
\end{equation}
\end{lemma}
\begin{proof}
For a centered random variable $Z$ of range length $b$, the second derivative
of its log moment-generating function is a tilted variance, at most $b^2/4$.
This variance bound follows from $(Z-a)(a+b-Z)\ge0$ and completing a square
around the tilted mean.  Since the log moment-generating function and its
first derivative vanish at zero, integration gives
$\log\E e^{tZ}\le t^2b^2/8$.  Applying this to independent
$c_i\xi_i$, multiplying their moment-generating functions and optimizing in
$t$ gives
\[
 \PP\{|\sum_i c_i\xi_i|\ge z\}
       \le2\exp\left(-\frac{2z^2}{\sum_i c_i^2}\right).
\]
For rank $d>0$, the unit sphere of $\Ran P$ has a $1/2$-net with at most
$5^d$ points.  Indeed, a maximal $1/2$-separated set has disjoint radius-$1/4$
balls inside the radius-$5/4$ ball in that $d$-dimensional space; comparison
of their volumes proves the count, and maximality proves the covering.
If $\|P\xi\|_2\ge u$, some net vector $c$ satisfies
$|c\cdot\xi|\ge u/2$.  Use the last tail estimate with $\|c\|_2=1$ and
sum over the net.  The rank-zero case is immediate.

For the second claim, the independent variables $\xi_i^2$ lie in $[0,1]$
and have mean $v\ge v_*$.  The same exponential argument gives
\[
 \PP\{n^{-1}\|\xi\|_2^2<v_*/2\}\le e^{-nv_*^2/2}.
\]
If the total energy is at least $v_*/2$ but the complementary energy is at
most $v_*/4$, then $n^{-1}\|P\xi\|_2^2\ge v_*/4$.
Insert $u^2=nv_*/4$ in \eqref{v3:eq:projection-concentration} and add the
total-energy failure probability.
\end{proof}

\begin{theorem}[Selection-cost obstruction for the actual full phase reader]
\label[theorem]{v3:thm:selection-cost}
Use the source of \cref{v3:thm:transient-product} with $r_0=0$, fix
$0<t_0<T$, and let
$v_*:=\min_{t\in[t_0,T]}r(t)(1-r(t))>0$.
Let $1\le K_N$, $K_N/N\to0$, and let $t_{N,1},\ldots,t_{N,J_N}$ be any predetermined
physical observation times in $[t_0,T]$.  An index may be selected by any
rule depending on all the observed records.  For sufficiently large $N$,
\begin{equation}
 \boxed{
 \PP\!\left\{\exists j\le J_N:
       \|(I-P_{K_N})u_N(t_{N,j})\|_2^2\le v_*/4\right\}
       \le 3J_N e^{-c_*N^3},}
 \label{v3:eq:selection-obstruction}
\end{equation}
where $c_*:=\min\{v_*^2/2,v_*/16\}>0$.
Consequently, if $\log J_N=o(N^3)$, no selector on these observations produces
a small-tail record with a probability bounded away from zero.  To obtain
success probability at least a fixed $\varepsilon>0$ requires
$J_N\ge (\varepsilon/3)e^{c_*N^3}$ under this bound.
\end{theorem}
\begin{proof}
The discrete real projection corresponding to the symmetric Fourier cube
has rank $d_{N,K_N}=o(N^3)$, and its Euclidean squared norm divided by $n$
is the trigonometric $L^2$ norm.  Each deterministic-time marginal is the
product law of \cref{v3:thm:transient-product}.  Apply
\eqref{v3:eq:small-tail-probability} at each time.  For sufficiently large
$N$, $d_{N,K_N}\log5\le N^3v_*/16$, giving an upper bound
$3e^{-c_*N^3}$ for each event.  The union bound proves
\eqref{v3:eq:selection-obstruction}; correlations between the times play no
role.  Any selector can succeed only when this union event occurs.  The
remaining assertions follow by rearranging the same bound.
\end{proof}

\begin{corollary}[Failure of the weighted-tail entrance at polynomial comparison counts]
\label[corollary]{v3:cor:weighted-failure}
For this point reader define the spatial weighted tail
\[
 \tau^{\rm sp}_{N,m}(K)
 =\sum_{\substack{k\in\mathcal B_N\\|k|_\infty>K}}
          (1+|k|_1/K)^m|\widehat\xi_N(k)|,\qquad m\ge0.
\]
Since $\tau^{\rm sp}_{N,m}(K)\ge\|(I-P_K)u_N\|_2$, the probability of
finding an observed record with
$\tau^{\rm sp}_{N,m}(K_N)\le\sqrt{v_*}/2$ satisfies
\eqref{v3:eq:selection-obstruction}.  In particular, a tail threshold tending
to zero cannot be met with nonvanishing probability by any polynomial-size
predetermined observation bank after time $t_0$.
The same assertion holds, with the threshold multiplied by $c_0$, for any
common two-value local reader $F_N:\{H,P\}\to\R^d$ whose contrast satisfies
$|F_N(P)-F_N(H)|\ge c_0>0$.
\end{corollary}
\begin{proof}
The weighted $\ell^1$ norm dominates the unweighted $\ell^2$ norm of the tail
coefficients; Parseval identifies the latter.  For a two-value local reader,
\[
 F_N(\eta_x)=F_N(H)+[F_N(P)-F_N(H)]\eta_x.
\]
Every nonzero Fourier coefficient is therefore the phase coefficient
multiplied by this same contrast vector.  Its high-frequency $L^2$ norm is
at least $c_0$ times the phase norm.  Apply
\cref{v3:thm:selection-cost}.
\end{proof}

\begin{theorem}[Adaptive observation and physical waiting-time obstruction]
\label[theorem]{v3:thm:hitting-cost}
For the same source with $r_0=0$, suppose in addition that there are at most
$C_EN^3$ swap edges and $c_{N,e}(\eta)\le C_SN^2$ uniformly.  These bounds
hold for the supplied finite-range interchange generator with fixed primitive
rates and its allowed bounded modulation.  Fix $t_0>0$ and put
\[
 \underline v=\min_{r\in[r(t_0),p]}r(1-r)>0,
 \qquad c_*:=\min\{\underline v^2/2,\underline v/16\}.
\]
Let $1\le K_N$, $K_N/N\to0$.  There is $C_{\rm rate}<\infty$, independent
of $N$ and the deterministic horizon $T_N\ge t_0$, such that for sufficiently
large $N$,
\begin{equation}
 \boxed{
 \begin{aligned}
 &\PP\!\left\{\exists t\in[t_0,T_N]:
       \|(I-P_{K_N})u_N(t)\|_2^2\le\underline v/4\right\}\\
 &\hspace{2em}\le
       3\bigl[1+C_{\rm rate}N^5(T_N-t_0)\bigr]e^{-c_*N^3}.
 \end{aligned}}
 \label{v3:eq:hitting-bound}
\end{equation}
This bounds any selection based on the whole observed trajectory, including
adapted stopping times and jump-time observations, with success understood
unconditionally.  In particular any fixed or polynomial physical horizon
has success probability tending to zero.  A fixed success probability
$\varepsilon>0$ requires
\begin{equation}
 T_N-t_0\ge
 \frac{\varepsilon e^{c_*N^3}/3-1}{C_{\rm rate}N^5}
 \label{v3:eq:waiting-lower}
\end{equation}
whenever the numerator is positive.  This is a necessary waiting-time bound,
not a sufficiency statement about eventual sampling efficiency.
\end{theorem}
\begin{proof}
Let $\mathscr B_N$ be the set of bit configurations whose nonzero high-mode
energy is at most $\underline v/4$.  It is a fixed set: changing $r(t)$
changes only the zero mode, which $P_{K_N}$ retains.  For every deterministic
$t\ge t_0$, \cref{v3:thm:transient-product,v3:lem:projection-concentration}
give
\[
 \mu_{N,t}(\mathscr B_N)\le3e^{-c_*N^3}
\]
for all sufficiently large $N$, uniformly in $t$.

Write $q_N(x,y)$ for the off-diagonal rates.  At any target configuration
$y$ and any $t\ge t_0$, consider its incoming intensity divided by its
positive product probability:
\[
 R^-_{N,t}(y):=
 \frac{1}{\mu_{N,t}(y)}\sum_{x\ne y}\mu_{N,t}(x)q_N(x,y).
\]
For a swap predecessor the product probabilities agree, and
\eqref{v3:eq:swap-symmetry} identifies the forward and reverse swap rate.
Their total is at most $C_EC_SN^5$.  For a flip predecessor $x=y^i$,
the ratio of product probabilities is either $r(t)/(1-r(t))$ or its
reciprocal.  Both are bounded uniformly for $t\ge t_0$, because
$r(t)\in[r(t_0),p]\Subset(0,1)$.  There are $N^3$ such predecessors,
with flip rates bounded by $\max(a,b)$.  Hence
\[
 R^-_{N,t}(y)\le C_{\rm rate}N^5
 \qquad(t\ge t_0)
\]
with fixed physical rate constants.  No reversibility of the full transient
law is assumed in this estimate.

Let $N_{\mathscr B}$ count all jumps landing in $\mathscr B_N$ during
$(t_0,T_N]$, including jumps whose starting state was already there.
The elementary finite-state jump intensity formula gives
\[
 \begin{aligned}
 \E N_{\mathscr B}
 &=\int_{t_0}^{T_N}\sum_x\mu_{N,t}(x)
                    \sum_{y\in\mathscr B_N,\ y\ne x}q_N(x,y)\,\dd t\\
 &=\int_{t_0}^{T_N}\sum_{y\in\mathscr B_N}
                 \mu_{N,t}(y)R^-_{N,t}(y)\,\dd t\\
 &\le3 C_{\rm rate}N^5(T_N-t_0)e^{-c_*N^3}.
 \end{aligned}
\]
The intensity formula follows by conditioning over a time increment and
integrating; all state sums and rates are finite at fixed $N$.
If the trajectory visits $\mathscr B_N$, it either starts there at $t_0$
or has at least one such jump.  A union bound and
$\PP(N_{\mathscr B}\ge1)\le\E N_{\mathscr B}$ prove
\eqref{v3:eq:hitting-bound}.  Every rule selecting a visited state can
succeed only on this hitting event, even when it uses the complete future
trajectory.  Rearrangement gives \eqref{v3:eq:waiting-lower}.
\end{proof}

\begin{remark}[The precise scope of the obstruction]
The statement concerns the actual supplied phase law and its full local
point reader.  It is not an assertion about an unmeasured coupled
Einstein--SM decoder or a reader that physically averages several cells.
Such changes require their own same-history specification.  The deterministic
observation bound alone does not address stopping times; the actual-rate
calculation in \cref{v3:thm:hitting-cost} supplies that extension for this
specified phase source, retaining its unconditional failure probability.  The initial constant record at $t=0$ is also
not excluded.  What is excluded is a claim that a regular initial point
reader automatically remains regular, or that a polynomially long
positive-time selection can repair its leakage, merely because the native
source has mixing, action-neutral swaps, or positive source covariances.
\end{remark}

\section{An exact physical-reader spectrum instead of a generic energy hypothesis}
\label{v3:sec:reader-spectrum}

The energy obstruction can also be evaluated directly at every Sobolev order.
Use the equivalent torus norm
$\|u\|_{H^q}^2=\sum_k(1+|k|_2^2)^q|\widehat u(k)|^2$.

\begin{proposition}[Exact reader occupation energy for the phase source]
\label[proposition]{v3:prop:reader-energy}
Let a \emph{specified} translation-invariant linear physical reader have
Fourier multiplier $M_N(k)$, scalar or vector valued, and let $\alpha_N$
be its fixed real amplitude calibration.  For
\[
 \widehat w_N(k,t)=\alpha_N M_N(k)\widehat\xi_N(k,t),
\]
one has the exact identity
\begin{equation}
 \boxed{
 \E\|w_N(t)\|_{H^q}^2
 =\frac{\alpha_N^2v(t)}{N^3}
       \sum_{k\in\mathcal B_N}(1+|k|_2^2)^q|M_N(k)|^2.}
 \label{v3:eq:reader-spectrum}
\end{equation}
On a fixed positive-time interval this formula is also the exact
occupation-energy test after integrating $v(t)$.  If $\alpha_N=1$ and
$|M_N(k)|\ge c>0$ on a fixed positive fraction of modes with
$|k|_\infty\ge c'N$, the mean $H^q$ energy grows at least as $N^{2q}$
for $q>0$.  Hence a uniformly noncollapsed full-band reader cannot provide
an $N$-uniform mean $H^q$ bound for this source.
\end{proposition}
\begin{proof}
Multiply \eqref{v3:eq:Fourier-covariance} at $k=l$ by the displayed
Sobolev weight and squared multiplier and sum.  On the specified fraction
of high modes, there are at least $c''N^3$ terms with weight at least
$c'''N^{2q}$.  Since $v(t)>0$ at positive time, the lower bound follows.
Integration is justified by the finite sums and nonnegativity.
\end{proof}

\begin{theorem}[Raw Sobolev energy diverges with high probability]
\label[theorem]{v3:thm:raw-energy}
For the full point reader and every fixed $t>0$, $q>0$, there are constants
$c_{t,q}>0$, $C_t<\infty$ such that for all sufficiently large odd $N$,
\begin{equation}
 \PP\{\|u_N(t)\|_{H^q}^2\ge c_{t,q}N^{2q}\}
                 \ge1-C_tN^{-3}.
 \label{v3:eq:energy-probability}
\end{equation}
The constants can be chosen uniformly for $t\in[t_0,T]$.
Thus the failure of a uniform high-order energy entrance is not merely
an expectation produced by rare outliers.
\end{theorem}
\begin{proof}
Let $P_N^{\rm hi}$ be the real discrete projection onto
$|k|_\infty>N/4$.  Its rank divided by $n$ tends to $7/8$.
For independent centered real coordinates of variance $v$ and fourth moment
$m_4$, expansion of the four indices gives, for any real symmetric matrix $P$,
\[
 \Var(\xi^TP\xi)
 =2v^2\tr(P^2)+(m_4-3v^2)\sum_xP_{xx}^2.
\]
Indeed the only nonzero terms after subtracting the mean square are the two
pairwise contractions and the coincident-index fourth-cumulant correction.
For an orthogonal projection, $\tr(P^2)\le n$ and
$\sum_xP_{xx}^2\le n$; Bernoulli fourth moments are bounded by one.
It follows that
\[
 X_N:=n^{-1}\xi^TP_N^{\rm hi}\xi,
 \qquad \E X_N\longrightarrow\tfrac78v(t),
 \qquad \Var(X_N)\le C/n.
\]
Chebyshev therefore gives $X_N\ge v(t)/4$ with probability at least
$1-C_t/n$ for large $N$.  Every mode in this shell has weight at least
$(N/4)^{2q}$, so
$\|u_N(t)\|_{H^q}^2\ge(N/4)^{2q}X_N$.
This proves \eqref{v3:eq:energy-probability}; compact positive-time variance
margins give uniform constants.
\end{proof}

\begin{corollary}[Calibrated smoothing demand, when that reader is actually present]
\label[corollary]{v3:cor:smoothing}
For the scalar point reader, uniform mean $H^q$ control of
$\alpha_N u_N$ requires $\alpha_N=O(N^{-q})$; it then removes the point
reader's order-one $L^2$ noise when $q>0$.
For a separately specified fluctuation-normalized reader with
$\alpha_N=N^{3/2}$ and
$M_N(k)=(1+|k|_2^2)^{-s/2}$, uniform mean $H^q$ control holds precisely when
\begin{equation}
 s>q+\tfrac32.
 \label{v3:eq:smoothing-threshold}
\end{equation}
At equality its mean energy grows logarithmically; below equality it grows
as $N^{3+2(q-s)}$.
These are spatial reader conditions, not sufficient conditions for the
full time--space tails of \cref{v1:in:closure}.
\end{corollary}
\begin{proof}
With $M_N=1$, the sum in \eqref{v3:eq:reader-spectrum} is bounded above
and below by positive constants times $N^{3+2q}$, by counting a high
frequency cube and the entire band.  The first claim follows, together with
$\E\|\alpha_Nu_N\|_2^2=\alpha_N^2v(t)$.
For the fluctuation normalization, \eqref{v3:eq:reader-spectrum} becomes
$v(t)\sum_{k\in\mathcal B_N}(1+|k|^2)^{q-s}$.
The dyadic shell of radius $2^j$ contains order $2^{3j}$ lattice points and
contributes order $2^{[3+2(q-s)]j}$.  Summing the geometric series gives
\eqref{v3:eq:smoothing-threshold} and the two divergence rates.
Nothing in this sum estimates a time derivative.
\end{proof}

\paragraph{Relation to the regenerative Hodge statements.}
The upstream Hodge estimates control specified source energies and the
spectral law of their declared operators.  Equation
\eqref{v3:eq:reader-spectrum} instead evaluates the energy of the \emph{actual
field reader} under an actual transient law.  These are not interchangeable
notions of compactness.  A new coarse reader or contrast rescaling may be
physically legitimate, but it must be independently declared and its common
action identified again.  Neither is inferred from a source gap, and neither
is performed on the records in this iteration.

\section{From resolved action-sector data to the full physical drift}
\label{v3:sec:full-bank}

The exact Higgs moment test reduces a real observation burden, but does not
control the independent Einstein and gauge rows.  The following finite
composition spells out that remaining incidence condition.  Its algebra is
standard; the point is to prevent a restricted sector gradient from being
reported as the full source covector.

\begin{proposition}[Resolved-sector conorm for the common-action drift]
\label[proposition]{v3:prop:full-bank}
At a current full field $x$, suppose an actually recorded transition label
$b$ has probability $\pi_b(x)$, and let $P_b(x)$ be a physical-mass
orthogonal projection onto its tested variation directions.  Include
rejection or stopped-exit labels with zero increment and zero projection
when needed, so $\sum_b\pi_b=1$.  If the conditional increment of the
\emph{same total action} on each label obeys
\begin{equation}
 d_b(x)\le-\delta_h\kappa_b(x)\|P_b(x)E_h(x)\|_{0,h}^2
                        +\delta_h r_b(x),
 \quad r_b\ge0,\quad\kappa_b\ge0,
 \label{v3:eq:block-work}
\end{equation}
and the finite operator
\begin{equation}
 \mathcal M_h(x)=\sum_b\pi_b(x)\kappa_b(x)P_b(x)
 \label{v3:eq:block-conorm}
\end{equation}
satisfies $\mathcal M_h(x)\succeq\kappa_*I$ on the full retained physical
variation bank, then the actual stopped drift satisfies
\begin{equation}
 \mathfrak d_h(x)\le-\delta_h\kappa_*\|E_h(x)\|_{0,h}^2
                  +\delta_h\sum_b\pi_b(x)r_b(x).
 \label{v3:eq:full-drift}
\end{equation}
The largest constant obtainable from these sector estimates is
$\lambda_{\min}(\mathcal M_h(x))$.  If the tested blocks have a common
unseen direction, they supply no positive uniform full-bank constant.
\end{proposition}
\begin{proof}
Condition on the actually reported label and sum
\eqref{v3:eq:block-work}.  Orthogonality gives
$\|P_bE\|^2=(E,P_bE)_h$, so the negative term is
$-\delta_h(E,\mathcal M_hE)_h$.  Its Rayleigh quotient proves the claimed
bound and the optimal constant.  If a nonzero vector lies in every
$\Ker P_b$ with $\pi_b\kappa_b>0$, the quadratic form vanishes on it, so
its least eigenvalue is zero.  This is a failure of inference from the
sector estimates; it is not a claim that a fully evaluated total drift
could never be favorable at one particular field.
\end{proof}

\begin{corollary}[What the action-moment audit would discharge]
\label[corollary]{v3:cor:conditional-handoff}
Suppose the actual resolved source supplies the sector action increments,
including \eqref{v3:eq:Higgs-work} when a Higgs-only label is present,
with cutoff-uniform positive conorm in \eqref{v3:eq:block-conorm}.
If their certified remainders $r_h=\sum_b\pi_b r_b$ have the occupation
bound of \eqref{v2:eq:work-occupation}, then
\cref{v2:thm:occupation,v2:thm:joint,v2:cor:closure} apply with
$\kappa_A=\kappa_*$.  The reader-energy, alignment, action-span, and exit
conditions in those results remain unchanged.  In particular a full-bank
stationarity conclusion cannot be deduced from the Higgs moment test alone.
\end{corollary}
\begin{proof}
Equation \eqref{v3:eq:full-drift} is exactly the row hypothesis
\eqref{v2:eq:action-drift}.  Invoke the already proved occupation and joint
selection theorems, and the inherited closure only when its remaining
physical-reader and gauge hypotheses hold.  No new transition kernel is
defined in this composition.
\end{proof}

\section{Outcome of the native calculation and continuation boundary}
\label{v3:sec:ledger}

This iteration resolves two concrete questions left by Version 2.
First, the signed work of the minimal Higgs sector cannot be read from
conditional covariance alone: \cref{v3:prop:quartic-witness} exhibits a
strictly positive same-support ambiguity, while
\cref{v3:thm:Higgs-work} gives an exact finite action-aware replacement.
Second, the explicitly supplied phase source does not propagate smoothness
of its full point reader from constant initial data.  Its failure is
quantitative at the field level, at the energy level, and at the cost of
selecting a rare small-tail record.  Fast or modulated swaps preserve the
transient product law and hence do not repair this specific defect.

\paragraph{Why this is not another abstract closure loop.}
The raw reader and the generator in the leakage calculation are fixed source
data.  Its drift, product entrance, Fourier covariance, and rarity bound
are evaluated, not assumed.  The two moment-matched Higgs rows are explicitly
labelled alternative compatible rows, not falsely reported measurements of
the missing coupled source.  Their role is to identify exactly why an
outgoing Gram cannot decide that action-level question.  The common-action
PDE theorem, the native Hodge short, general target-extension duality, the
phase-count obstruction, and the Fourier--midpoint population were not
reproved.

\paragraph{The first remaining physical identification.}
For the intended locked source one needs its \emph{actual} map from the
retained phase/current/holonomy record to the physical coframe, gauge and
Higgs fields.  If it is merely a noncollapsed local point reader of the
explicit phase regulator, the preceding estimates exclude the required
regularity at polynomial observation counts after a positive physical time.
The actual jump-rate calculation further excludes any trajectory-based
selection on a polynomial physical horizon, without assuming deterministic
observation times.  If it is a distinct covariant multi-cell reader, its spectrum or energy
comparison must be evaluated on that same source; the monograph's missing
joint field row is not supplied by changing its name.  For a fixed-fields
Higgs transition, the first missing action-aware columns can be reduced to
$q_1,q_2,c_3,c_4$ rather than the full successor table.  Metric, gauge,
fermion and mixed updates still require their own same-action increments
and a noncollapsing full-bank conorm.

\paragraph{What is not concluded.}
We have neither disproved a classical limit for an independently specified
coupled Renewal model nor verified such a model's complete source law.
The actual interchange regulator is already known to have a stochastic
fluctuation continuum; absence of strong classical point-reader compactness
is compatible with that result.  The present obstruction also leaves open
a physically provided coarse field reader, exponentially long waiting,
a differently specified source law, and small-noise scalings, but these are
different specified problems.  Conditioning on a rare successful event does
not remove its unconditional probability cost.  No entropy correction, noise short, Fourier projection, action
change, or comparison-time/physical-time identification has been applied to
make the evaluated source pass.

\paragraph{Verification and append-only record.}
The diagnostic script verifies the quartic expansion and both positive
moment-matched laws, computes the sharp finite action interval, checks the
mass-normalized vector Higgs work identity, and tests the transient product
forward equation and matrix exponential on a small modulated ring.  The
ring checks the graph-independent product-law identity; it is not presented
as the three-dimensional physical source.  Its exact killed semigroup also
checks the trajectory-entrance inequality and the incoming-rate bound.
Separate three-dimensional
Fourier calculations check the exact reader-energy sum and finite samples
of the leakage statistics.  Every theorem above has an analytic proof;
these tests are diagnostics only.  The full Version 2 source before its
final end-document command is preserved byte for byte, including both
prior mathematical bodies.  No earlier theorem is silently amended.

\begingroup
\renewcommand{\refname}{Additional references for Version 3}
\begin{thebibliography}{9}
\small
\setlength{\itemsep}{1pt}
\raggedright
\bibitem{v3:Hoeffding}
W.~Hoeffding, \emph{Probability inequalities for sums of bounded random
variables}, Journal of the American Statistical Association \textbf{58}
(1963), no.~301, 13--30.
\href{https://doi.org/10.1080/01621459.1963.10500830}{doi:10.1080/01621459.1963.10500830}.
\bibitem{v3:Funaki}
T.~Funaki, \emph{Interface motion from Glauber--Kawasaki dynamics of
non-gradient type}, 2024.
\href{https://arxiv.org/abs/2404.18364}{arXiv:2404.18364}.
\end{thebibliography}
\endgroup
% END IMMUTABLE VERSION 3 BODY
% APPEND VERSION 4 HERE, BEFORE THE FINAL END-DOCUMENT COMMAND.


% BEGIN IMMUTABLE VERSION 4 BODY
\clearpage
\begin{center}
{\Large\bfseries Version 4: The physical decoder modulo gauge}\\[0.5em]
{\large Mixed link--Higgs incidence, nonlinear local amplitudes, and a metric-source test}
\end{center}
\addcontentsline{toc}{part}{Version 4: Gauge-invariant decoder and metric-source audit}

\section{Upstream target: identify physical fluctuations before suppressing noise}
\label{v4:sec:scope}

Version 3 proves a roughness obstruction for the complete point reader of the
specified phase process.  It does not prove that every nonlinear physical
reader is rough modulo gauge.  In a gauge theory a pointwise change of Higgs
orientation can be exactly compensated by its incident links.  Conversely,
flat links and a constant Higgs modulus do not guarantee that this compensation
occurs on the same record.  The question studied here is therefore the
\emph{mixed physical decoder}, not another covariance or compactness criterion.

\paragraph{Inherited statements and nonduplication.}
The finite gauge covariance of the common action, its invariant internal
vacuum, and its complete metric-variation convention are supplied in
\cite{GT}, at \nolinkurl{thm:SMST-universal-coupled-action},
\nolinkurl{thm:SM-active-SM-II}, and
\nolinkurl{thm:SMST-action-rigidity}.  They are used, not claimed as new.
The same-history marginal-identification obstruction is already present at
\nolinkurl{thm:GT-one-history-projection}; we do not reprove a general
identifiability theorem.  The phase law in
\cref{v3:thm:transient-product} is also inherited without repeating its
forward-equation argument.  Neither the Hodge short, the general quartic
work formula, nor the rare-record waiting bound is reconstructed here.

The new results are an exact radial--angular finite energy decomposition
with a quantitative weak-frame readout; a pair of full-history couplings
which agree on every separate Higgs and link observation but have different
covariant energies; a variance--energy bound for \emph{nonlinear local radial
decoders}; and a positive metric-source test which cannot be canceled by a
volume potential.  These calculations decide whether an apparent reader
failure is a removable gauge-coordinate fluctuation or a physical source
obstruction.  The elementary gauge identities have standard precedents in
lattice gauge theory; see \cite{v4:CH,v4:Bazavov}.  No general novelty claim is
made for gauge covariance or unitary gauge.

\paragraph{Scope of the source evaluation.}
The locked Einstein source and the separately specified renewal--interchange
process are still distinct, as stated in \cite{GT} at
\nolinkurl{fire:renewal-interchange-not-Einstein}.  Where we evaluate a proposed
Higgs decoder on the latter process, the proposal is explicitly an input to
that calculation, not a newly discovered identification of the locked source.
The two joint completions below are exact mathematical witnesses; neither is
reported as a measurement of an unavailable native joint row.  All prior bodies
are preserved verbatim, and all additions have label prefix \texttt{v4:}.

\section{An exact finite gauge-invariant decoder test}
\label{v4:sec:polar}

Let $V_N=(\Z/N\Z)^3$, $h=N^{-1}$, $n=N^3$, and let $H_x\in\C^d$.
An oriented link $U_{x,i}$ is unitary and transports the fibre at $x+e_i$
back to $x$.  All fields and links in a formula belong to the same record.
Define the positive spatial Higgs energy
\begin{equation}
 \mathcal K_h(H,U)
 =h^3\sum_{x\in V_N}\sum_{i=1}^3
       \left|\frac{U_{x,i}H_{x+e_i}-H_x}{h}\right|^2.
 \label{v4:eq:kinetic}
\end{equation}
This is a positive comparison quantity, not a claim that the Lorentzian
common action is positive.  Write $r_x=|H_x|$.  When $r_x>0$, put $v_x=H_x/r_x$;
when $r_x=0$, choose any unit vector $v_x$.  Expressions below with the factor
$r_xr_{x+e_i}$ are independent of this choice.

\begin{proposition}[Radial--angular splitting on the actual link]
\label[proposition]{v4:prop:polar}
For each oriented edge $xy$ one has the exact identity
\begin{equation}
 |U_{xy}H_y-H_x|^2
  =(r_y-r_x)^2+r_xr_y|U_{xy}v_y-v_x|^2.
 \label{v4:eq:polar}
\end{equation}
Both terms are invariant under simultaneous site-frame changes
$H_x\mapsto g_xH_x$, $U_{xy}\mapsto g_xU_{xy}g_y^{-1}$.
In particular
\begin{equation}
 \mathcal K_h(H,U)\ge
 \mathcal K_h^{\rm rad}(r):=
 h^3\sum_{x,i}\left|\frac{r_{x+e_i}-r_x}{h}\right|^2
 \label{v4:eq:radial-lower}
\end{equation}
for every unitary link choice, even if the links depend arbitrarily on the
complete phase record.  If $0<r_-\le r_x\le r_+$, then
\begin{equation}
 \mathcal K_h^{\rm rad}+r_-^2\mathcal K_h^{\rm ang}
 \le\mathcal K_h(H,U)
 \le\mathcal K_h^{\rm rad}+r_+^2\mathcal K_h^{\rm ang},\qquad
 \mathcal K_h^{\rm ang}:=h^3\sum_{x,i}
      h^{-2}|U_{x,i}v_{x+e_i}-v_x|^2.
 \label{v4:eq:polar-equivalence}
\end{equation}
\end{proposition}
\begin{proof}
Expand the left side as
$r_y^2+r_x^2-2r_xr_y\Re(v_x^*U_{xy}v_y)$ and use
$|U_{xy}v_y-v_x|^2=2-2\Re(v_x^*U_{xy}v_y)$.
This proves \eqref{v4:eq:polar}, including zero radii.  Unitarity of $g_x$
proves the two invariances.  Nonnegativity and summation give
\eqref{v4:eq:radial-lower}; bounding the factor $r_xr_y$ proves
\eqref{v4:eq:polar-equivalence}.
\end{proof}

\begin{proposition}[A finite mixed moment and its physical accuracy scale]
\label[proposition]{v4:prop:edge-audit}
For any joint field/link law, conditional or unconditional, put
\begin{equation}
 B_h=h^3\sum_{x,i}\E\bigl[r_x^2+r_{x+e_i}^2
                 -2\Re(H_x^*U_{x,i}H_{x+e_i})\bigr].
 \label{v4:eq:mixed-bank}
\end{equation}
Then $B_h\ge0$ and $\E\mathcal K_h=h^{-2}B_h$.  If a finite audit gives
$|\widehat B_h-B_h|\le e_h$, then
\begin{equation}
 \max\{0,h^{-2}(\widehat B_h-e_h)\}
 \le\E\mathcal K_h
 \le h^{-2}(\widehat B_h+e_h).
 \label{v4:eq:edge-uncertainty}
\end{equation}
Consequently an $O(1)$ energy bound through this audit needs an $O(h^2)$
upper bound for the unnormalized mixed deficit, including the certified
uncertainty.  Bounds on the two separate marginal laws do not determine
$B_h$.
\end{proposition}
\begin{proof}
Expand the square in \eqref{v4:eq:kinetic}, take its expectation, and insert
the two enclosure endpoints.  The final statement follows concretely from
\cref{v4:thm:two-completions} below.  No independence assumption or new
observation rule has been used to justify the enclosure.
\end{proof}

\subsection{What an observed nonzero weak doublet determines}

For $v=(v_1,v_2)^T\in\C^2$ with $|v|=1$, set
\begin{equation}
 Q(v)=\begin{pmatrix}v_1&-\overline v_2\\v_2&\overline v_1\end{pmatrix}.
 \label{v4:eq:weak-frame}
\end{equation}
Thus $Q(v)\in SU(2)$ and $Q(v)e_1=v$.  On a nonzero-Higgs chart this is an
algebraic frame computed from the actual field, not a Fourier filter.
For the next theorem restrict the link representation to $SU(2)$.

\begin{theorem}[Quantitative weak-frame normalization of the mixed reader]
\label[theorem]{v4:thm:weak-frame}
Suppose $0<r_-\le|H_x|\le r_+$ and $U_{xy}\in SU(2)$.  Define
\[
 W_{xy}=Q(v_x)^*U_{xy}Q(v_y).
\]
In these frames the Higgs value is $r_xe_1$, and
\begin{equation}
 |U_{xy}v_y-v_x|^2=|W_{xy}e_1-e_1|^2
                          =\tfrac12\|W_{xy}-I\|_{\rm F}^2.
 \label{v4:eq:SU2-incidence}
\end{equation}
For links on a fixed principal-logarithm chart with eigenangles in
$[-\theta_*,\theta_*]$, $\theta_*<\pi$, put
$A^{\rm rel}_{xy}=h^{-1}\log W_{xy}$.  Then
\begin{equation}
 \frac{2r_-^2}{\pi^2}\|A^{\rm rel}_{xy}\|_{\rm F}^2
 \le r_xr_yh^{-2}|U_{xy}v_y-v_x|^2
 \le\frac{r_+^2}{2}\|A^{\rm rel}_{xy}\|_{\rm F}^2.
 \label{v4:eq:SU2-log}
\end{equation}
Thus the measured radial and mixed angular energies control the normalized
weak links at this first-order level, independently of how rapidly the
original site frames fluctuate.  Under an $SU(2)$ site-frame change these $W_{xy}$ are unchanged.
The frame normalization transports, rather than changes, the joint field/link record.
\end{theorem}
\begin{proof}
Direct multiplication gives $Q(v)^*Q(v)=I$ and $\det Q(v)=1$.  Every
$W\in SU(2)$ is of the form
$\left(\begin{smallmatrix}a&-\overline b\\b&\overline a\end{smallmatrix}\right)$,
with $|a|^2+|b|^2=1$.  Hence
$\|W-I\|_{\rm F}^2=2(|a-1|^2+|b|^2)=2|We_1-e_1|^2$.
Conjugating the edge difference gives the first equality in
\eqref{v4:eq:SU2-incidence}.  If the eigenvalues are $e^{\pm\ii\theta}$,
then
\[
 \|\log W\|_{\rm F}=\sqrt2|\theta|,\qquad
 \|W-I\|_{\rm F}=2\sqrt2\sin(|\theta|/2).
\]
The elementary inequalities
$(2/\pi)|\theta|\le2\sin(|\theta|/2)\le|\theta|$ for
$|\theta|\le\pi$ give \eqref{v4:eq:SU2-log} after multiplication by
$r_xr_y/(2h^2)$ and the radial bounds.  Uniqueness of an $SU(2)$ matrix with prescribed first column gives
$Q(gv)=gQ(v)$ for $g\in SU(2)$, so the two endpoint changes cancel
in $W_{xy}$.  No derivative of $Q(v_x)$ has been estimated or assumed small.
\end{proof}

\begin{remark}[The electroweak stabilizer and the remaining sectors]
For $U(2)$ links, $W=\operatorname{diag}(1,e^{\ii\alpha})$ satisfies
$We_1=e_1$ without being the identity.  The Higgs edge word therefore does
not control the stabilizer connection.  In the full Standard-Model carrier
one must retain the corresponding gauge curvature and every other sector;
the $SU(2)$ estimate above is not a full-bank coercivity claim.  It also gives
no high-order spacetime tail, constraint propagation, or coframe estimate.
At a Higgs zero the normalized frame is undefined; the radial lower bound
\eqref{v4:eq:radial-lower} remains valid without a lower-radius assumption.
\end{remark}

\section{Identical marginal histories and flat links do not determine physical energy}
\label{v4:sec:couplings}

This section is a concrete witness for the mixed quantity
\eqref{v4:eq:mixed-bank}.  It does not posit a new coupled native law.
Let $\sigma_x(t)\in\{-1,1\}$ be the sign of the supplied phase process,
so that $\sigma=2\eta-1$, and let $\zeta(t)$ be an independent process
with the same complete law.  The process may be stationary or may have
the product entrance of \cref{v3:thm:transient-product}.  At each finite
spacetime vertex $a=(t_j,x)$ and oriented spacetime edge $ab$, fix a
unit doublet $e_1$ and a radius $r_*>0$.  Compare the following joint
completions:
\begin{align}
 \mathsf P:\quad&H_a=r_*\sigma_a e_1,&
      U_{ab}&=\sigma_a\sigma_b I_2,\\
 \mathsf Q:\quad&H_a=r_*\sigma_a e_1,&
      U_{ab}&=\zeta_a\zeta_b I_2.
 \label{v4:eq:two-laws}
\end{align}
The center $\{\pm I_2\}$ lies in $SU(2)$.  Edges can include temporal
edges; no independence of distinct physical times is required.

\begin{theorem}[A fully marginally matched mixed-decoder witness]
\label[theorem]{v4:thm:two-completions}
The two completions in \eqref{v4:eq:two-laws} have the same complete Higgs
history law and the same complete link history law, separately.  Every
closed link word equals $I_2$ under either law, so all plaquette curvatures
and all unmarked closed-loop gauge observables agree exactly.  Their
covariant Higgs energies, however, differ:
\begin{equation}
 \mathcal K_h^{\mathsf P}=0\quad\hbox{pathwise},\qquad
 \E\mathcal K_h^{\mathsf Q}(t)
   =\frac{6r_*^2}{h^2}\bigl[1-m(t)^4\bigr],\qquad
 m(t)=2r(t)-1.
 \label{v4:eq:mixed-gap}
\end{equation}
At a fixed time with $0<r(t)<1$,
\begin{equation}
 h^2\mathcal K_h^{\mathsf Q}(t)
 \longrightarrow6r_*^2[1-m(t)^4]
 \quad\hbox{in }L^2(\PP).
 \label{v4:eq:mixed-gap-concentration}
\end{equation}
The distinguishing one-edge observation is
$\Re(H_a^*U_{ab}H_b)$: it is $r_*^2$ under $\mathsf P$, whereas its
spatial-edge expectation under $\mathsf Q$ is $r_*^2m(t)^4$.
Thus even the complete \emph{separate} histories and zero curvature do not
supply the mixed kinetic or metric-source information.
\end{theorem}
\begin{proof}
The Higgs history in both cases is the same function of $\sigma$.  The link
history is the same deterministic function of either $\sigma$ or an
identically distributed $\zeta$, proving equality of its marginal law.
Products around a closed word cancel every vertex sign, giving $I_2$.
Under $\mathsf P$, $U_{ab}H_b=H_a$ at every edge.

Under $\mathsf Q$, transform to the site frames $\zeta_aI_2$.
The links become the identity and the Higgs field becomes
$r_*\xi_a e_1$, where $\xi_a=\sigma_a\zeta_a$.  At a fixed time the
$\xi_x$ are independent signs with mean $m(t)^2$.  On a spatial edge,
\[
 \E|\xi_{x+e_i}-\xi_x|^2=2[1-m(t)^4].
\]
There are $3n$ such edges and $h^3n=1$, giving
\eqref{v4:eq:mixed-gap}.  The random variables
$|\xi_{x+e_i}-\xi_x|^2\le4$ are independent on disjoint edges.
Each edge shares an endpoint with at most a fixed number of other edges.
The variance of their spatial average is therefore at most $C/n$:
expand its variance into edge pairs and discard the independent pairs.
This proves \eqref{v4:eq:mixed-gap-concentration}.  The mixed-word
expectation is $r_*^2\E(\sigma_x\zeta_x\sigma_y\zeta_y)=r_*^2m^4$;
the pathwise equality under $\mathsf P$ was already computed.
\end{proof}

\paragraph{Positive interpretation, with its proper boundary.}
At $r_*=v_H$, zero fermions, a flat coframe, and $\Lambda=0$, the matched
completion is an actual finite gauge transform of the internal vacuum
already supplied by \cite{GT}.  Its complete finite covariant matter
residual and matter stress vanish, however rough the site signs are.
This is an application of inherited finite gauge covariance, not a new
vacuum-existence theorem or a model of nontrivial physical Einstein--SM
evolution.  In the scrambled completion, separate marginal audits would
still pass, but the mixed kinetic energy does not.  This is precisely why a
source noise-short or a statement that the phase is ``only gauge'' needs a
verified same-history link incidence.

The construction shows possibility and ambiguity, not occurrence of either
completion in the locked source.  The independent copy used for $\mathsf Q$
is stated explicitly.  Its presence is explicit and is not inferred for the actual one-source law.
In particular, $\mathsf P$ is not selected because it has the desired limit.
An actual joint measurement would decide which, if either, is relevant.
Neither completion identifies the phase process's bare product action with
the common Einstein--Standard-Model action.  The latter remains the
independently declared physical function used in this test.

\section{Nonlinear local radial decoders cannot hide finite-amplitude disorder}
\label{v4:sec:local-radial}

The matched gauge example has constant radius.  We now test the opposite
case without choosing a gauge or assuming that the decoder is linear.
Let $\eta$ have an equal-time product Bernoulli law on $V_N$, as supplied
by \cref{v3:thm:transient-product}.  Let $R_N\ge0$ be an integer,
$N\ge8(R_N+1)$, and put
\[
 B_{R_N}=\{-R_N,\ldots,R_N\}^3,\qquad L_N=2R_N+1.
\]
A predeclared translation-covariant radial decoder has the form
\begin{equation}
 r_x=F_N\bigl((\eta_{x+a})_{a\in B_{R_N}}\bigr),\qquad
             0\le F_N\le M,
 \quad v_N=\Var(r_0).
 \label{v4:eq:local-decoder}
\end{equation}
The function $F_N$ can be arbitrary, nonlinear, and cutoff dependent.
Only its support, amplitude, and actual variance enter the bounds.
Take any complex fields $H_x$ with $|H_x|=r_x$ and any unitary links,
possibly depending on the whole configuration and additional recorded data.

\begin{theorem}[Local radial variance forces covariant energy]
\label[theorem]{v4:thm:radial-obstruction}
For every such decoder and link choice,
\begin{equation}
 \E\mathcal K_h(H,U)\ge
                 \frac{6v_N}{h^2L_N^2}.
 \label{v4:eq:local-mean-lower}
\end{equation}
If $v_N>0$, then
\begin{equation}
 \PP\!\left(\mathcal K_h(H,U)<\frac{v_N}{h^2L_N^2}\right)
 \le\frac{3(4R_N+1)^3M^4}{n v_N^2}.
 \label{v4:eq:local-prob-lower}
\end{equation}
Consequently, if $v_N\ge v_*>0$ and $R_N/N\to0$, the covariant spatial
energy diverges in probability.  This conclusion holds uniformly over all
unitary gauge completions.  Conversely a uniform \emph{mean} energy bound
$\E\mathcal K_h\le C$ forces
\begin{equation}
                    v_N\le\frac{C}{6}(hL_N)^2.
 \label{v4:eq:radial-collapse-needed}
\end{equation}
It does not forbid a local averaging decoder whose radial variance actually
vanishes at this scale.
\end{theorem}
\begin{proof}
The supports of $r_x$ and $r_{x+L_Ne_i}$ are disjoint, including on the
periodic grid under the displayed size condition.  Thus they are independent
and identically distributed, and
$\E(r_{x+L_Ne_i}-r_x)^2=2v_N$.
Telescoping along $L_N$ nearest-neighbor edges and applying Cauchy--Schwarz
gives pathwise
\[
 \frac1n\sum_x(r_{x+L_Ne_i}-r_x)^2
 \le L_N^2\frac1n\sum_x(r_{x+e_i}-r_x)^2.
\]
Taking expectations and using \eqref{v4:eq:radial-lower}, separately for
each of the three directions, proves \eqref{v4:eq:local-mean-lower}.

For the probability assertion use just $i=1$ and put
$Z_x=(r_{x+L_Ne_1}-r_x)^2$, $A_N=n^{-1}\sum_xZ_x$.
Each $Z_x$ lies in $[0,M^2]$ and depends on the input set
$x+S_N$, where $S_N=B_{R_N}\cup(L_Ne_1+B_{R_N})$.
Two such variables are independent unless their translated input sets
intersect.  The difference set $S_N-S_N$ is contained in three translates
of $\{-2R_N,\ldots,2R_N\}^3$, hence has at most
$3(4R_N+1)^3$ elements.  Expanding the variance gives
\[
 \Var(A_N)\le\frac{3(4R_N+1)^3M^4}{n},\qquad \E A_N=2v_N.
\]
The preceding pathwise telescope gives
$\mathcal K_h\ge A_N/(h^2L_N^2)$.  Chebyshev's inequality at distance
$v_N$ proves \eqref{v4:eq:local-prob-lower}.  If $R_N/N\to0$ and
$v_N\ge v_*$, the error probability tends to zero and the lower energy
threshold tends to infinity.  Rearranging
\eqref{v4:eq:local-mean-lower} proves \eqref{v4:eq:radial-collapse-needed}.
\end{proof}

\begin{corollary}[Observation windows without a time-independence assumption]
\label[corollary]{v4:cor:radial-time}
At finitely many deterministic times $t_a$, let the phase marginals be the
product laws of \cref{v3:thm:transient-product}, and allow arbitrary dependence
between times.  Suppose the same support and amplitude bounds hold and
$v_N(t_a)\ge v_*>0$.  Let $w_a\ge0$ and $W=\sum_aw_a>0$.
Then, for the same numerical constant as above,
\begin{equation}
 \PP\!\left(\sum_aw_a\mathcal K_h(t_a)
                 <\frac{Wv_*}{h^2L_N^2}\right)
 \le\frac{3(4R_N+1)^3M^4}{n v_*^2}.
 \label{v4:eq:radial-time}
\end{equation}
The bound is independent of the number of observation times.
\end{corollary}
\begin{proof}
Construct $A_N(t_a)$ as in the preceding proof.  Its mean is at least
$2v_*$ and its variance has the same uniform bound.  By
$|\Cov(X,Y)|\le\sqrt{\Var X\Var Y}$,
\[
 \Var\!\left(\sum_aw_a A_N(t_a)\right)
 \le W^2\frac{3(4R_N+1)^3M^4}{n}.
\]
Its mean is at least $2Wv_*$.  Chebyshev and the pathwise energy lower
bound prove \eqref{v4:eq:radial-time}.  No mixing or independence of the
time marginals has been inferred.
\end{proof}

\begin{proposition}[The exact point-amplitude scale]
\label[proposition]{v4:prop:point-radius}
For the proposed nonnegative point-radius reader
$r_x=c_N+d_N\eta_x$, $c_N\ge0$, $d_N\ge0$, one has at a fixed time
\begin{equation}
 \E\mathcal K_h^{\rm rad}
       =6d_N^2 h^{-2}r(t)(1-r(t)).
 \label{v4:eq:Bernoulli-energy}
\end{equation}
If $d_N>0$, the random variable
$h^2\mathcal K_h^{\rm rad}/d_N^2$ converges in $L^2$ to
$6r(t)(1-r(t))$ as $N\to\infty$.  Thus contrast bounded away from zero forces
an $h^{-2}$ energy, while $d_N\asymp h$ allows finite nonvanishing energy.
For $c_N=0$, $d_N=h$, the Higgs amplitude tends uniformly to zero but its
radial covariant-energy lower bound does not tend to zero.  Reducing the
amplitude to make the field converge is therefore not by itself a zero-stress
or zero-defect argument.
\end{proposition}
\begin{proof}
Each squared nearest-neighbor radial difference equals
$d_N^2\one_{\{\eta_x\ne\eta_{x+e_i}\}}$ and has expectation
$2d_N^2r(t)(1-r(t))$.  Summing the $3n$ terms proves
\eqref{v4:eq:Bernoulli-energy}.  For the normalized sum only pairs of edges
sharing an endpoint can have nonzero covariance; there are at most $Cn$
such pairs and every summand is bounded by one.  Its variance is $O(n^{-1})$,
proving the stated $L^2$ convergence and the scaling assertions.
\end{proof}

\section{A metric-source test insensitive to the volume potential}
\label{v4:sec:metric}

We now connect the reader calculation to a physical action variation, rather
than treating energy divergence as only a coordinate problem.  This section
uses the minimal Higgs sector on a flat spacetime cell chart with
$g_{\mu\nu}=\eta_{\mu\nu}=\operatorname{diag}(-1,1,1,1)$.  Define the actual
forward covariant differences $K_\mu=h^{-1}(U_\mu H_{+\mu}-H)$ on the
spacetime grid.  Their link and field variables are held fixed in the
inverse-metric variation, as in the independent physical variation bank.
The local density is
\[
 \mathcal L_H=-g^{\mu\nu}\Re(K_\mu^*K_\nu)-V(H),
\]
with its common volume factor.  The function $V$ may be any real algebraic
potential for which this variation is defined.  This restriction is enough
to test the claimed minimal action; it is not a replacement for its
fermionic or gravitational equations.

\begin{proposition}[Positive null-average metric source]
\label[proposition]{v4:prop:null-stress}
The complete Higgs Hilbert stress on this finite density satisfies
\begin{equation}
 \mathfrak n(T^H):=T^H_{00}+\frac13\sum_{i=1}^3T^H_{ii}
      =2|K_0|^2+\frac23\sum_{i=1}^3|K_i|^2.
 \label{v4:eq:null-stress}
\end{equation}
It is independent of $V$ and is unchanged by adding any tensor
$c\,\eta_{\mu\nu}$, even with a record-dependent scalar $c$.
Consequently
\begin{equation}
 h^3\sum_x\mathfrak n(T^H)(t,x)
             \ge\frac23\mathcal K_h(H(t),U(t)).
 \label{v4:eq:null-energy}
\end{equation}
These identities are valid with arbitrary temporal link differences and
arbitrary correlations between times.
\end{proposition}
\begin{proof}
Varying the common volume factor and the inverse metric gives
\[
 T^H_{\mu\nu}=2\Re(K_\mu^*K_\nu)
       -\eta_{\mu\nu}\left(-|K_0|^2+\sum_i|K_i|^2+V\right).
\]
Thus $T^H_{00}=|K_0|^2+\sum_i|K_i|^2+V$ and
$\sum_iT^H_{ii}=3|K_0|^2-\sum_i|K_i|^2-3V$.
Their displayed combination proves \eqref{v4:eq:null-stress}.
For $c\eta$ the same combination is $-c+c=0$.
Nonnegativity and spatial summation give \eqref{v4:eq:null-energy}.
Equivalently this is the average of the contractions with the six null
vectors $(1,\pm e_i)$; no continuum energy-condition theorem is invoked.
\end{proof}

\begin{theorem}[Local invariant disorder forces a nonvanishing metric row]
\label[theorem]{v4:thm:metric-obstruction}
Evaluate the local decoder of \cref{v4:thm:radial-obstruction} on finitely
sampled source histories with times in a fixed positive-time cylinder and
$v_N(t_a)\ge v_*>0$.  Take nonnegative time weights $w_a$ with
$\sum_aw_a=W>0$.  Under any unitary link completion,
\begin{equation}
 \PP\!\left(
 \sum_aw_ah^3\sum_x\mathfrak n(T^H)(t_a,x)
          <\frac{2Wv_*}{3h^2L_N^2}\right)
 \le\frac{3(4R_N+1)^3M^4}{n v_*^2}.
 \label{v4:eq:metric-radial}
\end{equation}
For $R_N/N\to0$ and $v_*>0$ this gives divergence of a positive physical
metric pairing in probability.  Neither a Higgs volume-potential change
nor a cosmological-constant counterterm cancels this pairing.

For the full common action, specialize to a constant flat coframe, flat
internal links and zero fermions, with a fixed nonzero gravitational
normalization and arbitrary constant $\Lambda$.  Use a fixed nonnegative
collar-supported time test $\chi$, sampled on the mesh, and the inverse-metric
variation
\begin{equation}
 k^{00}=\chi,\qquad k^{11}=k^{22}=k^{33}=\chi/3,
           \qquad k^{\mu\nu}=0\ (\mu\ne\nu).
 \label{v4:eq:metric-test}
\end{equation}
On this restriction,
\begin{equation}
 |D\mathcal A_h[k]|=\frac12h^4\sum_{j,x}
               \chi(jh)\mathfrak n(T^H)(jh,x).
 \label{v4:eq:metric-gradient}
\end{equation}
Hence the full physically normalized finite Euler row cannot tend to zero
for a decoder covered by \eqref{v4:eq:metric-radial}.  This is a conclusion
about this specified evaluation class, not about arbitrary nonflat solutions
or possible cancellations with unconstrained fermion stresses.
\end{theorem}
\begin{proof}
Apply \eqref{v4:eq:null-energy} inside
\cref{v4:cor:radial-time}.  This proves \eqref{v4:eq:metric-radial} and
its divergence claim.  The pointwise cancellation for $c\eta$ and $V$
is exactly \cref{v4:prop:null-stress}.

For the specialization, zero plaquette curvature gives zero Yang--Mills
metric variation and zero fermions give zero fermionic variation.  The
first-derivative gravitational bulk density has zero first variation at a
constant flat coframe, for compactly supported tests; its factors contain
vanishing first coframe derivatives.  The cosmological volume variation
is zero because $\eta_{\mu\nu}k^{\mu\nu}=0$ pointwise.  The scalar Hilbert
variation is $-\frac12h^4\sum T^H_{\mu\nu}k^{\mu\nu}$, giving
\eqref{v4:eq:metric-gradient}.  The mass norm of the fixed smooth test is
bounded independently of $h$, and the corresponding coframe lift on this
flat chart is uniformly bounded as well.  Duality therefore bounds the full
Euler-row norm below by a fixed positive multiple of the right side of
\eqref{v4:eq:metric-gradient}.  Use \eqref{v4:eq:metric-radial} with
$w_j=h\chi(jh)$ and $W_h\to\int\chi>0$.
\end{proof}

\begin{corollary}[The scrambled flat-link witness fails a full metric test]
\label[corollary]{v4:cor:scrambled-metric}
On the flat coframe and zero-fermion restriction, the scrambled completion
$\mathsf Q$ of \cref{v4:thm:two-completions} has a full finite metric row
of order at least $h^{-2}$ in probability on any fixed positive-time test
cylinder on which $1-m(t)^4$ is bounded below.  The matched completion
$\mathsf P$, at the inherited vacuum radius and $\Lambda=0$, has zero
matter stress and zero matter variation.  Both have the same separate
field and link histories, and both have identically zero gauge curvature.
\end{corollary}
\begin{proof}
The covariance-counting bound in
\eqref{v4:eq:mixed-gap-concentration} is uniform on the cylinder.
For the time-weighted energy sum, bound cross-time covariances by the
product of the corresponding standard deviations, as in
\cref{v4:cor:radial-time}.  Its normalized variance is $O(n^{-1})$
and its normalized mean is bounded below.  Equations
\eqref{v4:eq:null-energy}--\eqref{v4:eq:metric-gradient} give the claimed
Euler-row lower bound.  Flatness includes temporal plaquettes because
every spacetime link is an endpoint sign product.  The matched statement
is the inherited finite covariance and vacuum calculation described after
\cref{v4:thm:two-completions}.
\end{proof}

\section{A same-record decoder classification and its use in the closure programme}
\label{v4:sec:ledger}

The results distinguish three source situations which neither a separate
outgoing Gram nor ordinary componentwise roughness can distinguish.
\begin{enumerate}[label=\textup{(\roman*)},leftmargin=2.5em]
\item \emph{Pure orientation fluctuation with matched links.}
The radius is controlled and the actual mixed edge deficit is small.
In the weak $SU(2)$ sector, \cref{v4:thm:weak-frame} gives a physical
frame-normalized first-order reader without discarding any recorded noise.
The exact matched witness has zero deficit at every edge.
\item \emph{Constant radius but unverified mixed transport.}
Even all Higgs marginals, all gauge marginals and all closed-loop curvatures
can agree while the Higgs kinetic and metric-source bounds fail.
The joint word $\Re(H_x^*U_{xy}H_y)$ is load bearing.
\item \emph{Noncollapsed local invariant amplitude.}
For the specified product phase source, \cref{v4:thm:radial-obstruction}
excludes bounded covariant energy whenever the decoder's physical support
shrinks and its radial variance stays positive.  No gauge choice repairs
this case.  The metric theorem supplies a direct physical-row obstruction
on the stated flat evaluation, independent of the volume potential.
\end{enumerate}

\paragraph{The positive finite test.}
For an independently supplied physical decoder, evaluate the radii and the
mixed edge deficits on the same accepted records.  Their mass-normalized
sum is the exact finite Higgs energy, with the enclosure
\eqref{v4:eq:edge-uncertainty}.  On the nonzero weak-doublet chart the algebraic
frame \eqref{v4:eq:weak-frame} then exposes the actual relative links.
This is more informative than deciding to suppress every phase innovation:
matched orientation noise requires no suppression, whereas radial or
relative-link disorder must be paid for.  Higher spacetime differences,
curvature in the stabilizer and other gauge directions, metric noncollapse,
physical initial/gauge alignment and full same-action stationarity remain
separate.  The present first-order identity is not renamed as the high-order
tail certificate of \cref{v1:in:closure}.

\paragraph{What this adds to Versions 1--3.}
The raw phase reader's leakage obstruction was a statement about its specified
coordinates.  The new radial theorem is invariant under every unitary
completion and allows arbitrary nonlinear local amplitude decoding.  The
marginally matched witness shows why even curvature information cannot supply
the missing mixed reader.  The trace-free metric test then moves the issue
from abstract Sobolev regularity to a concrete common-action Euler direction.
These are distinct from re-proving Fourier tails, prefix descent, or a new
numerical population.  The matched gauge example also prevents the opposite
overstatement that every rough phase record is physically incompatible with
a classical field.

\paragraph{Remaining native obligation.}
The supplied locked source still needs its actual joint phase--Higgs--link
incidence.  If that incidence is matched, the action-neutral-noise route of
Version 2 is compatible with a controlled Higgs reader; if it is scrambled
or radially noisy as above, the corresponding bounds cannot be supplied by
marginal reconstruction alone.  The witness does not determine which law the
locked source uses.  One cannot select the matched completion after seeing
that it has good energy, or replace the same-action work by a newly fitted
potential.  A proof about one given coupled source must retain the complete
metric, gauge, Higgs and fermion variation bank and verify their joint rows.
No universal microscopic Einstein--SM emergence is asserted here.

\paragraph{Verification and append-only record.}
The diagnostic script checks the exact polar identity for random unitary
links, the $SU(2)$ frame and logarithm inequalities, both marginally matched
finite laws by exhaustive enumeration on a small ring, the Bernoulli and
nonlinear local radial energy bounds, and the metric-potential cancellation.
A finite Hilbert variation is checked by a centered numerical metric
difference.  Enumeration and floating-point checks are diagnostics, not
machine-checked proofs or an Einstein--SM simulation.  The complete V3 file
before its final end-document command is preserved byte for byte.

\begingroup
\renewcommand{\refname}{Additional references for Version 4}
\begin{thebibliography}{9}
\small\raggedright
\bibitem{v4:CH}
S.~H.~Christiansen and T.~G.~Halvorsen,
\emph{A simplicial gauge theory}, Journal of Mathematical Physics
\textbf{53} (2012), 033501.
\href{https://arxiv.org/abs/1006.2059}{arXiv:1006.2059}.
Background for exact discrete gauge covariance and consistency; no result
from this reference supplies the native joint source law used in these notes.
\bibitem{v4:Bazavov}
A.~Bazavov, Y.~Meurice, S.-W.~Tsai, J.~Unmuth-Yockey, and J.~Zhang,
\emph{Gauge-invariant implementation of the Abelian Higgs model on optical
lattices}, Physical Review D \textbf{92} (2015), 076003.
\href{https://arxiv.org/abs/1503.08354}{arXiv:1503.08354}.
Background for gauge-invariant lattice Higgs formulations; its model and
quantum limits are not substituted for the source studied here.
\end{thebibliography}
\endgroup
% END IMMUTABLE VERSION 4 BODY
% APPEND VERSION 5 HERE, BEFORE THE FINAL END-DOCUMENT COMMAND.


% BEGIN IMMUTABLE VERSION 5 BODY
\clearpage
\begin{center}
{\Large\bfseries Version 5: Covariant jets before coordinate regularity}\\[0.5em]
{\large A finite mixed-path reader, temporal normalization, and recovery of the full action covector}
\end{center}
\addcontentsline{toc}{part}{Version 5: Covariant-jet readers and full-variation gauge transport}

\section{Upstream audit and the precise positive target}
\label{v5:sec:audit}

Version 4 identifies the joint Higgs--link deficit but deliberately stops at
first order.  A gauge-equivalent representative with small link values does
not yet have controlled high spacetime derivatives, and a derivative of a
restricted gauge-fixed action need not have the norm of the full physical
Euler row.  These are the two interfaces treated here.  The source law and its
signed common action are not changed.

\paragraph{Inherited material.}
We use the radial--angular identity and the algebraic weak frame of
\cref{v4:prop:polar,v4:thm:weak-frame}; neither is a new result of this
iteration.  Finite gauge covariance and the complete Ward identity belong to
\cite{GT}, at \nolinkurl{thm:SMST-universal-coupled-action} and
\nolinkurl{thm:SMST-action-rigidity}.  The scalar Fourier-tail estimate and the
same-record selector are \cref{v1:eq:Fourier-tail,v1:thm:joint}.  The prefix
balance, regenerated-state estimates, coupled hyperbolic theorem, and
numerical population are not rederived.  The previous source audit does not
supply an independently verified high covariant-jet moment bound for the
locked Einstein source.

\paragraph{New statements and scope.}
For a nowhere-zero weak doublet, a finite tower of \emph{actual covariant}
differences controls high differences of a Higgs-normalized link reader,
without any derivative bound on the original site frames.  Its positive
energy is an explicit linear contraction of mixed path Grams of bounded word
length.  A causal finite temporal normalization retains these estimates with
one declared derivative loss.  We then prove an exact relation between the
full source covector and the derivative of the restricted action.  An
oscillatory example shows why the latter cannot silently replace the former.
The \(U(2)\) determinant channel, color, geometry, Higgs zeros, and the actual
propagation of the moment bank remain explicit boundaries.  There is no
claim of a new elementary unitary-gauge or dressing-field construction:
related continuum formulations are discussed in \cite{v5:ILM,v5:Zajac}.
The additions are the finite ordered-jet estimates, their same-record
normalization, and the complete-source accounting needed here.

\section{A source-readable positive covariant-jet energy}
\label{v5:sec:words}

Use a periodic grid \(\Lambda_h=(\Z/N\Z)^d\), \(h=N^{-1}\),
\(1\le d\le4\), with physical mass \(h^d\).  A link \(U_i(x)\in SU(2)\)
transports \(x+e_i\) to \(x\).  For a section \(F\) in a fixed unitary
representation \(\rho\), define
\begin{equation}
 T_i^UF(x)=\rho(U_i(x))F(x+e_i),\qquad
 D_i^UF=h^{-1}(T_i^UF-F),\qquad
 \delta_i f=h^{-1}(f(x+e_i)-f(x)).
 \label{v5:eq:covdiff}
\end{equation}
For an ordered word \(I=(i_1,\ldots,i_\ell)\), set
\(D_I^U=D_{i_1}^U\cdots D_{i_\ell}^U\), including the empty word.
The operators need not commute.  Put
\begin{equation}
 \mathcal J_{h,L}(F;U)^2
 =\sum_{\ell=0}^{L}\ \sum_{I\in\{1,\ldots,d\}^{\ell}}
       \|D_I^UF\|_{2,h}^2,\qquad
 \|F\|_{p,h}^p=h^d\sum_x|F(x)|^p.
 \label{v5:eq:jet-energy}
\end{equation}
This is a positive comparison energy of classical realified variables, not a
positivity assertion about the Lorentzian action.  All link and section
factors in a word are on the same physical record.

\begin{proposition}[Exact mixed-path Gram representation]
\label[proposition]{v5:prop:path-Gram}
For \(S\subset\{1,\ldots,\ell\}\), let \(p_S\) traverse the letters of
\(I\) in \(S\), in their original order, and write
\(F_S(x)=T_{p_S}^UF(x)\), \(c_S=(-1)^{\ell-|S|}\).  Then
\begin{equation}
 D_I^UF(x)=h^{-\ell}\sum_Sc_SF_S(x),\qquad
 |D_I^UF(x)|^2=h^{-2\ell}c^*G_I(x)c,
 \label{v5:eq:word-Gram}
\end{equation}
where \(G_I(x)_{ST}=\langle F_S(x),F_T(x)\rangle\succeq0\).
Every entry is gauge invariant.  Hence \(\mathcal J_{h,L}^2\) is an explicit
positive contraction of same-record transported endpoint Grams of paths of
length at most \(L\).  If certified entries obey
\(\max_{S,T}|\widehat G_{ST}-G_{ST}|\le e_{I,x}\), their contribution to
the energy uncertainty is at most
\begin{equation}
 h^d\sum_{\ell\le L}\sum_{|I|=\ell}\sum_x
                 h^{-2\ell}4^\ell e_{I,x}.
 \label{v5:eq:Gram-error}
\end{equation}
A direct enclosure of \(c^*G_Ic\) can replace this conservative entrywise
bound.
\end{proposition}
\begin{proof}
Expand the ordered product
\(h^{-\ell}(T_{i_1}^U-I)\cdots(T_{i_\ell}^U-I)\).
Each subset selects one path with exactly the stated sign; no interchange
of two shifts is made.  Squaring the resulting vector gives the Gram
identity.  Under site changes \(g_x\), every \(F_S(x)\) transforms by the
same unitary \(\rho(g_x)\), which cancels in its inner product.  There are
\(2^\ell\) coefficients of modulus one, so the absolute error in the
contraction is at most \(4^\ell e_{I,x}\).  Summing with the physical
weights proves \eqref{v5:eq:Gram-error}.
\end{proof}

\begin{remark}[What the moment bank does not assert]
Neither marginal Higgs data nor plaquette traces alone supply
\eqref{v5:eq:word-Gram}.  The positive bank contains mixed endpoint and open
transport information, including temporal words when \(d=4\).  Its expected
energy and its expected increment under an accepted row are respectively the
same displayed linear contraction of the expected Gram and of the
terminal-minus-initial Gram.  This identifies a finite source audit.  It does
not establish a negative Lyapunov drift or bounded occupation moments until
those inequalities have actually been checked.  The cancellation factors
\(h^{-2\ell}\) are retained, not absorbed into an unspecified constant.
\end{remark}

\begin{lemma}[A cutoff-uniform covariant embedding]
\label[lemma]{v5:lem:cov-embed}
For arbitrary unitary links and a fixed finite-dimensional unitary
representation, there is a constant independent of \(h\) and \(U\) such that
\begin{equation}
 \|F\|_{\infty,h}\le C_d\mathcal J_{h,3}(F;U),\qquad
 \max_{|I|\le L-3}\|D_I^UF\|_{\infty,h}
             \le C_{d,L}\mathcal J_{h,L}(F;U).
 \label{v5:eq:cov-embed}
\end{equation}
On a fixed rectangular slab the same assertion holds locally on a smaller
slab, with a fixed collar and the corresponding extra cutoff constants.
There is no derivative or small-link hypothesis on \(U\).
\end{lemma}
\begin{proof}
The reverse triangle inequality and unitarity give the exact discrete Kato
inequality
\[
 |\delta_i|F|(x)|\le|D_i^UF(x)|.
\]
For a scalar grid function, tensor-product piecewise multilinear
interpolation gives uniform norm comparisons with its grid \(L^p\) norms
and bounds its continuum first derivatives by the grid first differences.
These follow by scaling each cell to the fixed unit cube and finite-dimensional
norm equivalence; its derivative in direction \(i\) is a convex interpolation
of the corresponding edge differences.  The ordinary Sobolev inequalities
on the fixed torus therefore give, for \(d\le4\),
\[
 \|f\|_{4,h}\le C\|f\|_{W_h^{1,2}},\quad
 \|f\|_{8,h}\le C\|f\|_{W_h^{1,8/3}},\quad
 \|f\|_{\infty,h}\le C\|f\|_{W_h^{1,8}}.
\]
The zero-order term is included, so no mean-zero or boundary hypothesis is
needed.  Apply these inequalities to \(|F|\).  The first gives
\(\|F\|_{4,h}\le C\mathcal J_{h,1}(F;U)\).  Since physical volume is
one, \(\|\cdot\|_{8/3,h}\le\|\cdot\|_{4,h}\); applying the first
bound also to every \(D_i^UF\) gives
\(\|F\|_{8,h}\le C\mathcal J_{h,2}(F;U)\).
The last inequality, applied in turn to those sections, proves the first
claim.  Apply it to \(D_I^UF\) to obtain the second, retaining all ordered
prefix derivatives.  A fixed smooth cutoff proves the collar version:
\(D_i^U(\chi F)=\chi D_i^UF+(\delta_i\chi)T_i^UF\), and its iterates
contain only bounded cutoff differences and the declared covariant words.
\end{proof}

\section{High-order weak-frame control without derivatives of the original gauge}
\label{v5:sec:frame}

Let \(H_x\in\C^2\), \(r_x=|H_x|\ge\rho_*>0\), and
\(Q_x=Q(H_x/r_x)\) be the frame from \eqref{v4:eq:weak-frame}.
Extend the same formula to the real-linear quaternionic map
\(\mathsf Q(z)=\left(\begin{smallmatrix}z_1&-\bar z_2\\z_2&\bar z_1\end{smallmatrix}\right)\)
on all \(\C^2\).  Thus \(\mathsf Q(e_1)=I\).
Define actual normalized variables
\begin{equation}
 W_i(x)=Q_x^*U_i(x)Q_{x+e_i},\quad
 a_i(x)=h^{-1}(W_i(x)-I),\quad
 J_I(x)=Q_x^*D_I^UH(x).
 \label{v5:eq:normalized-jets}
\end{equation}
For the empty word \(J_\varnothing=re_1\).  These variables are invariant
under the weak \(SU(2)\) site action; no independent regularity of \(Q_x\)
is assumed.

\begin{lemma}[Exact closed difference identities]
\label[lemma]{v5:lem:closed-identities}
For every ordered word \(I\),
\begin{equation}
 \delta_iJ_I=J_{iI}-a_iT_iJ_I,
 \label{v5:eq:jet-recursion}
\end{equation}
where \(T_i\) now denotes the ordinary shift.  Put, for real \(r>0\) and
\(z\in\C^2\),
\begin{align}
 b_h(r,z)&=|re_1+hz|,\\
 R_h(r,z)&=\frac{2r\Re z_1+h|z|^2}{b_h(r,z)+r},\\
 A_h(r,z)&=\frac{\mathsf Q(z)-R_h(r,z)I}{b_h(r,z)}.
 \label{v5:eq:closed-maps}
\end{align}
Then the data satisfy the exact identities
\begin{equation}
 T_ir=b_h(r,J_i),\quad
 \delta_ir=R_h(r,J_i),\quad
 a_i=A_h(r,J_i).
 \label{v5:eq:closed-reader}
\end{equation}
On any set \(\rho_*\le r\le M\), \(|z|\le M\), \(hM\le\rho_*/2\),
the maps \(R_h,A_h\) and every fixed number of their derivatives in
\((r,z)\) are bounded independently of \(h\).  They extend smoothly to
\(h=0\); no inverse power of \(h\) is hidden in these derivative bounds.
\end{lemma}
\begin{proof}
Conjugating \(D_i^U(D_I^UH)\) gives
\(J_{iI}=h^{-1}(W_iT_iJ_I-J_I)\), proving
\eqref{v5:eq:jet-recursion}.  The first-jet identity is
\(re_1+hJ_i=(T_ir)W_ie_1\).
Taking norms gives the first part of \eqref{v5:eq:closed-reader}.
Rationalizing the norm difference gives the expression for \(R_h\).
An \(SU(2)\) matrix is determined by its first column, so
\[
 W_i=\frac{\mathsf Q(re_1+hJ_i)}{T_ir}
    =\frac{rI+h\mathsf Q(J_i)}{T_ir}.
\]
Subtract \(I\), divide by \(h\), and use \(T_ir-r=hR_h\).
The denominators satisfy \(b_h\ge\rho_*/2\) and \(b_h+r\ge3\rho_*/2\).
The radicand in \(b_h\) is uniformly positive.  Smooth differentiation on
this fixed compact set proves every stated uniform derivative bound,
including the limiting values at \(h=0\).
\end{proof}

For ordinary grid differences, write
\(\|f\|_{C_h^s}=\max_{|\alpha|\le s}\|\delta^\alpha f\|_{\infty,h}\).

\begin{theorem}[Covariant jets control a high-order normalized reader]
\label[theorem]{v5:thm:jet-reader}
Fix an integer \(s\ge0\), a lower radius \(\rho_*>0\), and a finite
bound \(R\).  Suppose
\begin{equation}
 |H_x|\ge\rho_*,\qquad
 \mathcal J_{h,s+4}(H;U)
       +\sum_b\mathcal J_{h,s+3}(F_b;U)\le R,
 \label{v5:eq:covariant-budget}
\end{equation}
where the \(F_b\) are any fixed finite collection of fields in fixed
unitary weak representations.  For sufficiently small \(h\), depending
only on the displayed data, every \(W_i\) belongs to a common principal
logarithm chart.  Set
\(\mathsf A_i=h^{-1}\log W_i\) and
\(f_b=\rho_b(Q_x)^*F_b\).  Then
\begin{equation}
 \|r\|_{C_h^{s+1}}+\sum_i\bigl(\|a_i\|_{C_h^s}
              +\|\mathsf A_i\|_{C_h^s}\bigr)
              +\sum_b\|f_b\|_{C_h^s}
       \le C_{s,\rho_*,R}.
 \label{v5:eq:high-reader}
\end{equation}
The constant is independent of every ordinary derivative or Fourier
frequency of the original \(Q_x\).  It does not require the original
\(U_i\) to be close to the identity.  A locally bounded version holds on
fixed nested collars.
\end{theorem}
\begin{proof}
By \cref{v5:lem:cov-embed}, all \(J_I\), \(|I|\le s+1\), have a
uniform pointwise bound \(M\), as does \(r\).  Choose \(hM\le\rho_*/2\).
Equations \eqref{v5:eq:closed-reader} give uniform bounds for \(a_i\) and
\(\delta_ir\).

We use the following elementary finite-difference chain rule.  For a smooth
map \(F\) on a fixed convex compact neighbourhood of the grid values,
\[
 \delta_iF(u_x)=\int_0^1DF(u_x+t(u_{x+e_i}-u_x))[\delta_iu_x]dt.
\]
Iteration and \(\delta_i(fg)=(\delta_if)T_ig+f\delta_ig\) bound each
fixed-order difference of \(F(u)\) by a polynomial in the differences of
\(u\) through the same order.  The coefficients use only finitely many
bounded derivatives of \(F\).  Induction using this identity proves the
claim without a factor \(h^{-1}\) not already in a declared difference.
For \eqref{v5:eq:closed-maps}, the neighbourhood may be chosen with
\(r>\rho_*/2\), bounded \(z\), and small \(h\), independently of the
record.

Apply a simultaneous induction in \(n\).  If
\(|I|+n\le s+1\), applying \(n-1\) ordinary differences to
\eqref{v5:eq:jet-recursion} bounds \(\delta^nJ_I\) using
\(\delta^{n-1}J_{iI}\), differences of \(a_i\) through order \(n-1\),
and lower differences of \(J_I\).  Next
\(\delta_ir=R_h(r,J_i)\) bounds the order-\(n\) differences of \(r\).
The chain rule for \(a_i=A_h(r,J_i)\) then bounds its order-\(n\)
differences for \(n\le s\).  At the last step,
\(\delta_ir=R_h(r,J_i)\) gives order \(s+1\) for \(r\).
The base case was established above, so this induction proves the asserted
bounds for \(r\) and \(a\).

For \(h\|a_i\|_{\rm op}\le1/2\),
\[
 \mathsf A_i=\int_0^1(I+th a_i)^{-1}a_i\,dt.
\]
This is the principal matrix logarithm divided by \(h\), by its convergent
power series.  Its derivatives as a function of \(a_i\) are uniformly
bounded on the fixed set, so the same chain rule proves its \(C_h^s\) bound.
For a fixed representation put
\(a_i^{(b)}=h^{-1}(\rho_b(W_i)-I)\).
Using \(W_i=e^{h\mathsf A_i}\) and integrating the exponential derivative
shows that \(a_i^{(b)}\) has the same fixed-order bounds.  The normalized
covariant words of \(F_b\) satisfy
\[
 \delta_i\{\rho_b(Q)^*D_I^UF_b\}
 =\rho_b(Q)^*D_{iI}^UF_b
       -a_i^{(b)}T_i\{\rho_b(Q)^*D_I^UF_b\}.
\]
Their pointwise bounds through order \(s\) follow from
\cref{v5:lem:cov-embed}; induction in this last identity gives
\(\|f_b\|_{C_h^s}\le C\).  The collar conclusion follows using the fixed
nested regions needed for the finitely many stencils.
\end{proof}

\begin{corollary}[Higgs nonvanishing from an initial margin on a short slab]
\label[corollary]{v5:cor:radius}
On a physical-time grid, suppose \(|H_{0,x}|\ge2\rho_*\) and
\(\sup_{j,x}|D_0^UH_{j,x}|\le M_0\).  Then
\[
 |H_{j,x}|\ge2\rho_*-jhM_0.
\]
In particular the lower-radius premise of \cref{v5:thm:jet-reader} follows
on \(0\le jh\le\rho_*/M_0\); when \(M_0=0\) it persists throughout
the supplied slab.  A covariant-jet bound of order four on a fixed larger
slab supplies \(M_0\) by \cref{v5:lem:cov-embed}, without first assuming
that the Higgs frame is regular.
\end{corollary}
\begin{proof}
Unitarity and the reverse triangle inequality give
\(\big||H_{j+1,x}|-|H_{j,x}|\big|\le h|D_0^UH_{j,x}|\).
Telescope along the actual temporal links.  Apply the embedding to
\(D_0^UH\) for the last statement.  The energy is taken on a fixed supplied
slab before the smaller interval is chosen, so its embedding constant does
not implicitly depend on the conclusion.
\end{proof}

\begin{proposition}[A higher-order curvature readout on this normalized branch]
\label[proposition]{v5:prop:curvature}
For \(s\ge1\), the normalized plaquette
\(P_{ij}=W_i(T_iW_j)(T_jW_i)^{-1}W_j^{-1}\) satisfies
\begin{equation}
 \frac{P_{ij}-I}{h^2}
 =\bigl(\delta_i a_j-\delta_j a_i
           +a_iT_i a_j-a_jT_j a_i\bigr)
                (W_jT_jW_i)^{-1}.
 \label{v5:eq:curvature-word}
\end{equation}
Consequently both \((P_{ij}-I)/h^2\) and, for small \(h\),
\(h^{-2}\log P_{ij}\) have uniformly bounded \(C_h^{s-1}\) norms under
\eqref{v5:eq:covariant-budget}.  They are the original plaquette transported
to \(Q_x\), not a replacement curvature.
\end{proposition}
\begin{proof}
Put \(X=W_iT_iW_j\), \(Y=W_jT_jW_i\).  Direct multiplication of
\(W_i=I+ha_i\) gives
\(X-Y=h^2(\delta_i a_j-\delta_j a_i+a_iT_ia_j-a_jT_ja_i)\).
Since \(P_{ij}=XY^{-1}\), this is \eqref{v5:eq:curvature-word}.
The product rule, unitarity, and \cref{v5:thm:jet-reader} give the first
bound.  Apply the preceding logarithm formula with \(h^2\) in place of
\(h\) to obtain the second.  Endpoint cancellations in the four links give
the exact conjugacy to the original plaquette.
\end{proof}

\section{Temporal normalization and all-mode reader tails}
\label{v5:sec:temporal}

The algebraic Higgs frame is not generally temporal gauge.  This distinction
matters because \cref{v1:in:closure} has a specified temporal normalization.
The following finite calculation keeps that step explicit.

\begin{theorem}[Finite causal temporal gauge with controlled derivatives]
\label[theorem]{v5:thm:temporal}
On \(\{0,h,\ldots,Mh\}\times\T_h^3\), \(Mh\le T\), suppose the
Higgs-normalized links obey
\(\max_\mu\|a_\mu\|_{C_h^{q+1}}\le B\), where
\(W_\mu=I+ha_\mu\), and \(q\ge0\).  Define
\begin{equation}
 g_{0,x}=I,\qquad g_{j+1,x}=g_{j,x}W_0(j,x).
 \label{v5:eq:temporal-product}
\end{equation}
Then \(g_j\in SU(2)\), \(\|g\|_{C_h^{q+1}}\le C_{q,B,T}\), and
\(W_\mu^{\rm T}=gW_\mu(T_\mu g)^{-1}\) satisfies
\begin{align}
 W_0^{\rm T}&=I,\\
 a_i^{\rm T}:=h^{-1}(W_i^{\rm T}-I)
       &=(ga_i-\delta_i g)(T_i g)^{-1},\\
 \max_i\|a_i^{\rm T}\|_{C_h^q}
       +\max_i\|h^{-1}\log W_i^{\rm T}\|_{C_h^q}
       &\le C_{q,B,T}
 \label{v5:eq:temporal-estimate}
\end{align}
for sufficiently small \(h\).  The same transformation preserves controlled
\(C_h^q\) norms of the already normalized Higgs and matter fields.  These
constants do not depend on the frequencies of the original Higgs frame.
No periodicity of \(g\) in physical time is claimed or needed.
\end{theorem}
\begin{proof}
Unitarity is preserved by the product and gives the zeroth-order bound.
The temporal-link equality follows directly from
\eqref{v5:eq:temporal-product}.  For an order-\(r\) spatial difference,
the iterated product rule gives a leading term
\((\delta^\alpha g_j)(T^\alpha W_{0,j})\), plus terms with a lower-order
difference of \(g_j\) and a positive-order difference of \(W_{0,j}\).
The leading multiplier is unitary.  Every other term contains
\(\delta^\beta W_0=h\delta^\beta a_0\), \(|\beta|>0\).
Taking maxima and inducting on \(r\le q+1\), then summing at most \(T/h\)
steps, bounds all spatial differences of \(g\).  Equivalently their finite
triangular recurrences are bounded by a discrete Gronwall estimate with
factor \(1+Ch\), not \(1+C\).  Mixed temporal differences follow from the
exact identity \(\delta_0g=ga_0\), the product rule and these bounds.

Expand \(g(I+ha_i)(T_ig)^{-1}-I\) and use
\(g-T_ig=-h\delta_ig\) to prove the middle formula.  It uses one more
spatial difference of \(g\) than of the connection being estimated; this
is the declared derivative loss.  Unitarity and the product rule now bound
\(a_i^{\rm T}\) through order \(q\), and the integral logarithm formula
proves the same bound for its logarithmic coordinate.  Finally multiply
\(re_1\) and the normalized matter fields by their respective unitary
\(g\)-representations and use the product rule.  No bound on a derivative
of the original rough frame enters any of these identities.
\end{proof}

\begin{proposition}[From the finite difference estimate to the declared Fourier reader]
\label[proposition]{v5:prop:Fourier-reader}
For odd \(N\), let \(I_hf\) be the all-mode trigonometric interpolation of
a grid array on a fixed \(d\)-torus.  For an integer \(q\ge0\),
\begin{equation}
 \|I_hf\|_{H^q}\le C_{d,q}
       \left(\sum_{|\alpha|\le q}\|\delta^\alpha f\|_{2,h}^2\right)^{1/2}.
 \label{v5:eq:Fourier-compare}
\end{equation}
Thus \cref{v5:thm:jet-reader} gives a uniform \(H^q\) bound for its
normalized primitive reader when \(s\ge q\).  After temporal normalization
one may take \(s=q+1\).  On a local slab, multiplication by a fixed smooth
cutoff equal to one on a smaller evaluation cylinder and supported in the
supplied collar gives the same bound for a zero-extended periodic reader.
The local statement asserts exact agreement at the physical nodes in that
smaller cylinder, not equality of two interpolants between nodes.
\end{proposition}
\begin{proof}
The symbol of \(\delta_i\) is
\(h^{-1}(e^{2\pi\ii k_i/N}-1)\).  For the represented frequencies it has
modulus between \(4|k_i|\) and \(2\pi|k_i|\).  Parseval and the equivalence
between \((1+|k|^2)^q\) and the finite sum of monomials in \(|k_i|^2\)
prove \eqref{v5:eq:Fourier-compare}.  The \(C_h^q\) bounds imply the
right side is bounded on the fixed physical volume.  On a slab, the product
rule bounds differences of the cutoff product.  For small \(h\), its support
is separated from the artificial periodic seam by more than \(qh\), so
zero extension creates no unestimated difference.  The interpolation retains
every Fourier coefficient of these specified nodal coordinates.
\end{proof}

\begin{corollary}[Gauge-robust quantitative tail transfer]
\label[corollary]{v5:cor:tail}
Let the weak-sector arrays satisfy \eqref{v5:eq:covariant-budget} with
\(s=q+1\), and let the preceding all-mode, collar-compatible reader be the
one used in the source-to-continuum comparison.  In four spacetime dimensions,
for \(q>m+2\), its normalized coordinates obey
\begin{equation}
 \tau_{h,m}(K)\le C_{q,m,\rho_*,R,T}K^{2-q}.
 \label{v5:eq:dressed-tail}
\end{equation}
This uses no cutoff-uniform bound on ordinary derivatives of the original
Higgs orientations.  It does not bound unobserved color, determinant-line,
coframe, or other reader components.
\end{corollary}
\begin{proof}
Apply \cref{v5:thm:temporal,v5:prop:Fourier-reader} and then the already
proved four-dimensional weighted Fourier estimate
\eqref{v1:eq:Fourier-tail}.  The normalized arrays are computed from the
actual joint field/link record.  The periodic extension, when used, is a
declared collar realization of that reader; an unrelated or unspecified
physical decoder cannot be substituted without a separate comparison.
\end{proof}

\paragraph{A quantitative conditional application.}
For example, take \(k=4\), \(m=6\), \(q=14\).
The conservative source orders above are \(19\) for the Higgs and \(18\)
for the other weak fields, including the one derivative used by temporal
normalization.  Suppose these have a fixed bound and lower-radius chart,
and the remaining physical sectors and collars independently supply the
same reader regularity.  Use the full link-tangent source norm of
\cref{v5:prop:full-variation}, or an independently verified equivalent
physical norm; \cref{v5:prop:log-source} converts it to the normalized
log-link row required by the inherited action comparison.  If that
\emph{full} source row obeys \(\sigma_h=O(h)\), choose \(K_h\asymp h^{-1/13}\).  Then
\[
 \tau_{h,2}+\tau_{h,6}=O(h^{12/13}),\qquad
 \varepsilon_{c,h}=O(h^{10/13}),\qquad
 F_{h,4}=O\bigl(h^{5/13}(1+|\log h|)^5\bigr).
\]
If the initial and harmonic budgets in these same normalized coordinates are
\(O(h^{1/3})\), \cref{v1:in:closure} gives an \(O(h^{1/3})\)
state/curvature/stress error.  This is a consequence of the inherited
closure theorem with its original geometric and action-identification
hypotheses, not a new proof of that PDE theorem.  The powers follow directly:
\(hK^3\) and \(K^2K^{2-q}\) are both \(h^{10/13}\), multiplication
by \(K^{k+1}\) gives \(h^{5/13}\), and the restored tail is
\(K^{k+2}K^{2-q}=h^{6/13}\).  The spare power absorbs the logarithm
relative to \(h^{1/3}\).  High covariant-jet energies and source rows have
not been proved small for an unspecified native law by this substitution.

\section{An exact audit of the full physical covector after normalization}
\label{v5:sec:variation}

Let \(\mathcal A_h\) be a differentiable, exactly weak-gauge-invariant
finite action on an invariant open configuration set.  This subsection is
finite and applies also to signed actions.  Give link variations the left
physical coordinate
\(a_i=h^{-1}(\dot U_i)U_i^{-1}\in\mathfrak{su}(2)\), with
\(\langle X,Y\rangle_{\mathfrak g}=-\tfrac12\Re\tr(XY)\).
The full tangent norm is the physical-mass sum of these link coordinates,
all Higgs and realified matter variations, and any retained independent
metric and other-sector coordinates.  Site gauge transformations act
unitarily on this norm.  Let \(E_h\) be the Riesz gradient in \emph{this}
full norm.  A different source Gram must be transported with the same
coordinate map, or compared by an explicit conorm.

\begin{proposition}[Field-dependent normalization loses no full first variation]
\label[proposition]{v5:prop:full-variation}
For a fixed site transformation \(g\),
\begin{equation}
 \mathcal A_h(g\cdot z)=\mathcal A_h(z),\qquad
 E_h(g\cdot z)=T_g E_h(z),\qquad
 \|E_h(g\cdot z)\|_{0,h}=\|E_h(z)\|_{0,h},
 \label{v5:eq:full-covector}
\end{equation}
where \(T_g\) is the isometric fixed-gauge tangent map.  If \(g=g(z)\)
is a differentiable, record-dependent normalization on a nonzero-Higgs
chart, then for every full variation \(v\)
\begin{equation}
 D\mathcal A_h(z)[v]
 =D\mathcal A_h(g(z)\cdot z)[T_{g(z)}v]
 =D\mathcal A_h(g(z)\cdot z)[D(g(\cdot)\cdot\cdot)_zv].
 \label{v5:eq:dependent-gauge}
\end{equation}
No derivative norm of \(g(z)\) is needed for the first equality or for
\eqref{v5:eq:full-covector}.  These statements concern the full variation
bank at the normalized configuration, not the gradient restricted to the
radial Higgs slice.
\end{proposition}
\begin{proof}
Differentiate invariance with \(g\) held fixed.  In the chosen coordinates
its tangent map sends \(a_i(x)\) to \(\operatorname{Ad}_{g_x}a_i(x)\),
\(\dot H_x\) to \(g_x\dot H_x\), and the other matter variations to
their unitary transforms.  Riesz representation proves
\eqref{v5:eq:full-covector}.  Differentiating a record-dependent \(g\) adds
to \(T_gv\) the tangent of a gauge orbit at \(g\cdot z\).  The exact
finite Ward identity annihilates that tangent, proving
\eqref{v5:eq:dependent-gauge}.  This cancellation does not imply that the
restricted slice has the full ambient metric or the full gradient norm.
\end{proof}

\begin{proposition}[Logarithmic-coordinate source conversion on the normalized chart]
\label[proposition]{v5:prop:log-source}
For a normalized link $W=e^{h\mathsf A}$ with $h\|\mathsf A\|\le c$,
the map from its logarithmic variation $\dot{\mathsf A}$ to its left
physical link variation is
\[
 J_h(\mathsf A)\dot{\mathsf A}
 :=h^{-1}(D e^{h\mathsf A}[h\dot{\mathsf A}])e^{-h\mathsf A}
 =\int_0^1\operatorname{Ad}_{e^{t h\mathsf A}}
                                  \dot{\mathsf A}\,dt.
\]
Here $D e^{h\mathsf A}[h\dot{\mathsf A}]$ denotes the derivative of the
matrix exponential at $h\mathsf A$ in direction $h\dot{\mathsf A}$.
For a sufficiently small fixed $c$, $\|J_h-I\|\le Cc<1/2$.
Consequently $E_{\mathsf A}=J_h^*E_W$ and the two full source norms
are uniformly equivalent on this normalized chart.  This assertion makes
no such equivalence on an uncontrolled original logarithm chart.
\end{proposition}
\begin{proof}
Differentiate the exponential by its convergent series, or its integral
identity, and multiply on the right by $e^{-h\mathsf A}$.  The difference
of the adjoint action from the identity has norm at most
$Cth\|\mathsf A\|$ on the compact matrix chart.  Integration proves the
operator bound; the Neumann series gives its inverse.  Pairing a link
covector with this tangent map gives the adjoint relation.  Sum over links
with the same physical cell masses; the unchanged field and metric rows
complete the full norm comparison.
\end{proof}

\begin{theorem}[The missing angular Higgs row is a divergence of retained sources]
\label[theorem]{v5:thm:Ward-recovery}
At a normalized configuration \(H_x=r_xe_1\), decompose
\(E_H=E_r e_1+E_H^{\rm ang}\), with \(E_r=\Re(e_1^*E_H)\).
Let \(E_i\) be the link-source components and put
\begin{equation}
 \operatorname{Div}_W E(x)=h^{-1}\sum_i
 \bigl[E_i(x)-\operatorname{Ad}_{W_i(x-e_i)}^{-1}E_i(x-e_i)\bigr].
 \label{v5:eq:Ward-divergence}
\end{equation}
For the other charged fields define \(\mathcal J_x(E)\in\mathfrak{su}(2)\)
by
\(\langle\mathcal J_x(E),\xi\rangle_{\mathfrak g}
=\sum_b\Re\langle E_b(x),\rho_{b*}(\xi)f_b(x)\rangle\).
Then the complete finite Ward identity gives
\begin{align}
 E_H^{\rm ang}(x)
   &=-r_x^{-1}\bigl(\operatorname{Div}_WE(x)+\mathcal J_x(E)\bigr)e_1,
       \label{v5:eq:angular-recovery}\\
 \|E_h\|_{0,h}^2
   &=\|E_r\|_{2,h}^2+\sum_i\|E_i\|_{2,h}^2
     +\sum_b\|E_b\|_{2,h}^2+\|E_{\rm other}\|_{2,h}^2\notag\\
   &\quad+\bigl\|r^{-1}
             (\operatorname{Div}_WE+\mathcal J(E))\bigr\|_{2,h}^2.
       \label{v5:eq:full-recovery}
\end{align}
The formula supplies the full norm from the slice derivatives only when the
additional divergence/current term is retained.  Its inverse-difference
cost is in general of order \(h^{-1}\).
\end{theorem}
\begin{proof}
For an arbitrary site parameter \(\xi_x\), the vertical tangent has
\(a_i(x)=h^{-1}(\xi_x-\operatorname{Ad}_{W_i(x)}\xi_{x+e_i})\),
\(\dot H_x=r_x\xi_xe_1\), and
\(\dot f_b=\rho_{b*}(\xi_x)f_b\).
Summation by parts on the finite grid gives
\[
 \operatorname{Div}_WE+T_r^*E_H+\mathcal J(E)=0,
 \qquad T_r\xi=r\xi e_1.
\]
The real-linear map \(T_r\) maps \(\mathfrak{su}(2)\) onto the
three-dimensional orthogonal complement of \(\R e_1\), and direct
\(2\times2\) multiplication gives \(T_r^*T_r=r^2I\).
Consequently
\(E_H^{\rm ang}=T_r(T_r^*E_H)/r^2\), proving
\eqref{v5:eq:angular-recovery}.  Orthogonality of radial and angular
components gives \eqref{v5:eq:full-recovery}.  Finally each unitary shifted
term in \eqref{v5:eq:Ward-divergence} has the same \(L^2\) norm as its
unshifted version, giving an upper bound with \(Ch^{-1}\).
The following example attains this order.
\end{proof}

\begin{proposition}[Small slice slopes do not control the full physical row]
\label[proposition]{v5:prop:gradient-witness}
On an even periodic grid, let \(H_x=e_1\), let
\(T=\operatorname{diag}(\ii,-\ii)\), and set
\[
 W_1(x)=\exp(h^2(-1)^{x_1}T),\qquad W_i(x)=I\ (i\ne1).
\]
Consider the gauge-invariant positive kinetic test action
\(\mathcal A_h=\tfrac12h^d\sum_{x,i}|D_i^WH(x)|^2\).
In the physical coordinates above its only nonzero link source is
\begin{equation}
 E_1(x)=h^{-1}\sin(h^2)(-1)^{x_1}T,
 \quad E_r=2h^{-2}(1-\cos(h^2)),
 \quad E_H^{\rm ang}=-2h^{-2}\sin(h^2)(-1)^{x_1}\ii e_1.
 \label{v5:eq:gradient-witness}
\end{equation}
Thus the naive product-norm gradient of the action restricted to
\(H=re_1\) is \(O(h)\), while the full gradient norm tends to \(2\).
All plaquettes are the identity.  The factor \(h^{-1}\) in the Ward
recovery cannot be removed by a cutoff-independent equivalence of these
naive full and slice norms.
\end{proposition}
\begin{proof}
The edge density is \(h^{-2}(1-\cos(h a_1))\) when
\(W_1=e^{ha_1T}\).  Its derivative in \(a_1\) is
\(h^{-1}\sin(ha_1)\), giving the link formula at
\(a_1=h(-1)^{x_1}\).  Differentiating in the independent real and
imaginary Higgs components gives
\[
 E_H(x)=h^{-2}\{2H_x-W_1(x)H_{x+e_1}
                       -W_1(x-e_1)^*H_{x-e_1}\}
       =2h^{-2}(1-e^{\ii h^2(-1)^{x_1}})e_1.
\]
This proves the other formulas.  With unit total mass, their norms have
orders \(h\), \(h^2\), and a limit \(2\), respectively.
The links depend only on \(x_1\) and lie in one commuting subgroup, so every
plaquette cancels.  The Ward divergence is
\(2h^{-2}\sin(h^2)(-1)^{x_1}T\), verifying
\eqref{v5:eq:angular-recovery} directly.
\end{proof}

\begin{remark}[What this witness rules out]
The example rules out replacement of a full \emph{mass-norm} source row by a
naively restricted gradient.  It does not rule out distributional
stationarity or a classical vacuum limit: the oscillatory angular row may
vanish against fixed smooth tests, and the covariant kinetic energy here
actually tends to zero.  The source norm required by the inherited
quantitative theorem is a sufficient, stronger entrance.  Neither changing
that norm nor discarding the recovered angular equation is justified by
calling the operation a gauge choice.
\end{remark}

\section{Residual electroweak information and the same-source continuation}
\label{v5:sec:boundary}

\begin{proposition}[The exact additional link coordinate for the \(U(2)\) sector]
\label[proposition]{v5:prop:U2}
For \(W\in U(2)\), its first column \(w=We_1\) and determinant
\(d=\det W\) determine it exactly:
\begin{equation}
 W=\mathsf Q(w)\operatorname{diag}(1,d).
 \label{v5:eq:U2-factor}
\end{equation}
For an arbitrary \(U(2)\) gauge change \(g_x\), the Higgs-normalized link
transforms by the residual determinant-line action
\begin{equation}
 W_{xy}\longmapsto
 \operatorname{diag}(1,\det g_x)\,W_{xy}\,
                         \operatorname{diag}(1,\det g_y)^{-1}.
 \label{v5:eq:U2-residual}
\end{equation}
Consequently no number of covariant derivatives of a constant nonzero Higgs
alone controls all \(U(2)\) link directions: links
\(W_i(x)=\operatorname{diag}(1,d_i(x))\) annihilate every positive-order
covariant Higgs word for arbitrary \(d_i(x)\).
If the same record independently supplies normalized determinant links
\(d_i=e^{h b_i}\) with bounded \(C_h^s\) coefficients \(b_i\), the
reconstruction argument of \cref{v5:thm:jet-reader} extends to this sector
with those determinant estimates retained as additional inputs.
\end{proposition}
\begin{proof}
The columns of \(W\) are orthonormal, so its second column is a unit
phase times \((-\bar w_2,\bar w_1)^T\).  Taking determinants identifies
that phase as \(d\), proving \eqref{v5:eq:U2-factor}.
The same formula gives
\(Q(gv)=gQ(v)\operatorname{diag}(1,\det g)^{-1}\); apply it at the
endpoints to obtain \eqref{v5:eq:U2-residual}.
For the stated diagonal links, \(T_i^UH=H\), and every positive-order
covariant word is zero.  This includes flat determinant connections with
uncontrolled period holonomy as well as nonflat ones.
For the final assertion put
\(a_i^{\rm SU}=h^{-1}(\mathsf Q(w_i)-I)\).
Then
\[
 h^{-1}(W_i-I)=a_i^{\rm SU}
                \operatorname{diag}(1,e^{hb_i})
               +\operatorname{diag}(0,(e^{hb_i}-1)/h).
\]
The first-column reconstruction and the same induction apply with this
extra known coefficient.  The product and exponential integral rules give
all stated finite-order bounds.  No new physical Higgs or determinant
probe is introduced by the argument.
\end{proof}

\paragraph{What has been advanced.}
The physical reader problem is no longer an unspecified claim that a Hodge
screen somehow controls a gauge potential.  On the nonzero weak-doublet
chart, an explicit finite bank of covariant path Grams gives a positive jet
energy; that energy produces a high-order frame-normalized reader and its
actual curvature.  The transformation can be put into temporal gauge with a
known derivative count.  The complete finite Euler row is unchanged in its
isometric full physical norm, and any attempt to use a reduced derivative
must restore the divergence/current term in \eqref{v5:eq:full-recovery}.
These are positive transfer results, not another roughness example in a
newly chosen sampler.

\paragraph{The native calculation still to perform.}
The next finite source test is an occupation or drift bound for the actual
mixed path Grams of \cref{v5:prop:path-Gram}.  The regenerative Lyapunov theorem
applies only if its energy dominates these words with uniform constants;
spatial control alone does not cover temporal words.  The supplied phase
marginals do not provide that identification for the locked coupled source.
Determinant/color control, the full action covector, geometric noncollapse,
harmonic and initial alignment, and the declared collar reader remain the
separate inputs of \cref{v1:in:closure}.  Gauge normalization cannot create
missing physical correlation, occurrence, or stationarity.

\paragraph{Relation to the earlier diagnostics.}
In the matched completion of \cref{v4:thm:two-completions}, all positive-order
covariant Higgs words vanish and the normalized links are the identity.
This is an application, not a new source-generation theorem.  The scrambled
completion fails already at first order, and the radial obstruction remains
gauge invariant.  The positive transfer therefore respects the previous
actual-source obstructions.

\paragraph{Verification and append-only record.}
The accompanying script checks ordered noncommuting path expansion, the
closed normalized-jet identities, covariance under rough independent site
frames, exact plaquette recovery, temporal gauge products, the \(U(2)\)
factorization and residual law, and the cutoff-dependent full/slice gradient
witness.  It also verifies the example's exponents with rational arithmetic.
These are finite numerical diagnostics, not machine-checked proofs or an
Einstein--SM simulation.  The complete V4 source before its final
end-document command is preserved byte-for-byte.  All new labels have prefix
\texttt{v5:}.

\begingroup
\renewcommand{\refname}{Additional references for Version 5}
\begin{thebibliography}{9}
\small\raggedright
\bibitem{v5:ILM}
A.~Ilderton, M.~Lavelle, and D.~McMullan,
\emph{Symmetry breaking, conformal geometry and gauge invariance}, 2010.
\href{https://arxiv.org/abs/1002.1170}{arXiv:1002.1170}.
Background for Higgs-dependent invariant variables and their gauge-chart
limitations; no cutoff-uniform discrete estimate is imported from this work.
\bibitem{v5:Zajac}
M.~Zaj\k{a}c,
\emph{The dressing field method in gauge theories---geometric approach},
2021 preprint; published 2023.
\href{https://arxiv.org/abs/2105.09919}{arXiv:2105.09919}.
Background for geometric symmetry reduction; the physical source and
finite-reader bounds of this iteration are stated and proved separately.
\end{thebibliography}
\endgroup
% END IMMUTABLE VERSION 5 BODY
% APPEND VERSION 6 HERE, BEFORE THE FINAL END-DOCUMENT COMMAND.


% BEGIN IMMUTABLE VERSION 6 BODY
\clearpage
\begin{center}
{\Large\bfseries Version 6: Chronological innovations before temporal regularity}\\[0.5em]
{\large Exact covariant time-difference spectra, a native count calculation,
and an interpolation obstruction}
\end{center}
\addcontentsline{toc}{part}{Version 6: Chronological source tests and temporal reader regularity}

\section{Upstream audit: the clock and the reader must be identified together}
\label{v6:sec:scope}

Version 5 identifies a finite, gauge-invariant bank of mixed path Grams whose
positive covariant-jet energy controls a normalized physical reader.  Its
hypotheses concern spacetime words, including temporal words.  The next
question is not another Sobolev embedding: it is whether an actual
chronological source can supply those temporal moments.  A source may have
bounded amplitudes and even spatially constant readouts while its individual
physical records have large temporal differences.

\paragraph{Nonduplication and inherited inputs.}
We use \cref{v5:prop:path-Gram,v5:thm:jet-reader} for the spatial/covariant
reader construction.  The elementary occupation selector, action-drift
entrance, high-mode spatial obstruction, and gauge normalization from
Versions 1--5 are not proved again.  The supplied Grand Tensor \cite{GT}, at
\nolinkurl{thm:renewal-interchange-actual-audit}, already establishes the
birth--death count quotient, and at
\nolinkurl{thm:renewal-interchange-OU} establishes its separate fluctuation
interpretation.  Its regenerative Lyapunov and covariance theorems concern
specified source forms, not arbitrarily high time derivatives of a decoded
path.  The physical source norm and the time/space forcing requirements of
\cite{CC}, at \nolinkurl{v27:thm:source}, remain inherited inputs.

\paragraph{The additional calculation.}
We compute the exact drift of the mixed-jet energy under a joint transition,
and then an exact temporal finite-difference norm for a reversible finite
source equipped with its actual unitary transport.  The resulting positive
spectral polynomial gives sharp constants and a necessary and sufficient
finite temporal-energy test.  A transient one-jump formula separates the
role of reversibility from the role of genuine innovations.  Evaluation on
the supplied count process proves a high-temporal-derivative obstruction in
probability, including for every interpolant retaining the observed nodes.
A final conditional result records exactly what temporal integration by an
already specified reader could supply.  No such reader is added to the
native source during this calculation.

\paragraph{Two clocks are not interchangeable.}
The variable $t$ below is the calibrated chronological clock of a field
reader when that reader and clock are part of the actual source.  It is
\emph{not} automatically the accepted comparison index used in
\cref{v1:in:prefix,v2:thm:occupation}.  Successive candidates in an action
search can each encode a complete smooth spacetime history; jumps between
such candidates are not temporal jumps within a history.  Every application
of the temporal theorems must specify which of these objects is being read.
Spatial mesh size $h$ and physical observation spacing $\delta$ are also
kept independent until an explicit specialization.

\paragraph{Attribution and scope.}
The quadratic generator identity, finite spectral theorem, and characteristic
function argument below use standard tools.  Covariant Markov semigroup
representations on graph vector bundles are established theory; see
\cite{v6:GMT,v6:KL}.  We give all finite proofs needed here and make no
priority claim for Feynman--Kac representations or finite-difference
filtering.  The contribution within this lineage is their same-record,
physical-mass normalization, the exact temporal audit, and its evaluated
consequences for the supplied source.  No assertion about an unspecified
coupled Einstein--SM generator follows from these calculations.

\section{The actual transition's work on the positive covariant-jet energy}
\label{v6:sec:jet-drift}

At one fixed cutoff let $\mathscr S$ be a finite source state space and
$q(x,y)\ge0$ its actual transition rates, $x\ne y$.
Its scalar generator is
\[
 (\mathcal L\phi)(x)=\sum_{y\ne x}q(x,y)[\phi(y)-\phi(x)].
\]
A state contains the complete record required to make these rates valid.
Let $\mathscr H_x$ be isometric finite-dimensional physical reader fibres,
with their declared mass inner products, and let $\mathsf U_{xy}$ transport
$\mathscr H_y$ unitarily to $\mathscr H_x$ on that transition.  Distinct
resolved transition labels can be retained as parallel edges with their own
rates and transports; all sums below then include the label.  For readability
we write one edge per ordered pair.

Let $f(x)\in\mathscr H_x$ be an actual decoded vector.  In the application to
Version 5 it is the direct sum of the retained ordered covariant spatial jets,
including all their $h^{-\ell}$ factors and physical spatial masses.  The
transition transport on this direct sum must be the actual induced one; it
cannot be chosen after observing the vectors to minimize their difference.
Define
\begin{align}
 e_{xy}&=\mathsf U_{xy}f(y)-f(x),&
 b_f(x)&=\sum_yq(x,y)e_{xy},\\
 \Gamma_f(x)&=\sum_yq(x,y)\|e_{xy}\|^2,&
 V_f(x)&=\|f(x)\|^2.
 \label{v6:eq:source-increments}
\end{align}

\begin{proposition}[Exact same-transition energy work]
\label[proposition]{v6:prop:energy-work}
Without stationarity or reversibility,
\begin{equation}
 \boxed{\mathcal L V_f(x)
 =2\Re\langle f(x),b_f(x)\rangle+\Gamma_f(x).}
 \label{v6:eq:energy-work}
\end{equation}
All terms are contractions of the initial, transported terminal, and mixed
initial--terminal Grams of the same transition.  In particular, the
conditional mean term alone does not determine the energy drift.
Moreover $\Gamma_f(x)=0$ if and only if $e_{xy}=0$ on every positive-rate
outgoing edge.  Arbitrarily frequent exact joint gauge changes therefore
produce no energy or temporal innovation for a covariantly unchanged reader.
\end{proposition}
\begin{proof}
Unitarity and expansion of one squared norm give
\[
 \|f(y)\|^2-\|f(x)\|^2
 =2\Re\langle f(x),e_{xy}\rangle+\|e_{xy}\|^2.
\]
Multiply by the actual rate and sum.  The zero criterion follows because
all summands of $\Gamma_f$ are nonnegative.  Exact joint gauge covariance
identifies the terminal jets with the initial ones under the supplied
transport, so every such $e_{xy}$ is zero.  This does not erase a transition
which changes a gauge-invariant amplitude or mixed path Gram.
\end{proof}

For example, an independently verified row bound
$2\Re\langle f,b_f\rangle+\Gamma_f\le C_0+C_1V_f$ implies
\[
 \mathbb E V_f(X_t)\le e^{C_1t}\bigl(\mathbb E V_f(X_0)+C_0t\bigr)
 \qquad(C_0,C_1\ge0)
\]
by differentiating the finite forward equation and integrating the scalar
inequality.  This is an explicit means of checking the occupation-energy
input of Version 1, not a proof that an unknown source satisfies that bound.
Exits must still be recorded as failures; killing mass is not negative
energy work.  The finite stopped-state convention of Version 2 applies
without changing this accounting.

\section{A finite covariant compiler for chronological correlations}
\label{v6:sec:semigroup}

Assume for this section a strictly positive reversible law $\pi$:
\begin{equation}
 \pi(x)q(x,y)=\pi(y)q(y,x),\qquad
 \mathsf U_{yx}=\mathsf U_{xy}^{*}.
 \label{v6:eq:reversibility}
\end{equation}
The connection is part of the source.  Put
\begin{equation}
 (\mathcal L^{\mathsf U}f)(x)
 =\sum_yq(x,y)(\mathsf U_{xy}f(y)-f(x)),\qquad
 \mathsf H=-\mathcal L^{\mathsf U},\qquad
 \langle f,g\rangle_\pi=\sum_x\pi(x)\langle f(x),g(x)\rangle.
 \label{v6:eq:connection-generator}
\end{equation}

\begin{lemma}[Positive source form and actual path transport]
\label[lemma]{v6:lem:semigroup}
The operator $\mathsf H$ is self-adjoint and nonnegative, and
\begin{equation}
 \langle f,\mathsf Hf\rangle_\pi
 =\frac12\sum_{x,y}\pi(x)q(x,y)
                 \|\mathsf U_{xy}f(y)-f(x)\|^2.
 \label{v6:eq:Dirichlet}
\end{equation}
Let $\mathsf U_{s,t}$ be the ordered product of all actual jump transports
from time $t$ back to time $s$, equal to the identity on a jump-free interval.
Then
\begin{equation}
 (e^{-t\mathsf H}f)(x)
 =\mathbb E_x[\mathsf U_{0,t}f(X_t)].
 \label{v6:eq:FK}
\end{equation}
For $X_0\sim\pi$, define the random vector
$Y_t=\mathsf U_{0,t}f(X_t)\in\mathscr H_{X_0}$.  For $s\le t$,
\begin{equation}
 \mathbb E\Re\langle Y_s,Y_t\rangle
 =\langle f,e^{-(t-s)\mathsf H}f\rangle_\pi.
 \label{v6:eq:covariance}
\end{equation}
The norm $\mathbb E\|Y_t\|^2$ equals $\|f\|_\pi^2$.  No replacement of
$Y_t$ by its conditional mean is made.
\end{lemma}
\begin{proof}
Pair the $x,y$ and $y,x$ terms using \eqref{v6:eq:reversibility}.  Expansion
of their squared differences gives \eqref{v6:eq:Dirichlet}; the corresponding
polarization proves self-adjointness.  Over a short interval of length $u$,
the probabilities of no jump and of one $x\to y$ jump are respectively
$1-u\sum_yq(x,y)+O(u^2)$ and $u q(x,y)+O(u^2)$.  The transported conditional
expectation consequently has generator $\mathcal L^{\mathsf U}$.
The Markov property and path concatenation give the semigroup property;
uniqueness of the finite matrix initial-value problem proves
\eqref{v6:eq:FK}.  In the inner product at two times, the common unitary
factor $\mathsf U_{0,s}$ cancels.  Condition at $X_s$, apply
\eqref{v6:eq:FK}, and use its law $\pi$ to obtain
\eqref{v6:eq:covariance}.  The equal-time norm assertion follows directly
from unitarity and stationarity.
\end{proof}

There is an exact nonstationary version.  If $\mu_t$ is the actual law of
$X_t$, then the right side of \eqref{v6:eq:covariance} is
\begin{equation}
 \sum_x\mu_s(x)\Re\langle f(x),(e^{(t-s)\mathcal L^{\mathsf U}}f)(x)\rangle.
 \label{v6:eq:nonstationary-correlation}
\end{equation}
The proof by conditioning is unchanged and does not require reversibility.
Reversibility is used later for the positive spectral comparison, not for
this finite correlation formula.  A non-Markov reader is permitted when
the finite source record retained in $\mathscr S$ is Markov-sufficient;
no claim that the reader by itself is Markov is needed.

\section{The exact temporal covariant-difference test}
\label{v6:sec:filters}

For $m\ge1$ set $c_j=(-1)^{m-j}\binom mj$ and define the actual chronological
covariant difference at time $t$ by
\begin{equation}
 \mathfrak D_{\delta,m}f(t)
 =\delta^{-m}\sum_{j=0}^m c_j\mathsf U_{t,t+j\delta}f(X_{t+j\delta}),
 \qquad
 \mathscr T_m(\delta;f)=\mathbb E_\pi\|\mathfrak D_{\delta,m}f(t)\|^2.
 \label{v6:eq:temporal-difference}
\end{equation}
The products in this formula are the actual links across the observation
interval, including all intermediate jumps.  It is precisely an ordered
pure-time word when this is the supplied temporal link reader.

\begin{lemma}[A positive finite-difference spectral polynomial]
\label[lemma]{v6:lem:polynomial}
For $0\le r\le1$ let
\begin{align}
 p_m(r)&=\sum_{a,b=0}^m c_ac_b r^{|a-b|}\\
 &=\binom{2m}{m}+2\sum_{\ell=1}^m(-1)^\ell
                                  \binom{2m}{m+\ell}r^\ell.
 \label{v6:eq:filter-polynomial}
\end{align}
Then
\begin{equation}
 \boxed{p_m(r)=2(1-r)\sum_{j=0}^{m-1}
             \binom{2m-2-j}{m-1}(1-r)^j.}
 \label{v6:eq:positive-polynomial}
\end{equation}
In particular, with $B_m=\binom{2m-2}{m-1}$,
\begin{equation}
 2B_m(1-r)\le p_m(r)\le\binom{2m}{m}(1-r).
 \label{v6:eq:sharp-filter-bounds}
\end{equation}
Both constants are sharp on this interval.
\end{lemma}
\begin{proof}
The coefficient at lag $\ell$ is
$\sum_{a=0}^{m-\ell}c_ac_{a+\ell}
=(-1)^\ell\binom{2m}{m+\ell}$ by the binomial convolution identity,
which proves \eqref{v6:eq:filter-polynomial}.
For completeness, a useful proof of the positive form uses only a geometric
series.  For $0\le r<1$ put $X(\theta)=2-2\cos\theta$.  The Fourier series
\[
 \frac{1-r^2}{1-2r\cos\theta+r^2}
 =\sum_{\ell\in\Z}r^{|\ell|}e^{\ii\ell\theta}
\]
gives
\[
 p_m(r)=\frac1{2\pi}\int_{-\pi}^{\pi}
 X(\theta)^m\frac{1-r^2}{(1-r)^2+rX(\theta)}\,\dd\theta.
\]
For $r>0$, polynomial division in $X$ therefore gives the recurrence
\begin{equation}
 r p_m(r)=(1-r^2)\binom{2m-2}{m-1}-(1-r)^2p_{m-1}(r),
 \quad p_0(r)=1.
 \label{v6:eq:polynomial-recurrence}
\end{equation}
The proposed expression has $p_1(r)=2(1-r)$.  To verify its induction step,
put $s=1-r$ and
$R_m(s)=\sum_{j=0}^{m-1}\binom{2m-2-j}{m-1}s^j$.
Pascal's identity, coefficient by coefficient, gives
\[
 (1-s)R_m(s)=(1-s/2)\binom{2m-2}{m-1}-s^2R_{m-1}(s)
 \qquad(m\ge2).
\]
This is \eqref{v6:eq:polynomial-recurrence} after division by $2s$.
Continuity includes $r=0,1$, proving \eqref{v6:eq:positive-polynomial}.
All coefficients of $R_m$ are positive.  Its minimum on $[0,1]$ is
$B_m$ and its maximum is
$\sum_{j=0}^{m-1}\binom{2m-2-j}{m-1}=\binom{2m-1}{m}$.
This proves the bounds, with equality at $r=0$ for the upper constant and
the limiting ratio at $r\uparrow1$ for the lower constant.
\end{proof}

\begin{theorem}[Sharp same-source temporal-energy criterion]
\label[theorem]{v6:thm:temporal-spectrum}
Under \eqref{v6:eq:reversibility},
\begin{equation}
 \boxed{\mathscr T_m(\delta;f)
 =\delta^{-2m}\langle f,p_m(e^{-\delta\mathsf H})f\rangle_\pi.}
 \label{v6:eq:exact-temporal-energy}
\end{equation}
Consequently
\begin{equation}
 \boxed{
 2B_m\delta^{-2m}\mathscr Q_\delta(f)
 \le\mathscr T_m(\delta;f)
 \le\binom{2m}{m}\delta^{-2m}\mathscr Q_\delta(f),\qquad
 \mathscr Q_\delta(f)=\langle f,(I-e^{-\delta\mathsf H})f\rangle_\pi.}
 \label{v6:eq:temporal-comparison}
\end{equation}
For a cutoff family and fixed $m$, bounded temporal $m$-difference energy
is therefore equivalent to $\mathscr Q_{\delta_h}(f_h)=O(\delta_h^{2m})$.
If $\mu_f$ is the finite spectral measure of $\mathsf H$ at $f$, this is
also equivalent, up to fixed constants, to
\begin{equation}
 \int \min\{\delta_h\lambda,1\}\,\dd\mu_{f_h}(\lambda)
                   =O(\delta_h^{2m}).
 \label{v6:eq:spectral-temporal-test}
\end{equation}
The assertions keep every source mode and every chronological outcome.
\end{theorem}
\begin{proof}
Expansion of the square in \eqref{v6:eq:temporal-difference}, followed by
\eqref{v6:eq:covariance}, gives the first equality.  Since the finite
self-adjoint operator $\mathsf H$ is nonnegative, its exponential has
spectrum in $[0,1]$.  Apply \eqref{v6:eq:sharp-filter-bounds} eigenvalue by
eigenvalue to get \eqref{v6:eq:temporal-comparison}.
Finally
\[
 (1-e^{-1})\min\{x,1\}\le1-e^{-x}\le\min\{x,1\}\quad(x\ge0),
\]
which follows from monotonicity and concavity on $[0,1]$ and the elementary
exponential bound outside it.  Integrate this inequality against $\mu_f$.
\end{proof}

\begin{corollary}[Visible innovation and a noncollapsing chronological rate]
\label[corollary]{v6:cor:innovation-scale}
If $\delta\|\mathsf H\|\le M$, then
\begin{equation}
 \mathscr T_m(\delta;f)\asymp_{m,M}
 \delta^{1-2m}\langle f,\mathsf H f\rangle_\pi.
 \label{v6:eq:small-clock}
\end{equation}
If, instead, only a spectral gap $\mathsf H\ge\lambda_*\Pi$ on the
orthogonal complement $\Pi$ of its kernel is known, then
\begin{equation}
 \mathscr T_m(\delta;f)\ge
 2B_m\delta^{-2m}(1-e^{-\delta\lambda_*})\|\Pi f\|_\pi^2.
 \label{v6:eq:gap-time}
\end{equation}
For fixed $\lambda_*>0$ and $\delta_h\downarrow0$, a uniform temporal
$m$-difference bound forces
$\|\Pi f_h\|_\pi^2=O(\delta_h^{2m-1})$.
\end{corollary}
\begin{proof}
For $0\le\delta\lambda\le M$,
$e^{-M}\delta\lambda\le1-e^{-\delta\lambda}\le\delta\lambda$.
Use \eqref{v6:eq:temporal-comparison}.  For the gap statement, use
$1-e^{-\delta\lambda}\ge1-e^{-\delta\lambda_*}$ on $\Ran\Pi$.
\end{proof}

The variance in \eqref{v6:eq:gap-time} is that of the \emph{specified
physical reader modulo its covariantly parallel kernel}.  A noncollapse
bound for a differently normalized source Gram cannot be substituted for it.
Nor does a reversible jump-source conclusion apply without proof to a causal
hyperbolic evolution or to a nonreversible deterministic-rate limit.

\begin{proposition}[A short-lag observation bank and its audit precision]
\label[proposition]{v6:prop:lag-audit}
The $m+1$ same-path quantities
$\rho_\ell=\langle f,e^{-\ell\delta\mathsf H}f\rangle_\pi$,
$0\le\ell\le m$, determine the exact temporal energy by
\begin{equation}
 \mathscr T_m=\delta^{-2m}
 \left[\binom{2m}{m}\rho_0+
 2\sum_{\ell=1}^m(-1)^\ell\binom{2m}{m+\ell}\rho_\ell\right].
 \label{v6:eq:lag-bank}
\end{equation}
If every supplied $\widehat\rho_\ell$ has absolute error at most $e$, the
resulting energy error is at most $4^m\delta^{-2m}e$.
No full successor table is required when these correlations are directly
provided by the source.  Separate initial and terminal marginals do not
provide them.
\end{proposition}
\begin{proof}
Use \eqref{v6:eq:filter-polynomial} in
\eqref{v6:eq:exact-temporal-energy}.  The sum of the absolute coefficients
is $\sum_{a,b}|c_ac_b|=(\sum_a\binom ma)^2=4^m$, proving the error bound.
The factors $\mathsf U_{0,\ell\delta}$ in
\eqref{v6:eq:covariance} retain the mixed path observation.
\end{proof}

\section{What can be hidden by transport, and what remains without equilibrium}
\label{v6:sec:radial-transient}

\begin{proposition}[Parallel noise and a transport-independent radial lower bound]
\label[proposition]{v6:prop:radial-time}
The kernel of $\mathsf H$ consists precisely of sections satisfying
$\mathsf U_{xy}f(y)=f(x)$ on every positive-rate edge.  Such sections have
$\mathfrak D_{\delta,m}f=0$ pathwise for all $m$ and $\delta$.
For general $f$ put $r(x)=\|f(x)\|$ and let $\mathsf H_0=-\mathcal L$
be the scalar reversible source operator.  Every transport satisfying
\eqref{v6:eq:reversibility} obeys
\begin{equation}
 \boxed{
 \mathscr T_m(\delta;f)\ge
 2B_m\delta^{-2m}\langle r,(I-e^{-\delta\mathsf H_0})r\rangle_\pi.}
 \label{v6:eq:radial-time}
\end{equation}
Thus a change of temporal gauge cannot suppress a nonzero rate of invariant
radial fluctuations.  For $m=1$, aligned pure-gauge transports attain the
scalar lower bound whenever such transports are admissible.  No optimality
claim for all higher $m$ is needed.
\end{proposition}
\begin{proof}
The kernel description follows from the sum of squares
\eqref{v6:eq:Dirichlet}.  Its edge identities telescope along every actual
path.  For the radial comparison, pointwise Cauchy--Schwarz gives
\[
 \Re\langle f(X_0),\mathsf U_{0,t}f(X_t)\rangle
                  \le r(X_0)r(X_t).
\]
Take expectations and subtract from the equal squared norms to get
$\mathscr Q_\delta(f)\ge
\langle r,(I-e^{-\delta\mathsf H_0})r\rangle_\pi$.
Now apply \eqref{v6:eq:temporal-comparison}.
For the last claim write $f(x)=g_xr(x)v$ with a fixed unit vector $v$ and
unitaries $g_x$; where allowed choose
$\mathsf U_{xy}=g_xg_y^*$.  The transported difference then equals the
scalar radius difference times $g_xv$.  This construction demonstrates the
bound's meaning; it does not replace the actual source transport.
\end{proof}

\begin{theorem}[Transient single-jump production of every temporal order]
\label[theorem]{v6:thm:single-jump}
For an arbitrary finite generator and actual unitary jump transports, let
$X_t$ have any prescribed law $\mu$.  Set
$q_{\max}=\max_x\sum_yq(x,y)$ and $M_f=\max_x\|f(x)\|$.
For fixed $m\ge1$,
\begin{equation}
 \left|\mathbb E_\mu
 \left\|\sum_{j=0}^m c_j\mathsf U_{0,j\delta}f(X_{j\delta})\right\|^2
 -B_m\delta\sum_x\mu(x)\Gamma_f(x)\right|
 \le C_m M_f^2 q_{\max}^2\delta^2.
 \label{v6:eq:one-jump}
\end{equation}
There is no equilibrium or reversibility premise.  The error retains its
rate dependence and must not be called uniform when
$q_{\max,h}\delta_h$ is large.
\end{theorem}
\begin{proof}
With no jump the unscaled difference is zero because $\sum c_j=0$.
If the only jump is $x\to y$ in $(\ell\delta,(\ell+1)\delta)$, its
unscaled value is
\[
 \left(\sum_{j=\ell+1}^m c_j\right)e_{xy}
 =(-1)^{m-1-\ell}\binom{m-1}{\ell}e_{xy}.
\]
The partial binomial sum follows by Pascal's identity.  The probability
of that jump and no others is
\[
 \mu(x)q(x,y)\int_{\ell\delta}^{(\ell+1)\delta}
 e^{-q_xs-q_y(m\delta-s)}\,\dd s,
 \qquad q_x=\sum_yq(x,y).
\]
Replacing its exponential by $1$ changes the integral by at most
$m q_{\max}\delta^2$.  Sum the squared coefficients and use
$\sum_{\ell=0}^{m-1}\binom{m-1}{\ell}^2=B_m$ and
$\|e_{xy}\|\le2M_f$; the total one-jump error is bounded by a constant
$C_mM_f^2q_{\max}^2\delta^2$.  Uniformization by a rate-$q_{\max}$
Poisson process bounds the probability of at least two actual jumps by
$(m q_{\max}\delta)^2/2$.  On this event the unscaled difference has norm
at most $2^mM_f$.  Adding its contribution proves
\eqref{v6:eq:one-jump}.  This is an estimate at the actual source rate,
not an identification of physical time with an action-comparison clock.
\end{proof}

\section{An exact temporal audit of the supplied phase-count reader}
\label{v6:sec:count}

Retain the explicit source of \cref{v3:thm:transient-product}, including its
allowed nonlinear swap modulation.  Write $n=N^3$, $h=N^{-1}$, and
\begin{equation}
 a=\frac{4\lambda_0}{5},\quad b=\frac{2\lambda_0}{3},\quad
 g=a+b=\frac{22\lambda_0}{15},\quad p=\frac6{11},\quad
 R_N(t)=\frac{K_N(t)}n.
 \label{v6:eq:actual-count}
\end{equation}
Here $\Delta_\delta^m u(t)=\sum_{j=0}^m
(-1)^{m-j}\binom mj u(t+j\delta)$ denotes the unnormalized forward
difference.  This is the actual count readout already supplied by \cite{GT}; no Higgs
semantic or Einstein action is assigned to it here.  It is spatially constant,
so it tests the independent temporal requirement even on a reader with no
positive-order spatial differences.  The count quotient, not the full field,
has generator $a(n-k)$ upward and $bk$ downward.  Swaps leave it unchanged.

\begin{theorem}[Exact count time-differences, with and without stationarity]
\label[theorem]{v6:thm:count-energy}
At the supplied stationary count law,
\begin{equation}
 \boxed{
 \mathbb E_\pi|\delta^{-m}\Delta_\delta^mR_N(t)|^2
 =\frac{p(1-p)}{n\delta^{2m}}p_m(e^{-g\delta}).}
 \label{v6:eq:stationary-count-energy}
\end{equation}
Starting instead from the deterministic all-$H$ configuration, put
$r(t)=p(1-e^{-gt})$ and $v(t)=r(t)(1-r(t))$.  For $s\le t$,
\begin{equation}
 \mathbb E R_N(t)=r(t),\qquad
 \operatorname{Cov}(R_N(s),R_N(t))=n^{-1}v(s)e^{-g(t-s)}.
 \label{v6:eq:count-two-time}
\end{equation}
Hence the exact nonstationary difference variance is
\begin{equation}
 \operatorname{Var}(\delta^{-m}\Delta_\delta^mR_N(t))
 =\frac1{n\delta^{2m}}\sum_{a,b=0}^m
 c_ac_b e^{-g|a-b|\delta}v(t+\min(a,b)\delta).
 \label{v6:eq:count-variance}
\end{equation}
For fixed $m,t$, uniformly in $N$ as $\delta\downarrow0$,
\begin{equation}
 \operatorname{Var}(\delta^{-m}\Delta_\delta^mR_N(t))
 =\frac{B_m\lambda(t)}{n\delta^{2m-1}}
       +O\!\left(\frac1{n\delta^{2m-2}}\right),
 \quad \lambda(t)=a(1-r(t))+br(t)>0.
 \label{v6:eq:count-variance-asymptotic}
\end{equation}
The deterministic mean difference tends to $r^{(m)}(t)$.
\end{theorem}
\begin{proof}
The inherited count generator gives the exact conditional mean
$\mathbb E[R_N(t+u)\mid K_N(t)]
=p+(R_N(t)-p)e^{-gu}$ by applying its finite backward equation to $k/n$.
The stationary variance is $p(1-p)/n$.  This proves its exponential
covariance, and \eqref{v6:eq:stationary-count-energy} follows by expanding
the difference.  From the all-$H$ initial condition, the inherited transient
product law gives mean $r(t)$ and variance $v(t)/n$.  Conditioning as above
proves \eqref{v6:eq:count-two-time} and then
\eqref{v6:eq:count-variance}.

Expand each factor in the finite sum to first order in $\delta$.  The
constant term cancels because $\sum c_j=0$.  The identities
\[
 \sum_{a,b}c_ac_b|a-b|=-2B_m,\qquad
 \sum_{a,b}c_ac_b\min(a,b)=B_m
\]
follow respectively by differentiating
\eqref{v6:eq:positive-polynomial} at $r=1$ and by
$2\min(a,b)=a+b-|a-b|$.  The linear coefficient is therefore
$B_m(v'(t)+2gv(t))=B_m[a(1-r(t))+br(t)]$.
The second derivatives of $v$ and the fixed exponential are bounded on
compact time intervals independently of $N$, proving the stated remainder.
The mean is a fixed analytic function, so its normalized finite differences
converge to its derivatives.  In particular this expansion does not use the
small-total-rate condition of \cref{v6:thm:single-jump}; the total count rate
can grow with $n$.
\end{proof}

For the explicit observation choice $\delta=h$, stationary formulas give
\begin{equation}
 \mathbb E_\pi|h^{-m}\Delta_h^mR_N|^2
 =\frac{8\lambda_0}{11}B_m h^{4-2m}(1+O(h)).
 \label{v6:eq:count-h-rate}
\end{equation}
The first three leading orders are thus
\begin{equation}
 \frac{8\lambda_0}{11}h^2,\qquad
 \frac{16\lambda_0}{11},\qquad
 \frac{48\lambda_0}{11}h^{-2}.
 \label{v6:eq:count-123}
\end{equation}
At a general spacing $\delta_h=h^\zeta$, the noise term has order
$h^{3-\zeta(2m-1)}$.  A physical-clock or reader change must be declared;
these powers cannot be changed by relabeling the accepted comparison count.

\begin{theorem}[Temporal roughness in probability from deterministic initial data]
\label[theorem]{v6:thm:count-CLT}
For the same all-$H$ initial source, every fixed $t\ge0$ and integer $m\ge1$
satisfy, at $\delta=h$,
\begin{equation}
 \boxed{
 h^{m-2}\left(h^{-m}\Delta_h^mR_N(t)
                   -h^{-m}\Delta_h^m r(t)\right)
 \ \Longrightarrow\ \mathcal N(0,B_m\lambda(t)).}
 \label{v6:eq:count-CLT}
\end{equation}
For $m\ge3$, $|h^{-m}\Delta_h^mR_N(t)|\to\infty$ in probability.
For $m=2$ the error about the analytic mean has a nondegenerate Gaussian
limit rather than convergence to zero.  This is not an assertion of temporal
independence.
\end{theorem}
\begin{proof}
The count process has the same path law as the sum of $n$ independent
$0\leftrightarrow1$ chains with rates $a,b$ started at zero.  Indeed that sum
has precisely the inherited birth--death generator and initial condition;
uniqueness of the finite Markov law identifies the count path, also in the
presence of the permitted count-preserving swaps.  This argument does not
assert independence of the original spatial sites at different times.

For one independent chain let
$Z_h=\sum_{j=0}^m c_j(\eta(t+jh)-r(t+jh))$.
Its mean is zero, $|Z_h|\le2^{m+1}$, and
\eqref{v6:eq:count-variance-asymptotic} at $n=1$ gives
$w_h=\mathbb E Z_h^2=B_m\lambda(t)h+O(h^2)$.
Consequently $\mathbb E|Z_h|^3\le2^{m+1}w_h$.
For fixed real $u$, Taylor's formula for the scalar exponential yields
\[
 \mathbb E\exp\!\left(\frac{\ii u Z_h}{\sqrt{n w_h}}\right)
 =1-\frac{u^2}{2n}
   +O_m\!\left(\frac{|u|^3}{n^{3/2}\sqrt{w_h}}\right).
\]
Raising this expression to the $n$th power gives $e^{-u^2/2}$, since
$n w_h\asymp N^3/N=N^2\to\infty$.  The continuity theorem for
characteristic functions proves convergence of the normalized sum to a
standard Gaussian.  Multiplying its variance normalization by
$h^{m-2}h^{-m}/n=h^{-2}/n$ gives limiting variance
$h^{-4}w_h/n\to B_m\lambda(t)$, proving
\eqref{v6:eq:count-CLT}.

For $m\ge3$ the centered scaled variable has a continuous nondegenerate
limit and the scaled deterministic mean tends to zero.  For any fixed $M$,
the event $|h^{-m}\Delta_h^mR_N|\le M$ forces its scaled value into an
interval of radius $Mh^{m-2}\to0$.  First bound by any fixed interval
about zero, use weak convergence, and then shrink that interval.  The
Gaussian has no atom at zero, so the probability tends to zero.  The $m=2$
statement follows directly from \eqref{v6:eq:count-CLT}.
\end{proof}

\section{Interpolation cannot manufacture the missing temporal jet bound}
\label{v6:sec:interpolation}

\begin{lemma}[A nodal divided difference costs actual derivative energy]
\label[lemma]{v6:lem:interpolation-cost}
If $y\in H^m([t,t+m\delta])$ has the indicated nodal values, then
\begin{equation}
 \boxed{
 \int_t^{t+m\delta}|y^{(m)}(s)|^2\,\dd s
 \ge\delta\left|\delta^{-m}\sum_{j=0}^m c_jy(t+j\delta)\right|^2.}
 \label{v6:eq:interpolation-cost}
\end{equation}
The inequality is valid pathwise for a randomly chosen interpolant as well.
\end{lemma}
\begin{proof}
Repeated fundamental theorem of calculus gives
\[
 \delta^{-m}\Delta_\delta^m y(t)
 =\delta^{-m}\int_{[0,\delta]^m}
         y^{(m)}(t+s_1+\cdots+s_m)\,\dd s_1\cdots\dd s_m.
\]
Apply Jensen's inequality.  The sum of $m$ independent uniform variables on
$[0,\delta]$ has density bounded by $1/\delta$: convolution with a
probability density does not increase its supremum.  The squared right side
is consequently at most
$\delta^{-1}\int_t^{t+m\delta}|y^{(m)}|^2$.
Approximation or the absolutely continuous Sobolev representatives proves the
formula for $H^m$ functions.  No assumption on how the interpolant was chosen
enters this deterministic inequality.
\end{proof}

\begin{corollary}[No node-preserving uniform temporal $H^3$ reconstruction]
\label[corollary]{v6:cor:interpolation-obstruction}
Fix an interior time $t$ and $m\ge3$.  Let $y_N$ be any, possibly
path-dependent, $H^m$ interpolation on a fixed containing interval which
retains the actual count samples $R_N(t+jh)$, $0\le j\le m$.
For the all-$H$ source of \cref{v6:thm:count-CLT},
\begin{equation}
 \|y_N^{(m)}\|_{L^2}^2\longrightarrow\infty
       \quad\hbox{in probability}.
 \label{v6:eq:interpolation-obstruction}
\end{equation}
This includes every exact smooth nodal interpolant, not just a particular
piecewise polynomial construction.
\end{corollary}
\begin{proof}
Put $Z_N=h^{m-2}h^{-m}\Delta_h^mR_N(t)$.
By \cref{v6:thm:count-CLT}, $Z_N$ converges in law to a nondegenerate
Gaussian.  Equation \eqref{v6:eq:interpolation-cost} gives
\[
 \|y_N^{(m)}\|_{L^2}^2\ge h^{5-2m}|Z_N|^2.
\]
For fixed $M$, a left side at most $M$ would require
$|Z_N|^2\le Mh^{2m-5}\to0$.  The same shrinking-interval argument as in
\cref{v6:thm:count-CLT} proves \eqref{v6:eq:interpolation-obstruction}.
\end{proof}

This result does \emph{not} exclude a deterministic limit of the count
values.  Indeed \eqref{v6:eq:count-two-time} gives variance $O(N^{-3})$
and a smooth mean.  It excludes the particular high-temporal-jet entrance
for an exact chronological count reader at the displayed scales.
The count has discarded spatial information and is not identified with an
Einstein--SM field.  Nor is the result a failure theorem for the distinct
locked spacetime source, whose chronological field/link law is not supplied
by this count process.  For a genuinely different temporal observation rule,
its actual samples and source normalization must be evaluated anew.

\section{What a source-supplied temporal response can and cannot repair}
\label{v6:sec:response}

There is a precise positive distinction between interpolating existing
samples and reading a physical response to them.  The latter is admissible
only if the response is independently part of the source.  We record its
regularity cost rather than silently adding it.

\begin{proposition}[Temporal integration order of a declared causal reader]
\label[proposition]{v6:prop:causal-reader}
Let an actually specified finite-dimensional reader satisfy, on $[0,T]$,
\begin{equation}
  (\partial_t+\kappa)^r y=x,\qquad\kappa>0,\quad r\ge1,
 \label{v6:eq:response}
\end{equation}
with declared initial derivatives through order $r-1$ and $x\in L^2$.
Then
\begin{equation}
 \|y\|_{H^r(0,T)}
 \le C_{r,\kappa,T}
  \left(\|x\|_{L^2(0,T)}+\sum_{j=0}^{r-1}|y^{(j)}(0)|\right).
 \label{v6:eq:response-bound}
\end{equation}
If $x$ is a finite-jump path with a nonzero interior jump, this same reader
is not in $H^{r+1}$ on a neighborhood of that jump.  Thus an order-$r$
response supplies order $r$ temporal regularity, not arbitrary smoothness.
\end{proposition}
\begin{proof}
Set $z_j=(\partial_t+\kappa)^j y$ for $0\le j<r$.
The equations $\dot z_j+\kappa z_j=z_{j+1}$, with $z_r=x$, have the
explicit causal solution
\[
 z_j(t)=e^{-\kappa t}z_j(0)
       +\int_0^te^{-\kappa(t-s)}z_{j+1}(s)\,\dd s.
\]
Cauchy--Schwarz on the fixed interval and the equation bound each $z_j$
in $H^1$, descending from $j=r-1$; equivalently the finite first-order
system yields $y\in H^r$ with \eqref{v6:eq:response-bound}.
The initial $z_j(0)$ are fixed linear combinations of the declared initial
derivatives.  Expanding \eqref{v6:eq:response} gives
\[
 y^{(r)}=x-\sum_{j<r}\binom rj\kappa^{r-j}y^{(j)}.
\]
Every term in the sum is continuous because $y\in H^r$ in one dimension.
Thus $y^{(r)}$ has the same nonzero jump as $x$.
An $H^{r+1}$ function would have $y^{(r)}\in H^1$, with an absolutely
continuous representative, a contradiction.
\end{proof}

The state $y$ in \eqref{v6:eq:response} is a different physical reader from
$x$, and need not agree with its nodes.  No common-action or field-semantic
identity between them is proved by \eqref{v6:eq:response-bound}.  A reader
which additionally has gauge transport must state that transport and the
covariant response equation; unobserved filtering is not justified by the
positive estimate.  The source audit must distinguish a genuine accumulated
response from a convenient reconstruction applied after the fact.

\section{Consequence for the same-source closure programme}
\label{v6:sec:ledger}

The occupation/drift problem left by Version 5 has a chronological component
which is now explicitly measurable.  On a known reversible finite source,
\eqref{v6:eq:exact-temporal-energy} or the $m+1$ correlation bank
\eqref{v6:eq:lag-bank} evaluates each pure-time word without a continuum
solution, analytic interpolation premise, or invented source law.  Applying
this to the direct sum of spatial covariant jets tests the corresponding
ordered temporal-prefix words, when their chronological transports are the
same supplied ones.  Other mixed orderings, color/determinant information,
and the geometric reader still require their own actual identities.

There are three mathematically different mechanisms.  Covariantly parallel
noise has zero visible jump defect and needs no temporal smoothing.  A
nonparallel reversible source can satisfy a high-time-jet budget only with
the weighted spectral suppression in \eqref{v6:eq:spectral-temporal-test};
an equilibrium covariance window for another source metric does not imply
that suppression.  A separately specified integrated reader can have the
regularity \eqref{v6:eq:response-bound}, but must be identified with the
physical field and common action independently.

The exact count calculation is a test of a source already present in the
monograph.  It demonstrates failure of a strong temporal entrance even after
spatial averaging has removed the previous spatial obstruction.  It does
not license a claim that every classical limit requires arbitrarily high
temporal smoothness.  In particular the abstract low-regularity variational
closure and the high-order quantitative reader criterion have different
hypotheses.  A stochastic-to-variational route with weaker temporal control
would need its own actual-source and stress proof; it is not obtained here
by replacing a norm with a conditional mean.

The next source-specific test should therefore identify the actual
chronological transition/transport bank, its visible jump form, and any
retained temporal-response variables, before postulating the high
spacetime-energy drift of Version 5.  For the intended coupled source, none
of these joint observations is manufactured by the present calculation.
The full physical action row, same-record stationarity, legal-domain control,
initial and harmonic alignment, and remaining physical sectors retain their
previous status.  No universal microscopic selection or Einstein--SM
continuum theorem has been added under stronger quantifiers.

\paragraph{Verification and preservation.}
The diagnostic script checks the exact filter polynomial with integer
arithmetic, the finite covariant generator identity and semigroup spectrum,
the radial comparison, count covariance and asymptotic constants, the
triangular-array characteristic functions, and the deterministic
interpolation inequality.  These checks are finite diagnostics, not formal
verification or an Einstein--SM simulation.  The full V5 text preceding its
final end-document command is preserved byte-for-byte.  All appended labels
have the prefix \texttt{v6:}.

\begingroup
\renewcommand{\refname}{Additional references for Version 6}
\begin{thebibliography}{9}
\small\raggedright
\bibitem{v6:GMT}
B.~G\"uneysu, O.~Milatovic, and F.~Truc,
\emph{Generalized Schr\"odinger semigroups on infinite graphs},
Potential Analysis \textbf{41} (2014), 517--541.
\href{https://arxiv.org/abs/1306.6062}{arXiv:1306.6062};
\href{https://doi.org/10.1007/s11118-013-9381-6}{doi:10.1007/s11118-013-9381-6}.
Background for covariant graph semigroups; the finite formulas used here are
proved directly.
\bibitem{v6:KL}
A.~Kassel and T.~L\'evy,
\emph{Covariant Symanzik identities}, Probability and Mathematical Physics
\textbf{2} (2021), 419--475.
\href{https://arxiv.org/abs/1607.05201}{arXiv:1607.05201};
\href{https://doi.org/10.2140/pmp.2021.2.419}{doi:10.2140/pmp.2021.2.419}.
Background for same-path holonomy and vector-bundle Markov formulas; no
Gaussian-field construction is substituted for the native reader.
\end{thebibliography}
\endgroup
% END IMMUTABLE VERSION 6 BODY
% APPEND VERSION 7 HERE, BEFORE THE FINAL END-DOCUMENT COMMAND.


% BEGIN IMMUTABLE VERSION 7 BODY
\clearpage
\begin{center}
{\Large\bfseries Version 7: First variations before high temporal jets}\\[0.5em]
{\large A sharp kinetic-noise threshold, an explicit stress defect, and a fixed-decoder dynamical test}
\end{center}
\addcontentsline{toc}{part}{Version 7: Kinetic variation and the native mean-action interface}

\section{Upstream audit and the change of analytic entrance}
\label{v7:sec:audit}

Version 6 proves that the actual phase-count samples, observed at spacing
$h=N^{-1}$, admit no node-preserving uniform temporal $H^3$ reconstruction.
That conclusion concerns the high-order entrance to
\cref{v1:in:closure}.  It does not settle the first-derivative topology in
which a minimal classical matter action and its metric variation are defined.
The present iteration evaluates this weaker question on the \emph{same}
count source.  Neither the transition rates, the count normalization, nor the
recorded nodal values are replaced by their conditional means.

\paragraph{Nonduplication audit.}
The count generator, path-law equivalence, two-time covariance, higher
finite-difference estimates, and the $H^3$ obstruction are inherited from
\cref{v6:thm:count-energy,v6:thm:count-CLT,v6:cor:interpolation-obstruction}.
The source-specific rates and lumpability are already in \cite{GT}, at
\nolinkurl{thm:renewal-interchange-actual-audit}.  The abstract quadratic
Higgs-defect and weak fermionic variation statements are already in
\cite{CC}; no such general closure theorem is reproved here.  In particular,
the source-to-action identification requirement at
\nolinkurl{v16:action} in that file is retained.  Searches of these inputs
and the preceding bodies located the moment and interpolation results just
listed, but not the integrated first-jet concentration and the joint
scalar/metric variational limit calculated below.

\paragraph{What is new in this body.}
We compute the exact conditional innovation of successive count samples and
prove a quantitative concentration formula for their \emph{integrated
kinetic energy}.  It has a sharp transition at $n\delta\asymp1$, where
$n=N^3$ and $\delta$ is the specified physical sampling interval.  At
$\delta=h$ a node-preserving affine reconstruction converges strongly in
$H^1$, despite the inherited high-jet obstruction.  At the critical sampling
scale a positive kinetic stress defect survives.  Every node-preserving
interpolant pays at least that kinetic cost.  Finally, the actual limiting
count motion is tested against a fixed minimal real-Higgs action and against
an independently given nonlinear scalar decoder.  Regularity and
variational stationarity turn out to be different obligations even on the
positive $H^1$ branch.

\paragraph{Attribution and scope.}
Fluid limits of density-dependent Markov chains and the theory of realized
quadratic variations are established subjects; see \cite{v7:Kurtz,v7:Jacod}.
No priority for those principles is asserted.  The finite innovation and
concentration estimates required here are proved directly, with both the
population and temporal-resolution parameters retained.  The count is an
inherited scalar reader, not an identified Einstein--SM field.  Evaluating a
fixed target action on it is a \emph{compatibility test}, not a new physical
identification or a change of its native bare phase action.  All theorems
below are proved locally in this body; none assumes the existence of a
coupled Einstein--SM limit.

\section{A first-jet identity on actual sampled records}
\label{v7:sec:innovation}

Let $\mathcal H_h$ be a real finite-dimensional Hilbert space with its
declared physical mass inner product.  Its dimension may grow.  Fix
$t_j=j\delta$, $0\le j\le J$, $J\delta=T$, and an adapted square-integrable
record $X_j$.  Its piecewise affine reconstruction $Y$ keeps every node:
$Y(t_j)=X_j$.  Write
\begin{equation}
 B_j=\delta^{-1}\E[X_{j+1}-X_j\mid\mathcal F_j],\qquad
 \xi_j=X_{j+1}-X_j-\delta B_j,\qquad
 Y'=B_j+\delta^{-1}\xi_j\quad(t_j<t<t_{j+1}).
 \label{v7:eq:innovation-split}
\end{equation}
The conditional expectation is on the actual history.  No Markov or
independence hypothesis is needed.  The interpolation is an explicit
reconstruction, not the original jump path; its physical use must be declared.

\begin{proposition}[Exact first-derivative error and its noise cost]
\label[proposition]{v7:prop:first-jet}
For any deterministic $v\in L^2(0,T;\mathcal H_h)$,
\begin{equation}
 \boxed{
 \E\int_0^T\|Y'-v\|^2\,\dd t
 =\sum_{j<J}\E\int_{t_j}^{t_{j+1}}\|B_j-v(t)\|^2\,\dd t
   +\delta^{-1}\sum_{j<J}\E\|\xi_j\|^2.}
 \label{v7:eq:first-jet-identity}
\end{equation}
Consequently, once the predictable velocities converge to $v$ in this
mean-square norm, strong first-derivative convergence holds if and only if
$\delta^{-1}\sum_j\E\|\xi_j\|^2\to0$.  If also the initial values converge,
then $Y$ converges in mean square in $H^1$ to the primitive of $v$ with that
initial value.  The constants are independent of $\dim\mathcal H_h$.
\end{proposition}
\begin{proof}
Expand the square on each interval.  The mixed term has expectation zero:
$B_j$ is $\mathcal F_j$-measurable, $v$ is deterministic, and
$\E[\xi_j\mid\mathcal F_j]=0$.  The squared innovation contributes
$\delta^{-1}\E\|\xi_j\|^2$.  Sum the exact identities.  The two terms on the
right are nonnegative, proving necessity and sufficiency under the stated
predictable convergence.  For $y(t)=y(0)+\int_0^t v(s)\dd s$,
\[
 \sup_{t\le T}\|Y(t)-y(t)\|^2
 \le2\|X_0-y(0)\|^2+2T\int_0^T\|Y'-v\|^2\dd t.
\]
Take expectations.  No finite-dimensional norm equivalence enters.
\end{proof}

The division by $\delta$ in \eqref{v7:eq:first-jet-identity} is a physical
\emph{kinetic} cost.  It is not the $\upsilon_h^{-2}$ cost of the
Bregman comparison discrepancy in Version 1.  Nor is the deterministic
$v$ a substitute for the reader: the entire innovation remains in $Y$.
For random comparison velocities one must keep the mixed conditional term
unless their appropriate interval averages are $\mathcal F_j$-measurable.

\section{The exact sampled innovation of the supplied count source}
\label{v7:sec:count}

Use the source and all-$H$ initial state of \cref{v6:sec:count}.  To keep
sampling and population visible independently, write
\begin{equation}
 X_n(t)=K_N(t)/n,\quad n=N^3,\quad
 a=4\lambda_0/5,\quad b=2\lambda_0/3,\quad
 c=a+b,\quad p=a/c=6/11,
 \label{v7:eq:count-source}
\end{equation}
where $\lambda_0>0$ is fixed.  The generator on $x\in\{0,1/n,\ldots,1\}$ is
\begin{equation}
 (\mathcal L_nf)(x)=na(1-x)[f(x+n^{-1})-f(x)]
                  +nbx[f(x-n^{-1})-f(x)].
 \label{v7:eq:count-generator}
\end{equation}
Let
\begin{equation}
 r(t)=p(1-e^{-ct}),\qquad
 \mathfrak a(x)=a(1-x)+bx,\qquad \mathfrak a_*(t)=\mathfrak a(r(t)).
 \label{v7:eq:activity}
\end{equation}
The same process is obtained in count law when the permitted nonlinear
spatial swaps are present.  We use that inherited lumpability; no
independence of the original spatial histories is asserted.

\begin{lemma}[A population-uniform path estimate]
\label[lemma]{v7:lem:path}
For each fixed $T$ there is $C_T$ independent of $n$ such that
\begin{equation}
 \E\sup_{0\le t\le T}|X_n(t)-r(t)|^2\le C_T/n.
 \label{v7:eq:path-bound}
\end{equation}
If $Y_{n,\delta}$ is the affine interpolation of the \emph{actual} samples
$X_n(t_j)$, $J\delta=T$, then
\begin{equation}
 \E\|Y_{n,\delta}-r\|_\infty^2\le C_T(n^{-1}+\delta^4).
 \label{v7:eq:value-bound}
\end{equation}
\end{lemma}
\begin{proof}
Applying the finite generator to $x$ and $x^2$ shows that
\[
 M_n(t)=X_n(t)-\int_0^t(a-cX_n(s))\dd s
\]
is a martingale with predictable quadratic variation
$\langle M_n\rangle_t=n^{-1}\int_0^t\mathfrak a(X_n(s))\dd s$.
These identities follow directly by conditioning each possible jump in an
infinitesimal interval; equivalently they follow from the finite forward
system.  Hence
\[
 X_n(t)-r(t)=e^{-ct}\int_0^t e^{cs}\dd M_n(s).
\]
The martingale isometry bounds the squared terminal stochastic integral by
$C_T/n$, since $\mathfrak a$ is bounded on $[0,1]$.  Doob's $L^2$ inequality
then gives \eqref{v7:eq:path-bound}.  For completeness, that inequality
follows by applying the stopping-time maximal inequality to the
nonnegative submartingale $|Z_t|$, integrating its level bound, and using
Cauchy--Schwarz to obtain $\E\sup|Z|^2\le4\E|Z_T|^2$.
The affine interpolant of the deterministic $C^2$ function $r$ differs
from $r$ by at most $\delta^2\|r''\|_\infty/8$.  The interpolation of the
nodal errors is bounded by their maximum.  Squaring proves
\eqref{v7:eq:value-bound}.
\end{proof}

Define $\theta_\delta=1-e^{-c\delta}$,
$\alpha_\delta=p\theta_\delta$, and
$\beta_\delta=(1-p)\theta_\delta$.  Given $X_n(t_j)=x=k/n$, the endpoint
count has the law of the sum of independent variables
\begin{equation}
 K_N(t_{j+1})\ \overset{\rm law}{=}\
 \operatorname{Bin}(n-k,\alpha_\delta)
 +\operatorname{Bin}(k,1-\beta_\delta).
 \label{v7:eq:binomial-row}
\end{equation}
This is an exact statement about the count transition: it follows from
its inherited representation as a sum of two-state chains.  It remains
valid conditional on the full past because the quotient is strongly
lumpable.

Put $x_j=X_n(t_j)$ and define
\begin{align}
 b_\delta(x)&=\frac{\theta_\delta}{\delta}(p-x),\qquad
 \xi_j=x_{j+1}-x_j-\delta b_\delta(x_j),\\
 q_\delta(x)&=\delta^{-1}\bigl[(1-x)\alpha_\delta(1-\alpha_\delta)
                            +x\beta_\delta(1-\beta_\delta)\bigr].
 \label{v7:eq:conditional-coefficients}
\end{align}

\begin{lemma}[Conditional moments with no small-total-rate requirement]
\label[lemma]{v7:lem:moments}
Uniformly for $0<\delta\le\delta_0$, $n\ge1$, and all current counts,
\begin{align}
 \E[\xi_j\mid\mathcal F_j]&=0,&
 \E[\xi_j^2\mid\mathcal F_j]&=\frac{\delta}{n}q_\delta(x_j),
 \label{v7:eq:conditional-variance}\\
 \sup_{0\le x\le1}|q_\delta(x)-\mathfrak a(x)|&\le C\delta,&
 \sup_{0\le x\le1}|b_\delta(x)-(a-cx)|&\le C\delta,
 \label{v7:eq:coefficient-errors}\\
 \E[\xi_j^4\mid\mathcal F_j]&\le
 C\left(\frac{\delta}{n^3}+\frac{\delta^2}{n^2}\right).
 \label{v7:eq:fourth-bound}
\end{align}
In particular, $q_\delta$ has a positive lower bound for sufficiently small
$\delta$.  These estimates do not require $n\delta\to0$.
\end{lemma}
\begin{proof}
Use \eqref{v7:eq:binomial-row}.  Its conditional mean gives $b_\delta$, and
summing the Bernoulli variances gives \eqref{v7:eq:conditional-variance}.
The Taylor formula $\theta_\delta=c\delta+O(\delta^2)$ gives
\eqref{v7:eq:coefficient-errors}; $\min(a,b)>0$ gives the lower bound.
For the fourth moment write $n\xi_j=\sum_{\ell=1}^n Z_\ell$ with
conditionally independent centered Bernoulli variables.  Each has
$\E Z_\ell^4\le\E Z_\ell^2\le C\delta$.  In expanding the fourth power,
only a single fourth moment or a pair of second moments survives.  Thus
\[
 \E\left(\sum_\ell Z_\ell\right)^4
 \le Cn\delta+3(Cn\delta)^2.
\]
Divide by $n^4$.  This proof uses the exact endpoint distribution, not an
approximation allowing at most one jump of the entire count.
\end{proof}

\section{A sharp integrated kinetic-energy law}
\label{v7:sec:kinetic}

Set
\begin{equation}
 \varepsilon_{n,\delta}=\frac1{n\delta},\qquad
 \mathcal K_{n,\delta}(\varphi)
 =\int_0^T\varphi(t)|Y_{n,\delta}'(t)|^2\dd t,
 \label{v7:eq:kinetic-def}
\end{equation}
for deterministic bounded measurable $\varphi$.  The reconstruction is
piecewise affine, so its derivative is defined almost everywhere.  No
classical second derivative is assigned at a knot.

\begin{theorem}[Quantitative concentration of actual kinetic variation]
\label[theorem]{v7:thm:kinetic}
For all $n$ and $0<\delta\le\delta_0$ with $T/\delta$ integral,
\begin{equation}
 \boxed{
 \begin{aligned}
 &\left\|\mathcal K_{n,\delta}(\varphi)
    -\int_0^T\varphi r'^2\dd t
    -\varepsilon_{n,\delta}\int_0^T\varphi\mathfrak a_*(t)\dd t
       \right\|_{L^2(\PP)}\\
 &\qquad\le C_T\|\varphi\|_\infty
 \left[\delta+n^{-1/2}
       +\varepsilon_{n,\delta}(\sqrt\delta+n^{-1/2})\right].
 \end{aligned}}
 \label{v7:eq:kinetic-concentration}
\end{equation}
Moreover
\begin{equation}
 \E\|Y_{n,\delta}-r\|_{H^1}^2
 \le C_T\left(n^{-1}+\delta^2+\varepsilon_{n,\delta}\right).
 \label{v7:eq:H1-rate}
\end{equation}
The constants depend on $a,b,T,\delta_0$ but not on $n,\delta$ or on
whether the original phase source has nonlinear count-preserving swaps.
\end{theorem}
\begin{proof}
Write $\bar\varphi_j=\delta^{-1}\int_{t_j}^{t_{j+1}}\varphi$ and use
$Y'=b_\delta(x_j)+\xi_j/\delta$.  The kinetic functional is the sum of
\begin{align*}
 A&=\sum_j\delta\bar\varphi_j b_\delta(x_j)^2,\\
 B&=2\sum_j\bar\varphi_j b_\delta(x_j)\xi_j,\\
 C&=\sum_j\bar\varphi_j\xi_j^2/\delta.
\end{align*}
The path estimate and the uniform coefficient bounds give
\[
 \|A-\textstyle\int\varphi r'^2\|_{L^2(\PP)}
 \le C_T\|\varphi\|_\infty(\delta+n^{-1/2}).
\]
The terms of $B$ are martingale differences with bounded predictable
coefficients.  Orthogonality at different indices and
\eqref{v7:eq:conditional-variance} yield
$\E|B|^2\le C_T\|\varphi\|_\infty^2/n$.
Decompose the final term as
\[
 C=Z+\varepsilon_{n,\delta}
           \sum_j\delta\bar\varphi_j q_\delta(x_j),\qquad
 Z=\sum_j\frac{\bar\varphi_j}{\delta}
       (\xi_j^2-\E[\xi_j^2\mid\mathcal F_j]).
\]
Again the centered summands are martingale differences.  The conditional
fourth-moment estimate proves
\begin{align*}
 \E|Z|^2
 &\le C_T\|\varphi\|_\infty^2
       \left(\frac1{n^3\delta^2}+\frac1{n^2\delta}\right)\\
 &=C_T\|\varphi\|_\infty^2\varepsilon_{n,\delta}^2(n^{-1}+\delta).
\end{align*}
The remaining predictable sum differs from
$\int\varphi\mathfrak a(r(t))\dd t$ in $L^2(\PP)$ by at most
$C_T\|\varphi\|_\infty(\delta+n^{-1/2})$.
This follows since $q_\delta$ is uniformly Lipschitz, its uniform error
from $\mathfrak a$ is $O(\delta)$, and $r$ is Lipschitz.  Using cell
averages means that no modulus of continuity of $\varphi$ is needed.
Combine the estimates, increasing constants to absorb $\delta\le C\sqrt\delta$.

For \eqref{v7:eq:H1-rate}, apply \cref{v7:prop:first-jet} to $v=r'$.
Its predictable term is at most $C_T(n^{-1}+\delta^2)$ by the path bound.
Its innovation term is
$\delta^{-1}\sum_j\E\xi_j^2\le C_T/(n\delta)$.
Combine with \eqref{v7:eq:value-bound}.
\end{proof}

\begin{corollary}[Three regimes of the same recorded source]
\label[corollary]{v7:cor:trichotomy}
Let $n\to\infty$, $\delta\to0$, and write $\varepsilon=1/(n\delta)$.
The values converge uniformly to $r$ in probability in every regime.
\begin{enumerate}[label=\textup{(\roman*)},leftmargin=2.7em]
\item If $\varepsilon\to0$, then $Y_{n,\delta}\to r$ in mean square
      in $H^1$, and $|Y'|^2\to r'^2$ in mean $L^1$.
\item If $\varepsilon\to\ell\in(0,\infty)$, then the derivatives
      converge weakly in $L^2$ in probability and
      \begin{equation}
        |Y'|^2\dd t\ \rightharpoonup\
            (r'^2+\ell\mathfrak a_*(t))\dd t
        \quad\hbox{in probability}.
        \label{v7:eq:critical-defect}
      \end{equation}
      In particular the first derivatives do not converge strongly.
\item If $\varepsilon\to\infty$, then for every bounded deterministic
      $\varphi$,
      \begin{equation}
       \varepsilon^{-1}\mathcal K_{n,\delta}(\varphi)
           \longrightarrow\int\varphi\mathfrak a_*(t)\dd t
           \quad\hbox{in }L^2(\PP).
       \label{v7:eq:divergent-kinetic}
      \end{equation}
      On any fixed nonempty time interval the kinetic energy diverges in
      probability.
\end{enumerate}
For $n=h^{-3}$ and $\delta=h$, the first regime holds with
$\E\|Y-r\|_{H^1}^2=O(h^2)$, simultaneously with the inherited failure of
uniform node-preserving $H^3$ bounds.
\end{corollary}
\begin{proof}
The value claim is \eqref{v7:eq:value-bound}.  The first case follows from
\eqref{v7:eq:H1-rate} and
$\||u|^2-|v|^2\|_1\le\|u-v\|_2(\|u\|_2+\|v\|_2)$ with
Cauchy--Schwarz in probability.  In the second case,
\eqref{v7:eq:kinetic-concentration} tends to zero.  Uniform boundedness of
the expected $L^2$ derivative norm follows from \eqref{v7:eq:H1-rate}.
Integration by parts against compactly supported smooth tests and the
value convergence give weak convergence of $Y'$ to $r'$; density and the
uniform second-moment bound extend this to every fixed $L^2$ test.
Here ``weakly in probability'' means convergence of each such pairing,
not an assertion of a particular cross-cutoff almost-sure coupling.
Taking $\varphi=1$ in \eqref{v7:eq:critical-defect} excludes strong
convergence since $\mathfrak a_*(t)>0$.
Divide \eqref{v7:eq:kinetic-concentration} by $\varepsilon$ in the third
case.  The deterministic drift part then vanishes and every error tends
to zero.  Finally $\varepsilon=h^2$ at $\delta=h$.
\end{proof}

\begin{proposition}[Minimum kinetic cost among all node-preserving readers]
\label[proposition]{v7:prop:minimum-interpolant}
Let $\widetilde Y\in H^1(0,T)$ be any, possibly path-dependent,
interpolant of the same samples.  Then, pathwise,
\begin{equation}
 \boxed{
 \int_0^T|\widetilde Y'|^2\dd t
 =\int_0^T|Y_{n,\delta}'|^2\dd t
       +\int_0^T|(\widetilde Y-Y_{n,\delta})'|^2\dd t.}
 \label{v7:eq:kinetic-Pythagoras}
\end{equation}
In the critical regime of \cref{v7:cor:trichotomy}, for every nonnegative
continuous $\varphi$ and every $\eta>0$,
\begin{equation}
 \PP\left\{\int\varphi|\widetilde Y'|^2\dd t
      <\int\varphi(r'^2+\ell\mathfrak a_*)\dd t-\eta\right\}\longrightarrow0.
 \label{v7:eq:weighted-minimum}
\end{equation}
There exists a node-preserving family converging strongly in $H^1$ in
probability to $r$ if and only if $n\delta\to\infty$.  In the divergent
regime every such interpolant has diverging first-derivative energy.
\end{proposition}
\begin{proof}
On a cell $Y_{n,\delta}'$ is constant and
$\int_{t_j}^{t_{j+1}}(\widetilde Y-Y_{n,\delta})'\dd t=0$.
The cross term in the squared norm therefore vanishes cell by cell,
proving \eqref{v7:eq:kinetic-Pythagoras}.
For a positive weight, cellwise Cauchy--Schwarz gives a lower bound by
$\sum_j(\inf_{[t_j,t_{j+1}]}\varphi)(x_{j+1}-x_j)^2/\delta$.
Its difference from $\mathcal K_{n,\delta}(\varphi)$ is bounded by the
modulus of continuity of $\varphi$ at $\delta$ times
$\mathcal K_{n,\delta}(1)$, which is tight in the critical regime.
This proves \eqref{v7:eq:weighted-minimum}.  Sufficiency of
$n\delta\to\infty$ follows from the affine construction.  If it fails,
a subsequence has $\varepsilon\ge\eta_0>0$, and a further subsequence
has a finite positive limit or tends to infinity.  The preceding lower
bounds contradict strong $H^1$ convergence in either case.
\end{proof}

This is a statement about the declared observation interval.  For
$\delta=h^\zeta$, the kinetic-noise scale is $h^{3-\zeta}$.
Coarser sampling is not silently adopted to improve a rate.  At the
original choice $\delta=h$, all actual nodal values are retained and the
positive first-variation conclusion requires no change of clock.

\section{The same samples tested against a fixed scalar and metric action}
\label{v7:sec:action}

The following calculation does not identify the count with a physical Higgs
field.  It tests that proposed identification against an action fixed
\emph{before} the stochastic record is evaluated.  Take a unit spatial
torus, homogeneous fields, identity internal links, zero spinors, and the
canonical real-Higgs coordinate $H=y e_1/\sqrt2$.  On the minimal quartic
branch put
\begin{equation}
 \mathcal V(y)=\lambda_H(y^2/2-v_H^2)^2,\qquad\lambda_H>0.
 \label{v7:eq:Higgs-potential}
\end{equation}
The scalar metric is written
$\dd s^2=-L(t)^2\dd t^2+s(t)^2\dd\mathbf x^2$.
On each temporal cell use the prescribed forward-difference restriction
of its first-derivative matter action,
\begin{equation}
 S_{n,\delta}(y,L,s)
 =\sum_{j<J}\delta\left\{
      \frac{s_j^3}{2L_j}\left(\frac{y_{j+1}-y_j}{\delta}\right)^2
                  -L_js_j^3\mathcal V(y_j)\right\}.
 \label{v7:eq:finite-scalar-action}
\end{equation}
The lapse and scale entries are independent test coordinates, not variables
determined from the count law.  The first variation is taken at $L_j=s_j=1$.
This is the scalar restriction of the indicated flat-chart local density,
not the bare product phase action $-\sum_x\log\pi_{\eta_x}$ of \cite{GT}.
All claims about coupling it to the source remain conditional on this
reader/action identification.

For $v_j=(y_{j+1}-y_j)/\delta$, its exact rows are
\begin{align}
 D_y S[\zeta]&=\sum_j\delta[v_jD_\delta\zeta_j
                                    -\mathcal V'(y_j)\zeta_j],
 \label{v7:eq:scalar-variation}\\
 D_L S[\ell]&=-\sum_j\delta(\tfrac12v_j^2+\mathcal V(y_j))\ell_j,
 \label{v7:eq:lapse-variation}\\
 D_s S[b]&=3\sum_j\delta(\tfrac12v_j^2-\mathcal V(y_j))b_j.
 \label{v7:eq:scale-variation}
\end{align}
Here $b_j$ in \eqref{v7:eq:scale-variation} denotes the scale test only,
not the predictable coefficient $b_\delta$.

\begin{theorem}[Scalar Euler convergence and metric stress have different thresholds]
\label[theorem]{v7:thm:action-limit}
Evaluate \eqref{v7:eq:finite-scalar-action} at $y_j=X_n(t_j)$.
For every fixed compactly supported smooth scalar test $\zeta$, as
$n\to\infty$ and $\delta\to0$, with no restriction on $n\delta$,
\begin{equation}
 D_yS_{n,\delta}[\zeta(t_j)]
 \longrightarrow
 \int_0^T[r'\zeta'-\mathcal V'(r)\zeta]\dd t
 =-\int_0^T[r''+\mathcal V'(r)]\zeta\dd t
 \quad\hbox{in }L^2(\PP).
 \label{v7:eq:scalar-limit}
\end{equation}
If $\varepsilon=1/(n\delta)\to\ell\in[0,\infty)$, the lapse and scale
rows converge, against fixed smooth tests, to those determined by
\begin{equation}
 \boxed{
 \rho_\ell(t)=\tfrac12r'^2+\mathcal V(r)+\tfrac\ell2\mathfrak a_*(t),
 \qquad
 P_\ell(t)=\tfrac12r'^2-\mathcal V(r)+\tfrac\ell2\mathfrak a_*(t).}
 \label{v7:eq:stress-limit}
\end{equation}
In particular, at $\delta=h$, the first-derivative action, its scalar
variation, and its matter metric variations converge without any uniform
temporal $H^3$ bound.  At finite positive $\ell$, the additional flat-chart
stress tensor is
\begin{equation}
 \mathfrak S_{\mu\nu}(t)
 =\ell\mathfrak a_*(t)
   \left(n_\mu n_\nu+\tfrac12\eta_{\mu\nu}\right),\qquad
 n_\mu=(-1,0,0,0),\quad\eta=\operatorname{diag}(-1,1,1,1).
 \label{v7:eq:timelike-defect}
\end{equation}
It has equal positive energy density and isotropic pressure, not a
cosmological-constant form.  Its conservation is not asserted on these
off-shell records.
\end{theorem}
\begin{proof}
The exact nodal derivatives give
\eqref{v7:eq:scalar-variation}--\eqref{v7:eq:scale-variation}.
For compactly supported $\zeta$ and sufficiently small $\delta$, discrete
summation by parts in the kinetic scalar row gives
\[
 \sum_j\delta v_jD_\delta\zeta_j
 =-\sum_{j=1}^{J-1}\delta x_j
       \frac{\zeta(t_{j+1})-2\zeta(t_j)+\zeta(t_{j-1})}{\delta^2}.
\]
The boundary terms vanish because the test vanishes on fixed endpoint
collars.  The central quotient converges uniformly to $\zeta''$; the
uniform value estimate \eqref{v7:eq:path-bound} and deterministic
quadrature therefore prove convergence to $-\int r\zeta''$ in
$L^2(\PP)$.  The potential row converges by its Lipschitz bound on $[0,1]$.
Integration by parts of the deterministic limit proves
\eqref{v7:eq:scalar-limit}, even when the stochastic kinetic energy diverges.

For the metric rows, regard each nodal smooth test as a piecewise constant
weight.  Its uniform error from the continuous test is $O(\delta)$.
When $\varepsilon$ is bounded the kinetic total has bounded $L^2(\PP)$
by \cref{v7:thm:kinetic}; this replacement therefore has vanishing error.
That theorem and convergence of the potential sums give
\eqref{v7:eq:stress-limit}.  Isotropy and spatial homogeneity give zero
mixed components and the tensor \eqref{v7:eq:timelike-defect}.
The action values themselves converge to
$\int(\tfrac12r'^2-\mathcal V(r))\dd t+\tfrac\ell2\int\mathfrak a_*\dd t$.
For $\delta=h$, \cref{v7:cor:trichotomy} has $\ell=0$.
\end{proof}

\begin{proposition}[A dual first-variation estimate on the defect-free branch]
\label[proposition]{v7:prop:dual}
Let $r^\delta_j=r(t_j)$ and give scalar nodal tests vanishing at the two
endpoints the norm
\[
 \|\zeta\|_{1,\delta}^2
 =\sum_{j<J}\delta\bigl(|\zeta_j|^2+|D_\delta\zeta_j|^2\bigr).
\]
Then
\begin{equation}
 \left\|\sup_{\|\zeta\|_{1,\delta}\le1}
  |D_yS(x)[\zeta]-D_yS(r^\delta)[\zeta]|\right\|_{L^2(\PP)}
 \le C_T\bigl(n^{-1/2}+\delta+(n\delta)^{-1/2}\bigr).
 \label{v7:eq:dual-rate}
\end{equation}
For the same records the expected $L^1$ difference of the two kinetic
energy densities is bounded by
$C_T e_{n,\delta}(1+e_{n,\delta})$, where
$e_{n,\delta}=n^{-1/2}+\delta+(n\delta)^{-1/2}$.
\end{proposition}
\begin{proof}
Subtract \eqref{v7:eq:scalar-variation}, use Cauchy--Schwarz on the
kinetic term and the uniform Lipschitz bound for $\mathcal V'$ on $[0,1]$
on the potential term.  The supremum is bounded by the discrete
first-derivative and value error.  Equations
\eqref{v7:eq:H1-rate} and \eqref{v7:eq:value-bound}, and the deterministic
$O(\delta)$ derivative interpolation error, prove the estimate.
For the quadratic density use
$\||u|^2-|v|^2\|_1\le\|u-v\|_2(2\|v\|_2+\|u-v\|_2)$ and take
expectations.  This is a dual first-variation norm, not the stronger
nodal mass norm of the second-order Euler row.
\end{proof}

\paragraph{Action values do not define a new effective Euler equation.}
At critical resolution the extra scalar action value is
$\tfrac\ell2\int\mathfrak a(r(t))\dd t$, but
\eqref{v7:eq:scalar-limit} has no derivative of that term.  There is no
contradiction.  The kinetic concentration formula is proved along one
specified stochastic family, not as a locally differentiable identity on a
neighborhood of arbitrary varied field histories.  In the finite first
variation the source law is held fixed and the nodal values are varied
independently.  Differentiating an inferred family-dependent action value
would interchange operations not justified by the theorem.  Likewise, a
field-independent subtraction of that value does not remove the metric
source \eqref{v7:eq:timelike-defect}.

\section{Go further upstream: the fixed decoder must match the action dynamics}
\label{v7:sec:compatibility}

The positive $H^1$ result does not make the count solve a Higgs equation.
The following test isolates that issue without introducing a new transition
law or fitting a potential to a realized trajectory.

\begin{proposition}[The direct count reader is off shell for the fixed quartic]
\label[proposition]{v7:prop:off-shell}
For the all-$H$ source, the direct scalar assignment $y=r(t)$ does not
solve the flat homogeneous Euler equation of
\eqref{v7:eq:Higgs-potential} on any nonempty open time interval.
There is a fixed compactly supported smooth test $\zeta$ on that interval
and a number $c_\zeta>0$ such that
\begin{equation}
 \PP\{|D_yS_{n,\delta}[\zeta(t_j)]|\ge c_\zeta\}\longrightarrow1.
 \label{v7:eq:off-shell-test}
\end{equation}
Thus even the defect-free $\delta=h$ branch does not give vanishing
physical scalar variation for this target action.
\end{proposition}
\begin{proof}
The mean obeys $r'=c(p-r)$ and $r''=c^2(r-p)$, while
$\mathcal V'(r)=\lambda_H r(r^2-2v_H^2)$.
Consequently its residual is the polynomial
\begin{equation}
 r''+\mathcal V'(r)
 =\lambda_H r^3+(c^2-2\lambda_Hv_H^2)r-c^2p.
 \label{v7:eq:quartic-residual}
\end{equation}
It cannot vanish identically on an interval of $r$ values because its
cubic coefficient is positive and its constant term is nonzero.  The
nonconstant analytic function $r(t)$ takes an interval of values on every
open time interval.  A smaller interval therefore has a residual of one
strict sign.  Choose a nonnegative smooth bump supported there.
Equation \eqref{v7:eq:scalar-limit} has a nonzero limit on that test;
convergence in probability gives \eqref{v7:eq:off-shell-test} after taking
$c_\zeta$ smaller than its absolute limiting value.
\end{proof}

This assertion is restricted to the stated flat metric, identity links,
zero spinors, and direct scalar assignment.  A backreacting metric supplies
additional terms, including the cosmological friction term in a homogeneous
chart.  It is not legitimate to infer its absence or presence from the count
law.  Constant equilibrium data are also different: a constant scalar can
be critical if the independently calibrated potential has an extremum at
that value.  No such special tuning is used above.

\begin{theorem}[A fixed nonlinear decoder transports both the defect and the compatibility test]
\label[theorem]{v7:thm:decoder}
Let $\Phi\in C^3$ on a neighborhood of $[0,1]$ be independently specified,
with bounded derivatives, and retain the actual decoded samples
$z_j=\Phi(x_j)$.  Let $Z_{n,\delta}$ be their affine interpolation and
$z(t)=\Phi(r(t))$.  If $\varepsilon\to\ell\in[0,\infty)$, then
\begin{equation}
 |Z_{n,\delta}'|^2\dd t\ \rightharpoonup\
 \bigl(z'^2+\ell\Phi'(r)^2\mathfrak a_*(t)\bigr)\dd t
 \quad\hbox{in probability}.
 \label{v7:eq:decoded-defect}
\end{equation}
The values converge uniformly to $z$ in probability.  At $\ell=0$,
$Z_{n,\delta}\to z$ strongly in $H^1$ in probability.  At $\ell>0$, any
nonvanishing $\Phi'$ on a time interval retains a positive kinetic defect
there.  In particular a fixed nondegenerate scalar change of coordinates
does not erase the critical kinetic noise.

For a fixed $C^2$ potential $V$ on the decoded range, the weak scalar
Euler limit vanishes on a nonconstant interval if and only if
\begin{equation}
 \boxed{
 c^2(p-r)^2\Phi''(r)-c^2(p-r)\Phi'(r)+V'(\Phi(r))=0
 }
 \label{v7:eq:decoder-compatibility}
\end{equation}
throughout the corresponding interval of count values.  If this identity
holds up to $r=p$, $\Phi'(p)\ne0$, and $V$ is $C^2$ there, then necessarily
\begin{equation}
 V'(\Phi(p))=0,\qquad V''(\Phi(p))=-c^2.
 \label{v7:eq:equilibrium-compatibility}
\end{equation}
Thus no such nondegenerate smooth decoder identifies the relaxing count
trajectory with a nonconstant flat conservative scalar trajectory
approaching a strict nondegenerate potential minimum.
\end{theorem}
\begin{proof}
The uniform value estimate and bounded $\Phi'$ give convergence of the
decoded values and their affine interpolants.  Define the divided derivative
\[
 A_j=\int_0^1\Phi'(x_j+u(x_{j+1}-x_j))\dd u.
\]
Exactly $D_\delta z_j=A_jD_\delta x_j$.  The uniform path estimate and
Lipschitz continuity of $r$ imply
$\sup_j|A_j-\Phi'(r(t_j))|\to0$ in probability.  The original kinetic
total is tight when $\varepsilon$ is bounded.  Hence, against a bounded
continuous test, replacing $A_j^2$ by $\Phi'(r(t_j))^2$ in the kinetic
sum has vanishing error in probability.  Apply
\cref{v7:thm:kinetic} with the deterministic weight
$\varphi(t)\Phi'(r(t))^2$ to obtain \eqref{v7:eq:decoded-defect}.
At $\ell=0$, use the first-derivative convergence of the count and the
same divided-derivative estimate to prove strong $H^1$ convergence in
probability.  Positivity of $\mathfrak a$ proves the retained defect claim.

Discrete summation by parts as in \cref{v7:thm:action-limit} gives the
scalar Euler limit $-\int(z''+V'(z))\zeta$ from value convergence alone.
The chain rule gives
$z''=\Phi''(r)c^2(p-r)^2-\Phi'(r)c^2(p-r)$, proving
\eqref{v7:eq:decoder-compatibility}.  At $r=p$ the first two terms vanish,
so $V'(\Phi(p))=0$.  Write $r=p+u$.  The first term is $O(u^2)$,
the second is $c^2u\Phi'(p)+o(u)$, and the potential term is
$V''(\Phi(p))\Phi'(p)u+o(u)$.  Divide by $u$ and let $u\to0$ to get
$(c^2+V''(\Phi(p)))\Phi'(p)=0$.  Nondegeneracy proves
\eqref{v7:eq:equilibrium-compatibility}.  A positive $V''$ is incompatible
with that necessary identity.
\end{proof}

Equation \eqref{v7:eq:decoder-compatibility} is an audit of a \emph{fixed}
source-to-field map and a \emph{fixed} action.  Solving it afterwards for
$\Phi$ or changing $V$ can construct another model, but does not verify
the intended native reader.  The obstruction at the equilibrium does not
apply to general dissipative matter equations, to a nontrivial expanding
Einstein geometry, or to the distinct locked spacetime source without its
joint decoder being supplied.

\section{What this removes, and the remaining source-level task}
\label{v7:sec:ledger}

The high-time-jet obstruction in Version 6 is not a necessary obstruction to
classical variational passage.  At its own sampling scale $\delta=h$,
the same count nodes have a defect-free affine first-derivative limit with
root-mean-square $H^1$ error $O(h)$.  This is a positive exact-source
calculation, not a new sampler or a filtered conditional-mean field.
At finer critical resolution $n\delta\asymp1$ the same source supplies a
nonzero timelike kinetic stress, even though all fixed scalar Euler tests
have the smooth mean-field limit.  The optimal interpolation identity shows
that this extra kinetic cost cannot be removed by a more elaborate
node-preserving interpolation.

The main quantities are now distinct:
\begin{center}\small
\begin{tabularx}{\textwidth}{@{}L{0.29\textwidth}X@{}}
\toprule
Question & Quantity actually tested here\\
\midrule
Uniform values & Native count martingale, with $O(n^{-1})$ maximal mean-square error.\\[0.3em]
First-derivative convergence & Predictable velocity error plus
$\delta^{-1}\sum_j\E|\xi_j|^2$, not a high-jet bound.\\[0.3em]
Metric stress & Concentrated kinetic variation with exact defect scale $(n\delta)^{-1}$.\\[0.3em]
Scalar stationarity & The independently fixed decoder/action identity
\eqref{v7:eq:decoder-compatibility}; regularity does not supply it.\\[0.3em]
Full Einstein--SM closure & Still requires the physical geometry, all other sectors,
common-action identification, compatible variations, and their actual source estimates.\\
\bottomrule
\end{tabularx}
\end{center}

For a supplied native reader whose first differences have small realized
kinetic innovations, the low-regularity variational route may be more
appropriate than forcing its exact values into the high-temporal-jet route.
The present results do not provide the spatial covariant regularity or the
metric compactness of such a coupled reader.  They do remove a false
necessity: failure of the high-order entrance alone is not a reason to
abandon that first-derivative route.

The action mismatch is then exposed rather than hidden.  The explicit count
process is regular enough at $\delta=h$, but its relaxing mean is not a
solution of the fixed flat quartic scalar action.  Its bare product action,
its count decoder, and the physical Einstein--SM action remain different
objects.  The next useful upstream observation is therefore the joint
\emph{field and conjugate-momentum} reader, or the corresponding identified
first-order action flow, for the intended source.  Its predictable velocity
must match the independently fixed physical equations, and its innovation
budget must be paid in the topology of the actual action.  Another generic
$H^3$ estimate, a fitted potential, or a relabelled comparison clock would
not establish that identity.

\paragraph{Verification and preservation.}
The diagnostic script verifies the exact binomial count row against a finite
birth--death matrix exponential, the conditional second/fourth moments, the
kinetic expectation from two-time covariance, the interpolation
Pythagorean identity, and the independent scalar/lapse/scale derivatives.
It also samples exact endpoint count transitions in the three kinetic
regimes.  The simulations are diagnostics; the concentration theorem is
proved by the conditional fourth-moment estimate and martingale
orthogonality, not inferred from those simulations.  The complete V6 source
preceding its final end-document command is retained verbatim.  New labels
have prefix \texttt{v7:}.  There is no machine-checked proof or simulation
of the full Einstein--SM system.

\begingroup
\renewcommand{\refname}{Additional references for Version 7}
\begin{thebibliography}{9}
\small\raggedright
\bibitem{v7:Kurtz}
T.~G.~Kurtz,
\emph{Solutions of ordinary differential equations as limits of pure jump
Markov processes}, Journal of Applied Probability \textbf{7} (1970), 49--58.
\href{https://doi.org/10.2307/3212147}{doi:10.2307/3212147}.
Background for the deterministic limit of density-dependent chains; the
count path and sampling estimates needed here are proved directly.
\bibitem{v7:Jacod}
J.~Jacod,
\emph{Asymptotic properties of realized power variations and related
functionals of semimartingales}, 2006 preprint.
\href{https://arxiv.org/abs/math/0604450}{arXiv:math/0604450}.
Background for quadratic-variation asymptotics; the two-parameter count
calculation and its physical first-variation test above are separate.
\end{thebibliography}
\endgroup
% END IMMUTABLE VERSION 7 BODY
% APPEND VERSION 8 HERE, BEFORE THE FINAL END-DOCUMENT COMMAND.



% BEGIN IMMUTABLE VERSION 8 BODY
\clearpage
\begin{center}
{\Large\bfseries Version 8: The canonical pair before second-order equations}\\[0.5em]
{\large A source-generator Hamilton test, Legendre incidence, and a phase-space obstruction}
\end{center}
\addcontentsline{toc}{part}{Version 8: Native canonical-pair dynamics and Hamilton compatibility}

\section{Upstream audit: what Legendre reciprocity does and does not give}
\label{v8:sec:audit}

Version 7 shows that an actual source may pass a low-regularity first-variation
limit while its deterministic mean still fails the independently fixed scalar
Euler equation.  The next upstream object is therefore not another temporal
Sobolev norm.  It is the joint field--momentum record and the first-order flow
that the accepted source assigns to it.

\paragraph{Inherited structural inputs.}
The Grand Tensor monograph \cite{GT}, at
\nolinkurl{thm:SMST-matter-Legendre}, already proves finite Legendre
reciprocity for a matter Lagrangian with positive velocity Hessian.  In its
notation
\[
 p=D_v\mathscr L(z,v),\qquad
 D_p\mathscr H=v,\qquad D_z\mathscr H=-D_z\mathscr L,
 \qquad D_p^2\mathscr H=(D_v^2\mathscr L)^{-1}.
\]
The same theorem identifies the finite typed matter bracket and its reciprocal
kinetic coefficients.  We use these facts and do not reprove them.  Likewise,
\cite{GT}, at \nolinkurl{thm:SMST-positive-screen} and
\nolinkurl{cor:SMST-regenerative-screen}, already gives the positive-screen
mechanism for the metric-velocity, matter-normal, and matter-gradient packet
when its hypotheses hold.  Those are compactness/identification statements,
not a theorem that an accepted chronological transition obeys Hamilton's
equations.  The conditional Einstein--matter handoff in \cite{GT} explicitly
retains an actual-cycle/action defect as a separate row.

The present iteration isolates that missing chronological row.  It gives a
finite-table residual which can be evaluated directly from one accepted
transition law and one independently fixed Hamiltonian.  A martingale estimate
then passes this residual to the continuum first variation.  No Fourier--midpoint
writer, fitted potential, or new probability law is introduced.

\paragraph{Scope.}
All phase spaces below are finite-dimensional real Hilbert spaces, possibly
after the already declared source alignment to a fixed physical comparison
bank.  The Hamiltonian can depend on the cutoff, but every convergence theorem
states the required uniform chart and Lipschitz bounds.  Gauge and constraint
reductions are not silently included: the canonical pair must be the same
physical pair on which the finite common action and its Legendre map are
represented.

\section{The native generator residual of a canonical pair}
\label{v8:sec:generator}

Let $\Omega_h$ be a finite state space with continuous-time jump rates
$q_h(x,y)$ and generator
\[
 (\mathcal L_h f)(x)=\sum_yq_h(x,y)(f(y)-f(x)).
\]
Let $\Phi_h:\Omega_h\to E$ be the \emph{declared physical canonical-pair
reader} into one real Hilbert space $E$.  Put $Z_h(t)=\Phi_h(X_h(t))$ for the
actual accepted source path.  Let $\mathbb J:E\to E$ be a fixed skew-adjoint
isometry, $\mathbb J^*=-\mathbb J$, $\mathbb J^2=-I$.  For a $C^2$ Hamiltonian
$\mathscr H_h$ on a common compact chart set
\begin{equation}
 X_{\mathscr H_h}(z)=\mathbb J\nabla\mathscr H_h(z).
 \label{v8:eq:H-vector}
\end{equation}
Define the source-readable drift and visible jump form
\begin{align}
 b_h(x)&=(\mathcal L_h\Phi_h)(x)
       =\sum_yq_h(x,y)(\Phi_h(y)-\Phi_h(x)),\\
 \Gamma_h(x)&=\sum_yq_h(x,y)\|\Phi_h(y)-\Phi_h(x)\|^2,\\
 \mathfrak r_h(x)&=b_h(x)-X_{\mathscr H_h}(\Phi_h(x)).
 \label{v8:eq:H-residual}
\end{align}
Every quantity is a same-history finite transition/action quantity.  In
particular $\mathfrak r_h$ is not obtained by differentiating a prospective
continuum solution.

\begin{theorem}[Finite-source Hamilton residual and pathwise comparison]
\label[theorem]{v8:thm:generator-handoff}
Let $z_h$ solve the deterministic Hamilton equation
\begin{equation}
 \dot z_h=X_{\mathscr H_h}(z_h),\qquad z_h(0)=z_{h,0},
 \label{v8:eq:reference-H}
\end{equation}
on $[0,T]$.  Suppose all decoded and reference paths stay in a common compact
chart on which $X_{\mathscr H_h}$ is $L$-Lipschitz, uniformly in $h$.  Then
there is $C_{L,T}$ independent of the number of source states such that
\begin{equation}
\boxed{
 \E\sup_{t\le T}\|Z_h(t)-z_h(t)\|^2
 \le C_{L,T}\left[
   \E\|Z_h(0)-z_{h,0}\|^2
  +\E\int_0^T\|\mathfrak r_h(X_h(t))\|^2\dd t
  +\E\int_0^T\Gamma_h(X_h(t))\dd t
 \right].}
 \label{v8:eq:generator-comparison}
\end{equation}
Moreover, for every deterministic $\psi\in C_c^1((0,T);E)$, the weak
Hamilton residual
\begin{equation}
 \mathcal R_h(\psi)
 :=-\int_0^T\ip{Z_h}{\dot\psi}\dd t
    -\int_0^T\ip{X_{\mathscr H_h}(Z_h)}{\psi}\dd t
 \label{v8:eq:weak-residual}
\end{equation}
satisfies
\begin{equation}
\boxed{
 \|\mathcal R_h(\psi)\|_{L^2(\PP)}
 \le \|\psi\|_{L^2_t}
       \left(\E\int_0^T\|\mathfrak r_h(X_h(t))\|^2\dd t\right)^{1/2}
   +\|\psi\|_{L^\infty_t}
       \left(\E\int_0^T\Gamma_h(X_h(t))\dd t\right)^{1/2}.}
 \label{v8:eq:weak-residual-bound}
\end{equation}
Thus a vanishing generator residual and a vanishing total visible bracket are
sufficient for weak Hamiltonian stationarity.  No factor equal to the inverse
physical observation mesh occurs in \eqref{v8:eq:weak-residual-bound}.
\end{theorem}

\begin{proof}
Dynkin's formula for the vector-valued finite reader gives the exact martingale
\begin{equation}
 M_h(t)=Z_h(t)-Z_h(0)-\int_0^tb_h(X_h(s))\dd s.
 \label{v8:eq:Dynkin}
\end{equation}
Choose an orthonormal basis of $E$ and apply the scalar jump-martingale
isometry componentwise.  Summing the components gives
\begin{equation}
 \E\|M_h(T)\|^2
 =\E\int_0^T\Gamma_h(X_h(s))\dd s,
 \qquad
 \E\sup_{t\le T}\|M_h(t)\|^2
 \le4\E\int_0^T\Gamma_h(X_h(s))\dd s.
 \label{v8:eq:martingale-isometry}
\end{equation}
Subtract \eqref{v8:eq:reference-H} from \eqref{v8:eq:Dynkin}.  With
$e_h=Z_h-z_h$,
\[
 e_h(t)=e_h(0)+M_h(t)+\int_0^t\mathfrak r_h(X_h(s))\dd s
 +\int_0^t\{X_{\mathscr H_h}(Z_h(s))-X_{\mathscr H_h}(z_h(s))\}\dd s.
\]
Take the running supremum, square, use Cauchy--Schwarz on the two time
integrals and \eqref{v8:eq:martingale-isometry}, then apply Gronwall's
inequality.  This proves \eqref{v8:eq:generator-comparison}.

For the weak statement, integration by parts for a càdlàg finite-jump path and
a compactly supported $C^1$ test gives
\[
 -\int\ip{Z_h}{\dot\psi}\dd t
 =\int\ip{\psi}{\dd Z_h}
 =\int\ip{\mathfrak r_h(X_h)}{\psi}\dd t
  +\int\ip{X_{\mathscr H_h}(Z_h)}{\psi}\dd t
  +\int\ip{\psi}{\dd M_h}.
\]
The deterministic integral is bounded by Cauchy--Schwarz.  The stochastic
integral is a square-integrable martingale with second moment at most
$\|\psi\|_\infty^2\E\int\Gamma_h$.  Subtract the Hamilton term and take square
roots.  No time discretization was introduced, which proves the final claim.
\end{proof}

\subsection{The same residual is the first variation of the canonical action}

On an absolutely continuous phase-space path define the symmetric canonical
action
\begin{equation}
 \mathscr S_{\mathscr H}(z)
 =\int_0^T\left\{
    -\tfrac12\ip{\mathbb Jz}{\dot z}-\mathscr H(z)
                  \right\}\dd t.
 \label{v8:eq:phase-action}
\end{equation}
For compactly supported variations $\eta$, direct integration by parts gives
\begin{equation}
 D\mathscr S_{\mathscr H}(z)[\eta]
 =-\int_0^T\ip{\mathbb J(\dot z-X_{\mathscr H}(z))}{\eta}\dd t
 =\mathcal R(\mathbb J\eta),
 \label{v8:eq:action-residual}
\end{equation}
where $\mathcal R$ is the distributional Hamilton residual of
\eqref{v8:eq:weak-residual}.  Hence \eqref{v8:eq:weak-residual-bound} is already
a bound on the canonical first variation on every fixed smooth test bank.  A
separate high-temporal-jet estimate is not needed for this first-order action.

\begin{corollary}[Weak native canonical-action handoff]
\label[corollary]{v8:cor:weak-action}
Suppose $\mathscr H_h\to\mathscr H$ in $C^1$ on the common chart,
$Z_h(0)\to z_0$ in mean square, and
\begin{equation}
 \E\int_0^T\|\mathfrak r_h(X_h(t))\|^2\dd t\longrightarrow0,
 \qquad
 \E\int_0^T\Gamma_h(X_h(t))\dd t\longrightarrow0.
 \label{v8:eq:weak-H-hyp}
\end{equation}
Then $Z_h\to z$ in mean square uniformly on $[0,T]$, where
$\dot z=X_{\mathscr H}(z)$ and $z(0)=z_0$.  On every fixed compactly supported
$C^1$ variation bank, the canonical first variations converge to zero.
If, in addition, the inherited positive-screen hypotheses identify the
canonical momentum, matter gradients and metric coefficients strongly in the
stated local $L^2$ topologies, all quadratic Hamiltonian and stress insertions
built continuously from that packet pass to their classical limits.
\end{corollary}

\begin{proof}
Uniform $C^1$ convergence of the Hamiltonians gives uniform convergence of
the vector fields and a common local Lipschitz constant.  Apply
\cref{v8:thm:generator-handoff}, absorbing the vector-field difference into
$\mathfrak r_h$, to obtain the uniform mean-square convergence.  Equation
\eqref{v8:eq:action-residual} and \eqref{v8:eq:weak-residual-bound} give the
first-variation statement.  The last assertion is exactly the elementary
strong-$L^2$ product passage already used by the inherited positive-screen
theorem: if $Y_h\to Y$ in $L^2$ and the coefficients converge uniformly on the
common chart, then a bounded quadratic contraction differs in $L^1$ by at
most $C\|Y_h-Y\|_2(\|Y_h\|_2+\|Y\|_2)$.
\end{proof}

\section{Strong chronological control is a strictly stronger entrance}
\label{v8:sec:strong}

The weak theorem charges the total martingale bracket
$\int\Gamma_h$.  An $H^1$ reconstruction of the exact nodes pays an additional
inverse time scale, as Versions 6--7 already exhibit in the scalar count.  The
following general statement separates the two demands.

Fix observation times $t_j=j\delta$.  Put $Z_j=\Phi_h(X_h(t_j))$ and
\begin{equation}
 B_j=\delta^{-1}\E[Z_{j+1}-Z_j\mid\mathcal F_j],\qquad
 \xi_j=Z_{j+1}-Z_j-\delta B_j.
 \label{v8:eq:sampled-pair}
\end{equation}
Let $\widetilde Z$ be the affine interpolation of the actual nodes.

\begin{proposition}[Exact strong Hamilton residual of sampled canonical pairs]
\label[proposition]{v8:prop:strong-H}
For the piecewise constant Hamilton field $F_j=X_{\mathscr H_h}(Z_j)$,
\begin{equation}
\boxed{
 \E\int_0^T\|\dot{\widetilde Z}-F_j\|^2\dd t
 =\sum_j\delta\,\E\|B_j-F_j\|^2
   +\delta^{-1}\sum_j\E\|\xi_j\|^2.}
 \label{v8:eq:strong-H-exact}
\end{equation}
If $X_{\mathscr H_h}$ is uniformly Lipschitz on the chart, then
\begin{align}
 \E\int_0^T\|\dot{\widetilde Z}
      -X_{\mathscr H_h}(\widetilde Z)\|^2\dd t
 &\le 2\sum_j\delta\E\|B_j-F_j\|^2
      +2\delta^{-1}\sum_j\E\|\xi_j\|^2\notag\\
 &\quad+C L^2\left[
       \delta^2\sum_j\delta\E\|B_j\|^2
       +\delta\sum_j\E\|\xi_j\|^2\right].
 \label{v8:eq:strong-H-full}
\end{align}
Thus strong first-order residual control requires the realized innovation cost
$\delta^{-1}\sum_j\E\|\xi_j\|^2\to0$, whereas weak action stationarity in
\cref{v8:thm:generator-handoff} only charges the unscaled total bracket.
\end{proposition}

\begin{proof}
On a cell,
$\dot{\widetilde Z}=B_j+\xi_j/\delta$.  The predictable vector
$B_j-F_j$ is $\mathcal F_j$-measurable and the conditional mean of $\xi_j$ is
zero.  Expanding the square, integrating over the cell and taking expectation
gives exactly
\[
 \delta\E\|B_j-F_j\|^2+\delta^{-1}\E\|\xi_j\|^2.
\]
Summation proves \eqref{v8:eq:strong-H-exact}.
For $t=t_j+s$,
$\|X_{\mathscr H_h}(\widetilde Z(t))-F_j\|
 \le L(s/\delta)\|Z_{j+1}-Z_j\|$.
Integration of $(s/\delta)^2$ gives a factor $\delta/3$.
Conditioning gives
$\E\|Z_{j+1}-Z_j\|^2
 =\delta^2\E\|B_j\|^2+\E\|\xi_j\|^2$.
Use $(a+b)^2\le2a^2+2b^2$ and sum.
\end{proof}

\section{The missing chronological Legendre incidence}
\label{v8:sec:Legendre}

Static reciprocity of two Hessians does not identify the momentum carried by
an accepted chronology with the Legendre transform of its actual field
velocity.  This section gives the exact finite residual for that identification.

Consider a bosonic sector with
\begin{equation}
 \mathscr L(q,v)=\tfrac12\ip{v}{\mathbb K(q)v}
                  +\ip{c(q)}{v}-V(q),
 \qquad mI\preceq\mathbb K(q)\preceq MI,
 \label{v8:eq:quadratic-L}
\end{equation}
so that
\begin{equation}
 p=D_v\mathscr L(q,v)=\mathbb K(q)v+c(q).
 \label{v8:eq:Legendre-map}
\end{equation}
Let $Q_j$ be the actual sampled field, $P_j$ the independently supplied
canonical-momentum reader on the same record, and write
\[
 B_j^q=\delta^{-1}\E[Q_{j+1}-Q_j\mid\mathcal F_j],\qquad
 \xi_j^q=Q_{j+1}-Q_j-\delta B_j^q.
\]
Define the predictable Legendre-incidence residual
\begin{equation}
 \iota_j=P_j-c(Q_j)-\mathbb K(Q_j)B_j^q.
 \label{v8:eq:incidence}
\end{equation}

\begin{theorem}[Exact momentum--velocity incidence identity]
\label[theorem]{v8:thm:incidence}
On each observation cell use the actual affine velocity
$v_j=(Q_{j+1}-Q_j)/\delta$.  Then
\begin{equation}
\boxed{
 \begin{aligned}
 &\E\sum_j\delta\,
 \bigl\|P_j-D_v\mathscr L(Q_j,v_j)\bigr\|_{\mathbb K(Q_j)^{-1}}^2\\
 &\qquad=\E\sum_j\delta\,\|\iota_j\|_{\mathbb K(Q_j)^{-1}}^2
       +\delta^{-1}\sum_j\E\|\xi_j^q\|_{\mathbb K(Q_j)}^2.
 \end{aligned}}
 \label{v8:eq:incidence-identity}
\end{equation}
Consequently a source-supplied momentum can agree strongly with the Legendre
momentum of an exact node-preserving $H^1$ field reconstruction only if both
the predictable incidence and the kinetic innovation cost vanish.  Uniform
invertibility of the kinetic Hessian, by itself, supplies neither condition.
\end{theorem}

\begin{proof}
By \eqref{v8:eq:Legendre-map},
\[
 P_j-D_v\mathscr L(Q_j,v_j)
 =\iota_j-\mathbb K(Q_j)\xi_j^q/\delta.
\]
In the $\mathbb K^{-1}$ norm the squared cross term is
$-2\delta^{-1}\ip{\iota_j}{\xi_j^q}$.  Its conditional expectation is zero
because $\iota_j$ is $\mathcal F_j$-measurable and
$\E[\xi_j^q\mid\mathcal F_j]=0$.  The final square is
$\delta^{-2}\|\xi_j^q\|_{\mathbb K}^2$.  Multiply by $\delta$, sum and take
expectation.
\end{proof}

\begin{remark}[Predictable momentum is a different semantic]
The weaker relation
$P_j\approx D_v\mathscr L(Q_j,B_j^q)$ can hold even when the second term in
\eqref{v8:eq:incidence-identity} survives.  It identifies $P_j$ with the
Legendre transform of the \emph{predictable drift}, not with the velocity of
an exact node-preserving path.  Either semantic can be legitimate only when it
is the one actually declared by the source and by the common-action compiler;
they cannot be interchanged to remove a kinetic defect.
\end{remark}

\section{A direct native test on the supplied count process}
\label{v8:sec:count}

Return to the inherited density process $X_n=K_N/n$ of Version 7, with generator
\eqref{v7:eq:count-generator}.  Instead of assigning it directly to a scalar
configuration, let
\begin{equation}
 \Phi:[0,1]\longrightarrow E
 \label{v8:eq:phase-decoder}
\end{equation}
be any independently specified $C^3$ decoder into a fixed canonical phase
space.  This covers, in particular, any proposal in which both a field and a
momentum are fixed smooth functions of the same count variable.  It does not
cover an additional independently fluctuating momentum source.

\begin{proposition}[Exact generator and bracket asymptotics of a fixed phase decoder]
\label[proposition]{v8:prop:count-generator}
Uniformly for $x\in[0,1]$,
\begin{align}
 (\mathcal L_n\Phi)(x)
 &=c(p-x)\Phi'(x)
   +\frac{\mathfrak a(x)}{2n}\Phi''(x)+O(n^{-2}),
 \label{v8:eq:count-generator-expansion}\\
 \Gamma_n^\Phi(x)
 &:=na(1-x)\|\Phi(x+n^{-1})-\Phi(x)\|^2
   +nbx\|\Phi(x-n^{-1})-\Phi(x)\|^2\notag\\
 &=\frac{\mathfrak a(x)}{n}\|\Phi'(x)\|^2+O(n^{-2}),
 \label{v8:eq:count-bracket-expansion}
\end{align}
where $c=a+b$, $p=a/c$ and $\mathfrak a(x)=a(1-x)+bx$.  The remainders are
bounded by fixed $C^3$ norms of $\Phi$ and the source rates.
\end{proposition}

\begin{proof}
Taylor's formula with integral remainder gives
\begin{align*}
 \Phi(x+n^{-1})-\Phi(x)
 &=n^{-1}\Phi'(x)+\tfrac12n^{-2}\Phi''(x)+O(n^{-3}),\\
 \Phi(x-n^{-1})-\Phi(x)
 &=-n^{-1}\Phi'(x)+\tfrac12n^{-2}\Phi''(x)+O(n^{-3}),
\end{align*}
uniformly on the compact interval, using one-sided expansions at the endpoints
where the corresponding jump rate vanishes.  Multiply by the two rates
$na(1-x)$ and $nbx$.  The first-order coefficient is
$a(1-x)-bx=a-cx=c(p-x)$, while the sum of the second-order coefficients is
$\mathfrak a(x)$.  This proves \eqref{v8:eq:count-generator-expansion}.
Squaring the same increments gives
$n^{-2}\|\Phi'(x)\|^2+O(n^{-3})$; multiply by the rates and sum.
\end{proof}

\begin{theorem}[Hamilton compatibility equation of every fixed count decoder]
\label[theorem]{v8:thm:count-Hamilton}
Let $\mathscr H\in C^2$ on a neighborhood of $\Phi([0,p])$ and put
$X_{\mathscr H}=\mathbb J\nabla\mathscr H$.  Starting from the all-$H$ state,
the decoded source $Z_n(t)=\Phi(X_n(t))$ has vanishing total visible bracket,
\begin{equation}
 \E\int_0^T\Gamma_n^\Phi(X_n(t))\dd t=O(n^{-1}),
 \label{v8:eq:count-bracket-vanish}
\end{equation}
and converges uniformly in mean square to $z(t)=\Phi(r(t))$.
Its weak canonical-action residual tends to zero for every smooth compact test
if and only if
\begin{equation}
\boxed{
 c(p-r)\Phi'(r)=\mathbb J\nabla\mathscr H(\Phi(r))
 }
 \label{v8:eq:count-H-compatibility}
\end{equation}
for every $r$ in the count-value interval traversed by the limiting trajectory.
If the identity fails on a nonempty subinterval, there is a fixed smooth test
bank on which the limiting canonical first variation is nonzero.
\end{theorem}

\begin{proof}
The path estimate of \cref{v7:lem:path} and Lipschitz continuity of $\Phi$
give uniform mean-square convergence to $\Phi(r(t))$.  Equation
\eqref{v8:eq:count-bracket-vanish} follows from
\eqref{v8:eq:count-bracket-expansion} and boundedness of $\Phi'$.  The source
Hamilton residual is, by \eqref{v8:eq:count-generator-expansion},
\[
 \mathfrak r_n(x)
 =c(p-x)\Phi'(x)-\mathbb J\nabla\mathscr H(\Phi(x))+O(n^{-1}).
\]
If \eqref{v8:eq:count-H-compatibility} holds, its occupation $L^2$ norm tends
to zero, and \cref{v8:thm:generator-handoff,v8:cor:weak-action} give weak
stationarity.

Conversely, suppose the continuous residual
$R(r)=c(p-r)\Phi'(r)-\mathbb J\nabla\mathscr H(\Phi(r))$ is nonzero at one
interior value.  By continuity there are a unit vector $e$, a subinterval
$I$, and $\epsilon>0$ such that $\ip{R(r)}e\ge\epsilon$ on $I$ after possibly
reversing $e$.  Since the deterministic $r(t)$ traverses $I$, choose a
nonnegative smooth time bump $\chi$ supported where $r(t)\in I$ and put
$\psi(t)=\chi(t)e$.  Uniform convergence of the source values and the generator
expansion show that the predictable part of
$\mathcal R_n(\psi)$ converges in probability to
$\int\chi\ip{R(r(t))}e\dd t>0$, while the martingale part tends to zero in
$L^2$ by \eqref{v8:eq:count-bracket-vanish}.  Equation
\eqref{v8:eq:action-residual} gives the same conclusion for the canonical first
variation.
\end{proof}

\section{A phase-space no-go for a single relaxing source coordinate}
\label{v8:sec:no-go}

Version 7 proved a scalar endpoint obstruction for a flat quartic potential.
The next theorem is its genuinely phase-space form.  It applies to an arbitrary
autonomous Hamiltonian and therefore shows what adding a momentum coordinate
that is merely another function of the same relaxing count can and cannot do.

\begin{theorem}[A one-coordinate relaxation cannot approach a strict Hamiltonian vacuum nontrivially]
\label[theorem]{v8:thm:strict-vacuum}
Let $\Phi\in C^1([r_0,p];E)$ satisfy the compatibility identity
\eqref{v8:eq:count-H-compatibility} on $[r_0,p)$, and suppose
$z_*=\Phi(p)$.  Then
\begin{equation}
 \mathscr H(\Phi(r))\equiv\mathscr H(z_*)
 \quad\text{on }[r_0,p].
 \label{v8:eq:energy-level}
\end{equation}
Consequently, if $z_*$ is a strict local minimum or a strict local maximum of
$\mathscr H$, then $\Phi$ is constant on some interval $(p-\epsilon,p]$.
In particular no decoder which is nonconstant arbitrarily close to $p$ can
identify the relaxing count with an autonomous canonical trajectory approaching
a strict Hamiltonian vacuum.

If additionally $\Phi'(p)\ne0$ and $\mathscr H\in C^2$, then
\begin{equation}
 \mathbb J D^2\mathscr H(z_*)\Phi'(p)=-c\Phi'(p),
 \label{v8:eq:endpoint-linear}
\end{equation}
and hence
\begin{equation}
 \ip{\Phi'(p)}{D^2\mathscr H(z_*)\Phi'(p)}=0.
 \label{v8:eq:endpoint-null}
\end{equation}
Thus a positive- or negative-definite Hamiltonian Hessian is incompatible with
$\Phi'(p)\ne0$.
\end{theorem}

\begin{proof}
Along the deterministic limiting trajectory $r'=c(p-r)$ and
$z=\Phi(r)$, compatibility gives $\dot z=\mathbb J\nabla\mathscr H(z)$.
Therefore
\[
 \frac{\dd}{\dd t}\mathscr H(z(t))
 =\ip{\nabla\mathscr H(z)}{\mathbb J\nabla\mathscr H(z)}=0
\]
because $\mathbb J$ is skew-adjoint.  As $t\to\infty$, $r(t)\to p$ and
continuity gives $z(t)\to z_*$.  Hence the constant energy equals
$\mathscr H(z_*)$, proving \eqref{v8:eq:energy-level}.  If $z_*$ is a strict
local extremum, the nearby level set
$\{z:\mathscr H(z)=\mathscr H(z_*)\}$ contains only $z_*$ in some
neighborhood.  Since $\Phi(r)\to z_*$, it must equal $z_*$ for all $r$ close
enough to $p$.

At $r=p$, continuity of the compatibility identity first gives
$\nabla\mathscr H(z_*)=0$.  For $r<p$, subtract this zero value on the
right-hand side of \eqref{v8:eq:count-H-compatibility} and divide by
$r-p$.  Since $\Phi$ is differentiable at $p$ and $\mathscr H$ is $C^2$,
the right-hand difference quotient tends to
$\mathbb J D^2\mathscr H(z_*)\Phi'(p)$, whereas the left-hand quotient is
$-c\Phi'(r)$.  Letting $r\to p$ gives
\[
 -c\Phi'(p)=\mathbb J D^2\mathscr H(z_*)\Phi'(p),
\]
which is \eqref{v8:eq:endpoint-linear}.  Multiply by $-\mathbb J$ and pair with
$\Phi'(p)$.  Since $\ip v{\mathbb Jv}=0$,
\[
 \ip{\Phi'(p)}{D^2\mathscr H(z_*)\Phi'(p)}=0.
\]
Definiteness excludes a nonzero vector with this property.
\end{proof}

\begin{remark}[What can evade the no-go]
The theorem concerns an autonomous Hamiltonian pair which is a fixed function
of one relaxing scalar source coordinate.  It does not exclude an independently
fluctuating momentum source, an explicitly time-dependent Hamiltonian, an
expanding gravitational background, additional retained bath degrees of
freedom, or a larger source whose projection onto the count is dissipative.
Those possibilities require their own same-history canonical reader and must
be tested by \eqref{v8:eq:H-residual}; they are not reconstructed from the count
marginal by the theorem.
\end{remark}

\section{Composition with the inherited positive screen}
\label{v8:sec:composition}

The new generator residual and the inherited positive screen solve different
problems.  The former identifies the chronological equation; the latter passes
quadratic stress insertions.  Their hypotheses must hold on the same physical
record.

\begin{corollary}[Canonical-pair replacement of a high-time-jet entrance]
\label[corollary]{v8:cor:screen-composition}
Consider one loaded matter sector of the Grand Tensor finite common action.
Assume:
\begin{enumerate}[label=\textup{(H\arabic*)},leftmargin=3em]
\item the finite Legendre certificate of
      \nolinkurl{thm:SMST-matter-Legendre} has a uniform kinetic-Hessian margin
      and vanishing reciprocity residual on the retained canonical pair;
\item the actual accepted source exports that same canonical pair through a
      decoder $\Phi_h$, and its finite generator residual and total visible
      bracket satisfy \eqref{v8:eq:weak-H-hyp};
\item the positive-screen packet of \nolinkurl{thm:SMST-positive-screen}, or its
      regenerative discharge, holds on the canonical momentum, matter gradient,
      and required metric coefficients;
\item the remaining common-action, gauge, potential, boundary, and gravitational
      hypotheses of the intended weak Einstein--matter handoff are satisfied on
      the same cofinal realization.
\end{enumerate}
Then the matter chronological first variation is supplied by the actual source
generator and converges to the Hamilton equations of the same common action;
its quadratic kinetic and stress insertions converge through the positive
screen.  No independent high-order temporal interpolation bound is required
for this matter sector.  This corollary does not remove any of the remaining
gravity, constraint, gauge, or source-occurrence hypotheses.
\end{corollary}

\begin{proof}
Item \textup{(H1)} identifies the source momentum coordinates with the
Hamiltonian coordinates of the same finite action at each cutoff.  Item
\textup{(H2)} and \cref{v8:cor:weak-action} pass the chronological canonical
first variation.  Item \textup{(H3)} gives strong local $L^2$ convergence of
the quadratic packet and therefore $L^1$ convergence of its stress insertions.
The remaining terms are precisely those left to item \textup{(H4)}.  None of
these steps uses a node-preserving high-time-jet reconstruction.
\end{proof}

\section{Outcome and next native calculation}
\label{v8:sec:ledger}

This iteration moves the programme one level upstream from the second-order
field equation.  The finite quantity that must now be acquired from the native
source is
\begin{equation}
\boxed{
 \mathfrak r_h(x)=\sum_yq_h(x,y)(\Phi_h(y)-\Phi_h(x))
                  -\mathbb J\nabla\mathscr H_h(\Phi_h(x)),
 \qquad
 \Gamma_h(x)=\sum_yq_h(x,y)\|\Phi_h(y)-\Phi_h(x)\|^2.}
 \label{v8:eq:final-audit}
\end{equation}
The first is the actual predictable Hamilton mismatch; the second is the
visible martingale bracket of the same canonical reader.  Both are finite table
quantities.  They test chronology directly and do not require a guessed
continuum acceleration.

The distinction from static Legendre reciprocity is now exact.  Reciprocal
kinetic Hessians show that the Lagrangian and Hamiltonian coordinates are
compatible.  Equation \eqref{v8:eq:incidence-identity} separately tests whether
the momentum on an accepted chronology is the Legendre momentum of its actual
node-preserving field velocity.  Equation \eqref{v8:eq:H-residual} tests whether
the chronological drift obeys the Hamiltonian vector field.  None follows from
the other by notation.

For the one concrete count primitive already available, the visible bracket is
small enough for weak first-order passage, but every fixed phase-space decoder
must satisfy the explicit ODE \eqref{v8:eq:count-H-compatibility}.  A decoder
which is merely a function of that relaxing count cannot approach a strict
autonomous Hamiltonian vacuum nontrivially.  Therefore the next noncircular
native calculation should seek the independently retained \emph{joint}
field--momentum source (or the equivalent normal-velocity/current packet) of the
locked spacetime law, rather than deriving momentum from the count marginal.
Its transition table should be tested directly by \eqref{v8:eq:final-audit} and
its same-record Legendre incidence by \eqref{v8:eq:incidence}.

\paragraph{Preservation and verification.}
The complete Version 7 body before its final document command is preserved
verbatim.  All new labels have prefix \texttt{v8:}.  The accompanying finite
diagnostics check Dynkin's generator/bracket identities on sample finite chains,
the exact Legendre-incidence decomposition, the count-generator Taylor
coefficients, the phase-action first variation, and the strict-vacuum endpoint
linearization.  They are diagnostics only; no machine-checked proof or full
Einstein--SM simulation is asserted.

% END IMMUTABLE VERSION 8 BODY
% APPEND VERSION 9 HERE, BEFORE THE FINAL END-DOCUMENT COMMAND.


% BEGIN IMMUTABLE VERSION 9 BODY
\clearpage
\begin{center}
{\Large\bfseries Version 9: Directed canonical correlations before source closure}\\[0.5em]
{\large A finite Hamiltonian audit, a circulation--noise inequality, and a parity-resolved positive branch}
\end{center}
\addcontentsline{toc}{part}{Version 9: Directed canonical correlations and reversal parity}

\section{Source audit, corrections, and the nonduplicating target}
\label{v9:sec:audit}

Version 8 identifies the actual generator residual of a declared canonical
reader.  Merely asking again for that residual to vanish would not advance the
source problem.  The present iteration asks which \emph{directed correlations}
can supply it, whether a reversible regenerative source can pass the test on a
noncollapsed canonical sector, and how a finite stochastic source can retain
Hamiltonian circulation while its visible jump noise tends to zero.

\paragraph{What is inherited.}
The graph spectral universality fibre in \cite{GT}, at
\nolinkurl{thm:GT-NCG-graph-fibre}, already proves that positive graph
spectralization retains symmetric conductances but forgets divergence-free
cycle currents.  Its \nolinkurl{thm:SMST-classical-shuffle-obstruction}
already separates self-adjoint permutation channels from skew Hamiltonian
\emph{operator} derivations.  Its future-word and hypocoercive results, at
\nolinkurl{thm:GT-future-word-OS} and
\nolinkurl{thm:GT-hypocoercive-memory}, distinguish reflected dissipation
from forward circulation.  None of those results is reproved as a new
structural theorem here.  The generator-to-path estimate is
\cref{v8:thm:generator-handoff}; the finite Legendre map, covariant reader
estimates, and common-action identification requirements are unchanged.

The additions below concern \emph{classical phase-space readouts on the same
finite source}: an exact normalized residual decomposition, its ordinary and
momentum-reversed symmetry tests, an explicit current--noise likelihood bound,
and a finite positive-rate family with identical symmetric data but opposite
canonical limits.  The exact audited phase source already in \cite{GT} is
tested separately from that diagnostic family.  Neither is substituted for an
unprovided locked Einstein--SM source.

\subsection{Two explicit corrections to Version 8}
\label{v9:sec:corrections}

The append-only convention requires that the following corrections be stated
here rather than silently altering the preceding text.

\begin{proposition}[Canonical sign and the independent Lipschitz hypothesis]
\label[proposition]{v9:prop:corrections}
With the real inner product and the convention
$\mathbb J^*=-\mathbb J$, $\mathbb J^2=-I$,
$X_{\mathscr H}=\mathbb J\nabla\mathscr H$ used in Version 8, the compatible
symmetric canonical action is
\begin{equation}
 \boxed{\mathscr S^+_{\mathscr H}(z)
 =\int_0^T\left\{\tfrac12\ip{\mathbb Jz}{\dot z}
                    -\mathscr H(z)\right\}\dd t.}
 \label{v9:eq:correct-action}
\end{equation}
For compactly supported variations,
\begin{equation}
 D\mathscr S^+_{\mathscr H}(z)[\eta]
 =-\int_0^T\ip{\mathbb J(\dot z-X_{\mathscr H}(z))}{\eta}\dd t
 =\mathcal R(\mathbb J\eta).
 \label{v9:eq:correct-variation}
\end{equation}
The minus sign in front of the kinetic term in
\eqref{v8:eq:phase-action} is therefore inconsistent with
\eqref{v8:eq:action-residual}.  If that old kinetic sign is kept, the Euler
equation is instead $\dot z=-\mathbb J\nabla\mathscr H$.

In \cref{v8:cor:weak-action} one must separately retain the common chart and
cutoff-uniform Lipschitz constant of
\cref{v8:thm:generator-handoff}.  Convergence of the Hamiltonians in $C^1$
only gives uniform convergence of their vector fields, not that Lipschitz
constant.  A common bound on $\|D^2\mathscr H_h\|$ is one sufficient way to
supply it.
\end{proposition}
\begin{proof}
The variation of the kinetic term in \eqref{v9:eq:correct-action} is
\[
 \tfrac12\int\{\ip{\mathbb J\eta}{\dot z}
                         +\ip{\mathbb Jz}{\dot\eta}\}\dd t
 =-\int\ip{\mathbb J\dot z}{\eta}\dd t.
\]
Subtracting the potential variation and using $\mathbb J^2=-I$ proves
\eqref{v9:eq:correct-variation}.  Changing the kinetic sign gives
$\mathbb J\dot z-\nabla\mathscr H=0$, hence the stated reversed flow.
For the second claim the one-coordinate functions
$H_n(x)=n^{-3}\sin(n^2x)$ converge to zero in $C^1$ on compact intervals,
while $H_n''(x)=-n\sin(n^2x)$ has unbounded supremum.  This example can be
included in one coordinate of a canonical space.  The comparison proof in
Version 8 works with the separately assumed common Lipschitz bound; it does
not require a stronger convergence of Hamiltonian Hessians.
\end{proof}

The generator, martingale, incidence, and count-decoder calculations in Version
8 use the declared vector field directly and are unaffected.  Its canonical
first-variation interpretations are to use \eqref{v9:eq:correct-action} and
the stated common Lipschitz bound.  The accompanying Version 8 numerical
variation check used the \emph{positive} kinetic sign; it consequently did
not check the sign printed in that version's action.  The present diagnostic
checks both signs explicitly.

\paragraph{Attribution.}
Orthogonal regression, stationary covariance identities, and symmetric/skew
splitting are standard linear algebra and Markov-process tools.  Precision
versus current-asymmetry tradeoffs have extensive precedents, including
\cite{v9:BS}.  The finite inequality below is proved directly for the visible
reader and is not an invocation of a long-time uncertainty relation.  Reversal
parity in phase space must also be specified; see \cite{v9:FS} for its
importance in stochastic dynamics.  No general priority claim is made for
these principles or for random walks approximating circular motion.

\section{An exact directed-correlation audit of a nonlinear Hamiltonian target}
\label{v9:sec:bank}

Let $\Omega$ be a finite state space, $k(x,y)\ge0$ its actual continuous-time
rates for $x\ne y$, and $\pi(x)>0$ an invariant probability.  Neither
reversibility nor irreducibility is needed in this section.  The generator is
\[
 (\mathcal Lf)(x)=\sum_{y\ne x}k(x,y)(f(y)-f(x)).
\]
All finite-dimensional coordinates use an orthonormal basis of the declared
physical mass Hilbert space.  Let $z:\Omega\to\R^d$ be the actual canonical
reader, with $d$ even, and let the target
$F(z)=\mathbb J\nabla\mathscr H(z)$ be fixed independently of the source
statistics.  Define
\begin{equation}
 \mu=\E_\pi z,\qquad
 C=\E_\pi[(z-\mu)(z-\mu)^{\mathsf T}]\succ0,\qquad
 x=C^{-1/2}(z-\mu).
 \label{v9:eq:whitening}
\end{equation}
Thus $\E_\pi x=0$, $\E_\pi xx^{\mathsf T}=I$.  Whitening is an analysis of
the measured covariance; it neither changes the physical reader nor defines
the Hamiltonian.  Without a uniform positive covariance window it is not a
cutoff-uniform equivalence of physical norms.

Put
\begin{equation}
 \begin{gathered}
 \beta=\mathcal Lx=C^{-1/2}\mathcal Lz,\qquad f=C^{-1/2}F(z),\\
 G=\E_\pi[\beta x^{\mathsf T}],\qquad T=\E_\pi[f x^{\mathsf T}],
 \qquad\Omega_c=\tfrac12(G-G^{\mathsf T}),\\
 Q=\E_\pi\sum_y k(X,y)
       (x(y)-x(X))(x(y)-x(X))^{\mathsf T}.
 \end{gathered}
 \label{v9:eq:bank}
\end{equation}
Here $Q\succeq0$ is the visible jump covariance \emph{per unit physical time}.
The symbol $\Omega_c$ is a correlation matrix, not the source state space.
For a matrix $A$ set
$\operatorname{sym}A=(A+A^{\mathsf T})/2$ and
$\operatorname{skew}A=(A-A^{\mathsf T})/2$.  Define the orthogonal residual
\begin{equation}
 \chi=\E_\pi\|\beta-f-(G-T)x\|^2\ge0.
 \label{v9:eq:unresolved}
\end{equation}
This contains, in particular, any constant-mean target mismatch and any
nonlinear or hidden-source component outside the span of the current
canonical coordinates.

\begin{theorem}[Three exact parts of the physical Hamilton mismatch]
\label[theorem]{v9:thm:decomposition}
For the actual source and independently fixed, possibly nonlinear target,
\begin{equation}
 G+G^{\mathsf T}=-Q,
 \label{v9:eq:balance-matrix}
\end{equation}
and the full normalized generator residual obeys
\begin{equation}
 \boxed{
 \begin{aligned}
 e^2&:=\E_\pi\|C^{-1/2}(\mathcal Lz-F(z))\|^2\\
 &=\chi+\|\Omega_c-\operatorname{skew}T\|_{\rm F}^2
              +\|\tfrac12Q+\operatorname{sym}T\|_{\rm F}^2.
 \end{aligned}}
 \label{v9:eq:decomposition}
\end{equation}
The identity is exact at every cutoff.  It separates unresolved conditional
motion, incorrect directed canonical circulation, and the symmetric
covariance/noise mismatch.
\end{theorem}
\begin{proof}
Apply stationarity to the scalar functions $x_ax_b$.  The finite product rule
is
\[
 \mathcal L(x_ax_b)=x_a\mathcal Lx_b+x_b\mathcal Lx_a
 +\sum_y k(\cdot,y)(x_a(y)-x_a)(x_b(y)-x_b).
\]
Taking $\pi$ expectation gives \eqref{v9:eq:balance-matrix} entry by entry.
For any vector-valued function $v$, its orthogonal projection onto the span
of the coordinates of $x$ is $(\E_\pi vx^{\mathsf T})x$, since these
coordinates are orthonormal in $L^2(\pi)$.  Apply Pythagoras to
$v=\beta-f$ to obtain
\[
 e^2=\chi+\|G-T\|_{\rm F}^2.
\]
Equation \eqref{v9:eq:balance-matrix} gives $G=\Omega_c-Q/2$.
Symmetric and skew matrices are orthogonal in the real Frobenius inner
product, which proves \eqref{v9:eq:decomposition}.
\end{proof}

\begin{corollary}[A closed linear Hamiltonian sector]
\label[corollary]{v9:cor:linear}
Suppose $F(z)=A(z-\mu)$, where $A=\mathbb J K$, $K=K^{\mathsf T}$ is the
fixed Hamiltonian Hessian.  Then
\[
 T=C^{-1/2}AC^{1/2},\qquad
 \chi=\E_\pi\|\beta-Gx\|^2.
\]
If the measured covariance satisfies $AC+CA^{\mathsf T}=0$, then
\begin{equation}
 e^2=\chi+\|\Omega_c-C^{-1/2}AC^{1/2}\|_{\rm F}^2
                         +\tfrac14\|Q\|_{\rm F}^2.
 \label{v9:eq:linear-audit}
\end{equation}
If the coordinate span is invariant under $\mathcal L$, then $\chi=0$.
Otherwise equality of the two displayed matrices does not remove the
unresolved residual.
\end{corollary}
\begin{proof}
The formula for $T$ follows by inserting $F=A(z-\mu)$ and
$\E xx^{\mathsf T}=I$.  Its symmetric part is
$C^{-1/2}(AC+CA^{\mathsf T})C^{-1/2}/2$.  The target is in the coordinate
span, so the orthogonal residual is the one stated.  Invariance of the
coordinate span is exactly $\beta=Gx$.
\end{proof}

\subsection{What a short chronological panel measures}

\begin{proposition}[Directed lag extraction with the error scale retained]
\label[proposition]{v9:prop:lag}
For the actual stationary chronological coupling define
\[
 R(t)=\E_\pi[x(X_t)x(X_0)^{\mathsf T}].
\]
Then $R(0)=I$, $R'(0)=G$, and, with
$M_2=\|\mathcal L^2x\|_{L^2(\pi)}$,
\begin{align}
 \left\|\frac{R(t)-R(t)^{\mathsf T}}{2t}-\Omega_c\right\|_{\rm F}
      &\le\tfrac12tM_2,\label{v9:eq:lag-skew}\\
 \left\|\frac{2I-R(t)-R(t)^{\mathsf T}}{t}-Q\right\|_{\rm F}
      &\le tM_2.\label{v9:eq:lag-noise}
\end{align}
If a measured lag matrix has Frobenius error at most $\epsilon_R$, its
additional errors in these estimates are respectively
$\epsilon_R/t$ and $2\epsilon_R/t$.
\end{proposition}
\begin{proof}
The finite semigroup gives
$R(t)=\E_\pi[(e^{t\mathcal L}x)x^{\mathsf T}]$.
An invariant Markov semigroup is contractive in $L^2(\pi)$, by Jensen's
inequality followed by stationarity.  Twice integrating its finite
matrix differential equation gives
\[
 e^{t\mathcal L}x=x+t\mathcal Lx+
             \int_0^t(t-s)e^{s\mathcal L}\mathcal L^2x\,\dd s.
\]
The last term has $L^2(\pi)$ norm at most $t^2M_2/2$.  For every vector
function $v$, Bessel's inequality for the orthonormal coordinates $x_a$
gives $\|\E_\pi vx^{\mathsf T}\|_{\rm F}\le\|v\|_{L^2(\pi)}$.
Hence $\|R(t)-I-tG\|_{\rm F}\le t^2M_2/2$.  Take the skew and symmetric
parts and use \eqref{v9:eq:balance-matrix}.  The measurement bounds are the
triangle inequality applied to $R$ and its transpose.
\end{proof}

There is no claimed cutoff-uniform bound on $M_2$ until it is computed.
Moreover, this lag panel alone does not determine $\chi$.  That quantity
uses the actual conditional drift squared (or an equivalent conditional
successor-replica observation), together with the target contractions.  A
second future draw conditional on the same current record is not the same
observation as an uncontrolled second time step.

\section{Ordinary reversibility, physical momentum reversal, and the actual phase source}
\label{v9:sec:reversal}

\begin{proposition}[Ordinary reversible sources cannot supply a directed canonical block]
\label[proposition]{v9:prop:reversible}
If $\pi(x)k(x,y)=\pi(y)k(y,x)$, then $\Omega_c=0$, so
\begin{equation}
 e^2\ge\|\operatorname{skew}T\|_{\rm F}^2.
 \label{v9:eq:reversible-bound}
\end{equation}
There is also a covariance-free form.  For any decoder $z$, any fixed
$z_*$, and a positive definite quadratic Hamiltonian
$\mathscr H(z)=\tfrac12(z-z_*)^{\mathsf T}K(z-z_*)$,
\begin{equation}
 \boxed{
 \E_\pi\|\mathcal Lz-\mathbb JK(z-z_*)\|_K^2
 =\E_\pi\|\mathcal Lz\|_K^2
       +\E_\pi\|\mathbb JK(z-z_*)\|_K^2.}
 \label{v9:eq:quadratic-reversible}
\end{equation}
Thus a nonvanishing quadratic Hamiltonian motion cannot be approximated in
stationary mean square by such an ordinarily reversible instantaneous source.
\end{proposition}
\begin{proof}
Detailed balance makes $\mathcal L$ self-adjoint in $L^2(\pi)$.
Thus $G_{ab}=\langle\mathcal Lx_a,x_b\rangle_\pi$ is symmetric,
proving \eqref{v9:eq:reversible-bound}.
For the second identity put $w=K^{1/2}(z-z_*)$ and
$A_K=K^{1/2}\mathbb JK^{1/2}$, which is skew-symmetric.  Since
$\mathcal Lw=K^{1/2}\mathcal Lz$, the left side is
$\E\|\mathcal Lw-A_Kw\|^2$.  The matrix
$\E[(\mathcal Lw)w^{\mathsf T}]$ is symmetric.  Its contraction with
$A_K$ vanishes, proving the identity by expansion of the square.
\end{proof}

\begin{corollary}[What is ruled out for the supplied renewal--interchange source]
\label[corollary]{v9:cor:actual-phase}
At its supplied stationary product law, the phase-flip and permitted
modulated-swap generator of \cref{v3:thm:transient-product} satisfies
\eqref{v9:eq:quadratic-reversible} for \emph{every} instantaneous
canonical decoder of the complete phase configuration, not only for a
function of the count.  If the second term on its right has a positive
cutoff-uniform lower bound, the Hamiltonian generator residual cannot tend
to zero on that sector.
\end{corollary}
\begin{proof}
The inherited actual-update audit of \cite{GT} gives detailed balance for
each phase flip with the product law.  Its allowed swaps preserve that law
and have equal forward and reverse rates because the modulation site is
unchanged by the swap.  Their sum is self-adjoint.  Apply
\cref{v9:prop:reversible} to the complete instantaneous decoder.
\end{proof}

This does not exclude the small-bracket count limit of Version 8: its
stationary physical variance collapses, and a nonstationary count trajectory
was treated there separately.  Nor does the corollary cover a history-based
momentum reader without auditing its enlarged chronological state and law.
The signed, gauge-constrained gravitational Hamiltonian is not positive
definite and is not covered by \eqref{v9:eq:quadratic-reversible} merely by
calling it another oscillator.

\subsection{Odd momenta change the reversal test}

\begin{proposition}[Parity-resolved detailed balance permits canonical circulation]
\label[proposition]{v9:prop:parity}
Let $\vartheta:\Omega\to\Omega$ be an actual involution preserving $\pi$,
and let $(\mathcal Pf)(x)=f(\vartheta x)$.  Suppose
\begin{equation}
 \mathcal L^*=\mathcal P\mathcal L\mathcal P,
 \qquad z(\vartheta x)-\mu=\mathsf R(z(x)-\mu),
 \qquad \mathsf R^*\mathsf R=I,\quad\mathsf R^2=I.
 \label{v9:eq:parity-input}
\end{equation}
Then $C$ commutes with $\mathsf R$ and
\begin{equation}
 G^{\mathsf T}=\mathsf R G\mathsf R,\qquad
 \mathsf RQ\mathsf R=Q,\qquad
 \mathsf R\Omega_c\mathsf R=-\Omega_c.
 \label{v9:eq:parity-blocks}
\end{equation}
Thus an even-coordinate/odd-momentum circulation is allowed.  In particular,
if $\mathsf R\mathbb J\mathsf R=-\mathbb J$ and a quadratic target has
$\mathsf RK\mathsf R=K$, its target drift has precisely that odd parity.
\end{proposition}
\begin{proof}
Invariance of $\pi$ under $\vartheta$ gives
$C=\mathsf RC\mathsf R$; spectral calculus then implies commutation of
$C^{-1/2}$ with $\mathsf R$.  Consequently $x\circ\vartheta=\mathsf Rx$.
Use the adjoint identity in \eqref{v9:eq:parity-input} in the scalar pairings
defining $G$ to obtain $G^{\mathsf T}=\mathsf RG\mathsf R$.  Taking its
symmetric and skew parts gives the remaining identities.  The last assertion
follows by multiplying the two declared parity identities for
$\mathbb J$ and $K$.
\end{proof}

Ordinary detailed balance and physical time reversal with an odd momentum
are different statements.  Neither a sign flip nor a parity operation may be
assigned to a source label after the fact solely to pass this test.  The
actual reflected experiment must provide that identification.  Even the
correct parity does not fix the \emph{magnitude or orientation} of the
circulation; it must still be compared with the independently fixed action.

\section{A quantitative circulation--noise cost on the same source}
\label{v9:sec:tradeoff}

For an unordered edge choose an orientation $x\to y$, and put
\[
 w_{xy}=\pi(x)k(x,y),\quad
 c_{xy}=\tfrac12(w_{xy}+w_{yx}),\quad
 j_{xy}=\tfrac12(w_{xy}-w_{yx}).
\]
Assume both orientations have positive rates on every retained edge.
Stationarity says that $j$ is divergence-free.  Define the ordinary
stationary-reversal relative-entropy rate
\begin{equation}
 \mathcal S_{\rm ord}
   =\sum_{\{x,y\}}(w_{xy}-w_{yx})\log\frac{w_{xy}}{w_{yx}}
   =2\sum_{\{x,y\}}j_{xy}\log\frac{c_{xy}+j_{xy}}{c_{xy}-j_{xy}}.
 \label{v9:eq:entropy}
\end{equation}
This is a statistical path-asymmetry quantity for \emph{ordinary} reversal.
It is not automatically thermodynamic heat or entropy production when the
physical reversal flips momenta.  The distinction will be exact in the
finite family below.

\begin{theorem}[Visible canonical circulation needs a noise--asymmetry budget]
\label[theorem]{v9:thm:tradeoff}
Let $M=\max_x\|x(x)\|$ for the normalized reader of
\eqref{v9:eq:whitening}.  Then
\begin{equation}
 \boxed{\mathcal S_{\rm ord}\,\tr Q
          \ge\frac{4}{M^2}\|\Omega_c\|_{\rm F}^2.}
 \label{v9:eq:tradeoff}
\end{equation}
For a target with residual $e$ from \eqref{v9:eq:decomposition}, this gives
\begin{equation}
 \mathcal S_{\rm ord}\,\tr Q
 \ge\frac4{M^2}
       \bigl(\|\operatorname{skew}T\|_{\rm F}-e\bigr)_+^2.
 \label{v9:eq:target-tradeoff}
\end{equation}
In particular, with a uniformly bounded normalized reader and a nonzero
limiting target circulation, simultaneous $e\to0$ and $\tr Q\to0$
force $\mathcal S_{\rm ord}\to\infty$.
\end{theorem}
\begin{proof}
Write $u=x(x)$, $v=x(y)$, $d=v-u$ on each oriented unordered edge.
Pairing the two directed contributions in $G$ gives
\[
 -c_{xy}dd^{\mathsf T}+j_{xy}d(u+v)^{\mathsf T}.
\]
Its skew part is $j_{xy}(vu^{\mathsf T}-uv^{\mathsf T})$, so
\begin{equation}
 \Omega_c=\sum_{\{x,y\}}j_{xy}(vu^{\mathsf T}-uv^{\mathsf T}),
 \qquad Q=2\sum_{\{x,y\}}c_{xy}dd^{\mathsf T}.
 \label{v9:eq:flux-bank}
\end{equation}
The elementary wedge identity gives
\[
 \|vu^{\mathsf T}-uv^{\mathsf T}\|_{\rm F}
 =\left\|\tfrac12\{d(u+v)^{\mathsf T}-(u+v)d^{\mathsf T}\}\right\|_{\rm F}
 \le\sqrt2 M\|d\|.
\]
For $|a|<1$, $a\log((1+a)/(1-a))\ge2a^2$, by integrating
$2/(1-a^2)\ge2$ and using the sign of $a$.  Therefore
$\mathcal S_{\rm ord}\ge4\sum j_{xy}^2/c_{xy}$.  Cauchy--Schwarz now gives
\[
 \|\Omega_c\|_{\rm F}^2
 \le2M^2\left(\sum\frac{j_{xy}^2}{c_{xy}}\right)
              \left(\sum c_{xy}\|d\|^2\right)
 \le\tfrac14M^2\mathcal S_{\rm ord}\tr Q.
\]
This proves \eqref{v9:eq:tradeoff}.  Equation
\eqref{v9:eq:decomposition} implies
$\|\Omega_c-\operatorname{skew}T\|_{\rm F}\le e$;
the reverse triangle inequality proves \eqref{v9:eq:target-tradeoff}.
\end{proof}

The maximum-reader bound is an explicit hypothesis, not a consequence of
unit covariance.  When a covariance is collapsed and its whitening magnifies
rare source states, $M$ may diverge.  The exact cutoff-dependent bound remains
valid, but the uniform conclusion must not then be asserted.

\begin{corollary}[What the inherited spectral fibre can and cannot change]
\label[corollary]{v9:cor:spectral-fibre}
Hold fixed the masses $\pi$, symmetric fluxes $c_{xy}$, and the physical
reader.  Vary only the stationary circulation within the inherited graph
fibre.  Then $Q$ and the positive graph spectralization are unchanged,
whereas $\Omega_c$ changes by the explicit current contraction in
\eqref{v9:eq:flux-bank}.  Thus circulation cannot reduce the stationary
visible bracket at fixed symmetric data; it can change whether the
independently fixed Hamiltonian is matched.
\end{corollary}
\begin{proof}
The two statements follow respectively from the formulas for $Q$ and
$\Omega_c$ in \eqref{v9:eq:flux-bank}.  The invariance of graph
spectralization under that circulation is the inherited graph-fibre theorem
of \cite{GT}.
\end{proof}

\section{A positive finite chronological family with opposite canonical limits}
\label{v9:sec:cycle}

This section is a declared diagnostic family, not the native locked
Einstein--SM source.  Its purpose is positive: finite stochastic dynamics
\emph{can} pass the canonical generator test without action-gradient descent
or future continuum samples, but symmetric source data do not determine which
Hamiltonian motion occurs.

Fix $R_*>0$, $a_*>|\omega|>0$, and let $n\ge5$,
$\alpha_n=2\pi/n$.  Use the canonical matrix
\[
 \mathbb J=\begin{pmatrix}0&1\\-1&0\end{pmatrix}
\]
and, on $\Z/n\Z$, the reader and positive rates
\begin{equation}
 z_n(j)=R_*e^{\alpha_n j\mathbb J}e_1,\qquad
 k_+(v)=\frac{a_*+v}{2\sin\alpha_n},\qquad
 k_-(v)=\frac{a_*-v}{2\sin\alpha_n},\qquad |v|<a_*.
 \label{v9:eq:cycle-rates}
\end{equation}
The actual jump rule adds or subtracts one at those constant rates.  Its
stationary law is uniform.  The Hamiltonian against which it is tested is
fixed separately:
\begin{equation}
 \mathscr H_\omega(z)=\tfrac\omega2\|z\|^2,
 \qquad X_{\mathscr H_\omega}=\omega\mathbb Jz.
 \label{v9:eq:cycle-H}
\end{equation}

\begin{theorem}[Exact canonical audit and path limit of the finite cycle]
\label[theorem]{v9:thm:cycle}
Put $\nu_n=a_*\tan(\alpha_n/2)$.  The source has
\begin{align}
 \mathcal L_n^{(v)}z_n&=-\nu_n z_n+v\mathbb Jz_n,
 &\Gamma_n(j)&=2R_*^2\nu_n,\label{v9:eq:cycle-drift}\\
 C_n&=\tfrac12R_*^2I,
 &Q_n&=2\nu_n I,\qquad \Omega_{c,n}=v\mathbb J,\qquad\chi_n=0.
 \label{v9:eq:cycle-bank}
\end{align}
Its exact normalized residual for \eqref{v9:eq:cycle-H} is
\begin{equation}
 \boxed{e_n^2=2\nu_n^2+2(v-\omega)^2.}
 \label{v9:eq:cycle-error}
\end{equation}
For every initial state, or any initial distribution, its actual path obeys
\begin{equation}
 \E\sup_{t\le T}
 \|z_n(X_t)-e^{tv\mathbb J}z_n(X_0)\|^2
 \le C_{T,R_*,a_*,v}\alpha_n.
 \label{v9:eq:cycle-path}
\end{equation}
At any fixed terminal time, the exact error is
\begin{equation}
 \E\|z_n(X_t)-e^{tv\mathbb J}z_n(X_0)\|^2
          =2R_*^2(1-e^{-\nu_nt}).
 \label{v9:eq:cycle-terminal}
\end{equation}
In particular, the family $v=\omega$ converges to the fixed Hamiltonian
motion with visible total bracket $O(n^{-1})$ and weak canonical residual
$O(n^{-1/2})$.  This holds for actual generated paths without sampling a
prescribed future orbit.
\end{theorem}
\begin{proof}
Use $e^{\pm\alpha\mathbb J}=\cos\alpha\,I\pm\sin\alpha\,\mathbb J$.
The sum and difference of the rates are $a_*/\sin\alpha$ and
$v/\sin\alpha$.  Their action on the reader gives
$-a_*(1-\cos\alpha)z/\sin\alpha+v\mathbb Jz$,
which is \eqref{v9:eq:cycle-drift}.  Each jump has squared reader length
$2R_*^2(1-\cos\alpha)$.  Multiplying by the total rate proves the bracket
formula.  The finite root-of-unity sums give mean zero and covariance
$R_*^2I/2$.  The coordinate span is invariant, so its normalized drift is
$-\nu_nI+v\mathbb J$, and \eqref{v9:eq:balance-matrix} gives the stated
$Q_n$.  This proves \eqref{v9:eq:cycle-bank} and
\eqref{v9:eq:cycle-error}.

For a direct path proof realize the two types of jumps by independent
Poisson counts $N_+,N_-$ of the displayed rates.  This is exactly the finite
cycle law after reduction modulo $n$.  The lifted angle increment is
$A_t=\alpha(N_+(t)-N_-(t))$, with mean
$t\alpha v/\sin\alpha$.  Its centered martingale has variance
$t a_*\alpha^2/\sin\alpha$.  Doob's inequality therefore gives
\[
 \E\sup_{t\le T}|A_t-vt|^2
 \le 8T a_*\alpha^2/\sin\alpha
       +2T^2v^2(\alpha/\sin\alpha-1)^2\le C_T\alpha.
\]
The map $\theta\mapsto R_*e^{\theta\mathbb J}e_1$ is $R_*$-Lipschitz.
This proves \eqref{v9:eq:cycle-path}, uniformly in the initial angle.
The closed reader drift also gives
$\E[z_n(X_t)\mid X_0]=e^{-\nu_nt}e^{tv\mathbb J}z_n(X_0)$.
Both the actual and reference vectors have norm $R_*$, so expansion of the
squared difference proves \eqref{v9:eq:cycle-terminal}.
Alternatively the generator comparison theorem applies with residual
$-\nu_nz$ on the matched branch.  Its weak test estimate and the corrected
canonical action \eqref{v9:eq:correct-action} give the final residual rate.
\end{proof}

\begin{proposition}[Identical symmetric data do not identify the sign of Hamiltonian motion]
\label[proposition]{v9:prop:opposite}
The families $v=\omega$ and $v=-\omega$ have the same invariant law,
symmetric generator, positive graph spectral data, and visible jump
covariance at every $n$.  Their entire symmetrized chronological semigroups
also coincide.  Nevertheless their limits are opposite Hamiltonian flows,
and the second family's normalized residual for the independently fixed
$\mathscr H_\omega$ tends to $8\omega^2$ in squared norm.
On the canonical reader the complete lag matrix is
\begin{equation}
 R_n^{(v)}(t)=e^{-\nu_nt}
           \{\cos(vt)I+\sin(vt)\mathbb J\}.
 \label{v9:eq:cycle-lag}
\end{equation}
Its symmetric part cannot distinguish these two sources; its directed
position--momentum cross correlation can.
\end{proposition}
\begin{proof}
The total rate is independent of $v$ and the two directed rates are
interchanged by $v\mapsto-v$.  Hence
$(\mathcal L_n^{(v)})^*=\mathcal L_n^{(-v)}$ in the uniform law and their
symmetric parts agree.  The cyclic shift commutes with its inverse, so the
same statement at positive times also follows by exponentiation (or directly
from taking adjoints): the symmetrizations of the two semigroups coincide.
The covariance assertion is \eqref{v9:eq:cycle-bank}.
The invariant two-dimensional coordinate span has generator
$-\nu_nI+v\mathbb J$, proving \eqref{v9:eq:cycle-lag}.
Finally \eqref{v9:eq:cycle-path} and \eqref{v9:eq:cycle-error} give the
opposite limits and the nonzero target residual.
\end{proof}

\begin{proposition}[The exact reversal cost and a physical parity check]
\label[proposition]{v9:prop:cycle-parity}
For $v\ne0$, the ordinary reversal rate and normalized bracket satisfy
\begin{equation}
 \mathcal S_{{\rm ord},n}
 =\frac{v}{\sin\alpha_n}\log\frac{a_*+v}{a_*-v},\qquad
 \mathcal S_{{\rm ord},n}\tr Q_n
 \longrightarrow 2a_*v\log\frac{a_*+v}{a_*-v}\ge4v^2.
 \label{v9:eq:cycle-cost}
\end{equation}
The constant in \eqref{v9:eq:tradeoff} is asymptotically sharp within this
family when $n\to\infty$ and then $v/a_*\to0$.

For the declared involution $\vartheta j=-j$ and canonical parity
$\mathsf R=\operatorname{diag}(1,-1)$, one has
$z_n(\vartheta j)=\mathsf Rz_n(j)$ and
$(\mathcal L_n^{(v)})^*=\mathcal P\mathcal L_n^{(v)}\mathcal P$.
The stationary path law is exactly invariant under reversal of time
\emph{together with} this parity.  Its relative entropy from that
parity-reversed law is zero, despite
$\mathcal S_{{\rm ord},n}\asymp n$.
\end{proposition}
\begin{proof}
There are $n$ unoriented cycle edges, each with forward flux $k_+/n$
and reverse flux $k_-/n$.  Substitution into
\eqref{v9:eq:entropy} gives the first formula.  Since
$\tr Q_n=4a_*\tan(\alpha_n/2)$ and
$\tan(\alpha/2)/\sin\alpha=1/(1+\cos\alpha)$, the product tends to
$2a_*v\log((a_*+v)/(a_*-v))$.  The elementary logarithmic inequality used
above proves the lower bound.  Here $M^2=2$ and
$\|\Omega_c\|_{\rm F}^2=2v^2$, so the right side of
\eqref{v9:eq:tradeoff} is $4v^2$.  Expansion of the logarithm at small
$v/a_*$ proves sharpness in the stated successive limit.

Reflection of the cycle interchanges its two directed shifts and leaves
the uniform law invariant, proving the adjoint identity and reader parity.
For a stationary finite-jump path, reverse its event order and reflect
every state.  Each forward jump becomes a forward jump under this combined
operation, and similarly for each backward jump.  The initial stationary
weight is unchanged, the constant escape rates give the same holding-time
factors, and the products of jump rates agree.  Thus the complete path
probabilities, including their jump times, are identical.  Their relative
entropy is zero.  This proves the asserted distinction without a
thermodynamic interpretation of the ordinary-reversal cost.
\end{proof}

The target energy $\mathscr H_\omega$ is constant on this fixed-radius
reader, but its Hamiltonian vector field is nonzero.  This is consistent:
canonical evolution follows energy levels, not the negative gradient of the
instantaneous energy.  The result concerns the \emph{canonical history}
action \eqref{v9:eq:correct-action}, not smallness of
$\nabla\mathscr H_\omega$ itself.  It neither establishes an Einstein
geometry nor selects the target orientation from the unoriented cycle data.

\section{From a stationary directed audit to actual nonstationary paths}
\label{v9:sec:handoff}

The stationary matrix audit need not force an equilibrium initial law in
its application.  The following conclusion keeps its precise price.

\begin{theorem}[A same-source canonical handoff for dominated initial populations]
\label[theorem]{v9:thm:handoff}
For a cutoff family as above suppose
$c_*I\preceq C_h\preceq C_*I$, the actual reader and reference Hamiltonian
orbits remain in a common chart, and their target vector fields have a common
Lipschitz constant there.  Let the actual initial law satisfy
$\mu_{h,0}\le D_h\pi_h$ pointwise.  No independence between cutoffs is
required.  Let $e_h^2$ be the sum in \eqref{v9:eq:decomposition}, and let
$z_h^{\rm ref}(t)$ be the target Hamiltonian trajectory started at the
\emph{actual random} initial reader $Z_h(0)$.  Then
\begin{equation}
 \E\sup_{t\le T}\|Z_h(t)-z_h^{\rm ref}(t)\|^2
 \le C_{T,L,C_*}D_h\bigl(e_h^2+\tr Q_h\bigr).
 \label{v9:eq:dominated-path}
\end{equation}
For any fixed compactly supported $C^1$ canonical variation $\eta$,
its corrected weak first variation has the bound
\begin{equation}
 \|\mathcal R_h(\mathbb J\eta)\|_{L^2(\PP)}
 \le \sqrt{TC_*D_h}\left(
       \|\eta\|_{L^2}e_h+\|\eta\|_\infty\sqrt{\tr Q_h}\right).
 \label{v9:eq:dominated-variation}
\end{equation}
Consequently $D_h(e_h^2+\tr Q_h)\to0$ gives actual-path canonical
closure relative to those random initial data.  A deterministic common
limit additionally requires convergence of the initial reader and the
Hamiltonian vector fields in the hypotheses of the corrected
\cref{v8:cor:weak-action}.
\end{theorem}
\begin{proof}
Positivity and stationarity preserve domination:
$\mu_{h,t}=\mu_{h,0}e^{t\mathcal L_h}\le D_h\pi_he^{t\mathcal L_h}
=D_h\pi_h$.  The physical residual obeys
\[
 \E_{\pi_h}\|\mathcal L_hz_h-F_h(z_h)\|^2\le C_*e_h^2,
 \qquad
 \E_{\pi_h}\Gamma_h^{z}\le C_*\tr Q_h.
\]
Integrate these inequalities under the dominated law on $[0,T]$.
The proof of \cref{v8:thm:generator-handoff} applies with a random
initial reference measurable at time zero: its deterministic Gronwall
comparison is pathwise and the martingale estimates are unchanged.
This proves \eqref{v9:eq:dominated-path}.  The weak residual estimate
of that theorem gives \eqref{v9:eq:dominated-variation}, with
\eqref{v9:eq:correct-variation} identifying the canonical sign.  Initial
and Hamiltonian convergence are separate from stationarity of a population,
which proves the final qualification.
\end{proof}

This theorem does not supply confinement, spatial stress compactness,
constraint propagation, or gauge fixing.  It converts an independently
verified \emph{same-source} directed audit into the chronological row of
Version 8.  The remaining positive-screen and Einstein--SM hypotheses are
exactly those retained in \cref{v8:cor:screen-composition}.  In a large
state space a single prescribed initial atom may have a large domination
cost $D_h$; it cannot be replaced by one without proof.  The cycle calculation
above avoids that cost by its separate uniform-in-initial-state path proof.

\section{Outcome: acquire the directed canonical panel, not another static screen}
\label{v9:sec:ledger}

The new positive construction shows that persistent Hamiltonian motion and
vanishing visible martingale noise are compatible with genuinely stochastic
finite transition rules.  A stationary law can carry circulating trajectories;
stationarity of a law is not static fieldwise criticality.  Its canonical
orientation is not supplied by the positive spectralization, by the invariant
measure, or even by all symmetrized lag correlations.  The actual directed
position--momentum correlation is load bearing.

For the only explicitly audited phase source in the supplied monograph,
ordinary detailed balance gives a stronger negative conclusion than a test of
the scalar count alone.  No instantaneous decoder of its whole stationary
configuration can match a noncollapsed positive quadratic Hamiltonian motion
in the physical mean-square norm.  A history-dependent, momentum-odd, or
coherent extension is a different source object and must be acquired with its
actual dynamics.  The diagnostic cycle proves that such directionality can
matter, but does not identify that cycle with the locked opportunity source.

The finite upstream audit is now
\begin{equation}
 \boxed{
 (C_h,G_h,Q_h,T_h,\chi_h;\ \vartheta_h,\mathsf R_h),
 \qquad
 e_h^2=\chi_h+\|\operatorname{skew}G_h-\operatorname{skew}T_h\|_{\rm F}^2
       +\|Q_h/2+\operatorname{sym}T_h\|_{\rm F}^2.
 }
 \label{v9:eq:final-bank}
\end{equation}
The parenthesized parity data are used only when actually supplied.  The
Hamiltonian in $T_h$ must still be the independently reconstructed common
one, not a fitted rotation inferred from $G_h$.  The residual $\chi_h$ must
be evaluated or controlled; closing the visible linear bank does not certify
unobserved conditional motion.  A stationary audit also pays the actual
initial-population and chart costs before it becomes a path theorem.

The next useful native observation is the \emph{ordered} field--momentum
cross correlation and its physically reflected partner, together with the
conditional-drift residual outside that bank.  If it is ordinarily
self-adjoint on a retained noncollapsed stable matter block, the present
test identifies an obstruction immediately.  If it has the required
momentum-odd circulation, its magnitude and orientation can be compared
with the same action, and its visible bracket can be propagated by the
lower-regularity route.  This is a specific missing chronological source
row, not another assumption that the Einstein--SM limit already exists.

\paragraph{Verification and preservation.}
All pre-existing Version 8 bytes before its final end-document command are
preserved.  The two corrections in \cref{v9:prop:corrections} govern the
subsequent use of that text.  The diagnostic script checks both canonical
signs, exact stationary matrix balances and residual decompositions,
directed-lag bounds, the reversal and parity identities, the flux/noise
inequality on finite directed graphs, and the explicit cycle rates and
limits.  These are numerical and symbolic diagnostics, not machine-checked
proofs or a simulation of the coupled Einstein--SM system.  All new labels
have prefix \texttt{v9:}.

\begingroup
\renewcommand{\refname}{Additional references for Version 9}
\begin{thebibliography}{9}
\small\raggedright
\bibitem{v9:BS}
A.~C.~Barato and U.~Seifert,
\emph{Thermodynamic uncertainty relation for biomolecular processes},
Physical Review Letters \textbf{114} (2015), 158101.
\href{https://arxiv.org/abs/1502.05944}{arXiv:1502.05944};
\href{https://doi.org/10.1103/PhysRevLett.114.158101}{doi:10.1103/PhysRevLett.114.158101}.
Background for current-precision tradeoffs; no asymptotic uncertainty theorem
from this reference is used to prove the finite reader inequality above.
\bibitem{v9:FS}
I.~J.~Ford and R.~E.~Spinney,
\emph{Entropy production from stochastic dynamics in discrete full phase space},
Physical Review E \textbf{86} (2012), 021127.
\href{https://arxiv.org/abs/1204.4822}{arXiv:1204.4822};
\href{https://doi.org/10.1103/PhysRevE.86.021127}{doi:10.1103/PhysRevE.86.021127}.
Background for parity-sensitive reversal.  The finite parity identities and
cycle path-law comparison in this body are proved directly.
\end{thebibliography}
\endgroup
% END IMMUTABLE VERSION 9 BODY
% APPEND VERSION 10 HERE, BEFORE THE FINAL END-DOCUMENT COMMAND.


% BEGIN IMMUTABLE VERSION 10 BODY
\clearpage
\begin{center}
{\Large\bfseries Version 10: The physical renewal clock before instantaneous drift}\\[0.5em]
{\large An explicit phase corrector, accepted-event Hamiltonian transfer, and exact chronological fluctuations}
\end{center}
\addcontentsline{toc}{part}{Version 10: The actual renewal clock and a phase-corrected Hamiltonian entrance}

\section{Upstream audit: distinguish a fast clock from a wrong physical equation}
\label{v10:sec:audit}

Version 9 tests the directed canonical current of a continuous-time Markov
source. Its small occupation-norm drift residual is sufficient, not
necessary: an actual microscopic clock can produce a rapidly alternating
conditional velocity whose accumulated effect vanishes. Before declaring
such a velocity to be a physical Hamiltonian mismatch, its retained clock
phase must be examined. This iteration does that calculation for the exact
completed-private timing source already supplied in \cite{GT}.

\paragraph{Inherited inputs and the nonduplication boundary.}
The following source results are used with their stated scope:
\nolinkurl{thm:two-state-active-renewal} gives the protected-tick phase matrix;
\nolinkurl{thm:accepted-response-renewal} gives the accepted waiting-time
transform; \nolinkurl{thm:accepted-response-pressure} calibrates the mean
accepted-response rate; and \nolinkurl{thm:accepted-phase-feedback} treats the
phase feedback for an idempotent payload. Their proofs, the general returned
memory calculus, and the generic renewal-reward pressure identities are not
new results here. The locked operator normal form at
\nolinkurl{thm:GT76-locked-normalform} and the semantic boundary at
\nolinkurl{thm:renewal-semantic-stripping} are also retained: a timing law
alone does not specify an accepted field operation.

What is added is a nonlinear, first-moment and second-moment transfer for an
\emph{independently supplied} accepted field kernel on this actual clock.
An explicit two-state Poisson corrector proves chronological convergence
without a small raw phase-resolved drift residual. A finite circular reader
then shows that the latter residual can stay strictly positive while the
paths converge. For that reader the complete phase correction, its
martingale variance, and its first-mode dispersion are calculated exactly.
No exponential replacement of the waiting law, resampling of field values,
or modification of the source's protected ticks is made.

\paragraph{Attribution and logical scope.}
Martingale decompositions, Poisson correctors, and averaging over a fast
finite-state variable are standard methods. Random-time-change and
semi-Markov representations have extensive precedents, including
\cite{v10:Kurtz,v10:MS}. The contribution of this body is the explicit
same-clock calculation, the cutoff-uniform constants even for rare
acceptance, the retained fluctuation coefficient, and the resulting
weakening of the source entrance used in Versions 8--9. Every assertion below
is proved here from finite rows and elementary estimates. The circular
accepted operation is a diagnostic extension of the actual timing source,
not the acquired locked Einstein--SM field operation. The general theorem
requires that operation's actual event moments; it does not manufacture them.
The canonical action convention is the corrected one in
\cref{v9:prop:corrections}.

\section{The actual timing source and its accepted-event moment bank}
\label{v10:sec:clock}

Let $s_j\in\{H,P\}$ be the retained active phase at protected tick $j$.
The inherited phase matrix and invariant law are
\begin{equation}
 M=\begin{pmatrix}1/5&4/5\\2/3&1/3\end{pmatrix},
 \qquad \pi=(5/11,6/11),\qquad \lambda_{\rm ph}=-7/15.
 \label{v10:eq:phase}
\end{equation}
At a $P\to H$ completion an accepted mark occurs with probability
$0<\vartheta\le1$. Rejection leaves the field unchanged. On acceptance
use the independently specified stochastic kernel $K_\epsilon(z,\dd y)$
for the \emph{actual physical canonical reader} $z$ in a finite-dimensional
real mass Hilbert space $E$. At each cutoff the attainable state set is
finite; its size may grow. The full tick table is
\begin{equation}
\begin{array}{c|c|c}
\text{current phase}&\text{next phase and field}&\text{conditional probability}\\\hline
H&(H,z)&1/5\\
H&(P,z)&4/5\\
P&(P,z)&1/3\\
P&(H,z)\text{ with rejected mark}&(2/3)(1-\vartheta)\\
P&(H,y)\text{ with accepted mark}&(2/3)\vartheta K_\epsilon(z,\dd y).
\end{array}
\label{v10:eq:full-table}
\end{equation}
Different marked outcomes can have the same terminal field and are still
retained as distinct source records. In particular the rejection mark is not
inferred from whether $y=z$. The acceptance parameter is independent of the
current field on this declared branch. State-dependent acceptance or a
history-dependent change of this table is not silently included.

Let $d>0$ be the independently calibrated physical duration of one tick.
Define
\begin{equation}
 \epsilon:=\frac{11d}{4\vartheta},\qquad
 d=\frac{4\vartheta}{11}\epsilon.
 \label{v10:eq:physical-calibration}
\end{equation}
Thus $\epsilon$ is the \emph{mean physical time between accepted responses},
not a new choice of physical units. The dense small-event branch considered
below has $\epsilon\to0$. Keeping $d/\vartheta$ nonzero instead is a
different timing regime. A family cannot be made to pass the mean-event
calibration by relabelling the action-comparison index as physical time.

\begin{proposition}[The retained variance of the actual accepted clock]
\label[proposition]{v10:prop:clock-moments}
For a fresh completed-boundary entrance the accepted raw waiting time $T$
has the inherited transform
\begin{equation}
 \E z^T=\frac{8\vartheta z^2}
 {15-8z+(8\vartheta-7)z^2}.
 \label{v10:eq:waiting-transform}
\end{equation}
Its first two moments, in the physical normalization
\eqref{v10:eq:physical-calibration}, satisfy
\begin{align}
 \E T&=\frac{11}{4\vartheta},&
 \operatorname{Var}T&=\frac{121-104\vartheta}{16\vartheta^2},\\
 \E(dT)&=\epsilon,&
 \frac{\operatorname{Var}(dT)}{\epsilon^2}
   &=c_{\rm clk}^2:=1-\frac{104}{121}\vartheta.
 \label{v10:eq:clock-moments}
\end{align}
In particular $17/121\le c_{\rm clk}^2<1$. The actual renewal clock is not
a deterministic sequence of duration-$\epsilon$ events, even at
$\vartheta=1$.
\end{proposition}
\begin{proof}
The transform is the input just cited, not a new reconstruction theorem.
One inter-opportunity time $W$ is the sum of independent geometric variables
on $\{1,2,\ldots\}$ with success probabilities $4/5$ and $2/3$, as is also
seen from its transform $8z^2/((5-z)(3-z))$. Hence
$\E W=11/4$ and $\operatorname{Var}W=17/16$.
The accepted wait is the sum of $G$ independent copies of $W$, where
$G$ is geometric with success probability $\vartheta$. Conditioning on $G$
gives
\[
 \operatorname{Var}T=\frac{17}{16\vartheta}
        +\frac{121(1-\vartheta)}{16\vartheta^2}.
\]
This is the second formula; the remaining assertions are multiplication by
$d$ and substitution of \eqref{v10:eq:physical-calibration}.
\end{proof}

Fix an independently specified Hamiltonian $\mathscr H$ and the fixed
canonical skew isometry $\mathbb J$. Write $F=\mathbb J\nabla\mathscr H$.
For the actually observed accepted kernel put
\begin{equation}
 \begin{aligned}
 m_\epsilon(z)&=\epsilon^{-1}\int(y-z)K_\epsilon(z,\dd y),\\
 a_\epsilon(z)&=m_\epsilon(z)-F(z),\\
 v_\epsilon(z)&=\epsilon^{-2}\int\|y-z\|^2K_\epsilon(z,\dd y).
 \end{aligned}
 \label{v10:eq:event-bank}
\end{equation}
These are event-level finite sums with physical normalization. $F$ is not
fitted to their observed values. $v_\epsilon$ is a second moment, not only a
centered covariance; its drift contribution is intentionally included.

Put $f_H=0$, $f_P=11/6$, and $g_s=f_s-1$. Directly from
\eqref{v10:eq:full-table}, with $Z_j$ denoting the actual reader,
\begin{equation}
 d^{-1}\E[Z_{j+1}-Z_j\mid s_j=s,Z_j=z]
       =f_s\{F(z)+a_\epsilon(z)\}.
 \label{v10:eq:raw-drift}
\end{equation}
Although $\pi f=1$, $g_H=-1$ and $g_P=5/6$ do not shrink with the cutoff.
There is consequently no reason for the raw conditional drift minus $F$ to
be small in an occupation $L^2$ norm. Averaging $f$ without controlling its
correlation with $Z$ would not be a proof. The next exact identity supplies
that control.

\section{An explicit phase corrector for the unchanged physical field}
\label{v10:sec:corrector}

The centered phase Poisson equation is solved once and for all:
\begin{equation}
 (I-M)\chi=g,\qquad \pi\chi=0,
 \qquad \chi_H=-\frac{15}{22},\quad \chi_P=\frac{25}{44}.
 \label{v10:eq:chi}
\end{equation}
Indeed $g$ is an eigenvector of $M$ with eigenvalue $-7/15$, so
$\chi=(15/22)g$. In particular $\chi_P-\chi_H=5/4$.
Define the proof-only observable
\begin{equation}
 W(s,z)=z+d\chi_sF(z),\qquad W_j=W(s_j,Z_j).
 \label{v10:eq:corrected-observable}
\end{equation}
Neither the actual field $Z_j$, the action, nor the tick law is replaced by
$W_j$ in any conclusion. Its difference from $Z_j$ is bounded by
$d\max|\chi|\,\|F(Z_j)\|$. The protected phase is not deleted.

\begin{theorem}[Exact corrected drift and a uniform visible-noise bound]
\label[theorem]{v10:thm:corrector}
Suppose $F$ is $L$-Lipschitz and bounded by $B$ on the field chart and on all
successors under consideration. Suppose
\begin{equation}
 \sup_z\|a_\epsilon(z)\|\le A_\epsilon,
 \qquad \sup_z v_\epsilon(z)\le V_\epsilon.
 \label{v10:eq:event-bounds}
\end{equation}
On the full conditional history, the corrected mean increment is exactly
\begin{equation}
 \begin{aligned}
 d^{-1}\E[W_{j+1}-W_j\mid s,z]
 ={}&F(z)+f_s a_\epsilon(z)\\
 &+\mathbf1_{\{s=P\}}\frac{2\vartheta}{3}\chi_H
        \int[F(y)-F(z)]K_\epsilon(z,\dd y).
 \end{aligned}
 \label{v10:eq:corrected-drift}
\end{equation}
If $F$ is extended to a common neighborhood containing $W(s,z)$ with the
same bounds, the residual from $F(W_j)$ is at most
\begin{equation}
 \mathfrak b_\epsilon
 :=C\{A_\epsilon+\vartheta\epsilon L\sqrt{V_\epsilon}+dLB\}.
 \label{v10:eq:corrected-bias}
\end{equation}
For $dL\le1$, the centered corrected increment
$\zeta_j=W_{j+1}-W_j-\E[W_{j+1}-W_j\mid\mathcal F_j]$ obeys
\begin{equation}
 \E[\|\zeta_j\|^2\mid\mathcal F_j]
 \le C d\{\epsilon V_\epsilon+dB^2\}.
 \label{v10:eq:corrected-noise}
\end{equation}
The constants are independent of $\vartheta\in(0,1]$, of the number of
attainable records, and of the dimension in the declared Hilbert norm.
\end{theorem}
\begin{proof}
Write $s',y$ for the next phase and field. Then
\[
 W(s',y)-W(s,z)=(y-z)
  +d(\chi_{s'}-\chi_s)F(z)
  +d\chi_{s'}[F(y)-F(z)].
\]
The mean of the first term divided by $d$ is
$f_s(F+a_\epsilon)$. The mean of the second term divided by $d$ is
$(M\chi-\chi)_sF=-g_sF$. Only an accepted $P\to H$ transition changes $z$,
so the mean of the last term divided by $d$ is the second line of
\eqref{v10:eq:corrected-drift}. This proves the exact identity. Cauchy--Schwarz
on the event kernel gives
$\int\|y-z\|K_\epsilon\le\epsilon\sqrt{V_\epsilon}$.
Subtract $F(W(s,z))$ and use
$\|F(W)-F(z)\|\le Ld|\chi_s|B$ to prove
\eqref{v10:eq:corrected-bias}.

For the variance, use the same decomposition and Lipschitz continuity:
\[
 \|W(s',y)-W(s,z)\|^2
 \le 2(1+d\|\chi\|_\infty L)^2\|y-z\|^2
       +2d^2|\chi_{s'}-\chi_s|^2B^2.
\]
The conditional second moment of $y-z$ is at most
$(2\vartheta/3)\epsilon^2 V_\epsilon$.
Since $(2\vartheta/3)\epsilon^2=(11/6)d\epsilon$ and
$|\chi_{s'}-\chi_s|\le5/4$, the right side has the bound in
\eqref{v10:eq:corrected-noise}. Centering cannot increase the second moment.
No factor $1/\vartheta$ remains after the physical mean-duration calibration.
\end{proof}

The corrector is determined by the \emph{known clock phase}, not by choosing
a potential or conditioning on a favorable future. The term involving
$F(y)-F(z)$ is retained precisely because the field changes on completion;
ignoring it would be an unjustified independence assumption. The smallness
of this term is a consequence of the actual accepted jump bound.

\section{Physical-time canonical convergence from accepted-event estimates}
\label{v10:sec:handoff}

\begin{theorem}[Same-clock nonlinear Hamiltonian transfer]
\label[theorem]{v10:thm:handoff}
Consider cutoff families of \eqref{v10:eq:full-table}. Retain
\eqref{v10:eq:physical-calibration}, the common bounds of
\cref{v10:thm:corrector}, and a common globally Lipschitz bounded extension
of $F$ from the declared chart. Let $z_*(t)$ solve
$\dot z_*=F(z_*)$ with specified $z_*(0)$ on $[0,T+1]$.
For the unchanged piecewise-constant actual reader
$Z^d(t)=Z_{\lfloor t/d\rfloor}$, there is a constant depending only on
$T,L,B$ such that
\begin{equation}
 \boxed{
 \E\sup_{t\le T}\|Z^d(t)-z_*(t)\|^2
 \le C_{T,L,B}\left[
 \E\|Z_0-z_*(0)\|^2+\mathfrak b_\epsilon^2
       +\epsilon V_\epsilon+dB^2+d^2\right].}
 \label{v10:eq:path}
\end{equation}
No stationary initial phase or equilibrium field population is assumed.
In particular, uniformly bounded $V_\epsilon$, $A_\epsilon\to0$, and
$\epsilon\to0$ imply chronological convergence from convergent initial data.
If also $A_\epsilon=O(\epsilon)$ and the initial mean-square error is
$O(\epsilon)$, the squared error is $O(\epsilon)$, uniformly even when
$\vartheta\to0$.
\end{theorem}
\begin{proof}
For $n\le\lceil T/d\rceil$, the exact corrected recursion is
\[
 W_n=W_0+d\sum_{j<n}F(W_j)+d\sum_{j<n}r_j+M_n,
 \qquad \|r_j\|\le\mathfrak b_\epsilon,
 \qquad M_n=\sum_{j<n}\zeta_j.
\]
Orthogonality of martingale differences and the Hilbert-valued $L^2$ maximal
inequality give
\[
 \E\max_{n\le\lceil T/d\rceil}\|M_n\|^2
 \le4\sum_{j<\lceil T/d\rceil}\E\|\zeta_j\|^2
 \le C(T+d)(\epsilon V_\epsilon+dB^2).
\]
The deterministic reference satisfies
\[
 z_*((j+1)d)-z_*(jd)=dF(z_*(jd))+r_j^*,\qquad
 \|r_j^*\|\le LBd^2/2,
\]
by integrating its equation and using $\|\dot z_*\|\le B$.
Subtract the two recursions. If
$u_n=\max_{i\le n}\|W_i-z_*(id)\|$, then
\[
 u_n\le \|W_0-z_*(0)\|+\max_{i\le n}\|M_i\|
 +(T+d)(\mathfrak b_\epsilon+LBd/2)+Ld\sum_{j<n}u_j.
\]
The elementary discrete Gronwall bound multiplies the first terms by at
most $e^{L(T+d)}$. Square and take expectations. Finally
$\|W_j-Z_j\|\le d\|\chi\|_\infty B$ and
$\|z_*(t)-z_*(\lfloor t/d\rfloor d)\|\le Bd$.
These bounds also control the corrected initial value and prove
\eqref{v10:eq:path}. Because $d\le4\epsilon/11$, the rate assertions follow
from \eqref{v10:eq:corrected-bias}. No asymptotic replacement of the clock
by a Poisson process is used.
\end{proof}

\begin{corollary}[A local chart is preserved with a quantified failure probability]
\label[corollary]{v10:cor:exit}
Suppose the event bounds and the Lipschitz extension hold for every departure
from an open safe chart $D$, including the possible terminal successors, and
$\operatorname{dist}(z_*(t),E\setminus D)\ge\rho>0$ for $t\le T$.
Let $\tau$ be the first tick whose physical reader leaves $D$.
The right side of \eqref{v10:eq:path} bounds
\[
 \E\max_{n\le\lfloor T/d\rfloor}
 \|Z_{n\wedge\tau}-z_*((n\wedge\tau)d)\|^2.
\]
Consequently the probability of exit by time $T$ is at most that bound divided
by $\rho^2$. Exiting paths remain failures; they are not conditioned away.
\end{corollary}
\begin{proof}
In the preceding proof multiply every increment by
$\mathbf1_{\{j<\tau\}}$, which is measurable before tick $j+1$.
The stopped martingale has no larger quadratic-variation bound, and the
reference recursion stops at the same random index. The same running-supremum
Gronwall proof applies. At an exit tick the reader is outside $D$ whereas
its reference is at distance at least $\rho$ from the complement; the error
is therefore at least $\rho$. Markov's inequality gives the claim.
\end{proof}

\begin{theorem}[Weak canonical first variation without the raw drift norm]
\label[theorem]{v10:thm:weak}
Under the global hypotheses above, let
$\psi\in C_c^1((0,T);E)$ and use its zero extension when a terminal tick
straddles $T$. For all sufficiently small $d$ the unchanged record has
\begin{equation}
 \mathcal R_d(\psi)=-\int_0^T\langle Z^d,\dot\psi\rangle\dd t
                  -\int_0^T\langle F(Z^d),\psi\rangle\dd t
 \label{v10:eq:weak-residual}
\end{equation}
and
\begin{equation}
 \boxed{
 \begin{aligned}
 \|\mathcal R_d(\psi)\|_{L^2(\PP)}
 \le C_T\{&[A_\epsilon+\vartheta\epsilon L\sqrt{V_\epsilon}
                  +\sqrt{\epsilon V_\epsilon+dB^2}]\|\psi\|_\infty\\
         &+dB\|\dot\psi\|_{L^1}\}.
 \end{aligned}}
 \label{v10:eq:weak-bound}
\end{equation}
For $F=\mathbb J\nabla\mathscr H$ this is the corrected canonical
first-variation estimate on the test $\eta=-\mathbb J\psi$.
It does not require
$\sum_jd\,\E\|d^{-1}\E[\Delta Z_j\mid\mathcal F_j]-F(Z_j)\|^2\to0$.
\end{theorem}
\begin{proof}
Set $C_j=d\chi_{s_j}F(Z_j)$, so $W_j=Z_j+C_j$.
Equation \eqref{v10:eq:corrected-drift} and the centered increments give
\[
 \Delta Z_j=dF(Z_j)+d r_j^0+\zeta_j-\Delta C_j,
 \qquad \|r_j^0\|\le C(A_\epsilon+\vartheta\epsilon L\sqrt{V_\epsilon}).
\]
The distributional time derivative of the unchanged step reader is the sum
of its actual jumps at $(j+1)d$. Pair this identity with
$\psi((j+1)d)$. The martingale part has $L^2$ norm bounded by
$C_T\|\psi\|_\infty\sqrt{\epsilon V_\epsilon+dB^2}$, by
\eqref{v10:eq:corrected-noise}. The drift error is bounded by its supremum
times $T\|\psi\|_\infty$. Summation by parts of the correction has no
boundary term because of the test collars, and is bounded by
$d\|\chi\|_\infty B\|\dot\psi\|_1$.
Replacing $d\psi((j+1)d)$ in the $F(Z_j)$ sum by
$\int_{jd}^{(j+1)d}\psi(t)\dd t$ costs at most
$dB\|\dot\psi\|_1$. These are exactly the two integrals in
\eqref{v10:eq:weak-residual}. The canonical interpretation follows from
\eqref{v9:eq:correct-variation}, since $\psi=\mathbb J\eta$.
\end{proof}

This is a proof-only coboundary removal, not a physical smoothing operation.
The large raw conditional drift is integrated against the actual fast phase,
and the remainder is bounded. Its cancellation cannot be obtained by simply
substituting the stationary mean phase into a field-dependent equation.
Uniform Lipschitz and event-moment bounds on a growing full PDE state are
additional requirements; finite dimensionality alone does not provide them.

\section{A finite canonical reader on the supplied non-Poisson clock}
\label{v10:sec:rotor}

Reuse the circular physical reader of Version 9, but not its exponential
jump scheduler. Let $n\ge5$, $\epsilon=2\pi/n$, and
\begin{equation}
 z_n(\ell)=R_*e^{\epsilon\ell\mathbb J}e_1,
 \qquad \mathbb J=\begin{pmatrix}0&1\\-1&0\end{pmatrix},
 \qquad \ell\in\Z/n\Z.
 \label{v10:eq:rotor-reader}
\end{equation}
On an accepted completion move to $\ell+1$ with probability $(1+v)/2$
and to $\ell-1$ with probability $(1-v)/2$, where $|v|\le1$ is fixed.
Use exactly \eqref{v10:eq:full-table} and tick duration
$d=4\vartheta\epsilon/11$. For $|v|<1$ both accepted directions have
positive probability. The signed target is fixed as
\begin{equation}
 \mathscr H_v(z)=\tfrac v2\|z\|^2,\qquad F(z)=v\mathbb Jz.
 \label{v10:eq:rotor-target}
\end{equation}
This is an explicit finite marked extension of the inherited timing law.
It tests timing and canonical motion. Its accepted rotation is \emph{not}
asserted to be the locked field kernel or its common Einstein--SM action.
In particular the symbol $v$ is not identified with the locked sign Read.

\begin{proposition}[Convergence with a nonvanishing instantaneous phase residual]
\label[proposition]{v10:prop:raw-obstruction}
Put
\begin{equation}
 a_\epsilon=\frac{\cos\epsilon-1}{\epsilon},\qquad
 b_\epsilon=v\frac{\sin\epsilon}{\epsilon}.
 \label{v10:eq:rotor-coefficients}
\end{equation}
The accepted mean is
$m_\epsilon(z)=a_\epsilon z+b_\epsilon\mathbb Jz$, and
$v_\epsilon(z)=2R_*^2(1-\cos\epsilon)/\epsilon^2\le R_*^2$.
For every initial phase and rotor value,
\begin{equation}
 \E\sup_{t\le T}\|Z^d(t)-e^{tv\mathbb J}Z_0\|^2\le C_{T,R_*,v}\epsilon,
 \label{v10:eq:rotor-path}
\end{equation}
and the weak canonical variation in \cref{v10:thm:weak} vanishes.
Nevertheless, in the invariant product law of phase $\pi$ and uniform rotor,
\begin{equation}
 \boxed{
 \E\left\|d^{-1}\E[\Delta Z_j\mid\mathcal F_j]-v\mathbb JZ_j\right\|^2
 =R_*^2\left[\frac{11}{6}(a_\epsilon^2+b_\epsilon^2)
                    -2vb_\epsilon+v^2\right]
 \longrightarrow\frac56R_*^2v^2.}
 \label{v10:eq:nonvanishing-raw}
\end{equation}
For $v\ne0$, the strong raw occupation-residual entrance from Version 8
therefore fails on this converging family.
\end{proposition}
\begin{proof}
Average the two rotation matrices. Each accepted increment has squared norm
$2R_*^2(1-\cos\epsilon)$, giving the moment formulas. The event drift error
from $v\mathbb Jz$ is $O(\epsilon)$ uniformly on the circle.
The circle is invariant and the target orbit remains on it. A bounded
Lipschitz extension on a containing compact chart permits
\cref{v10:thm:handoff,v10:thm:weak} to apply, proving the positive conclusions.
The phase transition does not depend on the rotor, and each accepted
rotation preserves its uniform law, so the displayed product law is
invariant. In that law the raw drift equals
$f_s(a_\epsilon z+b_\epsilon\mathbb Jz)$.
Since $z\perp\mathbb Jz$, $\pi f=1$, and $\pi f^2=11/6$, expanding the
square proves \eqref{v10:eq:nonvanishing-raw}.
This does not contradict the sufficient theorem of Version 8; it shows that
its small strong residual is not necessary before the clock phase is removed
from the integrated error.
\end{proof}

\section{An exact chronological martingale and its retained fluctuation coefficient}
\label{v10:sec:angle}

Let $I_j$ be the accepted-completion indicator at tick $j+1$, and let
$\sigma_j\in\{-1,1\}$ be its resolved direction, with mean $v$ conditional
on acceptance. Write the actual unwrapped accepted angular reward
\begin{equation}
 A_m=A_0+\epsilon\sum_{j<m}I_j\sigma_j,
 \qquad Z_m=R_*e^{A_m\mathbb J}e_1.
 \label{v10:eq:angle}
\end{equation}
The unwrapped reward is used only for the proof; the physical rotor remains
finite. Define
\begin{equation}
 \mathcal M_m=A_m-A_0-mdv+dv(\chi_{s_m}-\chi_{s_0}).
 \label{v10:eq:angle-martingale}
\end{equation}

\begin{theorem}[Exact clock-corrected current variance]
\label[theorem]{v10:thm:angle}
For the stated finite table $\mathcal M_m$ is an exact martingale. Its
conditional one-tick variances are
\begin{align}
 q_H&=\tfrac14d^2v^2,\\
 q_P&=\tfrac{2\vartheta}{3}\epsilon^2
                    -3\vartheta\epsilon d v^2+\tfrac{89}{24}d^2v^2.
 \label{v10:eq:angle-q}
\end{align}
Their stationary mean is
\begin{equation}
 \boxed{\pi_Hq_H+\pi_Pq_P=d\epsilon\,D_{\vartheta,v},\qquad
 D_{\vartheta,v}=1-\frac{104}{121}\vartheta v^2
            =(1-v^2)+v^2c_{\rm clk}^2.}
 \label{v10:eq:diffusion-coefficient}
\end{equation}
Hence, from the stationary phase entrance,
\begin{equation}
 \E\mathcal M_m^2=md\epsilon D_{\vartheta,v}.
 \label{v10:eq:angle-variance}
\end{equation}
From a fixed phase $s_0$, let $\widehat q_s=q_s/(d\epsilon)$. Then exactly
\begin{equation}
 \E\mathcal M_m^2=d\epsilon\left[
 mD_{\vartheta,v}
 +(\widehat q_{s_0}-D_{\vartheta,v})
       \frac{1-(-7/15)^m}{1+7/15}\right].
 \label{v10:eq:angle-initial}
\end{equation}
In particular $17/121\le D_{\vartheta,v}\le1$ and
\begin{equation}
 \E\max_{m\le\lfloor T/d\rfloor}|A_m-A_0-mdv|^2
 \le C(T\epsilon+d^2+d\epsilon).
 \label{v10:eq:angle-max}
\end{equation}
All statements retain the actual renewal law rather than its mean or a
Poisson surrogate.
\end{theorem}
\begin{proof}
The conditional mean angular increment is $df_sv$. The conditional mean of
$dv(\chi_{s_{j+1}}-\chi_{s_j})$ is $-dvg_s$ by
\eqref{v10:eq:chi}. Subtracting $dv$ therefore leaves zero, proving the
martingale assertion. At an $H$ tick its possible martingale increments are
$-dv$ with probability $1/5$ and $dv/4$ with probability $4/5$, giving $q_H$.
At a $P$ tick they are $-dv$ on a $P$ survivor, $-9dv/4$ on a rejected
completion, and $\epsilon\sigma-9dv/4$ on an accepted completion. Their
mean is zero, and their second moment is
\[
 \tfrac13d^2v^2+\tfrac23\left[
 (1-\vartheta)\tfrac{81}{16}d^2v^2
 +\vartheta\{\epsilon^2-\tfrac92\epsilon dv^2
                             +\tfrac{81}{16}d^2v^2\}\right],
\]
which is $q_P$. Substitute $d/\epsilon=4\vartheta/11$ and
$\pi=(5/11,6/11)$ to obtain \eqref{v10:eq:diffusion-coefficient}.
The final expression in that equation uses
\eqref{v10:eq:clock-moments}; it is an algebraic decomposition into mark
variance and timing variance, not a replacement of the joint path law.

Martingale orthogonality sums the conditional variances. Every centered
function on the two-state phase has eigenvalue $-7/15$, so
\[
 \E_{s_0}\widehat q_{s_j}
 =D_{\vartheta,v}+(-7/15)^j(\widehat q_{s_0}-D_{\vartheta,v}).
\]
Sum this geometric sequence to prove both variance formulas. The normalized
$q_s$ are uniformly bounded for $0<\vartheta\le1$, $|v|\le1$.
Doob's inequality and
$|dv(\chi_{s_m}-\chi_{s_0})|\le5d/4$ now prove
\eqref{v10:eq:angle-max}. The lower bound on $D_{\vartheta,v}$ follows
from \eqref{v10:eq:diffusion-coefficient}.
\end{proof}

Even a deterministic accepted direction ($|v|=1$) retains timing variance.
At $\vartheta=1$ its coefficient is $17/121$, not zero. For rare acceptance
$\vartheta\to0$ the coefficient tends to one, consistently with the
inherited rare-response limit, but no such limit is needed in the proof.
The physical path error vanishes because it is multiplied by $\epsilon$.
The accepted direction and the mean physical cycle time remain independent
source data whose compatibility with a specified Hamiltonian must be tested.

\section{Exact slow dispersion of the canonical reader on the same clock}
\label{v10:sec:dispersion}

For the first circular character use
\begin{equation}
 \eta_\epsilon=\cos\epsilon+\ii v\sin\epsilon,\qquad
 B_{\vartheta,\epsilon}=
 \begin{pmatrix}
 1/5&4/5\\
 \frac23(1-\vartheta+\vartheta\eta_\epsilon)&1/3
 \end{pmatrix}.
 \label{v10:eq:mode-matrix}
\end{equation}
The sign $+\ii$ here encodes the formal positive accepted angle in
$e^{A\mathbb J}$; real reader correlations are obtained by its real and
imaginary coefficients, not by changing the convention for $\mathbb J$.

\begin{theorem}[Source-specific slow eigenvalue and leading path error]
\label[theorem]{v10:thm:dispersion}
The exact eigenvalues of \eqref{v10:eq:mode-matrix} are
\begin{equation}
 \lambda_\pm=\frac{4\pm\sqrt{121+120\vartheta(\eta_\epsilon-1)}}{15},
 \label{v10:eq:mode-eigenvalues}
\end{equation}
where the square root is the analytic branch equal to $11$ at
$\epsilon=0$. Uniformly for $\vartheta\in(0,1]$, $|v|\le1$, as
$\epsilon\to0$,
\begin{equation}
 \boxed{d^{-1}\log\lambda_+
 =\ii v-\frac{\epsilon}{2}D_{\vartheta,v}+O(\epsilon^2).}
 \label{v10:eq:slow-dispersion}
\end{equation}
Let the initial phase have law $\pi$, independently of a fixed initial rotor.
For $t_m=md\le T$, its angular characteristic factor is
\begin{equation}
 g_m=\pi B_{\vartheta,\epsilon}^{\,m}\one
       =c_+\lambda_+^m+c_-\lambda_-^m,
 \quad
 c_-=\frac{\lambda_+-g_1}{\lambda_+-\lambda_-},
 \quad c_+=1-c_-,
 \label{v10:eq:mode-coefficients}
\end{equation}
where $g_1=1+(4\vartheta/11)(\eta_\epsilon-1)$ and
$c_-=O(\vartheta^2\epsilon^2)$. Consequently
\begin{equation}
 \sup_{md\le T}
 \left|g_m-\exp\{t_m(\ii v-\epsilon D_{\vartheta,v}/2)\}\right|
 \le C_T\epsilon^2,
 \label{v10:eq:mode-uniform}
\end{equation}
and the terminal physical error has the expansion
\begin{equation}
 \boxed{
 \E\|Z_m-e^{t_mv\mathbb J}Z_0\|^2
 =R_*^2t_m\epsilon D_{\vartheta,v}+O_T(\epsilon^2).}
 \label{v10:eq:terminal-leading}
\end{equation}
The constants do not deteriorate as acceptance becomes rare.
\end{theorem}
\begin{proof}
Conditioning one tick of the marked source on its phase multiplies the
first character by the matrix \eqref{v10:eq:mode-matrix}; this proves its
meaning. Its characteristic polynomial is
\[
 \lambda^2-\tfrac8{15}\lambda-\tfrac7{15}
                      -\tfrac8{15}\vartheta(\eta_\epsilon-1),
\]
which gives \eqref{v10:eq:mode-eigenvalues}.
Put $x=\vartheta(\eta_\epsilon-1)$. On a fixed disk about zero the
analytic expansions are
\[
 \lambda_+=1+\tfrac4{11}x-\tfrac{120}{1331}x^2+O(x^3),\qquad
 \log\lambda_+=\tfrac4{11}x-\tfrac{208}{1331}x^2+O(x^3).
\]
Now
$\eta_\epsilon-1=\ii v\epsilon-\epsilon^2/2+O(\epsilon^3)$.
Divide by $d=4\vartheta\epsilon/11$ to get
\[
 d^{-1}\log\lambda_+
 =\ii v-\tfrac\epsilon2
                 +\tfrac{52}{121}\vartheta v^2\epsilon+O(\epsilon^2).
\]
All remainders have at least one factor $\vartheta$ before division, so the
bound is uniform down to $\vartheta=0$. This proves
\eqref{v10:eq:slow-dispersion}.

For small $\epsilon$, the eigenvalue gap is uniformly positive and
$|\lambda_-|$ is uniformly below one. The scalar sequence $g_m$ obeys the
same second-order characteristic recurrence, with $g_0=1$ and the stated
$g_1$; solving for its two coefficients gives
\eqref{v10:eq:mode-coefficients}. Since $g_1$ is exactly the linear part of
$\lambda_+$ in $x$, $c_-=O(x^2)=O(\vartheta^2\epsilon^2)$.
Exponentiating \eqref{v10:eq:slow-dispersion} for $md\le T$ gives a uniform
$O_T(\epsilon^2)$ error. The fast-mode coefficient is already of that
order; this proves \eqref{v10:eq:mode-uniform} even at $m=0$.
Finally $\|Z_m\|=\|Z_0\|=R_*$ and
\[
 \E\|Z_m-e^{t_mv\mathbb J}Z_0\|^2
   =2R_*^2\{1-\Re(e^{-\ii vt_m}g_m)\}.
\]
Insert \eqref{v10:eq:mode-uniform} and expand the remaining real exponential.
This gives \eqref{v10:eq:terminal-leading}.
\end{proof}

This dispersion formula concerns the supplied two-phase timing mechanism
with a non-idempotent circular accepted operation. It is not the scalar
survivor or the idempotent payload formula already proved in \cite{GT}.
It independently verifies the fluctuation coefficient obtained from the
pathwise martingale. Replacing the clock by exponential waiting would set
that coefficient to one and would miss its finite-$\vartheta$ correction.

\section{What is discharged, and what remains an actual field-source question}
\label{v10:sec:ledger}

The strong raw drift mismatch of Versions 8--9 is an unnecessarily restrictive
entrance when the microscopic phase is explicitly retained. For the actual
completed-private clock its centered phase equation is solved by
\eqref{v10:eq:chi}; the corrected identity then proves chronological and
weak canonical convergence under the accepted-event bank
\eqref{v10:eq:event-bank}. The finite circular example makes this strict:
its raw drift error has the positive limit \eqref{v10:eq:nonvanishing-raw},
while its actual physical path error tends to zero at the explicitly proved
rate. This is not achieved by assuming that a nonexistent instantaneous
physical generator already agrees with the target.

The source distinction is now:
\begin{center}\small
\begin{tabularx}{\textwidth}{@{}L{0.31\textwidth}X@{}}
\toprule
Quantity & Status in this iteration\\
\midrule
Protected-tick phase and accepted waiting law & Inherited, actually specified
source; its finite moments and Poisson corrector are evaluated explicitly.\\[0.3em]
Physical tick and event normalization & $\epsilon=11d/(4\vartheta)$ is retained
as a measured calibration; no comparison-index substitution.\\[0.3em]
Accepted field increment & Must be supplied on the same completed event;
its mean discrepancy and second moment are \eqref{v10:eq:event-bank}.\\[0.3em]
Fast-phase contribution to the drift & Removed from the integrated error by an
explicit proof-only coboundary, without altering physical records.\\[0.3em]
Canonical direction and Hamiltonian & Still independently identified. The
clock cannot choose an unmeasured field kernel or fit its Hamiltonian.\\[0.3em]
Einstein--SM stress and constraints & Spatial/geometric positive screens,
full physical variation, constraint propagation, and source/action incidence
remain separate.\\
\bottomrule
\end{tabularx}
\end{center}

The actual phase transition has therefore been used rather than replaced by
a convenient continuous-time sampler. The circular field kernel merely
establishes that this exact timing law is compatible with a nondissipative
classical limit when its independently supplied accepted updates have the
right directed tangent. It does not identify the locked opportunity
instrument's operator-valued outputs with those classical rotations, nor
turn the accepted count anchor into the complete common action.

The next noncircular source calculation is the accepted-event canonical
increment conditioned on its \emph{resolved completed history}, with the
physical mean duration retained. One should compare its mean with the fixed
common Hamiltonian and its second moment with the physical action topology;
there is no longer a need to force the known alternating clock phase itself
into a vanishing raw $L^2$ drift residual. Any genuine state dependence or
memory of the accepted operation must be retained in that joint event bank.
A full Einstein--SM theorem still needs the remaining sector and geometric
hypotheses of \cref{v8:cor:screen-composition}; they have not been inferred
from the phase corrector.

\paragraph{Verification and preservation.}
All V9 source bytes preceding its final document command are preserved.
All new labels have prefix \texttt{v10:}. The accompanying script checks the
inherited clock transform and its differentiated moments, the exact Poisson
solution, the nonlinear corrected-drift identity on finite rows, the
phase-resolved martingale variances, invariant full-table matrices,
nonvanishing raw drift, and both the exact dispersion and its leading-error
coefficient. Symbolic equalities and finite numerical tests are diagnostics,
not machine-checked proofs or simulations of the complete Einstein--SM model.

\begingroup
\renewcommand{\refname}{Additional references for Version 10}
\begin{thebibliography}{9}
\small\raggedright
\bibitem{v10:Kurtz}
T.~G.~Kurtz,
\emph{Strong approximation theorems for density dependent Markov chains},
Stochastic Processes and their Applications \textbf{6} (1978), 223--240.
\href{https://doi.org/10.1016/0304-4149(78)90020-0}{doi:10.1016/0304-4149(78)90020-0}.
Background for martingale and deterministic-limit comparisons; the exact
clock-specific formulas above are proved directly.
\bibitem{v10:MS}
M.~M.~Meerschaert and P.~Straka,
\emph{Semi-Markov approach to continuous time random walk limit processes},
Annals of Probability \textbf{42} (2014), 1699--1723.
\href{https://arxiv.org/abs/1206.1960}{arXiv:1206.1960};
\href{https://doi.org/10.1214/13-AOP905}{doi:10.1214/13-AOP905}.
Background for retaining chronological state variables in non-Poisson
processes. No anomalous-diffusion or physical field identification is
imported from this reference.
\end{thebibliography}
\endgroup
% END IMMUTABLE VERSION 10 BODY
% APPEND VERSION 11 HERE, BEFORE THE FINAL END-DOCUMENT COMMAND.


% BEGIN IMMUTABLE VERSION 11 BODY
\clearpage
\begin{center}
{\Large\bfseries Version 11: The resolved locked instrument before its field semantics}\\[0.5em]
{\large Exact event moments, a sharp locking noise floor, and coherent small-event realizations}
\end{center}
\addcontentsline{toc}{part}{Version 11: Resolved locked events and coherent trajectory limits}

\section{Upstream audit: use the instrument, not an invented event table}
\label{v11:sec:audit}

Version 10 treats the actual protected-tick timing mechanism, but its accepted
field kernel is independently supplied. Merely asking that kernel to have
small moments would repeat an entrance condition. Here we return to the
\emph{resolved six-edge Kraus operators} of the locked opportunity instrument.
Their operator entries are not numerically specified by the outcome identities
alone. Nevertheless those identities allow an exact, useful calculation of
what their resolved trajectories must do and of what an admissible coherent
small-event family can do.

\paragraph{Inherited results, not new claims.}
The supplied Grand Tensor monograph \cite{GT}, at
\nolinkurl{thm:GT76-locked-normalform}, gives
\begin{equation}
 \sum_{e=1}^6K_e^*K_e=I,\qquad
 \sum_{e=1}^6K_e=\sqrt2 I,\qquad
 \mathcal E=\tfrac13\operatorname{id}+\tfrac23\Psi,
 \label{v11:eq:inherited-lock}
\end{equation}
where \(\mathcal E(\rho)=\sum_eK_e\rho K_e^*\) and \(\Psi\) is a channel.
Its normal form parametrizes \(K_e\) by five Kraus operators for \(\Psi\).
The theorem \nolinkurl{thm:GT76-locked-pullback} identifies the additional
operator-valued panels that reconstruct these matrices. The separate
\nolinkurl{thm:SMST-channel-native-tangent} characterizes a Hamiltonian
channel tangent on an identified logical carrier. None of these structural
results, nor the timing proof of \cref{v10:thm:handoff}, is claimed anew.

The additions are a resolved-event moment identity in the \emph{original}
Kraus outcomes, a sharp noise constraint imposed by the locked identity
fraction, a trajectory consequence on the supplied non-Poisson clock, and an
explicit full-rank family satisfying the exact locking identities. Its
coherent tangent may be prescribed independently, but is not selected by
locking. We also retain the cost in the Kraus source Gram and the initial
pinching of different protected record sectors.

\paragraph{Which state is being read.}
The reader in this body is the conditional pure-state projector of a supplied
finite-dimensional instrument, with a specified pure preparation. It is a
well-defined function of the preparation and complete resolved record. It is
\emph{not} automatically the classical coframe, Higgs, canonical momentum,
or Einstein--SM phase-space reader. A final proposition gives the separate
intertwining test for a proposed classical decoder. No quantum continuum,
Born-to-classical field identification, or microscopic attraction theorem is
inferred from a finite matrix trajectory.

\paragraph{Attribution.}
Resolved quantum trajectories and their continuous-time limits are established
subjects; see \cite{v11:AP,v11:BBB}. The calculations here are elementary
finite Kraus, conditional-moment, and renewal-clock calculations. No priority
is asserted for the general trajectory formalism. Their purpose is the
particular \(1/3\)-identity locking constraint, the unchanged resolved
outcomes, and the source-norm and sector conditions required in this lineage.

\section{Exact event moments of the original resolved outcomes}
\label{v11:sec:moments}

Work in \(\mathscr H=\C^r\), \(r<\infty\). The physical matrix-reader norm
in this body is the \emph{unnormalized} Hilbert--Schmidt norm
\(\|X\|_2^2=\tr(X^*X)\); on Hermitian matrices its real inner product is
\(\tr(XY)\). This normalization is different from the coefficient Gram
normalization explicitly introduced in \cref{v11:sec:Gram}.
For a pure projector \(P=|\psi\rangle\langle\psi|\), \(\|\psi\|=1\), set
\begin{equation}
 w_e(P)=\tr(K_ePK_e^*),\qquad
 P_e=\frac{K_ePK_e^*}{w_e(P)}\quad(w_e(P)>0).
 \label{v11:eq:instrument-row}
\end{equation}
Zero-probability outcomes contribute zero to every sum. Every positive-weight
\(P_e\) is a pure projector. The label \(e\) remains the original edge
outcome: no rotation of the Kraus outcome basis is performed.

\begin{proposition}[The pure-state event's first and second displacement moments]
\label[proposition]{v11:prop:event-moments}
For every trace-preserving single-Kraus-per-outcome instrument,
\begin{align}
 m(P)&:=\sum_ew_e(P)(P_e-P)=\mathcal E(P)-P,
       \label{v11:eq:event-mean}\\
 s_2(P)&:=\sum_ew_e(P)\|P_e-P\|_2^2
       =2\{1-\tr(P\mathcal E(P))\},
       \label{v11:eq:event-square}\\
 v_2(P)&:=\sum_ew_e(P)\|P_e-P-m(P)\|_2^2
       =s_2(P)-\|m(P)\|_2^2.
       \label{v11:eq:event-var}
\end{align}
In particular the coarse channel determines these two scalar displacement
moments of the resolved pure-state reader. It need not determine the full
resolved trajectory law, its labels, or the distribution of other nonlinear
readers.
\end{proposition}
\begin{proof}
Trace preservation gives \(\sum_ew_e=1\), and summing \(w_eP_e\) gives
\(\mathcal E(P)\). Since \(P^2=P\), \(P_e^2=P_e\), and both have trace
one, \(\|P_e-P\|_2^2=2-2\tr(PP_e)\). Multiply by \(w_e\) and sum.
The final identity is the finite Hilbert variance identity. Purity is used
in the second formula; it is not an assertion for arbitrary mixed conditional
states or instruments whose resolved outcomes contain several unrecorded
Kraus terms.
\end{proof}

Fix a bounded Hermitian matrix \(H\), independently of the observed channel,
and put
\begin{equation}
 F_H(P)=-\ii[H,P],\qquad
 a_\epsilon=\sup_{P\ \mathrm{pure}}
 \left\|\frac{\mathcal E_\epsilon(P)-P}{\epsilon}-F_H(P)\right\|_2.
 \label{v11:eq:tangent-audit}
\end{equation}
The positive number \(\epsilon=11d/(4\vartheta)\) is the mean physical
accepted-event duration of Version 10. It is not chosen again to fit \(H\).

\begin{theorem}[A coherent channel tangent supplies the resolved noise budget]
\label[theorem]{v11:thm:tangent-moments}
For the actual resolved pure-state event law,
\begin{equation}
 \left\|\epsilon^{-1}m(P)-F_H(P)\right\|_2\le a_\epsilon,
 \qquad
 0\le s_2(P)\le2\epsilon a_\epsilon.
 \label{v11:eq:coherent-moments}
\end{equation}
Consequently \(a_\epsilon\to0\) makes the visible second moment per mean
physical time tend to zero. An \(O(\epsilon^2)\) channel remainder gives
\(s_2(P)=O(\epsilon^2)\), uniformly in the initial pure state, without a
separate noise or lower-outcome-probability assumption.
\end{theorem}
\begin{proof}
The first statement is the definition. Write
\(\mathcal E_\epsilon(P)-P=\epsilon F_H(P)+R_\epsilon(P)\).
Cyclicity and \(P^2=P\) give \(\tr(PF_H(P))=0\). Thus
\(s_2(P)=-2\tr(PR_\epsilon(P))\), which is nonnegative by
\cref{v11:prop:event-moments} and at most \(2\|R_\epsilon(P)\|_2\), since
\(\|P\|_2=1\). No division by an individual \(w_e\) enters the bound.
\end{proof}

The statement does not say that the unknown locked channel has this tangent.
It says that, \emph{once the same operator channel passes the tangent audit},
its original resolved pure-state events cannot carry an additional hidden
order-one displacement variance. This implication is stronger than an
independent first-moment and second-moment hypothesis for this reader.

\section{A sharp noise floor forced by the locked identity fraction}
\label{v11:sec:floor}

The upper estimate above is compatible with a nonzero lower noise bound at
second order. The locking identity forces such a bound.

\begin{theorem}[Exact locked resolved-variance identity]
\label[theorem]{v11:thm:noise-floor}
Suppose \(\mathcal E=\lambda\operatorname{id}+(1-\lambda)\Psi\), where
\(0<\lambda<1\) and \(\Psi\) is a channel. For a pure input \(P\), put
\(S=\Psi(P)\) and use the original instrument's \(m,s_2,v_2\). Then
\begin{equation}
 \boxed{
 v_2(P)=\frac{\lambda}{1-\lambda}\|m(P)\|_2^2
             +(1-\lambda)\{1-\tr(S^2)\}.}
 \label{v11:eq:noise-floor}
\end{equation}
Hence the locked value \(\lambda=1/3\) gives
\begin{equation}
 v_2(P)\ge\tfrac12\|m(P)\|_2^2,\qquad
 s_2(P)\ge\tfrac32\|m(P)\|_2^2.
 \label{v11:eq:locked-floor}
\end{equation}
If \(m_\epsilon(P)/\epsilon\to F_H(P)\), then
\begin{equation}
 \liminf_{\epsilon\to0}\epsilon^{-2}v_2(P)
                  \ge\tfrac12\|F_H(P)\|_2^2.
 \label{v11:eq:asymptotic-floor}
\end{equation}
Thus a nonzero coherent displacement cannot have an arbitrarily smaller
resolved conditional variance. The lower bound is per event: it does not
force a nonvanishing diffusion on a time interval with \(O(\epsilon^{-1})\)
accepted events.
\end{theorem}
\begin{proof}
Here \(m=(1-\lambda)(S-P)\) and
\(s_2=2(1-\lambda)(1-\tr(PS))\). Direct subtraction gives
\[
 s_2-\frac{\|m\|_2^2}{1-\lambda}
 =(1-\lambda)\{1-\tr(S^2)\}.
\]
Subtract \(\|m\|_2^2\) once more to obtain
\eqref{v11:eq:noise-floor}. A density matrix has \(\tr(S^2)\le1\).
The other assertions follow by substitution and division by \(\epsilon^2\).
The decomposition of \(\mathcal E\) is used algebraically, not as a new
resolved random choice between an identity event and a \(\Psi\) event.
\end{proof}

\begin{corollary}[The corresponding operator-valued multiplicative defect]
\label[corollary]{v11:cor:Schwarz}
Let \(\mathcal E^\dagger\) and \(\Psi^\dagger\) be the unital Heisenberg
duals. For every matrix \(A\),
\begin{align}
 &\mathcal E^\dagger(A^*A)
          -\mathcal E^\dagger(A)^*\mathcal E^\dagger(A)\notag\\
 &\quad\succeq
 \frac{\lambda}{1-\lambda}
 (\mathcal E^\dagger(A)-A)^*(\mathcal E^\dagger(A)-A).
 \label{v11:eq:Schwarz-floor}
\end{align}
In particular a locked channel cannot equal a nontrivial unitary channel on
the full matrix algebra at a fixed event scale.
\end{corollary}
\begin{proof}
Expand the left side as the sum of \((1-\lambda)\) times the
multiplicative defect of \(\Psi^\dagger\) and
\(\lambda(1-\lambda)(\Psi^\dagger(A)-A)^*(\Psi^\dagger(A)-A)\).
The first term is positive. Directly, with its Kraus isometry \(V\), it is
\(V^*(A^*\otimes I)(I-VV^*)(A\otimes I)V\).
This proves the inequality. A unitary channel has zero multiplicative defect;
then \(\mathcal E^\dagger(A)=A\) for every \(A\), so that channel is the
identity. This finite-event obstruction allows a nontrivial first-order
unitary limit.
\end{proof}

\section{Resolved trajectories on the supplied non-Poisson clock}
\label{v11:sec:clock}

Initially restrict to one protected sign--type sector, so that the accepted
operators on the edge system are exactly \(K_{e,\epsilon}\), and rejection
leaves the conditional state unchanged. At each accepted response retain its
actual label \(e\), probability \(w_e(P)\), and terminal projector \(P_e\).
Use the full two-phase table of \eqref{v10:eq:full-table} with that event
kernel and with \(d=4\vartheta\epsilon/11\). The resulting trajectory
\(P^d(t)=P_{\lfloor t/d\rfloor}\) is the original resolved instrument's
conditional state, not the averaged state \(\mathcal E(P)\).

\begin{theorem}[Coherent locked-event trajectory limit with no separate event-noise hypothesis]
\label[theorem]{v11:thm:trajectory}
Suppose \(\|H\|_{\rm op}\le M\), \(\epsilon\to0\), and the actual channels
obey \(a_\epsilon\to0\) in \eqref{v11:eq:tangent-audit}. For any initial
phase and pure preparation \(P_0\),
\begin{equation}
 \boxed{
 \E\sup_{t\le T}
 \|P^d(t)-e^{-\ii tH}P_0e^{\ii tH}\|_2^2
 \le C_{T,M}(a_\epsilon+a_\epsilon^2+\epsilon).}
 \label{v11:eq:trajectory-bound}
\end{equation}
The constants are independent of \(\vartheta\in(0,1]\), the number of
resolved paths, and \(r\), with the displayed norm convention and common
bound on \(H\). In particular \(a_\epsilon=O(\epsilon)\) gives mean-square
uniform error \(O(\epsilon)\).
For every fixed compactly supported \(C^1\) Hermitian test \(A(t)\),
\begin{equation}
 -\int\tr(P^d\dot A)\dd t
 -\int\tr(-\ii[H,P^d]A)\dd t\longrightarrow0
 \quad\hbox{in }L^2(\PP).
 \label{v11:eq:weak-von-Neumann}
\end{equation}
This is a weak von Neumann equation for the conditional-state reader, not yet
the classical Einstein--SM first variation.
\end{theorem}
\begin{proof}
Use the real Hilbert space of Hermitian matrices. On the pure-state orbit,
\(\|F_H(P)\|_2\le2M\) and
\(\|F_H(X)-F_H(Y)\|_2\le2M\|X-Y\|_2\).
A fixed Hilbert-ball cutoff supplies a bounded globally Lipschitz extension
agreeing with this field on the orbit and the proof-only corrected states;
its constants depend on \(M\), not \(r\).

By \cref{v11:thm:tangent-moments}, the event hypotheses of
\cref{v10:thm:corrector} hold with
\(A_\epsilon=a_\epsilon\), \(V_\epsilon=2a_\epsilon/\epsilon\).
The latter need not stay bounded. Its actual contribution in
\eqref{v10:eq:path} is \(\epsilon V_\epsilon=2a_\epsilon\), and the
corrected bias is bounded by
\[
 C_M\{a_\epsilon+\vartheta\sqrt{\epsilon a_\epsilon}+d\}.
\]
The proofs of \cref{v10:thm:corrector,v10:thm:handoff} use only a bounded
Lipschitz vector field; their stochastic estimates do not use the fixed
canonical-matrix representation of that field. They therefore apply to this
matrix reader. Substitution in \eqref{v10:eq:path}, with
\(d\le4\epsilon/11\), proves \eqref{v11:eq:trajectory-bound}.
This application uses the same conditional-state law and every original
resolved edge outcome.

Finally subtract the weak equation of the deterministic unitary path.
Cauchy--Schwarz and the Lipschitz estimate bound the difference by
\[
 \sup_{t\le T}\|P^d(t)-P_*(t)\|_2
       \bigl(\|\dot A\|_{L^1_t}+2M\|A\|_{L^1_t}\bigr).
\]
Equation \eqref{v11:eq:trajectory-bound} proves the weak assertion. Each
finite physical time horizon contains finitely many protected ticks; for a
fixed preparation its attainable resolved tree is finite, even if the union
over all horizons has infinitely many distinct conditional states.
\end{proof}

This result turns a verified \emph{operator} tangent into the two event
moment estimates needed by the actual clock theorem. It does not require
replacing the original unraveling by a specially selected Kraus basis.
Purity, complete outcome resolution, preparation, and the mean event duration
are load-bearing inputs. Mixed preparations require their own moment argument
or a declared purification with control of the amplified channel; they are
not covered by asserting that \eqref{v11:eq:event-square} still holds.

\section{An exact six-outcome small-event realization for a prescribed Hamiltonian}
\label{v11:sec:family}

The locking identities might conceivably prevent the preceding tangent
entrance altogether. The following explicit family answers that algebraic
question, without identifying it with the unspecified physical edge matrices.

Choose \(r\ge3\), a nonzero traceless Hermitian target \(H\), and four fixed
matrices \(C_2,\ldots,C_5\) such that
\(I,H,C_2,\ldots,C_5\) are linearly independent over \(\C\). Put
\(R=\sum_{\mu=2}^5C_\mu^*C_\mu\). Let \(O\) be a real \(6\times5\)
isometry onto \(\one^\perp\), with first column
\begin{equation}
 O_{\cdot1}=6^{-1/2}(1,1,1,-1,-1,-1)^{\mathsf T}.
 \label{v11:eq:O}
\end{equation}
Complete this column by any fixed orthonormal basis of its orthogonal
complement inside \(\one^\perp\). For sufficiently small positive
\(\epsilon\), define
\begin{align}
 S_\epsilon&=(I-\epsilon^4R)^{1/2},&
 U_\epsilon&=\exp(-\tfrac32\ii\epsilon H),\notag\\
 B_{1,\epsilon}&=U_\epsilon S_\epsilon,&
 B_{\mu,\epsilon}&=\epsilon^2C_\mu\quad(2\le\mu\le5),
       \label{v11:eq:B-family}\\
 K_{e,\epsilon}&=\frac{I}{\sqrt{18}}
        +\sqrt{\frac23}\sum_{\mu=1}^5O_{e\mu}B_{\mu,\epsilon}.
       \label{v11:eq:K-family}
\end{align}
All these objects are specified before a trajectory is generated.

\begin{theorem}[Exact locking, full finite rank, and an arbitrary coherent tangent]
\label[theorem]{v11:thm:family}
For all sufficiently small \(\epsilon>0\), \eqref{v11:eq:K-family} satisfies
all three identities in \eqref{v11:eq:inherited-lock}, including independence
of its six Kraus operators. Its accepted channel is
\begin{equation}
 \mathcal E_\epsilon(\rho)
 =\tfrac13\rho+\tfrac23\left[
 U_\epsilon S_\epsilon\rho S_\epsilon U_\epsilon^*
       +\epsilon^4\sum_{\mu=2}^5C_\mu\rho C_\mu^*\right].
 \label{v11:eq:family-channel}
\end{equation}
In the Hilbert--Schmidt induced superoperator norm,
\begin{equation}
 \mathcal E_\epsilon
 =\operatorname{id}+\epsilon\mathcal F_H
       +\tfrac34\epsilon^2\mathcal F_H^2+O(\epsilon^3),
 \qquad \mathcal F_H(X)=-\ii[H,X].
 \label{v11:eq:family-expansion}
\end{equation}
Every original edge outcome has probability bounded away from zero, uniformly
on pure preparations for small \(\epsilon\). Nevertheless its actual
resolved trajectory on the supplied clock converges to
\(e^{-\ii tH}P_0e^{\ii tH}\) with the error
\eqref{v11:eq:trajectory-bound} of order \(O(\epsilon)\).

After tensoring with the four nonzero protected record projections \(Q_a\),
the \(24\) accepted operators \(K_{e,\epsilon}\otimes Q_a\) remain linearly
independent. The additional no-response operator is the original scalar
identity operator. No accepted outcome has been replaced by a fictitious
identity branch.
\end{theorem}
\begin{proof}
We have \(B_{1,\epsilon}^*B_{1,\epsilon}=I-\epsilon^4R\), so
\(\sum_{\mu=1}^5B_{\mu,\epsilon}^*B_{\mu,\epsilon}=I\).
The inherited normal form gives the two exact sum identities and
\eqref{v11:eq:family-channel}. To check the rank condition, consider
\[
 I,\quad\frac{B_{1,\epsilon}-I}{\epsilon},\quad
 C_2,\ldots,C_5.
\]
This list converges to
\(I,-\tfrac32\ii H,C_2,\ldots,C_5\), which is linearly independent.
A nonzero minor remains nonzero for sufficiently small \(\epsilon\).
Invertible column operations and nonzero rescalings therefore show that
\(I,B_{1,\epsilon},\ldots,B_{5,\epsilon}\) is independent. The fixed
orthogonal Kraus-coordinate transformation in the inherited normal form then
proves independence of the six \(K_e\).

The spectral square root obeys \(S_\epsilon=I+O(\epsilon^4)\). Thus the
bracketed channel in \eqref{v11:eq:family-channel} differs from
\(\exp(\tfrac32\epsilon\mathcal F_H)\) by \(O(\epsilon^4)\). Its
exponential series gives \eqref{v11:eq:family-expansion}. The constants are
finite functions of the declared matrix norms; no inverse time scale is
suppressed. Hence \(a_\epsilon=O(\epsilon)\), and
\cref{v11:thm:trajectory} applies.

For the outcome probabilities, the leading scalar Kraus coefficients are
\begin{equation}
 c_e=\frac1{\sqrt{18}}+\sqrt{\frac23}O_{e1}
     =\frac1{\sqrt{18}}\mathbin{\pm}\frac13.
 \label{v11:eq:scalar-coefficients}
\end{equation}
Both values are nonzero. The first-order perturbation is
\(-\tfrac32\ii\epsilon\sqrt{2/3}\,O_{e1}H\), so
\(K_e^*K_e=c_e^2I+O(\epsilon^2)\). Therefore
\(w_e(P)\ge\tfrac12\min_ec_e^2>0\) for small \(\epsilon\).
Compression to each record block proves independence of the tensor-product
accepted operators. Their sum is \(\sqrt2 I\), so the no-response identity
lies in their span, as in the inherited locked instrument.
\end{proof}

\begin{corollary}[The locking noise floor is attained to leading order]
\label[corollary]{v11:cor:sharp}
For the family above and each pure \(P\),
\begin{align}
 m_\epsilon(P)&=\epsilon F_H(P)+O(\epsilon^2),\notag\\
 s_{2,\epsilon}(P)&=\tfrac32\epsilon^2\|F_H(P)\|_2^2+O(\epsilon^3),\notag\\
 v_{2,\epsilon}(P)&=\tfrac12\epsilon^2\|F_H(P)\|_2^2+O(\epsilon^3).
 \label{v11:eq:sharp-family}
\end{align}
Thus \(1/2\) in \eqref{v11:eq:locked-floor} is sharp even within
full-six-rank locking families at every nonzero cutoff.
\end{corollary}
\begin{proof}
The commutator superoperator \(\mathcal F_H\) is skew-adjoint in the
Hilbert--Schmidt inner product, so
\(\tr(P\mathcal F_H^2(P))=-\|F_H(P)\|_2^2\).
Insert \eqref{v11:eq:family-expansion} in
\cref{v11:prop:event-moments}. The final variance subtracts
\(\epsilon^2\|F_H(P)\|_2^2+O(\epsilon^3)\).
For any \(P\) not commuting with \(H\), the ratio to
\(\|m_\epsilon(P)\|_2^2\) tends to \(1/2\).
\end{proof}

This is a positive \emph{identity-class realization}. It uses the same
six-edge locking equations and the same protected-tick timing law, unlike a
freely chosen classical rotation kernel. But its \(B_\mu\) are explicitly
chosen to test feasibility, not acquired from the locked Einstein--SM source.
In particular \(H\) and \(-H\) can both be prescribed in the construction;
the acceptance probability and the number of records do not select the
physical Hamiltonian or its sign.


\begin{theorem}[The resolved locking variance and the actual clock variance add]
\label[theorem]{v11:thm:clock-coefficient}
For the family of \cref{v11:thm:family}, start at phase \(H\) with a pure
state \(P_0\), and use the original resolved Kraus outcomes and protected
ticks. Uniformly for \(0<t_0\le t_n=nd\le T\) and
\(0<\vartheta\le1\),
\begin{equation}
 \boxed{
 \E\|P_n-e^{-\ii t_nH}P_0e^{\ii t_nH}\|_2^2
 =\epsilon t_n\left(\frac32-\frac{104}{121}\vartheta\right)
        \|[H,P_0]\|_2^2+O_{t_0,T}(\epsilon^2).}
 \label{v11:eq:clock-coefficient}
\end{equation}
The coefficient splits exactly at leading order as
\begin{equation}
 \frac32-\frac{104}{121}\vartheta
 =\underbrace{\frac12}_{\text{locking event variance}}
  +\underbrace{1-\frac{104}{121}\vartheta}_{\text{actual clock variance}}.
 \label{v11:eq:variance-split}
\end{equation}
It is positive even at \(\vartheta=1\). This describes the actual resolved
trajectory error, not merely the deviation of its averaged state.
\end{theorem}
\begin{proof}
First replace \(\mathcal E_\epsilon\), only for this error calculation,
by \(\mathcal E_\epsilon^0=\tfrac13\operatorname{id}
+\tfrac23\exp(\tfrac32\epsilon\mathcal F_H)\).
The difference has trace-norm induced bound \(C\epsilon^4\) on Hermitian
inputs, directly from \eqref{v11:eq:family-channel}. The two-phase mean-state
tick maps are trace-preserving positive maps on block-diagonal phase--system
matrices, hence contractions in the sum trace norm. Their one-tick difference
is at most \(C\vartheta\epsilon^4\). Telescoping \(n\) ticks gives
mean-state error at most \(Cn\vartheta\epsilon^4\le C_T\epsilon^3\).
This is a comparison of averages for the proof, not a replacement of the
physical instrument or its resolved record.

For \(\mathcal E_\epsilon^0\), every eigendirection of the skew-adjoint
\(\mathcal F_H\) is propagated by the exact scalar two-phase matrix of
Version 10 with accepted multiplier
\[
 \eta_\epsilon=\tfrac13+\tfrac23\exp(\tfrac32\epsilon f)
 =1+\epsilon f+\tfrac34\epsilon^2f^2+O(\epsilon^3),
 \qquad f\in\ii[-2\|H\|,2\|H\|].
\]
Its slow root is
\(\lambda_+=(4+\sqrt{121+120\vartheta(\eta_\epsilon-1)})/15\).
Expanding the square root and then the logarithm gives
\begin{equation}
 \frac{\log\lambda_+}{d}
 =f+\epsilon\left(\tfrac34-\tfrac{52}{121}\vartheta\right)f^2
       +O(\epsilon^2).
 \label{v11:eq:operator-dispersion}
\end{equation}
The remainders are uniform in \(\vartheta\): before division by
\(d=4\vartheta\epsilon/11\), every nonconstant omitted term carries a
factor \(\vartheta\epsilon^3\).
For an initial phase \(H\), the first two scalar mean multipliers are
\(u_0=u_1=1\), and therefore
\[
 u_n=A_+\lambda_+^n+A_-\lambda_-^n,\quad
 A_-=(\lambda_+-1)/(\lambda_+-\lambda_-)
       =\tfrac{30}{121}\vartheta\epsilon f+O(\epsilon^2),\quad
 A_+=1-A_-.
\]
The fast root has modulus at most a fixed \(q<1\) for small \(\epsilon\).
Since \(d\le4\epsilon/11\), its contribution is \(O_{t_0}(\epsilon^2)\)
on the stated positive-time interval. Spectral calculus for the normal
operator \(\mathcal F_H\) consequently gives
\[
 \E P_n=e^{t_n\mathcal F_H}
 \left[P_0+\epsilon t_n c_\vartheta\mathcal F_H^2P_0
             -\tfrac{30}{121}\vartheta\epsilon\mathcal F_HP_0\right]
             +O_{t_0,T}(\epsilon^2),\quad
 c_\vartheta=\tfrac34-\tfrac{52}{121}\vartheta.
\]
The same formula applies to the original mean up to the already bounded
\(O(\epsilon^3)\) difference. Because every original resolved state and the
target are pure,
\(\E\|P_n-P_*(t_n)\|_2^2=2-2\tr(P_*(t_n)\E P_n)\).
Skew-adjointness gives
\(\langle P_0,\mathcal F_HP_0\rangle=0\) and
\(\langle P_0,\mathcal F_H^2P_0\rangle=-\|[H,P_0]\|_2^2\).
Substitution proves \eqref{v11:eq:clock-coefficient}.
The splitting \eqref{v11:eq:variance-split} then uses the sharp event
coefficient from \cref{v11:cor:sharp} and the clock coefficient from
\cref{v10:prop:clock-moments}. No exponential waiting-time substitution is
made.
\end{proof}

\section{What becomes singular: the dynamical Kraus Gram, not the outcome labels}
\label{v11:sec:Gram}

A full-rank family at every cutoff need not have a cutoff-uniform source
floor. This distinction cannot be suppressed when the construction is
composed with a physical positive-screen theorem. Define the normalized
coefficient Hilbert product
\(\langle A,B\rangle_{\rm c}=r^{-1}\tr(A^*B)\) and the \(6\times6\) Gram
\begin{equation}
 G_{ef}=r^{-1}\tr(K_e^*K_f),\qquad
 k_e=r^{-1}\tr K_e,\qquad
 \eta=\sum_e\|K_e-k_eI\|_{\rm c}^2.
 \label{v11:eq:Kraus-Gram}
\end{equation}
Trace preservation gives \(\tr G=1\) and
\(\eta=1-\sum_e|k_e|^2\ge0\).

\begin{theorem}[A coherent full-algebra tangent necessarily collapses unscaled Kraus contrasts]
\label[theorem]{v11:thm:Gram-collapse}
Let the six eigenvalues of \(G\) be
\(\lambda_1\ge\cdots\ge\lambda_6\ge0\). Then
\begin{equation}
 \sum_{j=2}^6\lambda_j\le\eta.
 \label{v11:eq:Gram-tail}
\end{equation}
If, on the whole matrix algebra,
\begin{equation}
 \|\mathcal E_\epsilon-\operatorname{id}-\epsilon\mathcal F_H\|_{2\to2}
       \le\epsilon b_\epsilon,
 \label{v11:eq:whole-tangent}
\end{equation}
then \(0\le\eta\le\epsilon b_\epsilon\) and
\(\lambda_6\le\epsilon b_\epsilon/5\).
In particular an \(O(\epsilon^2)\) coherent remainder is incompatible
with a fixed positive lower bound on every unscaled Kraus-contrast eigenvalue.
For the explicit family in \cref{v11:thm:family},
\begin{equation}
 \lambda_1\asymp1,\qquad\lambda_2\asymp\epsilon^2,\qquad
 \lambda_3,\ldots,\lambda_6\asymp\epsilon^4.
 \label{v11:eq:Gram-scales}
\end{equation}
\end{theorem}
\begin{proof}
Let \(B:\C^6\to M_r\) synthesize the six \(K_e\) in the normalized
coefficient Hilbert space. The unit vector \(I\) has Rayleigh quotient
\(\langle I,BB^*I\rangle_{\rm c}=\sum_e|k_e|^2=1-\eta\).
Therefore the largest eigenvalue of \(BB^*\), equivalently of \(G=B^*B\),
is at least \(1-\eta\). Subtract from \(\tr G=1\) to obtain
\eqref{v11:eq:Gram-tail}.

In the ordinary matrix-unit superoperator representation,
\(\operatorname{Tr}_{\rm sup}\mathcal E=\sum_e|\tr K_e|^2\), whereas
\(\operatorname{Tr}_{\rm sup}\mathcal F_H=0\).
Taking the superoperator trace of the remainder in
\eqref{v11:eq:whole-tangent} gives
\[
 \eta=-r^{-2}\operatorname{Tr}_{\rm sup}
 (\mathcal E_\epsilon-\operatorname{id}-\epsilon\mathcal F_H)
       \le\epsilon b_\epsilon,
\]
since the superoperator dimension is \(r^2\).

For \eqref{v11:eq:Gram-scales}, the columns
\(I,(B_{1,\epsilon}-I)/\epsilon,C_2,\ldots,C_5\) have a uniformly
positive Gram for small \(\epsilon\), by the rank argument in
\cref{v11:thm:family}. Multiplication by
\(\operatorname{diag}(1,\epsilon,\epsilon^2,\ldots,\epsilon^2)\)
and a fixed uniformly invertible column operation gives
\(I,B_{1,\epsilon},\ldots,B_{5,\epsilon}\). The fixed Kraus-coordinate
transformation gives the original \(K_e\). Singular-value min--max bounds
under uniformly invertible factors give singular-value scales
\(1,\epsilon,\epsilon^2,\ldots,\epsilon^2\); their squares are
\eqref{v11:eq:Gram-scales}.
\end{proof}

\begin{remark}[Which noncollapse is, and is not, ruled out]
This theorem concerns the full-algebra operator Gram of the \(K_e\). It does
not contradict the inherited positive accepted \emph{outcome} likelihood
metric or the fact that there are \(24\) nonzero resolved outcomes at each
cutoff. A large environment can retain distinguishable labels whose system
actions become close to scalar multiples of the identity. Nor does the
whole-algebra estimate apply when only a proper physical reader subalgebra
has a coherent tangent and other Kraus directions act on a retained
spectator. Rescaling small contrast directions can define a tangent-source
bank, but its physical normalization, error amplification, and equality with
the source metric have to be specified; it is not a free restoration of a
missing uniform margin.
\end{remark}

\section{The protected sign--type pinching has an initial-layer cost}
\label{v11:sec:sectors}

The sector qualification in \cref{v11:thm:trajectory} is essential.
Let \(Q_a\), \(1\le a\le4\), be the fixed orthogonal nonzero record
projections. The complete accepted operators are
\(A_{e,a}=K_e\otimes Q_a\). For a pure initial projector \(P_0\) on the
joint system define
\begin{equation}
 p_a=\tr((I\otimes Q_a)P_0),\qquad
 P_{0,a}=(I\otimes Q_a)P_0(I\otimes Q_a)/p_a
       \quad(p_a>0).
 \label{v11:eq:sector-data}
\end{equation}

\begin{proposition}[Exact first-acceptance sector selection]
\label[proposition]{v11:prop:sector}
The first accepted sign--type label has probability \(p_a\), and after its
occurrence the actual conditional state remains in that sector. If
\(p_{\max}:=\max_ap_a<1\), the first accepted state \(P_{\rm acc}\) satisfies
\begin{equation}
 \|P_{\rm acc}-P_0\|_2^2\ge2(1-p_{\max}).
 \label{v11:eq:sector-jump}
\end{equation}
For a fresh completed-boundary entrance the physical first-acceptance time
\(\tau_{\rm acc}\) has expectation \(\epsilon\). Consequently, as
\(\epsilon\to0\), the actual conditional paths cannot converge uniformly
in probability on \([0,T]\) to any deterministic continuous path starting
at \(P_0\).
\end{proposition}
\begin{proof}
Summing over \(e\) gives
\(\sum_e A_{e,a}^*A_{e,a}=I\otimes Q_a\), proving the label probability.
Every subsequent Kraus outcome either preserves that sector or has zero
probability. If \(P_{\rm acc}=|\phi\rangle\langle\phi|\) lies in sector
\(a\), then
\(\tr(P_0P_{\rm acc})=|\langle\psi_0,\phi\rangle|^2\le p_a\).
This gives \eqref{v11:eq:sector-jump}. The waiting mean is the inherited
clock identity. Markov's inequality gives
\(\PP(\tau_{\rm acc}>s)\le\epsilon/s\) for fixed \(s>0\).
A deterministic continuous path remains arbitrarily close to \(P_0\) for
all sufficiently small \(s\), whereas the accepted state is the stated
fixed distance away with probability tending to one before that \(s\).
The triangle inequality proves failure of uniform convergence at the initial
time. No mean-state averaging or conditioning away a sector is used.
\end{proof}

\begin{corollary}[Positive-time random-sector limit]
\label[corollary]{v11:cor:random-sector}
For the family of \cref{v11:thm:family} and a pure initial joint state, let
\(A\) be the actual first accepted sector label. On the same probability
space, for each fixed \(0<s<T\),
\begin{equation}
 \E\sup_{s\le t\le T}
 \left\|P^d(t)-
 (e^{-\ii tH}\otimes I)P_{0,A}(e^{\ii tH}\otimes I)\right\|_2^2
       \le C_{T,H,C}(\epsilon+\epsilon/s).
 \label{v11:eq:sector-limit}
\end{equation}
The initial layer produces a random sector with weights \(p_a\); it is not
a deterministic unpinched initial condition.
\end{corollary}
\begin{proof}
The probability of no acceptance by \(s\) is at most \(\epsilon/s\), and
the squared distance of pure projectors is at most two. Conditional on
\(A=a\), the first accepted edge outcome is exactly the edge instrument
applied to \(P_{0,a}\). The family expansion also holds after tensoring
with the fixed identity record factor, so its mean-square displacement from
\(P_{0,a}\) is \(O(\epsilon^2)\). At that event the clock returns to phase
\(H\). Apply \cref{v11:thm:trajectory} to the subsequent trajectory, with
that conditional initial state. The difference between the reference started
at \(\tau_{\rm acc}\) and the reference started at time zero is bounded by
\(C\tau_{\rm acc}\). By \cref{v10:prop:clock-moments},
\(\E\tau_{\rm acc}^2\le2\epsilon^2\). Conditional edge probabilities do
not affect the acceptance-time law, since total acceptance is scalar.
The triangle inequality, the finite number of record sectors, and the
preceding failure event prove \eqref{v11:eq:sector-limit}.
\end{proof}

\section{From reconstructed operator panels to a proposed classical field reader}
\label{v11:sec:physical}

The operator-valued pullbacks required by the inherited locked-source theorem
are stronger data than outcome counts. Their accuracy must resolve the
physical event scale.

\begin{proposition}[Sufficient operator-panel precision at the event scale]
\label[proposition]{v11:prop:precision}
Suppose two legitimate locked instruments, represented in the same admitted
Stinespring and edge frames, have their ten Hermitian pullback panels within
\(e_{\rm pan}\) in operator norm. Put
\(s_\vartheta=\sqrt{\vartheta(12-11\vartheta)}\).
Then their accepted channels satisfy
\begin{equation}
 \|\mathcal E-\widehat{\mathcal E}\|_{2\to2}
       \le C\,e_{\rm pan}/s_\vartheta.
 \label{v11:eq:panel-error}
\end{equation}
Thus \(e_{\rm pan}=o(\epsilon s_\vartheta)\) is a sufficient worst-case
precision to certify a vanishing tangent mismatch against a fixed \(H\),
and \(e_{\rm pan}=O(\epsilon^2s_\vartheta)\) is sufficient to retain an
\(O(\epsilon^2)\) channel remainder and event second-moment bound.
These are sufficient full-panel budgets, not lower bounds on every possible
target-specific measurement strategy.
\end{proposition}
\begin{proof}
The inherited reconstruction reads
\[
 D_\mu\otimes I=\frac{2\sqrt3}{s_\vartheta}
      \bigl(\mathcal P_\vartheta^*(X_\mu)
                  +\ii\mathcal P_\vartheta^*(Y_\mu)\bigr),
 \qquad K_e=I/\sqrt{18}+\sum_\mu O_{e\mu}D_\mu.
\]
Hence \(\|D_\mu-\widehat D_\mu\|\le4\sqrt3 e_{\rm pan}/s_\vartheta\).
The fixed finite matrix \(O\) bounds the sum of the six Kraus errors by
\(Ce_{\rm pan}/s_\vartheta\). Trace preservation gives
\(\|K_e\|_{\rm op},\|\widehat K_e\|_{\rm op}\le1\). Subtract each
\(K_e X K_e^*\) term in two steps and use
\(\|AXB\|_2\le\|A\|_{\rm op}\|X\|_2\|B\|_{\rm op}\).
This gives \eqref{v11:eq:panel-error}. Dividing by \(\epsilon\) yields
the tangent-error budget. Formula \eqref{v11:eq:event-square} bounds the
additional second-moment error by twice the channel error on a pure state.
\end{proof}

\begin{proposition}[The separate classical-decoder incidence]
\label[proposition]{v11:prop:decoder}
Let \(\Phi\) be an independently specified \(C^2\) decoder from a
neighborhood of the density matrices to a real physical Hilbert space,
with bounded first and second derivatives on connecting segments. Retain the
actual decoded outcomes \(z=\Phi(P)\), \(z_e=\Phi(P_e)\). Then
\begin{align}
 &\left\|\epsilon^{-1}\sum_ew_e(z_e-z)
                   -D\Phi(P)F_H(P)\right\|
      \le(\|D\Phi\|_\infty+\|D^2\Phi\|_\infty)a_\epsilon,
       \label{v11:eq:decoder-drift}\\
 &\sum_ew_e\|z_e-z\|^2
      \le2\|D\Phi\|_\infty^2\epsilon a_\epsilon.
       \label{v11:eq:decoder-noise}
\end{align}
For a separately fixed classical Hamiltonian vector field
\(F_{\rm phys}\), the remaining leading mismatch is exactly
\begin{equation}
 \mathfrak I(P)=D\Phi(P)(-\ii[H,P])-F_{\rm phys}(\Phi(P)).
 \label{v11:eq:classical-incidence}
\end{equation}
For a linear observable decoder \(\Phi_a(P)=\tr(A_aP)\), its first term is
\(\tr(\ii[H,A_a]P)\). Neither locking nor the coherent-state trajectory
limit proves that \(\mathfrak I\) vanishes for a chosen Einstein--SM reader.
\end{proposition}
\begin{proof}
Taylor's integral formula on the segment from \(P\) to \(P_e\) gives
\[
 \|\Phi(P_e)-\Phi(P)-D\Phi(P)(P_e-P)\|
       \le\tfrac12\|D^2\Phi\|_\infty\|P_e-P\|_2^2.
\]
Average, divide by \(\epsilon\), and use
\cref{v11:thm:tangent-moments}. Lipschitz continuity gives the second
inequality. Subtracting the independently fixed \(F_{\rm phys}\) proves
the stated incidence decomposition. The linear-reader expression follows
by cyclicity of the trace. The decoder need not be Markov-sufficient on its
own; all conditional expectations here remain on the complete state and
resolved history.
\end{proof}

\section{Outcome and the remaining actual-source calculation}
\label{v11:sec:ledger}

The accepted-event moment entrance has now been partially discharged for a
precisely declared reader. A Hamiltonian tangent of the actual coarse edge
channel gives the original pure-state outcome law's small second moment
\emph{automatically}; the exact renewal clock then gives chronological
convergence. The locking constraint forces a sharp, but harmless at leading
continuum order, variance of at least half the squared accepted displacement.
An explicit six-rank family realizes that scale for any independently chosen
bounded non-scalar Hamiltonian. Thus locking does not prohibit coherent
small-event dynamics, but it does not determine the physical generator.

The price is explicit. Coherent convergence on the whole matrix algebra
collapses unscaled Kraus-contrast directions. The protected sign--type
pinching has an initial layer unless the input is already in one sector.
A matrix conditional state is not a classical Einstein--SM field. Its
classical reader must pass \eqref{v11:eq:classical-incidence}, and its
physical source norm, locality, gravitational and matter stress, constraints,
and cofinal estimates remain the corresponding obligations of the earlier
notes. This construction does not populate all those conditions by assigning
new semantics to its matrix coordinates.

The next native calculation is therefore not to choose another event kernel.
It is to evaluate the supplied operator-valued locked-source pullbacks at the
correct event precision, extract their \emph{actual} tangent relative to the
independently reconstructed action, and test that tangent on the declared
field--momentum decoder. If a reader subalgebra rather than the whole edge
algebra carries the physical fields, its reducing property and source norm
must be retained before applying the full-algebra contrast obstruction.
The current sources identify these as additional joint operator observations;
they do not provide their numerical values for the intended coupled law.

\paragraph{Verification and preservation.}
The complete Version 10 source before its final document command is retained
byte-for-byte. All new labels have prefix \texttt{v11:}. The diagnostic
script checks the exact original-outcome moment and variance identities,
the operator Schwarz bound, the explicit six-edge normalization and rank,
the sharp noise coefficient, the Kraus-Gram scales, the protected-sector
jump, the decoder Taylor bound, and the mean trajectory on the literal
two-phase tick matrix, and the summed locking/clock fluctuation coefficient. It also checks the \(O(\epsilon)\) terminal error
on the original resolved instrument by a finite superoperator computation.
These finite checks do not replace the written proofs, are not
machine-checked formalizations, and do not simulate an Einstein--SM field.

\begingroup
\renewcommand{\refname}{Additional references for Version 11}
\begin{thebibliography}{9}
\small\raggedright
\bibitem{v11:AP}
S.~Attal and C.~Pellegrini,
\emph{Return to equilibrium for some quantum trajectories}, 2010.
\href{https://arxiv.org/abs/1004.3356}{arXiv:1004.3356}.
Background for resolved discrete instruments and their continuous-time
trajectory interpretation; no equilibrium or stochastic-Schr\"odinger theorem
is imported into the proofs above.
\bibitem{v11:BBB}
M.~Bauer, T.~Benoist, and D.~Bernard,
\emph{Repeated quantum non-demolition measurements: convergence and
continuous time limit}, Annales Henri Poincar\'e \textbf{14} (2013), 639--679.
\href{https://arxiv.org/abs/1206.6045}{arXiv:1206.6045};
\href{https://doi.org/10.1007/s00023-012-0204-x}{doi:10.1007/s00023-012-0204-x}.
Background for sector selection under repeated measurements; the exact
one-acceptance projection and renewal-clock boundary in this body are proved
from the displayed locked branches.
\end{thebibliography}
\endgroup
% END IMMUTABLE VERSION 11 BODY
% APPEND VERSION 12 HERE, BEFORE THE FINAL END-DOCUMENT COMMAND.


% BEGIN IMMUTABLE VERSION 12 BODY
\clearpage
\begin{center}
{\Large\bfseries Version 12: Nonlinear classical readers on the resolved clock}\\[0.5em]
{\large Finite-observable concentration, same-action polynomial sources, and an unbounded-generator branch}
\end{center}
\addcontentsline{toc}{part}{Version 12: Nonlinear observable closure on actual instrument histories}

\section{Audit: the next gap is nonlinear concentration, not another unitary limit}
\label{v12:sec:audit}

Version 11 derives event moments and coherent conditional-state trajectories
from a whole-channel tangent with a bounded generator.  Its final decoder
incidence \eqref{v11:eq:classical-incidence} remains a separate condition.
There are two related difficulties.  An expectation of a nonlinear operator
polynomial need not equal that polynomial of the expectations, and a physically
normalized collective Hamiltonian can have a norm growing with the cutoff.
Consequently a bounded-generator matrix-state theorem cannot simply be
relabelled as a nonlinear classical field theorem.

\paragraph{The substantial overlap is retained as an input, not claimed again.}
The SM--Einstein monograph \cite{SM}, at
\nolinkurl{thm:many-copy-transport} and
\nolinkurl{eq:many-copy-exact-algebra}, already proves transport of empirical
polynomial moments for its specified fixed-mode parent--fermion Hamiltonian.
It also distinguishes a concentrated initial point from a transported
probability measure and explicitly separates the many-copy parameter from
physical species.  We do not reprove that model or its many-copy theorem.
The locked Kraus normalization and complementary pullbacks are already in
\cite{GT}, at \nolinkurl{thm:GT76-locked-normalform} and
\nolinkurl{thm:GT76-locked-pullback}.  The clock, its phase Poisson equation,
and the fact that the corrector is not a physical filtering operation are
\cref{v10:sec:clock,v10:sec:corrector}.  All are used with their original scope.

\paragraph{New target.}
We prove a maximal-in-time concentration estimate for nonlinear macroscopic
observables of \emph{the actual resolved conditional states}, on the supplied
non-Poisson renewal clock.  A finite action-aware bank replaces a tangent
bound on the entire matrix algebra.  Neither purity nor small Hilbert--Schmidt
displacement of the full conditional state is assumed.  A positive
second-moment comparison controls every outcome-conditioned reader, including
its internal quantum variance.  Polynomial action derivatives then factorize
on those same records.  An explicit six-outcome locked family realizes this
entrance with an unbounded full Hamiltonian.  The additional semiclassical
scale and its preparation are declared; they are not asserted to have been
selected by the locked Einstein--SM source.

\paragraph{Attribution and interpretation.}
Classical limits of quantum spin dynamics and collective empirical moments
are established mathematics; see \cite{v12:FKL} and the preceding monograph
result.  Polynomial telescoping, commutator derivations, and positive
supermartingale estimates below are standard tools, proved in the form used.
The contribution of this body is their resolved-instrument, physical-clock,
and finite-observable composition, including constants independent of the
full Hamiltonian norm.  It is not a new claim of priority for a general
mean-field or large-spin limit.  The nonlinear spin realization is a
\emph{diagnostic action-compatible branch}, not the Standard Model or a
newly identified microscopic law of spacetime.

\section{A finite positive comparison for nonlinear operator readers}
\label{v12:sec:polynomial}

Let $A_1,\ldots,A_r$ be self-adjoint matrices on a finite Hilbert space, with
$\|A_a\|_{\rm op}\le M$.  Matrix dimension may vary; $r$ is fixed here.
For $x\in\R^r$, $|x_a|\le M$, define
\begin{equation}
 V_x=\sum_{a=1}^r(A_a-x_aI)^2\succeq0.
 \label{v12:eq:V}
\end{equation}
For a real commutative polynomial $p$, $\operatorname{Sym}p(A)$ denotes its
complete symmetric ordering.  Constants map to scalar multiples of $I$.
The order convention is fixed as part of the action representation.

\begin{lemma}[Polynomial concentration without an assumed moment factorization]
\label[lemma]{v12:lem:poly}
For each fixed polynomial $p$ there is $L_{p,M}<\infty$, independent of
matrix dimension and of the commutators, such that
\begin{equation}
 (\operatorname{Sym}p(A)-p(x)I)^2\preceq L_{p,M}^2 V_x.
 \label{v12:eq:poly-order}
\end{equation}
If $p(x)=\sum_w c_w x_{w_1}\cdots x_{w_{|w|}}$ is any fixed monomial
expansion, one may take
\[L_{p,M}=\sum_{|w|\ge1}|c_w||w|\max(1,M)^{|w|-1}.\]
For every density matrix $\rho$ and $z_a=\tr(\rho A_a)$,
\begin{align}
 |\tr(\rho\operatorname{Sym}p(A))-p(x)|^2
   &\le L_{p,M}^2\tr(\rho V_x),\\
 \tr(\rho V_x)&=|z-x|^2+
                 \sum_a\tr\bigl(\rho(A_a-z_aI)^2\bigr).
 \label{v12:eq:variance-split}
\end{align}
\end{lemma}
\begin{proof}
For an ordered word use the exact identity
\[
 A_{w_1}\cdots A_{w_\ell}-x_{w_1}\cdots x_{w_\ell}I
 =\sum_{j=1}^{\ell}A_{w_1}\cdots A_{w_{j-1}}
       (A_{w_j}-x_{w_j}I)\prod_{k=j+1}^{\ell}x_{w_k}.
\]
All factors to the right of the difference are scalars.  Thus for every
vector $u$ the norm of this word difference applied to $u$ is at most
$\ell\max(1,M)^{\ell-1}\langle u,V_xu\rangle^{1/2}$.
Average over orderings and sum the polynomial coefficients.  The resulting
operator is self-adjoint for real $p$, so this norm estimate proves
\eqref{v12:eq:poly-order} as a quadratic-form inequality.
The density-matrix Cauchy--Schwarz inequality gives
$|\tr\rho B|^2\le\tr\rho B^2$ for self-adjoint $B$.
Expanding each square around $z_a$ proves \eqref{v12:eq:variance-split}.
No products of expectations have been inserted into an operator equation.
\end{proof}

Let $\widehat{\mathcal H}$ be the independently fixed self-adjoint matrix
Hamiltonian and write
\begin{equation}
 \mathscr D O=\ii[\widehat{\mathcal H},O].
 \label{v12:eq:D}
\end{equation}
Let $F:\R^r\to\R^r$ be a fixed polynomial vector field.  On the actual
operator bank suppose
\begin{equation}
 \mathscr D A_a=\operatorname{Sym}F_a(A)+R_a,
 \qquad \kappa^2:=\sum_a\|R_a\|_{\rm op}^2.
 \label{v12:eq:action-incidence}
\end{equation}
This is an action/reader identity, or an explicitly bounded operator defect;
it is not obtained by fitting $F$ to observed trajectories.

\begin{proposition}[One positive observable controls nonlinear action incidence]
\label[proposition]{v12:prop:Lyapunov}
Let $x(t)$ solve $\dot x=F(x)$ on a fixed interval, with $|x_a(t)|\le M$.
Put $V(t)=V_{x(t)}$.  Then
\begin{equation}
 \dot V(t)+\mathscr D V(t)\preceq C_F V(t)+\kappa^2 I,
 \qquad C_F=2\Bigl(\sum_a L_{F_a,M}^2\Bigr)^{1/2}+1.
 \label{v12:eq:operator-Lyapunov}
\end{equation}
For every density matrix $\rho$ and $z_a=\tr(\rho A_a)$, the nonlinear
reader incidence satisfies
\begin{equation}
 \sum_a|\tr(\rho\mathscr D A_a)-F_a(z)|^2
 \le C_F'\sum_a\tr\bigl(\rho(A_a-z_aI)^2\bigr)+2\kappa^2.
 \label{v12:eq:conditional-incidence}
\end{equation}
\end{proposition}
\begin{proof}
Write $e_a=A_a-x_aI$ and
$B_a=\operatorname{Sym}F_a(A)-F_a(x)I$.  The derivation rule gives
\[
 \dot V+\mathscr D V=\sum_a\{e_aB_a+B_ae_a+e_aR_a+R_ae_a\}.
\]
For each vector $u$, \cref{v12:lem:poly} implies
$\sum_a\|B_au\|^2\le (\sum_aL_{F_a,M}^2)\langle u,Vu\rangle$.
Cauchy--Schwarz bounds the first two terms by
$2(\sum_aL_{F_a,M}^2)^{1/2}\langle u,Vu\rangle$.
The remaining terms are bounded by
$2\kappa\|u\|\langle u,Vu\rangle^{1/2}
 \le\langle u,Vu\rangle+\kappa^2\|u\|^2$.
This proves the operator inequality.  Apply the same polynomial lemma at
$x=z$, add the $R_a$ terms, and square with $(a+b)^2\le2a^2+2b^2$ to prove
\eqref{v12:eq:conditional-incidence}.
\end{proof}

\subsection{A same-action algebra that verifies, rather than assumes, the incidence}
Let $\nu\downarrow0$ and suppose a fixed finite set of real structure constants
satisfies the Jacobi identity and
\begin{equation}
 [A_a^{(\nu)},A_b^{(\nu)}]
   =\ii\nu\sum_c c_{ab}^{\ c}A_c^{(\nu)},
 \qquad \|A_a^{(\nu)}\|\le M.
 \label{v12:eq:Lie}
\end{equation}
For a fixed real polynomial $\mathcal H$ set
\begin{equation}
 \widehat{\mathcal H}_\nu=\nu^{-1}\operatorname{Sym}\mathcal H(A^{(\nu)}),
 \qquad
 F_a(x)=\sum_{b,c}c_{ab}^{\ c}x_c\partial_b\mathcal H(x).
 \label{v12:eq:Lie-action}
\end{equation}
The associated Poisson bracket is
$\{x_a,x_b\}=\sum_c c_{ab}^{\ c}x_c$, with $\dot x_a=\{x_a,\mathcal H\}$.

\begin{proposition}[Exact coordinate commutator and uniform fixed-word bounds]
\label[proposition]{v12:prop:Lie}
For \eqref{v12:eq:Lie-action},
\begin{equation}
 \mathscr D_\nu A_a^{(\nu)}=\operatorname{Sym}F_a(A^{(\nu)})
 \label{v12:eq:exact-Lie}
\end{equation}
exactly.  Moreover, for every fixed ordered polynomial $O(A^{(\nu)})$ and
fixed integer $j$, $\|\mathscr D_\nu^j O(A^{(\nu)})\|$ is bounded
independently of $\nu$.  No bound on $\|\widehat{\mathcal H}_\nu\|$ follows
or is needed.
\end{proposition}
\begin{proof}
On one symmetrized monomial, the commutator with $A_a$ replaces each factor
$A_b$ by $\ii\nu\sum_c c_{ba}^{\ c}A_c$.  Multiplication by $\ii/\nu$
and antisymmetry of $c_{ba}^{\ c}$ give exactly the symmetrized polynomial
$\sum_{b,c}c_{ab}^{\ c}x_c\partial_b\mathcal H$.
For a general ordered word, apply the derivation rule and replace a factor
$A_a$ by the fixed polynomial in \eqref{v12:eq:exact-Lie}.  One commutator
therefore produces a finite sum of bounded words with coefficients
independent of $\nu$.  Iterating a fixed number of times proves the assertion.
This argument does not assert the exact Poisson rule for two arbitrary
symmetrized polynomials; only the coordinate identity is exact.
\end{proof}

\section{The actual instrument requires only a finite action-aware bank}
\label{v12:sec:bank}

At each cutoff let $\{\mathscr K_e\}$ be the actual finite-outcome completely
positive instrument on density matrices, with trace-preserving sum
$\mathscr E_*$.  Its adjoint $\mathscr E^*$ is unital.  At an accepted event,
from the actual current state $\rho$ retain the original label $e$ with
probability $p_e=\tr\mathscr K_e(\rho)$ and update to
$\rho_e=\mathscr K_e(\rho)/p_e$ when $p_e>0$.
No new outcome resolution is imposed.  Conditional states may be mixed.
For any matrix $O$ the exact conditional identity is
\begin{equation}
 \E[\tr(\rho_e O)\mid\rho]=\tr(\rho\mathscr E^*O).
 \label{v12:eq:instrument-dual}
\end{equation}

Use precisely the clock in \eqref{v10:eq:full-table}:
$M=\left(\begin{smallmatrix}1/5&4/5\\2/3&1/3\end{smallmatrix}\right)$,
acceptance on a completion with probability $\vartheta$, and
$d=4\vartheta\epsilon/11$.  Rejection does not change $\rho$.
Set $f_H=0,f_P=11/6$ and retain
$\chi_H=-15/22$, $\chi_P=25/44$, so
$(I-M)\chi=f-1$.  This is the inherited physical timing law, not an
exponential replacement.

Write
\begin{equation}
 \mathcal O=\{A_a,A_a^2:1\le a\le r\}.
 \label{v12:eq:basic-bank}
\end{equation}
Assume the following \emph{same-action} operator estimates, with constants
independent of the cutoff:
\begin{align}
 &\max_{O\in\mathcal O}
   \|\epsilon^{-1}(\mathscr E^*O-O)-\mathscr D O\|\le\eta,
       &&0\le\eta\le1,\\
 &\max_{O\in\mathcal O}\|\mathscr E^*(\mathscr D O)-\mathscr D O\|
       \le C_1\epsilon,
       &&\max_{O\in\mathcal O}\|\mathscr D O\|\le C_1.
 \label{v12:eq:bank-bounds}
\end{align}
These involve at most $4r$ fixed operator arguments:
$A_a,A_a^2,\mathscr D A_a,\mathscr D(A_a^2)$.
The operators can be nonlinear polynomials of the physical reader.  Their
actual pullbacks, not separate scalar means, are tested.  The initial
preparation will be tested separately by $\tr(\rho_0V(0))$.

\begin{lemma}[A positive phase-corrected error on every resolved record]
\label[lemma]{v12:lem:corrector}
Retain \eqref{v12:eq:action-incidence}, \eqref{v12:eq:bank-bounds}, and a
bounded target solution $x$ on $[0,T+1]$.  Put
$C(t)=\mathscr D V(t)$.  There is $K_0$ independent of matrix dimension,
$\vartheta$, and $\|\widehat{\mathcal H}\|$, such that
\begin{equation}
 \widehat V_s(t)=V(t)+d\chi_sC(t)+dK_0I
 \quad\hbox{satisfies}\quad
 V(t)\preceq\widehat V_s(t)\preceq V(t)+2dK_0I.
 \label{v12:eq:positive-corrector}
\end{equation}
For the actual conditional state $\rho_j$ at tick $j$, set
$W_j=\tr(\rho_j\widehat V_{s_j}(jd))$.  Then
\begin{equation}
 \E[W_{j+1}\mid\mathcal F_j]
 \le(1+C_2d)W_j+C_2d(\kappa^2+\eta+d)
 \label{v12:eq:one-step}
\end{equation}
for $jd\le T$, where $\mathcal F_j$ includes the entire resolved history.
\end{lemma}
\begin{proof}
The identity
$C(t)=\sum_a\mathscr D(A_a^2)-2\sum_ax_a(t)\mathscr D A_a$
and \eqref{v12:eq:bank-bounds} bound $C$ and $\dot C$ uniformly.
The polynomial ODE on a bounded range bounds $\dot x$ and $\ddot x$,
hence $\dot V$ and $\ddot V$.  Choose
$K_0\ge\max_s|\chi_s|\sup_t\|C(t)\|$; this proves positivity and the upper
bound in \eqref{v12:eq:positive-corrector}.

For clarity the conditional calculation can be done in operator form before
pairing with $\rho_j$.  Put $p_H=0,p_P=2\vartheta/3$, so that
$p_s\epsilon=df_s$.  The uncorrected increment is
\[
 V(t+d)-V(t)+p_s(\mathscr E^*-I)V(t+d)
 =d\{\dot V(t)+f_s C(t)\}+\mathcal R_1,
 \qquad \|\mathcal R_1\|\le C(d^2+d\eta).
\]
Here the scalar $|x(t+d)|^2I$ has zero channel increment because the adjoint
is unital.  The other terms use only the first line of
\eqref{v12:eq:bank-bounds}.  The correction increment is exactly
\[
 d\{(M\chi)_s C(t+d)-\chi_sC(t)\}
 +dp_s\chi_H(\mathscr E^*-I)C(t+d).
\]
The last term has norm at most $Cdp_s\epsilon\le Cd^2$, by the second
line of \eqref{v12:eq:bank-bounds}; the time change of $C$ costs $Cd^2$.
The remaining term is $-d(f_s-1)C(t)$.  Thus the entire corrected expected
increment equals $d(\dot V+\mathscr D V)+\mathcal R$, with
$\|\mathcal R\|\le C(d^2+d\eta)$.
Apply \cref{v12:prop:Lyapunov}, pair with the actual density matrix, and use
$\tr\rho V\le W_j$.  This proves \eqref{v12:eq:one-step}.  No equilibrium
of $s_j$ and no independence between the phase and conditional state was used.
\end{proof}

\section{Nonlinear classical concentration along the unchanged resolved trajectory}
\label{v12:sec:main}

\begin{theorem}[Finite-bank nonlinear closure on the actual renewal clock]
\label[theorem]{v12:thm:main}
Under the hypotheses of \cref{v12:lem:corrector}, let
\begin{equation}
 b=\E\tr(\rho_0V(0))+\kappa^2+\eta+d.
 \label{v12:eq:budget}
\end{equation}
There is $C_T$ depending only on the displayed finite-bank bounds, the
polynomial coefficients, $r,M,T$, but not on matrix dimension,
$\vartheta\in(0,1]$, or the full Hamiltonian norm, such that
\begin{align}
 \max_{0\le j\le\lfloor T/d\rfloor}\E\tr(\rho_jV(jd))&\le C_Tb,\\
 \PP\left\{\max_{0\le j\le\lfloor T/d\rfloor}
                  \tr(\rho_jV(jd))>u\right\}&\le C_Tb/u
                  \qquad(u>0).
 \label{v12:eq:maximal}
\end{align}
In particular, when $b\to0$, the actual decoded means
$z_{j,a}=\tr(\rho_jA_a)$ converge uniformly in physical time in probability
to $x(t)$, and their conditional total variance
\begin{equation}
 \Sigma_j=\sum_a\tr\bigl(\rho_j(A_a-z_{j,a}I)^2\bigr)
 \label{v12:eq:Sigma}
\end{equation}
converges to zero uniformly in probability.  For $0<b\le1/2$ one also has
\begin{equation}
 \E\max_j\tr(\rho_jV(jd))\le C_Tb(1+|\log b|).
 \label{v12:eq:max-mean}
\end{equation}
No factor equal to the number of ticks or observed outcomes occurs in
\eqref{v12:eq:maximal}.
\end{theorem}
\begin{proof}
Put $a=1+C_2d$, $\beta=\kappa^2+\eta+d$, and
$J=\lfloor T/d\rfloor$.  Iterating \eqref{v12:eq:one-step} gives the
first inequality since $a^J\le e^{C_2T}$ and
$\E W_0\le\E\tr\rho_0V(0)+2dK_0$.
For the maximal bound define the nonnegative process
\[
 S_j=a^{-j}W_j+C_2\beta\sum_{\ell=j}^{J-1}d\,a^{-(\ell+1)},
 \qquad 0\le j\le J.
\]
Equation \eqref{v12:eq:one-step} gives
$\E[S_{j+1}\mid\mathcal F_j]\le S_j$.
Stop this nonnegative supermartingale at the first index with $W_j>u$,
or at $J$ if no such index occurs.  On the crossing event its value is at
least $a^{-J}u$.  Successive conditional expectations up to this bounded
stopping index give
$\PP(\max_jW_j>u)\le a^J\E S_0/u\le C_Tb/u$.
Since $\tr\rho_jV(jd)\le W_j$, this proves the second inequality.

The exact identity \eqref{v12:eq:variance-split} gives
$|z_j-x(jd)|^2+\Sigma_j=\tr\rho_jV(jd)$.
For the piecewise-constant reader between ticks, the reference moves by at
most $d\sup|F(x)|$, so the same convergence holds in physical time without
altering any recorded state.  Finally $0\le\tr\rho_jV(jd)\le4rM^2$.
Integrate the bound $\min(1,C_Tb/u)$ over this bounded range to obtain
\eqref{v12:eq:max-mean}, enlarging the constant when necessary.
The logarithmic loss is retained; no sharper mean-square maximal rate is
asserted by this argument.
\end{proof}

\begin{corollary}[Nonlinear moments, source insertions, and the decoder equation]
\label[corollary]{v12:cor:moments}
For every fixed real polynomial $p$,
\begin{align}
 \max_j\E\left|\tr(\rho_j\operatorname{Sym}p(A))-p(x(jd))\right|^2
       &\le C_{p,T}b,\\
 \PP\left\{\max_j
   \left|\tr(\rho_j\operatorname{Sym}p(A))-p(x(jd))\right|>u\right\}
       &\le C_{p,T}b/u^2.
 \label{v12:eq:moments}
\end{align}
The same estimates hold for any fixed finite bank of these polynomials.
Moreover, \eqref{v12:eq:conditional-incidence} converges to zero in mean square
at fixed physical times and uniformly in probability.  Thus the nonlinear
operator-to-reader incidence is a \emph{conclusion} on this branch, not an
assumed factorization of conditional moments.
\end{corollary}
\begin{proof}
Apply \cref{v12:lem:poly} to each actual conditional state and use
\cref{v12:thm:main}.  A finite bank is handled by summing its constants.
Equation \eqref{v12:eq:conditional-incidence} is bounded by the same positive
error, plus $2\kappa^2$.  This gives the final claim.  No outcome-conditioned
state was replaced by its ensemble mean.
\end{proof}

The theorem applies to mixed preparations and to any actual finite-outcome
instrument satisfying the same observable estimates.  It does not state
that different unravellings have identical state laws.  It states that the
positive reader error controls each such specified law.  A nonconcentrated
initial macroscopic distribution requires a corresponding random initial
classical reference; replacing it by its mean is not licensed.

\section{A locked six-edge realization without a full-generator norm bound}
\label{v12:sec:locked}

The following supplies the finite bank rather than postulating it.  It is a
constructed diagnostic realization of the inherited locking equations, not
a claim about unprovided operator panels of the intended source.
Let $\widehat{\mathcal H}_N$ act on a space of dimension at least five, with
an orthonormal eigenbasis $e_0,e_1,\ldots$.  Set
\[
 C_\mu=|e_0\rangle\langle e_{\mu-1}|\quad(2\le\mu\le5),
 \qquad R=\sum_{\mu=2}^5 C_\mu^*C_\mu.
\]
Thus $R$ is a projection, $\|C_\mu\|=1$, and $R$ commutes with the Hamiltonian.
For $0<\epsilon<1$ define
\begin{equation}
 U_\epsilon=e^{-\frac32\ii\epsilon\widehat{\mathcal H}_N},\quad
 S_\epsilon=(I-\epsilon^4R)^{1/2},\quad
 B_1=U_\epsilon S_\epsilon,\quad B_\mu=\epsilon^2 C_\mu\ (\mu\ge2).
 \label{v12:eq:B-family}
\end{equation}
Choose a real $6\times5$ isometry $O$ onto $\one^\perp$ whose first column
is $(1,1,1,-1,-1,-1)^{\mathsf T}/\sqrt6$, and put
\begin{equation}
 K_e=I/\sqrt{18}+\sqrt{2/3}\sum_{\mu=1}^5O_{e\mu}B_\mu.
 \label{v12:eq:K-family}
\end{equation}
The six original labels are retained.  This choice of auxiliary directions
is useful for proving rank uniformly; it is not a claim of spatial locality
of these additional matrices.

\begin{theorem}[Exact locking and a uniform finite-observable tangent]
\label[theorem]{v12:thm:locked-bank}
The family \eqref{v12:eq:K-family} satisfies
\begin{equation}
 \sum_eK_e^*K_e=I,\qquad \sum_eK_e=\sqrt2 I.
 \label{v12:eq:locking}
\end{equation}
All six $K_e$ are linearly independent for every $0<\epsilon<1$.
For sufficiently small $\epsilon$, independently of $N$ and the spectrum
of $\widehat{\mathcal H}_N$, each $K_e$ has a common positive least
singular-value bound.  In particular all six outcomes have positive
probability on every nonzero density matrix.
For any operator $O_0$,
\begin{align}
 \|\mathscr E_\epsilon^*O_0-O_0-\epsilon\mathscr D_NO_0\|
 &\le\tfrac34\epsilon^2\|\mathscr D_N^2O_0\|
                         +C\epsilon^4\|O_0\|,\\
 \|\mathscr E_\epsilon^*(\mathscr D_NO_0)-\mathscr D_NO_0\|
 &\le\epsilon\|\mathscr D_N^2O_0\|
                         +C\epsilon^4\|\mathscr D_NO_0\|.
 \label{v12:eq:local-tangent}
\end{align}
Hence uniform bounds on $O_0,\mathscr D_NO_0,\mathscr D_N^2O_0$ for
$O_0\in\mathcal O$ prove \eqref{v12:eq:bank-bounds} with
$\eta\le C\epsilon$, without a restriction $\epsilon\|\widehat{\mathcal H}_N\|\to0$.
\end{theorem}
\begin{proof}
Here $\sum_\mu B_\mu^*B_\mu=I$.  The orthogonality of $O$ and
$O^{\mathsf T}\one=0$ give \eqref{v12:eq:locking} by direct multiplication,
which is the inherited normal-form calculation.
The operators $I,B_1$ are diagonal in the chosen eigenbasis, while the four
$C_\mu$ occupy distinct off-diagonal matrix units.  On the $R$ and
$R^\perp$ diagonal entries the moduli of $B_1$ are respectively
$\sqrt{1-\epsilon^4}$ and one.  Thus $B_1$ is not scalar; the six operators
$I,B_1,C_2,\ldots,C_5$ are independent.  The fixed invertible column
transformation to the $K_e$ preserves independence.

The first-column choice gives
$K_e=I/\sqrt{18}\pm B_1/3+E_e$, $\|E_e\|\le C\epsilon^2$.
For every unit vector $u$,
\[
 \|(I/\sqrt{18}\pm B_1/3)u\|
 \ge \sqrt{1-\epsilon^4}/3-1/\sqrt{18}.
\]
The limit on the right is positive, proving the uniform singular-value
claim after subtracting $\|E_e\|$.

The exact adjoint channel is
\[
 \mathscr E_\epsilon^*O_0=\tfrac13O_0+	frac23\left[
 S_\epsilon U_\epsilon^*O_0U_\epsilon S_\epsilon
              +\epsilon^4\sum_{\mu=2}^5C_\mu^*O_0C_\mu\right].
\]
Since $\|S_\epsilon-I\|\le C\epsilon^4$, it differs in operator norm by
at most $C\epsilon^4\|O_0\|$ from
$\frac13O_0+\frac23 e^{(3\epsilon/2)\mathscr D_N}O_0$.
Unitary conjugation is an operator-norm isometry.  Twice integrating its
exact derivative gives
\[
 \left\|e^{t\mathscr D_N}O_0-O_0-t\mathscr D_NO_0\right\|
 \le t^2\|\mathscr D_N^2O_0\|/2.
\]
Use $t=3\epsilon/2$ to obtain the first estimate in
\eqref{v12:eq:local-tangent}.  Once integrating the conjugation of
$\mathscr D_NO_0$ gives the second.  No series expansion in
$\epsilon\|\widehat{\mathcal H}_N\|$ is taken.
\end{proof}

The coefficients are a specialization of Version 11's family, not a new
six-outcome normal form.  The additional results are its uniform rank and
positivity for an unbounded Hamiltonian, and the estimate on a fixed
observable bank rather than the entire algebra.  A tensor product with the
four inherited protected record sectors retains twenty-four accepted
operators.  Work in one prepared sector, or retain the random-sector initial
layer of \cref{v11:prop:sector,v11:cor:random-sector}.

\section{A nonlinear classical Hamiltonian realized on every resolved history}
\label{v12:sec:spin}

For $N\ge4$ let $J_1,J_2,J_3$ be the spin-$N/2$ irreducible matrices on
$\C^{N+1}$, with $[J_a,J_b]=\ii\epsilon_{abc}J_c$.  The declared reader is
$A_a=J_a/N$, so $\|A_a\|=1/2$ and
$[A_a,A_b]=\ii\epsilon_{abc}A_c/N$.
Fix independently $\Omega\ne0$ and $g\ne0$, and take the same-action pair
\begin{equation}
 \mathcal H(x)=\Omega x_1+\tfrac g2 x_3^2,\qquad
 \widehat{\mathcal H}_N=\Omega J_1+\frac{g}{2N}J_3^2
                       =N\operatorname{Sym}\mathcal H(A).
 \label{v12:eq:spin-H}
\end{equation}
This is a nonlinear classical spin Hamiltonian, not a Higgs potential or the
Einstein--SM Hamiltonian.

\begin{theorem}[A simultaneous nonlinear and renewal-clock limit]
\label[theorem]{v12:thm:spin}
Use the actual six-outcome instrument of \cref{v12:thm:locked-bank} with
\eqref{v12:eq:spin-H} and the inherited clock, for any sequence
$\epsilon_N\to0$, $0<\vartheta_N\le1$,
$d_N=4\vartheta_N\epsilon_N/11$.
Prepare the spin-coherent pure state directed along a prescribed unit vector
$n_0$, and set $x(0)=n_0/2$.  Then the target trajectory solves
\begin{equation}
 \dot x_1=-g x_2x_3,\qquad
 \dot x_2=g x_1x_3-\Omega x_3,\qquad
 \dot x_3=\Omega x_2.
 \label{v12:eq:spin-ODE}
\end{equation}
It exists globally on $|x|=1/2$.  For every finite $T$ the actual resolved
conditional means and variances obey \eqref{v12:eq:maximal} with
\begin{equation}
 b_N\le C(N^{-1}+\epsilon_N).
 \label{v12:eq:spin-budget}
\end{equation}
Every fixed polynomial moment and every polynomial parameter derivative of
the intensive same action have the bounds \eqref{v12:eq:moments} at this rate.
No condition $N\epsilon_N\to0$ is imposed.
\end{theorem}
\begin{proof}
Equation \eqref{v12:eq:exact-Lie} gives \eqref{v12:eq:spin-ODE} exactly at
the operator-coordinate level, with symmetric products where needed and
$\kappa=0$.  The scalar identity $x\cdot F(x)=0$ confines the target to
the compact sphere and gives global existence by ordinary finite-dimensional
ODE continuation.  All its fixed derivatives and the first two action
commutators of $A_a,A_a^2$ are uniformly bounded by
\cref{v12:prop:Lie}.  \Cref{v12:thm:locked-bank} therefore gives
$\eta=O(\epsilon_N)$ independently of $N$.

The spin-coherent state is a rotation of the highest-weight vector.  Its
mean is $x(0)=n_0/2$, and the Casimir identity gives
\[
 \tr(\rho_0V(0))
 =\frac{(N/2)(N/2+1)}{N^2}-\frac14=\frac1{2N}.
\]
Insert this, $\kappa=0$, and $d_N\le4\epsilon_N/11$ in
\cref{v12:thm:main,v12:cor:moments}.  The parameter derivatives are covered
because the operator coordinates do not depend on $\Omega,g$:
$N^{-1}\partial_\Omega\widehat{\mathcal H}_N=A_1$ and
$N^{-1}\partial_g\widehat{\mathcal H}_N=A_3^2/2$.
\end{proof}

\begin{proposition}[The full-algebra tangent can diverge while the classical reader converges]
\label[proposition]{v12:prop:unbounded}
For the same family,
\begin{equation}
 \left\|\epsilon_N^{-1}(\mathscr E_{\epsilon_N}^*-I)-\mathscr D_N
                                      \right\|_{\rm op\to op}
 \ge |\Omega|N-2/\epsilon_N.
 \label{v12:eq:whole-failure}
\end{equation}
In particular, choosing $\epsilon_N=N^{-1/2}$ makes this whole-algebra
tangent defect diverge while \cref{v12:thm:spin} gives a nonlinear classical
limit of the unchanged resolved reader.
\end{proposition}
\begin{proof}
The two spin-coherent states along $\pm e_1$ have the same expectation of
$J_3^2$ and expectations $\pm N/2$ of $J_1$.  Their Hamiltonian
expectations differ by $|\Omega|N$, so the spectral diameter of
$\widehat{\mathcal H}_N$ is at least that number.  A matrix unit between
extreme eigenvectors has norm one and commutator norm equal to the spectral
diameter, hence $\|\mathscr D_N\|_{\rm op\to op}\ge|\Omega|N$.
A unital completely positive map is operator-norm contractive: for example,
its Schwarz inequality bounds $\|\mathscr E^*O\|^2$ by $\|O\|^2$.
Thus $\|\mathscr E^*-I\|\le2$.  The reverse triangle inequality proves
\eqref{v12:eq:whole-failure}.  This does not contradict the full-algebra
Kraus-Gram statement in \cref{v11:thm:Gram-collapse}; its small whole-algebra
premise fails in this regime.
\end{proof}

\subsection{Why preparation variance is a real source condition}
\begin{proposition}[Equal means can retain opposite nonlinear drifts]
\label[proposition]{v12:prop:means}
For $N\ge3$ in the $N$-spin tensor realization, there are pure symmetric
states with identical means $\tr\rho A_a=0$ but opposite nonzero values of
$\tr(\rho\mathscr D_N A_1)$ for \eqref{v12:eq:spin-H}.
Specifically those values are $\mp g(N-1)/(8N)$.
\end{proposition}
\begin{proof}
Let $v_\pm=(0,1,\pm1)/\sqrt2$, and let $|v\rangle$ denote a spin-half
pure state of Bloch direction $v$.  Define
\[
 |\Psi_\pm\rangle=2^{-1/2}
      \bigl(|v_\pm\rangle^{\otimes N}+|-v_\pm\rangle^{\otimes N}\bigr).
\]
The two local states are orthogonal.  For an observable supported on at
most two sites the cross terms vanish when $N\ge3$.  Each mean is therefore
zero, and direct expansion of the collective product gives
\[
 \tr\bigl(\rho_\pm\operatorname{Sym}(A_2A_3)\bigr)
          =\pm\frac{N-1}{8N}.
\]
The on-site symmetrized Pauli product is zero; the $N(N-1)$ distinct-site
terms give the displayed value.  Since
$\mathscr D_N A_1=-g\operatorname{Sym}(A_2A_3)$, the claim follows.
The states lie in the symmetric spin-$N/2$ sector.  Their nonvanishing
macroscopic variance is not covered by the concentrated initial entrance.
\end{proof}

\section{Same-action sources and the remaining classical field identity}
\label{v12:sec:sources}

\begin{corollary}[Polynomial source derivatives pass on the same resolved record]
\label[corollary]{v12:cor:sources}
Suppose an independently specified finite family of external source
parameters $\lambda$ enters
$\widehat{\mathcal H}_\nu(\lambda)
=\nu^{-1}\operatorname{Sym}\mathcal H(A^{(\nu)};\lambda)$,
with the observable basis independent of $\lambda$.  Assume the preceding
bounds hold at the evaluated parameter.  Then the actual intensive source
insertions
\begin{equation}
 S_{\alpha,j}=\nu\tr\bigl(\rho_j\partial_{\lambda_\alpha}
                      \widehat{\mathcal H}_\nu\bigr)
 \label{v12:eq:source}
\end{equation}
converge at fixed times in mean square, and uniformly on compact times in
probability, to $\partial_{\lambda_\alpha}\mathcal H(x(t);\lambda)$ with
the polynomial bounds \eqref{v12:eq:moments}.  The same holds for any fixed
quadratic or higher polynomial insertion actually represented by this bank.
\end{corollary}
\begin{proof}
Differentiation of this specified action at fixed $A$ gives
$\nu\partial_{\lambda_\alpha}\widehat{\mathcal H}_\nu
=\operatorname{Sym}\partial_{\lambda_\alpha}\mathcal H(A;\lambda)$.
Apply \cref{v12:cor:moments}.  This proves convergence of the independently
varied action insertion, not a derivative of an inferred trajectory-dependent
limit value.  If $A$ itself depends on $\lambda$, its additional chain-rule
terms must be included and controlled separately.
\end{proof}

\paragraph{Poisson versus canonical variables.}
The spin realization is a classical Lie--Poisson system.  On a coordinate
patch of its sphere avoiding a pole and an azimuthal cut, set
$q=\arg(x_1+\ii x_2)$ and $p=x_3$.  Direct differentiation of
$\{x_a,x_b\}=\epsilon_{abc}x_c$ gives $\{q,p\}=1$.
On the leaf $|x|=s=1/2$ the \emph{same} Hamiltonian is
\begin{equation}
 H_{\rm can}(q,p)=\Omega\sqrt{s^2-p^2}\cos q+\tfrac g2p^2.
 \label{v12:eq:canonical-spin}
\end{equation}
Its equations $\dot q=\partial_pH_{\rm can}$,
$\dot p=-\partial_qH_{\rm can}$ are exactly \eqref{v12:eq:spin-ODE}
in that chart.  This calculation demonstrates a genuinely nonlinear
canonical local limit, not just a linear rotation of matrix coordinates.
For any fixed smooth tests on a subslab where the reference remains within
that chart, uniform convergence of the actual decoded variables passes the
weak first-order equation.  No strong derivative of the jump path is asserted.

A classical Einstein--SM identification still requires its \emph{own}
source-to-Poisson or source-to-canonical map, physical normalization, and
same-action parameter insertions.  A spin orbit does not provide gravitational
metric variables, spacetime locality, the Higgs radial mode, the chiral matter
content, or the multiplier constraints.  The parameter $N$ here is a declared
large-spin normalization, not the number of Standard Model generations.
The theorem fixes the finite observable bank and time interval; it has no
unproved spatial ultraviolet or increasing-bank uniformity.

\section{Outcome, precision budget, and the next native test}
\label{v12:sec:ledger}

The positive result is now nonlinear at the physical-reader level:
\[
 \begin{gathered}
 \text{same-action finite observable bank + concentrated preparation}\\[0.3em]
 \Downarrow\\[0.3em]
 \text{nonlinear resolved trajectories + same-action source convergence}.
 \end{gathered}
\]
The known renewal phase is handled by a proof-only correction of a positive
error observable.  There is no repeated-observation union-bound cost, no
purity hypothesis in the general theorem, and no required uniform bound on
the norm of the full quantum Hamiltonian.  These are the distinctions from
an ensemble-only moment theorem and from Version 11's whole-state limit.
The explicit locked realization shows that the finite-bank premise is
nonempty even where a whole-algebra tangent estimate fails badly.

\paragraph{What must actually be acquired.}
For a proposed native classical reader, the next source data are its
normalized coordinate operators, their independently represented action
commutators, and the original accepted-channel pullbacks on
$A_a,A_a^2,\mathscr D A_a,\mathscr D(A_a^2)$.  The preparation must provide
the variance in \eqref{v12:eq:budget}.  If a pullback on $A_a$ or $A_a^2$
is enclosed with operator-norm error $e_{\rm obs}$, its additional tangent
error is $e_{\rm obs}/\epsilon$.  Thus $e_{\rm obs}=o(\epsilon)$ suffices
for this bank's vanishing defect, and $e_{\rm obs}=O(\epsilon^2)$ retains
an $O(\epsilon)$ rate; the action-commutator uncertainty is a separate term.
This statement concerns direct observable pullbacks.  Acquiring them from
the rare-acceptance complementary panels must additionally pay the
$s_\vartheta^{-1}$ reconstruction factor in \cref{v11:prop:precision}.

Neither the chosen semiclassical scaling nor the fixed polynomial action is
selected by the bare locking identities.  The constructed branch can realize
other Hamiltonians as well.  For the intended Einstein--SM source, the
field/action incidence, metric and spatial derivative topology, gauge and
constraint rows, legal-domain and cofinal controls are not discharged by
this finite-dimensional result.  The advance is that a nontrivial nonlinear
reader incidence and its polynomial source factorization can now be
\emph{proved as outputs} on a resolved native-clock branch, rather than
simply appended as extra assumptions to a von Neumann limit.

\paragraph{Verification and preservation.}
The complete Version 11 text preceding its final end-document command is
preserved byte-for-byte.  All new labels have prefix \texttt{v12:}.
The accompanying diagnostic script checks the coordinate commutators,
polynomial quadratic-form bounds, positive corrected one-step estimate,
exact locking, rank and outcome floors, the finite-observable remainder,
nonlinear conditional-moment witness, source derivatives, and the coexistence
of a good finite bank with a bad whole-algebra tangent.  A finite resolved
trajectory experiment is diagnostic only.  The proofs are the operator
identities and conditional estimates given above, not numerical evidence or
a machine-checked formalization of the Einstein--SM system.

\begingroup
\renewcommand{\refname}{Additional references for Version 12}
\begin{thebibliography}{9}
\small\raggedright
\bibitem{v12:FKL}
J.~Fr\"ohlich, A.~Knowles, and E.~Lenzmann,
\emph{Semi-classical dynamics in quantum spin systems},
Letters in Mathematical Physics \textbf{82} (2007), 275--296.
\href{https://arxiv.org/abs/0709.4561}{arXiv:0709.4561};
\href{https://doi.org/10.1007/s11005-007-0202-y}{doi:10.1007/s11005-007-0202-y}.
Background for large-spin and mean-field Hamiltonian limits.  No result from
this reference is used to replace the finite resolved-clock proof above.
\end{thebibliography}
\endgroup
% END IMMUTABLE VERSION 12 BODY
% APPEND VERSION 13 HERE, BEFORE THE FINAL END-DOCUMENT COMMAND.


% BEGIN IMMUTABLE VERSION 13 BODY
\clearpage
\begin{center}
{\Large\bfseries Version 13: A spatial scalar action before an increasing-bank bound}\\[0.5em]
{\large Exact relative-energy cancellation, finite spectral records, and a joint classical field limit}
\end{center}
\addcontentsline{toc}{part}{Version 13: Spatial real-Higgs action, finite records, and scalar stress}

\section{Audit: from a fixed spin reader to a spatial action topology}
\label{v13:sec:audit}

Version 12 fixes the number of reader coordinates.  Its polynomial constants
are not uniform merely because those coordinates are relabelled by the sites
of a refining spatial lattice.  In particular, an ordinary Lipschitz bound
for a lattice wave vector field charges the norm of its discrete Laplacian,
which is of order $h^{-2}$.  An exponential of that bound would not be a useful
continuum estimate.  The present calculation instead uses the kinetic and
spatial-gradient terms of the \emph{same} scalar action to cancel this cost.

\paragraph{Nonduplication and inherited inputs.}
The fixed-mode many-copy theorem in \cite{SM}, at
\nolinkurl{thm:many-copy-transport}, remains an input with fixed spatial modes;
we do not reprove it.  The canonical real-Higgs normalization is
\cref{v7:eq:Higgs-potential}, and the distinction between canonical momentum
and a slope of recorded field values remains
\cref{v8:thm:incidence}.  The locked six-outcome normalization, its
observable-specific estimate, and the physical phase correction are
\cref{v12:thm:locked-bank,v10:eq:chi}.  We reuse those exact identities.
The new results below establish a spatially uniform action estimate, its
finite spectral realization and preparation, and an explicit joint limit.
Search and exact-term checks of the supplied monographs and the cumulative
classical-closure file did not locate this operator relative-energy argument.
This is a targeted nonduplication audit, not an assertion that every line of
all supplied sources has been independently verified.

\paragraph{Scope of the positive construction.}
The action treated here is the canonical real-Higgs restriction, with fixed
flat coframe, identity internal transport, and zero spinors.  Spatially varying
scalar fields are allowed.  Its quantization, semiclassical parameter, energy
cutoff, and accepted instrument are explicitly declared.  They provide a
finite action-compatible branch; they are \emph{not} asserted to be the
unacquired operator panels of the locked Einstein--SM source.  In particular,
a positive scalar stress on a flat coframe does not solve the Einstein
metric equation.  Neither gravitational backreaction nor locality of the
accepted \emph{instrument} is established by locality of its Hamiltonian.

\paragraph{A typographical clarification in the immutable prefix.}
In the proof of \cref{v12:thm:locked-bank}, the adjoint-channel formula has a
missing backslash in the printed coefficient following $O_0/3$.  The identity
used there and below is
\[
 \mathscr E_\epsilon^*O_0=\frac13O_0+\frac23
 \left[S_\epsilon U_\epsilon^*O_0U_\epsilon S_\epsilon
              +\epsilon^4\sum_{\mu=2}^5 C_\mu^*O_0C_\mu\right].
\]
No earlier source bytes are changed.

\paragraph{Attribution.}
Coherent-state classical limits, quantum comparison costs, and relative-energy
arguments have established precedents, including \cite{v13:DT,v13:GP}.
No priority claim for those methods is made.  The specific addition is the
spatial mass normalization, exact spectral compression without false finite
canonical commutation relations, and the resolved locked-clock estimate for
this quartic scalar field.  The proofs below give their own estimates and do
not import a continuum field limit from those references.

\section{The fixed spatial scalar action and its finite spectral realization}
\label{v13:sec:action}

Let $\Lambda_h=(\Z/N\Z)^3$, $h=N^{-1}$, $N\ge3$, and set $w=h^3$.
There are $n=N^3$ sites and $wn=1$.  Ordinary forward differences are
$\delta_i u_x=(u_{x+e_i}-u_x)/h$, and
$\Delta_h=-\sum_i\delta_i^*\delta_i$ in the mass inner product
$(u,v)_h=w\sum_xu_xv_x$.
Use the actual quartic normalization of \cref{v7:eq:Higgs-potential}:
\begin{equation}
 V(u)=\frac\lambda4(u^2-b)^2,
 \qquad \lambda=\lambda_H>0,\quad b=2v_H^2\ge0.
 \label{v13:eq:potential}
\end{equation}
Thus the doublet restriction is $H=u e_1/\sqrt2$; $b$ is a fixed potential
coefficient and is not an error budget.  The spatially discrete classical
action and energy are
\begin{align}
 S_h(q)&=\int w\sum_x\left\{\frac12\dot q_x^2
              -\frac12\sum_i(\delta_iq_x)^2-V(q_x)\right\}\dd t,
       \label{v13:eq:classical-action}\\
 E_h(q,p)&=w\sum_x\left\{\frac12p_x^2
              +\frac12\sum_i(\delta_iq_x)^2+V(q_x)\right\}.
       \label{v13:eq:classical-energy}
\end{align}
The canonical momentum density is $p_x$, with symplectic form
$w\sum_x\dd q_x\wedge\dd p_x$.  Consequently the classical equations are
$\dot q=p$, $\dot p=\Delta_hq-V'(q)$.

For $\nu>0$, on $L^2(\R^n,\dd u)$ put
\begin{equation}
 Q_x=u_x,\qquad P_x=-\ii\nu\partial_{u_x},\qquad
 [Q_x,P_y]=\ii\nu\delta_{xy}I,
 \qquad \hbar_{h,\nu}=\nu w.
 \label{v13:eq:CCR}
\end{equation}
The physical Planck parameter for the mass-normalized action is
$\hbar_{h,\nu}$, not $\nu$ alone.  Define the nonnegative Schr\"odinger energy
and its action generator by
\begin{equation}
 \boldsymbol E=w\sum_x\left\{\frac12P_x^2
           +\frac12\sum_i(\delta_iQ_x)^2+V(Q_x)\right\},\qquad
 \mathscr D O=\frac{\ii}{\nu w}[\boldsymbol E,O].
 \label{v13:eq:quantum-energy}
\end{equation}
All products in this equation are on the uncompressed canonical space.

\begin{lemma}[Confining spectral cutoff and exact compression of a derivation]
\label[lemma]{v13:lem:spectral}
For fixed $h,\nu$, the Friedrichs realization of $\boldsymbol E$ has compact
resolvent.  Its eigenvectors lie in the domain of every polynomial in the
$Q_x,P_x$.  For $\Lambda>0$ let
\[
 \Pi=\one_{[0,\Lambda]}(\boldsymbol E),\qquad
 \mathcal K=\Ran\Pi,\qquad
 \boldsymbol E_\Lambda=\boldsymbol E|_{\mathcal K},\qquad
 \widehat{\mathcal H}=\boldsymbol E_\Lambda/(\nu w).
\]
Then $\mathcal K$ is finite dimensional.  For any polynomial observable $B$
and $B_\Lambda=\Pi B\Pi|_{\mathcal K}$,
\begin{equation}
 \ii[\widehat{\mathcal H},B_\Lambda]
       =\Pi\mathscr D B\Pi|_{\mathcal K}.
 \label{v13:eq:compression}
\end{equation}
In general $\Pi Q_x^2\Pi\ne(\Pi Q_x\Pi)^2$.  No finite-dimensional exact
canonical commutation relation is claimed.
\end{lemma}
\begin{proof}
The closed energy form is the sum of a positive multiple of
$\|\nabla\psi\|_2^2$ and a nonnegative multiplication form.  Its potential
$W_h(u)$ tends to infinity as $|u|\to\infty$, since
$W_h(u)\ge c_h\sum_xu_x^4-C_h$ for fixed $h$.
A form-bounded family has a uniformly small $L^2$ tail outside large balls;
inside each ball its $H^1$ bound gives compactness in $L^2$.  Combining these
two facts proves compactness of the form-domain embedding and hence of the
resolvent.

For completeness, the polynomial-domain assertion can be seen from the
eigenvalue equation.  Test it with a cutoff approximation to
$\langle u\rangle^{2k}\psi$ and integrate by parts.  With
$a_h=w\nu^2/2$, the weighted identity is
\[
 a_h\|\nabla(\omega\psi)\|_2^2+
 \int W_h|\omega\psi|^2
 =E\|\omega\psi\|_2^2+a_h\int|\nabla\omega|^2|\psi|^2.
\]
Choose $\omega=\langle u\rangle^k$ with bounded smooth truncations whose
first derivatives are at most $C_k\langle u\rangle^{k-1}$.
The quartic lower bound, induction in $k$, and then monotone convergence give
all polynomial position moments and weighted first derivatives.  The
identity $-a_h\Delta\psi=(E-W_h)\psi$ gives weighted second derivatives:
apply the Fourier $H^2$ estimate to a cutoff times
$\langle u\rangle^k\psi$, whose commutator terms have just been bounded.
Differentiate the equation and repeat; derivatives of $W_h$ are polynomials,
so this proves every weighted derivative.  Thus all polynomial operations
used below are legitimate on the finite eigenspaces.  Finally $\Pi$ commutes
with $\boldsymbol E$.  Taking its matrix elements between eigenvectors proves
\eqref{v13:eq:compression} exactly.  Compressing an observable is not an
algebra homomorphism, which explains the final qualification.
\end{proof}

The finite source is defined on this spectral space from the outset.  The
projection is part of its declared regulator, not a projection applied after
observing an unfavorable trajectory.  In particular, every quadratic or
quartic insertion below means compression of the original insertion.

\section{An exact quartic relative-energy cancellation with no spatial stiffness loss}
\label{v13:sec:relative}

Let $q_x(t),p_x(t)$ be real comparison data on a fixed time interval, with
\begin{equation}
 \dot q=p,\qquad
 r=\dot p-\Delta_hq+V'(q).
 \label{v13:eq:reference}
\end{equation}
Their relevant values and smooth spatial/time differences are bounded by a
fixed constant.  In the continuum application below they are samples of a
smooth solution and $\|r\|_{2,h}=O(h^2)$.  The source itself will use only
$q(0),p(0)$, not these future values.
Put
\begin{equation}
 \alpha=1+\lambda b,\quad W(u)=V(u)+\frac\alpha2u^2,\quad
 \mathcal B(u\mid q)=W(u)-W(q)-W'(q)(u-q).
 \label{v13:eq:Bregman}
\end{equation}
The positive comparison observable is
\begin{equation}
 \mathcal V(t)=w\sum_x\left\{\frac12(P_x-p_x)^2
       +\frac12\sum_i\bigl(\delta_i(Q-q)_x\bigr)^2
       +\mathcal B(Q_x\mid q_x)\right\}.
 \label{v13:eq:relative-energy}
\end{equation}
The added convexification is a comparison device.  It is not added to the
physical action or to $\boldsymbol E$.

\begin{lemma}[Global coercivity of the comparison potential]
\label[lemma]{v13:lem:quartic}
For every real $q,e$,
\begin{equation}
 \mathcal B(q+e\mid q)
 =\frac12e^2+\frac\lambda{12}e^4
             +\frac{3\lambda}{2}e^2(q+e/3)^2
 \ge\frac12e^2+\frac\lambda{12}e^4.
 \label{v13:eq:quartic-coercivity}
\end{equation}
For $|q|\le M$ it is also bounded above by $C_M(e^2+e^4)$.
These are multiplication-operator inequalities for $e=Q_x-q_x$.
\end{lemma}
\begin{proof}
Expand the quartic and cancel its affine part.  The result before completing
the square is
$\lambda e^4/4+\lambda qe^3+(3\lambda q^2+1)e^2/2$.
Completing the square gives \eqref{v13:eq:quartic-coercivity}.
The upper bound follows from $|e|^3\le(e^2+e^4)/2$.
\end{proof}

\begin{theorem}[Exact operator cancellation for the spatial scalar action]
\label[theorem]{v13:thm:cancellation}
Set $e_x=Q_x-q_x$ and $v_x=P_x-p_x$.  On the common polynomial domain,
\begin{equation}
 \boxed{
 (\partial_t+\mathscr D)\mathcal V
 =w\sum_x\left\{
 \lambda p_x(3q_xe_x^2+e_x^3)
       +\frac\alpha2(e_xv_x+v_xe_x)-r_xv_x\right\}.}
 \label{v13:eq:exact-cancellation}
\end{equation}
Consequently, if $|q_x|+|p_x|\le M$, then as quadratic forms
\begin{equation}
 (\partial_t+\mathscr D)\mathcal V
       \preceq C_M\mathcal V+\frac12\|r\|_{2,h}^2 I.
 \label{v13:eq:energy-inequality}
\end{equation}
The constant depends on $M,\lambda,b$, not on $h,\nu$, the number of sites,
or the norm of the lattice Laplacian.  The same identity and inequality hold
after spectral compression with $\mathcal V_\Lambda=\Pi\mathcal V\Pi$ and
$\mathscr D_\Lambda=\ii[\widehat{\mathcal H},\,\cdot\,]$.
\end{theorem}
\begin{proof}
The exact canonical equations are
$\mathscr DQ_x=P_x$ and
$\mathscr DP_x=\Delta_hQ_x-V'(Q_x)$.
For a one-coordinate polynomial $f$, the canonical commutator gives
$\mathscr Df(Q_x)=\{f'(Q_x),P_x\}/2$, with $\{A,B\}=AB+BA$.
This identity follows by the product rule and $[Q_x,P_x]=\ii\nu I$;
it applies in particular to the quartic and to its affine subtraction.

The kinetic comparison contributes
\[
 \frac12w\sum_x\{v_x,\Delta_he_x-[V'(Q_x)-V'(q_x)]-r_x\}.
\]
The gradient comparison contributes
$w\sum_{x,i}\{\delta_ie_x,\delta_iv_x\}/2$.
Finite summation by parts cancels the Laplacian contribution \emph{exactly},
including both operator orderings.  No commutation between $e_x$ and $v_x$
is assumed in this cancellation.
The nonconvex relative potential contributes
$\{V'(Q_x)-V'(q_x),P_x\}/2-V''(q_x)p_xe_x$.
Combining it with the remaining kinetic term leaves
$p_x[V'(Q_x)-V'(q_x)-V''(q_x)e_x]
=\lambda p_x(3q_xe_x^2+e_x^3)$.
The convexifying square contributes $\alpha\{e_x,v_x\}/2$.
These terms prove \eqref{v13:eq:exact-cancellation}.

The quartic lemma bounds the first term in absolute value by
$C_M\mathcal B(Q_x\mid q_x)$.
For self-adjoint $A,B$, positivity of $(A\pm B)^2$ gives
$\pm(AB+BA)\preceq A^2+B^2$.
Use this for the second term and
$-r_xv_x\preceq v_x^2/2+r_x^2I/2$ for the last.
Sum with mass $w$.  This proves \eqref{v13:eq:energy-inequality}.
Compression is exact by \cref{v13:lem:spectral}; it does not replace
$\Pi Q_x^k\Pi$ by powers of a compressed coordinate.
\end{proof}

\begin{proposition}[Explicit norms needed by the renewal step]
\label[proposition]{v13:prop:bank-norms}
Suppose the comparison data arise from a fixed smooth function on the unit
torus and interval, with the time/spatial derivatives used below uniformly
bounded.  Put $A=1+\Lambda$ and $\omega=\nu h^3$.
For $0<\nu\le h\le1$ and $\Lambda\ge1$,
\begin{align}
 \|\partial_t^j\mathcal V_\Lambda(t)\|&\le C A\quad(j=0,1,2),
       \label{v13:eq:bank-low}\\
 \|\mathscr D_\Lambda\mathcal V_\Lambda(t)\|&\le C A,\nonumber\\
 \|\mathscr D_\Lambda^2\mathcal V_\Lambda(t)\|
 +\|\mathscr D_\Lambda\partial_t\mathcal V_\Lambda(t)\|
       &\le C A^2/\omega.
       \label{v13:eq:bank-high}
\end{align}
The constants contain no number-of-sites factor.  The last displayed bound
is conservative and pays its inverse physical quantum scale explicitly.
\end{proposition}
\begin{proof}
The inequalities $(A-B)^2\preceq2A^2+2B^2$ for self-adjoint $A,B$,
$u^2\le C(V(u)+1)$, and the bounded comparison values imply
$0\preceq\mathcal V\preceq C(\boldsymbol E+I)$ as forms.
Here $w\sum_x1=1$, and the reference gradient has a uniform bound.
A time derivative of $\mathcal V$ is a sum of linear $P,Q,\delta Q$
terms and scalar terms with bounded coefficients.  In the potential piece
$\partial_t\mathcal B(Q\mid q)=-W''(q)p(Q-q)$.
A second time derivative has the same kind of terms.  Weighted
Cauchy--Schwarz and $2|ab|\le a^2+b^2$ give the two-sided form bounds
$\pm\partial_t^j\mathcal V\preceq C(\boldsymbol E+I)$, $j=1,2$.
The identity \eqref{v13:eq:exact-cancellation} has the same two-sided form
bound, by its proof, so $\pm\mathscr D\mathcal V\preceq
C(\boldsymbol E+I)$ as well.  Compressing gives \eqref{v13:eq:bank-low}.
Finally $\|\widehat{\mathcal H}\|\le\Lambda/\omega$ and
$\|[H,B]\|\le2\|H\|\|B\|$ prove \eqref{v13:eq:bank-high}.
The large norm is used only in this polynomial error budget, not in an
exponential stability constant.
\end{proof}

\section{A genuinely finite initial preparation with a quantitative cutoff cost}
\label{v13:sec:preparation}

Take prescribed smooth initial values $q_x^0,p_x^0$ and the normalized product
Gaussian on the uncompressed canonical space
\begin{equation}
 \phi_{h,\nu}(u)=(\pi\nu h)^{-n/4}
 \exp\left[-\frac1{2\nu h}\sum_x(u_x-q_x^0)^2
               +\frac\ii\nu\sum_xp_x^0(u_x-q_x^0)\right].
 \label{v13:eq:Gaussian}
\end{equation}
This is a declared squeezed preparation, not a consequence of the locking
identities.  Its position variance at a site is $\nu h/2$ and its momentum
variance is $\nu/(2h)$.  For sufficiently large $\Lambda$ define the
\emph{finite} initial vector
\begin{equation}
 \psi_0=\Pi\phi_{h,\nu}/\|\Pi\phi_{h,\nu}\|,
 \qquad \rho_0=|\psi_0\rangle\langle\psi_0|.
 \label{v13:eq:projected-preparation}
\end{equation}
Only initial data are used by this construction.

\begin{theorem}[Preparation and compression cost in the physical energy norm]
\label[theorem]{v13:thm:preparation}
For $0<\nu\le h$ and uniformly bounded initial values and gradients,
\begin{align}
 \langle\phi_{h,\nu},\mathcal V(0)\phi_{h,\nu}\rangle
 &=\frac{7\nu}{4h}
  +\frac{\nu h}{4}\,w\sum_x(3\lambda(q_x^0)^2+1)
  +\frac{3\lambda\nu^2h^2}{16},
       \label{v13:eq:exact-prep}\\
 \langle\phi_{h,\nu},(\boldsymbol E+I)^2\phi_{h,\nu}\rangle&\le C,
       \label{v13:eq:E2-prep}\\
 \tr(\rho_0\mathcal V_\Lambda(0))
       &\le C\left(\frac\nu h+\frac1{1+\Lambda}\right).
       \label{v13:eq:finite-prep}
\end{align}
The last assertion holds for $\Lambda\ge\Lambda_0$, with $\Lambda_0,C$
independent of $h,\nu$.  The finite spectral space has dimension at least
five for all sufficiently small $h$ and all sufficiently large fixed
$\Lambda_0$.
\end{theorem}
\begin{proof}
Write $\xi_x=u_x-q_x^0$, $s^2=\nu h/2$.
Exactly $(P_x-p_x^0)\phi=\ii\xi_x\phi/h$, so its squared norm is
$\nu/(2h)$.  Independent adjacent centered positions give
$\E(\delta_i\xi_x)^2=\nu/h$.  Multiplication by $1/2$, summation over three
directions, and $wn=1$ give the first contribution $7\nu/(4h)$.
The quartic expansion in \cref{v13:lem:quartic}, together with
$\E\xi^3=0$, $\E\xi^4=3s^4$, gives the other terms of
\eqref{v13:eq:exact-prep}.

To verify the cutoff cost rather than assume an energy moment, calculate
\[
 P_x^2\phi=
 \left[(p_x^0)^2+2\ii p_x^0\xi_x/h-\xi_x^2/h^2+\nu/h\right]\phi.
\]
Its $L^2$ norm is at most $C(1+\nu/h)$ for bounded $p_x^0$.
Each gradient-square multiplication has $L^2$ norm at most
$C(1+\nu/h)$, by its fourth Gaussian moment and bounded reference gradient.
Each $V(Q_x)\phi$ has bounded norm, since its eighth Gaussian moment is
bounded when $\nu h\le1$.  The triangle inequality after multiplication by
$w$, summed over $n$ sites, proves
$\|\boldsymbol E\phi\|\le C(1+\nu/h)$ and hence
\eqref{v13:eq:E2-prep}.

Set $u=(I-\Pi)\phi$.  Spectral calculus gives
\[
 \|u\|^2\le C/(1+\Lambda)^2,\qquad
 \langle u,(\boldsymbol E+I)u\rangle\le C/(1+\Lambda).
\]
The form domination in the proof of \cref{v13:prop:bank-norms} bounds
$\|\mathcal V(0)^{1/2}u\|^2$ by the latter quantity.  Apply the triangle
inequality to $\mathcal V(0)^{1/2}(\phi-u)$, divide by
$\|\Pi\phi\|^2\ge1/2$, and use \eqref{v13:eq:exact-prep}.
This proves \eqref{v13:eq:finite-prep}.

Finally take five such Gaussians with distinct fixed constant shifts of the
initial position center and the same initial momentum.  Their Gram matrix
tends to the identity as $h\to0$: the absolute overlap of two of them is
$\exp[-n c^2/(4\nu h)]$ for their nonzero center separation $c$.
Each has uniformly bounded energy expectation.  The energy form matrix of
these five vectors has bounded norm by form Cauchy--Schwarz, while their
Gram has a uniform lower bound.  The variational min--max principle bounds
the fifth eigenvalue by a constant.  Increasing $\Lambda_0$ proves the
rank assertion without any future continuum data.
\end{proof}

The local uncertainty scale $\nu$ and the physical scale $\nu h^3$ are not
inferred from the number of physical species.  In this explicit preparation,
$\nu/h\to0$ is a sufficient simultaneous classical/spatial scale; it is
not asserted to be necessary for all correlated preparations or renormalized
quantum limits.

\section{Resolved scalar trajectories on the unchanged two-phase clock}
\label{v13:sec:clock}

On $\mathcal K$, use the exact six-outcome construction of
\cref{v12:thm:locked-bank}, now with
$\widehat{\mathcal H}=\boldsymbol E_\Lambda/(\nu h^3)$.
Its eigenbasis defines $C_\mu=|e_0\rangle\langle e_{\mu-1}|$,
$R=\sum_{\mu=2}^5C_\mu^*C_\mu$, and
\[
 B_1=e^{-3\ii\epsilon\widehat{\mathcal H}/2}
                       (I-\epsilon^4R)^{1/2},\qquad
 B_\mu=\epsilon^2C_\mu\quad(2\le\mu\le5).
\]
The original six Kraus labels are given by the fixed normal-form matrix in
\cref{v12:eq:K-family}.  Their completeness, linear independence, and positive
outcome probabilities are inherited exact statements, not new lemmas.
One may tensor the four protected record sectors and prepare one such sector;
the initial-layer qualification from Version 11 is unchanged.

Use the \emph{same} protected-tick phase matrix and acceptance law as before:
\[
 M=\begin{pmatrix}1/5&4/5\\2/3&1/3\end{pmatrix},\qquad
 d=\frac{4\vartheta}{11}\epsilon,\quad0<\vartheta\le1,\qquad
 \chi_H=-15/22,\quad\chi_P=25/44.
\]
The complete resolved conditional density matrix at tick $j$ is $\rho_j$.
All its actual outcomes, not an averaged replacement path, enter the theorem.
The constructed generator is spatially local in the scalar variables;
the spectral cutoff and the accepted operators above are global.  Spatial
locality of those operators is \emph{not} a conclusion.

\begin{theorem}[Spatially uniform resolved relative-energy bound]
\label[theorem]{v13:thm:main}
Let $q(t,x)$ be a fixed smooth real solution on $[0,T+1]\times\T^3$ of
\begin{equation}
 q_{tt}-\Delta q+\lambda q(q^2-b)=0.
 \label{v13:eq:wave}
\end{equation}
Set $q_x(t)=q(t,hx)$, $p_x(t)=q_t(t,hx)$ in the comparison observable.
Assume the finitely many derivatives used in
\cref{v13:prop:bank-norms} and the fourth spatial derivatives are bounded;
for example $q\in C^4([0,T+1];C^6(\T^3))$ suffices.
For any finite initial density matrix on $\mathcal K$, put
\begin{equation}
 \mathfrak b=\tr(\rho_0\mathcal V_\Lambda(0))
      +h^4+\frac{\epsilon(1+\Lambda)^2}{\nu h^3}.
 \label{v13:eq:budget}
\end{equation}
For $0<\nu\le h$, $\Lambda\ge\Lambda_0$, $0<\epsilon\le\epsilon_0$,
and arbitrary initial phase,
\begin{align}
 \max_{j\le T/d}\E\tr(\rho_j\mathcal V_\Lambda(jd))&\le C_T\mathfrak b,
      \label{v13:eq:mean-bound}\\
 \PP\left\{\max_{j\le T/d}\tr(\rho_j\mathcal V_\Lambda(jd))>u\right\}
       &\le C_T\mathfrak b/u\quad(u>0).
      \label{v13:eq:max-bound}
\end{align}
The constants are independent of $h,\nu,\Lambda,\epsilon,\vartheta$, and
spectral-space dimension.  In particular they contain neither
$\exp(CT/h^2)$ nor $\exp(CT\|\widehat{\mathcal H}\|)$.
For $0<\mathfrak b\le1/2$,
\begin{equation}
 \E\max_{j\le T/d}\tr(\rho_j\mathcal V_\Lambda(jd))
 \le C_T\mathfrak b\left[1+\log\left(2+
                        \frac{1+\Lambda}{\mathfrak b}\right)\right].
 \label{v13:eq:log-max}
\end{equation}
\end{theorem}
\begin{proof}
Sampling \eqref{v13:eq:wave} gives
$r_x=(\Delta q)(t,hx)-\Delta_hq_x$, hence
$\|r\|_{2,h}\le Ch^2$ by the integral remainder for centered differences.
The operator relative-energy inequality therefore has error $Ch^4I$.

Here it is important not to apply Version 12 with a number of coordinates
that is then silently allowed to grow.  Work directly with the single
mass-weighted observable $V=\mathcal V_\Lambda$ and $C=\mathscr D_\Lambda V$.
Write $A=1+\Lambda$, $\omega=\nu h^3$, and $B=A^2/\omega$.
\Cref{v13:prop:bank-norms} gives
\[
 \|V\|+\|\dot V\|+\|\ddot V\|+\|C\|\le C_0A,
 \qquad \|\dot C\|+\|\mathscr D_\Lambda C\|\le C_0B.
\]
The inherited observable channel bound, applied at each fixed $t$, gives
\[
 \|\mathscr E^*V-V-\epsilon C\|
       \le C(\epsilon^2 B+\epsilon^4 A),\quad
 \|\mathscr E^*C-C\|\le C(\epsilon B+\epsilon^4 A).
\]
Define $\widehat V_s=V+d\chi_s C+dK_0I$ with
$K_0=\max_s|\chi_s|\sup_t\|C(t)\|\le CA$.
Then $V\preceq\widehat V_s\preceq V+2dK_0I$.

For clarity retain all scale-dependent errors in the conditional tick
calculation.  With $p_H=0,p_P=2\vartheta/3$ and $p_s\epsilon=df_s$, its
uncorrected leading term is $d(\dot V+f_sC)$.
Its error is bounded by
$C[d^2 A+d^2B+d\epsilon B+d\epsilon^3A]$.
The correction contributes $d(M\chi-\chi)_sC=-d(f_s-1)C$.
Changing the time in this term costs $Cd^2B$, and its accepted channel part
costs
$Cd p_s(\epsilon B+\epsilon^4A)\le C(d^2B+d^2\epsilon^3A)$.
The identity offset has zero conditional increment.
Since $d\le4\epsilon/11$, $\omega\le1$, and $A\ge1$, the complete
remainder is at most $Cd\epsilon B$.
Pairing with the actual $\rho_j$ proves
\[
 \E[W_{j+1}\mid\mathcal F_j]
       \le(1+Cd)W_j+Cd(h^4+\epsilon B),
 \qquad W_j=\tr(\rho_j\widehat V_{s_j}(jd)).
\]
This is the same positive supermartingale entrance as in the proof of
\cref{v12:thm:main}, now with its constants proved independent of the spatial
bank.  Its initial offset obeys $dK_0\le C\epsilon B$.
That stopped-supermartingale argument gives
\eqref{v13:eq:mean-bound}--\eqref{v13:eq:max-bound} without a union bound over
ticks.  Finally $0\le\tr\rho V\le CA$ on the finite spectral space.
Integrate $\min(1,C_T\mathfrak b/u)$ up to $CA$ to obtain
\eqref{v13:eq:log-max}.  The field is not replaced by the corrected observable
in any conclusion.
\end{proof}

The comparison solution in this theorem specifies the regular time interval;
no global existence theorem for all data is asserted.  It is used only in
the proof.  The finite energy, cutoff, initial preparation, and instrument
are determined by the potential, lattice, and initial data without consulting
any future value of $q$.

\section{Classical field, nonlinear insertions, and the scalar metric source}
\label{v13:sec:stress}

For each actual conditional state let
\[
 \bar q_{j,x}=\tr(\rho_j\Pi Q_x\Pi),\qquad
 \bar p_{j,x}=\tr(\rho_j\Pi P_x\Pi).
\]
Compression signs are suppressed inside an expectation when the state is
already supported on $\mathcal K$; products always precede compression.
Define the full quantum positive comparison error
$e_j=\tr(\rho_j\mathcal V_\Lambda(jd))$.

\begin{proposition}[A field topology and quartic convergence from the same error]
\label[proposition]{v13:prop:field}
At every tick, for a fixed constant $C$,
\begin{align}
 &\|\bar p_j-p(jd)\|_{2,h}^2+
   \|\bar q_j-q(jd)\|_{2,h}^2+
   \sum_i\|\delta_i(\bar q_j-q(jd))\|_{2,h}^2
       \le C e_j,\label{v13:eq:field-error}\\
 &w\sum_x\tr\rho_j(Q_x-q_x)^4\le C e_j.\label{v13:eq:quartic-error}
\end{align}
The total variances of position, canonical momentum, and spatial gradients,
with these same weights, also tend to zero whenever $e_j\to0$.
Periodic multilinear interpolation of $\bar q_j$ and piecewise constant
interpolation of $\bar p_j$ consequently converge in
$H^1(\T^3)\times L^2(\T^3)$ to $(q,q_t)$, uniformly in probability on
$[0,T]$, whenever $\mathfrak b\to0$.
All spatial nodes are retained; no Fourier truncation of the attained fields
is performed.
\end{proposition}
\begin{proof}
For a self-adjoint $A$ and scalar $a$,
$\tr\rho(A-a)^2=(\tr\rho A-a)^2+
\tr\rho(A-\tr\rho A)^2$.
Use this identity for $Q,P,\delta_iQ$ and the coercivity of
\cref{v13:lem:quartic}.  Jensen also bounds the fourth power of the mean
position error by its fourth moment.  This proves all discrete statements.
The standard cellwise multilinear interpolation estimate follows by scaling
each cell to the unit cube: its $H^1$ norm is bounded by the grid value and
first-difference norm with cutoff-independent constants.
Interpolation of a smooth reference has $H^1$ error $O(h)$; the momentum
interpolation has $L^2$ error $O(h)$.
Between ticks the reference changes by $O(d)$ in these norms.
Now apply \eqref{v13:eq:max-bound}.  This argument asserts no strong time
derivative of the piecewise-constant quantum trajectory.
\end{proof}

To identify the source insertions independently, vary lapse $L_x$, positive
spatial metric $\gamma_x$, and shift $\beta_x$, keeping the original canonical
variables fixed.  Put $J_x=\sqrt{\det\gamma_x}$ and define
\begin{equation}
 \begin{split}
 \boldsymbol E(L,\gamma,\beta)=w\sum_x\biggl[
 L_x\biggl\{\frac{P_x^2}{2J_x}
      +\frac{J_x}{2}\gamma_x^{ij}(\delta_iQ_x)(\delta_jQ_x)
      +J_xV(Q_x)\biggr\}
 +\frac{\beta_x^i}{2}\{P_x,\delta_iQ_x\}\biggr].
 \end{split}
 \label{v13:eq:metric-energy}
\end{equation}
This is the canonical discretization of the same scalar action.  Use the
projection fixed at the evaluated flat parameter when differentiating
$\Pi\boldsymbol E(L,\gamma,\beta)\Pi$.  Recomputing a source-dependent
spectral projection would be a different convention and would require its
additional derivatives.
At $L=1$, $\gamma=I$, $\beta=0$ the resulting energy and stress insertions are
\begin{align}
 \boldsymbol T_{00,x}&=\tfrac12P_x^2+\tfrac12\sum_i(\delta_iQ_x)^2+V(Q_x),\\
 \boldsymbol T_{0i,x}&=\tfrac12\{P_x,\delta_iQ_x\},\\
 \boldsymbol T_{ij,x}&=(\delta_iQ_x)(\delta_jQ_x)
       +\delta_{ij}\left[\tfrac12P_x^2
                 -\tfrac12\sum_k(\delta_kQ_x)^2-V(Q_x)\right].
 \label{v13:eq:stress}
\end{align}
In particular, $\partial_{\gamma_{ij}}\boldsymbol E=-w\boldsymbol T_{ij}/2$
at the flat parameter, with the symmetric-coordinate pairing understood.

\begin{theorem}[Defect-free scalar stress and same-action first variations]
\label[theorem]{v13:thm:stress}
Let $T^h(q,p)$ denote \eqref{v13:eq:stress} evaluated on the classical sampled
fields, replacing $P,Q$ by $p,q$.  For every attained conditional state,
\begin{equation}
 \boxed{\quad
 w\sum_x\left|\tr(\rho_j\boldsymbol T_{\mu\nu,x})
                    -T^h_{\mu\nu,x}(q(jd),p(jd))\right|
       \le C\bigl(\sqrt{e_j}+e_j\bigr).
 \quad}
 \label{v13:eq:stress-error}
\end{equation}
The same estimate holds for the fixed quartic force $V'(Q)$ in weighted
$L^1$ and for the independently differentiated potential-parameter insertions
$\partial_\lambda V(Q)$ and $\partial_bV(Q)$.
Consequently all these scalar source insertions converge, uniformly in
probability on the regular slab and in local spacetime $L^1$, if
$\mathfrak b\to0$.  For $0<u\le1$, the supremum-in-time failure probability
for a source error larger than $u$ is bounded by
$C_T\mathfrak b/u^2$, apart from the deterministic $O(h+d)$ reconstruction
error.  Their expected fixed-time $L^1$ errors are at most
$C_T(\sqrt{\mathfrak b}+\mathfrak b)$.
\end{theorem}
\begin{proof}
Write each linear observable $A=P_x$ or $\delta_iQ_x$ as its scalar reference
$a$ plus $E_A$.  Expanding its square gives
\[
 |\tr\rho A^2-a^2|
 \le2|a|\sqrt{\tr\rho E_A^2}+\tr\rho E_A^2.
\]
For a symmetric cross product use
$|\tr\rho E_AE_B|\le
(\tr\rho E_A^2)^{1/2}(\tr\rho E_B^2)^{1/2}$.
Weighted Cauchy--Schwarz and the bounded reference energy control their
spatial sums by $C(\sqrt{e_j}+e_j)$.
For the potential and force use the exact expansions in $e=Q_x-q_x$.
The linear term is bounded after summation by $C\sqrt{e_j}$;
all terms of degrees two through four are bounded by
$C\tr\rho(e^2+e^4)$, using $|e|^3\le(e^2+e^4)/2$ in multiplication
functional calculus.  The quartic coercivity bounds the sum by $Ce_j$.
The two parameter derivatives are polynomials of degree at most four, so the
same proof applies.  Differentiating \eqref{v13:eq:metric-energy} directly
verifies that these are the declared independent action insertions.

The bound $C(\sqrt e+e)>u$, $u\le1$, implies $e>c u^2$ after a fixed
choice of $c>0$.  Apply \eqref{v13:eq:max-bound}.  At fixed time apply Jensen
to \eqref{v13:eq:mean-bound}.  Classical spatial quadrature and reference time
regularity give the deterministic reconstruction errors.  Integration on a
finite time interval gives the spacetime statement without a temporal
independence assumption.
\end{proof}

\paragraph{Canonical chronology and its limits.}
The convergent canonical mean fields satisfy the weak Hamilton equations of
the scalar action: pair $(\bar q_j,\bar p_j)$ with fixed smooth spacetime
tests, integrate by parts in time, and use the established convergence of
position, momentum, gradients, and $V'(Q)$.  The limit is
$q_t=p$, $p_t=\Delta q-V'(q)$, with the same coefficients.
This is not an assertion that the finite step reader has ordinary velocities,
or that its finite-difference kinetic density equals $\tr\rho P^2$.
The momentum is independently represented by the canonical operator and its
limit equals $q_t$; the stronger finite-record Legendre-incidence criterion of
Version 8 remains a different test.  Likewise these are scalar matter first
variations, not the vanishing of an independently varied Einstein metric row.

\section{An explicit simultaneous limit and the remaining native boundary}
\label{v13:sec:limit}

\begin{corollary}[A fully finite, spatially coupled scalar branch]
\label[corollary]{v13:cor:joint}
For the preparation in \cref{v13:thm:preparation}, one has
\begin{equation}
 \mathfrak b\le C\left[
 \frac\nu h+\frac1{1+\Lambda}+h^4
            +\frac{\epsilon(1+\Lambda)^2}{\nu h^3}\right].
 \label{v13:eq:full-budget}
\end{equation}
In particular choose
\begin{equation}
 \boxed{\nu_h=h^3,\qquad \Lambda_h=h^{-2},\qquad
        \epsilon_h=h^{12},\qquad
        d_h=\frac{4\vartheta_h}{11}h^{12},\quad0<\vartheta_h\le1.}
 \label{v13:eq:scales}
\end{equation}
Then $\mathfrak b_h=O(h^2)$, uniformly in $\vartheta_h$.
Every cutoff has a finite spectral carrier, six independent original accepted
outcomes, and the unchanged two-phase clock.  On each regular classical time
slab, their actual conditional scalar readers converge in
$H^1_x\times L^2_x$, uniformly in probability in time, and their scalar stress
has no remaining kinetic, gradient, or quartic defect.
The mean of the maximal relative error is $O(h^2(1+|\log h|))$.
\end{corollary}
\begin{proof}
Insert \eqref{v13:eq:finite-prep} into \eqref{v13:eq:budget}.
In \eqref{v13:eq:scales}, the four terms have orders
$h^2,h^2,h^4,h^2$ respectively, since
$\nu_hh^3=h^6$ and $(1+\Lambda_h)^2=O(h^{-4})$.
The positive finite-rank branch applies for small $h$ by
\cref{v13:thm:preparation}.  The field and stress conclusions are
\cref{v13:prop:field,v13:thm:stress}.
The logarithm in \eqref{v13:eq:log-max} has order $1+|\log h|$.
This gives the final estimate.  The exponents are sufficient, not optimized;
in particular no claim of computational efficiency is made.
\end{proof}

For any dyadic spatial sequence, the failure probabilities for a fixed
positive error threshold are summable.  The first Borel--Cantelli implication
therefore also gives almost-sure convergence along that sequence on any joint
probability space carrying these cutoff processes.  No cross-cutoff
independence is needed.  This is not, by itself, a summable adjacent-cutoff
\emph{norm} estimate or a source-native coupling of different cutoffs.

\paragraph{What has actually been removed.}
The spatially refining reader no longer pays a stability constant growing
like the norm of the lattice Laplacian.  The positive comparison contains
the complete scalar canonical momentum, gradient, and quartic content;
its bound is derived from the same Hamiltonian and the accepted channel.
The source does not begin by sampling a prescribed future solution.
It begins with finitely represented initial data and evolves by its displayed
instrument.  A smooth reference is used for a local-in-time convergence
proof, not as a list of future updates.  Finite spectral compression does
not create spurious exact canonical commutation relations, because every
composite insertion is compressed only after it has been formed.

\paragraph{What remains independent.}
This is a \emph{constructed real-Higgs action-compatible branch on a prescribed
flat geometry}.  It is stronger than the fixed spin example as a spatial
field and scalar-stress result, but it does not identify the actual locked
source with this instrument.  The latter has global spectral and unitary
operations whose microscopic locality has not been established.  The
semiclassical scale and squeezed preparation are declared, not derived from
primitive renewal probabilities.  Gauge/color links, chiral spinor variables,
Higgs zeros in a dressing construction, Einstein constraints, metric
backreaction, and the native common-action incidence are not supplied by the
scalar argument.  The scalar stress is generally nonzero, so keeping the
metric flat cannot be advertised as a coupled Einstein solution.

The next action-specific target is consequently the same cancellation for
an independently identified joint metric--gauge--matter observable energy,
or a source-local realization of the scalar accepted-channel estimates above.
Repeating the fixed-bank polynomial lemma, the clock Poisson equation, or a
newly labelled compactness assumption would not discharge either target.

\paragraph{Verification and preservation.}
The full Version 12 prefix before its final document command is retained
byte-for-byte.  The accompanying finite diagnostics check quartic coercivity,
classical and low-energy operator cancellations, Gaussian preparation costs,
nonmultiplicativity of spectral compression, the exact locked channel on a
small spectral test carrier, and the explicit scale budget.  The numerical
spectral test uses a convergent finite oscillator-basis approximation and is
labelled as such; it is not an exact diagonalization of the full spatial
Schr\"odinger operator.  These are diagnostics of the written proofs, not
machine-checked formal proofs or an Einstein--SM simulation.  Every added
label has prefix \texttt{v13:}.

\begingroup
\renewcommand{\refname}{Additional references for Version 13}
\begin{thebibliography}{9}
\small\raggedright
\bibitem{v13:DT}
B.~K.~Driver and P.~W.~Tong,
\emph{On the classical limit of quantum mechanics I}, 2015 preprint.
\href{https://arxiv.org/abs/1511.02447}{arXiv:1511.02447}.
Background for coherent-state limits and polynomial position/momentum
observables; no result is used in place of the relative-energy calculation.
\bibitem{v13:GP}
F.~Golse and T.~Paul,
\emph{The Schr\"odinger equation in the mean-field and semiclassical regime},
2015 preprint; Archive for Rational Mechanics and Analysis
\textbf{223} (2017), 57--94.
\href{https://arxiv.org/abs/1510.06681}{arXiv:1510.06681}.
Background for quantum--classical comparison estimates.  The present quartic
lattice, spectral-cutoff, and resolved-renewal-clock statements are proved
separately above.
\end{thebibliography}
\endgroup
% END IMMUTABLE VERSION 13 BODY
% APPEND VERSION 14 HERE, BEFORE THE FINAL END-DOCUMENT COMMAND.


% BEGIN IMMUTABLE VERSION 14 BODY
\clearpage
\begin{center}
{\Large\bfseries Version 14: Local tensor records before global spectral updates}\\[0.5em]
{\large A quantitative product-cutoff limit, local locked instruments, and scalar stress on the actual renewal clock}
\end{center}
% Three-digit section numbers require a wider TOC number box.
% This formatting change is scoped to the appended TOC entries.
\makeatletter
\addtocontents{toc}{\protect\begingroup}
\addtocontents{toc}{\protect\makeatletter}
\addtocontents{toc}{\protect\def\protect\l@section{\protect\bfseries\protect\@dottedtocline{1}{0em}{2.5em}}}
\makeatother
\addcontentsline{toc}{part}{Version 14: Local finite instruments and tensor-cutoff scalar closure}

\section{Audit and the precise locality problem}
\label{v14:sec:audit}

Version 13 proves spatial scalar and stress convergence, but its global
energy projection and global accepted unitary are not local operations.
Locality of a Hamiltonian does not imply locality of an instrument obtained
by exponentiating the whole Hamiltonian.  This iteration replaces both of
those regulator operations.  The source will have a tensor product of
finite one-site spaces, a product initial preparation, and accepted Kraus
operators supported on one site or one nearest-neighbor edge.

\paragraph{What is inherited rather than repeated.}
The physical normalization, quartic scalar action, uncompressed
Schr\"odinger energy, relative-energy cancellation, and scalar stress
insertions are \cref{v13:eq:CCR,v13:eq:quantum-energy,v13:thm:cancellation,v13:thm:stress}.
The six-outcome locked coefficient identity and the elementary
observable-specific conjugation estimate are
\cref{v12:thm:locked-bank}.  The actual two-state clock and its Poisson
corrector are \cref{v10:eq:phase,v10:eq:chi}.  Their general proofs are not
new results here.  Targeted searches and literal-term checks in the supplied
Grand Tensor, classical closure, and SM--Einstein monographs located the
existing channel product formulas and classical Galerkin results, but not
this tensor-cutoff, local-instrument, energy-norm composition.  This is a
limited nonduplication audit, not a claim of priority for Galerkin or
randomized Hamiltonian approximation.

\paragraph{The extra spatial selection rule is declared.}
At an accepted completion, a site or edge label is drawn uniformly and then
one of its six locked outcomes is recorded.  Each local context satisfies
its original six-outcome locking equations.  The entire spatial source has
\emph{six times the number of local contexts} as its accepted alphabet; it
is not falsely described as the original single six-outcome global
instrument.  This spatial selection rule, the scalar quantization, and the
initial preparation are specified parts of the constructed branch.  They
have not been derived from the bare locked opportunity table.

\paragraph{Attribution and scope.}
Random selection of local Hamiltonian terms is an established simulation
method; see \cite{v14:Campbell}.  The present estimates are direct proofs for
local locked instruments, resolved conditional records, and the retained
non-Poisson clock.  The positive result removes the global-operation
requirement from the constructed scalar branch.  It does not identify the
intended primitive Einstein--SM law, prove a relativistic speed bound, or
supply dynamical geometry, color, or chiral matter.  The physical metric
continues to be prescribed and flat.

\section{A tensor-product cutoff with explicit energy control}
\label{v14:sec:tensor}

Retain the three-torus lattice of Version 13, with $w=h^3$, $n=h^{-3}$,
$0<\nu\le h$, and $\omega=\nu w$.  Write $E=\boldsymbol E$ for its
nonnegative scalar Schr\"odinger energy on $L^2(\R^n)$.  At one site define
\begin{equation}
 \ell_x=\tfrac12P_x^2+\tfrac\lambda4Q_x^4+I,\qquad
 A=\sum_x\ell_x,\qquad B=E+I.
 \label{v14:eq:reference-oscillator}
\end{equation}
All operators and forms in this section are uncompressed.  The one-site
quartic operator has discrete spectrum by the same confining-form argument
as in \cref{v13:lem:spectral}.  For $R\ge1$ use the \emph{local} projection
\begin{equation}
 P_R=\bigotimes_x\one_{[0,R]}(\ell_x),\qquad
 \mathcal K_R=\Ran P_R,\qquad E_R=P_REP_R|_{\mathcal K_R}.
 \label{v14:eq:local-cutoff}
\end{equation}
The projection is fixed as part of the regulator, not applied to each
attained record.  Its tensor structure, unlike a cutoff of $E$, preserves
supports of local operators.  Composite insertions are still formed before
compression.

Here are explicit constants sufficient for the approximation estimates:
\begin{equation}
 \begin{gathered}
 \beta_h=6h^{-2}+\lambda b/2,\qquad
 c_h=n(\lambda b^2/4-1),\qquad
 b_h=2\beta_h\sqrt{n/\lambda},\\
 C_h^{\rm g}=2w^{-1}+b_h^2+2|c_h|+1,\qquad
 C_h^{\rm f}=2w+1+w(b_h^2+|c_h|),\qquad
 D_h=C_h^{\rm f}(C_h^{\rm g})^2.
 \label{v14:eq:graph-constants}
 \end{gathered}
\end{equation}
The symbol $b_h$ here is a quadratic-potential bound, not an error budget.
For fixed $\lambda>0,b\ge0$, these constants obey
$D_h\le C h^{-18}$ for $0<h\le1$.

\begin{lemma}[Local cutoff approximation in the physical energy form]
\label[lemma]{v14:lem:form}
The domains of $E$ and $A$ agree, and for $u$ in this domain,
\begin{equation}
 \|Au\|\le C_h^{\rm g}\|Bu\|,\qquad
 \langle u,Bu\rangle\le C_h^{\rm f}\langle u,Au\rangle.
 \label{v14:eq:graph-form}
\end{equation}
Consequently, with $\|u\|_B^2=\langle u,Bu\rangle$,
\begin{equation}
 \boxed{\quad
 \|(I-P_R)u\|_B^2\le \frac{D_h}{R}\|Bu\|^2.
 \quad}
 \label{v14:eq:tail}
\end{equation}
The spaces $\mathcal K_R$ form an increasing form core for $E$ as
$R\to\infty$ at fixed $h,\nu$.  None of these assertions says that $P_R$
commutes with the interacting energy $E$.
\end{lemma}
\begin{proof}
Expand the prescribed potential and gradients to write
\[
 E/w=A+W_2+c_hI,
 \qquad
 W_2=\frac1{2h^2}\sum_{x,i}(Q_{x+e_i}-Q_x)^2
                  -\frac{\lambda b}{2}\sum_xQ_x^2.
\]
As multiplication functions, $|W_2|\le\beta_h\sum_xQ_x^2$.
Since $(\sum_xQ_x^2)^2\le n\sum_xQ_x^4$ and
$A\succeq(\lambda/4)\sum_xQ_x^4$ as forms,
\begin{equation}
 \|W_2u\|\le b_h\|A^{1/2}u\|.
 \label{v14:eq:lower-order}
\end{equation}
In particular $W_2$ is infinitesimally $A$-bounded: use
$\|A^{1/2}u\|^2\le\|Au\|\|u\|$ and the scalar Young inequality.
For clarity, the usual perturbation conclusion follows here from a
resolvent Neumann series: $\|W_2(A-\ii t)^{-1}\|$ tends to zero as
$|t|\to\infty$ by \eqref{v14:eq:lower-order} and the spectral calculus of
$A$.  The same holds at $-\ii t$.  The symmetric operator
$A+W_2+c_h$ on $D(A)$ therefore has surjective resolvents at both signs
and is self-adjoint.  Its form is the scalar energy form on a core, so it
is the Friedrichs realization already retained for $E$.

Now
\[
 \|Au\|\le w^{-1}\|Eu\|+b_h\|A^{1/2}u\|+|c_h|\|u\|
 \le w^{-1}\|Eu\|+\tfrac12\|Au\|
                  +(b_h^2/2+|c_h|)\|u\|.
\]
Because $E\ge0$, both $\|Eu\|$ and $\|u\|$ are at most $\|Bu\|$.
This proves the graph bound with the displayed constant.  Similarly,
\[
 \langle u,Bu\rangle
 \le w\langle u,Au\rangle
       +wb_h\|u\|\|A^{1/2}u\|+(1+w|c_h|)\|u\|^2
 \le C_h^{\rm f}\langle u,Au\rangle,
\]
using $A\ge I$ and Young's inequality.  The reverse form comparison with
some finite fixed-$h$ constant follows from the same calculation after
absorbing half of $\langle A\rangle$; hence the two form domains coincide.

All $\ell_x$ strongly commute and have a joint tensor eigenbasis.  On
$\Ran(I-P_R)$ at least one local eigenvalue exceeds $R$, so the corresponding
eigenvalue of $A$ exceeds $R$.  Since $P_R$ commutes with $A$,
\[
 \|A^{1/2}(I-P_R)u\|^2\le R^{-1}\|Au\|^2.
\]
The two preceding bounds give \eqref{v14:eq:tail}.  The same joint spectral
expansion gives convergence in the $A$ form norm for every vector in its
form domain.  Form equivalence proves the final core assertion.
\end{proof}

\begin{proposition}[Compression is local but not multiplicative]
\label[proposition]{v14:prop:local-composites}
Write $Q_x^R=P_RQ_xP_R$ on $\mathcal K_R$ and
\[
 D_x^R=P_RQ_x^2P_R-(Q_x^R)^2
       =P_RQ_x(I-P_R)Q_xP_R\succeq0.
\]
Both $Q_x^R$ and $D_x^R$ are supported on site $x$.  For distinct neighboring
sites $x,y$,
\begin{equation}
 P_R(Q_y-Q_x)^2P_R=(Q_y^R-Q_x^R)^2+D_x^R+D_y^R.
 \label{v14:eq:edge-compression}
\end{equation}
Thus $E_R$ is exactly a sum of $n$ one-site positive terms and $3n$
nearest-neighbor positive terms.  An observable supported away from a given
term commutes with it.  Discarding $D_x^R,D_y^R$ would change that energy.
\end{proposition}
\begin{proof}
Factor the tensor projection into the factors at $x,y$ and their complement.
Compression of a one-site observable has that same one-site support.
Compression of $Q_xQ_y$ factors as $Q_x^RQ_y^R$ because the sites differ.
Expand the square, insert the two diagonal defects, and use positivity of
$P_RQ_x(I-P_R)Q_xP_R$.  The uncompressed action is precisely the sum of the
stated one-site and edge terms, each positive, and compression preserves
positivity.  This proves the support and energy assertions.
\end{proof}

\section{Quantitative tensor Galerkin transport in energy norm}
\label{v14:sec:Galerkin}

The exact-compression identity of Version 13 used a spectral projection of
$E$.  It is not available for $P_R$.  The following estimate pays the
resulting error rather than pretending that identity persists.
Let $\phi$ be a unit vector in $D(B^2)$, set
\[
 M_\phi=\|B^2\phi\|,\quad
 \phi_R=P_R\phi/\|P_R\phi\|,\quad
 \psi(t)=e^{-\ii tE/\omega}\phi,\quad
 \psi_R(t)=e^{-\ii tE_R/\omega}\phi_R.
\]
Vectors in $\mathcal K_R$ are identified with their original tensor
subspace; no global energy filter is introduced.

\begin{theorem}[Energy-form approximation with a polynomial time/quantum cost]
\label[theorem]{v14:thm:Galerkin}
If $D_hM_\phi^2/R\le1/4$, then
\begin{equation}
 \boxed{
 \sup_{0\le t\le T}\|\psi_R(t)-\psi(t)\|_B^2
 \le C\frac{D_h}{R}M_\phi^2(1+M_\phi)^2
                         (1+T/\omega)^2.}
 \label{v14:eq:Galerkin}
\end{equation}
The constant $C$ is absolute.  In particular there is no exponential of
$\|E_R/\omega\|$ or of the spatial stiffness.  The projection used by the
source remains $P_R$; the Ritz map in the proof is only an estimate.
\end{theorem}
\begin{proof}
Let $\mathfrak R_R$ be the orthogonal projection onto $\mathcal K_R$ in the
$B$ form inner product.  The best-approximation property and
\cref{v14:lem:form} imply
\[
 \|u-\mathfrak R_Ru\|_B\le\|u-P_Ru\|_B
       \le\sqrt{D_h/R}\|Bu\|,
 \quad
 \|(P_R-\mathfrak R_R)u\|_B\le2\sqrt{D_h/R}\|Bu\|.
\]
For $u\in D(B^2)$, its defining form identity gives
\[
 E_R\mathfrak R_Ru=P_REu+P_Ru-\mathfrak R_Ru,
 \qquad
 (E_R\mathfrak R_R-\mathfrak R_RE)u
                 =(P_R-\mathfrak R_R)Bu.
\]
Consequently $e_R=\psi_R-\mathfrak R_R\psi$ satisfies
\[
 \ii\omega\dot e_R=E_Re_R+(P_R-\mathfrak R_R)B\psi.
\]
On $\mathcal K_R$ the $B$ form is $E_R+I$.  The unitary generated by
$E_R$ is therefore an isometry in this form norm.  Duhamel's formula gives
\[
 \|e_R(t)\|_B
 \le\|e_R(0)\|_B+\frac{2T}{\omega}\sqrt{D_h/R}\,M_\phi,
\]
since $\|B^2\psi(t)\|=M_\phi$.

Put $t_R=\|(I-P_R)\phi\|_B\le\sqrt{D_h/R}M_\phi\le1/2$.
Then $\|P_R\phi\|^2\ge1-t_R^2\ge3/4$ and
$|\|P_R\phi\|^{-1}-1|\le2t_R^2$.
Also $\|P_R\phi\|_B\le M_\phi+t_R$.
It follows by the triangle inequality that
$\|e_R(0)\|_B\le C(1+M_\phi)t_R$.
Finally $\|\psi(t)-\mathfrak R_R\psi(t)\|_B
\le\sqrt{D_h/R}M_\phi$.  Combine the three bounds and square.
The argument can first be made for vectors with bounded spectral support
for $E$, then passed to $D(B^2)$; the displayed estimates are continuous
in its graph norm.  No commuting tensor-cutoff assumption is used.
\end{proof}

\begin{lemma}[The needed higher moment is supplied by the local Gaussian preparation]
\label[lemma]{v14:lem:Gaussian}
For the product Gaussian $\phi_{h,\nu}$ of
\cref{v13:eq:Gaussian}, with bounded initial values and first spatial
differences, $0<\nu\le h\le1$ implies
\begin{equation}
 \|(E+I)^2\phi_{h,\nu}\|\le C.
 \label{v14:eq:Gaussian-high}
\end{equation}
The constant depends only on these initial bounds and on $\lambda,b$.
The normalized finite preparation $P_R\phi/\|P_R\phi\|$ is a product of
normalized one-site projected Gaussians.
\end{lemma}
\begin{proof}
Every local density in $E$ is a fixed polynomial of degree at most four in
$P_x,Q_x,\delta_iQ_x$, with coefficients depending only on $\lambda,b$.
On the Gaussian, all these elementary linear observables have bounded means
and bounded ordered two-point centered expectations, uniformly in $h,\nu$:
position variance is $\nu h/2$, momentum variance $\nu/(2h)$, and each
spatial-difference variance is $\nu/h$.  Mixed contractions are bounded by
Cauchy--Schwarz.  Their commutators are scalar with size at most $\nu/h$.
For completeness, differentiating the Gaussian exponential generating
function, with the scalar commutator factor retained when ordering two
exponentials, gives the ordered Gaussian moment rule: each centered even
word is the sum of pair contractions in its stated order, and odd centered
words have zero expectation.  It follows that every word of length at most
sixteen in these observables has uniformly bounded expectation.  Nonzero
means are handled by expanding each factor around its mean.

The expression $\langle\phi,E^4\phi\rangle$ is the sum of four local-density
products multiplied by $w^4$.  Each summand is a fixed number of such words,
and there are $n^4$ site choices.  Since $wn=1$, the sum is uniformly
bounded.  The same argument handles $E^j$ for $j\le4$, proving
\eqref{v14:eq:Gaussian-high}.  The Gaussian belongs to every polynomial
operator domain, so these calculations justify the moment rather than
presuppose it.  Lastly $P_R$ is a tensor product and $\phi$ is a product;
its norm and the normalization factor factorize site by site.
\end{proof}

\section{A local six-outcome instrument in each spatial context}
\label{v14:sec:local-instrument}

Write the exact local decomposition from \cref{v14:prop:local-composites} as
\begin{equation}
 E_R=\sum_{\alpha=1}^{M}E_{\alpha,R},\qquad
 M=4n,\qquad H_\alpha=E_{\alpha,R}/\omega,\qquad
 H_R=\sum_\alpha H_\alpha.
 \label{v14:eq:local-sum}
\end{equation}
Each $H_\alpha$ acts on one site or one neighboring pair.  Put
\begin{equation}
 g_\alpha=\|H_\alpha\|,\quad L=\sum_\alpha g_\alpha,\quad
 \mathfrak C=1+L^2+M\sum_\alpha g_\alpha^2+M^2,\quad
 A_R=1+\|E_R\|.
 \label{v14:eq:local-cost}
\end{equation}
These are finite action-derived norms, not unrecorded stability constants.
Their size will be paid by the physical accepted-event scale.

At an accepted completion choose $\alpha$ uniformly, independently of the
current state.  Set $\tau=M\epsilon$ and use the six-outcome family of
\cref{v12:thm:locked-bank} \emph{on that local support}, with Hamiltonian
$H_\alpha$ and small parameter $\tau$.  More explicitly, in an eigenbasis
$u_0,u_1,\ldots$ of its local matrix choose
\[
 C_{\alpha,\mu}=|u_0\rangle\langle u_{\mu-1}|\ (2\le\mu\le5),\quad
 R_\alpha=\sum_{\mu=2}^5C_{\alpha,\mu}^*C_{\alpha,\mu},
\]
\[
 B_{\alpha,1}=e^{-3\ii\tau H_\alpha/2}(I-\tau^4R_\alpha)^{1/2},
 \qquad B_{\alpha,\mu}=\tau^2C_{\alpha,\mu}.
\]
The same fixed real normal-form matrix $O$ as in Version 12 gives
\begin{equation}
 K_{\alpha,e}=I/\sqrt{18}+\sqrt{2/3}
                       \sum_{\mu=1}^5O_{e\mu}B_{\alpha,\mu}.
 \label{v14:eq:local-K}
\end{equation}
All these matrices are tensored with the identity outside their local
support.  A sufficiently large one-site cutoff supplies local dimension at
least five.  This follows, for $\nu\le1$, by using five fixed compactly
supported smooth trial functions in the one-site variational principle;
their $\ell_x$ form matrix is bounded independently of $\nu$.

\begin{proposition}[Exact contextual locking and a finite local channel budget]
\label[proposition]{v14:prop:channel}
For each context $\alpha$, the original six outcomes satisfy
\[
 \sum_eK_{\alpha,e}^*K_{\alpha,e}=I,\qquad
 \sum_eK_{\alpha,e}=\sqrt2 I.
\]
They are independent and have positive outcome probabilities when
$0<\tau\le\tau_0$, for a fixed sufficiently small $\tau_0$.
The actual global accepted Kraus operators are
$M^{-1/2}K_{\alpha,e}$, with both $\alpha$ and $e$ retained.
Their channel is
$\mathscr E_* =M^{-1}\sum_\alpha\mathscr E_{\alpha,*}$.
For $\mathscr D_\alpha O=\ii[H_\alpha,O]$ and
$\mathscr D O=\ii[H_R,O]$, one has
\begin{align}
 \|\mathscr E^*O-O-\epsilon\mathscr DO\|
      &\le C\epsilon^2\mathfrak C\|O\|,\label{v14:eq:channel-second}\\
 \|\mathscr E^*O-O\|
      &\le C\epsilon(L+M)\|O\|.\label{v14:eq:channel-first}
\end{align}
A local channel acts as the identity on an observable with disjoint support.
No global energy eigenbasis or global unitary is used by an accepted update.
\end{proposition}
\begin{proof}
The inherited locked construction applies to each finite local matrix,
including a scalar $H_\alpha$: the auxiliary diagonal modulus separates
$B_{\alpha,1}$ from $I$, and the four off-diagonal units provide the other
independent directions.  Its least singular-value bound is independent of
the spectrum.  Tensoring with an outside identity preserves these facts.
Averaging the local trace-preserving maps proves global normalization.
The local observable estimate is
\[
 \|\mathscr E_\alpha^*O-O-\tau\mathscr D_\alpha O\|
       \le\tfrac34\tau^2\|\mathscr D_\alpha^2O\|
                       +C\tau^4\|O\|.
\]
Since $\|\mathscr D_\alpha^2O\|\le4g_\alpha^2\|O\|$, averaging and
using $\tau=M\epsilon$ give an error at most
$C[\epsilon^2M\sum g_\alpha^2+\epsilon^4M^4]\|O\|$.
For $\tau\le1$, $\epsilon^4M^4\le\epsilon^2M^2$.
This proves \eqref{v14:eq:channel-second} without expanding an exponential
under a small-$\tau g_\alpha$ assumption.  Once integrating unitary
conjugation instead gives
$\|\mathscr E_\alpha^*O-O\|\le2\tau g_\alpha\|O\|+C\tau^4\|O\|$.
Averaging, and using $\tau^4\le\tau$, proves
\eqref{v14:eq:channel-first}.  Finally local support makes a disjoint
observable commute with every $K_{\alpha,e}$, so completeness gives the
identity on that observable.
\end{proof}

\paragraph{Precisely which locality has been proved.}
Every conditional operation has one-site or nearest-neighbor operator
support.  The product preparation and tensor cutoff use no global spectral
selection.  The context label is additional classical source data, selected
by the displayed spatial rule; it is not derived from the original edge
label.  A conditional state of a distant subsystem may still change through
pre-existing correlations when an outcome is learned.  The disjoint-support
statement is the exact nonselective identity, not a denial of this
conditional effect.  Neither bounded microscopic speed nor a relativistic
continuum light cone is inferred solely from operator support.

\section{Resolved quantum transport without an exponential generator cost}
\label{v14:sec:clock}

Use the literal two-state clock, with no field change on rejection,
$d=4\vartheta\epsilon/11$, and the averaged accepted instrument of
\cref{v14:prop:channel}.  Its resolved outcomes are still those just
specified, not a Kraus representation chosen after averaging.  Let
$\rho_0=|\phi_R\rangle\langle\phi_R|$ and retain every actual conditional
state $\rho_j$.  Define the proof-only reference
\[
 \Pi_*(t)=|\psi_R(t)\rangle\langle\psi_R(t)|,\qquad
 V_*(t)=I-\Pi_*(t),\qquad
 \xi_j=\tr(\rho_jV_*(jd)).
\]
This rank-one reference projection is not a physical measurement or update.

\begin{theorem}[Uniform-in-record local-instrument comparison]
\label[theorem]{v14:thm:trajectory}
For $M\epsilon\le\tau_0$ and every initial clock phase,
\begin{align}
 \max_{j\le T/d}\E\xi_j&\le C_T\epsilon\mathfrak C,\\
 \PP\{\max_{j\le T/d}\xi_j>u\}&\le
                         C_T\epsilon\mathfrak C/u\quad(u>0).
 \label{v14:eq:trajectory}
\end{align}
The constant $C_T$ is independent of the number of sites, local Hilbert
dimensions, $\vartheta$, and the Hamiltonian norm.  The scale dependence is
entirely in $\mathfrak C$.  In particular no factor
$\exp(T\|H_R\|)$ occurs.
\end{theorem}
\begin{proof}
The exact unitary reference obeys
$\dot V_*+\mathscr DV_*=0$.  Also
\[
 \|V_*\|\le1,\quad
 \|\dot V_*\|+\|\mathscr DV_*\|\le4L,\quad
 \|\ddot V_*\|+\|\mathscr D\dot V_*\|\le8L^2.
\]
Put $C_*(t)=\mathscr DV_*(t)$ and retain the explicit phase corrector
$\chi_H=-15/22$, $\chi_P=25/44$.
With $K_0=\|\chi\|_\infty\sup_t\|C_*(t)\|$, define
\[
 \widehat V_s(t)=V_*(t)+d\chi_sC_*(t)+dK_0I.
\]
Then $0\preceq V_*\preceq\widehat V_s\preceq V_*+2dK_0I$.

Here are the scale-dependent bounds in the conditional calculation.
The uncorrected clock step equals
$d(\dot V_*+f_s\mathscr DV_*)$ plus an error bounded by
$C[d^2L^2+d\epsilon\mathfrak C]$:
use \eqref{v14:eq:channel-second}, the time Taylor remainder, and
$\|\mathscr D(V_*(t+d)-V_*(t))\|\le4dL^2$.
The correction contributes $-d(f_s-1)C_*$, a time error $Cd^2L^2$,
and an accepted part bounded by
\[
 C d p_s\epsilon(L+M)\|C_*\|
       \le C d^2(L^2+M^2),\qquad
 p_H=0,\ p_P=2\vartheta/3,\ p_s\epsilon=df_s.
\]
The last estimate uses \eqref{v14:eq:channel-first}.
Since $d\le4\epsilon/11$, all the remainders are at most
$C d\epsilon\mathfrak C$.  The leading term is exactly zero.
Pair with the actual conditional state and set
$W_j=\tr(\rho_j\widehat V_{s_j}(jd))$.  One obtains
\[
 \E[W_{j+1}\mid\mathcal F_j]\le W_j+C d\epsilon\mathfrak C,
 \qquad \E W_0\le C d L\le C\epsilon\mathfrak C.
\]
No growth term proportional to $\|H_R\|W_j$ has been introduced.
The nonnegative process
$W_j+C\epsilon\mathfrak C(T-jd+d)$ is a supermartingale up to the last
tick, after increasing the fixed constant.  Its first-threshold stopping
argument gives \eqref{v14:eq:trajectory}; iteration gives the mean bound.
This repeats the elementary maximal principle already proved in Version 12,
but the zero reference drift and the \emph{local} channel estimates are the
new inputs.  All conditioning is on the complete resolved history.
\end{proof}

\section{Transfer back to the physical scalar energy and stress}
\label{v14:sec:field}

Let $q$ be a smooth solution of \cref{v13:eq:wave} on its declared regular
slab, and use its sampled values only in the comparison observable
$\mathcal V(t)$ of \cref{v13:eq:relative-energy}.  Initial values determine
the product Gaussian and all finite source operations.  The constants below
depend on the reference bounds in Version 13, but the source never consults
future values of $q$.

\begin{lemma}[Continuous quantum comparison from the inherited cancellation]
\label[lemma]{v14:lem:quantum-relative}
For the full scalar quantum trajectory
$\psi(t)=e^{-\ii tE/\omega}\phi_{h,\nu}$,
\begin{equation}
 \sup_{t\le T}\langle\psi(t),\mathcal V(t)\psi(t)\rangle
              \le C_T(\nu/h+h^4).
 \label{v14:eq:full-relative}
\end{equation}
For its tensor Galerkin approximation,
\begin{equation}
 \sup_{t\le T}\langle\psi_R(t),\mathcal V(t)\psi_R(t)\rangle
 \le C_T\left[\frac\nu h+h^4+\frac{D_h}{R}(1+T/\omega)^2\right].
 \label{v14:eq:tensor-relative}
\end{equation}
\end{lemma}
\begin{proof}
The Gaussian preparation formula of Version 13 gives an initial relative
energy at most $C\nu/h$.  Its operator cancellation gives
$\frac{d}{dt}\langle\mathcal V\rangle
\le C\langle\mathcal V\rangle+Ch^4$ for the quantum evolution.
This can be justified first on global energy spectral truncations, which
commute with $E$ exactly, and then passed to the uncompressed evolution.
Indeed the truncated initial vectors converge in the $B$ form norm,
unitary evolution preserves that convergence uniformly in $t$, and
$0\preceq\mathcal V(t)\preceq C B$ uniformly.  Thus the expectation
converges uniformly as well.  These global projections occur only in this
proof, not in the source construction.  Integrating the differential
inequality gives \eqref{v14:eq:full-relative} with $C_T$ independent of
$h,\nu$.

For any two vectors $u,v$ and positive $\mathcal V$,
$\|\mathcal V^{1/2}u\|^2
\le2\|\mathcal V^{1/2}v\|^2+2\|\mathcal V^{1/2}(u-v)\|^2$.
Use $u=\psi_R$, $v=\psi$, form domination by $B$, and
\cref{v14:thm:Galerkin,v14:lem:Gaussian}.  This proves
\eqref{v14:eq:tensor-relative} for the stated sufficiently large cutoff.
\end{proof}

\begin{theorem}[Local finite scalar and stress closure]
\label[theorem]{v14:thm:scalar}
For the finite local source just constructed define its actual comparison
error and budget by
\begin{equation}
 e_j=\tr(\rho_jP_R\mathcal V(jd)P_R),\qquad
 \mathfrak b_{h,\nu,R,\epsilon}
 =\frac\nu h+h^4+\frac{D_h}{R}(1+T/\omega)^2
                       +\epsilon A_R\mathfrak C.
 \label{v14:eq:budget}
\end{equation}
Retain the Gaussian, cutoff, and local-step hypotheses above.  Then
\begin{align}
 \max_{j\le T/d}\E e_j&\le C_T\mathfrak b_{h,\nu,R,\epsilon},\\
 \PP\{\max_{j\le T/d}e_j>u\}&\le
                      C_T\mathfrak b_{h,\nu,R,\epsilon}/u.
 \label{v14:eq:scalar-bound}
\end{align}
When $\mathfrak b\to0$, the actual conditional scalar means converge in
$H_x^1\times L_x^2$ uniformly in probability on the regular time slab.
Canonical momentum, gradient, and quartic conditional variances vanish in
their stated mass-weighted topologies.  The scalar stress insertions formed
before tensor compression satisfy
\begin{equation}
 w\sum_x|\tr(\rho_jP_R\boldsymbol T_{\mu\nu,x}P_R)
                         -T^h_{\mu\nu,x}(q,p)|
                \le C(\sqrt{e_j}+e_j).
 \label{v14:eq:stress}
\end{equation}
The same holds for the quartic force and potential-parameter sources.
\end{theorem}
\begin{proof}
Put $V_R=P_R\mathcal V P_R$ and $P=\Pi_*(t)$, $Q=I-P$ on the tensor
carrier.  Positivity and the vector triangle inequality give the operator
bound
\[
 V_R\preceq2PV_RP+2QV_RQ
 \preceq2\langle\psi_R,V_R\psi_R\rangle P+2\|V_R\|Q.
\]
Because $\mathcal V\preceq C B$, $\|V_R\|\le C A_R$.
Pairing with any actual conditional state therefore gives
\begin{equation}
 e_j\le C_T a_{h,\nu,R}+C A_R\xi_j,\qquad
 a_{h,\nu,R}=\nu/h+h^4+D_hR^{-1}(1+T/\omega)^2.
 \label{v14:eq:positive-transfer}
\end{equation}
The first term is deterministic and bounded by
\cref{v14:lem:quantum-relative}.  The mean bound follows from
\cref{v14:thm:trajectory}.  If $u\le2C_Ta_{h,\nu,R}$, the probability
bound follows from the trivial bound one; otherwise the event forces
$\max\xi_j>u/(2CA_R)$, and the same trajectory theorem applies.
This proves \eqref{v14:eq:scalar-bound} without a per-tick union bound.

Every vector in the tensor carrier lies in the domain of the original
polynomial insertions.  The pointwise coercivity, density-matrix
Cauchy--Schwarz, and quartic expansion of
\cref{v13:prop:field,v13:thm:stress} apply to it without requiring an energy
spectral projection.  They therefore give the same field and stress bounds,
now with the local tensor compression retained throughout.  The stress
operators are the independent lapse, shift, and spatial-metric derivatives
of $P_RE(L,\gamma,\beta)P_R$ with $P_R$ fixed at the evaluated parameter.
They are not defined by differentiating a limiting mean trajectory.
The standard cellwise interpolation estimates and the smooth reference
quadrature give the stated continuum topologies.  This still does not
identify $\tr\rho P^2$ with the finite differences of the stepwise reader;
canonical momentum is the independently represented action variable.
\end{proof}

For completeness, $e_j\le C A_R$ deterministically.  Integrating the maximal
probability bound gives, for $0<\mathfrak b\le1/2$,
\begin{equation}
 \E\max_{j\le T/d}e_j
 \le C_T\mathfrak b\left[1+\log\left(2+\frac{A_R}{\mathfrak b}\right)\right].
 \label{v14:eq:log}
\end{equation}
In particular large state dimension has not been hidden in a compactness
constant, although the explicit cutoff and step costs can be very large.

\section{An explicit simultaneous local limit and its remaining costs}
\label{v14:sec:scales}

\begin{proposition}[Norm costs of the local scalar decomposition]
\label[proposition]{v14:prop:scales}
For $R\ge C\max(1,h^{-4})$ and $0<\nu\le h$, the local construction obeys
\begin{equation}
 A_R\le C R,\qquad
 \mathfrak C\le C R^2/\omega^2,\qquad D_h\le C h^{-18}.
 \label{v14:eq:costs}
\end{equation}
The constants are independent of local Hilbert dimension.  Thus one sufficient
bound on the scalar budget is
\begin{equation}
 \boxed{
 \mathfrak b\le C_T\left[
 \frac\nu h+h^4+\frac{h^{-18}}R(1+T/\omega)^2
                         +\frac{\epsilon R^3}{\omega^2}\right].}
 \label{v14:eq:power-budget}
\end{equation}
\end{proposition}
\begin{proof}
On a one-site cutoff, $\ell_x\le R$ and
$\|P_RQ_x^2P_R\|\le2\sqrt{R/\lambda}$, by the form inequality
$\lambda Q_x^4/4\preceq\ell_x$ and Cauchy--Schwarz.
The onsite energy term consequently has norm at most $Cw(R+1)$, and each
edge term has norm at most $Cw h^{-2}\sqrt R$.
Summing $n$ onsite and $3n$ edge terms gives
\[
 \|E_R\|\le C(R+h^{-2}\sqrt R+1).
\]
Dividing the individual bounds by $\omega=\nu w$ gives
$g_{\rm site}\le C(R+1)/\nu$ and
$g_{\rm edge}\le Ch^{-2}\sqrt R/\nu$.
There are $M=4n$ contexts, so
\[
 L\le\frac C\omega(R+h^{-2}\sqrt R+1),\qquad
 M\sum_\alpha g_\alpha^2
       \le\frac C{\omega^2}(R^2+h^{-4}R+1).
\]
For $R\ge C\max(1,h^{-4})$ these prove the first two bounds, including
the $M^2$ term since $\omega=\nu h^3$ and $\nu\le1$.
Finally $b_h^2=O(h^{-7})$, $C_h^{\rm g}=O(h^{-7})$,
$C_h^{\rm f}=O(h^{-4})$, proving the displayed estimate for $D_h$.
Substitute in \eqref{v14:eq:budget}.
\end{proof}

\begin{corollary}[A fully finite branch with one-site and nearest-neighbor accepted operations]
\label[corollary]{v14:cor:joint}
On any fixed smooth scalar comparison slab, choose
\begin{equation}
 \boxed{
 \nu_h=h^3,\qquad R_h=h^{-32},\qquad
 \epsilon_h=h^{110},\qquad
 d_h=\frac{4\vartheta_h}{11}h^{110},\quad0<\vartheta_h\le1.
 }
 \label{v14:eq:scales}
\end{equation}
For all sufficiently small $h$, the local-dimension, projection, and
$M\epsilon\le\tau_0$ conditions hold and
$\mathfrak b_h=O(h^2)$.  The source has a finite tensor-product carrier, a
product initial preparation, local six-outcome instruments in its recorded
spatial contexts, and the unchanged two-state timing law.  Its actual scalar
readers converge in $H_x^1\times L_x^2$ uniformly in probability in time,
and its scalar metric source has no residual kinetic, gradient, or quartic
defect.  The expected maximal relative energy is
$O(h^2(1+|\log h|))$.
\end{corollary}
\begin{proof}
Here $\omega_h=h^6$.  The four terms in
\eqref{v14:eq:power-budget} have orders respectively
$h^2,h^4,h^{-18+32-12}=h^2$, and
$h^{110-96-12}=h^2$.  Also $M\epsilon=4h^{107}\to0$.
The higher Gaussian moment is uniformly bounded; hence the prerequisite
$D_hM_\phi^2/R\le1/4$ holds for small $h$.
All finite local spaces have dimension at least five once their one-site
cutoff exceeds the uniform trial-space bound noted above.
Apply \cref{v14:thm:scalar} and \eqref{v14:eq:log}; since
$A_R=O(h^{-32})$, the logarithm is $O(1+|\log h|)$.
\end{proof}

\paragraph{What has been removed, and what has not.}
The displayed source does not use a global interacting-energy projection,
a global accepted Hamiltonian exponential, or a nonlocal auxiliary Kraus
matrix.  The cutoff is a product of one-site cutoffs; accepted operators
have one-site or nearest-neighbor support.  The global Ritz map, exact
quantum trajectory, and rank-one reference projector are proof devices
only.  Their use neither conditions away outcomes nor changes an attained
record.  The quantitative Galerkin estimate makes the diagonal choice of
local cutoff explicit rather than hiding it behind an existence assertion.

The exponents are deliberately conservative and make no efficiency claim.
The context label and uniform spatial selection rule are \emph{new declared
source data}, not consequences of the original single locked instrument.
The physical preparation and semiclassical scale are likewise not selected
by the primitive law.  The result remains a real-Higgs matter branch on a
prescribed flat metric.  It supplies neither the Einstein constraints and
backreaction nor a cutoff-uniform physical light-cone bound for the random
local operations.  The latter would require additional control of their
rates and the relevant local observable norms; small operator support alone
is not that theorem.

\paragraph{Next nonduplicating target.}
For the independently specified native law the concrete missing comparison
is now between its \emph{actual local contexts and their operator pullbacks}
and the same-action local terms above.  If the original source does not
export those contexts, the present construction is only an existence
realization.  Within the coupled programme, the next analytic task is an
energy cancellation including dynamical gauge and geometric variables with
their constraint rows retained, not another scalar compactness criterion.
The scalar result should not be relabelled as a coupled Einstein solution
while its generally nonzero stress is placed on a fixed flat geometry.

\paragraph{Verification and append-only record.}
All bytes of Version 13 preceding its final end-document command are kept.
The diagnostic script checks support and nonmultiplicativity of tensor
compression, the Ritz commutator identity and energy-norm estimate on a
finite test system, exact contextual locking, local-channel remainder
scaling, the phase-corrected one-step identity, and the explicit exponent
budget.  The small oscillator test is itself an oscillator-basis numerical
approximation, not an exact diagonalization of the spatial Schr\"odinger
operator.  These are diagnostics of the proofs, not machine-checked
formalization, primitive-source verification, or an Einstein--SM simulation.
All new labels have prefix \texttt{v14:}.

\begingroup
\renewcommand{\refname}{Additional references for Version 14}
\begin{thebibliography}{9}
\small\raggedright
\bibitem{v14:Campbell}
E.~Campbell,
\emph{Random Compiler for Fast Hamiltonian Simulation},
Physical Review Letters \textbf{123} (2019), 070503.
\href{https://arxiv.org/abs/1811.08017}{arXiv:1811.08017};
\href{https://doi.org/10.1103/PhysRevLett.123.070503}{doi:10.1103/PhysRevLett.123.070503}.
Background for random selection of Hamiltonian terms.  The local locked
instrument, tensor-cutoff energy estimate, and non-Poisson resolved-clock
composition above are proved separately, with no efficiency comparison.
\end{thebibliography}
\endgroup
\addtocontents{toc}{\protect\endgroup}
% END IMMUTABLE VERSION 14 BODY
% APPEND VERSION 15 HERE, BEFORE THE FINAL END-DOCUMENT COMMAND.


% BEGIN IMMUTABLE VERSION 15 BODY
\clearpage
\begin{center}
{\Large\bfseries Version 15: Gauss constraints before a gauge-coupled instrument}\\[0.5em]
{\large Exact physical-sector preservation, relational canonical variables, and local locked gauge--Higgs records}
\end{center}
\makeatletter
\addtocontents{toc}{\protect\begingroup}
\addtocontents{toc}{\protect\makeatletter}
\addtocontents{toc}{\protect\def\protect\l@section{\protect\bfseries\protect\@dottedtocline{1}{0em}{2.5em}}}
\makeatother
\addcontentsline{toc}{part}{Version 15: Gauss-preserving local instruments and relational gauge--Higgs variables}

\section{Audit: a local scalar instrument is not yet a constrained gauge instrument}
\label{v15:sec:audit}

Version 14 supplies local finite scalar operations, but its arbitrary local
energy-eigenbasis matrix units were not required to commute with gauge
transformations.  Applying that construction unchanged to charged matter
would not establish preservation of Gauss's law on resolved records.
There is a second issue: an exactly neutral gauge state cannot have a
nonzero expectation of a charged Higgs coordinate.  Consequently the
fixed-frame scalar comparison observable cannot simply be renamed as a
physical comparison for a gauge-invariant preparation.  This body treats
these two issues before attempting another continuum energy estimate.

\paragraph{Inherited and external material.}
The finite Ward identities and classical subsidiary/Gauss equations in
\cite{GT,CC}, in particular \nolinkurl{thm:SMST-finite-action-Einstein} and
\nolinkurl{v2:thm:subsidiary}, are not reproved.  Versions 4--5 already
establish configuration-level Higgs dressing and retain its stabilizer and
full-source requirements.  The local locked normal form, its channel error,
and the clock comparison are
\cref{v14:prop:channel,v14:thm:trajectory}.  The new construction chooses
its auxiliary matrix units in a local gauge commutant rather than in an
unrestricted matrix algebra.  It also solves the \emph{momentum} Gauss
constraint and identifies the resulting electric--radial coupling.
Targeted searches and literal-term checks in the three supplied monographs
and the cumulative lineage did not locate this resolved-instrument and
relational-preparation composition.  This is a limited audit, not a
claim of exhaustive verification.

Gauge-invariant lattice Hamiltonians and physical Gauss sectors are
established subjects; see \cite{v15:ZCR,v15:Strocchi}.  The polar coordinate
reduction below is not claimed as a new general theory of gauge reduction.
Its purpose is to supply every term, local support, operator ordering,
preparation cost, and resolved-outcome condition needed in this lineage.
All finite assertions used in the new construction are proved below.

\paragraph{Exact scope.}
The explicit action is a compact Abelian gauge field coupled to one charged
complex scalar with the same quartic radial potential as the real-Higgs
restriction.  It is a declared diagnostic gauge--matter action on a fixed
finite spatial graph, not an identification of the complete Standard Model.
The abstract instrument lemma applies to any compact gauge group with the
stated local multiplicity.  There is no dynamical Einstein metric, no new
spatial-cutoff-uniform gauge--matter relative-energy theorem, and no claim
that the primitive source selects the constructed contexts.  In particular,
this body proves a finite constrained extension and an exact reader
reduction, not the full Einstein--SM continuum conclusion.

\section{The exact constraint test is on each resolved outcome}
\label{v15:sec:constraint}

Let $\mathcal K$ be a finite Hilbert space and let $\mathsf P$ be the
orthogonal projection onto the declared physical constraint subspace.  For
an actual instrument with Kraus operators $K_{e,a}$, retain both the
recorded label $e$ and any unobserved within-outcome Kraus index $a$ in the
following sums.  Completeness is $\sum_{e,a}K_{e,a}^*K_{e,a}=I$.
Put $\mathsf Q=I-\mathsf P$.

\begin{proposition}[Exact resolved leakage criterion]
\label[proposition]{v15:prop:leakage}
For every density matrix $\rho=\mathsf P\rho\mathsf P$,
\begin{equation}
 \tr\!\left(\mathsf Q\sum_{e,a}K_{e,a}\rho K_{e,a}^*\right)
 =\sum_{e,a}\|\mathsf QK_{e,a}\mathsf P\rho^{1/2}\|_{\rm HS}^2.
 \label{v15:eq:leakage}
\end{equation}
All positive-probability conditional outcomes preserve the physical
subspace for every physical input if and only if
\begin{equation}
 \mathsf QK_{e,a}\mathsf P=0\quad\text{for every }e,a.
 \label{v15:eq:branch-constraint}
\end{equation}
In particular, $[K_{e,a},U(g)]=0$ for every gauge transformation $U(g)$
is sufficient when $\mathsf P$ is the invariant-vector projection.
Preservation then holds on every finite resolved history, including on
histories of arbitrarily small positive probability.
\end{proposition}
\begin{proof}
Insert $\rho=\mathsf P\rho\mathsf P$ and use cyclicity of the trace.
Each summand equals the square in \eqref{v15:eq:leakage} and is
nonnegative.  If the sum vanishes for every rank-one physical input,
every $\mathsf QK_{e,a}\mathsf P$ vanishes.  Conversely these zeros
place every unnormalized outcome in $\Ran\mathsf P$; normalization
does not change its support.  Commutation with every $U(g)$ implies
commutation with its compact-group average $\mathsf P$.  Induction over
successive recorded outcomes proves the last assertion.  No projection
or postselection is performed after an outcome.
\end{proof}

\begin{example}[Covariance and a zero mean constraint do not suffice]
\label[example]{v15:ex:covariance}
On the charge basis $|{-}1\rangle,|0\rangle,|1\rangle$, put
$G=\operatorname{diag}(-1,0,1)$ and
\[
 K_+=2^{-1/2}|1\rangle\langle0|,\qquad
 K_-=2^{-1/2}|{-}1\rangle\langle0|,\qquad
 K_0=|{-}1\rangle\langle{-}1|+|1\rangle\langle1|.
\]
Their channel is trace preserving and covariant under $e^{\ii tG}$:
the first two Kraus operators acquire only phases and the third commutes.
The physical input $|0\rangle\langle0|$ is sent to the equal mixture of
charges $1$ and $-1$.  It has $\tr(G\rho')=0$, but
$\tr(G^2\rho')=1$, and every attained outcome violates $G\psi=0$.
This is a diagnostic of the distinction, not a proposed locked source.
\end{example}

\begin{proposition}[Small charge mean does not certify a physical preparation]
\label[proposition]{v15:prop:gap}
If commuting finite Gauss generators satisfy
$\sum_xG_x^2\succeq g_*^2(I-\mathsf P)$ with $g_*>0$, then
\[
 1-\tr(\mathsf P\rho)\le g_*^{-2}\sum_x\tr(\rho G_x^2).
\]
No corresponding bound follows just from the numbers $\tr(\rho G_x)$.
For the integer charge regulators below one can take $g_* =\hbar$ in
unscaled canonical charge units.  Thus that conversion has an explicit
$\hbar^{-2}$ cost; it is not a cutoff-uniform replacement for exact support.
\end{proposition}
\begin{proof}
Pair the stated positive operator inequality with $\rho$.  The preceding
example proves the failure of mean constraints.  In a simultaneous charge
basis each $G_x/\hbar$ is an integer, so every nonzero charge tuple has
squared length at least one.
\end{proof}

\section{A gauge-equivariant tensor regulator for the same local action}
\label{v15:sec:action}

Let $(\mathcal V,\mathcal E,\mathcal F)$ be a finite oriented graph with
chosen oriented plaquette cycles.  There are no external charges.
Write $B_{xe}=1$ at the tail of $e$, $-1$ at its head, and zero otherwise;
let $C_{fe}$ record oriented face incidence, so $BC^{\mathsf T}=0$.
One complex scalar is $z_x=q_{x,1}+\ii q_{x,2}$ and a compact link is
$u_e=e^{\ii\theta_e}$.  For $e:x\to y$, it transports $z_y$ to $x$.
Gauge transformations act by
\begin{equation}
 z_x\mapsto e^{\ii\beta_x}z_x,\qquad
 \theta\mapsto\theta+B^{\mathsf T}\beta.
 \label{v15:eq:gauge}
\end{equation}
The kinematical space is
\[
 \mathcal H_{\rm kin}
 =L^2\!\left(\C^{\mathcal V}\times\T^{\mathcal E},
      \prod_x\frac{\dd q_{x,1}\dd q_{x,2}}{2\pi}
      \prod_e\frac{\dd\theta_e}{2\pi}\right).
\]
Fix positive masses $m_x$, electric coefficients $c_e$, hopping weights
$k_e$, face coefficients $\beta_f^{\rm mag}$, and $\lambda>0$, $b\ge0$.
These are independently declared action coefficients.  Put
$\Pi_e=-\ii\hbar\partial_{\theta_e}$ and
\begin{equation}
\begin{split}
 E={}&\sum_x\left[-\frac{\hbar^2}{2m_x}\Delta_{q_x}
                     +m_x\frac\lambda4(|z_x|^2-b)^2\right]
       +\frac12\sum_ec_e\Pi_e^2\\
 &+\frac12\sum_{e:x\to y}k_e|e^{\ii\theta_e}z_y-z_x|^2
       +\sum_f\beta_f^{\rm mag}\{1-\cos(C\theta)_f\}.
 \label{v15:eq:energy}
\end{split}
\end{equation}
Use its nonnegative Friedrichs realization.  The physical action generator
is $E/\hbar$.  For comparison with Versions 13--14, setting $m_x=h^3$
and $\hbar=\nu h^3$ recovers their matter-density normalization after
writing the canonical Cartesian momentum as $h^3$ times the momentum
density.  No universal Planck parameter is derived here.

In polar coordinates $z_x=r_xe^{\ii\varphi_x}$ let
$L_x=-\ii\hbar\partial_{\varphi_x}$.  The Gauss generators and physical
space are
\begin{equation}
 G_x=L_x+(B\Pi)_x,\qquad
 \mathcal H_{\rm phys}=\bigcap_x\Ker G_x.
 \label{v15:eq:Gauss}
\end{equation}
These generators have the sign dictated by \eqref{v15:eq:gauge}.
They are the infinitesimal constraints, not merely mean charges.

\begin{theorem}[Local tensor cutoffs preserve the exact gauge action]
\label[theorem]{v15:thm:cutoff}
Set
\[
 \ell_x=-\frac{\hbar^2}{2m_x}\Delta_{q_x}
                        +\frac{m_x\lambda}{4}|z_x|^4+I,
 \qquad \ell_e=\frac{c_e}{2}\Pi_e^2+I,
 \qquad A=\sum_x\ell_x+\sum_e\ell_e,
\]
and define the product projection
\begin{equation}
 P_R=\bigotimes_x\one_{[0,R]}(\ell_x)
                     \otimes\bigotimes_e\one_{[0,R]}(\ell_e),
 \qquad E_R=P_REP_R|_{\Ran P_R}.
 \label{v15:eq:cutoff}
\end{equation}
At fixed graph and $\hbar>0$ its range is finite dimensional.  The projection
commutes with every gauge transformation and Gauss generator.  Each term
$P_RE_\alpha P_R$ of \eqref{v15:eq:energy} has the original bounded local
support and commutes separately with the full gauge group.  Thus
$E_R$ preserves $P_R\mathcal H_{\rm phys}$ exactly.

The physical products are compressed \emph{after} they are formed.  In
particular neither an exactly unitary compressed link nor an exact finite
Cartesian canonical commutation relation is asserted.  At fixed graph there
is a coefficient-dependent $D_\Gamma<\infty$, independent of
$0<\hbar\le1$, such that
\begin{equation}
 \|(I-P_R)v\|_{E+I}^2
 \le\frac{D_\Gamma}{R}\|(E+I)v\|^2,
 \qquad v\in D(E).
 \label{v15:eq:tail}
\end{equation}
The tensor spaces are a form core as $R\to\infty$ at fixed graph and
$\hbar$.  No spatial-refinement bound on $D_\Gamma$ is asserted.
\end{theorem}
\begin{proof}
The rotationally invariant two-dimensional quartic oscillator has compact
resolvent: a bounded energy-form family has small tails at large $|z|$
and is $H^1$ bounded on each fixed ball, giving compact $L^2$ embedding.
The link cutoff retains finitely many integer Fourier modes.  The local
oscillators commute with scalar rotations and the link kinetic operators
with both endpoint translations.  Their spectral projections consequently
commute with the product gauge action.  The hopping norm and the face
cosine are gauge invariant, since $BC^{\mathsf T}=0$; the other terms
are invariant individually.  Compression by the tensor projection preserves
support and this commutation.  This also proves the Gauss assertion.

For the estimate, expand $E=A+W_2+W_0+cI$, where $W_0$ is a bounded face
potential, and $W_2$ consists of the hopping quadratics and the negative
quadratic part of the radial quartic.  As multiplication functions,
$|W_2|\le C_\Gamma\sum_xr_x^2$ uniformly in the link angles.  Hence
\[
 \|W_2v\|\le C_\Gamma\|A^{1/2}v\|,
 \qquad A\succeq I.
\]
The proof of \cref{v14:lem:form}, with these constants, gives
$D(E)=D(A)$, $\|Av\|\le C_\Gamma\|(E+I)v\|$, and
$E+I\preceq C_\Gamma A$ as forms.  To make the domain step explicit,
$\|(W_2+W_0+c)(A\mp\ii t)^{-1}\|\to0$ as $t\to\infty$,
so the symmetric perturbation on $D(A)$ has both resolvents and agrees
with the Friedrichs form operator.  Constants in these estimates involve
only the fixed coefficients and number of factors, not $\hbar$.
On the complement of $P_R$ at least one local oscillator eigenvalue
exceeds $R$.  The joint spectral theorem gives
$\|A^{1/2}(I-P_R)v\|^2\le R^{-1}\|Av\|^2$.
The form and graph comparisons prove \eqref{v15:eq:tail}; the same joint
spectral expansion proves the form-core claim.  The support statement does
not turn compression into an algebra homomorphism.
\end{proof}

\section{Exact Gauss reduction exposes a local electric--radial term}
\label{v15:sec:reduction}

Define the invariant link angles, modulo $2\pi$, by
\begin{equation}
 \alpha_e=\theta_e+\varphi_y-\varphi_x
                 =\theta_e-(B^{\mathsf T}\varphi)_e.
 \label{v15:eq:dressed}
\end{equation}
No principal argument is chosen; these are circle-valued variables.  All
edge angles, including loop and noncontractible period information, remain.
For the radial measure use $\prod_xr_x\dd r_x$.

\begin{theorem}[Unitary physical-coordinate reduction and its complete Hamiltonian]
\label[theorem]{v15:thm:reduction}
The map
\begin{equation}
 (\mathcal U f)(r,\varphi,\theta)
       =f(r,\theta-B^{\mathsf T}\varphi)
 \label{v15:eq:isometry}
\end{equation}
is a unitary from
$\mathcal H_{\rm rel}=L^2((0,\infty)^{\mathcal V}\times\T^{\mathcal E},
\prod_xr_x\dd r_x\prod_e\dd\alpha_e/(2\pi))$
onto $\mathcal H_{\rm phys}$.  On its natural Friedrichs form domain,
$E_{\rm rel}=\mathcal U^*E\mathcal U$ is
\begin{equation}
\boxed{\begin{aligned}
 E_{\rm rel}={}&\sum_x\left[-\frac{\hbar^2}{2m_x}
           (\partial_{r_x}^2+r_x^{-1}\partial_{r_x})
          +\frac{(B\Pi)_x^2}{2m_xr_x^2}+m_xV(r_x)\right]
          +\frac12\sum_ec_e\Pi_e^2\\
 &+\frac12\sum_{e:x\to y}k_e
       (r_x^2+r_y^2-2r_xr_y\cos\alpha_e)
          +\sum_f\beta_f^{\rm mag}\{1-\cos(C\alpha)_f\},
\end{aligned}}
 \label{v15:eq:reduced-energy}
\end{equation}
where $V(r)=\lambda(r^2-b)^2/4$ and now
$\Pi_e=-\ii\hbar\partial_{\alpha_e}$.
The charge has been solved exactly as $L_x=-(B\Pi)_x$, not set to zero
separately.  The additional electric term has support on the star of $x$.
It is positive and is not an optional gauge-fixing penalty.
\end{theorem}
\begin{proof}
At each fixed $\varphi$, translation from $\theta$ to $\alpha$ preserves
normalized Haar measure.  Integrating over $\varphi$ proves isometry.
Gauge transformations translate $\varphi$ arbitrarily while leaving
$r,\alpha$ fixed, so their invariant $L^2$ functions are exactly the
functions in \eqref{v15:eq:isometry}.  This proves surjectivity, also
by the Fourier expansion in $\varphi$ after the measure-preserving change
of variables.  The set where some $r_x=0$ has measure zero and does not
affect this Hilbert-space statement.

For a smooth function supported where all radii are positive,
\[
 L_x\mathcal U f=-\mathcal U(B\Pi)_xf,\qquad
 \Pi_e\mathcal U f=\mathcal U\Pi_ef.
\]
The polar Laplacian is
$\partial_r^2+r^{-1}\partial_r+r^{-2}\partial_\varphi^2$;
its angular term therefore becomes $(B\Pi)_x^2/(2m_xr_x^2)$.
The hopping norm is the displayed cosine expression, and
$C\theta=C\alpha$ modulo $2\pi$ by $CB^{\mathsf T}=0$.
These calculations prove equality of the forms on this core.

Here the singular coefficient does not license a new boundary condition.
One may define the reduced form as the unitary image of the original
closed physical form.  Equivalently, the indicated punctured smooth core
is a form core: for a smooth compactly supported Cartesian test, remove
$r_x<\delta^2$ using a radial cutoff that changes logarithmically to one
on $[\delta^2,\delta]$.  Its two-dimensional Dirichlet energy is
$O(1/|\log\delta|)$.  Boundedness of the test and compact support control
the corresponding error in the other terms.  A finite product of such
cutoffs removes all zero radii.  Gauge averaging, which commutes with these
cutoffs and is contractive for the invariant energy form, gives the same
core in the physical subspace.  Closing the identical nonnegative forms
proves the operator claim.
\end{proof}

\begin{proposition}[The same classical action retains the complete charge feedback]
\label[proposition]{v15:prop:classical}
On the nonzero-radius classical chart, let $p_x$ be radial canonical
momentum, $L_x$ angular matter momentum, and $\pi_e$ link momentum.
The canonical one-form has the exact identity
\begin{equation}
 \sum_x(p_x\dd r_x+L_x\dd\varphi_x)+\sum_e\pi_e\dd\theta_e
 =\sum_xp_x\dd r_x+\sum_e\pi_e\dd\alpha_e
                       +\sum_x\{L_x+(B\pi)_x\}\dd\varphi_x.
 \label{v15:eq:one-form}
\end{equation}
Thus $G=0$ gives genuine canonical pairs $(r,p),(\alpha,\pi)$, and the
classical Hamiltonian is the symbol of \eqref{v15:eq:reduced-energy},
including $\sum_x(B\pi)_x^2/(2m_xr_x^2)$.  Writing
$s_x=(B\pi)_x$ and $a_x=s_x/(m_xr_x^2)$, its equations are
\begin{align}
 \dot r_x&=p_x/m_x,\nonumber\\
 \dot p_x&=s_x^2/(m_xr_x^3)-m_xV'(r_x)
                 -\sum_{e\ni x}k_e(r_x-r_{y(e,x)}\cos\alpha_e),\nonumber\\
 \dot\alpha_e&=c_e\pi_e+(B^{\mathsf T}a)_e,\nonumber\\
 \dot\pi_e&=-k_e r_xr_y\sin\alpha_e
                 -\sum_f C_{fe}\beta_f^{\rm mag}\sin(C\alpha)_f.
 \label{v15:eq:classical-flow}
\end{align}
They preserve Gauss's law upon reconstruction and include dynamical
matter--electric and magnetic interactions.  No Einstein equation follows
from this gauge--matter reduction.
\end{proposition}
\begin{proof}
Substitute $\theta=\alpha+B^{\mathsf T}\varphi$ into the original
one-form, giving \eqref{v15:eq:one-form}.  Substitute
$L=-B\pi$ into the Cartesian kinetic energy in polar form.  Hamilton
differentiation gives all four equations.  In the unreduced coordinates
$\dot L=B\,\partial_\alpha E$ and
$\dot\pi=-\partial_\alpha E$, so
$\dot L+B\dot\pi=0$ identically.  The face contribution has zero
incidence by $BC^{\mathsf T}=0$.  This verifies the reconstructed
constraint without removing either electric or matter charge.
\end{proof}

\begin{remark}[Radial measure and action derivatives are retained]
The quantum radial kinetic term acts in $r\dd r$.  If that measure is
changed to $\dd r$ by $f\mapsto r^{1/2}f$, it becomes
$-\hbar^2\partial_r^2/(2m)-\hbar^2/(8mr^2)$.
The last term must not be silently dropped.  Direct differentiation of
$r^{-1/2}u$ proves the formula on the same punctured core.  It is an
ordering/measure statement, not a negative contribution inserted into the
original nonnegative Cartesian energy.

If an external source parameter varies the displayed action coefficients
but not the gauge action or coordinate map $\mathcal U$, then
$\partial_\lambda E_{\rm rel}
 =\mathcal U^*(\partial_\lambda E)\mathcal U$ as forms on the common
core.  The analogous identity holds after any fixed regulator projection.
In particular the derivative of the electric--radial term remains in
mass/lapse coefficient variations.  Differentiating a parameter-dependent
cutoff would be a different convention and requires its extra terms.
\end{remark}

\begin{example}[The omitted term can be order one on exact Gauss data]
\label[example]{v15:ex:charge-feedback}
For one oriented edge $x\to y$, take any radii $r_x,r_y>0$ and electric
momentum $\pi_e=p\ne0$.  Exact Gauss data have $L_x=-p,L_y=p$.
Even with dressed link $\alpha_e=0$ and no magnetic curvature, their
angular matter energy is
\[
 \frac{p^2}{2}\left(\frac1{m_xr_x^2}+\frac1{m_yr_y^2}\right).
\]
The corresponding radial forces are $p^2/(m_xr_x^3)$ and
$p^2/(m_yr_y^3)$.  A quantum link Fourier mode with
$\hbar k\to p$ realizes the same limiting contribution on smooth
radial preparations.  Therefore dropping this term is not justified by
zero curvature or by imposing Gauss's law.  It changes both the action and
its physical source derivatives at leading order.
\end{example}

\section{Physical preparations must concentrate relational readers}
\label{v15:sec:preparation}

\begin{proposition}[A charged mean is the wrong comparison on the exact physical sector]
\label[proposition]{v15:prop:charged-mean}
For any physical density matrix with finite second position moments,
$\tr(\rho z_x)=0$.  Hence for every prescribed complex number $z_x^*$,
\begin{equation}
 \tr\rho\,|z_x-z_x^*I|^2
      =\tr\rho\,|z_x|^2+|z_x^*|^2\ge |z_x^*|^2.
 \label{v15:eq:charged-floor}
\end{equation}
This also holds in the gauge-equivariant tensor regulator, with products
formed before compression.  It does not prohibit nonzero Higgs radius,
nonzero invariant hopping expectation, or a classical limit modulo gauge.
\end{proposition}
\begin{proof}
A physical state is unchanged by every site rotation, whereas $z_x$
acquires its nontrivial character.  Its expectation must therefore vanish.
For unbounded $z_x$, first use bounded radial truncations, then pass by
Cauchy--Schwarz and the second moment.  Expand the square.  The cutoff
commutes with the rotations and preserves this character identity, giving
the regulated version.
\end{proof}

This is not an error in the neutral real-scalar theorem.  It shows why its
fixed-frame comparison cannot be used unchanged on charged singlet states.
Nor does gauge averaging an arbitrary density matrix project it onto
$G=0$: an incoherent mixture of nonzero charge sectors is already invariant
as a density matrix.  The preparation below supplies physical \emph{vector}
support, not merely invariance under conjugation.

Fix initial relational data $(r_x^0,p_x^0,\alpha_e^0,\pi_e^0)$ on the
finite graph, with $0<r_*\le r_x^0\le r^*$ and bounded momenta.  Choose
smooth radial cutoffs $\chi_x$ supported in a fixed positive-radius
interval and equal to one near $r_x^0$.  For $0<\hbar\le1$ let
\begin{align}
 f_x(r)&=Z_x^{-1}r^{-1/2}\chi_x(r)
   \exp\left[-\frac{(r-r_x^0)^2}{2\hbar}
                      +\frac{\ii p_x^0(r-r_x^0)}{\hbar}\right],\nonumber\\
 g_e(\alpha)&=Z_e^{-1}e^{\ii k_e^0(\alpha-\alpha_e^0)}
    \sum_{n\in\Z}\exp\left[-\frac{(\alpha-\alpha_e^0-2\pi n)^2}{2\hbar}\right],
 \qquad k_e^0=\operatorname{round}(\pi_e^0/\hbar),
 \label{v15:eq:preparation}
\end{align}
where normalization uses $r\dd r$ and normalized Haar measure respectively.
Let $f^0=\prod_xf_x\prod_eg_e$ and $\psi^0=\mathcal U f^0$.
The integer $k_e^0$ in this formula is not the hopping weight $k_e$.

Define a positive relational comparison form at time zero by
\begin{equation}
\begin{split}
 \mathcal C_0={}&\sum_x\left[(r_x-r_x^0)^2
       +(-\ii\hbar\partial_{r_x}-p_x^0)^*
                    (-\ii\hbar\partial_{r_x}-p_x^0)\right]\\
 &+\sum_e\left[2\{1-\cos(\alpha_e-\alpha_e^0)\}
                              +(\Pi_e-\pi_e^0)^2\right].
 \label{v15:eq:initial-cost}
\end{split}
\end{equation}
The radial derivative is defined through this positive form in $r\dd r$;
no self-adjoint half-line momentum is assumed.

\begin{theorem}[Nonempty exactly constrained finite preparations]
\label[theorem]{v15:thm:preparation}
The vector $\psi^0$ satisfies every $G_x\psi^0=0$ exactly.  Its
relational moments obey, with fixed-graph constants,
\begin{equation}
 \langle f^0,\mathcal C_0f^0\rangle\le C_\Gamma\hbar,
 \qquad \|(E+I)\psi^0\|\le C_\Gamma.
 \label{v15:eq:initial-bounds}
\end{equation}
For the cutoff of \cref{v15:thm:cutoff}, define
$\psi_R^0=P_R\psi^0/\|P_R\psi^0\|$ when $R$ is sufficiently large.
This vector is finite, physical, and has
\begin{equation}
 \left\langle\psi_R^0,
     P_R\mathcal U\mathcal C_0\mathcal U^*P_R\psi_R^0\right\rangle
                  \le C_\Gamma(\hbar+R^{-1}).
 \label{v15:eq:finite-preparation}
\end{equation}
The physical source has neither a posterior Gauss projection nor a
postselection of outcomes.  The uncompressed preparation is a product in
relational variables, generally \emph{not} a product in the original site
and link factors.  All displayed initial data are used before evolution.
\end{theorem}
\begin{proof}
Physical support is \cref{v15:thm:reduction}.  Under $r\dd r$ the factor
$r^{-1/2}$ reduces radial position expectations to those of a cut-off
one-dimensional Gaussian.  Gaussian tails outside the region where
$\chi_x=1$ are exponentially small in $1/\hbar$; their polynomial
inverse powers from differentiation remain exponentially small.  Thus
$\langle(r-r^0)^2\rangle\le C\hbar$ and
\[
 \|(-\ii\hbar\partial_r-p^0)f_x\|^2
 =\left\|\left(\ii(r-r^0)+\frac{\ii\hbar}{2r}
                       -\ii\hbar\chi_x'/\chi_x\right)f_x\right\|^2
 \le C\hbar.
\]
The last expression is interpreted after multiplying by $\chi_x$; there
is no singular division at a zero of that cutoff.  For a link, the
periodized real Gaussian envelope has exponentially small overlap between
its distinct images near its chosen center.  Gaussian integration therefore
gives $\langle1-\cos(\alpha-\alpha^0)\rangle\le C\hbar$ and
$\|(-\ii\hbar\partial_\alpha-\hbar k_e^0)g_e\|^2\le C\hbar$.
Its momentum mean is exactly $\hbar k_e^0$, since the envelope is real
and periodic; the rounding error is at most $\hbar/2$.  Summing proves
the first estimate.

The explicit reduced Hamiltonian on $f^0$ contains at most two scaled
radial or angular derivatives.  Their norms are bounded uniformly in
$\hbar$ by the same Gaussian calculation (fourth Gaussian moments for
a second derivative).  All inverse radii and potential coefficients are
bounded on the radial supports.  The star derivatives are fixed finite
sums.  Consequently $\|E_{\rm rel}f^0\|\le C_\Gamma$, which proves
the second estimate by unitary equivalence.  This argument makes no
spatial-cutoff-uniform claim about the number or coefficients of those sums.

The positive form $\mathcal U\mathcal C_0\mathcal U^*$ on the physical
subspace is bounded above by $C_\Gamma(E+I)$: radial derivative energy is
part of the Cartesian kinetic energy, link momenta have positive electric
coefficients, $r^2$ is bounded by the quartic potential plus a constant,
and the cosine terms are bounded.  Apply \eqref{v15:eq:tail} to
$\psi^0$.  The norm of $(I-P_R)\psi^0$ is $O(R^{-1/2})$, and its
comparison-form norm is of the same order.  The triangle inequality in
that form, followed by normalization $\|P_R\psi^0\|\ge1/2$, gives
\eqref{v15:eq:finite-preparation}.  Commutation of $P_R$ with the gauge
group proves exact physical support after the cutoff.  No future
classical trajectory was used.
\end{proof}

\section{Local locked outcomes can be chosen inside the gauge commutant}
\label{v15:sec:instrument}

The following finite algebraic lemma permits non-Abelian gauge groups as
well, but does not construct their physical preparation or continuum limit.
For a local context $\alpha$ let the compact product gauge group on its
incident vertices act on its support by $U_\alpha(g)$.  Its representation
decomposes as
\[
 \mathcal K_\alpha=\bigoplus_\lambda V_\lambda\otimes M_\lambda,
 \qquad U_\alpha(g)=\bigoplus_\lambda U_\lambda(g)\otimes I.
\]
Assume the local action generator commutes with this representation, so
$H_\alpha=\bigoplus_\lambda I\otimes h_{\alpha,\lambda}$.
Suppose some multiplicity $M_{\lambda_0}$ has dimension at least five.
Choose five eigenvectors $u_0,\ldots,u_4$ of $h_{\alpha,\lambda_0}$ and
set, extended by zero to other isotypic blocks,
\[
 C_{\alpha,\mu}=I_{V_{\lambda_0}}\otimes|u_0\rangle\langle u_{\mu-1}|,
 \quad 2\le\mu\le5,\qquad
 R_\alpha=\sum_{\mu=2}^5 C_{\alpha,\mu}^*C_{\alpha,\mu}.
\]
For $0<\tau<1$ use
\begin{equation}
\begin{aligned}
 B_{\alpha,1}&=e^{-3\ii\tau H_\alpha/2}(I-\tau^4R_\alpha)^{1/2},
 &B_{\alpha,\mu}&=\tau^2C_{\alpha,\mu},\\
 K_{\alpha,e}&=I/\sqrt{18}+\sqrt{2/3}\sum_{\mu=1}^5O_{e\mu}B_{\alpha,\mu},
 &1&\le e\le6.
\end{aligned}
 \label{v15:eq:locked}
\end{equation}
where the fixed real isometry $O$ has column one
$(1,1,1,-1,-1,-1)^{\mathsf T}/\sqrt6$, as in Version 12.

\begin{theorem}[Exact local locking and outcome-by-outcome Gauss preservation]
\label[theorem]{v15:thm:locked}
Each $K_{\alpha,e}$ is supported on the declared local context and commutes
with every local gauge transformation.  It satisfies
\[
 \sum_eK_{\alpha,e}^*K_{\alpha,e}=I,\qquad
 \sum_eK_{\alpha,e}=\sqrt2 I.
\]
All six are linearly independent for $0<\tau<1$ and have a common positive
least singular-value lower bound for sufficiently small fixed $\tau$,
independently of $\|H_\alpha\|$ and the dimensions.  Consequently they
preserve the global physical sector on every resolved outcome.  Their
adjoint channel has exactly the estimates of
\cref{v12:thm:locked-bank} for the local $H_\alpha$.
\end{theorem}
\begin{proof}
All $C_{\alpha,\mu}$ are intertwiners on one isotypic block.  The operator
$R_\alpha$ is an orthogonal projection commuting with $H_\alpha$ and
with the gauge group.  Every factor in \eqref{v15:eq:locked} therefore
commutes with the gauge group.  No commutation of $C_{\alpha,\mu}$ with
$H_\alpha$ is used.  Local gauge commutation implies commutation with the
full gauge action after tensoring with outside identities.
The identity $\sum_{\mu=1}^5B_{\alpha,\mu}^*B_{\alpha,\mu}=I$
and the orthogonality of $O$ give the two locking equations.

Inside the selected isotypic block the four $C_{\alpha,\mu}$ have distinct
off-diagonal multiplicity entries, while $I$ and $B_{\alpha,1}$ are
diagonal in a common eigenbasis.  The latter has modulus
$\sqrt{1-\tau^4}$ on the $u_1,\ldots,u_4$ entries and modulus one on
$u_0$; it cannot be scalar.  Thus these six columns are independent, and
the fixed invertible normal-form column change gives independence of
the $K_{\alpha,e}$.  Moreover
\[
 K_{\alpha,e}=I/\sqrt{18}\pm B_{\alpha,1}/3+O_{\rm op}(\tau^2),
\]
so its least singular value is at least
$\sqrt{1-\tau^4}/3-1/\sqrt{18}-C\tau^2$, positive for sufficiently
small $\tau$.  This argument does not depend on the Hamiltonian phases.
The conjugation estimate uses only $\|R_\alpha\|=\|C_{\alpha,\mu}\|=1$
and is exactly the integration proof of \cref{v12:thm:locked-bank}.
Finally apply \cref{v15:prop:leakage}; there is no later repair of an
unphysical outcome.
\end{proof}

\begin{proposition}[The required multiplicity is present in the charged scalar regulator]
\label[proposition]{v15:prop:multiplicity}
On a bounded-degree graph with bounded-size plaquettes, all local terms of
\eqref{v15:eq:energy} admit the preceding construction for sufficiently
large $R$.  If necessary, include one incident matter site in an electric-link
or plaquette context on which its Hamiltonian acts as the identity.
The local operator support then remains uniformly bounded in graph radius.
No extra matter species or gauge-breaking reference vector is added.
\end{proposition}
\begin{proof}
The site oscillator has a neutral angular-momentum-zero radial subspace.
Choose five smooth radial test functions with disjoint compact supports away
from zero.  Their Rayleigh quotients are bounded uniformly for
$0<\hbar\le1$ and fixed coefficients, so its neutral spectral subspace
has dimension at least five once $R$ exceeds a fixed-graph threshold.
The link cutoff contains the zero-flux vector.  In each enlarged local
context, take those five neutral radial states, one fixed neutral state at
other included sites, and zero flux on every included link.  They lie in
one trivial representation of all incident local gauges.  The context's
invariant subspace consequently has multiplicity at least five.
Its action term preserves that whole subspace and can be diagonalized there;
its eigenvectors need not be the five trial products.  Adding an identity
factor on an existing incident matter site does not change the action term.
The resulting Kraus correction may act there and is openly part of the
constructed local instrument.  No distant degree of freedom is used.
\end{proof}

\section{A finite constrained gauge--matter source on the actual clock}
\label{v15:sec:clock}

Write $E_R/\hbar=\sum_{\alpha=1}^{M}H_\alpha$ using the gauge-invariant
site, electric-link, hopping, and plaquette terms just constructed.  At an
accepted completion choose a context uniformly, record it, and then record
one of its six original outcomes with $\tau=M\epsilon$.  This spatial
choice is declared just as in Version 14, not inferred from the original
single-context opportunity record.  The accepted global Kraus operators
are $M^{-1/2}K_{\alpha,e}$.  Use the unchanged two-phase clock and
$d=4\vartheta\epsilon/11$, with arbitrary initial clock phase.

Set $g_\alpha=\|H_\alpha\|$, $L=\sum_\alpha g_\alpha$, and
\begin{equation}
 \mathfrak C_\Gamma=1+L^2+M\sum_\alpha g_\alpha^2+M^2.
 \label{v15:eq:clock-cost}
\end{equation}
All norms use the actually regulated local matrices.

\begin{theorem}[Exactly constrained resolved gauge--matter evolution]
\label[theorem]{v15:thm:source}
Start from $\rho_0=|\psi_R^0\rangle\langle\psi_R^0|$ of
\cref{v15:thm:preparation}.  Every attained conditional state satisfies
\begin{equation}
 \rho_j=\mathsf P\rho_j\mathsf P,\qquad
 \tr(\rho_jG_x^2)=0\quad\text{for every }x,j
 \label{v15:eq:exact-constraint}
\end{equation}
exactly.  Let $\psi_R(t)=e^{-\ii tE_R/\hbar}\psi_R^0$ and
$\xi_j=1-\langle\psi_R(jd),\rho_j\psi_R(jd)\rangle$.
For $M\epsilon\le\tau_0$,
\begin{equation}
 \max_{j\le T/d}\E\xi_j\le C_T\epsilon\mathfrak C_\Gamma,
 \qquad
 \PP\{\max_{j\le T/d}\xi_j>u\}
       \le C_T\epsilon\mathfrak C_\Gamma/u.
 \label{v15:eq:trajectory}
\end{equation}
The constant is independent of $\vartheta$, matrix dimensions, and
Hamiltonian norms; the latter costs are in $\mathfrak C_\Gamma$.
No constraint-violation term or projection-success probability is added.

For any fixed regulated self-adjoint physical source insertion $O_R$,
\begin{equation}
 |\tr(\rho_jO_R)-\langle\psi_R(jd),O_R\psi_R(jd)\rangle|
               \le2\|O_R\|\sqrt{\xi_j}.
 \label{v15:eq:insertion}
\end{equation}
Thus $\epsilon\mathfrak C_\Gamma\|O_R\|^2\to0$ is sufficient for
its resolved-versus-quantum convergence in probability uniformly on a fixed
time interval.  This includes derivatives of the same regulated action
with respect to fixed external coefficients; it is not yet a
quantum-to-classical or Einstein-stress convergence theorem.
\end{theorem}
\begin{proof}
The finite preparation is physical and all Kraus operators commute with
the global gauge group.  \Cref{v15:prop:leakage} proves
\eqref{v15:eq:exact-constraint} by induction; rejected ticks act as the
identity.  Every term of $E_R$ also commutes with the group, so the
reference remains physical.
The local channel estimates and their averaging are exactly those of
\cref{v14:prop:channel}, now verified with gauge-commuting auxiliary
operators.  Their scale budget is \eqref{v15:eq:clock-cost}.  Apply
\cref{v14:thm:trajectory}; its proof is a pure-reference positive-observable
calculation for an arbitrary finite local Hamiltonian and does not use the
real-scalar field equation.  This gives \eqref{v15:eq:trajectory} with
no new proof of the clock Poisson equation or maximal principle.

For completeness the trace-norm estimate against a pure reference follows
by decomposing $\rho$ into pure states: the trace distance of two pure
projectors is $2\sqrt{1-|\langle u,v\rangle|^2}$, as is seen on their
two-dimensional span.  Convexity of trace norm and concavity of the square
root give $\|\rho-|\psi\rangle\langle\psi|\|_1\le2\sqrt{\xi}$.
Pair with $O_R$ to obtain \eqref{v15:eq:insertion}.  The event that its
maximum exceeds $u$ forces $\max\xi_j>u^2/(4\|O_R\|^2)$, proving
the asserted sufficient condition.  All recorded histories, not just
favorable constraints or a nonselective mean state, remain in the result.
\end{proof}

\section{Outcome and the next coupled estimate}
\label{v15:sec:ledger}

The new finite branch has dynamical compact Abelian links, charged scalar
matter, a gauge-equivariant tensor cutoff, nonempty exactly physical
preparations, and local locked operations preserving every Gauss constraint
on every resolved record.  The actual renewal-clock comparison continues to
hold.  Its constraint control is exact, not another small-error assumption.
There is no claim that a gauge-covariant averaged channel alone suffices.

The reader correction is equally consequential.  Raw charged means vanish
on the physical sector, but the Higgs radius and the dressed link variables
retain physical information.  Exact momentum reduction produces the local
term
\[
 \boxed{\sum_x\frac{(B\Pi)_x^2}{2m_xr_x^2},}
\]
which couples electric divergence to the Higgs radius and contributes to
action-source derivatives.  Setting angular matter momentum to zero would
remove actual charged configurations rather than solve their Gauss law.
Higgs zeros do not obstruct the Hilbert-space isometry, but they do make
inverse-radius estimates important for any subsequent classical chart
estimate.  No all-time inverse-radius bound has been derived here.

The classical gauge--matter equations in \eqref{v15:eq:classical-flow} are
an exact reduction of a fixed action, not the result of fitting a decoder to
a recorded curve.  Nevertheless, the present initial concentration estimate
has constants for a fixed graph, and \eqref{v15:eq:trajectory} compares
the instrument to a finite quantum gauge--matter flow.  Combining those two
facts alone does \emph{not} prove a spatially uniform semiclassical limit.
The scalar relative-energy inequality of Version 13 does not already
control the new inverse-radius electric energy, magnetic terms, and their
commutators.  That joint estimate, including an actual propagated
nonvanishing-radius margin or another treatment of zeros, is now the precise
analytic target.  Its source derivatives must be differentiated before
compression and before discarding any constraint term.

Gravity, non-Abelian color dynamics, chiral fermions, physical context
selection, semiclassical normalization, and native action identification
remain separate.  The abstract local commutant lemma is applicable to other
compact gauge groups when its representation multiplicity is supplied, but
is not a proof of those missing sectors.  The main advance is a
\emph{constraint-consistent finite gauge--matter extension and exact
relational starting point}, not universal microscopic Einstein--SM emergence.

\paragraph{Verification and append-only record.}
The full Version 14 prefix preceding its final document command is
preserved byte-for-byte.  New labels have prefix \texttt{v15:}.
The companion diagnostics verify the Gauss-incidence signs, the canonical
one-form, Hamilton equations and charge feedback, polar kinetic reduction,
gauge-equivariant cutoff blocks, exact contextual locking, conditional
Gauss preservation, and a finite gauge-invariant preparation.  They also
check the covariance-only counterexample.  These are finite symbolic and
numerical checks of the written arguments, not machine-checked proofs,
a primitive-source identification, or an Einstein--SM simulation.

\begingroup
\renewcommand{\refname}{Additional references for Version 15}
\begin{thebibliography}{9}
\small\raggedright
\bibitem{v15:ZCR}
E.~Zohar, J.~I.~Cirac, and B.~Reznik,
\emph{Quantum simulations of gauge theories with ultracold atoms:
local gauge invariance from angular momentum conservation}, 2013.
\href{https://arxiv.org/abs/1303.5040}{arXiv:1303.5040}.
Background for exact local gauge-invariant interactions; its microscopic
implementation is not substituted for the declared locked instrument.
\bibitem{v15:Strocchi}
F.~Strocchi, \emph{Local Gauss law and local gauge symmetries in QFT}, 2023.
\href{https://arxiv.org/abs/2301.06543}{arXiv:2301.06543}.
Background for distinguishing gauge-invariant observables, physical states,
and local Gauss constraints.  All finite formulas in this body are proved
in their declared representation and charge sector.
\end{thebibliography}
\endgroup
\addtocontents{toc}{\protect\endgroup}
% END IMMUTABLE VERSION 15 BODY
% APPEND VERSION 16 HERE, BEFORE THE FINAL END-DOCUMENT COMMAND.


% BEGIN IMMUTABLE VERSION 16 BODY
\clearpage
\begin{center}
{\Large\bfseries Version 16: A reference amplitude before an inverse-radius bound}\\[0.5em]
{\large Zero-safe gauge--Higgs comparison, an exact quantum dilation identity, and a local constrained dynamical branch}
\end{center}
\makeatletter
\addtocontents{toc}{\protect\begingroup}
\addtocontents{toc}{\protect\makeatletter}
\addtocontents{toc}{\protect\def\protect\l@section{\protect\bfseries\protect\@dottedtocline{1}{0em}{2.5em}}}
\makeatother
\addcontentsline{toc}{part}{Version 16: Zero-safe coupled comparison and a local physical dynamical branch}

\section{Audit: change the comparison, not the physical source}
\label{v16:sec:audit}

Version 15 solves Gauss's law on every resolved record, but its polar
Hamiltonian contains an inverse-radius electric term. A bound on the inverse
of every attained Higgs radius would be a strong additional requirement.
It is not necessary for the comparisons below. The reference amplitude,
not the fluctuating amplitude, is divided out. All microscopic configurations,
including Higgs zeros, remain in the source and in the comparison energy.

\paragraph{Inherited results and new scope.}
We retain \cref{v15:thm:reduction,v15:thm:cutoff,v15:thm:locked,v15:thm:source}:
the exact physical-coordinate reduction, the equivariant tensor regulator,
the local gauge-commuting locked outcomes, and their resolved-clock comparison.
The clock calculation is not repeated. The tensor Galerkin estimate is
\cref{v14:thm:Galerkin}; only its action-specific constants and preparation
are evaluated here. The neutral scalar cancellation of
\cref{v13:thm:cancellation} does not contain the comparison constructed below.
Targeted searches for relative or modulated energy and reference-amplitude
arguments in \cite{GT,CC} located no duplicate of this argument. This is a
limited nonduplication audit, not a verification of all supplied manuscripts.

There are two separate positive statements. First, a continuum comparison
identity holds between general smooth Abelian gauge--Higgs solutions, with
no lower bound on the candidate amplitude. Second, an exact quantum identity
holds for the literal compact-link lattice action of Version 15 when the
reference is a spatially homogeneous, nonconstant radial solution with zero
reference gauge field. This second identity gives a local finite,
Gauss-preserving quantum-to-classical branch with explicit joint scales.
We do not claim the general continuum identity to be an already proved
quantum identity for inhomogeneous references.

\paragraph{Attribution and physical boundary.}
Modulated energies for Klein--Gordon--Maxwell systems have precedents, for
example \cite{v16:Salvi}, whose limit is to relativistic Euler--Maxwell.
That result is not used as a substitute for any estimate here. Our continuum
comparison has the same quartic field equation on both sides, and the quantum
result has its own finite compact-link action and resolved renewal clock.
Neither the action coefficients, the preparation, the small quantum scale,
nor the local-context selection law is derived from bare renewal probabilities.
The metric is flat and prescribed. Nonzero scalar stress is not asserted to
satisfy a flat Einstein equation.

\section{A coupled continuum identity which is defined at Higgs zeros}
\label{v16:sec:continuum}

On $[0,T]\times\T^3$ use the convention
$D_\mu=\partial_\mu+\ii A_\mu$,
$E_i=\partial_tA_i-\partial_iA_0$, and $B=\operatorname{curl}A$.
A gauge change is $z\mapsto e^{\ii\beta}z$, $A\mapsto A-\dd\beta$.
The action has density
\[
 \tfrac12|D_0z|^2-\tfrac12\sum_i|D_iz|^2
 -\tfrac\lambda4(|z|^2-b)^2+\tfrac12(|E|^2-|B|^2).
\]
Its equations, with optional actual forcing, are
\begin{align}
 D_0^2z-\sum_iD_i^2z+\lambda(|z|^2-b)z&=f,
 & j_\mu&=\operatorname{Im}(\bar zD_\mu z),\\
 \partial_t E&=-\operatorname{curl}B-j+g,
 &\partial_t B&=\operatorname{curl}E.
 \label{v16:eq:continuum}
\end{align}
The physical Gauss equation is $\operatorname{div}E+j_0=0$.
We impose it when calling a solution physical. The energy comparison itself
does not use it as a cancellation in place of a field equation.

Let $(z_*,A_*,E_*,B_*)$ be an unforced smooth reference with
$r_*=|z_*|\ge\rho_*>0$. Define the global, gauge-invariant reference functions
\begin{equation}
 \mu_\mu=\frac{D_\mu^*z_*}{z_*}=\beta_\mu+\ii a_\mu,
 \qquad \beta_\mu=\partial_\mu\log r_*.
 \label{v16:eq:reference-ratio}
\end{equation}
The star on $D_\mu^*$ denotes the reference connection, not an adjoint.
No global choice of the phase of $z_*$ is needed for this quotient. In particular
$\partial_t\mu_i-\partial_i\mu_0=\ii E_{*,i}$.
For the candidate put
\begin{equation}
 w_\mu=D_\mu z-\mu_\mu z,\quad s=|z|^2,\quad s_*=r_*^2,
 \quad d=s-s_*,\quad e=E-E_*,\quad b_1=B-B_*.
 \label{v16:eq:errors}
\end{equation}
The symbol $b_1$ is a magnetic error and is distinct from the potential
parameter $b$. Every $w_\mu$ transforms covariantly under the candidate gauge;
its squared modulus is invariant. None of these definitions divides by $z$.
Set
\begin{equation}
 \mathcal R(t)=\int_{\T^3}\left[
 \tfrac12\sum_{\mu=0}^3|w_\mu|^2+\tfrac\lambda4d^2
                +\tfrac12(|e|^2+|b_1|^2)\right]\dd x.
 \label{v16:eq:continuum-cost}
\end{equation}

\begin{theorem}[Exact zero-safe gauge--Higgs comparison]
\label[theorem]{v16:thm:continuum}
For smooth fields satisfying \eqref{v16:eq:continuum} and the reference
conditions above,
\begin{align}
 \frac{\dd}{\dd t}\mathcal R
 =\int\biggl[&\beta_0\left(-|w_0|^2+\sum_i|w_i|^2+\lambda d^2\right)
 -2\sum_i a_i\operatorname{Im}(\overline{w_0}w_i)
 -d\,a\cdot e\notag\\
 &\hspace{18mm}+\operatorname{Re}(\overline{w_0}f)+e\cdot g\biggr]\dd x.
 \label{v16:eq:continuum-exact}
\end{align}
Consequently
\begin{equation}
 \mathcal R(t)\le C_T\left[\mathcal R(0)
            +\int_0^t(\|f\|_2^2+\|g\|_2^2)\dd s\right],
 \label{v16:eq:continuum-bound}
\end{equation}
where $C_T$ depends only on $T,\lambda$ and bounds on $\beta_0,a$.
It is independent of a candidate lower-radius bound, candidate derivative
bounds, and any spatial discretization norm. Smooth candidates are allowed
to have arbitrary zero sets.
\end{theorem}
\begin{proof}
The reference scalar equation divided by $z_*$ gives
\[
 \partial_t\mu_0+\mu_0^2-\sum_i(\partial_i\mu_i+\mu_i^2)
                      +\lambda(s_*-b)=0.
\]
Substitution and the commutator $[D_0,D_i]=\ii E_i$ give the exact equations
\begin{align*}
 D_0w_0-\sum_iD_iw_i&=-\mu_0w_0+\sum_i\mu_iw_i-\lambda d z+f,\\
 D_0w_i-D_iw_0&=\mu_0w_i-\mu_iw_0+\ii e_i z,\\
 \partial_t d&=2\operatorname{Re}(\bar zw_0)+2\beta_0d,\\
 j_i-j_{*,i}&=\operatorname{Im}(\bar zw_i)+a_i d.
\end{align*}
The real part of a connection multiplication contributes nothing to the
time derivative of $|w_\mu|^2$. Spatial integration by parts cancels
$\operatorname{Re}(\overline{w_0}D_iw_i+\overline{w_i}D_iw_0)$.
The remaining mixed scalar term is
\[
 \operatorname{Re}\{\mu_i(\overline{w_0}w_i-\overline{w_i}w_0)\}
                 =-2a_i\operatorname{Im}(\overline{w_0}w_i).
\]
The derivative of $\lambda d^2/4$ cancels
$-\lambda d\operatorname{Re}(\overline{w_0}z)$ and leaves
$\lambda\beta_0d^2$. The Maxwell curl terms cancel because curl is
self-adjoint on the torus. Its matter term
$-e_i\operatorname{Im}(\bar zw_i)$ cancels
$\operatorname{Re}(\overline{w_i}\ii e_i z)$ exactly.
The remaining Maxwell current is $-d\,a\cdot e$. This proves the identity.
All calculations were in Cartesian complex variables, including at $z=0$.

Cauchy--Schwarz and $2uv\le u^2+v^2$ bound the first three terms by
$C_{\lambda,\|\beta_0\|_\infty,\|a\|_\infty}\mathcal R$.
The forcings are at most $\mathcal R+(\|f\|_2^2+\|g\|_2^2)/2$.
Integration of the resulting scalar differential inequality proves
\eqref{v16:eq:continuum-bound}. No strong-weak solution extension is being
assumed: the theorem concerns smooth solutions and approximations for which
the displayed integration by parts is justified.
\end{proof}

\begin{proposition}[Radius, currents and complete continuum stress]
\label[proposition]{v16:prop:continuum-sources}
On the same reference chart, with the reference $\mu_\mu,E_*,B_*$ bounded,
\begin{align}
 \||z|-r_*\|_{H_x^1}^2
       +\|\partial_t|z|-\partial_t r_*\|_2^2&\le C\mathcal R,\\
 \|j-j_*\|_1+\|j_0-j_{*,0}\|_1
                 +\|T-T_*\|_1&\le C(\sqrt{\mathcal R}+\mathcal R).
 \label{v16:eq:continuum-sources}
\end{align}
Here $T$ is the stress tensor of the stated classical scalar--Maxwell action;
all its scalar, electromagnetic, and potential terms are retained.
These are gauge-invariant comparisons, not convergence of charged field means.
\end{proposition}
\begin{proof}
For $r=|z|$, $|r-r_*|\le|d|/\rho_*$. The elementary modulus inequality
implies, almost everywhere including across zeros,
\[
 |\partial_\mu r-\beta_\mu r|\le|w_\mu|.
\]
It follows first where $z\ne0$ by taking the real part after division by
its unit phase; it extends using the Lipschitz modulus and approximation
by $(|z|^2+\delta^2)^{1/2}$. Subtract
$\partial_\mu r_* =\beta_\mu r_*$ to prove the first estimate.

Write $D_\mu z=w_\mu+\mu_\mu z$ in every scalar stress contraction.
The difference from its reference consists of products of two $w$'s,
terms bounded by $C|z||w|$, and terms bounded by $C|d|$.
Also $\int |z|^2\le\int s_*+\|d\|_1\le C+C\sqrt{\mathcal R}$.
After integration all these terms are bounded by
$C(\sqrt{\mathcal R}+\mathcal R)$; use
$\mathcal R^{3/4}\le\sqrt{\mathcal R}+\mathcal R$ if needed.
The potential difference is
$\lambda(s_*-b)d/2+\lambda d^2/4$. Maxwell stresses are quadratic in $E,B$
and satisfy the same bound by expansion around $E_*,B_*$. The current
identities in the preceding proof, including their $\mu=0$ version,
complete the argument. This is a comparison of the action-derived tensor,
not an inferred derivative of a trajectory-dependent energy value.
\end{proof}

The theorem is useful even when a candidate develops zeros: only the
\emph{reference} requires a nonzero amplitude. It does not yet prove that
Version 15's inhomogeneous quantum states obey this classical identity.
We next give an exact quantum realization on a nonstatic reference branch
for which there are no unproved quantum remainders.

\section{An exact quantum identity for the literal compact-link action}
\label{v16:sec:quantum}

Take the cubic three-torus, $h=N^{-1}$, $N\ge3$, with $n=h^{-3}$ vertices,
$3n$ oriented nearest-neighbor links, and $3n$ elementary plaquettes. Put
$w=h^3$ and specialize the coefficients of \cref{v15:eq:energy} to
\begin{equation}
 m_x=w,\qquad c_e=h^{-1},\qquad k_e=h,\qquad
 \beta_f^{\rm mag}=h^{-1}.
 \label{v16:eq:physical-coefficients}
\end{equation}
Thus $\Pi_e/h^2$ is the electric density and
$h^{-2}(e^{\ii(C\theta)_f}-1)$ is the complex plaquette difference.
The full uncompressed physical energy is
\begin{equation}
 \boldsymbol E=\boldsymbol T+\boldsymbol W+
 \frac\lambda4w\sum_x(|Z_x|^2-b)^2
                         +\boldsymbol E_{\rm el}+\boldsymbol E_{\rm mag},
 \label{v16:eq:energy}
\end{equation}
where
\[
 \boldsymbol T=\sum_{x,a}\frac{p_{x,a}^2}{2w},\quad
 \boldsymbol W=\frac h2\sum_{e:x\to y}|U_eZ_y-Z_x|^2,
 \quad\boldsymbol E_{\rm el}=\sum_e\frac{\Pi_e^2}{2h},\quad
 \boldsymbol E_{\rm mag}=\frac1h\sum_f(1-\cos(C\theta)_f).
\]
Here $Z_x=q_{x,1}+\ii q_{x,2}$, $p_{x,a}=-\ii\hbar\partial_{q_{x,a}}$,
$\Pi_e=-\ii\hbar\partial_{\theta_e}$, and
$\mathscr D=\ii[\boldsymbol E,\,\cdot\,]/\hbar$.
Products precede every regulator compression. Gauss constraints, operator
normalizations, and the gauge group are exactly those of Version 15.

Choose $r_0>0$, $v_0\ne0$, independently of the cutoff. Let
\begin{equation}
 \ddot r+\lambda r(r^2-b)=0,\qquad r(0)=r_0,\quad\dot r(0)=v_0,
 \qquad \kappa(t)=\dot r(t)/r(t).
 \label{v16:eq:reference-ODE}
\end{equation}
Use a fixed nonzero time interval $[0,T]$ on which $r\ge\rho_*>0$.
Such an interval follows directly from the smooth ODE and $r_0>0$;
energy conservation $\dot r^2/2+\lambda(r^2-b)^2/4=\mathrm{constant}$
bounds its velocity and gives an interval chosen from the initial data.
This reference is an exact spatially homogeneous scalar solution with zero
reference electric, magnetic, and covariant spatial gradient fields.

Define the gauge-invariant positive comparison
\begin{align}
 \boldsymbol C(t)&=\boldsymbol T_\kappa+\boldsymbol W
                +\boldsymbol V_r+\boldsymbol E_{\rm el}+\boldsymbol E_{\rm mag},
 \label{v16:eq:quantum-cost}\\
 \boldsymbol T_\kappa&=\sum_{x,a}
          \frac{(p_{x,a}-w\kappa q_{x,a})^2}{2w},
 &\boldsymbol V_r&=\frac\lambda4w\sum_x(|Z_x|^2-r(t)^2)^2.\nonumber
\end{align}
This contains no inverse power of any quantum radius.

\begin{theorem}[Exact coupled quantum dilation cancellation]
\label[theorem]{v16:thm:quantum}
On the common smooth polynomial core, and as a quadratic-form identity,
\begin{equation}
 \boxed{\quad
 (\partial_t+\mathscr D)\boldsymbol C
   =\kappa\{-2\boldsymbol T_\kappa+2\boldsymbol W+4\boldsymbol V_r\}.
 \quad}
 \label{v16:eq:quantum-exact}
\end{equation}
Consequently
\begin{equation}
 (\partial_t+\mathscr D)\boldsymbol C
       \preceq4|\kappa|\boldsymbol C,\qquad
 0\preceq\boldsymbol C(t)\preceq C_T(\boldsymbol E+I).
 \label{v16:eq:quantum-order}
\end{equation}
These constants are independent of $h$, $\hbar$, all microscopic radii,
and the number of sites. The identity holds before imposing Gauss's law
and on the physical subspace. Electric and magnetic exchanges have not
been set to zero in the candidate dynamics.
\end{theorem}
\begin{proof}
Let
\[
 \boldsymbol D=\tfrac12\sum_{x,a}(q_{x,a}p_{x,a}+p_{x,a}q_{x,a}),
 \qquad\boldsymbol M_2=w\sum_x|Z_x|^2.
\]
The exact Cartesian commutation relations give
\begin{align*}
 \mathscr D\boldsymbol M_2&=2\boldsymbol D,\\
 \mathscr D\boldsymbol D&=2\boldsymbol T-2\boldsymbol W
       -\lambda w\sum_x(|Z_x|^2-b)|Z_x|^2.
\end{align*}
The hopping energy is homogeneous of degree two in the Cartesian scalar
coordinates. The electric and magnetic energies commute with the scalar
dilation. These facts prove the second identity without a link expansion
or any assumption that neighboring links are close to the identity.
The comparison expands exactly as
\[
 \boldsymbol C=\boldsymbol E-\kappa\boldsymbol D
 +\tfrac12\{\kappa^2+\lambda(b-r^2)\}\boldsymbol M_2
                        +\tfrac\lambda4(r^4-b^2)I.
\]
Use $\dot\kappa=-\kappa^2-\lambda(r^2-b)$, the two commutators,
and $w\sum_x1=1$. The coefficient of $\boldsymbol D$ is $2\kappa^2$;
combining it with $-2\kappa\boldsymbol T$ produces
$-2\kappa\boldsymbol T_\kappa+\kappa^3\boldsymbol M_2$.
The remaining scalar multiplication terms are
$\lambda\kappa w\sum_x(|Z_x|^2-r^2)^2$.
The hopping term is $2\kappa\boldsymbol W$. This proves the exact identity,
with all quantum ordering terms already included in $\boldsymbol D$.

Each summand of $\boldsymbol C$ is nonnegative, proving the first inequality.
For the second, use $(A-B)^2\preceq2A^2+2B^2$ on self-adjoint Cartesian
coordinates and $s\le C[1+(s-b)^2]$, $s\ge0$, as well as bounded $r,\kappa$.
All spatial sums carry their displayed masses. There is no extra sites factor.
The core and spectral-truncation argument in \cref{v14:lem:quantum-relative}
justifies the same inequality for energy-form quantum evolutions.
\end{proof}

\begin{proposition}[The charge energy is retained, even though no inverse-radius bound is used]
\label[proposition]{v16:prop:charge}
Under the exact physical-coordinate unitary of \cref{v15:thm:reduction},
\begin{equation}
 \boldsymbol T_\kappa
 =\sum_x\left[
 \frac{(-\ii\hbar\partial_{r_x}-w\kappa r_x)^*
       (-\ii\hbar\partial_{r_x}-w\kappa r_x)}{2w}
                 +\frac{(B\Pi)_x^2}{2w r_x^2}\right]
 \label{v16:eq:radial-charge}
\end{equation}
as positive forms in $\prod r_x\dd r_x\prod\dd\alpha_e/(2\pi)$.
The adjoint is in this radial measure. This identity asserts no self-adjoint
half-line momentum and adds no boundary condition at zero.
\end{proposition}
\begin{proof}
Resolve the Cartesian vector $p_x-w\kappa q_x$ into its radial and angular
components. Their Euclidean norm squares sum without a cross term.
The radial component is $-\ii\hbar\partial_{r_x}-w\kappa r_x$;
the angular component is $L_x/r_x$. On the physical subspace
$L_x=-(B\Pi)_x$ by the inherited exact reduction.
The equality on the punctured smooth core extends by the inherited form
closure. Thus a finite form value pays precisely the charge energy at zeros;
we have not removed that energy by deleting the set of small radii.
\end{proof}

\section{An exactly physical preparation with spatially uniform moments}
\label{v16:sec:preparation}

Use the relational variables of Version 15, not a gauge average of charged
product vectors. Fix a smooth $\chi$ supported inside
$(r_0/2,3r_0/2)$ and equal to one near $r_0$. Put
\begin{equation}
 \sigma_r^2=\hbar h^{-2},\qquad \sigma_a^2=\hbar h^{-1},\qquad
                         0<\hbar\le h^7.
 \label{v16:eq:widths}
\end{equation}
Define the normalized factors
\begin{align}
 f_h(r)&=Z_r^{-1}r^{-1/2}\chi(r)
       \exp\left[-\frac{(r-r_0)^2}{2\sigma_r^2}
                          +\frac{\ii wv_0(r-r_0)}\hbar\right],\nonumber\\
 g_h(a)&=Z_a^{-1}\sum_{k\in\Z}
       \exp\left[-\frac{(a-2\pi k)^2}{2\sigma_a^2}\right],\nonumber\\
 f_h^0&=\prod_x f_h(r_x)\prod_e g_h(\alpha_e),\qquad
                    \psi_h^0=\mathcal U f_h^0.
 \label{v16:eq:initial}
\end{align}
The electric reference momentum is zero. The preparation is a product in
relational coordinates, generally correlated in the original site/link factors.
Its center and phase use only $r_0,v_0$ and the lattice.

\begin{theorem}[Uniform preparation in the coupled physical energy]
\label[theorem]{v16:thm:preparation}
The preparation satisfies every Gauss constraint exactly and
\begin{equation}
 \langle\psi_h^0,\boldsymbol C(0)\psi_h^0\rangle
                    \le C\hbar h^{-5},\qquad
 \| (\boldsymbol E+I)^2\psi_h^0\|\le C,
 \label{v16:eq:preparation}
\end{equation}
for all sufficiently small $h$ and $0<\hbar\le h^7$.
The constants depend on $r_0,v_0,\chi,\lambda,b$, not on the number of sites.
No all-time amplitude restriction is imposed on the quantum dynamics.
\end{theorem}
\begin{proof}
Physical support follows from the inherited unitary $\mathcal U$.
Gaussian integration, with exponentially small radial-cutoff and wrapped-image
errors, gives
\[
 \E(r-r_0)^2\le C\sigma_r^2,\quad
 \E(1-\cos\alpha)\le C\sigma_a^2,\quad
 \E\Pi_e=0,\quad \E\Pi_e^2\le C\hbar^2/\sigma_a^2.
\]
All inverse radial factors are bounded on this \emph{initial} support.
Using \eqref{v16:eq:radial-charge}, the contributions to the initial comparison
are bounded, respectively, by
\begin{align*}
 \text{radial kinetic: }&C(\hbar^2h^{-6}/\sigma_r^2+\sigma_r^2
                                     +\hbar^2h^{-6}),\\
 \text{Gauss angular kinetic: }&C\hbar^2h^{-6}/\sigma_a^2,\\
 \text{hopping: }&Ch^{-2}(\sigma_r^2+\sigma_a^2),\\
 \text{density quartic: }&C(\sigma_r^2+\sigma_r^4),\\
 \text{electric: }&C\hbar^2h^{-4}/\sigma_a^2,\\
 \text{magnetic: }&Ch^{-4}\sigma_a^2.
\end{align*}
There are six incident links at a vertex and four on a plaquette; their
independent initial angular factors give these bounds also for the incidence
sums. Substitution of \eqref{v16:eq:widths} proves the first estimate.
The $h^{-5}$ term includes the \emph{retained} electric-divergence kinetic
cost. It is not obtained by pretending that $B\Pi=0$.

Here are details for the higher moment. Transfer the radial measure to flat
measure by multiplication by $\prod r_x^{1/2}$, keeping its potential
$-\hbar^2/(8wr_x^2)$. Conjugate by the initial radial phase and use
$X_x=(r_x-r_0)/\sigma_r$, $Y_e=\alpha_e/\sigma_a$ on a central angular chart.
Write the energy as $w$ times a sum of $O(n)$ local differential densities.
Their differentiated coefficients, of every fixed order needed for four
density products, are uniformly bounded by polynomials in the finitely many
local $X,Y$, on the radial cutoff support. The relevant scaled derivative
coefficients are
\[
 \frac\hbar{w\sigma_r}=\sqrt\hbar\,h^{-2},\qquad
 \frac\hbar{w\sigma_a}=\sqrt\hbar\,h^{-5/2},\qquad
 \frac\hbar{h^2\sigma_a}=\sqrt\hbar\,h^{-3/2}.
\]
They are bounded when $\hbar\le h^7$. The hopping multiplication is a
square of quantities bounded by
$C[(\sigma_r/h)(|X_x|+|X_y|)+(\sigma_a/h)|Y_e|]$.
The magnetic density is bounded by
$C\sigma_a^2h^{-4}\sum_{e\in f}Y_e^2$; the same polynomial bounds apply
after any fixed number of scaled derivatives. The inverse-radius coefficients
and their scaled derivatives are bounded on the radial support. Terms where
such derivatives hit $\chi$ are supported a fixed distance from $r_0$ and
are bounded by a polynomial times $\exp(-c/\sigma_r^2)$, uniformly after any
fixed number of inverse powers of $h,\hbar$. Wrapped-image overlaps have the
analogous $\exp(-c/\sigma_a^2)$ bound. These statements follow by
$|e^{it}-1|\le|t|$, $1-\cos t\le t^2/2$, and direct differentiation.

It follows that the absolute expectation of any product of at most four
normalized local energy densities is uniformly bounded: each uses a fixed
number of independent Gaussian variables, even when supports overlap.
Expanding $\langle\psi_h^0,\boldsymbol E^j\psi_h^0\rangle$, $j\le4$,
gives $O(n^j)$ such terms, each with weight $w^j$, and $wn=1$.
Thus these moments are uniformly bounded. The smooth cut-off radial and
periodic angular preparation belongs to every operator domain at fixed
$h,\hbar$, so the calculation proves rather than presupposes the second
estimate of \eqref{v16:eq:preparation}.
\end{proof}

The initial nonzero-radius support only specifies a preparation. Schr\"odinger
evolution and the attained tensor-cutoff states are not required to retain
that support. The comparison and its propagation remain zero-safe.

\section{A fully finite local constrained dynamical branch}
\label{v16:sec:finite}

Use exactly the gauge-equivariant tensor cutoff and local gauge-commuting
locked instruments of Version 15, with coefficients
\eqref{v16:eq:physical-coefficients}. Write $E_R=P_R\boldsymbol E P_R$,
$A_R=1+\|E_R\|$, and
\[
 E_R/\hbar=\sum_{\alpha=1}^{M}H_\alpha,\quad M=10n,\quad
 \mathfrak C_h=1+\Bigl(\sum_\alpha\|H_\alpha\|\Bigr)^2
                       +M\sum_\alpha\|H_\alpha\|^2+M^2.
\]
A support enlargement by an already present incident matter site, as in
\cref{v15:prop:multiplicity}, does not change a term's norm or add an action
term. At acceptance a context is chosen uniformly and one of its six
outcomes is recorded; its local parameter is $M\epsilon$. This context
rule remains declared, not derived from the primitive opportunity labels.
The clock remains $d=4\vartheta\epsilon/11$.

\begin{proposition}[Explicit spatial costs for the gauge-compatible tensor cutoff]
\label[proposition]{v16:prop:costs}
For all sufficiently large $R\ge C h^{-1}$,
\begin{equation}
 D_h\le Ch^{-12},\qquad A_R\le Ch^{-3}R,\qquad
                   \mathfrak C_h\le C h^{-6}R^2/\hbar^2,
 \label{v16:eq:costs}
\end{equation}
where $D_h$ is a valid constant in the tensor approximation bound
\cref{v15:eq:tail}.
\end{proposition}
\begin{proof}
Let $A=\sum_x\ell_x+\sum_e\ell_e$ be the reference product energy of
\cref{v15:thm:cutoff}. Expanding the scalar potential gives
$\boldsymbol E=A+W_2+W_0+cI$, with
\[
 |W_2|\le(6h+w\lambda b/2)\sum_x|Z_x|^2,
 \quad \|W_0\|\le6h^{-4},\quad c=\lambda b^2/4-4h^{-3}.
\]
Since $A\succeq(w\lambda/4)\sum_x|Z_x|^4$,
\[
 \|W_2u\|\le q_h\|A^{1/2}u\|,
 \qquad q_h=2(6h+w\lambda b/2)\sqrt{n/(w\lambda)}=O(h^{-2}).
\]
Absorbing half of $\|Au\|$ gives
$\|Au\|\le C h^{-4}\|(\boldsymbol E+I)u\|$.
The form comparison gives $\boldsymbol E+I\preceq Ch^{-4}A$.
The inherited product-tail proof therefore gives
$D_h\le Ch^{-4}(h^{-4})^2=Ch^{-12}$.

On a site cutoff $\ell_x\le R$, the form inequality for $|Z_x|^4$ gives
$\|P_R|Z_x|^2P_R\|\le2\sqrt{R/(w\lambda)}$.
Onsite and electric terms are at most $C(R+1)$, hopping terms at most
$Ch^{-1/2}\sqrt R$, and each magnetic term at most $2h^{-1}$.
For $R\ge Ch^{-1}$ every one is at most $CR$.
There are $10n$ terms. Summing their norms, then dividing by $\hbar$,
proves all remaining assertions in \eqref{v16:eq:costs}. Local Hilbert-space
dimension has not entered any bound.
\end{proof}

Let $\phi_R=P_R\psi_h^0/\|P_R\psi_h^0\|$ and let $\rho_j$ be the actual
resolved state, initially $|\phi_R\rangle\langle\phi_R|$.
Define its positive physical comparison and budget by
\begin{align}
 e_j&=\tr(\rho_jP_R\boldsymbol C(jd)P_R),\nonumber\\
 \mathfrak b_h&=\hbar h^{-5}
       +\frac{h^{-12}}R(1+T/\hbar)^2
                          +\epsilon h^{-9}R^3/\hbar^2.
 \label{v16:eq:budget}
\end{align}

\begin{theorem}[Local resolved gauge--Higgs closure about a moving radial reference]
\label[theorem]{v16:thm:finite}
On the interval specified by \eqref{v16:eq:reference-ODE}, suppose
$\hbar\le h^7$, $R$ is large enough for the inherited tensor Galerkin
estimate and the gauge-commutant multiplicity, and $M\epsilon\le\tau_0$.
Then, for every initial clock phase,
\begin{align}
 \max_{j\le T/d}\E e_j&\le C_T\mathfrak b_h,
 &\PP\{\max_{j\le T/d}e_j>u\}&\le C_T\mathfrak b_h/u,\label{v16:eq:maximal}\\
 \E\max_{j\le T/d}e_j&\le C_T\mathfrak b_h
       \left[1+\log\left(2+\frac{h^{-3}R}{\mathfrak b_h}\right)\right]
       \quad(0<\mathfrak b_h\le1/2).\nonumber
\end{align}
Every attained state has exact Gauss support. Neither the paths nor the
physical regulator exclude small radii. Constants are independent of
$h,\hbar,R,\epsilon,\vartheta$ and of matrix dimension.
\end{theorem}
\begin{proof}
For the uncompressed unitary state
$\psi(t)=e^{-\ii t\boldsymbol E/\hbar}\psi_h^0$,
\cref{v16:thm:quantum,v16:thm:preparation} give
\[
 \sup_{t\le T}\langle\psi(t),\boldsymbol C(t)\psi(t)\rangle
                                      \le C_T\hbar h^{-5}.
\]
Justify this on commuting global energy cutoffs and pass in the energy form
norm; those cutoffs are proof devices, not source operations.
The abstract Ritz/Duhamel proof of \cref{v14:thm:Galerkin} applies to the
present positive energy and product projection, by \cref{v15:thm:cutoff}.
Its higher preparation moment is supplied by \cref{v16:thm:preparation}.
Together with \eqref{v16:eq:costs} it gives
\[
 \sup_{t\le T}\|e^{-\ii tE_R/\hbar}\phi_R-\psi(t)\|_{\boldsymbol E+I}^2
                       \le C_T h^{-12}R^{-1}(1+T/\hbar)^2.
\]
Form domination in \eqref{v16:eq:quantum-order} therefore bounds the quantum
reference comparison by $C_Ta_h$, where
$a_h=\hbar h^{-5}+h^{-12}R^{-1}(1+T/\hbar)^2$.
There is no lattice consistency term: this homogeneous reference is exactly
a solution on every spatial lattice.

Let $P_*(t)$ be the rank-one projector onto the tensor quantum reference,
and $\xi_j=\tr\rho_j(I-P_*(jd))$. The inherited Gauss-preserving local
trajectory theorem \cref{v15:thm:source} gives
$\max\E\xi_j\le C_T\epsilon\mathfrak C_h$ and
$\PP(\max\xi_j>u)\le C_T\epsilon\mathfrak C_h/u$.
For $C_R=P_R\boldsymbol C P_R\succeq0$, the positive block inequality
$C_R\preceq2P_*C_RP_*+2(I-P_*)C_R(I-P_*)$ implies
\[
 e_j\le C_Ta_h+C A_R\xi_j.
\]
Using \eqref{v16:eq:costs} gives the mean estimate and the maximal probability
estimate exactly as in \cref{v14:thm:scalar}; thresholds below $2C_Ta_h$
use the trivial probability bound one. No union bound over sites or ticks
is taken. Since $e_j\le C_TA_R$, integration of
$\min(1,C_T\mathfrak b_h/u)$ gives the logarithmic mean estimate.
Exact Gauss support follows from the inherited outcome-by-outcome
commutation, not from this comparison estimate.
\end{proof}

\section{Physical readers, source insertions, and the radius-chart probability}
\label{v16:sec:sources}

For an attained physical state, transfer it to the relational Hilbert space
when defining radial derivatives. Let $\bar r_x=\tr(\rho_j r_x)$ and
let $\bar v_x$ be the real expectation of
$-\ii\hbar\partial_{r_x}/w$, defined through its kinetic form.
We do not assert that this half-line derivative is self-adjoint.
All physical products are formed before tensor compression.

\begin{proposition}[Amplitude, radial momentum and gauge-field energy]
\label[proposition]{v16:prop:readers}
At each attained state,
\begin{align}
 &\|\bar r-r(t)\|_{2,h}^2+\sum_i\|\delta_i\bar r\|_{2,h}^2
                    +\|\bar v-\dot r(t)\|_{2,h}^2\le C_T e_j,\\
 &w\sum_x\tr\rho_j(|Z_x|^2-r(t)^2)^2
 +\sum_e\frac{\tr\rho_j\Pi_e^2}{h}
 +\frac1h\sum_f\tr\rho_j|e^{\ii(C\theta)_f}-1|^2\le C e_j.
 \label{v16:eq:reader-bounds}
\end{align}
The conditional variances of the radius, electric, and plaquette readers,
and the squared radial-derivative form errors, satisfy the same bounds. Thus the radius and radial momentum means
converge to $(r,\dot r)$ in $H_x^1\times L_x^2$, and the electric and
plaquette-curvature readers converge to zero in the displayed physical
energy norms, uniformly in probability when $\mathfrak b_h\to0$.
\end{proposition}
\begin{proof}
The pointwise inequality
$(r_x-r)^2\le(|Z_x|^2-r^2)^2/\rho_*^2$ controls the radius.
The exact diamagnetic inequality
$|r_y-r_x|\le|U_eZ_y-Z_x|$ controls its spatial differences, including
zeros. Jensen's inequality transfers both to the conditional means.
On the physical subspace \eqref{v16:eq:radial-charge} controls
$\sum_xw\|(-\ii\hbar\partial_{r_x}/w-\kappa r_x)f\|^2$.
Subtract $\dot r=\kappa r$ and use the radius estimate. Quadratic-form
Cauchy--Schwarz bounds the squared error of its real expectation by that
form error. The last three terms are summands of the positive comparison,
using $|e^{it}-1|^2=2(1-\cos t)$.
Variance decompositions prove the corresponding assertions for conditional
fluctuations; the same radius statement uses its multiplication spectral
measure. Cellwise interpolation and the constant spatial reference prove
the final topology. No charged expectation is substituted for a physical
radius, and no finite time-difference velocity is substituted for momentum.
\end{proof}

\begin{proposition}[Same-action lapse, diagonal metric and potential sources]
\label[proposition]{v16:prop:sources}
Use independent lapse $L_x$ and diagonal spatial scale variations $s_{i,x}$,
with $J_x=s_{1,x}s_{2,x}s_{3,x}$, in the canonical matter energy. The scalar
terms are weighted by $L/J$, $LJ/s_i^2$, and $LJ$ for kinetic, $i$-hopping,
and potential terms respectively. Weight electric $i$-terms by $Ls_i^2/J$
and $ij$-plaquettes by $LJ/(s_i^2s_j^2)$. At the evaluated flat coefficients
these give exactly \eqref{v16:eq:energy}. Hold $P_R$ fixed under variations.
For every bounded test of these local coefficients, and for potential
parameters $\lambda,b$, the difference of the attained action insertion
from its homogeneous classical value is at most
\begin{equation}
 C_{T,\mathrm{test}}(\sqrt{e_j}+e_j).
 \label{v16:eq:source-bound}
\end{equation}
In particular the lapse and diagonal spatial stresses tend to
\[
 \rho_*(t)=\tfrac12\dot r(t)^2+V(r(t)),\qquad
 P_*(t)=\tfrac12\dot r(t)^2-V(r(t)),
\]
with zero limiting electromagnetic contribution. Scalar momentum-flux
and quadratic electric--magnetic insertions also tend to zero under their
usual symmetric product convention. No vanishing Einstein metric row is
claimed.
\end{proposition}
\begin{proof}
Write $p_x/w=\kappa q_x+v_x$, with
$\sum_xw\tr\rho |v_x|^2\le2e_j$.
Then the kinetic-density error from $\dot r^2$ is the sum of a $v^2$ term,
a symmetric $qv$ term, and $\kappa^2(|Z_x|^2-r^2)$.
The radius density error is $O(\sqrt{e_j})$ in weighted $L^1$;
$\sum_xw\tr\rho|Z_x|^2\le C_T+C\sqrt{e_j}$.
State Cauchy--Schwarz bounds the cross term by
$C(\sqrt{e_j}+e_j)$, with the operator product symmetrized.
The quartic identity
\[
 V(|Z_x|)-V(r)=\tfrac\lambda2(r^2-b)(|Z_x|^2-r^2)
                          +\tfrac\lambda4(|Z_x|^2-r^2)^2
\]
proves the same bound for the potential and its parameter derivatives.
Every hopping, electric, and plaquette source in the stated coefficient
bank is bounded in absolute expectation by its positive energy, hence by
$Ce_j$. Differentiation of the displayed independent metric weights proves
that these are actual action derivatives. Products between momentum and
covariant gradient are $O(\sqrt{e_j}+e_j)$ by Cauchy--Schwarz; gauge quadratic
cross insertions are $O(e_j)$. For a sine-curvature convention use
$\sin^2\theta\le2(1-\cos\theta)$.
The argument does not specify an otherwise unprovided off-diagonal compact
metric regulator. It covers the stated action-source bank exactly.
\end{proof}

\begin{theorem}[A quantitative radius chart is an output, not a prior exclusion]
\label[theorem]{v16:thm:radius}
For the same cubic lattice and reference $r(t)\ge\rho_*$,
let $\mathsf Z_h$ be multiplication by
$\mathbf1_{\{\min_x r_x<\rho_*/2\}}$ in the full configuration space.
At every attained state,
\begin{equation}
 \tr(\rho_jP_R\mathsf Z_hP_R)\le C_T h^{-1}e_j,
 \qquad
 \max_x|\bar r_x-r(t)|^2\le C_T h^{-1}e_j.
 \label{v16:eq:radius}
\end{equation}
Hence
\begin{equation}
 \PP\left\{\max_{j\le T/d}\tr(\rho_jP_R\mathsf Z_hP_R)>u\right\}
                   \le C_T\mathfrak b_h/(h u).
 \label{v16:eq:radius-prob}
\end{equation}
This is a bound on the low-radius Born weight of every attained conditional
state; it does not perform repeated position measurements or assert a joint
unmeasured classical radius path. The mean-radius chart has the corresponding
actual-record probability bound.
\end{theorem}
\begin{proof}
For any scalar grid function $v$, normalized discrete Fourier expansion and
Cauchy--Schwarz give
\[
 \|v\|_\infty^2\le
 \left(\sum_{k\in\mathcal B_N}(1+\lambda_h(k))^{-1}\right)
 \left(\|v\|_{2,h}^2+\sum_i\|\delta_i v\|_{2,h}^2\right)
 \le C h^{-1}\|v\|_{H_h^1}^2.
\]
Here $\lambda_h(k)=4h^{-2}\sum_i\sin^2(\pi k_i/N)\ge c|k|^2$ on the
centered Brillouin zone, and summing $(1+|k|^2)^{-1}$ in three dimensions
gives $CN$. Apply this to $v_x=r_x-r(t)$ on each configuration.
The density and diamagnetic inequalities in \cref{v16:prop:readers} show
$\|v\|_{H_h^1}^2\le C_T(\boldsymbol W+\boldsymbol V_r)$ pointwise.
On the event defining $\mathsf Z_h$, $\|v\|_\infty\ge\rho_*/2$.
Thus $\mathsf Z_h\preceq C_Th^{-1}(\boldsymbol W+\boldsymbol V_r)$ as
multiplication operators. Compress and pair with the state. The mean bound
also follows by applying the grid inequality to the actual conditional means.
Use \cref{v16:thm:finite} with threshold $c h u$ for the last assertion.
No inverse of a fluctuating radius has been bounded or inserted.
\end{proof}

\section{An explicit simultaneous limit and its precise boundary}
\label{v16:sec:limit}

\begin{corollary}[A nonstatic local gauge-constrained classical branch]
\label[corollary]{v16:cor:scales}
Choose
\begin{equation}
 \boxed{\hbar_h=h^7,\qquad R_h=h^{-28},\qquad
 \epsilon_h=h^{109},\qquad
 d_h=\frac{4\vartheta_h}{11}h^{109},\quad0<\vartheta_h\le1.}
 \label{v16:eq:scales}
\end{equation}
Then $\mathfrak b_h=O_T(h^2)$, the hypotheses of the finite construction
hold for small $h$, and every attained record remains in the exact Gauss
sector. The conclusions of \cref{v16:prop:readers,v16:prop:sources} hold
uniformly in probability on the nonstatic reference interval.
The expected maximum comparison is $O_T(h^2(1+|\log h|))$; the probability
in \eqref{v16:eq:radius-prob} is $O_T(h/u)$ for fixed $u>0$.
\end{corollary}
\begin{proof}
The three contributions in \eqref{v16:eq:budget} have orders
$h^{7-5}=h^2$, $h^{-12+28-14}=h^2$, and
$h^{109-9-84-14}=h^2$. The uniform higher initial moment and
$D_h/R=O(h^{16})$ supply the inherited tensor-approximation threshold.
The local neutral multiplicity is present for sufficiently large $R$ by the
same fixed-support radial trial functions as in \cref{v15:prop:multiplicity};
for the present coefficients their kinetic weights $\hbar^2/w\le h^{11}$
and quartic weights $w$ do not increase those trial bounds.
Also $M\epsilon=10h^{106}\to0$.
Since $A_R\le Ch^{-31}$, the logarithm in \cref{v16:thm:finite} is
$O(1+|\log h|)$. The radius probability follows by substitution.
\end{proof}

\paragraph{What has actually advanced.}
The comparison does not demand that the fluctuating Higgs radius be bounded
away from zero. A gauge-invariant observable formed in the original Cartesian
variables controls the charge, radial kinetic, hopping, electric, magnetic,
and density-quartic energies simultaneously. For the homogeneous moving
reference it has an exact quantum drift identity on the unchanged compact-link
action. This supplies a genuinely dynamical, locally implemented and exactly
Gauss-preserving branch, with a spatially uniform estimate and a posterior
radius-chart bound. Its source operations use the action and initial data,
not future reference fields. The nonconstant radial motion and its pressure
are retained; this is not a static vacuum record.

\paragraph{What has not advanced under stronger quantifiers.}
The limiting gauge field of the explicit quantum branch is zero. General
inhomogeneous reference fields with nonzero gauge curvature are covered by
\cref{v16:thm:continuum} only at the classical PDE level, not by a new
cutoff-uniform quantum lattice propagation theorem. Transferring that
zero-safe comparison to the compact-link quantum operator, including its
spatial reference ratios and noncommuting electric terms, is the next analytic
step. No modification of the locked source can be inferred from the comparison
formula alone. The context rule, physical preparation, and semiclassical
normalization remain declared; native operator/action identification,
non-Abelian and chiral sectors, and the Einstein constraints and backreaction
remain independent. The exponents prove existence with explicit costs,
not computational efficiency.

\paragraph{Verification and preservation.}
All pre-existing bytes before the final document command in Version 15 are
preserved. New labels have prefix \texttt{v16:}. The diagnostics check the
continuum cancellation, exact dilation algebra, radial-charge square,
finite compact-link classical identity, source-weight derivatives, zero-radius
cases, discrete radius-capacity estimate, Gaussian scale powers, and the
simultaneous budget. An oscillator-basis matrix test checks the quantum
identity on an interior core where its finite canonical commutators agree
with the uncompressed ones. It is a diagnostic of the written algebra, not an
exact full-lattice diagonalization, a machine-checked proof, or a verification
of the intended primitive Einstein--SM source.

\begingroup
\renewcommand{\refname}{Additional reference for Version 16}
\begin{thebibliography}{9}
\small\raggedright
\bibitem{v16:Salvi}
T.~Salvi, \emph{Semi-classical limit of the massive Klein--Gordon--Maxwell
system toward the relativistic Euler--Maxwell system via an adapted
modulated energy method}, 2025 preprint.
\href{https://arxiv.org/abs/2502.06622}{arXiv:2502.06622}.
Background for modulated stress--energy in coupled scalar and Maxwell systems;
its target equations and limiting procedure are different from the two
comparisons proved in this body.
\end{thebibliography}
\endgroup
\addtocontents{toc}{\protect\endgroup}
% END IMMUTABLE VERSION 16 BODY
% APPEND VERSION 17 HERE, BEFORE THE FINAL END-DOCUMENT COMMAND.


% BEGIN IMMUTABLE VERSION 17 BODY
\clearpage
\begin{center}
{\Large\bfseries Version 17: Electric curvature before a general inhomogeneous estimate}\\[0.5em]
{\large Exact shifted-link quantum comparison, a positive mixed-product identity,
and a local constrained branch with electromagnetic stress}
\end{center}
\makeatletter
\addtocontents{toc}{\protect\begingroup}
\addtocontents{toc}{\protect\makeatletter}
\addtocontents{toc}{\protect\def\protect\l@section{\protect\bfseries\protect\@dottedtocline{1}{0em}{2.5em}}}
\makeatother
\addcontentsline{toc}{part}{Version 17: Nonzero electric field strength on the local resolved branch}

\section{Audit and the precise extension beyond a zero gauge reference}
\label{v17:sec:audit}

The quantum conclusion of Version 16 has zero limiting electromagnetic field,
although its candidate gauge variables are dynamical. Its classical continuum
comparison permits a nonzero reference curvature, but that comparison is not
itself an operator proof for the compact-link lattice. We take a definite step
between those statements: a spatially homogeneous reference with a genuinely
nonzero electric field and nonlinear Higgs--electric feedback. The candidate
quantum states remain fully spatial, the magnetic energy remains in their
action, and candidate Higgs zeros are not excluded. The limiting magnetic
field of this construction is still zero.

\paragraph{Nonduplication and inherited inputs.}
The action, physical masses, Gauss representation, equivariant tensor cutoff,
local commuting Kraus construction, and renewal clock are exactly
\cref{v15:eq:Gauss,v15:thm:cutoff,v15:thm:locked,v16:eq:physical-coefficients,v16:eq:energy}.
The tensor Galerkin estimate is \cref{v14:thm:Galerkin}; the local resolved-state
comparison is \cref{v15:thm:source}. Neither is proved anew. The zero-reference
dilation identity and its original preparation are
\cref{v16:thm:quantum,v16:thm:preparation}. Searches followed by literal checks
of the supplied Grand Tensor, classical closure, and SM--Einstein monographs
located the previously retained gauge, constraint, and approximation results,
but not the shifted-link operator identity below. This is a targeted audit,
not a proof of bibliographic novelty or an independent verification of those
entire sources.

\paragraph{What is added.}
A momentum-degree observation removes a possible quantum-ordering obstacle
before the lattice calculation. An exact mixed scalar--link product identity
then bounds the new electric force without a microscopic radius floor or an
uncontrolled cubic energy term. These facts give a spatially uniform operator
comparison for nonzero electric references. A shifted, exactly physical
initial preparation populates the estimate. The existing local instrument
then gives actual resolved trajectories with nonzero limiting electromagnetic
energy, current, and anisotropic metric sources.

\paragraph{Scope and attribution.}
Weyl commutator calculus and homogeneous gauge--scalar reductions are standard
methods. The distinction between exact Poisson and Moyal brackets is important;
see \cite{v17:Robert}. The restricted differential-operator identity needed
here is proved directly. No general exact quantization of nonlinear observables
is claimed. The positive realization is an extension of the \emph{constructed}
Abelian gauge--Higgs branch, not an identification of the intended native
Einstein--SM law. Its initial preparation, semiclassical normalization, and
local-context scheduler remain declared data. The background metric is flat
and is not solved from the resulting nonzero stress.

\section{The operator promotion has a restricted exact calculus}
\label{v17:sec:calculus}

Use Cartesian scalar coordinates and circle link coordinates jointly denoted
by $x$, with their canonical momenta $p$. Coefficients are smooth and periodic
in circle coordinates. All assertions below are on smooth functions compactly
supported in the Cartesian coordinates, with the corresponding periodic core.
The Poisson convention is $\{c,e\}=c_x\cdot e_p-c_p\cdot e_x$.

\begin{lemma}[Exact commutator on the comparison class]
\label[lemma]{v17:lem:Weyl}
Let
\[
 e(x,p)=\tfrac12p^{\mathsf T}Mp+W(x),\qquad M=M^{\mathsf T}
\]
with constant $M$, and let $c(t,x,p)$ be polynomial of degree at most two
in $p$, with smooth coefficient functions of $x$. Let $\operatorname{Op}^W_\hbar$
denote symmetric Weyl quantization. Then
\begin{equation}
 \left(\partial_t+\frac\ii\hbar
       [\operatorname{Op}^W_\hbar(e),\,\cdot\,]\right)
       \operatorname{Op}^W_\hbar(c)
 =\operatorname{Op}^W_\hbar(\partial_t c+\{c,e\}).
 \label{v17:eq:Weyl}
\end{equation}
This is an exact identity, not an expansion with an unspecified remainder.
It does not assert that quantization preserves every positive symbol.
\end{lemma}
\begin{proof}
Decompose $c$ into $c_0(x)$, $c_j(x)p_j$, and $c_{jk}(x)p_jp_k$.
Composition with $W(x)$ is the finite product rule for a differential
operator of order at most two. In its antisymmetrization, the even-order
terms cancel and only the first-order Poisson term remains. Composition
with the constant quadratic kinetic polynomial has the same property:
it has at most two momentum derivatives, and antisymmetrization again
leaves only the first-order term. This may also be checked using
$[p_j,f]=-\ii\hbar\partial_jf$ and the Weyl formulas
\[
 \operatorname{Op}^W(fp_j)=\tfrac12(fp_j+p_jf),\qquad
 \operatorname{Op}^W(fp_jp_k)
 =fp_jp_k-\frac{\ii\hbar}{2}(f_jp_k+f_kp_j)
                          -\frac{\hbar^2}{4}f_{jk}.
\]
Substitution gives \eqref{v17:eq:Weyl} term by term. Periodic coefficients
obey the identical product rule on the circle, so no global angle operator
is needed. Differentiating the coefficient functions accounts for $\partial_t$.
\end{proof}

For the comparison used below, its highest momentum coefficients are constant
and its linear terms have affine Cartesian coefficients or constant link
coefficients. Its Weyl operator is therefore precisely the displayed sum of
positive squares and multiplication operators. Positivity will be proved
through those squares, not inferred from symbol positivity.

\begin{example}[The restriction on the observable matters]
For $e=p^2/2+\lambda q^4/4$ and $c=p^3$, the quantum symbol of the derivative
contains, in addition to $\{p^3,e\}=-3\lambda q^3p^2$, the term
$3\hbar^2\lambda q/2$. This follows by applying the product rule three times
to $[q^4,p^3]$ in Weyl order. Thus the preceding lemma does not extend to
all observables for a quartic Hamiltonian. Nor does it give an exact
commutator after a noncommuting tensor cutoff; that cutoff error is paid
by the inherited Galerkin theorem.
\end{example}

\section{An electric reference generated by the same compact-link action}
\label{v17:sec:reference}

Retain $\Lambda_h=(\Z/N\Z)^3$, $h=N^{-1}$, $N\ge3$, $w=h^3$, the energy
$\boldsymbol E$ of \cref{v16:eq:energy}, and its potential parameter $b$.
An edge $e=(x,i)$ runs from $x$ to $x+e_i$. An incoming edge is $(x-e_i,i)$.
For directionwise real functions $a_i(t)$ set
\begin{equation}
 \delta_i=ha_i,\qquad s_i=\sin\delta_i,\qquad c_i=1-\cos\delta_i,
 \qquad \ell_h(a)=\frac2{h^2}\sum_{i=1}^3c_i.
 \label{v17:eq:trig}
\end{equation}
The letters $c_i,s_i$ in this body are trigonometric coefficients, not clock
phases or scalar densities. Consider the finite reference ODE
\begin{equation}
 \begin{split}
 \dot r_h=v_h,&\qquad
 \dot v_h=-r_h\{\ell_h(a_h)+\lambda(r_h^2-b)\},\\
 \dot a_{h,i}=\mathcal E_{h,i},&\qquad
 \dot{\mathcal E}_{h,i}=-r_h^2\frac{\sin(ha_{h,i})}{h}.
 \end{split}
 \label{v17:eq:reference-ODE}
\end{equation}
Initial values $r_0>0,v_0,a_0,\mathcal E_0$ are fixed independently of $h$.
We take $\mathcal E_0\ne0$; $a_0$ may also be nonzero. Define
$\kappa_h=v_h/r_h$ wherever $r_h>0$.

\begin{proposition}[Exact homogeneous reference and its regular limit]
\label[proposition]{v17:prop:reference}
The configuration
\begin{equation}
 Z_x^*=r_h,\quad p_{x,1}^*=wv_h,\quad p_{x,2}^*=0,\quad
 U_{x,i}^*=e^{\ii ha_{h,i}},\quad \Pi_{x,i}^*=h^2\mathcal E_{h,i}
 \label{v17:eq:reference-fields}
\end{equation}
is an exact classical trajectory of the same lattice action. Its Gauss
constraint holds, and its spatial plaquette angles vanish. It has conserved
energy per unit volume
\begin{equation}
 \mathcal H_h^{\rm ref}
 =\tfrac12v_h^2+\tfrac12|\mathcal E_h|^2
 +\frac{r_h^2}{h^2}\sum_i(1-\cos(ha_{h,i}))
 +\frac\lambda4(r_h^2-b)^2.
 \label{v17:eq:ref-energy}
\end{equation}
There is $T>0$, independent of sufficiently small $h$, on which $r_h\ge\rho_*>0$
and all reference coefficients and the time derivatives used below are uniformly
bounded. On a possibly shorter such interval,
$|\mathcal E_h(t)|\ge|\mathcal E_0|/2$.
The reference converges uniformly at rate $O(h^2)$, including its first
time derivative, to
\begin{equation}
 \ddot r+\{ |a|^2+\lambda(r^2-b)\}r=0,\qquad
 \ddot a_i+r^2a_i=0.
 \label{v17:eq:limit-ODE}
\end{equation}
The continuum reference has $F_{0i}=\dot a_i$, zero spatial magnetic field,
and scalar current $j_i=r^2a_i$, $j_0=0$.
\end{proposition}
\begin{proof}
At the homogeneous real scalar, the two incident links in direction $i$
contribute $2h(1-\cos(ha_i))r$ to the Cartesian force. Dividing by $w$
gives the term in $\ell_h$. The link equations are
$\dot\theta=\Pi/h$, $\dot\Pi=-h r^2\sin\theta$; the magnetic force is zero
on zero plaquette angle. These are exactly \eqref{v17:eq:reference-ODE}.
Electric momentum is constant over all edges of each direction, so $B\Pi=0$;
angular scalar momentum is zero. This proves Gauss's law without deleting
it. Differentiating \eqref{v17:eq:ref-energy} and inserting the ODE cancels
both scalar--electric exchange terms, proving conservation.

On each bounded set of $(r,v,a,\mathcal E)$ the ODE vector fields and their
first derivatives are bounded uniformly for $0<h\le1$; use
$|\sin(ha)|/h\le|a|$ and $2(1-\cos(ha))/h^2\le a^2$.
Picard iteration in a fixed ball about the initial point gives a common
existence interval. Reducing it controls $r$ and $\mathcal E$ by continuity
with the same bounds. The integral Taylor remainders for sine and cosine
give a uniform $O(h^2)$ difference of the finite and limiting vector fields
on that ball. Subtract their integral equations and apply the ordinary
finite-dimensional Gronwall inequality. This proves the stated rate; substituting
back into the ODE gives the derivative rate. The final field and current
identifications use $D_i r=\ii a_i r$ and $D_0r=\dot r$ in the conventions
of \cref{v16:eq:continuum}.
\end{proof}

This reference is not a gauge removal of a physical electric field. The
spatial cycle holonomy is $e^{\ii a_i}$, and the time derivative of $a_i$
gives nonzero field strength. Its shifts below occur only inside a
comparison observable. No source operation consults its future values.

\section{An exact shifted-link comparison without a candidate radius floor}
\label{v17:sec:identity}

Suppress the subscript $h$ on reference functions within this section.
For every candidate configuration, including configurations with $Z_x=0$,
put
\begin{equation}
 \begin{gathered}
 \widetilde U_{x,i}=e^{-\ii\delta_i}U_{x,i},\qquad
 W_{x,i}=\frac{\widetilde U_{x,i}Z_{x+e_i}-Z_x}{h},\\
 V_{x,a}=p_{x,a}/w-\kappa q_{x,a},\qquad
 D_x=|Z_x|^2-r^2,\qquad
 E_{x,i}=\Pi_{x,i}/h^2-\mathcal E_i,\\
 J_{x,i}=\overline Z_x\widetilde U_{x,i}Z_{x+e_i}.
 \end{gathered}
 \label{v17:eq:errors}
\end{equation}
Here $V_x$ may be identified with its two-component real vector, or with
$V_{x,1}+\ii V_{x,2}$. $W,D,J$ are multiplication operators. The quantities
called $E_{x,i}$ here are electric \emph{errors}, not the full Hamiltonian.
Define positive operators/forms
\begin{equation}
 \begin{gathered}
 T_\kappa=\frac w2\sum_{x,a}V_{x,a}^2,\qquad
 W_\delta=\frac w2\sum_{x,i}|W_{x,i}|^2,\qquad
 V_r=\frac{\lambda w}{4}\sum_xD_x^2,\qquad
 E_{\rm diff}=\frac w2\sum_{x,i}E_{x,i}^2,\\
 \boldsymbol C_h=T_\kappa+W_\delta+V_r+E_{\rm diff}
                                      +\boldsymbol E_{\rm mag}.
 \end{gathered}
 \label{v17:eq:cost}
\end{equation}
All actual magnetic plaquettes remain. Since the reference shift is constant
in each direction, it has zero plaquette curl. Every term is gauge invariant.
No inverse candidate radius and no principal argument of an actual link occurs.

Define the multiplication remainders
\begin{align}
 F_x=-\frac1h\sum_i\bigl[&(c_i-\ii s_i)W_{x,i}
 -(c_i+\ii s_i)\widetilde U_{x-e_i,i}^{*}W_{x-e_i,i}\bigr],
 \label{v17:eq:scalar-force}\\
 G_{x,i}&=\frac1h\bigl[c_i\operatorname{Im}J_{x,i}
                         -s_i(\operatorname{Re}J_{x,i}-r^2)\bigr].
 \label{v17:eq:electric-force}
\end{align}
For real self-adjoint components define
$\operatorname{Sym}_{\R}(V,F)=\frac12\sum_{a=1}^2(V_aF_a+F_aV_a)$;
for a real electric pair use $\operatorname{Sym}(E,G)=(EG+GE)/2$.

\begin{lemma}[The mixed product pays its electric force quadratically]
\label[lemma]{v17:lem:product}
For every edge and every candidate, one has the exact identity
\begin{equation}
 \boxed{\quad |J_{x,i}-r^2|^2
 =D_xD_{x+e_i}+r^2h^2|W_{x,i}|^2.\quad}
 \label{v17:eq:product}
\end{equation}
In particular its right side is nonnegative. If
$A_* =\max_i|a_i|$, with $\|\cdot\|_{2,h}^2=w\sum|\cdot|^2$, then
\begin{align}
 \|F\|_{2,h}^2&\le12A_*^2\|W\|_{2,h}^2,\\
 \|G\|_{2,h}^2&\le A_*^2
                 \bigl(3\|D\|_{2,h}^2+r^2h^2\|W\|_{2,h}^2\bigr).
 \label{v17:eq:force-bounds}
\end{align}
These are also pointwise multiplication-operator inequalities. There is no
bound on the size or derivative of the candidate in their constants.
\end{lemma}
\begin{proof}
Write $z=Z_x$, $t=\widetilde U_{x,i}Z_{x+e_i}$ and expand
\[
 |\bar zt-r^2|^2=(|z|^2-r^2)(|t|^2-r^2)+r^2|t-z|^2.
\]
Unitarity of $\widetilde U$ proves \eqref{v17:eq:product}, also at zeros.
Moreover $\sqrt{c_i^2+s_i^2}=|1-e^{\ii ha_i}|\le h|a_i|$.
The square of the sum of the six terms in \eqref{v17:eq:scalar-force}
is at most six times their squared sum. Summing over sites counts each
direction twice and gives the factor twelve. For $G$, Cauchy--Schwarz in
the real and imaginary parts gives $|G_{x,i}|^2\le A_*^2|J_{x,i}-r^2|^2$.
Use $D_xD_y\le(D_x^2+D_y^2)/2$ and sum over three positive directions.
\end{proof}

The mixed-product identity is important: estimating
$\operatorname{Re}J-r^2$ only by $D_x+D_y-h^2|W|^2$ would introduce an
apparently uncontrolled electric-error times $|W|^2$ term. Equation
\eqref{v17:eq:product} bounds the complete complex product directly in the
existing density and gradient energies.

\begin{theorem}[Exact electric-reference quantum comparison]
\label[theorem]{v17:thm:identity}
For the unchanged energy $\boldsymbol E$ of \cref{v16:eq:energy}, the
reference \eqref{v17:eq:reference-ODE}, and
$\mathscr D=\ii[\boldsymbol E,\cdot]/\hbar$, one has on the common core
and as quadratic forms
\begin{equation}
 \boxed{\begin{aligned}
 (\partial_t+\mathscr D)\boldsymbol C_h
 ={}&\kappa(-2T_\kappa+2W_\delta+4V_r)\\
   &+w\sum_x\operatorname{Sym}_{\R}(V_x,F_x)
    +w\sum_{x,i}\operatorname{Sym}(E_{x,i},G_{x,i}).
 \end{aligned}}
 \label{v17:eq:identity}
\end{equation}
The corresponding classical identity is the same formula with ordinary real
products. For $h\le1$ one may take
\begin{equation}
 (\partial_t+\mathscr D)\boldsymbol C_h\preceq K_T\boldsymbol C_h,
 \quad K_T=\sup_{t\le T}
 \left[4|\kappa|+2+12A_*^2+6A_*^2/\lambda+r^2A_*^2\right],
 \label{v17:eq:bound}
\end{equation}
and $0\preceq\boldsymbol C_h(t)\preceq C_T(\boldsymbol E+I)$ with
$C_T$ independent of $h,\hbar$, and the number of sites. In particular no
candidate lower-radius hypothesis or $h^{-2}$ stability constant is used.
\end{theorem}
\begin{proof}
We first compute the classical identity with all lattice terms retained.
In the translated link coordinate $\theta_{x,i}-\delta_i$ the velocity is
$hE_{x,i}$. Put $\Delta_{\widetilde U}$ for the covariant lattice Laplacian.
The scalar equations and
$\dot\kappa+\kappa^2=-\ell_h(a)-\lambda(r^2-b)$ give
\[
 \dot Z_x=V_x+\kappa Z_x,\qquad
 \dot V_x=\Delta_{\widetilde U}Z_x-\lambda D_xZ_x-\kappa V_x+F_x.
\]
To verify the last term explicitly, the difference of the original and
shifted hopping energies is
\[
 \boldsymbol W-W_\delta
      =h\sum_{x,i}(c_i\operatorname{Re}J_{x,i}
                                      +s_i\operatorname{Im}J_{x,i}).
\]
Its derivatives at the tail and head have complex coefficients $c_i-\ii s_i$
and $c_i+\ii s_i$. Subtract their scalar zero-order term
$2h\sum_i c_iZ_x$, which is exactly $w\ell_h(a)Z_x$.
The remaining outgoing and incoming differences are
\eqref{v17:eq:scalar-force}. This also verifies the incidence signs.

The shifted gradient and the density obey
\[
 \dot W_{x,i}=\kappa W_{x,i}
       +\frac{\widetilde U_{x,i}V_{x+e_i}-V_x}{h}
       +\ii E_{x,i}\widetilde U_{x,i}Z_{x+e_i},\qquad
 \dot D_x=2\operatorname{Re}(\overline Z_xV_x)+2\kappa D_x.
\]
The electric equation, after subtracting the reference equation, is
\[
 \dot E_{x,i}=-\frac{\operatorname{Im}J_{x,i}}h
             +G_{x,i}+\text{the original magnetic force}/h^2.
\]
Indeed $\operatorname{Im}(e^{\ii\delta_i}J)=
\cos\delta_i\operatorname{Im}J+\sin\delta_i\operatorname{Re}J$ and
$\dot{\mathcal E}_i=-r^2\sin\delta_i/h$.
The scalar Laplacian cancels the shifted-gradient derivative by finite
summation by parts. The quartic force cancels the first term in $\dot D$.
The remaining electric exchange cancels because
\[
 \operatorname{Re}\left(\overline W_{x,i}\,\ii E_{x,i}
                        \widetilde U_{x,i}Z_{x+e_i}\right)
       = E_{x,i}\operatorname{Im}J_{x,i}/h.
\]
Finally the magnetic derivative cancels its electric force. The reference
$h^2\mathcal E_i$ has zero discrete curl, so this cancellation is unchanged
by subtracting it. The dilation terms left over are exactly
$\kappa(-2T_\kappa+2W_\delta+4V_r)$. This proves the classical identity.

The Hamiltonian has constant diagonal quadratic kinetic energy in the
Cartesian and link momenta. The comparison symbol is of momentum degree
two, with the same constant highest-order coefficients. Its Weyl
quantization is \eqref{v17:eq:cost}: the shifted squares already include the
correct symmetric Cartesian dilation term. The right side just computed
is at most quadratic in momenta; products linear in momenta become precisely
the displayed symmetrizations. Apply \cref{v17:lem:Weyl}. This proves the
exact operator identity, with no omitted $\hbar^2$ remainder.

For the inequality use, on each vector, the elementary bounds
$\operatorname{Sym}_{\R}(V,F)\preceq(|V|^2+|F|^2)/2$ and
$\operatorname{Sym}(E,G)\preceq(E^2+G^2)/2$. Inserting
\eqref{v17:eq:force-bounds}, $\|W\|_{2,h}^2\le2\boldsymbol C_h$,
$\|D\|_{2,h}^2\le4\boldsymbol C_h/\lambda$, and the other positive
squares gives \eqref{v17:eq:bound}.

The energy-form upper bound follows termwise. The shifted scalar momentum
is bounded by twice its original kinetic square plus
$C_Tw\sum|Z|^2$. The shifted electric energy is bounded by twice its original
energy plus $C_T I$. The shifted hopping difference is bounded by twice
the original hopping plus
$h\sum|1-e^{-\ii ha_i}|^2|Z_{x+e_i}|^2\le C_Tw\sum|Z|^2$.
The density quartic is bounded by the original quartic plus a constant.
Finally $s\le C[1+(s-b)^2]$, $s\ge0$, and $w\sum_x1=1$ control all
remaining terms. These estimates are independent of the candidate amplitude.
The common-core identity and its energy-relative form bound extend to
energy-form quantum evolutions by the spectral approximation argument
already used in \cref{v14:lem:quantum-relative}: the auxiliary global spectral
projections appear only in that proof, not in the physical tensor regulator.
\end{proof}

\begin{corollary}[The complete Gauss charge energy remains inside the comparison]
\label[corollary]{v17:cor:charge}
On the exact physical space the scalar kinetic comparison has the same
positive decomposition as \cref{v16:eq:radial-charge}, with $\kappa=\kappa_h$.
Moreover $B\Pi=B(\Pi-h^2\mathcal E)$, since the reference electric field is
constant in each direction. Thus the retained term
$\sum_x(B\Pi)_x^2/(2wr_x^2)$ has not been dropped by either reference shift.
There is no need to define a self-adjoint half-line radial momentum.
\end{corollary}
\begin{proof}
The scalar square is unchanged except for its reference coefficient, so
\cref{v16:prop:charge} applies. Incoming and outgoing copies of each
constant directional reference momentum cancel under $B$.
\end{proof}

\section{A shifted physical preparation with the same spatial cost}
\label{v17:sec:preparation}

Keep the relational unitary $\mathcal U$ and the radial factor of
\cref{v16:eq:initial}, centered at $r_0$ with phase $wv_0(r-r_0)/\hbar$.
Use the same widths
$\sigma_r^2=\hbar h^{-2}$ and $\sigma_a^2=\hbar h^{-1}$.
For an edge in direction $i$, replace only its initial angular factor by
\begin{equation}
 g_{h,i}(\alpha)=Z_{h,i}^{-1}e^{\ii k_{h,i}(\alpha-ha_{0,i})}
  \sum_{m\in\Z}\exp\left[-\frac{(\alpha-ha_{0,i}-2\pi m)^2}
                                    {2\sigma_a^2}\right],\qquad
 k_{h,i}=\operatorname{round}(h^2\mathcal E_{0,i}/\hbar).
 \label{v17:eq:prep}
\end{equation}
All edges of the same direction have the same integer. The preparation is
still a product in relational variables, lifted by $\mathcal U$ into the
exact Gauss sector. It is not a gauge average of charged product vectors.

\begin{theorem}[Uniform preparation with nonzero electric center]
\label[theorem]{v17:thm:preparation}
For bounded fixed initial data as above, $0<\hbar\le h^7$ and sufficiently
small $h$, this preparation $\psi_h^0$ satisfies
\begin{equation}
 G_x\psi_h^0=0,\qquad
 \langle\psi_h^0,\boldsymbol C_h(0)\psi_h^0\rangle
                         \le C\hbar h^{-5},\qquad
 \|(\boldsymbol E+I)^2\psi_h^0\|\le C.
 \label{v17:eq:prep-bound}
\end{equation}
The constants are independent of the number of sites. Neither an all-time
support condition away from zero nor a zero electric preparation is imposed.
\end{theorem}
\begin{proof}
Exact physical support follows from \cref{v15:thm:reduction}. Translate each
angular integration variable by $ha_{0,i}$; its envelope moments are the
same as in \cref{v16:thm:preparation}. Its momentum mean is exactly
$\hbar k_{h,i}$, and
$|\hbar k_{h,i}-h^2\mathcal E_{0,i}|\le\hbar/2$.
The rounded mean has exactly zero incidence at every vertex. Thus there
is no new deterministic term in the Gauss angular kinetic energy.
The added electric error is at most $C\hbar^2h^{-4}$, while the angular
variance contribution remains $C\hbar^2h^{-4}/\sigma_a^2$.
The shifted hopping depends on the translated envelopes, and the reference
shift has zero plaquette curl. Its hopping and magnetic moments are therefore
unchanged. The radial kinetic and density terms are the ones evaluated in
\cref{v16:thm:preparation}. Together they give
\[
 \langle\boldsymbol C_h(0)\rangle
       \le C(\hbar h^{-5}+\hbar^2h^{-4})\le C\hbar h^{-5}.
\]

For completeness the higher energy-moment argument also survives these
shifts with no additional inverse lattice power. Conjugating by the angular
phase changes $\Pi_e/h^2$ by the bounded number
$\hbar k_{h,i}/h^2=\mathcal E_{0,i}+O(\hbar h^{-2})$.
Its incidence is zero. The hopping density now has the additional bounded
reference coefficient $(e^{\ii ha_{0,i}}-1)/h$; the plaquette density has no
additional reference term. In the scaled Gaussian variables, each original
local energy density and each of the finitely many derivatives needed for
four density products is still a polynomially bounded function with uniform
coefficients. The derivative scales are exactly those listed in the proof of
\cref{v16:thm:preparation}. Radial cutoff and wrapped-image errors remain
exponentially small. Expanding the fourth energy moment gives $O(n^4)$
local products, each with weight $w^4$, and $wn=1$. Their uniform Gaussian
moment bounds prove $\|(\boldsymbol E+I)^2\psi_h^0\|\le C$.
All operations are performed on the same smooth initial form core; no future
reference field enters the preparation.
\end{proof}

\section{Local resolved evolution and nonzero electromagnetic sources}
\label{v17:sec:finite}

Use exactly the gauge-equivariant local tensor cutoff and commuting local
locked operations of Version 15 with the physical coefficients of Version 16.
The Hamiltonian, context supports, and uniform context probabilities are not
changed in this iteration. Only the prescribed initial preparation has the
new centers. The physical tick duration is still
$d=4\vartheta\epsilon/11$, $0<\vartheta\le1$. Let
$\phi_R=P_R\psi_h^0/\|P_R\psi_h^0\|$ and start the actual conditional
trajectory at $|\phi_R\rangle\langle\phi_R|$.

\begin{theorem}[A local quantum-to-classical branch with nonzero electric field]
\label[theorem]{v17:thm:finite}
On the common reference interval of \cref{v17:prop:reference}, set
\begin{equation}
 e_j=\tr\bigl(\rho_jP_R\boldsymbol C_h(jd)P_R\bigr),\qquad
 \mathfrak b_h=\hbar h^{-5}
 +\frac{h^{-12}}{R}(1+T/\hbar)^2
 +\frac{\epsilon h^{-9}R^3}{\hbar^2}.
 \label{v17:eq:budget}
\end{equation}
Retain $0<\hbar\le h^7$, the sufficiently large local cutoff,
$10h^{-3}\epsilon\le\tau_0$, and the tensor-approximation threshold of
\cref{v14:thm:Galerkin}. Then
\begin{align}
 \max_{j\le T/d}\E e_j&\le C_T\mathfrak b_h,\\
 \PP\{\max_{j\le T/d}e_j>u\}&\le C_T\mathfrak b_h/u\qquad(u>0).
 \label{v17:eq:maximal}
\end{align}
Every attained state satisfies all Gauss constraints exactly.
For $0<\mathfrak b_h\le1/2$,
\begin{equation}
 \E\max_{j\le T/d}e_j\le C_T\mathfrak b_h
     \left[1+\log\left(2+\frac{h^{-3}R}{\mathfrak b_h}\right)\right].
 \label{v17:eq:log}
\end{equation}
Constants depend on the fixed reference data and time interval, not on the
spatial or matrix dimensions, $\vartheta$, or the microscopic radii.
\end{theorem}
\begin{proof}
The uncompressed quantum flow has comparison error at most
$C_T\hbar h^{-5}$ by \cref{v17:thm:identity,v17:thm:preparation}.
The tensor Galerkin theorem, the new uniform fourth energy moment, and
$\boldsymbol C_h\preceq C_T(\boldsymbol E+I)$ give, for the finite
unitary reference $\psi_R(t)$,
\[
 \sup_{t\le T}\langle\psi_R(t),\boldsymbol C_h(t)\psi_R(t)\rangle
 \le C_T\left[\hbar h^{-5}+h^{-12}R^{-1}(1+T/\hbar)^2\right].
\]
The action is unchanged, so the spatial norm costs are exactly
\cref{v16:prop:costs}:
$A_R\le Ch^{-3}R$ and $\mathfrak C_h\le Ch^{-6}R^2/\hbar^2$.
The existing local resolved-trajectory theorem gives a nonnegative fidelity
error $\xi_j$ with mean and maximal probability budget
$C_T\epsilon\mathfrak C_h$. For any positive observable $C_R$ and the
rank-one reference projection $P_*(t)$,
\[
 C_R\preceq2P_*C_RP_*+2(I-P_*)C_R(I-P_*).
\]
Since $\|P_R\boldsymbol C_hP_R\|\le C_T A_R$, its expectation is at most
twice the preceding deterministic reference error plus $C_TA_R\xi_j$.
The mean bound and first-threshold probability bound follow, splitting off
thresholds already smaller than the deterministic error. This is precisely
the positive transfer established in Version 14, now with the new coupled
reference estimate; there is no per-tick union bound. Gauge commutation of
each actual outcome gives exact physical support throughout.
Finally $0\le e_j\le C_TA_R$; integration of
$\min(1,C_T\mathfrak b_h/u)$ proves \eqref{v17:eq:log}.
\end{proof}

All expectations below mean the original insertion before tensor compression,
evaluated in a state on its tensor subspace. In particular, a squared
compressed coordinate is not substituted for a compressed square.
Write
\[
 \mathsf E_{x,i}=\Pi_{x,i}/h^2,\quad
 \mathsf D_{x,i}=\frac{U_{x,i}Z_{x+e_i}-Z_x}{h},\quad
 \mathsf B_{x,ij}=\frac{e^{\ii(C\theta)_{x,ij}}-1}{h^2}\quad(i<j),
 \qquad \mathsf P_{x,a}=p_{x,a}/w.
\]
For an actual physical density matrix, the radial momentum mean is understood
as the real expectation of $-\ii\hbar w^{-1}\partial_{r_x}$ in the radial
form domain, not as an assertion of a self-adjoint half-line momentum.

\begin{proposition}[Relational fields, current, and independent metric insertions]
\label[proposition]{v17:prop:sources}
At every attained record with comparison $e_j$, the mean radius and radial
momentum density obey
\begin{equation}
 \|\bar r_j-r_h\|_{H_h^1}^2+
 \|\bar v_{r,j}-v_h\|_{2,h}^2\le C_Te_j.
 \label{v17:eq:radial-readers}
\end{equation}
The mass-weighted second moments of $\mathsf E-\mathcal E_h$, of
$\mathsf B$, and of the shifted scalar gradients are at most $C_Te_j$.
For the actual gauge-invariant link current
\[
 \mathsf j_{x,i}=h^{-1}\operatorname{Im}(\overline Z_xU_{x,i}Z_{x+e_i}),
 \qquad j_{h,i}^*=r_h^2\sin(ha_{h,i})/h,
\]
one has
\begin{equation}
 w\sum_{x,i}|\tr(\rho_j\mathsf j_{x,i})-j_{h,i}^*|
                                      \le C_T(\sqrt{e_j}+e_j).
 \label{v17:eq:current}
\end{equation}
The energy density and diagonal spatial metric sources of the same action,
as defined below, have the identical weighted $L^1$ error bound relative to
their lattice reference values. Their continuum limits are
\begin{align}
 \rho_*={}&\tfrac12\dot r^2+\tfrac12r^2|a|^2+V(r)
                                      +\tfrac12|\dot a|^2,\\
 T_{ii,*}={}&r^2a_i^2+\tfrac12\dot r^2-\tfrac12r^2|a|^2-V(r)
                        +\tfrac12|\dot a|^2-\dot a_i^2.
 \label{v17:eq:stress-limit}
\end{align}
Potential-parameter insertions satisfy the same bounds. Thus the electric
energy and its generally anisotropic stress do not vanish in this limit.
\end{proposition}
\begin{proof}
The pointwise inequalities
$||Z_x|-r_h|\le |D_x|/\rho_*$ and
$||Z_{x+e_i}|-|Z_x||\le h|W_{x,i}|$ give the radius bound by Jensen.
The radial decomposition in \cref{v17:cor:charge} controls the quadratic
form of $-\ii\hbar w^{-1}\partial_{r_x}-\kappa_hr_x$.
Cauchy--Schwarz for that form controls its real mean; combining with the
radius bound and $\kappa_hr_h=v_h$ proves the momentum statement.
The electric, magnetic, and gradient assertions are direct positive
summands of \eqref{v17:eq:cost}.

The exact identity
\begin{equation}
 \mathsf D_{x,i}=e^{\ii\delta_i}W_{x,i}+\mu_i^hZ_x,
 \qquad \mu_i^h=(e^{\ii\delta_i}-1)/h.
 \label{v17:eq:gradient-decomp}
\end{equation}
has $|\mu_i^h|\le|a_i|$. Similarly
$\mathsf P_{x,a}=V_{x,a}+\kappa_hq_{x,a}$.
The current is
$\mathsf j_{x,i}=(\sin\delta_i/h)|Z_x|^2+
\operatorname{Im}(e^{\ii\delta_i}\overline Z_xW_{x,i})$.
The density error is bounded by the quartic square, while
$w\sum\tr\rho_j|Z|^2\le C_T+C\sqrt{e_j}$. Weighted Cauchy--Schwarz
bounds the last product by
$C_T\sqrt{e_j}(1+\sqrt{e_j})^{1/2}\le C_T(\sqrt{e_j}+e_j)$.
This proves \eqref{v17:eq:current} without a pointwise radius upper bound.

To specify the independent sources, introduce lapse $L_x$ and diagonal
spatial metric $\gamma_{x,i}>0$, with $J_x=(\prod_i\gamma_{x,i})^{1/2}$,
and hold the original canonical variables fixed. Define
\begin{equation}
 \begin{split}
 E(L,\gamma)=w\sum_x L_x\biggl[&
 \frac{\sum_a\mathsf P_{x,a}^2}{2J_x}
 +\frac{J_x}{2}\sum_i\gamma_{x,i}^{-1}|\mathsf D_{x,i}|^2
 +J_xV(|Z_x|)\\
 &+\frac1{2J_x}\sum_i\gamma_{x,i}\mathsf E_{x,i}^2
 +\frac{J_x}{2}\sum_{i<k}\gamma_{x,i}^{-1}\gamma_{x,k}^{-1}
                                       |\mathsf B_{x,ik}|^2\biggr].
 \end{split}
 \label{v17:eq:metric-action}
\end{equation}
At $L=1,\gamma_i=1$ this is exactly $\boldsymbol E$, including the compact
magnetic cosine. Its lapse derivative is $w\mathsf T_{00,x}$, with
\[
 \mathsf T_{00,x}=\tfrac12\sum_a\mathsf P_{x,a}^2
 +\tfrac12\sum_i|\mathsf D_{x,i}|^2+V(|Z_x|)
 +\tfrac12\sum_i\mathsf E_{x,i}^2
 +\tfrac12\sum_{i<k}|\mathsf B_{x,ik}|^2.
\]
Its independent diagonal metric derivative is $-w\mathsf T_{ii,x}/2$, where
\begin{align*}
 \mathsf T_{ii,x}={}&|\mathsf D_{x,i}|^2
 +\tfrac12\sum_a\mathsf P_{x,a}^2
 -\tfrac12\sum_k|\mathsf D_{x,k}|^2-V(|Z_x|)\\
 &-\mathsf E_{x,i}^2+\tfrac12\sum_k\mathsf E_{x,k}^2
 +\sum_{k\ne i}|\mathsf B_{x,ik}|^2
                    -\tfrac12\sum_{k<\ell}|\mathsf B_{x,k\ell}|^2.
\end{align*}
For $k<i$ the square of $\mathsf B_{x,ik}$ denotes that of the corresponding
oriented $ki$ face. This is an explicitly declared diagonal-metric extension;
no unproved uniqueness of arbitrary off-diagonal lattice metric couplings is
inferred from the flat action. In differentiating $P_RE(L,\gamma)P_R$ the
regulator is held fixed at the evaluated physical parameter, as in Versions
13--16.

Expand each scalar square using \eqref{v17:eq:gradient-decomp} and the momentum
identity. Its error consists of a reference-bounded multiple of $D_x$,
a quadratic error, or a mixed term bounded by density-matrix Cauchy--Schwarz.
Their spatial sums are at most $C_T(\sqrt{e_j}+e_j)$. Expand electric squares
around $\mathcal E_h$ in the same way. All magnetic squares are already
bounded by $C e_j$. Finally
\[
 V(|Z|)-V(r_h)=\frac\lambda2(r_h^2-b)D_x+\frac\lambda4D_x^2.
\]
This proves the source bound. Differentiation in $\lambda,b$ gives polynomials
in $|Z|^2$ of degree at most two and has the same estimate.
At the reference,
$|\mathsf D_i^*|^2=2r_h^2(1-\cos(ha_{h,i}))/h^2$ and
$\mathsf E_i^*=\mathcal E_{h,i}$, while the magnetic square is zero.
\Cref{v17:prop:reference} and the trigonometric remainders give
\eqref{v17:eq:stress-limit}, with deterministic $O(h^2)$ error on the
reference interval. The stipulated nonzero electric reference makes its
energy a nonzero retained contribution.
\end{proof}

Consequently a fixed threshold $u\le1$ for any of these integrated source
errors has maximal-in-time failure probability at most
$C_T\mathfrak b_h/u^2$, apart from its deterministic reference-discretization
error. Fixed-time expected source errors are at most
$C_T(\sqrt{\mathfrak b_h}+\mathfrak b_h)$. These follow directly from
\cref{v17:thm:finite}; they do not assume temporal independence.

\section{Explicit scale, physical interpretation, and the next remaining term}
\label{v17:sec:ledger}

\begin{corollary}[A nonzero-field-strength limit on local physical records]
\label[corollary]{v17:cor:scales}
Take
\begin{equation}
 \hbar_h=h^7,\qquad R_h=h^{-28},\qquad\epsilon_h=h^{109},\qquad
 d_h=\frac{4\vartheta_h}{11}h^{109},\quad 0<\vartheta_h\le1.
 \label{v17:eq:scales}
\end{equation}
Then $\mathfrak b_h=O(h^2)$ and
$\E\max_j e_j=O(h^2(1+|\log h|))$.
The unchanged local resolved instrument, initialized by \eqref{v17:eq:prep},
has exact Gauss support at every record. Its relational scalar means converge
in spatial $H^1\times L^2$, its electric means converge in $L^2$, and its
magnetic energy tends to zero, uniformly in probability on the reference
interval. Its lapse, diagonal spatial metric, potential, and current
insertions converge to those of \eqref{v17:eq:limit-ODE}, including the
nonzero electric energy and anisotropic pressure in
\eqref{v17:eq:stress-limit}. No candidate-radius floor is imposed.
\end{corollary}
\begin{proof}
The three powers in \eqref{v17:eq:budget} are
$h^{7-5}$, $h^{-12+28-14}$ and $h^{109-9-84-14}$, all $h^2$.
The local-dimension and tensor thresholds are the same as Version 16 and
hold for small $h$ by \cref{v17:thm:preparation}. Also
$10h^{-3}\epsilon_h=10h^{106}\to0$. Since $A_R\le Ch^{-31}$,
\eqref{v17:eq:log} gives the stated logarithm. Apply
\cref{v17:prop:sources} and the usual cellwise interpolation bounds already
used in this lineage. The reference radius is spatially constant, so its
lattice/continuum discrepancy is just its $O(h^2)$ ODE error.
\end{proof}

\paragraph{What is now genuinely nonzero.}
The reference solves both the homogeneous scalar equation and the Maxwell
evolution, with conserved total scalar--electric energy
\[
 \tfrac12\dot r^2+\tfrac12|\dot a|^2+\tfrac12r^2|a|^2+V(r).
\]
The scalar sector has energy derivative $r^2a\cdot\dot a$, and the electric
sector has its negative; differentiating the displayed expression verifies
the cancellation. Thus the field is not a test electric background prescribed
independently of the matter. It is an interacting matter--electric reference
selected by finite initial data. The electromagnetic field strength is
nonzero even though the spatial magnetic curvature is zero. This distinction
prevents calling the branch gauge-flat simply because every reference spatial
plaquette is the identity. The generally nonzero total stress still does not
solve a flat Einstein metric equation.

\paragraph{What remains outside the theorem.}
The reference is homogeneous, not a general inhomogeneous gauge--Higgs solution.
Its magnetic field is zero. The candidate has no such spatial restriction,
but concentration drives it to this particular reference. The general
inhomogeneous extension must retain spatially varying reference ratios and
nonzero reference plaquette curvature. The kinetic Weyl lemma shows that,
for a comparison of momentum degree at most two in these \emph{unreduced}
Cartesian/link coordinates, a new Moyal remainder is not automatically the
obstacle. The remaining work is the classical compact-link cancellation and
its positive operator-form bound: coefficient commutators, reference magnetic
transport, and force estimates must be controlled uniformly in $h$.
Quantizing a proposed positive classical inequality without an explicit square
argument would not supply that bound.

The primitive source identification has not changed. The accepted operators,
context scheduler, preparation, and semiclassical scaling in this positive
realization are declared, not determined by bare renewal probabilities.
Non-Abelian color, chiral matter, Einstein constraints, and metric backreaction
remain separate. The useful advance is a same-action, locally implemented,
exactly Gauss-preserving quantum branch with \emph{nonzero} limiting electric
curvature and its independently varied metric sources, rather than another
zero-field scalar estimate.

\paragraph{Verification and preservation.}
Every byte of Version 16 before its final end-document command is retained.
New labels have prefix \texttt{v17:}. The diagnostic script checks the mixed
product identity symbolically; the full compact-link Poisson identity on
$3^3$, $4^3$, and $5^3$ lattices, including exact candidate zeros; and the
quantum identity on an interior core of a 576-dimensional oscillator/circle
test. The latter is a finite basis diagnostic, not an exact finite-dimensional
canonical commutation relation. It also checks reference energy conservation,
the $O(h^2)$ reference discretization rate, the strict momentum-degree boundary
of the Weyl lemma, and the simultaneous scale powers. These checks complement
the written proofs and are not a machine-checked formalization or an
Einstein--SM simulation.

\begingroup
\renewcommand{\refname}{Additional reference for Version 17}
\begin{thebibliography}{9}
\small\raggedright
\bibitem{v17:Robert}
D.~Robert, \emph{When Poisson and Moyal brackets are equal?}, 2022 preprint.
\href{https://arxiv.org/abs/2207.08798}{arXiv:2207.08798}.
Background for the distinction between exact Moyal and Poisson brackets.
The restricted momentum-degree calculation and its lattice use are proved
in this body, not inferred from a general nonlinear exact-quantization claim.
\end{thebibliography}
\endgroup
\addtocontents{toc}{\protect\endgroup}
% END IMMUTABLE VERSION 17 BODY
% APPEND VERSION 18 HERE, BEFORE THE FINAL END-DOCUMENT COMMAND.


% BEGIN IMMUTABLE VERSION 18 BODY
\clearpage
\begin{center}
{\Large\bfseries Version 18: Reference-weighted edges before magnetic closure}\\[0.5em]
{\large An inhomogeneous compact-link identity, nonzero magnetic curvature,
and a local exactly constrained classical field limit}
\end{center}
\makeatletter
\addtocontents{toc}{\protect\begingroup}
\addtocontents{toc}{\protect\makeatletter}
\addtocontents{toc}{\protect\def\protect\l@section{\protect\bfseries\protect\@dottedtocline{1}{0em}{2.5em}}}
\makeatother
\addcontentsline{toc}{part}{Version 18: Inhomogeneous reference-weighted gauge--Higgs closure}

\section{Audit: the missing estimate is spatial and magnetic}
\label{v18:sec:audit}

The limiting magnetic field in \cref{v17:cor:scales} vanishes. Merely replacing
its homogeneous coefficients by functions of position would not prove the
same estimate: reference-amplitude differences change the discrete adjoint,
and a shifted magnetic energy no longer has the original magnetic force.
This body computes both terms on the unchanged compact Abelian action.
The candidate field need not have a positive radius or small link angles.
Only the smooth comparison solution has a positive radius.

\paragraph{Inherited inputs and nonduplication.}
The energy and mass normalization are \cref{v16:eq:energy}; Gauss reduction,
the equivariant tensor regulator, and the local gauge-commuting instruments
are \cref{v15:thm:reduction,v15:thm:cutoff,v15:thm:locked}.
The quantum promotion uses \cref{v17:lem:Weyl}, not a newly asserted exact
quantization principle. Tensor approximation and the resolved renewal-clock
comparison are \cref{v14:thm:Galerkin,v15:thm:source}; their proofs are not
repeated. The classical high-regularity local matter existence result in
\cite{CC}, at \nolinkurl{v2:thm:local}, supplies a comparison slab when applied
to this Abelian scalar subcase. We do not claim another local-existence
or Fourier--midpoint theorem. Targeted file searches and inspection of the
relevant final sections did not locate the reference-weighted compact-link
identity below in the supplied material. This is a limited overlap audit,
not an exhaustive verification or a bibliographic priority claim.

\paragraph{Scope.}
The result is for the same constructed compact Abelian gauge--Higgs model,
on a prescribed flat three-torus. It now allows inhomogeneous scalar fields,
nonzero electric and magnetic curvature, and nonzero local charge satisfying
Gauss's law. The proof may compare with samples of a smooth continuum solution;
the source uses only its initial data, the action, and the declared local
instrument, never future comparison values. The source-selection rule,
semiclassical normalization, and preparation are still declared, not derived
from bare renewal probabilities. Einstein backreaction, non-Abelian color,
and chiral matter are not conclusions.

\paragraph{Relation to established methods.}
Smooth gauge--scalar Cauchy theory, modulated energies, and semiclassical
comparison are established subjects; see the earlier references and
\cite{v18:Pecher} for temporal-gauge Cauchy theory. That external result is
background, not a replacement for a quartic compact-link estimate here.
The new work is the explicit lattice identity, its positive operator-form
bound, the reference-consistency and exactly physical preparation estimates,
and their composition on the original resolved renewal clock.

\section{The reference-weighted edge and the complete error bank}
\label{v18:sec:bank}

Use the cubic lattice of \cref{v16:eq:physical-coefficients}, with $h=N^{-1}$,
$N\ge3$, $w=h^3$. Write $e:x\to y$ for a positive edge. Vertex--edge
incidence $B$ has $+1$ at the tail and $-1$ at the head; face incidence is $C$,
so $CB^{\mathsf T}=0$. The uncompressed physical energy is
\begin{equation}
 \begin{split}
 \boldsymbol E={}&\sum_{x,a}\frac{p_{x,a}^2}{2w}
 +\frac h2\sum_{e:x\to y}|U_eZ_y-Z_x|^2
 +\frac{\lambda w}{4}\sum_x(|Z_x|^2-b)^2\\
 &+\sum_e\frac{\Pi_e^2}{2h}
 +\frac1h\sum_f(1-\cos(C\theta)_f),\qquad U_e=e^{\ii\theta_e}.
 \end{split}
 \label{v18:eq:energy}
\end{equation}
Here $Z_x=q_{x,1}+\ii q_{x,2}$, $p_{x,a}=-\ii\hbar\partial_{q_{x,a}}$,
$\Pi_e=-\ii\hbar\partial_{\theta_e}$ and
$\mathscr D=\ii[\boldsymbol E,\,\cdot\,]/\hbar$.
We write $p_x=p_{x,1}+\ii p_{x,2}$ only as a notation for two real
components. All physical products below are formed before regulator compression.

Let $z_x^*(t)\ne0$ and $\theta_e^*(t)$ be smooth reference data. Define
\begin{equation}
 r_x=|z_x^*|,\qquad \mu_x=\dot z_x^*/z_x^*=\kappa_x+\ii\omega_x,
 \qquad \Pi_e^*=h\dot\theta_e^*,\qquad E_e^*=\Pi_e^*/h^2.
 \label{v18:eq:reference}
\end{equation}
Here $U_e^*=e^{\ii\theta_e^*}$ is a reference link; a star on
$\widetilde U_e$ below denotes its adjoint instead. The symbol $\omega_x$
is a reference phase velocity, not the quantum parameter.
The reference may have residuals:
\begin{align}
 f_x^*&=\ddot z_x^*-\Delta_{U^*}z_x^*
                         +\lambda(r_x^2-b)z_x^*,\\
 g_e^*&=\dot E_e^*+h^{-1}\operatorname{Im}
                    (\overline{z_x^*}U_e^*z_y^*)
                         +h^{-3}(C^{\mathsf T}\sin\beta^*)_e,
 \qquad \beta_f^*=(C\theta^*)_f,
 \label{v18:eq:ref-residuals}
\end{align}
where $\Delta_U$ is the usual symmetric link Laplacian. No discrete
reference field equation is assumed when these residuals are retained.

The invariant reference edge phase and the radius weight are
\begin{equation}
 e^{\ii\delta_e}=U_e^*\frac{z_y^*/r_y}{z_x^*/r_x},\qquad
 \gamma_e=(r_x/r_y)^{1/2},\qquad
 \widetilde U_e=e^{-\ii\delta_e}U_e.
 \label{v18:eq:weights}
\end{equation}
Only the circle value of $\delta_e$ is used. Local differentiable lifts in
time give
\begin{equation}
 \dot\delta_e=hE_e^*+\omega_y-\omega_x,
 \qquad e^{\ii(C\delta)_f}=e^{\ii\beta_f^*}.
 \label{v18:eq:reference-transport}
\end{equation}
Consequently no global principal argument, removal of periods, or choice of
the candidate's phase is made.

Define
\begin{equation}
 \begin{gathered}
 v_x=p_x/w-\mu_x Z_x,\qquad D_x=|Z_x|^2-r_x^2,\\
 W_e=h^{-1}(\gamma_e\widetilde U_eZ_y-\gamma_e^{-1}Z_x),\\
 e_e=\Pi_e/h^2-E_e^*,\qquad
 \varphi_f=(C\theta)_f-\beta_f^*,\qquad
 J_e=\overline Z_x\widetilde U_eZ_y.
 \end{gathered}
 \label{v18:eq:errors}
\end{equation}
The real components of $v_x$ are self-adjoint, and its kinetic square always
means $v_{x,1}^2+v_{x,2}^2$. In particular it does \emph{not} mean $v_x^*v_x$:
when $\omega_x\ne0$, the real components need not commute. This convention
is essential for the exact Weyl promotion and the retained angular charge.

The comparison is
\begin{equation}
 \begin{gathered}
 \boldsymbol{\mathcal C}_h=T+W+V+E_d+M_d,\\
 T=\frac w2\sum_{x,a}v_{x,a}^2,\qquad
 W=\frac w2\sum_e|W_e|^2,\qquad
 V=\frac{\lambda w}{4}\sum_xD_x^2,\\
 E_d=\frac w2\sum_e e_e^2,\qquad
 M_d=\frac1h\sum_f(1-\cos\varphi_f).
 \end{gathered}
 \label{v18:eq:cost}
\end{equation}
Every term is gauge invariant and nonnegative on its form domain. Reference
weights occur only in this comparison, not in the source Hamiltonian.

\begin{lemma}[Weighted mixed-product identity]
\label[lemma]{v18:lem:product}
At every edge, including at candidate zeros,
\begin{equation}
 \boxed{|J_e-r_xr_y|^2=D_xD_y+r_xr_yh^2|W_e|^2.}
 \label{v18:eq:product}
\end{equation}
Let $\rho\le r_x\le R_*$, $Q_*=R_*/\rho$, and assume
$|1-e^{\ii\delta_e}|\le A_*h$. Put
\begin{align}
 F_x={}&\frac1h\left[
 \sum_{e:x\to y}(e^{\ii\delta_e}-1)\gamma_e^{-1}W_e
 -\sum_{e:y\to x}(e^{-\ii\delta_e}-1)\gamma_e\widetilde U_e^*W_e\right],\\
 G_e={}&h^{-1}\left[(1-\cos\delta_e)\operatorname{Im}J_e
                   -\sin\delta_e(\operatorname{Re}J_e-r_xr_y)\right].
 \label{v18:eq:forces}
\end{align}
With $\|\cdot\|_{2,h}^2=w\sum|\cdot|^2$ on the indicated sites or edges,
\begin{align}
 \|F\|_{2,h}^2&\le12Q_*A_*^2\|W_e\|_{2,h}^2,\\
 \|G\|_{2,h}^2&\le A_*^2\left(3\|D\|_{2,h}^2
                                      +R_*^2h^2\|W_e\|_{2,h}^2\right).
 \label{v18:eq:force-bounds}
\end{align}
These are pointwise multiplication inequalities and therefore operator-form
inequalities as well.
\end{lemma}
\begin{proof}
Expand both sides of
\[
 |\bar zt-ab|^2=(|z|^2-a^2)(|t|^2-b^2)
                   +ab|\sqrt{a/b}\,t-\sqrt{b/a}\,z|^2
\]
for $a,b>0$. Set $a=r_x,b=r_y,t=\widetilde U_eZ_y$ to obtain
\eqref{v18:eq:product}. Each vertex has six incident edges; the squared norm
of their sum is at most six times the sum of the squared norms.
Each edge is counted twice and $\gamma_e^{\pm2}\le Q_*$. This proves the
first force bound. For the second,
$\sqrt{(1-\cos\delta)^2+\sin^2\delta}=|1-e^{\ii\delta}|$ gives
$|G_e|^2\le A_*^2|J_e-r_xr_y|^2$.
Use $D_xD_y\le(D_x^2+D_y^2)/2$ and count six incidences at a vertex.
No candidate radius is inverted and no estimate of $|Z_x|$ is inserted.
\end{proof}

\section{The exact inhomogeneous identity and its magnetic remainder}
\label{v18:sec:identity}

Define the real magnetic multiplication term
\begin{equation}
 \mathcal M_f=(1-\cos\beta_f^*)\sin\varphi_f
                        +\sin\beta_f^*(1-\cos\varphi_f).
 \label{v18:eq:mag-remainder}
\end{equation}
For self-adjoint real components write
$\operatorname{Sym}(A,B)=(AB+BA)/2$, and sum this over two real components
for $\operatorname{Sym}_{\R}(v,F)$.

\begin{theorem}[Exact inhomogeneous compact-link quantum comparison]
\label[theorem]{v18:thm:identity}
On the common Cartesian/periodic smooth core the unchanged Hamiltonian
\eqref{v18:eq:energy} satisfies the exact quadratic-form identity
\begin{equation}
\boxed{\begin{aligned}
 (\partial_t+\mathscr D)\boldsymbol{\mathcal C}_h
 ={}&-w\sum_x\kappa_x\sum_a v_{x,a}^2
 +\frac w2\sum_{e:x\to y}(\kappa_x+\kappa_y)|W_e|^2
 +\lambda w\sum_x\kappa_xD_x^2\\
 &+w\sum_x\operatorname{Sym}_{\R}
                  \left(v_x,F_x-\frac{f_x^*}{z_x^*}Z_x\right)
 +w\sum_e\operatorname{Sym}(e_e,G_e-g_e^*)\\
 &+\sum_f\operatorname{Sym}\bigl((Ce)_f,\mathcal M_f\bigr).
\end{aligned}}
\label{v18:eq:exact}
\end{equation}
For classical fields the same identity has ordinary products.
Assume the radius and edge bounds of \cref{v18:lem:product}, bounded
$\kappa$, and $|1-e^{\ii\beta_f^*}|\le B_*h^2$. If
$\|f^*/z^*\|_\infty\le1$, then
\begin{equation}
 (\partial_t+\mathscr D)\boldsymbol{\mathcal C}_h
 \preceq K\boldsymbol{\mathcal C}_h+C\mathfrak d_h I,
 \qquad
 \mathfrak d_h=\|f^*/z^*\|_\infty^2+\|g^*\|_{2,h}^2.
 \label{v18:eq:energy-bound}
\end{equation}
The constants depend only on the displayed reference bounds and $\lambda$,
not on $h$, $\hbar$, the number of sites, candidate radii, or candidate
plaquette angles. If in addition $\mu,E^*$ and the first reference spatial
differences are uniformly bounded, then
\begin{equation}
 0\preceq\boldsymbol{\mathcal C}_h(t)\preceq C_T(\boldsymbol E+I).
 \label{v18:eq:form-domination}
\end{equation}
\end{theorem}
\begin{proof}
First compute classically, in temporal gauge. Write
\[
 (L_rZ)_x=h^{-2}\left[
 \sum_{e:x\to y}(\widetilde U_eZ_y-r_yZ_x/r_x)
 +\sum_{e:y\to x}(\widetilde U_e^*Z_y-r_yZ_x/r_x)\right].
\]
This is the negative first variation, divided by $w$, of the weighted
hopping energy $W$. Its off-diagonal coefficient is still the actual
transport; its diagonal coefficient is the reference radius ratio.
The reference equation gives
\[
 \dot\mu_x+\mu_x^2
 =\Delta_{U^*}z_x^*/z_x^*-\lambda(r_x^2-b)+f_x^*/z_x^*.
\]
Insert it in the actual scalar equation to obtain
\begin{align*}
 \dot Z_x&=v_x+\mu_xZ_x,\\
 \dot v_x&=(L_rZ)_x-\lambda D_xZ_x-\mu_xv_x
                       +F_x-(f_x^*/z_x^*)Z_x,\\
 \dot D_x&=2\operatorname{Re}(\overline Z_xv_x)+2\kappa_xD_x.
\end{align*}
For the force formula, an outgoing term after subtracting the reference
zero-order term is
$h^{-2}e^{\ii\delta_e}(\widetilde U_eZ_y-r_yZ_x/r_x)$.
Split its coefficient into $1+(e^{\ii\delta_e}-1)$.
The incoming term is the conjugate-oriented expression. Since
\[
 \widetilde U_eZ_y-r_yZ_x/r_x=h\gamma_e^{-1}W_e,
 \quad
 \widetilde U_e^*Z_x-r_xZ_y/r_y=-h\gamma_e\widetilde U_e^*W_e,
\]
these are precisely $L_rZ$ and \eqref{v18:eq:forces}.

The time derivatives of the weight and shifted link are
\[
 \dot\gamma_e/\gamma_e=(\kappa_x-\kappa_y)/2,
 \qquad \dot{\widetilde U}_e
   =\ii(he_e+\omega_x-\omega_y)\widetilde U_e.
\]
They yield the exact edge identity
\begin{equation}
 \dot W_e=\frac{\gamma_e\widetilde U_ev_y-\gamma_e^{-1}v_x}{h}
 +\left(\frac{\kappa_x+\kappa_y}{2}+\ii\omega_x\right)W_e
 +\ii e_e\gamma_e\widetilde U_eZ_y.
 \label{v18:eq:W-flow}
\end{equation}
The terms with $L_r$ cancel the first term here after summing over both
endpoints. This is an exact discrete adjoint cancellation, not an estimate
of the $h^{-2}$ Laplacian. Imaginary multiples of the squared scalar or
edge norms have zero real part. The density-quartic derivative cancels
$-\lambda D_x\operatorname{Re}(\overline v_xZ_x)$.
The terms left from these three parts are the first line of
\eqref{v18:eq:exact} and its scalar force pairing.

The electric error equation is
\[
 \dot e_e=-h^{-1}\operatorname{Im}J_e+G_e-g_e^*
 -h^{-3}\{C^{\mathsf T}[\sin(\beta^*+\varphi)-\sin\beta^*]\}_e.
\]
The edge--electric exchange cancels because
\[
 \operatorname{Re}(\overline W_e\,\ii e_e\gamma_e\widetilde U_eZ_y)
                         =e_e\operatorname{Im}J_e/h.
\]
Also $\dot\varphi=hCe$: the reference phase-velocity incidence vanishes
under $C$. The derivative of $M_d$ is $\sum_f(Ce)_f\sin\varphi_f$.
Combining it with the electric magnetic force leaves
\[
 \sum_f(Ce)_f\{\sin\varphi_f-\sin(\beta_f^*+\varphi_f)
                                      +\sin\beta_f^*\},
\]
which is \eqref{v18:eq:mag-remainder}. This proves the full classical
identity, including both reference residuals.

For quantum promotion, the Hamiltonian has constant quadratic kinetic
coefficients in the original Cartesian and link momenta. The comparison
has momentum degree two; its real linear momentum shifts have affine
Cartesian coefficients or constant link coefficients. Its Weyl quantization
is exactly the positive real-component squares in \eqref{v18:eq:cost}.
Apply \cref{v17:lem:Weyl}. Linear momentum pairings become the displayed
symmetric products. The spatial weights and all trigonometric terms are
multiplication operators. This gives \eqref{v18:eq:exact} with no omitted
ordering remainder. Positivity below is proved independently of this
symbol calculation.

The scalar and electric force terms are bounded on every core vector by
Cauchy--Schwarz, $2ab\le a^2+b^2$, and \eqref{v18:eq:force-bounds}.
The coefficient of the first line is at most $4\|\kappa\|_\infty$
times the comparison. For the magnetic part, pointwise
\[
 |\mathcal M_f|\le2B_*h^2|e^{\ii\varphi_f}-1|.
\]
Every face has four edges and every edge belongs to four faces, so
$\|C\|_{\ell^2\to\ell^2}\le4$. For a test vector $u$,
\begin{align*}
 \left|\sum_f\operatorname{Re}\langle(Ce)_fu,\mathcal M_fu\rangle\right|
 &\le16B_*h\,
       \langle u,E_du\rangle^{1/2}\langle u,M_du\rangle^{1/2}\\
 &\le8B_*h\langle u,(E_d+M_d)u\rangle.
\end{align*}
Here the powers follow from $w=h^3$ and
$M_d=(2h)^{-1}\sum_f|e^{\ii\varphi_f}-1|^2$.
This bound holds for arbitrary candidate plaquettes. The small quantity is
the \emph{reference} plaquette, not a restriction of the candidate support.

Finally $w\sum_x|Z_x|^2\preceq C(I+V)$, since
$s\le R_*^2+1+(s-r_x^2)^2$ for $s\ge0$.
Young's inequality bounds the $f^*$ pairing by a kinetic square plus
$C\|f^*/z^*\|_\infty^2(I+V)$; the $g^*$ pairing is bounded by an electric
square plus $\|g^*\|_{2,h}^2I/2$. This proves \eqref{v18:eq:energy-bound}.
For \eqref{v18:eq:form-domination}, shifted momenta are bounded by twice
the original squares plus bounded reference multiples of $w\sum|Z|^2+I$.
For edges, $\gamma_e=1+O(h)$ and $e^{-\ii\delta_e}=1+O(h)$ show that
$W_e$ is a bounded multiple of $(U_eZ_y-Z_x)/h$ plus a bounded multiple
of $Z_x$. The shifted magnetic chord is bounded by the original chord
plus $B_*h^2$. Sum with the physical weights. The quartic controls the
remaining scalar square. These are form inequalities with uniform constants.
The inherited energy-spectral approximation argument extends the core
identity and integrated bound to the stated energy-form evolutions; those
auxiliary global projections are proof devices, not changes to the tensor
regulator used by the source.
\end{proof}

\begin{corollary}[The inhomogeneous Gauss contribution is retained]
\label[corollary]{v18:cor:charge}
On the physical space of \cref{v15:thm:reduction}, the scalar kinetic
comparison is
\begin{equation}
 T=\sum_x\left[
 \frac{(-\ii\hbar\partial_{r_c}-w\kappa_xr_c)^*
             (-\ii\hbar\partial_{r_c}-w\kappa_xr_c)}{2w}
 +\frac{((B\Pi)_x+w\omega_xr_c^2)^2}{2wr_c^2}\right].
 \label{v18:eq:charge}
\end{equation}
Here $r_c=|Z_x|$ is the candidate radius and the adjoint uses $r_c\dd r_c$.
Let $\mathcal G_x^*=(B\Pi^*)_x+wr_x^2\omega_x$. The second numerator equals
\begin{equation}
 (B(\Pi-\Pi^*))_x+w\omega_xD_x+\mathcal G_x^*.
 \label{v18:eq:charge-center}
\end{equation}
Thus neither local charge nor a reference Gauss residual is erased by the
comparison. No self-adjoint half-line momentum or inverse-radius bound is
asserted.
\end{corollary}
\begin{proof}
Resolve the real Cartesian vector $p_x-w(\kappa_x+\ii\omega_x)Z_x$
into radial and angular components in its positive quadratic form.
The angular component is $L_x/r_c-w\omega_xr_c$. Substitute the exact
Gauss relation $L_x=-(B\Pi)_x$ and expand around $\Pi^*,r_x^2$.
The inherited Friedrichs form convention extends the equality from the
punctured core. It is not a new boundary condition at zero.
\end{proof}

\section{Continuum samples have explicit force and Gauss consistency}
\label{v18:sec:consistency}

Let $(z,A,E)$ be a fixed smooth solution on $[0,T]\times\T^3$ in temporal
gauge $A_0=0$ of the equations in \cref{v16:eq:continuum}, with the same
quartic coefficient, and $|z|\ge\rho>0$. All derivatives through a fixed
sufficient order, for example six spatial and four temporal derivatives,
are bounded. Assume a global periodic potential $A$; nonzero local magnetic
curvature is allowed, but nontrivial line-bundle topology is not claimed.
At an edge $e=(x,i)$ set
\begin{equation}
 z_x^*=z(t,hx),\qquad
 \theta_{x,i}^*=\int_0^h A_i(t,hx+s e_i)\dd s,\qquad
 \Pi_{x,i}^*=h\int_0^h E_i(t,hx+s e_i)\dd s.
 \label{v18:eq:sampling}
\end{equation}
These specify comparison quantities, not a schedule of source updates.

\begin{proposition}[Uniform second-order reference consistency]
\label[proposition]{v18:prop:consistency}
The samples satisfy $\dot\theta_e^*=\Pi_e^*/h$, all uniform reference
bounds of \cref{v18:thm:identity}, and
\begin{equation}
 \max_x|f_x^*|+\max_e|g_e^*|
       +\max_x|\mathcal G_x^*|/w\le C_T h^2.
 \label{v18:eq:consistency}
\end{equation}
In particular $\mathfrak d_h\le C_T h^4$ and the reference magnetic
phase has $|1-e^{\ii\beta_f^*}|\le C_T h^2$.
\end{proposition}
\begin{proof}
For a line based at $x$ let
$F_i(s)=\exp(\ii\int_0^s A_i(x+t e_i)\dd t)z(x+s e_i)$.
Its derivatives at zero are $D_i^kz(x)$. The incoming parallel transport
is $F_i(-h)$, so the symmetric second difference gives
$\Delta_{U^*}z^*=\sum_iD_i^2z+O(h^2)$ uniformly. The scalar equation
proves the first bound.

At an edge midpoint, the hopping product equals
$\overline{F_i(-h/2)}F_i(h/2)$. Its imaginary part is odd in $h$,
and its linear coefficient is $\operatorname{Im}(\bar zD_i z)$.
Thus the reference current divided by $h$ is the midpoint current plus
$O(h^2)$. The reference electric density and its time derivative are
line averages, also equal to their midpoint values up to $O(h^2)$.

Discrete Stokes gives $\beta_f^*=\int_f F_{ij}$ exactly. Expanding this
face integral about its center gives
$h^2F_{ij}+h^4(\partial_i^2+\partial_j^2)F_{ij}/24+O(h^6)$.
Two incident opposite faces are centered symmetrically around the edge;
their oriented difference, divided by $h^3$, is the corresponding
$\partial_jF_{ij}$ at the edge midpoint plus $O(h^2)$.
The sine correction contributes at most $C h^6/h^3=O(h^3)$ per
incident face, since $\beta_f^*=O(h^2)$. Summing the fixed number of faces
therefore gives
$h^{-3}C^{\mathsf T}\sin\beta^*=(\operatorname{curl}B)_{
m midpoint}
+O(h^2)$. The Maxwell equation proves the $g^*$ bound with the stated
orientation convention.

Finally, at a vertex,
\[
 (B\Pi^*)_x/w
 =\sum_i h^{-2}\left[\int_0^h E_i(x+s e_i)\dd s
                         -\int_{-h}^0 E_i(x+s e_i)\dd s\right]
 =\operatorname{div}E(x)+O(h^2).
\]
Since $r_x^2\omega_x=\operatorname{Im}(\bar z\dot z)(x)$, continuum
Gauss law proves the last assertion in \eqref{v18:eq:consistency}.
Bounded covariant first derivatives and the radius margin give
$|1-e^{\ii\delta_e}|\le Ch$ and $|r_y-r_x|\le Ch$.
The remaining time and radius bounds follow from the assumed smooth slab.
All remainders are uniform integral Taylor remainders with a fixed number
of derivatives; no estimate of a candidate field enters them.
\end{proof}

\section{Exactly physical preparations for variable radius, phase velocity, and flux}
\label{v18:sec:preparation}

Only time-zero samples are used in the preparation. In the exact relational
coordinates of Version 15 choose radial cutoffs equal to one near $r_x(0)$,
with uniformly bounded derivatives and support inside a fixed positive-radius
interval. Set
\[
 \sigma_r^2=\hbar h^{-2},\qquad \sigma_a^2=\hbar h^{-1},\qquad
 0<\hbar\le h^7,\qquad
 k_e=\operatorname{round}(\Pi_e^*(0)/\hbar).
\]
Use normalized factors
\begin{align}
 f_x(r_c)&=Z_x^{-1}r_c^{-1/2}\chi_x(r_c)
 \exp\left[-\frac{(r_c-r_x(0))^2}{2\sigma_r^2}
            +\frac{\ii w\dot r_x(0)(r_c-r_x(0))}{\hbar}\right],\\
 g_e(a)&=Z_e^{-1}e^{\ii k_e(a-\delta_e(0))}
   \sum_{m\in\Z}\exp\left[-\frac{(a-\delta_e(0)-2\pi m)^2}{2\sigma_a^2}\right].
 \label{v18:eq:preparation}
\end{align}
Lift their product by the unchanged unitary $\mathcal U$ to obtain $\psi_h^0$
in the original Gauss subspace. The radial and edge factors are independent
in relational coordinates, not in the original Cartesian/link coordinates.

\begin{theorem}[Uniform preparation and retained Gauss mismatch]
\label[theorem]{v18:thm:preparation}
For the smooth initial samples above, every $G_x\psi_h^0=0$ exactly, and
\begin{equation}
 \langle\psi_h^0,\boldsymbol{\mathcal C}_h(0)\psi_h^0\rangle
       \le C(\hbar h^{-5}+h^4),\qquad
 \|(\boldsymbol E+I)^2\psi_h^0\|\le C.
 \label{v18:eq:prep-bounds}
\end{equation}
The constants are uniform in the number of sites. No all-time support
restriction on the candidate radii is imposed. In particular the generally
nonzero rounded electric divergence is not discarded.
\end{theorem}
\begin{proof}
The physical-coordinate unitary proves exact Gauss support. Its positive
radial/angular decomposition is \eqref{v18:eq:charge}. Gaussian integration
and periodic-image estimates give position variances $O(\sigma_r^2)$,
angle chord variances $O(\sigma_a^2)$, momentum variances
$O(\hbar^2/\sigma_r^2)$ and $O(\hbar^2/\sigma_a^2)$, and
$|\hbar k_e-\Pi_e^*|\le\hbar/2$. Constants are uniform because all radial
centers and supports lie in one positive compact interval.

There are six links at a vertex, so
$|B(\hbar k-\Pi^*)|\le3\hbar$. In the angular numerator
\eqref{v18:eq:charge-center}, the random momentum part has second moment
$O(\hbar^2/\sigma_a^2)$, the rounding part is $O(\hbar^2)$,
the radial-density contribution is $O(w^2\sigma_r^2)$, and the reference
Gauss residual is $O(w^2h^4)$. Since the initial radii are bounded away
from zero, summing its energy gives
\[
 C\left(\hbar^2h^{-6}/\sigma_a^2+\hbar^2h^{-6}
                                  +\sigma_r^2+h^4\right).
\]
This calculation is where continuum Gauss consistency is used. Requiring
zero divergence of each rounded electric center would be wrong for these
charged data.

The other comparison terms have the bounds
\begin{align*}
 T_{\rm radial}&:\ C(\hbar^2h^{-6}/\sigma_r^2+\sigma_r^2+\hbar^2h^{-6}),\\
 W&:\ Ch^{-2}(\sigma_r^2+\sigma_a^2),\\
 V&:\ C(\sigma_r^2+\sigma_r^4),\\
 E_d&:\ C(\hbar^2h^{-4}/\sigma_a^2+\hbar^2h^{-4}),\\
 M_d&:\ Ch^{-4}\sigma_a^2.
\end{align*}
For $W$ the reference-weighted difference vanishes at the two centers;
its coefficients and ratios are uniformly bounded. For $M_d$ the four
center angles sum to the reference plaquette modulo $2\pi$, so only
the centered envelope angles occur. Substitute the widths and
$\hbar\le h^7$ in this list and the angular bound. This proves the first
estimate in \eqref{v18:eq:prep-bounds}.

Here are the additional details needed for the fourth energy moment.
Transfer the radial measure to flat measure and retain its
$-\hbar^2/(8wr_c^2)$ term. Conjugate by the initial radial and electric
phases. The classical parts of each normalized local momentum are bounded:
$w^{-1}p_r^*=\dot r_x(0)$, $h^{-2}\hbar k_e=E_e^*+O(\hbar h^{-2})$,
and, critically,
\[
 w^{-1}(B\hbar k)_x=-r_x(0)^2\omega_x(0)+O(h^2+\hbar h^{-3}).
\]
Use centered variables $X_x=(r_c-r_x(0))/\sigma_r$ and
$Y_e=(a-\delta_e(0))/\sigma_a$ on central angular charts. The scaled
first-derivative coefficients are
$\hbar/(w\sigma_r)=\sqrt\hbar\,h^{-2}$,
$\hbar/(w\sigma_a)=\sqrt\hbar\,h^{-5/2}$ and
$\hbar/(h^2\sigma_a)=\sqrt\hbar\,h^{-3/2}$, all bounded.
The reference hopping divided by $h$ and reference face angle divided by
$h^2$ are bounded by the consistency proposition. Centered hopping and
magnetic factors have coefficients controlled by
$\sigma_r/h$, $\sigma_a/h$ and $\sigma_a/h^2$, also bounded.
Their derivatives of every fixed order needed for four energy-density
products are bounded by a fixed polynomial in the finitely many local
$X,Y$. Initial inverse radii and their such derivatives are bounded on the
radial cutoff support. Derivatives of radial cutoffs and overlaps of periodic
images are bounded by polynomial factors times
$e^{-c/\sigma_r^2}$ or $e^{-c/\sigma_a^2}$; these remain bounded after the
fixed inverse powers of $h,\hbar$ in the calculation.

The energy is $w$ times a sum of $O(h^{-3})$ local densities with bounded
support size. Any product of at most four has a uniform Gaussian integral,
including overlapping supports. Expanding its expectation gives at most
$C h^{-3j}$ terms with weight $w^j$, $j\le4$. Thus
$\langle(\boldsymbol E+I)^4\rangle\le C$. At fixed cutoff the preparation
is in all required operator domains, by its smooth radial support and
periodic smooth factors. This proves the second estimate rather than
assuming it. No future values of a reference field were used.
\end{proof}

\section{Local resolved histories converge to the inhomogeneous gauge--Higgs solution}
\label{v18:sec:finite}

Retain the source's local gauge-equivariant tensor cutoff $P_R$, gauge-commuting
six-outcome operators in each recorded local context, uniform context
selection, and protected-tick clock, with $d=4\vartheta\epsilon/11$.
There are $M=10h^{-3}$ action contexts. The finite initial vector is
$\phi_R=P_R\psi_h^0/\|P_R\psi_h^0\|$; start the actual resolved state
at $|\phi_R\rangle\langle\phi_R|$.
The action and all updates are the same as in Versions 15--17. Only the
initial centers now vary in space. In particular there is no reference-driven
time-dependent update kernel.

\begin{theorem}[Inhomogeneous local physical-record limit]
\label[theorem]{v18:thm:finite}
Let the reference be a smooth positive-radius continuum solution on the stated
slab and let its samples and preparation be as above. Under the inherited
large-$R$ tensor threshold and $10h^{-3}\epsilon\le\tau_0$, define
\begin{equation}
 e_j=\tr(\rho_jP_R\boldsymbol{\mathcal C}_h(jd)P_R),\qquad
 \mathfrak b_h=\hbar h^{-5}+h^4
 +\frac{h^{-12}}R(1+T/\hbar)^2
 +\frac{\epsilon h^{-9}R^3}{\hbar^2}.
 \label{v18:eq:budget}
\end{equation}
For $0<\hbar\le h^7$ and any initial clock phase,
\begin{align}
 \max_{j\le T/d}\E e_j&\le C_T\mathfrak b_h,\\
 \PP\{\max_{j\le T/d}e_j>u\}&\le C_T\mathfrak b_h/u\quad(u>0).
 \label{v18:eq:maximal}
\end{align}
Every attained state is in the exact Gauss subspace. Constants depend on the
regular reference slab, not on spatial or matrix dimension, $\vartheta$,
or candidate radii. For $0<\mathfrak b_h\le1/2$,
\begin{equation}
 \E\max_j e_j\le C_T\mathfrak b_h
       \left[1+\log\left(2+\frac{h^{-3}R}{\mathfrak b_h}\right)\right].
 \label{v18:eq:log}
\end{equation}
\end{theorem}
\begin{proof}
The uncompressed quantum evolution of $\psi_h^0$ satisfies
\[
 \sup_{t\le T}\langle\boldsymbol{\mathcal C}_h(t)\rangle
                    \le C_T(\hbar h^{-5}+h^4)
\]
by \cref{v18:thm:identity,v18:prop:consistency,v18:thm:preparation} and
integration of \eqref{v18:eq:energy-bound}. The form bound
\eqref{v18:eq:form-domination}, the uniform fourth energy moment, and the
inherited tensor Galerkin theorem give for the finite unitary reference
\[
 \sup_{t\le T}\langle\psi_R(t),\boldsymbol{\mathcal C}_h(t)\psi_R(t)\rangle
 \le C_T\left[\hbar h^{-5}+h^4+\frac{h^{-12}}R(1+T/\hbar)^2\right].
\]
The Hamiltonian and tensor regulator are unchanged, so their costs are still
$A_R\le Ch^{-3}R$ and $\mathfrak C_h\le Ch^{-6}R^2/\hbar^2$ from
\cref{v16:prop:costs}. The local resolved-source theorem supplies fidelity
error $\xi_j$ with mean and maximal probability bounds
$C_T\epsilon\mathfrak C_h$. For $C_R=P_R\boldsymbol{\mathcal C}_hP_R$
and the unitary-reference rank-one projector $P_*$, positivity gives
\[
 C_R\preceq2P_*C_RP_*+2(I-P_*)C_R(I-P_*),\qquad \|C_R\|\le C_T A_R.
\]
Consequently $e_j$ is at most twice the deterministic comparison bound plus
$C_T A_R\xi_j$. Apply the inherited mean and first-threshold estimates,
separating thresholds already below the deterministic bound. This proves
\eqref{v18:eq:maximal} without summing probabilities over ticks.
Gauge commutation of every actual outcome preserves all constraints exactly.
Finally $e_j\le C_TA_R$ deterministically; integration of the bound
$\min(1,C_T\mathfrak b_h/u)$ gives \eqref{v18:eq:log}.
No whole-generator norm enters the exponential comparison constant.
\end{proof}

\section{Magnetic fields, local charge, and the independently varied sources}
\label{v18:sec:sources}

Suppress fixed tensor projections inside expectations, but continue to form
all composite insertions on the original space. Define physical readers
\[
 \mathsf E_e=\Pi_e/h^2,\qquad
 \mathsf B_f=(e^{\ii(C\theta)_f}-1)/h^2,\qquad
 \mathsf D_e=(U_eZ_y-Z_x)/h,\qquad \mathsf P_x=p_x/w.
\]
For the magnetic reader the real and imaginary parts are two commuting
Hermitian multiplication observables. Its comparison value is
$\mathsf B_f^*=(e^{\ii\beta_f^*}-1)/h^2$; the imaginary part tends to
the oriented continuum curvature.

\begin{proposition}[Field topology and nonzero magnetic second moments]
\label[proposition]{v18:prop:readers}
For every physical conditional state with comparison value $e_j$,
\begin{align}
 \|\bar r_j-r\|_{H_h^1}^2+\|\bar v_{r,j}-\dot r\|_{2,h}^2&\le C_T e_j,\\
 w\sum_e\tr\rho_j(\mathsf E_e-E_e^*)^2
 +w\sum_f\tr\rho_j|\mathsf B_f-\mathsf B_f^*|^2&\le2e_j.
 \label{v18:eq:reader-bounds}
\end{align}
The mean radial momentum density is the real expectation of
$-\ii\hbar w^{-1}\partial_{r_c}$ in its radial form domain.
The corresponding conditional variances and the weighted density-quartic
and covariant-gradient errors are also controlled by $C_Te_j$.
No candidate phase choice or principal-angle magnetic reader is needed.
\end{proposition}
\begin{proof}
$||Z_x|-r_x|\le |D_x|/\rho$ controls the zero-order radius error.
Taking moduli in the edge difference gives
\[
 \left|\gamma_e|Z_y|-\gamma_e^{-1}|Z_x|\right|\le h|W_e|.
\]
For $a_x=|Z_x|-r_x$, this implies
$|a_y-a_x|/h\le\gamma_e^{-1}|W_e|
 +h^{-1}|r_y/r_x-1||a_x|$.
The reference first-difference bound controls the last coefficient.
Square, sum, and use Jensen for the means. The radial term in
\eqref{v18:eq:charge} controls the form error in momentum around
$\kappa_x|Z_x|$; its real mean and $\dot r_x=\kappa_xr_x$ give the
momentum estimate. The magnetic identity
\[
 \mathsf B_f-\mathsf B_f^*
       =e^{\ii\beta_f^*}(e^{\ii\varphi_f}-1)/h^2
\]
shows that its squared weighted norm is $2M_d$ exactly. Electric error
has squared norm $2E_d$. The variance assertions follow by expanding each
positive square around its mean. Radial derivative statements remain
quadratic-form statements, not an assertion of a self-adjoint half-line
momentum.
\end{proof}

\begin{proposition}[Currents and action-derived metric sources]
\label[proposition]{v18:prop:sources}
Let the finite spatial and temporal scalar currents be
\[
 \mathsf j_e=h^{-1}\operatorname{Im}(\overline Z_xU_eZ_y),\qquad
 \mathsf j_{0,x}=w^{-1}(q_{x,1}p_{x,2}-q_{x,2}p_{x,1}).
\]
Their reference values are
$j_e^*=r_xr_y\sin\delta_e/h$ and $j_{0,x}^*=r_x^2\omega_x$.
Their weighted $L^1$ errors are bounded by $C_T(\sqrt{e_j}+e_j)$.
The same bound holds for the independently varied lapse, three diagonal
spatial metrics, and potential parameters of \cref{v17:eq:metric-action},
now evaluated at these inhomogeneous reference samples. In particular it
includes the complete electric and magnetic contributions to those sources.
\end{proposition}
\begin{proof}
The exact decomposition at an outgoing edge is
\begin{equation}
 \mathsf D_e=e^{\ii\delta_e}\gamma_e^{-1}W_e+\mu_e^h Z_x,
 \qquad
 \mu_e^h=\frac{(r_y/r_x)e^{\ii\delta_e}-1}{h}.
 \label{v18:eq:physical-gradient}
\end{equation}
Its coefficient $\mu_e^h$ is uniformly bounded and equals
$(U_e^*z_y^*/z_x^*-1)/h$. Also
$\mathsf P_x=v_x+\mu_xZ_x$, interpreted in real components.
The spatial current is consequently a bounded reference coefficient times
$|Z_x|^2$ plus a mixed $Z_x,W_e$ term. The temporal current is
$\operatorname{Im}(\overline Z_xv_x)+\omega_x|Z_x|^2$; the real
Cartesian cross-component factors in its definition commute, so no
unrecorded current ordering is needed. Weighted Cauchy--Schwarz,
$w\sum\tr\rho_j|Z_x|^2\le C+C\sqrt{e_j}$, and the density square
prove the current estimates.

Differentiate the already declared finite metric action before compression,
holding its canonical variables and regulator fixed. Its densities are the
ones in \cref{v17:prop:sources}, but neither electric nor magnetic reference
terms are now zero. Expand each scalar square using
\eqref{v18:eq:physical-gradient}, each Cartesian kinetic square using
$\mathsf P=v+\mu Z$, and electric and magnetic squares around $E^*$ and
$\mathsf B^*$. Every difference is a quadratic error, a mixed product
bounded by Cauchy--Schwarz, or a bounded coefficient times $D_x$.
The potential difference is exactly
$\lambda(r_x^2-b)D_x/2+\lambda D_x^2/4$.
Summing gives $C_T(\sqrt{e_j}+e_j)$, including the potential derivatives.
The finite metric convention is the explicitly declared diagonal extension;
no unique off-diagonal lattice coupling is inferred from it.
\end{proof}

The same expansion controls any separately declared bounded-coefficient
quadratic cross insertion in the transported scalar derivatives and the
Hermitian electric/magnetic components. This algebraic statement is not a
claim that the independently varied full Einstein--SM metric action has
already been constructed. The retained continuum source limits are the
standard scalar--Maxwell expressions, including
\[
 T_{00}=\tfrac12|D_0z|^2+\tfrac12\sum_i|D_iz|^2+V(|z|)
                        +\tfrac12(|E|^2+|B|^2),
\]
\[
 T_{ii}=|D_iz|^2+\tfrac12|D_0z|^2-\tfrac12\sum_k|D_kz|^2-V(|z|)
                      -E_i^2-B_i^2+\tfrac12(|E|^2+|B|^2).
\]
Staggered line/face sampling and vertex reconstruction have deterministic
$O(h)$ source error on the smooth reference slab. Field means are interpolated
with every node retained, as in the previous field estimates. Thus
\eqref{v18:eq:maximal} gives spatial $H^1\times L^2$ radius/momentum and
$L^2$ electric/magnetic convergence uniformly in probability, together with
weighted spacetime $L^1$ source convergence. For a fixed source-error threshold
$u\le1$, the stochastic maximal failure bound is $C_T\mathfrak b_h/u^2$,
apart from that deterministic reconstruction error. Fixed-time expected
source errors are at most $C_T(\sqrt{\mathfrak b_h}+\mathfrak b_h)$.
The charged mean remains zero on physical states; it is never substituted
for a classical charged field.

\section{A nonempty magnetic class, simultaneous scale, and remaining boundary}
\label{v18:sec:ledger}

\begin{corollary}[A local exactly constrained branch with spatial magnetic curvature]
\label[corollary]{v18:cor:scales}
For any smooth constrained initial data whose continuum solution has a
positive scalar radius on a regular slab as above, choose
\begin{equation}
 \boxed{\hbar_h=h^7,\qquad R_h=h^{-28},\qquad
 \epsilon_h=h^{109},\qquad d_h=4\vartheta_hh^{109}/11,
 \quad0<\vartheta_h\le1.}
 \label{v18:eq:scales}
\end{equation}
The prepared local finite instrument then has $\mathfrak b_h=O(h^2)$,
$\E\max_j e_j=O(h^2(1+|\log h|))$, exact Gauss support on every outcome,
and all field/current/source conclusions above. Nonzero limiting magnetic
curvature and local charge are permitted. The operations and context
probabilities are unchanged from the constructed local gauge branch.
\end{corollary}
\begin{proof}
The four terms in \eqref{v18:eq:budget} have powers
$h^2,h^4,h^{-12+28-14}=h^2$, and
$h^{109-9-84-14}=h^2$. The uniform fourth energy moment verifies the
inherited tensor threshold for small $h$, and $M\epsilon=10h^{106}\to0$.
Local neutral multiplicity is supplied by the unchanged radial trial-space
argument in Version 15. All remaining assertions follow from
\cref{v18:thm:finite,v18:prop:readers,v18:prop:sources}; the logarithm has
order $1+|\log h|$ because $h^{-3}R=h^{-31}$.
\end{proof}

\paragraph{Explicit nonempty nonzero-magnetic initial data.}
For example, on the unit torus take
\[
 z_0(x)=r_0+\rho_1\cos(2\pi x_3),\quad r_0>|\rho_1|>0,
 \qquad A_0=(0,a\sin(2\pi x_1),0),\quad a\ne0,
\]
\[
 E_0=(e\sin(2\pi x_1),0,e_c),\qquad
 \dot z_0=v_0-\ii\frac{2\pi e\cos(2\pi x_1)}{z_0(x)}.
\]
Here $A_0$ denotes the \emph{initial spatial potential}, not a temporal
component. Temporal gauge remains imposed. Directly
$j_{0,0}=-2\pi e\cos(2\pi x_1)$, so
$\operatorname{div}E_0+j_{0,0}=0$. The magnetic field is
$(0,0,2\pi a\cos(2\pi x_1))$, with nonzero $L^2$ norm. All data are
smooth, periodic, and exactly continuum constrained. The inherited smooth
classical local-existence result cited at the start supplies a regular slab;
continuity retains the positive scalar radius and nonzero magnetic energy
on a shorter one. This supplies nonempty comparison data without proving
another Cauchy theorem. The finite source uses only the initial samples and
has exact \emph{finite} Gauss support through its preparation, even though
its sampled reference Gauss row has the $O(h^2)$ consistency defect.

\paragraph{What has been closed.}
Spatially varying reference amplitudes and nonzero reference plaquettes no
longer require a new candidate compactness assumption. The weighted edge
keeps the exact discrete adjoint cancellation, its mixed-product identity
pays the electric force, and the magnetic sine mismatch is bounded by
$8B_*h(E_d+M_d)$ as a form. The entire calculation has an exact operator
promotion on the unchanged Cartesian/link action. Physical preparation,
cutoff, event, and clock costs are retained, and the resulting estimates
concern actual resolved states, not just an ensemble mean. The magnetic
field and its stress are no longer driven to zero by the comparison.

\paragraph{What has not been closed.}
This is a local-in-time, constructed Abelian gauge--Higgs classical limit on
a prescribed flat metric, on a trivial bundle and a reference branch without
Higgs zeros. Candidate zeros are allowed. The proof does not establish
an arbitrary-reference weak solution theorem, remove reference zeros or
topological obstructions, or identify the primitive law's actual operator
panels with the constructed local contexts. Semiclassical scaling and
preparation remain declared. Non-Abelian curvature, chiral matter,
Einstein constraints, and metric backreaction remain independent. In
particular nonzero limiting matter and electromagnetic stress do not solve
a flat Einstein equation.

The next analytic target can therefore move beyond a further Abelian scalar
variation: either the same positive comparison with genuinely dynamical
geometry, or the non-Abelian curvature/current terms that are absent from
this commutative-link calculation. The source-identification target remains
the actual local context/action pullbacks. Neither is discharged by
repeating the clock, local Kraus, or Gauss-preservation lemmas.

\paragraph{Preservation and diagnostics.}
Every pre-existing byte before Version 17's closing document command is
preserved. All new labels use \texttt{v18:}. The companion script checks
the exact inhomogeneous identity with nonzero reference magnetic curvature,
local charge, forcing residuals, and candidate zeros; the weighted
mixed-product identity and magnetic positive bound; second-order covariant,
Maxwell and Gauss sampling; the magnetic quantum commutator on an interior
Fourier core; and the simultaneous exponent budget. These are finite
symbolic/numerical diagnostics, not formal proof checking, primitive-source
verification, or an Einstein--SM simulation.

\begingroup
\renewcommand{\refname}{Additional reference for Version 18}
\begin{thebibliography}{9}
\small\raggedright
\bibitem{v18:Pecher}
H.~Pecher, \emph{Unconditional global well-posedness in energy space for the
Maxwell--Klein--Gordon system in temporal gauge}, 2015 preprint.
\href{https://arxiv.org/abs/1504.00506}{arXiv:1504.00506}.
Background for temporal-gauge Cauchy theory; no global or low-regularity
result from this reference is claimed for the quartic compact-link limit.
\end{thebibliography}
\endgroup
\addtocontents{toc}{\protect\endgroup}
% END IMMUTABLE VERSION 18 BODY
% APPEND VERSION 19 HERE, BEFORE THE FINAL END-DOCUMENT COMMAND.


% BEGIN IMMUTABLE VERSION 19 BODY
\clearpage
\begin{center}
{\Large\bfseries Version 19: Non-Abelian cancellation before a quantum orbit comparison}\\[0.5em]
{\large Reference-zero-safe energy transfer, genuinely noncommuting compact-link data,
and an exact Casimir calculation}
\end{center}
\makeatletter
\addtocontents{toc}{\protect\begingroup}
\addtocontents{toc}{\protect\makeatletter}
\addtocontents{toc}{\protect\def\protect\l@section{\protect\bfseries\protect\@dottedtocline{1}{0em}{2.5em}}}
\makeatother
\addcontentsline{toc}{part}{Version 19: Non-Abelian energy cancellation and the quantum orbit boundary}

\section{Audit: do not replace noncommuting curvature by an Abelian copy}
\label{v19:sec:audit}

Version 18 establishes a constructed Abelian quantum-to-classical branch.
Its magnetic calculation uses commuting circle links and its reference ratios
require a nonzero reference scalar. Neither fact licenses substitution of
matrix-valued phases. This iteration goes upstream of that substitution.
It first identifies a comparison for the complete classical non-Abelian
first-order equations, with no division by either scalar field. It then
exhibits exact compact $SU(2)$ reference trajectories with noncommuting
magnetic curvature and isolates an operator identity appropriate to the
compact-group kinetic energy.

\paragraph{What is inherited.}
The displayed temporal-gauge augmented matter system, its currents,
constraint identities, and regular local existence are already in
\cite{CC}, at \nolinkurl{v2:eq:gauge-evolution},
\nolinkurl{v2:eq:higgs-evolution}, \nolinkurl{v2:thm:subsidiary}, and
\nolinkurl{v2:thm:local}. They are not new equations or a new Cauchy theorem
here. That source's polynomial $H^s$ estimates use a bounded ball of both
states. The additional result below instead bounds a relative energy by
\emph{reference} amplitudes only: the candidate connection, curvature, and
matter first derivatives have no imposed pointwise upper bounds in its
constant. The scalar convexification from Version 13 is extended to a real
unitary representation, and the non-Abelian exchange terms are calculated
explicitly.

The local gauge-commuting Kraus theorem and its physical-sector conclusion
remain \cref{v15:thm:locked}. No clock, local-Kraus, or Gauss-preservation
proof is repeated. Targeted semantic searches followed by literal checks of
the cumulative lineage, the Grand Tensor, the classical-closure file, and
the SM--Einstein monograph located their existing gauge equations and
Abelian comparisons, but not the combined relative identity, compact
reference, and operator calculation proved here. This is a limited overlap
audit, not an exhaustive review or a priority claim.

\paragraph{Precisely what is, and is not, proved.}
The principal propagation theorem in this body is \emph{classical}, for
smooth temporal-gauge fields or smooth forced augmented records. Its
quotient statement is an infimum over time-independent classical gauge
alignments. It is not an operator inequality on a gauge-singlet quantum
state. In particular the new result does not, by itself, extend Version
18's resolved-instrument theorem to non-Abelian links or through reference
Higgs zeros. The compact-group commutator calculation at the end is exact,
but a positive, gauge-invariant quantum \emph{relative} comparison must
still be constructed. This boundary is stated before, rather than after,
the new estimates.

\paragraph{Attribution.}
Temporal-gauge energy methods and symmetric hyperbolic formulations of
Yang--Mills systems are established mathematics; see
\cite{v19:LiuWang,v19:Lada}. The proofs here use standard integration by
parts, the invariant Lie-algebra inner product, quartic algebra, and compact
Lie-group differentiation. Their role is to expose the exact non-Abelian
terms and the reference-zero and source-topology boundaries in this
lineage. They are not a claim of a new general existence theory.

\section{The same non-Abelian action and a positive vector quartic comparison}
\label{v19:sec:setup}

Work on $[0,T]\times\T^3$ with its prescribed flat metric. Let $K$ be a
compact Lie group with Lie algebra $\mathfrak k$, endowed with a positive
$\operatorname{Ad}$-invariant inner product. Let $W$ be a finite-dimensional
real orthogonal representation, with infinitesimal action $T_X^*=-T_X$.
A complex unitary representation is treated with the real part of its
Hermitian product. Fix constants $c_{\rm ad},c_T$ such that
\[
 |[X,Y]|\le c_{\rm ad}|X||Y|,\qquad
 |T_Xu|\le c_T|X||u|.
\]
Neither faithfulness nor a scalar-stabilizer condition is needed. In
particular a gauge factor acting trivially on the Higgs must still be
included in the connection and curvature variables.

Put $D_i=\partial_i+[A_i,\cdot]$ on adjoint fields and
$D_i=\partial_i+T_{A_i}$ on $W$. Define the real bilinear current by
\begin{equation}
 \langle J(u,v),X\rangle=(v,T_Xu)_W.
 \label{v19:eq:current}
\end{equation}
Thus $J(v,u)=-J(u,v)$. Use the fixed potential
\begin{equation}
 V(H)=\frac\lambda4(|H|^2-b)^2,\qquad \lambda>0,\quad b\ge0,
 \qquad \alpha=1+\lambda b.
 \label{v19:eq:potential}
\end{equation}
For physical fields write
\[
 P=\partial_tH,\quad Q_i=D_iH,\quad E_i=\partial_tA_i,\quad
 F_{ij}=\partial_iA_j-\partial_jA_i+[A_i,A_j].
\]
The action density is
\[
 \frac12|P|^2-\frac12\sum_i|Q_i|^2-V(H)
 +\frac12\sum_i|E_i|^2-\frac12\sum_{i<j}|F_{ij}|^2.
\]
For the comparison allow the forced augmented equations
\begin{align}
 \dot A_i&=E_i,&\dot H&=P,\\
 \dot E_i&=D_jF_{ji}-J(H,Q_i)+r_i^E,
 &\dot F_{ij}&=D_iE_j-D_jE_i+r_{ij}^F,\\
 \dot P&=D_iQ_i-\nabla V(H)+r^P,
 &\dot Q_i&=D_iP+T_{E_i}H+r_i^Q.
 \label{v19:eq:system}
\end{align}
Here $F_{ji}=-F_{ij}$ and the independent curvature rows have $i<j$.
The two head equations are exact. A record with independent $Q,F$ is
called physical only when their defining compatibility conditions and
\begin{equation}
 D_iE_i+J(H,P)=0
 \label{v19:eq:Gauss}
\end{equation}
hold. The estimate below does not substitute a Gauss equation for an
unproved evolution equation. Fermions and dynamical geometry are absent.

For $S,H\in W$ let $\phi=H-S$ and define
\begin{equation}
 \mathcal B(H\mid S)=V(H)-V(S)-(\nabla V(S),\phi)
                                  +\frac\alpha2|\phi|^2.
 \label{v19:eq:B}
\end{equation}
This is a comparison potential, not a modification of the physical action.

\begin{lemma}[A zero-safe vector quartic identity]
\label[lemma]{v19:lem:quartic}
For every $S,\phi\in W$,
\begin{equation}
\begin{split}
 \mathcal B(S+\phi\mid S)
 ={}&\frac12|\phi|^2+\frac\lambda{12}|\phi|^4
   +\frac\lambda2|\phi|^2|S+\phi/3|^2\\
   &+\lambda\bigl((S,\phi)+|\phi|^2/3\bigr)^2.
\end{split}
\label{v19:eq:quartic}
\end{equation}
Consequently it controls $|\phi|^2/2+\lambda|\phi|^4/12$ globally.
For bounded $|S|$ it is also at most $C(|\phi|^2+|\phi|^4)$.
The exact nonlinear force remainder is
\begin{equation}
 N(S,\phi):=\nabla V(S+\phi)-\nabla V(S)-D^2V(S)\phi
 =\lambda\{ |\phi|^2S+2(S,\phi)\phi+|\phi|^2\phi\}.
 \label{v19:eq:N}
\end{equation}
All formulas remain valid at $S=0$ and at $H=0$.
\end{lemma}
\begin{proof}
The quadratic part after convexification is $|\phi|^2/2$. The other
terms before completing squares are
\[
 \frac\lambda2|S|^2|\phi|^2+\lambda(S,\phi)^2
       +\lambda(S,\phi)|\phi|^2+\frac\lambda4|\phi|^4.
\]
Expanding the last two positive terms of \eqref{v19:eq:quartic} gives
these terms after adding $\lambda|\phi|^4/12$. This proves the identity
without division by a field norm. The upper bound follows from bounded
$S$ and $|\phi|^3\le(|\phi|^2+|\phi|^4)/2$.
Finally $\nabla V(H)=\lambda(|H|^2-b)H$; expansion and removal of its
constant and linear parts give \eqref{v19:eq:N}.
\end{proof}

\section{A full non-Abelian energy identity with reference-only constants}
\label{v19:sec:identity}

Let starred variables be a smooth unforced solution of
\eqref{v19:eq:system}. They are in the same temporal trivialization as the
candidate. Define
\[
 a_i=A_i-A_i^*,\quad e_i=E_i-E_i^*,\quad f_{ij}=F_{ij}-F_{ij}^*,\quad
 \phi=H-H^*,\quad p=P-P^*,\quad y_i=Q_i-Q_i^*.
\]
The positive comparison is
\begin{equation}
 \mathcal R(t)=\int_{\T^3}\left[
 \frac12\sum_i(|a_i|^2+|e_i|^2+|y_i|^2)
 +\frac12|p|^2+\frac12\sum_{i<j}|f_{ij}|^2
 +\mathcal B(H\mid H^*)\right]\dd x.
 \label{v19:eq:R}
\end{equation}
The $a$ term is a lower-order comparison, not a mass term added to the
gauge action. Its necessity is demonstrated below. Set
\[
 \mathscr F(t)=\|r^P\|_2^2+\sum_i(\|r_i^Q\|_2^2+\|r_i^E\|_2^2)
                                 +\sum_{i<j}\|r_{ij}^F\|_2^2.
\]

\begin{theorem}[Exact reference-zero-safe non-Abelian relative energy]
\label[theorem]{v19:thm:relative}
For smooth solutions of the displayed forced and reference systems,
\begin{align}
 \dot{\mathcal R}=\int\biggl[&
 \sum_i\langle a_i,e_i\rangle+\alpha(\phi,p)
                     +(P^*,N(H^*,\phi))\notag\\
 &+\sum_i\{(p,T_{a_i}Q_i^*)
       +(y_i,T_{a_i}P^*+T_{E_i^*}\phi)
       -\langle e_i,J(\phi,Q_i^*)\rangle\}\notag\\
 &+\sum_{i,j}\langle e_i,[a_j,F_{ji}^*]\rangle
   +\sum_{i<j}\langle f_{ij},[a_i,E_j^*]-[a_j,E_i^*]\rangle\notag\\
 &+(p,r^P)+\sum_i\{(y_i,r_i^Q)+\langle e_i,r_i^E\rangle\}
                         +\sum_{i<j}\langle f_{ij},r_{ij}^F\rangle
 \biggr]\dd x.
 \label{v19:eq:identity}
\end{align}
If the reference fields $H^*,P^*,Q^*,E^*,F^*$ are bounded on the slab,
then
\begin{equation}
 \mathcal R(t)\le e^{Kt}\left[
     \mathcal R(0)+\frac12\int_0^t\mathscr F(s)\dd s\right],
 \qquad 0\le t\le T,
 \label{v19:eq:bound}
\end{equation}
where $K$ depends only on $\lambda,b,c_T,c_{\rm ad}$ and those reference
bounds. No candidate pointwise field or derivative bounds enter $K$.
No lower bound on either Higgs radius is required. No norm of a spatial
derivative operator enters the estimate.
\end{theorem}
\begin{proof}
Use the \emph{candidate} covariant derivative in every principal term of
the difference system. Direct subtraction gives
\begin{align*}
 \dot a_i&=e_i,&\dot\phi&=p,\\
 \dot p&=D_i y_i+T_{a_i}Q_i^*
              -[\nabla V(H)-\nabla V(H^*)]+r^P,\\
 \dot y_i&=D_i p+T_{a_i}P^*+T_{e_i}H+T_{E_i^*}\phi+r_i^Q,\\
 \dot e_i&=D_j f_{ji}+[a_j,F_{ji}^*]
                   -J(H,y_i)-J(\phi,Q_i^*)+r_i^E,\\
 \dot f_{ij}&=D_i e_j-D_j e_i
                    +[a_i,E_j^*]-[a_j,E_i^*]+r_{ij}^F.
\end{align*}
This choice is essential: expanding every $D_i$ about $A_i^*$ too early
would produce unnecessary appearances of the candidate connection norm.

Skew adjointness of $T_{A_i}$ and of $\operatorname{ad}_{A_i}$, together
with periodic integration by parts, gives
\[
 \int\{(p,D_i y_i)+(y_i,D_i p)\}=0,
\]
\[
 \int\left\{\sum_i\langle e_i,D_jf_{ji}\rangle
       +\sum_{i<j}\langle f_{ij},D_i e_j-D_j e_i\rangle\right\}=0.
\]
There is no commutativity assumption on the $A_i$. Invariance of the
Lie-algebra metric is exactly what makes the second cancellation hold.
The two apparently dangerous matter--electric products cancel pointwise:
\[
 (y_i,T_{e_i}H)-\langle e_i,J(H,y_i)\rangle=0.
\]
This cancels the complete candidate $H$, not only its reference value.

Differentiation of \eqref{v19:eq:B}, using $\dot H=P$ and
$\dot H^*=P^*$, gives
\[
 \partial_t\mathcal B
 =(\nabla V(H)-\nabla V(H^*),p)
       +(N(H^*,\phi),P^*)+\alpha(\phi,p).
\]
Its first term cancels the potential force in the kinetic derivative.
Differentiating $|a|^2/2$ gives the first term of
\eqref{v19:eq:identity}; the remaining reference products and forcings
are exactly those displayed there. This proves the identity.

To bound it, write $M$ for a bound on the sum of the absolute values of the
reference fields in the statement. The current estimate
$|J(u,v)|\le c_T|u||v|$ follows from its definition. Every bilinear
reference term is consequently bounded by
$C(M,c_T,c_{\rm ad})$ times the sum of the squared errors in
\eqref{v19:eq:R}. For the only higher-order term,
\[
 |(P^*,N(H^*,\phi))|
 \le\lambda M(3M|\phi|^2+|\phi|^3)
 \le C_{M,\lambda}\mathcal B(H\mid H^*),
\]
by \cref{v19:lem:quartic}. The $\alpha$ term has the same bound by
Young's inequality and the quadratic coercivity. Each forcing product is
at most half the corresponding squared error plus half its forcing square.
Thus $\dot{\mathcal R}\le K\mathcal R+\mathscr F/2$. Multiply by
$e^{-Kt}$, integrate, and enlarge the exponential outside the forcing
integral to obtain \eqref{v19:eq:bound}.
All divisions in the proof are by fixed constants, never by a scalar
field. Smoothness justifies the integrations; this is not an existence or
an automatic weak-solution approximation theorem.
\end{proof}

\begin{remark}[What low-regularity topology has actually been controlled]
The estimate controls $a,p,e,y,f$ in $L^2$ and $\phi$ in $L^2\cap L^4$.
When $Q,Q^*$ are the actual covariant gradients, $D_A\phi=y-T_aH^*$
also tends to zero in $L^2$. It does not assert ordinary $H^1$ convergence
of the gauge potential or of the scalar in an arbitrary fixed gauge.
No such conclusion follows by forgetting the connection in a covariant
gradient. The finite inner-product calculation remains valid with skew
adjoint spatial difference operators, but that fact alone would not supply
local gauge covariance or compatibility of a proposed lattice regulator.
\end{remark}

\section{Gauge alignment and the connection information missing from curvature}
\label{v19:sec:gauge}

For a smooth time-independent $g:\T^3\to K$ transform a candidate by
\[
 A_i^g=gA_ig^{-1}-(\partial_i g)g^{-1},\qquad
 E_i^g=gE_ig^{-1},\qquad F_{ij}^g=gF_{ij}g^{-1},
\]
and by its representation on $H,P,Q$. Such transformations preserve
$A_0=0$. Let $\mathcal R_g(t)$ be the preceding comparison of the transformed
candidate with the fixed reference. Define the nonnegative orbit comparison
\begin{equation}
 \mathcal R_{\rm orb}(t)=\inf_{g\in C^\infty(\T^3,K)}\mathcal R_g(t).
 \label{v19:eq:orbit}
\end{equation}
It is not claimed to be a symmetric metric; the relative potential is
reference based. Nor is a minimizing gauge assumed to exist.

\begin{corollary}[One initial alignment controls the classical orbit]
\label[corollary]{v19:cor:orbit}
Under the hypotheses of \cref{v19:thm:relative},
\begin{equation}
 \mathcal R_{\rm orb}(t)\le e^{Kt}
 \left[\mathcal R_{\rm orb}(0)+\frac12\int_0^t\mathscr F(s)\dd s\right].
 \label{v19:eq:orbit-bound}
\end{equation}
The constant is independent of the chosen gauge and its spatial derivatives.
A time-dependent gauge minimizing the comparison separately at every time
is not used in the proof.
\end{corollary}
\begin{proof}
The transformed candidate solves the same forced equations with transformed
forcings, whose norms are unchanged. The reference is fixed, and the
constant in \cref{v19:thm:relative} contains no candidate bounds. Apply that
theorem for each single time-independent $g$ and then take the infimum of
its initial comparison. For each such $g$, $\mathcal R_{\rm orb}(t)$ is at
most $\mathcal R_g(t)$, so approximation to the initial infimum proves the
claim without compactness or attainment of that infimum.
\end{proof}

\begin{proposition}[Field strength alone is not a complete non-Abelian initial comparison]
\label[proposition]{v19:prop:connection-needed}
Removing $\sum_i|a_i|^2/2$ from \eqref{v19:eq:R} invalidates a bound of
its zero-initial-error form, even for smooth homogeneous, exactly
Gauss-compatible $SU(2)$ solutions. This remains true with identically
zero scalar fields and nonzero initial electric field.
\end{proposition}
\begin{proof}
Use $T_a=-\ii\sigma_a/2$ and $\langle X,Y\rangle=-2\tr(XY)$, so
$[T_a,T_b]=\epsilon_{abc}T_c$ and $|T_a|=1$. At time zero take the
reference
\[
 A_i^*=0,\qquad E_2^*=T_2,\quad E_1^*=E_3^*=0,\qquad H^*=P^*=0,
\]
and candidate
\[
 A_1=cT_1,\quad A_2=A_3=0,\qquad E_i=E_i^*,\qquad H=P=0,
 \quad c\ne0.
\]
Both have $F=0$, $Q=0$, zero matter charge, and $\sum_i[A_i,E_i]=0$.
They produce smooth homogeneous solutions of the polynomial Yang--Mills
ODE, and the inherited constraint identity preserves Gauss's law. Initially
all terms of the comparison except $|a|^2/2$ vanish. But the curvature
identity gives
\[
 \partial_t(F_{12}-F_{12}^*)\big|_{t=0}
 =[A_1,E_2]-[A_2,E_1]=cT_3.
\]
Thus the comparison with the connection term deleted is
$c^2t^2/2+O(t^3)>0$ for small positive $t$. No estimate multiplying its
zero initial value can hold. Taking, for example, $c\notin4\pi\Z$ also
keeps the initial first-cycle holonomy distinct from the reference. The
example concerns missing initial connection/transport information, not a
violation of gauge covariance.
\end{proof}

The connection term is therefore not dispensable bookkeeping when passing
to an unbroken non-Abelian sector. A Higgs singlet cannot supply the missing
colour alignment. Gauge-invariant traces of curvatures at one instant are
not substitutes for this information.

\section{The complete classical stress and actual source tests}
\label{v19:sec:stress}

For a physical candidate let $\mathsf T$ be the stress tensor obtained by
independent metric variation of the action in \cref{v19:sec:setup}. In the
specified flat coordinates its entries are
\begin{align}
 \mathsf T_{00}&=\tfrac12\bigl(|P|^2+\sum_i|Q_i|^2
                 +\sum_i|E_i|^2+\sum_{i<j}|F_{ij}|^2\bigr)+V(H),\\
 \mathsf T_{0i}&=(P,Q_i)+\sum_j\langle E_j,F_{ij}\rangle,\\
 \mathsf T_{ij}&=(Q_i,Q_j)+\delta_{ij}
             \bigl(\tfrac12|P|^2-\tfrac12\sum_k|Q_k|^2-V(H)\bigr)\notag\\
 &\quad-\langle E_i,E_j\rangle+\sum_k\langle F_{ik},F_{jk}\rangle
    +\delta_{ij}\bigl(\tfrac12\sum_k|E_k|^2
                         -\tfrac12\sum_{k<\ell}|F_{k\ell}|^2\bigr).
 \label{v19:eq:stress}
\end{align}
These include the off-diagonal and momentum sources. They are continuum
action derivatives, not a claimed uniqueness theorem for an off-diagonal
compact-lattice metric extension.

\begin{theorem}[All classical quadratic and quartic source errors are controlled]
\label[theorem]{v19:thm:stress}
At each time under the reference bounds of \cref{v19:thm:relative},
\begin{equation}
 \sum_{\mu\le\nu}\|\mathsf T_{\mu\nu}-\mathsf T_{\mu\nu}^*\|_1
  +\sum_{\mu=0}^3\|J_\mu-J_\mu^*\|_1
 \le C\bigl(\sqrt{\mathcal R}+\mathcal R\bigr),
 \label{v19:eq:stress-bound}
\end{equation}
where $J_0=J(H,P)$, $J_i=J(H,Q_i)$ and the current comparison is in the
chosen common gauge. The same bound holds for $\nabla V(H)-\nabla V(H^*)$
in $L^1$ and for the independent $\lambda,b$ derivatives of the potential.
For the gauge-invariant stress and potential insertions the right side may
be replaced by $C(\sqrt{\mathcal R_{\rm orb}}+\mathcal R_{\rm orb})$.
\end{theorem}
\begin{proof}
For every quadratic pairing, write each candidate factor as its reference
plus its error. Products of two errors have $L^1$ norm at most the product
of their $L^2$ norms. A product of a reference and an error is bounded by
its reference $L^2$ norm times the error $L^2$ norm. These give
$C(\sqrt{\mathcal R}+\mathcal R)$ for every quadratic term in
\eqref{v19:eq:stress}, without requiring that two Lie-algebra matrices
commute. The invariant inner product is the action's fixed contraction.
For the currents, use the bilinear identity
\[
 J(H,Q_i)-J(H^*,Q_i^*)=J(H^*,y_i)+J(\phi,Q_i^*)+J(\phi,y_i),
\]
and its temporal counterpart. The same $L^2$ estimates apply.

Expansion of a quartic or cubic polynomial about bounded $H^*$ produces
terms of degree one through four in $\phi$. The degree-one spatial
integral is bounded by $C\|\phi\|_2$. The remaining terms are bounded
by $C\int(|\phi|^2+|\phi|^4)$, including degree three by Young's
inequality. \Cref{v19:lem:quartic} proves the required bound. The potential
parameter derivatives have degrees at most four, so the same reasoning
applies. The stress formulas themselves follow by contracting the scalar
and Yang--Mills kinetic tensors under an arbitrary symmetric metric
variation; $F_{0i}=E_i$ and antisymmetry give the signs displayed in
\eqref{v19:eq:stress}. Thus they are independent insertions of the stated
action.

For each fixed gauge alignment $g$, the candidate stress and scalar
potential insertions are unchanged, while the preceding bound holds with
$\mathcal R_g$. Approximate the infimum in \eqref{v19:eq:orbit} and use
continuity of $\sqrt u+u$. Current components are adjoint covariants, not
invariant scalars, so their comparison retains its specified alignment.
\end{proof}

\begin{corollary}[A classical low-order entrance for non-Abelian residuals]
\label[corollary]{v19:cor:entrance}
Let a sequence of smooth forced augmented records have the same smooth
physical reference and
\[
 \mathcal R_n(0)+\int_0^T\mathscr F_n(t)\dd t\longrightarrow0.
\]
Then $a_n,e_n,f_n,p_n,y_n\to0$ in $L_t^\infty L_x^2$ and
$\phi_n\to0$ in $L_t^\infty(L_x^2\cap L_x^4)$, with the fourth power
controlled by the comparison. The full stress and potential source
insertions converge in $L_t^\infty L_x^1$. No independent all-time
candidate Sobolev bound is needed by this implication. An analogous
stress conclusion holds when the initial error is only small after one
classical temporal-gauge alignment.
\end{corollary}
\begin{proof}
Apply \eqref{v19:eq:bound}, quartic coercivity, and
\cref{v19:thm:stress}. For the last statement use
\cref{v19:cor:orbit}. The statement supplies neither the existence of a
proposed record nor its residual estimates. If its augmented variables
are not initially known to be physical, their defining constraints remain
a separate obligation. Those qualifications prevent relabelling this
classical transfer as a new primitive-source generation theorem.
\end{proof}

\section{An exact noncommuting SU(2) family, including a reference Higgs zero}
\label{v19:sec:SU2}

Use again $T_a=-\ii\sigma_a/2$, $\langle X,Y\rangle=-2\tr(XY)$,
and the fundamental doublet with real product $\operatorname{Re}(u^*v)$.
Fix a unit doublet $u$. Take the spatially homogeneous ansatz
\begin{equation}
 A_i=a_iT_i,\quad E_i=e_iT_i,\quad H=r u,\quad P=v u,
 \quad Q_i=a_i rT_i u,\quad F_{ij}=a_i a_j\epsilon_{ijk}T_k.
 \label{v19:eq:SU2-data}
\end{equation}
The scalar coefficient $r$ is signed. There is no polar-coordinate
singularity when it crosses zero.

\begin{theorem}[A genuinely non-Abelian classical reference]
\label[theorem]{v19:thm:SU2}
The ansatz is an exact physical solution of the continuum equations if
\begin{equation}
 \begin{split}
 \dot r=v,\qquad
 \dot v&=-r\left(\tfrac14\sum_i a_i^2+\lambda(r^2-b)\right),\\
 \dot a_i=e_i,\qquad
 \dot e_i&=-a_i\sum_{j\ne i}a_j^2-\tfrac14r^2a_i.
 \end{split}
 \label{v19:eq:SU2-flow}
\end{equation}
These equations conserve
\begin{equation}
 \mathcal H_*=	frac12v^2+\tfrac12\sum_i e_i^2
  +\tfrac12\sum_{i<j}a_i^2a_j^2
  +\tfrac18r^2\sum_i a_i^2+\tfrac\lambda4(r^2-b)^2.
 \label{v19:eq:SU2-energy}
\end{equation}
Every finite initial datum has a global smooth ODE solution. On every
finite time interval its reference fields in \cref{v19:thm:relative}
are bounded. Initial data with $r(0)=0$, $v(0)\ne0$, and all $a_i(0)\ne0$
provide a reference with a Higgs zero and genuinely noncommuting magnetic
curvature. In particular
\begin{equation}
 |[F_{12},F_{23}]|^2=a_1^2a_2^4a_3^2>0
 \label{v19:eq:noncommuting}
\end{equation}
on a nonempty interval adjoining that initial time.
\end{theorem}
\begin{proof}
The Pauli product gives $T_i^2=-I/4$ and
$\operatorname{Re}(T_i u,T_j u)=\delta_{ij}/4$.
Hence $D_iQ_i=-r(\sum_i a_i^2)u/4$ and
$J(H,Q_i)=r^2a_iT_i/4$. Also
\[
 \sum_j[A_j,F_{ji}]=-a_i\sum_{j\ne i}a_j^2 T_i.
\]
These identities give \eqref{v19:eq:SU2-flow} from the two acceleration
equations. The defining first-order identities then give the $Q,F$
equations. Each $[A_i,E_i]$ is zero, and
$J(H,P)=rvJ(u,u)=0$, so Gauss's law holds exactly.

Differentiation of \eqref{v19:eq:SU2-energy} cancels its gauge--scalar
and magnetic work terms. Nonnegativity bounds $v,e_i$ on every solution
and bounds $r$ by the quartic. Although the magnetic quartic does not
bound each $a_i$ separately when the others vanish, bounded $e_i$ gives
$|a_i(t)|\le|a_i(0)|+C|t|$. Thus every component remains bounded on
finite intervals. The polynomial ODE continuation criterion proves global
existence, not a global existence theorem for arbitrary spatially varying
fields. The commutator in \eqref{v19:eq:noncommuting} is
$a_1a_2^2a_3T_2$, up to its immaterial orientation sign. Continuity of
nonzero initial $a_i$ gives the stated interval. Smoothness at $r=0$
is explicit in the polynomial equations.
\end{proof}

The relative-energy theorem compares this reference with fully spatial
classical candidates; it does not require them to share the homogeneous
ansatz. The comparison remains valid across reference scalar zeros.
It still does not supply a physical quantum preparation concentrated in
its common-gauge coordinates.

\subsection{The compact action, not a Lie-algebra replacement of the links}
On the cubic torus with $w=h^3$ define the compact $SU(2)$ energy
\begin{equation}
\begin{split}
 E_h^{(2)}={}&\sum_x\frac{|p_x|^2}{2w}
  +\frac h2\sum_{x,i}|U_{x,i}H_{x+e_i}-H_x|^2
  +\frac{\lambda w}{4}\sum_x(|H_x|^2-b)^2\\
 &+\frac1{2h}\sum_{x,i}|\Pi_{x,i}|^2
  +\frac4h\sum_f\left(1-\frac12\operatorname{Re}\tr U_f\right).
\end{split}
\label{v19:eq:compact-action}
\end{equation}
The link kinetic convention is
$\dot U_{x,i}=h^{-1}\Pi_{x,i}U_{x,i}$. Its electric norm uses the
orthonormal $T_a$ basis above. The Wilson coefficient $4/h$ is declared
and normalized to the continuum kinetic term; it is not the Abelian
coefficient copied without its representation trace factor.

\begin{proposition}[An exact compact reference and its $O(h^2)$ limit]
\label[proposition]{v19:prop:compact}
Put
\[
 H_x=r_hu,\quad p_x=wv_hu,\quad
 U_{x,i}=e^{h a_{h,i}T_i},\quad \Pi_{x,i}=h^2e_{h,i}T_i.
\]
For $s_i=\sin(ha_{h,i}/2)$, $c_i=\cos(ha_{h,i}/2)$, these are exact
homogeneous classical trajectories of \eqref{v19:eq:compact-action} when
\begin{equation}
\begin{split}
 \dot r_h=v_h,\qquad
 \dot v_h&=-r_h\left(\frac2{h^2}\sum_i(1-c_i)
                                  +\lambda(r_h^2-b)\right),\\
 \dot a_{h,i}=e_{h,i},\qquad
 \dot e_{h,i}&=-\frac{r_h^2}{2h}s_i
                    -\frac8{h^3}s_i c_i\sum_{j\ne i}s_j^2.
\end{split}
\label{v19:eq:compact-flow}
\end{equation}
They satisfy the exact non-Abelian Gauss constraint. For the same fixed
initial $r,v,a,e$ as in \cref{v19:thm:SU2}, their reduced variables and
first time derivatives converge uniformly at rate $O_T(h^2)$ on each
fixed finite interval. Their reference plaquettes have nonzero
commutator curvature when the corresponding $a_i$ are nonzero.
\end{proposition}
\begin{proof}
For unit vectors $n,m\in\R^3$, direct Pauli multiplication gives
\begin{equation}
 1-\tfrac12\operatorname{Re}\tr
   \left(e^{s n\cdot T}e^{t m\cdot T}
                e^{-s n\cdot T}e^{-t m\cdot T}\right)
 =2\sin^2(s/2)\sin^2(t/2)\bigl(1-(n\cdot m)^2\bigr).
 \label{v19:eq:Wilson-trace}
\end{equation}
At the orthogonal axes $n=e_i,m=e_j$ this is $2s_i^2s_j^2$.
Moreover $\operatorname{Re}(u^*Uu)=\tr U/2$ for every $U\in SU(2)$,
so the homogeneous hopping energy does not select a direction of the unit
doublet $u$. The reduced energy per unit volume is therefore
\begin{equation}
 \tfrac12v_h^2+\tfrac12\sum_i e_{h,i}^2
 +\frac{r_h^2}{h^2}\sum_i(1-c_i)
 +\frac8{h^4}\sum_{i<j}s_i^2s_j^2
 +\tfrac\lambda4(r_h^2-b)^2.
 \label{v19:eq:compact-energy}
\end{equation}
To verify that reduction has not omitted transverse link forces, vary
an axis in \eqref{v19:eq:Wilson-trace}. Its first derivative at
$n\cdot m=0$ is zero. The hopping trace also has zero first variation
under transverse changes of the axis. Translation invariance makes each
single-link gradient equal to its homogeneous variation divided by the
number of links in that direction. Consequently the full local gauge
forces are along $T_i$, not just their projection onto the ansatz.
The scalar force is along $u$, since incoming and outgoing homogeneous
links have sum $2c_iI$. Differentiating \eqref{v19:eq:compact-energy}
now gives exactly \eqref{v19:eq:compact-flow}. Indeed the original
canonical one-form restricts to $v_h\dd r_h+\sum_i e_{h,i}\dd a_{h,i}$:
each link contributes $h^3e_{h,i}\dd a_{h,i}$, and there are $h^{-3}$
links in each direction; the scalar mass has the same cancellation.
The electric coadjoint
term from its quadratic norm is zero. Each incoming flux is conjugated
by a link commuting with that flux, so the outgoing and incoming Gauss
contributions cancel. The scalar angular charge is zero as before.

The reduced energy is conserved. It bounds $v_h,e_{h,i},r_h$ uniformly
for fixed initial data and small $h$; integration of $\dot a_h=e_h$
gives uniform $a_h$ bounds on each finite interval. Thus the ODE stays
in a common bounded set. The integral sine and cosine remainders show
that its vector field and its first derivatives are uniformly bounded
there and differ by $O(h^2)$ from \eqref{v19:eq:SU2-flow}. Comparing
integral equations and applying finite-dimensional Gronwall proves the
stated rate; substitution proves the first-derivative rate. This uses
no gauge--Higgs PDE approximation theorem. Finally expansion of the
ordered commutator gives
$U_{ij}=I+h^2a_ia_j[T_i,T_j]+O(h^3)$ on that same bounded set.
The Wilson trace formula verifies its positive magnetic energy directly.
\end{proof}

This compact family is a source-independent reference calculation, not a
claim that the local resolved quantum instrument concentrates on it. It
specifies noncommuting same-action data on which the next quantum
comparison can be tested, with correct electric and Wilson normalizations.

\section{An exact compact-group kinetic calculation before quantum promotion}
\label{v19:sec:Casimir}

The constant-coordinate kinetic hypothesis of \cref{v17:lem:Weyl} does
not hold for a non-Abelian compact-group Laplacian in arbitrary coordinates.
The following invariant calculation is an alternative for the pieces it
actually covers.

On a finite product of compact groups with normalized Haar measure, choose
orthonormal invariant vector fields $X_{e,a}$ for each link $e$, skew
adjoint on the smooth core. The Lie-algebra metric is bi-invariant. Set
\[
 L_{e,a}=-\ii\hbar X_{e,a},\qquad
 T=\sum_{e,a}\frac{L_{e,a}^2}{2I_e},\qquad I_e>0.
\]
The sign of the structure constants depends on the invariant-field
convention; the quadratic Casimir commutes with each $L_{e,a}$ in either
convention. All coefficients in the next statement are smooth real
multiplication functions. The $I_e$ and the reference components $p_{e,a}$
are scalar coefficients, not configuration-dependent inertia tensors.

\begin{theorem}[Exact Casimir drift of a shifted electric comparison]
\label[theorem]{v19:thm:Casimir}
Let $\mathcal H=T+V(U)$ and
\[
 C(t)=\sum_{e,a}\frac{(L_{e,a}-p_{e,a}(t)I)^2}{2I_e}+f(t,U).
\]
Then on the smooth group core, with
$\operatorname{Sym}(A,B)=(AB+BA)/2$,
\begin{align}
 \left(\partial_t+\frac\ii\hbar[\mathcal H,\cdot]\right)C
 ={}&\dot f+
 \sum_{e,a}\frac1{I_e}
   \operatorname{Sym}(L_{e,a}-p_{e,a},X_{e,a}(f-V))\notag\\
 &+\sum_{e,a}\frac{p_{e,a}}{I_e}X_{e,a}f
 -\sum_{e,a}\frac{\dot p_{e,a}}{I_e}(L_{e,a}-p_{e,a}).
 \label{v19:eq:Casimir}
\end{align}
There is no unlisted $\hbar^2$ remainder. When $f\ge0$, the displayed
comparison is positive. Positivity bounds on its derivative still require
bounds on the actual multiplication functions on the right.
\end{theorem}
\begin{proof}
Invariant differentiation and Haar integration give
$[L_{e,a},f]=-\ii\hbar X_{e,a}f$. It follows immediately that
\[
 \frac\ii\hbar[T,f]=\sum_{e,a}\frac1{I_e}
                         \operatorname{Sym}(L_{e,a},X_{e,a}f),
 \qquad \frac\ii\hbar[V,L_{e,a}]=-X_{e,a}V.
\]
To check centrality, write the commutator of $\sum_aL_a^2$ with $L_b$
as a constant multiple of $\sum_{a,c}c_{ab}^{\ c}(L_aL_c+L_cL_a)$.
The structure constants are antisymmetric in $a,c$ for an invariant
orthonormal metric, whereas the operator expression is symmetric, so the
sum is zero. Different links commute. Thus the kinetic part does not
contribute to the derivative of a shifted electric square. Differentiate
that square in time and use the second commutator identity. Add the first
identity for $f$, split $L=(L-p)+p$, and collect terms to obtain
\eqref{v19:eq:Casimir}. All symmetrizations are retained, so no commutation
of a momentum with a nonconstant force is presumed. The sum of squares
proves the final positivity statement. These are identities before any
noncommuting tensor compression; its error would still have to be paid.
\end{proof}

\begin{proposition}[A normalized $SU(2)$ Wilson-force square]
\label[proposition]{v19:prop:Wilson-gradient}
For $U\in SU(2)$ put $v(U)=1-\tr U/2$, and use the invariant derivatives
associated with $T_a=-\ii\sigma_a/2$. Then
\begin{equation}
 \sum_a|X_av(U)|^2=\tfrac14v(U)(2-v(U))\le\tfrac12v(U).
 \label{v19:eq:Wilson-gradient}
\end{equation}
For a plaquette word $P$ using four distinct links, each once with either
orientation, put $v_f=1-\tr P/2$. Its complete link derivative satisfies
\begin{equation}
 \sum_{e\in f,a}|X_{e,a}v_f|^2=v_f(2-v_f)\le2v_f.
 \label{v19:eq:plaquette-gradient}
\end{equation}
These are multiplication identities for arbitrary candidate links, not a
small-angle expansion.
\end{proposition}
\begin{proof}
Write $U=x_0I+2\sum_ax_aT_a$ with $x_0^2+\sum_ax_a^2=1$.
The normalized trace gives $v=1-x_0$ and the left derivative gives
$X_av=x_a/2$, up to a convention-dependent sign for right derivatives.
Squaring and summing proves \eqref{v19:eq:Wilson-gradient}. Variation
of one factor of a plaquette inserts an orthonormal Lie basis conjugated
by a partial product, with an additional sign for an inverted factor.
Adjoint invariance preserves the sum of squared derivatives. Each of
its four distinct links therefore contributes the preceding single-group
value with $U=P$. This proves \eqref{v19:eq:plaquette-gradient}.
\end{proof}

The two exact calculations identify where a compact non-Abelian operator
proof should start. They do not identify the relative Wilson force for
arbitrarily interleaved candidate and reference holonomies. Its order,
transport, and positivity are additional mathematical content; replacing
that word by an Abelian difference of angles would discard it.

\section{The remaining quantum and native-source boundary}
\label{v19:sec:ledger}

The new classical implication is
\[
 \begin{gathered}
 \text{initial alignment + actual low-order forcing budget}\\
 \Downarrow\\
 \text{non-Abelian field/current/stress comparison, including reference zeros}.
 \end{gathered}
\]
Its stability constant depends on the reference covariant fields rather
than the candidate's Sobolev norm. The magnetic and scalar--electric
cancellations are exact and do not require commuting gauge generators.
The compact $SU(2)$ family also gives explicit finite reference data with
noncommuting magnetic curvature, including a scalar zero. These statements
are substantive extensions of the \emph{classical analytic interface}.

They must not be confused with a positive quantum-to-classical extension.
The common-gauge subtraction in \eqref{v19:eq:R} is not a gauge-invariant
operator on the exact physical Hilbert space. The already-proved charged
mean obstruction \cref{v15:prop:charged-mean} prevents preparing a physical
singlet with arbitrarily small expectation of a raw charged-coordinate
square about a nonzero fixed reference. Haar averaging that square is not
minimization over an aligned orbit. Likewise the classical infimum
\eqref{v19:eq:orbit} is nonlinear; no operator with the required propagation
inequality has been obtained by writing that infimum inside a trace.
At scalar zeros, the earlier Higgs-frame construction also cannot supply
an invertible field frame. These are distinct and visible restrictions,
not corrections silently made to Version 18.

The next non-Abelian source calculation therefore has a concrete target:
a positive \emph{gauge-invariant relative} observable controlling aligned
connection or holonomy, electric, curvature, and matter errors on actual
physical records. Its compact-group drift can use \eqref{v19:eq:Casimir},
but must preserve the ordered reference/candidate Wilson products and pay
their force terms. A bounded source expectation and a small initial
preparation cost for that observable would then be genuine new inputs to
the inherited local instrument transfer. The present body neither assumes
those inputs to be already available nor labels them as proved.

The original primitive-law question also persists: the actual local contexts,
operator pullbacks, physical preparation, and semiclassical normalization
have not been identified with the constructed family. No Einstein metric is
solved, and the nonzero Yang--Mills--Higgs stress is not a flat Einstein
solution. Chiral fermions and gravitational constraints are outside all
statements in this body. No claim of global weak-solution theory is made.

\paragraph{Verification and append-only preservation.}
Every byte of Version 18 before its final end-document command is retained.
New labels have prefix \texttt{v19:}. The diagnostic script checks the
vector quartic identity in several representation dimensions; the complete
non-Abelian relative derivative on finite skew-adjoint periodic test grids,
with forcing and scalar zeros; the connection-only initial obstruction;
full local compact-$SU(2)$ force differentiation rather than only a restricted
energy; the Wilson trace and force-square normalizations; Casimir centrality and
the exact invariant differential-operator drift;
and the $O(h^2)$ compact reference rate. The periodic difference tests check
algebraic summation by parts, not a claimed gauge-invariant spatial regulator.
All are finite symbolic/numerical diagnostics, not machine-checked proofs,
primitive-source identification, or an Einstein--SM simulation.

\begingroup
\renewcommand{\refname}{Additional references for Version 19}
\begin{thebibliography}{9}
\small\raggedright
\bibitem{v19:LiuWang}
C.~Liu and J.~Wang,
\emph{A new symmetric hyperbolic formulation and the local Cauchy problem
for the Einstein--Yang--Mills system in the temporal gauge}, 2021 preprint.
\href{https://arxiv.org/abs/2111.04540}{arXiv:2111.04540}.
Background for tensorial symmetric hyperbolic and constraint formulations;
no local or global existence result is imported as a new theorem here.
\bibitem{v19:Lada}
A.~\L ada,
\emph{Global solutions of Yang--Mills--Higgs systems on four-dimensional
Minkowski space: I. Existence}, Mathematical Methods in the Applied Sciences
\textbf{12} (1990), 153--165.
\href{https://doi.org/10.1002/mma.1670120206}{doi:10.1002/mma.1670120206}.
Background for classical temporal-gauge Yang--Mills--Higgs theory, distinct
from the new comparison and compact-reference calculations.
\end{thebibliography}
\endgroup
\addtocontents{toc}{\protect\endgroup}
% END IMMUTABLE VERSION 19 BODY
% APPEND VERSION 20 HERE, BEFORE THE FINAL END-DOCUMENT COMMAND.


% BEGIN IMMUTABLE VERSION 20 BODY
\clearpage
\begin{center}
{\Large\bfseries Version 20: Invariant phases before a quantum orbit comparison}\\[0.5em]
{\large A finite Wilson-word norm, exact physical preparation, and positive
non-Abelian quantum comparison on the resolved renewal clock}
\end{center}
\makeatletter
\addtocontents{toc}{\protect\begingroup}
\addtocontents{toc}{\protect\makeatletter}
\addtocontents{toc}{\protect\def\protect\l@section{\protect\bfseries\protect\@dottedtocline{1}{0em}{2.5em}}}
\makeatother
\addcontentsline{toc}{part}{Version 20: Positive gauge-invariant quantum orbit comparison}

\section{Audit: separate configuration alignment from momentum alignment}
\label{v20:sec:audit}

Version 19 leaves a precise operator problem. Its classical infimum over
aligned field energies is not itself an operator on a physical singlet.
Nor can averaging a charged-coordinate square make that square concentrate
about a nonzero charged reference. The construction below does something
different: a nonnegative \emph{configuration} function separates gauge
orbits, and a gauge-invariant phase supplies the momentum shift. They are
then quantized separately as a multiplication operator and an explicit sum
of differential squares. No infimum of noncommuting operators is used.

\paragraph{Inherited results and the limited overlap audit.}
The maximal-tree holonomy reduction and finite invariant-loop saturation
already occur in \cite{GT}, in particular at
\nolinkurl{thm:YM-relative-tree-gauge} and
\nolinkurl{thm:YM-primitive-loop-saturation}. We use their root-loop
coordinates, not a new tree-gauge theorem. The exact compact-group calculus
is \cref{v19:sec:Casimir}, and the tensor Galerkin, local gauge-commuting
instrument, and physical-clock estimates are
\cref{v14:thm:Galerkin,v15:thm:locked,v15:thm:source}.
Their proofs are not claimed again. Targeted searches of those sources and
the cumulative lineage did not locate the invariant-phase comparison,
finite separating norm, or its physical preparation proved here. This is
not an exhaustive bibliographic or proof audit.

\paragraph{Exact scope of the new positive result.}
The configuration manifold is a \emph{fixed finite} product of compact
$SU(2)$ link groups. We work in the unbroken pure-gauge sector, with its
independently specified Wilson action, not with an Abelian replacement of
its links. No Higgs frame, additional reference matter, or nonzero scalar
is needed. This does \emph{not} prove a quantum Yang--Mills--Higgs limit
through an evolving scalar zero. The new conclusion is a non-Abelian
\emph{fixed-lattice} semiclassical and resolved-instrument theorem. Its
constants and initial Hamilton--Jacobi interval depend on the graph and
its physical coefficients. No spatial-refinement-uniform bound is asserted.
The contrast between this theorem and the spatially uniform Abelian
conclusions is retained throughout.

\paragraph{Attribution.}
Gauge-field coherent states, semiclassical concentration, and
Hamilton--Jacobi/WKB methods are established subjects; see
\cite{v20:TW2,v20:TW3}. The contributions within this lineage are an actual
positive gauge-invariant relative operator, its exact geometric drift,
an explicitly physical preparation, and its local locked-clock realization.
All estimates used below are proved here or identified as inherited. The
chosen action and preparation are not inferred from primitive outcome counts.

\section{A finite Wilson-word norm for a noncommuting configuration orbit}
\label{v20:sec:words}

Let $\Gamma$ be a finite connected graph and
$\mathcal Q=SU(2)^{\mathcal E}$. Vertex gauges act by
$U_e\mapsto g_xU_eg_y^{-1}$ for $e:x\to y$.
Choose a root and a spanning tree. The complete tuple of fundamental
rooted chord holonomies determines the configuration modulo the based
gauge group; the remaining root gauge acts by simultaneous conjugation.
Retain all these holonomies, including noncontractible cycles.
We may add two specified rooted loop words as anchors. Denote the resulting
finite list by $U_1,\ldots,U_m$, with anchors numbered 1 and 2. The list
still contains every fundamental chord holonomy. It is not replaced by
its separate marginal trace distributions.

Use $T_a=-\ii\sigma_a/2$ and write
\[
 U_j=s_j I+2\sum_a x_{j,a}T_a,\qquad s_j^2+|x_j|^2=1.
\]
Conjugation induces a common rotation of all $x_j\in\R^3$. Put
\[
 g_{ij}=x_i\cdot x_j,\qquad t_{12j}=x_1\cdot(x_2\times x_j).
\]
The following real bank uses only finitely many specified Wilson words:
\begin{equation}
 \mathfrak F(U)=\bigl((s_j)_{j=1}^m,g_{11},g_{12},g_{22},
                    (g_{1j},g_{2j},t_{12j})_{j=3}^m\bigr).
 \label{v20:eq:bank}
\end{equation}
Indeed, with $s(V)=\tr(V)/2$,
\begin{equation}
 g_{ij}=s_i s_j-s(U_iU_j),\qquad
 t_{12j}=-\tfrac12\{s(U_1U_2U_j)-s(U_2U_1U_j)\}.
 \label{v20:eq:word-formulas}
\end{equation}
The order of the three factors is load bearing.
For a fixed initial configuration $U^0$, assume
\begin{equation}
 \delta_0=|x_1^0\times x_2^0|>0,
 \qquad c_0(U)=|\mathfrak F(U)-\mathfrak F(U^0)|^2.
 \label{v20:eq:c0}
\end{equation}
This is a bounded smooth, nonnegative, gauge-invariant multiplication
function. Let $\mathcal O_0$ be the full vertex-gauge orbit of $U^0$.
Give $\mathcal Q$ any fixed product bi-invariant Riemannian metric.

\begin{theorem}[Finite ordered words control the configuration orbit]
\label[theorem]{v20:thm:orbit-norm}
There are constants $c_*,C_* >0$, depending on the specified graph,
loop words, metric, and reference anchors, such that on all of $\mathcal Q$,
\begin{equation}
 c_*\operatorname{dist}(U,\mathcal O_0)^2
 \le c_0(U)\le C_*\operatorname{dist}(U,\mathcal O_0)^2.
 \label{v20:eq:orbit-norm}
\end{equation}
Moreover $|\nabla c_0|^2\le L_0c_0$ for a finite constant $L_0$.
The zero set of $c_0$ is precisely $\mathcal O_0$, and its Hessian is
strictly positive on the normal bundle of that orbit. No candidate anchor
margin or candidate small-angle condition is a hypothesis.
\end{theorem}
\begin{proof}
First consider a tuple sufficiently close in bank values to the reference.
Its first two vectors have positive length and nonzero cross product.
In their positively oriented Gram--Schmidt frame set
\[
 a=\sqrt{g_{11}},\quad b=g_{12}/a,\quad
 c=\sqrt{g_{22}-b^2},\quad
 x_1=(a,0,0),\quad x_2=(b,c,0).
\]
For $j\ge3$ the same frame gives exactly
\begin{equation}
 x_j=\left(\frac{g_{1j}}a,
       \frac{g_{2j}-b g_{1j}/a}{c},\frac{t_{12j}}{ac}\right).
 \label{v20:eq:reconstruction}
\end{equation}
The denominators remain bounded away from zero in a fixed neighborhood of
the \emph{reference} bank. These formulas and the scalar entries give a
Lipschitz inverse there, modulo one common rotation. Every proper rotation
lifts to $SU(2)$ conjugation. Thus the squared orbit distance of the tuple
is at most a fixed multiple of $c_0$. In particular equal bank values give
one simultaneous conjugation, and the presence of every fundamental chord
holonomy identifies the original vertex-gauge orbit.

To relate the tuple to the original links, gauge both configurations to
the chosen tree. All tree links then agree exactly, while the remaining
links are the chord holonomies. Conversely every rooted word is a product
of finitely many original links or inverses. Telescoping that product
bounds its chord distance by the sum of the constituent link distances.
Bi-invariance and the equivalence of chord and geodesic distances on
$SU(2)$ give both comparisons, with explicit finite word-length costs.
These are uses of the inherited tree reduction, not additional physical
operations on a record.

Each component of $\mathfrak F$ is uniformly Lipschitz on the compact
configuration space and is constant on $\mathcal O_0$. This proves the
upper bound. The preceding inverse formulas prove the lower bound when
$c_0$ is small. Elsewhere $c_0\ge\eta_0>0$, while the orbit distance is
at most the diameter of $\mathcal Q$, proving the global lower bound.
Since $c_0=\sum_\ell f_\ell^2$, where
$f_\ell=\mathfrak F_\ell-\mathfrak F_\ell(U^0)$,
\[
 |\nabla c_0|^2\le4c_0\sum_\ell|\nabla f_\ell|^2.
\]
This proves the gradient assertion. Finally $c_0$ and its first derivative
vanish on its zero set. Along a normal geodesic the lower bound in
\eqref{v20:eq:orbit-norm} forces its quadratic Taylor coefficient to be
positive. Along tangent vectors that coefficient is zero, since the
function vanishes on the orbit. Compactness gives a uniform positive
normal Hessian. The orbit is an embedded compact submanifold; its
stabilizer consists of the common center, because its anchors do not commute.
\end{proof}

\begin{example}[Pairwise traces lose an actual orientation]
Take three nonzero orthogonal imaginary quaternion vectors, and keep
$s_j$ fixed. Replacing $x_3$ by $-x_3$ preserves every $s_j$ and every
$g_{ij}$, but not $t_{123}$. A rotation fixing $x_1$ and $x_2$ fixes
$x_1\times x_2$, so the two tuples are not simultaneously conjugate.
Equation \eqref{v20:eq:word-formulas} separates them using two differently
ordered three-loop words. Separate trace or pairwise-Gram data would not
supply this initial comparison.
\end{example}

The theorem is a quantitative finite comparison, stronger for this purpose
than density of an invariant algebra. It is not a claim that every loop
word is a microscopic local Read. The long rooted words and their joint
values are explicitly retained data of the constructed preparation.

\section{An invariant phase from finite Gauss-compatible initial data}
\label{v20:sec:phase}

Equip each link with the metric for which $T_1,T_2,T_3$ are orthonormal,
rescaled by a fixed inertia $I_e>0$. Let $g$ be the resulting product
metric, $\dd\mu$ its normalized Haar volume, and $\Delta_g$ its
nonpositive Laplace--Beltrami operator. Fix the same compact gauge action
and the mechanical Hamiltonian
\begin{equation}
 \mathsf h(U,p)=\tfrac12|p|_{g^{-1}}^2+W(U),\qquad
 W(U)=\sum_f\beta_f\left(1-\tfrac12\operatorname{Re}\tr P_f(U)\right),
 \quad\beta_f>0.
 \label{v20:eq:classical-H}
\end{equation}
Every $P_f$ is its actual ordered loop product. In particular $W\ge0$ is
smooth and gauge invariant. The proof permits any such smooth nonnegative
invariant potential on this fixed compact manifold. It does not replace
a Wilson loop by an Abelian angle difference.

Let $p^0\in T^*_{U^0}\mathcal Q$ obey the classical Gauss constraint:
it annihilates every tangent to $\mathcal O_0$. This is the moment-map
zero condition for the cotangent gauge action.

\begin{lemma}[Gauss data extend to an invariant phase]
\label[lemma]{v20:lem:initial-phase}
Under the noncommuting-anchor condition in \eqref{v20:eq:c0}, there is a
real smooth gauge-invariant function $S_0$ on $\mathcal Q$ with
$\dd S_0(U^0)=p^0$. If $p^0=0$, one may take $S_0=0$.
\end{lemma}
\begin{proof}
The gauge stabilizer of $U^0$ is the common center; it acts identically
on the entire pure-link configuration manifold and on its cotangent space.
Consequently the translates of $p^0$ define a well-defined smooth
covector section $p_q^0$ on $q\in\mathcal O_0$. The Gauss condition
makes it normal to the orbit. Choose a gauge-equivariant tubular
neighborhood using the fixed invariant metric. In its normal coordinates
$q,v$ define $S_0(\exp_qv)=\chi(|v|^2)p_q^0(v)$, where $\chi$ is one
near zero and vanishes before the tubular boundary. Extension by zero is
smooth. Invariance of the metric, the covector section, and $\chi$ makes
this function gauge invariant. Its differential on the orbit has normal
part $p_q^0$ and zero tangential part. This proves the assertion without
adding a gauge-breaking reference field.
\end{proof}

\begin{proposition}[A common short-time invariant Hamilton--Jacobi family]
\label[proposition]{v20:prop:HJ}
There is $T_\Gamma>0$ and a real smooth gauge-invariant $S(t,U)$ on
$[0,T_\Gamma]\times\mathcal Q$ satisfying
\begin{equation}
 \partial_tS+\tfrac12|\nabla S|^2+W=0,\qquad S(0)=S_0.
 \label{v20:eq:HJ}
\end{equation}
Writing $b_t=\nabla S(t)$, its configuration flow $\Phi_t$ is a smooth
gauge-equivariant diffeomorphism. If
\begin{equation}
 c_t=c_0\circ\Phi_t^{-1},\qquad
 \mathcal O_t=\Phi_t(\mathcal O_0),
 \label{v20:eq:transport-cost}
\end{equation}
then $\partial_tc_t+b_t\cdot\nabla c_t=0$ and, uniformly on this interval,
\begin{equation}
 c_*'\operatorname{dist}(U,\mathcal O_t)^2\le c_t(U)
 \le C_*'\operatorname{dist}(U,\mathcal O_t)^2,
 \qquad |\nabla c_t|^2\le L_c c_t.
 \label{v20:eq:transport-coercivity}
\end{equation}
The Hamiltonian reference starting at $(U^0,p^0)$ has
$U_*(t)=\Phi_t(U^0)$ and $p_*(t)=\dd S(t,U_*(t))$ and obeys Gauss's
law. Constants and the interval depend on the graph, inertias, potential,
and initial phase, not on a quantum parameter.
\end{proposition}
\begin{proof}
Start the smooth Hamiltonian flow of \eqref{v20:eq:classical-H} on the
graph of $\dd S_0$. Compactness of its base and bounded initial momentum
give a uniform short-time smooth flow. Its projected configuration map
has derivative equal to the identity at time zero. Compactness and the
inverse-function theorem make it a diffeomorphism for a common sufficiently
small time. One direct global argument uses uniform local inverse
neighborhoods for near pairs of points and uniform closeness to the identity
for the complement of those neighborhoods. Thus distinct starting points
cannot acquire the same image for these times.

For a characteristic $(U(t),p(t))$ from $U_0$ define
\[
 S(t,U(t))=S_0(U_0)+\int_0^t
       \{p(s)\dot U(s)-\mathsf h(U(s),p(s))\}\,\dd s.
\]
Vary its starting point. Integration by parts and Hamilton's equations
cancel all interior terms, leaving $p(t)\delta U(t)$; the initial term
cancels $\dd S_0$. Hence $\dd S=p(t)$. Varying the terminal time gives
\eqref{v20:eq:HJ}. This is a direct local construction, not an assumption
of a classical continuum field solution. Gauge invariance of the initial
phase, metric, and potential gives equivariance of the characteristics
and invariance of $S$. In particular its differential annihilates gauge
orbit tangents at every time, which is Gauss's law.

The projected Hamiltonian velocity is $\nabla S$, so its map is also the
flow $\Phi_t$. Transport gives \eqref{v20:eq:transport-cost}.
Uniform bounds on the derivatives of $\Phi_t$ and its inverse transport
the two distance bounds. Moreover
$c_t=\sum_\ell(f_\ell\circ\Phi_t^{-1})^2$; differentiating this finite
sum of squares proves the last inequality, just as before.
\end{proof}

The family $S,c_t$ is a \emph{proof-only comparison}. The finite source
below depends on $W$, its inertias, the initial finite data, and its declared
regulator. It never applies $e^{\ii S(t)/\hbar}$ as a corrective gate and
never consults future reference values when choosing a recorded outcome.

\section{An exact positive gauge-invariant quantum relative observable}
\label{v20:sec:operator}

Use $0<\eta\le1$ for the physical quantum parameter in this fixed-graph
body. Define on $L^2(\mathcal Q,\dd\mu)$
\begin{equation}
 T_\eta=-\frac{\eta^2}{2}\Delta_g,
 \qquad E_\eta=T_\eta+W,
 \qquad L_{e,a}=-\ii\eta X_{e,a}.
 \label{v20:eq:quantum-H}
\end{equation}
Here $X_{e,a}$ differentiates the actual $e$th group factor in an invariant
orthonormal Lie direction before multiplication by $I_e^{-1/2}$.
These operators and all following products precede finite compression.
Let $U_S$ be multiplication by $e^{\ii S/\eta}$ and set
\begin{equation}
 \begin{split}
 \mathcal C_\eta(t)&=T_{S,\eta}(t)+c_t,\\
 T_{S,\eta}&=U_ST_\eta U_S^*
 =\sum_{e,a}\frac{(L_{e,a}-X_{e,a}S)^2}{2I_e}.
 \end{split}
 \label{v20:eq:positive-operator}
\end{equation}
This is a nonnegative differential-form operator. Gauge invariance of
$S,c_t$ and of $T_\eta$ makes it commute with the complete gauge action.
It is therefore an actual positive observable on the exact Gauss subspace.
The multiplication function $c_t$ is not an infimum of quantum kinetic
forms; the kinetic comparison is the displayed, separately specified square.

\begin{theorem}[Exact invariant-phase drift and a quantum-order remainder]
\label[theorem]{v20:thm:quantum}
For every smooth $u$ on $\mathcal Q$, put $\psi=U_Su$. On the interval of
\cref{v20:prop:HJ},
\begin{equation}
\begin{split}
 &\left\langle\psi,
 \left(\partial_t+\frac\ii\eta[E_\eta,\cdot]\right)
             \mathcal C_\eta\psi\right\rangle\\
 &\quad=-\eta^2\int\operatorname{Hess}_gS(\nabla\bar u,\nabla u)\,\dd\mu
       +\frac{\eta^2}{4}\int\Delta_g^2S\,|u|^2\,\dd\mu\\
 &\qquad\quad+\eta\operatorname{Im}\int
                 \bar u\,\langle\nabla c_t,\nabla u\rangle_g\,\dd\mu.
\end{split}
\label{v20:eq:exact-drift}
\end{equation}
Consequently, for finite constants independent of $\eta$,
\begin{equation}
 \left(\partial_t+\frac\ii\eta[E_\eta,\cdot]\right)\mathcal C_\eta
       \preceq K\mathcal C_\eta+K_2\eta^2 I,
 \qquad 0\preceq\mathcal C_\eta(t)\preceq C(E_\eta+I).
 \label{v20:eq:operator-bound}
\end{equation}
One may take
$K=\max\{2\|\operatorname{Hess}S\|_\infty+1,L_c/2\}$ and
$K_2=\|\Delta_g^2S\|_\infty/4$, with suprema also over time.
For the uncompressed quantum evolution,
\begin{equation}
 \langle\psi(t),\mathcal C_\eta(t)\psi(t)\rangle
 \le e^{Kt}\bigl(\langle\psi(0),\mathcal C_\eta(0)\psi(0)\rangle
                                      +K_2t\eta^2\bigr).
 \label{v20:eq:quantum-estimate}
\end{equation}
No restriction on the candidate configuration, its curvature, or a Higgs
radius enters these inequalities. They are not spatially uniform assertions.
\end{theorem}
\begin{proof}
All calculations can first be made on the smooth compact-manifold core.
Conjugate a Schr\"odinger solution by $U_S^*$ and use
\eqref{v20:eq:HJ}. Its exact transformed generator is
\[
 T_\eta-\ii\eta D_b,\qquad
 D_b=b\cdot\nabla+\tfrac12\operatorname{div}b,\qquad b=\nabla S.
\]
The divergence term is essential; $D_b$ is skew adjoint in the fixed
volume. The transformed comparison is $T_\eta+c_t$. The transported
multiplication part cancels because
$\partial_tc_t+[D_b,c_t]=\partial_tc_t+b\cdot\nabla c_t=0$.
Its remaining quantum commutator contributes exactly
$\eta\operatorname{Im}\int\bar u\nabla c_t\cdot\nabla u$.

For the kinetic part compute the form of $[D_b,T_\eta]$ by integration
by parts. Differentiating $b\cdot\nabla u$ contributes
$-\eta^2\int\operatorname{sym}\nabla b(\nabla\bar u,\nabla u)$.
The transport of $|\nabla u|^2$ cancels the multiplication by
$\operatorname{div}b$. The remaining term is
\[
 -\frac{\eta^2}{2}\operatorname{Re}\int
      u\,\nabla\bar u\cdot\nabla\operatorname{div}b
 =\frac{\eta^2}{4}\int\Delta_g\operatorname{div}b\,|u|^2.
\]
Since $b=\nabla S$, these are the first two terms of
\eqref{v20:eq:exact-drift}. This derivation is intrinsic: the metric
Hessian and Laplacian retain the compact-group geometry. No Cartesian
canonical commutation relation for group angles or absent Ricci remainder
has been assumed.

The absolute Hessian term is at most
$2\|\operatorname{Hess}S\|_\infty\langle u,T_\eta u\rangle$.
The last term is at most
\[
 \sqrt{2L_c\langle u,T_\eta u\rangle\langle u,c_tu\rangle}
 \le\langle u,T_\eta u\rangle+(L_c/2)\langle u,c_tu\rangle.
\]
This proves the form inequality. The second bound follows by the squared
triangle inequality for $-\ii\eta\dd-\dd S$, boundedness of $\dd S,c_t$,
and $W\ge0$. Integrating the differential inequality proves
\eqref{v20:eq:quantum-estimate}. For energy-form data, approximate in the
form norm by smooth data; on this compact manifold the elliptic kinetic
operator has the usual smooth spectral core, also obtained directly from
its discrete group spectrum and bounded smooth potential. The uniform
form domination and the integrated inequality pass to the limit. Mixed
states follow by positive sums and monotone approximation.
\end{proof}

\begin{remark}[This is not a missing-force ansatz]
The phase equation contains the complete ordered Wilson potential $W$.
The potential commutator cancels its actual derivative through
\eqref{v20:eq:HJ}; no current has been fitted to a recorded path. The price
is control of derivatives of a phase on the \emph{entire finite
configuration manifold}, rather than only of continuum reference fields.
That price is finite here, but must be estimated anew under spatial
refinement. The new construction closes an operator-existence gap without
claiming that the spatial estimate of Version 19 has already been promoted.
\end{remark}

\section{A nonempty exactly physical preparation, with no posterior projection}
\label{v20:sec:preparation}

Define the normalized vector
\begin{equation}
 \psi_{\eta,0}(U)=Z_\eta^{-1/2}
       \exp[-c_0(U)/(2\eta)]\exp[\ii S_0(U)/\eta],
 \qquad Z_\eta=\int e^{-c_0/\eta}\,\dd\mu.
 \label{v20:eq:preparation}
\end{equation}
All its inputs are the specified graph, action, initial holonomies, and
Gauss-compatible initial momenta. Both factors are gauge invariant.
There is no Haar average of density matrices and no rejection of an
unphysical trajectory.

\begin{theorem}[Initial orbit concentration and a uniform quantum energy moment]
\label[theorem]{v20:thm:preparation}
Every Gauss generator annihilates $\psi_{\eta,0}$ exactly. With
$k=\operatorname{codim}\mathcal O_0$ in $\mathcal Q$, there are
fixed-graph positive constants such that, for small $\eta$,
\begin{equation}
 c\eta^{k/2}\le Z_\eta\le C\eta^{k/2},\qquad
 \langle\psi_{\eta,0},\mathcal C_\eta(0)\psi_{\eta,0}\rangle\le C\eta,
 \qquad \|(E_\eta+I)^2\psi_{\eta,0}\|\le C.
 \label{v20:eq:initial-cost}
\end{equation}
Thus \eqref{v20:eq:quantum-estimate} gives an $O(\eta)$ comparison
uniformly on the stated time interval, in an exactly physical state.
\end{theorem}
\begin{proof}
Invariance proves the Gauss equations on the smooth core. In a tubular
neighborhood of the compact orbit, \cref{v20:thm:orbit-norm} and its normal
Hessian give $c|v|^2\le c_0\le C|v|^2$ in the $k$ normal coordinates.
The volume Jacobians have positive upper and lower bounds. Outside that
neighborhood $c_0$ has a positive minimum. Gaussian integration therefore
gives the two normalization bounds and
\[
 Z_\eta^{-1}\int c_0^j e^{-c_0/\eta}\,\dd\mu\le C_j\eta^j
 \quad\text{for every fixed nonnegative integer }j.
\]
Conjugation by the initial phase removes its momentum shift. The real
amplitude $a_\eta=Z_\eta^{-1/2}e^{-c_0/(2\eta)}$ has exactly
\[
 \langle a_\eta,T_\eta a_\eta\rangle
       =\frac18\int|\nabla c_0|^2|a_\eta|^2\,\dd\mu.
\]
The gradient bound and the $j=1$ moment prove the initial comparison cost.

For the higher energy moment set $F=-c_0/2+\ii S_0$; all its derivatives
are bounded on the fixed compact manifold. Direct differentiation gives
\[
 E_\eta e^{F/\eta}
  =\left[W-\tfrac12\langle\dd F,\dd F\rangle
                         -\tfrac\eta2\Delta_gF\right]e^{F/\eta}.
\]
The contraction here is complex bilinear, as dictated by differentiation.
The multiplier on the right and its derivatives up to order two are
bounded independently of $0<\eta\le1$. Applying $E_\eta$ once more
uses the same identity and the product rule; its additional terms are
$-\eta\nabla F\cdot\nabla$ of that multiplier and
$-\eta^2\Delta_g/2$ of it. They are uniformly bounded. Normalization
therefore gives $\|E_\eta^j\psi_{\eta,0}\|\le C$ for $j=1,2$,
proving the last bound. This argument does not hide a negative power of
the normalization or assume a higher-energy moment.
\end{proof}

\section{Holonomy, kinetic sources, and phase changes are controlled}
\label{v20:sec:sources}

Let $e_\eta(t)=\tr(\rho\mathcal C_\eta(t))$ for any physical density
matrix of finite energy. The configuration Born distribution satisfies
\begin{equation}
 \tr\!\left(\rho\one_{\{\operatorname{dist}(U,\mathcal O_t)>a\}}\right)
        \le C e_\eta(t)/a^2\quad(a>0).
 \label{v20:eq:Born-orbit}
\end{equation}
This follows directly from \eqref{v20:eq:transport-coercivity}. It is not
a simultaneous measurement of configurations and electric momenta.
For every fixed smooth gauge-invariant configuration function $f$,
\begin{equation}
 \tr\rho\,|f-f(U_*(t))|^2\le C_f e_\eta(t).
 \label{v20:eq:configuration-sources}
\end{equation}
In particular the full named Wilson-word bank, not just separate plaquette
traces, concentrates. The shifted kinetic summand independently controls
the covector error relative to $\dd S(t,U)$.

\begin{proposition}[Actual action insertions of momentum degree at most two]
\label[proposition]{v20:prop:sources}
Consider a declared, gauge-invariant symmetric source insertion defined
in the fixed Haar space by the quadratic form
\begin{equation}
 \mathfrak o_\eta[\psi]
 =\tfrac12\int\langle-\ii\eta\dd\psi,
                      A(U)(-\ii\eta\dd\psi)\rangle
  +\operatorname{Re}\int\bar\psi\,B(U)\cdot(-\ii\eta\dd\psi)
  +\int V_0(U)|\psi|^2.
 \label{v20:eq:source-form}
\end{equation}
Here $A$ is a smooth real symmetric tensor and $B,V_0$ are smooth real
coefficients; invariance is for the complete contracted form. No positivity
of this source is assumed. Its classical symbol is
$o(U,p)=\langle p,A p\rangle/2+B\cdot p+V_0$.
Uniformly on the reference interval,
\begin{equation}
 |\tr(\rho O_\eta)-o(U_*(t),p_*(t))|
       \le C_O(\sqrt{e_\eta(t)}+e_\eta(t)).
 \label{v20:eq:source-bound}
\end{equation}
This applies, in particular, to independently varied Wilson coefficients
and link electric-energy weights. More general metric insertions are
included only when they are independently defined by the displayed form.
\end{proposition}
\begin{proof}
For a vector write $\psi=e^{\ii S/\eta}u$. Then
$-\ii\eta\dd\psi=e^{\ii S/\eta}(-\ii\eta\dd u+\dd S\,u)$.
Expanding the \emph{form} gives a quadratic term bounded by
$C\langle\psi,T_{S,\eta}\psi\rangle$, mixed terms bounded by its square
root times $C\|\psi\|$, and the multiplication
$o(U,\dd S(t,U))|u|^2$. This last smooth invariant function equals the
classical reference source on $\mathcal O_t$. Its difference is bounded
pointwise by $C\operatorname{dist}(U,\mathcal O_t)\le C\sqrt{c_t}$.
Cauchy--Schwarz proves \eqref{v20:eq:source-bound}. For density matrices
use spectral decomposition and concavity of the square root. The proof
expands the original quadratic form, so ordering/divergence terms have
not been deleted by replacing it with a symbol. Differentiating an
independently specified action at fixed regulator supplies exactly such
forms for the stated examples. A parameter-dependent regulator would
have additional derivatives and is not covered by that convention.
\end{proof}

The classical symbol $|p-\dd S|^2/2+c_t$ also controls the distance to the
reference \emph{phase-space} gauge orbit on bounded reference sets. Choose
a configuration alignment, use smoothness and equivariance of $\dd S$,
and bound the aligned momentum difference by
$|p-\dd S(U)|+C\operatorname{dist}(U,\mathcal O_t)$.
The corresponding quantum assertion is the pair of positive form and
Born estimates above, not a joint probability law for noncommuting observables.

\begin{proposition}[Comparison phases can be restarted without changing the source]
\label[proposition]{v20:prop:patch}
Suppose at a time $t_0$ two invariant comparison phases have the same
differential on the reference orbit, and their nonnegative configuration
costs are each comparable to its squared distance. Then their operators
at that time satisfy
$\mathcal C_\eta^{(2)}\preceq C\mathcal C_\eta^{(1)}$
with $C$ independent of $\eta$. Consequently the short-time comparison
can be iterated over finitely many such phase charts along a reference
whose connections remain in the regular noncommuting stratum. No phase
change is applied to any attained physical state.
\end{proposition}
\begin{proof}
The smooth covector difference $\dd S_2-\dd S_1$ vanishes on the orbit,
so its squared norm is bounded by the first configuration cost. The
squared triangle inequality for the two shifted differential operators
gives $T_{S_2,\eta}\preceq2T_{S_1,\eta}+C c^{(1)}$ as forms.
The two distance comparisons give $c^{(2)}\le Cc^{(1)}$.
For a new chart, \cref{v20:lem:initial-phase} constructs an invariant
phase with the current Gauss covector. Compactness of a fixed regular
reference segment permits a finite subdivision with these phase charts;
finite multiplication of their constants preserves independence from
$\eta$. This is not a bound uniform in graph size or in a reference
approaching a singular orbit stratum.
\end{proof}

\section{Finite link cutoffs and gauge-commuting local outcomes without matter spectators}
\label{v20:sec:finite}

On every link retain the Peter--Weyl representations of spin $j\le J$,
where $J$ is an integer or half-integer, and let $P_J$ be the product
projection. It commutes with the left and right endpoint gauge actions,
with every link Casimir, and with $T_\eta$. Let
\[
 \mathcal K_J=\Ran P_J,\qquad E_{\eta,J}=P_JE_\eta P_J|_{\mathcal K_J}.
\]
The original Wilson products and source insertions are formed before
compression. This tensor cutoff does not commute with $W$ in general.

\begin{lemma}[Explicit compact-link tensor-tail and unitary approximation costs]
\label[lemma]{v20:lem:cutoff}
With $B_\eta=E_\eta+I$ and constants depending on the fixed graph,
\begin{equation}
 \|(I-P_J)v\|_{B_\eta}^2
       \le\frac{C}{1+\eta^2(J+1)^2}\|B_\eta v\|^2.
 \label{v20:eq:cutoff-tail}
\end{equation}
For the preparation of \cref{v20:thm:preparation}, let
$\psi_{\eta,J,0}=P_J\psi_{\eta,0}/\|P_J\psi_{\eta,0}\|$ and let
$\psi_J(t)$ and $\psi(t)$ be the finite and uncompressed quantum
trajectories. For sufficiently large $J$,
\begin{equation}
 \sup_{t\le T}\|\psi_J(t)-\psi(t)\|_{B_\eta}^2
       \le C_T\frac{(1+T/\eta)^2}{1+\eta^2(J+1)^2}.
 \label{v20:eq:cutoff-evolution}
\end{equation}
The normalized finite preparation remains exactly physical.
\end{lemma}
\begin{proof}
Set $A_\eta=I+T_\eta$. Boundedness and nonnegativity of $W$ imply
$A_\eta\preceq B_\eta\preceq(1+\|W\|_\infty)A_\eta$ and
$\|A_\eta v\|\le(1+\|W\|_\infty)\|B_\eta v\|$.
On the complement of $P_J$ at least one spin exceeds $J$; the corresponding
Casimir contributes at least $c_\Gamma\eta^2(J+1)^2$ to $T_\eta$.
Its joint spectral expansion with $P_J$ gives
\[
 \|A_\eta^{1/2}(I-P_J)v\|^2
 \le\frac{C}{1+\eta^2(J+1)^2}\|A_\eta v\|^2.
\]
This proves \eqref{v20:eq:cutoff-tail}. Apply the inherited Ritz/Duhamel
proof of \cref{v14:thm:Galerkin}, replacing its $D_h/R$ by the displayed
tail constant and its quantum scale by $\eta$. The uniform bound on
$\|B_\eta^2\psi_{\eta,0}\|$ is supplied by
\cref{v20:thm:preparation}; no additional initial moment is assumed.
For large $J$ the normalization threshold is satisfied and this proves
\eqref{v20:eq:cutoff-evolution}. Commutation of $P_J$ with every gauge
transformation proves its physical-sector statement. No exact commutator
identity is incorrectly asserted for the noncommuting tensor cutoff.
\end{proof}

For local instruments suppose, as on a finite periodic cubic lattice, each
link is contained in a simple plaquette of bounded length. Decompose
$E_{\eta,J}$ into its positive electric-link and Wilson-plaquette terms.
An electric context may be enlarged to one incident plaquette, on whose
other links its action term is the identity. This is a bounded-radius
support enlargement, not a new action term.

\begin{proposition}[Pure-gauge contexts supply the required neutral multiplicity]
\label[proposition]{v20:prop:multiplicity}
For $J\ge2$ a simple plaquette context has at least five linearly
independent vectors invariant under every incident vertex gauge.
Consequently every preceding local action term admits the six-outcome,
gauge-commuting construction of \cref{v15:thm:locked}, with no extra
matter site or matter species.
\end{proposition}
\begin{proof}
For $j=0,1/2,1,3/2,2$ take the five functions $\chi_j(P_f)$ on the
plaquette links. Each is invariant under all endpoint gauges, since their
interior factors telescope and a trace removes conjugation at the base.
The spin-$j$ matrix-product expression has representation $j$ on each
link, so it lies in the link tensor cutoff when $J\ge2$.
Integrating all but one link by Haar invariance reduces their inner
products to character orthogonality on $SU(2)$; hence these five functions
are orthonormal. The trivial representation therefore has multiplicity
at least five in this context. Each local action term preserves that
invariant subspace. Choose the auxiliary matrix units of
\cref{v15:thm:locked} in its multiplicity space, not in a gauge-charged
basis. The inherited proof gives exact locking, independence of all six
operators, and Gauss preservation on every outcome. The auxiliary
operations may act on the additional plaquette links; this is openly part
of the constructed local context.
\end{proof}

Let $M$ be the number of contexts and write
$E_{\eta,J}/\eta=\sum_{\alpha=1}^M H_\alpha$.
At acceptance choose a context uniformly and retain its original six
locked outcomes with local step $\tau=M\epsilon$.
The complete accepted label is $(\alpha,e)$, not six global labels.
Set
\begin{equation}
 A_{\eta,J}=1+\|E_{\eta,J}\|,\qquad
 \mathfrak C_{\eta,J}=1+
   \left(\sum_\alpha\|H_\alpha\|\right)^2
     +M\sum_\alpha\|H_\alpha\|^2+M^2.
 \label{v20:eq:local-cost}
\end{equation}
The supplied non-Poisson timing law and its physical units are unchanged:
$d=4\vartheta\epsilon/11$, $0<\vartheta\le1$.
All context choices, initial preparation, and semiclassical parameters
remain declared source data, not consequences of the bare locking identities.

\section{Resolved physical records have a non-Abelian orbit and source limit}
\label{v20:sec:trajectory}

\begin{theorem}[Fixed-lattice non-Abelian quantum orbit closure on the actual clock]
\label[theorem]{v20:thm:resolved}
Initialize the preceding finite local source at
$|\psi_{\eta,J,0}\rangle\langle\psi_{\eta,J,0}|$ and retain each actual
conditional state $\rho_j$. For $M\epsilon\le\tau_0$ put
\begin{equation}
 \begin{split}
 e_j&=\tr(\rho_jP_J\mathcal C_\eta(jd)P_J),\\
 \mathfrak b_{\eta,J,\epsilon}
 &=\eta+\frac{(1+T/\eta)^2}{1+\eta^2(J+1)^2}
                     +\epsilon A_{\eta,J}\mathfrak C_{\eta,J}.
 \end{split}
 \label{v20:eq:budget}
\end{equation}
On the fixed comparison interval, with arbitrary initial clock phase,
\begin{equation}
 \max_{j\le T/d}\E e_j\le C_T\mathfrak b_{\eta,J,\epsilon},\qquad
 \PP\{\max_{j\le T/d}e_j>u\}
       \le C_T\mathfrak b_{\eta,J,\epsilon}/u\quad(u>0).
 \label{v20:eq:resolved}
\end{equation}
The constants are independent of $\eta,J,\epsilon,\vartheta$ and finite
matrix dimension, but depend on the fixed graph, its action coefficients,
and the comparison phases. Every recorded outcome preserves the exact
Gauss subspace. Configuration-orbit concentration, all fixed smooth
invariant Wilson observables, and every independently defined source of
\cref{v20:prop:sources} converge when the displayed budget vanishes.
\end{theorem}
\begin{proof}
The exact uncompressed trajectory has comparison at most $C_T\eta$ by
\cref{v20:thm:quantum,v20:thm:preparation}. Form domination of
$\mathcal C_\eta$ by $B_\eta$ and the squared triangle inequality,
together with \cref{v20:lem:cutoff}, give
\[
 \sup_{t\le T}\langle\psi_J(t),\mathcal C_\eta(t)\psi_J(t)\rangle
 \le C_T\left[\eta+
           \frac{(1+T/\eta)^2}{1+\eta^2(J+1)^2}\right].
\]
The inherited local locked-clock comparison
\cref{v14:thm:trajectory,v15:thm:source} bounds the fidelity error
$\xi_j=1-\langle\psi_J(jd),\rho_j\psi_J(jd)\rangle$
in expectation and maximal probability by
$C_T\epsilon\mathfrak C_{\eta,J}$.
For the positive compressed comparison, the rank-one projection $P_*(t)$
onto $\psi_J(t)$ gives
\[
 P_J\mathcal C_\eta P_J
 \preceq2P_*\mathcal C_\eta P_*
     +2(I-P_*)P_J\mathcal C_\eta P_J(I-P_*).
\]
The second term has norm at most $C_TA_{\eta,J}$ on its support.
Pairing with $\rho_j$ bounds $e_j$ by twice the preceding deterministic
comparison plus $C_TA_{\eta,J}\xi_j$. The inherited maximal bound gives
\eqref{v20:eq:resolved}, splitting off thresholds smaller than the
deterministic term. No union bound over ticks is introduced.
The initial state and every local Kraus operator are physical, proving
outcome-by-outcome Gauss support. Finally the pointwise configuration and
source-form estimates apply to any vector in the tensor subspace, provided
products precede compression. They therefore apply to every attained state.
\end{proof}

For source errors larger than $0<u\le1$, the corresponding maximal-in-time
probability is at most $C_{T,O}\mathfrak b/u^2$, by
\eqref{v20:eq:source-bound}. At fixed times their expected absolute errors
are at most $C_{T,O}(\sqrt{\mathfrak b}+\mathfrak b)$.
Since $e_j\le C_TA_{\eta,J}$, integration of the maximal tail yields
\begin{equation}
 \E\max_{j\le T/d}e_j
 \le C_T\mathfrak b\left[1+
                \log\left(2+\frac{A_{\eta,J}}{\mathfrak b}\right)\right]
 \quad(0<\mathfrak b\le1/2).
 \label{v20:eq:max-mean}
\end{equation}
All Born bounds refer to each attained conditional state. No unobserved
classical configuration trajectory or repeated position measurement is
introduced by this statement.

\begin{corollary}[An explicit finite local scale choice]
\label[corollary]{v20:cor:scales}
For this fixed graph one may choose
\begin{equation}
 J_\eta=\lceil\eta^{-3}\rceil,\qquad
 \epsilon_\eta=\eta^{15},\qquad
 d_\eta=\frac{4\vartheta_\eta}{11}\eta^{15}.
 \label{v20:eq:scales}
\end{equation}
Then $\mathfrak b_{\eta,J_\eta,\epsilon_\eta}=O(\eta)$ and
$\E\max_j e_j=O(\eta(1+|\log\eta|))$.
Every cutoff is a finite gauge-equivariant tensor carrier with original
local locked outcomes and exact physical support.
\end{corollary}
\begin{proof}
The fixed-graph Casimir bound and bounded Wilson multiplication give
\[
 A_{\eta,J}\le C_\Gamma(1+\eta^2J^2),\qquad
 \mathfrak C_{\eta,J}\le
                 C_\Gamma\left[1+\frac{(1+\eta^2J^2)^2}{\eta^2}\right].
\]
With $J=\lceil\eta^{-3}\rceil$, these have orders $\eta^{-4}$ and
$\eta^{-10}$. The tensor-evolution term is $O(\eta^2)$, while
$\epsilon A\mathfrak C=O(\eta^{15-4-10})=O(\eta)$.
The first term of the budget is $\eta$. The normalization and local-step
thresholds hold for small $\eta$, since the graph and $M$ are fixed.
Substitution into \eqref{v20:eq:max-mean} gives the logarithm. These are
sufficient existence powers, not a computational efficiency result.
\end{proof}

\section{Noncommuting references and the spatial-conditioning obstruction}
\label{v20:sec:boundary}

The initial population need not be flat or static. For example on a fixed
cubic lattice take the homogeneous compact links
$U_i^0=e^{h a_iT_i}$ with all $a_i\ne0$, and initially zero electric
momenta. Choose the axes and small nonzero $ha_i$ so that the rooted
$12$ and $23$ plaquettes have nonparallel imaginary quaternion vectors.
Their leading vectors are proportional to $h^2a_1a_2T_3$ and
$h^2a_2a_3T_1$, respectively. Thus the anchor condition holds. Gauss data
hold because the electric momenta are zero, and $S_0=0$ suffices.
For generic such values the Wilson force is nonzero, so the subsequent
same-action reference develops electric momentum. The exact compact
force calculation is already supplied by \cref{v19:prop:compact}; setting
its scalar to zero gives this pure-gauge reference. It is not a new
homogeneous ODE theorem. For a sufficiently short interval its two anchor
plaquettes remain noncommuting. The preceding construction consequently
has genuinely non-Abelian magnetic data and nonstatic dynamics, not an
embedded $U(1)$ example.

\begin{proposition}[Unscaled small-loop data can lose an eighth spatial power]
\label[proposition]{v20:prop:conditioning}
The constant controlling orbit distance by the bank error need not remain
uniform when the reference anchors shrink. On a three-loop tuple take
$x_1^0=\sin\alpha\,e_1$, $x_2^0=\sin\beta\,e_2$, $U_3^0=I$,
and vary only $U_3=e^{tT_3}$. With
$\delta_0=|\sin\alpha\sin\beta|$, one has
\begin{equation}
 c_0(t)=(\cos(t/2)-1)^2+\delta_0^2\sin^2(t/2),\qquad
 d_{\rm tuple}(t)^2=4(1-\cos(t/2)),
 \label{v20:eq:conditioning-example}
\end{equation}
where $d_{\rm tuple}^2$ is the minimum sum of squared Hilbert--Schmidt
link distances under simultaneous conjugation. Hence
\begin{equation}
 \lim_{t\to0}\frac{d_{\rm tuple}(t)^2}{c_0(t)}=\frac2{\delta_0^2}.
 \label{v20:eq:conditioning-floor}
\end{equation}
When the two reference angles have size comparable to $h^2$, this inverse
comparison costs at least a constant times $h^{-8}$ on this family.
\end{proposition}
\begin{proof}
All anchor data remain unchanged. The third scalar trace changes by
$\cos(t/2)-1$, both anchor dot products with $x_3$ vanish, and its
oriented triple is $\delta_0\sin(t/2)$ up to a sign. This proves the
first formula. The third reference loop is the identity and hence is
fixed by every conjugation. Its squared distance is always
$4(1-\cos(t/2))$; the identity conjugation makes the two other contributions
zero, proving the second formula. Taylor expansion proves the ratio and
the stated scaling. This is an explicit conditioning example, not a
universal optimal lower bound for every differently normalized source bank.
\end{proof}

\paragraph{What has been closed.}
There is now a positive gauge-invariant operator controlling configuration
\emph{orbit} information and an independently aligned momentum square,
with a proved propagation inequality, exact Gauss-compatible preparation,
finite tensor cutoff, and local resolved-clock realization. It does not
require a Higgs frame. The former quantum-orbit gap is therefore not an
absolute obstruction on a fixed non-Abelian gauge graph.

\paragraph{What has not been closed.}
The spatially uniform non-Abelian continuum theorem is not proved by this
fixed-graph construction. The phase derivatives, the Hamilton--Jacobi
interval or finite patch constants, the loop-word metric, and the
preparation volume constants can all depend on refinement. Equation
\eqref{v20:eq:conditioning-floor} shows a concrete source of that dependence.
Dividing small Wilson quantities by powers of $h$ would change the declared
reader normalization and its precision budget; it is not a free repair.
A useful next calculation is to control these constants in a local,
physically normalized covariant comparison, or to construct a spatially
uniform positive quantum comparison directly. The full Yang--Mills--Higgs
matter extension, Einstein backreaction and constraints, and chiral
fermions remain separate. Likewise no theorem here shows that the
independently given primitive law selects the declared operator panels,
contexts, phase-space preparation, or quantum scale.

\paragraph{Verification and preservation.}
Every byte of Version 19 before its final end-document command is retained.
New labels have prefix \texttt{v20:}. The companion diagnostics check the
ordered quaternion trace identities and frame reconstruction; random
simultaneous-gauge invariance; the sharp shrinking-anchor example; the
intrinsic drift on an $SU(2)$ class-function test with its Haar measure;
the physical orbit-packet preparation scale; the pure-gauge character
multiplicity; and the finite cutoff/clock power budget. These are symbolic
and finite numerical checks of the written arguments, not machine-checked
proofs, a complete native-source identification, or an Einstein--SM simulation.

\begingroup
\renewcommand{\refname}{Additional references for Version 20}
\begin{thebibliography}{9}
\small\raggedright
\bibitem{v20:TW2}
T.~Thiemann and O.~Winkler,
\emph{Gauge Field Theory Coherent States (GCS): II. Peakedness Properties},
2000 preprint. \href{https://arxiv.org/abs/hep-th/0005237}{arXiv:hep-th/0005237}.
Background for gauge-field coherent preparation and concentration. The
invariant Wilson-bank packet and its kinetic cost above are proved directly.
\bibitem{v20:TW3}
T.~Thiemann and O.~Winkler,
\emph{Gauge Field Theory Coherent States (GCS): III. Ehrenfest Theorems},
2000 preprint. \href{https://arxiv.org/abs/hep-th/0005234}{arXiv:hep-th/0005234}.
Background for semiclassical gauge observables. No ensemble-only or
infinitesimal theorem from this reference substitutes for the resolved-clock
maximal comparison proved here.
\end{thebibliography}
\endgroup
\addtocontents{toc}{\protect\endgroup}
% END IMMUTABLE VERSION 20 BODY
% APPEND VERSION 21 HERE, BEFORE THE FINAL END-DOCUMENT COMMAND.


% BEGIN IMMUTABLE VERSION 21 BODY
\clearpage
\begin{center}
{\Large\bfseries Version 21: Spatial normalization and the real-phase caustic}\\[0.5em]
{\large A scale-aware invariant bank, an exact Wilson obstruction, and a positive phase-space alternative}
\end{center}
\makeatletter
\addtocontents{toc}{\protect\begingroup}
\addtocontents{toc}{\protect\makeatletter}
\addtocontents{toc}{\protect\def\protect\l@section{\protect\bfseries\protect\@dottedtocline{1}{0em}{2.5em}}}
\makeatother
\addcontentsline{toc}{part}{Version 21: Scale-aware Wilson data and caustic-free quadratic comparison}

\section{Audit: two different obstructions to spatially uniform phase comparison}
\label{v21:sec:audit}

Version 20 leaves the normalization of shrinking loop data and the time interval
of a real Hamilton--Jacobi phase dependent on the spatial graph. These are not
one problem. A normalization can repair a small signal without preventing a
configuration-space caustic. This iteration separates the two issues and
computes both costs before trying to extend that phase argument.

\paragraph{Inputs and nonduplication.}
The invariant-word formulas and finite-frame reconstruction are
\cref{v20:eq:word-formulas,v20:eq:reconstruction}; their fixed-graph result is
not claimed anew. The invariant phase, its quantum identity, and its local
resolved implementation remain \cref{v20:prop:HJ,v20:thm:quantum,v20:sec:trajectory}.
The classical non-Abelian comparison of Version 19 is a different,
reference-covariant energy estimate. The supplied classical closure manuscript
\cite{CC}, at \nolinkurl{v9:lem:HJ}, already constructs a homogeneous
scalar--gravity Hamilton--Jacobi branch; this is not a result about the
projected characteristics of the refining Wilson configuration manifold.
Targeted text searches of these files found no caustic or Riccati calculation
of the kind below. This is a limited overlap check, not an exhaustive audit of
all manuscripts or a claim of general mathematical novelty.

\paragraph{New conclusions and exact scopes.}
First, a weighted, scale-declared ordered-loop bank controls normalized tuple
orbit distance with constants independent of the number of loops and their
small amplitudes. Its raw-trace precision and derivative costs are retained.
Second, for the physically normalized \emph{nonlinear compact Wilson action},
the global zero-initial-phase Hamilton--Jacobi solution cannot remain smooth
past a time proportional to the lattice spacing. This is an exact
characteristic obstruction, not an estimate of numerical stiffness.
Third, a positive full phase-space comparison remains regular through these
caustics. It is proved for quadratic Hamiltonians and evaluated on every
radiative and harmonic mode of the actual Wilson Hessian. The resulting sharp
quantum cost is of order $\eta h^{-4}$ in three spatial dimensions.

The last calculation is a theorem about the \emph{quadratic diagnostic}, not
an identification of that diagnostic with the full non-Abelian source. It does
not extend the nonlinear, exactly Gauss-preserving instrument conclusion of
Version 20 to a common spatial continuum interval. No source transition,
physical field, or nonlinear action is changed to evade the obstruction.

\paragraph{Attribution.}
Canonical lattice gauge Hamiltonians, caustics of real generating functions,
and positive Gaussian propagation are established mathematics; see
\cite{v21:KS,v21:CR}. The proofs below establish their precise normalization
and consequences for this lineage. In particular, the finite oscillator
calculation is not advertised as a new theory of free gauge fields. Its role
is to decide which proposed upstream estimate is possible and to identify the
correct energy and uncertainty cost for a replacement.

\section{A uniformly conditioned bank with its physical scales declared}
\label{v21:sec:bank}

Take a finite tuple of the actual rooted $SU(2)$ holonomies used in Version 20,
including all fundamental chord holonomies whenever the whole graph orbit is
the target. Write
\[
 U_j=s_jI+2x_j^aT_a,\qquad T_a=-\ii\sigma_a/2,\qquad
 s_j^2+|x_j|^2=1,\qquad 1\le j\le m.
\]
Fix \emph{before evaluating the record} scales $0<\alpha_j\le1$ and positive
weights $w_j$, with $w_1=w_2=1$ and $\sum_{j=3}^m w_j\le1$. Put
$y_j=x_j/\alpha_j$. Simultaneous conjugation rotates all $y_j$ by the same
proper rotation. For a specified reference tuple suppose
\begin{equation}
 |y_j^*|\le M,\qquad |y_1^*\times y_2^*|\ge\delta>0.
 \label{v21:eq:reference-window}
\end{equation}
There is \emph{no} boundedness or anchor condition on the candidate $y_j$.
Define
\[
 a=(|y_1|^2,y_1\cdot y_2,|y_2|^2),\quad r_j=|y_j|^2,\quad
 v_j=(y_1\cdot y_j,y_2\cdot y_j,(y_1\times y_2)\cdot y_j)\quad(j\ge3).
\]
The extra diagonal data $r_j$ will control candidates far from the reference
frame. The scale-aware squared error is
\begin{equation}
 \begin{split}
 c_{\alpha,w}(U)={}&|a-a^*|^2+
 \sum_{j=1}^m w_j\bigl[(s_j-s_j^*)^2+(r_j-r_j^*)^2\bigr]\\
 &+\sum_{j=3}^m w_j|v_j-v_j^*|^2,
 \end{split}
 \label{v21:eq:bank-cost}
\end{equation}
and its target distance is
\begin{equation}
 d_{\alpha,w}(U,U^*)^2
 =\inf_{R\in SO(3)}\sum_{j=1}^m w_j
       \bigl[(s_j-s_j^*)^2+|y_j-Ry_j^*|^2\bigr].
 \label{v21:eq:tuple-distance}
\end{equation}
These are functions of one joint record. Separate loop marginals are not
substituted for the mixed observations.

\begin{theorem}[Scale-aware global control with no candidate anchor hypothesis]
\label[theorem]{v21:thm:normalized}
Under \eqref{v21:eq:reference-window}, a constant depending only on $M,\delta$
satisfies
\begin{equation}
 \boxed{d_{\alpha,w}(U,U^*)^2\le C_{M,\delta}\,c_{\alpha,w}(U)}
 \label{v21:eq:normalized-bound}
\end{equation}
for every candidate tuple. It is independent of $m$, the $w_j$, and all
$\alpha_j$. On any declared candidate window $|y_j|\le M_1$ there is also
$c_{\alpha,w}\le C_{M,M_1}d_{\alpha,w}^2$. The zero set is the simultaneous
conjugacy orbit. If the tuple contains every fundamental chord, its zero set
on the original configuration manifold is the full vertex-gauge orbit.
No uniform conversion to an original link metric through long tree paths is
asserted by this theorem.
\end{theorem}
\begin{proof}
Choose $\varepsilon_0>0$, depending only on $M,\delta$, so that
$|a-a^*|^2<\varepsilon_0$ forces both ordered anchor frames to be
nondegenerate with a common margin. This follows from continuity of
$a_1a_3-a_2^2$ and its reference lower bound $\delta^2$. In their respective
oriented Gram--Schmidt coordinates write
\[
 y_1=(A,0,0),\quad y_2=(B,C,0),\quad
 A=\sqrt{a_1},\quad B=a_2/A,\quad C=\sqrt{a_3-B^2}.
\]
All denominators $A,C,AC$ have positive bounds depending only on $M,\delta$.
Every remaining vector has coordinates
\[
 y_j=\left(\frac{v_{j,1}}A,
       \frac{v_{j,2}-Bv_{j,1}/A}{C},\frac{v_{j,3}}{AC}\right).
\]
The linear map from $v_j$ to these coordinates, and its derivative in $a$,
are uniformly bounded on the anchor neighborhood. Subtract reference values
by writing the difference as the new inverse applied to $v_j-v_j^*$ plus
the change of inverse applied to $v_j^*$. The latter vector is bounded using
only the reference bound $M$. Thus one common proper rotation gives
\[
 |y_j-Ry_j^*|^2\le C\bigl(|v_j-v_j^*|^2+|a-a^*|^2\bigr)
 \quad(j\ge3),
\]
and the corresponding anchor error is at most $C|a-a^*|^2$.
Sum with the weights; their total is at most three. This proves
\eqref{v21:eq:normalized-bound} in this neighborhood, without an upper
bound on any candidate $v_j$ or $y_j$.

Outside it, $c_{\alpha,w}\ge\varepsilon_0$. The retained diagonal bank gives
\[
 \sum_jw_j|y_j|^2\le3M^2+\sqrt{3c_{\alpha,w}}.
\]
Using any rotation in \eqref{v21:eq:tuple-distance}, the vector errors are
therefore at most $C_M(1+\sqrt{c_{\alpha,w}})$; the scalar errors are already
in $c_{\alpha,w}$. For $c_{\alpha,w}\ge\varepsilon_0$ this is bounded by
$C_{M,\delta}c_{\alpha,w}$. This establishes the global assertion. Notice
why diagonal data for every loop were retained: small weights on remote
candidate vectors do not otherwise give this uniform far-field estimate.

For the reverse bound on the specified bounded window, subtract each dot
product and ordered triple by telescoping one factor at a time. All remaining
factors are bounded by $M+M_1$, and the sum of nonanchor weights is at most one.
The sum of squared bank errors is bounded by the weighted sum of squared
vector and scalar errors, for each common rotation. Take its minimum. Finally,
zero bank error gives equality of the oriented reconstructed vectors and
scalar parts. Proper rotations lift to $SU(2)$ conjugations. The graph-orbit
assertion is exactly the inherited tree reduction, with no claim that its
metric constants are uniform in word length.
\end{proof}

All entries are measured functions of ordinary ordered Wilson traces. With
$s(V)=\tfrac12\operatorname{tr}V$, the exact formulas are
\begin{equation}
 \begin{split}
 y_i\cdot y_j&=\frac{s_i s_j-s(U_iU_j)}{\alpha_i\alpha_j},\\
 |y_j|^2&=\frac{1-s_j^2}{\alpha_j^2},\\
 (y_1\times y_2)\cdot y_j
 &=-\frac{s(U_1U_2U_j)-s(U_2U_1U_j)}{2\alpha_1\alpha_2\alpha_j}.
 \end{split}
 \label{v21:eq:normalized-traces}
\end{equation}
These reuse the inherited quaternion identities; the new estimate is their
scale- and population-uniform weighted use. For plaquettes of size $h^2$ one
may declare $\alpha_j=h^2$; macroscopic holonomies need not have that scale
and may instead retain $\alpha_j=1$. A normalized curvature floor such as
\eqref{v21:eq:reference-window} is still a reference condition, not a result
about arbitrary continuum gauge fields.

\section{Normalization repairs a signal only when precision and derivatives are paid}
\label{v21:sec:precision}

\begin{proposition}[Raw-trace uncertainty at the small-loop scale]
\label[proposition]{v21:prop:precision}
Suppose all $\alpha_j=h^2$, and every raw scalar trace in
\eqref{v21:eq:normalized-traces}, for candidate and reference, is enclosed
with absolute error at most $e\le1$. Let $\widehat c_{h,w}$ be the cost
computed from those approximations. Then
\begin{equation}
 d_{h^2,w}(U,U^*)\le
 C_{M,\delta}\bigl(\sqrt{\widehat c_{h,w}}+e h^{-6}\bigr).
 \label{v21:eq:precision}
\end{equation}
In particular, $e=o(h^6)$ is a sufficient worst-case precision for certifying
vanishing normalized orbit error from a vanishing measured cost. This is not
a lower bound for all possible specialized measurement strategies.
\end{proposition}
\begin{proof}
Since every true scalar trace has absolute value at most one, an approximate
product $s_i s_j$ differs by at most $3e$, after enlarging constants to cover
the additional single trace. The diagonal data have the same bound before
division. Thus each normalized quadratic entry has error $Ce h^{-4}$;
each normalized triple has error at most $e h^{-6}$. Scalar errors are at
most $e$. There is a fixed number of entries per loop, and total weight at
most three. The weighted bank error is consequently at most $Ce h^{-6}$,
for each of the two tuples. The triangle inequality in that weighted
Euclidean space bounds $\sqrt{c_{h,w}}$ by
$\sqrt{\widehat c_{h,w}}+Ce h^{-6}$. Apply
\cref{v21:thm:normalized}.
\end{proof}

\begin{proposition}[A conservative derivative cost in the physical kinetic metric]
\label[proposition]{v21:prop:gradient}
Give every original link inertia $I_e=h$, and suppose each rooted loop word
has length at most $L_h$, counting repetitions and inverses. With all
$\alpha_j=h^2$, the same cost satisfies on the entire compact configuration
space
\begin{equation}
 |\nabla c_{h,w}|_{g_h}^2\le
 C L_h^2h^{-13} c_{h,w}.
 \label{v21:eq:gradient-cost}
\end{equation}
The constant is independent of the number of retained loops. This bound is
not asserted to be sharp on a small-field region.
\end{proposition}
\begin{proof}
For a length-$\ell$ word, a derivative in one link direction is the sum of
its occurrences in that word. Each quaternion scalar-trace derivative has
Euclidean Lie-gradient norm at most $1/2$ per occurrence. Hence its product
metric gradient is bounded by $\ell/(2\sqrt h)$, including repeated or
inverted letters. The ordinary product rule applies to products of traces.
The longest words in \eqref{v21:eq:normalized-traces} have length at most
$3L_h$. The worst normalization is $h^{-6}$, so the sum of squared gradients
of the weighted bank components is at most $CL_h^2h^{-13}$; the weights
sum to a fixed constant. Finally, for a squared bank difference,
$|\nabla\sum_a w_a f_a^2|^2\le
4(\sum_aw_af_a^2)(\sum_aw_a|\nabla f_a|^2)$.
This proves the assertion on every candidate, not only a bounded normalized
window.
\end{proof}

The normalized metric has removed the eighth-power inverse signal loss in
\cref{v20:prop:conditioning}, for the newly declared norm. It has not removed
the cost from a Hamilton--Jacobi proof: the gradient coefficient in its
positive operator estimate can still grow. Nor does it identify small local
plaquettes with all global connection information. These distinctions are
independent of the dynamical obstruction proved next.

\section{The actual Wilson action forces a real-phase caustic at time of order h}
\label{v21:sec:caustic}

On the periodic cubic lattice with $N\ge4$ even and $h=N^{-1}$, take the
pure-gauge mechanical Hamiltonian in the normalization of Version 19,
\begin{equation}
 H_h(U,P)=\sum_{e,a}\frac{(P_e^a)^2}{2h}
     +\frac4h\sum_f\left(1-\frac12\operatorname{Re}\tr P_f(U)\right).
 \label{v21:eq:Wilson}
\end{equation}
Here $P_f(U)$ is the ordered plaquette holonomy and $P_e^a$ is canonical
cotangent momentum; the notation distinguishes these different uses of $P$.
Every link has the bi-invariant metric with inertia $h$.
Let $B$ be the tail-minus-head vertex incidence and $C$ the oriented
edge-to-plaquette incidence, so $BC^{\mathsf T}=0$.

\begin{lemma}[The normalized transverse Hessian of the compact action]
\label[lemma]{v21:lem:Hessian}
At the equilibrium $U_e=I$, $P_e=0$, write
$U_e=\exp(u_e^aT_a)$ and use orthonormal canonical variables
$q_e^a=\sqrt h\,u_e^a$, $p_e^a=P_e^a/\sqrt h$. The quadratic Taylor
Hamiltonian is exactly
\begin{equation}
 H_h^{(2)}(q,p)=\frac12|p|^2+
                         \frac1{2h^2}|Cq|^2.
 \label{v21:eq:quadratic}
\end{equation}
On each nonzero spatial Fourier mode $k\in(\Z/N\Z)^3$, restricted to
$Bq=Bp=0$, its frequency is
\begin{equation}
 \Omega_h(k)=\frac2h
     \left(\sum_{i=1}^3\sin^2(\pi k_i/N)\right)^{1/2},
 \label{v21:eq:frequency}
\end{equation}
with two transverse polarizations and three color components. There are also
nine constant harmonic canonical pairs, which have zero frequency. In
particular $\Omega_{\max}=2\sqrt3/h$ for even $N$.
\end{lemma}
\begin{proof}
The first variation of a plaquette at identity is $\sum_eC_{fe}u_e^aT_a$.
For $X=x^aT_a$,
$1-\tfrac12\tr e^X=|x|^2/8+O(|x|^4)$. Since the trace has zero first
variation at identity, the quadratic plaquette term is
$|Cu|^2/(2h)$, with no contribution from the second-order noncommuting
part of its logarithm. The kinetic Taylor term is $|P|^2/(2h)$.
The declared canonical rescaling proves \eqref{v21:eq:quadratic}.

On a Fourier mode set $d_i=e^{2\pi\ii k_i/N}-1$. The exact incidence
identity is
\[
 |C(k)v|^2=|d|^2|v|^2-|d^*v|^2.
\]
The co-closed subspace is $d^*v=0$. Its dimension is two for $k\ne0$,
and the eigenvalue is $|d|^2=4\sum_i\sin^2(\pi k_i/N)$. The third
polarization is a gauge gradient and is not a physical oscillatory mode.
At $k=0$ all three spatial directions survive for each color; these are the
nine harmonic modes. Fourier diagonalization over the real field gives the
same mode count after pairing conjugate frequencies. The checkerboard mode
$k=(N/2,N/2,N/2)$ attains the displayed maximum with transverse polarization
orthogonal to $(1,1,1)$.
\end{proof}

\begin{theorem}[Exact obstruction to a global, spatially uniform zero-initial phase]
\label[theorem]{v21:thm:caustic}
For \eqref{v21:eq:Wilson}, suppose a real $C^2$ solution of
\[
 \partial_t S_h+\tfrac12|\dd S_h|_{g_h^{-1}}^2+W_h=0,
 \qquad S_h(0,U)=0
\]
is defined on $[0,T]\times SU(2)^{\mathcal E}$. Then necessarily
\begin{equation}
 \boxed{T<t_h:=\frac{\pi h}{4\sqrt3}.}
 \label{v21:eq:caustic-time}
\end{equation}
More precisely, as long as the phase is smooth, its Hessian at the identity
configuration on a transverse mode of frequency $\Omega$ is
\begin{equation}
 -\Omega\tan(\Omega t).
 \label{v21:eq:HJ-Hessian}
\end{equation}
Thus this full-manifold phase cannot provide the fixed positive time interval
needed for spatial refinement, even when the selected initial reference
configuration elsewhere on that manifold is noncommuting.
\end{theorem}
\begin{proof}
The characteristic from $(U=I,p=0)$ is stationary: $\dd W_h(I)=0$.
Differentiate the \emph{exact nonlinear} characteristic flow with respect
to its initial configuration on the graph $p=\dd S_h(0)=0$. Along this
stationary characteristic its variational equation is constant and equals
\eqref{v21:eq:quadratic}. In a transverse normal coordinate the projected
configuration derivative is $\cos(\Omega t)$ and the momentum derivative is
$-\Omega\sin(\Omega t)$. No truncation error appears in this variational
equation: the Hessian at the stationary configuration is the exact linearization
of the compact nonlinear flow.

If a global $C^2$ phase exists, its $C^1$ configuration velocity
$\nabla S_h$ has a differentiable flow with invertible derivative on each
finite time interval. Its characteristics coincide with those just computed,
by uniqueness of the smooth Hamilton equations. At
$t=\pi/(2\Omega_{\max})$ their projected derivative annihilates a nonzero
checkerboard transverse variation. This contradicts invertibility, proving
\eqref{v21:eq:caustic-time}. Before that crossing, differentiating
$p(t)=\dd S_h(t,U(t))$ gives the Hessian as momentum derivative times the
inverse configuration derivative, which proves \eqref{v21:eq:HJ-Hessian}.
The checkerboard direction is co-closed and not a gauge tangent, so gauge
redundancy does not remove the crossing.

Version 20 permits $S_0=0$ whenever its specified reference momentum is zero.
Its global phase is defined on the entire configuration manifold, not only
near the selected noncommuting orbit. The stationary identity configuration
is therefore in its required domain. Noncommutativity of a different selected
orbit does not avoid the obstruction to this \emph{global} construction.
\end{proof}

\begin{proposition}[Changing a real slope does not remove oscillator crossings]
\label[proposition]{v21:prop:real-slope}
For one oscillator of frequency $\Omega>0$ and initial real quadratic phase
$S_0(q)=s_0q^2/2$, the projected graph factor and Hessian are
\[
 A(t)=\cos(\Omega t)+(s_0/\Omega)\sin(\Omega t),\qquad
 s(t)=\frac{s_0\cos(\Omega t)-\Omega\sin(\Omega t)}{A(t)}.
\]
For every finite real $s_0$, $A(t)$ has a zero in $(0,\pi/\Omega)$.
\end{proposition}
\begin{proof}
Solve the two linear Hamilton equations from $(q_0,s_0q_0)$.
The endpoint values of $A$ at $0$ and $\pi/\Omega$ are $1$ and $-1$.
Its derivative and the numerator cannot both vanish at a zero, since the
Hamiltonian fundamental matrix is invertible. Hence the graph Hessian has
a genuine pole there.
\end{proof}

The theorem does not invalidate the fixed-graph short-time statement in
Version 20. It identifies an impossible proposed \emph{uniform upgrade} for
its permitted zero-initial global phase. Local phase charts, multiple
momentum charts, or full phase-space comparisons are different possibilities.
No general no-go for a classical Yang--Mills limit follows.

\section{A positive full phase-space observable stays regular through a caustic}
\label{v21:sec:phase-space}

Let $Z=(Q_1,\ldots,Q_d,P_1,\ldots,P_d)^{\mathsf T}$ on the Schwartz
core of $L^2(\R^d)$, with
$[Z_a,Z_b]=\ii\eta\mathbb J_{ab}I$ and
$\mathbb J=\left(\begin{smallmatrix}0&I\\-I&0\end{smallmatrix}\right)$.
These are uncompressed canonical operators, not link-angle operators on a
finite compact-group cutoff. Let
\[
 H(t)=\tfrac12 Z^{\mathsf T}K(t)Z+\ell(t)^{\mathsf T}Z,
 \qquad K(t)=K(t)^{\mathsf T},\quad A(t)=\mathbb J K(t),
\]
where quadratic expressions are symmetrically ordered, and let
$\dot z=A(t)z+\mathbb J\ell(t)$.
For a real positive matrix $G(t)$ set
\begin{equation}
 \mathcal C_G(t)=\frac12
   \sum_{a,b}G_{ab}(t)\operatorname{Sym}(Z_a-z_a,Z_b-z_b).
 \label{v21:eq:quadratic-cost}
\end{equation}

\begin{theorem}[Exact positive transport in the full canonical space]
\label[theorem]{v21:thm:quadratic}
The operator \eqref{v21:eq:quadratic-cost} is nonnegative, and
\begin{equation}
 \boxed{(\partial_t+\ii[H,\cdot]/\eta)\mathcal C_G
 =\tfrac12\operatorname{Sym}
       \bigl((Z-z)^{\mathsf T}(\dot G+A^{\mathsf T}G+GA)(Z-z)\bigr).}
 \label{v21:eq:quadratic-drift}
\end{equation}
If $\dot G+A^{\mathsf T}G+GA\preceq c(t)G$, its expectation is bounded by
$e^{\int_0^t c}\langle\mathcal C_G(0)\rangle$. If $\Phi$ is the full
Hamiltonian fundamental matrix and
$G(t)=\Phi(t)^{-\mathsf T}G(0)\Phi(t)^{-1}$, the drift is exactly zero.
This $G(t)$ stays positive through singularities of any configuration block
of $\Phi$. Uniform comparison of $G(t)$ with a physically selected norm is
a separate question; it is verified below for the Wilson quadratic modes.
\end{theorem}
\begin{proof}
Factor $G=L^{\mathsf T}L$. The cost is the sum of squares
$\frac12\sum_a(L(Z-z))_a^2$, proving positivity without an inference from a
general Weyl-positive symbol. Direct canonical commutation gives
$\ii[H,Z]/\eta=AZ+\mathbb J\ell$; subtract the reference and apply the
operator product rule. The result is exactly \eqref{v21:eq:quadratic-drift}.
A positive matrix inequality becomes a sum-of-squares quadratic-form
inequality by the same factorization, proving the expectation bound.
The fundamental matrix solves $\dot\Phi=A\Phi$ and is symplectic because
$A\mathbb J+\mathbb J A^{\mathsf T}=0$. In particular it is invertible
for every finite time where its coefficients are defined. Differentiating
$\Phi^{-\mathsf T}G(0)\Phi^{-1}$ gives zero in the drift matrix. Its
positivity follows from congruence. A configuration block may have zero
determinant while the full symplectic matrix remains invertible.
All identities first hold on the Schwartz core and extend to the corresponding
finite second-moment states by quadratic-form approximation.
\end{proof}

\begin{proposition}[Positive complex Gaussian widths have no real-graph pole]
\label[proposition]{v21:prop:complex}
Let a real symplectic matrix have blocks
$\Phi=\left(\begin{smallmatrix}A&B\\C&D\end{smallmatrix}\right)$.
For any complex symmetric $Z_0$ with $\operatorname{Im}Z_0\succ0$,
$X=A+BZ_0$ is invertible, and
\begin{equation}
 Z_t=(C+DZ_0)(A+BZ_0)^{-1},\qquad
 \operatorname{Im}Z_t=X^{-*}(\operatorname{Im}Z_0)X^{-1}\succ0.
 \label{v21:eq:complex}
\end{equation}
For an oscillator matrix $\Omega\succ0$, choosing $Z_0=\ii\Omega$ gives
$Z_t=\ii\Omega$ under its autonomous flow, even at every real-phase caustic.
\end{proposition}
\begin{proof}
Put $Y=C+DZ_0$. The symplectic identities give
$X^{\mathsf T}Y-Y^{\mathsf T}X=0$ and
$X^*Y-Y^*X=2\ii\operatorname{Im}Z_0$.
If $Xv=0$ for a nonzero $v$, pairing the latter identity with $v$ contradicts
strict positivity. Thus $X$ is invertible. Conjugate the two identities
by its inverse to prove symmetry and the imaginary-part formula.
For the oscillator, $A=D=\cos(t\Omega)$,
$B=\Omega^{-1}\sin(t\Omega)$ and $C=-\Omega\sin(t\Omega)$.
Then $X=e^{\ii t\Omega}$ and $Y=\ii\Omega e^{\ii t\Omega}$, proving
the last assertion.
\end{proof}

This positive width is an amplitude/phase description of a Gaussian
comparison, not a complex-valued replacement for the real physical action.
The directly positive operator \eqref{v21:eq:quadratic-cost} is enough for
all conclusions below; no Gaussian interpretation of the actual nonlinear
instrument is assumed.

\section{All Wilson quadratic modes: uniform energy control and a sharp quantum cost}
\label{v21:sec:all-modes}

Retain the exact Hessian \eqref{v21:eq:quadratic} on the linearized Gauss
quotient. Let $\Omega_h$ be the positive square root of $h^{-2}C^*C$ on the
mean-zero co-closed sector. Its real dimension is $6(N^3-1)$. The nine
constant harmonic pairs are retained separately, not erased as gauge modes.
In an orthonormal modal basis write
\[
 H_h^{(2)}=\tfrac12|P_{\rm r}|^2+
     \tfrac12|\Omega_h Q_{\rm r}|^2+\tfrac12|P_0|^2.
\]
Let $q_{\rm r},p_{\rm r},q_0,p_0$ solve these same classical equations. Define
\begin{equation}
 \begin{split}
 \mathcal R_h(t)={}&\tfrac12|P_{\rm r}-p_{\rm r}|^2
      +\tfrac12|\Omega_h(Q_{\rm r}-q_{\rm r})|^2\\
 &+\tfrac12|P_0-p_0|^2
      +\tfrac12|Q_0-q_0(t)-t(P_0-p_0)|^2.
 \end{split}
 \label{v21:eq:all-mode-cost}
\end{equation}

\begin{theorem}[No-caustic comparison in the physical quadratic energy]
\label[theorem]{v21:thm:all-modes}
For every $h$ the cost \eqref{v21:eq:all-mode-cost} has exactly zero drift
under $H_h^{(2)}$. On each fixed interval $[0,T]$ it controls the radiative
momentum and curl-energy error, harmonic momentum and configuration error,
and their conditional variances, with constants independent of $h$. It
also controls the mean-zero co-closed configuration error because
$\inf_k\Omega_h(k)\ge4$ for $N\ge2$.

Every normalized state satisfies the sharp lower bound
\begin{equation}
 \boxed{\langle\mathcal R_h(t)\rangle\ge
        \frac\eta2\bigl(\operatorname{tr}\Omega_h+9\bigr),\qquad
        c h^{-4}\le\operatorname{tr}\Omega_h\le C h^{-4}.}
 \label{v21:eq:quantum-floor}
\end{equation}
Displaced oscillator Gaussians, with unit-width Gaussian harmonic modes at
time zero, attain equality initially and under their exact quadratic evolution.
In particular, existence of states with vanishing expectation of this
\emph{unsubtracted positive all-mode energy comparison} is equivalent to
$\eta_h h^{-4}\to0$. No such necessity is asserted for a different,
renormalized observable or for the unproved nonlinear gauge limit.
\end{theorem}
\begin{proof}
The radiative cost has matrix $G=\operatorname{diag}(\Omega_h^2,I)$ and
$A=\left(\begin{smallmatrix}0&I\\-\Omega_h^2&0\end{smallmatrix}\right)$.
Thus $A^{\mathsf T}G+GA=0$, independently of the largest eigenvalue.
For a free harmonic pair, both $P_0-p_0$ and
$Q_0-q_0(t)-t(P_0-p_0)$ have zero total derivative. Apply
\cref{v21:thm:quadratic}, or the product rule, to get exactly zero drift
of every summand. The squared harmonic configuration error is at most
$2|Q_0-q_0(t)-t(P_0-p_0)|^2+2T^2|P_0-p_0|^2$.
For nonzero $k$, $\Omega_h(k)\ge(2/h)\sin(\pi/N)\ge4$,
using $\sin x\ge2x/\pi$ on $[0,\pi/2]$.
The identity
$\tr\rho(A-a)^2=(\tr\rho A-a)^2+\tr\rho(A-\tr\rho A)^2$
proves the statements about means and variances in these energy norms.

For one radiative oscillator of frequency $\Omega$, set
$X=Q-q$, $Y=P-p$. Since $[X,Y]=\ii\eta I$,
\[
 (Y-\ii\Omega X)^*(Y-\ii\Omega X)
       =Y^2+\Omega^2X^2-\eta\Omega I\succeq0.
\]
Its contribution is at least $\eta\Omega/2$.
For each harmonic pair replace $X$ by
$Q_0-q_0(t)-t(P_0-p_0)$; its commutator with $Y=P_0-p_0$ is still
$\ii\eta I$, giving the lower bound $\eta/2$.
These are operator lower bounds, not just uncertainty estimates on a special
family of states.

The exact real-mode trace is
\begin{equation}
 \operatorname{tr}\Omega_h
  =\frac{12}{h}\sum_{k\in(\Z/N\Z)^3\setminus\{0\}}
               \sqrt{\sum_i\sin^2(\pi k_i/N)}.
 \label{v21:eq:trace}
\end{equation}
The upper bound is immediate from $N^3=h^{-3}$ and the bound $\sqrt3$.
For the lower bound restrict $k_1$ to the middle half of its period, where
$|\sin(\pi k_1/N)|\ge1/\sqrt2$; at least a fixed fraction of the $N^3$
modes lie there, for all $N\ge4$. This proves the two-sided $h^{-4}$
estimate. More precisely the Riemann sum gives
\[
 h^4\operatorname{tr}\Omega_h\longrightarrow
 12\int_{[0,1]^3}\sqrt{\sum_i\sin^2(\pi t_i)}\,\dd t>0.
\]

At time zero prepare the product of displaced normalized Gaussians
$\exp[-\Omega_j(Q_j-q_j)^2/(2\eta)+\ii p_j(Q_j-q_j)/\eta]$
in the radiative eigenmodes, and the corresponding frequency-one Gaussians
in the harmonic modes. They have variances $\eta/(2\Omega_j)$ and
$\eta\Omega_j/2$, respectively. They saturate all the completed-square
bounds. Their exact quadratic evolution preserves the cost expectation, so
equality persists for all finite times. These statements prove both
necessity and sufficiency for the specified positive expectation to vanish.
They also supply energy-norm convergence of the means and disappearance of
all-mode fluctuations when that scale condition holds.
\end{proof}

In the original link-linearization variables, $q_e=\sqrt h\,u_e$ and
$p_e=P_e/\sqrt h$. Since $u_e$ represents $hA_e$ to first order, the two
radiative squares are precisely the mass-weighted electric and linear
magnetic energy errors. There is no number-of-sites normalization missing
from \eqref{v21:eq:quantum-floor}. At $\eta_h=h^7$ the optimal diagnostic
cost is of order $h^3$, while its real-phase caustic still occurs at order
$h$. Shrinking a quantum scale therefore suppresses fluctuations but does
not change the real classical caustic time.

\paragraph{The physical-sector distinction is retained.}
The constraint $Bp=0$ and the quotient by linear gradient coordinates are
\emph{linearized} Gauss reduction at the flat configuration. Global color
rotations act on the remaining modes; choosing nonzero colored reference
coordinates is an alignment in this quadratic model. The construction is
not a density matrix on the nonlinear compact-link singlet space, and its
quadratic cost is not claimed to descend as-is to that space. In particular,
this calculation cannot be substituted for the invariant nonlinear operator
of Version 20. It tests and repairs the high-frequency comparison mechanism,
not the missing nonlinear gauge-orbit construction.

\section{Why merely subtracting the fluctuation floor loses the source information}
\label{v21:sec:source-cost}

\begin{proposition}[Centered means do not certify a vanishing physical quadratic source]
\label[proposition]{v21:prop:source-floor}
For the centered Gaussian state in \cref{v21:thm:all-modes}, all radiative
canonical means vanish but
\[
 \langle\tfrac12|P_{\rm r}|^2\rangle
 =\langle\tfrac12|\Omega_hQ_{\rm r}|^2\rangle
 =\frac\eta4\operatorname{tr}\Omega_h.
\]
Thus vanishing coordinate means alone does not make the raw kinetic or
magnetic action insertion vanish. Subtracting
$\eta\operatorname{tr}\Omega_h/2$ makes a different energy observable.
If the independently varied action changes the frequencies, its subtracted
source differs by the derivative of that same term; it is not licensed to
leave the original source unchanged.
\end{proposition}
\begin{proof}
The two Gaussian variances computed in the preceding proof give the equal
energy shares. Subtracting a parameter-dependent scalar from a Hamiltonian
subtracts its scalar derivative from the corresponding independent first
variation. For an explicit check scale the stiffness to $s\Omega_h^2$, with
$s>0$. The ground energy is
$\eta\sqrt s\operatorname{tr}\Omega_h/2$, and its derivative is
$\eta\operatorname{tr}\Omega_h/(4\sqrt s)$.
Directly, the unsubtracted insertion is
$\tfrac12Q^{\mathsf T}\Omega_h^2Q$; its expectation in that ground state
is the same nonzero derivative. A subtraction removes exactly this term,
not only an irrelevant statement about the coordinate means.
\end{proof}

This is an elementary source-accounting test, not a prohibition on
renormalization. A declared subtraction with its complete source derivatives
can be studied separately. The present programme uses an unsubtracted positive
comparison to prove disappearance of stress defects, so its normalization
must remain explicit.

\section{Outcome: the next comparison must live in invariant phase space}
\label{v21:sec:ledger}

The shrinking-loop inverse constant is not an absolute obstruction. With
scales and weights declared, \cref{v21:thm:normalized} supplies a uniform
normalized orbit estimate, at the price of explicit raw-trace precision,
reference noncommutativity, and derivative costs. This is an observation
norm result; it does not provide all spatial connection control from
plaquettes or infer a native operator panel from its existence.

The global real-phase obstruction is stronger. For an initial phase expressly
permitted in Version 20, the actual nonlinear compact Wilson characteristics
have a transverse configuration caustic by $\pi h/(4\sqrt3)$. There can
therefore be no fixed positive interval for that global phase under refinement.
Making the accepted time step or quantum parameter smaller does not change
this classical time. Proving a uniform Hessian bound for that phase would be
a circular target, because the bound is false on the displayed family.

The positive alternative is also exact within its stated scope. The full
canonical phase-space cost passes the caustic with zero drift for the Wilson
quadratic modes, with a constant independent of the largest frequency. Its
all-mode uncertainty floor is quantified and attained, with harmonic modes
retained. These are diagnostics and rigorous design requirements, not an
already proved nonlinear Yang--Mills approximation theorem.

A substantive continuation should construct a \emph{gauge-invariant positive
cotangent-orbit comparison} on the nonlinear physical space, whose metric is
comparable to the covariant electric--curvature--connection energy and does
not require a single real configuration generating function. It must retain
ordered holonomy information, nonlinear Gauss constraints, group-curvature
commutators, and the differentiated action sources. A multiple-chart or
positive Gaussian construction would additionally have to pay its chart
transitions and spatial constants rather than restart a real phase for free.
The new calculations neither provide those missing estimates nor identify the
primitive source with a chosen instrument. Dynamical gravity, chiral matter,
and selection of physical preparations and scales remain separate.

\paragraph{Verification and preservation.}
All Version 20 bytes preceding its last end-document command are retained.
New labels use \texttt{v21:}. The diagnostic script checks scaled quaternion
invariants, weighted reconstruction, raw-trace error budgets, the exact Wilson
second variation, transverse mode frequencies, the projected-characteristic
zero, the Riccati graph and complex positive width, full symplectic energy
transport, mode counting, and the uncertainty/source scales. Finite matrices
and floating-point checks are diagnostics of these written proofs. They are
not machine-checked formalizations, nonlinear continuum simulations, or an
identification of the intended Einstein--SM source.

\begingroup
\renewcommand{\refname}{Additional references for Version 21}
\begin{thebibliography}{9}
\small\raggedright
\bibitem{v21:KS}
J.~Kogut and L.~Susskind, \emph{Hamiltonian formulation of Wilson's lattice
gauge theories}, Physical Review D \textbf{11} (1975), 395--408.
\href{https://doi.org/10.1103/PhysRevD.11.395}{doi:10.1103/PhysRevD.11.395}.
Background for the canonical compact-link action. The coefficient convention,
Hessian, and characteristic calculation above are derived explicitly.
\bibitem{v21:CR}
M.~Combescure and D.~Robert, \emph{Semiclassical spreading of quantum wave
packets and applications near unstable fixed points of the classical flow},
Asymptotic Analysis \textbf{14} (1997), 377--404.
\href{https://doi.org/10.3233/ASY-1997-14405}{doi:10.3233/ASY-1997-14405}.
Background for propagation of Gaussian packets by the full stability matrix.
No external propagation estimate replaces the finite quadratic proofs here.
\end{thebibliography}
\endgroup
\addtocontents{toc}{\protect\endgroup}
% END IMMUTABLE VERSION 21 BODY
% APPEND VERSION 22 HERE, BEFORE THE FINAL END-DOCUMENT COMMAND.


% BEGIN IMMUTABLE VERSION 22 BODY
\clearpage
\begin{center}
{\Large\bfseries Version 22: A distance form before a smooth comparison phase}\\[0.5em]
{\large Spatially uniform nonlinear vacuum-orbit control, exact tree-packet energy,
and local resolved pure-gauge records}
\end{center}
\makeatletter
\addtocontents{toc}{\protect\begingroup}
\addtocontents{toc}{\protect\makeatletter}
\addtocontents{toc}{\protect\def\protect\l@section{\protect\bfseries\protect\@dottedtocline{1}{0em}{2.5em}}}
\makeatother
\addcontentsline{toc}{part}{Version 22: Nonlinear vacuum-orbit comparison without Hamilton--Jacobi phases}

\section{Audit: a weaker observable can avoid the phase obstruction}
\label{v22:sec:audit}

The real-phase caustic in \cref{v21:thm:caustic} is not a singularity of
Hamiltonian evolution. A new fixed-graph phase-patching theorem would also
be redundant: \cref{v20:prop:patch} already supplies that conclusion, with
its graph-dependent constants. We instead ask whether a positive physical
comparison must have a globally smooth phase or a smoothly conditioned
invariant-word bank at all. On a zero-energy reference orbit the answer is
no. A bounded Lipschitz configuration cost and the \emph{unchanged nonlinear
energy} suffice. Its commutator is controlled as an energy form using only
one weak derivative.

\paragraph{New scope, and an important restriction.}
The positive result below is spatially uniform for the full compact $SU(2)$
Wilson Hamiltonian, but its classical target is the \emph{static,
trivial-holonomy vacuum gauge orbit with zero electric momentum}. All
candidate links, non-Abelian plaquette products, electric fluctuations, and
nonlinear Gauss constraints are retained. No quadratic action replaces the
Wilson action, no small-angle condition is imposed on candidates, and no
noncommuting reference anchor is required. In particular the reference orbit
may have its full diagonal $SU(2)$ stabilizer. The limiting field strengths
and the unsubtracted gauge action sources vanish. This is not a theorem
about nonzero Yang--Mills reference fields, a mass gap, a fixed-Planck-constant
quantum continuum, or the coupled Einstein--SM system.

\paragraph{Inherited results and overlap check.}
The kinetic and Wilson normalizations, the vacuum Hessian and its sharp
quadratic uncertainty cost are retained from Version 21. The normalized
$SU(2)$ force identity is \cref{v19:prop:Wilson-gradient}. Tree reduction
and Haar disintegration are already used in Version 20; only the short
measure argument needed for the new packet calculation is recalled below.
The tensor Ritz estimate, pure-gauge neutral multiplicity, local locked
outcomes, and actual clock comparison remain
\cref{v14:thm:Galerkin,v20:prop:multiplicity,v15:thm:locked,v14:thm:trajectory}.
Targeted semantic searches followed by literal checks in the cumulative
lineage, Grand Tensor, and SM--Einstein monograph did not locate the
Lipschitz vacuum-orbit estimate or the tree-packet Dirichlet identity below.
This is a limited nonduplication audit, not an exhaustive proof review.

Gauge-invariant semiclassical preparation and quantum comparison costs are
established subjects; see \cite{v20:TW2,v13:GP}. The elementary conserved-energy
argument is not claimed as a new general stability method. The additional
content is its nonsmooth but positive invariant observable on the nonlinear
physical space, the spatial normalization, the exact correlated preparation
cost, and the explicit tensor-cutoff/local-clock composition. No previously
unstated value of a primitive operator panel is inferred.

\section{The physically normalized configuration distance is an admissible form}
\label{v22:sec:distance}

Let $N\ge3$, $h=N^{-1}$, $n=N^3$, and let $\mathcal V_h,\mathcal E_h,
\mathcal F_h$ be the periodic cubic vertices, positive links, and elementary
plaquettes. Thus $|\mathcal E_h|=|\mathcal F_h|=3n$.
Set $\mathcal Q_h=SU(2)^{\mathcal E_h}$ with normalized product Haar measure
and metric
\[
 g_h=h\sum_{e,a}(\vartheta_{e,a})^2,
 \qquad T_a=-\ii\sigma_a/2.
\]
Here $X_{e,a}$ is the invariant vector field in the $T_a$ direction and
$\vartheta_{e,a}$ is its dual coframe. Left versus right conventions are
fixed consistently with the endpoint gauge action; the Casimir and every
squared norm below are independent of this choice. Write
\begin{equation}
 \begin{aligned}
 v(U)&=1-\tfrac12\operatorname{Re}\tr U,
       \qquad L_{e,a}=-\ii\eta X_{e,a},\\
 T_{h,\eta}&=-\frac{\eta^2}{2h}\sum_{e,a}X_{e,a}^2,
       \qquad W_h=\frac4h\sum_fv(U_f),\\
 E_{h,\eta}&=T_{h,\eta}+W_h.
 \end{aligned}
 \label{v22:eq:energy}
\end{equation}
The physical quantum parameter is $\eta>0$, and $U_f$ is the actual ordered
plaquette product. All expressions precede finite compression. The smooth
bounded potential is nonnegative and
\begin{equation}
 0\le W_h\le24h^{-4}I.
 \label{v22:eq:Wsup}
\end{equation}
The Friedrichs energy has the usual compact-manifold self-adjoint realization;
its form domain is $H^1(\mathcal Q_h)$ at every fixed $h,\eta$.

The vertex gauge action is $(g\cdot U)_{xy}=g_xU_{xy}g_y^{-1}$. Its vacuum
orbit is $\mathcal O_h=\{(g_xg_y^{-1})_{xy}:g\in SU(2)^{\mathcal V_h}\}$.
Define the multiplication function
\begin{equation}
 \boxed{\quad c_h(U)=h\min_g\sum_{e:x\to y}v(g_x^{-1}U_eg_y)
 =\frac h4\min_g\sum_e\|U_e-g_xg_y^{-1}\|_{\rm HS}^2.\quad}
 \label{v22:eq:cost}
\end{equation}
The minimum exists by compactness. Only the \emph{commuting configuration
variables} are minimized pointwise. No infimum of noncommuting kinetic
operators is taken, and no minimizing gauge is inserted into a physical
state. The constant comparison time unit is fixed to one; adding $c_h$ to
the comparison does not add a mass term to the action.

\begin{lemma}[Uniform weak-gradient bound without a smooth gauge slice]
\label[lemma]{v22:lem:distance}
The function $c_h$ is gauge invariant, vanishes exactly on $\mathcal O_h$,
and belongs to $W^{1,\infty}(\mathcal Q_h)$ at fixed $h$. It satisfies
\begin{equation}
 0\le c_h\le6h^{-2},\qquad
 |\nabla c_h|_{g_h}^2\le\tfrac12c_h\quad\hbox{almost everywhere}.
 \label{v22:eq:gradient}
\end{equation}
The gradient constant is independent of $h$, the number of links, and the
stabilizer of the reference orbit. Moreover,
\begin{equation}
 4c_h(U)=\inf_g h^3\sum_e
       \left\|\frac{g_x^{-1}U_eg_y-I}{h}\right\|_{\rm HS}^2.
 \label{v22:eq:connection-norm}
\end{equation}
\end{lemma}
\begin{proof}
The second equality in \eqref{v22:eq:cost} follows from
$\|A-B\|_{\rm HS}^2=4v(A^*B)$ for $A,B\in SU(2)$.
Embed $\mathcal Q_h$ in the real space of link matrices with squared norm
$h\sum\|\cdot\|_{\rm HS}^2$. Its differential has norm $1/\sqrt2$
relative to $g_h$, because $\tr(T_a^*T_b)=\delta_{ab}/2$.
Distance to any closed subset of a Euclidean space is one-Lipschitz, by
the triangle inequality. Hence $\sqrt{c_h}$ has weak gradient at most
$1/(2\sqrt2)$ in the manifold metric. This weak-gradient assertion follows
in local charts from the difference quotients of a Lipschitz function;
the almost-everywhere chain rule applied to its square gives
$|\nabla c_h|^2\le c_h/2$. Compactness gives the fixed-$h$ Lipschitz
bound for $c_h$ itself. Gauge transformations act by ambient and intrinsic
isometries and preserve the orbit, proving invariance. Since the orbit is
compact, zero distance characterizes it. Finally $v\le2$ proves the
supremum bound, and rearranging the powers of $h$ proves
\eqref{v22:eq:connection-norm}. No differentiability of a minimizing gauge
or regularity of the orbit quotient was used.
\end{proof}

\begin{proposition}[Macroscopic periods are retained, not inferred from curvature]
\label[proposition]{v22:prop:periods}
For direction $i$, let $P_{i,\ell}(U)$ be the ordered product around a
periodic coordinate line, with $N^2$ such lines. Then pointwise
\begin{equation}
 \frac1{N^2}\sum_\ell v(P_{i,\ell}(U))\le c_h(U).
 \label{v22:eq:periods}
\end{equation}
For $U_{x,1}=e^{haT_3}$ and $U_{x,2}=U_{x,3}=I$,
\begin{equation}
 W_h(U)=0,\qquad c_h(U)\ge v(e^{aT_3}).
 \label{v22:eq:flat-period}
\end{equation}
Thus a nontrivial fixed period cannot be identified with the declared
vacuum merely from zero plaquette curvature.
\end{proposition}
\begin{proof}
For unitaries, telescoping gives
$\|V_1\cdots V_N-I\|_{\rm HS}^2\le
 N\sum_k\|V_k-I\|_{\rm HS}^2$. Apply this after any vertex gauge.
The trace of a closed period is invariant. Divide by four, sum over the
$N^2$ disjoint direction-$i$ lines and divide by $N^2$. Since $N^{-1}=h$,
the result is at most $h\sum_ev(g_x^{-1}U_eg_y)$. Minimize over $g$.
The displayed constant links have identity plaquettes and the stated period,
which proves the final assertion. This controls the average of the named
period contrasts, not a supremum over individual lines.
\end{proof}

\section{A nonlinear quantum estimate with no Hamilton--Jacobi time interval}
\label{v22:sec:propagation}

Set
\begin{equation}
 \mathcal C_{h,\eta}=E_{h,\eta}+c_h\succeq0.
 \label{v22:eq:comparison}
\end{equation}
It is a closed positive form on the energy domain and commutes with every
gauge transformation. It therefore restricts to the exact physical Gauss
subspace. It is not the quantization of a classical kinetic infimum.

\begin{theorem}[Conserved energy controls a Lipschitz orbit comparison]
\label[theorem]{v22:thm:quantum}
Let $\rho(t)=e^{-\ii tE_{h,\eta}/\eta}\rho_0
 e^{\ii tE_{h,\eta}/\eta}$ have finite energy, and write
$\mathsf E_0=\tr(\rho_0 E_{h,\eta})$ and
$\mathsf c(t)=\tr(\rho(t)c_h)$. For every finite $T$, with no restriction
$T=O(h)$,
\begin{equation}
 \boxed{\quad
 \sqrt{\mathsf c(t)}\le\sqrt{\mathsf c(0)}+	frac t2\sqrt{\mathsf E_0},
 \qquad 0\le t\le T.\quad}
 \label{v22:eq:transport}
\end{equation}
In particular,
\begin{equation}
 \sup_{t\le T}\tr(\rho(t)\mathcal C_{h,\eta})
 \le (2+T^2/2)\tr(\rho_0\mathcal C_{h,\eta}).
 \label{v22:eq:quantum-bound}
\end{equation}
The constants are independent of $h,\eta$, state-space dimension and
candidate link angles. The entire nonlinear ordered Wilson potential is
retained. If the initial state is physical, it stays physical exactly.
\end{theorem}
\begin{proof}
For a normalized smooth vector, the bounded Lipschitz multiplier $c_h$
maps $H^1$ into itself. Integration by parts in the energy form gives
\begin{equation}
 \frac{\ii}{\eta}\langle\psi,[E_{h,\eta},c_h]\psi\rangle
 =\eta\operatorname{Im}\int\bar\psi
                  \langle\nabla c_h,\nabla\psi\rangle_{g_h}\,\dd\mu.
 \label{v22:eq:form-commutator}
\end{equation}
The notation on the left denotes the form commutator when $c_h\psi$ is
not in the operator domain. The potential commutes with the multiplier.
By \eqref{v22:eq:gradient}, the absolute value is at most
\[
 \left(\eta^2\int|\nabla\psi|^2\right)^{1/2}
 \left(\int|\nabla c_h|^2|\psi|^2\right)^{1/2}
 \le \sqrt{\langle T_{h,\eta}\rangle\langle c_h\rangle}.
\]
Energy conservation and $W_h\ge0$ give
$|\mathsf c'(t)|\le\sqrt{\mathsf E_0\mathsf c(t)}$.
Integrate the inequality for $\sqrt{\mathsf c+\delta}$ and let
$\delta\downarrow0$ to obtain \eqref{v22:eq:transport}. Squaring with
$(a+b)^2\le2a^2+2b^2$ gives \eqref{v22:eq:quantum-bound}.

For domain justification begin with initial vectors in $D(E_{h,\eta})$.
Their evolutions are differentiable in $L^2$ and continuous in the energy
form norm; the expectation of bounded $c_h$ is differentiable with its
form expression above. Approximate arbitrary energy-form data by energy
spectral cutoffs. Unitary evolution preserves that form approximation
uniformly on compact times, while $c_h$ is bounded at each fixed $h$.
The integrated inequality passes to the limit with the same displayed
constants. Positive sums of vector states and Cauchy--Schwarz for those
sums prove the mixed-state assertion. Gauge invariance of both terms of
$E_{h,\eta}$ gives the physical-sector preservation.
\end{proof}

The proof has no real generating function, no phase-chart restart, no
Hessian of a reconstructed orbit coordinate, and no second derivative of
$c_h$. It therefore does not encounter the caustic in Version 21. The
reference energy is exactly zero; subtracting the energy of a nonzero
reference would not produce the nonnegative operator used here.

\section{An exactly physical tree packet and its exact kinetic cost}
\label{v22:sec:packet}

Choose a rooted spanning tree $\mathsf T_h$ with root-to-vertex paths of
length at most $3(N-1)$. For example, successively use the first nonzero
coordinate to connect a vertex to a vertex one step nearer the origin,
without periodic wrap edges in the tree. There are
$m=|\mathcal E_h|-|\mathcal V_h|+1=2n+1$ chords. Let $H_j(U)$ be their
rooted fundamental holonomies. Denote the corresponding simple fundamental
cycle by $\mathsf C_j$, and put
\begin{equation}
 \ell_j=|\mathsf C_j|,\qquad
 \mathfrak L_h=\sum_{j=1}^m\ell_j\le C h^{-4},\qquad
 n_e=\#\{j:e\in\mathsf C_j\},\quad \sum_en_e=\mathfrak L_h.
 \label{v22:eq:tree-size}
\end{equation}
The common root-path segments disappear inside every central function of
$H_j$, leaving the simple cycle; thus each relevant edge occurs only once.

For $0<a\le a_0<1$ define a smooth positive central one-loop factor
\begin{equation}
 f_a(V)=Z_a^{-1/2}e^{-v(V)/(2a)},\qquad
 Z_a=\int_{SU(2)}e^{-v(V)/a}\,\dd V,
 \qquad \Phi_{h,a}(U)=\prod_{j=1}^m f_a(H_j(U)).
 \label{v22:eq:packet}
\end{equation}
This is a declared correlated preparation, not a product of original link
states. It is smooth and strictly positive everywhere; no configuration is
removed by its definition. The tree is used to prepare the state and to
estimate it, not to replace the local Hamiltonian or update supports.

\begin{lemma}[Loop-factor moments and the conditional product measure]
\label[lemma]{v22:lem:factor}
The vector $\Phi_{h,a}$ has norm one and obeys every Gauss equation exactly.
Under its configuration probability measure the rooted chord holonomies
are independent with densities $f_a^2$, conditional on the original tree
links. For each fixed positive integer $r$,
\begin{equation}
 c a^{3/2}\le Z_a\le C a^{3/2},\qquad
 \int v(V)^r f_a(V)^2\,\dd V\le C_r a^r,
 \qquad I_a:=\sum_{b=1}^3\int|X_bf_a|^2\,\dd V\le C/a.
 \label{v22:eq:one-loop}
\end{equation}
Every component of $\nabla\log f_a$ has zero mean in this measure.
\end{lemma}
\begin{proof}
For fixed tree links, replacing each chord link by $H_j$ is a left-right
Haar translation, independently for each chord. The Jacobian is one.
The remaining tree links have product Haar measure; this proves the norm
and conditional product assertions. A gauge change conjugates the rooted
holonomies by the root gauge, so centrality proves invariance of the vector,
not merely invariance of a density matrix. Differentiating gives Gauss zero.

Write $V=\cos\chi I-\ii\sin\chi\,\mathbf n\cdot\sigma$, with
$0\le\chi\le\pi$. Normalized class integration has weight
$(2/\pi)\sin^2\chi\,\dd\chi$, and $v=1-\cos\chi$.
On a fixed small interval both $v\asymp\chi^2$ and
$\sin^2\chi\asymp\chi^2$; outside it the exponential is bounded by
$e^{-c/a}$. Rescale $\chi=\sqrt a\,s$ and bound the resulting Gaussian
integrals to prove the first two estimates, including all fixed moments.
The retained exact force identity
$\sum_b|X_bv|^2=v(2-v)/4$ gives
\begin{equation}
 I_a=\frac1{16a^2}\int v(2-v)f_a^2\,\dd V.
 \label{v22:eq:Iexact}
\end{equation}
It implies the stated bound. Finally
$\int f_a^2X_b\log f_a=\frac12\int X_b(f_a^2)=0$
by Haar invariance. The same moment calculation and $|X_bv|\le\sqrt{v/2}$
bound the $p$th moment of $X_b\log f_a$ by $C_p a^{-p/2}$.
\end{proof}

\begin{theorem}[Exact tree Dirichlet identity and preparation error]
\label[theorem]{v22:thm:preparation}
The full physical kinetic energy of the original link variables is exactly
\begin{equation}
 \boxed{\quad
 \langle\Phi_{h,a},T_{h,\eta}\Phi_{h,a}\rangle
            =\frac{\eta^2}{2h}\mathfrak L_h I_a.\quad}
 \label{v22:eq:kinetic-exact}
\end{equation}
Furthermore,
\begin{equation}
 \langle\mathcal C_{h,\eta}\rangle_{\Phi_{h,a}}
 \le C\left[\frac{\eta^2\mathfrak L_h}{ha}
                            +a h^{-4}+a h^{-2}\right].
 \label{v22:eq:preparation-general}
\end{equation}
In particular, if $0<\eta\le h^7$ and $a=\eta/h$, then
\begin{equation}
 \langle\mathcal C_{h,\eta}\rangle_{\Phi_{h,\eta/h}}
                          \le C\eta h^{-5}.
 \label{v22:eq:preparation}
\end{equation}
The constants do not grow with the number of loops or links.
\end{theorem}
\begin{proof}
Let $\varepsilon_{ej}\in\{0,\pm1\}$ be the oriented simple-cycle
incidence. In the representative with all tree links equal to the identity,
differentiating a central factor at a tree edge inserts its generator to
the left or right of $H_j$, with sign $\varepsilon_{ej}$. The left and
right derivatives of a central function agree. Thus, up to one common
orthogonal change of Lie basis at that edge,
\[
 (X_{e,b}\Phi)/\Phi
       =\sum_j\varepsilon_{ej}(X_b\log f_a)(H_j).
\]
At an arbitrary tree configuration the same assertion for the summed square
follows by the fixed gauge transformation to that representative. The electric
Dirichlet density summed over $b$ is gauge invariant. Conditional on the tree
links, the different $H_j$ are independent and every displayed logarithmic
derivative is centered by \cref{v22:lem:factor}. All distinct-loop cross
terms therefore have \emph{exactly zero integral}. The remaining terms give
$n_e I_a$ for edge $e$. Sum over edges and insert $\eta^2/(2h)$ to prove
\eqref{v22:eq:kinetic-exact}. A bound by the square of the total number of
loops would lose this cancellation.

In tree gauge, a plaquette is a product of at most four chord matrices or
their inverses; tree-edge factors are identities. Unitary telescoping gives
$v(V_1\cdots V_r)\le r\sum_kv(V_k)$ for $r\le4$, also with inverse
factors. Its expectation is at most $Ca$, without requiring the factors to
be distinct. There are $3n$ plaquettes, so
$\langle W_h\rangle\le Ca h^{-4}$. The same explicit gauge is admissible
in the minimum defining $c_h$. Therefore
$c_h\le h\sum_jv(H_j)$ and $\langle c_h\rangle\le Cha m\le Ca h^{-2}$.
Combine these estimates with \eqref{v22:eq:Iexact} and
\eqref{v22:eq:tree-size}. Substituting $a=\eta/h$ bounds the kinetic term
by $C\eta h^{-4}$ and the other terms by
$C\eta h^{-5}+C\eta h^{-3}$, proving \eqref{v22:eq:preparation}.
\end{proof}

The equality is for the \emph{original electric Dirichlet form}, not a
kinetic energy chosen on the reduced chord coordinates. The independent
variables in the Haar calculation do not make the preparation a product
in the physical link factors. The optional choice
$a\asymp\eta h^{-1/2}$ balances the two leading upper bounds and gives
$C\eta h^{-9/2}$ on the displayed tree family; no optimality is asserted.
We retain $a=\eta/h$ below for a simple simultaneous scale. The sharper
quadratic lower cost $\eta h^{-4}$ in Version 21 is a distinct result;
it is not claimed to be attained by this tree packet.

\section{The higher energy moment is derived with no hidden loop-count factor}
\label{v22:sec:moments}

The tensor approximation needs an operator graph moment, not only a small
mean energy. The next bound supplies it for the same preparation.
Let $\mu_{h,a}=|\Phi_{h,a}|^2\dd\mu$, and put
\[
 \tau=\frac{\eta^2\mathfrak L_h}{ha},\qquad \omega=a h^{-4}.
\]
These are preparation scales, not a time step or an action generator.

\begin{lemma}[Fixed derivative words of the correlated preparation]
\label[lemma]{v22:lem:words}
For each fixed word $I=((e_1,b_1),\ldots,(e_r,b_r))$ of invariant
link derivatives, and each fixed $1\le p<\infty$, one has
\begin{equation}
 \left\|\frac{X_I\Phi_{h,a}}{\Phi_{h,a}}\right\|_{L^p(\mu_{h,a})}
 \le C_{r,p}a^{-r/2}\prod_{k=1}^r\sqrt{n_{e_k}}.
 \label{v22:eq:word-bound}
\end{equation}
A word containing an edge with $n_e=0$ is zero. Constants depend on the
fixed derivative order, not on graph size.
\end{lemma}
\begin{proof}
For one derivative use the centered independent sum in the preceding proof.
For an even integer $2s\ge p$, expand its $2s$th moment. Any term with an
index appearing only once has zero expectation. The remaining terms have
at most $s$ distinct indices, so their number is at most $C_s n_e^s$.
The single-loop moment bound in \cref{v22:lem:factor} bounds each product
by $C_s a^{-s}$, using H\"older. Passing down to $p$ proves the first-order
bound. Orthogonal basis changes conditional on the tree do not change these
bounds or their centering.

For an ordered derivative of $\log\Phi$ of order $k\ge2$, each contributing
central Wilson word must contain \emph{every} differentiated edge.
Each edge occurs once in its simple cycle, and insertion of $k$ fixed Lie
generators has absolute trace bounded by a constant depending only on $k$,
not on word length. Repeated derivatives at one edge have the same bound.
Consequently its pointwise magnitude is at most $C_k n_I/a$, where $n_I$
counts such common cycles. Since
$n_I\le\min_j n_{e_j}\le\prod_j\sqrt{n_{e_j}}$ for $k\ge2$,
and $a^{-1}\le a^{-k/2}$, this is the required bound for that log derivative.
The statement is evaluated before any coordinate-dependent gauge
transformation is differentiated; it is simply differentiation of each
original closed word. Finally, repeated use of the product rule expresses
$X_I\Phi/\Phi$ as a finite sum of products of ordered derivatives of
$\log\Phi$, with total order $r$. Apply H\"older to its first-order factors
and the pointwise estimates to all other factors. The number of terms depends
only on $r$. This proves \eqref{v22:eq:word-bound}.
\end{proof}

\begin{theorem}[Uniform preparation graph moment]
\label[theorem]{v22:thm:moment}
For $0<\eta\le h^7$, $a=\eta/h$, and all sufficiently small $h$,
\begin{equation}
 \boxed{\quad \|(E_{h,\eta}+I)^2\Phi_{h,a}\|\le C.\quad}
 \label{v22:eq:graph-moment}
\end{equation}
No negative power of the many-loop normalization occurs in this estimate.
\end{theorem}
\begin{proof}
The word lemma with orders two and four, followed by the triangle inequality
in the electric sums, gives
\[
 \|T\Phi\|\le C\tau,\qquad
 \|(T\Phi)/\Phi\|_{L^4(\mu_{h,a})}\le C\tau,\qquad
 \|T^2\Phi\|\le C\tau^2.
\]
For example the last sum is bounded by
$C\eta^4h^{-2}a^{-2}(\sum_en_e)^2$; the three Lie directions per
link only multiply its fixed constant. Plaquette telescoping and the
single-loop moment bound give
$\|W_h\|_{L^p(\mu_{h,a})}\le C_p\omega$ for every fixed $p$.
Thus $\|W_hT\Phi\|\le C\omega\tau$ and
$\|W_h^2\Phi\|\le C\omega^2$.

Each link occurs in four elementary plaquettes. The exact group force bound
and the fixed second derivative of a normalized matrix trace imply
\[
 \|X_{e,b}W_h\|_{L^4(\mu_{h,a})}\le C\sqrt a/h,
 \qquad \|X_{e,b}^2W_h\|_\infty\le C/h.
\]
For the first bound use $|Xv(U_f)|^2\le Cv(U_f)$ and
$\|v(U_f)\|_{L^2}\le Ca$. Expanding $T(W_h\Phi)$ by its actual product
rule, its difference from $W_hT\Phi$ has norm at most
\[
 \frac{C\eta^2}{h}\sum_e
 \left(h^{-1}+h^{-1}\sqrt{n_e}\right)
 \le C\eta^2h^{-6}.
\]
Here the zero-incidence edges contribute only the first term, and
$\sum_{n_e>0}\sqrt{n_e}\le\sum_en_e=\mathfrak L_h\le Ch^{-4}$.
It follows that
\[
 \|E^2\Phi\|\le C(\tau^2+\tau\omega+\omega^2+\eta^2h^{-6}),
 \qquad \|E\Phi\|\le C(\tau+\omega).
\]
For the chosen width, $\tau\le C\eta h^{-4}$ and
$\omega=\eta h^{-5}$; all displayed quantities are bounded when
$\eta\le h^7$. This proves \eqref{v22:eq:graph-moment}.
The vector is smooth on the compact group product at every fixed cutoff,
so every differentiation lies in the required operator domain. The estimates
supply the graph moment rather than assuming it.
\end{proof}

\section{Spatially explicit tensor cutoff and local instrument costs}
\label{v22:sec:finite}

Let $P_J$ be the product of the link Peter--Weyl cutoffs $j\le J$, $J\ge2$,
and $E_J=P_JE_{h,\eta}P_J$ on its finite tensor range. The source preparation
is the normalized $\Phi_J=P_J\Phi/\|P_J\Phi\|$, not a projection applied
repeatedly to observed outcomes. These projections commute with every
vertex gauge transformation and with the kinetic Casimirs. They need not
commute with the Wilson potential or with $c_h$.
Put
\begin{equation}
 b_J=1+\eta^2(J+1)^2,\qquad
 \Lambda_J=1+\eta^2(J+1)^2/h.
 \label{v22:eq:Jscales}
\end{equation}

\begin{proposition}[No fixed-graph constant is hidden in the cutoff step]
\label[proposition]{v22:prop:cutoff}
For $0<h\le1$, $0<\eta\le1$, with $B=E_{h,\eta}+I$,
\begin{equation}
 \|(I-P_J)u\|_B^2\le\delta_{h,\eta,J}\|Bu\|^2,
 \qquad \delta_{h,\eta,J}=C h^{-12}/\Lambda_J.
 \label{v22:eq:tail}
\end{equation}
For the preparation of \cref{v22:thm:moment}, and sufficiently small
$\delta_{h,\eta,J}$, the finite and uncompressed unitary evolutions obey
\begin{equation}
 \sup_{t\le T}\|\psi_J(t)-\psi(t)\|_B^2
 \le C_T\frac{h^{-12}}{\Lambda_J}(1+T/\eta)^2.
 \label{v22:eq:Galerkin}
\end{equation}
Consequently,
\begin{equation}
 \sup_{t\le T}\langle\psi_J(t),\mathcal C_{h,\eta}\psi_J(t)\rangle
 \le C_T\left[\eta h^{-5}
       +\frac{h^{-14}}{\Lambda_J}(1+T/\eta)^2\right].
 \label{v22:eq:finite-quantum}
\end{equation}
\end{proposition}
\begin{proof}
Set $A=I+T_{h,\eta}$. By \eqref{v22:eq:Wsup},
$A\preceq B\preceq Ch^{-4}A$ and
$\|Au\|\le Ch^{-4}\|Bu\|$. These statements concern both the form and
operator domains: a bounded multiplication perturbation preserves $D(T)$.
On $\Ran(I-P_J)$ one link Casimir has spin greater than $J$ and therefore
$A\ge c\Lambda_J$ there. Its joint spectral decomposition with $P_J$ gives
\[
 \|A^{1/2}(I-P_J)u\|^2\le C\Lambda_J^{-1}\|Au\|^2.
\]
The two comparisons prove \eqref{v22:eq:tail}. Insert this tail constant,
quantum scale $\eta$, and the \emph{proved} uniform graph moment into the
Ritz/Duhamel estimate \cref{v14:thm:Galerkin}. This gives
\eqref{v22:eq:Galerkin}, with its normalization threshold.

The new comparison has $0\preceq\mathcal C_{h,\eta}\preceq6h^{-2}B$
by \eqref{v22:eq:gradient}. The squared triangle inequality for its positive
square root, \cref{v22:thm:quantum,v22:thm:preparation}, and
\eqref{v22:eq:Galerkin} prove \eqref{v22:eq:finite-quantum}. The additional
$h^{-2}$ in this last estimate is intentional: its omission would incorrectly
claim a uniform upper norm for the global configuration cost.
\end{proof}

Decompose $E_J/\eta=\sum_{\alpha=1}^{M}H_\alpha$ into electric-link and
Wilson-plaquette terms, $M=6n$. An electric context can be enlarged to an
incident plaquette without altering its Hamiltonian term. The five neutral
character vectors of \cref{v20:prop:multiplicity} supply the local multiplicity
for the exact gauge-commuting six-outcome construction. Use that existing
construction with local parameter $M\epsilon$. The complete accepted label
is $(\alpha,e)$, with the context chosen uniformly. This scheduler and the
preparation remain declared, not derived from the primitive opportunity bit.
The physical timing law remains $d=4\vartheta\epsilon/11$.

\begin{proposition}[Dimension-independent local norm budget]
\label[proposition]{v22:prop:local}
For $J\ge2$, the explicit costs satisfy
\begin{equation}
 \begin{split}
 \|P_J\mathcal C_{h,\eta}P_J\|&\le C h^{-4}b_J=:A_C,\\
 \mathfrak C_{h,\eta,J}
 :=1+\Big(\sum_\alpha\|H_\alpha\|\Big)^2
       +M\sum_\alpha\|H_\alpha\|^2+M^2
       &\le C h^{-8}\eta^{-2}b_J^2.
 \end{split}
 \label{v22:eq:local-costs}
\end{equation}
Every actual local outcome preserves Gauss's law exactly.
\end{proposition}
\begin{proof}
An electric-link generator has norm at most $C\eta(J+1)^2/h$; a plaquette
generator has norm at most $C/(\eta h)$. There are $3n$ of each.
Summing and squaring gives the second bound, including the $M^2$ term for
$h,\eta\le1$. Their energy norms, together with $c_h\le6h^{-2}$, give
the first. No local or global matrix dimension enters these bounds.
The outcome statement is precisely the inherited commutant construction,
not a new inference from covariance of its averaged channel. Source products
are formed before compression throughout.
\end{proof}

\section{A spatially uniform nonlinear local vacuum branch on actual records}
\label{v22:sec:resolved}

Let $\rho_j$ be the actual conditional state at tick $j$, initialized by
$|\Phi_J\rangle\langle\Phi_J|$. Define
\begin{equation}
 \begin{split}
 e_j&=\tr(\rho_jP_J\mathcal C_{h,\eta}P_J),\\
 \mathfrak b_h&=\eta h^{-5}
 +\frac{h^{-14}}{1+\eta^2(J+1)^2/h}(1+T/\eta)^2
 +\epsilon h^{-12}\eta^{-2}\{1+\eta^2(J+1)^2\}^{3}.
 \end{split}
 \label{v22:eq:budget}
\end{equation}

\begin{theorem}[Full Wilson action, exact Gauss support, and a uniform time interval]
\label[theorem]{v22:thm:resolved}
For $0<\eta\le h^7$, $a=\eta/h$, $J\ge2$, the preceding tensor
normalization threshold, and $6h^{-3}\epsilon\le\tau_0$, one has, on any
fixed finite physical interval and with any initial clock phase,
\begin{equation}
 \max_{j\le T/d}\E e_j\le C_T\mathfrak b_h,\qquad
 \PP\{\max_{j\le T/d}e_j>u\}\le C_T\mathfrak b_h/u\quad(u>0).
 \label{v22:eq:resolved}
\end{equation}
The constants are independent of spatial and matrix dimensions,
$\eta,J,\epsilon$, and $0<\vartheta\le1$. Every attained state has exact
Gauss support. If $0<\mathfrak b_h\le1/2$, then
\begin{equation}
 \E\max_{j\le T/d}e_j\le C_T\mathfrak b_h
       \left[1+\log\left(2+\frac{h^{-4}b_J}{\mathfrak b_h}\right)\right].
 \label{v22:eq:log}
\end{equation}
The target is the zero-electric, trivial-holonomy vacuum orbit. These
estimates do not assert convergence to a nonzero classical gauge field.
\end{theorem}
\begin{proof}
The actual local locked-clock theorem \cref{v14:thm:trajectory} applies with
the explicit costs of \cref{v22:prop:local}. Its pure finite-unitary reference
$\psi_J(t)$ gives a nonnegative fidelity error $\xi_j$ whose expectation and
maximal tail are bounded by $C_T\epsilon\mathfrak C_{h,\eta,J}$.
No new proof of the clock correction or of its maximal inequality is needed.
For $V_J=P_J\mathcal C_{h,\eta}P_J$ and the rank-one projection $P_*$ onto
$\psi_J(jd)$, positivity gives
\[
 V_J\preceq2P_*V_JP_*+2(I-P_*)V_J(I-P_*).
\]
Pairing with the actual conditional state bounds $e_j$ by twice the
\emph{deterministic} bound in \eqref{v22:eq:finite-quantum} plus
$2A_C\xi_j$. Substitution of \eqref{v22:eq:local-costs} proves the mean
bound. For the maximal probability, thresholds below twice the deterministic
term use the bound one; all other thresholds force a corresponding crossing
of $\xi_j$. This gives \eqref{v22:eq:resolved} without a union bound over
ticks. All preparation and outcome operators preserve physical support.
Finally $0\le e_j\le A_C$. Integrate the tail
$\min(1,C_T\mathfrak b_h/u)$ over this bounded interval to obtain
\eqref{v22:eq:log}. The source never applies the global comparison operator,
a phase correction, or a reference-state projection as an update.
\end{proof}

\begin{corollary}[One explicit simultaneous choice]
\label[corollary]{v22:cor:scales}
Take
\begin{equation}
 \boxed{\quad \eta_h=h^7,\quad a_h=h^6,\quad
 J_h=\lceil h^{-22}\rceil,\quad \epsilon_h=h^{118},\quad
 d_h=\frac{4\vartheta_h}{11}h^{118}.\quad}
 \label{v22:eq:scales}
\end{equation}
Then $\mathfrak b_h=O_T(h^2)$ and
$\E\max_j e_j=O_T(h^2(1+|\log h|))$ on every fixed finite interval.
\end{corollary}
\begin{proof}
Here $b_J=O(h^{-30})$ and $\Lambda_J\asymp h^{-31}$.
The three terms in \eqref{v22:eq:budget} have orders $h^2,h^3,h^2$,
respectively. The tensor tail constant is $O(h^{19})$, so its uniform
preparation-moment threshold holds. Also $M\epsilon=6h^{115}\to0$ and
$J\ge2$. The comparison norm is $O(h^{-34})$, hence the logarithm has
order $1+|\log h|$. These are sufficient, deliberately conservative powers;
no computational efficiency or sharp semiclassical threshold is claimed.
\end{proof}

\section{Connection alignment, physical sources, and the exact remaining boundary}
\label{v22:sec:outcome}

On a physical state, products inside expectations always refer to original
operators followed by the fixed tensor projection. The configuration cost
controls a Born-weight tail:
\begin{equation}
 \tr\bigl(\rho_jP_J\one_{\{c_h>r^2\}}P_J\bigr)\le e_j/r^2,
 \qquad r>0.
 \label{v22:eq:Born}
\end{equation}
Together with \eqref{v22:eq:connection-norm}, this is a physically normalized
connection-orbit comparison. It is not a simultaneous classical trajectory
of noncommuting configurations and electric momenta. By
\cref{v22:prop:periods}, each averaged macroscopic period contrast is also
bounded by $e_j$. The positive electric and magnetic energies obey
\begin{equation}
 \begin{split}
 \frac{h^3}{2}\sum_{e,a}\tr\rho_j(L_{e,a}/h^2)^2&\le e_j,\\
 h^3\sum_f\tr\rho_j
       \left\|(U_f-I)/h^2\right\|_{\rm HS}^2&\le e_j .
 \end{split}
 \label{v22:eq:fields}
\end{equation}
The equalities between these expressions and the action energies follow
from \eqref{v22:eq:energy} and $\|U-I\|_{\rm HS}^2=4v(U)$.
This retains every group-valued plaquette, not only a quadratic expansion
near the identity. The vanishing expressions are invariant squared field
norms and their conditional quantum contributions, not merely small means.

For independently varied relative electric and Wilson weights, define
\[
 E(\alpha,\beta)=\sum_e\alpha_e\sum_a L_{e,a}^2/(2h)
                         +(4/h)\sum_f\beta_f v(U_f).
\]
At $\alpha_e=\beta_f=1$, a direction with
$|\dot\alpha_e|,|\dot\beta_f|\le B_0$ has its exact insertion $O$ bounded
as forms by
\begin{equation}
 -B_0 E_{h,\eta}\preceq O\preceq B_0 E_{h,\eta},\qquad
 |\tr(\rho_jP_JOP_J)|\le B_0e_j.
 \label{v22:eq:sources}
\end{equation}
These follow termwise from positivity before compression. Thus the
unsubtracted action insertions tend to their vacuum value zero, with the
same maximal bound $C_T\mathfrak b_h/u$ for a fixed source threshold.
This covers the displayed independently defined coefficient sources. A full
lattice metric or shift variation must first be specified; it is not
inferred uniquely from the flat Wilson energy. No vacuum-energy subtraction
or omitted derivative of such a subtraction is used.

\paragraph{What is genuinely new.}
The spatially uniform estimate uses the full nonlinear compact action and
exact Gauss constraints at a singular vacuum orbit. It requires neither
linearized reduction, noncommuting anchors, a real-phase interval, nor
smoothing of attained records. The higher preparation moment and the local
finite-source costs are derived. The finite stochastic state evolves through
all its local outcomes; its macroscopic invariant errors converge to the
\emph{static vacuum}, not to a nonzero classical field.

\paragraph{What remains independent.}
For nonzero reference fields, $E-E_{\rm ref}$ need not be positive and
conserved energy alone does not control trajectory separation. The moving
covariant force comparison therefore remains unproved; fixed-graph phase
restarts do not supply its spatial constants. The preparation is correlated
along a possibly macroscopic spanning tree. Locality is proved for the
subsequent Kraus supports, not for its preparation circuit or a relativistic
speed bound. The tree, initial packet, quantum scale, contexts and their
selection law remain declared, not obtained from the independently specified
primitive Renewal source. Chiral matter, Einstein constraints and backreaction
are absent. No generic microscopic attraction or Yang--Mills mass gap follows.

\paragraph{Verification and preservation.}
All Version 21 bytes before its final document command are retained. New
labels use \texttt{v22:}. The diagnostics test tree-incidence derivatives,
Haar moments, Dirichlet normalization, gauge invariance, the form commutator,
period information and scale powers. These are finite checks, not
machine-checked proofs or an Einstein--SM simulation.
\addtocontents{toc}{\protect\endgroup}
% END IMMUTABLE VERSION 22 BODY
% APPEND VERSION 23 HERE, BEFORE THE FINAL END-DOCUMENT COMMAND.


% BEGIN IMMUTABLE VERSION 23 BODY
\clearpage
\begin{center}
{\Large\bfseries Version 23: Actual source reduction before another constructed limit}\\[.5em]
{\large Record-sensitive covariance, decoded positive-energy transport, and the kinetic price of conditional centering}
\end{center}
\makeatletter
\addtocontents{toc}{\protect\begingroup}
\addtocontents{toc}{\protect\makeatletter}
\addtocontents{toc}{\protect\def\protect\l@section{\protect\bfseries\protect\@dottedtocline{1}{0em}{2.5em}}}
\makeatother
\addcontentsline{toc}{part}{Version 23: Same-source decoded transport and a quantitative full-record obstruction}

\section{Four-source audit and a change in the immediate target}
\label{v23:sec:audit}

Version 22 proves a nonlinear, spatially uniform vacuum theorem for a
constructed local instrument. Its actual-source identification remains
independent. The four newly supplied companions sharpen that distinction.
Their latest theorem blocks and the earlier blocks needed by those statements
were read, together with targeted overlap checks against this lineage. This is
not an independent verification of every proof in their cumulative histories.
The following are inherited results, not new claims of this continuation.

\begin{center}\small
\begin{tabularx}{\textwidth}{@{}L{.21\textwidth}X@{}}
\toprule
Companion & Usable result and restriction retained here\\
\midrule
YM mass-gap V43 \cite{v23:YM43} &
\nolinkurl{thm:v43-triangular} and \nolinkurl{thm:v43-free-history} identify
unit-Jacobian conditional innovations and a uniform free Hessian at the zero
history. \nolinkurl{prop:v43-norm-price}, \nolinkurl{prop:v43-actual-vertex},
and \nolinkurl{prop:v43-wall} retain, respectively, the inverse-coordinate
cost, a nonzero actual Wilson cubic, and a transported-wall contribution.
This is not nonlinear all-depth convexity or a physical-time gap.\\[.4em]
SM source selection V17 \cite{v23:SEL17} &
\nolinkurl{sel12:thm:coarse-descent} proves an averaged multiplicity
reduction on all inputs. \nolinkurl{sel12:thm:record-descent} requires more
to reproduce original labelled histories. \nolinkurl{sel17:prop:old-defect}
evaluates a nonzero fixed-label covariance defect of the unchanged full-record
witness. The internal invariant algebra is not a symmetry of that entire
instrument.\\[.4em]
Exact-action Einstein V16 \cite{v23:EB16} &
\nolinkurl{v16:lift-theorem} lifts a specified bare soft compression, but
\nolinkurl{v16:actual-domain} retains boundedness of the actual lapse as a
premise. The full observed inverse, including complementary flux cancellation,
remains unresolved. Small matter stress is not a substitute for that inverse.\\[.4em]
YM shell mixing V16 \cite{v23:SH16} &
\nolinkurl{thm:v16-flat} and \nolinkurl{thm:v16-limit} identify the complete
measured one-loop coefficient and its differentiated refinement limit on a
flat one-holonomy branch. The statement is about that coefficient, not a
nonperturbative joint measure, a preparation, or Lorentzian trajectories.\\
\bottomrule
\end{tabularx}
\end{center}

The most immediate native issue is therefore not another freely chosen
Hamiltonian realization. It is the relation between the supplied source and
the proposed physical reader. We prove two different conclusions. First, the
already evaluated fixed-label covariance defect gives a strictly positive
operational separation from every covariant full-record model; rarity does
not remove it at the first accepted observation. Second, \emph{full-record
autonomous descent is not needed for positive-energy concentration of a
retained decoded reader}. A completely positive state map and an averaged
intertwining estimate on the full conditional input suffice. No source outcome
is erased, and no intermediate reset is performed.

The final calculation transports the kinetic term under the YM conditional
coordinate change. It explains why the companion's repaired potential Hessian
cannot by itself replace a physical-clock estimate. The shell measure and
Einstein inverse results above are used as boundaries on composition, not
claimed to discharge missing hypotheses.

\paragraph{Attribution.}
Choi operators, completely bounded trace norms, and their elementary
composition inequalities are established tools; see \cite{v23:Watrous}.
The stopped positive-observable argument is inherited from Versions 12--14.
The additions are the quantitative application to the supplied source witness,
the physical-clock and decoded-history estimates, and their explicit
Version 22 costs. The evaluated source-selection witness is a specified
normalized source in that companion, not a verified unique primitive
Einstein--SM law.

\section{A quantitative separation of the evaluated full-record source}
\label{v23:sec:covariance}

Let $\mathcal I=(\mathcal I_e)_e$ be a normalized instrument on $M_r$ and
retain its original labels through the channel
\[
 \mathbf I(X)=\sum_e|e\rangle\langle e|\otimes\mathcal I_e(X).
\]
The Choi convention is
$J(\mathcal L)=\sum_{a,b}\mathcal L(E_{ab})\otimes E_{ab}$, so a channel
has Choi trace $r$. Output tensor reorderings do not change the norms below.
We use the diamond norm with no factor $1/2$; two perfectly distinguishable
channels have distance two. A self-adjoint $T$ acts on input and output by
$U_t=e^{\ii tT}$, and
$\widehat T=T\otimes I-I\otimes\overline T$ acts on a branch Choi space.
Labels transform trivially. Write
\[
 q_T=\operatorname{diam}(\operatorname{spec}\widehat T)
      =2\operatorname{diam}(\operatorname{spec}T).
\]

\begin{theorem}[Covariance residual bounds the distance to every covariant instrument]
\label[theorem]{v23:thm:distance}
Suppose $q_T>0$. Let $J_e=J(\mathcal I_e)$ and
$\mathfrak d_T=\sum_e\|[\widehat T,J_e]\|_{\rm HS}^2$.
For every normalized instrument $\mathcal J$ on the same carrier and original
label space, whose individual branches are $U_t$-covariant for all real $t$,
\begin{equation}
 \boxed{\quad
 \|\mathbf I-\mathbf J\|_\diamond
       \ge \frac{\sqrt{\mathfrak d_T}}{r q_T}.
 \quad}
 \label{v23:eq:distance}
\end{equation}
It suffices on the right to retain any subset of the nonnegative branch
summands. The comparison branches may have arbitrary Choi rank.
\end{theorem}
\begin{proof}
Let $\mathcal P_T$ be the HS orthogonal projection onto the commutant of
$\widehat T$. In its spectral decomposition,
\[
 \|[\widehat T,J_e]\|_{\rm HS}^2
  =\sum_{\lambda,\mu}(\lambda-\mu)^2
                 \|P_\lambda J_eP_\mu\|_{\rm HS}^2
  \le q_T^2\|(I-\mathcal P_T)J_e\|_{\rm HS}^2.
\]
A covariant branch Choi operator belongs to $\Ran\mathcal P_T$. Orthogonal
projection and summation over the original flags therefore give
$\|J(\mathbf I)-J(\mathbf J)\|_{\rm HS}
 \ge\sqrt{\mathfrak d_T}/q_T$.
Apply the channel difference to half of the normalized maximally entangled
state of input dimension $r$. Its output is the Choi difference divided by
$r$. Hence the diamond norm is at least its trace norm divided by $r$,
which is at least its HS norm divided by $r$. This proves the assertion.
Nothing in this argument selects a Kraus representation of the comparison.
\end{proof}

\begin{corollary}[A numerical separation certified by the supplied exact source calculation]
\label[corollary]{v23:cor:actual-distance}
For the unchanged source-selection V13/V15 witness at
$(u,v,\zeta)=(3/5,4/5,\ii)$, use the actual generator $T$ which is
$\operatorname{diag}(1,-1,0)$ on its three-dimensional trivial block and
zero on the other blocks. Then every fixed-label instrument covariant under
this internal colour subgroup has
\begin{equation}
 \boxed{\quad
 \|\mathbf I-\mathbf J\|_\diamond
 \ge c_{\rm src}:=\frac{\sqrt{87688}}{12600}
 >\frac1{43}.
 \quad}
 \label{v23:eq:actual-distance}
\end{equation}
This also excludes approximation by a full-record instrument covariant under
the proposed entire colour group, in this same represented identification.
\end{corollary}
\begin{proof}
The inherited exact computation
\nolinkurl{sel17:prop:old-defect} gives, for one original branch $K$,
\[
 \|[\widehat T,J_K]\|_{\rm HS}^2=87688/50625.
\]
Here $r=14$ and $q_T=4$. Apply \cref{v23:thm:distance} to that single
branch. Since $\sqrt{50625}=225$, its denominator is $14\cdot4\cdot225$.
The strict rational comparison follows from
$87688\cdot43^2-12600^2=3375112>0$.
The branch coefficient and its commutator are inherited, not regenerated
from unprovided auxiliary scripts in this continuation.
\end{proof}

This is a separation on the identified fourteen-dimensional full source.
It is not a comparison between that carrier and Version 22's growing tensor
carrier without a map between them. It also does not say that forgetting the
labels, restricting inputs, or applying a noninvertible physical decoder
preserves the lower bound. Those are different experiments.

For noisy operator data, the same proof gives the following one-sided test.
If $\widehat J$ has direct-sum HS error at most $e$ from $J$, then
\[
 \inf_{\mathbf J\ \mathrm{covariant}}\|\mathbf I-\mathbf J\|_\diamond
 \ge \frac{(\|[\widehat T,\widehat J]\|_{\rm HS}-q_Te)_+}{rq_T}.
\]
Indeed the commutator map has HS norm at most $q_T$ by the same spectral
calculation. A small observed residual is not a certificate of exact symmetry.

\section{Rarity does not turn a full-record mismatch into a physical identification}
\label{v23:sec:clock}

Keep the inherited phase matrix, accepted probability, and physical units:
\[
 M=\begin{pmatrix}1/5&4/5\\2/3&1/3\end{pmatrix},\qquad
 d=\frac{4\vartheta}{11}\epsilon,\qquad0<\vartheta\le1.
\]
The mean physical waiting time from a fresh completed boundary to acceptance
is $\epsilon$. Until that first acceptance the compared field is unchanged.
The clock and the retained waiting record are the same in the two experiments.

\begin{proposition}[Exact opportunity and first-acceptance normalization]
\label[proposition]{v23:prop:clock}
Put $\delta_{\rm rec}=\|\mathbf I-\mathbf J\|_\diamond$. At one opportunity,
with identical no-response branch, the flagged channel distance is
$\vartheta\delta_{\rm rec}$. In the $P$ phase of one protected tick it is
$(2\vartheta/3)\delta_{\rm rec}$, so its norm divided by the tick duration is
\begin{equation}
 \frac{11}{6\epsilon}\delta_{\rm rec}.
 \label{v23:eq:rate-distance}
\end{equation}
The complete experiment retaining the first accepted label, field output,
and waiting time has distance exactly $\delta_{\rm rec}$. If stopped at a
physical deadline $T>0$, recording ``no acceptance'' as a separate outcome,
its distance is
\begin{equation}
 F_{\vartheta,d}(T)\delta_{\rm rec}
       \ge(1-\epsilon/T)_+\delta_{\rm rec},
 \label{v23:eq:first}
\end{equation}
where $F_{\vartheta,d}(T)$ is the actual clock's first-acceptance probability.
\end{proposition}
\begin{proof}
In the first two maps the only nonzero difference is the accepted flagged
block with its displayed scalar probability. Homogeneity of the diamond norm
gives the first statements and the calibration gives \eqref{v23:eq:rate-distance}.
At first acceptance the output map is the tensor product of the unchanged
waiting-time distribution with $\mathbf I$, or with $\mathbf J$.
Tensoring a normalized classical state preserves the norm: preparation gives
the upper bound and tracing that state gives the lower bound. At finite
deadline the difference has total weight $F_{\vartheta,d}(T)$ and the failed
block is identical. The mean waiting identity and Markov's inequality give
$1-F_{\vartheta,d}(T)\le\epsilon/T$. Infinite waiting records are handled by
finite truncations and their vanishing tail.
\end{proof}

Thus \cref{v23:cor:actual-distance} has an order-one first-accepted-record
separation even if $\vartheta$ tends to zero, when $\epsilon\to0$.
Equation \eqref{v23:eq:rate-distance} is a norm of \emph{instrument-map}
increments, not a lower bound for accumulated field drift. Version 10 shows
why a large fast-phase drift alone would not imply such a path obstruction.
The first-acceptance assertion is a distinct, fully recorded operational test.

\begin{proposition}[A sufficient all-history error bound pays accepted count, not raw ticks]
\label[proposition]{v23:prop:hybrid}
Suppose two instruments on a common state/record carrier differ by at most
$\delta$ in flagged diamond norm at each accepted event, uniformly on every
possible classical past. They have the same independent clock and the same
initial state. Their full finite-history output channels through $J$ ticks
obey
\begin{equation}
 \|\mathsf H_J-\widetilde{\mathsf H}_J\|_\diamond
       \le\min\{2,\delta\,\E A_J\},
 \qquad \E A_J\le Jd/\epsilon+\vartheta,
 \label{v23:eq:hybrid}
\end{equation}
where $A_J$ is the accepted-event count. No independence between attained
quantum states is required.
\end{proposition}
\begin{proof}
Condition on the entire common clock and acceptance sequence, which is
independent of the field. For a sequence with $a$ acceptances, replace the
accepted instruments one at a time. Each preceding and following normalized
quantum/classical channel is a trace-norm contraction at every amplification;
thus telescoping costs at most $a\delta$. A classical feedback rule is a
direct sum over past flags, so the uniform conditional norm bound gives the
same estimate. Average the clock and use convexity. The bound two is the
universal channel-distance bound.
For $p_0=\PP(s_0=P)$ and $\lambda_0=-7/15$, the exact phase recurrence gives
\[
 \E A_J=\frac{4\vartheta}{11}J+
 \frac{5\vartheta}{11}(p_0-6/11)(1-\lambda_0^J).
\]
The correction has absolute value at most $60\vartheta/121<\vartheta$.
This proves the stated bound. This is an upper bound; no linear-in-$J$
lower accumulation claim is made.
\end{proof}

\section{A retained physical reader need not be an autonomous labelled process}
\label{v23:sec:reduction}

Let $\mathsf R_*$ be a CPTP map from an actual finite source algebra
$\mathfrak A$ to a physical matrix algebra $\mathfrak B$; finite direct sums
with their usual sum trace are allowed. Let $\sigma_j$ be the source's actual
conditional state and $\rho_j=\mathsf R_*(\sigma_j)$. An actual accepted
instrument $\{\mathcal A_{y,*}\}$ has sum $\mathcal A_*$.
For every physical observable $O$ the exact conditional identity is
\begin{equation}
 \E[\tr(\rho_{j+1}O)\mid\sigma_j,\text{ acceptance}]
       =\tr\bigl(\sigma_j\mathcal A^*\mathsf R^*O\bigr).
 \label{v23:eq:conditional}
\end{equation}
Normalization by an outcome probability cancels against that probability
in this identity. There is no division by a small lower bound for a recorded
outcome. The reduced state need not determine its next conditional law.

\begin{proposition}[An exact locked witness separates averaged reduction from record autonomy]
\label[proposition]{v23:prop:nonautonomous}
Let the source carrier be $\C^6\otimes\mathcal K$, with $\mathcal K$ any
finite physical space, and let $Q_e$ be the six coordinate projections on
$\C^6$. Put
\[
 b=\frac2{\sqrt5},\qquad a=\frac{\sqrt2-b}{6}>0,
 \qquad K_e=(aI+bQ_e)\otimes I_{\mathcal K}.
\]
These six original coefficients are independent, have positive probabilities
on every nonzero state, and satisfy
\[
 \sum_eK_e^*K_e=I,\qquad \sum_eK_e=\sqrt2 I.
\]
For $\mathsf R_* =\operatorname{Tr}_{\C^6}$, their summed channel obeys
$\mathsf R_*\mathcal A_* =\mathsf R_*$ on all inputs, including ancillary
extensions. Nevertheless no collection of reduced branch maps reproduces
$\mathsf R_*\mathcal A_{e,*}=\mathcal B_{e,*}\mathsf R_*$ for all inputs.
For every initial state with pure physical marginal, the physical marginal
remains that same pure state on every outcome.
\end{proposition}
\begin{proof}
The sum is $(6a+b)I=\sqrt2 I$. The squared sum has coefficient
$6a^2+2ab+b^2=(2+5b^2)/6=1$. Independence follows because the coefficient
matrix of the six diagonal operators is $bI+a\one\one^{\mathsf T}$, with
nonzero eigenvalues $b$ and $\sqrt2$. Their least singular value is $a$.
A trace-preserving operation on the traced factor leaves its other marginal
unchanged, proving the averaged identity.
The inputs $|1\rangle\langle1|\otimes\tau$ and
$|2\rangle\langle2|\otimes\tau$ have identical reduction but give, at label
one, the different unnormalized reductions $(a+b)^2\tau$ and $a^2\tau$.
This excludes branch autonomy. Finally a positive joint state with pure
marginal $|v\rangle\langle v|$ is supported on $\C^6\otimes\C v$:
pair it with the positive complementary projection to get zero, then use
its positive square root. Every $K_e$ preserves that support, proving the
last statement. This is a diagnostic identity-action reader, not the
Einstein--SM source.
\end{proof}

The source-selection companion's actual external-symmetry reduction gives a
more general useful entrance. Its maps satisfy
$\mathsf R_*\mathsf S_* =\mathrm{id}$ and, when its stated covariance
hypothesis is verified,
\[
 \mathsf R_*\mathcal A_* =\Psi_*\mathsf R_*.
\]
For an \emph{independently fixed} target action channel $\mathscr E_*$, put
\begin{equation}
 \delta_h=\|\mathsf R_*\mathcal A_*-
                      \mathscr E_*\mathsf R_*\|_\diamond.
 \label{v23:eq:intertwining}
\end{equation}
On that exact averaged reduction,
\begin{equation}
 \delta_h=\|\Psi_*-\mathscr E_*\|_\diamond.
 \label{v23:eq:small-algebra}
\end{equation}
Indeed composition with $\mathsf R_*$ gives the upper bound, and composition
with its CPTP right inverse $\mathsf S_*$ gives the reverse bound. This is an
audit on the smaller physical algebra. It neither sets $\mathscr E_*:=\Psi_*$
by fitting an action nor executes the resetting map $\mathsf S_*\mathsf R_*$
between source events. Failed labelled descent in
\nolinkurl{sel12:thm:record-descent} is entirely compatible with this identity.

\section{A same-source decoded-energy theorem on the actual renewal clock}
\label{v23:sec:transfer}

At each cutoff let the target physical carrier, Hamiltonian
$H=\sum_\alpha H_\alpha$, and constructed channel $\mathscr E_*$ be those of
\cref{v14:prop:channel}. Put
\[
 L=\sum_\alpha\|H_\alpha\|,\qquad
 \mathfrak C=1+L^2+M\sum_\alpha\|H_\alpha\|^2+M^2.
\]
The target is a comparison channel, not a replacement of the actual source.
Use $M\epsilon\le\tau_0$ and the actual clock
$d=4\vartheta\epsilon/11$. Assume $dL\le1$.

The actual source may have a different and larger finite memory. It has a
fixed CPTP reader $\mathsf R_*$ and its original accepted instrument. Allow
that instrument to depend on its complete prior classical history, but require
\eqref{v23:eq:intertwining} with a common bound $\delta_h$ at every such
history. No source-to-physical branch autonomy is assumed. On nonaccepted
ticks the decoded state is unchanged; an invisible native update is allowed
only when this \emph{conditional} reader identity is verified. The displayed
clock probabilities and acceptance rate are independent of that source state.
These are explicit source hypotheses, not consequences of the word renewal.

For a physical unit vector $\phi$, set
\[
 \psi_*(t)=e^{-\ii tH}\phi,\quad P_*(t)=|\psi_*(t)\rangle\langle\psi_*(t)|,
 \quad \xi_j=1-\tr(\rho_jP_*(jd)),\quad
 \xi_0=1-\tr(\mathsf R_*\sigma_0\,|\phi\rangle\langle\phi|).
\]
A random initial source state is permitted; its expectation is used below.

\begin{theorem}[Approximate averaged intertwining controls the actual decoded histories]
\label[theorem]{v23:thm:decoded}
Under the preceding hypotheses put
\begin{equation}
 B_h=\E\xi_0+\epsilon\mathfrak C+\delta_h/\epsilon.
 \label{v23:eq:B}
\end{equation}
For every fixed $T$, independently of source and physical matrix dimensions,
the number of observed ticks, and $0<\vartheta\le1$,
\begin{equation}
 \max_{j\le T/d}\E\xi_j\le C_T B_h,
 \qquad \PP\{\max_{j\le T/d}\xi_j>u\}\le C_TB_h/u.
 \label{v23:eq:decoded}
\end{equation}
All states in this conclusion are decoded from the \emph{unchanged original
source histories}. If $H$ preserves a physical Gauss subspace with projector
$\mathsf P$ and $\phi\in\Ran\mathsf P$, then the same bounds hold for
$\tr(\rho_j(I-\mathsf P))$. Exact Gauss support requires the separate exact
branch condition; it is not inferred from a small $B_h$.
\end{theorem}
\begin{proof}
Use the positive physical reference observable $V_*=I-P_*$ and
$C_*=\ii[H,V_*]$. Its exact drift is
$\dot V_*+\ii[H,V_*]=0$. The inherited phase coefficients are
$\chi_H=-15/22$, $\chi_P=25/44$, with $(I-M)\chi=f-1$ and
$f_H=0,f_P=11/6$. As in \cref{v14:thm:trajectory}, put
\[
 \widehat V_s(t)=V_*(t)+d\chi_s C_*(t)+dK_0I,
 \qquad K_0=\|\chi\|_\infty\sup_t\|C_*(t)\|\le C L.
\]
Then $0\preceq V_*\preceq\widehat V_s\preceq V_*+2dK_0I$.
The same-clock target calculation gives, for every physical input state,
the conditional increment at most $Cd\epsilon\mathfrak C$.
This is an operator inequality and not an assertion about a distribution of
target outcomes.

Pair instead with the actual source state using \eqref{v23:eq:conditional}.
At a current phase $s$ the only difference from that target calculation is
\[
 p_s\tr\left[\sigma_j
 (\mathcal A^*\mathsf R^*-\mathsf R^*\mathscr E^*)
                         \widehat V_H(t+d)\right],
 \qquad p_H=0,\quad p_P=2\vartheta/3.
\]
The dual of an induced trace-norm bound controls its pairing by
$p_s\delta_h\|\widehat V_H(t+d)\|$. The diamond bound is stronger than
that induced bound. Since $dL\le1$, the last norm is bounded by an absolute
constant. Also $p_s\le11d/(6\epsilon)$. Thus the added increment is at most
$Cd\delta_h/\epsilon$. This argument remains conditional on the entire
history; a non-Markov decoded state introduces no extra term.

Set $W_j=\tr(\rho_j\widehat V_{s_j}(jd))\ge\xi_j$. We have proved
\[
 \E[W_{j+1}\mid\mathcal F_j]\le W_j+Cd\beta_h,
 \quad \beta_h=\epsilon\mathfrak C+\delta_h/\epsilon,
 \quad \E W_0\le\E\xi_0+CdL\le C B_h.
\]
Up to $J=\lfloor T/d\rfloor$, the process
$W_j+C\beta_h(T-jd+d)$ is nonnegative and a supermartingale after enlarging
the fixed constant. Stopping at its first $W_j>u$ gives the maximal bound,
and iteration gives the expectation bound. This is the same elementary
maximal principle already proved in Version 12, now on the \emph{native}
conditional states through a CP reader. Finally
$I-\mathsf P\preceq I-P_*(t)$ whenever $\psi_*(t)$ is physical, proving
the leakage statement.
\end{proof}

\begin{corollary}[Positive physical action observables transfer without record autonomy]
\label[corollary]{v23:cor:energy}
Let $0\preceq V_h(t)\preceq A_hI$ on the physical carrier and suppose
$\sup_{t\le T}\langle\psi_*(t),V_h(t)\psi_*(t)\rangle\le a_h$.
For the actual decoded $e_j=\tr(\rho_jV_h(jd))$, put
$b_h=a_h+A_hB_h$. Then
\begin{equation}
 \max_j\E e_j\le C_Tb_h,\qquad
 \PP(\max_j e_j>u)\le C_Tb_h/u.
 \label{v23:eq:energy}
\end{equation}
For $b_h>0$ the maximal expectation is at most
$C_Tb_h[1+\log(2+A_h/b_h)]$. Independently specified insertions with
$-cV_h\preceq O_h\preceq cV_h$ have
$|\tr(\rho_jO_h)|\le c e_j$ on those same records.
\end{corollary}
\begin{proof}
For $P=P_*(t)$ and $Q=I-P$, the squared triangle inequality gives
$V_h\preceq2PV_hP+2QV_hQ\preceq2a_hP+2A_hQ$.
Therefore $e_j\le2a_h+2A_h\xi_j$. Apply \cref{v23:thm:decoded}, using the
trivial bound one for thresholds below $4a_h$. Since $0\le e_j\le A_h$,
integration of $\min(1,C_Tb_h/u)$ proves the logarithmic bound. Pairing the
two stated insertion inequalities with the positive state proves the last
claim. These are bounds for actual differentiated insertions, not a definition
of an action by its desired limit.
\end{proof}

\paragraph{What the theorem does and does not bypass.}
The required defect is a single-step averaged map on \emph{every full
conditional source input}, not agreement of ensemble means for one initial
population. It does not reproduce the source's branch probabilities from its
physical marginal. It keeps those original probabilities and states and
bounds their physical positive errors directly. A source may therefore fail
\nolinkurl{sel12:thm:record-descent} and still satisfy this theorem. Conversely,
knowing only a fitted averaged trajectory or a source Gram gives no bound on
\eqref{v23:eq:intertwining}.

\section{An explicit cost for transferring the nonlinear vacuum branch}
\label{v23:sec:scales}

Specialize the preceding theorem to Version 22's finite Wilson system.
There $V_h=P_J\mathcal C_{h,\eta}P_J$, and the proved deterministic
comparison and norm bounds are
\[
 \begin{gathered}
 a_h\le C_T\left[\eta h^{-5}+\frac{h^{-14}}{\Lambda_J}(1+T/\eta)^2\right],
 \qquad A_h\le Ch^{-4}b_J,\\
 b_J=1+\eta^2(J+1)^2,\qquad
 \Lambda_J=1+\eta^2(J+1)^2/h.
 \end{gathered}
\]
The source reader, its preparation, and its actual averaged intertwining with
\emph{that same independently fixed action channel} still have to be acquired.

\begin{corollary}[A quantitative native-to-constructed handoff, with its explicit precision]
\label[corollary]{v23:cor:vacuum}
Under the actual-source hypotheses of \cref{v23:thm:decoded}, the entire
Version 22 positive vacuum and coefficient-source bound holds with budget
\begin{equation}
 \widetilde{\mathfrak b}_h
   =C_T\left[\mathfrak b_h^{(22)}
        +h^{-4}b_J\left(\E\xi_0+\frac{\delta_h}{\epsilon}\right)\right].
 \label{v23:eq:vacuum}
\end{equation}
For its declared scales $\eta=h^7$, $J=\lceil h^{-22}\rceil$,
$\epsilon=h^{118}$, it suffices to have
\begin{equation}
 \E\xi_0=O(h^{36}),\qquad \delta_h=O(h^{154})
 \label{v23:eq:precision}
\end{equation}
to retain $\widetilde{\mathfrak b}_h=O_T(h^2)$ and the
$O_T(h^2(1+|\log h|))$ maximal mean bound. Exact physical support is
retained only if the decoded original branches preserve it exactly;
otherwise the theorem supplies the explicit small leakage bound.
\end{corollary}
\begin{proof}
Insert \cref{v22:prop:cutoff,v22:prop:local} into
\cref{v23:cor:energy}. Under the displayed powers,
$A_h=O(h^{-34})$ and $\mathfrak C=O(h^{-82})$. Hence
$A_h\epsilon\mathfrak C=O(h^2)$,
$A_h\E\xi_0=O(h^2)$ and
$A_h\delta_h/\epsilon=O(h^2)$.
Also $L\le C h^{-41}$, so $dL\to0$ uniformly in $\vartheta$ and the
extra hypothesis $dL\le1$ holds. Version 22 already verifies its remaining
local and tensor thresholds. Its positive electric, Wilson, orbit-distance,
and specified coefficient-source inequalities hold on the finite carrier
before restricting to physical states, so they transfer through this same
positive observable. The Gauss qualification is
\cref{v23:thm:decoded}, not an added repair operation.
\end{proof}

These powers are sufficient bounds for an intentionally conservative
full-state-fidelity route. They are \emph{not} necessary precision thresholds
for every physical reader. A direct positive-observable intertwining estimate
may cost much less. A flagged full-record distance is a sufficient upper bound
for a coarse-channel defect on the same carrier, but the fixed positive lower
bound in \cref{v23:cor:actual-distance} is only for the flagged experiment.
It cannot be used to rule out a smaller averaged defect after a verified CP
reduction. Neither an arbitrarily small acceptance probability nor an unnamed
coarse-graining establishes such a reduction.

\section{The YM conditional coordinates also transport the kinetic term}
\label{v23:sec:kinetic}

The mass-gap companion repairs the free \emph{potential} Hessian on a
history fibre. It does not identify that fibre with a physical canonical
state. The following exact calculation states what is required before one
could make that additional identification.

\begin{proposition}[Unit-Jacobian centering is not a change of physical frequencies]
\label[proposition]{v23:prop:kinetic}
On one fixed final-source fibre suppose the inherited linear conditional
centering is $\xi=U_m\eta$, with
\[
 \det U_m=1,\qquad H_m^{\rm raw}=U_m^*D_mU_m,
 \qquad D_m=\operatorname{diag}(H_1,\ldots,H_m).
\]
If a canonical kinetic matrix $M\succ0$ is \emph{independently specified} in
raw coordinates, the cotangent transformation is
$p_\xi=U_m^{-\mathsf T}p_\eta$. It takes the quadratic Hamiltonian to
\begin{equation}
 \frac12p_\xi^{\mathsf T}(U_mM^{-1}U_m^{\mathsf T})p_\xi
                      +\frac12\xi^{\mathsf T}D_m\xi.
 \label{v23:eq:kinetic}
\end{equation}
Its squared frequencies are exactly those of $M^{-1}H_m^{\rm raw}$.
In particular a depth-independent lower bound for $D_m$ alone is not a
lower bound for these frequencies. For the \emph{diagnostic} choice $M=I$
and the actual raw test vector supplied in YM V42--V43,
\begin{equation}
 \omega_{\min,m}^{\,2}\le\frac{16}{7(m-3)}\qquad(m\ge4).
 \label{v23:eq:frequency}
\end{equation}
No physical YM mass-gap inference from that diagnostic kinetic choice is made.
\end{proposition}
\begin{proof}
Equality of the canonical one-form gives
$p_\eta^{\mathsf T}d\eta=p_\xi^{\mathsf T}d\xi$ and hence
$p_\eta=U_m^{\mathsf T}p_\xi$. Substitution proves
\eqref{v23:eq:kinetic}. The second-order configuration equation has matrix
$U_mM^{-1}U_m^{\mathsf T}D_m
 =U_m(M^{-1}H_m^{\rm raw})U_m^{-1}$,
which proves equality of spectra, including multiplicities. The inherited
actual vector $z_m$ has
$z_m^*H_m^{\rm raw}z_m=N^4$ and
$\|z_m\|^2\ge(7/16)(m-3)N^4$; its Rayleigh quotient proves
\eqref{v23:eq:frequency}. The potential floor and this test concern different
metrics, so there is no contradiction.
\end{proof}

There is an exact nonlinear version at the operator level. Let
$\xi=F(\eta)$ be the companion's smooth unit-Jacobian centering on its
actual open branch, and initially work on compactly supported smooth tests.
For a raw kinetic energy $-\eta_q^2\operatorname{div}_\eta
(M^{-1}\nabla_\eta)/2$, with $M$ constant and $\eta_q$ a separately specified
quantum scale, the unitary coordinate map
$(\mathcal U\psi)(\xi)=\psi(F^{-1}\xi)$ gives
\begin{equation}
 \mathcal U T\mathcal U^{-1}
   =-\frac{\eta_q^2}{2}\operatorname{div}_\xi
       \left(K_F(\xi)\nabla_\xi\right),\qquad
 K_F=DF\,M^{-1}DF^{\mathsf T}\big|_{F^{-1}\xi}.
 \label{v23:eq:nonlinear-kinetic}
\end{equation}
Indeed the change of variables in the Dirichlet form has Jacobian one and
transforms each gradient by $DF^{\mathsf T}$; integration by parts proves
this differential expression and its induced form realization. The potential,
Haar/pivot densities, and boundaries are transported with the same map.
The divergence in \eqref{v23:eq:nonlinear-kinetic} includes its variable
coefficient first-order terms. Deleting them or replacing the transported
domain by a new rectangular one changes the operator. The supplied source
has not specified this diagnostic canonical kinetic energy, so this result
prevents an unjustified transfer; it does not assign one to the YM history.

\section{Outcome: the next acquisition is a physical map with its action test}
\label{v23:sec:outcome}

There are now two source questions, and the new notes show why they should
not be conflated. Reproducing the complete efficient source as a covariant
full-record process is obstructed in the evaluated witness by
\eqref{v23:eq:actual-distance}. Proving positive-energy concentration of a
physical reader requires only the weaker but still all-input identity
\eqref{v23:eq:intertwining}, with its physical-time error normalization.
The latter may survive even when the source's labels do not descend
as an autonomous process. The theorem keeps those labels and the full
conditional memory instead of resetting or silently discarding either.

The concrete next bank is therefore
\[
 \boxed{\quad
 \mathsf R_*,\quad
 \mathsf R_*\mathcal A_*-\mathscr E_*\mathsf R_*,\quad
 \E\xi_0,\quad
 \text{decoded branch support and physical clock}.
 \quad}
\]
The original source must supply these maps in its actual representation.
The target $\mathscr E_*$ must be fixed by the same action before testing.
When the source-selection averaged multiplicity reduction is verified, the
nontrivial action test reduces exactly to \eqref{v23:eq:small-algebra}; its
labelled nonautonomy is not a reason to abandon that route. When only internal
algebra dimensions or low-order probability marginals are known, the bank
has not been acquired. The companion's five-total-slot determinant-mode
criterion remains a separate structural test and is not replaced by the
present dynamical theorem.

The other two upstream boundaries remain equally real. The shell file's
complete signed determinant coefficient is not a bound for the actual
nonequilibrium source law, and YM V43's free innovation floor does not supply
that law or its physical kinetic metric. The Einstein bridge's lifted soft
compression is not the full observed inverse after complementary variables
relax. Accordingly no Einstein response estimate, nonzero-reference
non-Abelian continuum theorem, mass gap, or universal microscopic selection
is claimed here. The change of direction is to verify a potentially smaller
\emph{actual source-to-physical channel} rather than keep manufacturing new
locally supported comparison dynamics.

\paragraph{Verification and preservation.}
All Version 22 bytes before its final document command are preserved.
The accompanying script checks the covariance projection inequality, its
rational witness constant, exact locked nonautonomous reduction, clock-count
normalization, the decoded conditional drift comparison, symplectic spectrum
invariance, and the explicit powers. Companion rational coefficients are
identified as inherited inputs. These diagnostics are not a regeneration of
all companion certificates, a machine-checked proof, or an Einstein--SM
simulation. New labels carry the prefix \texttt{v23:}.

\begingroup
\renewcommand{\refname}{Additional sources for Version 23}
\begin{thebibliography}{9}
\small\raggedright
\bibitem{v23:YM43} A.~P\'elissier,
\emph{Renewal Yang--Mills Mass-Gap Research Notes}, supplied cumulative V43,
file \nolinkurl{Renewal_Yang_Mills_Mass_Gap_Research_Notes_V43(1).tex}.
\bibitem{v23:SEL17} A.~P\'elissier,
\emph{Renewal SM Source-Selection Research Notes}, supplied cumulative V17,
file \nolinkurl{Renewal_SM_Source_Selection_Research_Notes_V17(3).tex}.
\bibitem{v23:EB16} A.~P\'elissier,
\emph{Renewal Exact-Action Einstein Bridge Research Notes}, supplied cumulative V16,
file \nolinkurl{Renewal_Exact_Action_Einstein_Bridge_Research_Notes_V16(3).tex}.
\bibitem{v23:SH16} A.~P\'elissier,
\emph{Renewal Yang--Mills Shell-Mixing Research Notes}, supplied cumulative V16,
file \nolinkurl{Renewal_Yang_Mills_Shell_Mixing_Research_Notes_V16(2).tex}.
\bibitem{v23:Watrous} J.~Watrous,
\emph{The Theory of Quantum Information}, Cambridge University Press, 2018,
Chapter~3. Author's source and errata:
\href{https://cs.uwaterloo.ca/~watrous/TQI/}{\nolinkurl{cs.uwaterloo.ca/~watrous/TQI/}}.
Used for background on channel norms; the specialized inequalities above
are proved directly.
\end{thebibliography}
\endgroup
\addtocontents{toc}{\protect\endgroup}
% END IMMUTABLE VERSION 23 BODY
% APPEND VERSION 24 HERE, BEFORE THE FINAL END-DOCUMENT COMMAND.


% BEGIN IMMUTABLE VERSION 24 BODY
\clearpage
\begin{center}
{\Large\bfseries Version 24: Test physical heating before full-channel proximity}\\[.5em]
{\large One-sided action-observable transport, a strict locked example, and the actual reduced Pauli source on the renewal clock}
\end{center}
\makeatletter
\addtocontents{toc}{\protect\begingroup}
\addtocontents{toc}{\protect\makeatletter}
\addtocontents{toc}{\protect\def\protect\l@section{\protect\bfseries\protect\@dottedtocline{1}{0em}{2.5em}}}
\makeatother
\addcontentsline{toc}{part}{Version 24: Direct action-energy audits and a reduced-source heating test}

\section{Audit: a smaller test must still be a test of the actual source}
\label{v24:sec:audit}

Version 23 transfers physical energy through a pure-reference fidelity estimate.
Its sufficient error contains the full physical-observable norm multiplied by a
source channel defect divided by the mean accepted interval. It explicitly leaves
open whether direct positive-observable control can be cheaper. We develop that
alternative, and then apply a separate necessary test to an already specified
reduced source. Merely renaming a channel defect as an energy error would not
establish either conclusion.

\paragraph{Inherited material and exact nonduplication boundary.}
We retain the completely positive reader and original-history identity in
\cref{v23:eq:conditional}, the phase correction of Version 10, and the positive
maximal principle of Versions 12--14. The source-selection companion
\cite{v23:SEL17}, at \nolinkurl{sel12:thm:coarse-descent}, supplies an exact averaged
multiplicity reduction under its stated covariance hypothesis. Its
\nolinkurl{sel8:eq:reduced-channel} specifies
\[
 \Psi_*(\rho)=\tfrac13\rho+\tfrac12D\rho D^*+\tfrac16F\rho F^*,
\]
and \nolinkurl{sel8:thm:quadratic-blindness} supplies the two normalized Pauli
panels $(D,F)=(Z,Z)$ and $(X,Z)$. Their original coefficients, covariance,
independence, and static internal-algebra calculation are not new results here.
This Pauli benchmark is \emph{not} the fourteen-dimensional V13 witness used in
\cref{v23:cor:actual-distance}. No coefficients of that other witness are changed
or silently assigned to this calculation.

The new assertions are: an all-input, one-sided energy transfer on the unchanged
conditional source histories; an exact finite matrix certificate and cold-energy
obstruction; a genuinely moving locked example with order-one full-channel
mismatch; and the physical-time heating, contraction, and recovery costs of the
supplied reduced Pauli panel. Searches for order-sensitive intertwining and reduced
trace contraction in the cumulative lineage and the source-selection companion
located no duplicate of this composition. This is a targeted audit, not exhaustive
verification of those files or a claim of priority for general Lyapunov methods.

\paragraph{Attribution and interpretation.}
Positive-observable supermartingales, trace-norm contraction, and protected quantum
subspaces are established tools; see \cite{v23:Watrous,v24:TV,v24:RKW}. The proofs
below establish their particular action normalization and renewal-clock use.
An operator inequality must be verified on the actual source inputs; it is not
established by naming a slack variable. Conversely, an algebraically full internal
sector need not carry conservative physical-time dynamics. The exact Pauli
calculation below concerns one supplied benchmark, not every admitted Renewal
source. No Einstein inverse, nonzero-reference non-Abelian continuum estimate,
or primitive selection of a physical action is inferred here.

\section{A one-sided energy transfer with no pure-state or full-channel norm cost}
\label{v24:sec:order}

All algebras in this body are finite matrix algebras or finite direct sums, with
the ordinary trace. Let $\mathsf R_*$ be a fixed CPTP map from the actual source
algebra to the physical algebra. Denote its unital positive adjoint by $\mathsf i$;
it need not be multiplicative. The actual conditional states are $\sigma_j$ and
$\rho_j=\mathsf R_*(\sigma_j)$. Every original outcome is retained. The accepted
channel $\mathcal A_{j,*}$ may depend on the entire prior classical history.
On nonaccepted ticks the reader values used below are unchanged conditionally;
hidden updates require precisely that identity, not just equality in equilibrium.

Retain the actual clock and its independently calibrated physical tick:
\begin{equation}
 M=\begin{pmatrix}1/5&4/5\\2/3&1/3\end{pmatrix},\qquad
 d=\frac{4\vartheta}{11}\epsilon,\qquad
 p_H=0,\quad p_P=2\vartheta/3,\qquad
 p_P\epsilon=\frac{11}{6}d.
 \label{v24:eq:clock}
\end{equation}
The mean accepted interval is $\epsilon$. The clock is independent of the source
state, as in the specified branch of Version 23.

Let $C(t)\succeq0$ be a \emph{separately fixed physical action comparison}. It
must control the desired field or source topology by an independently proved
estimate. Suppose there are physical corrected observables $\widehat C_s(t)$
with
\begin{equation}
 C(t)\preceq\widehat C_s(t)\preceq C(t)+\zeta I,\qquad \zeta\ge0,
 \label{v24:eq:comparison-window}
\end{equation}
and a separately fixed target accepted channel $\mathscr E_*$ whose same-clock
one-step operator calculation gives
\begin{equation}
 \mathcal T_s^*\widehat C(t+d)-\widehat C_s(t)
 \preceq d\lambda_0\widehat C_s(t)+d r_0I.
 \label{v24:eq:target-order}
\end{equation}
Here $\lambda_0,r_0\ge0$, and the precise meaning of $\mathcal T_s^*$ is
\begin{align}
 \mathcal T_H^*\widehat C(t+d)
   &=\tfrac15\widehat C_H(t+d)+\tfrac45\widehat C_P(t+d),\nonumber\\
 \mathcal T_P^*\widehat C(t+d)
   &=\tfrac13\widehat C_P(t+d)+\tfrac23(1-\vartheta)\widehat C_H(t+d)
                        +\tfrac23\vartheta\mathscr E^*\widehat C_H(t+d).
 \label{v24:eq:target-tick}
\end{align}
This is an operator estimate, not a claim about reproducing target labels.
For the usual phase correction one takes
\[
 \widehat C_s=C+d\chi_s\mathscr DC+dK_0I,\qquad
 \chi_H=-15/22,\quad\chi_P=25/44,
\]
with $K_0\ge\max_s|\chi_s|\sup_t\|\mathscr DC(t)\|$ and
$\zeta=2dK_0$. The hypotheses above also allow another already proved positive
correction. In particular they do not assert that Version 22's statewise tensor
comparison automatically supplies \eqref{v24:eq:target-order}.

Define the actual Heisenberg intertwining error
\begin{equation}
 \Xi_j(O)=\mathcal A_j^*\mathsf i(O)-\mathsf i(\mathscr E^*O).
 \label{v24:eq:Xi}
\end{equation}
Only the following \emph{one-sided} test is required, on every possible
conditional source input and at every relevant tick:
\begin{equation}
 \boxed{\quad
 \Xi_j\bigl(\widehat C_H(t+d)\bigr)
 \preceq\epsilon a\,\mathsf i(\widehat C_P(t))+
                      \epsilon r I,\qquad a,r\ge0.
 \quad}
 \label{v24:eq:source-order}
\end{equation}
Negative energy work needs no small lower bound. Errors on unrelated observables
or on the full channel may be large. The number $r$ has physical energy-per-time
units; $a$ is a relative growth rate.

\begin{theorem}[Direct action-observable transport on the original source histories]
\label[theorem]{v24:thm:order}
Assume \eqref{v24:eq:comparison-window}--\eqref{v24:eq:source-order}. Put
\begin{equation}
 \lambda=\lambda_0+\tfrac{11}{6}a,\qquad
 R_0=r_0+\tfrac{11}{6}r,\qquad
 b=\E\tr(\rho_0C(0))+\zeta+TR_0.
 \label{v24:eq:budget}
\end{equation}
For $J=\lfloor T/d\rfloor$ and $e_j=\tr(\rho_jC(jd))$,
\begin{equation}
 \boxed{\quad
 \max_{j\le J}\E e_j\le e^{\lambda T}b,\qquad
 \PP\{\max_{j\le J}e_j>u\}\le e^{\lambda T}b/u\quad(u>0).
 \quad}
 \label{v24:eq:maximal}
\end{equation}
No pure physical preparation, labelled autonomy, or diamond-norm proximity is
assumed. In particular the budget has no factor $\sup_t\|C(t)\|$ multiplying an
initial fidelity or a channel error. If $C(t)\preceq A I$ and $b>0$, then
\begin{equation}
 \E\max_{j\le J}e_j\le C_{\lambda,T}b
                    \bigl[1+\log(2+A/b)\bigr].
 \label{v24:eq:log}
\end{equation}
The factor $A$ enters only this final integration of the tail bound. All
constants displayed in the hypotheses must remain controlled in any claimed
cutoff-uniform application.
\end{theorem}
\begin{proof}
Set $W_j=\tr(\sigma_j\mathsf i(\widehat C_{s_j}(jd)))\ge e_j$. Conditional
expectation of an original accepted outcome gives
$\tr(\sigma_j\mathcal A_j^*\mathsf i(\widehat C_H(t+d)))$; its normalization
probability cancels exactly as in \cref{v23:eq:conditional}. Lift
\eqref{v24:eq:target-order} by the positive unital map $\mathsf i$.
The only additional term is $p_s\Xi_j(\widehat C_H(t+d))$. It is zero in phase
$H$. In phase $P$, \eqref{v24:eq:source-order} and \eqref{v24:eq:clock} bound it
by $d(11/6)(a\mathsf i(\widehat C_P(t))+rI)$. Therefore
\[
 \E[W_{j+1}\mid\mathcal F_j]\le(1+\lambda d)W_j+dR_0,
 \qquad\E W_0\le\E e_0+\zeta.
\]
This is conditional on the entire actual history; no Markov property of the
physical marginal is used.

Write $q=1+\lambda d$. The process
\[
 S_j=q^{-j}W_j+R_0\sum_{\ell=j}^{J-1}d\,q^{-(\ell+1)}
\]
is nonnegative and a supermartingale. Iteration proves the mean estimate.
Stopping at the first $W_j>u$, or at $J$ when no crossing occurs, gives
$\PP(\max W_j>u)\le q^J\E S_0/u\le e^{\lambda T}b/u$. Since $e_j\le W_j$,
this proves \eqref{v24:eq:maximal}. Integrate
$\min(1,e^{\lambda T}b/u)$ up to $A$ to obtain \eqref{v24:eq:log}.
The argument never multiplies a full-state distance by $A$.
\end{proof}

\begin{corollary}[Stationary action-energy test without any phase correction]
\label[corollary]{v24:cor:static}
For a fixed $C\succeq0$, suppose the actual accepted source satisfies
\[
 \mathcal A_j^*\mathsf i(C)-\mathsf i(C)
       \preceq\epsilon a\mathsf i(C)+\epsilon r I
\]
on all conditional inputs. Then \eqref{v24:eq:maximal} holds with
$\lambda=11a/6$, $b=\E\tr(\rho_0C)+11Tr/6$. This proves concentration or
stability in that action energy, not by itself the equations of a nonstationary
reference. The comparison must encode those equations separately when needed.
\end{corollary}
\begin{proof}
On nonaccepted ticks the tested value is unchanged. The exact accepted
probability gives the same scalar conditional inequality with no correction
and no target channel. Apply the preceding supermartingale calculation.
\end{proof}

\section{A finite spectral certificate and the unavoidable cost of heating cold states}
\label{v24:sec:certificate}

\begin{proposition}[Exact one-sided certificate and its robustness]
\label[proposition]{v24:prop:certificate}
For fixed time and history, let
\[
 D_j=\epsilon^{-1}\Xi_j(\widehat C_H(t+d)),\qquad
 B_j=\mathsf i(\widehat C_P(t))\succeq0.
\]
At any fixed relative rate $a\ge0$, the smallest nonnegative additive rate in
\eqref{v24:eq:source-order} is exactly
\begin{equation}
 \boxed{\quad r_j(a)=\bigl[\lambda_{\max}(D_j-aB_j)\bigr]_+.\quad}
 \label{v24:eq:spectral}
\end{equation}
It is a convex, nonincreasing function of $a$. If $\widehat\Xi_j(\widehat C_H)$
has certified operator-norm error at most $e_{\rm obs}$, replacing the exact
largest eigenvalue by the measured one and adding $e_{\rm obs}/\epsilon$ is a
valid upper certificate for $r_j(a)$. A further error $e_B$ in $B_j$ costs
$a e_B$. Supremums over all relevant histories and times remain necessary.

If the reader has a CPTP section $\mathsf S_*$ and exact averaged descent
$\mathsf R_*\mathcal A_{j,*}=\Psi_{j,*}\mathsf R_*$, then this matrix test can
be performed entirely in the physical algebra, with
$D_j=\epsilon^{-1}(\Psi_j^*-\mathscr E^*)\widehat C_H$ and
$B_j=\widehat C_P$. The optimal $r_j(a)$ is identical in the two algebras.
\end{proposition}
\begin{proof}
For a Hermitian matrix $D-aB$, the inequality $D-aB\preceq rI$ is equivalent
to $r\ge\lambda_{\max}(D-aB)$. A largest eigenvalue is the supremum of its
unit-vector Rayleigh quotients, each affine and nonincreasing in $a$.
This proves the formula and its stated shape. The variational formula also
gives the norm perturbation estimates.

Under descent, the source matrix is $\mathsf i(D-aB)$. Unital positivity
preserves an upper order bound. Conversely, the adjoint of the CPTP section
is positive, unital, and obeys $\mathsf S^*\mathsf i=\mathrm{id}$, so it reflects
that same bound. Thus the admissible $r$ values, and hence their minima,
agree exactly. This applies in particular to the companion's normalized
multiplicity reduction; no reset $\mathsf S_*\mathsf R_*$ is executed.
\end{proof}

This is a finite linear matrix inequality, not a proof that its slack is small.
It is smaller than full-channel tomography once the physical action insertion is
known. Estimating it from a diamond norm still gives
$e_{\rm obs}\le\delta\|\widehat C_H\|$, which recovers the old large-norm cost.
The improvement requires \emph{actual action-sensitive information}, not a
change of names or units.

\begin{proposition}[A cold-energy necessary test, including a nonzero quantum floor]
\label[proposition]{v24:prop:heating}
Let $C\succeq0$ on a physical carrier, and let $\Psi_*$ be its actual accepted
channel. If
\begin{equation}
 \Psi^*C-C\preceq\epsilon a C+\epsilon rI,
 \label{v24:eq:static-order}
\end{equation}
then, for the spectral projection $P_s=\one_{[0,s]}(C)\ne0$,
\begin{equation}
 r\ge\epsilon^{-1}
 \left[\lambda_{\max}\bigl(P_s\Psi^*C P_s|_{\Ran P_s}\bigr)
                         -(1+\epsilon a)s\right]_+.
 \label{v24:eq:low-energy}
\end{equation}
For $P_0$ onto $\Ker C$, the compressed heating is
\begin{equation}
 P_0\Psi^*C P_0
    =\sum_\ell(C^{1/2}K_\ell P_0)^*(C^{1/2}K_\ell P_0)\succeq0.
 \label{v24:eq:kernel-heat}
\end{equation}
In finite dimension, a zero additive rate $r=0$ with some finite $a\ge0$ exists
if and only if every original Kraus operator preserves $\Ker C$.
This is an energy-specific consequence of the resolved-support criterion of
\cref{v15:prop:leakage}, not a new claim of gauge invariance from mean constraints.
\end{proposition}
\begin{proof}
Compress \eqref{v24:eq:static-order} to $\Ran P_s$ and use
$P_s C P_s\preceq sP_s$. Its largest Rayleigh quotient gives
\eqref{v24:eq:low-energy}. The Kraus formula proves \eqref{v24:eq:kernel-heat}.
If $r=0$, every positive summand vanishes, so $K_\ell\Ker C\subseteq\Ker C$.
Conversely that property makes the positive operator $\Psi^*C$ annihilate
$\Ker C$, including its mixed blocks. On its orthogonal complement $C_+$
is invertible and
\[
 \Psi^*C\preceq\beta C,\qquad
 \beta=\lambda_{\max}(C_+^{-1/2}(\Psi^*C)_+C_+^{-1/2}).
\]
Take $a=(\beta-1)_+/\epsilon$. This is finite but need not be cutoff uniform.
The case $C=0$ is immediate. The spectral projection is a test subspace, not a
postselection inserted in the source process.
\end{proof}

\subsection{A concrete direct transfer already supported by this lineage}
\begin{corollary}[Decoded spatial-scalar transfer in the action norm]
\label[corollary]{v24:cor:scalar}
Use Version 13's \emph{global spectral} scalar regulator and its positive
comparison, with $A=1+\Lambda$ and $\omega=\nu h^3$. The proof of
\cref{v13:thm:main} supplies \eqref{v24:eq:target-order} with
\[
 \lambda_0\le C_T,\quad
 r_0\le C_T(h^4+\epsilon A^2/\omega),\quad
 \zeta\le C_T dA.
\]
Let an actual source have the fixed CP reader into that carrier and the
one-sided test \eqref{v24:eq:source-order} with $a\le a_0$ and rate $r_h$.
If its initial decoded \emph{energy}, not fidelity, is at most $s_h$ in
expectation, its complete scalar field and stress comparison obeys
\eqref{v24:eq:maximal} with a constant times
\begin{equation}
 b_h=s_h+h^4+\epsilon A^2/\omega+r_h.
 \label{v24:eq:scalar-budget}
\end{equation}
In particular the displayed sufficient conditions
\[
 \nu=h^3,\quad\Lambda=h^{-2},\quad\epsilon=h^{12},\qquad
 s_h=O(h^2),\quad r_h=O(h^2)
\]
give an $O(h^2)$ energy budget. Absolute pullback errors on the two tested
insertions can certify $r_h=O(h^2)$ if
\begin{equation}
 e_C+d|\chi_H|e_{\mathscr DC}=O(\epsilon h^2)=O(h^{14}).
 \label{v24:eq:insertion-precision}
\end{equation}
These are insertion errors in their physical units, \emph{not} a new necessary
precision threshold on the raw channel. No uniform purity of the actual decoded
initial state is required. The scalar action and the source-to-carrier map remain
independently specified inputs.
\end{corollary}
\begin{proof}
The quoted one-step estimate is the explicitly calculated operator inequality
in the proof of \cref{v13:thm:main}, not an inferred inequality from its averaged
conclusion. Since $d\le4\epsilon/11$, $A\ge1$ and $\omega\le1$, the positive
offset is absorbed by $C\epsilon A^2/\omega$. Apply
\cref{v24:thm:order}. The energy controls on each attained physical density
matrix in \cref{v13:prop:field,v13:thm:stress} need no purity and therefore
apply to the actual decoded states.

Both maps in \eqref{v24:eq:Xi} are unital, so $\Xi_j(I)=0$. Thus the tested
error is $\Xi_j(C(t+d))+d\chi_H\Xi_j(\mathscr DC(t+d))$, with no cost for
the identity offset. Its operator norm is at most the left side of
\eqref{v24:eq:insertion-precision}. Use $a=0$ to certify the claimed additive
rate. Substitution of the powers in \eqref{v24:eq:scalar-budget} completes the
calculation. An actual prepared source must separately supply $s_h$; the
companion's CP section makes this preparation feasible as a declared
construction but does not prove its native selection.
\end{proof}

This corollary deliberately uses the spectral branch whose direct operator
estimate is available. It does \emph{not} claim that the noncommuting tensor
cutoff in Version 22 has already passed the same local order test, or that its
$h^{154}$ sufficient diamond budget can simply be replaced by $h^{120}$.
Such a transfer needs a direct same-action certificate on that carrier.

\section{A strictly weaker locked entrance with moving fields and arbitrarily costly spectators}
\label{v24:sec:strict}

A concrete example proves that the new test can pass while a full channel
comparison fails by a fixed amount. The Hamiltonian is fixed before the source
is evaluated, and no operation uses a future reference state.

Let $\mathcal K=\C^2\oplus\C^4$, with projections $P$ and $Q=\sum_{\mu=2}^5Q_\mu$,
where the $Q_\mu$ are the four rank-one coordinate projections of $\C^4$.
Fix a nonscalar self-adjoint $H_c$ on $\Ran P$, a positive number $M\ge1$, and
\[
 H_M=H_c\oplus M I_4,\qquad
 U_\epsilon=e^{-3\ii\epsilon H_c/2}\quad\hbox{on }\Ran P.
\]
In the inherited locked normal form choose
\begin{equation}
 B_1=U_\epsilon P,\qquad B_\mu=Q_\mu\ (2\le\mu\le5),\qquad
 K_e=I/\sqrt{18}+\sqrt{2/3}\sum_{\mu=1}^5O_{e\mu}B_\mu,
 \label{v24:eq:strict-K}
\end{equation}
where $O$ is a real isometry onto $\one^\perp$ with first column
$(1,1,1,-1,-1,-1)^{\mathsf T}/\sqrt6$. This uses six original outcomes, not
an executed Kraus rotation. The event channel is exactly
\begin{equation}
 \mathcal A_*(\rho)=\tfrac13\rho+\tfrac23\left[
 U_\epsilon P\rho P U_\epsilon^*+\sum_{\mu=2}^5Q_\mu\rho Q_\mu\right].
 \label{v24:eq:strict-channel}
\end{equation}
Let $\psi_c(t)=e^{-\ii tH_c}\psi_c(0)$ for a fixed unit vector which is not an
eigenvector. The physical comparison is
\begin{equation}
 C_M(t)=P-|\psi_c(t)\rangle\langle\psi_c(t)|+M Q\succeq0.
 \label{v24:eq:strict-C}
\end{equation}
It includes the same Hamiltonian's expensive $MQ$ energy, rather than declaring
that component absent after an outcome.

\begin{theorem}[Nontrivial locked motion with order-one channel discrepancy]
\label[theorem]{v24:thm:strict}
For all sufficiently small $\epsilon>0$, the six $K_e$ in
\eqref{v24:eq:strict-K} are linearly independent and satisfy exactly
\[
 \sum_eK_e^*K_e=I,\qquad \sum_eK_e=\sqrt2 I.
\]
Each preserves both $P$ and $Q$, and every outcome has positive probability
on every nonzero state supported in $P$. On the unchanged renewal clock,
for any initial conditional state, the actual comparison $e_j=\tr\rho_jC_M(jd)$
satisfies
\begin{equation}
 \max_{j\le T/d}\E e_j\le C_T b,\qquad
 \PP(\max_{j\le T/d}e_j>u)\le C_Tb/u,\qquad
 b=\E e_0+\epsilon.
 \label{v24:eq:strict-bound}
\end{equation}
The constants are independent of $M$, even if $M\to\infty$. In contrast,
for the target unitary channel $\mathcal U_{\epsilon,*}=\operatorname{Ad}_{e^{-\ii\epsilon H_M}}$,
\begin{equation}
 \boxed{\quad\|\mathcal A_*-\mathcal U_{\epsilon,*}\|_\diamond\ge2/3.\quad}
 \label{v24:eq:strict-diamond}
\end{equation}
Thus the direct comparison can converge while the full-channel error divided
by $\epsilon$ diverges. The example is a declared locked realization, not an
identification of any unprovided native field operation.
\end{theorem}
\begin{proof}
The five $B$ operators obey $\sum B_\mu^*B_\mu=P+Q=I$. The fixed normal-form
algebra gives both locking identities and \eqref{v24:eq:strict-channel}.
If $aI+bB_1+\sum c_\mu Q_\mu=0$, restriction to $P$ gives
$aI_2+bU_\epsilon=0$. For small nonzero $\epsilon$, $U_\epsilon$ is nonscalar;
therefore $a=b=0$, and then all $c_\mu=0$. The normal-form column transformation
is invertible, proving six-coefficient independence. On $P$ each $K_e$ equals
$I_2/\sqrt{18}\pm U_\epsilon/3$, with least singular value at least
$1/3-1/\sqrt{18}>0$. On $Q$ some label probabilities may vanish; zero-probability
labels are handled in the ordinary instrument convention, not postselected.

Write $\mathscr D O=\ii[H_M,O]$. The identity
$\dot C_M+\mathscr DC_M=0$ is exact. Every positive-order time derivative
or action commutator of $C_M$ is supported on $P$ and bounded independently of
$M$. Moreover $\mathcal A^*(MQ)=MQ$. Twice integrating the $P$-unitary gives
\[
 \|\mathcal A^*C_M-C_M-\epsilon\mathscr DC_M\|
       \le\tfrac34\epsilon^2\|\mathscr D^2C_M\|\le C\epsilon^2,
 \qquad
 \|\mathcal A^*\mathscr DC_M-\mathscr DC_M\|\le C\epsilon.
\]
The same bounds hold after a time shift by $d$. The already established
phase correction $C_M+d\chi_s\mathscr DC_M+dK_0I$, with $K_0$ independent of
$M$, has expected one-tick increase at most $C d\epsilon I$; the time
Taylor terms cost $Cd^2$, and the accepted correction costs
$Cd p_P\epsilon\le Cd^2$. The $MQ$ term has exactly zero increment, so
its norm is never estimated by $M$. Applying the proof of
\cref{v24:thm:order} gives \eqref{v24:eq:strict-bound}.

Finally input the pure state $(|2\rangle+|3\rangle)/\sqrt2$ inside $Q$.
The target unitary leaves its density matrix unchanged, whereas
\eqref{v24:eq:strict-channel} multiplies its off-diagonal entries by $1/3$.
The difference has eigenvalues $\pm1/3$ and trace norm $2/3$, proving
\eqref{v24:eq:strict-diamond}. The original resolved channel is unchanged in
this test. For $M\ge1$, $I-|\psi_c(t)\rangle\langle\psi_c(t)|\preceq C_M(t)$,
so the estimate also controls full-state fidelity to the moving code state.
It does so without demanding uniform channel proximity on unused inputs.
\end{proof}

The preparation in this example is not artificially rescued by a large initial
cost: $\rho_0=|\psi_c(0)\rangle\langle\psi_c(0)|$ has $e_0=0$, and each actual
outcome preserves the code. Near-code mixed preparations are also covered by
\eqref{v24:eq:strict-bound} with their actual $e_0$. The large $MQ$ term can be
an independently varied energy coefficient, and its attained insertion is
bounded by $e_j$. No spatial locality or classical field identification is
inferred from this six-dimensional example.

\section{The companion's rigidity form can measure dissipation rather than Hamiltonian motion}
\label{v24:sec:Dirichlet}

The source-selection companion assigns a positive commutator form to its actual
reduced coefficient bank. Before interpreting its positive eigenvalues as a
physical dynamical scale, the corresponding channel must be inspected.

\begin{proposition}[Exact dissipative meaning of a calibrated unital commutator form]
\label[proposition]{v24:prop:Dirichlet}
Let $\Psi_*$ be both trace preserving and unital, with Kraus operators $L_\ell$.
For any matrix $X$ let
\[
 q_\Psi(X)=\sum_\ell\bigl(\|[L_\ell,X]\|_{\rm HS}^2+
                           \|[L_\ell^*,X]\|_{\rm HS}^2\bigr).
\]
Then, in the ordinary Hilbert--Schmidt product,
\begin{equation}
 \boxed{\quad q_\Psi(X)=4\operatorname{Re}
                 \langle X,(I-\Psi^*)X\rangle_{\rm HS}.\quad}
 \label{v24:eq:Dirichlet}
\end{equation}
For any independently fixed Hermitian $H$, let $\mathscr DX=\ii[H,X]$ and
$\mathcal R_\epsilon X=\epsilon^{-1}(\Psi^*X-X)-\mathscr DX$. For $X\ne0$,
\begin{equation}
 \|\mathcal R_\epsilon X\|_{\rm HS}
     \ge\frac{q_\Psi(X)}{4\epsilon\|X\|_{\rm HS}}.
 \label{v24:eq:tangent-floor}
\end{equation}
In particular a coherent tangent with vanishing normalized remainder on a
specified observable bank requires $q_{\Psi_h}(X_h)=o(\epsilon_h)$ for bounded,
nonvanishing normalized tests. The identity is restricted to bistochastic
channels; a nonunital faithful-state problem needs its own weighted calculation.
\end{proposition}
\begin{proof}
Expand the two commutator squares. Trace preservation and unitality make the
two positive terms in each sum equal to $2\|X\|_{\rm HS}^2$. The two cross
terms give twice the real part of $\langle X,\Psi^*X\rangle$, and for the
adjoint Kraus bank the same real part with $\Psi_*$. These real parts agree
because the two superoperators are adjoints. This gives
\eqref{v24:eq:Dirichlet}. The derivation $\mathscr D$ is skew-adjoint in HS,
so $\operatorname{Re}\langle X,\mathscr DX\rangle=0$. Consequently
\[
 \operatorname{Re}\langle X,\mathcal R_\epsilon X\rangle
                 =-q_\Psi(X)/(4\epsilon).
\]
Cauchy--Schwarz proves the bound. This does not prohibit finite-rank channels
with nearly coherent action on a protected bank; their relevant commutator
form must shrink at the displayed event scale.
\end{proof}

For the supplied panel $(D,F)=(X,Z)$, \nolinkurl{sel8:thm:calibrated-form}
identifies its original source form with $q_\Psi$ in this proposition.
Thus its positive internal-algebra rigidity gap has an exact dissipative
interpretation on the reduced bank. The structural algebra statement remains
correct. A static positive gap is not automatically the frequency of a
conservative Hamiltonian reader.

\section{An exact reduced-source heating and physical-time calculation}
\label{v24:sec:Pauli}

Retain, without changing the companion's coefficients, its supplied Pauli panel
$(D,F)=(X,Z)$. By \nolinkurl{sel8:thm:reduced-channel} its averaged physical
channel on the multiplicity qubit is
\begin{equation}
 \Psi_*(\rho)=\tfrac13\rho+\tfrac12X\rho X+\tfrac16Z\rho Z.
 \label{v24:eq:Pauli}
\end{equation}
This is a partial-trace identity on every full source input. The three summands
are a channel representation; they are not substituted for the original six
resolved edge outcomes.

\begin{theorem}[Exact contraction and minimum action-energy injection]
\label[theorem]{v24:thm:Pauli}
For \eqref{v24:eq:Pauli}, the Bloch vector is mapped by
\begin{equation}
 (x,y,z)\longmapsto(2x/3,-y/3,0).
 \label{v24:eq:Bloch}
\end{equation}
The sharp trace-distance contraction coefficient is $2/3$, the unique fixed
state is $I_2/2$, and the calibrated commutator form on the normalized $X,Y,Z$
basis has eigenvalues $4/3,16/3,4$, respectively. For every rank-one projector
$P$ and every physical density matrix $\rho$,
\begin{equation}
 \boxed{\quad
 \tr\bigl((I-P)\Psi_*(\rho)\bigr)\ge1/6.
 \quad}
 \label{v24:eq:heating-sixth}
\end{equation}
The constant is attained. With the no-factor-$1/2$ diamond convention of
Version 23 one also has the exact distances
\begin{equation}
 \inf_{U\in U(2)}\|\Psi_*-\operatorname{Ad}_U\|_\diamond=1,
 \qquad \|\Psi_*-\mathrm{id}\|_\diamond=4/3.
 \label{v24:eq:Pauli-diamond}
\end{equation}
For a target action cost $C=\Delta(I-P)$ with a
fixed $\Delta>0$, \eqref{v24:eq:static-order} therefore requires
$r\ge\Delta/(6\epsilon)$ when tested on its exact zero-energy input $P$.
No choice of finite relative rate $a$ removes this injection.
\end{theorem}
\begin{proof}
Conjugation by $X$ preserves its Bloch component and negates the other two;
conjugation by $Z$ has the analogous property. Their given weights yield
\eqref{v24:eq:Bloch}. For a traceless Hermitian qubit matrix $u\cdot\sigma$,
its trace norm is $2|u|$. The largest singular value of this diagonal Bloch
map is $2/3$, attained on the $X$ axis. Its three nonconstant eigenvalues
have modulus less than one, proving uniqueness of $I_2/2$.
\Cref{v24:prop:Dirichlet} gives the stated calibrated eigenvalues, agreeing
with the inherited static calculation rather than replacing it.

Every output Bloch vector has length at most $2/3$, hence its least eigenvalue
is at least $(1-2/3)/2=1/6$. Pairing with $I-P$ proves the energy injection.
Take $P=\rho=(I+X)/2$ for equality.
For the channel distances, its normalized Choi state is diagonal in the
orthogonal Pauli Bell basis with weights $1/3,1/2,0,1/6$. Its largest
eigenvalue is $1/2$, whereas a unitary channel has a pure normalized Choi
state. The projector onto that pure state gives a trace-distance lower
bound $2(1-1/2)=1$. The choice $U=X$ attains it: the triangle inequality
for the Pauli-channel difference gives the sum of the absolute coefficient
differences, equal to one, and the Bell input gives the same lower bound.
The identical argument against the identity gives $4/3$.
On $\rho=P$, the right side of
\eqref{v24:eq:static-order} has expectation $\epsilon r$, while the left side
is at least $\Delta/6$. This proves the last claim. It is a dynamical test of
one supplied source, not a change to its calibrated internal algebra.
\end{proof}

The other supplied panel $(D,F)=(Z,Z)$ has channel
$\rho\mapsto\rho/3+2Z\rho Z/3$, Bloch multipliers $(-1/3,-1/3,1)$, and hence
a retained diagonal classical sector. The differing classical sector is not
silently promoted to an autonomous coherent qubit. These formulas follow from
the same one-line Pauli calculation and delimit the negative result above.

\begin{theorem}[The supplied renewal clock amplifies the reduced contraction]
\label[theorem]{v24:thm:mixing}
Apply the unchanged original source instrument whose reduction is
\eqref{v24:eq:Pauli} on the actual clock \eqref{v24:eq:clock}. Assume that the
physical reduction is unchanged on nonaccepted ticks. For any initial clock
phase and source state, its actual decoded states obey, at $t_j=jd$,
\begin{equation}
 \boxed{\quad
 \bigl\|\E\rho_j-I_2/2\bigr\|_1
          \le\frac65\exp\!\left(-\frac{t_j}{3\epsilon}\right).
 \quad}
 \label{v24:eq:physical-mixing}
\end{equation}
For any deterministic rank-one target $P_*(t)$, including a smooth trajectory
of an independently specified bounded Hamiltonian,
\begin{equation}
 \E\{1-\tr(\rho_jP_*(t_j))\}
       \ge\frac12-\frac35\exp\!\left(-\frac{t_j}{3\epsilon}\right).
 \label{v24:eq:moving-failure}
\end{equation}
At any fixed positive physical time under $\epsilon\to0$, the actual decoded
states therefore cannot converge in probability to that pure target. More
quantitatively, for $\xi_j=1-\tr(\rho_jP_*(t_j))$,
\begin{equation}
 \PP(\xi_j>1/4)\ge
 \left[\frac13-\frac45e^{-t_j/(3\epsilon)}\right]_+.
 \label{v24:eq:failure-prob}
\end{equation}
This conclusion does not assume that $\rho_j$ is an autonomous labelled
process or that different times are independent.
\end{theorem}
\begin{proof}
Condition on a complete clock and acceptance sequence. Repeated use of the
all-input averaged reduction gives averaged output $\Psi_*^{A_j}\rho_0$, where
$A_j$ is its accepted count. If the initial source is random or correlated
with its initial phase, apply this identity conditionally; the uniform bound
below is unchanged. The Bloch length contracts by at most $q^{A_j}$, with
$q=2/3$. The exact generating function for $\E_s q^{A_j}$ is
\[
 B_\vartheta(q)^j\one,\qquad
 B_\vartheta(q)=
 \begin{pmatrix}1/5&4/5\\(2/3)(1-\vartheta+\vartheta q)&1/3\end{pmatrix}.
\]
Its positive eigenvalue at $q=2/3$ is
\[
 \lambda_\vartheta=\frac{4+\sqrt{121-40\vartheta}}{15},\qquad
 v_\vartheta=\bigl(1,(5\lambda_\vartheta-1)/4\bigr)^{\mathsf T}.
\]
For $0<\vartheta\le1$, $13/15\le\lambda_\vartheta<1$ and
$5/6\le(v_\vartheta)_P\le1$. Positivity gives
$B_\vartheta(q)^j\one\le(6/5)\lambda_\vartheta^jv_\vartheta
\le(6/5)\lambda_\vartheta^j\one$. Also
\[
 1-\lambda_\vartheta
  =\frac{8\vartheta/3}{11+\sqrt{121-40\vartheta}}
  \ge\frac{4\vartheta}{33}=\frac d{3\epsilon}.
\]
Thus $\E q^{A_j}\le(6/5)e^{-jd/(3\epsilon)}$, uniformly even for rare
acceptance. This proves \eqref{v24:eq:physical-mixing}.

For a traceless Hermitian qubit difference $X$, pairing with a rank-one
projector has magnitude at most $\|X\|_1/2$. Subtract the maximally mixed
value to obtain \eqref{v24:eq:moving-failure}. Finally $0\le\xi_j\le1$ gives
$\E\xi_j\le1/4+(3/4)\PP(\xi_j>1/4)$, proving
\eqref{v24:eq:failure-prob}. In particular a bounded random variable
converging in probability to zero could not have this positive expectation
limit. No statement that each individual conditional state tends to $I_2/2$
is made; only its ensemble mean does so, and that is enough for the obstruction.
\end{proof}

\begin{corollary}[First-acceptance heating also survives arbitrary rarity]
\label[corollary]{v24:cor:first}
From a fresh completed-boundary entrance let $\tau$ be the first accepted time.
For any predetermined rank-one comparison $P_*(t)$, and energy
$e_\tau=\Delta\{1-\tr(\rho_\tau P_*(\tau))\}$,
\[
 \PP\{\tau\le T,\ e_\tau>\Delta/12\}
        \ge \frac1{11}(1-\epsilon/T)_+.
\]
The event concerns the conditional decoded energy after the actual first
record, not a new measurement. The clock and source must retain the same
independence and no-response reader conditions as above.
\end{corollary}
\begin{proof}
At any fixed waiting history the actual first-outcome mean has
$\E(e_\tau\mid\tau)\ge\Delta/6$ by \cref{v24:thm:Pauli}. Since
$0\le e_\tau\le\Delta$, its conditional probability of exceeding
$\Delta/12$ is at least $1/11$. Multiply by
$\PP(\tau\le T)\ge(1-\epsilon/T)_+$ from the inherited mean waiting time.
The last inequality is Markov's inequality for the same waiting law.
\end{proof}

\section{A faithful instantaneous decoder cannot restore the lost coherent distinction}
\label{v24:sec:recovery}

Version 23 correctly notes that arbitrary noninvertible readers need not preserve
its full-record lower bound. The next statement pays precisely for a reader's
ability to retain distinguishable physical inputs. It is not a no-go theorem
for all classical coarse variables or history-dependent readers.

\begin{theorem}[Contraction obstruction after a recoverable CP reduction]
\label[theorem]{v24:thm:recovery}
Let an actual source channel $\Psi_*$ have trace-distance contraction coefficient
at most $q<1$ on Hermitian state differences. Let $\mathsf R_*,\mathsf S_*$ be
CPTP maps to and from a physical matrix algebra of dimension at least two,
with
\[
 \|\mathsf R_*\mathsf S_*-\mathrm{id}\|_\diamond\le\kappa.
\]
For every physical unitary channel $\mathcal U_*$,
\begin{equation}
 \boxed{\quad
 \|\mathsf R_*\Psi_*-\mathcal U_*\mathsf R_*\|_\diamond
                       \ge(1-q-\kappa)_+.
 \quad}
 \label{v24:eq:recovery}
\end{equation}
In particular the reduced Pauli panel has a bound $1/3$ for every CP reader
with an exact CPTP section and a nontrivial quantum output. If a separately
fixed target action channel is at diamond distance $\chi_\epsilon$ from its
unitary step, the lower bound becomes
$(1-q-\kappa-\chi_\epsilon)_+$.
\end{theorem}
\begin{proof}
Take two orthogonal pure physical states $\rho_+,\rho_-$ and put
$\tau_\pm=\mathsf R_*\mathsf S_*\rho_\pm$. The recovery estimate gives
$\|\tau_+-\tau_-\|_1\ge2-2\kappa$. Let the left side of
\eqref{v24:eq:recovery} be $\delta$. Applying its definition to the two
source states $\mathsf S_*\rho_\pm$ and taking a difference gives
\[
 \|\mathsf R_*\Psi_*\mathsf S_*(\rho_+-\rho_-)\|_1
       \ge2-2\kappa-2\delta.
\]
The same norm is at most $q\|\mathsf S_*\rho_+-\mathsf S_*\rho_-\|_1\le2q$,
by positivity, trace preservation, and the assumed contraction. Rearranging
proves the bound. Only ordinary input states were used, so the result holds
for the diamond norm as well. Replacing $\mathcal U_*$ by a channel within
$\chi_\epsilon$ costs at most that amount after composition with the CPTP
reader. This gives the final assertion.
\end{proof}

The section is an information-retention condition, not a demanded physical
reset between events. A constant reader has no such section and is not ruled
out. Nor is a growing tensor source, a reader of the full recorded history,
a new momentum-odd extension, or a classical variable which deliberately loses
most quantum inputs covered without further argument. For the original
$W\otimes\C^2$ source, the bound applies to reductions \emph{factoring through
this identified multiplicity qubit}; it does not eliminate all other possible
readers of the larger source. A fixed positive internal algebra gap therefore
supplies a useful rejection test for one proposed chronology, not a universal
rejection of Standard-Model emergence.

\section{Outcome: distinguish a cold action bank from a merely rigid internal algebra}
\label{v24:sec:outcome}

The positive improvement is a strict one. A same-action comparison can be
propagated on every original source history by an all-input upper order bound
on its own insertions. It need not pay for source errors on the entire quantum
algebra, an initial pure-state fidelity, or a large spectator energy norm.
The explicit six-outcome example carries a nonconstant physical reference and
passes this test while remaining a fixed distance from its target full channel.
The scalar corollary shows where an existing field-and-stress theorem really
supplies the needed operator input, and where a tensor-cutoff extension still
needs its own calculation.

The actual benchmark gives the complementary message. The supplied $(X,Z)$
multiplicity panel is a full internal matrix algebra, but its retained channel
heats every rank-one zero-energy state by at least one sixth of a fixed gap per
acceptance. On the calibrated dense renewal clock its ensemble mean approaches
$I_2/2$ at the explicit physical rate in \eqref{v24:eq:physical-mixing}.
It cannot be declared a noncollapsed conservative qubit reader by forgetting
its original labels or by a recoverable instantaneous CP re-encoding. The
companion's algebraic source-selection conclusion is unchanged; its dynamical
interpretation requires this extra test.

The next upstream acquisition is consequently concrete:
\[
 \boxed{\begin{gathered}
 \text{actual reduced channel on the physical action bank,}\\
 \text{its cold-state heating, and its normalized action drift}
 \end{gathered}}
\]
The next actual source must pass these tests with its own preparation and
clock. Smoothing failed records does not establish this same-action entrance.
A passing one-sided certificate supplies the more economical transfer. The
intended Einstein--SM source bank and field/action identification remain
unverified, as do the independent geometric and continuum requirements.
In particular the full observed Einstein inverse and the nonzero-reference
non-Abelian spatial quantum theorem remain unresolved within this lineage.

\paragraph{Verification and preservation.}
All bytes of Version 23 before its final end-document command are preserved.
New labels have prefix \texttt{v24:}. The companion script checks the one-sided
spectral certificates, a direct conditional tick calculation with a CP reader,
the exact locked moving-code example, its order-one channel-distance witness,
Pauli Bloch and commutator spectra, the actual non-Poisson clock eigenvector
bound, cold-energy injection, the trace-contraction recovery inequality, and
the scalar scale budget. Finite tests are diagnostics of the written proofs,
not machine-checked formalizations or a simulation of the Einstein--SM system.
The source benchmark and the constructed locked example are kept separate in
both the script and this text.

\begingroup
\renewcommand{\refname}{Additional references for Version 24}
\begin{thebibliography}{9}
\small\raggedright
\bibitem{v24:TV} F.~Ticozzi and L.~Viola,
\emph{Quantum information encoding, protection, and correction from trace-norm
isometries}, Physical Review A \textbf{81} (2010), 032313.
\href{https://doi.org/10.1103/PhysRevA.81.032313}{doi:10.1103/PhysRevA.81.032313}.
Background for recoverable quantum encodings and protected sectors. The finite
contraction and locked-clock bounds here are proved directly.
\bibitem{v24:RKW} D.~Reeb, M.~J.~Kastoryano, and M.~M.~Wolf,
\emph{Hilbert's projective metric in quantum information theory}, 2011 preprint.
\href{https://arxiv.org/abs/1102.5170}{arXiv:1102.5170}.
Background for order and distinguishability contraction. No projective-metric
theorem is imported as a substitute for the explicit Pauli calculation.
\end{thebibliography}
\endgroup
\addtocontents{toc}{\protect\endgroup}
% END IMMUTABLE VERSION 24 BODY
% APPEND VERSION 25 HERE, BEFORE THE FINAL END-DOCUMENT COMMAND.


% BEGIN IMMUTABLE VERSION 25 BODY
\clearpage
\begin{center}
{\Large\bfseries Version 25: Original histories before an assumed recovery}\\[.5em]
{\large Exact quantum-retention costs of the supplied six-edge source,
passive word effects, and the distinction between a record and its purification}
\end{center}
\makeatletter
\addtocontents{toc}{\protect\begingroup}
\addtocontents{toc}{\protect\makeatletter}
\addtocontents{toc}{\protect\def\protect\l@section{\protect\bfseries\protect\@dottedtocline{1}{0em}{2.5em}}}
\makeatother
\addcontentsline{toc}{part}{Version 25: Original-history recovery and a sharp two-acceptance source test}

\section{Audit: the history-reader alternative has to use the original instrument}
\label{v25:sec:audit}

Version 24 proves a dissipative obstruction for the supplied multiplicity
Pauli channel and explicitly leaves a reader of the full original history
outside its instantaneous-recovery theorem. That distinction is real: a
random-unitary channel can lose coherence when its labels are forgotten and
retain it perfectly when its actual unitary labels are known. It would be
incorrect, however, to replace the present source's six recorded coefficients
by a more favorable Kraus decomposition of its reduced channel. This body
evaluates the original coefficients, with all six labels and the complete
terminal source system available to the reader.

\paragraph{Inherited inputs and nonduplication.}
The source-selection companion \cite{v23:SEL17}, at
\nolinkurl{sel7:eq:matrix-family},
\nolinkurl{sel8:eq:contrast-calibration}, and
\nolinkurl{sel8:thm:quadratic-blindness}, supplies the six-edge family and its
$(D,F)=(X,Z)$ benchmark. Its external covariance, exact normalization,
independence, and partial-trace reduction are inherited, not new results.
The source is the six-dimensional triplet--qubit benchmark, not the
fourteen-dimensional witness of \cref{v23:cor:actual-distance}.
The calibrated two-phase clock and its accepted waiting mean are retained
from Versions 10 and 23. No clock proof, internal-algebra saturation theorem,
or reduced Pauli mixing theorem is repeated.
Targeted semantic searches and exact-term checks in the cumulative notes and
the source-selection companion found the earlier rank, record-descent, and
information-disturbance results, but not the explicit recovery values or the
two-acceptance matrices calculated here. This is a limited overlap audit,
not a claim to have independently checked every supplied proof.

\paragraph{Established recovery mathematics.}
Optimal recovery from a specified measured environment, the scalar-effect
criterion, and qubit determinant formulas have established treatments in
Gregoratti--Werner \cite{v25:GW02,v25:GW04}. Their polar-recovery argument is
included below with a rectangular encoded input so that the retained source
memory, diamond-norm convention, and passive-history use are explicit. No
priority claim is made for that general recovery principle. The additions
are its exact evaluation and preparation optimization for the unchanged
source, its two-acceptance obstruction even with an arbitrary retained
purification, and the corresponding original-clock and action-bank boundaries.

\paragraph{Which question is being answered.}
The input is an \emph{unknown coherent qubit} in the independently identified
multiplicity factor. A reader is a fixed CPTP map, possibly depending on all
recorded labels and the known time, but not on the unknown input state. The
source itself is not corrected or reset between events. Allowing every such
terminal reader makes the negative result stronger; it does not assert that
the optimizing reader is an admitted native operation. Ancillas carrying a
second copy or a bypass encoding of the unknown logical state are not allowed.
Independent preparation memory is considered explicitly and generously.
A classical field reader need not retain an arbitrary coherent qubit, so the
results do not establish a no-go theorem for classical Einstein--SM closure.

\section{The effect bank gives an exact qubit recovery norm}
\label{v25:sec:recovery}

Let $V:\C^2\to\mathcal H$ be an isometry, and let $K_w$ be the actual
coefficients of a finite efficient instrument, or actual products along a
finite recorded history. Put
\[
 L_w=K_wV,\qquad F_w=L_w^*L_w,\qquad \sum_wF_w=I_2,
 \qquad
 \mathcal T_*(\rho)=\bigoplus_w L_w\rho L_w^*.
\]
The direct sum is the classical record flag together with the full terminal
quantum output. Each output block may be larger than the logical input.
Let $|\Omega\rangle=(|00\rangle+|11\rangle)/\sqrt2$, and for a qubit
channel $\Lambda_*$ define
\[
 F_{\rm e}(\Lambda_*)=
 \langle\Omega|(\Lambda_*\otimes\mathrm{id})(|\Omega\rangle\langle\Omega|)
 |\Omega\rangle.
\]
The diamond norm has no factor $1/2$, as in Version 23.

\begin{lemma}[Polar recovery with the original flags and a rectangular code]
\label[lemma]{v25:lem:polar}
For CPTP terminal readers $\mathcal R_*$ of the entire flagged output,
\begin{align}
 \sup_{\mathcal R_*}F_{\rm e}(\mathcal R_*\mathcal T_*)
 &=\frac14\sum_w(\tr\sqrt{F_w})^2
   =\frac{1+s(\mathcal T)}2,\qquad
 s(\mathcal T)=\sum_w\sqrt{\det F_w},
 \label{v25:eq:fidelity}\\
 \inf_{\mathcal R_*}
 \|\mathcal R_*\mathcal T_*-\mathrm{id}\|_\diamond
 &=1-s(\mathcal T).
 \label{v25:eq:diamond}
\end{align}
The same infimum holds against any prescribed logical unitary channel.
One optimal reader performs the polar inverse on each original output block.
In particular, perfect recovery is possible precisely when every $F_w$ is
a scalar multiple of $I_2$.
\end{lemma}
\begin{proof}
A CPTP map on a classical direct sum restricts to a CPTP map on each block.
Write its block Kraus operators as $R_{w,a}$, and extend the polar partial
isometry of $L_w$ to an isometry $W_w:\C^2\to\mathcal H$ so that
$L_w=W_w\sqrt{F_w}$. The output dimension permits the extension, also when
$F_w$ is singular. In an eigenbasis of $\sqrt{F_w}$, the vector with entries
$\langle i|R_{w,a}W_w|i\rangle$ indexed by $a$ has Euclidean norm at most one:
completeness bounds the sum of its squared entries by
$\sum_a\|R_{w,a}W_w|i\rangle\|^2=1$.
The vector triangle inequality therefore gives
\[
 \sum_a|\tr(R_{w,a}L_w)|^2\le(\tr\sqrt{F_w})^2.
\]
The channel fidelity is one quarter of the sum on the left, proving the
upper bound. The CPTP map
\[
 R_{w,*}(Y)=W_w^*YW_w+
          \tr\bigl((I-W_wW_w^*)Y\bigr)\tau_0
\]
for any fixed qubit state $\tau_0$ is an allowed extension of the polar
inverse and gives the corrected channel
$\Lambda_*(\rho)=\sum_w\sqrt{F_w}\rho\sqrt{F_w}$.
This attains the bound. The identity
$(\tr\sqrt F)^2=\tr F+2\sqrt{\det F}$ and $\sum_w\tr F_w=2$
prove \eqref{v25:eq:fidelity}.

For the exact diamond statement write
$\sqrt{F_w}=a_wI+b_w\cdot\sigma$, with real $a_w\ge0,b_w\in\R^3$.
Completeness gives
\[
 \sum_w(a_w^2+|b_w|^2)=1,\qquad \sum_w a_wb_w=0.
\]
Hence
\[
 \Lambda_*=f\,\mathrm{id}+\sum_w\operatorname{Ad}_{b_w\cdot\sigma},
 \qquad f=\sum_wa_w^2=F_{\rm e}(\Lambda_*).
\]
When $f<1$, the last sum divided by $1-f$ is CPTP since
$(b_w\cdot\sigma)^2=|b_w|^2I$. Thus
$\|\Lambda_*-\mathrm{id}\|_\diamond\le2(1-f)$.
For \emph{every} corrected channel, the Bell input and the two-outcome test
$|\Omega\rangle\langle\Omega|,I-|\Omega\rangle\langle\Omega|$ give
$\|\mathcal R_*\mathcal T_*-\mathrm{id}\|_\diamond
 \ge2(1-F_{\rm e}(\mathcal R_*\mathcal T_*))$.
The fidelity maximum makes this lower bound $1-s$, and the polar reader
attains it. Composing a reader with a fixed logical unitary proves the
unitary-target assertion. Finally
$\sqrt{\det F_w}\le\tr F_w/2$, with equality exactly for scalar effects.
Every deficit is nonnegative, proving the last assertion.
\end{proof}

This is a \emph{terminal} recovery calculation. The matrices $W_w$ are
functions of known source coefficients, encoding, and word, not of an unknown
input state. No polar inverse is inserted into the actual source chronology.
The result also does not say that the corrected logical action equals a
previously specified physical action; it measures preservation of quantum
information before that separate incidence test.

\subsection{A negative test needs only original word probabilities}
\begin{proposition}[A record-probability lower bound for any terminal quantum reader]
\label[proposition]{v25:prop:probability}
Let $\mathcal T_w$ be arbitrary CP branches of a trace-preserving flagged
channel with qubit input; branches need not be efficient. For two orthogonal
pure input states set
$p_w=\tr\mathcal T_w(|0\rangle\langle0|)$ and
$q_w=\tr\mathcal T_w(|1\rangle\langle1|)$.
Then
\begin{equation}
 \inf_{\mathcal R_*}
 \|\mathcal R_*\mathcal T_*-\mathrm{id}\|_\diamond
 \ge 1-\sum_w\sqrt{p_wq_w}.
 \label{v25:eq:probability-bound}
\end{equation}
If normalized measured distributions $\widehat p,\widehat q$ have certified
$\ell^1$ errors at most $a,b$, respectively, the right side may be replaced by
\begin{equation}
 \left[1-\sum_w\sqrt{\widehat p_w\widehat q_w}-\sqrt a-\sqrt b\right]_+.
 \label{v25:eq:probability-precision}
\end{equation}
\end{proposition}
\begin{proof}
After a blockwise CPTP reader, let $M_{w,k}$ be the composite Kraus matrices.
Their squared output norms on $|0\rangle,|1\rangle$ still sum to $p_w,q_w$
within each flag. Consequently
\[
 \sum_k|\tr M_{w,k}|^2
 \le p_w+q_w+2\sqrt{p_wq_w},
\]
by Cauchy--Schwarz on the two vectors of diagonal matrix elements.
Summing and dividing by four bounds entanglement fidelity by
$(1+\sum_w\sqrt{p_wq_w})/2$. The Bell-input argument from the preceding
lemma proves \eqref{v25:eq:probability-bound}. For the robust statement use
$\|\sqrt p-\sqrt{\widehat p}\|_2^2\le\|p-\widehat p\|_1$ and the
unit $\ell^2$ norms of square roots of probability vectors. Two applications
of Cauchy--Schwarz bound the affinity error by $\sqrt a+\sqrt b$.
No positive lower bound on any particular word probability is required.
\end{proof}

\section{A passive-history determinant identity and its physical-time criterion}
\label{v25:sec:words}

Allow history-dependent actual coefficients $K_{e|w}$, with
$\sum_eK_{e|w}^*K_{e|w}=I$ on each current source carrier. Start from an
isometric logical preparation $V$. For a word $w$, let $L_w$ be the actual
product times $V$, with $L_\varnothing=V$, and put $F_w=L_w^*L_w$.
The source may have more than two quantum dimensions. The following identity
therefore uses the current image of the logical input, not determinants of
whole-source matrices.

\begin{theorem}[Exact loss accumulation without intermediate recovery]
\label[theorem]{v25:thm:word}
For a polar isometry $V_w$ of $L_w$, define
\[
 Q_{e|w}=V_w^*K_{e|w}^*K_{e|w}V_w,\qquad
 \ell_w=1-\sum_e\sqrt{\det Q_{e|w}},\qquad
 s_n=\sum_{|w|=n}\sqrt{\det F_w}.
\]
One has $0\le\ell_w\le1$, $s_0=1$, and
\begin{equation}
 \boxed{\quad
 s_{n+1}=s_n-\sum_{|w|=n}\sqrt{\det F_w}\,\ell_w.
 \quad}
 \label{v25:eq:word-ledger}
\end{equation}
The best terminal recovery error in diamond norm from all $n$ original labels
and the terminal source equals $1-s_n$. In particular it cannot decrease
when further source operations are performed.
If $\ell_w\le a$ on every attained code image then
$s_n\ge(1-a)^n$; if $\ell_w\ge a$ whenever $\det F_w>0$, then
$s_n\le(1-a)^n$.
\end{theorem}
\begin{proof}
The polar decomposition gives
\[
 F_{we}=\sqrt{F_w}\,Q_{e|w}\sqrt{F_w},\qquad
 \det F_{we}=\det F_w\det Q_{e|w}.
\]
These formulas remain true when $F_w$ is singular, and all such terms have
zero determinant. Summing over descendants proves
\eqref{v25:eq:word-ledger}. Completeness gives
$\sum_eQ_{e|w}=I_2$; the determinant--trace inequality proves
$0\le\ell_w\le1$. Apply \cref{v25:lem:polar} to the full word instrument
for the optimal error, and iterate the two scalar inequalities for the
bounds. The polar isometries appear only in the calculation; the products
$L_w$ remain the unchanged source products. The weights $\sqrt{\det F_w}$
are retention weights, \emph{not} asserted to be the Born probabilities of a
particular input state.
\end{proof}

\begin{corollary}[An accepted-event retention scale, separate from action incidence]
\label[corollary]{v25:cor:word-clock}
On the independent renewal clock of Version 23, suppose the previous upper
bound $\ell_w\le a_\epsilon$ holds for all words and current code images.
With all waiting and outcome records retained up to physical time $T$, a
terminal recovery exists with diamond error at most
\begin{equation}
 a_\epsilon\,\E A_T\le
 a_\epsilon(T/\epsilon+1),
 \label{v25:eq:retention-scale}
\end{equation}
where $A_T$ is the actual accepted count. Thus
$a_\epsilon=o(\epsilon)$ is a sufficient all-history quantum-retention
condition. This is not a derivation of a desired Hamiltonian from the source.
\end{corollary}
\begin{proof}
Condition on the complete clock and acceptance record, which is independent
of the unknown input. For $n$ acceptances, \cref{v25:thm:word} and
$1-(1-a)^n\le na$ give a blockwise terminal reader of error at most $na$.
Average using the convexity of the diamond norm. The accepted-count estimate
is the already proved \cref{v23:prop:hybrid}. No independence between attained
quantum states is required. Nor is intermediate re-encoding performed.
\end{proof}

The exact loss identity is often the useful test: a single prefix with
nonvanishing retention loss already prevents a uniform all-input coherent
limit, even if its later physical one-sided energy test happens to hold on a
smaller observable class. Retention is necessary for a faithful coherent
qubit identification; it is not a prerequisite for every classical decoder.

\section{One original acceptance: sharp recovery cost for the standard section}
\label{v25:sec:standard}

We now reconstruct the supplied source in a fixed real orthonormal triplet
frame. Let $X,Z$ be Pauli matrices and set
\begin{equation}
 A=\frac{2I+Z+3X}{6},\qquad
 B=\frac{2I+Z-3X}{6},\qquad C=\frac{I-Z}{3}.
 \label{v25:eq:ABC}
\end{equation}
Thus $D=A-B=X$ and $F=A+B-2C=Z$ exactly as in the companion. For each
unordered coordinate pair $i,j$ in $\{1,2,3\}$, let
\[
 P_{ij,\pm}=\frac12|e_i\pm e_j\rangle\langle e_i\pm e_j|,
 \qquad Q_{ij}=I_3-P_{ij,+}-P_{ij,-}.
\]
These are the six tetrahedral edge projectors in an equivalent orthonormal
triplet frame. For $e=(ij,\pm)$ let $\bar e=(ij,\mp)$ and define
\begin{equation}
 K_e=\frac1{\sqrt2}
 (P_e\otimes A+P_{\bar e}\otimes B+Q_e\otimes C).
 \label{v25:eq:source}
\end{equation}
The inherited identities give
$\sum_eK_e^*K_e=I_6$ and $\sum_eK_e=\sqrt2I_6$.
All coefficients are real Hermitian in this frame. The source's averaged
partial-trace channel is \cref{v24:eq:Pauli}; its three Pauli summands are
not the original recorded outcomes.

Use the companion's standard CPTP section of the physical partial trace,
\[
 \mathsf S_{0,*}(\rho)=\frac{I_3}{3}\otimes\rho,
 \qquad
 \mathcal T_{0,*}(\rho)=\bigoplus_{e=1}^6
                         K_e\mathsf S_{0,*}(\rho)K_e^*.
\]
The entire $\C^3\otimes\C^2$ terminal state, not just its reduced qubit,
is available to the optimizing reader.

\begin{theorem}[Sharp recovery deficit of the unchanged six-edge source]
\label[theorem]{v25:thm:standard}
For the preceding standard preparation and original first accepted instrument,
\begin{equation}
 \boxed{\quad
 \sup_{\mathcal R_*}F_{\rm e}(\mathcal R_*\mathcal T_{0,*})=\frac23,
 \qquad
 \inf_{\mathcal R_*}\|\mathcal R_*\mathcal T_{0,*}-\mathrm{id}\|_\diamond
 =\frac23.
 \quad}
 \label{v25:eq:standard}
\end{equation}
The same optimal distance holds to any prescribed logical unitary. An optimal
reader yields the corrected Pauli channel
\begin{equation}
 \Lambda_{\rm pol,*}(\rho)
     =\frac23\rho+\frac15X\rho X+\frac{2}{15}Z\rho Z.
 \label{v25:eq:standard-corrected}
\end{equation}
\end{theorem}
\begin{proof}
Orthogonality of the three external projections gives exactly, for each $e$,
\[
 K_e\mathsf S_{0,*}(\rho)K_e^*
 =\frac16\left(P_e\otimes A\rho A+
                P_{\bar e}\otimes B\rho B+Q_e\otimes C\rho C\right).
\]
The flagged channel is CPTP-equivalent, for this single experiment, to
$\rho\mapsto(A\rho A)\oplus(B\rho B)\oplus(C\rho C)$:
measure the already block-diagonal external projections conditional on $e$
and trace the triplet to obtain the latter; conversely draw $e$ uniformly
and prepare the corresponding rank-one external projection. Both are fixed
CPTP postprocessings and reproduce the channels on every logical input.
This proves equality of their optimal recovery figures. It is an analytical
equivalence of output experiments, not an added intermediate measurement in
the original source.

Directly,
\[
 \tr|A|=\tr|B|=\frac{\sqrt{10}}3,\qquad
 \tr|C|=\frac23,\qquad
 |\det A|=|\det B|=\frac16,\quad\det C=0.
\]
The polar lemma therefore gives
$F_{\rm e}^{\max}=(10/9+10/9+4/9)/4=2/3$ and the exact diamond error
$1-(1/6+1/6)=2/3$.
For the explicit corrected map use
\[
 |A|=\frac{\sqrt{10}}6I+\frac X{\sqrt{10}}
                      +\frac Z{3\sqrt{10}},\qquad
 |B|=\frac{\sqrt{10}}6I-\frac X{\sqrt{10}}
                      +\frac Z{3\sqrt{10}},\qquad |C|=C.
\]
All scalar--Pauli and $XZ$ cross terms cancel in
$\sum_{L=A,B,C}|L|\rho|L|$, leaving
\eqref{v25:eq:standard-corrected}. The original operators, their normalization,
and all flags have been retained.
\end{proof}

For comparison only, an instrument which actually recorded the three
coefficients $I/\sqrt3$, $X/\sqrt2$, and $Z/\sqrt6$ would have the same averaged
multiplicity channel and would be perfectly corrected by those recorded
Pauli inverses. It is a \emph{different instrument}. The source computation
above proves that its perfect recovery cannot be assigned to the original
six-edge experiment merely because the averaged channels agree.

\section{Changing only the unpurified external preparation cannot make the first step exact}
\label{v25:sec:section}

\begin{lemma}[Sections of the identified partial trace have no hidden logical bypass]
\label[lemma]{v25:lem:section}
Every CPTP map $\mathsf S_*:M_2\to\mathcal B(\C^3\otimes\C^2)$
with $\operatorname{Tr}_{\C^3}\mathsf S_*=\mathrm{id}$ is
$\mathsf S_*(\rho)=\tau\otimes\rho$ for a fixed density matrix $\tau$.
The same assertion holds with any additional preparation-memory factor on
the traced side.
\end{lemma}
\begin{proof}
Apply $\mathsf S_*\otimes\mathrm{id}$ to the normalized Bell state.
Its logical-output/reference marginal is the pure Bell state. A positive
joint state with a pure marginal has product support on that marginal: pair
with its complementary projection and use the positive square root to remove
all complementary rows and columns. The Choi state therefore equals
$\tau\otimes|\Omega\rangle\langle\Omega|$. Equality of Choi matrices proves
the channel formula. This also proves the additional-memory assertion.
\end{proof}

\begin{proposition}[A uniform nonzero first-step deficit without access to a purifier]
\label[proposition]{v25:prop:unpurified}
For every section $\tau\otimes\rho$ of the identified physical partial trace,
when no purifier of $\tau$ is available at the terminal cut, the original
flagged first-acceptance channel has
\begin{equation}
 \sup_{\mathcal R_*}F_{\rm e}(\mathcal R_*\mathcal T_{\tau,*})\le\frac{83}{84},
 \qquad
 \inf_{\mathcal R_*}\|\mathcal R_*\mathcal T_{\tau,*}-\mathrm{id}\|_\diamond
 \ge\frac1{42}.
 \label{v25:eq:all-unpurified}
\end{equation}
These constants are sufficient uniform bounds, not claimed optimal. For the
particular pure external vector $(1,\ii,0)/\sqrt2$, the exact optimal fidelity is
$(6+\sqrt{13}+\sqrt5)/12$.
\end{proposition}
\begin{proof}
First let $\tau=|w\rangle\langle w|$, $\|w\|=1$. At a complementary
pair of edge labels, with missing coordinate $k$ and other coordinates $i,j$,
put $r_k=|w_k|^2$ and $g_{ij}=\operatorname{Re}(\bar w_iw_j)$.
Compression of the original effect to $w\otimes\C^2$ gives
\begin{equation}
 F_{ij,\pm}=a_kI\ \pm\frac{g_{ij}}3X+
                    \frac{1-3r_k}{18}Z,
 \qquad a_k=\frac{7-3r_k}{36}\le\frac7{36}.
 \label{v25:eq:pure-effects}
\end{equation}
Indeed $A^2=7I/18+X/3+Z/9$, $B^2=7I/18-X/3+Z/9$,
and $C^2=2(I-Z)/9$. The root probabilities satisfy
$p-q=\pm2g_{ij}$ and $p+q=1-r_k$, proving the formula.
Let $F_e=a_eI+b_e\cdot\sigma$.
The real matrix $G=(\operatorname{Re}(\bar w_iw_j))$ is positive with trace
one and rank at most two, since it is the sum of the outer products of the
real and imaginary parts of $w$. Write $t_2=\tr G^2$ and
$s_2=\sum_kr_k^2$. Then $t_2\ge1/2$ and $s_2\le t_2$, and direct summation
over the three complementary pairs gives
\[
 \sum_e|b_e|^2
 =\frac{6t_2-3s_2-1}{54}\ge\frac1{108}.
\]
Since $\sum_ea_e=1$ and
\[
 a_e-\sqrt{a_e^2-|b_e|^2}
 =\frac{|b_e|^2}{a_e+\sqrt{a_e^2-|b_e|^2}}
 \ge\frac{18}{7}|b_e|^2,
\]
the optimal fidelity loss is at least $1/84$ by
\cref{v25:lem:polar}. For mixed $\tau$, any fixed reader's entanglement
fidelity is affine in $\tau$. Decompose $\tau$ into pure states and apply
the same upper bound to that reader on each component. Taking the supremum
preserves the bound. The Bell lower bound gives the stated diamond estimate.
For $w=(1,\ii,0)/\sqrt2$, all $g_{ij}=0$, and the three pairwise determinant
radicands in \eqref{v25:eq:pure-effects} give
$s=(\sqrt{13}+\sqrt5)/6$. This proves the displayed exact value without
claiming that this preparation is globally optimal.
\end{proof}

\section{A retained purification repairs one acceptance, but the second has an exact loss}
\label{v25:sec:purifier}

It is important to allow genuinely retained source memory before drawing a
strong conclusion. Give the terminal reader an untouched finite memory
$\mathcal M$ and a pure preparation $|\xi\rangle\in\C^3\otimes\mathcal M$,
with external marginal $\tau$. Encode the logical qubit by
$V_\xi|v\rangle=|\xi\rangle\otimes|v\rangle$ and execute exactly the
original $K_e$ on the triplet and logical factors, with the identity on
$\mathcal M$. This does not grant the reader a coherent copy of the measured
outcome register. It grants only a purifier of the independent preparation.

For a word $w$, its logical effect is
\begin{equation}
 F_w(\tau)=V_\xi^*(K_w^*K_w\otimes I_{\mathcal M})V_\xi
       =\operatorname{Tr}_{\C^3}[(\tau\otimes I_2)K_w^*K_w].
 \label{v25:eq:purified-effect}
\end{equation}
The last expression is defined by its matrix elements, so its positivity is
also immediate from the first. Effects depend only on $\tau$, not on the
choice of purification. The full source output and all preparation memory
are available to the optimizing terminal reader.

\begin{lemma}[The symmetric external marginal is optimal when its purifier is retained]
\label[lemma]{v25:lem:purified-optimum}
For every fixed word length $n$, let
$s_n(\tau)=\sum_{|w|=n}\sqrt{\det F_w(\tau)}$.
Among all independent external preparations and all their retained finite
purifications, the largest optimal entanglement fidelity is achieved at
$\tau=I_3/3$, and the smallest optimal diamond recovery error is
$1-s_n(I_3/3)$.
The conclusion also bounds arbitrary mixed preparation memory: allowing its
purification cannot make recovery worse.
\end{lemma}
\begin{proof}
The function $A\mapsto\sqrt{\det A}$ is concave on positive $2\times2$
matrices. One elementary proof is the superadditivity
$\sqrt{\det(A+B)}\ge\sqrt{\det A}+\sqrt{\det B}$:
for invertible $A$, diagonalize $A^{-1/2}BA^{-1/2}$ and use
$\lambda_1+\lambda_2\ge2\sqrt{\lambda_1\lambda_2}$; approximation handles
singular $A$. Positive homogeneity then gives concavity. The effects in
\eqref{v25:eq:purified-effect} are linear in $\tau$, so $s_n$ is concave.

The inherited external $S_4$ covariance of the actual coefficients says
$K_{g e}=(u_g\otimes I_2)K_e(u_g^*\otimes I_2)$.
It permutes entire original words and gives
$F_{g w}(u_g\tau u_g^*)=F_w(\tau)$. Hence $s_n$ is invariant under the
external triplet action. That representation is irreducible, so its finite
group average is $I_3/3$. Concavity yields
$s_n(I_3/3)\ge |S_4|^{-1}\sum_gs_n(u_g\tau u_g^*)=s_n(\tau)$.
The polar lemma proves both recovery assertions. Every mixed preparation
memory admits a retained purification with the same marginal $\tau$;
tracing the additional purifier is CPTP. The optimum with that additional
access is an upper bound for the original experiment. By
\cref{v25:lem:section}, a CPTP encoding that leaves the unknown logical
marginal unchanged can contain only such independent preparation memory,
not an unrecorded logical bypass.
\end{proof}

\begin{theorem}[Sharp two-acceptance recovery obstruction with arbitrary retained preparation memory]
\label[theorem]{v25:thm:two}
Retain the original coefficients \eqref{v25:eq:source}, with no intervening
operation except those coefficients. With arbitrary independent preparation
memory and arbitrary terminal CPTP recovery, one acceptance can be recovered
perfectly. After two acceptances the \emph{best possible} diamond error,
also optimized over that preparation and memory, is
\begin{equation}
 \boxed{\quad
 \delta_2=1-\frac{\sqrt{57}+\sqrt{1209}+2\sqrt{957}}{108}
          =0.0352653247\ldots>\frac1{30}.
 \quad}
 \label{v25:eq:two-defect}
\end{equation}
The optimal entanglement fidelity is $1-\delta_2/2$.
The following table gives the exact logical effects at the optimal external
marginal $I_3/3$; $e$ is the first and $f$ the second original label:
\begin{center}
\begin{tabular}{@{}lcc@{}}
\toprule
Relation of $(e,f)$ & Number of ordered pairs & $F_{ef}(I_3/3)$\\
\midrule
$f=e$ & $6$ & $\operatorname{diag}(13/216,\,31/648)$\\[.35em]
$f=\bar e$ & $6$ & $\operatorname{diag}(1/216,\,19/648)$\\[.35em]
$f\ne e,\bar e$ & $24$ & $\operatorname{diag}(11/432,\,29/1296)$\\
\bottomrule
\end{tabular}
\end{center}
All one-label probabilities are $1/6$ independently of the logical input.
Nevertheless the two-label distributions of $|0\rangle$ and $|1\rangle$
have exact total variation $4/27$.
\end{theorem}
\begin{proof}
For one acceptance,
\[
 F_e(I_3/3)=\frac16(A^2+B^2+C^2)=I_2/6.
\]
A maximally entangled triplet--memory preparation realizes this marginal.
Every branch is therefore a scaled isometry of the unknown input, and
\cref{v25:lem:polar} gives exact first-step recovery. This is only a
terminal recovery statement; none of its inverses is executed between the
two acceptances considered next.

Here is an exact calculation of the table, rather than a fit to sampled
outcome frequencies. At $\tau=I_3/3$,
\[
 F_{ef}=\frac13\operatorname{Tr}_{\C^3}(K_eK_f^2K_e).
\]
For a repeated edge and its complement, respectively, the three orthogonal
external lines give
\[
 F_{ee}=\frac1{12}(A^4+B^4+C^4),\qquad
 F_{e\bar e}=\frac1{12}(AB^2A+BA^2B+C^4).
\]
Direct multiplication using \eqref{v25:eq:ABC} gives
\[
 A^4+B^4+C^4=\operatorname{diag}(13/18,31/54),\quad
 AB^2A+BA^2B+C^4=\operatorname{diag}(1/18,19/54),
\]
which proves the first two rows.
The external tetrahedral action is transitive on ordered adjacent edge pairs.
It has no logical action, so all the remaining $24$ effects are identical.
For fixed $e$, completeness gives
$\sum_fF_{ef}=F_e=I_2/6$. Thus the common adjacent effect is
$(I_2/6-F_{ee}-F_{e\bar e})/4$, proving the third row.
One can equivalently verify this row from the explicit root projectors;
no independence of the two recorded labels is assumed.

The three determinants are
\[
 \frac{403}{139968},\qquad\frac{19}{139968},\qquad
 \frac{319}{559872}.
\]
Their square roots, multiplied by $6,6,24$, sum to
$(\sqrt{1209}+\sqrt{57}+2\sqrt{957})/108$.
Apply \cref{v25:lem:polar,v25:lem:purified-optimum} for the exact optimal
error over all allowed preparations. The simple bounds
$\sqrt{57}<38/5$, $\sqrt{1209}<174/5$, and $\sqrt{957}<31$
give $s_2<29/30$, proving the strict rational lower bound.
All effects are diagonal; their polar-corrected logical channel is exactly
\[
 \rho\longmapsto(1-\delta_2/2)\rho+(\delta_2/2)Z\rho Z.
\]
This explicitly attains both recovery optima. Finally the two input
probability differences are $1/81$, $-2/81$, and $1/324$ in the three rows.
Half their absolute values, with the displayed multiplicities, sum to $4/27$.
\end{proof}

The theorem identifies a precise memory boundary. Retaining the purification
repairs the information lost by the \emph{first} averaged external mixture,
but the second unchanged acceptance writes input-dependent information into
the \emph{original classical labels}. The optimal logical channel then
retains the classical $Z$ distinction and attenuates its conjugate coherences.
A procedure that recovers and re-prepares the triplet--memory code after each
acceptance changes the source between events. It is not the passive-history
reader analyzed here. The two-acceptance result does not forbid that different
feedback law; it shows why it must be supplied rather than inferred.

The table gives an experimentally smaller rejection test as well. At the
symmetric preparation, its classical affinity is exactly $s_2$, so
\cref{v25:prop:probability} reproduces the optimal obstruction from just the
$36$ original two-label probabilities for two specified inputs. The positive
construction of the maximizing recovery still uses terminal quantum access.
A negative probability test and a positive recovery construction have different
data requirements.

\section{No later passive history can undo the certified loss on the original clock}
\label{v25:sec:physical-time}

Let $\mathsf H_{T,*}$ be the full output experiment up to physical time $T$,
retaining waiting records, every original accepted label, the terminal source,
and any allowed untouched preparation memory. Assume the exact clock in
\cref{v24:eq:clock}, a fresh completed-boundary entrance, and no hidden quantum
changes between accepted events. This is the literal repeated benchmark;
nonaccepted reader-invisible but quantum-active updates would define another
full-history experiment and must be analyzed separately.
Let $\tau_k$ be the time of its $k$th acceptance. The inherited regenerative
waiting identity gives
\[
 \E\tau_k=k\epsilon,\qquad
 F_k(T):=\PP(\tau_k\le T)\ge(1-k\epsilon/T)_+.
\]

\begin{theorem}[A two-event loss persists at every positive physical time]
\label[theorem]{v25:thm:deadline}
For the original repeated source, every CPTP terminal reader of its full
history, and every prescribed unitary target $\mathcal U_T$ on the unknown
logical qubit, the best possible error with arbitrary allowed preparation
memory satisfies
\begin{equation}
 \boxed{\quad
 \inf\|\mathcal R_{T,*}\mathsf H_{T,*}-\mathcal U_T\|_\diamond
 \ge F_2(T)\delta_2
 \ge(1-2\epsilon/T)_+\delta_2.
 \quad}
 \label{v25:eq:deadline}
\end{equation}
For the standard unpurified section $I_3/3\otimes\rho$, the stronger first-event
bound is
\begin{equation}
 \inf\|\mathcal R_{T,*}\mathsf H_{T,*}-\mathcal U_T\|_\diamond
 \ge\frac23F_1(T)\ge\frac23(1-\epsilon/T)_+.
 \label{v25:eq:standard-deadline}
\end{equation}
The infimum in the first assertion includes the independent preparation;
no recovery is executed during the native evolution.
\end{theorem}
\begin{proof}
First condition on the complete clock and acceptance schedule. It is independent
of the unknown input and of the quantum outcomes in this scalar-acceptance
benchmark. If the schedule has at least $k$ acceptances, its full quantum/history
output is a CPTP postprocessing of the flagged output after the first $k$
original acceptances: execute the remaining original operations, retain their
flags, and apply the terminal reader. Thus its achievable entanglement fidelity
cannot exceed the first-$k$ optimum. Waiting records are independent input-free
side information and do not change that optimum.
For schedules with fewer than $k$ acceptances, fidelity is at most one.
Entanglement fidelity is affine in the overall channel, so
\[
 F_{\rm e}(\mathcal U_T^{-1}\mathcal R_{T,*}\mathsf H_{T,*})
 \le 1-F_k(T)\delta_k/2,
\]
where $\delta_2$ is the universal two-event optimum and $\delta_1=2/3$
is the standard-section optimum. The Bell test yields the stated diamond
bounds. Conditional scheduling of additional source operations after the
first $k$ events cannot improve their recovery figure either: it is still
CPTP postprocessing of that full recorded prefix. Finally the waiting-mean
bound is Markov's inequality applied to $\tau_k$.
No subtraction or renormalization of a failed clock history is used.
\end{proof}

\begin{corollary}[A pure-input and positive-error interpretation]
\label[corollary]{v25:cor:pure-input}
For the standard section, any such terminal procedure has an unknown pure
logical input whose expected overlap error relative to its prescribed unitary
target is at least $2F_1(T)/9$. For that input the attained conditional error
$X\in[0,1]$ satisfies
\begin{equation}
 \PP(X>1/9)\ge[(2F_1(T)-1)/8]_+.
 \label{v25:eq:probability-failure}
\end{equation}
For arbitrary independent preparation memory, the corresponding maximal
pure-input mean error is at least $F_2(T)\delta_2/3$.
These are uniform all-input obstructions; the witnessing input need not be
identical for different terminal readers or cutoffs.
\end{corollary}
\begin{proof}
For any qubit CPTP channel $\Lambda$, the uniform pure-state average overlap
is $(2F_{\rm e}(\Lambda)+1)/3$. To verify it directly, write a Kraus
representation and use
$\int |v\rangle\langle v|^{\otimes2}\dd v=(I+\mathrm{Swap})/6$;
this identity follows from the rotational second moments
$\int n_an_b\dd n=\delta_{ab}/3$ on the Bloch sphere.
Insert the entanglement-fidelity bounds in the previous proof. Some pure
input has error at least that average. For the standard-section input,
$\E X\le1/9+(8/9)\PP(X>1/9)$ gives
\eqref{v25:eq:probability-failure}. Any internal randomness of the terminal
reader can be included in the expectation; it is not assumed independent of
the original history.
\end{proof}

For an inert logical reference qubit, the fixed positive comparison
$\Delta(I-|\Omega_{U_T}\rangle\langle\Omega_{U_T}|)$ has expectation at
least $\Delta F_2(T)\delta_2/2$ after every allowed terminal procedure, and
at least $\Delta F_1(T)/3$ for the standard section. This is a precise
coherent-code energy test, not the assertion that this Bell-projector cost
is an Einstein--SM action observable. It cannot be transferred to an unrelated
classical action merely by relabelling its variables.

\section{Outcome: the original short-word effects, not a more favorable unravelling}
\label{v25:sec:outcome}

The history-reader alternative is now quantitatively resolved for the supplied
benchmark under a generous information policy. With the companion's standard
unpurified external preparation, the optimal first-acceptance recovery error
is $2/3$ even when all original labels and the terminal triplet are available.
Changing the unpurified section cannot make the first step exact. Granting
an arbitrary retained purification does permit perfect first-step recovery,
but its best possible two-step recovery still has the fixed deficit
\eqref{v25:eq:two-defect}. Finite-group symmetry proves that this latter value
is already optimized over the independent external preparation and memory.
The calibrated dense clock turns that two-event loss into an order-one
fixed-positive-time obstruction, rather than suppressing it by rarity.

The positive general tool is an exact passive-word calculation:
\[
 \boxed{\quad
 F_w=V^*K_w^*K_wV,\qquad
 \text{best qubit recovery error}=1-\sum_w\sqrt{\det F_w}.
 \quad}
\]
Its descendant identity \eqref{v25:eq:word-ledger} gives a nonnegative,
source-normalized accumulation test without changing an outcome or inserting
a reset. The original word probabilities already give a robust negative
certificate; the full word effects give the exact recovery optimum on an
efficient code. An internal algebra, a reduced channel decomposition, and
one-letter label independence do not supply these history effects for free.

\paragraph{What has not been claimed.}
This benchmark does not become the uniquely selected primitive source by
being fully specified. The result is about retention of an unknown coherent
multiplicity qubit, not about every classical low-energy reader. It does not
rule out the one-sided action-energy route of Version 24 on a smaller bank,
a growing source with a different normalization, or an independently supplied
feedback/quantum-memory law. It does rule out assigning perfect reversible
history dynamics to this unchanged repeated instrument on the strength of
its Pauli decomposition or its available first-step purification.
Recovering a code, where possible, still does not select its Hamiltonian:
a terminal unitary can be appended by fiat, so the separate same-action
chronological incidence must remain explicit.

The next native acquisition is therefore the intended physical code or
classical action bank together with its \emph{original finite-word effects},
its actually retained preparation memory, and its action generator. A
hypothetical ability to recover erased amplitudes is not a substitute for
those source data. No new nonzero-reference non-Abelian spatial estimate,
Einstein observed inverse, chiral continuation, or universal microscopic
selection theorem is claimed in this body.

\paragraph{Verification and append-only preservation.}
All bytes of Version 24 before the final end-document command are preserved.
The companion diagnostics reconstruct the original six coefficients in exact
arithmetic, verify the three two-letter effect classes and their probabilities,
check polar recovery and the corrected channels, test the preparation bounds,
and verify the passive-word determinant identity and clock probabilities.
The finite calculations are checks of the written arguments, not a
machine-checked formalization or an Einstein--SM simulation. New labels have
prefix \texttt{v25:}. The external general recovery result is attributed
rather than presented as a new theorem of the Renewal programme.

\begingroup
\renewcommand{\refname}{Additional references for Version 25}
\begin{thebibliography}{9}
\small\raggedright
\bibitem{v25:GW02} M.~Gregoratti and R.~F.~Werner,
\emph{Quantum Lost and Found}, Journal of Modern Optics \textbf{50} (2003),
915--933. \href{https://arxiv.org/abs/quant-ph/0209025}{arXiv:quant-ph/0209025}.
Background for correction with a specified classical environment record and
the scalar-effect criterion; no freedom to change the source's measured
instrument is assumed here.
\bibitem{v25:GW04} M.~Gregoratti and R.~F.~Werner,
\emph{On quantum error-correction by classical feedback in discrete time},
Journal of Mathematical Physics \textbf{45} (2004), 2600--2612.
\href{https://arxiv.org/abs/quant-ph/0403092}{arXiv:quant-ph/0403092};
\href{https://doi.org/10.1063/1.1758320}{doi:10.1063/1.1758320}.
Proposition~4 and the qubit multi-step analysis supply the established
polar-fidelity framework. The rectangular-code normalization, exact diamond
calculation, preparation optimization, and original six-edge/renewal-clock
evaluations used above are proved explicitly in this body.
\end{thebibliography}
\endgroup
\addtocontents{toc}{\protect\endgroup}
% END IMMUTABLE VERSION 25 BODY
% APPEND VERSION 26 HERE, BEFORE THE FINAL END-DOCUMENT COMMAND.


% BEGIN IMMUTABLE VERSION 26 BODY
\clearpage
\begin{center}
{\Large\bfseries Version 26: Classical action information before another recovery assumption}\\[.5em]
{\large A sharp three-acceptance bit loss, the actual record's observable span,
and a quantified independent-preparation alternative}
\end{center}
\makeatletter
\addtocontents{toc}{\protect\begingroup}
\addtocontents{toc}{\protect\makeatletter}
\addtocontents{toc}{\protect\def\protect\l@section{\protect\bfseries\protect\@dottedtocline{1}{0em}{2.5em}}}
\makeatother
\addcontentsline{toc}{part}{Version 26: Classical distinguishability and the original finite action-information bank}

\section{Audit: a preserved classical distinction is not a preserved qubit}
\label{v26:sec:audit}

Version 25 leaves a precise possibility rather than a general licence to claim
classical dynamics: its optimal two-acceptance corrected channel preserves the
classical $Z$ distinction although it loses conjugate coherence. It explicitly
retains smaller classical action banks outside its coherent-qubit obstruction.
We test that possibility on the \emph{next original acceptance}, not by choosing
a different decomposition of the averaged Pauli channel. We then determine
what the original classical labels can actually measure, and what an explicitly
larger preparation resource can achieve.

\paragraph{Inherited inputs and targeted overlap audit.}
The six original coefficients, their tetrahedral covariance and exact locking,
the physical partial trace, and the independent retained-purification policy
are \cref{v25:eq:ABC,v25:eq:source,v25:eq:purified-effect}. The two-event effects
and their remaining diagonal distinction are \cref{v25:thm:two}. The accepted
waiting mean and its original-clock interpretation are retained from
\cref{v25:thm:deadline}. We do not reprove the optimal coherent-qubit recovery
formula, source-algebra saturation, Pauli mixing, or the general one-sided
energy theorem. Exact-term checks of this cumulative source and the supplied
source-selection companion \cite{v23:SEL17} found related rank and three-letter
\emph{occurrence} tests, but not the three-event classical bit optimum or the
record-information bank below. This is a limited overlap audit, not a claim
of exhaustive verification or bibliographic priority.

\paragraph{Information policy.}
The input is an unknown classical sign encoded as either of two orthogonal
states of the same multiplicity qubit. Its basis may be chosen in advance.
The triplet and any preparation memory are independent of that sign. All
original source labels, the complete terminal source system, and any retained
purifier of the triplet are available to a terminal reader. No correction,
reset, or measurement is inserted between acceptances. No copy of the input
sign is stored before the first event. This last restriction is substantive:
an unknown \emph{classical} sign can be copied in its known basis, unlike an
unknown coherent qubit. A supplied classical archive would therefore evade
this experiment and must not be silently counted as independent preparation
memory. We do not exclude that different resource policy.

\paragraph{Attribution and scope.}
Binary trace-norm discrimination and correction of observable algebras are
established mathematics; see \cite{v23:Watrous,v26:BKK}. Their elementary
specializations are proved below to fix the original flags and normalization.
The new results are the exact depth-three source calculation, its optimization
over all independent preparation memory and all bit axes, and the positive
record-statistic and preparation-count estimates. This is not a universal
classical no-go, a new physical Hamiltonian, or an Einstein--SM continuum
construction. The source remains the supplied six-dimensional triplet--qubit
benchmark, not the fourteen-dimensional witness of Version 23.

\section{The exact classical bit norm on the original effects}
\label{v26:sec:bit}

Let $V:\C^2\to\mathcal H$ be an isometry and retain an efficient flagged
experiment
\[
 L_w=K_wV,\qquad F_w=L_w^*L_w,\qquad \sum_w F_w=I_2,
 \qquad \mathcal T_*(\rho)=\bigoplus_wL_w\rho L_w^*.
\]
The words are the actual passive source products. For a unit vector
$n\in S^2$, let $P_{\pm n}=(I\pm n\cdot\sigma)/2$ and write
\[
 F_w=a_wI+b_w\cdot\sigma,\qquad a_w\ge|b_w|,\qquad
 \sum_wa_w=1,\quad\sum_wb_w=0.
\]
The symbols $n$ and $a_w$ here are not a number of sites or clock rates.

\begin{lemma}[Exact terminal bit discrimination with complete flags]
\label[lemma]{v26:lem:bit}
With equal prior probabilities on the two signs, the minimum error of any
terminal quantum measurement, including arbitrary preceding CPTP processing,
is
\begin{equation}
 p_{\rm bit}(n)=\frac{1-D(n)}2,\qquad
 D(n)=\sum_w\sqrt{a_w^2-|b_w\times n|^2}.
 \label{v26:eq:bit}
\end{equation}
The bit can be read perfectly if and only if $[F_w,n\cdot\sigma]=0$ for every
word. More generally, for an isometric $d$-dimensional input, a prescribed
classical basis is perfectly recoverable precisely when every $F_w$ is
diagonal in that basis. The reader need not preserve superpositions.
\end{lemma}
\begin{proof}
For any two density matrices $\rho_+,\rho_-$ with equal priors, a binary
measurement $0\preceq M\preceq I$ has success probability
$1/2+\tr M(\rho_+-\rho_-)/2$. Diagonalization of the traceless Hermitian
difference shows that its positive spectral projection attains
$1/2+\|\rho_+-\rho_-\|_1/4$. Every terminal CPTP map followed by a measurement
is already a measurement on its input through the adjoint map; it cannot
improve this optimum.

In block $w$ the two unnormalized outputs are $|u_w\rangle\langle u_w|$ and
$|v_w\rangle\langle v_w|$, where $u_w=L_w|+n\rangle$ and
$v_w=L_w|-n\rangle$. On their span their difference has trace norm
\[
 \sqrt{(\|u_w\|^2+\|v_w\|^2)^2-4|\langle u_w,v_w\rangle|^2}.
\]
This follows by solving the quadratic characteristic equation; it also holds
when their span has dimension one. The two quantities in the radicand are
$2a_w$ and $|\langle+n|F_w|-n\rangle|^2=|b_w\times n|^2$.
Trace norm adds over the retained orthogonal flags, proving
\eqref{v26:eq:bit}. Each summand is at most $a_w$. Their sum equals one
exactly when every off-diagonal matrix element vanishes.

For a general classical basis $|i\rangle$, perfect discrimination requires
orthogonal supports of the outputs of different basis states. Within each
flag these supports are the lines spanned by $L_w|i\rangle$. Their pairwise
inner products are $\langle i|F_w|j\rangle$. Vanishing of all off-diagonal
entries is therefore necessary and sufficient: project onto the mutually
orthogonal lines in each flag, and assign the unused complement arbitrarily.
This constructs a classical recovery on every basis mixture, not a coherent
quantum recovery.
\end{proof}

\begin{lemma}[Independent preparation memory is optimized by the symmetric marginal]
\label[lemma]{v26:lem:preparation}
Use the original six source coefficients of \cref{v25:eq:source}. Let an
independent pure triplet--memory preparation have triplet marginal $\tau$,
and retain its entire purifier. For words of any fixed length $k$, define
\[
 F_w(\tau)=\operatorname{Tr}_{\C^3}[(\tau\otimes I)K_w^*K_w].
\]
For every \emph{fixed} bit axis $n$, the discrimination quantity
$D_k(n;\tau)$ of \eqref{v26:eq:bit} is maximized at $\tau=I_3/3$.
Consequently optimization over all independent preparations, all finite
preparation memories, and all bit axes reduces to
$\max_{n\in S^2}D_k(n;I_3/3)$. An arbitrary mixed preparation memory cannot
improve that maximum.
\end{lemma}
\begin{proof}
For a fixed basis, the function
\[
 F\longmapsto\sqrt{a(F)^2-|b(F)\times n|^2}
\]
is concave on positive $2\times2$ matrices. To verify this without a
purification-dependent formula, put $u=b\times n$ and use
\[
 \sqrt{a^2-|u|^2}
 =\inf_{|v|<1}\frac{a-v\cdot u}{\sqrt{1-|v|^2}}\quad(a\ge|u|).
\]
The reverse Lorentzian Cauchy--Schwarz inequality proves the lower bound for
each $v$, and $v=u/a$ gives equality in the interior; continuity gives the
boundary cases. The infimum of these linear functions is concave. Each
$F_w(\tau)$ is linear in $\tau$, so their sum is concave as well.

The original external tetrahedral action permutes entire words and has no
logical-qubit action. Thus $D_k(n;u_g\tau u_g^*)=D_k(n;\tau)$, with $n$
unchanged. Its triplet representation is irreducible and its average sends
$\tau$ to $I_3/3$. Concavity gives the claimed maximum. A maximally entangled
triplet--memory vector realizes that marginal, so the bound is attained
within the allowed resource class. Every mixed independent preparation can
be purified into larger untouched memory; discarding that purifier is a
terminal CPTP map and cannot improve discrimination. No averaging gate or
reset is executed on the actual source.
\end{proof}

\begin{proposition}[Passive classical loss is monotone and has a finite quadratic test]
\label[proposition]{v26:prop:ledger}
For any fixed input encoding and axis, $D_{k+1}(n)\le D_k(n)$ under an
additional unchanged, fully recorded source operation. For the effects at
one depth put
\[
 M_k=\sum_{w:a_w>0}\frac{b_wb_w^{\mathsf T}}{a_w}.
\]
Then
\begin{equation}
 p_{\rm bit}(n)\ge\frac14 n^{\mathsf T}
       \{(\tr M_k)I-M_k\}n.
 \label{v26:eq:quadratic-test}
\end{equation}
In particular two nonparallel nonzero effect Bloch vectors exclude perfect
recovery of every bit basis. This is a \emph{classical} discrimination test,
not the determinant sum for coherent-qubit recovery in Version 25.
\end{proposition}
\begin{proof}
For each original flag, the next trace-preserving flagged instrument is a
CPTP map applied to its Hermitian signed output. Such a map contracts trace
norm on Hermitian matrices: write the positive and negative parts, use
positivity and trace preservation, and apply the triangle inequality.
Summing proves monotonicity and gives an exact nonnegative loss at each
parent flag. For the second assertion,
\[
 a_w-\sqrt{a_w^2-|b_w\times n|^2}
 =\frac{|b_w\times n|^2}{a_w+\sqrt{a_w^2-|b_w\times n|^2}}
 \ge\frac{|b_w\times n|^2}{2a_w}.
\]
Sum and divide by two. The kernel of the right-hand matrix consists of the
vectors parallel to every nonzero $b_w$, proving the last assertion.
\end{proof}

\section{A complete three-acceptance table and a sharp bit obstruction}
\label{v26:sec:three}

Number the original labels in the triplet frame as
\[
 1=(23,+),\ 2=(23,-),\ 3=(13,+),\ 4=(13,-),\
 5=(12,+),\ 6=(12,-).
\]
A word $(e,f,g)$ means $K_gK_fK_e$, in chronological order. For the symmetric
external marginal with retained purification, set
\[
 F_{efg}=\frac13\operatorname{Tr}_{\C^3}
              (K_eK_fK_g^2K_fK_e).
\]
The original matrices are real Hermitian, but their products need not commute.

\begin{theorem}[Exact original third-event effects]
\label[theorem]{v26:thm:table}
The $216$ effects split into the twelve external tetrahedral word orbits in
\cref{v26:tab:effects}. Each row lists
$93312F_{efg}=\left(\begin{smallmatrix}u&t\\t&v\end{smallmatrix}\right)$.
There is no change of the original outcome basis. The effects sum to $I_2$.
\begin{table}[htbp]
\centering\small
\begin{tabular}{@{}ccrrr@{}}
\toprule
Representative $(e,f,g)$ & Orbit size & $u$ & $t$ & $v$\\
\midrule
$(1,1,1)$ & $6$  & $2076$ & $0$ & $1420$\\
$(1,1,2)$ & $6$  & $348$ & $0$ & $460$\\
$(1,1,3)$ & $24$ & $798$ & $0$ & $646$\\
$(1,2,1)$ & $6$  & $60$ & $0$ & $364$\\
$(1,2,2)$ & $6$  & $60$ & $0$ & $748$\\
$(1,2,3)$ & $24$ & $78$ & $0$ & $406$\\
$(1,3,1)$ & $24$ & $573$ & $0$ & $553$\\
$(1,3,2)$ & $24$ & $141$ & $0$ & $409$\\
$(1,3,3)$ & $24$ & $726$ & $0$ & $718$\\
$(1,3,4)$ & $24$ & $294$ & $0$ & $190$\\
$(1,3,5)$ & $24$ & $321$ & $72$ & $109$\\
$(1,3,6)$ & $24$ & $321$ & $-72$ & $109$\\
\bottomrule
\end{tabular}
\caption{Three original acceptances; common denominator $93312$.}
\label{v26:tab:effects}
\end{table}
\end{theorem}
\begin{proof}
Here is a rational multiplication procedure fixing every entry. On the pair
$ij$ and its missing coordinate $k$, let
$S_{ij}=|i\rangle\langle i|+|j\rangle\langle j|$,
$T_{ij}=|i\rangle\langle j|+|j\rangle\langle i|$, and
$Q_k=|k\rangle\langle k|$. The integer matrices
\begin{equation}
 N_{ij,\pm}=6\sqrt2K_{ij,\pm}
 =S_{ij}\otimes(2I+Z)\ \pm3T_{ij}\otimes X
                         +2Q_k\otimes(I-Z)
 \label{v26:eq:integer-source}
\end{equation}
have only integer entries. If $N_w=N_gN_fN_e$, then
\begin{equation}
 F_w=\frac{\operatorname{Tr}_{\C^3}(N_w^{\mathsf T}N_w)}{1119744}.
 \label{v26:eq:integer-effects}
\end{equation}
Thus the entries in the table are the entries of that partial trace divided
by $12$. Multiplication for its twelve representatives gives the displayed
integers; for example the last two representatives give respectively
$12\left(\begin{smallmatrix}321&\pm72\\\pm72&109\end{smallmatrix}\right)$.
For the first representative the partial trace is
$12\operatorname{diag}(2076,1420)$. These computations use the fixed matrices
\eqref{v26:eq:integer-source}, not reconstructed outcome frequencies.

For completeness the orbit reduction is finite and explicit. The $24$ signed
coordinate-permutation matrices whose coordinate signs have product $+1$
permute the six root lines. Conjugation carries $K_e$ to the coefficient with
the permuted label and leaves the logical factor unchanged. It therefore
leaves the partial-trace effect of each permuted word unchanged. Fixing its
first label as $1$, its second is $1$, $2$, or one of $3,4,5,6$; the stabilizer
is transitive on the latter four. If the second is $1$ or $2$, the third has
the same three alternatives. If it is $3$, its third has the six separately
listed alternatives. Their orbit sizes are those in the table and sum to
$216$, giving all words with no assumed temporal independence. Finally the
weighted sum of the integer matrices is $93312 I_2$, also following from
successive source completeness. This proves the assertion.
\end{proof}

\begin{theorem}[Sharp classical bit loss after three unchanged acceptances]
\label[theorem]{v26:thm:sharp}
Optimize jointly over every bit basis in the multiplicity qubit, every
independent external preparation and retained finite purification, and every
terminal CPTP reader and binary measurement. With no intervening reset or
correction, the minimum equal-prior bit error after three original acceptances
is exactly
\begin{equation}
 \boxed{\quad p_3^{\rm cl}
   =\frac{215-\sqrt{41041}}{3888}
   =0.0031929567\ldots>\frac1{320}.\quad}
 \label{v26:eq:sharp}
\end{equation}
The symmetric triplet marginal with its retained purification, the $Z$ input
basis, and the flagwise optimal binary measurement attain it. This is a
classical decision-error probability, not the coherent diamond error
$\delta_2$ of \cref{v25:eq:two-defect}.
\end{theorem}
\begin{proof}
By \cref{v26:lem:preparation} it suffices to use the symmetric marginal and
optimize the axis $n$. The ten diagonal classes have $b_w$ parallel to $Z$,
so each summand of \eqref{v26:eq:bit} increases with $n_z^2$.
For the two other classes put
\[
 a=\frac{215}{93312},\qquad x=\frac{72}{93312},\qquad
 z=\frac{106}{93312},\qquad b_\pm=(\pm x,0,z).
\]
Their paired contributions, before their multiplicity $24$, satisfy
\begin{align*}
 &\sqrt{a^2-x^2-z^2+(x n_x+z n_z)^2}
 +\sqrt{a^2-x^2-z^2+(-x n_x+z n_z)^2}\\
 &\qquad\le2\sqrt{a^2-x^2-z^2+x^2n_x^2+z^2n_z^2}
 \le2\sqrt{a^2-z^2+(z^2-x^2)n_z^2}.
\end{align*}
The first inequality is concavity of the square root, and the second uses
$n_x^2\le1-n_z^2$. Since $z^2>x^2$, this bound is maximized at $n_z^2=1$,
where both inequalities are equalities. The diagonal contributions are
simultaneously maximized there. Thus $Z$ is an optimal axis; the presence
of nonzero diagonal contrasts also makes it the only axis up to sign.

At this axis all diagonal words have zero decision loss. The $48$ tilted
words have $\sqrt{a^2-x^2}=\sqrt{41041}/93312$, so
\[
 D_3^{\max}=1-\frac{48(215-\sqrt{41041})}{93312}
           =1-\frac{215-\sqrt{41041}}{1944}.
\]
Substitution into \eqref{v26:eq:bit} proves the exact value and constructs
an attaining flagwise spectral measurement. Its preparation is allowed, so
the bound is not merely a relaxation. To check the displayed strict rational
bound, $\sqrt{41041}<4057/20$ because $41041<(4057/20)^2$; hence
$215-\sqrt{41041}>243/20=3888/320$.
\end{proof}

\begin{corollary}[The surviving classical algebra disappears at depth three]
\label[corollary]{v26:cor:algebra}
For the symmetric purified preparation, every bit basis is perfectly readable
after one acceptance; after two acceptances exactly the $Z$ basis is perfectly
readable; after three acceptances no nonconstant classical qubit observable
is perfectly readable. The last assertion holds for all independent
preparation memories by \cref{v26:thm:sharp}. No later passive history restores
perfect readability.
\end{corollary}
\begin{proof}
The inherited first effects are $I/6$ and the inherited second effects are
diagonal with at least one nonscalar diagonal. Apply \cref{v26:lem:bit}.
At depth three, the effect for $(1,1,1)$ is nonscalar diagonal and that for
$(1,3,5)$ has a nonzero off-diagonal entry. Their common Hermitian commutant
is scalar. The sharp positive error also proves the statement for every
other independent preparation. Passive monotonicity proves the final claim.
This does not identify physical fields with arbitrary qubit observables;
it gives the exact observable boundary of this encoding experiment.
\end{proof}

\section{The full record algebra and the measurable linear span are different}
\label{v26:sec:record}

Now retain only the original classical word labels and waiting records, with
the symmetric triplet marginal. A retained triplet purifier does not change
these marginal probabilities. For an input
\[
 \rho(u)=\frac12(I+xX+yY+zZ),\qquad x^2+y^2+z^2\le1,
\]
the probability of word $w$ is $\tr(\rho(u)F_w)$. The following result is
about this label-only readout, not the larger terminal quantum experiment
used for \cref{v26:thm:sharp}.

\begin{theorem}[All-depth label blindness and the minimal informative prefix]
\label[theorem]{v26:thm:span}
At every word length every effect $F_w(I_3/3)$ is real symmetric. Thus all
finite or infinite original label histories have identical laws for
$(x,y,z)$ and $(x,-y,z)$. On the same input qubit the real linear effect spans
at depths one, two, and every depth at least three are, respectively,
\begin{equation}
 \R I,\qquad \operatorname{span}_{\R}\{I,Z\},\qquad
 \operatorname{span}_{\R}\{I,X,Z\}.
 \label{v26:eq:span}
\end{equation}
At depth three the unital complex star algebra generated by the effects is
nevertheless $M_2(\C)$. That algebra equality does not make the classical
record informationally complete for the input qubit.
\end{theorem}
\begin{proof}
Every $K_e$ is real, so $K_w^{\mathsf T}K_w$ and its symmetric partial trace
are real symmetric. Their trace pairing with the imaginary antisymmetric
Pauli $Y$ is zero. This proves equality of each finite cylinder probability.
The timing law is input-independent, so adding it changes neither equality.
Cylinder sets generate the sigma field of infinite histories; the equal
finite-dimensional laws consequently give identical infinite laws as well.
One may also see this by a monotone-class argument.

The one- and two-event spans follow from the inherited exact effects. The
three-event table supplies both a nonscalar $Z$ component and the two signs
of a nonzero $X$ component; together with completeness it spans $I,X,Z$.
For every prefix $w$, $F_w=\sum_eF_{we}$. The spans are therefore increasing
with depth, while reality confines all of them to the three-dimensional real
symmetric space. This proves \eqref{v26:eq:span}. The complex algebra
containing $X,Z$ also contains $iXZ=Y$, and hence is $M_2(\C)$. But products
and imaginary combinations of effect matrices are not automatically effects
of recorded events. The linear probability span, not its generated algebra,
determines this classical observation experiment.
\end{proof}

\begin{proposition}[Which linear action insertions can be estimated from the labels]
\label[proposition]{v26:prop:unbiased}
At a fixed prefix length, an input-independent real statistic $t(w)$ has
expectation $\tr(\rho A)$ for every input density matrix precisely when
\begin{equation}
 A=\sum_w t(w)F_w.
 \label{v26:eq:unbiased}
\end{equation}
Consequently the three-event labels can estimate every Hermitian insertion
in $\operatorname{span}_{\R}\{I,X,Z\}$ without bias, but no such statistic,
even using an arbitrary passive history or any number of identical
independent histories, can estimate the $Y$ expectation for all inputs.
For the two inputs with $y=\pm1$, every label-only binary decision has error
$1/2$ and every real estimator $\widehat y$ has averaged squared error at
least one.
\end{proposition}
\begin{proof}
Linearity of the Born probabilities gives the displayed operator condition;
density matrices separate Hermitian matrices. The span theorem gives the
finite-prefix assertion. Equality of the entire laws for $y=\pm1$ persists
under independent products and under any classical postprocessing. Thus a
binary decision has no advantage over a prior guess. The two risks have
average $\E[(\widehat y-1)^2+(\widehat y+1)^2]/2
=\E\widehat y^2+1\ge1$. No statement is made here about terminal quantum
measurements, other external preparations, or a newly supplied complex
instrument.
\end{proof}

\section{Physical-time and separating-action consequences of the three-event loss}
\label{v26:sec:time}

Retain the original renewal clock and the fresh completed-boundary entrance,
with $d=4\vartheta\epsilon/11$ and mean physical accepted interval $\epsilon$.
Let $\tau_k$ be the $k$th accepted time and
$F_k(T)=\PP(\tau_k\le T)$. The previously proved timing identity gives
$\E\tau_k=k\epsilon$, so $F_k(T)\ge(1-k\epsilon/T)_+$.
The full output below includes all waiting records, all original labels, the
terminal source system, and the allowed independent preparation memory.

\begin{theorem}[Classical loss persists under the original dense clock]
\label[theorem]{v26:thm:time}
For every fixed $T>0$, every bit basis, every allowed independent preparation,
and every terminal reader of the unchanged source, its equal-prior decision
error satisfies
\begin{equation}
 \boxed{\quad p_{
m err}(T)\ge F_3(T)p_3^{\rm cl}
        \ge(1-3\epsilon/T)_+p_3^{\rm cl}.\quad}
 \label{v26:eq:time}
\end{equation}
The bound also permits arbitrary passive outcome-dependent processing after
the first three original events. Its first lower bound uses the exact
calibrated waiting probability, not a Poisson replacement.
\end{theorem}
\begin{proof}
Condition on the input-independent complete clock schedule. For a schedule
with at least three acceptances, its later flagged source output is a CPTP
postprocessing of its first three flagged outputs. Trace-norm contraction
and \cref{v26:thm:sharp} give error at least $p_3^{\rm cl}$. Any reader
conditioned on that schedule is still among the readers bounded there.
Schedules with fewer than three acceptances contribute nonnegative error.
Averaging proves the first bound, and the waiting-mean inequality proves
the second. This argument charges incomplete histories instead of removing
or renormalizing them. The record retains the same unmodified source
operations before the third acceptance.
\end{proof}

\begin{corollary}[A quantitative test for a separating classical action bank]
\label[corollary]{v26:cor:action}
Suppose two independently specified classical initial data are encoded as the
two multiplicity inputs above. Their target classical trajectories give a
finite vector of declared action/readout values $a_+(T),a_-(T)\in\R^m$,
with $\Delta_T=|a_+(T)-a_-(T)|>0$. The target equations and bank are fixed
before choosing a terminal reader. For every terminal quantum-to-classical
estimate $\widehat a$ independent of the unknown sign,
\begin{align}
 \frac12\sum_{s=\pm}\PP_s\!\bigl(
      |\widehat a-a_s(T)|\ge\Delta_T/2\bigr)
       &\ge F_3(T)p_3^{\rm cl},\
 \frac12\sum_{s=\pm}\E_s|\widehat a-a_s(T)|^2
       &\ge\frac{\Delta_T^2}{4}F_3(T)p_3^{\rm cl}.
 \label{v26:eq:action-risk}
\end{align}
Thus this unchanged single-source encoding cannot supply vanishing
all-input error on a bank which keeps these two target values separated,
even though no coherent qubit recovery was requested.
\end{corollary}
\begin{proof}
From an estimate choose the nearer of the two target vectors, with a fixed
tie convention. On a wrong decision its distance from the true target is at
least $\Delta_T/2$, by the triangle inequality. This is an allowed terminal
binary measurement and decision, so \cref{v26:thm:time} bounds its averaged
error below. The event and squared-risk bounds follow. Independent internal
randomness of the reader may be included in the output measurement.
\end{proof}

\begin{corollary}[Positive action costs must pay either mean error or variance]
\label[corollary]{v26:cor:positive-cost}
Let a fixed terminal CPTP reader of the same history produce physical states
$\rho_s$, with $s=\pm$ the unknown input sign. Let $A_1,\ldots,A_m$ be
independently declared self-adjoint physical insertions, not necessarily
commuting. Suppose their reference vectors $a_s\in\R^m$ have separation
$\Delta=|a_+-a_-|>0$, and positive comparison operators satisfy
\[
 C_s\succeq c_0\sum_{\ell=1}^m(A_\ell-a_{s,\ell}I)^2,
 \qquad c_0>0.
\]
Then, with expectations over all original histories and reader outcomes,
\begin{equation}
 \frac12\sum_{s=\pm}\E_s\tr(\rho_s C_s)
 \ge\frac{c_0\Delta^2}{4}F_3(T)p_3^{\rm cl}.
 \label{v26:eq:positive-cost}
\end{equation}
This bound is on the complete positive comparison, not on the conditional
means alone. In particular an exact deterministic interpretation of the
conditional means does not remove the retained variance contribution.
\end{corollary}
\begin{proof}
Put $u=(a_+-a_-)/\Delta$, $A_u=\sum u_\ell A_\ell$, and
$a_{s,u}=u\cdot a_s$. For every vector $v$, Cauchy--Schwarz gives
\[
 \|(A_u-a_{s,u})v\|^2
 \le\sum_\ell\|(A_\ell-a_{s,\ell})v\|^2.
\]
Thus the stated comparison controls $(A_u-a_{s,u})^2$ even though the
individual $A_\ell$ need not commute. A spectral measurement of the single
operator $A_u$ gives a classical estimator whose two targets differ by
$\Delta$. Its mean squared error is exactly
$\E_s\tr\rho_s(A_u-a_{s,u})^2$. Apply
\cref{v26:cor:action}. Finally
\[
 \tr\rho(A-a)^2=(\tr\rho A-a)^2+
                          \tr\rho(A-\tr\rho A)^2
\]
separates mean error and conditional variance. The measurement is a
mathematical distinguishability test, not an extra operation executed by the
source. A general nonlinear functional of an unknown density matrix is not
a terminal CPTP measurement and is not covered by the earlier estimator
claim; the present positive-square statement is the applicable action test.
\end{proof}

The separation premise is essential. A constant observable, a bank assigning
the same values to both inputs, a supplied archive of the original classical
input, or a construction valid only for one known initial point is not
excluded. Nor does this corollary assert that an Einstein--SM action bank has
already been identified with these two qubit preparations. If a smooth
Hamiltonian flow on a regular finite-dimensional phase space is used as the
target, distinct initial points remain distinct at fixed times by uniqueness;
a bank must still be shown to separate their physical states. This elementary
observation is not a substitute for the missing field/action identification.

\section{A positive larger-source alternative with an exact preparation cost}
\label{v26:sec:replicas}

The obstruction is not a claim that the labels contain no classical
information. We now compute that information exactly, while declaring the
extra resource. Take $N$ \emph{independently prepared} copies of the same
source, each with initial state $I_3/3\otimes\rho(x,y,z)$. The triplet purifier
is unnecessary for label statistics. No unknown qubit is cloned and no
original single trajectory is reset: these are $N$ supplied preparations.
Their source copies and clocks are independent conditional on the common
classical preparation parameters. This resource has not been derived from
the primitive law or identified with a field population.

For one completed three-event prefix define $S=+1$ on the orbit represented
by $(1,3,5)$, $S=-1$ on the orbit represented by $(1,3,6)$, and $S=0$ on
all other original words. Let $B$ be the indicator that the second original
label is complementary to the first, regardless of the third.

\begin{proposition}[Two exact classical statistics of the original labels]
\label[proposition]{v26:prop:statistics}
For every qubit input,
\begin{align}
 \E S&=x/27,&\quad \E S^2&=(215+106z)/1944,\
 \E B&=(11-8z)/108.
 \label{v26:eq:statistics}
\end{align}
Hence the estimators
\begin{equation}
 \widehat x=\frac{27}{N}\sum_{j=1}^NS_j,
 \qquad
 \widehat z=\frac{11}{8}-\frac{27}{2N}\sum_{j=1}^NB_j
 \label{v26:eq:estimators}
\end{equation}
are globally unbiased, and
\begin{align}
 \Var(\widehat x)
   &=\frac1N\left[\frac38(215+106z)-x^2\right]
       \le\frac{963}{8N},\
 \Var(\widehat z)
   &=\frac{(11-8z)(97+8z)}{64N}
       \le\frac{1691}{64N}.
 \label{v26:eq:variances}
\end{align}
No iteration of the source after these prefixes is required by these
statistics; keeping their original classical records retains them.
\end{proposition}
\begin{proof}
Summing the last two table rows with their signs and multiplicities gives
\[
 \sum_wS(w)F_w=X/27,\qquad
 \sum_wS(w)^2F_w=(215I+106Z)/1944.
\]
Summing all third outcomes with a complementary first pair gives, from the
inherited two-event table,
\[
 \sum_wB(w)F_w=\operatorname{diag}(1/36,19/108)
                        =(11I-8Z)/108.
\]
Pair with $\rho$ and use independence of the $N$ preparations. The exact
variances follow by expanding the squares. For $-1\le z\le1$ the Bernoulli
probability lies in $[1/36,19/108]\subset[0,1/2]$, so its maximum variance
is at $z=-1$. This proves the displayed uniform bounds, including the
negative sign of that endpoint. There is no independence assumption between
successive labels of one source.
\end{proof}

\begin{theorem}[The $X$ preparation-noise coefficient is optimal for this record policy]
\label[theorem]{v26:thm:Fisher}
For the complete first three labels, at the maximally mixed input $x=y=z=0$,
the classical Fisher information for the $x$ parameter is exactly
\begin{equation}
 \mathcal I_x=\sum_w\frac{b_{w,x}^2}{a_w}=\frac8{645},
 \qquad \mathcal I_{xz}=0.
 \label{v26:eq:Fisher}
\end{equation}
For $N$ independent such records, every estimator locally unbiased for $x$
at that point has variance at least $645/(8N)$. The estimator $\widehat x$
in \eqref{v26:eq:estimators} attains this lower bound, and is unbiased on the
entire Bloch ball. In particular a root-mean-square accuracy $Ch$ at that
point requires $N\ge645/(8C^2h^2)$ for a locally unbiased reader of this
specified three-event record bank.
\end{theorem}
\begin{proof}
The probability is $p_w=a_w+b_{w,x}x+b_{w,z}z$ and is strictly positive at
the center. Only the $48$ tilted words have $b_{w,x}\ne0$, each with
$a_w=215/93312$ and $b_{w,x}=\pm72/93312$. Therefore
\[
 \mathcal I_x=48\frac{72^2}{93312\cdot215}=8/645.
\]
The two signs cancel the mixed Fisher entry. The score of one record is
$b_{w,x}/a_w$, with mean zero, and the $N$-record score is their sum, with
variance $N\mathcal I_x$. Local unbiasedness gives
$\E(\widehat x\,\mathrm{score})=1$ by differentiating the finite sum of
record probabilities. Cauchy--Schwarz gives
$1\le\Var(\widehat x)N\mathcal I_x$.
For the displayed estimator, the one-record coefficient $27S(w)$ is exactly
$(b_{w,x}/a_w)/\mathcal I_x$, so equality holds. This is a classical
record-information bound, not a quantum Fisher-information bound or an
optimality statement over additional terminal quantum measurements or longer
prefixes. The final preparation-count inequality follows directly.
\end{proof}

\begin{proposition}[Deadlines and incomplete clocks can be retained without postselection]
\label[proposition]{v26:prop:deadline-estimation}
At a fixed physical deadline $T$, put $q_k=F_k(T)$ for the known original
clock. Use $S^{(T)}=S$ if the first three acceptances are complete, and zero
otherwise. Use $B^{(T)}=B$ if the first two are complete, and zero otherwise.
Retain the incomplete flags as ordinary outcomes. If $q_2,q_3>0$, the
unbiased estimators are
\begin{equation}
 \widehat x_T=\frac{27}{Nq_3}\sum_jS_j^{(T)},\qquad
 \widehat z_T=\frac{11}{8}-\frac{27}{2Nq_2}\sum_jB_j^{(T)}.
 \label{v26:eq:deadline-estimators}
\end{equation}
They obey
\begin{align}
 \Var(\widehat x_T)&=\frac1N
       \left[\frac{3(215+106z)}{8q_3}-x^2\right]
                  \le\frac{963}{8Nq_3},\
 \Var(\widehat z_T)&\le\frac{513}{16Nq_2}.
 \label{v26:eq:deadline-variance}
\end{align}
Moreover for $u>0$,
\begin{align}
 \PP(|\widehat x_T-x|\ge u)&\le2e^{-Nq_3^2u^2/1458},\
 \PP(|\widehat z_T-z|\ge u)&\le2e^{-8Nq_2^2u^2/729}.
 \label{v26:eq:tails}
\end{align}
No information about the sign or state is contained in the independent
waiting flags.
\end{proposition}
\begin{proof}
The waiting event is independent of the source labels and input. Thus
$\E S^{(T)}=q_3x/27$,
$\E(S^{(T)})^2=q_3(215+106z)/1944$, and
$B^{(T)}$ is Bernoulli with probability $q_2p$, where
$p=(11-8z)/108\le19/108$.
The first variance formula follows. Directly the second variance is
\[
 \Var(\widehat z_T)=\frac{729}{4Nq_2}p(1-q_2p)
                  \le\frac{513}{16Nq_2}.
\]
This is a conservative bound valid for every $q_2\in(0,1]$; at $q_2=1$
the sharper bound in \eqref{v26:eq:variances} applies.

Here are the probability bounds. If $Y$ lies in an interval of length $L$,
then $\log\E e^{t(Y-\E Y)}\le t^2L^2/8$: the second derivative of the
log moment generating function is the variance under its tilted law, at most
$L^2/4$, and its value and first derivative at zero are zero. Multiply these
bounds for independent preparations and optimize the exponential Markov
bound. Apply it with $L=2$ to $S^{(T)}$, threshold $q_3u/27$, and with
$L=1$ to $B^{(T)}$, threshold $2q_2u/27$. This gives
\eqref{v26:eq:tails} without a union bound over individual clock completions.
Normalization uses the known predeclared clock probability, not a
state-dependent success selection. For fixed $T$ and $\epsilon\to0$,
$q_k\ge(1-k\epsilon/T)_+\to1$, so the preparation cost remains of order
$N^{-1}$ in mean square.
\end{proof}

The independent-replica result does not contradict the single-preparation
bit floor. Its input consists of repeated preparations from a classical
parameter source, not a CPTP cloning of one unknown logical qubit. It gives
initial-parameter statistics on two real observable directions; it supplies
neither a Hamiltonian time evolution of those parameters nor the absent
$Y$ information. In particular, applying an arbitrary desired classical flow
to the estimates afterwards would still require an independent physical
reader/action identification and would not prove that the benchmark itself
executes that flow.

\section{Outcome: what the next native acquisition must distinguish}
\label{v26:sec:outcome}

The two-event diagonal survival in Version 25 is not a stable classical
pointer of the repeated source. At the third acceptance the exact optimized
bit error is \eqref{v26:eq:sharp}, and its physical-time obstruction is
\eqref{v26:eq:time}. This already tests a commutative input family, but
explicitly excludes a prior classical input copy, growing preparation
resources, and banks that do not separate the tested inputs. The
positive-square consequence \eqref{v26:eq:positive-cost} retains conditional
variance as well as mean error.

The positive alternative is equally specific. Original labels estimate the
$X,Z$ preparation bank with a sharp $X$ sampling coefficient, provided the
model supplies independent repeated preparations. Observing one source longer
is not that resource. Its entire label history is blind to $Y$ at the
symmetric preparation even though the effects generate $M_2(\C)$.
Algebraic saturation, measurable linear span, classical readability, and
physical action incidence are distinct statements.

The next native acquisition should therefore be the intended action
insertions and reader, together with the actual initial archive and replica
policy and the source's normalized work on that bank. A single-input,
separating-bank claim under the present policy already has a finite rejection
certificate. A growing empirical preparation, archive, or different operation
must instead be evaluated on its own terms before applying Version 24.
Another longer-prefix recovery obstruction on this benchmark would not
identify the intended law. No Einstein inverse, moving-reference nonlinear
Yang--Mills continuum estimate, or universal selection theorem is added here.

\paragraph{Verification and preservation.}
Every Version 25 byte before its final document command is preserved.
The new script checks all $216$ effects in exact arithmetic, all twelve
word orbits, the sharp axis and preparation bounds against full terminal
outputs, the record span, Fisher coefficient, original clock, and sampling
estimates. Numerical optimization tests diagnose, not replace, the written
proofs. New labels have prefix \texttt{v26:}. These checks are not
proof-assistant formalizations or an Einstein--SM simulation.

\begingroup
\renewcommand{\refname}{Additional reference for Version 26}
\begin{thebibliography}{9}
\small\raggedright
\bibitem{v26:BKK} C.~B\'eny, A.~Kempf, and D.~W.~Kribs,
\emph{Quantum error correction of observables}, Physical Review A
\textbf{76} (2007), 042303.
\href{https://arxiv.org/abs/0705.1574}{arXiv:0705.1574};
\href{https://doi.org/10.1103/PhysRevA.76.042303}{doi:10.1103/PhysRevA.76.042303}.
Background for the distinction between preserving a classical observable
algebra and preserving an unknown coherent state. The exact source table,
preparation optimization, bit error, and sampling costs are derived directly
in this body. Binary trace-norm discrimination is also treated in
\cite{v23:Watrous}; its finite form is proved in \cref{v26:lem:bit}.
\end{thebibliography}
\endgroup
\addtocontents{toc}{\protect\endgroup}
% END IMMUTABLE VERSION 26 BODY
% APPEND VERSION 27 HERE, BEFORE THE FINAL END-DOCUMENT COMMAND.


% BEGIN IMMUTABLE VERSION 27 BODY
\clearpage
\begin{center}
{\Large\bfseries Version 27: The nonlinear preparation scale before another source test}\\[.5em]
{\large An exact Wilson uncertainty floor, tree-free physical packets,
and same-action heat regularization}
\end{center}
\makeatletter
\addtocontents{toc}{\protect\begingroup}
\addtocontents{toc}{\protect\makeatletter}
\addtocontents{toc}{\protect\def\protect\l@section{\protect\bfseries\protect\@dottedtocline{1}{0em}{2.5em}}}
\makeatother
\addcontentsline{toc}{part}{Version 27: Sharp nonlinear vacuum preparation scale and local resolved transport}

\section{Audit: a preparation question that the record obstructions do not answer}
\label{v27:sec:audit}

Version 26 settles the useful short-prefix classical retention test for its
specified benchmark. Extending that calculation to longer words would not
identify a physical field law. We instead return to a different upstream
question: at a declared physical quantum scale, can \emph{any} exactly
Gauss-constrained state have the unsubtracted action energy and connection
comparison required by the nonlinear vacuum branch? A preparation theorem
and an obstruction valid for every preparation are the appropriate tests.

\paragraph{Inherited inputs, and why this is not the quadratic calculation again.}
The action, vacuum-orbit cost, weak-gradient estimate and real-time comparison
are \cref{v22:eq:energy,v22:eq:cost,v22:lem:distance,v22:thm:quantum}.
The tensor-tail, Ritz and local-instrument estimates are
\cref{v22:prop:cutoff,v14:thm:Galerkin,v22:prop:local,v14:thm:trajectory}.
Their proofs and the original renewal clock are not reconstructed here.
Version 21 proves an $\eta h^{-4}$ uncertainty scale for the \emph{quadratic}
Wilson Hessian on the linearized Gauss quotient. Version 22 instead constructs
states on the full nonlinear physical space, with the larger sufficient
preparation cost $\eta h^{-5}$. Its tree packet pays the sum of fundamental
cycle lengths. Neither result proves the two-sided nonlinear preparation
scale established below.

The companion exact-action Einstein bridge \cite{v23:EB16}, especially
\nolinkurl{v16:lift-theorem} and \nolinkurl{v16:actual-domain}, still does not
bound the full observed inverse after complementary relaxation. Its positive
compression cannot replace that inverse. No new Einstein inverse estimate is
inferred from the present preparation theorem. The source-selection results
\cite{v23:SEL17} also do not specify a physical spatial quantum scaling.
The completed algebra and history calculations therefore remain separate
from deciding whether a proposed normalization can even support a small
\emph{unsubtracted} physical action cost.

\paragraph{New results and exact scope.}
An electric--plaquette commutator gives a nonlinear operator lower bound at
arbitrary link angles. A geometric tube-volume estimate constructs a
\emph{tree-free} invariant packet with the matching power of $h$. No regular
orbit slice or small-angle expansion is used. A fixed imaginary-time
regularization by the same positive energy supplies the higher graph moment
needed for the finite regulator while preserving the preparation scale.
This gives a necessary and sufficient scale for existence of nonlinear
vacuum-concentrated physical preparations and a less restrictive explicit
local resolved branch. The constants in the lower and upper bounds are not
claimed optimal.

The target is still the zero-electric, trivial-holonomy vacuum. The packet
and its regularization are declared correlated preparation data, not an
identified primitive selection rule or a local preparation circuit. The
Hamiltonian, subsequent local Kraus operators and clock are unchanged.
No Yang--Mills mass gap, nonzero-reference limit or Einstein--SM emergence
claim follows from an absolute energy floor.

\paragraph{Attribution and audit limits.}
Hamiltonian lattice gauge theory and gauge-invariant coherent preparations
are established subjects \cite{v27:KS,v27:Chin,v27:BT}. The new estimates are
proved directly in the normalization of this lineage. Targeted checks of
the cumulative notes and the supplied companion endpoints did not locate
the nonlinear commutator floor, dimension-tracked orbit packet, or weighted
heat estimate used here. This is not an exhaustive novelty or independent
proof audit of all companion manuscripts. In particular, no claim is made
that a variational ansatz or imaginary-time preparation is itself new.

\section{An uncertainty floor for the complete nonlinear Wilson Hamiltonian}
\label{v27:sec:floor}

Retain a periodic cubic lattice with $N\ge3$, $h=N^{-1}$ and $n=h^{-3}$.
There are $3n$ positive edges and $3n$ elementary oriented plaquettes, indexed
by $(x,i,j)$ with $i<j$. On $\mathcal Q_h=SU(2)^{\mathcal E_h}$, use normalized
product Haar measure, the metric $g_h$ and generators $T_a=-\ii\sigma_a/2$
from Version 22. Write $\eta\in(0,1]$ for its physical quantum parameter and
\begin{equation}
 T=\frac1{2h}\sum_{e,a}L_{e,a}^2,\qquad
 W=\frac4h\sum_f(1-s_f),\qquad E=T+W,
 \quad s_f=\tfrac12\tr U_f,\quad L_{e,a}=-\ii\eta X_{e,a}.
 \label{v27:eq:action}
\end{equation}
The trace $s_f$ is real. All operators in this section precede compression.
Choose the left multiplication derivative $X_{e,a}$ at the first edge
$e=(x,i)$ of $U_f=U_{x,i}U_{x+e_i,j}U_{x+e_j,i}^{-1}U_{x,j}^{-1}$.
The kinetic Casimir is unchanged by this convention. Decompose
\[
 U_f=s_fI+2\sum_a x_{f,a}T_a,\qquad s_f^2+\sum_a x_{f,a}^2=1.
\]
The functions $x_{f,a}$ are real smooth multiplication functions; they need
not individually preserve the physical subspace.

\begin{lemma}[A first-edge derivative of the full ordered plaquette]
\label[lemma]{v27:lem:plaquette}
For the stated first edge and each $a$,
\begin{equation}
 X_{e,a}x_{f,a}=s_f/2,\qquad
 \sum_a x_{f,a}^2=(1-s_f)(1+s_f)\le2(1-s_f).
 \label{v27:eq:plaquette}
\end{equation}
These identities require no restriction on any link or plaquette angle.
\end{lemma}
\begin{proof}
Left multiplication on the first link changes $U_f$ to $e^{tT_a}U_f$.
The multiplication rule
$T_aT_b=-\delta_{ab}I/4+\epsilon_{abc}T_c/2$ gives
$X_{e,a}x_{f,c}=s_f\delta_{ac}/2+\sum_b\epsilon_{abc}x_{f,b}/2$.
Set $c=a$. The second assertion follows from the exact unit-quaternion
identity and $s_f\le1$.
\end{proof}

\begin{theorem}[A full nonlinear, unrenormalized energy floor]
\label[theorem]{v27:thm:floor}
As quadratic forms on $H^1(\mathcal Q_h)$,
\begin{equation}
 \boxed{\quad
 4T+\left(2+\frac{3\eta}{4}\right)W
       \succeq9\eta h^{-4}I,\qquad
 E\succeq\frac94\eta h^{-4}I\quad(0<\eta\le1).
 \quad}
 \label{v27:eq:floor}
\end{equation}
The bounds hold on the entire kinematical Hilbert space, hence on the exact
Gauss subspace and for every finite-energy density matrix. They also hold
after any finite orthogonal compression, provided the original action
products are compressed after being formed.
\end{theorem}
\begin{proof}
On the smooth core, for each $(f,a)$ set
$A_{f,a}=L_{e(f),a}/\sqrt h$ and
$B_{f,a}=2x_{f,a}/\sqrt h$.
The multiplication commutator and \cref{v27:lem:plaquette} give
\[
 \ii[A_{f,a},B_{f,a}]=\eta s_f/h.
\]
Positivity of $(A_{f,a}-\ii B_{f,a})^*(A_{f,a}-\ii B_{f,a})$ proves
$A_{f,a}^2+B_{f,a}^2\succeq\eta s_f/h$.
Sum over the three colours and every plaquette. In the chosen orientation
an edge is first in at most two plaquettes, so
\[
 \sum_{f,a}A_{f,a}^2\preceq4T,\qquad
 \sum_{f,a}B_{f,a}^2
   =\frac4h\sum_f(1-s_f^2)\preceq2W.
\]
It follows that
\[
 4T+2W\succeq\frac{3\eta}{h}\sum_fs_f
  =9\eta h^{-4}I-\frac{3\eta}{4}W.
\]
This is the first bound. For $0<\eta\le1$ the coefficient of $W$ on its
left is at most four, giving the second. Smooth functions are a form core;
all multiplication functions and their first derivatives are bounded at
fixed $h$, so the inequalities pass to $H^1$. Trace against positive sums
of vector states to obtain the density-matrix assertion. Finally an operator
form lower bound survives compression by pairing it with a vector in the
compressed subspace. No false compressed canonical commutator is used.
\end{proof}

The lower bound is not a gap above the quantum ground state. It bounds the
\emph{absolute} energy in the fixed convention \eqref{v27:eq:action}, whose
classical vacuum energy is zero. Subtracting a ground-energy or other
parameter-dependent scalar changes this comparison and requires the
subtraction's derivatives in independently varied action sources. The
argument also does not identify $\eta$ with a physical species count.

\section{A dimension-tracked Gibbs packet around an arbitrary closed orbit}
\label{v27:sec:tubes}

Let $d_h$ be the intrinsic distance of $g_h$, let $\mathcal O_h$ be the
trivial-holonomy gauge orbit and let $c_h$ be the original cost
\cref{v22:eq:cost}. Denote the real dimension of the configuration manifold by
\[
 \mathfrak d_h=3|\mathcal E_h|=9h^{-3}.
\]
This is the ambient dimension, not a regular-orbit codimension.
For a single $SU(2)$ factor, an eigenangle $\theta\in[0,\pi]$ has distance
$2\sqrt h\,\theta$. The elementary cosine bounds therefore imply
\begin{equation}
 \ell\,d_h(U,\mathcal O_h)^2\le c_h(U)
       \le u\,d_h(U,\mathcal O_h)^2,
 \qquad \ell=\frac1{2\pi^2},\quad u=\frac18.
 \label{v27:eq:distance-comparison}
\end{equation}
These inequalities are obtained before minimizing over the same orbit and
then passed through the minimum. They do not require a unique minimizing
gauge or a smooth quotient.

\begin{lemma}[Ball and tube volumes on the actual product group]
\label[lemma]{v27:lem:volume}
On $\mathcal Q_h$, balls of equal radius have equal Haar volume and, for
$0<r\le R$,
\begin{equation}
 \mu(B(R))\le(R/r)^{\mathfrak d_h}\mu(B(r)).
 \label{v27:eq:ball}
\end{equation}
For any nonempty closed subset $S\subseteq\mathcal Q_h$, put
$F_S(r)=\mu\{d_h(U,S)<r\}$. For $\lambda\ge1$,
\begin{equation}
 F_S(\lambda r)\le[2(\lambda+1)]^{\mathfrak d_h}F_S(r).
 \label{v27:eq:tube}
\end{equation}
Neither assertion assumes that $S$ is a submanifold or has a fixed stabilizer.
\end{lemma}
\begin{proof}
Homogeneity follows by left translations of the product group. Away from the
measure-zero antipodal cut sets, exponential coordinates based at a point
are products of three-dimensional balls $|v_e|<2\pi\sqrt h$. In metric
orthonormal coordinates their volume Jacobian, up to a common normalization,
is
\[
 j(v)=\prod_e\left[
   \frac{\sin(|v_e|/(2\sqrt h))}{|v_e|/(2\sqrt h)}\right]^2.
\]
The function $\sin t/t$ is decreasing on $[0,\pi]$: its derivative has
numerator $t\cos t-\sin t$, with derivative $-t\sin t\le0$.
The inverse exponential image of $B(R)$ is the star-shaped set
$\sum_e|v_e|^2<R^2$ inside the product cut balls. Scaling it by $r/R$
stays in the inverse image of $B(r)$ and increases $j$. The change of
variables gives \eqref{v27:eq:ball}, even when $R$ exceeds an individual
factor's injectivity radius or the full manifold's diameter.

Take a finite maximal $r$-separated family of centres in $S$. Compactness
guarantees existence. Their open balls of radius $r/2$ are disjoint and
belong to the $r$-tube; their balls of radius $(\lambda+1)r$ cover the
$\lambda r$-tube. If their number is $m$, homogeneity and \eqref{v27:eq:ball}
give
\[
 F_S(\lambda r)\le m\mu(B((\lambda+1)r))
 \le[2(\lambda+1)]^{\mathfrak d_h}m\mu(B(r/2))
 \le[2(\lambda+1)]^{\mathfrak d_h}F_S(r).
\]
No estimate of $m$, orbit curvature or reach is needed.
\end{proof}

\begin{theorem}[Uniform orbit-packet moments without a logarithmic volume loss]
\label[theorem]{v27:thm:Gibbs}
For $a>0$, define
\begin{equation}
 Z_h(a)=\int e^{-c_h/a}\,\dd\mu,\qquad
 \phi_{h,a}=Z_h(a)^{-1/2}e^{-c_h/(2a)},\qquad
 A_0=\pi^2/2+\log7.
 \label{v27:eq:Gibbs}
\end{equation}
Then $\phi_{h,a}$ is an exactly gauge-invariant unit vector in the energy
form domain and
\begin{align}
 \frac{Z_h(2a)}{Z_h(a)}&\le e^{A_0\mathfrak d_h},\\
 \langle c_h\rangle_{\phi_{h,a}}
   &\le2A_0\mathfrak d_h a,\\
 \PP_{|\phi_{h,a}|^2\mu}
  \{c_h\ge2a(A_0\mathfrak d_h+t)\}&\le e^{-t}\quad(t\ge0).
 \label{v27:eq:Gibbs-bounds}
\end{align}
For each fixed integer $k\ge1$,
$\langle c_h^k\rangle_{\phi_{h,a}}\le C_k(a\mathfrak d_h)^k$.
All constants are independent of $h$, the number of links and the geometry
of singular orbit strata.
\end{theorem}
\begin{proof}
Write $D=\mathfrak d_h$, $F=F_{\mathcal O_h}$ and set
$r^2=2aD/\ell$. From \eqref{v27:eq:distance-comparison},
\[
 Z_h(a)\ge e^{-ur^2/a}F(r)=e^{-2uD/\ell}F(r).
\]
Split the numerator integral into the tube $d_h(U,\mathcal O_h)<r$ and
the annuli $kr\le d_h(U,\mathcal O_h)<(k+1)r$, $k\ge1$.
Equations \eqref{v27:eq:distance-comparison} and \eqref{v27:eq:tube} give
\[
 Z_h(2a)\le F(r)\left[1+\sum_{k\ge1}
      \{4(k+1)e^{-k^2}\}^{D}\right]\le7^D F(r).
\]
For the last inequality, $D$ is a positive integer and
$1+\sum b_k^D\le(1+\sum b_k)^D$. Also $8/e<4$ and, for $k\ge2$,
$k^2\ge2k$, so
\[
 1+\sum_{k\ge1}4(k+1)e^{-k^2}
 <1+4+4\sum_{k\ge2}(k+1)4^{-k}=55/9<7.
\]
This proves the ratio bound, since $2u/\ell=\pi^2/2$.
Under the normalized packet,
$\E e^{c_h/(2a)}=Z_h(2a)/Z_h(a)$. Jensen gives the mean bound;
exponential Markov gives the stated tail. Integrating the shifted
exponential tail gives, for example,
\[
 \E c_h^k\le2^{k-1}
   \left[(2aA_0D)^k+(2a)^k k!\right]\le C_k(aD)^k.
\]
At fixed $h$, the inherited weak-gradient bound makes $c_h$ bounded
Lipschitz, so its exponential is in $H^1$ by the weak chain rule. It is
invariant because $c_h$ is invariant. Thus it is a vector in the exact
Gauss subspace, not merely an invariant density matrix with charged support.
\end{proof}

The ratio estimate is essential. A bound using only the absolute partition
normalization would typically leave $D\log(1/a)$. The tube argument cancels
that normalization loss without asserting a smooth gauge slice. The packet
uses the complete orbit distance; no minimizing gauge is measured or applied.

\section{A tree-free preparation attains the nonlinear quantum scale}
\label{v27:sec:packet}

\begin{proposition}[Magnetic action is bounded by the same orbit cost]
\label[proposition]{v27:prop:magnetic}
Pointwise on the whole configuration manifold,
\begin{equation}
                     W_h\le64h^{-2}c_h.
 \label{v27:eq:magnetic-cost}
\end{equation}
\end{proposition}
\begin{proof}
For four unitary factors, telescoping and Cauchy--Schwarz give
$\|V_1V_2V_3V_4-I\|_{\mathrm{HS}}^2
 \le4\sum_{j=1}^4\|V_j-I\|_{\mathrm{HS}}^2$.
Consequently $v(U_f)\le4\sum_{e\in f}v(g_x^{-1}U_eg_y)$ for every
vertex gauge, also for inverse edge factors. Each edge occurs in four
plaquettes. Sum and multiply by $4/h$, then minimize over the same gauge
to obtain \eqref{v27:eq:magnetic-cost}. This is an upper bound and does
not infer connection alignment from magnetic energy alone.
\end{proof}

\begin{theorem}[Matching-power upper bound on the full nonlinear physical space]
\label[theorem]{v27:thm:packet-energy}
Let $\mathcal C=E+c_h$ be the unchanged positive comparison. For every $a>0$,
\begin{equation}
 \langle\mathcal C\rangle_{\phi_{h,a}}
 \le2A_0\mathfrak d_h a
       \left(1+\frac{64}{h^2}+\frac{\eta^2}{16a^2}\right).
 \label{v27:eq:packet-general}
\end{equation}
In particular $a=\eta h$ gives, for $0<\eta\le1$ and $h\le1$,
\begin{equation}
 \boxed{\quad
 \langle\mathcal C\rangle_{\phi_{h,\eta h}}
                       \le C_*\eta h^{-4},
 \quad C_*=18A_0(65+1/16).
 \quad}
 \label{v27:eq:packet-energy}
\end{equation}
No tree, candidate-angle restriction, orbit-slice Hessian, or mode deletion
occurs in this preparation estimate.
\end{theorem}
\begin{proof}
The inherited $|\nabla c_h|_{g_h}^2\le c_h/2$ and the weak chain rule give
\[
 \langle T\rangle_{\phi_{h,a}}
  =\frac{\eta^2}{8a^2}\E_{\phi}|\nabla c_h|_{g_h}^2
  \le\frac{\eta^2}{16a^2}\E_{\phi}c_h.
\]
Add \cref{v27:prop:magnetic} and the comparison cost itself, and use
\cref{v27:thm:Gibbs}. For $a=\eta h$ and $\mathfrak d_h=9h^{-3}$,
$1+64h^{-2}+\eta^2/(16a^2)
 \le(65+1/16)h^{-2}$. This proves the stated constant.
The estimate is of the full group-valued potential, not its Hessian.
\end{proof}

The packet is not asserted to minimize the action and the constant $C_*$
is deliberately coarse. Its improvement is the \emph{power} of the physical
mesh: the preparation now pays the ambient number of modes and their energy
scale, not the extra tree-path length in Version 22. The function $c_h$ is
globally defined but need not be smooth, so a higher graph moment cannot be
assumed just from \eqref{v27:eq:packet-energy}. The next result supplies it.

\section{Same-action heat regularization preserves the small comparison}
\label{v27:sec:heat}

\begin{lemma}[A weighted heat estimate using only one derivative of the cost]
\label[lemma]{v27:lem:heat}
For a unit vector $\phi$ in the energy form domain, set
$u_s=e^{-sE}\phi$, $e_0=\langle\phi,E\phi\rangle$,
$c_0=\langle\phi,c_h\phi\rangle$ and $n_s=\|u_s\|^2$. Then for $s\ge0$,
\begin{align}
 n_s&\ge e^{-2se_0},\\
 \langle u_s,c_hu_s\rangle&\le c_0+\eta^2s/8,\\
 \frac{\langle u_s,E u_s\rangle}{n_s}&\le e_0.
 \label{v27:eq:heat}
\end{align}
The last inequality holds whenever the normalized expression is defined,
which the first guarantees. The constants have no spatial dimension factor.
\end{lemma}
\begin{proof}
The first bound is Jensen's inequality for the spectral measure of the
nonnegative $E$ in $\phi$. At $s>0$ the heat-evolved vector lies in every
power domain of $E$. The form identity for bounded Lipschitz $c_h$ is
\begin{align*}
 \frac{\dd}{\dd s}\langle u_s,c_hu_s\rangle
  ={}&-\eta^2\int c_h|\nabla u_s|^2
       -\eta^2\operatorname{Re}\int\overline u_s
                                \nabla c_h\cdot\nabla u_s
       -2\int c_hW_h|u_s|^2.
\end{align*}
Completing a square in $\sqrt{c_h}\nabla u_s$ bounds the first two
terms above by
$\eta^2\int |\nabla c_h|^2|u_s|^2/(4c_h)
 \le\eta^2\|u_s\|^2/8\le\eta^2/8$.
The quotient is interpreted as zero where $c_h=0$; alternatively use
$c_h+\delta$ in this completion and let $\delta\downarrow0$.
The last term is nonpositive. Integrate, using strong convergence at zero
and boundedness of $c_h$ at fixed $h$, to prove the second assertion.
This proof uses no second derivative of $c_h$.

For the third assertion, write $\lambda$ for the spectral variable. The
mean of $\lambda$ under the weight $e^{-2s\lambda}$ cannot exceed its
unweighted mean, because
\[
 \frac12\iint(\lambda-\mu)
      (e^{-2s\lambda}-e^{-2s\mu})\,\dd\nu(\lambda)\dd\nu(\mu)\le0.
\]
Rearranging gives the result. Truncation of the spectral integrals followed
by monotone bounds justifies it for finite $e_0$.
\end{proof}

\begin{theorem}[A smooth exactly physical preparation with a uniform fourth energy moment]
\label[theorem]{v27:thm:smooth-packet}
For $\phi=\phi_{h,\eta h}$ define the declared initial vector
\begin{equation}
             \psi_{h,\eta}=\frac{e^{-E}\phi}{\|e^{-E}\phi\|}.
 \label{v27:eq:smooth-packet}
\end{equation}
If $C_*\eta h^{-4}\le1$, it is smooth, exactly gauge invariant, and obeys
\begin{equation}
 \boxed{\quad
 \langle\psi_{h,\eta},\mathcal C\psi_{h,\eta}\rangle
             \le C\eta h^{-4},\qquad
 \|(E+I)^2\psi_{h,\eta}\|\le4.
 \quad}
 \label{v27:eq:smooth-bounds}
\end{equation}
Here $C$ is independent of $h,\eta$ and the number of links.
\end{theorem}
\begin{proof}
By \cref{v27:thm:packet-energy}, $e_0+c_0\le C_*\eta h^{-4}\le1$.
Use \cref{v27:lem:heat} at $s=1$ to get
\[
 \langle\mathcal C\rangle_{\psi}
       \le e_0+e^{2e_0}(c_0+\eta^2/8)
       \le C\eta h^{-4},
\]
since $\eta^2\le\eta h^{-4}$ for $0<\eta,h\le1$.
Moreover
\[
 \|(E+I)^2\psi\|
 \le e^{e_0}\sup_{\lambda\ge0}(1+\lambda)^2e^{-\lambda}
 \le e\,(4/e)=4.
\]
At each fixed $h,\eta>0$, $E$ is elliptic with smooth bounded potential
on the compact product manifold. The spectral bound puts $e^{-E}\phi$ in
every power domain; local elliptic regularity, applied successively to
$E$, gives smoothness. This regularity statement is only at each fixed
cutoff; the uniform quantity used later is the displayed graph moment.
Both $E$ and $c_h$ commute with the gauge action. Functional calculus
therefore preserves the invariant vector, proving every Gauss constraint
exactly without a subsequent physical projection.
\end{proof}

The parameter $s=1$ in \eqref{v27:eq:smooth-packet} is a fixed inverse-energy
regularization parameter in the stated units, \emph{not} physical renewal
time. Equation \eqref{v27:eq:smooth-packet} specifies a preparation before
any source history is generated. It is not inserted between accepted events.
It uses the full original energy and is generally nonlocal. No local
preparation circuit, preparation probability or native source selection of
this vector has been proved. Its role is an explicit existence construction
with a derived moment bound, not a change of the real-time generator.

\section{Necessary and sufficient nonlinear preparation scale}
\label{v27:sec:sharp}

Let $\mathcal H_h^{\mathrm{phys}}$ be the exact invariant subspace and put
\[
 m(h,\eta)=\inf\{\langle\psi,(E+c_h)\psi\rangle:
   \psi\in\mathcal H_h^{\mathrm{phys}}\cap H^1,\ \|\psi\|=1\}.
\]

\begin{theorem}[The full nonlinear unsubtracted vacuum scale]
\label[theorem]{v27:thm:sharp}
For $N\ge3$ and $0<\eta\le1$,
\begin{equation}
 \boxed{\qquad
 \frac94\eta h^{-4}\le m(h,\eta)\le C_*\eta h^{-4}.
 \qquad}
 \label{v27:eq:two-sided}
\end{equation}
Consequently, for any refining sequence $h\downarrow0$, $0<\eta_h\le1$,
the following are equivalent:
\begin{enumerate}[label=\textup{(\roman*)},itemsep=3pt]
\item $\eta_hh^{-4}\to0$;
\item there exist exact physical finite-energy states with
  $\langle E_{h,\eta_h}+c_h\rangle\to0$;
\item there exist smooth exact physical unit vectors with the same vanishing
  comparison and a uniform bound $\|(E_{h,\eta_h}+I)^2\psi_h\|\le4$ for
  all sufficiently small $h$.
\end{enumerate}
The sufficient preparations also have vanishing comparison under the
unchanged quantum evolution, uniformly on every fixed finite physical-time
interval. The necessity is valid at every time and for every state, not
just for the constructed packets.
\end{theorem}
\begin{proof}
The lower bound is \cref{v27:thm:floor}; $c_h\ge0$. The upper bound is
the exactly physical packet of \cref{v27:thm:packet-energy}. This proves
\eqref{v27:eq:two-sided} and the necessity. If (i) holds, the smallness
threshold of \cref{v27:thm:smooth-packet} eventually holds, proving (iii).
The implication (iii)$\Rightarrow$(ii) is immediate. Finally the already
proved \cref{v22:thm:quantum} multiplies the initial comparison by at most
$2+T^2/2$, with no $h$ in this time factor. Apply it to the new vectors.
\end{proof}

This is a statement about the full nonlinear group-valued Wilson action on
its exact physical sector. It upgrades the earlier \emph{quadratic-model}
scale, but not its sharp leading coefficient. It also does not say that a
fixed nonzero physical Planck parameter has a classical vacuum limit in
these unsubtracted norms: \eqref{v27:eq:floor} says the opposite. A
renormalized quantum continuum, a different physical normalization, and a
classical limit are different questions. Neither a small charged mean nor
extra independent preparations can evade this floor while keeping the same
positive action operator.

\section{Transfer to finite local resolved histories and improved explicit powers}
\label{v27:sec:resolved}

Use the same product Peter--Weyl cutoff $P_J$, $J\ge2$, local electric-link
and Wilson-plaquette contexts, gauge-commuting six-outcome instruments and
context probabilities as Version 22. The accepted alphabet still includes
its declared spatial context. Only the initial vector changes to
\[
 \psi_{h,\eta,J}=P_J\psi_{h,\eta}/\|P_J\psi_{h,\eta}\|.
\]
All products in $P_J\mathcal C P_J$ are formed before compression. Retain
$b_J=1+\eta^2(J+1)^2$ and
$\Lambda_J=1+\eta^2(J+1)^2/h$. The exact clock remains
$d=4\vartheta\epsilon/11$, $0<\vartheta\le1$.

\begin{theorem}[The improved preparation enters the unchanged local instrument]
\label[theorem]{v27:thm:resolved}
Assume $C_*\eta h^{-4}\le1$, the sufficiently small inherited tensor-tail
threshold, and $6h^{-3}\epsilon\le\tau_0$. For the actual conditional
states initialized by $\psi_{h,\eta,J}$, put
\begin{equation}
 \begin{split}
 e_j&=\tr(\rho_jP_J(E+c_h)P_J),\\
 \mathfrak b_h^{(27)}&=\eta h^{-4}
  +\frac{h^{-14}}{\Lambda_J}(1+T/\eta)^2
  +\epsilon h^{-12}\eta^{-2}b_J^3.
 \end{split}
 \label{v27:eq:budget}
\end{equation}
Then on every fixed finite physical interval,
\begin{align}
 \max_{j\le T/d}\E e_j&\le C_T\mathfrak b_h^{(27)},\\
 \PP\{\max_{j\le T/d}e_j>u\}&\le C_T\mathfrak b_h^{(27)}/u
                                              \quad(u>0).
 \label{v27:eq:resolved}
\end{align}
Every attained state satisfies the exact Gauss constraints. At \emph{every}
attained record one also has
\begin{equation}
                         e_j\ge\frac94\eta h^{-4}.
 \label{v27:eq:record-floor}
\end{equation}
For $0<\mathfrak b_h^{(27)}\le1/2$,
\begin{equation}
 \E\max_j e_j\le C_T\mathfrak b_h^{(27)}
       \left[1+\log\left(2+\frac{h^{-4}b_J}{\mathfrak b_h^{(27)}}\right)\right].
 \label{v27:eq:maximal-mean}
\end{equation}
\end{theorem}
\begin{proof}
The tail bound in \cref{v22:prop:cutoff} is stated for all $0<\eta\le1$;
its restriction $\eta\le h^7$ in the earlier resolved theorem came from
the \emph{old preparation}. The new vector supplies the required graph
moment uniformly by \cref{v27:thm:smooth-packet}. Apply the inherited
Ritz/Duhamel estimate \cref{v14:thm:Galerkin} with that moment and
$\delta=C h^{-12}/\Lambda_J$. Thus the finite and uncompressed unitary
vectors differ in the $E+I$ form norm by at most
$C_T\delta(1+T/\eta)^2$ after squaring.

The full quantum comparison is $O_T(\eta h^{-4})$ by
\cref{v27:thm:sharp}. The inherited domination
$0\preceq E+c_h\preceq6h^{-2}(E+I)$ and the squared triangle inequality
for its form square root therefore give
\[
 \sup_{t\le T}\langle\psi_J(t),(E+c_h)\psi_J(t)\rangle
 \le C_T\left[\eta h^{-4}
       +h^{-14}\Lambda_J^{-1}(1+T/\eta)^2\right].
\]
The local norm estimates of \cref{v22:prop:local} are also valid for every
$0<\eta\le1$: $A_C\le Ch^{-4}b_J$ and
$\mathfrak C\le Ch^{-8}\eta^{-2}b_J^2$. The unchanged local clock theorem
\cref{v14:thm:trajectory} bounds the finite-unitary fidelity error $\xi_j$
in mean and maximal probability by $C_T\epsilon\mathfrak C$.
The positive comparison transfer proved in Version 14 gives
\[
 e_j\le2\langle\psi_J(jd),\mathcal C\psi_J(jd)\rangle+2A_C\xi_j.
\]
This proves \eqref{v27:eq:resolved}, using the trivial probability bound
for thresholds below the deterministic term. No union bound over ticks
occurs. Gauge invariance of the preparation, cutoff and each local Kraus
operator gives the exact constraints. Compress \cref{v27:eq:floor} to
obtain \eqref{v27:eq:record-floor}; this is independent of any fidelity
argument. Finally $0\le e_j\le A_C$. Integrating
$\min(1,C_T\mathfrak b_h^{(27)}/u)$ gives
\eqref{v27:eq:maximal-mean}.
\end{proof}

\begin{corollary}[Explicit simultaneous local scales]
\label[corollary]{v27:cor:scales}
One sufficient choice is
\begin{equation}
 \boxed{\quad
 \eta_h=h^6,\quad a_h=\eta_hh=h^7,\quad
 J_h=\lceil h^{-20}\rceil,\quad\epsilon_h=h^{110},\quad
 d_h=\frac{4\vartheta_h}{11}h^{110}.
 \quad}
 \label{v27:eq:scales}
\end{equation}
Then $\mathfrak b_h^{(27)}=O_T(h^2)$, and the expected maximal comparison
is $O_T(h^2(1+|\log h|))$. Its fixed-time mean is $O_T(h^2)$ and its
pointwise lower bound is $9h^2/4$. Thus the energy power $h^2$ cannot be
improved to $o(h^2)$ at this same $\eta_h$ without changing or subtracting
the stated action. The instrument-step and cutoff exponents are not claimed
optimal.
\end{corollary}
\begin{proof}
Here $b_J\asymp h^{-28}$ and $\Lambda_J\asymp h^{-29}$.
The three terms of \eqref{v27:eq:budget} have orders
$h^{6-4}=h^2$, $h^{-14+29-12}=h^3$ and
$h^{110-12-12-84}=h^2$. The tensor-tail constant is $O(h^{17})$, the
uniform graph moment is at most four, and
$6h^{-3}\epsilon_h=6h^{107}\to0$, so all thresholds eventually hold.
The comparison norm is $A_C=O(h^{-32})$; consequently the logarithm is
$O(1+|\log h|)$. Apply \cref{v27:thm:resolved} and its recordwise lower
bound.
\end{proof}

The physical connection-orbit comparison, electric and full ordered magnetic
energies, averaged noncontractible period contrasts, and independently
specified electric/Wilson coefficient insertions retain exactly their
Version 22 bounds. Each is bounded by the same positive $e_j$ (up to its
stated fixed factor), so its vacuum conclusion follows from
\eqref{v27:eq:resolved}. This is not convergence to a nonzero field, an
invention of a simultaneous trajectory of noncommuting observables, or a
new complete metric-source coupling.

\section{Outcome and the remaining primitive-source obligations}
\label{v27:sec:outcome}

The new information is a two-sided statement about the full nonlinear
physical action, rather than a longer history of the Pauli benchmark:
\[
 \boxed{\quad
 \text{existence of exact physical vacuum-concentrated preparations}
       \quad\Longleftrightarrow\quad \eta_hh^{-4}\to0.
 \quad}
\]
Its lower bound is independent of the preparation or source law. Its upper
bound is attained in order by a declared gauge-invariant packet and a
same-action regularization, with a proved uniform fourth energy moment.
The nonlinear action and singular vacuum orbit are both retained. The
extra tree-length power and a possible logarithmic entropy loss have been
removed; this is not a claim that the preceding tree construction was
incorrect.

The resulting local resolved continuation needs a less restrictive explicit
quantum, tensor and event scale. The new packet, however, uses a global
configuration cost and a global initial heat regularization. It does not
provide a local preparation circuit or show that the primitive law chooses
these states. Subsequent local operator support remains exactly the inherited
property. The actual source-to-action map, selection of $\eta_h$, and
one-sided physical energy work from Version 24 still have to be verified on
the intended source. At a fixed nonvanishing $\eta$, this unsubtracted
vacuum comparison cannot vanish under spatial refinement; changing to a
renormalized quantity requires a new source accounting.

The full Einstein inverse audited at the start remains unresolved, as do
spatially uniform nonzero-reference non-Abelian quantum propagation, chiral
matter and Einstein backreaction. The next noncircular test is whether the
intended physical normalization and preparation satisfy the now sharp
nonlinear energy entrance, or whether a different explicitly differentiated
renormalized action is intended. Another improvement of the already rejected
single-input benchmark would not answer that question.

\paragraph{Verification and preservation.}
All Version 26 bytes before its final document command are preserved.
New labels have prefix \texttt{v27:}. The diagnostics check the exact
$SU(2)$ derivative and plaquette-incidence constants, the completed-square
floor, intrinsic/chord constants, ball-volume and partition-ratio bounds,
weighted heat identities, finite Galerkin tests with products formed before
compression, and the simultaneous powers. Small-group spectral tests are
labelled finite approximations, not simulations of the full spatial process.
They do not replace the proofs or constitute machine-checked verification.
The initial heat operator is evaluated numerically only in those small tests;
no computational-efficiency claim is made for the constructed field state.

\begingroup
\renewcommand{\refname}{Additional references for Version 27}
\begin{thebibliography}{9}
\small\raggedright
\bibitem{v27:KS} J.~Kogut and L.~Susskind,
\emph{Hamiltonian formulation of Wilson's lattice gauge theories},
Physical Review D \textbf{11} (1975), 395.
\href{https://doi.org/10.1103/PhysRevD.11.395}{doi:10.1103/PhysRevD.11.395}.
Background for the compact-link Hamiltonian framework; all scale and
commutator estimates here use the explicit normalization above.
\bibitem{v27:Chin} S.~A.~Chin, O.~S.~van Roosmalen, E.~A.~Umland,
and S.~E.~Koonin,
\emph{Exact ground-state properties of the $SU(2)$ Hamiltonian lattice gauge theory},
Physical Review D \textbf{31} (1985), 3201.
\href{https://doi.org/10.1103/PhysRevD.31.3201}{doi:10.1103/PhysRevD.31.3201}.
Background for the compact-sphere many-body and variational description;
no ground-state numerical result is imported as a bound here.
\bibitem{v27:BT} B.~Bahr and T.~Thiemann,
\emph{Gauge-invariant coherent states for Loop Quantum Gravity II:
Non-abelian gauge groups}, 2007 preprint.
\href{https://arxiv.org/abs/0709.4636}{arXiv:0709.4636}.
Background for invariant preparations and singular gauge orbits. The
product-volume, weak-gradient and weighted-heat arguments in this body
are proved independently and do not assume a regular-orbit asymptotic.
\end{thebibliography}
\endgroup
\addtocontents{toc}{\protect\endgroup}
% END IMMUTABLE VERSION 27 BODY
% APPEND VERSION 28 HERE, BEFORE THE FINAL END-DOCUMENT COMMAND.


% BEGIN IMMUTABLE VERSION 28 BODY
\clearpage
\begin{center}
{\Large\bfseries Version 28: The actual configuration law before a quantum preparation}\\[.5em]
{\large A sharp-cost heat lift, a physical-scale nonconcentration test,
and the distinction between Wilson measures and physical Hamiltonians}
\end{center}
\makeatletter
\addtocontents{toc}{\protect\begingroup}
\addtocontents{toc}{\protect\makeatletter}
\addtocontents{toc}{\protect\def\protect\l@section{\protect\bfseries\protect\@dottedtocline{1}{0em}{2.5em}}}
\makeatother
\addcontentsline{toc}{part}{Version 28: Actual-law preparation tests and a sharp-cost physical heat lift}

\section{Companion audit: three measures and two generators are not interchangeable}
\label{v28:sec:audit}

The supplied Yang--Mills V86 file \cite{v28:YM86} changes the immediate
question. Version 27 constructs low-energy physical states, but it does not
show that their configuration probabilities are supplied by a native law.
The companion now gives quantitative bounds for an \emph{actual} restricted
weak-coupling law. This permits a direct preparation test, rather than
another freely selected wave packet. It also supplies earlier physical-transfer
results that must not be confused with that restricted measure.

\paragraph{What was inspected and what is inherited.}
The continuation and theorem map of the cumulative companion was inspected,
with detailed reads of the physical-transfer and time-zero blocks V52--V53
and the relevant gauge-stability, normalization, inverse, transport and
nonconcentration blocks V79--V86. This is a dependency-directed audit, not
an independent verification of all sixty thousand source lines.
The following literal labels fix the inputs and avoid page-number ambiguity:
\begin{itemize}[leftmargin=1.6em,itemsep=3pt]
\item \nolinkurl{prop:v52-fine-transfer} and
\nolinkurl{thm:v53-static} identify the \emph{full compact} Wilson transfer,
its Perron time-zero law, and a complete static interaction in the stated
small-$\beta$ regime. Their transfer Hamiltonian is
$-a_t^{-1}\log\mathsf P$, with an independently specified physical time
spacing $a_t$. We do not rederive that Hamiltonian or set it equal to the
canonical $E_{h,\eta}$ used here.
\item \nolinkurl{thm:v84-amplitude},
\nolinkurl{prop:v84-inverse}, and \nolinkurl{thm:v85-TV}
concern the \emph{restricted} mean-zero Coulomb law, with its determinant,
Haar factor, positive gauge domain and original normalization retained.
They are not bounds for the full compact Perron law without another theorem.
\item \nolinkurl{thm:v86-joint} gives the actual short-interval plaquette
nonconcentration used below. \nolinkurl{prop:v86-thermal} and
\nolinkurl{lem:v86-doubling} retain a still-unproved dyadic normalizer
estimate. We do not replace it by Version 27's partition ratio for a
\emph{different} orbit-distance density.
\end{itemize}
The previously proved nonlinear energy floor, weak-gradient orbit estimate,
tensor approximation, local gauge-commuting instruments and renewal-clock
comparison remain \cref{v27:thm:floor,v22:lem:distance,v22:prop:cutoff,v22:prop:local,v14:thm:trajectory}. No new proof of those results is claimed.

\paragraph{New direction and exact scope.}
We prove a quantitative procedure taking \emph{any} spatial configuration
probability to an exactly physical pure state after a declared, linkwise
heat regularization of that probability. Its kinetic cost and Wilson bias
have the same $\eta h^{-4}$ order as the inherited lower bound. For this
specified lift, vanishing comparison is equivalent to two separate conditions:
a small original configuration cost and the small quantum scale.
We then test the latest restricted law under an \emph{explicit additional}
identification of its slice with physical spacing $h=M^{-1}$.
On logarithmic coupling sequences that law fails the raw classical magnetic
entrance in mean and in probability. This does not exclude a renormalized
quantum field or the companion's different transfer Hamiltonian.

The density-to-vector step is an existence construction, generally nonlinear
in the input law. It is not a CPTP decoder of one source sample, an executed
native preparation, or a proof that the source selects our operations.
Heat regularization of a density, heat regularization of a wavefunction, and
physical real-time evolution are three different operations throughout.
The density-gradient and heat-entropy tools have classical precedents
\cite{v28:HO,v28:Ni}; all constants and specialized estimates used here are
proved below.

\section{The actual weak-coupling input in its own normalization}
\label{v28:sec:input}

Fix the companion's $0<R\le1/4$. Its lattice has $M^4$ sites,
mean-zero Coulomb link coordinates $a$, and probability
\begin{equation}
 \dd\mu_{M,\beta}(a)=Z_{M,\beta}^{-1}
 \one_{\Omega_R}(a)\det F(a)
             e^{-\beta S(a)-\mathcal H(a)}\,\dd a.
 \label{v28:eq:actual-law}
\end{equation}
All notation in this equation is the companion's. In particular $F$ is its
relative gauge matrix, not a physical field-strength tensor. Its plaquette
observable is the literal nonlinear trace
\[
 s_p(a)=1-\tfrac12\operatorname{Re}\tr U_p(a)\in[0,2].
\]
Set $H_\beta=\log(2+\beta)$. The inherited \nolinkurl{thm:v86-joint}
states that if
\begin{equation}
 M\longrightarrow\infty,\quad\beta_M\longrightarrow\infty,
 \qquad \frac{\sqrt{\beta_M}}{M^2\sqrt{H_{\beta_M}}}\longrightarrow0,
 \label{v28:eq:window}
\end{equation}
then for a fixed explicitly positive $\sigma_R$ and all sufficiently large
$M$, uniformly in plaquettes,
\begin{equation}
 b_M^2=\frac{\sigma_R^2}{\beta_M H_{\beta_M}},\qquad
 \sup_{|I|\le2b_M^2}\mu_{M,\beta_M}(s_p\in I)\le\frac34,
 \qquad \E_\mu s_p\ge\frac{b_M^2}{2}.
 \label{v28:eq:NC-input}
\end{equation}
No independence of plaquettes is contained in this input. The companion
also proves $\E s_p\le C_R H_\beta/\beta$ by translation invariance and its
complete action bound. The lower and upper orders are not claimed to match.

To perform a physical test, choose one coordinate slice and keep its three
spatial link directions. \emph{Declare}, for this test only, $h=M^{-1}$,
and identify each source group element $U_e$ with that same physical link.
Let $\nu_h$ be the exact pushforward of \eqref{v28:eq:actual-law} under this
slice map. No conditional reweighting or three-dimensional Wilson reset is
made. The slice contains $3h^{-3}$ spatial plaquettes. The companion does
not itself establish that this declaration is its physical continuum map.

On $\mathcal Q_h=SU(2)^{3h^{-3}}$ use normalized product Haar measure,
$T_a=-\ii\sigma_a/2$, and the physical metric
\[
 g_h=h\sum_{e,a}\vartheta_{e,a}^2,\qquad
 D_h^{\rm dim}=\dim\mathcal Q_h=9h^{-3}.
\]
The superscript on $D_h^{\rm dim}$ distinguishes dimension from every earlier
operator named $D_h$. Retain exactly
\begin{equation}
 T_{h,\eta}=-\frac{\eta^2}{2}\Delta_{g_h},\quad
 W_h=\frac4h\sum_p s_p,\quad E_{h,\eta}=T_{h,\eta}+W_h,\quad
 \mathcal C_{h,\eta}=E_{h,\eta}+c_h.
 \label{v28:eq:canonical}
\end{equation}
Here $c_h$ is Version 22's trivial-holonomy orbit cost, not a new source
weight. $\beta$ and $\eta$ are independent parameters unless a separate
physical theorem identifies them.

\section{The minimum kinetic cost of a prescribed configuration density}
\label{v28:sec:Fisher}

\begin{proposition}[Exact density-constrained kinetic infimum]
\label[proposition]{v28:prop:minimal}
Let $p\ge0$, $\int p\,\dd\mu_H=1$, on $\mathcal Q_h$.
Among all density matrices with configuration marginal $p\,\dd\mu_H$, the
infimum of $\tr(\rho T_{h,\eta})$ is
\begin{equation}
 \frac{\eta^2}{2}\int|\nabla\sqrt p|_{g_h}^2\,\dd\mu_H
 =\frac{\eta^2}{8}\mathcal I_h(p),\qquad
 \mathcal I_h(p)=\int_{p>0}\frac{|\nabla p|_{g_h}^2}{p}\,\dd\mu_H,
 \label{v28:eq:minimal}
\end{equation}
with value $+\infty$ when $\sqrt p\notin H^1$. The formula on the right is
interpreted through the first expression for nonsmooth densities.
The real pure vector $\sqrt p$ attains the infimum. If $p$ is gauge invariant,
that vector is in the exact Gauss subspace, so the same infimum holds with
exact physical support required. Untouched finite auxiliary memory and
mixing cannot lower the bound.
\end{proposition}
\begin{proof}
Write a density matrix spectrally as
$\rho=\sum_j\lambda_j|\psi_j\rangle\langle\psi_j|$.
For finite kinetic energy, the Hilbert-valued map
$\Psi(U)=(\sqrt{\lambda_j}\psi_j(U))_j$ has square-integrable weak gradient,
and $\|\Psi(U)\|=\sqrt{p(U)}$. The norm map is one-Lipschitz; its weak
chain rule, obtained by finite truncation and replacing the norm by
$(\|\Psi\|^2+\delta^2)^{1/2}$, gives
\[
 |\nabla\sqrt p|^2\le\sum_j\lambda_j|\nabla\psi_j|^2
 \quad\hbox{almost everywhere}.
\]
Integration gives the lower bound. If the left side is infinite no
finite-energy density matrix with that marginal exists. For the real vector
$\sqrt p$ equality is immediate. Gauge invariance means invariance of that
vector, not merely of a density matrix under conjugation; hence every weak
Gauss generator annihilates it. Auxiliary components are simply additional
coordinates of the same Hilbert-valued map, so the proof is unchanged.
\end{proof}

For $p>0$ smooth, an additional same-action test is exact:
\begin{equation}
 \frac{E_{h,\eta}\sqrt p}{\sqrt p}
 =W_h-\frac{\eta^2}{4}\Delta_{g_h}\log p
           -\frac{\eta^2}{8}|\nabla\log p|_{g_h}^2.
 \label{v28:eq:eigen-test}
\end{equation}
This follows by differentiating $\sqrt p=e^{(\log p)/2}$. Thus $\sqrt p$ is
an eigenvector of the \emph{independently fixed} canonical energy precisely
when the right side is spatially constant on configuration space.
Conversely a positive eigenvector is the ground vector, by the positivity of
the compact Schr\"odinger heat kernel, or by the ground-state transform
$\langle u,(E-E_0)u\rangle=(\eta^2/2)\int p|\nabla(u/\sqrt p)|^2$.
No stationarity of a statistical law implies \eqref{v28:eq:eigen-test} for
an unrelated Hamiltonian. The companion's Perron eigenfunction already
satisfies its own \emph{transfer} eigenvalue equation; that is a different
identity, not a missing proof in that companion.

\section{Heat transport of an arbitrary law and its exact Wilson bias}
\label{v28:sec:heat}

Let $\nu$ be any probability measure on $\mathcal Q_h$, even singular.
Let $\overline\nu$ be its average over vertex gauge transformations. This
leaves the joint distribution of all gauge-invariant configuration
observables unchanged. For $t>0$ define
\begin{equation}
 p_t\,\dd\mu_H=(e^{t\Delta_{g_h}})_*\overline\nu,
 \qquad \phi_t=\sqrt{p_t},\qquad
 w_\nu=\int W_h\,\dd\nu,\quad c_\nu=\int c_h\,\dd\nu.
 \label{v28:eq:heat-law}
\end{equation}
The heat operator is the product of one-link Casimir heat operators; it
commutes with independent left and right translations at link endpoints.
Consequently $p_t$ is smooth, strictly positive and gauge invariant, and
$\phi_t$ is an exactly physical unit vector. The parameter $t$ is a
configuration-heat parameter, not a time in the accepted renewal history.

\begin{lemma}[Dimension-explicit density-gradient estimate]
\label[lemma]{v28:lem:Fisher}
For every initial probability in \eqref{v28:eq:heat-law},
\begin{equation}
 \mathcal I_h(p_t)\le\frac{D_h^{\rm dim}}{2t},\qquad
 \langle\phi_t,T_{h,\eta}\phi_t\rangle
                 \le\frac{9\eta^2h^{-3}}{16t}.
 \label{v28:eq:heat-Fisher}
\end{equation}
No derivative of the initial density, positive lower bound for that density,
or inverse gauge determinant is assumed.
\end{lemma}
\begin{proof}
The product of compact groups with a bi-invariant positive metric has
nonnegative Ricci tensor: the sectional form is one quarter of the squared
Lie bracket, with the metric scaling retained. For a positive heat solution,
put $f=\log p$ and $I=\int p|\nabla f|^2$.
The equations $\dot f=\Delta f+|\nabla f|^2$ and
$\frac12\Delta|\nabla f|^2=\langle\nabla f,\nabla\Delta f\rangle+
|\nabla^2f|^2+\operatorname{Ric}(\nabla f,\nabla f)$, followed by integration
by parts, give the exact identity
\[
 I'=-2\int p\bigl(|\nabla^2f|^2+
                     \operatorname{Ric}(\nabla f,\nabla f)\bigr).
\]
For clarity, differentiation of $I$ gives
$\int(\Delta p)|\nabla f|^2+2\int p\langle\nabla f,
\nabla(\Delta f+|\nabla f|^2)\rangle$.
Apply the displayed differential identity and use
$\int p\Delta|\nabla f|^2=-\int p\langle\nabla f,\nabla|\nabla f|^2\rangle$;
the result is the stated formula. Since
$|\nabla^2 f|^2\ge(\Delta f)^2/D_h^{\rm dim}$ and
$\int p\Delta f=-I$, weighted Cauchy--Schwarz yields
$I'\le-2I^2/D_h^{\rm dim}$. On any interval $[s,t]$ with $s>0$ this gives
$I(t)\le D_h^{\rm dim}/[2(t-s)]$; the case $I=0$ is immediate.
Send $s\downarrow0$. Compact heat-kernel smoothing justifies every calculation
at positive times for an arbitrary initial probability. The second assertion
is \cref{v28:prop:minimal}, with $D_h^{\rm dim}=9h^{-3}$.
\end{proof}

\begin{proposition}[Exact plaquette transport and a two-way orbit estimate]
\label[proposition]{v28:prop:bias}
Put $q_t=e^{-3t/(4h)}$. Then
\begin{align}
 \int W_h p_t\,\dd\mu_H
 &=e^{-3t/h}w_\nu+12h^{-4}(1-e^{-3t/h}),
                 \label{v28:eq:Wilson-heat}\\
 \int c_hp_t\,\dd\mu_H
 &\le2c_\nu+6h^{-2}(1-q_t),
                 \label{v28:eq:orbit-heat}\\
 c_\nu&\le2\int c_hp_t\,\dd\mu_H+6h^{-2}(1-q_t).
                 \label{v28:eq:orbit-reverse}
\end{align}
There is also a coupling of $U\sim\overline\nu$ and $U'\sim p_t\mu_H$ with,
for every elementary plaquette,
\begin{equation}
 \E\bigl|\sqrt{s_p(U')}-\sqrt{s_p(U)}\bigr|^2
                      \le12t/h.
 \label{v28:eq:coupling}
\end{equation}
All constants are independent of volume and of the initial law.
\end{proposition}
\begin{proof}
In the fundamental representation $\sum_aT_a^2=-3I/4$.
An elementary plaquette uses four distinct link variables for $N\ge3$.
Applying the link Casimir to its trace therefore gives
$\Delta_{g_h}(1-s_p)=-(3/h)(1-s_p)$, including inverse-oriented links.
Heat evolution of this exact eigenfunction gives
$\int s_pp_t=1-e^{-3t/h}+e^{-3t/h}\int s_p\dd\nu$.
Sum the $3h^{-3}$ plaquettes with coefficient $4/h$ to obtain
\eqref{v28:eq:Wilson-heat}. This is not a linear-curvature approximation.

Realize the product heat transition as $U'_e=U_eB_e$, with independent
one-link central heat increments satisfying
$\E\tfrac12\tr B_e=q_t$.
For any candidate vertex gauge, the triangle inequality for matrix chords
and $(a+b)^2\le2a^2+2b^2$ give
\[
 c_h(U')\le2c_h(U)+2h\sum_e v(B_e),\quad
 c_h(U)\le2c_h(U')+2h\sum_e v(B_e).
\]
These follow by first using a minimizer at one endpoint in the cost of the
other endpoint. Only existence is used; no derivative or measurable choice
of a minimizing gauge enters the integrated inequalities. Taking expectations
proves \eqref{v28:eq:orbit-heat}--\eqref{v28:eq:orbit-reverse}.

Telescoping four unitary factors gives
$\|U'_p-U_p\|_{\rm HS}^2\le4\sum_{e\in\partial p}\|B_e-I\|_{\rm HS}^2$.
The same bound holds for inverse occurrences by unitary invariance.
Each expected squared chord is $4(1-q_t)\le3t/h$.
Also $\sqrt{s_p}=\|U_p-I\|_{\rm HS}/2$, so its squared difference is at
most $\|U'_p-U_p\|_{\rm HS}^2/4$. This proves \eqref{v28:eq:coupling}.
Gauge averaging before the coupling changes none of the $s_p$ values.
\end{proof}

\section{A sharp-cost preparation lift and its two independent entrance conditions}
\label{v28:sec:lift}

\begin{theorem}[Uniform heat lift of a supplied configuration probability]
\label[theorem]{v28:thm:lift}
For any probability $\nu$ on $\mathcal Q_h$, $0<\eta,h\le1$, and $t>0$,
the exactly physical vector in \eqref{v28:eq:heat-law} satisfies
\begin{equation}
 \begin{split}
 \langle\phi_t,\mathcal C_{h,\eta}\phi_t\rangle
 \le{}&\frac{9\eta^2h^{-3}}{16t}
       +e^{-3t/h}w_\nu+12h^{-4}(1-e^{-3t/h})\\
      &+2c_\nu+6h^{-2}(1-e^{-3t/(4h)}).
 \end{split}
 \label{v28:eq:general-lift}
\end{equation}
In particular the predetermined choice
\begin{equation}
 t_{h,\eta}=\eta h/8
 \quad\hbox{gives}\quad
 \boxed{\langle\phi_{t_{h,\eta}},\mathcal C_{h,\eta}
                    \phi_{t_{h,\eta}}\rangle
 \le w_\nu+2c_\nu+\frac{153}{16}\eta h^{-4}.}
 \label{v28:eq:sharp-lift}
\end{equation}
For this lift along any refining sequence, vanishing comparison is
\emph{equivalent} to
\begin{equation}
             w_{\nu_h}+c_{\nu_h}\longrightarrow0,
                   \qquad \eta_hh^{-4}\longrightarrow0.
 \label{v28:eq:iff-input}
\end{equation}
The quantum power is sharp by the inherited nonlinear lower bound; neither
the heat choice nor its numerical coefficient is asserted optimal.
\end{theorem}
\begin{proof}
Add \eqref{v28:eq:heat-Fisher} and
\eqref{v28:eq:Wilson-heat}--\eqref{v28:eq:orbit-heat}.
Use $1-e^{-x}\le x$. The two leading added terms at $t=\eta h/8$ are
$9\eta h^{-4}/2$ each; the orbit addition is
$9\eta h^{-2}/16\le9\eta h^{-4}/16$. This proves
\eqref{v28:eq:sharp-lift} and sufficiency.
Conversely the inherited $\mathcal C_{h,\eta}\succeq(9/4)\eta h^{-4}I$
forces the second condition. In particular $\eta_h\to0$.
Equation \eqref{v28:eq:Wilson-heat} and positivity show
$e^{-3\eta_h/8}w_{\nu_h}\to0$, hence $w_{\nu_h}\to0$.
Equation \eqref{v28:eq:orbit-reverse} gives
$c_{\nu_h}\le2\langle c_h\rangle+9\eta_hh^{-2}/16\to0$.
This proves necessity for the specified lift, not for an arbitrary
reweighted preparation.
\end{proof}

The heat operation in this theorem is local in link variables as a
\emph{Markov smoothing of a classical measure}. The input law may have
arbitrary correlations, and the square-root amplitude generally inherits
those correlations. No local quantum preparation circuit follows.
In particular the procedure is not an affine state decoder: for two distinct
positive invariant heat densities $p_1,p_2$, the mixture of their
square-root pure states has rank two, whereas the square root of
$(p_1+p_2)/2$ is a pure state. Thus one cannot implement this whole law-to-vector
map as a CPTP map on a single unknown input probability state.
This rank argument also prevents interpreting the heat lift as the actual
source transition without an independent preparation mechanism.

\begin{corollary}[A derived high moment, with the additional preparation change exposed]
\label[corollary]{v28:cor:moment}
Set $b_h^0=w_{\nu_h}+2c_{\nu_h}+(153/16)\eta_hh^{-4}$ and assume $b_h^0\le1$.
With $\phi_h=\phi_{t_{h,\eta}}$, define
\begin{equation}
              \psi_h=\frac{e^{-E_{h,\eta}}\phi_h}
                            {\|e^{-E_{h,\eta}}\phi_h\|}.
 \label{v28:eq:amplitude-filter}
\end{equation}
Then $\psi_h$ is smooth and exactly physical,
\begin{equation}
 \langle\psi_h,\mathcal C_{h,\eta}\psi_h\rangle\le Cb_h^0,
 \qquad \|(E_{h,\eta}+I)^2\psi_h\|\le4.
 \label{v28:eq:moment}
\end{equation}
Its configuration marginal differs from $p_{t_{h,\eta}}\mu_H$ in total
variation by at most $2\sqrt{b_h^0}$. This \emph{extra} inverse-energy
filter is not required for the exact density lift, and it does change that
density. It is part of an optional declared initial preparation only.
\end{corollary}
\begin{proof}
Apply the weighted amplitude-heat inequalities of \cref{v27:lem:heat} at
parameter one, now to the arbitrary vector $\phi_h$ with its newly proved
small energy and orbit expectation. They give
$\langle E+c_h\rangle_{\psi_h}\le e_0+e^{2e_0}(c_0+\eta^2/8)\le Cb_h^0$.
The spectral estimate
$\|(E+I)^2\psi_h\|\le e^{e_0}\sup_{x\ge0}(1+x)^2e^{-x}\le4$
is exactly the inherited one. Gauge invariance is preserved.
Moreover $(1-e^{-x})^2\le x$ for $x\ge0$, so
$\|\phi_h-e^{-E}\phi_h\|\le\sqrt{e_0}$.
Normalization costs at most the same amount by the reverse triangle
inequality, giving $\|\psi_h-\phi_h\|\le2\sqrt{e_0}$.
Cauchy--Schwarz applied to
$|\psi_h|^2-|\phi_h|^2$ bounds its half-$L^1$ norm by that vector norm.
This proves the marginal claim. No energy estimate is inferred from total
variation alone; the weighted heat inequality supplied it separately.
\end{proof}

\section{Testing the companion law: a raw classical-field obstruction}
\label{v28:sec:actual-test}

The following consequences apply to the \emph{unchanged} slice probability
$\nu_h$ specified after \eqref{v28:eq:NC-input}. Write
\begin{equation}
                L_M=b_M^2h^{-4}
                   =\frac{\sigma_R^2M^4}{\beta_MH_{\beta_M}}.
 \label{v28:eq:physical-width}
\end{equation}
The physical squared chord of one plaquette is
$\|(U_p-I)/h^2\|_{\rm HS}^2=4h^{-4}s_p$.

\begin{theorem}[Native plaquette fluctuations cannot be relabelled as a small raw magnetic field]
\label[theorem]{v28:thm:source-test}
Under the inherited window \eqref{v28:eq:window}, eventually
\begin{equation}
 \boxed{\quad \E_{\nu_h}W_h\ge6L_M,
                  \qquad \nu_h\{W_h>3L_M\}\ge\frac17.\quad}
 \label{v28:eq:mean-prob}
\end{equation}
The probability inequality uses no independence or spatial mixing.
If $L_M\to\infty$, these raw magnetic energies are not tight and cannot
converge in probability to a finite classical magnetic energy.
No quantum state with that exact configuration marginal, regardless of phase,
mixture or untouched memory, can have a smaller mean magnetic insertion.

For any deterministic scalar subtraction $d_M$, the one-plaquette variables
$h^{-4}s_p-d_M$ are also not tight when $L_M\to\infty$: for every fixed
$A<\infty$, eventually
\begin{equation}
       \nu_h\{|h^{-4}s_p-d_M|\le A\}\le\frac34.
 \label{v28:eq:counterterm}
\end{equation}
This last assertion is pointwise, not a theorem excluding a smeared or
renormalized continuum field.
\end{theorem}
\begin{proof}
Sum the inherited $\E s_p\ge b_M^2/2$ over $3h^{-3}$ spatial plaquettes
with weight $4/h$. This gives the first inequality.
From \eqref{v28:eq:NC-input}, every plaquette has
$\nu_h(s_p>2b_M^2)\ge1/4$.
Let $Y$ be the fraction of spatial plaquettes with that property. Then
$0\le Y\le1$ and $\E Y\ge1/4$. For $\theta=1/8$,
\[
 \E Y\le\theta+(1-\theta)\PP(Y>\theta)
 \quad\Longrightarrow\quad \PP(Y>1/8)\ge1/7.
\]
On that event
$W_h>(4/h)(3h^{-3})(1/8)(2b_M^2)=3L_M$.
A probability at least $1/7$ above a diverging threshold rules out tightness.
Magnetic energy is a multiplication observable, whose expectation depends
only on the configuration marginal, proving the phase/memory statement.
Finally the event in \eqref{v28:eq:counterterm} restricts $s_p$ to an
interval of length $2Ah^4\le2b_M^2$ eventually. Apply the inherited
concentration bound at its arbitrary center. Constants do not alter the
length, so the estimate is uniform in the deterministic subtraction.
\end{proof}

For example, $\beta_M=c\log M$, $c>0$, satisfies the companion window and
has $L_M\to\infty$. Thus the new lower bound concerns a regime in which
\emph{unscaled} plaquettes approach identity in mean, yet their declared
classical curvature normalization is not controlled. A necessary condition
for $W_h\to0$ in probability within this window is
$\beta_MH_{\beta_M}h^4\to\infty$; this is not asserted sufficient.
The companion has not identified $M^{-1}$ as physical spacing, nor its
restricted law as a physical time-zero state. These conclusions test that
\emph{additional} proposed identification, not its unproved premise.

A bounded-energy smooth reference cannot repair the raw conclusion. With
$B_p=(U_p-I)/h^2$ and any deterministic reference $B_p^*$ with
$h^3\sum_p\|B_p^*\|_{\rm HS}^2\le C_*$,
\[
 h^3\sum_p\|B_p-B_p^*\|_{\rm HS}^2
 \ge\left(\sqrt{W_h}-\sqrt{C_*}\right)^2
 \quad\hbox{on }\{W_h\ge C_*\}.
\]
This is the reverse triangle inequality, since
$h^3\sum_p\|B_p\|_{\rm HS}^2=W_h$ exactly. It establishes the corresponding
failure of a raw strong $L^2$ curvature comparison in that experiment.
It does not bound a gauge-field distribution tested only after spatial
smearing or a specified nontrivial coarse reader.

\begin{proposition}[The small physical heat lift does not erase this obstruction]
\label[proposition]{v28:prop:heat-NC}
For the heat law of the same slice at any $t>0$, set
\[
 \delta_t=\left[\frac14-
       \frac{12t}{h(\sqrt2-1)^2b_M^2}\right]_+.
\]
Then every spatial plaquette obeys
$\PP_{p_t}(s_p>b_M^2)\ge\delta_t$, and
\begin{equation}
 \PP_{p_t}\{W_h>6\delta_t L_M\}
                       \ge\frac{\delta_t}{2-\delta_t}
 \qquad(\delta_t>0).
 \label{v28:eq:heat-NC}
\end{equation}
At $t=\eta h/8$, if $\eta/b_M^2\to0$ and $L_M\to\infty$, the heat-lifted
magnetic energy is still not tight. These two conditions hold for logarithmic
$\beta_M$ whenever the necessary quantum scale $\eta_hh^{-4}\to0$ holds.
\end{proposition}
\begin{proof}
Use the coupling in \cref{v28:prop:bias}. On $s_p(U)>2b_M^2$, if
$|\sqrt{s_p(U')}-\sqrt{s_p(U)}|\le(\sqrt2-1)b_M$, then
$s_p(U')>b_M^2$. Markov's inequality and
\eqref{v28:eq:coupling} give the stated one-plaquette lower bound.
The fraction $Y_t$ of these plaquettes has expectation at least $\delta_t$.
The same bounded-variable argument at threshold $\delta_t/2$ gives
$\PP(Y_t>\delta_t/2)\ge\delta_t/(2-\delta_t)$.
On this event the energy is greater than
$(4/h)(3h^{-3})(\delta_t/2)b_M^2=6\delta_tL_M$.
The asymptotic assertions follow directly. For logarithmic coupling,
$\eta/b_M^2=(\eta h^{-4})(\beta H_\beta h^4)/\sigma_R^2\to0$.
\end{proof}

The amplitude filter \eqref{v28:eq:amplitude-filter} is not covered by this
probability-preservation statement when the initial energy is large. It can
reweight a high-energy law substantially. Calling that reweighting an
unchanged native preparation would be incorrect. The small-$b_h^0$ premise
of \cref{v28:cor:moment} is precisely what makes its alteration quantitatively
controlled; it is unavailable for the logarithmic-coupling slice just tested.

\section{A source-probability entrance to the inherited local resolved dynamics}
\label{v28:sec:resolved}

\begin{theorem}[Conditional transfer from a verified configuration law]
\label[theorem]{v28:thm:resolved}
Suppose a family of spatial probabilities $\nu_h$ is independently supplied
and its gauge-invariant costs satisfy
$w_{\nu_h}+c_{\nu_h}+\eta_hh^{-4}\to0$.
Prepare the heat lift and optional amplitude filter of
\cref{v28:cor:moment}, followed by the original tensor Peter--Weyl projection
$P_J$ and normalization. Use the same local electric/plaquette locked
instruments, context rule, and clock as Version 27. Set
\[
 b_J=1+\eta^2(J+1)^2,\qquad \Lambda_J=1+\eta^2(J+1)^2/h,
 \qquad d=4\vartheta\epsilon/11.
\]
Under the inherited sufficiently large tensor cutoff and small local-step
conditions, the actual conditional records satisfy
\begin{align}
 \max_{j\le T/d}\E e_j&\le C_T\mathfrak b_h,
 &\PP\{\max_{j\le T/d}e_j>u\}&\le C_T\mathfrak b_h/u,\label{v28:eq:resolved}\\
 \mathfrak b_h&=w_{\nu_h}+c_{\nu_h}+\eta h^{-4}
       +\frac{h^{-14}}{\Lambda_J}(1+T/\eta)^2
       +\epsilon h^{-12}\eta^{-2}b_J^3,
 &&\label{v28:eq:budget}
\end{align}
where $e_j=\tr[\rho_jP_J(E+c_h)P_J]$.
All attained states have exact Gauss support. The electric, Wilson, orbit,
period and independently specified coefficient-source bounds of Version 22
follow with this budget. No assumption that the original native law executes
these accepted operations is made; that remains a separate source/action test.
\end{theorem}
\begin{proof}
The new preparation supplies both an initial comparison bounded by
$C(w_\nu+c_\nu+\eta h^{-4})$ and the uniform graph moment at most four.
Version 22's full nonlinear energy propagation multiplies the former by
$2+T^2/2$. Its tensor approximation, with the latter moment, adds
$C_Th^{-14}\Lambda_J^{-1}(1+T/\eta)^2$ to the comparison.
The unchanged local outcome theorem then adds
$C_T\epsilon h^{-12}\eta^{-2}b_J^3$.
More explicitly, for the pure finite-unitary reference projection $P_*(t)$,
positivity gives $V_J\preceq2P_*V_JP_*+2(I-P_*)V_J(I-P_*)$;
the deterministic first term has just been bounded and the second is at
most the inherited comparison norm $Ch^{-4}b_J$ times twice the fidelity
error. Its maximal estimate is the inherited stopped-positive-process
bound. This proves \eqref{v28:eq:resolved} without a per-tick union bound.
Every initial operation and every local Kraus coefficient commutes with the
gauge action. The initial vector is invariant, so all outcomes remain in
the exact physical subspace. All source insertions and composite products
are those already differentiated before compression, not definitions fitted
to this new preparation.
\end{proof}

For example, if the \emph{measured original} configuration cost is $O(h^2)$,
then $\eta=h^6$, $J=\lceil h^{-20}\rceil$, $\epsilon=h^{110}$ give
$\mathfrak b_h=O_T(h^2)$ and the same
$O_T(h^2[1+|\log h|])$ maximal expectation as in Version 27.
The supplied logarithmic-coupling restricted law does not supply that
configuration premise under the raw slice identification. This example is
an entrance criterion, not a claimed evaluation of missing native moments.

\section{The dyadic gate is not closed by changing the preparation}
\label{v28:sec:outcome}

The companion's latest open scalar estimate is
\[
 \log\frac{Z_{M,\beta/2}}{Z_{M,\beta}}\le CV,
 \qquad V=M^4,
\]
for its \emph{same} restricted measure. By its already proved dyadic identity,
this would imply $\E S/V=O(\beta^{-1})$ and its thermal-scale
nonconcentration result. Version 27 instead controls a ratio with exponent
$c_h/a$ and Haar background on a different, spatial compact configuration
space. Neither its tube-volume proof nor the present heat lift bounds the
numerator of the companion's dyadic ratio. No removal of that logarithm is
claimed here. Even its removal would strengthen, rather than repair, the
raw magnetic obstruction at logarithmic coupling: the fluctuation scale
would increase from $(\beta H_\beta)^{-1}$ to $\beta^{-1}$.

The new positive theorem is a source-probability-to-physical-state
\emph{construction with costs}. It separates a geometric requirement,
$\int(W_h+c_h)\dd\nu_h\to0$, from the necessary quantum requirement
$\eta_hh^{-4}\to0$. It does not turn either requirement into a consequence
of a measure's internal algebra, statistical stationarity or an auxiliary
mixing estimate. The normalized square-root map and the optional energy
filter are declared preparation operations, not native reader identities.

The new negative test concerns raw unsubtracted classical energy and strong
curvature norms under a stated slice/spacing identification. It is compatible
with ultraviolet fluctuations of a renormalized quantum theory. The actual
restricted weak-coupling law, the full compact Perron time-zero law, and a
new heat-lifted preparation are not silently substituted for one another.
Likewise the companion's $-a_t^{-1}\log\mathsf P$ is not identified with
$E_{h,\eta}$ without a generator comparison and source-derivative calibration.
A mean-zero Coulomb coordinate condition does not by itself identify trivial
global holonomies; the separate orbit cost remains necessary.

The next source-specific task is consequently to specify the actual spatial
or coarse physical reader and its measured magnetic/orbit costs, or to state
the independently differentiated renormalized comparison that is intended.
The present bounds give a finite test and a quantitative sufficient lift for
that task. No interacting quantum continuum, Hamiltonian mass gap,
nonzero-reference nonlinear non-Abelian propagation, Einstein inverse,
chiral sector, or Einstein backreaction is established by this iteration.

\paragraph{Verification and append-only scope.}
All Version 27 bytes preceding its final end-document command are preserved.
The accompanying diagnostics check the exact fundamental Casimir and
four-link heat eigenvalue, entropy-dissipation normalization on a smooth
$SU(2)$ density, the mixed-state density-gradient inequality, the square-root
map's nonaffinity, the pulse-to-energy constants without independence, the
heat coupling, and the inherited cutoff powers. These finite calculations
support the written proofs; they do not sample the complete restricted
Yang--Mills measure, recompute the companion's certified $\sigma_R$, or
constitute proof-assistant verification. Every new label has prefix
\texttt{v28:}.

\begingroup
\renewcommand{\refname}{Additional references for Version 28}
\begin{thebibliography}{9}
\small\raggedright
\bibitem{v28:YM86} A.~P\'elissier,
\emph{Renewal Yang--Mills Mass-Gap Research Notes}, supplied cumulative V86,
file \nolinkurl{Renewal_Yang_Mills_Mass_Gap_Research_Notes_V86(3).tex}.
The specific physical-transfer, restricted-law, and fluctuation results used
are named in the audit and in \cref{v28:sec:input}.
\bibitem{v28:HO} M.~Hoffmann-Ostenhof and T.~Hoffmann-Ostenhof,
\emph{``Schr\"odinger inequalities'' and asymptotic behavior of the electron
density of atoms and molecules}, Physical Review A \textbf{16} (1977), 1782.
\href{https://doi.org/10.1103/PhysRevA.16.1782}{doi:10.1103/PhysRevA.16.1782}.
Background for square-root density inequalities; the present whole-configuration
kinetic minimization is proved directly, including its physical-sector claim.
\bibitem{v28:Ni} L.~Ni,
\emph{The entropy formula for linear heat equation}, 2003 preprint.
\href{https://arxiv.org/abs/math/0306147}{arXiv:math/0306147}.
Background for heat entropy and nonnegative Ricci curvature. The compact-product
Fisher estimate and the literal Wilson and orbit-transport constants are
calculated independently above.
\end{thebibliography}
\endgroup
\addtocontents{toc}{\protect\endgroup}
% END IMMUTABLE VERSION 28 BODY
% APPEND VERSION 29 HERE, BEFORE THE FINAL END-DOCUMENT COMMAND.



% BEGIN IMMUTABLE VERSION 29 BODY
\clearpage
\begin{center}
{\Large\bfseries Version 29: A reader trilemma from the actual blocked Wilson history}\\[0.5em]
{\large Plaquette retention, physical-time continuity, and vacuum-energy entrance cannot be imposed simultaneously}
\end{center}
\makeatletter
\addtocontents{toc}{\protect\begingroup}
\addtocontents{toc}{\protect\makeatletter}
\addtocontents{toc}{\protect\def\protect\l@section{\protect\bfseries\protect\@dottedtocline{1}{0em}{2.5em}}}
\makeatother
\addcontentsline{toc}{part}{Version 29: Reader-level consequences of the actual blocked Wilson source}

\section{Audit: use the actual blocked history before designing another preparation}
\label{v29:sec:audit}

Version 28 separated three objects: a supplied configuration law, a declared quantum
preparation obtained from that law, and the independently fixed Hamiltonian action.  The
new Yang--Mills source supplies an additional object which is more directly relevant to
the native-reader problem: the \emph{actual Perron time-zero law} of the four-dimensional
Wilson transfer and its slice-local completed spatial history.  We do not reconstruct that
law here.  We use two literal conclusions of the supplied V53 block in
\cite{v29:YM86}:
\begin{enumerate}[label=\textup{(I\arabic*)},leftmargin=2.7em]
\item for every spatial blocking depth, every elementary blocked plaquette trace
      $F_r\in[-1,1]$ has
      \begin{equation}
       \Var(F_r)>c_0,\qquad c_0:=\frac{999}{4000};
       \label{v29:eq:var-input}
      \end{equation}
\item along the isotropic fine-time refinement $a_r\downarrow0$, for the centered
      variable $X_r=F_r-\E F_r$ in the same stationary history,
      \begin{equation}
       \E|X_r(1)-X_r(0)|^2>c_0,
       \label{v29:eq:increment-input}
      \end{equation}
      where the two observations are separated by one original fine time step of
      physical duration $a_r$.
\end{enumerate}
The companion proves these statements on the original physical history space; the blocked
snapshots need not form a new Markov chain.  It also proves that their physical transfer is
a reducing restriction of the original positive transfer.  These are inherited inputs,
not new claims below.

The present question is different from V28's raw-slice test.  Suppose an intended physical
reader is applied to this \emph{actual blocked history}.  Can it both retain the supplied
plaquette information and have the strongly continuous small-time behaviour demanded of a
classical or Hamiltonian continuum field?  The answer is quantitatively negative for any
reader whose plaquette distortion tends to zero.  The proof uses only
\eqref{v29:eq:var-input}--\eqref{v29:eq:increment-input}, so no Markov assumption,
independence, or reconstruction of the companion's transfer matrix is added.

This does not rule out a reader which deliberately erases the plaquette variable, a new
renormalized observable with a different normalization, or a different ultraviolet
coupling trajectory.  Those possibilities have to be specified and tested on their own
terms.  The point is to quantify the amount of information loss required before the
actual strong-coupling history can be called time-continuous or vacuum-concentrated.

\section{A source-level magnetic floor at every completed spatial depth}
\label{v29:sec:magnetic}

Let a blocked spatial slice contain $N^3$ sites and $3N^3$ positively oriented elementary
plaquettes, with declared coarse spatial spacing $\ell>0$.  For a plaquette put
$s=1-F\in[0,2]$ and use the same unsubtracted Wilson magnetic normalization as in the
constructed Hamiltonian branch,
\begin{equation}
 W_{N,\ell}=\frac4\ell\sum_{p=1}^{3N^3}s_p.
 \label{v29:eq:coarse-W}
\end{equation}
No independence of the plaquettes is assumed.

\begin{lemma}[Variance forces a positive mean plaquette action]
\label[lemma]{v29:lem:var-mean}
If $F\in[-1,1]$ and $\Var(F)\ge c$ with $0<c<1$, then
\begin{equation}
 \boxed{\quad \E(1-F)\ge \chi(c):=1-\sqrt{1-c}.\quad}
 \label{v29:eq:var-mean}
\end{equation}
The constant is sharp given only the support and variance information.
\end{lemma}
\begin{proof}
Write $m=\E F$.  Since $F^2\le1$,
$\Var(F)=\E F^2-m^2\le1-m^2$.  Hence
$m\le\sqrt{1-c}$ and \eqref{v29:eq:var-mean} follows.  Equality is possible for a
two-point $\{\pm1\}$ law with mean $\sqrt{1-c}$, proving sharpness at this level of
information.
\end{proof}

\begin{theorem}[All-depth coarse magnetic floor for the actual Perron history]
\label[theorem]{v29:thm:magnetic-floor}
Under the inherited strong-coupling hypotheses giving
\eqref{v29:eq:var-input}, every blocked law satisfies
\begin{equation}
 \boxed{\quad
  \E W_{N,\ell}\ge
    \frac{12N^3}{\ell}\,\chi(c_0),
  \qquad \chi(c_0)=1-\sqrt{\frac{3001}{4000}}>0.
 \quad}
 \label{v29:eq:magnetic-floor}
\end{equation}
Consequently these actual blocked laws do not satisfy the vacuum entrance
$\E W_{N,\ell}\to0$ at any fixed $(N,\ell)$, nor along a refinement on which
$N^3/\ell$ stays bounded below by a positive constant.  This is a statement about the
unsubtracted magnetic action \eqref{v29:eq:coarse-W}, not about a renormalized field.
\end{theorem}
\begin{proof}
Apply \cref{v29:lem:var-mean} separately to each of the $3N^3$ plaquettes and sum the
nonnegative expectations.  The uniform variance input gives the same lower bound for
every summand.  Multiply by $4/\ell$.
\end{proof}

Numerically,
\begin{equation}
 c_0=0.24975\ldots,\qquad
 \chi(c_0)=0.133830\ldots .
 \label{v29:eq:constants}
\end{equation}
These decimal values are for orientation only; all proofs use the exact rational $c_0$.

\section{A stable-reader lower bound for one-step physical increments}
\label{v29:sec:reader}

We now allow an arbitrary reader, including extra private randomness and complete source
history.  Work on one probability space carrying the original stationary source and the
reader output.  Let $Y_r(j)$ be a real decoded observable at fine time $j$ and let
$X_r(j)=F_r(j)-\E F_r$ be the centered actual blocked plaquette.  Stationarity is used only
to identify the same distortion at the two endpoints.

Define the root-mean-square reader distortion
\begin{equation}
 \delta_r^2=\E|Y_r(0)-X_r(0)|^2.
 \label{v29:eq:distortion}
\end{equation}
The output need not be a deterministic function of the current blocked snapshot.

\begin{theorem}[Retention--continuity tradeoff]
\label[theorem]{v29:thm:tradeoff}
For every such reader,
\begin{equation}
 \boxed{\quad
 \Bigl(\E|Y_r(1)-Y_r(0)|^2\Bigr)^{1/2}
       +2\delta_r>\sqrt{c_0}.
 \quad}
 \label{v29:eq:tradeoff}
\end{equation}
In particular:
\begin{enumerate}[label=\textup{(\roman*)},leftmargin=2.6em]
\item if $\delta_r\to0$, then
\begin{equation}
 \liminf_r\E|Y_r(1)-Y_r(0)|^2\ge c_0;
 \label{v29:eq:retained-jump}
\end{equation}
\item if the decoded observable is mean-square continuous on the original physical clock,
      in the sense that
      $\E|Y_r(1)-Y_r(0)|^2\to0$ while $a_r\to0$, then
\begin{equation}
 \boxed{\quad
       \liminf_r\delta_r^2\ge\frac{c_0}{4}
              =\frac{999}{16000}.
 \quad}
 \label{v29:eq:continuity-price}
\end{equation}
\end{enumerate}
The result is unchanged if the reader uses the full earlier history, terminal hidden
source state, or additional independent randomness.
\end{theorem}
\begin{proof}
In $L^2$ of the complete joint probability space,
\begin{align*}
 \|X_r(1)-X_r(0)\|_2
 &\le \|Y_r(1)-Y_r(0)\|_2\\
 &\quad+\|Y_r(1)-X_r(1)\|_2
       +\|Y_r(0)-X_r(0)\|_2.
\end{align*}
Stationarity gives the same bound $\delta_r$ at both endpoints.  Insert the inherited
strict lower bound \eqref{v29:eq:increment-input}.  The two consequences follow by taking
limits.  No conditional independence or Markov property is used.
\end{proof}

Thus a reader that accurately retains the actual coarse plaquette cannot repair the
companion's physical-time discontinuity by postprocessing.  Conversely a reader that
achieves mean-square continuity must discard an order-one amount of that observable in
the precise sense \eqref{v29:eq:continuity-price}.

\section{Vacuum concentration requires order-one plaquette information loss}
\label{v29:sec:vacuum-reader}

There is a separate, static version of the same issue.  Let $\widetilde F_r\in[-1,1]$ be
any decoded plaquette trace on the same joint probability space and set
$\widetilde s_r=1-\widetilde F_r$.  Define
\begin{equation}
 \varepsilon_r^2=\E|\widetilde F_r-F_r|^2.
 \label{v29:eq:static-distortion}
\end{equation}

\begin{proposition}[Vacuum-energy entrance has a fixed information price]
\label[proposition]{v29:prop:vacuum-price}
Under \eqref{v29:eq:var-input},
\begin{equation}
 \boxed{\quad
 \varepsilon_r\ge
       \sqrt{c_0}-\sqrt{2\E\widetilde s_r}.
 \quad}
 \label{v29:eq:vacuum-price}
\end{equation}
Consequently, if the decoded plaquette is vacuum-concentrated in the unsubtracted sense
$\E\widetilde s_r\to0$, then
\begin{equation}
 \boxed{\quad \liminf_r\varepsilon_r^2\ge c_0.\quad}
 \label{v29:eq:vacuum-price-limit}
\end{equation}
Thus an instantaneous reader cannot both reproduce this source plaquette in mean square
and turn it into a vacuum plaquette.
\end{proposition}
\begin{proof}
For the actual source,
\[
 \|F_r-1\|_2^2
   =\Var(F_r)+(1-\E F_r)^2\ge c_0.
\]
Since $0\le1-\widetilde F_r\le2$,
\[
 \|\widetilde F_r-1\|_2^2
  =\E(1-\widetilde F_r)^2\le2\E\widetilde s_r.
\]
The reverse triangle inequality gives
$\|F_r-\widetilde F_r\|_2\ge
 \|F_r-1\|_2-\|\widetilde F_r-1\|_2$, proving the result.
\end{proof}

The constants in \cref{v29:thm:tradeoff,v29:prop:vacuum-price} concern retention of the
\emph{actual blocked trace}.  They do not apply to a physically justified renormalized
observable which is not required to approximate that trace.  Such a reader, however, must
be specified as a new map with its own action incidence and physical-time calibration.

\section{A three-way incompatibility and its quantitative forms}
\label{v29:sec:trilemma}

\begin{theorem}[Reader trilemma for the actual blocked strong-coupling source]
\label[theorem]{v29:thm:trilemma}
Along the inherited isotropic fine-time refinement, no sequence of real reader observables
can satisfy all three properties:
\begin{enumerate}[label=\textup{(T\arabic*)},leftmargin=2.8em]
\item \emph{plaquette retention:}
      $\E|Y_r(j)-X_r(j)|^2\to0$ at $j=0,1$;
\item \emph{physical-time mean-square continuity:}
      $\E|Y_r(1)-Y_r(0)|^2\to0$ while $a_r\to0$;
\item \emph{vacuum retention of the same plaquette coordinate:} after adding back the
      source mean to form $\widetilde F_r$, one has
      $\E(1-\widetilde F_r)\to0$ and
      $\E|\widetilde F_r-F_r|^2\to0$.
\end{enumerate}
More quantitatively, enforcing \textup{(T2)} costs asymptotic squared distortion at least
$c_0/4$, while enforcing the vacuum statement in \textup{(T3)} costs at least $c_0$.
\end{theorem}
\begin{proof}
The incompatibility of \textup{(T1)} and \textup{(T2)} is
\cref{v29:thm:tradeoff}.  The two clauses inside \textup{(T3)} are incompatible by
\cref{v29:prop:vacuum-price}.  The quantitative constants are exactly
\eqref{v29:eq:continuity-price} and \eqref{v29:eq:vacuum-price-limit}.
\end{proof}

This theorem is deliberately reader-relative.  It does not say that the strong-coupling
source has no useful coarse description.  It says that a useful ultraviolet reader must
make an explicit choice: retain this order-one blocked plaquette, or smooth/renormalize it
enough to restore physical-time continuity or vacuum concentration.  That choice cannot
be hidden inside a statement that the reader is ``approximately the same'' observable.

\section{An inherited transfer gap becomes a reader-design condition}
\label{v29:sec:semigroup}

The companion also proves, on these same centered blocked history vectors $v_r$, the
transfer estimate
\begin{equation}
 0\le\langle v_r,e^{-tH_{a_r}}v_r\rangle
       \le2^{-t/a_r}\|v_r\|^2,
 \qquad \|v_r\|^2>c_0.
 \label{v29:eq:transfer-input}
\end{equation}
It follows that every fixed positive physical time has vanishing correlation while the
time-zero norm stays nonzero.  The companion already uses this to rule out a strongly
continuous nonzero semigroup limit for the plaquette itself.  We do not reprove that
statement.

The new reader estimate above shows that this conclusion is stable under vanishing
mean-square distortion.  In particular, if a reader $Y_r$ obeys
$\|Y_r-X_r\|_{L^2}\to0$ at each fixed source time, then its one-step modulus cannot tend
to zero.  Therefore any reader intended to generate a strongly continuous continuum
Hamiltonian must have a nonvanishing ultraviolet renormalization distance from this
plaquette bank.  A scalar amplitude rescaling that keeps a finite nonzero limiting norm is
not enough: the companion's spectral estimate already annihilates all positive-time
correlations at fixed norm.

This is not a claim that every possible nonlinear reader inherits the spectral bound
\eqref{v29:eq:transfer-input}; only the one-step mean-square statement is needed for the
quantitative trilemma.  A materially different reader must be analysed through its actual
joint laws and action insertions.

\section{Relation to the V28 energy entrance and the Einstein bridge}
\label{v29:sec:interfaces}

Version 28 proves a sufficient law-to-preparation entrance when a supplied configuration
law has
\[
             w_{\nu_h}+c_{\nu_h}+\eta_hh^{-4}\longrightarrow0.
\]
The current strong-coupling blocked source fails that particular vacuum entrance already
through \eqref{v29:eq:magnetic-floor}; spatial completion does not make the measured
plaquette variance vanish.  This statement is stronger in a different direction from
Version 28's logarithmic-coupling raw-slice obstruction: it concerns the companion's
\emph{actual Perron time-zero law} and every completed spatial depth, but only in its
explicit very-small-coupling basin.

There is no contradiction.  The two companion regimes are different, and neither has
been identified as the ultraviolet trajectory of a classical Yang--Mills field.  The
reader trilemma instead tells us what a successful trajectory must change: its coarse
plaquette bank, its physical clock, or both must scale in a way that restores continuity
without simply relabelling an order-one source observable as a small field.

The Exact-Action Einstein bridge remains a separate bottleneck.  Its supplied V16 analysis
proves that the retained positive scalar flux excludes the bare diagonal soft sector, but
still leaves the full low-source inverse response unresolved.  The present plaquette
reader inequalities do not bound that inverse or the geometric lapse output.  Conversely,
solving the Einstein inverse would not identify a Yang--Mills reader for the strong-coupling
source.  These two upstream obligations should remain separate.

\section{Outcome and next noncircular target}
\label{v29:sec:outcome}

The new information is not another constructed gauge trajectory.  It is a restriction on
what an \emph{actual-source reader} can legitimately claim:
\[
 \boxed{
 \begin{gathered}
 \text{accurately retain the actual blocked plaquette}\
 \Longrightarrow\text{retain an order-one fine-time increment};\\[0.2em]
 \text{force fine-time continuity}\
 \Longrightarrow\text{pay at least }c_0/4\text{ in squared plaquette distortion};\\[0.2em]
 \text{force vacuum concentration of that same coordinate}\
 \Longrightarrow\text{pay at least }c_0\text{ in squared distortion}.
 \end{gathered}}
\]
All constants are inherited from one proved source variance $c_0=999/4000$.

The next useful acquisition is therefore the \emph{actual intended coarse/renormalized
physical reader}, including its physical spatial normalization, its source-to-reader map,
its action insertion, and its one-step physical-time law.  Once that map is supplied, the
Version 24 one-sided action-work theorem and the Version 28 law-to-state entrance give
finite tests.  Without it, extending either the strong-coupling terminal basin or the
restricted weak-coupling raw coordinates would remain circular.

No physical Yang--Mills mass gap, continuum construction, Einstein inverse, chiral sector,
or primitive Renewal-to-Einstein--SM theorem is claimed in this body.

\paragraph{Verification and preservation.}
Every Version 28 byte before its final document command is preserved.  New labels have
prefix \texttt{v29:}.  The companion script verifies the exact source constant, the
variance-to-mean inequality on an extremizing two-point law, all reader-tradeoff constants,
the magnetic-floor coefficient, and randomized finite examples saturating the triangle
bounds.  These are diagnostics of the written arguments, not proof-assistant
formalization or a simulation of the companion's full Wilson source.

\begingroup
\renewcommand{\refname}{Additional source for Version 29}
\begin{thebibliography}{9}
\small\raggedright
\bibitem{v29:YM86} A.~P\'elissier,
\emph{Renewal Yang--Mills Mass-Gap Research Notes}, supplied cumulative V86,
file \nolinkurl{Renewal_Yang_Mills_Mass_Gap_Research_Notes_V86(3).tex}.
The inherited inputs used here are the V52--V53 actual Perron spatial history, its
all-depth plaquette-variance lower bound, and its fine-time increment/transfer obstruction.
No other theorem of that cumulative manuscript is silently imported.
\end{thebibliography}
\endgroup
\addtocontents{toc}{\protect\endgroup}
% END IMMUTABLE VERSION 29 BODY
% APPEND VERSION 30 HERE, BEFORE THE FINAL END-DOCUMENT COMMAND.


% BEGIN IMMUTABLE VERSION 30 BODY
\clearpage
\begin{center}
{\Large\bfseries Version 30: The observed Schur impedance of the Einstein low-source bank}\\[0.5em]
{\large A finite source-specific range criterion, certified high-block lifting, and an initial geometric closure test}
\end{center}
\makeatletter
\addtocontents{toc}{\protect\begingroup}
\addtocontents{toc}{\protect\makeatletter}
\addtocontents{toc}{\protect\def\protect\l@section{\protect\bfseries\protect\@dottedtocline{1}{0em}{2.5em}}}
\makeatother
\addcontentsline{toc}{part}{Version 30: Source-specific observed Schur impedance for the exact-action Einstein bridge}

\section{Audit: leave the Yang--Mills reader obstruction and attack the remaining Einstein range problem}
\label{v30:sec:audit}

Version 29 makes another blocking or reader-distortion estimate on the same Yang--Mills
source progressively less informative.  The Exact-Action Einstein companion identifies a
different bottleneck: after the shift variables are eliminated, the leading active lapse
satisfies one positive scalar equation, its exponentially small high-frequency \emph{source}
tail is already controlled through the complete inverse and the geometric observation, and
the unresolved part is the response of a logarithmically sized low-source bank through the
\emph{full} scalar inverse \cite{v23:EB16}.  We therefore change direction here.

The following four facts are inherited from that companion and are not reproved.
They are stated to fix the interfaces used below.
\begin{enumerate}[label=\textup{(E\arabic*)},leftmargin=2.7em]
\item On the mean-zero finite scalar space, the unchanged positive operator
      $\mathscr S_h$ and the lapse-to-curvature observation
      \begin{equation}
        \mathcal O_h=\mathfrak E_hT_h^\dagger V_h
        \label{v30:eq:O}
      \end{equation}
      are the actual operators obtained after the exact transverse elimination.
\item For $P=P_{\le R}$ and $Q=I-P$, with blocks
      \begin{equation}
        A=P\mathscr S_hP,\qquad C=Q\mathscr S_hP,\qquad D=Q\mathscr S_hQ,
        \label{v30:eq:blocks}
      \end{equation}
      one has $D\succ0$ and the exact low-source Schur objects
      \begin{equation}
        K=A-C^*D^{-1}C,\qquad Zx=x-D^{-1}Cx,
        \label{v30:eq:KZ}
      \end{equation}
      with $K\succeq c h^{12}I$ on $\Ran P$.
\item If $g_h$ is the actual scalar source, then for every fixed output Sobolev order
      $s$ and every prescribed $p>0$ there is
      $R_h=O(\log(1/h))$ such that
      \begin{equation}
        \|\mathcal O_h\mathscr S_h^{-1}Q_h^{\rm src}g_h\|_{H_h^s}
            \le C_s h^p,
        \qquad Q_h^{\rm src}:=I-P_{\le R_h}.
        \label{v30:eq:high-tail}
      \end{equation}
      The solution generated by the retained low source is \emph{not} frequency truncated.
\item The direct transverse part of the prepared geometric output is $O(h)$ in every
      fixed Sobolev norm.  What remains is exactly the lapse-mediated observation in
      \eqref{v30:eq:O}.
\end{enumerate}
The companion additionally proves that the bare diagonal soft sector is not an admissible
substitute for the full scalar inverse: the retained positive flux excludes its nonzero
limiting vectors from the finite-energy domain.  We keep that conclusion as a boundary.
Nothing below diagonalizes the bare product-defect operator or deletes the high block.

The new question is narrower.  Can the \emph{specific} low source be tested against the
complete inverse without demanding a cutoff-uniform observation estimate for every scalar
right side?  The answer is yes: the exact response is encoded by one positive matrix on the
low-source space.  Its dimension is only $O((\log N)^3)$ for the inherited choice of $R_h$,
while every high-mode relaxation remains inside its entries.

\section{The observed Schur impedance and its exact source spectral measure}
\label{v30:sec:impedance}

We first state the finite Hilbert-space result independently of the particular coefficients.
Let $\mathcal H=P\mathcal H\oplus Q\mathcal H$ be a finite Hilbert space,
$S=S^*\succ0$, and let $A,C,D,K,Z$ be defined by
\eqref{v30:eq:blocks}--\eqref{v30:eq:KZ} with $S$ in place of $\mathscr S_h$.
Let $\mathcal Y$ be another finite Hilbert space and $O:\mathcal H\to\mathcal Y$ any
linear observation.  Put
\begin{equation}
  B:=OZ:P\mathcal H\to\mathcal Y,\qquad
  G:=B^*B\succeq0,
  \label{v30:eq:BG}
\end{equation}
and define the \emph{observed Schur impedance}
\begin{equation}
  \boxed{\quad J:=K^{-1/2}GK^{-1/2}\succeq0.\quad}
  \label{v30:eq:J}
\end{equation}
This is not a new action.  It is a diagnostic built from the unchanged positive solve and
the independently specified observation.

\begin{theorem}[Exact source-specific observed Schur identity]
\label[theorem]{v30:thm:schur-identity}
For every $b\in P\mathcal H$, let $u=S^{-1}b$ and
$q=K^{-1/2}b$.  Then
\begin{align}
 u&=ZK^{-1}b,\label{v30:eq:u}\\
 \langle u,Su\rangle&=\langle b,K^{-1}b\rangle=\|q\|^2,\label{v30:eq:energy}\\
 \boxed{\quad\|Ou\|_{\mathcal Y}^2=\langle q,Jq\rangle.\quad}
 \label{v30:eq:output}
\end{align}
In particular, a large worst-case norm $\|J\|$ does not by itself imply a large response of
the specified source; only the spectral measure of $J$ in the vector $q$ enters.
\end{theorem}
\begin{proof}
For $x\in P\mathcal H$, block multiplication gives
\[
 S Zx=
 \begin{pmatrix}A&C^*\\ C&D\end{pmatrix}
 \binom{x}{-D^{-1}Cx}
 =\binom{Kx}{0}.
\]
Thus $SZ=\iota_PK$, where $\iota_P$ is the inclusion of $P\mathcal H$, and
$S^{-1}b=ZK^{-1}b$.  If $x=K^{-1}b$, then
\[
 \langle Zx,SZx\rangle=\langle x,Kx\rangle
       =\langle b,K^{-1}b\rangle.
\]
Finally $Ou=BK^{-1}b=BK^{-1/2}q$, so
\[
 \|Ou\|^2
 =\langle q,K^{-1/2}B^*BK^{-1/2}q\rangle
 =\langle q,Jq\rangle.
\]
No estimate of $D^{-1}$ or of $S^{-1}$ was used in these identities.
\end{proof}

The same information can be written without a matrix square root.

\begin{corollary}[Generalized-eigenvalue source decomposition]
\label[corollary]{v30:cor:pencil}
There is a $K$-orthonormal basis $v_1,\ldots,v_m$ of $P\mathcal H$ and numbers
$\lambda_j\ge0$ satisfying
\begin{equation}
  Gv_j=\lambda_jKv_j,\qquad
  \langle v_i,Kv_j\rangle=\delta_{ij}.
  \label{v30:eq:pencil}
\end{equation}
For $\alpha_j=\langle v_j,b\rangle$,
\begin{equation}
 \boxed{\quad
  \langle b,K^{-1}b\rangle=\sum_j|\alpha_j|^2,
  \qquad
  \|OS^{-1}b\|^2=\sum_j\lambda_j|\alpha_j|^2.
 \quad}
 \label{v30:eq:pencil-output}
\end{equation}
Hence for a cutoff family and a specified source $b_h$, the bound
$\|O_hS_h^{-1}b_h\|=O(h)$ is equivalent to
\begin{equation}
             \sum_j\lambda_{j,h}|\alpha_{j,h}|^2=O(h^2).
             \label{v30:eq:iff}
\end{equation}
\end{corollary}
\begin{proof}
Diagonalize $J$ and put $v_j=K^{-1/2}e_j$ for an orthonormal eigenbasis $e_j$.
Then \eqref{v30:eq:pencil} holds.  If $x=K^{-1}b$, its $K$-orthonormal expansion is
$x=\sum_j\alpha_jv_j$, because
$\langle v_j,Kx\rangle=\langle v_j,b\rangle$.  The two quadratic forms of
\cref{v30:thm:schur-identity} give \eqref{v30:eq:pencil-output}.
\end{proof}

For $\Lambda>0$ define the source energy above impedance $\Lambda$ by
\[
  M_{>\Lambda}(b)=\sum_{\lambda_j>\Lambda}|\alpha_j|^2.
\]
Then, writing $E_b=\langle b,K^{-1}b\rangle$ and
$\lambda_{\max}=\|J\|$,
\begin{equation}
  \boxed{\quad
  \Lambda M_{>\Lambda}(b)
  \le\|OS^{-1}b\|^2
  \le\Lambda E_b+\lambda_{\max}M_{>\Lambda}(b).
  \quad}
  \label{v30:eq:tail-impedance}
\end{equation}
Thus a divergent worst observed impedance is compatible with a small particular response
only if the normalized source loses its mass in those divergent directions.  This is a
source-specific range statement, not an ellipticity claim for arbitrary residuals.

\section{Only the low bank must be probed, but every high mode remains in the probe}
\label{v30:sec:finite-bank}

Return to the inherited Exact-Action blocks.  Let $e_1,\ldots,e_m$ be any orthonormal
basis of $\Ran P$, with $m=\operatorname{rank}P$.  Define the exact high lifts
\begin{equation}
  y_a=D^{-1}Ce_a,\qquad z_a=e_a-y_a=Ze_a.
  \label{v30:eq:lifts}
\end{equation}
Then
\begin{equation}
  K_{ab}=\langle e_a,Ae_b\rangle-\langle Ce_a,y_b\rangle,
  \qquad
  G_{ab}=\langle Oz_a,Oz_b\rangle_{\mathcal Y}.
  \label{v30:eq:matrices}
\end{equation}

\begin{theorem}[Exact finite low-source acquisition]
\label[theorem]{v30:thm:finite-acquisition}
The complete response of every source in $\Ran P$ is determined by the $m$ high-block
solves
\begin{equation}
                    Dy_a=Ce_a\qquad(1\le a\le m),
                    \label{v30:eq:high-solves}
\end{equation}
together with the finite matrices in \eqref{v30:eq:matrices}.  No principal-block inverse
of $S$ appears.  In particular, when the inherited radius $R_h$ is used,
\begin{equation}
                     m=O(R_h^3)=O((\log N)^3).
                     \label{v30:eq:dimension}
\end{equation}
Every high Fourier mode of the exact solution may contribute to the vectors $y_a$ and
$Oz_a$; none is deleted by this reduction.
\end{theorem}
\begin{proof}
Equations \eqref{v30:eq:lifts}--\eqref{v30:eq:matrices} are exactly the columns of $Z$,
$K$, and $G$.  Once they are known, \cref{v30:thm:schur-identity} computes the response
for any coefficient vector of $b\in\Ran P$.  The dimension count is the inherited count
of real nonconstant Fourier modes in the three-dimensional $\ell^1$ ball of radius $R_h$.
The equation for $y_a$ is solved on the original $Q$ space, so its output contains any
high-mode scattering generated by the full coefficients.
\end{proof}

This theorem is an exact reduction of the \emph{number of right sides}, not of the
high-block operator.  It is therefore compatible with the companion's warning that the
low-frequency principal block is not the full inverse.

\section{Residual-certified construction of the impedance matrix}
\label{v30:sec:certified}

The preceding theorem still contains exact high solves.  We now record an a posteriori
certificate which permits interval or rational residual bounds without mistaking a floating
solve for a proof.

Let $\widetilde y_a\in Q\mathcal H$ be approximate lifts and define
\begin{equation}
  r_a=D\widetilde y_a-Ce_a,
  \qquad
  \varepsilon_R:=\left(\sum_{a=1}^m\|r_a\|^2\right)^{1/2}.
  \label{v30:eq:residuals}
\end{equation}
Let $\widetilde Z e_a=e_a-\widetilde y_a$, and form
\begin{align}
 \widetilde K_{ab}&=\langle e_a,Ae_b\rangle-
                         \langle Ce_a,\widetilde y_b\rangle,\nonumber\\
 \widetilde B e_a&=O\widetilde Ze_a,
 \qquad \widetilde G=\widetilde B^*\widetilde B.
 \label{v30:eq:tilde}
\end{align}

\begin{proposition}[High-solve residuals give explicit matrix enclosures]
\label[proposition]{v30:prop:lift-error}
Put $d_D=\|D^{-1}\|$, $c_C=\|C\|$, and $c_O=\|O\|$.  Then
\begin{align}
 \|\widetilde Z-Z\|&\le d_D\varepsilon_R,\label{v30:eq:Zerr}\\
 \|\widetilde K-K\|&\le c_Cd_D\varepsilon_R,\label{v30:eq:Kerr}\\
 \|\widetilde B-B\|&\le c_Od_D\varepsilon_R,\label{v30:eq:Berr}\\
 \|\widetilde G-G\|&\le
   c_Od_D\varepsilon_R
   \bigl(2\|\widetilde B\|+c_Od_D\varepsilon_R\bigr).
   \label{v30:eq:Gerr}
\end{align}
Here $\varepsilon_R$ may be replaced by the operator norm of the residual column map if
that sharper quantity is certified.
\end{proposition}
\begin{proof}
Let $Y e_a=D^{-1}Ce_a$ and $\widetilde Y e_a=\widetilde y_a$.  The residual column map
$R= D\widetilde Y-C$ gives
$\widetilde Y-Y=D^{-1}R$.  Its operator norm is at most its Hilbert--Schmidt norm, which
is bounded by $d_D\varepsilon_R$.  Since $Z=I-Y$ and
$K=A-C^*Y$, the first two inequalities follow.  The third is obtained by applying $O$.
Finally
\[
 G-\widetilde G=(B-\widetilde B)^*\widetilde B+
 \widetilde B^*(B-\widetilde B)+(B-\widetilde B)^*(B-\widetilde B),
\]
which proves \eqref{v30:eq:Gerr}.
\end{proof}

For the companion's scalar family one may take the inherited coarse estimate
$d_D\le Ch^{-12}$ and, at output order $s=0$, the inherited full-band observation bound
$c_O\le Ch^{-1}$.  The resulting residual tolerance can therefore be severe.  The point
of \cref{v30:prop:lift-error} is not to claim numerical ease; it states exactly how a
finite verification must pay that severity.

The next lemma converts matrix enclosures into a certified output interval.

\begin{lemma}[Perturbation bound for the source-specific response]
\label[lemma]{v30:lem:perturb}
Suppose $K\succeq\kappa I$, $\|\widetilde K-K\|\le\varepsilon_K<\kappa$,
$\|\widetilde G-G\|\le\varepsilon_G$, and
$\|\widetilde b-b\|\le\varepsilon_b$.  Put
\begin{align*}
 x&=K^{-1}b,\qquad \widetilde x=\widetilde K^{-1}\widetilde b,\\
 R_x&=\frac{\varepsilon_b+\varepsilon_K\|b\|/\kappa}
                 {\kappa-\varepsilon_K},\\
 X&=\|b\|/\kappa,
 \qquad \widetilde X=(\|b\|+\varepsilon_b)/(\kappa-\varepsilon_K),\\
 M_G&=\|\widetilde G\|+\varepsilon_G.
\end{align*}
Then $\|x-\widetilde x\|\le R_x$ and
\begin{equation}
 \boxed{\quad
 \left|\,x^*Gx-\widetilde x^*\widetilde G\widetilde x\,\right|
 \le M_G(X+\widetilde X)R_x+\varepsilon_G\widetilde X^2.
 \quad}
 \label{v30:eq:perturb}
\end{equation}
Thus a strict certified upper bound on the right-hand finite quadratic form, plus the
explicit error in \eqref{v30:eq:perturb}, is a rigorous upper bound on the actual observed
low-source response.
\end{lemma}
\begin{proof}
The identity
\[
 \widetilde K(\widetilde x-x)
   =(\widetilde b-b)+(K-\widetilde K)x
\]
and $\|\widetilde K^{-1}\|\le(\kappa-\varepsilon_K)^{-1}$ prove the first bound.
Next split
\[
 x^*Gx-\widetilde x^*\widetilde G\widetilde x
 =(x-\widetilde x)^*Gx+
   \widetilde x^*G(x-\widetilde x)+
   \widetilde x^*(G-\widetilde G)\widetilde x.
\]
Use $\|G\|\le M_G$ and the displayed bounds for $x,\widetilde x$.
\end{proof}

This is intentionally an \emph{a posteriori} theorem.  It permits exact rational, interval,
or independently residual-certified solves.  It does not promote an unverified floating
matrix inversion to a theorem.

\section{A finite source-specific Einstein certificate}
\label{v30:sec:certificate}

Apply the preceding construction with the inherited operators
$S=\mathscr S_h$, $O=\mathcal O_h$, and with output Hilbert norm $H_h^s$.
For the inherited logarithmic radius $R_h$ write
\begin{equation}
 b_h=P_{\le R_h}g_h,
 \qquad
 \boxed{\quad
 \mathfrak C_{h,s}:=
   b_h^*K_{R_h}^{-1}G_{R_h,s}K_{R_h}^{-1}b_h.
 \quad}
 \label{v30:eq:certificate}
\end{equation}
Here $G_{R_h,s}$ is formed with the $H_h^s$ inner product on
$\mathcal O_hZ_{R_h}\Ran P_{\le R_h}$.  Every object is finite at each cutoff and the
high block is retained through $Z_{R_h}$.

\begin{theorem}[Necessary and sufficient low-source certificate for the leading geometric observation]
\label[theorem]{v30:thm:certificate}
For the actual prepared scalar source of the Exact-Action bridge,
\begin{equation}
 \boxed{\quad
 \|\mathcal O_h\mathscr S_h^{-1}P_{\le R_h}g_h\|_{H_h^s}^2
       =\mathfrak C_{h,s}.
 \quad}
 \label{v30:eq:certificate-exact}
\end{equation}
Consequently
\begin{equation}
 \|\mathcal O_h\mathscr S_h^{-1}P_{\le R_h}g_h\|_{H_h^s}=O(h)
 \quad\Longleftrightarrow\quad
 \mathfrak C_{h,s}=O(h^2).
 \label{v30:eq:certificate-iff}
\end{equation}
The condition is source-specific.  It is strictly weaker than a uniform
$\|\mathcal O_h\mathscr S_h^{-1}\|=O(1)$ estimate, which earlier quasimodes exclude in
related active spaces.
\end{theorem}
\begin{proof}
Equation \eqref{v30:eq:certificate-exact} is
\cref{v30:thm:schur-identity} applied to the exact low-source interface.  Taking square
roots gives the equivalence.  No statement about other right sides follows from it.
\end{proof}

The generalized pencil in \cref{v30:cor:pencil} gives an equivalent finite audit.  If
\[
 G_{R_h,s}v_{j,h}=\lambda_{j,h}K_{R_h}v_{j,h},
 \qquad \langle v_{i,h},K_{R_h}v_{j,h}\rangle=\delta_{ij},
\]
then
\begin{equation}
 \boxed{\quad
 \mathfrak C_{h,s}
   =\sum_j\lambda_{j,h}
       |\langle v_{j,h},b_h\rangle|^2.
 \quad}
 \label{v30:eq:source-spectrum}
\end{equation}
This explicitly separates a possible growing observed impedance from the actual source
weight carried by its corresponding directions.

\section{Composition with the inherited high-source tail}
\label{v30:sec:closure}

The previous section concerns only the low source.  The inherited inverse-tail theorem
already controls the complement through the complete inverse.  Combining the two gives a
clean sufficient condition for the full leading initial observation.

\begin{corollary}[From one finite certificate to an $O(h)$ full leading curvature output]
\label[corollary]{v30:cor:curvature}
Fix a Sobolev order $s\ge0$ and choose the inherited $R_h$ with high-source exponent
$p>1$.  If
\begin{equation}
                         \mathfrak C_{h,s}\le C^2 h^2,
                         \label{v30:eq:certificate-target}
\end{equation}
then the complete lapse-mediated geometric observation satisfies
\begin{equation}
 \|\mathcal O_h\mathscr S_h^{-1}g_h\|_{H_h^s}\le C'_s h.
 \label{v30:eq:full-low-high}
\end{equation}
Adding the already controlled direct transverse contribution gives the prepared leading
active geometric output
\begin{equation}
                         \|\mathfrak E_h(\beta_h)\|_{H_h^s}
                              \le C''_s h.
                         \label{v30:eq:Einstein-leading}
\end{equation}
Thus the unresolved initial Einstein-response row on this selected family is reduced to
verifying the finite positive number \eqref{v30:eq:certificate}; no uniform inverse theorem
for arbitrary scalar sources is required.
\end{corollary}
\begin{proof}
Split $g_h=P_{\le R_h}g_h+Q_h^{\rm src}g_h$.  The first term has observation at most
$Ch$ by \eqref{v30:eq:certificate-exact} and
\eqref{v30:eq:certificate-target}; the second is $O(h^p)=o(h)$ by
\eqref{v30:eq:high-tail}.  This proves \eqref{v30:eq:full-low-high}.
The exact active reconstruction is the inherited direct term minus
$T_h^\dagger V_h\mathscr S_h^{-1}g_h$.  Apply $\mathfrak E_h$ and use inherited item
\textup{(E4)}.
\end{proof}

The conclusion is deliberately limited.  It is a bound for the \emph{leading prepared
initial active observation}.  It does not establish the higher weak rows, a nonlinear
exact-action lifespan, or the full observed Einstein inverse away from this selected source
family.  In the bridge's finite-amplitude metric expansion it can be inserted into the
already derived curvature formula, but the independent amplitude and evolution remainders
remain exactly where that companion leaves them.

\section{Why low source dimension is not enough}
\label{v30:sec:counterexample}

The finite criterion is useful precisely because the dimension count by itself gives no
bound.  A one-dimensional example makes the quantifiers explicit.  Let
\[
 K_h=h^{12},\qquad G_h=h^{-10},\qquad b_h=h.
\]
Then the low-source space has dimension one and
\[
 \langle b_h,K_h^{-1}b_h\rangle=h^{-10},
 \qquad
 \|OS^{-1}b_h\|^2=b_h^2K_h^{-2}G_h=h^{-32}.
\]
Changing $G_h$ to $h^{24}$ while leaving the same dimension, source size and Schur floor
gives an output of order $h^2$.  The distinction is exactly the observed impedance.

This example is not a model of the Exact-Action coefficients.  It records a logical
point: neither $m_R=O((\log N)^3)$ nor $K_R\succeq ch^{12}I$ determines the required
source-specific output.  The finite matrix \eqref{v30:eq:certificate} is therefore a real
additional certificate, not notation for a consequence already proved in the companion.

\section{Outcome: the next calculation is finite and source-specific}
\label{v30:sec:outcome}

The Yang--Mills reader programme has now supplied several explicit rejection and entrance
criteria.  The Exact-Action Einstein bridge had a different remaining issue: a particular
analytic low source might still scatter through high scalar modes before being observed.
Version 30 reduces that question exactly to
\[
 \boxed{
   \mathfrak C_{h,s}
      =b_h^*K_{R_h}^{-1}G_{R_h,s}K_{R_h}^{-1}b_h,
   \qquad \dim K_{R_h}=O((\log N)^3).
 }
\]
The matrices $K_{R_h}$ and $G_{R_h,s}$ retain the full high-block inverse through the lifted
basis vectors.  A verified $O(h^2)$ bound for this one positive number gives the missing
$O(h)$ leading geometric response; a lower bound larger than that scale rejects the branch.
The generalized eigen-expansion \eqref{v30:eq:source-spectrum} identifies exactly which
observed high-impedance directions are actually populated by the prepared source.

This is where I would continue rather than return to another generic compactness theorem:
construct or certify the $O((\log N)^3)$ lifted columns for the companion's actual rational
family, using the residual bounds of \cref{v30:prop:lift-error,v30:lem:perturb}, and test
\eqref{v30:eq:certificate-target}.  That calculation can succeed even if the worst
observation norm diverges, and it can fail for a concrete source without making a false
claim about every residual.

No new nonlinear Einstein lifespan, primitive-source selection, Yang--Mills mass gap, or
universal microscopic Renewal implication is asserted.  The advance is an exact,
source-specific finite certificate for the upstream Einstein row which remained open after
the companion's high-source and flux analyses.

\paragraph{Verification and preservation.}
Every Version 29 byte before its final end-document command is retained.  The diagnostic
script checks the block Schur identity, generalized eigen expansion, spectral-tail bounds,
high-lift residual enclosures and perturbation inequality on deterministic positive matrix
banks.  It also checks the one-dimensional quantifier counterexample.  These checks are
regressions for the written algebra; they are not a proof-assistant formalization or a
numerical certificate for the companion's still-uncomputed actual low-source matrices.
New labels use the prefix \texttt{v30:}.

\addtocontents{toc}{\protect\endgroup}
% END IMMUTABLE VERSION 30 BODY
% APPEND VERSION 31 HERE, BEFORE THE FINAL END-DOCUMENT COMMAND.



% BEGIN IMMUTABLE VERSION 31 BODY
\clearpage
\begin{center}
{\Large\bfseries Version 31: The continuum flux-range gate behind the Einstein impedance}\\[0.5em]
{\large A source-angle criterion, an explicit Hodge compensation equation, and finite witness certificates without the singular scalar inverse}
\end{center}
\makeatletter
\addtocontents{toc}{\protect\begingroup}
\addtocontents{toc}{\protect\makeatletter}
\addtocontents{toc}{\protect\def\protect\l@section{\protect\bfseries\protect\@dottedtocline{1}{0em}{2.5em}}}
\makeatother
\addcontentsline{toc}{part}{Version 31: Continuum flux-range and source-angle gate for the exact-action Einstein lapse}

\section{Why the finite impedance certificate has a more primitive range gate}
\label{v31:sec:audit}

Version 30 reduces the unresolved leading Einstein observation to the finite positive number
\[
 \mathfrak C_{h,s}=b_h^*K_{R_h}^{-1}G_{R_h,s}K_{R_h}^{-1}b_h,
 \qquad \dim K_{R_h}=O((\log N)^3),
\]
while retaining the complete high-mode relaxation inside the lifted columns.  That is an exact
finite-cutoff criterion, but it does not by itself explain which continuum geometry makes the
actual source safe or dangerous.  We therefore go one level upstream and use the normalized
coefficient equation already isolated in the Exact-Action bridge \cite{v23:EB16}.

The inherited continuum data are as follows.  On the real off-diagonal coefficient Hilbert
space let $\overline T$ be the normalized transverse source and let $\overline\Pi$ be the
orthogonal projection onto $\overline{\Ran\overline T}$.  The normalized lapse source is
$\overline V$.  On smooth mean-zero scalar tests,
\begin{equation}
 \boxed{\quad
 \mathscr S^\circ=\tfrac12 R^{\circ *}R^\circ,
 \qquad R^\circ=(I-\overline\Pi)\overline V .\quad}
 \label{v31:eq:R}
\end{equation}
The actual leading forcing is
\begin{equation}
 \boxed{\quad
 g^\circ=2\overline V^*p^\circ,
 \qquad p^\circ=M^\circ w^\circ .\quad}
 \label{v31:eq:g-p}
\end{equation}
The inherited identity $\overline T\,\overline\beta=2M^\circ w^\circ$ gives
\begin{equation}
                 p^\circ\in\Ran\overline T=\Ran\overline\Pi .
 \label{v31:eq:p-transverse}
\end{equation}
Thus the source is generated by the component of $\overline V$ which lies in the transverse
range, whereas the limiting positive scalar energy measures the orthogonal residual
$(I-\overline\Pi)\overline V$.  This is the structural mismatch which a generic elliptic
estimate cannot remove.

The bridge already proves that if $u_h=h n_h$ is bounded and has a weak limit, then that limit
solves $\mathscr S^\circ n^\circ=g^\circ$ distributionally and has finite limiting flux.  The
new point below is to characterize that finite-energy solvability exactly, without first
inverting $\mathscr S^\circ$.  The criterion is a source-weighted angle between the two pieces
of $\overline V$.

\section{A general finite-energy range theorem for a positive square}
\label{v31:sec:range}

Let $X_0$ be a real vector space densely embedded in a Hilbert distribution pairing, let
$Y$ be a real Hilbert space, and let $R:X_0\to Y$ be linear.  We do not assume that $\Ran R$
is closed.  Give $X_0/\ker R$ the energy norm
\begin{equation}
      \|\phi\|_{\mathcal E}^2=\tfrac12\|R\phi\|_Y^2
 \label{v31:eq:energy-norm}
\end{equation}
and denote its completion by $\mathcal E_R$.  A linear functional $g$ on $X_0$ has extended
dual norm
\begin{equation}
 \|g\|_{R,*}:=
 \sqrt2\sup_{R\phi\ne0}\frac{|\langle g,\phi\rangle|}{\|R\phi\|_Y},
 \label{v31:eq:dual-norm}
\end{equation}
with value $+\infty$ if $g$ does not annihilate $\ker R$.

\begin{theorem}[Finite-energy range theorem without a closed-range hypothesis]
\label[theorem]{v31:thm:range}
For a functional $g$ on $X_0$, the following are equivalent.
\begin{enumerate}[label=\textup{(\roman*)},leftmargin=2.7em]
\item $\|g\|_{R,*}<\infty$.
\item There is a unique $y_g\in\overline{\Ran R}$ such that
\begin{equation}
        \langle g,\phi\rangle=\tfrac12\langle y_g,R\phi\rangle_Y
        \qquad(\phi\in X_0).
 \label{v31:eq:factor}
\end{equation}
\item The weak equation
\begin{equation}
          \tfrac12R^*Rn=g
 \label{v31:eq:weak-square}
\end{equation}
has a unique solution class $n\in\mathcal E_R$, modulo $\ker R$.
\end{enumerate}
In that case
\begin{equation}
 \boxed{\quad
     Rn=y_g,\qquad
     \tfrac12\|Rn\|_Y^2
       =\langle g,n\rangle
       =\|g\|_{R,*}^2
       =\tfrac12\|y_g\|_Y^2 .\quad}
 \label{v31:eq:range-energy}
\end{equation}
\end{theorem}
\begin{proof}
The map $X_0/\ker R\to\overline{\Ran R}$ induced by $R$ is an isometry after multiplying
the target norm by $1/\sqrt2$.  A functional is continuous for
\eqref{v31:eq:energy-norm} exactly when it descends through this map.  The Riesz theorem on
$\overline{\Ran R}$ then gives a unique $y_g$ and the factor $1/2$ in
\eqref{v31:eq:factor}.  Conversely that factorization immediately bounds the dual norm.
The weak equation means precisely
$\langle g,\phi\rangle=\langle Rn,R\phi\rangle/2$ for every test.  Uniqueness in the
completion gives $Rn=y_g$.  Substitution into the weak equation and the definition of the
dual norm prove all identities in \eqref{v31:eq:range-energy}.  No inverse of $R^*R$ and no
closed-range assertion is used.
\end{proof}

This theorem separates two logically different questions.  The finite-cutoff matrices are
invertible, whereas their continuum energy form can have nonclosed range.  Finite-energy
solvability is therefore a statement about the source functional on the energy completion,
not about a pointwise inverse of the limiting differential expression.

\section{The actual Einstein source is governed by one transverse-to-residual angle}
\label{v31:sec:angle}

Apply \cref{v31:thm:range} with $R=R^\circ$ from \eqref{v31:eq:R}.  Put
$Q^\circ_o=I-\overline\Pi$ on the off-diagonal coefficient space.  Because
$p^\circ\in\Ran\overline\Pi$, define the extended source-angle number
\begin{equation}
 \boxed{\quad
 \Gamma^\circ:=
 \sup_{Q_o^\circ\overline V\phi\ne0}
 \frac{|\langle p^\circ,\overline\Pi\overline V\phi\rangle|}
      {\|Q_o^\circ\overline V\phi\|}
 \quad}
 \label{v31:eq:Gamma}
\end{equation}
with $\Gamma^\circ=+\infty$ if some test has zero denominator and nonzero numerator.

\begin{theorem}[Exact source-angle criterion for the normalized Einstein lapse]
\label[theorem]{v31:thm:source-angle}
For the inherited continuum coefficients, the following are equivalent.
\begin{enumerate}[label=\textup{(\alph*)},leftmargin=2.7em]
\item The equation $\mathscr S^\circ n^\circ=g^\circ$ has a finite-energy solution in
      $\mathcal E_{R^\circ}$.
\item $\Gamma^\circ<\infty$.
\item There exists $q\in\Ran Q_o^\circ$ such that
\begin{equation}
                 \overline V^*(q-p^\circ)=0 .
 \label{v31:eq:compensation}
\end{equation}
Equivalently,
\begin{equation}
 \boxed{\quad
 (p^\circ+\ker\overline V^*)\cap\Ran Q_o^\circ\ne\varnothing .\quad}
 \label{v31:eq:affine-intersection}
\end{equation}
\end{enumerate}
When these conditions hold there is a unique minimum-norm
$q_*\in\overline{\Ran R^\circ}$ satisfying \eqref{v31:eq:compensation}, and
\begin{equation}
 \boxed{\quad
 \|q_*\|=\Gamma^\circ,\qquad
 R^\circ n^\circ=4q_*,\qquad
 \tfrac12\|R^\circ n^\circ\|^2=8(\Gamma^\circ)^2 .\quad}
 \label{v31:eq:angle-energy}
\end{equation}
\end{theorem}
\begin{proof}
For a smooth mean-zero test $\phi$, \eqref{v31:eq:g-p} and
\eqref{v31:eq:p-transverse} give
\begin{equation}
 \langle g^\circ,\phi\rangle
  =2\langle p^\circ,\overline V\phi\rangle
  =2\langle p^\circ,\overline\Pi\overline V\phi\rangle .
 \label{v31:eq:source-pairing}
\end{equation}
Hence \eqref{v31:eq:dual-norm} gives
$\|g^\circ\|_{R^\circ,*}=2\sqrt2\,\Gamma^\circ$.  This proves the equivalence of
(a) and (b) by \cref{v31:thm:range} and also gives the energy in
\eqref{v31:eq:angle-energy}.

Assume (b).  The functional
$R^\circ\phi\mapsto\langle p^\circ,\overline\Pi\overline V\phi\rangle$
is well defined and bounded on $\Ran R^\circ$.  Its Riesz representative is a unique
$q_*\in\overline{\Ran R^\circ}$ with norm $\Gamma^\circ$.  Since
$q_*\in\Ran Q_o^\circ$,
\[
 \langle q_*-p^\circ,\overline V\phi\rangle
 =\langle q_*,Q_o^\circ\overline V\phi\rangle
  -\langle p^\circ,\overline\Pi\overline V\phi\rangle=0,
\]
which is \eqref{v31:eq:compensation}.  Thus (c) holds.

Conversely, if $q\in\Ran Q_o^\circ$ satisfies \eqref{v31:eq:compensation}, then
\[
 |\langle p^\circ,\overline\Pi\overline V\phi\rangle|
   =|\langle q,Q_o^\circ\overline V\phi\rangle|
   \le\|q\|\,\|R^\circ\phi\|,
\]
so $\Gamma^\circ\le\|q\|$.  Projecting any such $q$ orthogonally onto
$\overline{\Ran R^\circ}$ preserves the represented functional and decreases its norm;
therefore $q_*$ is the unique minimum-norm compensator.  Finally
\cref{v31:thm:range} represents $g^\circ$ as
$\langle g^\circ,\phi\rangle=\langle 4q_*,R^\circ\phi\rangle/2$, proving
$R^\circ n^\circ=4q_*$.
\end{proof}

The theorem explains the characteristic issue without calling it an eigenvalue problem.  A
lapse direction is dangerous when $\overline V\phi$ lies very close to the transverse range,
so $Q_o^\circ\overline V\phi$ is small, while the actual transverse source vector $p^\circ$
still has appreciable pairing with $\overline\Pi\overline V\phi$.  Only that
\emph{source-weighted} angle enters.

\section{The gate reduces to one scalar Hodge potential and three means}
\label{v31:sec:hodge}

The preceding condition can be written with the explicit coefficient matrices already used in
the bridge.  Let $M_b=M^\circ$ be the invertible off-diagonal map formed from $b^\circ$ and
let $M_r$ be the same map formed from $r^\circ$.  Then
\begin{equation}
       \overline V=-\tfrac12 M_r\nabla^\circ,
       \qquad p^\circ=M_b w^\circ .
 \label{v31:eq:MrMb}
\end{equation}
Let $P^\circ$ be the mean-zero transverse vector projection and
$Q_H^\circ=I-P^\circ$.  Since
$\Ran\overline T=M_b\Ran P^\circ$ up to the inherited nonzero normalization,
\begin{equation}
     \Ran Q_o^\circ=M_b^{-*}\Ran Q_H^\circ .
 \label{v31:eq:complement-parameterization}
\end{equation}
The Hodge complement has the exact parameterization
\begin{equation}
 \Ran Q_H^\circ
   =\left\{c+\nabla^\circ(\Omega^\circ)^{-1}\varphi:
       c\in\mathbb R^3,\ \overline\varphi=0\right\}.
 \label{v31:eq:Hodge-param}
\end{equation}

\begin{corollary}[One-scalar compensation equation]
\label[corollary]{v31:cor:hodge-gate}
The finite-energy range condition of \cref{v31:thm:source-angle} holds if and only if there
are $c\in\mathbb R^3$ and a mean-zero scalar $\varphi$ with
$z=c+\nabla^\circ(\Omega^\circ)^{-1}\varphi$ such that
\begin{equation}
 \boxed{\quad
   \overline V^*M_b^{-*}z=\overline V^*M_bw^\circ
   \quad\hbox{distributionally}.\quad}
 \label{v31:eq:Hodge-compensation}
\end{equation}
Equivalently, with $A=(M_b)^{-1}M_r$ and with the common nonzero adjoint normalization
cancelled from both sides,
\begin{equation}
 \boxed{\quad
 \operatorname{div}^\circ\!\left[A^*
   \bigl(c+\nabla^\circ(\Omega^\circ)^{-1}\varphi\bigr)\right]
   =\operatorname{div}^\circ\!\left[M_r^*M_bw^\circ\right].\quad}
 \label{v31:eq:scalar-compensation}
\end{equation}
Moreover
\begin{equation}
 (\Gamma^\circ)^2=
 \inf\left\{\|M_b^{-*}z\|_{L^2}^2:
       z\ \hbox{satisfies \eqref{v31:eq:Hodge-compensation} and
       \eqref{v31:eq:Hodge-param}}\right\}.
 \label{v31:eq:Gamma-inf}
\end{equation}
\end{corollary}
\begin{proof}
The orthogonal complement of $M_b\Ran P^\circ$ is
$M_b^{-*}(\Ran P^\circ)^\perp=M_b^{-*}\Ran Q_H^\circ$, proving
\eqref{v31:eq:complement-parameterization}.  Insert
$q=M_b^{-*}z$ in \eqref{v31:eq:compensation}; this gives
\eqref{v31:eq:Hodge-compensation}.  The Hodge parameterization is the same constants-plus-
gradients decomposition retained in the finite flux coordinates of the bridge.  Since
$\overline V=-M_r\nabla^\circ/2$, taking the distributional adjoint on both sides gives
\eqref{v31:eq:scalar-compensation}.  Finally the minimum-norm compensator of
\cref{v31:thm:source-angle} gives \eqref{v31:eq:Gamma-inf}.
\end{proof}

This is a substantial reduction of the continuum question.  It is one scalar equation for one
mean-zero Hodge potential plus three constants.  It still is not automatically elliptic.  The
visible flux characteristic proved in the bridge is precisely why no unproved uniform
coercivity of \eqref{v31:eq:scalar-compensation} is inserted here.

\section{Finite source-angle banks give monotone obstruction certificates}
\label{v31:sec:finite-bank}

The range criterion can be attacked without the characteristic scalar inverse.  Choose smooth
mean-zero tests $\phi_1,\ldots,\phi_m$ and put
\begin{equation}
 A^{(m)}_{ij}=\langle R^\circ\phi_i,R^\circ\phi_j\rangle,
 \qquad
 c^{(m)}_i=\langle p^\circ,\overline\Pi\overline V\phi_i\rangle .
 \label{v31:eq:bank}
\end{equation}

\begin{proposition}[Finite bank formula and exact rejection test]
\label[proposition]{v31:prop:bank}
Let $\Gamma_m$ be the supremum in \eqref{v31:eq:Gamma} restricted to the span of the first
$m$ tests.  Then
\begin{equation}
 \boxed{\quad
 \Gamma_m^2=
 \begin{cases}
  (c^{(m)})^*(A^{(m)})^\dagger c^{(m)},&c^{(m)}\in\Ran A^{(m)},\\
  +\infty,&c^{(m)}\notin\Ran A^{(m)}.
 \end{cases}\quad}
 \label{v31:eq:bank-formula}
\end{equation}
For nested banks, $\Gamma_m$ is nondecreasing and
$\sup_m\Gamma_m=\Gamma^\circ$ whenever their union is dense in the test space for the energy
seminorm.  In particular:
\begin{enumerate}[label=\textup{(\roman*)},leftmargin=2.7em]
\item one finite bank with $c^{(m)}\notin\Ran A^{(m)}$ proves that no finite-energy
      continuum solution exists;
\item any sequence of individual tests satisfying
\begin{equation}
 \|R^\circ\phi_m\|\longrightarrow0,
 \qquad
 |\langle p^\circ,\overline\Pi\overline V\phi_m\rangle|\ge c>0
 \label{v31:eq:witness-sequence}
\end{equation}
forces $\Gamma^\circ=+\infty$.
\end{enumerate}
\end{proposition}
\begin{proof}
For a coefficient vector $z$, the squared denominator and numerator of the restricted quotient
are $z^*A^{(m)}z$ and $|(c^{(m)})^*z|^2$.  If $c^{(m)}$ has a component in
$\ker A^{(m)}$, that component gives a zero denominator and nonzero numerator.  Otherwise the
standard positive-semidefinite Rayleigh quotient gives
$(c^{(m)})^*(A^{(m)})^\dagger c^{(m)}$.  Enlarging the test span can only increase the
supremum.  Density gives the full supremum by definition.  The final assertion follows from
the one-dimensional quotient.
\end{proof}

The projection $\overline\Pi$ in \eqref{v31:eq:bank} is not a new hard inverse.  The inherited
transverse problem has
$\overline B=\overline T^*\overline T$ on its physical slice and a uniformly controlled
inverse in every fixed Sobolev order and in a small exponential Fourier weight.  Thus for a
finite trigonometric bank one obtains
\begin{equation}
 \overline\Pi y=\overline T\,\overline B^{-1}\overline T^*y
 \label{v31:eq:Pi-compute}
\end{equation}
using only the already solved transverse normal equation.  The singular scalar inverse
$\mathscr S_h^{-1}$ never enters this witness calculation.  This gives a concrete route to
an exact rational or residual-certified source-specific rejection test.

\section{The previously constructed characteristic packets are not the missing source witness}
\label{v31:sec:packets}

It is useful to dispose of one tempting circular direction.  At finite cutoff put
\begin{equation}
 \mathfrak D_h^2:=\langle g_h,\mathscr S_h^{-1}g_h\rangle
   =\sup_{\phi\ne0}\frac{|\langle g_h,\phi\rangle|^2}
                              {\langle\phi,\mathscr S_h\phi\rangle}.
 \label{v31:eq:finite-dual}
\end{equation}
The equality is the finite-dimensional version of \cref{v31:thm:range}.  If
$n_h=\mathscr S_h^{-1}g_h$ and $u_h=hn_h$, then
\begin{equation}
          \mathfrak D_h^2
          =\left\langle\frac{g_h}{h},u_h\right\rangle .
 \label{v31:eq:dual-normalized}
\end{equation}
Since the inherited source expansion bounds $g_h/h$ in every fixed Sobolev norm, boundedness
of $u_h$ in $L^2$ is sufficient for $\mathfrak D_h=O(1)$.  Conversely
$\mathfrak D_h\to\infty$ would rule out a bounded normalized lapse.

The actual source already has a nonzero finite-dual floor.  Let
$\psi(x)=\sin(\vartheta_2+\vartheta_3)$ be the fixed smooth test used in the inherited rational
source certificate.  Its pairing with $g^\circ$ is bounded away from zero, while smooth-test
consistency gives $\langle\mathcal S_h\psi,\mathscr S_h\mathcal S_h\psi\rangle=O(h^2)$.
Together with $g_h=h\lambda_*^4\mathcal S_hg^\circ+O(h^2)$, the variational formula yields
\begin{equation}
                 \boxed{\quad \mathfrak D_h^2\ge c_*>0\quad}
 \label{v31:eq:dual-floor}
\end{equation}
for all sufficiently fine selected cutoffs.  Thus a successful normalized continuum lapse,
if it exists, carries a genuinely nonzero scalar energy; it is not approaching the zero class.

The localized visible characteristic packets of the bridge do not prove the opposite, namely
divergence of \eqref{v31:eq:finite-dual}.  For their normalized lapse component
$\widehat n_h$, the inherited estimates give
\begin{equation}
 |\langle \widehat n_h,g_h\rangle|
       \le Ch\exp(-c h^{-2/3}),
 \qquad
 \mathscr S_h\succeq c h^{12}I,
 \qquad \|\widehat n_h\|=1.
 \label{v31:eq:packet-input}
\end{equation}
Therefore their individual contribution to the variational quotient obeys
\begin{equation}
 \boxed{\quad
 \frac{|\langle g_h,\widehat n_h\rangle|}
      {\langle\widehat n_h,\mathscr S_h\widehat n_h\rangle^{1/2}}
 \le C h^{-5}e^{-c h^{-2/3}}\longrightarrow0 .\quad}
 \label{v31:eq:packet-harmless}
\end{equation}
So the already known geometrically visible quasimodes are not the source-weighted obstruction
we are looking for.  Any genuine rejection of the bounded-lapse branch must find a different
family whose source overlap does not collapse faster than its residual energy.

\section{Relation to the finite V30 impedance and the next exact calculation}
\label{v31:sec:interface}

Versions 30 and 31 answer different layers of the same problem.  Version 30 is an exact
finite-cutoff statement:
\[
 \mathfrak C_{h,s}=\|\mathcal O_h\mathscr S_h^{-1}P_{\le R_h}g_h\|_{H_h^s}^2
\]
and $\mathfrak C_{h,s}=O(h^2)$ is exactly the required low-source curvature estimate.  The
present body identifies the more primitive continuum entrance which would be required by a
bounded normalized lapse:
\begin{equation}
 \boxed{\quad
 \Gamma^\circ<\infty
 \quad\Longleftrightarrow\quad
 \exists\ q\in\Ran(I-\overline\Pi):
       \overline V^*(q-p^\circ)=0 .\quad}
 \label{v31:eq:final-gate}
\end{equation}
Equivalently it is the one-scalar compensation equation
\eqref{v31:eq:scalar-compensation}.

A failure of \eqref{v31:eq:final-gate}, certified for example by
\eqref{v31:eq:witness-sequence}, rules out the bounded-$hn_h$ continuum branch.  It does not
by itself prove that the observed curvature in Version 30 diverges: a rough lapse could still
be partially hidden by the geometric observation.  Conversely, passing
\eqref{v31:eq:final-gate} gives a finite-energy continuum solution class but does not by
itself supply the quantitative finite-$h$ regularity needed to prove
$\mathfrak C_{h,s}=O(h^2)$.  A stability/Mosco-type passage for the actual source would still
be required.  These qualifications prevent the continuum factorization from being used as a
surrogate for the finite observed certificate.

The useful change of target is that the first source-specific continuum test no longer needs
the poorly conditioned scalar inverse.  One can now proceed in either direction:
\begin{enumerate}[label=\textup{(A\arabic*)},leftmargin=2.8em]
\item build nested finite Fourier banks and certify $\Gamma_m$ using only the already
      controlled transverse solve \eqref{v31:eq:Pi-compute}; divergence gives an obstruction;
\item construct an explicit Hodge compensator $(c,\varphi)$ for
      \eqref{v31:eq:scalar-compensation}; its norm gives a positive finite-energy certificate,
      after which one returns to Version 30 for the finite observed regularity.
\end{enumerate}
The first existing characteristic packets are now proved to be the wrong bank for (A1) by
\eqref{v31:eq:packet-harmless}; repeating them would be circular.

No claim is made here that the actual compensation equation passes or fails.  No nonlinear
Einstein lifespan, higher weak row, Yang--Mills mass gap, chiral sector, or primitive
Renewal-to-Einstein--SM theorem is inferred.  The advance is an exact source-angle/range gate
and a scalar-plus-three-means equation which can be attacked before the singular scalar
inverse.

\paragraph{Verification and preservation.}
Every Version 30 byte before its final document command is preserved.  The accompanying
script checks the finite-dimensional analogues of the energy-completion theorem, the source
angle and minimum-norm compensation identities, the affine-intersection criterion, the
finite-bank pseudoinverse formula including its infinite case, and the normalized packet
scaling.  These are regressions for the written algebra, not a proof-assistant formalization
or a numerical evaluation of the still-open coefficient equation.  New labels use the prefix
\texttt{v31:}.

\addtocontents{toc}{\protect\endgroup}
% END IMMUTABLE VERSION 31 BODY
% APPEND VERSION 32 HERE, BEFORE THE FINAL END-DOCUMENT COMMAND.


% BEGIN IMMUTABLE VERSION 32 BODY
\clearpage
\begin{center}
{\Large\bfseries Version 32: The explicit residual-stress angle after the V22 bridge audit}\\[0.5em]
{\large A corrected cofinal interface, a monotone Fourier certificate, and a bounded-compensator criterion at zero spectral parameter}
\end{center}
\makeatletter
\addtocontents{toc}{\protect\begingroup}
\addtocontents{toc}{\protect\makeatletter}
\addtocontents{toc}{\protect\def\protect\l@section{\protect\bfseries\protect\@dottedtocline{1}{0em}{2.5em}}}
\makeatother
\addcontentsline{toc}{part}{Version 32: Explicit residual-stress angle after the latest Exact-Action Einstein audit}

\section{Audit correction: the Exact-Action bridge has advanced beyond the V16 interface}
\label{v32:sec:audit}

Versions 30--31 were deliberately formulated against the supplied Exact-Action bridge V16,
because that was the latest bridge body audited there.  A library audit now identifies a later
cumulative source, Versions 1--22 \cite{v32:EB22}.  Its V17--V19 bodies materially sharpen the
cofinal scalar problem, while V20--V22 construct and then dynamically test a distinct tilted
initial-multiplier branch at the fixed three-site cutoff.  Continuing Version 31 as though those
results did not exist would duplicate stronger work, so the present body changes the interface.

The following later results are inherited, not reproved.
\begin{enumerate}[label=\textup{(L\arabic*)},leftmargin=2.7em,itemsep=3pt]
\item V17 defines the maximal projected-flux graph
\[
  Zu=QA\nabla^\circ u,\qquad
  \mathcal D=\{u\in L_0^2(\T^3):Zu\in L^2\},
\]
with smooth graph core and closed form
\begin{equation}
 a(u,v)=\frac14\ip{Zu}{K^{-1}Zv},\qquad
 324^{-1}I\preceq K\preceq\frac{16}{9}I .
 \label{v32:eq:form}
\end{equation}
The complete normalized finite scalar forms Mosco-converge to this form.  Fixed positive-parameter
resolvents for the actual normalized source converge strongly, but the zero-parameter inverse is
left open.  These are \nolinkurl{v17:graph}, \nolinkurl{v17:form-limit}, and \nolinkurl{v17:resolvent} of
\cite{v32:EB22}.
\item V18 splits the actual limiting forcing as
\begin{equation}
 \boxed{\quad g^\circ=g_{\rm b}+g_{\rm r},\qquad
 g_{\rm b}=-2Z^*Qy,\qquad
 g_{\rm r}=-2D^*f^\circ,\quad D=A\nabla^\circ,\quad}
 \label{v32:eq:split}
\end{equation}
where $g_{\rm b}$ is already a bounded functional in the energy norm, while $f^\circ$ is the
\emph{explicit finite Fourier field}
$C^\circ(\Omega^\circ)^{-2}q$ from the original source table.  Its exact squared norm is
\begin{equation}
                    \|f^\circ\|^2=\frac{594441}{65536}.
 \label{v32:eq:f-norm}
\end{equation}
V18 proves that finite susceptibility of the actual source is equivalent to an $L^2$ stress
$j\in\Ran Q$ solving
\begin{equation}
             \operatorname{div}^\circ\!\left[A^*(j+2f^\circ)\right]=0.
 \label{v32:eq:stress}
\end{equation}
These are \nolinkurl{v18:source-split} and \nolinkurl{v18:dual}.
\item V18 also finds a nonconstant chamber mode $\kappa$ with $Z\kappa=0$ and
$\mathscr A\kappa=0$, and proves exact orthogonality of the actual forcing to that mode.
V19 goes further: zero lies in the essential spectrum of the inversion-odd sector containing the
actual source.  Hence neither deleting finitely many modes nor asserting a Fredholm gap can settle
\eqref{v32:eq:stress}.  These are \nolinkurl{v18:kernel-theorem},
\nolinkurl{v18:orthogonality}, and \nolinkurl{v19:essential-theorem}.
\item For the actual source parity, V19 proves that any compensator may be chosen with no harmonic
part:
\begin{equation}
 \boxed{\quad j=\nabla^\circ\varphi,\qquad
 \varphi\in H^1\cap L_0^2,\quad \varphi(-\theta)=\varphi(\theta),\qquad
 \operatorname{div}^\circ\!\left[A^*(\nabla^\circ\varphi+2f^\circ)\right]=0.\quad}
 \label{v32:eq:potential}
\end{equation}
Thus Version 31's three harmonic compensation parameters are unnecessary for this particular
source.  The precise inherited result is \nolinkurl{v19:parity}.
\item V20--V22 concern a different branch: the initial multiplier values are tilted by order
amplitude at $N=3$.  V21 cancels the complete initial weak rate, while V22 proves a nonzero
first-jet tangency defect $\Delta_*$ at the selected tilt, an $a^{-2}$ weak-acceleration pole, and
an original-action layer of physical duration $O(a^3)$.  It also leaves an explicit
$11$-parameter quadratic first-tilt gate unsolved.  None of those fixed-cutoff tilted records
supplies a cofinal solution of \eqref{v32:eq:potential} for the old unit-initial-multiplier family.
\end{enumerate}

Accordingly, Version 31's source-angle language remains a useful abstract Hilbert-space idea, but
its V16 implementation is no longer the minimal cofinal test.  The bounded term $g_{\rm b}$ has
already been paid, parity removes the means, and the only unresolved energy-dual incidence is the
explicit finite stress $f^\circ$ in \eqref{v32:eq:potential}.  The rest of this body isolates that
incidence without introducing the singular scalar inverse.

\section{The corrected dangerous-source angle uses only the explicit finite stress}
\label{v32:sec:angle}

Let $P=P^\circ$ be the mean-zero transverse Hodge projection and $Q=I-P$.  The coefficient
matrix is
\[
 A=M_b^{-1}M_r,
\]
with the literal trigonometric $M_b,M_r$ of the bridge.  Since $f^\circ\in\Ran P$, for every
smooth scalar $u$,
\begin{equation}
 \ip{g_{\rm r}}u=-2\ip{f^\circ}{Du}
      =-2\ip{f^\circ}{PDu},\qquad Zu=QDu.
 \label{v32:eq:pairing}
\end{equation}
The source and the form are inversion odd/even in the sense of V19, so only odd scalar tests can
contribute.  Let $C^\infty_{-,0}$ denote the real smooth mean-zero functions satisfying
$u(-\theta)=-u(\theta)$.

\begin{definition}[Explicit residual-stress angle]
\label[definition]{v32:def:Gamma}
Define
\begin{equation}
 \boxed{\quad
 \Gamma_f:=\sup_{u\in C^\infty_{-,0}}
     \frac{|\ip{f^\circ}{PDu}|}{\|QDu\|_2} .\quad}
 \label{v32:eq:Gamma}
\end{equation}
A test with $QDu=0$ and nonzero numerator makes the quotient $+\infty$; a $0/0$ test is ignored.
\end{definition}

This quantity contains no scalar inverse and no unknown transverse solve: $A$ is the explicit
analytic coefficient matrix and $P,Q$ are ordinary Fourier Hodge projections.  The source
$f^\circ$ has only the nine positive-frequency rows inherited from the original action.

\begin{theorem}[Exact zero-parameter range criterion after the later source split]
\label[theorem]{v32:thm:angle}
Let
\begin{equation}
 \mathcal E_{\rm r}^*:=\sup_{u\in\mathcal D}
            \{2\ip{g_{\rm r}}u-a(u,u)\}\in[0,+\infty].
 \label{v32:eq:Er}
\end{equation}
Then the following are equivalent:
\begin{enumerate}[label=\textup{(\roman*)},leftmargin=2.5em,itemsep=3pt]
\item $\Gamma_f<\infty$;
\item $\mathcal E_{\rm r}^*<\infty$;
\item the actual limiting source has finite energy-dual susceptibility;
\item the stress equation \eqref{v32:eq:stress} has an $L^2$ solution;
\item the even-potential equation \eqref{v32:eq:potential} has a solution
$\varphi\in H^1\cap L_0^2$.
\end{enumerate}
Moreover the explicit metric bounds in \eqref{v32:eq:form} give
\begin{equation}
 \boxed{\qquad
       \frac4{81}\Gamma_f^2
        \le \mathcal E_{\rm r}^*
        \le \frac{256}{9}\Gamma_f^2 .
       \qquad}
 \label{v32:eq:Gamma-bounds}
\end{equation}
In particular the zero essential spectral threshold does not by itself make the actual source
singular; it is singular precisely when the source angle in \eqref{v32:eq:Gamma} diverges.
\end{theorem}
\begin{proof}
The inherited metric bounds imply
\begin{equation}
        \frac9{64}\|Zu\|_2^2\le a(u,u)\le81\|Zu\|_2^2.
 \label{v32:eq:metric-equivalence}
\end{equation}
For $L(u)=\ip{g_{\rm r}}u$, \eqref{v32:eq:pairing} gives
$|L(u)|=2|\ip{f^\circ}{PDu}|$.  Hence finiteness of $\Gamma_f$ is exactly boundedness of
$L$ with respect to the form seminorm, up to the fixed constants in
\eqref{v32:eq:metric-equivalence}.  The smooth graph core from V17 extends this statement from
smooth functions to all of $\mathcal D$.  Even tests pair to zero by the inherited inversion
parities, so restriction to $C^\infty_{-,0}$ loses nothing.

For a nonnegative quadratic form, direct optimization on each real line gives
\[
 \sup_u\{2L(u)-a(u,u)\}
      =\sup_{a(u,u)>0}\frac{|L(u)|^2}{a(u,u)},
\]
with value $+\infty$ if $L$ is nonzero on the form kernel.  Substitution of
\eqref{v32:eq:metric-equivalence} proves \eqref{v32:eq:Gamma-bounds}.

The bounded term $g_{\rm b}$ in \eqref{v32:eq:split} is already continuous in the form norm.
Therefore the class of finite-energy-dual functionals contains $g^\circ$ exactly when it contains
$g_{\rm r}$.  The equivalence with an $L^2$ stress is the inherited V18 dual-range theorem, and
V19's parity theorem gives the equivalence with the even potential in
\eqref{v32:eq:potential}.  No closed-range assertion is used.
\end{proof}

The comparison with Version 31 is now precise.  There we tested a full-source incidence before the
later split was audited.  Here the already bounded component is removed by a proved theorem, the
remaining source is the finite $f^\circ$, and the harmonic compensation is removed by a proved
parity theorem.  Thus \eqref{v32:eq:Gamma} is the smaller source-native gate.

\section{A monotone finite Fourier certificate which is equivalent to the full range question}
\label{v32:sec:finite}

For $L\ge1$, let $\mathcal V_L^-$ be the real inversion-odd trigonometric polynomials of mean zero
with $|k|_1\le L$.  Choose any real $L^2$-orthonormal basis
$\phi_1,\ldots,\phi_{m_L}$ and define
\begin{equation}
 (H_L)_{ab}=\ip{Z\phi_a}{Z\phi_b},
 \qquad
 (c_L)_a=\ip{f^\circ}{PD\phi_a}.
 \label{v32:eq:finite-matrices}
\end{equation}
The temporary notation $ZD\!{}^0\phi:=Z\phi$ only emphasizes that there is one first-order
coefficient map; no extra derivative is introduced.  Thus $H_L$ is simply the Gram matrix of the
vectors $Q A\nabla^\circ\phi_a$.

Define
\begin{equation}
 \Gamma_L:=\sup_{0\ne u\in\mathcal V_L^-}
       \frac{|\ip{f^\circ}{PDu}|}{\|Z u\|_2},
 \label{v32:eq:GammaL}
\end{equation}
with the same zero-denominator convention as in \cref{v32:def:Gamma}.

\begin{theorem}[Monotone Fourier exhaustion of the exact source angle]
\label[theorem]{v32:thm:Fourier}
For every $L$,
\begin{equation}
 \boxed{
 \Gamma_L^2=
 \begin{cases}
   c_L^*H_L^\dagger c_L,&c_L\in\Ran H_L,\\[1mm]
   +\infty,&c_L\notin\Ran H_L,
 \end{cases}}
 \label{v32:eq:pseudoinverse}
\end{equation}
and
\begin{equation}
             \boxed{\qquad \Gamma_L\nearrow\Gamma_f.\qquad}
 \label{v32:eq:monotone}
\end{equation}
Consequently the following finite-bank statement is equivalent to the full zero-parameter range
condition:
\begin{equation}
 \boxed{\qquad
 \sup_{L\ge1}\Gamma_L<\infty
 \quad\Longleftrightarrow\quad
 \text{the V18--V19 actual-source compensator exists.}
 \qquad}
 \label{v32:eq:finite-iff}
\end{equation}
Thus one finite bank with $c_L\notin\Ran H_L$, or any certified sequence
$\Gamma_L\to\infty$, proves failure.  Conversely a uniform certified bound on the nested banks
proves existence; no passage through the unregularized scalar inverse is needed.
\end{theorem}
\begin{proof}
For $u=\sum_ax_a\phi_a$, the denominator squared is $x^*H_Lx$ and the numerator is
$|c_L^*x|$.  A finite-dimensional linear functional is bounded in this seminorm exactly when it
annihilates $\ker H_L$, equivalently when $c_L\in\Ran H_L$.  Diagonalizing $H_L$ on its positive
range gives \eqref{v32:eq:pseudoinverse}.

The spaces $\mathcal V_L^-$ are nested, so $\Gamma_L$ is nondecreasing and at most
$\Gamma_f$.  V17 proves that smooth mean-zero functions are a graph core for $Z$.  Symmetrizing a
graph-core approximation under inversion gives an odd graph-core approximation for every odd
domain vector.  A smooth odd function is in turn approximated in every fixed Sobolev norm by its
odd Fourier partial sums; since $A$ is smooth and $Q$ is an order-zero Fourier projection, this
implies convergence of both $u$ and $Zu$.  The numerator can be written
$-\ip{g_{\rm r}}u/2$, so it also converges under the scalar $L^2$ convergence.  Every quotient in
\eqref{v32:eq:Gamma} is therefore the limit of quotients from the finite banks, proving
\eqref{v32:eq:monotone}.  Combine with \cref{v32:thm:angle}.
\end{proof}

This theorem is deliberately stronger than saying that a small Fourier ansatz ``did not find'' a
compensator.  The finite matrices provide monotone \emph{lower} certificates for the dangerous
source angle.  Failure at a small support is conclusive only through the exact range alternative in
\eqref{v32:eq:pseudoinverse}; a merely nonzero least-squares residual is not.

\section{A bounded-potential compactness certificate: what a constructive solve must prove}
\label{v32:sec:compactness}

The parity reduction in V19 permits a complementary primal test.  For an even mean-zero
trigonometric polynomial $\varphi$, define the stress residual
\begin{equation}
 \mathcal R(\varphi)=
   \operatorname{div}^\circ\!\left[A^*(\nabla^\circ\varphi+2f^\circ)\right]
          \in H^{-1}(\T^3).
 \label{v32:eq:residual}
\end{equation}

\begin{theorem}[Bounded finite-polynomial compensators are necessary and sufficient]
\label[theorem]{v32:thm:compactness}
The exact potential equation \eqref{v32:eq:potential} has a solution if and only if there is a
sequence of real even mean-zero trigonometric polynomials $\varphi_n$ such that
\begin{equation}
 \boxed{\qquad
     \sup_n\|\nabla^\circ\varphi_n\|_2<\infty,
     \qquad
     \|\mathcal R(\varphi_n)\|_{H^{-1}}\longrightarrow0.
 \qquad}
 \label{v32:eq:compactness-certificate}
\end{equation}
Equivalently, a computational stress search is a proof of existence only after it controls both its
residual and its $H^1$ size.  Residual decay with an unbounded compensating gradient is not, by
itself, an upper certificate at the essential zero threshold.
\end{theorem}
\begin{proof}
If $\varphi$ solves \eqref{v32:eq:potential}, take its Fourier partial sums after the even
mean-zero projection.  They converge to $\varphi$ in $H^1$.  Smoothness and boundedness of $A$
then imply
$A^*\nabla\varphi_n\to A^*\nabla\varphi$ in $L^2$, hence their divergences converge in $H^{-1}$.
This gives \eqref{v32:eq:compactness-certificate}.

Conversely the Poincare inequality and the gradient bound make $\varphi_n$ bounded in $H^1$.
Pass to a weakly convergent subsequence $\varphi_n\rightharpoonup\varphi$ in $H^1$; parity and
mean zero are weakly closed.  Multiplication by the smooth $A^*$ is bounded on $L^2$, so
$A^*\nabla\varphi_n\rightharpoonup A^*\nabla\varphi$.  The divergence map is continuous
$L^2\to H^{-1}$.  Since the residual converges strongly to zero in $H^{-1}$, its weak limit is
zero, proving \eqref{v32:eq:potential}.  No coercive estimate for the degenerate coefficient is
used.
\end{proof}

The theorem explains exactly why the numerical small-Fourier attempts reported in V19 do not
settle the question.  Their residuals alone are not a nonexistence proof, and an improving sequence
would still need a uniform gradient bound to become an existence proof.  This is the correct
compactness quantity in the presence of the essential zero threshold.

\section{The explicit source table makes the new gate independent of the singular transverse feedback}
\label{v32:sec:explicit}

The residual angle is unusually concrete.  The nine positive-frequency coefficients of the original
$q$ are the V15 table, satisfy $k\cdot q_k=0$, and determine
\begin{equation}
 \widehat f^\circ(k)=\frac{\mathrm i\,k\times q_k}{|k|^2}
 \qquad(k\ne0).
 \label{v32:eq:f-Fourier}
\end{equation}
Consequently Parseval gives
\begin{equation}
 \|f^\circ\|^2
   =2\sum_{k\in\mathcal B_+}\frac{|q_k|^2}{|k|^2}
   =\frac{594441}{65536},
 \label{v32:eq:f-replay}
\end{equation}
which replays the inherited V18 value directly from the printed source table.  The coefficient
$A=M_b^{-1}M_r$ is an explicit analytic trigonometric-rational matrix with a uniformly invertible
$M_b$.  Therefore every finite matrix in \eqref{v32:eq:finite-matrices} can be evaluated and
interval-certified using only:
\begin{enumerate}[label=(\alph*),leftmargin=2.3em,itemsep=2pt]
\item the printed finite source table for $q$;
\item the explicit $b^\circ,r^\circ$ coefficient functions;
\item ordinary Fourier differentiation and the Hodge projections $P,Q$;
\item quadrature or Fourier convolution with a rigorous enclosure for the analytic coefficient
$M_b^{-1}$.
\end{enumerate}
There is no call to $\mathscr A^{-1}$ and no unknown $w^\circ$ in the dangerous part.  The
transverse-normal solve that generated $w^\circ$ survives only inside $g_{\rm b}$, whose energy-dual
boundedness is already proved in V18.

This is the principal gain over the Version 31 acquisition plan.  A cofinal obstruction or entrance
can now be sought by nested finite banks without recomputing the singular scalar response at each
cutoff.  The result remains an analytical continuum gate; transferring a positive conclusion back
to the unregularized finite inverse still requires control of the zero spectral parameter under
refinement, beyond the fixed-$\rho$ Mosco/resolvent convergence of V17.

\section{The old and tilted branches must now be kept separate}
\label{v32:sec:branches}

The later bridge reveals two distinct open routes.

\paragraph{Unit-initial-multiplier cofinal route.}
The old complementary branch keeps initial multiplier values $(1,0)$.  V17--V19 prove its full
form limit, explicit source split, kernel compatibility, essential zero threshold, and source
suppression on a growing low-energy space.  Its remaining zero-parameter gate is exactly
\eqref{v32:eq:finite-iff}, equivalently \eqref{v32:eq:compactness-certificate}.  A proof that
$\Gamma_L$ stays uniformly bounded supplies the missing $L^2$ compensator; a certified divergence
supplies a source-specific obstruction.  Essential spectrum alone decides neither outcome.

\paragraph{Tilted fixed-cutoff route.}
V20--V22 change the initial multiplier values by order amplitude and therefore define different
original records of the same action.  At $N=3$, V21 constructs an $11$-parameter affine family of
first tilts after cancelling the quadratic rows and cubic means.  V22 evaluates one Poisson-null
tangency component at its selected point and proves it nonzero, then derives the $a^3$ initial
layer and the curvature-rate response.  The next finite gate is the common zero problem for the
quadratic polynomials
\begin{equation}
 \mathcal T_\nu(\eta)=
 Dd_{2,W}(Y)[F_1Y+G\ell(\eta),\nu]
 +\mathbb K_1(F_1Y+G\ell(\eta))[\nu,\ell(\eta)],
 \qquad \eta\in\mathcal N_{21},
 \label{v32:eq:tilt-gate}
\end{equation}
for all Poisson-null tests $\nu$.  V22 proves only a local nonzero neighborhood at the chosen
point; it does not decide the global $11$-parameter polynomial system.  Even a solution there
would still need refinement-uniform tilts and a propagated invariant-manifold or rate estimate.

The source $g_h$ and stress equation studied above belong to the old unit-initial branch.  They
must not be imported into the tilted branch after its initial values have changed.  Conversely the
fixed-$N=3$ tangency defect does not prove failure of the old cofinal stress equation.  The two
routes are complementary upstream tests, not competing descriptions of one record.

\section{Outcome and next exact computation}
\label{v32:sec:outcome}

The V17--V22 audit changes the next calculation substantially.  The full scalar-domain and
fixed-$\rho$ convergence theory already exist; repeating Version 31's abstract range theorem is no
longer useful.  For the old cofinal branch the unresolved content is now concentrated in the single
explicit number
\begin{equation}
 \boxed{\qquad
   \Gamma_f=\sup_{u\in C^\infty_{-,0}}
        \frac{|\ip{f^\circ}{P A\nabla^\circ u}|}
             {\|Q A\nabla^\circ u\|_2},
   \qquad A=M_b^{-1}M_r,
 \qquad}
 \label{v32:eq:final-gate}
\end{equation}
and its nested finite certificates
\begin{equation}
               \Gamma_L^2=c_L^*H_L^\dagger c_L
 \label{v32:eq:final-finite}
\end{equation}
whenever the finite range condition holds.  Uniform boundedness is equivalent to the actual V19
compensator; divergence is an obstruction.  The fully explicit $f^\circ$ means this calculation
can be performed without either the singular scalar inverse or the already discharged transverse
feedback.

I would therefore continue by assembling $H_L,c_L$ for increasing Fourier radii with certified
analytic-coefficient quadrature.  In parallel, the primal search should track the gradient norm of
its even potential as required by \cref{v32:thm:compactness}.  Those two calculations are dual:
a growing $\Gamma_L$ diagnoses source incidence with the closing flux range, while a bounded
potential sequence with vanishing $H^{-1}$ residual constructs the missing stress directly.

No claim is made here that either sequence has already been evaluated for enough radii to decide
the range.  No fixed positive physical Einstein lifespan, cofinal tilted family, higher weak
closure, Yang--Mills mass gap, or primitive Renewal-to-Einstein--SM theorem is inferred.  The
advance is that the latest bridge collapses the cofinal zero-parameter question to an explicit
finite-stress angle and supplies an exact monotone finite certificate for it.

\paragraph{Verification and preservation.}
Every Version 31 byte before its final document command is preserved.  The verification script
checks the metric constants in \eqref{v32:eq:Gamma-bounds}, the finite Gram/pseudoinverse formula
including the infinite alternative, monotonicity on nested toy banks, the bounded-compensator
compactness implication in finite analogues, inversion-parity bookkeeping, and the exact replay
of \eqref{v32:eq:f-replay} from the nine printed source coefficients.  These checks validate the
new algebra and constants; they do not numerically decide the actual infinite-dimensional range.
New labels use the prefix \texttt{v32:}.

\begingroup
\renewcommand{\refname}{Additional source for Version 32}
\begin{thebibliography}{9}
\small\raggedright
\bibitem{v32:EB22}
\emph{Renewal Exact-Action Einstein Bridge Research Notes}, supplied cumulative Versions 1--22,
file \nolinkurl{Renewal_Exact_Action_Einstein_Bridge_Research_Notes_V22.tex}.
The inherited interfaces used here are \nolinkurl{v17:graph}, \nolinkurl{v17:form-limit},
\nolinkurl{v17:resolvent}, \nolinkurl{v18:source-split}, \nolinkurl{v18:dual},
\nolinkurl{v18:kernel-theorem}, \nolinkurl{v18:orthogonality}, \nolinkurl{v19:parity},
\nolinkurl{v19:essential-theorem}, and the V20--V22 tilted-branch status results.  They are treated at
their stated scopes; no result of that cumulative source is represented as independently
reproved in this body.
\end{thebibliography}
\endgroup

\addtocontents{toc}{\protect\endgroup}
% END IMMUTABLE VERSION 32 BODY
% APPEND VERSION 33 HERE, BEFORE THE FINAL END-DOCUMENT COMMAND.



% BEGIN IMMUTABLE VERSION 33 BODY
\clearpage
\begin{center}
{\Large\bfseries Version 33: Escape from the V23 first-jet seed before another root search}\\[0.5em]
{\large A rank-stratum dichotomy, a bounded-tilt local no-go, and the joint amplitude--first-jet scale}
\end{center}
\makeatletter
\addtocontents{toc}{\protect\begingroup}
\addtocontents{toc}{\protect\makeatletter}
\addtocontents{toc}{\protect\def\protect\l@section{\protect\bfseries\protect\@dottedtocline{1}{0em}{2.5em}}}
\makeatother
\addcontentsline{toc}{part}{Version 33: Escape from the V23 first-jet seed and the bounded-tilt local no-go}

\section{Superseding audit: V23 closes the old eleven-tilt route}
\label{v33:sec:audit}

Version 32 audited the Exact-Action Einstein bridge only through cumulative V22.  The newer
cumulative V23 source \cite{v33:EB23} materially changes the fixed-cutoff tilted branch and must be
inserted before another calculation on that branch.  We retain the V32 unit-initial-multiplier
stress problem as a separate cofinal question; nothing below changes its explicit source angle or
its Fourier certificate.

The following V23 statements are inherited at their literal scope.
\begin{enumerate}[label=\textup{(J\arabic*)},leftmargin=2.7em,itemsep=3pt]
\item At the complementary three-site first jet, the leading tilts which cancel the complete
quadratic drift, satisfy the four multiplier normalizations, and cancel the three cubic constant
shift rows form an eleven-dimensional affine space
\[
                         \mathfrak L_{21}=\ell_{21}+\mathcal N_{21}.
\]
\item There is an exact weak Poisson-null test
\[
 \phi_\diamond=\frac23\{\cos(2\pi x_1)-\cos(2\pi x_2)\},\qquad
 \nu_\diamond=(0,\nabla_\delta\phi_\diamond)
\]
for which the complete V22 tangency polynomial is \emph{constant on the entire affine space}:
\begin{equation}
 \boxed{\quad \mathcal T_{\nu_\diamond}(Y_*,\ell)
             =\Delta_\diamond\ne0
             \qquad(\ell\in\mathfrak L_{21}).\quad}
 \label{v33:eq:V23-no-go}
\end{equation}
The source proves all eleven linear, eleven diagonal-quadratic and fifty-five mixed quadratic
coefficients vanish exactly over the rationals; the remaining constant is nonzero.
\item V23 therefore moves the attack upstream to the real-analytic five-parameter first-jet chart
\begin{equation}
 p=(\alpha_1,\alpha_2,\alpha_3,b_2,b_3),\qquad
 b_1(p)=\frac{(\alpha_1^2+\alpha_2^2+\alpha_3^2)/4-b_2b_3}{b_2+b_3},
 \label{v33:eq:five-chart}
\end{equation}
with base point
\begin{equation}
 p_*=(1,9/8,1,-8,-8),\qquad b_{1,*}=16175/4096.
 \label{v33:eq:pstar}
\end{equation}
Every first jet in a neighborhood of this point obeys the original linear constraints, exact
curl-free quadratic weak neutrality, the four quadratic mean conditions, and admits the original
nonlinear initial-constraint preparation.  V23 does not assert a tangency root on this chart.
\end{enumerate}

Thus Version 32's description of the global eleven-tilt polynomial as still open is superseded by
\eqref{v33:eq:V23-no-go}.  The new question is whether changing the canonical first jet can remove
the defect.  V23 proposes differentiating a tilt-invariant tangency quotient in the five variables.
Before doing that finite coefficient calculation, there is a structural consequence of
\eqref{v33:eq:V23-no-go}: a root cannot approach the old first jet with bounded leading tilt while
the lower-incidence ranks remain regular.  We prove that consequence here.

\section{The lower-incidence bundle on a regular first-jet stratum}
\label{v33:sec:bundle}

Let $\mathcal M$ denote the finite physical leading-multiplier space at the three-site cutoff.
For a first jet $Y(p)$ from \eqref{v33:eq:five-chart}, stack the three lower incidences of V23 as
\begin{equation}
 \mathscr L(p)\ell=d(p),
 \qquad
 \mathscr L(p)=
 \begin{pmatrix}
  \mathbb K_1(Y(p))\\[1mm]
  \mathsf m\\[1mm]
  \mathcal R_{Y(p)}
 \end{pmatrix},
 \qquad
 d(p)=
 \begin{pmatrix}
   -\widehat B_2(Y(p))\\[1mm]0\\[1mm]-t(Y(p))
 \end{pmatrix}.
 \label{v33:eq:lower-stack}
\end{equation}
Thus $\mathscr L(p)\ell=d(p)$ is exactly
$\mathfrak F_2=0$, $\mathsf m\ell=0$, and $\mathfrak F_c=0$ from the V23 upstream system.
All entries are real analytic in $p$ on the chart.

The weak Poisson-null space is
\begin{equation}
 \mathcal N_W(p)=\{\nu\in\mathcal W:
      \mathbb K_1(Y(p))[\nu,\mu]=0\text{ for every multiplier }\mu\}.
 \label{v33:eq:NW}
\end{equation}
We call a neighborhood $U$ of $p_*$ a \emph{regular solvable incidence stratum} when the ranks of
the first Poisson restriction defining \eqref{v33:eq:NW}, of the stacked lower map
$\mathscr L(p)$, and of its augmented matrix $[\mathscr L(p)\mid d(p)]$ are constant on $U$,
with
\[
 \operatorname{rank}[\mathscr L(p)\mid d(p)]=\operatorname{rank}\mathscr L(p)=\operatorname{rank}\mathscr L(p_*) .
\]
Thus the lower system remains solvable with its base kernel dimension, while the weak Poisson-null
bundle keeps its base dimension.  This is the precise rank-and-solvability branch on which the V23
null bundle and lower solution set can be transported.  We do not assume that all five first-jet
directions stay in such a stratum; loss of solvability or a rank change is one of the alternatives
in the theorem below.

\begin{lemma}[Analytic coordinates for the regular lower-incidence fibre]
\label[lemma]{v33:lem:bundle}
On every sufficiently small regular solvable incidence stratum $U$ through $p_*$ there are real-analytic
maps
\begin{equation}
 \ell_0:U\to\mathcal M,
 \qquad E:U\to\mathcal L(\mathbb R^{11},\mathcal M),
 \label{v33:eq:ellE}
\end{equation}
with $E(p)$ injective, such that
\begin{equation}
 \boxed{\quad
   \{\ell:\mathscr L(p)\ell=d(p)\}
      =\{\ell_0(p)+E(p)c:c\in\mathbb R^{11}\}.
 \quad}
 \label{v33:eq:fibre}
\end{equation}
They may be chosen with $\ell_0(p_*)=\ell_{21}$ and
$\Ran E(p_*)=\mathcal N_{21}$.
There is also an analytic nonzero section $\nu(p)\in\mathcal N_W(p)$ with
$\nu(p_*)=\nu_\diamond$.
\end{lemma}
\begin{proof}
At $p_*$ the inherited V20--V23 certificates give the stated lower solution dimension and the
weak-null vector.  On a constant-rank neighborhood, keep one nonzero maximal minor of the stacked
map and one nonzero maximal minor of the weak Poisson restriction.  Analytic Gaussian elimination
using those fixed minors gives an analytic particular solution and an analytic basis of the
kernel.  This proves \eqref{v33:eq:fibre}.  The same fixed-minor elimination transports
$\nu_\diamond$ analytically inside the weak nullspace.  Shrinking $U$ keeps $E(p)$ injective and
$\nu(p)$ nonzero.  This is ordinary finite-dimensional analytic linear algebra; no scalar inverse,
SVD threshold, or new source equation is introduced.
\end{proof}

For $p\in U$ and $c\in\mathbb R^{11}$ define the transported scalar tangency test
\begin{equation}
 \Delta(p,c):=
 \mathcal T_{\nu(p)}\!\left(Y(p),\ell_0(p)+E(p)c\right).
 \label{v33:eq:Delta-pc}
\end{equation}
The V22--V23 formula for $\mathcal T$ is analytic in the first jet and the test and is polynomial
of degree at most two in the leading tilt.  Hence
\begin{equation}
 \Delta(p,c)=d_0(p)+d_1(p)\cdot c+c^TD_2(p)c,
 \label{v33:eq:quadratic-c}
\end{equation}
with analytic coefficient functions (replace $D_2$ by its symmetric part without changing the
quadratic polynomial).  The inherited global no-go gives the much stronger base identity
\begin{equation}
 d_0(p_*)=\Delta_\diamond\ne0,
 \qquad d_1(p_*)=0,
 \qquad D_2(p_*)=0.
 \label{v33:eq:base-coefficients}
\end{equation}

\section{Every nearby tangency root either changes rank or escapes to infinite tilt}
\label{v33:sec:escape}

Fix any Euclidean norm on the five chart parameters and put
\begin{equation}
                         \delta(p)=|p-p_*|.
 \label{v33:eq:delta}
\end{equation}
All finite-dimensional norm choices are equivalent; the exponent below is therefore intrinsic to
this local statement.

\begin{theorem}[Square-root escape law from the V23 global no-go]
\label[theorem]{v33:thm:escape}
There are positive constants $\delta_0,C_0,c_0,C_1$, depending only on a sufficiently small
regular incidence stratum, such that for $0<\delta(p)<\delta_0$,
\begin{equation}
 \left|\Delta(p,c)-\Delta_\diamond\right|
       \le C_0\,\delta(p)\bigl(1+|c|+|c|^2\bigr).
 \label{v33:eq:Delta-Lipschitz}
\end{equation}
Consequently every solution of the complete lower incidences and this transported tangency
component,
\begin{equation}
 \mathscr L(p)\ell=d(p),\qquad
 \mathcal T_{\nu(p)}(Y(p),\ell)=0,
 \label{v33:eq:root-system}
\end{equation}
obeys
\begin{equation}
 \boxed{\qquad
 |c|\ge c_0\,\delta(p)^{-1/2},
 \qquad
 \|\ell\|\ge C_1^{-1}\delta(p)^{-1/2}-C_1 .
 \qquad}
 \label{v33:eq:escape}
\end{equation}
In particular there is no bounded-leading-tilt tangency root converging to $p_*$ while the
incidence ranks stay on their base stratum.
\end{theorem}
\begin{proof}
Analyticity of the three coefficient functions in \eqref{v33:eq:quadratic-c}, together with
\eqref{v33:eq:base-coefficients}, gives after shrinking the parameter ball
\[
 |d_0(p)-\Delta_\diamond|+\|d_1(p)\|+\|D_2(p)\|
                                      \le C_0\delta(p),
\]
which is \eqref{v33:eq:Delta-Lipschitz} after changing $C_0$ by a fixed factor.

Suppose $\Delta(p,c)=0$.  Then
\[
 |\Delta_\diamond|
 \le C_0\delta(p)(1+|c|+|c|^2).
\]
Choose $\delta_0$ so that $3C_0\delta_0<|\Delta_\diamond|$.  The preceding inequality then
forces $|c|>1$.  Since $1+|c|+|c|^2\le3|c|^2$ for $|c|\ge1$,
\[
 |c|\ge
 \left(\frac{|\Delta_\diamond|}{3C_0}\right)^{1/2}\delta(p)^{-1/2}.
\]
This is the first inequality in \eqref{v33:eq:escape}.

On the compact closure of a smaller $U$, injectivity of $E(p)$ gives a uniform lower singular
value $\sigma>0$, while $\ell_0(p)$ is uniformly bounded by some $M_0$.  Thus
\[
 \|\ell_0(p)+E(p)c\|\ge \sigma|c|-M_0.
\]
Absorb $\sigma,M_0$ and the preceding constant into $C_1$ to prove the second inequality.
\end{proof}

The theorem is deliberately local and fixed-cutoff.  It does not claim that the full five-parameter
chart is a regular solvable incidence stratum; if a candidate approaches $p_*$ through a locus on
which the lower system loses solvability, or on which the Poisson or lower-incidence rank changes,
it leaves the hypotheses before the escape conclusion can be applied.  That is not a loophole
hidden in the proof: it is the other upstream mechanism to be tested.

\begin{corollary}[Rank-change-or-escape dichotomy]
\label[corollary]{v33:cor:dichotomy}
Let $(p_n,\ell_n)$ be tangency-root candidates for the complete V23 upstream system with
$p_n\to p_*$.  Then, after passage to any subsequence on which the lower system remains solvable and all relevant
incidence ranks equal their base values, one has
\begin{equation}
                         \|\ell_n\|\longrightarrow\infty,
 \qquad
 \|\ell_n\|\ge c\,|p_n-p_*|^{-1/2}-C.
 \label{v33:eq:sequence-escape}
\end{equation}
Therefore every root sequence approaching the current first jet must either leave the base
solvability/rank stratum infinitely often or escape every bounded leading-tilt chart.
\end{corollary}
\begin{proof}
A constant-rank subsequence lies in finitely many local minor charts after shrinking around $p_*$;
apply \cref{v33:thm:escape} on one such chart and pass to a further subsequence.  If no such
base-rank solvable subsequence exists, the complementary alternative already holds.
\end{proof}

\begin{corollary}[Bounded-control local no-go]
\label[corollary]{v33:cor:bounded}
For every finite $M$ there is $\varepsilon_M>0$ such that, on every regular solvable incidence stratum,
\begin{equation}
 0<|p-p_*|<\varepsilon_M,
 \qquad \|\ell\|\le M,
 \qquad \mathscr L(p)\ell=d(p)
 \label{v33:eq:bounded-hyp}
\end{equation}
imply
\begin{equation}
                   \mathcal T_{\nu(p)}(Y(p),\ell)\ne0.
 \label{v33:eq:bounded-nogo}
\end{equation}
Thus the original analytic initializer on any fixed bounded multiplier-control neighborhood cannot
find a tangency root by an arbitrarily small perturbation of the V23 first jet.
\end{corollary}
\begin{proof}
Choose $\varepsilon_M$ so that the lower bound in \eqref{v33:eq:escape} exceeds $M$.
\end{proof}

This is stronger than merely observing that $\Delta_\diamond$ is nonzero at one point.  The global
flatness of the old eleven-tilt polynomial forces any regular-stratum cancellation generated by a
small first-jet perturbation to use a tilt whose magnitude diverges at least as the inverse square
root of that perturbation.

\section{The joint amplitude--first-jet scaling required by any escaping root}
\label{v33:sec:amplitude}

The V20--V23 tilted preparations have initial multiplier values of the form
\begin{equation}
                 \lambda(a,0)=e+a\ell+\text{higher amplitude terms}.
 \label{v33:eq:lambda-expansion}
\end{equation}
For bounded $\ell$ this is exactly the local analytic regime used in those bodies.  The escape
law shows that a root approaching $p_*$ cannot stay in that bounded-control regime.  One may still
ask whether a joint limit can let $\ell$ diverge while $a\ell$ remains small.  The following is a
necessary scale statement only; V23 does not provide the required initializer uniformly for
unbounded $\ell$.

\begin{corollary}[Necessary joint scale for a near-flat escaping family]
\label[corollary]{v33:cor:joint-scale}
Suppose a family of exact or candidate upstream roots lies on a regular incidence stratum,
$p(a)\to p_*$, and write $\delta_a=|p(a)-p_*|$.  If its leading tilt obeys the V23 lower
incidences and tangency equation, then
\begin{equation}
                    \|\ell(a)\|\ge c\,\delta_a^{-1/2}-C.
 \label{v33:eq:joint-lower}
\end{equation}
If, in addition, a genuinely constructed extension of the original initializer satisfies
\begin{equation}
 \|\lambda(a,0)-e-a\ell(a)\|=o(a\|\ell(a)\|)
 \quad\text{and}\quad
 \lambda(a,0)\longrightarrow e,
 \label{v33:eq:uniform-expansion}
\end{equation}
then necessarily
\begin{equation}
                    \boxed{\qquad \frac{a}{\sqrt{\delta_a}}\longrightarrow0,
                    \qquad\text{equivalently}\qquad
                    \frac{\delta_a}{a^2}\longrightarrow\infty .\qquad}
 \label{v33:eq:joint-scale}
\end{equation}
In particular a tangency-cancelling first-jet displacement of order $O(a^2)$ cannot remain in a
near-flat multiplier chart under \eqref{v33:eq:uniform-expansion}.  If
$\delta_a\asymp a^\kappa$, the necessary inequality is $\kappa<2$.
\end{corollary}
\begin{proof}
The first statement is \cref{v33:thm:escape}.  Under \eqref{v33:eq:uniform-expansion}, convergence
of $\lambda(a,0)$ to $e$ forces $a\|\ell(a)\|\to0$.  Combine this with
\eqref{v33:eq:joint-lower}.  The power-law restatement is immediate.
\end{proof}

The hypothesis \eqref{v33:eq:uniform-expansion} is intentionally exposed.  V23 proves analytic
initialization for every fixed finite control after shrinking the amplitude interval; it does
\emph{not} prove a chart radius uniform as $\|\ell\|\to\infty$.  Hence
\eqref{v33:eq:joint-scale} is a necessary condition on any future joint extension, not an existence
theorem for such escaping preparations.

\section{What to compute next after the escape theorem}
\label{v33:sec:next}

The updated bridge now separates three genuinely different routes.

\paragraph{First: test rank change on the five-dimensional chart.}
The escape theorem assumes that the lower system stays solvable and that the Poisson and incidence
ranks stay at their base values.  A candidate bounded root approaching $p_*$ therefore has only the
complementary alternative: reach a locus where solvability fails or one of those ranks changes.  This
is a finite determinantal/augmented-rank problem in the five explicit parameters of
\eqref{v33:eq:five-chart}.  It should be tested before any large Newton search, because the V23
null bundle itself changes there.

\paragraph{Second: search for finite-distance bounded-tilt roots.}
On a regular rank stratum, \cref{v33:cor:bounded} excludes roots in every sufficiently small
first-jet neighborhood with a prescribed control bound.  A bounded-control root, if it exists,
therefore occurs at a finite displacement from $p_*$; evaluating only the derivative at the old
seed cannot by itself certify such a root.  The V23 requested derivative remains useful for choosing
a search direction, but it is no longer a local existence argument.

\paragraph{Third: treat an escaping joint limit as a new cofinal problem.}
If a root is sought with $p\to p_*$ and $\|\ell\|\to\infty$, the required scale is at least
\eqref{v33:eq:joint-lower}, and a near-flat multiplier record must satisfy the separation
\eqref{v33:eq:joint-scale}.  The stationary chart, higher-preparation controls, signed Gram,
curvature reads and physical lifespan would all need estimates uniform in that growing tilt.  None
of V20--V23 supplies those estimates.

The old unit-initial-multiplier cofinal branch remains logically independent.  Version 32's
explicit stress angle $\Gamma_f$ and monotone Fourier certificates are still its correct
zero-parameter gate.  The V23 first-jet no-go and the present escape theorem concern the separate
order-amplitude tilted records and must not be used to infer a result for that old source.

No root of the five-parameter upstream system, cofinal tilted family, fixed positive Einstein
lifespan, higher weak closure, or primitive Renewal-to-Einstein--SM implication is asserted here.
The advance is the elimination of a whole class of local root searches: near the V23 seed, a
regular-stratum tangency root cannot remain in any bounded leading-tilt preparation chart.

\paragraph{Verification and preservation.}
Every Version 32 byte before its final document command is preserved.  The companion verifier for
this body checks the exact five-parameter mean surface and base point, the algebra of a generic
analytic quadratic fibre with vanishing base linear/quadratic coefficients, the square-root escape
inequality, the invariance of its exponent under uniformly invertible affine fibre-coordinate
changes, and the amplitude-scale implication.  It does not regenerate the V23 rational
$P_2,\Phi_3$ coefficient package or the exact value of $\Delta_\diamond$; those are inherited from
the newer bridge at their stated scope.  New labels use the prefix \texttt{v33:}.

\begingroup
\renewcommand{\refname}{Additional source for Version 33}
\begin{thebibliography}{9}
\small\raggedright
\bibitem{v33:EB23}
\emph{Renewal Exact-Action Einstein Bridge Research Notes}, supplied cumulative Versions 1--23,
file \nolinkurl{Renewal_Exact_Action_Einstein_Bridge_Research_Notes_V23.tex}.
The inherited interfaces used here are the V23 global eleven-tilt no-go, its five-dimensional
weak-neutral first-jet chart, and its finite upstream incidence system, together with the V20--V22
Poisson-rank, preparation and tangency formulas retained inside that cumulative source.  No new
rational coefficient value from that source is represented as independently regenerated here.
\end{thebibliography}
\endgroup

\addtocontents{toc}{\protect\endgroup}
% END IMMUTABLE VERSION 33 BODY
% APPEND VERSION 34 HERE, BEFORE THE FINAL END-DOCUMENT COMMAND.



% BEGIN IMMUTABLE VERSION 34 BODY
\clearpage
\begin{center}
{\Large\bfseries Version 34: Blow-up of the V23 tangency fibre}\\[0.5em]
{\large The first escape tensor, the sharp square-root cone, and an upgraded upstream gate}
\end{center}
\makeatletter
\addtocontents{toc}{\protect\begingroup}
\addtocontents{toc}{\protect\makeatletter}
\addtocontents{toc}{\protect\def\protect\l@section{\protect\bfseries\protect\@dottedtocline{1}{0em}{2.5em}}}
\makeatother
\addcontentsline{toc}{part}{Version 34: Blow-up of the V23 tangency fibre and the first escape tensor}

\section{Audit and target after the square-root escape theorem}
\label{v34:sec:audit}

A renewed Library audit found no supplied Exact-Action Einstein Bridge continuation beyond cumulative
V23.  We therefore continue from the literal V23 five-parameter upstream system and from Version
33's local regular-stratum escape theorem.  The latter gives a lower bound of order
$|p-p_*|^{-1/2}$ for any tangency root that approaches the old first jet while the complete lower
incidence remains solvable at the base ranks.  The present body asks the next sharper question:
\emph{what finite equation is seen after rescaling exactly at that minimal escape rate?}

Work on one of the regular solvable incidence charts of Version 33.  Its analytically transported
lower-incidence fibre can be written
\begin{equation}
             \ell=\ell_0(p)+E(p)c,\qquad c\in\mathbb R^{11},
 \label{v34:eq:fibre}
\end{equation}
with $E(p)$ injective and with the transported weak Poisson-null test denoted by $\nu(p)$.  Restrict
its scalar tangency observation to this fibre:
\begin{equation}
 \Delta(p,c):=\mathcal T_{\nu(p)}\bigl(Y(p),\ell_0(p)+E(p)c\bigr).
 \label{v34:eq:Delta}
\end{equation}
The V22 tangency expression is at most quadratic in the leading tilt, hence on this analytic fibre
there are analytic coefficient functions
\begin{equation}
 \Delta(p,c)=d_0(p)+d_1(p)^Tc+c^TD_2(p)c,
 \qquad D_2(p)=D_2(p)^T.
 \label{v34:eq:quadratic}
\end{equation}
Version 23's exact global eleven-tilt calculation gives, at the base point,
\begin{equation}
       d_0(p_*)=\Delta_\diamond\ne0,\qquad
       d_1(p_*)=0,\qquad D_2(p_*)=0.
 \label{v34:eq:base-flat}
\end{equation}
These identities, not merely the nonzero value of $\Delta_\diamond$, are the input which creates the
blow-up structure below.  No new value of the inherited rational coefficient is claimed here.

\section{The first escape tensor and the sharp square-root compactification}
\label{v34:sec:tensor}

\begin{definition}[First escape tensor]
\label[definition]{v34:def:tensor}
For $v\in T_{p_*}\mathbb R^5\cong\mathbb R^5$ define the symmetric form
\begin{equation}
     \mathfrak B(v):=D_pD_2(p_*)[v]\in\operatorname{Sym}(11),
 \qquad
     \mathfrak B(v)[y,y]=y^T\mathfrak B(v)y.
 \label{v34:eq:B}
\end{equation}
It is linear in $v$.  We call $\mathfrak B$ the \emph{first escape tensor} of the transported V23
tangency test on the solved lower-incidence fibre.
\end{definition}

The tensor is a derivative \emph{after} the lower equations, normalization rows and cubic constant
rows have been solved and the null test has been transported.  Differentiating the unsolved V23
formula while holding the old null basis fixed would therefore compute a different object.

\begin{theorem}[Sharp square-root blow-up equation]
\label[theorem]{v34:thm:blowup}
Let $(p_n,c_n)$ be roots of \eqref{v34:eq:Delta} on a fixed regular solvable incidence chart, with
$p_n\to p_*$.  Put
\[
 \delta_n=|p_n-p_*|,\qquad v_n=\frac{p_n-p_*}{\delta_n}.
\]
Assume, after passage to a subsequence,
\begin{equation}
                   v_n\longrightarrow v,\qquad
                   y_n:=\sqrt{\delta_n}\,c_n\longrightarrow y
 \label{v34:eq:scaled-conv}
\end{equation}
with $y$ finite.  Then $|v|=1$, $y\ne0$, and
\begin{equation}
            \boxed{\qquad
              \Delta_\diamond+\mathfrak B(v)[y,y]=0.
            \qquad}
 \label{v34:eq:cone}
\end{equation}
Thus every root which attains the minimal Version-33 escape scale has a limit point in the finite
algebraic \emph{first escape cone}
\begin{equation}
 \mathscr C_{\rm esc}
   =\{(v,y)\in S^4\times(\mathbb R^{11}\setminus\{0\}):
                   \Delta_\diamond+\mathfrak B(v)[y,y]=0\}.
 \label{v34:eq:escape-cone}
\end{equation}
\end{theorem}
\begin{proof}
Analyticity and \eqref{v34:eq:base-flat} give, uniformly for $v$ in the unit sphere,
\begin{align*}
 d_0(p_*+\delta v)&=\Delta_\diamond+O(\delta),\\
 d_1(p_*+\delta v)&=\delta\,Dd_1(p_*)[v]+o(\delta),\\
 D_2(p_*+\delta v)&=\delta\,\mathfrak B(v)+o(\delta).
\end{align*}
Insert $c_n=\delta_n^{-1/2}y_n$ in \eqref{v34:eq:quadratic}.  The constant correction is
$O(\delta_n)$ and the linear term is $O(\sqrt{\delta_n})$ because $y_n$ is bounded.  For the
quadratic term,
\[
 c_n^TD_2(p_n)c_n
   =y_n^T\mathfrak B(v_n)y_n+o(1).
\]
Passing to the limit in $\Delta(p_n,c_n)=0$ proves \eqref{v34:eq:cone}.  The limit cannot have
$y=0$ because $\Delta_\diamond\ne0$.
\end{proof}

The point of \eqref{v34:eq:cone} is that neither $Dd_0(p_*)$ nor $Dd_1(p_*)$ contributes at the
minimal escape scale.  The first finite quantity capable of cancelling the inherited defect is the
first $p$-derivative of the \emph{quadratic fibre coefficient}.  This is a stronger reduction than a
full derivative of the unsolved tangency quotient.

\begin{proposition}[The scalar escape cone is nonempty exactly when the tensor is nonzero]
\label[proposition]{v34:prop:cone}
For the single transported V23 test used above,
\begin{equation}
       \boxed{\quad \mathscr C_{\rm esc}\ne\varnothing
               \quad\Longleftrightarrow\quad
               \mathfrak B\not\equiv0.\quad}
 \label{v34:eq:cone-nonempty}
\end{equation}
If $\mathfrak B(v_0)[z,z]\ne0$, one of $v_0$ or $-v_0$ has quadratic value opposite in sign to
$\Delta_\diamond$, and a scalar rescaling of $z$ then gives a point of \eqref{v34:eq:escape-cone}.
\end{proposition}
\begin{proof}
The forward implication follows immediately from \eqref{v34:eq:cone}.  Conversely choose
$v_0,z$ with $q=\mathfrak B(v_0)[z,z]\ne0$.  Since $\mathfrak B(-v_0)=-\mathfrak B(v_0)$, replace
$v_0$ by $-v_0$ if necessary so that $q\Delta_\diamond<0$.  Then
$y=\sqrt{|\Delta_\diamond/q|}\,z$ satisfies \eqref{v34:eq:cone}.
\end{proof}

This proposition is only a statement about the one scalar necessary test inherited from V23.  A
nonempty scalar escape cone is not a solution of the complete vector-valued upstream tangency
system, does not restore the lower incidence at a new first jet, and does not establish an exact
history.

\section{Coordinate invariance and the finite derivative package}
\label{v34:sec:invariance}

The fibre coordinates in \eqref{v34:eq:fibre} are not canonical.  A useful certificate must not
depend on the particular rational basis chosen at the old seed.

\begin{proposition}[Intrinsic nature of the first escape tensor]
\label[proposition]{v34:prop:invariant}
Let
\begin{equation}
               c=A(p)\widetilde c+b(p)
 \label{v34:eq:fibre-change}
\end{equation}
be an analytic change of fibre coordinates with $A(p_*)$ invertible and $b$ bounded.  Let
$p=\chi(\widetilde p)$ be an analytic first-jet reparametrization with invertible derivative at the
base point.  Then the first escape tensor transforms by
\begin{equation}
 \widetilde{\mathfrak B}(\widetilde v)
  =A_*^T\,\mathfrak B(D\chi_*\widetilde v)\,A_*,
 \qquad A_*=A(p_*).
 \label{v34:eq:B-transform}
\end{equation}
Consequently the assertions $\mathfrak B\equiv0$ and $\mathfrak B\not\equiv0$, as well as the
existence of the square-root escape cone, are independent of every such analytic trivialization of
the lower-incidence fibre and of the five-parameter first-jet chart.
\end{proposition}
\begin{proof}
The quadratic coefficient after \eqref{v34:eq:fibre-change} is
$\widetilde D_2(p)=A(p)^TD_2(p)A(p)$.  Differentiate at $p_*$ and use $D_2(p_*)=0$: derivatives of
$A$ multiply the vanishing base coefficient and drop out.  Translation by $b$ changes only the
constant and linear coefficients and does not change the quadratic coefficient.  The parameter
change contributes $D\chi_*$ by the chain rule.  This proves \eqref{v34:eq:B-transform}; congruence
by an invertible matrix and an invertible base reparametrization preserve whether the tensor
vanishes identically.  The corresponding cone variables satisfy $y=A_*\widetilde y$.
\end{proof}

In any one analytic fibre chart, deciding the first escape tensor requires only
\begin{equation}
           B_a:=\partial_{p_a}D_2(p_*)\in\operatorname{Sym}(11),
           \qquad 1\le a\le5.
 \label{v34:eq:Ba}
\end{equation}
There are $5\times66=330$ symmetric scalar entries.  The relevant directional form is
$\mathfrak B(v)=\sum_av_aB_a$.  At the square-root scale no derivative of $d_0$ or $d_1$ enters the
limit equation.  Thus the next exact coefficient computation can be smaller than differentiating the
complete tangency map with all fibre variables left unsolved.

\section{If the first escape tensor vanishes, the control blow-up strengthens}
\label{v34:sec:higher}

The tensor may vanish.  In that case the square-root law of Version 33 is not sharp.

\begin{theorem}[Vanishing first escape tensor upgrades the lower bound]
\label[theorem]{v34:thm:linear-escape}
Suppose $\mathfrak B\equiv0$ on a regular solvable incidence neighborhood of $p_*$.  Then, after
shrinking the neighborhood, every root of \eqref{v34:eq:Delta} satisfies
\begin{equation}
                 \boxed{\qquad |c|\ge c_0|p-p_*|^{-1}.\qquad}
 \label{v34:eq:linear-escape}
\end{equation}
Consequently
\begin{equation}
                 \|\ell\|\ge c_1|p-p_*|^{-1}-C_1.
 \label{v34:eq:linear-escape-ell}
\end{equation}
\end{theorem}
\begin{proof}
Put $\delta=|p-p_*|$.  Analyticity, \eqref{v34:eq:base-flat}, and
$D_pD_2(p_*)=0$ give
\[
 |d_0(p)-\Delta_\diamond|\le C\delta,
 \qquad |d_1(p)|\le C\delta,
 \qquad \|D_2(p)\|\le C\delta^2.
\]
For a root and sufficiently small $\delta$,
\[
 \frac{|\Delta_\diamond|}{2}
   \le C\{\delta+\delta|c|+\delta^2|c|^2\}.
\]
If both $\delta|c|$ and $\delta^2|c|^2$ were smaller than a sufficiently small fixed constant, the
right side would contradict the left side.  Hence $\delta|c|\ge c_0$ after changing constants.
The physical-tilt estimate follows from the uniform injectivity of $E(p)$ exactly as in Version 33.
\end{proof}

\begin{corollary}[Stronger joint amplitude--first-jet separation when $\mathfrak B=0$]
\label[corollary]{v34:cor:joint}
Under the hypotheses of \cref{v34:thm:linear-escape}, suppose a future joint construction has
$p(a)\to p_*$, remains on the regular solvable incidence stratum, satisfies the tangency equation,
and has a uniform leading-multiplier expansion in the sense of Version 33,
\[
 \|\lambda(a,0)-e-a\ell(a)\|=o(a\|\ell(a)\|),
 \qquad \lambda(a,0)\to e.
\]
Then, with $\delta_a=|p(a)-p_*|$,
\begin{equation}
                    \boxed{\qquad \frac{a}{\delta_a}\longrightarrow0.\qquad}
 \label{v34:eq:strong-joint}
\end{equation}
In particular $\delta_a\asymp a^\kappa$ would require $\kappa<1$.
\end{corollary}
\begin{proof}
Near-flatness gives $a\|\ell(a)\|\to0$, while
\eqref{v34:eq:linear-escape-ell} gives $\|\ell(a)\|\gtrsim\delta_a^{-1}$.
\end{proof}

More generally, if along a specified analytic first-jet arc the fibre coefficients obey
$|d_1|=O(\delta^r)$ and $\|D_2\|=O(\delta^s)$ with integers $r,s\ge1$, the same one-line comparison
gives the necessary order
\begin{equation}
                  |c|\gtrsim \delta^{-\min\{r,s/2\}}.
 \label{v34:eq:newton-order}
\end{equation}
This is only a lower-order Newton-polygon diagnostic; equality requires the corresponding leading
coefficient equation to have a solution.

\section{The next exact calculation is now sharply localized}
\label{v34:sec:next}

The V23 five-parameter search and Version 33's rank-change-or-escape dichotomy can now be organized
without another broad root search.

\paragraph{1. Compute the five first-escape matrices.}
Regenerate the solved lower-incidence fibre and transported null test on the V23 five-parameter
chart and evaluate the $330$ entries in \eqref{v34:eq:Ba}.  The source already supplies exact
$P_2$ and $\Phi_3$ polarization machinery at the base cut; what is still needed is its derivative
through the \emph{moving solved fibre}.  A single nonzero entry proves
$\mathfrak B\not\equiv0$ and produces a finite square-root escape direction for the scalar test via
\cref{v34:prop:cone}.  Vanishing of all entries upgrades the necessary blow-up to
\eqref{v34:eq:linear-escape}.

\paragraph{2. Keep the complete upstream system separate from the scalar cone.}
Even when $\mathscr C_{\rm esc}$ is nonempty, the other transported weak-null tangency components,
Poisson rank, lower-incidence augmented rank, signed source Gram and preparation-control ranks must
all be imposed at the same first jet.  The cone only identifies the leading scaling at which the
already certified V23 scalar obstruction can first be cancelled.

\paragraph{3. Test the rank-change locus in parallel.}
A bounded root approaching $p_*$ is already excluded on every regular solvable stratum by Version
33.  The determinantal and augmented-rank locus of the five-parameter lower system therefore remains
the only local bounded-control alternative.  The first escape tensor does not replace that test.

\paragraph{4. Do not conflate the old cofinal stress branch.}
Version 32's source-angle certificate $\Gamma_f$ concerns the unit-initial-multiplier cofinal branch.
The blow-up tensor here concerns the different order-amplitude tilted branch of V20--V23.  Neither
criterion proves the other.

No root of the five-parameter V23 upstream system, no cofinal escaping family, no cutoff-uniform
stationary chart, no common positive Einstein lifespan, and no primitive Renewal-to-Einstein--SM
implication is claimed.  The advance is a compactification of the only regular near-seed escape
allowed by Version 33 and a reduction of its first nontrivial coefficient calculation to five
symmetric $11\times11$ matrices.

\paragraph{Verification and preservation.}
Every Version 33 byte before its final document command is preserved.  The companion verifier for
this body checks the square-root limit identity on generic analytic quadratic fibres, the
nonempty-cone equivalence, congruence invariance under analytic fibre-coordinate changes, the
$|p-p_*|^{-1}$ upgrade when the first escape tensor vanishes, and the associated amplitude-scale
implication.  It does not regenerate the V23 rational $P_2,\Phi_3$ package or claim values for the
$330$ derivative entries.  New labels use the prefix \texttt{v34:}.

\begingroup
\renewcommand{\refname}{Additional source for Version 34}
\begin{thebibliography}{9}
\small\raggedright
\bibitem{v34:EB23}
\emph{Renewal Exact-Action Einstein Bridge Research Notes}, supplied cumulative Versions 1--23,
file \nolinkurl{Renewal_Exact_Action_Einstein_Bridge_Research_Notes_V23.tex}.
The inherited interfaces are the global V23 eleven-tilt no-go, the five-dimensional weak-neutral
first-jet chart, and the complete finite upstream incidence problem.  Their exact coefficients are
used only at the scopes stated in that source.
\end{thebibliography}
\endgroup

\addtocontents{toc}{\protect\endgroup}
% END IMMUTABLE VERSION 34 BODY
% APPEND VERSION 35 HERE, BEFORE THE FINAL END-DOCUMENT COMMAND.



% BEGIN IMMUTABLE VERSION 35 BODY
\clearpage
\begin{center}
{\Large\bfseries Version 35: A surviving source test across rank changes}\\[0.5em]
{\large Exact original coefficients, a changing incidence rank, and rank-independent escape}
\end{center}
\makeatletter
\addtocontents{toc}{\protect\begingroup}
\addtocontents{toc}{\protect\makeatletter}
\addtocontents{toc}{\protect\def\protect\l@section{\protect\bfseries\protect\@dottedtocline{1}{0em}{2.5em}}}
\makeatother
\addcontentsline{toc}{part}{Version 35: A surviving source test across rank changes}

\section{The prerequisite for differentiating a solved fibre}
\label{v35:sec:context}

The five matrices in \eqref{v34:eq:Ba} presuppose a differentiable solved
lower-incidence fibre with eleven free coordinates and a transported weak
Poisson-null test. Neither premise follows merely from the five-dimensional
initial-data chart. The present calculation returns to the actual finite
coefficient functional before differentiating that fibre. There are two
source-specific findings. The coefficient ranks do increase along an explicit
arc through the old seed, so the eleven-coordinate fibre cannot be continued
over an open five-parameter neighborhood. Nevertheless the particular weak
null test used to exclude the old fibre survives on the entire first-jet
family. That survival makes it possible to prove the escape estimate without
any constant-rank or solvability-stratum hypothesis.

The result has a precise scope. It concerns the complementary three-site
first-jet chart of the supplied Exact-Action Einstein Bridge, not a cofinal
spatial family. It excludes bounded controls approaching its old seed even
through rank-changing strata. It does not produce a tangency root, prohibit
all finite-distance roots, or settle the distinct continuum stress equation
of Version 32. Selecting multipliers to solve finite constraints and
propagating an additional preparation condition are different questions;
the general constrained-discretization background is
\cite{v35:DBGP}, not a source for the coefficient certificates below.

\paragraph{Inherited facts used at their stated scope.}
The Bridge V20 defines the quadratic constraint coefficient $P_{\nu,2}$,
its first Poisson matrix $\mathbb K_1$, and the prepared quadratic drift.
Bridge V21 supplies the cubic constant-row map and the lower-incidence
rank $97$ at the base point. Bridge V23 proves that, on the complete base
lower-incidence fibre, one fixed weak-null tangency test has the constant
value $\Delta_\diamond\ne0$. These facts are retained in Versions 33--34.
The new exact computation reconstructs their displayed finite coefficients,
checks the base residues, and proves the fixed-test identity and nearby
rank opening stated below. It does not reprove the inherited nonlinear
stationary-chart or initialization theorems.

\section{The actual first-jet family and its fixed diagnostic}
\label{v35:sec:family}

Set $h=1/3$, use the 27 original sites in lexicographic order, and retain
$\delta_i=3(S_i-S_i^{-1})$, where $S_i$ shifts the $i$th coordinate forward
by one site. Let $c_i(x)=\cos(2\pi x_i)$, whose three sampled values are
$(1,-1/2,-1/2)$. On the two coordinates $j<k$ different from $i$, $T_i$
has off-diagonal entries $(T_i)_{jk}=(T_i)_{kj}=1$ and $P_i$ has diagonal
entries $(P_i)_{jj}=1$, $(P_i)_{kk}=-1$. All other entries vanish.

For six real coefficients $(\alpha,b)$ put
\begin{equation}
 \begin{split}
 U(\alpha)&=\sum_{i=1}^3\alpha_i T_i c_i,\\
 p_R(\alpha,b)&=3\{\operatorname{diag}(b)-(b_1+b_2+b_3)I\}
                    +\tfrac32\sum_{i=1}^3\alpha_iP_i c_i,\\
 Y(\alpha,b)&=(U(\alpha),\sqrt3\,p_R(\alpha,b)).
 \end{split}
 \label{v35:eq:six-family}
\end{equation}
The actual five-dimensional, mean-compatible chart has parameter
$\zeta=(\alpha_1,\alpha_2,\alpha_3,b_2,b_3)$ and
\begin{equation}
 b_1(\zeta)=\frac{(\alpha_1^2+\alpha_2^2+\alpha_3^2)/4-b_2b_3}{b_2+b_3},
 \qquad
 \zeta_*=(1,9/8,1,-8,-8).
 \label{v35:eq:chart}
\end{equation}
Thus $b_{1,*}=16175/4096$. The denominator is bounded away from zero on
a sufficiently small fixed neighborhood of $\zeta_*$. The inherited
linear constraints, quadratic weak neutrality and four quadratic mean
balances hold on this chart.

A multiplier is $\nu=(n,\beta)$, and its flat canonical vector is
\begin{equation}
 G\nu=(\operatorname{SymGrad}_\delta\beta,
       \nabla_\delta^2n-I\Delta_\delta n),\qquad
 \mathbb K_1(Y)[\nu,\mu]
 =DP_{\nu,2}(Y)[G\mu]-DP_{\mu,2}(Y)[G\nu].
 \label{v35:eq:K-G}
\end{equation}
Here $\operatorname{SymGrad}_\delta\beta$ has entries
$\delta_i\beta_j+\delta_j\beta_i$, without an additional factor $1/2$.
Define the fixed, nonconstant weak test
\begin{equation}
 \phi_\diamond=\tfrac23(c_1-c_2),\qquad
 \nu_\diamond=(0,\nabla_\delta\phi_\diamond).
 \label{v35:eq:diamond}
\end{equation}
Its original site values are rational; its shift components have values
$0,\pm3$. Its shift mean and curl vanish and
$\|\nu_\diamond\|_h^2=12$.

\begin{theorem}[The same weak-null test survives the entire first-jet family]
\label[theorem]{v35:thm:fixed-null}
For every real $(\alpha,b)$ in \eqref{v35:eq:six-family},
\begin{equation}
             \boxed{\mathbb K_1(Y(\alpha,b))\nu_\diamond=0.}
 \label{v35:eq:fixed-null}
\end{equation}
In particular this identity holds throughout \eqref{v35:eq:chart},
including where the Poisson rank changes. It does not require the mean
surface, a lower-incidence solution, or a chosen nullspace basis.
\end{theorem}
\begin{proof}
The quadratic constraint coefficient is homogeneous of field degree two,
so \eqref{v35:eq:K-G} is linear in $Y$. It suffices to evaluate it on
the six spanning fields. In rational coordinates these are
\begin{equation}
 \begin{array}{ll}
 Y^{\alpha_i}_R=(T_ic_i,\tfrac32P_ic_i),&1\le i\le3,\\
 Y^{b_i}_R=(0,3(e_ie_i^T-I)),&1\le i\le3.
 \end{array}
 \label{v35:eq:bank}
\end{equation}
The exact finite calculation specified in \cref{v35:sec:compiler}
gives zero on all 108 multiplier rows for each of these six fields:
\begin{equation}
 \bigl(\mathbb K_1(Y_R^a)\nu_\diamond\bigr)_j=0
 \quad(1\le a\le6,\ 0\le j<108).
 \label{v35:eq:648}
\end{equation}
These 648 equalities are evaluated over $\mathbb Q$, not inferred from
modular zeros or a floating rank tolerance. The compiler uses the explicit
$P_{\nu,2}$ below and its exact quadratic polarization.

The physical factor $\sqrt3$ is retained as follows. For a pure-shift
column, the lapse output block $K_{N\beta}$ depends only on $U$ and the
shift output block $K_{\beta\beta}$ depends only on the canonical
momentum, linearly. This follows directly from the field grading of
\eqref{v35:eq:P2}, and is the same grading used in Bridge V20's
rationalization. Consequently the zero lapse and shift outputs in
\eqref{v35:eq:648} remain zero when the momentum is multiplied by
$\sqrt3$. Linearity now proves \eqref{v35:eq:fixed-null} for all six
real coefficients. The norm and weak-space assertions follow by summing
the displayed rational grid values with weight $1/27$.
\end{proof}

\section{Coefficient regeneration and exact finite certificates}
\label{v35:sec:compiler}

For reproducibility, the quadratic functional used for all new zero and
rank tests is recorded explicitly. All expressions in this section use
spatial sums; the physical Frobenius pairing counts both off-diagonal
entries. Write a Lorentz coefficient as $(R,B)$, with bracket
\[
 [(R,B),(S,C)]=(-R\times S+B\times C,-R\times C-B\times S).
\]
For $Y=(U,p)$ set
\[
 z_i=(R_i,B_i),\qquad
 R_{ia}=-\tfrac12\epsilon_{abc}\delta_bU_{ic},\qquad
 B=p-\tfrac12(\tr p)I,
\]
\[
 \Pi_1=(\tr U)I-U,\qquad
 \Pi_2=\tfrac34U^2-\tfrac12(\tr U)U+
       \{\tfrac14(\tr U)^2-\tfrac12\tr(U^2)\}I.
\]
For $\nu=(n,\beta)$ retain $a=(-\operatorname{curl}_\delta\beta/2,
\nabla_\delta n)$ and $w=\operatorname{SymGrad}_\delta\beta$. For $i<j$
and $\epsilon=\epsilon_{ijk}$, let
\[
 \Sigma_{0,ij}=(2\epsilon ne_k,2\beta_je_i-2\beta_ie_j),\qquad
 \Sigma_{1,ij}=(\epsilon nU_k,\beta_j(\Pi_1)_i-\beta_i(\Pi_1)_j).
\]
The retained scalar coefficient is
\begin{equation}
 \begin{split}
 -P_{\nu,2}(Y)={}&-\langle D\Pi_2(U)[w],B\rangle
 +\sum_i\langle\delta_i(\Pi_2)_i,a_B\rangle
 +\sum_i\langle(\Pi_1)_i,[a,z_i]_B\rangle\\
 &-h\sum_i\langle e_i,[z_i,[z_i,S_ia]]_B\rangle\\
 &+\sum_{i<j}\bigl\{
 \langle\Sigma_{1,ij},\delta_i z_j-\delta_j z_i\rangle
 +\langle\Sigma_{0,ij},[z_i,z_j]\rangle\bigr\}.
 \end{split}
 \label{v35:eq:P2}
\end{equation}
In particular the shifted cubic-electric contribution in its second line
is retained. Here
\[
 D\Pi_2(U)[w]=\tfrac34(Uw+wU)-\tfrac12\{(\tr w)U+(\tr U)w\}
 +\{\tfrac12\tr U\tr w-\tr(Uw)\}I.
\]
No finite-difference derivative is used: for the quadratic field functional,
\begin{equation}
 DP_{\nu,2}(Y)[Z]
       =\tfrac12\{P_{\nu,2}(Y+Z)-P_{\nu,2}(Y-Z)\}.
 \label{v35:eq:polarization}
\end{equation}
For example, let $P_2(Y)$ denote its full 108-component multiplier
covector. The entire column in \eqref{v35:eq:648} is evaluated by
\begin{equation}
 DP_2(Y)[G\nu_\diamond]
       -\bigl(DP_{\nu_\diamond,2}(Y)[Ge_j]\bigr)_{j=0}^{107}.
 \label{v35:eq:column-evaluation}
\end{equation}
Equations \eqref{v35:eq:bank}--\eqref{v35:eq:column-evaluation}, ordinary
cyclic array shifts and rational arithmetic specify all 648 entries.
The accompanying verifier performs these sums with arbitrary-precision
fractions. It stores one exact zero verdict per field, with the number
of tested original rows, rather than a list of rounded small values.

The joint rank test also uses the cubic coefficients. They are not
replaced by the quadratic Poisson matrix. Put $V(p)=2p-(\tr p)I$. The
retained unit-multiplier cubic functional is
\begin{equation}
 \Phi_3(Y)=-\langle D\Pi_2(U)[V(p)],B\rangle+
 \sum_{i<j}\epsilon_{ijk}\bigl\{
 -\tfrac14\langle(U^2)_k,\delta_iR_j-\delta_jR_i\rangle
 +\langle U_k,[z_i,z_j]_R\rangle
 +2h\langle e_k,\mathcal B_{3,ij}(z)_R\rangle\bigr\}.
 \label{v35:eq:Phi3}
\end{equation}
The BCH term is that of the original ordered factors
$(z_i,S_i z_j,-S_jz_i,-z_j)$. With $l_1=l_2=l_3=0$, each successive
factor $v$ updates the old values by
\[
 l_3\leftarrow l_3+\tfrac12[l_2,v]
 +\tfrac1{12}\{[l_1,[l_1,v]]+[v,[v,l_1]]\},\quad
 l_2\leftarrow l_2+\tfrac12[l_1,v],\quad l_1\leftarrow l_1+v.
\]
Cubic differentiation uses
$DF(Y)[Z]=\{F(Y+Z)-F(Y-Z)\}/2-F(Z)$.
The vector $F_1Y=(V(p),\mathcal F(U))$ is
\[
 \mathcal F(U)_{ij}=\tfrac12\Delta_\delta U_{ij}
 -\tfrac12(\delta_i v_j+\delta_jv_i)
 +\tfrac12\delta_i\delta_j\tr U
 +\tfrac12\delta_{ij}(\delta_kv_k-\Delta_\delta\tr U),
 \qquad v_j=\delta_iU_{ij}.
\]
Hence the full prepared quadratic target is
\begin{equation}
 \widehat B_2(Y)[\nu]
 =D\Phi_3(Y)[G\nu]+DP_{\nu,2}(Y)[F_1Y]
                    +P_{(0,\nabla_\delta n),2}(Y).
 \label{v35:eq:B2}
\end{equation}
The last term is not discarded.

For a constant shift $e_c$, its original cubic constraint coefficient
$P_{c,3}$ is obtained by extracting minus the amplitude-three coefficient
of the same stationary scalar functional with
\[
 (z,A_0,w,u,p,N,\beta)
 =(az_1(Y),az_1(Y)_c,a\delta_cU,aU,ap,0,e_c).
\]
This includes the temporal auxiliary connection. Equivalently the verifier
expands the cofactor through $\Pi_2$, keeps $\Sigma_2$ against the linear
spatial curvature, $\Sigma_1$ against $[z_i,z_j]$, and $\Sigma_0$ against
the ordered cubic BCH term. The constant-row data are precisely
\begin{equation}
 \begin{split}
 t_c(Y)&=D\Phi_3(Y)[\delta_cY]+DP_{c,3}(Y)[F_1Y],\\
 \mathcal R_Y[c,\ell]&=DP_{c,3}(Y)[G\ell]
 -DP_{\ell,2}(Y)[\delta_cY]-\langle\delta_c\ell,P_2(Y)\rangle.
 \end{split}
 \label{v35:eq:Rc}
\end{equation}
The last preparation term is retained by exact cyclic summation by parts.
These formulas are inherited coefficient identities; their independent
implementation and the additional evaluations are the new finite work.

\section{An actual rank-opening arc through the old seed}
\label{v35:sec:rank}

Use rational multiplier coordinates $\ell=Q_3x$ with
$Q_3=\operatorname{diag}(I_{27},\sqrt3 I_{81})$. Put
\begin{equation}
 \begin{split}
 \widehat K(\zeta)&=Q_3\mathbb K_1(Y(\zeta))Q_3/\sqrt3,\\
 \widehat b(\zeta)&=Q_3\widehat B_2(Y(\zeta))/\sqrt3,\qquad
 \widehat R(\zeta)=\mathcal R_{Y(\zeta)}Q_3,\\
 \mathcal L(\zeta)&=
 \begin{pmatrix}\widehat K(\zeta)\\ \mathsf m\\ \widehat R(\zeta)\end{pmatrix},
 \qquad d(\zeta)=\begin{pmatrix}-\widehat b(\zeta)\\0_4\\-t(Y(\zeta))\end{pmatrix}.
 \end{split}
 \label{v35:eq:joint}
\end{equation}
The four rows of $\mathsf m$ take the component spatial sums. Dividing
these rows by 27 gives the equivalent zero-mean normalization. Thus
$\mathcal Lx=d$ is the complete lower-incidence problem, with 115 rows
and 108 columns, not just a subset of the Poisson equations.

Consider the real analytic arc
\begin{equation}
 \zeta(t)=(1+t,9/8,1,-8,-8),\qquad
 b_1(t)=\frac{16175}{4096}-\frac{t}{32}-\frac{t^2}{64}.
 \label{v35:eq:arc}
\end{equation}
Let $W$ synthesize the 29-dimensional original weak shift space: its first
three columns are constant shifts, followed by
$(0,\nabla_\delta(e_j-e_{26}))$, $0\le j<26$.

\begin{theorem}[The eleven-coordinate incidence fibre is not open in five parameters]
\label[theorem]{v35:thm:rank}
For all sufficiently small nonzero real $t$,
\begin{equation}
 \boxed{\operatorname{rank}\widehat K(\zeta(t))\ge92,\qquad
 \operatorname{rank}(\widehat K(\zeta(t))W)\ge20,\qquad
 \operatorname{rank}\mathcal L(\zeta(t))\ge99.}
 \label{v35:eq:rank-lower}
\end{equation}
At $t=0$ the respective ranks are $90,18,97$. Consequently, wherever
the lower equations on this punctured arc are compatible, their solution
fibre has dimension at most nine; the weak Poisson-null space also has
dimension at most nine. A rank-$97$ joint system with an eleven-dimensional
solution fibre therefore does not extend to an open five-parameter
neighborhood of the old seed.
\end{theorem}
\begin{proof}
All matrix entries along \eqref{v35:eq:arc} are rational polynomials in
$t$. The field grading makes \eqref{v35:eq:joint} rational; alternatively
one may keep $\sqrt3$ in its exact quadratic number field. Each fixed
minor is thus an analytic function of one real variable. It suffices to
show the three minors specified below are not identically zero.

Let
\[
 \mathcal F=(8,17,20,23,24,25,26,35,44,53,74,77,80,89,98,105,106,107),
 \qquad \mathcal J=\{0,\ldots,107\}\setminus\mathcal F.
\]
For the Poisson minor use the 92-index principal set
$\mathcal J\cup\{89,98\}$. For the joint minor use the following
99 rows and columns. An interval denotes all its integer indices.
\begin{equation}
 \begin{split}
 \mathcal I_{99}={}&(0{:}7,9{:}16,18{:}19,21{:}22,27{:}34,36{:}43,45{:}52,\\
 &54{:}73,75{:}76,78{:}79,81{:}104,108{:}114),\\
 \mathcal J_{99}={}&(0{:}23,27{:}43,45{:}52,54{:}76,78{:}79,81{:}105).
 \end{split}
 \label{v35:eq:joint-minor-indices}
\end{equation}
For the weak matrix use the row and column sets
\[
 \begin{aligned}
 \mathcal I_{20}&=(0{:}1,3{:}4,9{:}10,12{:}13,27{:}30,36{:}39,54{:}55,57{:}58),\\
 \mathcal J_{20}&=(3{:}19,21{:}22,24).
 \end{aligned}
\]
The original-coefficient compiler gives these exact nonzero residues at $t=1$:
\begin{center}\small
\begin{tabular}{@{}lrr@{}}\toprule
Actual minor & modulo 241 & modulo 1009\\\midrule
Poisson principal $92\times92$ & $40$ & $725$\\
Weak restriction $20\times20$ & $83$ & $991$\\
Complete lower incidence $99\times99$ & $230$ & $284$\\\bottomrule
\end{tabular}
\end{center}
All coefficient denominators are units at both primes. The values of
$\sqrt3$ used to check the original formulas are respectively $56$ and
$149$; their squares equal three in the corresponding fields. A nonzero
residue proves the exact minor is nonzero in characteristic zero. The
second prime is an independent regression, not an additional assumption.

A nonzero one-variable polynomial has an isolated zero, if any, at zero.
Thus each of these three minors is nonzero for all sufficiently small
$t\ne0$, proving \eqref{v35:eq:rank-lower}. The evaluation at $t=1$
serves only to prove that a polynomial is not identically zero; no physical
initialization or trajectory at that distant parameter is inferred.
The ranks at zero use the inherited exact nullspace and mean-row
certificates. The regeneration reproduces the old Poisson principal minor
$25$ and the old cubic-mean control minor $37$ modulo 241. Finite rank
then bounds the dimensions of the kernels and compatible affine fibres as
claimed.
\end{proof}

The displayed new statements are characteristic-zero \emph{lower} rank
bounds. The modular ranks at the two sampled fields are $92,20,99$, but
these modular values alone are not used to assert matching real upper
bounds or compatibility of $d(\zeta(t))$ for every small real $t$.
Neither the signs nor the order of vanishing of a minor at zero are needed.
In particular the proof makes no unverified inference from numerical
near-zero singular values.

\section{Projection onto one fixed base fibre avoids every moving inverse}
\label{v35:sec:projection}

The rank change is real, but it need not be a way around the tangency
obstruction. The fixed test in \cref{v35:thm:fixed-null} is the reason.
Here is the finite argument that replaces a moving nullspace or
pseudoinverse.

\begin{lemma}[A rank-independent estimate for a moving affine incidence]
\label[lemma]{v35:lem:projection}
Let $L(p),d(p)$ be continuously differentiable finite matrices and vectors
near $p_*$, and suppose $d_*\in\Ran L_*$. Let $T(p,x)$ be continuously
differentiable in $p$ and polynomial of degree at most two in $x$. Assume
\begin{equation}
              T(p_*,z)=\Delta\ne0\quad\hbox{whenever }L_*z=d_*.
 \label{v35:eq:base-constant}
\end{equation}
Then, on a smaller fixed neighborhood, there is $C<\infty$ such that for
all $x$, with $r=L(p)x-d(p)$ and $\delta=|p-p_*|$,
\begin{equation}
 \boxed{|T(p,x)-\Delta|
 \le C\{\delta(1+\|x\|)^2+
                 \|r\|(1+\|x\|+\|r\|)\}.}
 \label{v35:eq:robust}
\end{equation}
No rank or range hypothesis is imposed on $L(p)$ away from $p_*$.
\end{lemma}
\begin{proof}
Use only the fixed Moore--Penrose inverse $L_*^\dagger$. Define
\[
 e=L_*^\dagger(L_*x-d_*),\qquad z=x-e.
\]
The vector $L_*x-d_*$ belongs to $\Ran L_*$, hence $L_*z=d_*$ exactly.
Also
\[
 L_*x-d_*=(L_*-L(p))x+d(p)-d_*+r,
\]
so $\|e\|\le C_0\{\delta(1+\|x\|)+\|r\|\}$.
The coefficient bounds for a quadratic polynomial give
\[
 |T(p,x)-T(p_*,x)|\le C_1\delta(1+\|x\|)^2,
\]
\[
 |T(p_*,x)-T(p_*,z)|
       \le C_2\|e\|(1+\|x\|+\|e\|).
\]
Substitute the bound for $e$, restrict $\delta\le1$, and absorb mixed
terms using $2ab\le a^2+b^2$ or their direct linear bounds. The result
is \eqref{v35:eq:robust}, since $T(p_*,z)=\Delta$. All constants use
only this fixed base inverse and coefficient bounds on a fixed compact
parameter neighborhood. No inverse at $p$ appears.
\end{proof}

This proof is also an explicit polynomial consequence of the base affine
equations: $T(p_*,x)-\Delta$ lies in the ideal generated by
$L_*x-d_*$. The multiplier of that vector can be chosen affine in $x$.
No implicit parametrization of a varying-rank solution set is required.

\section{A source-specific escape theorem across all nearby rank strata}
\label{v35:sec:escape}

Let $\mathcal T_{\nu_\diamond}(Y,\ell)$ be the actual Bridge V22 tangency
coefficient
\begin{equation}
 \mathcal T_{\nu_\diamond}(Y,\ell)
 =Dd_{2,W}(Y)[F_1Y+G\ell,\nu_\diamond]
  +\mathbb K_1(F_1Y+G\ell)[\nu_\diamond,\ell].
 \label{v35:eq:T}
\end{equation}
Set $T(\zeta,x)=\mathcal T_{\nu_\diamond}(Y(\zeta),Q_3x)$.
It is analytic in $\zeta$ and at most quadratic in $x$. Bridge V23
proves \eqref{v35:eq:base-constant} for the full base fibre of
\eqref{v35:eq:joint}. As a new incidence regression, solving that base
system modulo 241 with a different free-coordinate choice and evaluating
\eqref{v35:eq:T} gives $T/\sqrt3=148$ modulo 241, in agreement with
its inherited nonzero coefficient. The global constancy on the base
fibre remains the inherited characteristic-zero theorem, not a conclusion
from this one modular check.

\begin{theorem}[Rank change cannot produce bounded tangency roots near the seed]
\label[theorem]{v35:thm:escape}
There are a neighborhood of $\zeta_*$ and constants $c,C>0$ such that
any solution of the \emph{complete} V23 upstream system in that
neighborhood, at $\zeta\ne\zeta_*$, satisfies
\begin{equation}
              \boxed{\|\ell\|\ge c|\zeta-\zeta_*|^{-1/2}-C.}
 \label{v35:eq:escape}
\end{equation}
This holds without fixing the Poisson rank, weak-null dimension, joint
incidence rank, or a fibre basis. In particular no sequence of complete
roots can have $\zeta_n\to\zeta_*$ and bounded $\ell_n$.
\end{theorem}
\begin{proof}
Every complete root satisfies $\mathcal L(\zeta)x=d(\zeta)$.
By \cref{v35:thm:fixed-null}, $\nu_\diamond$ belongs to its actual weak
Poisson-null space even where that space changes dimension. Therefore
every such root must also satisfy $T(\zeta,x)=0$. Apply
\cref{v35:lem:projection} with $r=0$ and the inherited nonzero
$\Delta_\diamond$. It gives
\[
        |\Delta_\diamond|\le C_0|\zeta-\zeta_*|(1+\|x\|)^2.
\]
Take square roots and use the fixed norm equivalence under $Q_3$.
There is no solution at $\zeta=\zeta_*$ by the base-fibre theorem.
\end{proof}

\begin{corollary}[Uniform exclusion at every prescribed finite control bound]
\label[corollary]{v35:cor:bounded}
For each fixed $M<\infty$ there is $\rho_M>0$ such that the complete
upstream system has no root with
$|\zeta-\zeta_*|<\rho_M$ and $\|\ell\|\le M$, including on
rank-changing loci. One may choose $\rho_M=c(1+M)^{-2}$ with a
sufficiently small fixed numerator determined by the constants above.
\end{corollary}
\begin{proof}
Use the last inequality in the proof with the fixed control bound and
choose $\rho_M$ strictly smaller than its necessary threshold. The
base point itself is excluded separately.
\end{proof}

The statement that a nearby bounded root might avoid the Version-33
estimate by changing rank was only an untested alternative. The new
source identity closes that alternative on this particular first-jet
chart. The rank-opening theorem does not contradict the escape theorem:
the fibre changes dimension, while a diagnostic in its weak test space
survives. The proof follows that fixed diagnostic rather than the entire
changing test bundle.

\section{Residual tolerances and the physical amplitude scale}
\label{v35:sec:tolerances}

The same estimate supplies a practical rejection condition that does not
mistake a nearly compatible floating solve for an exact root. Let
\[
 r=\mathcal L(\zeta)x-d(\zeta),\qquad
 \tau=|T(\zeta,x)|,\qquad R=\|x\|.
\]
Then
\begin{equation}
 \boxed{|\Delta_\diamond|
 \le \tau+C\{\delta(1+R)^2+\|r\|(1+R+\|r\|)\},
 \qquad \delta=|\zeta-\zeta_*|.}
 \label{v35:eq:residual-rejection}
\end{equation}
In particular, lower-incidence errors must be charged together with the
control norm. A small raw residual alone does not resolve an escaping
control problem. Equation \eqref{v35:eq:residual-rejection} is an
analytic bound with explicitly identifiable finite coefficient constants;
no numerical upper enclosure for $C$ is asserted in this iteration.

\begin{corollary}[Joint amplitude separation without a rank restriction]
\label[corollary]{v35:cor:amplitude}
Suppose a future family of complete roots has $\zeta(a)\to\zeta_*$ and
an actual initial multiplier expansion
\[
 \|\lambda(a,0)-e-a\ell(a)\|=o(a\|\ell(a)\|),\qquad
                         \lambda(a,0)\to e.
\]
Then, writing $\delta_a=|\zeta(a)-\zeta_*|$,
\begin{equation}
                  \boxed{a/\sqrt{\delta_a}\longrightarrow0,
                         \qquad \delta_a/a^2\longrightarrow\infty.}
 \label{v35:eq:amplitude}
\end{equation}
The conclusion allows arbitrary nearby rank changes.
\end{corollary}
\begin{proof}
The assumed remainder and near-flatness give $a\|\ell(a)\|\to0$ by
the reverse triangle inequality. Apply \eqref{v35:eq:escape} and use
$a\to0$. The two limits in \eqref{v35:eq:amplitude} are equivalent.
\end{proof}

This is not a uniform initializer for diverging $\ell$. The stationary
chart, signed Gram, higher-preparation controls and physical lifespan
still require estimates in that joint regime. No physical time is
rescaled and no reinitialization occurs during a trajectory.

\section{What the coefficient calculation settles}
\label{v35:sec:status}

\paragraph{Proved in this body.}
The original diamond test is weak Poisson-null on the entire six-dimensional
linear source span, hence on the mean-compatible five-parameter chart.
Explicit original-source minors show that the Poisson, weak-restriction
and full lower-incidence ranks increase along a punctured real arc through
the base point. The base-fibre tangency obstruction nevertheless enforces
square-root control escape across all nearby rank strata, with a robust
bound for approximate lower solves and the stated amplitude consequence.

\paragraph{Inherited, not independently proved here.}
The unchanged stationary action, the nonlinear initial-constraint
construction, the base rank upper bounds and full base-fibre constancy of
$\Delta_\diamond$ are the named Bridge inputs. The regeneration reproduces
its base Poisson, cubic-mean and tangency residues; this is not an independent
validation of every earlier analytical theorem or nonlinear history.

\paragraph{Unresolved and next direction.}
No full tangency root has been constructed or excluded at finite first-jet
displacement. An unbounded-control sequence approaching the seed also
remains possible at the necessary scale; its asymptotic equations must
respect the changing lower incidence. The full ambient 330-entry
Version-34 package is therefore not an automatically defined physical
derivative on this family. It remains a conditional construction on a
proved regular solvable subchart, using only that subchart's admissible
tangent directions. It is not appropriate to count every direction of
$\mathbb R^5$ as an available regular direction before testing rank.

The productive next finite target is a compatible root at finite
first-jet displacement, or the complete rank-aware unbounded-control
system; another search for bounded roots converging to this seed is
excluded by \cref{v35:thm:escape}. The Version-32 unit-initial-multiplier
stress equation remains an independent option. No cofinal Einstein
limit, common positive physical lifespan, chiral completion, nonzero-field
Yang--Mills quantum limit, or primitive Renewal-to-Einstein--SM theorem
is asserted.

\paragraph{Verification record.}
All 648 fixed-test coefficient identities were evaluated using exact
arbitrary-precision rational arithmetic. The original Poisson and complete
lower-incidence matrices, the weak restrictions and the stated minors were
regenerated independently at two good primes. Their arrays and index
lists are supplied. Nonzero modular residues certify nonzero
characteristic-zero minors; zero modular residues are never used to certify
an identity. The analytic projection and residual bounds have the written
proofs above, supplemented by finite regression examples. All bytes of
Version 34 before its final document command are preserved. The checks do
not constitute proof-assistant formalization or a numerical integration of
the original nonlinear spacetime.

\begingroup
\renewcommand{\refname}{Additional sources for Version 35}
\begin{thebibliography}{9}\small\raggedright
\bibitem{v35:EB23}
\emph{Renewal Exact-Action Einstein Bridge Research Notes}, supplied
cumulative Versions 1--23. Original coefficient interfaces:
\nolinkurl{v20:K1}, \nolinkurl{v20:P2}, \nolinkurl{v20:Phi3},
\nolinkurl{v20:B2}, \nolinkurl{v21:R}; base source facts:
\nolinkurl{v20:rank-certificate}, \nolinkurl{v21:joint-tilt},
\nolinkurl{v23:global-no-go}, \nolinkurl{v23:five-chart}.
The statements imported from this source are distinguished from the new
fixed-test and rank-opening calculations above.
\bibitem{v35:DBGP}
C. Di Bartolo, R. Gambini and J. Pullin,
\emph{Consistent and mimetic discretizations in general relativity},
J. Math. Phys. \textbf{46} (2005), 032501;
\url{https://arxiv.org/abs/gr-qc/0404052}.
Primary-source background for discrete multiplier selection and constraint
preservation; no coefficient value or estimate for the present source is
imported from this article.
\end{thebibliography}
\endgroup
\addtocontents{toc}{\protect\endgroup}
% END IMMUTABLE VERSION 35 BODY
% APPEND VERSION 36 HERE, BEFORE THE FINAL END-DOCUMENT COMMAND.



% BEGIN IMMUTABLE VERSION 36 BODY
\clearpage
\begin{center}
{\Large\bfseries Version 36: Solvable lower incidence and six intrinsic tangency functions}\\[0.5em]
{\large A fixed multiplier kernel, exact source compatibility, and elimination of the generic tilt fibre}
\end{center}
\makeatletter
\addtocontents{toc}{\protect\begingroup}
\addtocontents{toc}{\protect\makeatletter}
\addtocontents{toc}{\protect\def\protect\l@section{\protect\bfseries\protect\@dottedtocline{1}{0em}{2.5em}}}
\makeatother
\addcontentsline{toc}{part}{Version 36: Solvable lower incidence and six intrinsic tangency functions}

\section{The remaining computation after rank-independent escape}
\label{v36:sec:context}

Version 35 excludes bounded leading tilts approaching the old first jet even
through changing-rank loci. The next useful calculation is therefore not
another local escape inequality. We ask whether the complete lower incidence
is actually solvable away from that seed, and which of its free controls can
change the necessary tangency equations. Both questions can be answered for
an explicit nonempty open subset of the original five-parameter family.

There is a fixed sixteen-dimensional multiplier space annihilated by the
first Poisson matrix throughout the entire underlying six-dimensional linear
source span. The actual prepared quadratic source annihilates the same
space. A retained nonzero minor then determines the \emph{exact} generic
Poisson rank, rather than just a lower rank bound. Three explicit,
mean-zero lapse directions in that fixed kernel control the three cubic
constant-shift equations. Consequently the full lower system is solvable,
with nine free parameters, on a specified open family; compatibility is a
conclusion, not an augmented-rank hypothesis.

The second calculation is more decisive for the search. Six fixed axial
weak tests exhaust the nonconstant weak Poisson kernel on that family.
Their complete tangency polynomials are invariant under \emph{every}
translation by the fixed sixteen-dimensional multiplier kernel, including
translations which do not satisfy the cubic mean equations. Thus none of
the nine remaining admissible tilt controls changes those six tests. The
full finite first-jet incidence problem reduces to six rational functions
of the first-jet parameters, or four ratios after removing their common
homogeneous scale. No simultaneous zero of these functions is proved here.

\paragraph{Scope and inputs.}
All assertions in this body concern the original three-site cutoff
$h=1/3$. The stationary action and the coefficient formulas $P_{\nu,2}$,
$\Phi_3$, $\widehat B_2$, and the cubic constant-row map are the unchanged
Bridge V20--V23 formulas retained in \cite{v36:EB23} and reconstructed in
Version 35's executable coefficient code. The computations below use those
formulas, not a continuum stencil or a positive replacement of the signed
system. The inherited initializer and higher signed response remain
separate analytical inputs. The distinction between selecting discrete
multipliers and preserving a supplementary condition under their actual
flow is also present in the general constrained-discretization literature
\cite{v36:DBGP}; no source-specific identity below is imported from it.


\paragraph{Companion V24 audit.}
The newer Bridge V24 source \cite{v36:EB24} constructs one formal
$\alpha_2$ tangent of the complete lower incidence and proves a nonzero
first derivative of the diamond tangency test. It explicitly leaves
integration of that tangent and a full tangency root open. Those results
are not recounted here as new. The present open-set solution need not
extend with bounded controls to the old rank-90 seed. V24 also repairs an
omitted base summand in an auxiliary V23 recomputation; the missing value
is proved zero there. Our coefficient formula below retains the full
$Dd_{2,W}(Y)[F_1Y+G\ell]$ throughout, including the $G\ell$ term.
This targeted audit does not independently regenerate V24's derivative
certificate.

\section{A fixed kernel and an exact compatibility identity}
\label{v36:sec:kernel}

Use the six-dimensional family $Y(\alpha,b)$ in
\eqref{v35:eq:six-family}, with $p=\sqrt3\,p_R$. A multiplier coordinate
$x\in\R^{108}$ represents the physical multiplier $\ell=Q_3x$, where
$Q_3=\operatorname{diag}(I_{27},\sqrt3I_{81})$. Write
\begin{equation}
 \widehat K(Y)=\frac1{\sqrt3}Q_3\mathbb K_1(Y)Q_3,
 \qquad \widehat b(Y)=\frac1{\sqrt3}Q_3\widehat B_2(Y).
 \label{v36:eq:rational}
\end{equation}
These are the original rationalized quantities, with the spatial-sum
convention used consistently in every row. Define the mean map
$\mathsf m:\R^{108}\to\R^4$ using spatial \emph{means}.

Let $\mathcal N$ be the following fixed subspace of these multiplier
coordinates:
\begin{equation}
 \boxed{\quad
 \mathcal N=\{(n,\beta):
 n(x)=a_1(x_1)+a_2(x_2)+a_3(x_3),\quad
 \beta_i(x)=b_i(x_i)\ (i=1,2,3)\}.
 \quad}
 \label{v36:eq:N}
\end{equation}
The $a_i,b_i$ are arbitrary real functions on the three-point circle;
$b_i$ in this display are multiplier profiles, not the canonical
parameters $b$. Additive lapse profiles have dimension seven and the
axial shifts have dimension nine, hence $\dim\mathcal N=16$.

For an explicit integer synthesis $N:\R^{16}\to\R^{108}$ put
$c=(2,-1,-1)$ and $s=(0,1,-1)$ on the three-point circle. Its columns,
numbered from zero, are the constant lapse, the six lapse profiles
$c(x_i),s(x_i)$ in increasing $i$, and then the nine shifts
$e_i,e_i c(x_i),e_i s(x_i)$ in increasing $i$. Orthogonality of the
constant, cosine and sine profiles proves injectivity. Set
\begin{equation}
 \mathcal N_0=\mathcal N\cap\ker\mathsf m,
 \qquad \dim\mathcal N_0=12.
 \label{v36:eq:N0}
\end{equation}

\begin{theorem}[Universal Poisson kernel and prepared-source compatibility]
\label[theorem]{v36:thm:kernel}
For every real $(\alpha,b)\in\R^6$, without imposing the quadratic
scalar-mean surface,
\begin{equation}
 \boxed{\widehat K(Y)N=0,\qquad N^T\widehat b(Y)=0.}
 \label{v36:eq:identities}
\end{equation}
In particular $\operatorname{rank}\widehat K(Y)\le92$ on this whole
linear source span. These are identities of the original coefficients,
not statements about their principal symbols.
\end{theorem}
\begin{proof}
We give a finite exact coefficient proof and its exhaustion argument.
Let $Y_a^R=(U_a,p_{R,a})$, $1\le a\le6$, be the coefficient vectors
obtained by setting in turn one of $(\alpha_1,\alpha_2,\alpha_3,b_1,b_2,b_3)$
to one and the others to zero in \eqref{v35:eq:six-family}. Use
$\delta_i=3(S_i-S_i^{-1})$ on the original arrays. The defining identity
\eqref{v35:eq:K-G} and quadratic polarization evaluate a column against
$\nu=N_j$ by
\[
 DP_{\bullet,2}(Y_a^R)[G\nu]
       -DP_{\nu,2}(Y_a^R)[G\bullet].
\]
Here a bullet ranges over all 108 original multiplier point tests. Direct
rational contraction of the retained $P_2$ gives zero in all 108 rows,
for all sixteen $N_j$ and all six $Y_a^R$. The reversal grading in
\eqref{v36:eq:rational} multiplies blocks by nonzero constants and
preserves these zeros, since each column $N_j$ is either a pure lapse
or a pure shift. This supplies $6\cdot16\cdot108=10368$ exact scalar
identities. Linearity of $\mathbb K_1$ in $Y$ proves the first statement
for every real linear combination.

For the second statement the complete rationalized target is evaluated by
\begin{equation}
 \widehat b_N(U,p_R)=\widehat B_{2,N}(U,p_R),\qquad
 \widehat b_\beta(U,p_R)=\widehat B_{2,\beta}(U,p_R)
                         +2\widehat B_{2,\beta}(0,p_R).
 \label{v36:eq:b-grading}
\end{equation}
Its lapse component includes the original prepared correction
$P_{(0,\delta n),2}$; it is not an unprepared drift. The map is
homogeneous quadratic in the six real source coefficients. Evaluate
$N^T\widehat b$ at the six coordinate vectors and their fifteen pairwise
sums. All sixteen components are exactly zero at every node, giving
336 rational equalities. If $q$ is any scalar homogeneous quadratic
polynomial, its diagonal coefficients are $q(e_a)$ and its cross
coefficients are $q(e_a+e_b)-q(e_a)-q(e_b)$. The stated nodes therefore
exhaust every coefficient of every component and prove the identity.

The supplied verifier reconstructs these contractions from the original
quadratic constraint and cubic Hamiltonian formulas. Its arithmetic is
integer-array arithmetic with exact common rational denominators. No
floating zero or zero residue modulo a prime enters this proof. Finally
$N$ is injective, so the first identity implies the rank bound.
\end{proof}

\section{An explicit nonempty generic set with exact ranks}
\label{v36:sec:ranks}

Return to the five-parameter mean-compatible chart
$\zeta=(\alpha_1,\alpha_2,\alpha_3,b_2,b_3)$, with
\begin{equation}
 b_1(\zeta)=\frac{(\alpha_1^2+\alpha_2^2+\alpha_3^2)/4-b_2b_3}{b_2+b_3}.
 \label{v36:eq:mean}
\end{equation}
Let $\widehat R(\zeta)=\mathcal R_{Y(\zeta)}Q_3$ and let
$t(\zeta)$ be the actual cubic constant-shift target. Thus the full lower
incidence is
\begin{equation}
 \widehat Kx=-\widehat b,\qquad \mathsf mx=0,\qquad
                       \widehat Rx=-t.
 \label{v36:eq:lower}
\end{equation}
The right side $t$ is retained even where it has not been separately
simplified. The first and last coefficient maps are linear and quadratic,
respectively, in the canonical first jet; $t$ is cubic.

Define the index set
\begin{align*}
 \mathcal F&=(8,17,20,23,24,25,26,35,44,53,74,77,80,89,98,105,106,107),\\
 \mathcal I&=(\{0,\ldots,107\}\setminus\mathcal F)\cup\{89,98\}.
\end{align*}
It has 92 indices. Let $E_{\mathcal I}$ be coordinate insertion on this
set, and set $A=E_{\mathcal I}^T\widehat KE_{\mathcal I}$. Let $E_3$
contain columns 1, 2, and 3 of $N$: explicitly, the lapse profiles
$c(x_1),s(x_1),c(x_2)$, with zero shift. They belong to $\mathcal N_0$.
Write $H=\widehat R E_3$.

\begin{theorem}[A nonempty open family with ranks 92, 9 and 99]
\label[theorem]{v36:thm:generic}
The set
\begin{equation}
 \mathcal U=\{\zeta:(b_2+b_3)\alpha_1\alpha_2\alpha_3\ne0,
                         \ \det A(\zeta)\det H(\zeta)\ne0\}
 \label{v36:eq:U}
\end{equation}
is nonempty and open, and is dense in the displayed rational parameter
chart. For every $\zeta\in\mathcal U$,
\begin{equation}
 \boxed{\ker\widehat K=\mathcal N,\quad
 \operatorname{rank}\widehat K=92,\quad
 \dim(\ker\mathbb K_1\cap\mathcal W)=9,\quad
 \operatorname{rank}\begin{pmatrix}\widehat K\\\mathsf m\\\widehat R\end{pmatrix}=99.}
 \label{v36:eq:exact-ranks}
\end{equation}
Here $\mathcal W$ is the original curl-free physical shift space. Its
nine null directions are precisely the shifts $\beta_i=b_i(x_i)$;
six are nonconstant and three are constant.
\end{theorem}
\begin{proof}
Use the exact rational point
\begin{equation}
 \zeta^\dagger=(2,9/8,1,-8,-8),\qquad b_1^\dagger=15983/4096.
 \label{v36:eq:point}
\end{equation}
The original coefficient regeneration gives, at the good prime 241,
\begin{equation}
                     \det A(\zeta^\dagger)=40\pmod{241},\qquad
                     \det H(\zeta^\dagger)=226\pmod{241}.
 \label{v36:eq:residues}
\end{equation}
The first nonzero residue is the retained Version-35 rank witness,
recomputed here. The second uses the three explicit lapse profiles above.
All denominators at this rational point are units at 241. Nonzero residues
therefore prove nonzero characteristic-zero determinants. A second
prime provides an independent regeneration check recorded in the package.
No zero residue is used to establish an identity or an upper rank bound.

By \cref{v36:thm:kernel}, rank is at most 92, while invertibility of $A$
gives rank at least 92. Its kernel is therefore exactly $\mathcal N$.
Congruence by $Q_3$ gives the physical kernel; its intersection with pure
curl-free shifts is the nine-dimensional axial shift space in
\eqref{v36:eq:N}. The four constant multiplier columns show that
$\mathsf m|_{\mathcal N}$ is onto. The homogeneous kernel of the first
two rows in \eqref{v36:eq:lower} is $\mathcal N_0$, of dimension twelve.
Invertibility of $H$ makes $\widehat R|_{\mathcal N_0}$ onto $\R^3$,
so the full homogeneous kernel has dimension nine and the full row rank
is $108-9=99$.

Every coefficient is rational in $\zeta$, with denominators supported on
$b_2+b_3$ and fixed nonzero rational constants. Clearing those denominators
makes the determinant product a polynomial which is not identically zero,
as witnessed by \eqref{v36:eq:residues}. The zero set of a nonzero real
polynomial has empty interior. This proves openness, density on the
rational chart, and nonemptiness at the specified point.
\end{proof}

The old seed is not in $\mathcal U$: its Poisson rank is 90. Thus this
is not a return to an eleven-dimensional solved fibre through that seed.
There are real nearby rank-92 points, but the generic kernel is the fixed
sixteen-dimensional space, not its eighteen-dimensional enlargement at
the old seed.

\section{The full lower system is solved, not postulated compatible}
\label{v36:sec:solution}

Let $C_4:\R^4\to\R^{108}$ insert constant multiplier components, so
$\mathsf mC_4=I_4$. For $\zeta\in\mathcal U$ define
\begin{align}
 x_{\rm p}&=-E_{\mathcal I}A^{-1}E_{\mathcal I}^T\widehat b,
                                                        \label{v36:eq:xp}\\
 x_{\rm m}&=x_{\rm p}-C_4\mathsf m x_{\rm p},\notag\\
 x_0&=x_{\rm m}+E_3H^{-1}(-t-\widehat Rx_{\rm m}).       \label{v36:eq:x0}
\end{align}
Both inverses in these definitions are nonsingular by a proved open-set
condition with a nonempty exact witness.

\begin{theorem}[Generic exact lower solvability and all its remaining controls]
\label[theorem]{v36:thm:solve}
For every $\zeta\in\mathcal U$, $x_0(\zeta)$ solves all 108 quadratic
drift rows, the four mean normalizations, and the three cubic
constant-shift equations in \eqref{v36:eq:lower}. Every solution is
\begin{equation}
                \boxed{x=x_0+e,\qquad
                 e\in\mathcal N_0\cap\ker\widehat R,}
 \label{v36:eq:all-solutions}
\end{equation}
and the latter space has dimension nine. The selected $x_0$ is rational,
hence analytic on $\mathcal U$, and bounded with each fixed derivative on
compact subsets of $\mathcal U$.
\end{theorem}
\begin{proof}
The coordinate space $\operatorname{Ran}E_{\mathcal I}$ intersects
$\mathcal N$ only at zero: an intersection vector has coordinates $y$
with $Ay=0$. Dimension then gives
$\R^{108}=\operatorname{Ran}E_{\mathcal I}\oplus\mathcal N$.
The residual $r=\widehat Kx_{\rm p}+\widehat b$ annihilates the first
summand by \eqref{v36:eq:xp}. It annihilates $\mathcal N$ because
$\widehat K$ is skew and both identities in \eqref{v36:eq:identities}
hold. Thus $r=0$ on the whole space, not just on the retained 92 rows.

The mean correction lies in $\mathcal N$ and enforces $\mathsf mx_{\rm m}=0$.
The final correction lies in $\mathcal N_0$ and has image
$-t-\widehat Rx_{\rm m}$ under $\widehat R$, by $H=\widehat RE_3$.
This proves every original lower equation. Subtracting two solutions and
using \cref{v36:thm:generic} gives \eqref{v36:eq:all-solutions} and its
dimension. Rationality, analyticity and compact-subset bounds follow
from the two displayed finite inverse formulas.
\end{proof}

This theorem removes a lower-incidence compatibility assumption on an
actual full-dimensional open family. It does not remove tangency or
higher signed regularity assumptions. In particular, solving these
first-jet coefficient equations is not yet a proof of an invariant
preparation manifold or a common physical lifespan.

\section{All six generic tangency tests are independent of the multiplier kernel}
\label{v36:sec:invariance}

Let $\nu_1,\ldots,\nu_6$ be the fixed physical weak tests with zero lapse
and shifts $e_i c(x_i),e_i s(x_i)$, in increasing $i$. Their coordinate
columns are columns $8,9,11,12,14,15$ of $N$. Multiplication by the
common physical factor $\sqrt3$ would change only their normalization,
not their span or the following zeros. Define
\begin{equation}
 \Theta_a(Y,x)=\frac1{\sqrt3}\mathcal T_{\nu_a}(Y,Q_3x),
                   \qquad 1\le a\le6.
 \label{v36:eq:theta}
\end{equation}

\begin{theorem}[Exact translation invariance of the six tangency polynomials]
\label[theorem]{v36:thm:invariance}
For every first jet in the six-dimensional linear source span, every
$x\in\R^{108}$, and every $e\in\mathcal N$,
\begin{equation}
              \boxed{\Theta_a(Y,x+e)=\Theta_a(Y,x),
                         \qquad a=1,\ldots,6.}
 \label{v36:eq:invariance}
\end{equation}
Neither $x$ nor $x+e$ is required to solve a lower equation. In particular,
all nine admissible controls in \eqref{v36:eq:all-solutions} leave the
complete nonconstant weak-null tangency covector unchanged on $\mathcal U$.
\end{theorem}
\begin{proof}
We specify the rational polynomial and exhaust its coefficients, retaining
the physical grading. Write $Y=(U,\sqrt3p_R)$ and $x=(n,\beta_R)$.
Put $Y_R=(U,p_R)$, let $F_1^R Y_R=(2p_R-\operatorname{tr}(p_R)I,
\mathcal F(U))$, and use $G$ on the unscaled multiplier coordinates.
For a pure weak test $\nu$, let $d_\nu(Y_R)$ be the unscaled quadratic
weak drift obtained from the original $P_2,\Phi_3$. Its types are $U^2$
and $p_R^2$. The reversal grading gives the exact identity
\begin{equation}
 \Theta_\nu(Y,x)=
 Dd_\nu(Y_R)[F_1^R Y_R+Gx]
 +\mathbb K_1(F_1^R Y_R+Gx)[\nu,x],
 \label{v36:eq:rational-T}
\end{equation}
where the $\mathbb K_1$ on the right is evaluated in unscaled fields.
Indeed the physical first velocity is
$(\sqrt3(V(p_R)+\operatorname{SymGrad}\beta_R),
\mathcal F(U)+\nabla^2n-I\Delta n)$; both terms in the physical weak
tangency have a single factor $\sqrt3$ after this grading.

The polynomial in \eqref{v36:eq:rational-T} is the sum of a quadratic
function of $Y_R$, a bilinear function of $(Y_R,x)$, and
\[
 q_\nu(x)=\mathbb K_1(Gx)[\nu,x].
\]
For $e\in\mathcal N$, its change is therefore
\begin{equation}
 L_\nu(Y_R,e)+Dq_\nu(e)[x]+q_\nu(e),\qquad
 L_\nu(Y_R,e)=Dd_\nu(Y_R)[Ge]
                       +\mathbb K_1(F_1^RY_R)[\nu,e].
 \label{v36:eq:change}
\end{equation}
Here $L_\nu$ is bilinear and $q_\nu$ is homogeneous quadratic. All its
coefficients can consequently be tested on the same six source basis
vectors, the sixteen columns of $N$, and the 108 original multiplier
point columns.

For clarity the gradient of $q_\nu$ can be evaluated without differentiating
an inverse or a moving nullspace. If
$H_\nu(V,W)=D^2P_{\nu,2}[V,W]$, then
\begin{equation}
 Dq_\nu(e)[z]=2H_\nu(Ge,Gz)
                -H_z(Ge,G\nu)-H_e(Gz,G\nu).
 \label{v36:eq:q-gradient}
\end{equation}
Substitute the unchanged $P_2$ into this expression. Exact rational
contraction gives zero for all six choices of $\nu$, all sixteen $e=N_j$,
and every point test $z$, or 10368 exact scalar zeros. The retained cubic
$\Phi_3$ and quadratic $P_2$ likewise give
$L_{\nu_a}(Y_b^R,N_j)=0$ for all $6\cdot6\cdot16=576$ choices.
These finite identities are regenerated in the supplied script. Every
spatial shift and the prepared weak-drift convention is retained.

Bilinearity now gives $L_\nu=0$ for all the stated inputs and
$Dq_\nu(e)=0$ for every $e\in\mathcal N$. Euler's identity for a
homogeneous quadratic polynomial gives $2q_\nu(e)=Dq_\nu(e)[e]=0$.
Equation \eqref{v36:eq:change} proves \eqref{v36:eq:invariance}. The
exhaustive arithmetic is over $\mathbb Q$ (implemented as exact integer arrays
and denominators), not a modular or numerical inference of zero.
\end{proof}

The three constant weak-shift tests give no further tangency equation.
The inherited translation identity gives
$\mathbb K_1(V)[(0,c),\ell]=\langle\delta_c\ell,LV\rangle_h$,
while their quadratic weak drift is identically zero. For the present
family $LY=0$ and $L(F_1Y+G\ell)=0$. Thus their tangency is zero for
all $\ell$. This is a constraint identity, not deletion of constant
physical shifts from the action.

\section{Six rational functions exactly replace the generic upstream system}
\label{v36:sec:reduction}

On $\mathcal U$ define the six actual source functions
\begin{equation}
 \boxed{\tau_a(\zeta)=\Theta_a\bigl(Y(\zeta),x_{\rm p}(\zeta)\bigr),
                         \quad a=1,\ldots,6.}
 \label{v36:eq:tau}
\end{equation}
The particular vector $x_{\rm p}$ only solves the original Poisson rows.
It is legitimate in this definition because $x_0-x_{\rm p}\in\mathcal N$
and \cref{v36:thm:invariance} is stronger than invariance on the normalized
lower fibre. The full lower correction \eqref{v36:eq:x0} remains necessary
when constructing an actual admissible initial multiplier.

\begin{theorem}[Exact generic reduction of every lower and weak-null tangency row]
\label[theorem]{v36:thm:reduction}
For $\zeta\in\mathcal U$, the complete Bridge V23 first-jet system
\[
 \widehat B_2(Y)+\mathbb K_1(Y)\ell=0,\quad
 \mathsf m\ell=0,\quad t(Y)+\mathcal R_Y\ell=0,\quad
 \mathcal T_\nu(Y,\ell)=0\quad(\nu\in\mathcal W\cap\ker\mathbb K_1(Y))
\]
has a multiplier solution if and only if
\begin{equation}
                        \boxed{\tau_1(\zeta)=\cdots=\tau_6(\zeta)=0.}
 \label{v36:eq:six-zero}
\end{equation}
If these six functions vanish, every $\ell=Q_3x$ from
\eqref{v36:eq:all-solutions} solves that finite coefficient system.
Each $\tau_a$ is a rational function on $\mathcal U$, independent of
the choice of complement used to solve the Poisson equation.
\end{theorem}
\begin{proof}
The lower equations are solved by \cref{v36:thm:solve}. The weak
Poisson-null space on $\mathcal U$ is exactly the six displayed
nonconstant tests plus the three constant shifts, by
\cref{v36:thm:generic}. The latter tests vanish by the translation
identity. Every lower solution differs from $x_{\rm p}$ by an element
of $\mathcal N$; hence all other tangency components are precisely the
six values in \eqref{v36:eq:tau}. This proves both implications.

The construction uses rational coefficient functions and the inverse of
one rational 92 by 92 matrix with nonzero determinant, so each $\tau_a$
is rational. A different Poisson solution differs by $\ker\widehat K
=\mathcal N$, and cannot change the six values. This proves the
independence statement without a preferred gauge or a moving
pseudoinverse.
\end{proof}

The theorem is an equivalence for the full \emph{finite first-jet}
incidence, not the nonlinear Einstein equation. In particular it does not
replace the absolute higher weak targets, their signed rank, or their
actual flow by six scalar equations. The generic open set deliberately
excludes exceptional rank loci, where additional weak-null tests must be
included anew.

\section{Four scale ratios and differentiation of the actual reduced map}
\label{v36:sec:ratios}

\begin{proposition}[Homogeneity and four shape parameters]
\label[proposition]{v36:prop:scale}
Under the simultaneous scaling $(\alpha,b)\mapsto(s\alpha,sb)$,
$s\ne0$, the mean surface is preserved and
\begin{equation}
 \widehat K\mapsto s\widehat K,\quad
 \widehat b\mapsto s^2\widehat b,\quad
 \widehat R\mapsto s^2\widehat R,\quad t\mapsto s^3t,
 \quad x_{\rm p},x_0\mapsto s(x_{\rm p},x_0),\quad
 \boxed{\tau\mapsto s^2\tau.}
 \label{v36:eq:homogeneity}
\end{equation}
On the part with $\alpha_3\ne0$, the zero problem thus depends only on
\begin{equation}
 \left(\frac{\alpha_1}{\alpha_3},\frac{\alpha_2}{\alpha_3},
                         \frac{b_2}{\alpha_3},\frac{b_3}{\alpha_3}\right).
 \label{v36:eq:ratios}
\end{equation}
Fixing $\alpha_3=1$ in a search prevents a spurious improvement obtained
solely by shrinking the whole first jet to zero.
\end{proposition}
\begin{proof}
The first four scalings follow from the homogeneous degrees of the
original coefficient functionals; the mean condition is quadratic.
Equations \eqref{v36:eq:xp}--\eqref{v36:eq:x0} give the next scaling
since the fixed insertion matrices do not change. Both terms in
\eqref{v36:eq:rational-T} are homogeneous of total degree two in
$(Y_R,x)$, giving the last scaling. The determinant exclusions stay
nonzero for $s\ne0$. Division by $\alpha_3$ proves
\eqref{v36:eq:ratios}. This is a parameter reduction of the coefficient
problem, not a new normalization of physical time or deletion of a
physical amplitude from a reconstructed history.
\end{proof}

The new map can also be differentiated without transporting an old null
basis. Put $u=A^{-1}(-\widehat b_{\mathcal I})$. For any admissible
parameter variation,
\begin{equation}
 A\,\partial_j u=-\partial_j\widehat b_{\mathcal I}
                              -(\partial_j A)u,\qquad
 \partial_j\tau_a=
  (\partial_j\Theta_a)(Y,x_{\rm p})+
      D_x\Theta_a(Y,x_{\rm p})E_{\mathcal I}\partial_j u .
 \label{v36:eq:derivative}
\end{equation}
Here the first derivative includes the actual derivative of $Y(\zeta)$,
including $b_1(\zeta)$. These are differentiated linear equations for
the actual source solve, not derivatives of a fitted inverse. They give
a concrete way to assemble a certified Jacobian or an interval residual
bound on compact subsets of $\mathcal U$.

Six equations in four ratios are not a dimension-counting proof of
nonexistence. They may have identities, exceptional zeros, or no zeros.
Conversely a small least-squares value does not prove a zero. The
reduction specifies the exact remaining map without prejudging that
question.

\section{Executed finite tests and the limits of the exploratory search}
\label{v36:sec:verification}

The four exact identity banks have sizes
\begin{center}
\begin{tabular}{@{}lr@{}}
\toprule
Identity bank & Rational scalar equalities\\\midrule
Six source directions times sixteen Poisson-null columns & 10368\\
Prepared-source compatibility at 21 quadratic nodes & 336\\
Tangency quadratic gradients on sixteen kernel directions & 10368\\
Tangency source--kernel bilinear coefficients & 576\\\midrule
Total & 21648\\\bottomrule
\end{tabular}
\end{center}
Every equality was regenerated and checked in characteristic zero. The
shared-denominator array implementation is checked separately against
ordinary scalar \texttt{Fraction} arithmetic. The nonzero minors and an
original-source point solve are regenerated at the good primes 241 and
1009. These positive minor tests are distinct from the zero-identity
proofs. The supplied package contains the actual code, profile banks,
index sets, residue reports and input fingerprints.

A separate floating search used the four ratios with $\alpha_3=1$.
It did not produce a verified simultaneous root. Some reductions of the
six residuals were accompanied by increasingly large homogeneous
momenta; the original first jet was therefore moving far from the old
seed, not approaching a bounded control solution there. These searches
are archived as diagnostics only. They provide neither an all-parameter
nonexistence proof nor a certified finite-distance solution. No assertion
of an asymptotic residual power is made from their finite values.

The point \eqref{v36:eq:point} certifies a nonempty generic lower-solvable
family, not a tangency root. A nonzero rational tangency component at
that point is detected by a nonzero good-prime residue in the package.
Thus the new solvability theorem is not being presented as the missing
full tangency cancellation.

\section{Status and the next nonduplicating target}
\label{v36:sec:status}

\paragraph{Proved.}
The original Poisson matrix has a universal sixteen-dimensional kernel
on the six-dimensional source span, and the actual quadratic prepared
source annihilates it. On an explicitly nonempty open dense part of the
five-parameter mean surface, the Poisson rank is exactly 92, the weak
nullspace has dimension nine, and the complete lower rank is exactly 99.
Every lower right-hand side supplied by the source is solved there by
the displayed formulas. Six fixed nonconstant tangency functions are
invariant under the entire universal multiplier kernel. They exactly
replace the generic full first-jet incidence problem, and homogeneity
reduces their zero search to four shape ratios.

\paragraph{Inherited.}
The stationary action, flat canonical identities, physical grading,
original initialization mechanism, old-seed tangency obstruction and
Version-35 escape theorem retain their stated scopes. The coefficient
arithmetic above checks new identities in that action; it is not an
independent verification of all earlier analytical continuum claims.

\paragraph{Unresolved.}
No simultaneous root of the six reduced functions is proved, and
exceptional rank loci are not classified by the generic result. No
higher signed-rank certificate, invariant nonlinear preparation, common
positive physical lifespan or cofinal Einstein--SM limit follows. The
unit-initial-multiplier stress problem of Version 32 remains independent.

The next finite target is now the actual map
$\tau(\alpha_1,\alpha_2,1,b_2,b_3)$, with its determinant exclusions,
not an optimization over nine irrelevant generic tilt controls. A
certified common zero or a source-specific separating identity would
resolve something new. Another adjustment of those free tilt controls
cannot alter the six remaining tangency observations.

\begingroup
\renewcommand{\refname}{Additional sources for Version 36}
\begin{thebibliography}{9}\small\raggedright
\bibitem{v36:EB23}
\emph{Renewal Exact-Action Einstein Bridge Research Notes}, supplied
cumulative Versions 1--23. Original coefficient interfaces
\nolinkurl{v20:K1}, \nolinkurl{v20:P2}, \nolinkurl{v20:Phi3},
\nolinkurl{v20:B2}, \nolinkurl{v21:R}, and
\nolinkurl{v23:upstream-system}. The present body retains their
literal three-site conventions and distinguishes new coefficient
identities from inherited initialization and flow results.
\bibitem{v36:EB24}
\emph{Renewal Exact-Action Einstein Bridge Research Notes}, supplied cumulative
Versions 1--24. The targeted audit uses \nolinkurl{v24:missing-zero},
\nolinkurl{v24:lower-tangent}, \nolinkurl{v24:nonzero-upstream}, and
\nolinkurl{v24:status}; their certificates are inherited, not newly computed here.
\bibitem{v36:DBGP}
C. Di Bartolo, R. Gambini and J. Pullin,
\emph{Consistent and mimetic discretizations in general relativity},
J. Math. Phys. \textbf{46} (2005), 032501;
\url{https://arxiv.org/abs/gr-qc/0404052}.
\end{thebibliography}
\endgroup
\addtocontents{toc}{\protect\endgroup}
% END IMMUTABLE VERSION 36 BODY
% APPEND VERSION 37 HERE, BEFORE THE FINAL END-DOCUMENT COMMAND.



% BEGIN IMMUTABLE VERSION 37 BODY
\clearpage
\begin{center}
{\Large\bfseries Version 37: A universal lapse--shift tangency form}\\[0.5em]
{\large An exact two-profile boundary root and a certified anisotropic residual tail}
\end{center}
\makeatletter
\addtocontents{toc}{\protect\begingroup}
\addtocontents{toc}{\protect\makeatletter}
\addtocontents{toc}{\protect\def\protect\l@section{\protect\bfseries\protect\@dottedtocline{1}{0em}{2.5em}}}
\makeatother
\addcontentsline{toc}{part}{Version 37: A lapse--shift form, a boundary root, and an anisotropic tail}

\section{From the six rational functions to their actual coefficients}
\label{v37:sec:context}

The lower-solvability and kernel-translation theorems of Version 36 are retained,
not reproved. They leave six rational tangency functions on the three-profile
five-parameter family, or four ratios after homogeneous normalization. The
present calculation evaluates the coefficients behind those functions. There
is a stronger cancellation: their explicit canonical quadratic and
canonical--multiplier bilinear terms vanish. The remaining expression is a
fixed bilinear pairing of lapse and shift, evaluated on the multiplier
selected by the original Poisson equation.

Two consequences are established with original-coefficient arithmetic. An
explicit rational point on a two-profile boundary solves every lower and
weak-null tangency equation. It does not satisfy the nonzero-three-profile
entrance, and one of the mean-preparation derivatives vanishes there. Separately,
a ray with all three oscillatory profiles fixed and nonzero has tangency
residual of exact order $R^{-4}$ as its homogeneous momentum grows like $R$.
The leading residual is nonzero, so its small values are not finite roots.
Both observations are relevant to the remaining search; neither is an
Einstein evolution theorem.

\paragraph{Scope and dependencies.}
Throughout this body $h=1/3$. We use precisely the six-dimensional rationalized
source span, physical map $Q_3$, first Poisson matrix $\widehat K$, prepared
quadratic source $\widehat b$, and constant-shift rows of
\cref{v36:sec:kernel,v36:sec:invariance,v36:sec:reduction}. The original
quadratic constraint and cubic action coefficient are the ones reconstructed
by the retained executable compiler. All new zero identities below are
re-evaluated in characteristic zero; the earlier universal kernel and
source-compatibility theorem is an inherited input. No new companion
continuation beyond the V24 audit recorded in Version 36 is needed for these
calculations. Solving a finite constraint incidence and preserving it under
its own dynamics remain different tasks, as also emphasized in the general
constrained-discretization setting \cite{v37:DBGP}.

\section{The canonical terms cancel before the Poisson equation is solved}
\label{v37:sec:canonical}

Write $Y_R=(U,p_R)$ with physical momentum $\sqrt3p_R$, and let
$x=(n,\beta_R)\in\R^{108}$ represent $\ell=Q_3x$. Use the six fixed axial
weak tests $\nu_a$ of Version 36. Denote their quadratic weak drifts by
$d_a(Y_R)$ and put
\[
 F_1^R Y_R=(2p_R-\operatorname{tr}(p_R)I,\mathcal F(U)).
\]
Every occurrence of $G$ below is the original flat multiplier Hamiltonian
vector, with the same unscaled rational coordinates as in
\eqref{v36:eq:rational-T}.

\begin{theorem}[Canonical independence of the six tangency polynomials]
\label[theorem]{v37:thm:pure}
For every $Y_R$ in the entire six-dimensional source span and every
$x\in\R^{108}$,
\begin{align}
 Dd_a(Y_R)[F_1^R Y_R]&=0,\\
 Dd_a(Y_R)[Gx]+\mathbb K_1(F_1^R Y_R)[\nu_a,x]&=0.
 \label{v37:eq:cancellations}
\end{align}
Consequently the full tangency polynomial satisfies
\begin{equation}
 \boxed{\Theta_a(Y,x)=q_a(x):=\mathbb K_1(Gx)[\nu_a,x],
                         \qquad 1\le a\le6.}
 \label{v37:eq:pure}
\end{equation}
No mean-surface equation, lower-incidence equation, or generic-rank
hypothesis is required for this identity.
\end{theorem}
\begin{proof}
The retained complete expression is
\[
 \Theta_a(Y,x)=Dd_a(Y_R)[F_1^R Y_R+Gx]
                  +\mathbb K_1(F_1^R Y_R+Gx)[\nu_a,x].
\]
Here $d_a$ is homogeneous quadratic and $\mathbb K_1$ is linear in its
canonical input. Expanding therefore separates a canonical quadratic term,
a canonical--multiplier bilinear term, and $q_a$.

For an exhaustive exact coefficient evaluation let $Y_j^R$, $1\le j\le6$,
be the six literal source columns generated by unit values of
$(\alpha_1,\alpha_2,\alpha_3,b_1,b_2,b_3)$, and let $e_k$, $1\le k\le108$,
be the original multiplier point columns. Quadratic polarization of the
unchanged $P_{\nu,2}$ and cubic polarization of $\Phi_3$ give
\begin{align*}
 Dd_a(Y_j^R)[F_1^R Y_l^R]&=0
                          &&(1\le j,l\le6),\\
 Dd_a(Y_j^R)[Ge_k]
       +\mathbb K_1(F_1^R Y_j^R)[\nu_a,e_k]&=0
                          &&(1\le j\le6,\ 1\le k\le108).
\end{align*}
The companion script evaluates all $6\cdot36=216$ first entries and all
$6\cdot6\cdot108=3888$ second entries as exact rational arrays; all vanish.
It uses literal $3(S_i-S_i^{-1})$, both symmetric off-diagonal occurrences,
and the shifted cubic action term. Neither a finite-amplitude derivative nor
an inverse is used. Bilinearity exhausts the two displayed finite coefficient
banks and proves \eqref{v37:eq:cancellations}; substitution proves
\eqref{v37:eq:pure}. The cancellation does not discard the $Gx$ summand:
its derivative is explicitly included in the second bank.
\end{proof}

This result removes only the \emph{explicit} dependence on the canonical
jet. The solution $x=x(Y)$ of the original Poisson equation still depends on
that jet, and so do the resulting six rational source functions.

\section{Six fixed lapse--shift pairings}
\label{v37:sec:bilinear}

Let $n\in\R^{27}$ and $\beta_R\in\R^{81}$ retain point-coordinate ordering.
For the ordinary point columns define the rational matrix
\begin{equation}
 (C_a)_{jk}=\tfrac12 q_a(e^N_j+e^\beta_k),
                 \qquad C_a\in\R^{27\times81}.
 \label{v37:eq:C}
\end{equation}

\begin{proposition}[Bilinear form and its finite rank]
\label[proposition]{v37:prop:bilinear}
For all multiplier coordinates,
\begin{equation}
 \boxed{q_a(n,\beta_R)=2n^TC_a\beta_R.}
 \label{v37:eq:bilinear}
\end{equation}
Each $C_a$ has rank sixteen. Each pairing vanishes if its lapse is additive
axial, or its shift has the axial form in \eqref{v36:eq:N}. In particular a
pure lapse or pure shift makes every $q_a$ zero, although it need not solve
any lower equation.
\end{proposition}
\begin{proof}
For a pure shift test, the quadratic constraint has canonical type $Up_R$.
The same-type shift--shift block of $\mathbb K_1$ is linear in momentum,
and the mixed shift--lapse block is linear in configuration. Since
$G(n,\beta_R)$ has configuration depending only on $\beta_R$ and momentum
depending only on $n$, $q_a$ has type $n\beta_R$ and no pure square term.
Polarization proves \eqref{v37:eq:bilinear} with the normalization
\eqref{v37:eq:C}.

For reproducibility an exact evaluation of a row uses
$H_\nu(V,W)=D^2P_{\nu,2}[V,W]$ and
\[
 Dq_\nu(e)[z]=2H_\nu(Ge,Gz)-H_z(Ge,G\nu)-H_e(Gz,G\nu).
\]
For each of the 27 lapse columns this evaluates every shift column
simultaneously. The package regenerates all $6\cdot27\cdot81=13122$
rational entries. Transforming by the integer tensor basis
$\{1,c,s\}^{\otimes3}$, $c=(2,-1,-1)$ and $s=(0,1,-1)$, leaves exactly
forty nonzero cross entries per test. There are sixteen nonzero lapse rows,
and their submatrix on the shift component in the test's axial direction
has rank sixteen. Exact row reduction of this displayed sparse bank checks
both assertions, proving the stated upper and lower rank bounds. The
annihilation of the axial subspaces also follows directly from the stronger
translation identity \cref{v36:thm:invariance} and \eqref{v37:eq:bilinear}.
\end{proof}

\begin{corollary}[A positivity certificate must use the source equation]
\label[corollary]{v37:cor:indefinite}
A constant real linear combination of the six $q_a$ cannot be a nonzero
positive-semidefinite quadratic form on the full multiplier space.
\end{corollary}
\begin{proof}
Every such combination is $2n^TC\beta_R$. Replacing $\beta_R$ by
$-\beta_R$ reverses its sign. If it is nonzero, suitable $n,\beta_R$
therefore give both signs.
\end{proof}

Thus a possible separating identity for the remaining three-profile problem
must retain the original Poisson/source incidence or a restriction it implies.
There is no global positive multiplier form to be obtained merely by summing
these six tests. On the generic set of Version 36 the exact problem is now
\[
 \widehat K(Y)x=-\widehat b(Y),\qquad
                     n^TC_a\beta_R=0\quad(1\le a\le6),
\]
followed by the already proved kernel corrections for the means and cubic
rows. This polynomial representation does not permit deleting the generic
rank exclusions when it is used as an equivalence.

\section{An explicit exact root on a two-profile boundary}
\label{v37:sec:root}

Choose the rational first-jet data
\begin{equation}
 \boxed{\alpha=(0,2,1),\qquad b=(1/4,-1,-2).}
 \label{v37:eq:root-data}
\end{equation}
They obey the exact scalar mean relation
$b_1b_2+b_1b_3+b_2b_3=5/4=|\alpha|^2/4$.
Let $f(x)=(1,c(x),s(x))^T$. Define the rational multiplier coordinates,
all independent of $x_1$, by
\begin{equation}
 n(x)=f(x_2)^TA_Nf(x_3),\qquad
 (\beta_R)_i(x)=f(x_2)^TA_if(x_3),\qquad \ell=Q_3(n,\beta_R),
 \label{v37:eq:root-fields}
\end{equation}
where
\begin{align}
 A_N&=\frac1{2450304}
 \begin{pmatrix}
 0&10739400&295992\\
 -1605011&-15252&155376\\
 -115259&-302004&140616
 \end{pmatrix},\notag\\
 A_1&=\frac1{102096}
 \begin{pmatrix}
 0&-168&13352\\-41616&-1729&-19052\\37136&3571&-7854
 \end{pmatrix},\notag\\
 A_2&=\frac1{288}
 \begin{pmatrix}0&18&-18\\26&-1&27\\12&0&12\end{pmatrix},\qquad
 A_3=\frac1{288}
 \begin{pmatrix}0&23&3\\144&-4&0\\-144&27&3\end{pmatrix}.
 \label{v37:eq:root-matrices}
\end{align}

\begin{theorem}[An exact full first-jet coefficient root]
\label[theorem]{v37:thm:root}
The literal arrays \eqref{v37:eq:root-data}--\eqref{v37:eq:root-matrices}
satisfy
\begin{equation}
 \boxed{\widehat Kx=-\widehat b,\qquad
 \mathsf mx=0,\qquad \widehat Rx=-t,\qquad
 \mathcal T_\nu(Y,Q_3x)=0\quad
                   (\nu\in\mathcal W\cap\ker\mathbb K_1(Y)).}
 \label{v37:eq:root-all}
\end{equation}
At this point the Poisson rank is 92, the weak-null dimension is nine, and
the complete lower rank is 99. The positive physical multiplier norm is
\begin{equation}
 \boxed{\|\ell\|_h^2=\frac{96734411978389}{2251496134656},
                  \qquad 6.55<\|\ell\|_h<6.56.}
 \label{v37:eq:root-norm}
\end{equation}
These statements concern the finite first-jet equations, not a nonlinear
exact initial history or an Einstein spacetime.
\end{theorem}
\begin{proof}
Substitute the four integer numerator matrices into the original rational
arrays. The exact verifier independently re-evaluates all 108 Poisson rows,
the four means, and the three cubic constant-shift equations in characteristic
zero. It does not accept a reconstructed or rounded matrix equation as the
verification. Specifically, it applies the polarized original $P_2$ directly
to $Y_R,Gx$, retains the rationalized shift-block factor three, and recomputes
$\widehat b$ with its lapse preparation term.

The cubic check can be reproduced without a matrix inverse. For any
multiplier $x$, let $P_{x,2}$ denote the actual quadratic constraint. The
unscaled constant-row action is
\[
 R_Yx\big|_i=D P_{i,3}(Y)[Gx]
       -DP_{x,2}(Y)[\delta_iY]-P_{\delta_i x,2}(Y).
\]
The last sign is exact discrete integration by parts. The rational physical
conversion is $\widehat Rx=3R_{(U,p_R)}x-2R_{(U,0)}x_N$ and
$\widehat t=3t(U,p_R)-2t(U,0)$. Those direct rational values give all three
zeros in \eqref{v37:eq:root-all}. All six nonconstant tangency rows are also
checked from the complete original $Dd+\mathbb K_1$ expression, independently
of using only \eqref{v37:eq:bilinear}.

The inherited fixed kernel gives rank at most 92. A nonzero original
92-column Poisson minor at a good prime supplies the matching lower bound.
The weak kernel is consequently the same nine-dimensional axial shift
space as before. Its three constant tests vanish by the retained translation
and linear-constraint identity; the other six were checked directly. For the
cubic means, use columns $3,4,5$ of the integer kernel bank of Version 36,
namely $c(x_2),s(x_2),c(x_3)$ in the lapse slot. Direct rational contraction
gives the determinant
\[
 \det(\widehat R N_{[:,(3,4,5)]})=-\frac{847288609443}{16}\ne0.
\]
They have zero means; the four constant multiplier columns independently
control the means. Together with the rank-92 Poisson block this proves lower
rank $92+4+3=99$. This uses a different three-column minor from the one
selected in Version 36, which vanishes at the present point.
Finally sum the squared rational arrays with weight $1/27$, including the
factor three in the physical shift norm. The exact result is
\eqref{v37:eq:root-norm}; rational squaring proves its two numerical bounds.
\end{proof}

\section{Why the boundary root is not the desired three-profile entrance}
\label{v37:sec:boundary}

The first component of $\alpha$ in \eqref{v37:eq:root-data} is zero. Thus this
point is outside the three-profile open set of Version 36, even though its
Poisson and complete lower ranks retain their generic values. It cannot be
used as a root inside that open set by suppressing the nonzero-profile
condition.

There is also an analytical distinction at initialization. Both canonical
first fields are independent of $x_1$. For the original quadratic
constant-shift constraint the retained translation identity gives, up to the
fixed sign convention, its canonical Hamiltonian vector $\delta_1Y$.
Consequently
\begin{equation}
 \boxed{DP_{(0,e_1),2}(Y)[Z]=0
                    \quad\hbox{for every canonical variation }Z.}
 \label{v37:eq:mean-degenerate}
\end{equation}
Indeed $\delta_1U=\delta_1p=0$, and the symplectic pairing with $Z$ is zero.
Therefore the four quadratic mean-preparation derivatives cannot be onto
at this first jet. This is not merely the failure of one previously chosen
balancing direction; the entire indicated row vanishes.

The old four-mean implicit-function initializer consequently does not apply
at this root. A singular mean solve, an additional symmetry argument, or a
new source family would have to be supplied. No impossibility of every such
initializer is claimed. Similarly, the higher signed source Gram and the
preservation of the tangency condition by the original flow have not been
verified here. The exact root is a nonempty coefficient-level boundary case,
not a hidden replacement for a nondegenerate Einstein branch.

\section{A nondegenerate three-profile ray has a removable Poisson singularity}
\label{v37:sec:ray}

Keep all three oscillatory profiles fixed and nonzero, and put
\begin{equation}
 \boxed{\alpha=(2,9/8,1),\quad
 b_2=-R,\quad b_3=-2R,\quad
 b_1=\frac{2R}{3}-\frac{401}{768R},\qquad R>0.}
 \label{v37:eq:ray}
\end{equation}
The quadratic mean relation holds identically. The variable $z=R^{-1}$
below is an algebraic parameter, not physical time.

Let $V$ synthesize the integer tensor-Fourier complement of $\mathcal N$:
retain all columns of $I_4\otimes\{1,c,s\}^{\otimes3}$ except the sixteen
columns spanning \eqref{v36:eq:N}. Then $V$ has 92 columns and is a fixed
complement throughout the family. The actual source equation on it has
\begin{equation}
 M(z)y=-g(z),\qquad
 M(z)=V^T[z\widehat K(Y(z))]V=M_0+zM_1+z^2M_2,
 \quad g(z)=g_0+zg_1+z^2g_2.
 \label{v37:eq:pencil}
\end{equation}
Here $g(z)=V^T[z\widehat b(Y(z))]$. These exact polynomial expressions
follow by the linear and quadratic field degrees and the vanishing of the
pure homogeneous-momentum prepared source. They are independently
regenerated from the original rational coefficient functional in the package.

\begin{lemma}[Exact leading splitting and source cancellation]
\label[lemma]{v37:lem:split}
For \eqref{v37:eq:pencil}, $\operatorname{rank}M_0=72$. Let $E$ insert its
72 pivot coordinates in exact row-reduction order, and let $Z$ be the null
basis with identity on the complementary twenty coordinates. Put
$T=(E,Z)$ and $\widetilde M_j=T^TM_jT$, $\widetilde g_j=T^Tg_j$. Then
\begin{equation}
 \widetilde M_0=\begin{pmatrix}A_0&0\\0&0\end{pmatrix},\quad
 \widetilde M_1=\begin{pmatrix}A_1&B_1\\C_1&0\end{pmatrix},\quad
 \widetilde M_2=\begin{pmatrix}A_2&B_2\\C_2&D_2\end{pmatrix}.
 \label{v37:eq:blocks}
\end{equation}
Writing $u_0=-A_0^{-1}\widetilde g_{0,A}$, the actual source satisfies
\begin{equation}
 \widetilde g_{0,B}=0,\qquad C_1u_0=-\widetilde g_{1,B},
 \qquad S_2:=D_2-C_1A_0^{-1}B_1\text{ is invertible}.
 \label{v37:eq:split-cancel}
\end{equation}
In these exact coordinates
$\det A_0\equiv233$ and $\det S_2\equiv40\pmod{241}$;
all denominators are units at this prime.
\end{lemma}
\begin{proof}
The three matrices and the three source vectors are finite rational arrays
specified by \eqref{v37:eq:ray}--\eqref{v37:eq:pencil}. Exact elimination
of $M_0$ supplies the pivot/null decomposition. Direct rational multiplication
checks the zero blocks, $M_0Z=0$, and the two source cancellations. The
nonzero determinant residues, supplemented by the exact rational
invertibility checks, prove the two stated inverses. The zero claims are
not inferred from modular arithmetic. This is a calculation on the original
source, not an assumed solvability condition for a general pencil.
\end{proof}

Substitute $T^{-1}y=(u_0+z\widetilde u,v)$. The top equation of
\eqref{v37:eq:pencil} has a removable factor $z$ after its constant row is
cancelled; the bottom equation has a removable factor $z^2$ by
\eqref{v37:eq:split-cancel}. The exact resulting system is
\begin{equation}
 \boxed{\mathcal H(z)\binom{\widetilde u}{v}=h_0+zh_1,\qquad
 \mathcal H(z)=
 \begin{pmatrix}
 A_0+zA_1+z^2A_2&B_1+zB_2\\
 C_1+zC_2&D_2
 \end{pmatrix},}
 \label{v37:eq:desingularized}
\end{equation}
where
\[
 h_0=\binom{-\widetilde g_{1,A}-A_1u_0}
                 {-\widetilde g_{2,B}-C_2u_0},\qquad
 h_1=\binom{-\widetilde g_{2,A}-A_2u_0}{0}.
\]

\begin{theorem}[A bounded analytic Poisson representative at large momentum]
\label[theorem]{v37:thm:analytic}
For all sufficiently small $|z|$, \eqref{v37:eq:desingularized} has a unique
real-analytic solution. For $z\ne0$ it gives a solution
\begin{equation}
 x_{\rm q}(z)=VT\binom{u_0+z\widetilde u(z)}{v(z)}
 \label{v37:eq:xanalytic}
\end{equation}
of every original Poisson row, with $\|x_{\rm q}(z)\|=O(1)$ as $z\to0$.
The original mean and cubic lower equations are also solvable on the ray
for all sufficiently large $R$.
\end{theorem}
\begin{proof}
The Schur complement of $A_0$ in $\mathcal H(0)$ is exactly $S_2$, so
\cref{v37:lem:split} proves invertibility. Analytic inversion proves the
first assertion. The two algebraic divisions used above were justified by
identities, not by dropping small source rows. Reversing them for $z\ne0$
proves the full quotient equation, and the universal source compatibility
from Version 36 gives all original rows. Formula \eqref{v37:eq:xanalytic}
gives boundedness.

The quotient invertibility and the fixed kernel give Poisson rank 92 for
small nonzero $z$. For the Version-36 three-column cubic control minor on
this ray, direct original-coefficient evaluation at $R=8$ gives a nonzero
value at a good prime, recorded in the package. Its entries are Laurent
polynomials in $z$, so its determinant is not identically zero and has no
zeros in some sufficiently small punctured neighborhood of zero. The four
means and cubic rows are then solved by the retained kernel correction.
This last assertion is existence of the full lower solution; boundedness of
all of those additional correction coefficients is not used or asserted.
\end{proof}

The analyticity is effective. Writing
$\mathcal H(z)=\mathcal H_0+z\mathcal H_1+z^2\mathcal H_2$, define
$m_j=\|\mathcal H_0^{-1}\mathcal H_j\|_\infty$ and
$\rho=[4(1+m_1+m_2)]^{-1}$. Then for $|z|\le\rho$ the Neumann
residual is at most $1/4$ and
$\|\mathcal H(z)^{-1}\|_\infty\le(4/3)\|\mathcal H_0^{-1}\|_\infty$.
This supplies a rational lower bound on a valid parameter neighborhood;
no fixed physical evolution interval is inferred from it.

\section{Exact fourth-order tangency suppression, without a root}
\label{v37:sec:tail}

Let $\tau(R)$ use the six tests, the actual source \eqref{v37:eq:ray},
and any complete lower solution. By Version 36's kernel invariance it can
be evaluated on $x_{\rm q}(R^{-1})$ instead.

\begin{theorem}[A certified anisotropic residual tail]
\label[theorem]{v37:thm:tail}
On the ray \eqref{v37:eq:ray}, in the original spatial-sum rational
normalization,
\begin{equation}
 \boxed{\tau(R)=R^{-4}\gamma+O(R^{-5}),\qquad \gamma\ne0.}
 \label{v37:eq:tail}
\end{equation}
All six coefficients are nonzero. In particular
\begin{equation}
 \gamma_1=
 \frac{134919942606783961765870993920227582667}
 {34331898520299363323114467596697600},\qquad 3929<\gamma_1<3930.
 \label{v37:eq:gamma1}
\end{equation}
The others have exact rational bounds
\begin{align*}
 6850&<\gamma_2<6851, & 4010&<\gamma_3<4011,\\
 -3423&<\gamma_4<-3422, & 1522&<\gamma_5<1523,\\
 4237&<\gamma_6<4238. &&
\end{align*}
Thus for some finite $R_0$ and all $R\ge R_0$, the lower system is
solvable but the complete tangency system has no root at that first jet.
Nevertheless the infimum of the six-component residual norm over that ray
is zero.
\end{theorem}
\begin{proof}
For $w=(\widetilde u,v)$ expand the analytic system
\eqref{v37:eq:desingularized} as $w(z)=\sum_{k\ge0}w_kz^k$. Its exact
recurrence is
\begin{align*}
 \mathcal H_0w_0&=h_0,\qquad
 \mathcal H_0w_1=h_1-\mathcal H_1w_0,\\
 \mathcal H_0w_k&=-\mathcal H_1w_{k-1}-\mathcal H_2w_{k-2}
                                      \quad(k\ge2).
\end{align*}
All solves are rational. Let $x(z)=\sum_{k\ge0}x_kz^k$ be obtained from
\eqref{v37:eq:xanalytic}, and put
$Q_a=\left(\begin{smallmatrix}0&C_a\\C_a^T&0\end{smallmatrix}\right)$.
By \cref{v37:thm:pure,v37:prop:bilinear}, the coefficient of $z^k$ in the
actual $\tau_a$ is exactly
\begin{equation}
                 t_{a,k}=\sum_{i=0}^k x_i^TQ_ax_{k-i}.
 \label{v37:eq:recurrence-tau}
\end{equation}
The package constructs $w_0,\ldots,w_6$ from the exactly regenerated
source matrices, not from fitted large-$R$ values. Rational evaluation
of \eqref{v37:eq:recurrence-tau} gives $t_{a,k}=0$ for all six $a$ and
$0\le k\le3$, and $t_{a,4}=\gamma_a$ with the displayed values. The
24 lower-order cancellations are exact. The full fractions for all six
fourth- and fifth-order terms and all series vectors are retained in the
certificate. Rational numerator comparisons prove each interval above.
Analyticity supplies the stated remainder; the positive first component
then excludes a zero for all sufficiently large $R$.

A finite remainder enclosure is also available directly. For
$w^{[n]}=\sum_{k=0}^n w_kz^k$, its system residual is
\[
 z^{n+1}(\mathcal H_1w_n+\mathcal H_2w_{n-1})
                      +z^{n+2}\mathcal H_2w_n.
\]
Multiply by the inverse bound following \cref{v37:thm:analytic}, then
use \eqref{v37:eq:xanalytic} and the finite quadratic matrices $Q_a$.
This yields an explicit rational constant for the $O(z^5)$ remainder
if a numerical threshold is required. No floating sample is needed for
the asymptotic assertion or the existence of $R_0$.
\end{proof}

The four-order cancellation is not a consequence of the small lapse
alone. The exact series gives zero limiting lapse and lapse size $O(R^{-1})$,
but pairing that estimate with the bounded shift gives only $O(R^{-1})$.
The further cancellations through order three use the actual source solve.
This is why neither an empirical slope nor a crude bilinear norm bound
would prove \eqref{v37:eq:tail}.

\section{What the two limiting regimes do to the first jet}
\label{v37:sec:interpretation}

The boundary root and the anisotropic ray have different meanings. The first
is a finite exact coefficient root with one absent spatial profile and a
singular mean-preparation derivative. The second keeps every oscillatory
profile nonzero but increases the homogeneous canonical momentum without
bound. It proves that arbitrarily small generic tangency residuals exist;
it does not turn any finite point on its tail into an exact root.

On \eqref{v37:eq:ray}, $\|Y_R\|\asymp R$ while the oscillatory
configuration is fixed. Therefore after division by the whole first-jet
size the limit is a spatially homogeneous momentum jet:
\[
 R^{-1}U\longrightarrow0,\qquad
 R^{-1}p_R\longrightarrow
 3\{\operatorname{diag}(2/3,-1,-2)-(2/3-1-2)I\}.
\]
In a putative near-flat physical preparation $X=aY+\text{remainder}$
with a remainder smaller than its leading term, this regime would require
$aR\to0$. That observation is only a necessary normalization condition.
Neither uniform initializer estimates at growing $R$, the scaling of the
higher signed inverse, nor a common physical lifespan has been proved here.

Consequently a useful search for a nondegenerate finite root must distinguish
bounded ratio regions with all profiles separated from zero from both
boundary collapse and unbounded homogeneous momentum. On such a bounded
region the six equations remain a genuine simultaneous source problem.
The two constructions above do not prove its existence or nonexistence.
Version 32's unit-initial-multiplier compensation problem remains independent.

\section{Verification ledger and the remaining mathematical target}
\label{v37:sec:status}

\paragraph{Proved.}
The explicit canonical terms cancel from all six tangency polynomials on the
entire six-dimensional source span. The surviving forms are fixed rank-16
lapse--shift pairings. A fully specified rational two-profile point solves
all lower and weak-null first-jet equations, with finite multiplier norm;
its missing mean-preparation derivative is also identified. On a separate
nondegenerate anisotropic ray, an exact rank-72/20 splitting makes the
Poisson solve analytic and bounded in a fixed quotient. Its six residuals
have nonzero exact fourth-order leading coefficients, so the residual can
approach zero without a root anywhere on the sufficiently far tail.

\paragraph{Inherited.}
The universal multiplier kernel, actual-source compatibility, weak-space
identification at rank 92, and generic kernel corrections of Version 36
are used at their stated finite scopes. The original $P_2,\Phi_3$ compiler
and its physical conventions are unchanged. No independent verification of
every preceding analytic theorem is claimed.

\paragraph{Still unresolved.}
A root with all three profiles nonzero at finite first-jet ratios has not
been certified. The boundary root does not have the regular four-mean
initializer supplied earlier; a replacement has not been constructed.
Neither branch has a new higher signed-source certificate, propagated
nonlinear tangency theorem, common Einstein lifespan, or cofinal physical
limit. No arbitrary microscopic Renewal-to-Einstein--SM implication follows.

\paragraph{Executed arithmetic and preservation.}
The verifier checks 4104 exact canonical-cancellation entries, regenerates
all 13122 cross coefficients and their ranks, and directly verifies the
121 directly evaluated root equations of \cref{v37:thm:root}. It regenerates
the three full original ray Poisson matrices and three prepared source
coefficients in exact arithmetic, verifies the leading block and source
identities, and applies rational recurrences for the asymptotic coefficients.
Nonzero-minor tests are kept separate from zero-identity proofs. Numerical
searches and large-$R$ comparisons are diagnostics, not root or asymptotic
certificates. Every Version 36 byte before its final document command is
preserved. The package records source hashes and the executed stages; no
proof-assistant formalization or nonlinear spacetime simulation is asserted.

\begingroup
\renewcommand{\refname}{Additional background for Version 37}
\begin{thebibliography}{9}
\small\raggedright
\bibitem{v37:DBGP}
C. Di Bartolo, R. Gambini and J. Pullin,
\emph{Consistent and mimetic discretizations in general relativity},
J. Math. Phys. \textbf{46} (2005), 032501;
\url{https://arxiv.org/abs/gr-qc/0404052}.
Background for the distinction between discrete multiplier consistency and
constraint-preserving evolution. None of the present finite source values,
root certificates, or asymptotic estimates is imported from that work.
\end{thebibliography}
\endgroup
\addtocontents{toc}{\protect\endgroup}
% END IMMUTABLE VERSION 37 BODY
% APPEND VERSION 38 HERE, BEFORE THE FINAL END-DOCUMENT COMMAND.


% BEGIN IMMUTABLE VERSION 38 BODY
\clearpage
\begin{center}
{\Large\bfseries Version 38: Transverse isolation of the boundary root}\\[0.5em]
{\large Exact normal coefficients, arbitrary standing-wave phases, and preparation conditioning}
\end{center}
\makeatletter
\addtocontents{toc}{\protect\begingroup}
\addtocontents{toc}{\protect\makeatletter}
\addtocontents{toc}{\protect\def\protect\l@section{\protect\bfseries\protect\@dottedtocline{1}{0em}{2.5em}}}
\makeatother
\addcontentsline{toc}{part}{Version 38: Phase-independent transverse isolation of the boundary root}

\section{The continuation question and the scope of the new calculation}
\label{v38:sec:context}

Version 37 supplies an exact first-jet coefficient root with
$\alpha=(0,2,1)$ and $b=(1/4,-1,-2)$. It does not supply a three-profile
root or a regular four-mean initializer there. This body tests the proposed
transverse continuation, rather than re-proving the lower-solvability
criterion. The outcome is a source-specific exclusion: no sufficiently
nearby complete coefficient root can switch on the missing first profile.
The exclusion persists when arbitrary standing-wave phases are permitted
in all three directions.

The new ingredient is an exact nonzero quadratic normal map of two actual
weak-null tangency rows. Its weighted norm is independent of the phase
of the missing profile. Consequently a small tangency residual near this
root measures collapse of that profile, not proximity to a nearby
three-profile root. The same collapse makes the original quadratic
mean-preparation derivative ill-conditioned.

All statements below have the original spatial cutoff $h=1/3$ and original
physical normalizations. A neighborhood means a finite-dimensional
parameter neighborhood, not a spatial-refinement neighborhood.
The primitive update, action, multiplier equations, and physical clock
are unchanged. The wider phase family consists of different initial
coefficients of the same action. No nonlinear prepared history or common
physical lifespan is constructed.

\paragraph{Inherited inputs and new verification.}
We use the universal kernel and prepared-source compatibility of
\cref{v36:thm:kernel}, the fixed quotient and lower corrections of
\cref{v36:thm:solve}, and the canonical cancellation and two-profile
root of \cref{v37:thm:pure,v37:prop:bilinear,v37:thm:root}.
Their extension to the phase family is proved below by exact cyclic
translation and polynomial exhaustion. The new normal coefficients
are verified directly against the original coefficient functional;
no numerical slope is used as a derivative certificate.
The distinction between discrete multiplier consistency and subsequent
evolution is also standard background \cite{v37:DBGP}, not a conclusion
obtained by solving the present finite equations.

\section{An eight-parameter phase family with exactly solvable lower equations}
\label{v38:sec:family}

On the three-point circle set
\[
 c_0=(1,-1/2,-1/2),\qquad s_0=(0,1,-1),\qquad
 u_i(x_i)=\alpha_i c_0(x_i)+\sigma_i s_0(x_i).
\]
The actual sine amplitude is $2\sigma_i/\sqrt3$, since the sampled sine
is $(\sqrt3/2)s_0$. Retain the original transverse tensors $T_i,P_i$
and rational momentum convention:
\begin{align}
 U&=\sum_{i=1}^3u_iT_i,\notag\\
 p_R&=3\{\operatorname{diag}(b)-(\operatorname{tr}\operatorname{diag}(b))I\}
                     +\frac32\sum_{i=1}^3u_iP_i,\qquad p=\sqrt3p_R.
 \label{v38:eq:family}
\end{align}
The four quadratic means vanish when
\begin{equation}
 b_1=\frac{\frac14\sum_i(\alpha_i^2+\frac43\sigma_i^2)-b_2b_3}{b_2+b_3}.
 \label{v38:eq:mean}
\end{equation}
Indeed $\langle u_i^2\rangle=\alpha_i^2/2+2\sigma_i^2/3$;
the complementary tensors have equal Frobenius norms and zero inner
product. The linear constraints follow from their axial transversality,
and the constant-shift means vanish from the same inner-product identity.
Equation \eqref{v38:eq:mean} is the retained scalar mean, not a new source
selection rule.

Write
\[
 z=(\alpha_1,\sigma_1),\qquad
 q=(\alpha_2,\sigma_2,\alpha_3,\sigma_3,b_2,b_3),\qquad
 q_*=(2,0,1,0,-1,-2).
\]
The boundary root has $(z,q)=(0,q_*)$. Denote its transverse amplitude by
\begin{equation}
 \rho^2=\alpha_1^2+\frac43\sigma_1^2.
 \label{v38:eq:rho}
\end{equation}

\begin{lemma}[Polynomial exhaustion by the three cyclic phase lines]
\label[lemma]{v38:lem:lines}
Let a vector-valued polynomial have degree at most two in each of several
two-dimensional blocks. If it vanishes when each block is restricted to
any of three distinct lines through zero, independently of the other
blocks and of unrestricted additional variables, then it vanishes identically.
\end{lemma}
\begin{proof}
For one block, restrict to each line and vary its scalar parameter.
Every homogeneous piece of degrees zero, one, and two vanishes on those
lines. A nonzero homogeneous binary polynomial of degree at most two
cannot have three distinct linear factors. Its coefficients, which can
depend polynomially on the remaining variables, therefore vanish.
Repeat one block at a time. Applying the argument to each scalar
coordinate proves the vector statement.
\end{proof}

\begin{proposition}[The phase extension preserves the actual finite algebraic reductions]
\label[proposition]{v38:prop:extension}
On the nine-dimensional linear span of the three cosine profiles, three
sine profiles, and three homogeneous momenta, one has
\[
 \widehat K(Y)N=0,\qquad N^T\widehat b(Y)=0,\qquad
 \Theta_a(Y,x)=2n^TC_a\beta_R\quad(1\le a\le6).
\]
Here $N$ is exactly the sixteen-column kernel bank of Version 36 and
$C_a$ are exactly the six Version-37 matrices. On a neighborhood
$\mathcal U_*$ of $(0,q_*)$, the Poisson rank is 92 and the complete
lower rank is 99. Every parameter there admits an analytic solution of
all Poisson, mean, and cubic constant-shift equations.
\end{proposition}
\begin{proof}
Cyclic translation in each spatial coordinate is an exact symmetry of
the retained scalar coefficient functional: it commutes with
$\delta_i=3(S_i-S_i^{-1})$, the literal shifts, and the spatial sum.
It also commutes with $F_1$ and intertwines $G$. The subspaces $N$ and
the six axial nonconstant weak tests are invariant under these
translations, the latter with an invertible two-dimensional change
of test basis in each direction.

The three translates of the cosine span three distinct lines in its
cosine--sine plane. For example,
\[
 S c_0=-\tfrac12c_0-\tfrac34s_0,\qquad
 S^{-1}c_0=-\tfrac12c_0+\tfrac34s_0.
\]
The inherited kernel identity is linear in the canonical jet; translating
it and taking spans proves it on all nine columns. The source compatibility
is homogeneous quadratic. It holds for arbitrary cosine amplitudes and
homogeneous momenta and thus, by independent translations, on every product
of those three phase lines. \Cref{v38:lem:lines} proves it on the entire
phase span.

For the tangency identity, the canonical-only part is quadratic in the
canonical jet and the mixed part is bilinear in the jet and the arbitrary
multiplier. The two zero identities of \cref{v37:thm:pure} transform
covariantly with the test bank. The same three-line argument extends their
vanishing. The remaining multiplier-only form is unchanged. In particular
its translation invariance under addition of $N$ is retained.

At the root, \cref{v37:thm:root} supplies a nonzero Poisson minor of size 92.
It stays nonzero nearby, while the universal sixteen-dimensional kernel
gives the reverse bound. Its exact weak intersection consists of the same
nine axial shift profiles. The alternative cubic-control columns $3,4,5$
of $N$ have determinant $-847288609443/16$ at the root, and hence stay
independent nearby. Together with the four constant multiplier columns
they control all remaining seven lower rows. Prepared-source compatibility
proves solvability of the Poisson equation. Its fixed-complement inverse,
followed by the original mean and three cubic corrections, proves the
asserted analytic lower solution. The actual cubic right side is retained.
\end{proof}

This result supplies lower-incidence solvability for the phase extension.
It does not assert regularity of the four \emph{canonical mean-preparation}
derivatives at $z=0$, which is a different map.

\section{The first-axis tests see only first-axis nonconstant multiplier modes}
\label{v38:sec:selection}

Let $P_0$ average a scalar array over $x_1$, at fixed $(x_2,x_3)$,
and act componentwise on shift arrays. Use the fixed tensor-Fourier
complement $V$ from Version 37; its range is the orthogonal complement
of $N$, invariant under all spatial translations. Define the unique
quotient representative
\begin{equation}
 x(z,q)=-V\{V^T\widehat K(Y(z,q))V\}^{-1}V^T\widehat b(Y(z,q)).
 \label{v38:eq:quotient}
\end{equation}
It is analytic on a smaller $\mathcal U_*$. Kernel corrections to solve
the full lower system do not change its tangency values. Let
$\tau_c,\tau_s$ denote the tests $\nu_{1,c},\nu_{1,s}$, in that order,
with $c=(2,-1,-1)$ and $s=(0,1,-1)$ as in Version 36.

\begin{lemma}[An exact two-sided mode exclusion]
\label[lemma]{v38:lem:mode}
For each of the first two original cross matrices,
\begin{equation}
 P_0^TC_a=0,\qquad
 C_a\operatorname{diag}(P_0,P_0,P_0)=0,\qquad a=1,2.
 \label{v38:eq:mode}
\end{equation}
Consequently, if $x=(n,\beta_R)$,
\[
 q_a(x)=2((I-P_0)n)^TC_a((I-P_0)\beta_R).
\]
Also $x(0,q)=P_0x(0,q)$.
\end{lemma}
\begin{proof}
In the exact tensor basis $\{1,c,s\}^{\otimes3}$, the two supplied
Version-37 sparse cross matrices have no nonzero entry whose lapse
or shift coordinate is constant in $x_1$. The verifier reconstructs
these matrices from the eighty retained rational sparse entries and
checks both matrix products in \eqref{v38:eq:mode}. This is a new use of
the inherited exact cross bank, not a newly claimed regeneration of its
entire compiler.

At $z=0$ the canonical coefficients and prepared source are independent
of $x_1$. Translation of the quotient solution in that coordinate
therefore solves the same original equation and stays in $\Ran V$.
Uniqueness in the complement implies translation invariance, equivalent
to $P_0x=x$.
\end{proof}

Analyticity now gives
\[
 (I-P_0)x(z,q)
   =\alpha_1X_c(q)+\sigma_1X_s(q)+O(\rho^2).
\]
Because both arguments of each cross form must be first-axis nonconstant,
the first two residuals have no constant or linear term in $z$ for
\emph{every} nearby $q$. This is stronger than merely observing a
quadratic residual on one numerical parameter ray.

\section{Exact normal coefficients from the original source equation}
\label{v38:sec:coefficients}

At the boundary root, write $K_0=\widehat K(Y_*)$ and
$b_0=\widehat b(Y_*)$. Let $Y_c,Y_s$ be the first-axis cosine and
raw-sine columns in \eqref{v38:eq:family}. Their normal derivatives
do not change $b_1$ at $z=0$ in \eqref{v38:eq:mean}. Set
\[
 K_\epsilon=\widehat K(Y_\epsilon),\qquad
 b_\epsilon=D\widehat b(Y_*)[Y_\epsilon],
 \qquad \epsilon=c,s.
\]
The quotient vectors $x_0,X_c,X_s$ are specified by the original
full-row equations
\begin{equation}
 \begin{aligned}
 K_0x_0&=-b_0,&x_0&\in\Ran V,\\
 K_0X_c&=-b_c-K_cx_0,&X_c&\in\Ran V,\\
 K_0X_s&=-b_s-K_sx_0,&X_s&\in\Ran V.
 \end{aligned}
 \label{v38:eq:derivatives}
\end{equation}
These are well posed by the proved quotient inverse. They are exactly
the differentiated source equation, not differentiation of an approximate
pseudoinverse.

\begin{theorem}[The two rational quadratic normal forms]
\label[theorem]{v38:thm:coefficients}
Define $D,A,B$ by
\begin{align}
 D&=587693396116325617134218074626554406701361841,\notag\\
 A&=-\frac{2500363378752229226973293064905761635211973321676}{D},\notag\\
 B&=-\frac{2092474420167555721391383484185182870099150711594}{D}.
 \label{v38:eq:AB}
\end{align}
The quadratic normal terms of $(\tau_c,\tau_s)$ at $(0,q_*)$ are
\begin{align}
 Q_c(a,s)&=A a^2+4B as-\tfrac43 A s^2,\notag\\
 Q_s(a,s)&=B a^2-\tfrac43 A as-\tfrac43 B s^2.
 \label{v38:eq:quadratics}
\end{align}
In particular
\[
 -4255<A<-4254,\qquad -3561<B<-3560,\qquad
 7492<\sqrt{A^2+3B^2}<7493 .
\]
\end{theorem}
\begin{proof}
The certificate gives all 108 rational entries of the three vectors in
\eqref{v38:eq:derivatives}. It checks their membership in $N^\perp$
and checks all 324 equations directly using the original polarized
$P_2$ and prepared $B_2$, including its scalar-row correction and
the factor three in the rationalized shift--shift block. For example
\[
 b_\epsilon=
 \tfrac12\{\widehat b(Y_*+Y_\epsilon)-\widehat b(Y_*-Y_\epsilon)\}.
\]
This is exact homogeneous quadratic polarization. A proposed matrix
or floating derivative is not the final verification.

For each of the two tests, the original functional $q_\nu(x)$ is
evaluated on $X_c,X_s,X_c+X_s$. The resulting symmetric matrices are
exactly
\[
 \begin{pmatrix}A&2B\\2B&-4A/3\end{pmatrix},\qquad
 \begin{pmatrix}B&-2A/3\\-2A/3&-4B/3\end{pmatrix}.
\]
The first cosine values are also checked independently in the
zero-order/cosine-derivative stage. The two-sided mode exclusion removes
every term involving $x_0$ and a second derivative of $x$.
Hence these matrices are the full quadratic Taylor coefficient,
not merely part of it. Direct integer comparisons prove the intervals,
and comparing $A^2+3B^2$ with $7492^2,7493^2$ proves the modulus bound.
\end{proof}

The large rational numbers fix the inherited field, test, and spatial-sum
normalizations. They are not a new physical coupling. Decimal values
$A\simeq-4254.53713667$, $B\simeq-3560.48652919$ are explanatory only.

\section{A complex-square normal map excludes every phase of the missing profile}
\label{v38:sec:isolation}

Set
\[
 w=\alpha_1+\frac{2i}{\sqrt3}\sigma_1,\qquad
 \mathcal T(z,q)=\tau_c(z,q)+i\sqrt3\,\tau_s(z,q),\qquad
 \Gamma_*=A+i\sqrt3 B .
\]
Then $|w|=\rho$. The exact two real quadratic forms in
\eqref{v38:eq:quadratics} combine as
\begin{equation}
 Q_c+i\sqrt3 Q_s=\Gamma_*\overline w^{\,2}.
 \label{v38:eq:complex}
\end{equation}

\begin{theorem}[Phase-independent transverse isolation]
\label[theorem]{v38:thm:isolation}
There are a neighborhood of $(0,q_*)$ and a finite constant $C$ such that
\begin{equation}
 \left|\mathcal T(z,q)-\Gamma_*\overline w^{\,2}\right|
       \le C\{\rho+|q-q_*|\}\rho^2.
 \label{v38:eq:remainder}
\end{equation}
After reducing the neighborhood,
\begin{equation}
 \boxed{\frac{|\Gamma_*|}{2}\rho^2
       \le \sqrt{\tau_c^2+3\tau_s^2}
       \le\frac{3|\Gamma_*|}{2}\rho^2.}
 \label{v38:eq:lower}
\end{equation}
Therefore every complete lower-incidence and weak-null tangency root in
this neighborhood has $\alpha_1=\sigma_1=0$. No such root restores
the first profile, even with independent nearby changes in the other
two amplitudes, their phases, and the two homogeneous momentum parameters.
For any path approaching $(0,q_*)$ with $\rho>0$,
\[
 \frac{\sqrt{\tau_c^2+3\tau_s^2}}{\rho^2}\longrightarrow|\Gamma_*|.
\]
\end{theorem}
\begin{proof}
Analyticity of the quotient representative gives a uniformly bounded
second-order Taylor remainder for its first-axis nonconstant part.
The two-sided mode exclusion makes each tangency row quadratic in that
part. Its quadratic coefficients vary analytically with $q$. This proves
a remainder of size $O(\rho^3+|q-q_*|\rho^2)$; the coefficient at $q_*$
is \eqref{v38:eq:complex}. Since $\Gamma_*\ne0$ and
$|\overline w^2|=\rho^2$, choose the neighborhood so that
$C(\rho+|q-q_*|)\le|\Gamma_*|/2$. The two triangle inequalities give
\eqref{v38:eq:lower}. Every full tangency root annihilates these two
actual weak-null tests, giving $\rho=0$. Dividing
\eqref{v38:eq:remainder} by $\rho^2$ proves the last assertion.
\end{proof}

The neighborhood is proved to exist by the finite analytic bounds and
nonzero source coefficients. No numerical radius or interval enclosure
for that neighborhood is claimed. The theorem does not say that the
boundary root is isolated \emph{within} the boundary, nor that all
boundary parameters solve the other four tangency rows. It excludes a
transverse three-profile continuation near this specified root.
Already on the old cosine-only slice,
\[
 (\tau_c,\tau_s)=(A,B)\alpha_1^2+O(\alpha_1^3)
\]
with the other parameters fixed. Allowing a sine component rotates the
leading pair but cannot reduce its weighted norm to zero.

\section{Approximate source solves and the conditioning of mean preparation}
\label{v38:sec:conditioning}

The preceding exclusion is exact. It also gives a useful residual test
for computations that do not solve the Poisson equation exactly.

\begin{proposition}[A quotient-residual lower bound]
\label[proposition]{v38:prop:residual}
For any nearby parameter and any multiplier $x$, put
\[
 r=\widehat K(Y)x+\widehat b(Y),\qquad
 \varepsilon_T=\sqrt{q_c(x)^2+3q_s(x)^2}.
\]
There is a finite neighborhood-dependent $C_1$ such that
\begin{equation}
 \boxed{\frac{|\Gamma_*|}{2}\rho^2
       \le \varepsilon_T+C_1(\|r\|+\|r\|^2).}
 \label{v38:eq:residual}
\end{equation}
The constant does not grow with an arbitrary added vector in $N$.
\end{proposition}
\begin{proof}
Project $x$ onto the fixed complement $N^\perp$.
Its difference $e$ from \eqref{v38:eq:quotient} satisfies the quotient
equation with right side $r$, so $\|e\|\le C\|r\|$ by the bounded
local inverse. The two forms are unchanged by the removed kernel
component. Their quadratic expansion at the bounded analytic solution
therefore gives a pair error bounded by $C(\|r\|+\|r\|^2)$.
Combine this with \eqref{v38:eq:lower}.
\end{proof}

A second cost is imposed by the canonical mean equations. Use normalized
spatial mean, the Frobenius norm on symmetric fields, and the positive
canonical norm
$\|(Z_U,Z_p)\|^2=\|Z_U\|_h^2+\|Z_p\|_h^2$, with physical momentum $p$.
Let $\mathscr M_Y$ be the four quadratic canonical mean-constraint
derivatives. Its first constant-shift row satisfies the retained
translation identity
\[
 DP_{(0,e_1),2}(Y)[Z]
  =\langle\delta_1U,Z_p\rangle_h-\langle\delta_1p,Z_U\rangle_h .
\]

\begin{proposition}[Preparation loses conditioning at exactly the collapsing profile]
\label[proposition]{v38:prop:mean}
On \eqref{v38:eq:family}, the operator norm of that scalar row is exactly
\begin{equation}
 \|DP_{(0,e_1),2}(Y)\|=\frac{\sqrt{837}}2\,\rho.
 \label{v38:eq:row-norm}
\end{equation}
Every right inverse $\mathscr B_Y$ of the full four-row derivative,
when one exists, obeys
\begin{equation}
 \boxed{\|\mathscr B_Y\|\ge\frac{2}{\sqrt{837}\,\rho}.}
 \label{v38:eq:inverse}
\end{equation}
Restricting the allowable canonical variations cannot improve this bound.
For the exact lower solution, with
$\mu=\sqrt{\tau_c^2+3\tau_s^2}$, it follows that
\[
 \|\mathscr B_Y\|\ge
 \sqrt{\frac{2|\Gamma_*|}{837}}\,\mu^{-1/2}.
\]
\end{proposition}
\begin{proof}
Here $\omega_h=3\sqrt3$, $\|T_1\|_F^2=\|P_1\|_F^2=2$, and
$p_{\rm osc}=(\omega_h/2)\sum_i u_iP_i$.
Only the first profile contributes to $\delta_1Y$. Exact three-point
summation gives
\[
 \|\delta_1U\|_h^2=27\rho^2,\qquad
 \|\delta_1p\|_h^2=\frac{729}{4}\rho^2.
\]
The canonical symplectic map is an isometry in the stated norm.
Thus the norm of the displayed linear functional is
$\sqrt{27+729/4}\,\rho$, proving \eqref{v38:eq:row-norm}.
Apply a right inverse to the unit target vector in this particular shift
row. Its image $Z$ has row value one, so
$1\le(\sqrt{837}/2)\rho\|Z\|$, proving \eqref{v38:eq:inverse}.
The lower inequality in \eqref{v38:eq:lower} gives
$\rho^{-1}\ge\sqrt{|\Gamma_*|/(2\mu)}$, which yields the last estimate.
\end{proof}

This is a derivative conditioning statement, not a theorem about the
nonlinear preparation radius or its lifespan. It shows why reducing
the residual by collapsing the missing profile does not repair the
original regular-initialization requirement.

\section{Verification record and the next distinct source problem}
\label{v38:sec:status}

\paragraph{Proved.}
The lower algebraic reduction extends to arbitrary standing-wave phases
on the displayed eight-parameter mean-compatible family. Its exact
quotient is analytic near the Version-37 boundary point.
The two first-axis tangency rows have an explicitly verified nonzero
complex-square normal form. Their weighted norm is comparable to the
missing profile amplitude squared, uniformly in its phase and in nearby
changes of the other parameters. Every complete nearby coefficient root
therefore remains on the missing-profile boundary. Approximate lower
solves have a quotient-residual rejection bound, and the canonical
mean-preparation inverse has a quantified singular cost at that boundary.

\paragraph{What was actually checked.}
The source-level verifier checks all 108 base equations and both sets
of 108 differentiated Poisson equations by direct rational contraction
of the retained coefficient functional. It checks eight original
quadratic-form values across its two executed stages. The rational
vectors and the two normal matrices are printed in the executable
certificate. A separate arithmetic stage checks complement membership,
the four two-sided mode-exclusion matrix identities on the inherited
cross bank, the complex-square identity, its strict modulus interval,
the three-line interpolation rank, and the physical mean-row norm.
No finite-difference derivative, floating rank, or fitted asymptotic
coefficient is accepted as one of these equalities.

The supplied cross matrices and compiler are inherited Version-37 data,
with hashes recorded. The current calculation does not repeat all
Version-37 canonical-cancellation banks or every older analytic proof.
Exploratory floating solves suggested the normal scaling; the accepted
coefficient values were then verified against all original rows over
$\mathbb Q$. The analytic neighborhood is not assigned an unverified
numerical radius. These are executable arithmetic checks and written
proofs, not a proof-assistant formalization or an exact-action spacetime
integration.

\paragraph{Inherited boundaries.}
The physical quantum and renewal source construction, the separate
unit-initial-multiplier stress equation of Version 32, and the higher
signed-source and common-lifespan questions are unchanged. No
three-profile coefficient root, regular nonlinear preparation at this
boundary, propagated tangency, or primitive Einstein--Standard-Model
selection theorem follows from this iteration.

\paragraph{Direction.}
Repeating a local Newton search which approaches this boundary point,
even with phase freedom, is now excluded as a route to a three-profile
root. A genuinely different boundary bifurcation would first require
the two-component normal coefficient to vanish at another lower-regular
background. Alternatively one must seek a finite-distance interior root
or enlarge the canonical profile class. The present result does not
claim either exists. It determines the local zero-set obstruction
at the actual root rather than postulating another regularity certificate.

Every Version 37 byte before its final document command is preserved.
New labels use the prefix \texttt{v38:}.
\addtocontents{toc}{\protect\endgroup}
% END IMMUTABLE VERSION 38 BODY
% APPEND VERSION 39 HERE, BEFORE THE FINAL END-DOCUMENT COMMAND.


% BEGIN IMMUTABLE VERSION 39 BODY
\clearpage
\begin{center}
{\Large\bfseries Version 39: Mixed-mode cubic means and exact initial drift zeros}\\[.45em]
{\large A verified source root beyond the axial obstruction, with all profiles retained}
\end{center}
\hypersetup{pdftitle={Native Regenerative Stationarity Closure: cumulative V1--V39}}
\makeatletter
\addtocontents{toc}{\protect\begingroup}
\addtocontents{toc}{\protect\makeatletter}
\addtocontents{toc}{\protect\def\protect\l@section{\protect\bfseries\protect\@dottedtocline{1}{0em}{2.5em}}}
\makeatother
\addcontentsline{toc}{part}{Version 39: Mixed-mode cubic means and exact initial drift zeros}

\section{Companion audit and the changed finite problem}
\label{v39:sec:audit}

The four supplied companions change the useful next calculation. The exact-action
Einstein source has advanced to V28 \cite{v39:EB28}. Its mixed-frequency first
fields leave the axial Poisson stratum used in Versions 35--38: the first
Poisson bracket can have only its four constant null directions. In particular,
the old nonconstant weak-null tests, including the diamond test, are no longer
null on this enlarged class. The axial obstruction and the boundary normal form
proved earlier remain valid on their stated families; they are not obstructions
to this new mixed family.

The source proves that maximal first Poisson rank makes the complete prepared
quadratic drift compatible, with a unique four-mean-zero first multiplier tilt.
Second initial multiplier values can then control all 104 nonconstant cubic
drift rows. The unresolved coefficient is the three-component \emph{constant}
shift source, evaluated after that unique first tilt. The displayed V28 rank
witness has three nonzero cubic means. Neither its rank certificate nor its
quadratic cancellation is a solution of these three equations.

The other companions are retained without transporting conclusions across
inequivalent problems. Source-selection V35 \cite{v39:SS35} reconstructs a
stationary reverse relative to an actually specified faithful state and gives
an output-sensitive reverse test. It does not identify the historical matter
generator or make a reverse comparison an executed operation. Perturbative ESM
V23 \cite{v39:PE23} proves a derivative floor for its populated one-loop
counterterm and bounds its native cubic-artifact return; its first-background
identity still retains higher-antifield and two-background terms. This is not
a symmetry or propagation theorem for the uncompleted Lorentzian action here.
Yang--Mills V32 \cite{v39:YM32} controls a twice differentiable local logarithmic
remainder and normalized conditional laws at a common exterior. It explicitly
retains the leading moving minimum, determinant and rough factors, and does not
deduce unconditional mixing or a physical Hamiltonian gap. Those results supply
useful scope checks, not missing estimates for the deterministic Einstein flow.
This audit covers the latest bodies and the dependencies used below, not an
independent verification of every cumulative proof in the four sources.

The new result is an original-coefficient, computer-assisted zero of the three
cubic means, with a nonsingular three-parameter Jacobian. All three axial
profiles and all four added mixed profiles are nonzero there. A further analytic
argument uses that zero to construct exact initial records with both the original
constraints and their complete canonical drift zero. The latter is an initial-cut
statement: invariance of this selected set, regularity of the full multiplier
Hessian, and a common physical lifespan are not assumed or concluded.

\section{Independent mixed amplitudes and the original mean constraints}
\label{v39:sec:family}

Fix throughout this body the original three-site torus, $h=1/3$, the pairing
$\langle u,v\rangle_h=27^{-1}\sum_x\overline u(x)v(x)$, and
\[
 \delta_i=3(S_i-S_i^{-1}),\qquad \omega=3\sqrt3.
\]
All symmetric tensor contractions use the full Frobenius pairing. The physical
field amplitude is denoted $a$, separately from every first-field parameter.
For $\{i,j,k\}=\{1,2,3\}$ with $j<k$, put
\[
 T_i=e_je_k^T+e_ke_j^T,\qquad P_i=e_je_j^T-e_ke_k^T,
 \qquad f_i(x)=\cos(2\pi x_i).
\]
For the ordered mixed frequency list
\begin{equation}
 k_1=(1,1,0),\quad k_2=(1,0,1),\quad
 k_3=(0,1,1),\quad k_4=(1,-1,0),
 \label{v39:eq:modes}
\end{equation}
let $i<j$ be the nonzero indices of $k$ and $m$ its zero index. Define
\begin{align}
 v_k&=k_je_i-k_ie_j,& w_k&=e_m,\\
 T_k&=v_kw_k^T+w_kv_k^T,& Q_k&=v_kv_k^T-2w_kw_k^T,\\
 \phi_k(x)&=\cos(2\pi k\cdot x)+3^{-1/2}\sin(2\pi k\cdot x).
 \label{v39:eq:mixed-tensors}
\end{align}
The last profile takes exactly the rational values $(1,0,-1)$ according to
$k\cdot(3x)$ modulo three. No approximate trigonometric data enter the certificate.

For real parameters $\alpha=(\alpha_1,\alpha_2,\alpha_3)$,
$d_2,d_3$, and $\xi=(\xi_1,\ldots,\xi_4)$, define
\begin{equation}
 d_1=
 \frac{\frac14\sum_i\alpha_i^2+\frac43\sum_{r=1}^4\xi_r^2-d_2d_3}
      {d_2+d_3},\qquad d_2+d_3\ne0,
 \label{v39:eq:d1}
\end{equation}
and the canonical first field $Y=(U,p)$ by
\begin{align}
 U&=\sum_i\alpha_iT_if_i+\sum_{r=1}^4\xi_rT_{k_r}\phi_{k_r},\\
 p&=\omega\left\{\operatorname{diag}(d)-(\operatorname{tr}\operatorname{diag}(d))I
       +\frac12\sum_i\alpha_iP_if_i
       +\frac12\sum_{r=1}^4\xi_rQ_{k_r}\phi_{k_r}\right\}.
 \label{v39:eq:first-field}
\end{align}
This is a nine-parameter extension of the equal-mixed-amplitude curve in the
companion. Every parameter changes initial canonical coefficients, not the action.

\begin{lemma}[Compatibility and nondegenerate mean preparation]
\label{v39:lem:means}
Every field \eqref{v39:eq:first-field} satisfies the original linear constraints
$LY=0$ and all four quadratic mean constraints. If
$\alpha_1\alpha_2\alpha_3(d_2+d_3)\ne0$, the four quadratic mean derivatives have
an invertible restriction to four explicit canonical controls. This conclusion
does not require any mixed amplitude to be small.
\end{lemma}
\begin{proof}
All inhomogeneous tensors are trace-free and annihilate their respective phase
wavevectors. The seven nonzero frequency pairs are distinct on the actual finite
frequency group. Thus the configuration double-divergence/trace row and the
momentum divergence vanish. Frequency orthogonality and $T_i:P_i=T_k:Q_k=0$
give all three momentum means $\langle p,\delta_iU\rangle_h=0$.

The mixed tensor norms and profile mean are
\[
 \|T_k\|_F^2=4,\quad \|Q_k\|_F^2=8,\quad
 |k|^2=2,\quad \overline{\phi_k^2}=2/3.
\]
In the complete original quadratic Hamiltonian the trace and divergence
configuration terms are zero on these fields. Direct orthogonal summation gives
\begin{equation}
 \frac{H_2(Y)}{\omega^2}
  =-2\sum_{i<j}d_id_j+\frac12\sum_i\alpha_i^2
                               +\frac83\sum_r\xi_r^2=0.
 \label{v39:eq:mean-identity}
\end{equation}
The equality is exactly \eqref{v39:eq:d1}. Each mixed configuration and its
momentum partner contribute $4\omega^2\xi_r^2/3$ separately. In particular the
mean balance retains both contributions.

For the scalar mean use variation of the homogeneous parameter $d_1$ before
solving \eqref{v39:eq:d1}. Its derivative is
$-2\omega^2(d_2+d_3)$. For momentum mean $i$ use the physical momentum
variation $Z_{p,i}=\alpha_iT_i\sin(2\pi x_i)$, with zero configuration variation.
Its derivative on that mean is $-\omega\alpha_i^2$ and on the other two is zero.
These controls have zero derivative on $H_2$, by the same frequency and
polarization orthogonality. The homogeneous control has zero derivative on all
momentum means. The resulting four-by-four matrix is diagonal with the stated
nonzero entries, up to the retained harmless sign of the scalar constraint.
No continuum product rule or continuous translation symmetry is used.
\end{proof}

The flat analytic range-and-mean initializer therefore applies locally about
any finite field in this family with the displayed nonzero factors: after
reducing $|a|$, $aY$ lies in the original stationary action germ. Its original
nonconstant linear right inverse and the four controls in the lemma solve the
nonlinear range rows and the divided mean rows. This use of the initializer
does not require a regular multiplier-rate Hessian; that additional issue is
kept separate below.

\section{The three actual cubic means as a finite rational map}
\label{v39:sec:map}

Retain the original finite stationary Hamiltonian $H(X,\lambda)$,
$P=D_\lambda H$, and
\[
 F=\mathbb J\nabla_XH,\qquad \mathsf B=D_XP\,F,
 \qquad M=D_\lambda P.
\]
The first Poisson bracket $\mathbb K_1(Y)$ is skew-adjoint; $M$ is symmetric.
They are not interchangeable inverses. Put $p=\sqrt3p_R$ and
\[
 Q_3=\operatorname{diag}(I_{27},\sqrt3I_{81}),\quad
 \widehat K=3^{-1/2}Q_3\mathbb K_1Q_3,\quad
 \widehat b=3^{-1/2}Q_3\widehat B_2.
\]
The 108 multiplier coordinates are component-major: lapse followed by the
three shift components. Let $\mathcal E$ be the four-dimensional space of
component constants. On $LY=0$, the companion proves
\begin{equation}
 \mathcal E\subseteq\ker\widehat K,\qquad
 \widehat b\perp\mathcal E.
 \label{v39:eq:constants}
\end{equation}
Let $\mathcal J$ omit indices $26,53,80,107$, let $E$ insert those 104 remaining
point coordinates, and let $S=(I-P_{\mathcal E})E$. Set
\begin{equation}
 A=E^T\widehat K E,\qquad b=E^T\widehat b,\qquad
 y=-A^{-1}b,\qquad x=Sy,\qquad \ell=Q_3x.
 \label{v39:eq:solve}
\end{equation}
Whenever $A$ is invertible, these formulas solve \emph{all} 108 Poisson rows
and all four means. Indeed $\widehat KS=\widehat KE$; the remaining four row
sums follow from \eqref{v39:eq:constants}. The full rank is then exactly 104.

For a constant shift $c$, the original prepared cubic mean law is
\begin{equation}
 C_c(Y)=t_c(Y)+\mathcal R_Y[c,\ell(Y)].
 \label{v39:eq:actual-C}
\end{equation}
Here, using the literal original homogeneous coefficients,
\begin{align}
 t_c(Y)&=D\Phi_3(Y)[\delta_cY]+DP_{c,3}(Y)[F_1Y],\\
 \mathcal R_Y[c,\ell]&=DP_{c,3}(Y)[G\ell]
       -DP_{\ell,2}(Y)[\delta_cY]
       -\langle\delta_c\ell,P_2(Y)\rangle.
 \label{v39:eq:mean-functional}
\end{align}
The last term is the correction imposed by exact initial constraint
preparation. It is retained in every coefficient calculation. The source
functional $P_{\ell,2}$ includes the flat auxiliary response and shifted
cubic electric term; $\Phi_3=-H_3$ and $P_{c,3}$ retain the original ordered
magnetic BCH terms. The provided compiler is the literal one retained in
Versions 35--38 and in the companion's equations \texttt{v20:P2},
\texttt{v20:Phi3}, and \texttt{v21:constant-law}.

Write $\widehat R\,x$ for the three values
$\mathcal R_Y[e_i,Q_3x]$. Define
\begin{equation}
 B=\widehat R S,\qquad
 \mathbf C=t+By=t-BA^{-1}b\in\R^3.
 \label{v39:eq:reduced-C}
\end{equation}
This is the actual three-mean map. It is not the six axial tangency functions
of Versions 36--38, and no old axial weak-null frame is imposed on it.

\paragraph{Finite coefficient specification.}
For the root certificate fix $\alpha=(1,9/8,1)$ and $d_2=d_3=-8$. In rational
momentum coordinates the first field is linear in
\begin{equation}
 z=(1,\xi_1^2+\xi_2^2+\xi_3^2+\xi_4^2,\xi_1,\xi_2,\xi_3,\xi_4).
 \label{v39:eq:z}
\end{equation}
Its six basis fields are the old axial field, the homogeneous correction
$(0,\operatorname{diag}(0,1/4,1/4))$, and the four mixed pairs
$(T_{k_r}\phi_{k_r},\frac32Q_{k_r}\phi_{k_r})$.
Thus $\widehat K$ is linear in $z$, the concatenation
$(\widehat b,\widehat R)$ is homogeneous quadratic in $z$, and $t$ is homogeneous
cubic in $z$. The exact bank consists of six $108$-square matrices, 21
quadratic coefficient vectors of length 432, and 56 cubic coefficient vectors
of length three. All are rational.

These arrays are regenerated by polynomial exhaustion, not fitted on arbitrary
samples. For a quadratic map $Q$, its diagonal coefficient is $Q(e_i)$ and
its mixed coefficient is $Q(e_i+e_j)-Q(e_i)-Q(e_j)$. For a cubic map $T$, the
$iik$ coefficient is
$[T(e_k+e_i)+T(e_k-e_i)-2T(e_k)]/2$, the $iii$ coefficient is $T(e_i)$,
and three distinct indices use inclusion--exclusion on their three unit
vectors. These exact identities specify every retained coefficient from
\eqref{v39:eq:mean-functional}. The physical grading is applied before the
reduced solve: the $\beta\beta$ Poisson block is multiplied by three,
$\widehat b_\beta=B_{2,\beta}(U,p_R)+2B_{2,\beta}(0,p_R)$,
$\widehat R_N=3R_N(U,p_R)-2R_N(U,0)$,
$\widehat R_\beta=3R_\beta(U,p_R)$, and
$t=3t(U,p_R)-2t(U,0)$.

\section{A verified simultaneous zero with every spatial profile present}
\label{v39:sec:root}

Set $\xi_4=1/5$ and use $q=(\xi_1,\xi_2,\xi_3)$ as the three root variables.
Define the exact dyadic centre
\begin{equation}
 q_0=2^{-120}
 \begin{pmatrix}
 -415139097047449945527169767258117582\\
  405914304295007683928888200118499045\\
  235008126046521569502553452453475741
 \end{pmatrix},
 \qquad r=2^{-80}.
 \label{v39:eq:centre}
\end{equation}

\begin{theorem}[An original-action cubic-mean root]
\label{v39:thm:root}
The three equations $\mathbf C(q,1/5)=0$, with the fixed axial and homogeneous
parameters above, have exactly one solution $q_*$ in
$\|q-q_0\|_\infty\le r$. Throughout that cube the full original Poisson rank
is 104, and $D_q\mathbf C$ is nonsingular. For orientation,
\begin{equation}
 \xi_*\simeq
 (-0.3123159445662353833,\ 0.3053759818347139375,
       \ 0.1768004637216116406,\ 0.2).
 \label{v39:eq:root-decimal}
\end{equation}
All three axial and all four mixed profiles are nonzero. The uniquely
normalized physical first tilt obeys
\begin{equation}
 15.80<\|\ell(Y_*)\|_h<15.82,\qquad
                  \max_{x,\mu}|\ell^\mu(Y_*)(x)|<21.
 \label{v39:eq:tilt-bounds}
\end{equation}
The homogeneous coefficient is $d_1\simeq3.9271367640846675$ and $d_2=d_3=-8$.
The assertion is about exact roots and inequalities; the decimal values do
not define or certify them.
\end{theorem}

\begin{proof}
We give the full inclusion argument and the arithmetic conditions. The
certificate supplies fixed rational matrices $P$ and $W$, a rational vector
$y_0$, and three rational vectors $y_j$. Here $P$ approximates the inverse of
$A(q_0)$, $W$ approximates the inverse of the three-mean Jacobian, and the vectors
approximate the principal-coordinate solution and its three derivatives.
No approximate inverse is treated as an exact solve. The arithmetic verifier
regenerates the original coefficient bank and then uses rational operations
only in the following inequalities.

For any finite matrix, $\|\cdot\|_\infty$ is the induced maximum-row-sum norm;
for vectors it is the maximum norm. Put
\[
 q_A=\|I-PA_0\|_\infty,\quad
 a_A=\frac{\|P\|_\infty}{1-q_A},\quad
 e_y=\frac{\|P(A_0y_0+b_0)\|_\infty}{1-q_A}.
\]
The strict inequality $q_A<1$ proves invertibility and
$\|A_0^{-1}\|_\infty\le a_A$, as well as
$\|y(q_0)-y_0\|_\infty\le e_y$. With subscript $j$ denoting an exact
parameter derivative at the centre, the differentiated equation is
\begin{equation}
 A_0\partial_jy=-b_j-A_jy.
 \label{v39:eq:linearized-solve}
\end{equation}
Consequently
\[
 e_j=a_A\{\|A_0y_j+b_j+A_jy_0\|_\infty+
                         \|A_j\|_\infty e_y\}
\]
bounds $\|\partial_jy(q_0)-y_j\|_\infty$. Form the three columns
$J_{0,j}=t_j+B_jy_0+B_0y_j$. The exact centre bounds are
\begin{align}
 \eta&=\|W(t_0+B_0y_0)\|_\infty+\|WB_0\|_\infty e_y,\\
 q_C&=\|I-WJ_0\|_\infty+
 \|W\|_\infty\sum_{j=1}^3
                  (\|B_j\|_\infty e_y+\|B_0\|_\infty e_j).
 \label{v39:eq:centre-enclosures}
\end{align}
They bound respectively $\|W\mathbf C(q_0)\|_\infty$ and
$\|I-WD\mathbf C(q_0)\|_\infty$.

It remains to bound the whole cube, rather than only its centre. The cube is
inside $|q_i|<1$. In \eqref{v39:eq:z}, each coordinate has magnitude at most
four, each first derivative at most two, and each second derivative at most
two. For a degree-two monomial the corresponding value, first-derivative and
second-derivative bounds are $16,16,24$; for a degree-three monomial its
second-derivative bound is 192. Multiplying these constants by the sums of
exact coefficient norms gives bounds $A_1,A_2$ for any one first or second
derivative of $A$, $b_0^+,b_1^+,b_2^+$ for $b$, $B_0^+,B_1^+,B_2^+$ for
$B$, and $t_2^+$ for a second derivative of $t$ on $|q_i|<1$.

In particular $\|A(q)-A_0\|_\infty\le3A_1r$. The certified inequality
$3a_AA_1r<1$ gives
\[
 a_B=\frac{a_A}{1-3a_AA_1r},\quad
 u_0=a_Bb_0^+,\quad
 u_1=a_B(b_1^++A_1u_0),\quad
 u_2=a_B(b_2^++A_2u_0+2A_1u_1).
\]
Differentiating $Ay=-b$ proves that these bound $y$, each first derivative,
and each second derivative on the cube. Hence
\begin{equation}
 H_C=t_2^++B_2^+u_0+2B_1^+u_1+B_0^+u_2
 \label{v39:eq:hessian-bound}
\end{equation}
bounds the maximum norm of each $\partial_j\partial_k\mathbf C$. The row-sum
Jacobian variation is at most $9H_Cr$, so a valid whole-cube contraction bound is
$\kappa=q_C+9\|W\|_\infty H_Cr$.

The independently checked rational comparisons give the following deliberately
rounded-up dyadic bounds:
\begin{center}
\begin{tabular}{@{}lr@{}}
\toprule
Quantity & Certified upper bound\\\midrule
$\|I-PA_0\|_\infty$ & $2^{-37}$\\
$\|A_0^{-1}\|_\infty$ & $1$\\
$\|I-WD\mathbf C(q_0)\|_\infty$ & $2^{-100}$\\
$H_C$ & $2^{75}$\\
$\sup_{\|q-q_0\|_\infty\le r}\|I-WD\mathbf C(q)\|_\infty$ & $2^{-18}$\\
$\|W\mathbf C(q_0)\|_\infty$ & $2^{-115}$\\
\bottomrule
\end{tabular}
\end{center}
All statements concern the actual rational coefficient bank, not an interpolated
floating map. In particular
\[
 \eta+\kappa r\le2^{-115}+2^{-98}<2^{-80}=r.
\]
The map $q\mapsto q-W\mathbf C(q)$ is therefore a strict contraction of the
closed cube into itself. Its unique fixed point is a zero: the bound on
$I-WD\mathbf C$ makes both square factors nonsingular. Conversely every zero
in the cube is its fixed point. This also proves nonsingularity of the
Jacobian throughout the cube. The Neumann argument for $A$ and
\eqref{v39:eq:constants} prove the full rank statement.

Finally $x(q_*)=Sy(q_*)$ lies within
$\|S\|_\infty(e_y+3u_1r)$ of $Sy_0$. Exact comparisons of the weighted sum
$(\|x_N\|^2+3\|x_\beta\|^2)/27$ with $15.80^2$ and $15.82^2$, and of every
weighted component square with $21^2$, give \eqref{v39:eq:tilt-bounds}. The
profile nonvanishing and the stated mean chart factors are separated from zero
by their exact centre and radius. Standard verified nonlinear-inclusion
methodology is background \cite{v39:Rump}; the particular source bank and
rational inequalities above are the certificate used here.
\end{proof}

\section{A six-parameter family of simultaneous cubic cancellations}
\label{v39:sec:family-zero}

The numerical root is not used as a surrogate for a root curve. The verified
Jacobian supplies the latter by the ordinary analytic implicit-function theorem.
Let
\[
 \eta=(\alpha_1,\alpha_2,\alpha_3,d_2,d_3,\xi_4),\qquad
 \eta_*=(1,9/8,1,-8,-8,1/5).
\]

\begin{corollary}[Persistent mixed-mode coefficient zeros]
\label{v39:cor:six}
There is an open neighborhood of $\eta_*$ and a unique nearby real-analytic
map $q=q_*(\eta)$ for which all three original cubic means vanish. Along this
six-parameter family the Poisson rank is 104, the unique normalized first tilt
is analytic, all seven spatial profiles remain present, and the four quadratic
mean-preparation derivatives remain independent.
\end{corollary}
\begin{proof}
The first fields and their source coefficients are analytic wherever
$d_2+d_3\ne0$. Invertibility of $A$ persists near the certified root, so the
normalized tilt and $\mathbf C$ are analytic there. Theorem
\ref{v39:thm:root} proves invertibility of $D_q\mathbf C$ at the zero. Apply
the analytic implicit-function theorem and shrink the neighborhood to preserve
all finitely many nonzero margins. Lemma \ref{v39:lem:means} gives the mean
control assertion.
\end{proof}

One of these six parameters can represent a common first-field scaling; no
claim of six distinct scale-free physical moduli is needed. This is a
finite-dimensional family of selected initial coefficients, not an open basin
of arbitrary microscopic dynamics, and not a cofinal family in spatial cutoff.

\section{Exact original constraints and zero complete initial drift}
\label{v39:sec:exact-data}

The verified zero permits more than removal of the cubic coefficient. We now
use the original analytic initializer and the next initial multiplier values;
no update is performed during a history.

\begin{lemma}[Analytic family with its true cubic control]
\label{v39:lem:analytic-prep}
For $(q,\eta)$ close to the verified family and a fixed finite normalized
multiplier $m\in\mathcal E^\perp$, the original action has jointly analytic
initial records, for sufficiently small $a$, with
\begin{align}
 \lambda(a,q,\eta,m)&=e+a\ell(Y(q,\eta))+a^2m,\\
 P(X(a,q,\eta,m),\lambda)&=0,\\
 X(a,q,\eta,m)&=aY(q,\eta)+O(a^2),\\
 \mathsf B(a,q,\eta,m)&=a^3\{B_3(q,\eta,0)+\mathbb K_1(Y)m\}+O(a^4).
 \label{v39:eq:analytic-prep}
\end{align}
The three constant-shift components of the braces are exactly $\mathbf C(q,\eta)$,
and the constant-lapse component is zero. The construction is valid locally
about every specified finite value of $m$; uniformity as $\|m\|\to\infty$ is not
asserted.
\end{lemma}
\begin{proof}
Use the original analytic stationary germ at the flat field, its nonconstant
linear right inverse, and the four independent mean controls in Lemma
\ref{v39:lem:means}. The range equation starts at second amplitude order.
After it is solved, divide the four mean equations by $a^2$ and use their
analytic Taylor extension through $a=0$. Their derivative in the four mean
controls is the invertible matrix already proved. The analytic implicit-function
theorem gives exact constraint zeros, locally in all the displayed finite
parameters, with unchanged first coefficient. This argument is at a fixed
finite field $Y$; reducing $a$ places it inside the same original analytic
germ regardless of the size of that fixed coefficient.

The exact first-tilt law in the companion cancels every quadratic drift row,
so division of the drift by $a^3$ is analytic. For the second value $a^2m$, use
the original identity
\[
 D_\lambda\mathsf B[m]=\mathbb K m+D_X(Mm)[F],
 \qquad \mathbb K=a\mathbb K_1(Y)+O(a^2),\quad M=O(a^2).
\]
Integrating that value change gives $a^3\mathbb K_1(Y)m+O(a^4)$ at fixed fields.
Exact re-preparation can change a third field coefficient, since its four-mean
constraint discrepancy begins at order $a^4$. That coefficient change belongs
to $\ker L$. The linear drift factors through $L$ on lapse tests and is zero
on shift tests, so this change contributes nothing to the cubic drift. This
is precisely the retained mean-order bookkeeping, not a claim that no third
field correction occurs. The companion's full constant-shift formula gives
\eqref{v39:eq:actual-C}, independently of $m$ and of higher preparation.
Finally the exact homogeneity identity $\langle\lambda,\mathsf B\rangle=0$
and $\mathsf B=O(a^3)$ make the constant-lapse cubic coefficient zero.
\end{proof}

\begin{theorem}[Nonempty exact initial constraint-drift zero family]
\label{v39:thm:exact-zero}
For $\eta$ near $\eta_*$ there are analytic original initial records, for all
sufficiently small $a$, such that
\begin{equation}
 \boxed{P(X(a,\eta),\lambda(a,\eta))=0,\qquad
        \mathsf B(X(a,\eta),\lambda(a,\eta))=0.}
 \label{v39:eq:exact-zero}
\end{equation}
They have
\begin{align}
 X(a,\eta)&=aY(q_*(\eta),\eta)+O(a^2),\\
 \lambda(a,\eta)&=e+a\ell(Y(q_*(\eta),\eta))+O(a^2),\\
 \overline{\lambda^N}&=1,\qquad\overline{\lambda^\beta}=0.
 \label{v39:eq:initial-jets}
\end{align}
The lapse and spatial metric are positive for sufficiently small amplitude.
No invertibility of the symmetric multiplier-rate Hessian is used in this
initial-data theorem.
\end{theorem}
\begin{proof}
Take three constant-shift drift means followed by all 104 nonconstant drift
coordinates, and divide these 107 rows in \eqref{v39:eq:analytic-prep} by $a^3$.
Denote their analytic extension by $\mathcal F(a,q,\eta,m)$. At $a=0$ its first
three entries are $\mathbf C(q,\eta)$, independent of $m$. For the remaining
entries its derivative in the normalized second multiplier is the invertible
restriction of $\mathbb K_1(Y)$ to $\mathcal E^\perp$.

Fix first $\eta=\eta_*$. At the proved zero $q_*$, the actual finite vector
$B_3^\perp(q_*,\eta_*,0)$ has a unique cancelling second multiplier
\[
 m_*=-\bigl(\mathbb K_1(Y_*)|_{\mathcal E^\perp}\bigr)^{-1}
                                  B_3^\perp(q_*,\eta_*,0).
\]
No value of this unlisted higher baseline is assumed to vanish. It is the
actual coefficient of an analytic finite function and is therefore finite.
Choose the initializer in a neighborhood of this fixed $m_*$. At
$(0,q_*,\eta_*,m_*)$ all 107 rows of $\mathcal F$ vanish. Its derivative in
$(q,m)$ has block form
\begin{equation}
 \begin{pmatrix}
 D_q\mathbf C(q_*,\eta_*)&0\\
 *&\mathbb K_1(Y_*)|_{\mathcal E^\perp}
 \end{pmatrix}.
 \label{v39:eq:joint-jacobian}
\end{equation}
Both diagonal blocks are invertible. The analytic implicit-function theorem,
with $a,\eta$ as external parameters, supplies $q(a,\eta),m(a,\eta)$ making
these 107 original rows identically zero. At zero amplitude
$q(0,\eta)=q_*(\eta)$ by local uniqueness. The exact constraints have already
been solved throughout the initializer; no post-hoc projection is added.

The only remaining drift component could be a constant lapse covector.
The exact identity $\langle\lambda,\mathsf B\rangle=0$ and
$\overline{\lambda^N}=1$ force that last component to be zero as well. This
proves all 108 drift equations. Analyticity gives \eqref{v39:eq:initial-jets},
and the positive flat lapse and metric persist after reducing the amplitude.
The inverse in \eqref{v39:eq:joint-jacobian} is the first Poisson bracket for
\emph{initial-value controls}, not the symmetric Hessian governing time rates.
\end{proof}

The canonical field vector is not made static by this construction. At the
central first field its configuration component has
\begin{equation}
 \lim_{a\to0}a^{-1}\overline{F_U(X(a),\lambda(a))}
       =2\omega\operatorname{diag}(d_{1,*},-8,-8)\ne0.
 \label{v39:eq:nonstatic}
\end{equation}
Indeed $F_U/a\to2p-(\operatorname{tr}p)I+
\operatorname{SymGrad}_\delta\ell^\beta$; the last term has zero spatial mean,
and the first two have mean $2\omega\operatorname{diag}(d)$. The theorem thus
selects nontrivial, instantaneously constraint-preserving canonical data, not
a constant zero-field state.

\begin{corollary}[What follows on a separately regular normalized continuation]
\label{v39:cor:rate}
If the full original multiplier Hessian on these prepared records has only
its radial null direction and is invertible on the normalized complement,
then their own original normalized continuation has $\lambda_t(0)=0$.
Its initial canonical second derivative is $O(a)$, and the original metric
and literal link curvatures are $O(a)$ at this fixed cutoff.
\end{corollary}
\begin{proof}
The original consistency equation is $M\lambda_t=-\mathsf B$. Equation
\eqref{v39:eq:exact-zero} and the separately stated invertibility give zero
normalized rate. The analytic original canonical equation then gives
$X_t=O(a)$ and $X_{tt}=D_XF\,F+D_\lambda F\,\lambda_t=O(a)$.
The original ADM and stationary link reconstructions use these initial jets
and the first multiplier time jet, so their fixed-cutoff curvature estimates
follow from the same analytic flat expansions retained in the companion.
\end{proof}

The extra hypothesis of this corollary has \emph{not} been verified at the
new isolating box. The companion proves signed regularity on an unspecified
small neighborhood of its old seed; the present certificate does not show
that this finite-distance root lies there. The unconditional result of this
body is \eqref{v39:eq:exact-zero}, not the existence or lifespan of a fully
regular nonlinear history.

\section{Verification boundary and the next dynamical test}
\label{v39:sec:status}

The new computation does not repeat the companion's generic Poisson-rank
argument as a new theorem. It evaluates the remaining actual cubic means,
certifies a simultaneous zero and its Jacobian, and combines that zero with
the existing initial-value control mechanism to remove the complete initial
drift exactly. The four-mean derivative does not degenerate, and no axial or
mixed profile is sent to zero.

\paragraph{Executed verification.}
The package contains the full rational coefficient bank, its regeneration from
the unchanged $P_2,\Phi_3,P_{c,3}$ compiler, the dyadic centre and preconditioners,
and an exact-arithmetic inclusion checker. The compiler regenerates all six
Poisson matrices, all 21 quadratic coefficients of the prepared source and
cubic control map, and all 56 cubic target coefficients. The checker verifies
skewness, all four constant null directions and source means at coefficient
level, the complete mean-compatible field formula, all strict Neumann and
contraction inequalities, and the physical tilt bounds. Thus floating searches
only propose the centre; the accepted inclusion uses rational operations.
The independent arithmetic reconstruction is compared entry by entry, not
through floating singular-value tolerances. The portable checker can rerun the
coefficient regeneration as a longer mode, separately from checking the
included certificate. No proof-assistant formalization is asserted.

\paragraph{Proved and inherited.}
The independent-amplitude mean identity and nondegenerate mean controls, the
certified three-mean zero, the local six-parameter coefficient family, and the
exact initial constraint-drift zero family are established above. The analytic
flat stationary germ, original initializer framework, first-tilt law, prepared
cubic mean functional, homogeneity identity and second-value control law are
inherited at their stated finite-cutoff scope. The higher baseline defining
$m_*$ and the analytic nonlinear initial roots are not reported as numerically
tabulated; their existence follows from the proved onto blocks and the original
analytic functions, not from an assumed zero baseline.

\paragraph{Unresolved.}
A zero initial drift need not remain zero under the original action. On a
regular continuation the next equation is
\[
 M\lambda_{tt}=-D_X\mathsf B\,F
 \quad\hbox{at this prepared cut, since }\lambda_t=0.
\]
Its signed inverse, higher source coefficients and geometric observation must
be retained. Neither first Poisson rank nor a zero cubic mean bounds that
response. A fresh full signed-source certificate at this root, followed by
this actual differentiated drift, is the next nonduplicating dynamical task.
No repeated re-initialization, multiplier reset, alteration of the action or
renormalization of physical time is supplied here. Uniform spatial refinement,
a common positive physical lifespan, matter/chiral source identification and
the full Einstein--Standard-Model continuum remain outside the result. The
unit-initial-multiplier stress problem of Version 32 remains a separate branch.

Every Version 38 byte before its final document command is preserved. New
labels use the prefix \texttt{v39:}. The source hashes and the actually executed
arithmetic and compilation stages are recorded in the accompanying provenance.

\begingroup
\renewcommand{\refname}{Sources retained in Version 39}
\begin{thebibliography}{9}\small\raggedright
\bibitem{v39:EB28}
\emph{Renewal Exact-Action Einstein Bridge Research Notes}, supplied cumulative
V1--V28, file \nolinkurl{Renewal_Exact_Action_Einstein_Bridge_Research_Notes_V28(1).tex}.
In particular \texttt{v28:compatibility}, \texttt{v28:constants},
\texttt{v28:rank}, \texttt{v28:generic-theorem}, \texttt{v28:quadratic},
\texttt{v28:cubic-theorem} and their retained V20--V21 coefficient identities.
\bibitem{v39:SS35}
\emph{Renewal SM Source Selection Research Notes}, supplied V35,
file \nolinkurl{Renewal_SM_Source_Selection_Research_Notes_V35(2).tex},
stationary reverse reconstruction and its source-identification boundaries.
\bibitem{v39:PE23}
\emph{Renewal Perturbative ESM Continuum Research Notes}, supplied V23,
file \nolinkurl{Renewal_Perturbative_ESM_Continuum_Research_Notes_V23(2).tex},
derivative floor, native descendant transport and the retained first-background equation.
\bibitem{v39:YM32}
\emph{Renewal Yang--Mills Mass Gap Research Notes}, supplied restart V32,
file \nolinkurl{Renewal_Yang_Mills_Mass_Gap_Research_Notes_V32(1).tex},
\texttt{r32:derivatives}, \texttt{r32:native}, \texttt{r32:kernels} and their
explicit distinction from global retained response.
\bibitem{v39:Rump}
S.~M.~Rump, \emph{Solving nonlinear systems with least significant bit accuracy},
Computing \textbf{29} (1982), 183--200; institutional record
\url{https://doi.org/10.15480/882.305}. Verified inclusion is established methodology;
the source coefficients and rational certificate above are specific to this calculation.
\end{thebibliography}
\endgroup
\addtocontents{toc}{\protect\endgroup}
% END IMMUTABLE VERSION 39 BODY
% APPEND VERSION 40 HERE, BEFORE THE FINAL END-DOCUMENT COMMAND.



% BEGIN IMMUTABLE VERSION 40 BODY
\clearpage
\begin{center}
{\Large\bfseries Version 40: Differentiated drift and a transverse initial-rate correction}\\[.45em]
{\large An exact source test at the mixed-mode root, and compatible initial records beyond zero rate}
\end{center}
\hypersetup{pdftitle={Native Regenerative Stationarity Closure: cumulative V1--V40}}
\makeatletter
\addtocontents{toc}{\protect\begingroup}
\addtocontents{toc}{\protect\makeatletter}
\addtocontents{toc}{\protect\def\protect\l@section{\protect\bfseries\protect\@dottedtocline{1}{0em}{2.5em}}}
\makeatother
\addcontentsline{toc}{part}{Version 40: Differentiated drift and a transverse initial-rate correction}

\section{Recovery, inherited inputs, and the dynamical question}
\label{v40:sec:context}

The complete Version 39 source, embedded original-coefficient bank, isolating
cube and rational verification program were recovered. The exact inclusion and
physical-tilt enclosures were rerun successfully. No later unfinished mathematical
body was found. This continuation preserves Version 39 and starts from its
verified central mixed-mode first field $Y_*$ and normalized first tilt $\ell_*$.
The four companion files retain the scope recorded in \cref{v39:sec:audit};
no additional source-selection, perturbative or Yang--Mills hypothesis is inserted.

The initial identities $P=0$ and $\mathsf B=0$ in
\cref{v39:thm:exact-zero} do not imply that their zero set is invariant. Before
attempting another full signed inverse, we compute a necessary differentiated
source at that very root. It is nonzero. There is, however, a constructive
alternative to repeating a zero-rate initialization: a specifically determined,
mean-zero transverse initial rate cancels the whole leading weak differentiated
source. A new analytic selection of initial records realizes its exact first
consistency equation without changing the action or prescribing a rate during
an evolution.

All results below concern the original cutoff $h=1/3$. The full signed
multiplier-Hessian certificate remains unresolved; the new positive initial-data
theorem does not use it. Claims about an actual unique continuation are stated
separately and retain this hypothesis. The general distinction between consistent
multiplier selection and dynamical preservation is standard in constrained
discretization \cite{v40:DBGP}; the numerical constants and original-source
correction established here are not imported from that literature.

We use the original stationary Hamiltonian and definitions
\begin{equation}
 P=D_\lambda H,\qquad F=\mathbb J\nabla_XH,\qquad
 \mathsf B=D_XP\,F,\qquad M=D_\lambda P,
 \qquad \mathbb K=D_XP\,\mathbb J D_XP^* .
 \label{v40:eq:objects}
\end{equation}
The first Poisson coefficient $\mathbb K_1$ is not $M$. The inherited
homogeneity identities, flat germ and grading are
\begin{gather}
 M\lambda=0,\qquad \langle\lambda,\mathsf B\rangle_h=0,\qquad
 D_\lambda\mathsf B[v]=\mathbb K v+D_X(Mv)[F],\label{v40:eq:identities}\\
 F(aY+O(a^2),e+a\ell+O(a^2))
       =a(F_1Y+G\ell)+O(a^2),\notag\\
 M=O(a^2),\qquad D_XM=O(a),\qquad D_\lambda M=O(a^2).
 \label{v40:eq:flat-orders}
\end{gather}
Here $e=(1,0)$ and $G=\mathbb J L^*$, with $L=D_XP(0,e)$.
These are identities or Taylor bounds of the original finite analytic germ,
retained from the Einstein companion \cite{v40:EB28}, not assumptions that its
inverse is bounded. Constants may depend on this cutoff and on a fixed bounded
parameter neighborhood.

\section{The weak differentiated source is determined by the first records}
\label{v40:sec:derivative}

Let $\mathcal W$ be the full curl-free shift space, including the three
constants, and $\mathcal A$ the mean-zero lapse and mean-zero transverse-shift
space. On the mean-lapse-zero multiplier complement, these give the fixed
direct sum $\mathcal A\oplus\mathcal W$. Write $d_{2,W}$ for the original
quadratic weak drift. For fixed first records $Y,\ell$, put
\begin{equation}
 V=F_1Y+G\ell,\qquad
 \mathfrak t_{Y,\ell}[\nu]
   =Dd_{2,W}(Y)[V,\nu]+\mathbb K_1(V)[\nu,\ell],
 \quad \nu\in\mathcal W.
 \label{v40:eq:t}
\end{equation}
This definition is valid for every weak test; it does not require that the test
be Poisson-null. That distinction matters on the maximal-rank mixed family,
where the old nonconstant Poisson-null tests no longer have that property.

\begin{lemma}[Differentiated weak drift at a prescribed initial tangent]
\label{v40:lem:leading}
Suppose $P(X_a,\lambda_a)=0$, $X_a=aY+O(a^2)$,
$\lambda_a=e+a\ell+O(a^2)$, and $LY=0$. At the initial tangent
$(X_t,\lambda_t)=(F(X_a,\lambda_a),av)$, for a fixed finite multiplier $v$,
\begin{equation}
 (D_X\mathsf B\,F+D_\lambda\mathsf B[av])\big|_{\mathcal W}
   =a^2\{\mathfrak t_{Y,\ell}+
                     \mathbb K_1(Y)[\,\cdot\,,v]\big|_{\mathcal W}\}
       +O(a^3).
 \label{v40:eq:source-leading}
\end{equation}
Moreover $LV=0$, and $\mathfrak t_{Y,\ell}$ annihilates the three constant
shifts. None of the coefficients displayed in \eqref{v40:eq:source-leading}
depends on a second or higher initial preparation coefficient.
\end{lemma}
\begin{proof}
The original weak linear drift is identically zero. The joint Taylor formula is
\[
 \mathsf B_W(X,e+\zeta)=d_{2,W}(X)
                  +\mathbb K_1(X)[\,\cdot\,,\zeta]\big|_{\mathcal W}
                  +O((\|X\|+\|\zeta\|)^3).
\]
Its $X$ derivative at the stated records is
$a\{Dd_{2,W}(Y)+\mathbb K_1(\,\cdot\,)[\,\cdot\,,\ell]\}+O(a^2)$.
Its multiplier derivative is $a\mathbb K_1(Y)|_W+O(a^2)$.
Insert $F=aV+O(a^2)$ and $av$. These expansions prove
\eqref{v40:eq:source-leading} with all derivatives of the remainder controlled
on a fixed finite analytic neighborhood.

The exact flat identities $LF_1Y=0$ for $LY=0$ and $LG=0$ give $LV=0$.
For a constant shift $c$, the original $d_{2,W}(X)[c]$ is identically zero.
The original translation incidence gives
$\mathbb K_1(V)[c,\ell]=\langle\delta_c\ell,LV\rangle_h=0$.
This proves the mean assertion without discarding any constant row.
\end{proof}

For verification we retain the rational grading $p=\sqrt3p_R$,
$\ell=Q_3x$ and $Q_3=\operatorname{diag}(I_{27},\sqrt3I_{81})$.
With $x=(n,\beta_R)$, define
\[
 v_U=2p_R-(\operatorname{tr}p_R)I+
             \operatorname{SymGrad}_\delta\beta_R,\qquad
 v_p=\mathcal F(U)+\nabla_\delta^2n-I\Delta_\delta n.
\]
Thus $V=(\sqrt3v_U,v_p)$. For a pure shift test, the exact rational formula is
\begin{equation}
 \frac{\mathfrak t_{Y,\ell}[\nu]}{\sqrt3}
 =Dd_{2,W}(U,p_R)[(v_U,v_p),\nu]
                   +\mathbb K_1(v_U,v_p)[\nu,(n,\beta_R)].
 \label{v40:eq:rational-t}
\end{equation}
The right side uses the normalized mean pairing. It follows from the original
$U^2,p^2$ grading of the weak drift and the linear grading of the Poisson blocks.
In particular the $G\ell$ terms in both occurrences of $V$ remain present.

\section{An exact nonzero derivative at the recovered root}
\label{v40:sec:defect}

Take the actual spatial scalar and weak test
\begin{equation}
 \phi_0(x)=\cos(2\pi x_1),\qquad
 \nu_0=(0,\nabla_\delta\phi_0),\qquad
                    \|\nu_0\|_h^2=\frac{27}{2}.
 \label{v40:eq:test}
\end{equation}
The cosine has the rational values $(1,-1/2,-1/2)$ on the retained lattice.

\begin{theorem}[Noninvariance detected by an original-source coefficient]
\label{v40:thm:defect}
At the exact central root of \cref{v39:thm:root},
\begin{equation}
 \boxed{-6671<\frac{\mathfrak t_{Y_*,\ell_*}[\nu_0]}{\sqrt3}<-6670.}
 \label{v40:eq:defect-interval}
\end{equation}
For every analytic initial family of \cref{v39:thm:exact-zero}, evaluated at
these central first records, the frozen-multiplier canonical tangent obeys
\begin{equation}
 D_X\mathsf B(X_a,\lambda_a)[F(X_a,\lambda_a)][\nu_0]
       =a^2\mathfrak t_{Y_*,\ell_*}[\nu_0]+O(a^3).
 \label{v40:eq:initial-derivative}
\end{equation}
Consequently that tangent is not tangent to the exact initial set
$\{P=0,\mathsf B=0\}$ for all sufficiently small nonzero $a$.
\end{theorem}
\begin{proof}
We first describe the exact-arithmetic enclosure of the new coefficient.
The Version 39 isolating cube has radius $r=2^{-80}$ in the first three mixed
amplitudes, with the fourth equal to $1/5$. Its retained rational proposal
supplies an approximate normalized Poisson solution $\widetilde x=S\widetilde y$.
The original coefficient bounds in its inverse certificate give, on that whole
cube,
\begin{equation}
 \|x-\widetilde x\|_\infty<2^{-43},\qquad
 \|U-U_0\|_{\max}<2^{-78},\qquad
 \|p_R-p_{R,0}\|_{\max}<2^{-76}.
 \label{v40:eq:input-radii}
\end{equation}
These are recomputed rational bounds, not rounded decimal uncertainties.
For example, if $A_0,b_0$ are the centre Poisson system and $P_0$ its retained
rational preconditioner, then
\[
 e_y=\frac{\|P_0(A_0\widetilde y+b_0)\|_\infty}
                   {1-\|I-P_0A_0\|_\infty},\qquad
 e_x\le\|S\|_\infty(e_y+3r\sup\|D y\|_\infty).
\]
The derivative bound comes from the differentiated actual solve
$A\partial_j y=-\partial_jb-(\partial_jA)y$. The $U,p_R$ bounds follow
directly from the six basis fields and the variation of $\sum_r\xi_r^2$.

The new evaluator computes \eqref{v40:eq:rational-t} directly from the retained
$P_{\nu,2}$ and $\Phi_3$ functional, not from a fitted tangent. The identities
\[
 DP_{\nu,2}(Y)[Z]=\tfrac12\{P_{\nu,2}(Y+Z)-P_{\nu,2}(Y-Z)\},
\]
\[
 D\Phi_3(Y)[Z]=\tfrac12\{\Phi_3(Y+Z)-\Phi_3(Y-Z)\}-\Phi_3(Z)
\]
are exact homogeneous polarizations. The second polarization of the quadratic
weak drift is exact as well. Every spatial shift and phase derivative is taken
before evaluation.

For reproducibility, each array is stored as an exact rational centre and a
single nonnegative rational radius controlling the largest entry error.
Addition adds radii. For an entrywise product the new radius is
$\|A\|_{\max}e_B+\|B\|_{\max}e_A+e_Ae_B$; a contraction over $m$ indices
multiplies this bound by $m$. A cross product uses $m=2$.
Permutations preserve the radius, and sums use their literal number of terms.
These rules prove enclosure by induction over the complete arithmetic graph.
No floating operation is used in an accepted inequality. On the actual cube
the resulting rational interval is contained in $(-6671,-6670)$; its centre
is approximately $-6670.546546052$ and its radius is less than $0.00135$.
The exact numerator and denominator data are included in the certificate.
This proves \eqref{v40:eq:defect-interval}. Finally apply
\cref{v40:lem:leading} with $v=0$. Since $\mathsf B=0$ initially, its
nonzero derivative proves noninvariance.
\end{proof}

The unchanged $P=0,\mathsf B=0$ initial-data theorem remains valid. Nor is this
a nonexistence theorem for evolution: an original continuation need not keep
$\mathsf B=0$ or $\lambda_t=0$ at later times.

\begin{corollary}[A necessary graded acceleration cost for zero initial rate]
\label{v40:cor:necessary}
Suppose a twice differentiable normalized original continuation of these data
exists with $\lambda_t(0)=0$. Then
\begin{equation}
 a\|\lambda_{tt,A}(0)\|_h+a^2\|\lambda_{tt,W}(0)\|_h\ge c>0
 \label{v40:eq:acc-cost}
\end{equation}
for sufficiently small positive $a$. No inverse of the weak signed block is
needed for this necessary inequality.
\end{corollary}
\begin{proof}
Differentiate $M\lambda_t=-\mathsf B$ at the initial cut. Since
$\lambda_t(0)=0$, the equation is $M\lambda_{tt}=-D_X\mathsf B F$.
The inherited source grading gives weak--active and weak--weak blocks of
orders $a^3$ and $a^4$. Pair the equation with \eqref{v40:eq:test}, use
\eqref{v40:eq:initial-derivative}, and divide by $a^2$. Fixed-dimensional
norm equivalence gives \eqref{v40:eq:acc-cost}. This argument asserts a
necessary cost if a continuation with that initial rate exists, not its existence.
\end{proof}

\section{Twenty-six transverse rate directions have a verified right inverse}
\label{v40:sec:control}

Let $E_W$ be the $108$-by-$26$ matrix with columns
$(0,\nabla_\delta(e_j-e_{26}))$, $0\le j<26$. It is the full nonconstant
weak space, not a selected set of rows. Let $E_A$ consist of twenty-six lapse
columns $e_j-e_{26}$ followed by the eighty-one pure-shift columns
$\operatorname{curl}_\delta e_{i,x}$, in component-major order. Select the
following zero-based columns:
\begin{equation}
 \begin{split}
 \mathcal J_T={}&(59,103,69,91,99,36,32,41,56,63,101,27,76,\\
                &50,87,46,29,83,66,79,90,64,47,97,106,54),\\
 Z_T={}&(E_A)_{:,\mathcal J_T},\qquad
 \mathcal H(Y)=E_W^T\widehat K(Y)Z_T.
 \end{split}
 \label{v40:eq:H}
\end{equation}
Every selected column is a pure transverse shift with zero spatial mean; none
uses a lapse column. Matrix entries in \eqref{v40:eq:H} use spatial sums,
consistent with the retained rational $\widehat K$.

\begin{theorem}[Active control at every point of the root cube]
\label{v40:thm:control}
Throughout the Version 39 isolating cube, $\mathcal H$ is invertible and
\begin{equation}
 \boxed{\|\mathcal H^{-1}\|_\infty<\frac1{32}.}
 \label{v40:eq:H-inverse}
\end{equation}
In particular the original transverse-rate map onto the twenty-six nonconstant
weak covectors has full row rank at the exact root.
\end{theorem}
\begin{proof}
The selected columns are built from the literal integer curl and gradient
matrices. Their divergence, all four means and lapse entries are checked to be
zero exactly. Write the six inherited rational Poisson matrices as $K_j$ and
$h_j=E_W^TK_jZ_T$. Their coefficient vector is
$z=(1,\sum_r\xi_r^2,\xi_1,\ldots,\xi_4)$, as in
\cref{v39:eq:z}. Thus $\mathcal H=\sum_jz_jh_j$.
If $\mathcal H_0$ is its value at the rational centre, then on the whole cube
\[
 \|\mathcal H-\mathcal H_0\|_\infty
 \le \|h_1\|_\infty\sum_{j=1}^3(2|\xi_{0,j}|r+r^2)
                    +r\sum_{j=1}^3\|h_{j+1}\|_\infty=:\delta_H.
\]
The supplied $26$-by-$26$ dyadic preconditioner $P_T$, with denominator
$2^{48}$, is checked by exact multiplication to satisfy
\begin{equation}
 q_T:=\|I-P_T\mathcal H_0\|_\infty+\|P_T\|_\infty\delta_H<2^{-25},
 \qquad \frac{\|P_T\|_\infty}{1-q_T}<\frac1{32}.
 \label{v40:eq:H-neumann}
\end{equation}
The Neumann series proves invertibility and \eqref{v40:eq:H-inverse} on every
point of the cube, including its exact root. This is a real enclosure, not a
rank inference from a determinant modulo a prime. The new multiplication uses
the recovered original-source coefficient bank whose inclusion was rerun;
the entire older bank is not represented as newly regenerated in this iteration.
\end{proof}

\section{An explicit initial-rate coefficient cancels the leading weak source}
\label{v40:sec:rate}

For the central first records define the rationally graded covector
\begin{equation}
 t_j(Y,\ell)=\frac{27}{\sqrt3}\mathfrak t_{Y,\ell}[(E_W)_j],\qquad
 w_*=-\mathcal H(Y_*)^{-1}t(Y_*,\ell_*),\qquad
 v_*:=\sqrt3Q_3Z_Tw_*=3Z_Tw_*.
 \label{v40:eq:v}
\end{equation}
The factor $27$ restores the spatial-sum convention of $\mathcal H$.
The last equality holds because $Z_T$ has no lapse component. These formulas
define $v_*$ at the exact algebraic root, not merely at its printed decimal
approximation.

\begin{theorem}[Complete leading weak cancellation by a transverse rate]
\label{v40:thm:rate}
The coefficient $v_*$ has zero lapse, transverse mean-zero shift, and satisfies
\begin{equation}
 \boxed{\mathfrak t_{Y_*,\ell_*}[\nu]
                  +\mathbb K_1(Y_*)[\nu,v_*]=0
                  \quad\hbox{for every }\nu\in\mathcal W.}
 \label{v40:eq:cancel}
\end{equation}
At the central root its physical size is enclosed by
\begin{equation}
 \boxed{1785<\|v_*\|_h<1787,\qquad
                        \max_{x,\mu}|v_*^\mu(x)|<2500.}
 \label{v40:eq:v-bounds}
\end{equation}
The same construction is analytic along a sufficiently small part of the
six-parameter coefficient-root family of \cref{v39:cor:six}.
\end{theorem}
\begin{proof}
The algebraic change of grading is
$\widehat K=Q_3\mathbb K_1Q_3/\sqrt3$. On pure shift tests,
$\mathbb K_1[\nu,\sqrt3Q_3z]/\sqrt3=\nu^T\widehat Kz$ in the sum
convention. Thus \eqref{v40:eq:v} solves \eqref{v40:eq:cancel} on all
$26$ columns of $E_W$. The three constant tests have zero
$\mathfrak t$ by \cref{v40:lem:leading} and belong to the kernel of
$\mathbb K_1(Y_*)$, so their equations hold too. The geometric properties
of $v_*$ follow from those of $Z_T$.

For the physical bounds the new rational evaluator directly computes all
$26$ entries of $t$ from \eqref{v40:eq:rational-t}, with their joint root-cube
enclosure $\|t-\widetilde t\|_\infty\le e_t$. It gives $e_t<0.049405$ in
spatial-sum units. Put $\widetilde w=-P_T\widetilde t$ and
$\widetilde v=3Z_T\widetilde w$. From \eqref{v40:eq:H-neumann},
\[
 \|w_*-\widetilde w\|_\infty\le
 \frac{\|P_T\|_\infty}{1-q_T}
 \left(\|\mathcal H_0\widetilde w+\widetilde t\|_\infty
                  +\delta_H\|\widetilde w\|_\infty+e_t\right)=:e_w.
\]
Hence every physical component error is at most
$e_v=3\|Z_T\|_\infty e_w<0.02831$. The exact inequalities
\[
 1785^2<\frac1{27}\sum_i\max\{|\widetilde v_i|-e_v,0\}^2,
 \qquad
 \frac1{27}\sum_i(|\widetilde v_i|+e_v)^2<1787^2,
 \qquad \|\widetilde v\|_\infty+e_v<2500
\]
prove \eqref{v40:eq:v-bounds}. Their arithmetic is rational. Analyticity follows
from the independently verified coefficient-root chart and the nonzero
$\det\mathcal H$; no smooth continuation of a singular pseudoinverse is used.
\end{proof}

The coefficient is finite but appreciably larger than the first tilt.
Only $av_*$ tends to zero with amplitude. No minimal-norm, optimal-conditioning,
efficient-preparation or cutoff-uniform claim is made. Nor is
\eqref{v40:eq:v} an instruction to impose this rate throughout a history.
The next theorem selects initial records on which it satisfies the original
consistency equation exactly.

\section{Exact initial records realizing the corrected consistency equation}
\label{v40:sec:initial}

For nearby external parameters $\eta$ of \cref{v39:cor:six}, evaluate
$v(\eta)$ by \eqref{v40:eq:v} at that coefficient root. It is held fixed while
selecting the amplitude-dependent initial records. It is not recomputed from
a changing physical history.

\begin{theorem}[Nonempty exact initial consistency with the corrected rate]
\label{v40:thm:initial}
There are analytic initial records for the unchanged action, at all sufficiently
small amplitudes and $\eta$ near its central value, satisfying
\begin{equation}
 \boxed{P(X(a,\eta),\lambda(a,\eta))=0,\qquad
        \mathsf B(X(a,\eta),\lambda(a,\eta))
                       +aM(X(a,\eta),\lambda(a,\eta))v(\eta)=0.}
 \label{v40:eq:initial}
\end{equation}
They have the same leading canonical field and first multiplier tilt as
Version 39, with
\[
 X=aY_*(\eta)+O(a^2),\qquad \lambda=e+a\ell_*(\eta)+O(a^2),
 \qquad \overline{\lambda^N}=1,\quad \overline{\lambda^\beta}=0.
\]
Their metric and lapse are positive for sufficiently small amplitude. The
initial tangent $(F,av(\eta))$ satisfies every original first constraint
consistency equation exactly. This assertion requires no invertibility of $M$.
\end{theorem}
\begin{proof}
Use the actual analytic constraint initializer from
\cref{v39:lem:analytic-prep} with
$\lambda=e+a\ell(q,\eta)+a^2m$, where $m\in\mathcal E^\perp$.
It solves $P=0$ identically. Define
$\mathcal B_v=\mathsf B+aMv(\eta)$ on these initial records.
Quadratic drift cancellation and $M=O(a^2)$ make $a^{-3}\mathcal B_v$
analytic through $a=0$.

The original source grading has blocks
\begin{equation}
 M_{AA}=a^2(A_0+O(a)),\qquad
 M_{AW}=a^3(J_0+O(a)),\qquad
 M_{WW}=a^4(C_0+O(a)).
 \label{v40:eq:grading}
\end{equation}
Since $v(\eta)\in\mathcal A$, its contribution $aMv$ has order $a^3$
on active rows and at least $a^4$ on weak rows, including all constant shifts.
Thus the three divided constant-shift equations at $a=0$ are exactly the same
$\mathbf C(q,\eta)=0$ as in Version 39. Their derivative in $m$ is zero.
On all $104$ nonconstant rows the derivative in $m$ remains the invertible
restriction of $\mathbb K_1(Y)$ to $\mathcal E^\perp$: the leading term of
$M$ depends only on $Y$, while $m$ enters the multiplier at order $a^2$.
No numerical value or inverse of $A_0,J_0,C_0$ is used.

At $q=q_*(\eta_*)$, the actual nonconstant cubic baseline of
$\mathcal B_v$ is a finite vector. Choose the unique $m_*$ canceling it by
$\mathbb K_1|_{\mathcal E^\perp}$. The divided $107$-row map has derivative
in $(q,m)$
\[
 \begin{pmatrix}
 D_q\mathbf C&0\\ *&\mathbb K_1(Y_*)|_{\mathcal E^\perp}
 \end{pmatrix},
\]
whose diagonal blocks were verified in Version 39. The analytic
implicit-function theorem solves these $107$ exact rows for $q(a,\eta),m(a,\eta)$.
The remaining row is a constant lapse covector. Equation
\eqref{v40:eq:identities} gives
$\langle\lambda,\mathcal B_v\rangle_h=0$, and the mean lapse is one;
hence that last row also vanishes. This proves \eqref{v40:eq:initial}.
In particular $D_XP F+D_\lambda P[av]=\mathsf B+aMv=0$.
The leading records and positivity follow from the same analytic expansion
as in Version 39. Higher baseline values and nonlinear initial roots are
established by this onto argument, not reported as numerically tabulated.
\end{proof}

\begin{corollary}[Interpretation on an independently regular continuation]
\label{v40:cor:actual-rate}
If the true multiplier Hessian at these records has only its radial null
direction and is invertible on the mean-lapse-zero complement, their unique
normalized original continuation has $\lambda_t(0)=av(\eta)$.
Its initial canonical first and second derivatives and both retained initial
curvature reads are $O(a)$ at this fixed cutoff.
\end{corollary}
\begin{proof}
Subtract \eqref{v40:eq:initial} from $M\lambda_t=-\mathsf B$.
The difference lies in the fixed normalized complement and is killed by $M$,
so it is zero under the stated hypothesis. Then
$X_t=F=O(a)$ and
$X_{tt}=D_XF F+D_\lambda F[av]=O(a)$ on the fixed analytic chart.
The retained metric and stationary-link reconstruction uses these jets and
$\lambda_t=O(a)$, giving the stated initial curvature estimates.
\end{proof}

This theorem selects different second and higher initial coefficients from
Version 39. It does not preserve the condition $\mathsf B=0$, which was
unnecessarily restrictive for eliminating the leading differentiated weak
source. What is preserved exactly is the original first consistency equation.
No rate, force, constraint projection or reset is added to the physical evolution.

\section{The leading weak acceleration forcing is removed, not the full inverse}
\label{v40:sec:acceleration}

For the exact initial tangent of \cref{v40:thm:initial}, let
\begin{equation}
 \mathsf R_a=
 D_X\mathsf B[F]+D_\lambda\mathsf B[av]
                 +\{D_XM[F]+D_\lambda M[av]\}\,av.
 \label{v40:eq:acc-source}
\end{equation}
This is the actual differentiated consistency source at that initial tangent.
On any twice differentiable continuation realizing it, the equation is
$M\lambda_{tt}(0)=-\mathsf R_a$.

\begin{theorem}[One-order improvement of the complete weak differentiated source]
\label{v40:thm:improvement}
On the family \eqref{v40:eq:initial},
\begin{equation}
 \boxed{\mathsf R_a|_{\mathcal W}=O(a^3),\qquad
                        \mathsf R_a|_{\mathcal A}=O(a^2).}
 \label{v40:eq:improvement}
\end{equation}
The first bound includes all three constant weak rows. At the Version 39
zero-rate tangent, in contrast, the weak source has the nonzero $a^2$
coefficient in \eqref{v40:eq:defect-interval}.
\end{theorem}
\begin{proof}
By \cref{v40:lem:leading,v40:thm:rate}, the two derivatives of
$\mathsf B_W$ have canceling order-$a^2$ coefficients. Variations of the
first records induced by $q(a,\eta)=q_*(\eta)+O(a)$ change this cancellation
only at order $a^3$. Equations \eqref{v40:eq:flat-orders} give
$\{D_XM[F]+D_\lambda M[av]\}av=O(a^3)$ on every row.
For the active rows the possible order-$a$ term is $\mathsf B_1V$.
The original flat linear drift factors through $L$, and $LV=0$, so this
term vanishes. All remaining terms have order at least $a^2$.
\end{proof}

\begin{corollary}[Conditional acceleration orders under the full signed certificate]
\label{v40:cor:graded-orders}
Assume additionally that both the leading active block $A_0$ and the complete
leading graded signed matrix in \eqref{v40:eq:grading} are invertible at $Y_*$. Then the corrected family has
\begin{equation}
 \lambda_{tt,A}(0)=O(1),\qquad \lambda_{tt,W}(0)=O(a^{-1}).
 \label{v40:eq:conditional-acc}
\end{equation}
At the uncorrected Version 39 data the nonzero leading weak source instead
forces a weak acceleration of order at least $a^{-2}$, with the same
independently assumed regular graded inverse.
\end{corollary}
\begin{proof}
Put $D_a=\operatorname{diag}(aI_A,a^2I_W)$ and
$M=D_a\Gamma(a)D_a$, with $\Gamma(0)$ invertible. By
\eqref{v40:eq:improvement},
$D_a^{-1}\mathsf R_a=(O(a),O(a))$; multiplying by the bounded
$\Gamma(a)^{-1}$ and then $D_a^{-1}$ proves
\eqref{v40:eq:conditional-acc}.
For zero initial rate, the active right side has order $a^2$ and the weak
right side is $a^2\mathfrak t+O(a^3)$ with $\mathfrak t\ne0$.
The leading active block is invertible when the separately retained active
regularity holds. Eliminating it gives the weak Schur equation
$a^4(D_0+O(a))\lambda_{tt,W}=-a^2\mathfrak t+O(a^3)$,
where $D_0$ is invertible under the complete graded certificate.
Thus $\lambda_{tt,W}=-a^{-2}D_0^{-1}\mathfrak t+O(a^{-1})$.
For this last comparison we include active-block regularity in the phrase
``full signed certificate''; invertibility of an indefinite full matrix alone
would not imply invertibility of that principal block.
\end{proof}

The corollary is not a claim that the indicated signed matrices have been
certified here. In particular the correction cancels a leading forcing term;
it does not make every multiplier acceleration bounded or establish a common
physical lifespan. The residual $a^3$ weak source, the higher signed Hessian
and their geometric observation remain in the original equation.

\section{Verification record and the next unresolved calculation}
\label{v40:sec:status}

\paragraph{Proved.}
The new original-source interval \eqref{v40:eq:defect-interval} establishes
noninvariance of the recovered zero-rate initial family. An exact rational
Neumann certificate proves the transverse control map is onto throughout the
root cube. Its fully specified rate coefficient cancels every leading weak
differentiated row, with the physical size bounds \eqref{v40:eq:v-bounds}.
The original analytic initializer then produces exact constrained records
satisfying the corrected first consistency equation, and their complete weak
acceleration forcing is $O(a^3)$ rather than $a^2$. These last initial-data
statements do not assume an inverse of the rate Hessian.

\paragraph{Inherited and checked.}
The recovered root inclusion, source-bank structural identities and tilt bounds
were rerun. The new $27$ tested derivative values (one scalar witness and the
$26$-row control source) are evaluated afresh from the literal rational
$P_2,\Phi_3$ compiler, including both polarization subtractions, all shifted
terms and all Frobenius entries. Exact centre-radius arithmetic encloses the
algebraic root, rather than silently evaluating it at a decimal approximation.
The new $26$-square preconditioner is accepted only through rational residual
inequalities on the whole cube. No complete regeneration of all $79{,}224$
older bank entries, proof-assistant formalization, nonlinear initial root,
or physical-time simulation is represented as newly executed.

\paragraph{Still unresolved.}
The full signed multiplier-Hessian regularity at this finite-distance root
has not been established. It remains required for a unique original
continuation and for the conditional acceleration orders. The correction
changes higher initial records, not the leading canonical first field; the
leading signed-source certificate would therefore serve both initial families.
That certificate and the next remaining differentiated weak source are the
appropriate next calculation. Another attempt to propagate exact zero
multiplier rate at the unchanged first records would disregard the proved
nonzero source. The older unit-initial-multiplier cofinal stress problem of
Version 32, microscopic selection of these preparations, a common physical
lifespan and the Einstein--Standard-Model continuum remain separate.

\begingroup
\renewcommand{\refname}{Additional sources for Version 40}
\begin{thebibliography}{9}
\small\raggedright
\bibitem{v40:EB28}
\emph{Renewal Exact-Action Einstein Bridge Research Notes}, supplied cumulative
Versions 1--28, file
\nolinkurl{Renewal_Exact_Action_Einstein_Bridge_Research_Notes_V28(1).tex}.
Retained interfaces: the original $P_2,\Phi_3$ compiler; the flat identities,
\nolinkurl{v20:graded-source}, the weak joint Taylor formula and the mixed-mode
constraint initializer. No full signed certificate at the new V39 root is
imported from the old local neighborhood.
\bibitem{v40:DBGP}
C.~Di Bartolo, R.~Gambini and J.~Pullin,
\emph{Consistent and mimetic discretizations in general relativity},
J. Math. Phys. \textbf{46} (2005), 032501;
\url{https://arxiv.org/abs/gr-qc/0404052}.
Primary-source background only; all displayed source values and control bounds
in this body follow from the retained finite action and the stated certificate.
\end{thebibliography}
\endgroup
\addtocontents{toc}{\protect\endgroup}
% END IMMUTABLE VERSION 40 BODY
% APPEND VERSION 41 HERE, BEFORE THE FINAL END-DOCUMENT COMMAND.


% BEGIN IMMUTABLE VERSION 41 BODY
\clearpage
\begin{center}
{\Large\bfseries Version 41: The complete signed source at the mixed-mode root}\\[0.5em]
{\large Certified inertia, original-action continuation, and suppression of the first weak time layer}
\end{center}
\makeatletter
\addtocontents{toc}{\protect\begingroup}
\addtocontents{toc}{\protect\makeatletter}
\addtocontents{toc}{\protect\def\protect\l@section{\protect\bfseries\protect\@dottedtocline{1}{0em}{2.5em}}}
\makeatother
\addcontentsline{toc}{part}{Version 41: Complete signed regularity and the corrected original-flow layer}

\section{The signed inverse, rather than another Poisson control}
\label{v41:sec:context}

The mixed-mode root of \cref{v39:thm:root} and the transverse initial-rate
construction of \cref{v40:thm:initial} use the original constraint Poisson
matrix. Neither establishes regularity of the symmetric multiplier-rate
Hessian. This distinction is the remaining condition in
\cref{v40:cor:actual-rate,v40:cor:graded-orders}. We now evaluate that Hessian's
complete leading graded source at the \emph{same exact algebraic root},
including all higher weak columns and their signed pairings. The active block
is positive, but the complete matrix has ten negative directions. Its
nonsingularity cannot be certified by replacing it with an unsigned Gram.

The finite coefficient calculation removes the conditional regularity in the
V39 and V40 local-continuation statements. A further original-flow argument
shows that the V40 correction suppresses the leading weak response on the
physical interval $0\le t\le a^3S$. On this interval its weak rate is $O(a^2)$,
and both curvature observations change by $O(a^2)$ from their initial values.
The interval still shrinks with amplitude. No fixed positive physical
lifespan, spatial-refinement estimate, or Einstein--Standard-Model continuum
is inferred.

\paragraph{Retained objects and scope.}
We use $H,P,F,\mathsf B,M$ and the physical active/weak decomposition of
\cref{v40:eq:objects,v40:eq:grading}. All calculations are at $h=1/3$ on the
original $27$ sites. The four axial/homogeneous parameters and the fourth
mixed amplitude are initially fixed as in V39. Its isolating cube is
\begin{equation}
 \mathcal Q_* =\{q\in\R^3:\|q-q_0\|_\infty\le2^{-80}\},
 \qquad \xi=(q_1,q_2,q_3,1/5),
 \label{v41:eq:cube}
\end{equation}
where $q_0$ is the exact dyadic centre in \cref{v39:eq:centre}. The true root
$q_*$ is the unique zero isolated by that theorem, not its decimal display.
The canonical first field is $Y(q)=(U(q),p(q))$, with the physical momentum
$p=\sqrt3 p_R$ retained. The new matrix certificates hold on all of
$\mathcal Q_*$. Extension to a small part of V39's six-parameter root family
will follow by continuity; no numerical radius in those additional parameters
is claimed.

The original first/second source compiler, flat stationary tangent form, and
invariance of the leading graded source under order-amplitude initial tilts
are inherited from the Exact-Action bridge \cite{v41:EB28}, specifically its
\nolinkurl{v3:compiler}, \nolinkurl{v3:raw-K}, and
\nolinkurl{v20:graded-source}. They are restated sufficiently below to specify
the new executable calculation. The analytic stationary chart and normalized
continuation theorem are retained inputs. The root inclusion is replayed,
not proved again by a new coefficient bank. The additions are the populated
signed certificate at this root and its stated dynamical consequences.

\section{An independently regenerated complete graded source}
\label{v41:sec:compiler}

Use the original phase derivative $\delta_i=3(S_i-S_i^{-1})$, literal cyclic
shifts $S_i$, and Lorentz coordinates $(R,B)$ with
\[
 [(r,b),(s,c)]=(-r\times s+b\times c,-r\times c-b\times s).
\]
At each site the flat spatial normal map $T_0:\R^{18}\to\R^{12}$ and right
inverse $R_0=T_0^T(T_0T_0^T)^{-1}$ are fixed rational matrices. In the auxiliary
order $(A_0^R,A_0^B,w_{11},w_{22},w_{33},w_{12},w_{13},w_{23})$, the normal map
on $(R_i,B_i)_{i=1}^3$ has vector components
\[
 -2\sum_i B_i\times e_i,\qquad -2\sum_i R_i\times e_i,
\]
and symmetric components $B_{ii}-\tr B$ and $B_{ij}+B_{ji}$ for $i<j$.
Its kernel is the symmetric rotational tangent. Let $Z_0$ synthesize that
tangent in the raw order $(11,22,33,12,13,23)$. Its dual map adds the two
off-diagonal occurrences; it does not count them only once in a Frobenius
contraction. The corresponding signed inverse is
\begin{equation}
 K_{\rm raw}^{-1}
 =\frac14\begin{pmatrix}2I_3-\mathbf1\mathbf1^T&0\\0&I_3\end{pmatrix}.
 \label{v41:eq:cartan}
\end{equation}

For $Y=(U,p)$ the first stationary spatial coefficient is
\[
 R_{ia}=-\tfrac12\epsilon_{abc}\delta_bU_{ic},\qquad
 B=p-\tfrac12(\tr p)I.
\]
For a multiplier test $\nu=(n,\beta)$, its flat auxiliary response is
\[
 A_0^R=-\tfrac12\operatorname{curl}_\delta\beta,
 \qquad A_0^B=\nabla_\delta n,
 \qquad w_{ij}=\delta_i\beta_j+\delta_j\beta_i.
\]
The original stationary coefficient functional is denoted by
$\Phi(X,z,\eta;\nu)=\langle\eta,\mathcal T_X(z)\rangle+
\mathcal V_{X,\nu}(z)$, with its canonical $-\langle p,w\rangle$ slot
retained. It is affine in the auxiliary and multiplier entries. Define
\begin{align}
 g_1(\nu)&=[a]D_z\Phi(aY,az_1,\eta_0(\nu);\nu),
 &J_1\nu&=Z_0^*g_1(\nu),\label{v41:eq:first}\\
 \tau_2&=[a^2]\mathcal T_{aY}(az_1),
 &z_2^\perp&=-R_0\tau_2,\notag\\
 \eta_1(\nu)&=-R_0^*g_1(\nu).&&\notag
\end{align}
For weak tests the second source is
\begin{equation}
 J_2\nu=Z_0^*[a^2]D_z\Phi\bigl(aY,az_1+a^2z_2^\perp,
                         \eta_0+a\eta_1;\nu\bigr),
 \qquad\nu\in\mathcal W.
 \label{v41:eq:second}
\end{equation}
Here $[a^j]$ is a coefficient, without a factorial. Both corrections in
\eqref{v41:eq:second} are part of the actual source.

\begin{lemma}[Finite compiler and its action incidence]
\label[lemma]{v41:lem:compiler}
The accompanying rational arithmetic program evaluates
\eqref{v41:eq:first}--\eqref{v41:eq:second} from the original scalar
coefficient functional. Its active and weak outputs are respectively the
complete first and restricted second source columns. It does not require a
weak inverse or a second stationary tangent solve. Its midpoint-radius
version encloses these columns for every input in \eqref{v41:eq:cube}.
\end{lemma}
\begin{proof}
The compiler represents amplitude coefficients through degree two and takes
the spatial derivative by reverse differentiation of the finite arithmetic
circuit. Differentiation is in the independent degree-zero spatial slot;
the retained degree in the other amplitude slots is therefore the required
coefficient of $D_z\Phi$, rather than a truncation of its differentiated
value. The coframe uses $e=I+U/2-U^2/8$ at the necessary degrees, and
$\Pi=2\operatorname{cof}e$ uses
\[
 \Pi_1=(\tr U)I-U,\qquad
 \Pi_2=\tfrac34U^2-\tfrac12(\tr U)U+
       \bigl((\tr U)^2/4-\tr(U^2)/2\bigr)I.
\]
In the velocity slot its derivative additionally retains $D\Pi_3(U)[w]$,
where
\begin{align*}
 \Pi_3(U)={}&
 \bigl((\tr U)^3/24-(\tr U)\tr(U^2)/4+\tr(U^3)/3\bigr)I\\
 &+\bigl(- (\tr U)^2/8+\tr(U^2)/4\bigr)U
 +\tfrac38(\tr U)U^2-\tfrac58U^3.
\end{align*}
This term is not replaced by $D\Pi_2$ when a degree-zero velocity multiplies
it before spatial differentiation.

The electric term includes
$-(h/2)\sum_i\langle\Pi_i,[z_i,[z_i,S_iA_0]]_B\rangle$.
The magnetic term includes the cubic BCH coefficient of the four literal
ordered factors $(z_i,S_i z_j,-S_jz_i,-z_j)$, multiplied by the original
$h\Sigma_{ij}$. Its recursion, initialized with $l_1=l_2=l_3=0$, is
\[
 l_3\leftarrow l_3+\tfrac12[l_2,v]+
 \tfrac1{12}\bigl([l_1,[l_1,v]]+[v,[v,l_1]]\bigr),\quad
 l_2\leftarrow l_2+\tfrac12[l_1,v],\quad l_1\leftarrow l_1+v,
\]
with simultaneous old values. Every coframe, dual and shifted occurrence is
retained in these contractions. Higher logarithmic terms have too high a
degree in the active first source. In the weak second source the zeroth
auxiliary $A_0$ is zero, so the quartic electric term also has too high a
degree; this is the inherited restricted degree argument, not a deletion
from the action.

Differentiate the retained circuit by the product, sum, contraction and
transpose-shift rules. Induction over its operations proves that the output
is its exact spatial derivative. The normal equation supplies
$z_2^\perp$ and the first auxiliary equation supplies $\eta_1$. The inherited
flat tangent-Hessian cancellation makes the uncomputed second tangent
connection contribute zero on weak tests. Thus the result is exactly the
original restricted source, as in the bridge's compiler theorem.

For the enclosure, a rational centre array and a common nonnegative radius
represent componentwise closed intervals. Each arithmetic and reverse
operation propagates a valid inclusion by the triangle inequality. Centres
use arbitrary-precision integers with exact denominators. Radii are rounded
\emph{upwards} to multiples of $2^{-180}$ after each operation. The rational
endpoints for $\sqrt3$ are independently verified by squaring. The same
induction proves inclusion of all original values on the root cube, including
the derivative circuit. No floating differentiation enters this proof.
\end{proof}

The active columns initially use the point indices
\begin{align*}
\mathcal I_A=(&0{:}25,27{:}34,36{:}43,45{:}52,\\
 &54,55,56,57,58,60,61,63,64,65,66,67,69,70,72,73,75,76,78,79,\\
 &81,84,87,90,93,96,99,102).
\end{align*}
where intervals are inclusive and indices are zero based. Subtract their
constant lapse and curl-free shift parts to obtain physical active
representatives. Their $J_1$ columns do not change. The weak basis comprises
the three constant shifts followed by
$\nabla_\delta(e_j-e_{26})$, $0\le j<26$. Write
\begin{equation}
 C=J_1|_{\mathcal I_A}\in\R^{162\times78},\qquad
 W=J_2|_{\mathcal W}\in\R^{162\times29},\qquad Q=(C,W),
 \label{v41:eq:columns}
\end{equation}
where each of $27$ sites has six raw symmetric components. The two computed
matrices are
\begin{equation}
 A=C^T\mathcal K^{-1}C,\qquad
 \Gamma=Q^T\mathcal K^{-1}Q
     =\begin{pmatrix}A&L\\L^T&C_W\end{pmatrix},
 \qquad\mathcal K^{-1}=I_{27}\otimes K_{\rm raw}^{-1}.
 \label{v41:eq:Gram}
\end{equation}
These numerical matrices use spatial sums. Multiplying all covectors in the
consistency equation by $h^3=1/27$ restores the normalized spatial pairing.
It changes neither the solution, inertia nor rank. Stated numerical inverse
norms concern the displayed sum-coordinate matrices, not a
cutoff-independent physical norm.

\section{A signed inclusion on the whole algebraic root cube}
\label{v41:sec:certificate}

\begin{theorem}[Complete graded signed regularity]
\label[theorem]{v41:thm:signed}
For every $q\in\mathcal Q_*$ the matrices in \eqref{v41:eq:Gram} satisfy
\begin{equation}
 \boxed{A\succ0,\quad\|A^{-1}\|_\infty<1,\qquad
 \operatorname{Inertia}(\Gamma)=(97,10,0),\quad
 \|\Gamma^{-1}\|_\infty<13.}
 \label{v41:eq:certificate}
\end{equation}
Inertia is ordered as positive, negative, zero. Consequently the full source
$Q$ has rank $107$, and its weak Schur complement
\begin{equation}
 D=C_W-L^TA^{-1}L
 \label{v41:eq:D}
\end{equation}
has inertia $(19,10,0)$. These are statements at the exact V39 root as well
as at every other point of its isolating cube.
\end{theorem}
\begin{proof}
The package contains the rational centre arrays for all $17{,}334$ entries
of $C,W$ and fixed dyadic preconditioners $R_A,R_\Gamma$ and congruences
$S_A,S_\Gamma$. The verifier regenerates the centre arrays and enclosures
from \cref{v41:lem:compiler}; it does not accept stored floating eigenvalues.
The following strict bounds are proved by integer numerator comparisons:
\begin{center}\small
\begin{tabular}{@{}lcc@{}}
\toprule
Quantity & Active matrix $A$ & Complete matrix $\Gamma$\\\midrule
Uniform componentwise matrix radius & $2^{-48}$ & $2^{-28}$\\
$\|I-RX_c\|_\infty+\|R\|_\infty n r_X$ & $2^{-33}$ & $2^{-18}$\\
$\|R\|_\infty/(1-q_X)$ & $1$ & $13$\\
Congruence error from the specified sign matrix & $2^{-37}$ & $2^{-16}$\\
Specified positive signs & $78$ & $97$\\
Specified negative signs & $0$ & $10$\\\bottomrule
\end{tabular}
\end{center}
Here every entry in the first four rows is a strict upper bound;
$X_c$ is the exact rational centre, $r_X$ its common component radius,
$n$ its matrix dimension, and $q_X$ the preceding Neumann expression.
The source radii themselves are less than $2^{-65}$ for $C$ and $2^{-50}$
for $W$. All these inclusions concern the actual simultaneous root cube;
independent interval boxes may enlarge it but cannot invalidate inclusion.

For completeness, the matrix argument is elementary. If $q_X<1$, then
$RX=I-(I-RX)$ is invertible by its Neumann series; square matrices imply
$R$ and $X$ are both invertible, and
$\|X^{-1}\|_\infty\le\|R\|_\infty/(1-q_X)$.
For inertia the actual checked quantity is
\begin{equation}
 \|S^TX_cS-\mathcal D\|_\infty
 +\|S^T\|_\infty\|S\|_\infty n r_X<1,
 \label{v41:eq:inertia-test}
\end{equation}
where $\mathcal D$ is diagonal with the specified entries $\pm1$.
The perturbation is symmetric, so its Euclidean operator norm is at most
its infinity norm. The path from $\mathcal D$ to $S^TXS$ therefore never
has a zero eigenvalue. Its inertia is that of $\mathcal D$. In particular
$S^TXS$ is invertible, which also proves $S$ invertible. Sylvester's law
then gives the inertia of $X$ without any assumed accuracy of the proposed
congruence. The residual technique has established verification precedents
\cite{v41:Rump}; the source arrays and successful inequalities are supplied
here.

Applying these arguments to the two regenerated enclosures proves
\eqref{v41:eq:certificate}. A null column combination of $Q$ would be a
null vector of $\Gamma$, so $Q$ is injective. Block congruence separates
$A$ and $D$. Subtracting the active inertia $(78,0,0)$ gives $(19,10,0)$
for $D$. No sign has been replaced by its absolute value.
\end{proof}

\paragraph{Original-source cross-checks.}
Before evaluating the new root, the independent compiler was tested on the
older complementary seed of the bridge. With $\sqrt3\mapsto56$ modulo
$241$, it reproduces the inherited determinants
\[
 \det A=131,\qquad\det\Gamma=15\pmod{241},
\]
the first and full source ranks $78,107$, and the three constant-column
trace residues $(17,78,102)$ at the origin. These are regression checks,
not a transfer of regularity from an unquantified old neighborhood.
The verifier additionally checks $T_0R_0=I$, symmetry and idempotence of
$R_0T_0$, and first, weak and auxiliary derivative contractions by an
independent exact cubic value-polarization formula at primes $241,1009$.
For a cubic polynomial $f$, the formula is
\[
 f'(0)=\frac{8(f(1)-f(-1))-(f(2)-f(-2))}{12}.
\]
It is an exact coefficient identity, not a finite-difference approximation.
The actual cube inclusion is in characteristic zero throughout.

\section{The two prepared families now have original local continuations}
\label{v41:sec:continuation}

Let $\mathcal A$ be the physical mean-zero lapse/transverse-shift space,
$\mathcal W$ the complete curl-free shift space, and retain their fixed
bases above. The constant lapse is excluded by the normalized rate
condition, not deleted from the original constraints. Matrix covectors in
this section use the same sum convention as \eqref{v41:eq:Gram}.

\begin{theorem}[The finite-amplitude Hessian and its only null direction]
\label[theorem]{v41:thm:true-Hessian}
For either the V39 zero-drift family or the V40 corrected-rate family, at the
central root and on a sufficiently small part of its six-parameter family,
there is $a_0>0$ such that for $0<|a|<a_0$ the full original multiplier
Hessian has rank $107$ and inertia $(97,10,1)$. Its nullspace is exactly
$\R\lambda$. Its restriction to the fixed mean-lapse-zero complement is
invertible. Each such initial record has a unique local normalized
original-action continuation, preserving all original constraints.
\end{theorem}
\begin{proof}
The inherited graded-source identity, including its tilt and second-field
cancellations, gives on the fixed complement
\begin{equation}
 M_{\rm comp}(a)=D_a\bigl(\Gamma_*+O(a)\bigr)D_a,
 \qquad D_a=\operatorname{diag}(aI_{78},a^2I_{29}).
 \label{v41:eq:graded}
\end{equation}
The original stationary tangent inverse tends to
\eqref{v41:eq:cartan}. The leading $\Gamma_*$ is the matrix just certified,
independent of the different finite higher initial coefficients in the two
families. The remainder is analytic at this fixed cutoff. Its size need
not be inferred from a computed leading inverse. Nonsingularity and inertia
therefore persist for sufficiently small nonzero $a$. Congruence by the
invertible $D_a$ preserves the inertia $(97,10,0)$.

The exact homogeneity identity gives $M\lambda=0$, while
$\overline{\lambda^N}=1$ makes $\lambda$ transverse to the normalized
complement. Hence the full space splits as that complement plus
$\R\lambda$, and the form has exactly the stated one-dimensional radical.
There is no additional finite null direction. The original stationary chart
and its normalized vector field now satisfy the inherited continuation
hypotheses: solve $M\lambda_t=-\mathsf B$ on the complement, use
$\langle\lambda,\mathsf B\rangle=0$ to recover the omitted radial equation,
and apply local analytic ODE existence and uniqueness. The exact identity
$\dot P=\mathsf B+M\lambda_t=0$ propagates $P=0$.

All leading matrices depend continuously on the remaining root-family
parameters. Shrinking their neighborhood preserves the certified
nonsingularity. On a compact smaller parameter neighborhood the finite
analytic preparation and stationary charts also have common local bounds.
This proves the asserted local extension without claiming a numerical
amplitude radius, an all-cutoff estimate, or an amplitude-independent
physical interval.
\end{proof}

\begin{corollary}[Actual initial rates and initial geometric scale]
\label[corollary]{v41:cor:rates}
On these original continuations, the V39 family has $\lambda_t(0)=0$ and the
V40 family has $\lambda_t(0)=av(\eta)$ exactly. At their initial cuts,
\begin{equation}
 X,X_t,X_{tt},\lambda-e,\lambda_t=O(a),\qquad
 \|\Riem(g_h)(0)\|+\|E_h^{\rm link}(0)\|
                  +\|F_h^{{\rm sp},{\rm link}}(0)\|=O(|a|).
 \label{v41:eq:initial-geometry}
\end{equation}
The records retain the nonzero profiles and nontrivial leading canonical
motion of V39. These statements do not make them equilibria.
\end{corollary}
\begin{proof}
The respective exact initial equations are $\mathsf B=0$ and
$\mathsf B+aMv=0$. Both proposed rates belong to the normalized complement;
\cref{v41:thm:true-Hessian} makes them the unique actual rates. The original
canonical field is analytic with $F(0,e)=0$, so $X_t=O(a)$ and
$X_{tt}=D_XF[F]+D_\lambda F[\lambda_t]=O(a)$. At fixed grid the required
spatial derivatives have the same orders. The retained ADM curvature
formulas use these field derivatives, $\lambda_t$ and their spatial
derivatives, not $\lambda_{tt}$. The stationary connection and its first
time derivative have the same $O(a)$ scale, giving the literal ordered
curvature estimates. No equality of the first metric and link curvature
coefficients is asserted.
\end{proof}

\section{The acceleration comparison is now unconditional at this cutoff}
\label{v41:sec:acceleration}

Write the leading matrix as
$\Gamma_*=(\begin{smallmatrix}A_0&L_0\\L_0^T&C_0\end{smallmatrix})$
and $D_0=C_0-L_0^TA_0^{-1}L_0$. Both inverses now exist by
\cref{v41:thm:signed}. The fixed bases and the same covector convention are
used throughout.

\begin{theorem}[Zero-rate pole and corrected differentiated source]
\label[theorem]{v41:thm:acceleration}
For V39's zero-rate family at the central root, let $t_W$ be the complete
weak covector from \cref{v40:eq:t}. Then its actual continuation obeys
\begin{equation}
 \lambda_{tt,W}(0)=-a^{-2}D_0^{-1}t_W+O(a^{-1}),
 \qquad \|\lambda_{tt,W}(0)\|_h\ge c|a|^{-2}.
 \label{v41:eq:old-pole}
\end{equation}
For the V40 corrected family its actual continuation instead satisfies
\begin{equation}
 \boxed{\lambda_{tt,A}(0)=O(1),\qquad
        \lambda_{tt,W}(0)=O(|a|^{-1}).}
 \label{v41:eq:new-accel}
\end{equation}
The constants are finite at this cutoff. A nonzero leading coefficient of
the corrected weak acceleration is not asserted.
\end{theorem}
\begin{proof}
For the zero-rate family, differentiating consistency at the initial cut
gives $M\lambda_{tt}=-D_X\mathsf B[F]$. V40 supplies active target
$O(a^2)$ and weak target $a^2t_W+O(a^3)$. Eliminating the active equation
adds only $O(a^3)$ to the weak target; the weak Schur matrix is
$a^4(D_0+O(a))$. This proves the expansion. V40's exact interval
$-6671<t_W[\nu_0]/\sqrt3<-6670$ makes $t_W$ nonzero. Invertibility of
$D_0$, finite-dimensional norm equivalence and the reverse triangle
inequality give the lower bound.

For the corrected family the complete differentiated source is the original
\cref{v40:eq:acc-source}, including $\dot M\lambda_t$. By
\cref{v40:thm:improvement} its active part is $O(a^2)$ and its weak part
$O(a^3)$. The active-to-weak elimination again multiplies the active target
by $a$, leaving weak target $O(a^3)$. The Schur solve gives the weak
$O(a^{-1})$ rate; substitution into the active equation gives
$O(1)$, because $a^{-2}a^3a^{-1}=O(1)$. This discharges the regularity
hypotheses of \cref{v40:cor:graded-orders}, rather than treating the
source estimate alone as an inverse estimate.
\end{proof}

\section{An analytic scaled field for the corrected original histories}
\label{v41:sec:layer-field}

We now fix one of the corrected analytic initial families. Write its initial
records as $(X_a,\lambda_a)$, its actual initial rate as $(av,0)$ in the
active/weak splitting, and set
\[
 V=F_1Y+G\ell.
\]
The initial $v$ is chosen once; it is not recomputed during the trajectory.
Let $\mathsf B_1=D_X\mathsf B(0,e)$, let $R_A,R_W$ restrict covectors,
and let $I_A,I_W$ synthesize multiplier components. Define
\begin{align}
 \mathcal C_\ell Z[\nu]
    &=Dd_{2,W}(Y)[Z,\nu]+\mathbb K_1(Z)[\nu,\ell],\notag\\
 \mathcal K_W z[\nu]&=\mathbb K_1(Y)[\nu,z],\qquad\nu\in\mathcal W.
 \label{v41:eq:linear-maps}
\end{align}
These are the same weak Taylor maps as V40. Its exact cancellation is
\begin{equation}
 \mathcal C_\ell V+\mathcal K_W I_Av=0,
 \qquad \mathsf B_1V=0.
 \label{v41:eq:limit-cancel}
\end{equation}
The second identity is the original flat constraint identity: $LV=0$ and
$\mathsf B_1$ factors through the linear constraint map $L$.

Use the comparison variables
\begin{equation}
 t=a^3s,\qquad X(t)=X_a+a^4\xi(s),\qquad
 \lambda(t)=\lambda_a+a^4\zeta(s),\qquad a>0.
 \label{v41:eq:scales}
\end{equation}
The time $t$ remains the original physical time. The variable $s$ is not
substituted as a new physical clock.

\begin{lemma}[Removable singularities and the complete limiting field]
\label[lemma]{v41:lem:layer-field}
On every bounded set of $(\xi,\zeta)$, the normalized original vector field
in \eqref{v41:eq:scales} extends analytically to $a=0$. Its limiting system is
\begin{align}
 \xi'&=V,\qquad \zeta_A'=v,\notag\\
 \zeta_W'&=-D_0^{-1}\bigl\{
 \mathcal C_\ell\xi+\mathcal K_W(I_A\zeta_A+I_W\zeta_W)
          -L_0^TA_0^{-1}R_A\mathsf B_1\xi\bigr\}.
 \label{v41:eq:limit-field}
\end{align}
The convergence of the fields is $O(a)$ in every fixed finite $C^k$ norm
on a compact set. The extension uses the full original signed rate solve.
\end{lemma}
\begin{proof}
The weak joint Taylor expansion is
\[
 \mathsf B_W(X,e+\chi)=d_{2,W}(X)+
                 \mathbb K_1(X)|_W\chi+O((\|X\|+\|\chi\|)^3).
\]
Its initial field and multiplier derivatives are respectively
$a\mathcal C_\ell+O(a^2)$ and $a\mathcal K_W+O(a^2)$.
For the active rows the retained linear field derivative is $R_A\mathsf B_1$.
A multiplier derivative has order $a$. Hence on the scaled tube
\begin{align}
 \delta\mathsf B_A&=a^4R_A\mathsf B_1\xi+O(a^5),\notag\\
 \delta\mathsf B_W&=a^5(\mathcal C_\ell\xi+\mathcal K_W\zeta)+O(a^6).
 \label{v41:eq:delta-B}
\end{align}
The remainders and their fixed derivatives are analytic with these powers.

The original active source is of joint order one, and the weak source is
of joint order at least two near $(0,e)$. This uses the inherited universal
weak first-source cancellation and the zero flat source throughout a
multiplier neighborhood. The $a^4$ displacements therefore change those
sources by $O(a^4)$ and $O(a^5)$, respectively. In the original factorization
$M=J^*K^{-1}J$, with the analytic stationary tangent inverse retained, the
active, mixed and weak block changes have orders $a^5,a^6,a^7$. Consequently
\begin{equation}
 \Gamma(a,\xi,\zeta)-\Gamma(a,0,0)=O(a^3).
 \label{v41:eq:delta-Gamma}
\end{equation}
The newly certified $\Gamma_*$ makes its inverse analytic on each bounded
tube for all sufficiently small $a$.

At the actual initial state, $\mathsf B=-aMv$, so
$D_a^{-1}\mathsf B=O(a^2)$. Changing the inverse in
\eqref{v41:eq:delta-Gamma} thus contributes only $O(a^3)$ to the unscaled
weak rate. Using \eqref{v41:eq:delta-B}, the leading rate change is
\[
 \delta\lambda_{t,W}
 =-a\,R_W\Gamma_*^{-1}
 \binom{R_A\mathsf B_1\xi}{\mathcal C_\ell\xi+\mathcal K_W\zeta}
 +O(a^2),\qquad
 \delta\lambda_{t,A}=O(a^2).
\]
Block inversion gives the last line of \eqref{v41:eq:limit-field};
$d\lambda/dt=a\zeta'$ gives its stated normalization. Likewise
$F(X_a+a^4\xi,\lambda_a+a^4\zeta)/a=V+O(a)$.
The apparent negative powers in the rate have therefore removable analytic
singularities, not merely pointwise limits. Differentiating these finite
Taylor formulas proves the claimed $C^k$ convergence. In particular the
$\mathsf B_1\xi$ term is retained off the constraint tangent; it is not
deleted before the limiting trajectory is determined.
\end{proof}

\section{The leading weak layer vanishes on the corrected family}
\label{v41:sec:layer}

\begin{theorem}[Original-flow existence and the suppressed weak response]
\label[theorem]{v41:thm:layer}
For every fixed $S<\infty$ there exists $a_S>0$ such that the actual
normalized histories from the corrected initial family exist on
$0\le t\le a^3S$ for $0<a<a_S$, preserving every original constraint.
In \eqref{v41:eq:scales}, uniformly for $0\le s\le S$ and with their first
two $s$ derivatives,
\begin{equation}
 \boxed{\xi(s)=sV+O_S(a),\qquad
        \zeta_A(s)=sv+O_S(a),\qquad \zeta_W(s)=O_S(a).}
 \label{v41:eq:layer-limit}
\end{equation}
In particular,
\begin{align}
 X(a^3s)-X_a&=a^4sV+O_S(a^5),\notag\\
 \lambda_A(a^3s)-\lambda_{A,a}&=a^4sv+O_S(a^5),\notag\\
 \boxed{\lambda_W(a^3s)-\lambda_{W,a}=O_S(a^5),\qquad
        \lambda_{t,W}(a^3s)=O_S(a^2),}\notag\\
 \lambda_{t,A}(a^3s)&=av+O_S(a^2).
 \label{v41:eq:physical-layer}
\end{align}
All constants are for the fixed cutoff and initial family. The assertion is
not that $a^3$ is the maximal lifespan.
\end{theorem}
\begin{proof}
For zero initial scaled displacements the limiting field in
\eqref{v41:eq:limit-field} has $\xi_0=sV$, $\zeta_{A,0}=sv$. Substitution
and \eqref{v41:eq:limit-cancel} give
\[
 \zeta_{W,0}'=-D_0^{-1}\mathcal K_W I_W\zeta_{W,0},
 \qquad\zeta_{W,0}(0)=0.
\]
Thus $\zeta_{W,0}=0$ identically. This uses the exact rate correction;
the uncorrected V39 source does not have this cancellation.

The limiting trajectory and its derivatives are bounded on each fixed
$[0,S]$. Choose a compact tube with that trajectory in its interior.
\Cref{v41:lem:layer-field} supplies uniform Lipschitz bounds there and
$O(a)$ differences between the scaled original field and its limit.
The integral equations and Gronwall's inequality give an $O_S(a)$
trajectory difference until first exit. Taking $a$ sufficiently small
precludes that exit. Local existence and continuation on the tube then
supply the whole interval $[0,S]$; the $C^1$ field convergence and the ODE
supply convergence of first and second $s$ derivatives.

The original stationary chart, lapse positivity, spatial metric positivity
and signed regularity persist on this tube because its physical
displacements are $O_S(a^4)$ and the normalized leading matrix stays
nonsingular. Undoing the scaling therefore gives the actual original
histories, not an evolution defined by repeatedly solving the initial
conditions. Their exact equation propagates $P=0$.
Finally multiply the displacement estimates by $a^4$ and use
$\lambda_t=a\zeta'$ to obtain \eqref{v41:eq:physical-layer}.
\end{proof}

The signed weak Schur matrix is not positive; the proof makes no assertion
that generic errors have a contractive limiting flow. It only uses finite
ODE bounds on fixed $S$ and the exact zero forced response for the selected
limiting initial trajectory. Also, $a_S$ may deteriorate with $S$; one
cannot set $S=T/a^3$ in this theorem to obtain a fixed physical interval.

\section{The same suppression is visible in both curvature observations}
\label{v41:sec:geometry}

\begin{theorem}[Small curvature and no leading normalized layer change]
\label[theorem]{v41:thm:geometry}
On the original interval of \cref{v41:thm:layer}, the ordinary interpolated
ADM metric and the literal original link observations satisfy, in every
fixed spatial norm at this cutoff,
\begin{align}
 \sup_{t\le a^3S}\bigl(
 \|\Riem(g_h)(t)\|+\|E_h^{\rm link}(t)\|
                  +\|F_h^{{\rm sp},{\rm link}}(t)\|\bigr)&=O_S(a),
 \label{v41:eq:curvature-size}\\
 \sup_{t\le a^3S}\bigl(
 \|\Riem(g_h)(t)-\Riem(g_h)(0)\|
 +\|E_h^{\rm link}(t)-E_h^{\rm link}(0)\|\notag\\
 \hspace{30mm}
 +\|F_h^{{\rm sp},{\rm link}}(t)-F_h^{{\rm sp},{\rm link}}(0)\|\bigr)
 &=O_S(a^2).
 \label{v41:eq:curvature-change}
\end{align}
Thus each observation's change divided by $a$ tends uniformly to zero on
the scaled interval. No equality of the separate metric and link first
amplitude coefficients is asserted.
\end{theorem}
\begin{proof}
The initial fields and multiplier deviations are $O(a)$, their displacements
are $O_S(a^4)$, and their rates are $O_S(a)$ by
\cref{v41:cor:rates,v41:thm:layer}. The analytic canonical equation gives
$X_t=O_S(a)$ and $X_{tt}=D_XF[F]+D_\lambda F[\lambda_t]=O_S(a)$.
At fixed grid all required spatial derivatives satisfy the same estimates.
The retained coordinate formula
\[
 R_{0i0j}=\tfrac12(g_{0j,it}+g_{i0,tj}-g_{00,ij}-g_{ij,tt})
          +\text{Christoffel products}
\]
and the other independent Riemann blocks use at most $\lambda_t$, not
$\lambda_{tt}$. Expressed through the original canonical equation, the
metric curvature is therefore a smooth finite function of
$(X,\lambda,F,\lambda_t)$ on the positive chart. This proves its size bound.
The stationary connection and its velocity are likewise smooth functions
of these arguments, proving the literal size bounds without replacing the
link connection by a reconstructed Levi--Civita connection.

For the difference estimate, \cref{v41:eq:physical-layer} gives
$\lambda_t(t)-\lambda_t(0)=O_S(a^2)$, while the field and multiplier
displacements are $O_S(a^4)$. Analyticity gives $F(t)-F(0)=O_S(a^4)$.
The mean-value estimate for the same smooth finite curvature functions
therefore gives \eqref{v41:eq:curvature-change}. Their derivatives are
bounded on the fixed stationary chart and do not involve the singular
multiplier inverse: that inverse has already determined the separately
controlled rate. This does not bound a physical time derivative of
curvature uniformly in $a$.
\end{proof}

\section{Verification record and remaining dynamical question}
\label{v41:sec:status}

The recovered V39 root inclusion and its tilt bounds were replayed from the
unchanged embedded parent data. The new standalone verifier independently
regenerates all $12{,}636$ active and $4{,}698$ weak source entries on that
whole root cube. It checks the signed Gram symmetry, two inverse bounds,
two inertia congruences, the old-source signed determinant/trace regression,
and independent value-polarization derivatives. Its default run needs only
Python and NumPy; the original coefficient compiler, all exact rational
centres, root data and fixed dyadic proposals are embedded. The complete
$79{,}224$-entry older Poisson/cubic root bank is retained and replayed, not
claimed newly regenerated. The new $17{,}334$-entry stationary-source bank
is a different bank and is actually regenerated.

The arithmetic certificate is not a simulation of the nonlinear initial
roots. Their existence is inherited from the already proved analytic
initializers, and their unique original local continuations now follow
from the populated signed certificate. The scaled-flow and curvature
estimates are the written analytical arguments above, not conclusions
drawn from integrating a truncated amplitude polynomial. All V40 bytes
before its final document command are preserved.

The immediate improvement is substantive: the full leading signed inverse
at this mixed-mode root is no longer missing. The V40 transverse correction
is now the actual initial rate of a unique original history, and the
leading weak time-layer response is absent. The corrected weak acceleration
still has only the upper bound $O(a^{-1})$. Its next coefficient, the
corresponding original sources, and control beyond the shrinking interval
remain to be established. Negative directions of a constraint-rate
Hessian do not by themselves prove a physical instability, just as
nonsingularity does not prove stability of the evolved system.

The next question is the corrected flow's longer physical-time behavior or
its next invariant source, with this same signed inverse and both curvature
observations retained. Repeating first-Poisson rank or assuming positivity
of the weak Schur form would add nothing. Spatial refinement, a common
positive physical lifespan, identification of the primitive preparation and
operations, nonzero matter backreaction, and the full Einstein--Standard-Model
continuum remain distinct. The older unit-initial-multiplier stress problem
of Version 32 is not solved by a certificate for these different initial
records.

\begingroup
\renewcommand{\refname}{Additional sources for Version 41}
\begin{thebibliography}{9}
\small\raggedright
\bibitem{v41:EB28}
\emph{Renewal Exact-Action Einstein Bridge Research Notes}, supplied cumulative
Versions 1--28,
\nolinkurl{Renewal_Exact_Action_Einstein_Bridge_Research_Notes_V28(1).tex}.
Retained interfaces: the original stationary coefficient functional,
\nolinkurl{v3:compiler}, \nolinkurl{v3:raw-K},
\nolinkurl{v20:graded-source}, the normalized continuation theorem and the
V22 scaled-field/metric identities. No signed regularity at the new root is
imported from the older seed.
\bibitem{v41:Rump}
S.~M.~Rump, \emph{Verification of positive definiteness},
BIT Numerical Mathematics \textbf{46} (2006), 433--452.
\url{https://doi.org/10.1007/s10543-006-0056-1}.
Primary-source context for verified matrix inequalities. The signed
congruence criterion used here is proved above; the present source bank and
its numerical enclosures are independently supplied.
\bibitem{v41:DBGP}
C.~Di Bartolo, R.~Gambini and J.~Pullin,
\emph{Consistent and mimetic discretizations in general relativity},
J. Math. Phys. \textbf{46} (2005), 032501.
\url{https://arxiv.org/abs/gr-qc/0404052}.
Background distinguishing multiplier selection from propagation; not a
source of the mixed-mode certificate or corrected time-layer estimates.
\end{thebibliography}
\endgroup
\addtocontents{toc}{\protect\endgroup}
% END IMMUTABLE VERSION 41 BODY
% APPEND VERSION 42 HERE, BEFORE THE FINAL END-DOCUMENT COMMAND.



% BEGIN IMMUTABLE VERSION 42 BODY
\clearpage
\begin{center}
{\Large\bfseries Version 42: Growth of the weak linearized flow and a longer original-time interval}\\[0.5em]
{\large A verified real eigenmode, a positive comparison norm, and logarithmic-time propagation}
\end{center}
\makeatletter
\addtocontents{toc}{\protect\begingroup}
\addtocontents{toc}{\protect\makeatletter}
\addtocontents{toc}{\protect\def\protect\l@section{\protect\bfseries\protect\@dottedtocline{1}{0em}{2.5em}}}
\makeatother
\addcontentsline{toc}{part}{Version 42: Weak-flow spectrum and logarithmically longer propagation}

\section{The next question is dynamical, not another rank condition}
\label{v42:sec:context}

The signed certificate in \cref{v41:thm:signed} and the exact corrected
initializer of \cref{v40:thm:initial} already give unique original-action
histories and suppression of the leading weak response on every interval
$0\le t\le a^3S$ with $S$ fixed. They do not bound the growth of a small
remaining error when $S$ increases with $a^{-1}$. In particular, the ten
negative directions of the weak signed form neither prove instability nor
permit one to assume oscillatory dynamics. This body computes the actual
weak evolution matrix and uses its verified growth bound to enlarge the
proved interval.

The results are at the unchanged cutoff $h=1/3$, at the same exact mixed-mode
root of \cref{v39:thm:root}, and for one fixed analytic corrected initializer
from \cref{v40:thm:initial}. The physical time remains $t$. No higher initial
coefficient is assigned a new numerical value, and no time-dependent
preparation or feedback is executed. The parent root, original stationary
coefficient compiler, signed regularity, and analytic initializer are
inherited inputs. The new source-specific calculation is a real eigenpair
and a semigroup bound for the weak matrix. The new evolution result treats
an unbounded interval of the comparison time and retains both original
curvature observations.

The conclusion has two parts. There is a real weak eigenvalue
$0.3382576<\sigma_*<0.3382579$ in comparison time. Nevertheless the selected
corrected histories remain small on
$t\le c a^3\log(1/a)$ for every fixed $0<c<50/17$. A growing homogeneous mode
is not a proof that this particular inhomogeneous prepared history excites
it. Its first possible excitation is isolated at the end of this body and
is explicitly left unevaluated.

\section{The complete weak generator and its signed identity}
\label{v42:sec:generator}

Use the physical active/weak bases and spatial-sum convention of
\cref{v41:eq:Gram}. Denote the weak multiplier synthesis by
$\mathsf W:\R^{29}\to\R^{108}$; its first three columns are constant shifts
and its other columns are $\nabla_\delta(e_j-e_{26})$. This $\mathsf W$ is
not the stationary-source matrix called $W$ in \cref{v41:eq:columns}. Put
\begin{equation}
 \begin{gathered}
 D=C_W-L^TA^{-1}L,\
 \mathsf K=\mathsf W^T\mathbb K_1(Y_*)\mathsf W,
 \qquad \mathcal B=-D^{-1}\mathsf K.
 \end{gathered}
 \label{v42:eq:generator}
\end{equation}
The $29$ weak coordinates and all signs are retained. Multiplying both
$D$ and $\mathsf K$ by the common factor $1/27$ leaves $\mathcal B$
unchanged. In the rationalized parent coordinates,
\begin{equation}
 \mathsf K=\frac1{\sqrt3}\mathsf W^T\widehat K(Y_*)\mathsf W.
 \label{v42:eq:physical-factor}
\end{equation}
Indeed $\widehat K=3^{-1/2}Q_3\mathbb K_1Q_3$, and all columns of
$\mathsf W$ have zero lapse. The factor in \eqref{v42:eq:physical-factor}
is essential for the stated numerical rate.

\begin{proposition}[The exact signed conservation identity]
\label[proposition]{v42:prop:signed}
The matrix in \eqref{v42:eq:generator} is real and satisfies
\begin{equation}
 \boxed{\mathcal B^TD+D\mathcal B=0.}
 \label{v42:eq:signed-identity}
\end{equation}
Thus $w(s)^TDw(s)$ is constant when $w'=\mathcal Bw$. Nonzero real
eigenvalues occur in opposite pairs, and every real eigenvector for a
nonzero real eigenvalue is isotropic for $D$.
\end{proposition}
\begin{proof}
The signed certificate gives $D=D^T$ invertible, while the original Poisson
bracket gives $\mathsf K^T=-\mathsf K$. Hence
$\mathcal B^TD=\mathsf K=-D\mathcal B$. Differentiation proves conservation.
Also $\mathcal B^T=-D\mathcal BD^{-1}$, so $\mathcal B$ and
$-\mathcal B$ have the same characteristic polynomial. If
$\mathcal Br=\sigma r$ with $\sigma\in\R\setminus\{0\}$, the signed
identity gives $2\sigma r^TDr=0$.
\end{proof}

For the full layer variable $z=(u,w)$, where
$u=(\xi,\zeta_A)\in\R^{402}$ and $w=\zeta_W\in\R^{29}$, write the
limiting field from \cref{v41:eq:limit-field} as
\begin{equation}
 \mathscr F_0(z)=b+\mathbb Lz,\qquad
 b=(V,v_*,0),\qquad
 \mathbb L=\begin{pmatrix}0&0\\ \mathcal R&\mathcal B\end{pmatrix},
 \label{v42:eq:full-linear}
\end{equation}
where
\[
 \mathcal R(\xi,\zeta_A)
 =-D^{-1}\{\mathcal C_\ell\xi+\mathcal K_W I_A\zeta_A
                  -L^TA^{-1}R_A\mathsf B_1\xi\}.
\]
The verified correction and flat constraint identity say exactly
\begin{equation}
                  \mathbb L b=0.
 \label{v42:eq:Lb}
\end{equation}
Thus $z_0(s)=sb$ is the zero-initial-displacement limiting trajectory. The
spectral question concerns departures from this trajectory, not an
additional forcing on it.

\section{An exact spectral enclosure on a refined root cube}
\label{v42:sec:certificate}

A sharper root enclosure is obtained without finding a new root. The V39
contraction has centre displacement at most $2^{-115}$ and Lipschitz constant
at most $2^{-18}$ on its isolating cube. Therefore its fixed point obeys
\begin{equation}
 \|q_*-q_0\|_\infty
 \le\frac{2^{-115}}{1-2^{-18}}<2^{-114}<2^{-112}.
 \label{v42:eq:refined-root}
\end{equation}
The same stationary compiler is rerun on the cube
$\mathcal Q_{42}=\{\|q-q_0\|_\infty\le2^{-112}\}$. Its rational source
centres agree with V41; its radii are smaller. This is refinement of an
inherited enclosure, not a second independently selected mixed-mode root.

\begin{theorem}[A real growing mode and a verified positive growth norm]
\label[theorem]{v42:thm:spectrum}
For every $q\in\mathcal Q_{42}$, the actual matrices
\eqref{v42:eq:generator} have a simple real eigenpair $(\sigma,r)$ with
\begin{equation}
 \boxed{0.3382576<\sigma<0.3382579,\qquad r_3=1.}
 \label{v42:eq:eigen}
\end{equation}
The index is zero based, so $r_3$ is the first nonconstant weak coordinate.
There is one fixed rational invertible matrix $S\in\R^{29\times29}$,
printed as a rational array in the certificate, such that
\begin{equation}
 \boxed{\frac{S^{-1}\mathcal BS+(S^{-1}\mathcal BS)^T}{2}
                         \preceq\frac{17}{50}I.}
 \label{v42:eq:lognorm}
\end{equation}
In particular these statements hold at the exact root $q_*$.
\end{theorem}
\begin{proof}
We describe the finite inequalities completely; their rational proposals and
successful exact evaluations are in the accompanying executable certificate.
A centre and a common componentwise radius enclose each stationary-source
matrix. The compiler retains both weak corrections in
\cref{v41:eq:second}. The new radii satisfy
\[
 r_{J_1}<2^{-97},\qquad r_{J_2}<2^{-82}.
\]
The verifier also independently regenerates all six matrices
$\mathsf W^T\widehat K_j\mathsf W$ from quadratic constraint polarization:
$6\cdot29^2=5046$ rational entries. It verifies skew symmetry and the exact
three constant columns. For zero-lapse tests the rational block grading is
$\widehat K_{\beta\beta}=3K_{\beta\beta}(U,p_R)$; this is checked before
using \eqref{v42:eq:physical-factor}.

Here is the enclosure chain used instead of an unchecked inverse.
Let $A_c,L_c,C_c$ be the rational centre blocks, with common component radii
$r_A,r_L,r_C$. The inherited fixed preconditioner supplies
$\|A^{-1}\|_\infty\le\alpha$. For a fixed dyadic proposal $Y$ for
$A^{-1}L$, set
\[
 e_Y=\alpha\{\|A_cY-L_c\|_\infty
                   +78r_A\|Y\|_\infty+29r_L\}.
\]
Then $\|A^{-1}L-Y\|_\infty\le e_Y$. A symmetric rational centre for the
true Schur matrix is
\[
 D_c=C_c-\tfrac12(L_c^TY+Y^TL_c).
\]
Its error in induced infinity norm is bounded by
\begin{equation}
 \begin{split}
 e_D={}&29r_C+\|L_c^T\|_\infty e_Y
       +78r_L(\|Y\|_\infty+e_Y)\\
       &+\tfrac12\|L_c^TY-Y^TL_c\|_\infty.
 \end{split}
 \label{v42:eq:schur-error}
\end{equation}
For a fixed rational inverse proposal $R_D$, the condition
$\|I-R_DD_c\|_\infty+\|R_D\|_\infty e_D<1$ gives a rigorous inverse
bound $d_I$ by the Neumann series. Finally a fixed rational matrix
$B_c$ encloses $\mathcal B$ through the residual of
$DB_c+\mathsf K=0$. The cube uncertainty in the six Poisson coefficients
and a rational interval for $\sqrt3$ are included in this residual.
The resulting exact bounds are
\begin{equation}
 e_D<2^{-48},\qquad \|D^{-1}\|_\infty<3/500,
 \qquad \|\mathcal B-B_c\|_\infty<2^{-50}.
 \label{v42:eq:operator-bounds}
\end{equation}
They refer to the fixed raw weak coordinates, not a cofinal physical norm.

For \eqref{v42:eq:lognorm}, let $S_I$ be the fixed rational inverse proposal
for $S$, and put $e_S=\|I-S_IS\|_\infty<2^{-41}$. Thus $S$ is
invertible. With $T_c=S_IB_cS$, the exact similarity error has the bound
\[
 e_T=\frac{e_S}{1-e_S}\|T_c\|_\infty
    +\frac{\|S_I\|_\infty}{1-e_S}
                 \|\mathcal B-B_c\|_\infty\|S\|_\infty.
\]
Let $H_c=(T_c+T_c^T)/2$. The verified rational inequality is
\begin{equation}
 \max_i\left((H_c)_{ii}+\sum_{j\ne i}|(H_c)_{ij}|\right)
                   +15e_T<339/1000<17/50.
 \label{v42:eq:symmetric-test}
\end{equation}
The factor $15=(29+1)/2$ bounds the infinity norm of the symmetric part
of a matrix whose infinity norm is at most $e_T$. The largest eigenvalue
of a real symmetric matrix is bounded by its largest absolute-row disk
endpoint. This proves \eqref{v42:eq:lognorm}; positivity of the physical
signed form has not been assumed.

For the real eigenpair, impose $r_j=1$ at $j=3$. Let $E_j$ insert the
other 28 coordinates, and let $(r_c,\sigma_c)$ be the frozen dyadic
proposal. The nonlinear map is
\[
 \mathcal F(y,\sigma)=\mathcal B(e_j+E_jy)-\sigma(e_j+E_jy).
\]
Its rational centre derivative is
$M_c=((B_c-\sigma_cI)E_j,-r_c)$. For the fixed rational preconditioner
$R_e$, the verifier proves, on the infinity-norm cube of radius
$\rho_e=10^{-7}$ in $(y,\sigma)$,
\begin{align}
 d_e={}&\|R_e(B_c-\sigma_cI)r_c\|_\infty\notag\\
       &+\|R_e\|_\infty\|\mathcal B-B_c\|_\infty\|r_c\|_\infty
                                                     <2^{-44},\notag\\
 q_e={}&\|I-R_eM_c\|_\infty\notag\\
       &+\|R_e\|_\infty(\|\mathcal B-B_c\|_\infty+2\rho_e)
                                                     <2^{-16},\notag\\
 d_e+q_e\rho_e&<\rho_e.
 \label{v42:eq:eigen-contraction}
\end{align}
The last two lines give a strict contraction of that real cube into itself.
Its zero is an actual real normalized eigenpair. The entire interval
$[\sigma_c-\rho_e,\sigma_c+\rho_e]$ lies strictly between the endpoints
in \eqref{v42:eq:eigen}. The bordered derivative is invertible at the
zero. Consequently the eigenspace is one dimensional and its eigenvector
is not in the range of $\mathcal B-\sigma I$: otherwise that bordered
derivative would have a kernel. This proves algebraic simplicity.
All inequalities are exact rational comparisons and hold for every matrix
in the retained cube enclosure. Verified eigenpair methods have a broader
literature \cite{v42:Rump}; the particular inclusion above is proved by
these displayed residual inequalities.
\end{proof}

\section{What the growing mode does and does not imply}
\label{v42:sec:growth}

Fix the exact root. Define the positive norm
$\|w\|_S=\|S^{-1}w\|_2$ and put $\gamma=17/50$.

\begin{corollary}[Two-sided information about the homogeneous weak evolution]
\label[corollary]{v42:cor:semigroup}
For all $s\ge0$,
\begin{equation}
 \|e^{s\mathcal B}w\|_S\le e^{\gamma s}\|w\|_S.
 \label{v42:eq:semigroup}
\end{equation}
For the verified real eigenvector,
$\|e^{s\mathcal B}r\|_S=e^{\sigma_*s}\|r\|_S$.
There is no fixed positive definite matrix $G$ with
$\mathcal B^TG+G\mathcal B\preceq0$. The eigenvector is nonconstant and
has a nonzero image under the original weak shift-curvature observation
$\mathfrak E_h$.
\end{corollary}
\begin{proof}
For $y=S^{-1}w$, the derivative of $\|y\|_2^2$ is at most
$2\gamma\|y\|_2^2$ by \eqref{v42:eq:lognorm}. Integration proves the
upper bound. The eigenpair proves the equality and gives
$r^T(\mathcal B^TG+G\mathcal B)r=2\sigma_*r^TGr>0$ for $G\succ0$.
A constant shift is annihilated by $\mathsf K$, and hence by
$\mathcal B$, so the nonzero-eigenvalue vector cannot be constant.
The inherited all-frequency observation has exactly the constant shifts
as kernel on the weak space; therefore its image is nonzero.
\end{proof}

This is a spectral statement about the actual limiting weak block. The
selected corrected limiting trajectory has $w=0$ exactly, by
\eqref{v42:eq:Lb}, and a homogeneous growing mode need not be excited by
its next inhomogeneous source. Nor does this calculation construct an
arbitrary nearby exactly constrained initial perturbation. It therefore
is not presented as nonlinear instability of the selected physical
histories, a divergent curvature theorem, or a statement about ghosts.
The signed quantity $r^TDr$ is zero, consistently with
\cref{v42:prop:signed}.

The full linearized matrix in \eqref{v42:eq:full-linear} has the explicit
exponential
\[
 e^{s\mathbb L}=
 \begin{pmatrix}
 I&0\\ \displaystyle\int_0^s e^{(s-u)\mathcal B}\mathcal R\,\dd u
       &e^{s\mathcal B}
 \end{pmatrix}.
\]
Since $\gamma>0$, equivalence of fixed finite-dimensional norms and
\eqref{v42:eq:semigroup} give a constant $M_\gamma<\infty$ such that
\begin{equation}
                  \|e^{s\mathbb L}\|\le M_\gamma e^{\gamma s}
                         \qquad(s\ge0).
 \label{v42:eq:full-semigroup}
\end{equation}
All matrices in this bound are fixed at the same root. No assertion that
$M_\gamma$ is uniform under spatial refinement is made.

\section{Quantitative control on a growing scaled tube}
\label{v42:sec:tube}

The compact-set statement in \cref{v41:lem:layer-field} alone cannot be
used with a growing comparison-time interval. The following estimate
supplies the required extra uniformity. Write $\mathscr F_a$ for the
\emph{original} normalized field in the variables
\[
 X=X_a+a^4\xi,\qquad \lambda=\lambda_a+a^4\zeta,
 \qquad t=a^3s.
\]
All norms below are fixed finite-dimensional norms at $h=1/3$.

\begin{lemma}[A polynomial remainder on the expanding displacement domain]
\label[lemma]{v42:lem:tube}
There exist $a_0,r,C>0$ such that, for $0<a<a_0$ and $a^3\|z\|\le r$,
$\mathscr F_a(z)$ is defined on the original stationary chart and
\begin{align}
 \|\mathscr F_a(z)-b-\mathbb Lz\|
       &\le C\{a(1+\|z\|)+a^3\|z\|^2\},\label{v42:eq:remainder}\\
 \|D_z\mathscr F_a(z)-\mathbb L\|
       &\le C(a+a^3\|z\|),\qquad
 \|D_z^2\mathscr F_a(z)\|\le Ca^3.
 \label{v42:eq:derivative-remainder}
\end{align}
In particular the first line is at most $Ca(1+\|z\|)^2$ after enlarging
$C$ and taking $a\le1$. Signed regularity and positivity of the lapse and
spatial metric hold throughout a sufficiently small such normalized tube.
\end{lemma}
\begin{proof}
Write the analytic initial records as
$X_a=a\bar X(a)$ and $\lambda_a=e+a\bar\chi(a)$.
The displaced normalized arguments are
\[
 \bar X(a)+a^3\xi,\qquad \bar\chi(a)+a^3\zeta.
\]
The joint source grading retained in V41 gives analytic normalized sources
$J_A/a$, $J_W/a^2$ and an analytic normalized signed matrix $\Gamma$ in
these arguments. At $(a,z)=(0,0)$ its determinant is nonzero by
\cref{v41:thm:signed}. A fixed small neighborhood in the normalized
arguments therefore has bounded fixed derivatives of $\Gamma^{-1}$.
The stationary chart, positive lapse and positive metric persist there
for small $a$. This chooses $r,a_0$ independently of $z$ satisfying the
stated bound.

The weak drift has joint order at least two, so
$\mathsf B_W/a^2$ is analytic in these normalized arguments. The active
drift has order at least one, so $\mathsf B_A/a$ is analytic as well.
The original rate equation gives its active and weak scaled components as
\[
 -a^{-2}R_A\Gamma^{-1}
       \binom{\mathsf B_A/a}{\mathsf B_W/a^2},
 \qquad
 -a^{-3}R_W\Gamma^{-1}
       \binom{\mathsf B_A/a}{\mathsf B_W/a^2}.
\]
The canonical scaled component is $F/a$. Each $z$ derivative supplies
a factor $a^3$ through the normalized arguments. Therefore two such
derivatives give bounds $Ca^4$, $Ca^3$ and $Ca^7$ in the active,
weak and canonical components, respectively. They are uniformly bounded
as asserted in \eqref{v42:eq:derivative-remainder}.

At $z=0$ the base rate is exactly $(av_*,0)$ and the base canonical rate
is $aV+O(a^2)$. V41's exact first-derivative expansion gives
$D_z\mathscr F_a(0)=\mathbb L+O(a)$. Thus
$\mathscr F_a(0)=b+O(a)$ and its derivative has the claimed base error.
Taylor's formula on the segment from zero to $z$, using the uniform
second-derivative bound just proved, gives both remaining estimates.
This argument retains the off-constraint term $\mathsf B_1\xi$ and the
whole signed solve. The removable cancellations at the base are precisely
those already established for the exact corrected initializer; no new
vanishing higher coefficient is assumed.
\end{proof}

\section{Original histories on a logarithmically longer interval}
\label{v42:sec:longer}

\begin{theorem}[Logarithmic extension with vanishing normalized curvature change]
\label[theorem]{v42:thm:log-time}
Fix $0<c<50/17$ and put
\begin{equation}
 S_a=c\log(1/a),\qquad T_a=a^3S_a,\qquad\gamma=17/50.
 \label{v42:eq:time}
\end{equation}
For all sufficiently small $a>0$, the corrected initial records have their
unique normalized original-action histories on $0\le t\le T_a$.
Every original constraint is preserved. Uniformly for $0\le s\le S_a$,
\begin{equation}
 \|z(s)-sb\|+\|z'(s)-b\|+\|z''(s)\|
                         \le C_c a e^{\gamma s}.
 \label{v42:eq:long-layer}
\end{equation}
Consequently
\begin{align}
 X(a^3s)-X_a&=a^4sV+O_c(a^5e^{\gamma s}),\notag\\
 \lambda_A(a^3s)-\lambda_{A,a}
               &=a^4sv_*+O_c(a^5e^{\gamma s}),\notag\\
 \lambda_W(a^3s)-\lambda_{W,a}
               &=O_c(a^5e^{\gamma s}),\label{v42:eq:long-displacements}\\
 \lambda_{t,A}(a^3s)-av_*&=O_c(a^2e^{\gamma s}),\qquad
 \lambda_{t,W}(a^3s)=O_c(a^2e^{\gamma s}).
 \label{v42:eq:long-rates}
\end{align}
Both original curvature observations stay $O_c(a)$ and their changes from
the initial cut satisfy
\begin{equation}
 \begin{split}
 \sup_{0\le t\le T_a}\bigl(&\|\Riem(g_h)(t)-\Riem(g_h)(0)\|
             +\|E_h^{\rm link}(t)-E_h^{\rm link}(0)\|\\
             &+\|F_h^{\rm sp,link}(t)-F_h^{\rm sp,link}(0)\|\bigr)
                       =O_c(a^{2-\gamma c})=o(a).
 \end{split}
 \label{v42:eq:long-geometry}
\end{equation}
No coefficient equality between metric and link curvature is inferred.
\end{theorem}
\begin{proof}
Let $e(s)=z(s)-sb$. Equation \eqref{v42:eq:Lb} and variation of constants
give, as long as the solution is in the normalized tube,
\[
 e(s)=\int_0^s e^{(s-u)\mathbb L}
       \{\mathscr F_a(z(u))-b-\mathbb Lz(u)\}\,\dd u.
\]
Bootstrap $\|e(u)\|\le1$. Then $\|z(u)\|\le\|b\|u+1$ and
\cref{v42:lem:tube} bounds the braces by $Ca(1+u)^2$.
Using \eqref{v42:eq:full-semigroup} yields
\begin{equation}
 \|e(s)\|\le Ca e^{\gamma s}
              \int_0^s e^{-\gamma u}(1+u)^2\,\dd u
                    \le C' a e^{\gamma s}.
 \label{v42:eq:bootstrap}
\end{equation}
The last integral is bounded uniformly in $s$, since $\gamma>0$.
On $s\le S_a$ the right side is at most $C'a^{1-\gamma c}$, which
is less than $1/2$ for small $a$. Also
$a^3(\|b\|S_a+1)\to0$, so the growing trajectory stays inside the
normalized tube. A first-exit argument and finite ODE continuation give
existence on the entire interval. This is the extra argument that permits
$S$ to grow; it is not substitution of $S_a$ into a fixed-$S$ estimate.

The equation for $e'$ gives
$\|e'\|\le\|\mathbb L\|\|e\|+Ca(1+s)^2$.
Since $(1+s)^2\le C_\gamma e^{\gamma s}$, the same bound follows.
For the second derivative, use
$z''=D\mathscr F_a(z)z'$, $\mathbb Lb=0$ and
\eqref{v42:eq:derivative-remainder}; the terms are bounded by
$Cae^{\gamma s}$ on this bootstrap tube. This proves
\eqref{v42:eq:long-layer}. Undoing the original displacement and time
scales proves \eqref{v42:eq:long-displacements}--\eqref{v42:eq:long-rates}.
The original signed regularity ensures that this continuation is the
unchanged constrained flow, and its exact equation preserves $P=0$.

For curvature, the canonical and multiplier displacements are
$O(a^4s+a^5e^{\gamma s})$, while the change in the multiplier rate is
$O(a^2e^{\gamma s})$. The canonical field is analytic on the fixed
stationary chart, so its change has the displacement bound. As in
\cref{v41:thm:geometry}, each retained metric or link curvature is a smooth
finite function of $(X,\lambda,F,\lambda_t)$ with bounded derivatives
there; its formula uses no $\lambda_{tt}$. The mean-value bound is therefore
$C(a^4s+a^2e^{\gamma s})$. At $S_a$ this is
$O_c(a^{2-\gamma c})$, since $a^4\log(1/a)$ is smaller.
The initial $O(a)$ bounds and $2-\gamma c>1$ give the size statement.
\end{proof}

\begin{corollary}[One explicit longer interval]
\label[corollary]{v42:cor:particular}
For $t\le a^3\log(1/a)$, the weak multiplier rate and each curvature
change from the initial cut are $O(a^{83/50})$. The scaled trajectory
error is $O(a^{33/50})$.
\end{corollary}
\begin{proof}
Take $c=1$ in \cref{v42:thm:log-time}.
\end{proof}

The factor $\log(1/a)$ tends to infinity relative to the V41 scale, but
$T_a\to0$ in physical time. The upper restriction $c<50/17$ is sufficient
for this proof, not a claim about a maximal lifespan. No numerical value
for the sufficiently small amplitude threshold or the constants in the
analytic initializer is supplied. All of them may depend on the fixed
cutoff and the selected exact initial family.

\section{The first uncomputed excitation is a scalar source coefficient}
\label{v42:sec:excitation}

The growing eigenmode does not settle the next source. The new tube estimate
gives a more precise way to pose that calculation. Analyticity and the
$O(a^3)$ bound on $D_z^2\mathscr F_a$ imply that on each fixed bounded set
\begin{equation}
 \mathscr F_a(z)=b+\mathbb Lz+a(b_1+\mathbb L_1z)+O(a^2).
 \label{v42:eq:first-correction-field}
\end{equation}
The coefficient of $a$ is affine in $z$. The new symbols $b_1,\mathbb L_1$
are derivatives of the actual displaced original field of the selected
initializer, not additional free source parameters.

\begin{proposition}[Exact first-correction excitation formula]
\label[proposition]{v42:prop:excitation}
On every fixed comparison interval, the first amplitude derivative
$z_1=\partial_a z_a|_{a=0}$ exists and satisfies
\[
 z_1'=\mathbb Lz_1+b_1+s\mathbb L_1b,\qquad z_1(0)=0.
\]
Let $\ell_u$ be any nonzero real left eigenvector of $\mathbb L$ for the
verified $\sigma_*>0$. Put
\[
 \beta_0=\ell_ub_1,\qquad \beta_1=\ell_u\mathbb L_1b,
 \qquad \kappa_u=\frac{\beta_0}{\sigma_*}
                           +\frac{\beta_1}{\sigma_*^2}.
\]
Then
\begin{equation}
 \boxed{\ell_u z_1(s)=\kappa_u(e^{\sigma_*s}-1)
                                  -\frac{\beta_1}{\sigma_*}s.}
 \label{v42:eq:excitation}
\end{equation}
In particular $\kappa_u=0$ is necessary and sufficient to remove the
exponential term from this particular first-correction projection.
\end{proposition}
\begin{proof}
The coefficient of $a$ in the analytic field is affine because its second
$z$ derivative vanishes. Analytic parameter dependence of the finite ODE
on a fixed interval permits differentiation, yielding the displayed
inhomogeneous linear equation. The positive eigenvalue is simple for
$\mathbb L$ as well: its other diagonal block has only eigenvalue zero.
A left eigenvector exists; explicitly if $r_u^T\mathcal B=\sigma_*r_u^T$,
one may take $\ell_u=(r_u^T\mathcal R/\sigma_*,r_u^T)$.
Its projection obeys $\chi'=\sigma_*\chi+\beta_0+s\beta_1$ with zero
initial value. Integration gives \eqref{v42:eq:excitation}.
\end{proof}

Neither $\beta_0,\beta_1$ nor $\kappa_u$ has been evaluated for the actual
higher initializer in this body. The unknown higher initial coefficients
cannot be replaced by zero merely because the leading weak layer cancels.
Moreover, the first-correction expansion is proved on fixed comparison
intervals; it is not used as a lower bound at $s\sim\log(1/a)$ without
additional uniform remainder control. Even a verified $\kappa_u=0$ would
not remove other growing modes or higher-order forcing. The proposition
specifies a smaller next actual-source computation, not an assumed
stability condition or a new physical growth claim.

\section{Proved conclusions, retained inputs, and the remaining boundary}
\label{v42:sec:status}

\paragraph{Proved in this body.}
The weak signed generator has a simple real growing mode at the unchanged
algebraic root, with the stated rational enclosure. A fixed positive
comparison norm bounds its semigroup by $e^{17s/50}$. The original scaled
field has a polynomial remainder on a growing normalized tube. These
results extend the corrected original histories to
$t\le c a^3\log(1/a)$ for every $c<50/17$, with vanishing curvature change
after division by $a$. The first possible excitation of the certified
mode has the exact scalar formula \eqref{v42:eq:excitation}.

\paragraph{Inherited inputs.}
The V39 root and old Poisson/cubic bank, the V40 exact corrected initializer
and leading-source cancellation, and V41's full signed source, normalized
continuation, and metric/link observation formulas retain their stated
scopes. The two packaged verification stages replay the root inclusion,
recompute the complete $17334$-entry stationary-source enclosure on the smaller
cube, and independently recompute the $5046$ restricted Poisson entries. Both
stages were executed separately. The stored inverse, similarity, and eigenpair proposals are
validated by exact residuals, not accepted from floating calculations.
The older complete $79224$-entry root bank is replayed, not newly regenerated.

\paragraph{Still unresolved.}
The actual higher-order excitation coefficients, longer original-time
control beyond the logarithmic interval, a common positive physical
lifespan, spatial-refinement uniformity, and selection of these preparations
by the primitive source remain unproved. The verified growing homogeneous
mode is not proof that the selected exactly constrained family becomes
nonlinearly unstable. Conversely, its exactly zero leading response is not
proof of attraction or damping. The unit-initial-multiplier stress problem
of Version 32, dynamical matter backreaction and the full classical
Einstein--Standard-Model limit remain separate.

\paragraph{Verification and preservation.}
The exact verifier is standalone with Python and NumPy. It checks the
new matrix and eigenpair inequalities from rational data and original
coefficient polarization. The larger-time and curvature conclusions are
written analytic proofs, not a numerical spacetime integration. The
verification package records which regeneration stages were executed,
input/output hashes, and the preserved V41 prefix. New labels have prefix
\texttt{v42:}. No proof-assistant formalization is asserted.

\begingroup
\renewcommand{\refname}{Additional source for Version 42}
\begin{thebibliography}{9}
\small\raggedright
\bibitem{v42:Rump}
S.~M.~Rump, \emph{Verified error bounds for all eigenvalues and eigenvectors
of a matrix}, SIAM Journal on Matrix Analysis and Applications
\textbf{43} (2022), 1736--1754.
\url{https://doi.org/10.1137/21M1451440}.
Primary-source context for verified eigenpair bounds. The normalized real
contraction, the particular source matrices and their enclosures are
established in this body rather than imported from that work.
\end{thebibliography}
\endgroup
\addtocontents{toc}{\protect\endgroup}
% END IMMUTABLE VERSION 42 BODY
% APPEND VERSION 43 HERE, BEFORE THE FINAL END-DOCUMENT COMMAND.


% BEGIN IMMUTABLE VERSION 43 BODY
\clearpage
\begin{center}
{\Large\bfseries Version 43: Second-order preparation of the growing spectral sector}\\[0.5em]
{\large A verified nine-dimensional splitting, exact initial-rate access, and a longer logarithmic interval}
\end{center}
\makeatletter
\addtocontents{toc}{\protect\begingroup}
\addtocontents{toc}{\protect\makeatletter}
\addtocontents{toc}{\protect\def\protect\l@section{\protect\bfseries\protect\@dottedtocline{1}{0em}{2.5em}}}
\makeatother
\addcontentsline{toc}{part}{Version 43: Second-order preparation of the growing spectral sector}

\section{Changing initial selection rather than assuming a favorable source}
\label{v43:sec:context}

Work at the same cutoff $h=1/3$ and the same algebraic mixed-mode root as
\cref{v39:thm:root}. Fix one analytic corrected initializer supplied by
\cref{v40:thm:initial}. Its higher coefficients are well-defined by that
chosen analytic construction, but have not been tabulated. The exact
signed certificate of \cref{v41:thm:signed} and the actual weak matrix
$\mathcal B$ of \eqref{v42:eq:generator} are retained. No source matrix or
physical time normalization is changed.

The first excitation calculation in \cref{v42:prop:excitation} leaves an
uncomputed coefficient for the selected initializer. This body does not
assign it zero. Instead it proves that a one-time, second-order weak
\emph{initial-rate} coefficient can set the first growing response. The
initial records are then solved anew, so that this rate is selected by
the unique original-action continuation. It is not imposed during an
already evolving history.

There are three new steps. A rational block enclosure separates nine
rapidly growing directions from the other twenty weak directions. The
full exact initial consistency problem is shown to admit arbitrary fixed
second-order weak rates, with a triangular effect on the first displaced
field coefficient. That effect permits cancellation of the first
exponential response of the whole nine-dimensional sector, not only the
one real eigenvector isolated in V42. An error estimate uniform on growing
comparison intervals then enlarges the sufficient original-time range to
$t<c a^3\log(1/a)$ for every $c<100/17$.

The values of the old excitation coefficients, the new physical rate
coefficient, and its numerical norm are not evaluated here. The conclusion
is nevertheless an existence theorem, not a hypothesis that those
coefficients happen to vanish: the exact rate-access argument and an
explicit linear selection formula supply a nonempty new family. All
amplitude thresholds may depend on that finite selected coefficient.

\section{A certified nine-dimensional growing spectral sector}
\label{v43:sec:splitting}

Let $B_c,S,S_I$ be the fixed rational proposals replayed in V42. Its
original-source enclosure on the root cube $\mathcal Q_{42}$ proves
\[
 \|\mathcal B-B_c\|_\infty<2^{-50},\qquad
 e_S:=\|I-S_IS\|_\infty<2^{-41}.
\]
The matrix $S$ is invertible. Put $T_c=S_IB_cS$ and
\begin{equation}
 e_T=\frac{e_S}{1-e_S}\|T_c\|_\infty+
       \frac{\|S_I\|_\infty}{1-e_S}\,2^{-50}\|S\|_\infty.
 \label{v43:eq:similarity-error}
\end{equation}
Then $\|S^{-1}\mathcal BS-T_c\|_\infty\le e_T$. In the column ordering
of that \emph{fixed} $S$, retain the nine indices
\begin{equation}
 \mathcal I_{\rm f}=(0,2,3,9,10,11,16,17,20),
 \qquad \mathcal I_{\rm c}=\{0,\ldots,28\}\setminus\mathcal I_{\rm f}.
 \label{v43:eq:indices}
\end{equation}
Indices are zero based. After this permutation write the true similarity as
\begin{equation}
 T=\begin{pmatrix}A&E\\F&C\end{pmatrix},\qquad
 A\in\R^{9\times9},\quad C\in\R^{20\times20}.
 \label{v43:eq:blocks}
\end{equation}
The letters in this block split do not denote the active or signed Schur
blocks of V41.

\begin{lemma}[Exact block inequalities on the entire root cube]
\label[lemma]{v43:lem:blocks}
For every matrix in the retained root-cube enclosure,
\begin{equation}
 \boxed{
 \operatorname{sym}A\succeq\frac1{60}I,
 \quad \operatorname{sym}C\preceq10^{-6}I,
 \quad \|E\|_F,\|F\|_F<2^{-28}.}
 \label{v43:eq:block-bounds}
\end{equation}
Here $\operatorname{sym}A=(A+A^T)/2$ and $\|\cdot\|_F$ is the
Frobenius norm. No eigenvalue computed in floating arithmetic enters
these inequalities.
\end{lemma}
\begin{proof}
Every block error in Frobenius norm is at most $29e_T$: indeed
$\|Z\|_F\le\sqrt{29}\|Z\|_\infty\le29\|Z\|_\infty$ for a
$29$-row matrix, and restriction only decreases that bound. Compute the
Gershgorin lower bound of $\operatorname{sym}(T_c)_{\rm ff}$ and the
Gershgorin upper bound of $\operatorname{sym}(T_c)_{\rm cc}$, then
subtract or add $29e_T$. For each cross block use
$\|Z\|_F\le14\max_{ij}|Z_{ij}|$, since $9\cdot20<14^2$, and add
$29e_T$. The packaged rational comparisons give a fast lower bound greater
than $1/60$, a complementary upper bound less than $10^{-6}$, and cross
bounds less than $2^{-28}$. For scale only, their decimal displays are
approximately $0.018016348775$, $3.10\,10^{-11}$ and
$3.11\,10^{-11}$, respectively. The first, less restrictive rational
inequalities are the certificate. The parent replay regenerates the
$17334$ stationary-source entries and verifies the $2^{-50}$ input
before these new comparisons are made.
\end{proof}

\begin{theorem}[Invariant splitting with controlled forward and backward growth]
\label[theorem]{v43:thm:split}
At the exact root there is a real invariant decomposition
\begin{equation}
 \R^{29}=\mathcal E_{\rm f}\oplus\mathcal E_{\rm c},\qquad
 \dim\mathcal E_{\rm f}=9,
 \quad\dim\mathcal E_{\rm c}=20,
 \label{v43:eq:decomposition}
\end{equation}
with a commuting real projector $\Pi_{\rm f}$. In fixed positive norms on
the two spaces, the restrictions satisfy
\begin{equation}
 \boxed{
 \|e^{-s\mathcal B_{\rm f}}\|\le e^{-s/61},\qquad
 \|e^{s\mathcal B_{\rm c}}\|\le e^{s/1000},\qquad s\ge0.}
 \label{v43:eq:split-growth}
\end{equation}
In particular $\mathcal B_{\rm f}$ is invertible and
$\|\mathcal B_{\rm f}^{-1}\|\le61$. Exactly nine eigenvalues, with
algebraic multiplicity, have real part greater than $1/1000$. This is not
an assertion that every remaining eigenvalue has nonpositive real part.
The splitting depends analytically on the finite matrix throughout a
sufficient neighborhood of the root cube.
\end{theorem}
\begin{proof}
Put $\alpha=1/60$, $\delta=10^{-6}$, $\epsilon=2^{-28}$ and
$r=2^{-16}$. The map $X\mapsto AX-XC$ on $9$ by $20$ matrices has
inverse
\[
 Z\longmapsto\int_0^\infty e^{-At}Z e^{Ct}\,\dd t,
\]
whose Frobenius norm is at most $(\alpha-\delta)^{-1}<64$. This follows
by differentiating Euclidean squared norms using \eqref{v43:eq:block-bounds}.
An invariant complementary graph $(Xv,v)$ requires
\[
 AX-XC=-E+XFX.
\]
Its integral fixed-point map preserves $\|X\|_F\le r$ and is a strict
contraction, because the exact inequalities are
\begin{equation}
 64\epsilon(1+r^2)<r,\qquad 128\epsilon r<\frac12.
 \label{v43:eq:riccati-bounds}
\end{equation}
Likewise the graph $(u,Yu)$ is invariant precisely when
$CY-YA=-F+YEY$. The inverse Sylvester map has the same bound, using
$-\int_0^\infty e^{Ct}Ze^{-At}\,\dd t$. The identical contraction
argument gives $\|Y\|_F\le r$.

Since $\|XY\|<1$, the two graphs are complementary. The invertible
matrix $\left(\begin{smallmatrix}I&X\\Y&I\end{smallmatrix}\right)$
conjugates $T$ to the diagonal blocks $A+EY$ and $C+FX$. Their symmetric
parts obey
\[
 \operatorname{sym}(A+EY)\succeq(\alpha-\epsilon r)I\succ\frac1{61}I,
 \qquad
 \operatorname{sym}(C+FX)\preceq(\delta+\epsilon r)I\prec\frac1{1000}I.
\]
Pull the Euclidean norms on graph coordinates back through $S$ to obtain
\eqref{v43:eq:split-growth}. These inequalities also locate the real
parts of the two block spectra and give the inverse bound by integrating
the backward fast semigroup. The projector onto the first graph along
the second commutes with the generator. The contraction equations and
their uniformly invertible derivatives give analytic dependence.
\end{proof}

Invariant graph/Riccati methods are established operator theory; see, for
example, \cite{v43:Wyss}. The particular rational separation, the count
nine, and the constants above are obtained from the retained source, not
from that reference. The positive comparison norms in
\eqref{v43:eq:split-growth} do not replace the original signed action.

\section{An exact second-order weak rate is available by initial selection}
\label{v43:sec:rate-access}

Let $I_W:\R^{29}\to\R^{108}$ be the retained physical weak multiplier
synthesis. It includes all three constant shifts and has zero lapse. Let
$J_Ww=(0,0,w)$ denote its coordinate insertion in the $431$-dimensional
scaled state $z=(\xi,\zeta_A,\zeta_W)$. A constant shift rate is allowed:
only the mean lapse is fixed throughout the normalized evolution. The
initial shift means remain zero as in V40.

\begin{theorem}[Exact original initial consistency with any fixed weak second rate]
\label[theorem]{v43:thm:rate-access}
For every fixed finite $w\in\R^{29}$, there is an analytic initial family
$(X_a^w,\lambda_a^w)$, for sufficiently small amplitude, satisfying
\begin{equation}
 \boxed{
 P(X_a^w,\lambda_a^w)=0,\qquad
 \mathsf B(X_a^w,\lambda_a^w)+
       M(X_a^w,\lambda_a^w)(av_*+a^2I_Ww)=0.}
 \label{v43:eq:exact-initial}
\end{equation}
It has mean lapse one and initial shift means zero, positive lapse and
spatial metric, and the same canonical and multiplier coefficients
through order $a^2$ as the chosen $w=0$ family:
\begin{equation}
 X_a^w-X_a^0=O_w(a^3),\qquad
 \lambda_a^w-\lambda_a^0=O_w(a^3).
 \label{v43:eq:initial-comparison}
\end{equation}
The actual unique normalized continuation satisfies
\begin{equation}
                 \lambda_t(0)=av_*+a^2I_Ww.
 \label{v43:eq:actual-rate}
\end{equation}
The construction is jointly analytic locally in $w$; uniformity is only
claimed on bounded sufficiently small neighborhoods of a fixed finite $w$.
\end{theorem}
\begin{proof}
Use the same exact constraint initializer as
\cref{v39:lem:analytic-prep,v40:thm:initial}, with finite parameters
$(q,m)$ and multiplier $e+a\ell(q)+a^2m$. On these exactly constrained
records impose the $107$ nonredundant rows of
$\mathsf B+M(av_*+a^2I_Ww)$, divided by $a^3$.
The extra term has orders
\begin{equation}
 a^2MI_Ww=
       \binom{O_w(a^5)}{O_w(a^6)}
       \quad\hbox{on active and weak rows.}
 \label{v43:eq:extra-orders}
\end{equation}
This follows from the full source grading
$M_{AW}=O(a^3)$, $M_{WW}=O(a^4)$, before any inverse is applied.
Consequently the divided analytic system differs from V40's system by
$a^2$ times an analytic function; its three constant-shift rows differ
only at order $a^3$. Its derivative at $a=0$ in $(q,m)$ remains exactly
V40's verified invertible block-triangular matrix. The analytic
implicit-function theorem supplies $q(a,w),m(a,w)$, and uniqueness of
its local analytic branch gives
\[
 (q(a,w),m(a,w))-(q(a,0),m(a,0))=O_w(a^2).
\]
One may see this divisibility by differentiating the implicit equation in
$w$, whose inhomogeneous term has a factor $a^2$, and integrating on a
fixed bounded parameter segment. The initial map has $X=0,\lambda=e$
at $a=0$ for every $(q,m)$, so its parameter derivatives are $O(a)$.
This proves \eqref{v43:eq:initial-comparison}.

The last constant-lapse row follows without an additional equation:
$M\lambda=0$, $\langle\lambda,\mathsf B\rangle_h=0$ and mean lapse
one force its residual to vanish. The full leading signed matrix is the
unchanged matrix certified in V41. Thus, for sufficiently small nonzero
$a$, its normalized restriction is invertible and the unique original
rate equation gives \eqref{v43:eq:actual-rate}. Positivity and the
original stationary chart follow by the same leading expansion. The
sufficient amplitude depends on $w$, but existence for every fixed finite
$w$ does not require a bound uniform as $\|w\|\to\infty$.
\end{proof}

This is an initial-data theorem, not a feedback rule. The second rate is
chosen once; subsequent rates are obtained from $M\lambda_t=-\mathsf B$.
In particular no requirement of zero shift mean is imposed at later times,
and no constant weak mode is discarded.

\section{Triangular access to the first displaced source coefficient}
\label{v43:sec:triangular}

For each new initializer use its own initial cut in the original variables
\[
 t=a^3s,\qquad X=X_a^w+a^4\xi,\qquad
 \lambda=\lambda_a^w+a^4\zeta.
\]
Call the resulting exact scaled field $\mathscr F_a^w(z)$. For $w=0$,
write the coefficients already defined in V42 as
\[
 \mathscr F_a^0(z)=b+\mathbb Lz+a(b_1+\mathbb L_1z)+O(a^2).
\]
The vectors $b_1$ and $\mathbb L_1$ in this body always refer to this
\emph{fixed unmodified} initializer, including its actual unlisted
higher coefficients.

\begin{lemma}[The rate changes the affine source but not its first linear operator]
\label[lemma]{v43:lem:triangular}
On every fixed bounded set, and with fixed derivatives,
\begin{equation}
 \boxed{
 \mathscr F_a^w(z)=b+\mathbb Lz+
          a(b_1+J_Ww+\mathbb L_1z)+O_w(a^2).}
 \label{v43:eq:affine-access}
\end{equation}
In particular the coefficient $\mathbb L_1$ is unchanged by $w$.
\end{lemma}
\begin{proof}
By \eqref{v43:eq:initial-comparison}, the normalized base arguments
$(X_a^w/a,(\lambda_a^w-e)/a)$ differ from their $w=0$ values by
$O_w(a^2)$. In the analytic normalized inverse used in
\cref{v42:lem:tube}, a first $z$ derivative cancels the worst $a^{-3}$
rate prefactor. Thus $D_z\mathscr F_a^w(0)-D_z\mathscr F_a^0(0)=O_w(a^2)$.
The canonical component at zero changes by $O_w(a^2)$, since $F$ is
analytic and is divided by $a$. The two rate components at zero are
known \emph{exactly} from \eqref{v43:eq:actual-rate}: their changes
are zero in the active component and $aw$ in the weak component.
Finally $D_z^2\mathscr F_a^w=O_w(a^3)$ on the same normalized tube as
in V42. Taylor expansion in $z$ proves \eqref{v43:eq:affine-access}.
The initial-rate identity is needed here: a raw estimate of the singular
inverse on the change of initial records would lose this triangular
structure.
\end{proof}

This resolves a possible circular dependence. The coefficient selected
below is computed from $b_1,\mathbb L_1$ of the old initializer; changing
the initializer by that rate leaves $\mathbb L_1$ and the other order-$a$
source components unchanged. The required cancellation is therefore a
linear equation, not an unevaluated fixed-point condition on an unknown
new excitation coefficient.

\section{One initial selection removes the first response of all nine fast directions}
\label{v43:sec:cancel}

In \eqref{v42:eq:full-linear}, write $z=(u,w)$ with
$u\in\R^{402}$. Since the upper block of $\mathbb L$ is zero,
all its generalized eigenvectors for nonzero eigenvalues have zero $u$
component. The splitting in \cref{v43:thm:split} lifts to the full
state through the commuting projector
\begin{equation}
 \mathcal P_{\rm f}(u,w)=
 J_W\left(\Pi_{\rm f}w+
       \mathcal B_{\rm f}^{-1}\Pi_{\rm f}\mathcal Ru\right).
 \label{v43:eq:full-projector}
\end{equation}
Its range is $J_W\mathcal E_{\rm f}$, of dimension nine. Denote the
restriction of $\mathbb L$ to this range by $\mathbb L_{\rm f}$.
The inverse in \eqref{v43:eq:full-projector} only acts on the nine-dimensional
fast space; no inverse of a zero-frequency complementary block is used.

\begin{lemma}[Access and growth for the full linearized system]
\label[lemma]{v43:lem:full-split}
One has
\[
 \mathcal P_{\rm f}^2=\mathcal P_{\rm f},\quad
 [\mathcal P_{\rm f},\mathbb L]=0,\quad
 \mathcal P_{\rm f}J_W=J_W\Pi_{\rm f}.
\]
For $\nu=1/1000$ and a fixed constant $C_\nu$,
\begin{equation}
 \|e^{s\mathbb L}(I-\mathcal P_{\rm f})\|
       \le C_\nu e^{\nu s},\qquad s\ge0.
 \label{v43:eq:full-slow}
\end{equation}
\end{lemma}
\begin{proof}
Substitution in the block formula for $\mathbb L$ proves the three
identities. On the kernel of $\mathcal P_{\rm f}$, the fast weak component
is $-\mathcal B_{\rm f}^{-1}\Pi_{\rm f}\mathcal Ru$, while $u'=0$.
The complementary weak component solves
$w_c'=\mathcal B_cw_c+(I-\Pi_{\rm f})\mathcal Ru$.
Use \eqref{v43:eq:split-growth} and
$\int_0^s e^{\nu(s-r)}\,\dd r\le\nu^{-1}e^{\nu s}$.
Equivalence of the fixed graph norms proves \eqref{v43:eq:full-slow}.
\end{proof}

Define the unique $w^\sharp\in\mathcal E_{\rm f}$ by
\begin{equation}
 \boxed{
 J_Ww^\sharp=-\mathcal P_{\rm f}b_1
       -\mathbb L_{\rm f}^{-1}\mathcal P_{\rm f}\mathbb L_1b.}
 \label{v43:eq:wsharp}
\end{equation}
Both terms lie in $J_W\mathcal E_{\rm f}$, so this is well defined and
finite. Its norm is bounded by a fixed constant times
$\|b_1\|+\|\mathbb L_1b\|$. No numerical values for these two original
higher-source quantities are supplied.

\begin{theorem}[Exact preparation eliminating the first fast exponential response]
\label[theorem]{v43:thm:cancellation}
Take the exact initial family of \cref{v43:thm:rate-access} with
$w=w^\sharp$. Its first amplitude-correction profile $z_1$ solves
\begin{equation}
 z_1'=\mathbb Lz_1+b_1+J_Ww^\sharp+s\mathbb L_1b,
 \qquad z_1(0)=0,
 \label{v43:eq:first-profile}
\end{equation}
and has the exact fast projection
\begin{equation}
 \boxed{
 \mathcal P_{\rm f}z_1(s)=
       -s\mathbb L_{\rm f}^{-1}\mathcal P_{\rm f}\mathbb L_1b.}
 \label{v43:eq:fast-profile}
\end{equation}
In particular it has no $e^{s\mathbb L_{\rm f}}$ term. For $j=0,1,2$,
\begin{equation}
            \|\partial_s^jz_1(s)\|\le C e^{\nu s},
            \qquad s\ge0,\quad\nu=1/1000.
 \label{v43:eq:profile-bound}
\end{equation}
\end{theorem}
\begin{proof}
The first equation follows by differentiating the analytic initial-value
problem on fixed comparison intervals and using
\cref{v43:lem:triangular}, $\mathbb Lb=0$ and $z_0=sb$.
Projection and \eqref{v43:eq:wsharp} give an inhomogeneous fast equation
whose affine particular solution is exactly \eqref{v43:eq:fast-profile};
it also has the required zero initial value. Uniqueness proves the identity
for all $s\ge0$ as a statement about the finite linear profile equation.
On the complement, \eqref{v43:eq:full-slow} bounds the convolution with
$b_1+s\mathbb L_1b$ by
\[
 C e^{\nu s}\int_0^s e^{-\nu r}(1+r)\,\dd r\le C' e^{\nu s}.
\]
The fast affine profile has the same bound. Its first two derivatives are
controlled from the equation and its derivative. This does not extend a
formal amplitude expansion to large $s$ without a remainder estimate;
that estimate is proved next.
\end{proof}

The rate in \eqref{v43:eq:wsharp} cancels an exponential \emph{response},
not the spectrum of the action. The nine eigenvalues remain. The fast
profile is generally linear rather than zero, and the remaining twenty
weak directions are not asserted to be damped. This construction also
does not decide whether V42's unmodified initializer excites any growing
mode.

\section{A second-order remainder on a growing comparison interval}
\label{v43:sec:remainder}

Fix the selected finite $w^\sharp$ and write
$\widetilde b_1=b_1+J_Ww^\sharp$.
The exact same normalized-argument proof as V42, now with the initial
family of \cref{v43:thm:rate-access}, gives constants $a_0,r,C>0$ such
that on $a^3\|z\|\le r$,
\begin{align}
 \mathscr F_a^\sharp(z)&=b+\mathbb Lz+
                a(\widetilde b_1+\mathbb L_1z)+\mathscr Q_a(z),
                \label{v43:eq:second-expansion}\\
 \|\mathscr Q_a(z)\|&\le C\{a^2(1+\|z\|)+a^3\|z\|^2\},
                \label{v43:eq:Q}\\
 \|D_z\mathscr Q_a(z)\|&\le C(a^2+a^3\|z\|),\qquad
 \|D_z^2\mathscr F_a^\sharp(z)\|\le Ca^3.
                \label{v43:eq:DQ}
\end{align}
These inequalities follow by a second-order analytic expansion at $z=0$
in $a$, \cref{v43:lem:triangular}, and Taylor's formula in $z$ with the
uniform second-derivative bound. All expansions concern the original
inverse and original constraint source; they do not discard higher action
coefficients.

Let the comparison profile be
\begin{equation}
                   z_{\rm app}(s)=sb+az_1(s).
 \label{v43:eq:app}
\end{equation}
It is used only in estimates and never supplied to the source as an update.

\begin{lemma}[Uniform profile defect]
\label[lemma]{v43:lem:defect}
For $s\le c\log(1/a)$ with fixed $c<100/17$, all sufficiently small
$a$ satisfy
\begin{equation}
 \|\mathscr F_a^\sharp(z_{\rm app}(s))-z_{\rm app}'(s)\|
                 \le C_c a^2e^{\nu s}.
 \label{v43:eq:defect}
\end{equation}
\end{lemma}
\begin{proof}
The equations for $sb,z_1$ cancel the coefficients through order $a$.
The remaining defect is $a^2\mathbb L_1z_1+\mathscr Q_a(z_{\rm app})$.
Since $\nu c<1$, \eqref{v43:eq:profile-bound} gives
$a\|z_1(s)\|\le C a^{1-\nu c}=o(1)$ on this whole interval.
Thus $\|z_{\rm app}(s)\|\le C(1+s)$ and it lies in the normalized
tube. Use \eqref{v43:eq:Q} and
$(1+s)^2\le C_\nu e^{\nu s}$ to prove the estimate.
\end{proof}

\begin{theorem}[A second-order selected original family on a larger logarithmic interval]
\label[theorem]{v43:thm:longer}
Let $\gamma=17/50$ and $\nu=1/1000$. For every fixed
\begin{equation}
            0<c<100/17,\qquad S_a=c\log(1/a),\qquad T_a=a^3S_a,
 \label{v43:eq:time}
\end{equation}
the unique normalized original-action histories from the selected initial
family exist on $[0,T_a]$ for all sufficiently small positive $a$.
They preserve every original constraint. Their scaled displacements obey
\begin{equation}
 \boxed{
 \|z_a-z_{\rm app}\|+\|z_a'-z_{\rm app}'\|
                  +\|z_a''-z_{\rm app}''\|
          \le C_c a^2 e^{\gamma s},\qquad 0\le s\le S_a.}
 \label{v43:eq:comparison}
\end{equation}
Consequently
\begin{equation}
 \|z_a-sb\|+\|z_a'-b\|+\|z_a''\|
       \le C_c\{a e^{\nu s}+a^2e^{\gamma s}\}=o(1)
 \label{v43:eq:total-error}
\end{equation}
uniformly on the stated interval.
\end{theorem}
\begin{proof}
Put $e=z_a-z_{\rm app}$. Under the bootstrap $\|e\|\le1$, the
trajectory has $\|z_a\|\le C(1+s)$ and stays in the normalized tube.
The inherited full semigroup bound
$\|e^{s\mathbb L}\|\le M_\gamma e^{\gamma s}$ and
\eqref{v43:eq:DQ} give
\[
 \|e(s)\|\le M_\gamma\int_0^s e^{\gamma(s-u)}
 \{C[a+a^3(1+u)]\|e(u)\|+C_ca^2e^{\nu u}\}\,\dd u.
\]
Multiply by $e^{-\gamma s}$ and use Gronwall's inequality. Since
$\nu<\gamma$, the integral of $e^{-(\gamma-\nu)u}$ is bounded, while
the additional exponential factor is at most
$\exp(C[aS_a+a^3(S_a+S_a^2)])=O(1)$. Hence
$\|e(s)\|\le C_ca^2e^{\gamma s}$. At $S_a$ this is
$C_ca^{2-\gamma c}=o(1)$, so the bootstrap cannot fail. Also
$a^3(1+S_a)\to0$, which prevents exit from the original signed regular
chart or positive lapse/metric neighborhood. Finite-dimensional
continuation gives the whole asserted interval.

The error equation gives the same bound for $e'$. For $e''$, first
apply the chain rule to $z_a''=D\mathscr F_a^\sharp(z_a)z_a'$.
Equations \eqref{v43:eq:second-expansion}--\eqref{v43:eq:DQ} and the
profile equation show
\[
 \|D\mathscr F_a^\sharp(z_{\rm app})z_{\rm app}'-z_{\rm app}''\|
            \le C_ca^2e^{\nu s}.
\]
Here the order-$a$ terms are $\mathbb Lz_1'+\mathbb L_1b=z_1''$;
the remaining terms are bounded using $a e^{\nu S_a}=o(1)$ and
$\|z_{\rm app}\|=O(1+s)$. The first two derivative bounds already
proved, together with the uniform $D^2\mathscr F_a^\sharp$ estimate,
then give the claimed $e''$ bound. This avoids differentiating an
uncontrolled pointwise remainder in time. Finally combine with
\eqref{v43:eq:profile-bound} to obtain \eqref{v43:eq:total-error}.
The signed inverse and exact initial constraints identify this continuation
with the original action flow, whose equation $\dot P=0$ preserves $P$.
\end{proof}

\section{Physical rate, curvature, and preparation accuracy}
\label{v43:sec:observations}

\begin{corollary}[The original observations retain their leading small-amplitude values]
\label[corollary]{v43:cor:physical}
On the interval \eqref{v43:eq:time},
\begin{equation}
 \boxed{
 \|\lambda_{t,W}\|+\|\lambda_{t,A}-av_A\|
 \le C_c\{a^2e^{\nu s}+a^3e^{\gamma s}\}.}
 \label{v43:eq:rates}
\end{equation}
Each of the retained metric Riemann tensor, literal electric link curvature,
and literal spatial link curvature has size $O_c(a)$. If $\mathscr O_a(t)$
denotes any one of those observations, with its unchanged definition, then
\begin{equation}
 \boxed{
 \sup_{t\le T_a}\|\mathscr O_a(t)-\mathscr O_a(0)\|
       \le C_c\{a^{2-\nu c}+a^{3-\gamma c}\}=o(a).}
 \label{v43:eq:curvature}
\end{equation}
Here $\mathscr O_a(0)$ is the initial observation of the \emph{new} family.
The new and old initial metric and link curvatures differ by $O(a^2)$.
\end{corollary}
\begin{proof}
The physical rate equals $a\zeta'$ and the displacement equals $a^4z$.
Use \eqref{v43:eq:total-error}; $b_W=0$ gives \eqref{v43:eq:rates}.
Changes from the initial rate have the same bound after subtracting the
fixed $a^2I_Ww^\sharp$. The original curvature formulas are smooth
functions of $(X,\lambda,F,\lambda_t)$ on the retained chart, with
bounded fixed-cutoff derivatives, and involve no $\lambda_{tt}$.
The state displacement is bounded by
$C(a^4s+a^5e^{\nu s}+a^6e^{\gamma s})$; the rate change dominates it
on the stated logarithmic interval. This proves \eqref{v43:eq:curvature}.
Both exponents exceed one by \eqref{v43:eq:time}. At the initial cut,
\eqref{v43:eq:initial-comparison} and the rate difference $a^2I_Ww^\sharp$
give the $O(a^2)$ old/new observation difference. Initial size $O(a)$ and
the change estimate complete the proof.
\end{proof}

\begin{corollary}[An explicit interval beyond the previous coefficient bound]
\label[corollary]{v43:cor:example}
One may take
\begin{equation}
              \boxed{T_a=5a^3\log(1/a).}
 \label{v43:eq:example}
\end{equation}
The weak rate and all the displayed curvature changes are then
$O(a^{13/10})$, and $z_a-sb$ is $O(a^{3/10})$.
\end{corollary}
\begin{proof}
For $c=5$, $2-\nu c=399/200$, $3-\gamma c=13/10$,
$1-\nu c=199/200$, and $2-\gamma c=3/10$. The smaller exponent in
each sum determines the stated bound. Also $50/17<5<100/17$.
\end{proof}

The larger coefficient is a sufficient existence statement for a differently
selected initial family, not a maximal-lifespan theorem. In particular
$T_a\to0$ in physical time. The spectra have not been changed, and no
linear or nonlinear attracting basin has been proved.

\begin{proposition}[An error in the selected rate has a known first-order cost]
\label[proposition]{v43:prop:sensitivity}
Replace $w^\sharp$ by $w^\sharp+d$, with fixed finite $d$. The difference
between the two first profiles has fast projection
\begin{equation}
 \boxed{
 \mathcal P_{\rm f}(z_1^{\sharp+d}-z_1^\sharp)
       =\mathbb L_{\rm f}^{-1}
          (e^{s\mathbb L_{\rm f}}-I)J_W\Pi_{\rm f}d.}
 \label{v43:eq:sensitivity}
\end{equation}
An analytic amplitude-dependent implementation with $d(a)=O(a)$ retains
the sufficient interval and exponents of \cref{v43:thm:longer}.
\end{proposition}
\begin{proof}
Subtract the two linear profile equations and use
$\mathcal P_{\rm f}J_W=J_W\Pi_{\rm f}$. Integration of the constant
forcing gives \eqref{v43:eq:sensitivity}. When $d(a)=O(a)$, the added
term in the scaled field is $aJ_Wd(a)=O(a^2)$; the changed normalized
initial arguments have still higher order by the rate-access proof.
The same second-order remainder and bootstrap therefore apply.
\end{proof}

Thus the exact selection formula is not interchangeable with an arbitrary
fixed numerical approximation. Evaluating $b_1,\mathbb L_1b$ and
certifying the resulting rate would be necessary for a numerical
implementation with a quantified error. That calculation is distinct from
the existence proof given here.

\section{Proved conclusions, inherited inputs, and unresolved questions}
\label{v43:sec:status}

\paragraph{Proved in this body.}
The actual weak matrix admits a nine-dimensional growing sector and a
complementary growth bound $e^{s/1000}$, certified by rational block
inequalities and a proved invariant-graph construction. Every fixed
second-order weak rate can be realized by exact original constrained
initial records without changing their first two amplitude coefficients.
Its first displaced-source effect is a constant weak insertion, with the
first linear operator unchanged. The rate \eqref{v43:eq:wsharp}
therefore removes the first fast exponential response and gives a nonempty
original-action family with the longer interval
$t\le c a^3\log(1/a)$ for all $c<100/17$ and $o(a)$ changes of both
curvature observations.

\paragraph{Inherited inputs.}
The unchanged V39 root and Poisson/cubic bank, the analytic original
initializer, V40's leading rate correction, V41's complete signed matrix,
and V42's original-source enclosure and growth bound retain their stated
scopes. The default new verifier reruns V42's root and $17334$-entry
stationary-source stage; it does not regenerate the older $79224$-entry
root bank or the $5046$ weak-Poisson entries. All new splitting inequalities
are exact rational comparisons, including the similarity uncertainty on
the entire root cube. Finite model regressions check the projector,
selection identity and exponents separately from original-source data.

\paragraph{Still unresolved.}
The original higher coefficients $b_1,\mathbb L_1$, the numerical value
and norm of $w^\sharp$, and explicit amplitude thresholds have not been
computed. The new result is an analytic existence construction, not an
executed initial circuit or a tabulated nonlinear root. It does not prove
that the old family has nonzero excitation, remove the growing spectrum,
or exclude additional very slow growth on the complement. Higher-order
fast forcing, an invariant slow manifold, a fixed positive physical
lifespan, refinement-uniform bounds and primitive selection of these
records remain unproved. The independent V32 stress-compensation branch
and the full matter-coupled Einstein--Standard-Model limit are unchanged.

\paragraph{Verification and preservation.}
Every V42 byte before its final document command is preserved. The new
verifier is standalone with Python and NumPy and replays the parent
source certificate by default; a labelled option runs only the new
inequalities using the inherited matrix enclosure. The nonlinear
initial-data, remainder and continuation statements are written proofs,
not numerical spacetime integrations or proof-assistant formalizations.
New labels have prefix \texttt{v43:}.

\begingroup
\renewcommand{\refname}{Additional reference for Version 43}
\begin{thebibliography}{9}
\small\raggedright
\bibitem{v43:Wyss}
C.~Wyss, \emph{Hamiltonians with Riesz Bases of Generalised Eigenvectors
and Riccati Equations}, arXiv:1005.5336 (2010).
\url{https://arxiv.org/abs/1005.5336}.
Primary-source context for invariant graph/Riccati methods. The finite
contraction proof and original-source numerical enclosures used here are
given in full above, not imported from that work.
\end{thebibliography}
\endgroup
\addtocontents{toc}{\protect\endgroup}
% END IMMUTABLE VERSION 43 BODY
% APPEND VERSION 44 HERE, BEFORE THE FINAL END-DOCUMENT COMMAND.



% BEGIN IMMUTABLE VERSION 44 BODY
\clearpage
\begin{center}
{\Large\bfseries Version 44: Exact nonlinear selection of a locally invariant graph}\\[0.5em]
{\large A regular amplitude chart, transverse initial-rate access, and growth controlled by the complementary sector}
\end{center}
\makeatletter
\addtocontents{toc}{\protect\begingroup}
\addtocontents{toc}{\protect\makeatletter}
\addtocontents{toc}{\protect\def\protect\l@section{\protect\bfseries\protect\@dottedtocline{1}{0em}{2.5em}}}
\makeatother
\addcontentsline{toc}{part}{Version 44: Exact nonlinear selection of a locally invariant graph}

\section{Moving from a finite-order cancellation to an invariant initial selection}
\label{v44:sec:context}

Keep the cutoff $h=1/3$, the algebraic mixed-mode root of
\cref{v39:thm:root}, and one fixed analytic V40 initializer. The V41
signed certificate and V43 spectral splitting are inherited at exactly
that root. In particular, the weak growing space has dimension nine;
the rest of the full $431$-dimensional normalized state has dimension
$422$. No spectral or coefficient calculation on a new root is asserted.

Version 43 selects a fixed second-order weak rate to remove the first
exponential response. Repeating that procedure at every order would not
by itself construct an invariant set. The present body instead proves
that a locally invariant, finite-regularity graph exists for a regular
amplitude chart of the \emph{full original field}, and that exact
constrained initial records intersect this graph transversely through
the same nine initial-rate controls. The resulting selection is
amplitude dependent. It is not a truncation of a formal power series.

The general invariant-graph method is standard Lyapunov--Perron theory;
primary-source context for gap conditions is \cite{v44:LL}. A proof
with the needed differentiability is included below. The source-specific
steps are the regular extension, its graph jets, and the divided normal
velocity of the exact initial-record family. These establish a nonempty
intersection with the original constraints rather than merely assuming
that the constraints meet an abstract invariant manifold.

The graph is local, may depend on an auxiliary extension outside its
original chart, and is not asserted unique among all local invariant
graphs. It is normally repelling in the nine fast directions. The theorem
is therefore about selected initial records, not attraction of arbitrary
microscopic data. No positive physical lifespan independent of amplitude
or spatial-refinement estimate is inferred.

\section{A fixed normalized chart removes the apparent singular vector field}
\label{v44:sec:normalization}

Write the chosen $w=0$ base records as $(X_a^0,\lambda_a^0)$ and use the
fixed mean-lapse-zero multiplier complement. For $a\ne0$ define the
$431$-component coordinate
\begin{equation}
 u=\left(\frac{X-X_a^0}{a},\frac{\lambda-\lambda_a^0}{a}\right),
 \qquad t=a^3s.
 \label{v44:eq:coordinate}
\end{equation}
Multiplier components in this formula use their fixed physical synthesis.
The physical clock remains $t$; $s$ is only a comparison variable. The
V42 displacement is $z=u/a^3$. Denote its $w=0$ field by
$\mathscr F_a^0$ and put
\begin{equation}
                  \mathcal G(a,u)=a^3\mathscr F_a^0(u/a^3).
 \label{v44:eq:G-definition}
\end{equation}
Thus $u'=\mathcal G(a,u)$, where a prime means differentiation in $s$.

\begin{lemma}[Regular extension of the unchanged normalized field]
\label[lemma]{v44:lem:regular}
There is a fixed neighborhood of $(0,0)$ on which $\mathcal G$ extends
real analytically. With the V42 coefficients $b,b_1,\mathbb L,\mathbb L_1$
for the chosen initializer,
\begin{align}
 \mathcal G(a,0)&=a^3b+a^4b_1+O(a^5),
       &\mathbb Lb&=0,\label{v44:eq:base-field}\\
 D_u\mathcal G(a,0)&=\mathbb L+a\mathbb L_1+O(a^2),
       &D_u^2\mathcal G(a,u)&=O(1).\label{v44:eq:base-derivative}
\end{align}
In particular $\mathcal G(0,0)=0$, its full derivative at $(0,0)$ is
$(0,\mathbb L)$, and
\begin{equation}
 \|D_u\mathcal G(a,u)-\mathbb L\|\le C(|a|+\|u\|).
 \label{v44:eq:G-derivative}
\end{equation}
For small $a>0$, the field inside this neighborhood is precisely the
original normalized action field, with its complete signed solve.
\end{lemma}
\begin{proof}
Set $X_a^0=a\bar X(a)$ and $\lambda_a^0=e+a\bar\chi(a)$.
In the source formula used in \cref{v42:lem:tube}, replace the normalized
arguments by $(\bar X(a)+u_X,\bar\chi(a)+u_\lambda)$.
The divided sources $J_A/a,J_W/a^2$, the divided drifts
$\mathsf B_A/a,\mathsf B_W/a^2$, and the normalized signed matrix
$\Gamma$ are analytic on a fixed neighborhood of the retained root.
V41's nonsingularity makes $\Gamma^{-1}$ analytic there.

Let $\mathfrak d=(\mathsf B_A/a,\mathsf B_W/a^2)^T$.
After multiplication by $a^3$, the two multiplier components displayed
in the proof of \cref{v42:lem:tube} become
\[
              -aR_A\Gamma^{-1}\mathfrak d,
              \qquad -R_W\Gamma^{-1}\mathfrak d.
\]
There is no remaining negative power of $a$. The canonical component
is $a^2F$, also analytic. This proves an extension on the \emph{fixed}
normalized neighborhood, not just on bounded $z$ sets. Positivity of the
physical lapse and metric, the stationary chart and signed regularity
persist there for sufficiently small positive $a$.

The exact corrected base rate and V42's first field expansion give
\eqref{v44:eq:base-field}. For $a\ne0$,
$D_u\mathcal G(a,0)=D_z\mathscr F_a^0(0)$, yielding
\eqref{v44:eq:base-derivative} by analyticity. Bounded second derivatives
on a smaller fixed neighborhood give \eqref{v44:eq:G-derivative}.
The extension at $a=0$ is an analytical device; it is not a new physical
zero-amplitude rate law.
\end{proof}

Use the full commuting projector $\mathcal P_{\rm f}$ of
\cref{v43:lem:full-split}. Put
\[
 P_{\rm f}=\mathcal P_{\rm f},\quad P_{\rm c}=I-P_{\rm f},\quad
 E_{\rm f}=\Ran P_{\rm f},\quad E_{\rm c}=\Ran P_{\rm c},
 \quad L_{\rm f}=\mathbb L|_{E_{\rm f}},\quad
 L_{\rm c}=\mathbb L|_{E_{\rm c}}.
\]
Then $\dim E_{\rm f}=9$, $\dim E_{\rm c}=422$, and $b\in E_{\rm c}$.
The inherited norms give constants $M\ge1$ and
\begin{equation}
 \alpha=1/61,\qquad \nu=1/1000,\qquad
 \|e^{-sL_{\rm f}}\|\le e^{-\alpha s},\qquad
 \|e^{sL_{\rm c}}\|\le M e^{\nu s}\quad(s\ge0).
 \label{v44:eq:gap}
\end{equation}
The factor $M$ absorbs the full triangular block and the fixed norm
changes. It is finite, not assigned the value one. Also
$\operatorname{sym}L_{\rm f}\succeq\alpha I$ in its inherited graph norm.

\section{A locally invariant graph with sufficient differentiability}
\label{v44:sec:graph}

Write $u=x+y$ in $E_{\rm c}\oplus E_{\rm f}$ and let
$G_{\rm c}=P_{\rm c}\mathcal G$, $G_{\rm f}=P_{\rm f}\mathcal G$.
The argument includes $a'=0$ as an additional complementary coordinate.

\begin{theorem}[A local slow-growth graph for the original action field]
\label[theorem]{v44:thm:graph}
There is a $C^6$ map $h(a,x)$ from a neighborhood of $(0,0)$ in
$\R\times E_{\rm c}$ to $E_{\rm f}$, with
\begin{equation}
 h(0,0)=0,\qquad Dh(0,0)=0,
 \label{v44:eq:h-tangent}
\end{equation}
whose graph
\begin{equation}
                  \mathcal M_a=\{x+h(a,x):x\in E_{\rm c}\}
 \label{v44:eq:manifold}
\end{equation}
is locally invariant for $u'=\mathcal G(a,u)$.
It has codimension nine for each fixed $a$. On this graph,
\begin{equation}
 D_xh(a,x)G_{\rm c}(a,x,h(a,x))
           =G_{\rm f}(a,x,h(a,x)).
 \label{v44:eq:invariance}
\end{equation}
The theorem does not assert analyticity, attraction, global invariance
outside the original chart, or damping of every complementary direction.
\end{theorem}
\begin{proof}
Set $N(a,u)=\mathcal G(a,u)-\mathbb Lu$.
By \cref{v44:lem:regular}, $N(0,0)=DN(0,0)=0$.
Multiply $N$ by a smooth cutoff supported in a sufficiently small ball
in $(a,u)$ and equal to one on its concentric inner ball. Extend the
product by zero; call it $\widetilde N$. Its global Lipschitz constant
$\delta$ can be made arbitrarily small, and all its fixed higher
derivatives are bounded. This extension is used only to prove existence
of a local graph. It agrees with the original vector field in the inner
ball.

Choose $\beta=1/800$. The strict inequalities
\begin{equation}
              \nu<\beta,\qquad 7\beta<\alpha
 \label{v44:eq:regularity-gap}
\end{equation}
hold. For $\theta\in[\beta,7\beta]$ use the norm
$\|U\|_\theta=\sup_{s\ge0}e^{-\theta s}\|U(s)\|$.
For fixed $(a,x_0)$ solve
\begin{align}
 x(s)&=e^{sL_{\rm c}}x_0+
   \int_0^s e^{(s-r)L_{\rm c}}P_{\rm c}\widetilde N(a,x(r)+y(r))\,\dd r,
       \label{v44:eq:LP-c}\\
 y(s)&=-\int_s^\infty e^{(s-r)L_{\rm f}}
           P_{\rm f}\widetilde N(a,x(r)+y(r))\,\dd r.
       \label{v44:eq:LP-f}
\end{align}
Choose product norms so that the projection costs are incorporated in
$M$ and $\delta$. The Lipschitz constant of this map in the weighted
space is at most
\begin{equation}
 q_\theta=\delta\left\{\frac{M}{\theta-\nu}
                              +\frac1{\alpha-\theta}\right\}.
 \label{v44:eq:LP-contraction}
\end{equation}
Shrink the cutoff radius until $q_\theta<1/4$ on the displayed compact
weight interval. The value at zero has finite weighted norm, including
its parameter forcing. Banach's theorem gives a unique solution in the
$\beta$-weighted space. Define $h(a,x_0)=y(0)$. The same solution lies
in every larger weighted space and is its unique fixed point there.

Here are the differentiability details needed below. A formal first
parameter derivative solves the differentiated fixed-point equation;
its linear inverse is bounded by $(1-q_\theta)^{-1}$. A derivative of
order $j$ has a finite sum of products of lower derivatives as its
inhomogeneous term. If their orders sum to $j$, their weighted bounds
multiply to an $e^{j\beta s}$ bound. Bounded derivatives of
$\widetilde N$ therefore put the order-$j$ inhomogeneity in the
$j\beta$ space. Induction constructs these derivatives for $j\le6$.
They are genuine continuous derivatives, not just formal solutions:
subtract the finite Taylor formulas on $[0,S]$, where ordinary uniform
Taylor remainders apply, and bound the tail in the slightly larger
$(j+1)\beta$ norm by its extra factor $e^{-\beta S}$. First choose
$S$ large and then take the difference quotient to zero. The uniform
inverse bound closes the induction. Evaluation at $s=0$ is bounded in
all these spaces, proving $C^6$ regularity of $h$. The spare seventh
weight is why \eqref{v44:eq:regularity-gap} is stated in this form.

At $(a,x_0)=(0,0)$ the fixed point is zero. Since $DN(0,0)=0$, its
first variation has zero fast coordinate, proving
\eqref{v44:eq:h-tangent}. Shifting any fixed-point trajectory forward
in $s$ leaves it in the weighted class and gives the fixed point with
its new complementary initial coordinate. Uniqueness proves invariance
for the extension. Differentiating the graph along that trajectory gives
\eqref{v44:eq:invariance}. In the inner ball this is the original
field. Local backward invariance while the trajectory remains in that
ball follows by uniqueness of the finite-dimensional flow, or by adding
a finite backward segment to the same weighted future.
\end{proof}

The cutoff is not an action modification executed on a history. All
histories used below remain in the inner ball, where the two vector
fields coincide. Different admissible extensions can select different
local graphs; no canonical global graph or microscopic rule for choosing
one is inferred.

\section{Graph jets identify the initial-selection equation}
\label{v44:sec:jets}

The graph's existence alone does not establish that the original
constrained initializer reaches it. The needed transversality appears
at a higher amplitude order and must be calculated.

\begin{lemma}[The relevant graph jets are determined by the actual first field correction]
\label[lemma]{v44:lem:jets}
There is a linear map $H:E_{\rm c}\to E_{\rm f}$ such that
\begin{equation}
 h(a,0)=O(a^4),\qquad
 D_xh(a,x)=aH+O(a^2+\|x\|).
 \label{v44:eq:h-jets}
\end{equation}
It is the unique solution of
\begin{equation}
 L_{\rm f}H-HL_{\rm c}=-P_{\rm f}\mathbb L_1|_{E_{\rm c}},
 \qquad
 Hb=-L_{\rm f}^{-1}P_{\rm f}\mathbb L_1b.
 \label{v44:eq:H}
\end{equation}
\end{lemma}
\begin{proof}
Tangency initially gives $h(a,0)=O(a^2)$ and $D_xh(a,0)=O(a)$.
At $x=0$ the fast base force is $O(a^4)$ because $P_{\rm f}b=0$.
Insert these estimates into \eqref{v44:eq:invariance}. Taylor expansion
in the fast variable gives
\[
 \|L_{\rm f}h(a,0)\|
 \le C a^4+C(|a|+\|h(a,0)\|)\|h(a,0)\|.
\]
Invertibility of $L_{\rm f}$ absorbs the last term for small $a$ and
proves $h(a,0)=O(a^4)$.

Let $H=\partial_aD_xh(0,0)$. Differentiate invariance once in $x$
and once in $a$ at the origin. The term containing $D_x^2h$ multiplied
by the complementary base velocity is $O(a^3)$ and contributes zero at
this order. Equations \eqref{v44:eq:base-field}--\eqref{v44:eq:base-derivative}
give the Sylvester equation in \eqref{v44:eq:H}. Its unique solution is
\[
 H=-\int_0^\infty e^{-sL_{\rm f}}
                P_{\rm f}\mathbb L_1|_{E_{\rm c}}e^{sL_{\rm c}}\,\dd s,
\]
which converges by $\alpha>\nu$. Multiplying the Sylvester equation by
$b$, for which $L_{\rm c}b=0$, proves the last identity. Taylor's theorem
for the $C^6$ graph gives the remainder in \eqref{v44:eq:h-jets}.
\end{proof}

Let the exact V43 rate-access family be $(X_a^w,\lambda_a^w)$, and
express its initial point in the \emph{fixed} chart
\eqref{v44:eq:coordinate}:
\begin{equation}
 U(a,w)=\left(\frac{X_a^w-X_a^0}{a},
              \frac{\lambda_a^w-\lambda_a^0}{a}\right)=O_w(a^2).
 \label{v44:eq:initial-U}
\end{equation}
It is jointly analytic after removal of the factor $a$, locally near
any fixed finite $w$. The exact initial rate and V43's triangular
coefficient identity give
\begin{equation}
 \mathcal G(a,U(a,w))
       =a^3b+a^4(b_1+J_Ww)+O_w(a^5).
 \label{v44:eq:initial-G}
\end{equation}
This uses the initializer's actual rate identity, not a singular-inverse
estimate on $U(a,w)$.

Restrict the free coefficient to the fixed nine-dimensional weak space
$w\in\mathcal E_{\rm f}$, so $J_W:\mathcal E_{\rm f}\to E_{\rm f}$
is an isomorphism. Define the normal velocity of this initial record:
\begin{equation}
 \mathcal N(a,w)=P_{\rm f}\mathcal G(a,U(a,w))
   -D_xh(a,P_{\rm c}U(a,w))P_{\rm c}\mathcal G(a,U(a,w)).
 \label{v44:eq:normal-velocity}
\end{equation}

\begin{proposition}[Divided normal velocity is transverse in exactly nine controls]
\label[proposition]{v44:prop:normal}
The quotient $\mathcal N(a,w)/a^4$ has a $C^1$ extension to $a=0$,
and
\begin{equation}
 \boxed{
 a^{-4}\mathcal N(a,w)
       =P_{\rm f}b_1+J_Ww-Hb+O_w(a).}
 \label{v44:eq:normal-expansion}
\end{equation}
Its $w$ derivative at $a=0$ is the isomorphism $J_W$.
Moreover, in a sufficiently small fixed chart, zero normal velocity at
any point is equivalent to that point lying on $\mathcal M_a$.
\end{proposition}
\begin{proof}
Equations \eqref{v44:eq:initial-U} and \eqref{v44:eq:h-jets} give
$D_xh(a,P_{\rm c}U)=aH+O_w(a^2)$.
Use \eqref{v44:eq:initial-G}, $P_{\rm f}b=0$ and
$P_{\rm c}J_Ww=0$ to obtain \eqref{v44:eq:normal-expansion}.
The numerator is $C^5$, since $h$ is $C^6$. Its first four Taylor
coefficients below order four vanish identically for all nearby $w$.
The integral Taylor formula
\[
 \frac{\mathcal N(a,w)}{a^4}
   =\frac1{3!}\int_0^1(1-r)^3\partial_a^4\mathcal N(ra,w)\,\dd r
\]
therefore gives the claimed $C^1$ extension and its derivative.

For the equivalence, at fixed $(a,x)$ let
$\mathcal A(a,x,y)=G_{\rm f}(a,x,y)-D_xh(a,x)G_{\rm c}(a,x,y)$.
It vanishes at $y=h(a,x)$ by invariance. The mean-value integral in
$y$ gives
\begin{equation}
 \mathcal A(a,x,y)=[L_{\rm f}+E(a,x,y)](y-h(a,x)),\qquad
 \|E(a,x,y)\|\le C(|a|+\|x\|+\|y\|).
 \label{v44:eq:normal-factor}
\end{equation}
Here the linear cross blocks vanish because the projectors commute
with $\mathbb L$. Shrinking the chart makes
$\|L_{\rm f}^{-1}E\|<1/2$. The bracket is then invertible, proving
both directions. A vanishing normal \emph{velocity} is sufficient here
because of this proved factorization, not by definition of tangency at
an arbitrary off-graph point.
\end{proof}

\section{Exact constrained initial records on the nonlinear invariant graph}
\label{v44:sec:intersection}

\begin{theorem}[A one-time nonlinear initial selection]
\label[theorem]{v44:thm:selection}
For the chosen local graph there is a $C^1$ function $w(a)$, for
sufficiently small $a$, with
\begin{equation}
 \boxed{
 w(a)=w^\sharp+O(a),\qquad
 J_Ww^\sharp=-P_{\rm f}b_1-L_{\rm f}^{-1}P_{\rm f}\mathbb L_1b,}
 \label{v44:eq:selection}
\end{equation}
such that $U(a,w(a))\in\mathcal M_a$. It is locally unique among the
nine selected rate controls near $w^\sharp$ for this graph and initializer.
The corresponding original records satisfy exactly
\begin{equation}
 \boxed{
 P=0,\qquad
 \mathsf B+M(av_*+a^2I_Ww(a))=0,\qquad
 \lambda_t(0)=av_*+a^2I_Ww(a).}
 \label{v44:eq:physical-initial}
\end{equation}
Their unique normalized original histories stay on $\mathcal M_a$
while they remain in its local chart. Their canonical fields and
multiplier values through order $a^2$ are unchanged from the chosen
$w=0$ family.
\end{theorem}
\begin{proof}
At $a=0$, the divided normal equation in
\cref{v44:prop:normal} has the unique root determined by
$J_Ww=Hb-P_{\rm f}b_1$. Substitution of \eqref{v44:eq:H} identifies
this root with V43's $w^\sharp$ in \eqref{v43:eq:wsharp}.
The $C^1$ implicit-function theorem and the invertible $w$ derivative
give $w(a)$ and its $O(a)$ comparison. Proposition
\ref{v44:prop:normal} then gives membership in the graph, not only
cancellation of a finite jet.

The V43 initializer is jointly analytic on a neighborhood of the fixed
finite $w^\sharp$. Thus it can be composed with $w(a)$, with a common
sufficient amplitude after shrinking that neighborhood. It solves every
original constraint and exact initial consistency equation. V41's
unchanged signed regularity identifies the rate and the unique original
continuation. Local invariance proves the graph assertion. Finally
\eqref{v44:eq:initial-U} makes the physical initial difference $O(a^3)$.
The selected family is only claimed $C^1$ in $a$, not analytic: the
invariant graph was proved $C^6$, not analytic.
\end{proof}

This theorem replaces an infinite sequence of proposed fast-jet
cancellations by an exact nonlinear initial intersection. It does not
compute $b_1,\mathbb L_1,w^\sharp$ or the graph numerically. Their
finite, source-defined values suffice for existence. It also does not
prescribe a rate after the initial cut: the evolved rate is still
obtained from the original consistency equation. There is no repeated
projection onto $\mathcal M_a$.

\section{Propagation is controlled by the complementary growth bound}
\label{v44:sec:time}

\begin{theorem}[An invariantly selected family on the complementary logarithmic scale]
\label[theorem]{v44:thm:time}
Fix $c>0$ and $\mu>\nu=1/1000$ such that $c\mu<1$.
For all sufficiently small positive $a$, the histories in
\cref{v44:thm:selection} exist on
\begin{equation}
 S_a=c\log(1/a),\qquad T_a=a^3S_a,
 \label{v44:eq:time}
\end{equation}
preserve every original constraint, and remain on the local graph.
Relative to their \emph{own} initial cut, put
$z(s)=[u(s)-U(a,w(a))]/a^3$. Then
\begin{equation}
 \boxed{
 \|z(s)-sb\|+\|z'(s)-b\|+\|z''(s)\|
           \le C_{c,\mu}a e^{\mu s},\qquad 0\le s\le S_a.}
 \label{v44:eq:trajectory}
\end{equation}
Consequently every fixed $c<1000$ is admitted by a suitable $\mu$.
The same selected family works for all these fixed $c$, with thresholds
and constants permitted to depend on $c,\mu$.
\end{theorem}
\begin{proof}
On the graph define the reduced original field
$g_a(x)=G_{\rm c}(a,x,h(a,x))$.
Tangency and \eqref{v44:eq:G-derivative} imply
\begin{equation}
 D_xg_a(x)=L_{\rm c}+O(|a|+\|x\|),\qquad
 D_xh(a,x)=O(|a|+\|x\|).
 \label{v44:eq:reduced-derivatives}
\end{equation}
The initial point has $x(0)=O(a^2)$. Put $v=x'/a^3$.
Equation \eqref{v44:eq:initial-G} and boundedness of $w(a)$ give
$v(0)=b+O(a)$, with $b$ identified as an element of $E_{\rm c}$.
Along the reduced flow,
\[
                  v'=D_xg_a(x)v,\qquad x'=a^3v.
\]
Bootstrap $\|v-b\|\le1$ and remain in the original graph chart.
Then $\|x(s)\|\le C(a^2+a^3s)$, hence $\|x\|\le C_ca^2$
on the logarithmic interval for small $a$. In particular the
perturbation in the first equation of \eqref{v44:eq:reduced-derivatives}
is at most $C_ca$.

Set $r=v-b$ and use $L_{\rm c}b=0$. Variation of constants and
\eqref{v44:eq:gap} give
\[
 \|r(s)\|\le MC a e^{\nu s}
   +MC_ca\int_0^s e^{\nu(s-t)}(\|b\|+\|r(t)\|)\,\dd t.
\]
After multiplication by $e^{-\nu s}$, the constant-source integral is
bounded by $\|b\|/\nu$. Gronwall therefore yields
$\|r(s)\|\le C'a e^{(\nu+C''a)s}$. For small $a$ this is at most
$C'a e^{\mu s}$. At $S_a$ it is $O(a^{1-c\mu})=o(1)$.
This closes the bootstrap. Also $x=O(a^2)$ and $h(a,x)=O(a^3)$
there, so the original chart, positivity and signed regularity cannot
fail. Finite-dimensional continuation establishes the whole interval.

The full normalized velocity is
$u'/a^3=v+D_xh(a,x)v$. By
\eqref{v44:eq:reduced-derivatives}, its difference from $b$ is
$O(ae^{\mu s})$. Integration proves the position estimate in
\eqref{v44:eq:trajectory}. For the second derivative use
\[
 z''=D_u\mathcal G(a,u)z',\qquad \mathbb Lb=0.
\]
Here $u=O(a^2)$, $z'=b+O(ae^{\mu s})$ and
\eqref{v44:eq:G-derivative} gives the same bound. This does not
differentiate an uncontrolled time remainder. Constraint preservation
follows from the original exact identity $\dot P=0$.
For $c<1000$, the interval $(\nu,1/c)$ is nonempty, proving the
last assertion.
\end{proof}

\begin{corollary}[Original rate and both curvature observations]
\label[corollary]{v44:cor:observations}
On \eqref{v44:eq:time},
\begin{equation}
 \|\lambda_{t,W}\|+\|\lambda_{t,A}-av_A\|
                \le C_{c,\mu}a^2e^{\mu s}.
 \label{v44:eq:rates}
\end{equation}
Each retained metric Riemann tensor, literal electric link curvature,
and literal spatial link curvature has size $O(a)$. For each such
observation separately,
\begin{equation}
 \boxed{
 \sup_{t\le T_a}\|\mathscr O_a(t)-\mathscr O_a(0)\|
                \le C_{c,\mu}a^{2-\mu c}=o(a).}
 \label{v44:eq:curvature}
\end{equation}
For example one may take $c=100$, $\mu=1/500$. Then
\begin{equation}
 \boxed{
 T_a=100a^3\log(1/a),\qquad
 \lambda_{t,W}=O(a^{9/5}),\qquad
 \mathscr O_a(t)-\mathscr O_a(0)=O(a^{9/5}).}
 \label{v44:eq:example}
\end{equation}
The corresponding scaled trajectory error is $O(a^{4/5})$.
\end{corollary}
\begin{proof}
Physical displacements equal $a^4z$ and multiplier rates equal
$a\zeta'$. Apply \eqref{v44:eq:trajectory}, using $b_W=0$ and
$b_A=v_A$. Differences from the initial rate obey the same estimate;
the extra fixed initial term is $a^2I_Ww(a)=O(a^2)$.
Each original curvature formula is a smooth fixed-cutoff function of
$(X,\lambda,F,\lambda_t)$ and uses no $\lambda_{tt}$, as retained
in V41--V43. State displacements are
$O(a^4s+a^5e^{\mu s})$ and changes of the rate are
$O(a^2e^{\mu s})$. The latter dominate on the displayed interval.
Initial curvature size remains $O(a)$ by the same original formulas.
Finally $\mu c=1/5$ in the example.
\end{proof}

The cutoff is still fixed, $T_a\to0$, and no maximal-lifespan statement
is made. The constant $1000$ comes from the inherited complementary
upper growth estimate, not from an optimized spectrum. Local invariance
does not exclude accumulation in the other twenty weak directions or
the remaining canonical and active coordinates on longer times.

\section{The invariant graph is normally repelling, not a selection mechanism}
\label{v44:sec:repulsion}

\begin{proposition}[Exact transverse equation and local forward repulsion]
\label[proposition]{v44:prop:repulsion}
For any original normalized history staying in a sufficiently small
chart, define $d(s)=y(s)-h(a,x(s))$. Then
\begin{equation}
 d'=[L_{\rm f}+E(a,x,y)]d,
 \qquad
 \|d(s)\|_{\rm f}\ge e^{s/122}\|d(0)\|_{\rm f}.
 \label{v44:eq:repulsion}
\end{equation}
The second inequality holds for forward intervals staying in that chart.
For small $a>0$ and sufficiently small initial-rate errors $e\in\mathcal E_{\rm f}$,
the exact constrained initializer with $w=w(a)+e$ has
\begin{equation}
                  \|d(0)\|_{\rm f}\ge c_0a^4\|e\|.
 \label{v44:eq:rate-normal-error}
\end{equation}
\end{proposition}
\begin{proof}
The chain rule and \eqref{v44:eq:normal-factor} give the exact
transverse equation. In the V43 fast graph norm,
$\operatorname{sym}L_{\rm f}\succeq1/61$. Shrink the neighborhood
until the perturbation has norm at most $1/122$. Differentiating
$\|d\|_{\rm f}^2$ and integrating gives
\eqref{v44:eq:repulsion}, including the identically zero solution.

The divided normal-velocity map has an invertible $w$ derivative
uniformly near $(0,w^\sharp)$. Its mean-value formula, as a perturbation
of the fixed isomorphism $J_W$, gives
$\|\mathcal N(a,w(a)+e)\|\ge c a^4\|e\|$.
The operator in \eqref{v44:eq:normal-factor} has a bounded upper norm
on the same chart. Applying that formula at the initial point proves
\eqref{v44:eq:rate-normal-error}.
\end{proof}

The proposition concerns normalized-state distance to the selected graph,
not a lower bound on every physical curvature component. It does show
that the graph is not a proved attracting basin, even within the exactly
constrained initial-record family. No numerical amplitude threshold,
implementation accuracy, or microscopic procedure for selecting $w(a)$
is supplied. The exact selection must not be replaced by an arbitrary
fixed approximate rate while retaining its invariant conclusion.

\section{Proved conclusions, inherited inputs, and the remaining physical interval}
\label{v44:sec:status}

\paragraph{Proved in this body.}
The full original action field has a regular analytic amplitude chart.
Its certified spectral gap gives a locally invariant $C^6$ graph of
codimension nine. The divided normal velocity of the exact initial-rate
family is transverse, so a $C^1$ amplitude-dependent choice $w(a)$
places exactly constrained records on that graph. Its leading coefficient
is V43's $w^\sharp$, but graph membership is exact and is not limited
to first-order cancellation. Propagation along the graph yields original
histories on $t\le c a^3\log(1/a)$ for every fixed $c<1000$, with
$o(a)$ changes of both curvature observations. The graph is locally
repelling in the fast normal directions.

\paragraph{Inherited inputs.}
The exact V39 root, the original analytic initializer, V40's initial
rate correction, V41's complete signed regularity, V42's analytic
source grading and V43's spectral splitting and rate-access theorem
are used at their stated fixed-cutoff scope. The new verifier replays
V43, including its root inclusion and $17334$-entry stationary-source
reconstruction. It does not regenerate the older $79224$-entry root
bank or the $5046$ weak-Poisson bank. No new original-source matrix is
claimed to have been independently reconstructed beyond that replay.

\paragraph{Still unresolved.}
The graph and the coefficients $b_1,\mathbb L_1,w^\sharp,w(a)$ are
not numerically tabulated. No explicit local radius or sufficient
amplitude is enclosed. The graph may depend on the analytical cutoff
extension and is not asserted analytic or physically preferred. The
selected family is not an open attracting set. An invariant graph does
not establish a fixed positive physical lifespan: the complementary
motion, chart margins, higher curvature derivatives and spatial
refinement still require estimates. No primitive source-selection,
full matter coupling, or Einstein--Standard-Model continuum theorem
is inferred. The separate V32 stress-compensation branch is unchanged.

\paragraph{Verification and preservation.}
Every V43 byte before its final document command is preserved. The
standalone verifier checks the inherited spectral certificate, the seven
weighted spectral gaps, the graph-jet Sylvester and normal-velocity
identities, exact polynomial examples of invariant selection and its
repulsion, and the time/exponent conversions. These finite checks support
the written proof; they are not a computation of the native invariant
graph, a nonlinear spacetime integration, or proof-assistant formalization.
New labels have prefix \texttt{v44:}.

\begingroup
\renewcommand{\refname}{Additional reference for Version 44}
\begin{thebibliography}{9}
\small\raggedright
\bibitem{v44:LL}
Y.~Latushkin and B.~Layton, \emph{The optimal gap condition for invariant
manifolds}, Discrete and Continuous Dynamical Systems \textbf{5} (1999),
233--268. \url{https://doi.org/10.3934/dcds.1999.5.233}.
Primary-source context for spectral-gap invariant-manifold methods.
The finite-dimensional contraction, differentiability argument and
source-specific initial intersection used here are given above.
\end{thebibliography}
\endgroup
\addtocontents{toc}{\protect\endgroup}
% END IMMUTABLE VERSION 44 BODY
% APPEND VERSION 45 HERE, BEFORE THE FINAL END-DOCUMENT COMMAND.



% BEGIN IMMUTABLE VERSION 45 BODY
\clearpage
\begin{center}
{\Large\bfseries Version 45: Exact neutral spectrum and a polynomial physical interval}\\[0.5em]
{\large A cubic original-source kernel certificate, positive oscillatory energy, and continuation on the invariant family}
\end{center}
\makeatletter
\addtocontents{toc}{\protect\begingroup}
\addtocontents{toc}{\protect\makeatletter}
\addtocontents{toc}{\protect\def\protect\l@section{\protect\bfseries\protect\@dottedtocline{1}{0em}{2.5em}}}
\makeatother
\addcontentsline{toc}{part}{Version 45: Exact neutral spectrum and polynomial physical-time continuation}

\section{The complementary growth estimate, not another initial selection}
\label{v45:sec:context}

Keep the original cutoff $h=1/3$, the algebraic mixed-mode root of
\cref{v39:thm:root}, and one of the invariant initial families supplied by
\cref{v44:thm:selection}. All constants below are at this fixed cutoff.
There is no new choice of action, clock, first field, or initial-rate family.
The open point addressed here is the complementary linear growth estimate
used to continue this same family. The upper exponent $1/1000$ in V43--V44
is a valid enclosure, but does not determine whether the remaining weak
spectrum has a genuinely positive real part.

The actual weak matrices are
\begin{equation}
 D=C_W-L^TA^{-1}L,\qquad
 \widehat{\mathsf K}=\mathsf W^T\widehat K\mathsf W,\qquad
 \mathsf K=\widehat{\mathsf K}/\sqrt3,\qquad
 B=-D^{-1}\mathsf K.
 \label{v45:eq:weak}
\end{equation}
The $29$ columns of $\mathsf W$ retain the three constant shifts and the
$26$ nonconstant curl-free tests, in their original order. The matrix $D$
is the original signed Schur form, not a positive replacement. V41 proves
its inertia $(19,10,0)$; V42 proves $B^TD+DB=0$; V43 proves a
nine-dimensional fast invariant space and bounds the remaining real parts
above by $1/1000$. The new calculation proves exact nullity five, a
nondegenerate restriction of $D$ of inertia $(4,1,0)$ on that kernel,
and positivity of the remaining six oscillatory directions. It follows
that the weak nonfast evolution is bounded. The full complementary
state may still have linear growth, and this distinction is retained.

The indefinite-inner-product arguments below are finite-dimensional
linear algebra; Krein-signature theory gives broader context
\cite{v45:KM}. The new source content is the explicit cubic kernel
identity and its signed kernel certificate. The invariant graph,
initial intersection, stationary chart and curvature-read formulas
are inherited from V39--V44, not independently reproved here.

\section{An exact cubic null vector of the original weak Poisson matrix}
\label{v45:sec:kernel}

Let $\xi=(\xi_1,\xi_2,\xi_3,\xi_4)$ be the mixed amplitudes of V39,
with $\xi_4=1/5$ on its isolating cube. Restricting the six retained
rational Poisson coefficient matrices to the weak tests gives
\begin{equation}
 \widehat{\mathsf K}(\xi)=\sum_{j=1}^4\xi_jK_j,\qquad
 K_j^T=-K_j,\qquad K_je_i=0\quad(0\le i\le2).
 \label{v45:eq:Klinear}
\end{equation}
The two other coefficient matrices, corresponding to the constant and
quadratic-amplitude source coordinates, restrict to zero exactly.
Indices here and in the certificate are zero based.

Put $(p,q,r)=(\xi_2,\xi_3,\xi_4)$, and set the first three entries of
$k(p,q,r)\in\R^{29}$ equal to zero. Its remaining entries are specified
by the following table. Repeated indices have the same polynomial.
\begin{center}\small
\begin{tabular}{@{}ll@{}}
\toprule
Indices & Entry of $k$\\\midrule
$3,18$ & $-p(p-q)(q+r)$\\
$4,19$ & $q(-pq+pr-qr)$\\
$5,20$ & $-pq(-p+r)$\\
$6,27$ & $q^2(-p+r)$\\
$7,28$ & $p^2(q+r)$\\
$8$ & $(p-q)(-pq+pr-qr)$\\
$9$ & $q(p^2+qr)$\\
$10$ & $p(-pq+pr+q^2)$\\
$11$ & $p^2r-pq^2-2pqr+q^2r$\\
$12$ & $-2p^2r-pq^2+2pqr-q^2r$\\
$13$ & $p^2q-p^2r+2pqr-2q^2r$\\
$14$ & $-p^2q-p^2r+pq^2-q^2r$\\
$15,21$ & $-p(-pq+pr-qr)$\\
$16,22$ & $q(-p+r)(p-q)$\\
$17,23$ & $-pq(q+r)$\\
$24$ & $q(-p^2+pq+qr)$\\
$25$ & $p(pr-q^2)$\\
$26$ & $p^2q+p^2r-2pqr+q^2r$\\\bottomrule
\end{tabular}
\end{center}

\begin{theorem}[A structural upper rank bound and an exact root-cube rank]
\label[theorem]{v45:thm:kernel}
The original coefficient matrices obey the polynomial identity
\begin{equation}
 \boxed{\widehat{\mathsf K}(\xi)k(\xi_2,\xi_3,\xi_4)=0}
 \label{v45:eq:syzygy}
\end{equation}
for all four real amplitudes. On the V42 root cube
$\|(\xi_1,\xi_2,\xi_3)-q_0\|_\infty\le2^{-112}$, one has
\begin{equation}
 \boxed{\operatorname{rank}\mathsf K=24,\qquad
                  \dim\ker\mathsf K=5.}
 \label{v45:eq:rank}
\end{equation}
This kernel retains all three constant shifts and has two further
nonconstant directions.
\end{theorem}
\begin{proof}
The displayed vector is homogeneous cubic and independent of $\xi_1$.
The verifier contracts it with the four exact restricted Poisson
coefficient matrices and checks all $26\binom{7}{3}=910$ degree-four
coefficients, each over $\mathbb Q$. There are only $79$ nonzero coefficients
in the displayed vector itself. This is an exhaustive polynomial
identity, not a finite sample of parameter values. The first three
rows are zero by \eqref{v45:eq:Klinear}.

On the root cube, the retained isolating inequalities give
$p>3/10$, $q>17/100$, and $r=1/5$. In particular
$k_7=p^2(q+r)>333/10000$. The $26\times26$ skew block obtained by
removing the constant shifts is therefore singular. Its rank is even,
so it has rank at most $24$, not merely at most $25$.

For the reverse inequality use the principal index set
\begin{equation}
 \mathcal I=\{3,4,\ldots,28\}\setminus\{7,22\}.
 \label{v45:eq:indices}
\end{equation}
The fixed dyadic inverse proposal $P_{24}$ in the package satisfies,
throughout the same cube,
\begin{equation}
 \|I-P_{24}\widehat{\mathsf K}_{\mathcal I\mathcal I}\|_\infty
       <2^{-30},\qquad
 \|\widehat{\mathsf K}_{\mathcal I\mathcal I}^{-1}\|_\infty<1/500.
 \label{v45:eq:minor}
\end{equation}
The enclosure is detailed below. Thus the rank is at least $24$.
Combining both bounds proves \eqref{v45:eq:rank}. The uniform scalar
factor $1/\sqrt3$ changes neither rank nor kernel. The upper rank
argument does not assume the cubic mean equations beyond the stated
nonvanishing component; the lower bound is only claimed on the root
cube.
\end{proof}

\section{The signed form is nondegenerate on the five-dimensional kernel}
\label{v45:sec:inertia}

Let $\mathcal J=\{7,22\}$. Define a $29\times5$ kernel synthesis
$N$ by taking its first three columns to be $e_0,e_1,e_2$, and
its last two columns to have entries $I_2$ on $\mathcal J$, zero on
the constant coordinates, and
\begin{equation}
 N_{\mathcal I,\{3,4\}}
       =-\widehat{\mathsf K}_{\mathcal I\mathcal I}^{-1}
                          \widehat{\mathsf K}_{\mathcal I\mathcal J}.
 \label{v45:eq:N}
\end{equation}
The rank theorem makes the two remaining rows automatic: equivalently,
the residual skew Schur block is zero. Hence $\Ran N=\ker\mathsf K$.

\begin{theorem}[Certified signed kernel inertia]
\label[theorem]{v45:thm:kernel-inertia}
Throughout the retained root cube,
\begin{equation}
 \boxed{\operatorname{Inertia}(N^TDN)=(4,1,0).}
 \label{v45:eq:kernel-inertia}
\end{equation}
In particular the kernel of $B$ is a nondegenerate subspace for the
original signed form $D$.
\end{theorem}
\begin{proof}
We give the complete residual-to-enclosure argument. Denote the rational
weak midpoint by $K_0$ and its norm error by $e_K$. With
$P_{24}$ fixed, put
\[
 q_{24}=\|I-P_{24}(K_0)_{\mathcal I\mathcal I}\|_\infty
                  +\|P_{24}\|_\infty e_K,
 \qquad a_{24}=\frac{\|P_{24}\|_\infty}{1-q_{24}}.
\]
These give \eqref{v45:eq:minor} by a Neumann series. For the dyadic
kernel proposal $N_0$, whose free rows are exact, set
\[
 e_N=a_{24}\left[
 \|(K_0)_{\mathcal I\mathcal I}(N_0)_{\mathcal I,\{3,4\}}
                         +(K_0)_{\mathcal I\mathcal J}\|_\infty
 +e_K(1+\|(N_0)_{\mathcal I,\{3,4\}}\|_\infty)\right].
\]
Then $\|N-N_0\|_\infty\le e_N$. In particular
$\|(N-N_0)^T\|_\infty\le29e_N$.

The Schur midpoint is the same source midpoint as in V42,
\[
 D_m=(C_W)_m-\tfrac12(L_m^TY+Y^TL_m),
\]
where $Y$ is its fixed proposal for $A^{-1}L$. The retained source
balls and their residual calculation give
$\|D-D_m\|_\infty\le e_D<2^{-48}$. Write
$n=\|N_0\|_\infty$, $n_T=\|N_0^T\|_\infty$, and
$d=\|D_m\|_\infty$. The identity obtained by changing the three
factors one at a time gives
\begin{equation}
 \|N^TDN-N_0^TD_mN_0\|_\infty\le
 e_G:=29e_N(d+e_D)(n+e_N)+n_Te_D(n+e_N)+n_Tde_N.
 \label{v45:eq:eG}
\end{equation}
For the fixed rational $5\times5$ congruence proposal $S_0$, the
executed certificate proves
\begin{equation}
 \boxed{
 \|S_0^TN_0^TD_mN_0S_0-\operatorname{diag}(-1,1,1,1,1)\|_\infty
       +\|S_0^T\|_\infty e_G\|S_0\|_\infty<2^{-27}.}
 \label{v45:eq:congruence}
\end{equation}
Every actual congruence is symmetric. Its spectral norm distance from
the displayed sign matrix is less than one, so the straight symmetric
homotopy cannot cross zero. The same inequality proves $S_0$ invertible:
a singular congruence would be singular, contradicting its distance
from an invertible unit sign matrix. Sylvester inertia then proves
\eqref{v45:eq:kernel-inertia}.

For orientation only, the exact computed bounds are approximately
$q_{24}=5.014\,10^{-10}$, $e_N=1.051\,10^{-13}$,
$e_D=2.795\,10^{-15}$, and the left side of
\eqref{v45:eq:congruence} is $4.083\,10^{-9}$. The strict
rational inequalities, not these rounded displays, are the certificate.
The parent active and full graded signed congruences are also rechecked
on the same inherited source balls.
\end{proof}

The source balls in this calculation are the already established V42
entrywise enclosures. The current executable replays the root inclusion,
reconstructs all ensuing matrices from the exact centres and radii, and
checks the new kernel identities and residuals. Their provenance is
explicit: reusing those source balls is not a claim to have freshly
regenerated all stationary-source entries. A separate original weak
coefficient regeneration is recorded with the package.

\section{The entire nonfast weak evolution is bounded}
\label{v45:sec:spectrum}

\begin{theorem}[Exact weak spectral decomposition]
\label[theorem]{v45:thm:spectrum}
At the retained algebraic root, the weak space has an invariant direct
sum
\begin{equation}
 \R^{29}=E_+\oplus E_-\oplus E_0\oplus E_{\rm osc},
 \qquad (\dim E_+,\dim E_-,\dim E_0,\dim E_{\rm osc})=(9,9,5,6).
 \label{v45:eq:splitting}
\end{equation}
The space $E_+$ is exactly the fast space of V43. Its eigenvalues have
real parts at least $1/61$, and those on $E_-$ have real parts at most
$-1/61$. The space $E_0=\ker B$ is semisimple at eigenvalue zero.
On $E_{\rm osc}$ the form $D$ is positive definite and $B$ is skew-adjoint
for that positive form. Thus its six nonzero eigenvalues are purely
imaginary and semisimple over $\C$.

For the original nonfast weak restriction $B_{\rm c}$, there exists a
fixed finite constant $M_0$ such that
\begin{equation}
 \boxed{\|e^{sB_{\rm c}}\|\le M_0\quad(s\ge0).}
 \label{v45:eq:weak-bounded}
\end{equation}
No candidate nonlinear damping assertion is involved.
\end{theorem}
\begin{proof}
Since $B^TD+DB=0$ and $D$ is invertible, $B$ is similar to $-B^T$.
The spectrum and its algebraic multiplicities are invariant under
$\lambda\mapsto-\lambda$. V43 isolates exactly nine eigenvalues
with real part greater than $1/1000$, on a space whose backward
semigroup decays at rate $1/61$. Hence their negatives form an
invariant space $E_-$ of dimension nine. All remaining eigenvalues
have real parts between $-1/1000$ and $1/1000$.

A useful elementary orthogonality fact is the following. On two spectral
subspaces whose eigenvalue sums never vanish, the cross restriction
$D_{12}$ satisfies the Sylvester equation
$B_1^TD_{12}+D_{12}B_2=0$ and is zero. This follows by triangularizing
the two matrices over $\C$ and solving for the entries in triangular
order; the diagonal divisors are their eigenvalue sums. It applies
to $E_+\oplus E_-$ and the remaining spectral space. Their sum
$H=E_+\oplus E_-$ is therefore $D$-nondegenerate. Each of $E_+$
and $E_-$ is $D$-isotropic: preserve the $D$ pairing under evolution
and use backward decay on $E_+$ or forward decay on $E_-$. Thus
$D|_H$ has inertia $(9,9,0)$.

The kernel $E_0$ has dimension five and inertia $(4,1,0)$ by the two
new certificates. Its $D$-orthogonal complement is invariant. Also
\[
 \Ran B=(\ker B)^{\perp_D},
\]
since $D B$ is skew and both spaces have dimension 24. Their
intersection with $\ker B$ is zero by nondegeneracy on the kernel.
This proves semisimplicity at zero, without inferring it from rounded
eigenvalues.

The spaces $H$ and $E_0$ are $D$-orthogonal. Define
$E_{\rm osc}=(H\oplus E_0)^{\perp_D}$. Its dimension is six and
its inertia is
\[
       (19,10)-(9,9)-(4,1)=(6,0).
\]
It is invariant and contains no kernel vector. Conjugation by the
positive square root of $D|_{E_{\rm osc}}$ makes its generator an
ordinary real skew-symmetric matrix. This proves its asserted spectrum,
semisimplicity and bounded evolution. The evolution on $E_0$ is the
identity. On $E_-$, the nondegenerate cross pairing with $E_+$
identifies its generator with the negative transpose of the fast
restriction; its forward semigroup decays with a finite norm-conversion
constant. Combining these fixed norms proves \eqref{v45:eq:weak-bounded}.
The construction also works on the retained cube, with finite uniform
constants after fixing its spectral charts; no numerical value of
$M_0$ is required or supplied.
\end{proof}

The six imaginary eigenvalues are three conjugate pairs counted with
multiplicity. No assertion of distinct frequencies is needed. The one
negative signed direction in $E_0$ is retained. It is stationary at
this linear level and is not replaced by a positive kinetic term.
The original nine growing directions also remain present outside the
selected invariant graph.

\section{The full complementary state has at most linear growth}
\label{v45:sec:full-semigroup}

The full layer matrix is not just $B$: it is the inherited triangular
matrix on $402$ canonical/active and $29$ weak coordinates,
\begin{equation}
 \mathbb L=\begin{pmatrix}0&0\\ \mathcal R&B\end{pmatrix}.
 \label{v45:eq:full-L}
\end{equation}
Removing its nine-dimensional fast spectral space gives the same
$422$-dimensional complementary restriction $L_{\rm c}$ used in V44.
It is important not to promote \eqref{v45:eq:weak-bounded} to a bounded
full semigroup without checking the off-diagonal term.

\begin{proposition}[A polynomial bound with the zero-mode coupling retained]
\label[proposition]{v45:prop:polynomial-semigroup}
There is a fixed finite $M_1$ such that
\begin{equation}
 \boxed{\|e^{sL_{\rm c}}\|\le M_1(1+s)\quad(s\ge0).}
 \label{v45:eq:linear-growth}
\end{equation}
In spectral coordinates the only potentially linearly growing term is
$s\Pi_0\mathcal R_{\rm c}$, where $\Pi_0$ projects the nonfast weak
space onto its five-dimensional zero eigenspace. No vanishing of this
source coupling is assumed.
\end{proposition}
\begin{proof}
Write $\Pi_+$ for the fast projector on the weak space and
$\Pi_{\rm c}=I-\Pi_+$. The full complementary space is parametrized
by $(q,w)\in\R^{402}\oplus\Ran\Pi_{\rm c}$ through
\[
 (q,w)\longmapsto
 (q,w-B_+^{-1}\Pi_+\mathcal Rq).
\]
This is a fixed bounded linear isomorphism. In these coordinates the
linear equations are
$q'=0$, $w'=B_{\rm c}w+\mathcal R_{\rm c}q$, with
$\mathcal R_{\rm c}=\Pi_{\rm c}\mathcal R$. Thus
\[
 e^{sL_{\rm c}}
 \sim
 \begin{pmatrix}
 I&0\\
 (\int_0^s e^{rB_{\rm c}}\,\dd r)\mathcal R_{\rm c}
       &e^{sB_{\rm c}}
 \end{pmatrix}.
\]
The weak bound immediately gives \eqref{v45:eq:linear-growth}.
More precisely, the restriction $B_{\ne0}$ to $E_-\oplus E_{\rm osc}$
is invertible, so
\begin{equation}
 \int_0^s e^{rB_{\rm c}}\,\dd r
 =s\Pi_0+B_{\ne0}^{-1}(e^{sB_{\rm c}}-I)(I-\Pi_0).
 \label{v45:eq:primitive}
\end{equation}
The second term is bounded. This proves the last assertion and retains
all zero-mode forcing rather than deleting it. In particular the full
zero Jordan blocks have length at most two. Their exact number, and
whether the indicated coupling vanishes on further constraint tangents,
are not computed here.
\end{proof}

\section{Original histories on a polynomial physical-time interval}
\label{v45:sec:time}

Use exactly the V44 graph $u=x+h(a,x)$ and its selected initial
records; the graph need not be reconstructed with new spectral weights.
Its original reduced field is
$g_a(x)=P_{\rm c}\mathcal G(a,x+h(a,x))$. The retained regularity gives
\begin{equation}
 \begin{gathered}
 D_xg_a(x)=L_{\rm c}+O(a+\|x\|),\qquad
 D_xh(a,x)=O(a+\|x\|),\\
 x(0)=O(a^2),\qquad v(0)=b+O(a),\qquad
 v=x'/a^3,\qquad L_{\rm c}b=0.
 \end{gathered}
 \label{v45:eq:reduced-data}
\end{equation}
All these are properties of the previously selected original records.
No further derivative or preparation coefficient is assigned a value.

\begin{theorem}[A polynomial lower lifespan for the same invariant family]
\label[theorem]{v45:thm:time}
There exist constants $s_*>0$, $a_*>0$, and $C<\infty$ such that for
$0<a<a_*$ the selected unique original histories exist on
\begin{equation}
 \boxed{0\le t\le s_*a^{5/2},\qquad
              0\le s=t/a^3\le s_*a^{-1/2}.}
 \label{v45:eq:time}
\end{equation}
They remain in the original stationary and signed-regular chart, on
the selected invariant graph, and preserve all original constraints.
For $z(s)=(u(s)-u(0))/a^3$ one has
\begin{align}
 \|u(s)\|&\le C a^2,\label{v45:eq:u}\\
 \|z'(s)-b\|+\|z''(s)\|&\le C a(1+s)^2,\label{v45:eq:velocity}\\
 \|z(s)-sb\|&\le C a\bigl((1+s)^3-1\bigr).
 \label{v45:eq:position}
\end{align}
The constants are not numerically enclosed and are not claimed uniform
as the spatial cutoff changes.
\end{theorem}
\begin{proof}
On the graph, differentiation of the original reduced field gives
\[
                 x'=a^3v,\qquad v'=D_xg_a(x)v.
\]
Bootstrap $\|v-b\|\le1$ and remain inside the graph chart. For
$s\le s_*a^{-1/2}$ with $s_*\le1$,
\[
 \|x(s)\|\le C_0a^2+a^3s(\|b\|+1)\le C_1a^2.
\]
Consequently $E_a(s)=D_xg_a(x(s))-L_{\rm c}$ obeys
$\|E_a(s)\|\le K a$, with $K$ fixed on the bootstrap tube.
For $r=v-b$, variation of constants, \eqref{v45:eq:linear-growth},
and $L_{\rm c}b=0$ imply
\begin{align*}
 \|r(s)\|
 &\le M_1C_0a(1+s)
   +M_1Ka\int_0^s(1+s-t)(\|b\|+\|r(t)\|)\,\dd t\\
 &\le C_2a(1+s)^2.
\end{align*}
Here the second inequality uses only the bootstrap bound, not a
Gronwall exponential. Choose $s_*$ so that $C_2s_*^2<1/8$,
and then $a_*$ so that
$C_2(a+2s_*\sqrt a+s_*^2)<1/2$. This closes the bootstrap
strictly on the whole proposed interval.

The graph-jet estimate of V44 gives $h(a,x)=O(a^3)$ when
$x=O(a^2)$. Therefore $u=O(a^2)$, so the selected trajectory cannot
reach the boundary of the fixed original chart. Positivity of lapse
and metric and signed regularity persist, and finite-dimensional
continuation supplies the whole interval. Local invariance used here
is that of the original field; the auxiliary extension is never reached.
Constraint preservation is its exact identity $\dot P=0$.

The full normalized velocity is
$z'=v+D_xh(a,x)v$, so \eqref{v45:eq:reduced-data} and the bound
on $r$ prove the first part of \eqref{v45:eq:velocity}. Next,
$z''=D_u\mathcal G(a,u)z'$, $\mathbb Lb=0$, and
$D_u\mathcal G(a,u)-\mathbb L=O(a+\|u\|)=O(a)$ give its
second part. Finally integrate $z'-b$ from zero to obtain
\eqref{v45:eq:position}. No fixed-$s$ asymptotic has been substituted
at a growing endpoint.
\end{proof}

For every fixed $c>0$,
$s_*a^{5/2}/(ca^3\log(1/a))\to\infty$. Thus
\eqref{v45:eq:time} is longer than every fixed-coefficient V44
logarithmic interval, not merely a larger numerical prefactor on that
same interval. It still tends to zero in the original physical time.

\section{Physical observations and the cost of a longer selected interval}
\label{v45:sec:observations}

\begin{corollary}[Rate and curvature control]
\label[corollary]{v45:cor:observations}
On \eqref{v45:eq:time},
\begin{equation}
 \boxed{\|\lambda_{t,W}\|+\|\lambda_{t,A}-av_A\|
                  \le Ca^2(1+s)^2.}
 \label{v45:eq:rates}
\end{equation}
Each of the original metric Riemann tensor, literal electric link
curvature and literal spatial link curvature remains $O(a)$. For
each observation separately,
\begin{equation}
 \boxed{\|\mathscr O_a(a^3s)-\mathscr O_a(0)\|
                       \le Ca^2(1+s)^2.}
 \label{v45:eq:observations}
\end{equation}
If $0\le\varepsilon(a)\le s_*$ and $\varepsilon(a)\to0$, then
on $t\le a^{5/2}\varepsilon(a)$ the weak rate and every displayed
curvature change are $o(a)$. In particular, for sufficiently small $a$,
\begin{equation}
 \boxed{
 T_a=a^{11/4},\qquad
 \lambda_{t,W}=O(a^{3/2}),\qquad
 \mathscr O_a(t)-\mathscr O_a(0)=O(a^{3/2})
 }
 \label{v45:eq:example}
\end{equation}
uniformly on that interval. Its scaled position error is $O(a^{1/4})$.
\end{corollary}
\begin{proof}
The physical multiplier rate is $a$ times its normalized displacement
velocity. Since $b_W=0$ and $b_A=v_A$, \eqref{v45:eq:velocity}
gives \eqref{v45:eq:rates}, including the $O(a^2)$ initial weak
rate of the selected family. Original state changes have size
$O(a^4s+a^5(1+s)^3)$, by \eqref{v45:eq:position}; on the present
interval this is bounded by $Ca^2(1+s)^2$. Each retained curvature
formula is a smooth fixed-cutoff function of $(X,\lambda,F,\lambda_t)$
with no $\lambda_{tt}$ slot. Its initial size is $O(a)$, and the
same Lipschitz argument as in V44 proves \eqref{v45:eq:observations}.
At $s\le\varepsilon(a)/\sqrt a$, the bound divided by $a$ is at
most $C(\sqrt a+\varepsilon(a))^2$. This proves the $o(a)$ claim.
For the example, use $s\le a^{-1/4}$; this lies inside
\eqref{v45:eq:time} for all sufficiently small $a$.
\end{proof}

On the full endpoint $s_*a^{5/2}$ the conclusion is $O(a)$,
not necessarily $o(a)$, for the curvature change. This distinction is
not erased by renaming the comparison variable. The two curvature
observations are not identified coefficient by coefficient, and a
uniform physical derivative of curvature is not asserted.

\begin{proposition}[A necessary preparation-accuracy consequence]
\label[proposition]{v45:prop:accuracy}
For a perturbed member of the exact rate-access family, with fast
coefficient $w(a)+e(a)$ in the V44 local control chart, remaining
inside that same fixed normalized neighborhood up to
$t=s_*a^{5/2}$ necessarily requires
\begin{equation}
 \|e(a)\|\le C a^{-4}
           \exp\!\left(-\frac{s_*}{122\sqrt a}\right).
 \label{v45:eq:accuracy}
\end{equation}
This is a necessary local-state condition, not a sufficient accuracy
threshold or a curvature lower bound.
\end{proposition}
\begin{proof}
V44 gives initial normal distance at least $c_0a^4\|e(a)\|$
and growth at least $e^{s/122}$ while the history stays in that
neighborhood. Its normal distance is bounded there by a fixed
constant. Evaluate at $s=s_*a^{-1/2}$ and rearrange. The proposition
is a consequence of the inherited repulsion theorem; it is not an
additional source coefficient calculation.
\end{proof}

\section{The remaining zero-mode question and verification scope}
\label{v45:sec:status}

The small positive complementary exponent has been removed by an
original-source calculation. The result is not that the full linear
system has become coercive: five zero modes and their possible
coupling to the canonical/active variables remain. Formula
\eqref{v45:eq:primitive} specifies the finite map
$\Pi_0\mathcal R_{\rm c}$ whose vanishing would remove the linear
semigroup growth. It has not been shown to vanish, either on the full
state or on a further constrained tangent. The amplitude derivative of
the reduced field on this zero-mode coupling is likewise not computed.

A simple comparison shows why this next step matters. Set
\[
 L_0=\begin{pmatrix}0&1\\0&0\end{pmatrix},\qquad
 L_a=L_0+\begin{pmatrix}0&0\\a&0\end{pmatrix},\qquad b=(1,0)^T.
\]
Then $L_0b=0$, $\|L_a-L_0\|=a$, but the solution of
$v'=L_av$, $v(0)=b$ is
\[
 v(s)=(\cosh(\sqrt a\,s),\ \sqrt a\sinh(\sqrt a\,s))^T.
\]
Its first-component error is $as^2/2+O(a^2s^4)$ on bounded
$\sqrt a\,s$ sets. Thus the available $O(a)$ derivative
perturbation and linear semigroup bound alone do not justify a
longer $s\asymp a^{-1}$ interval with vanishing normalized velocity
error. This is an illustrative matrix, not a claim about an eigenvalue
of the actual reduced nonlinear field.

\paragraph{Proved here.}
An exact cubic weak-Poisson kernel identity, the root-cube rank $24$,
and the signed kernel inertia $(4,1,0)$ establish five semisimple zero
modes and six positive-energy oscillatory directions. The weak
nonfast semigroup is bounded and the full complementary semigroup
grows at most linearly. The unchanged V44 invariant initial family
therefore has original histories through $s_*a^{5/2}$ with $O(a)$
curvature, and $o(a)$ rate/curvature changes on the stated slightly
shorter windows. No new initial selection is needed.

\paragraph{Inherited inputs and executed checks.}
The exact root inclusion, original source normalization, signed total
inertia, fast-space certificate, analytic invariant graph and selected
initializer are retained at their stated scope. The new verifier
checks the cubic polynomial identity coefficientwise and the exact
kernel/inertia residual inequalities. It replays the V39 root inclusion
and rechecks the earlier signed and fast certificates using the
recovered V42 source-centre enclosures. It does not count those
inherited enclosures as a fresh $17334$-entry stationary-source
regeneration. The separately executed original weak-bank regeneration
and its equality checks are reported in the diagnostic record.
The complete nonlinear graph, higher initializer and original
spacetime histories are not numerically tabulated. The continuation
is a written finite-dimensional analytic proof, not an ODE experiment
or a proof-assistant formalization.

\paragraph{Unresolved.}
The constants $s_*,a_*,M_0,M_1$ and the local chart radius have no
numerical enclosure. No spectral damping of nonlinear center dynamics,
fixed positive physical lifespan, spatial-refinement estimate,
attracting microscopic selection, or full Einstein--Standard-Model
continuum theorem is asserted. The normal repulsion and the independent
V32 unit-initial-multiplier stress problem remain unchanged.

\begingroup
\renewcommand{\refname}{Additional reference for Version 45}
\begin{thebibliography}{9}
\small\raggedright
\bibitem{v45:KM}
R.~Koll\'ar and P.~D.~Miller,
\emph{Graphical Krein Signature Theory and Evans--Krein Functions},
arXiv:1209.3185 (2012). \url{https://arxiv.org/abs/1209.3185}.
Primary-source context for indefinite quadratic forms and Hamiltonian
spectral signatures. The finite signed orthogonality and inertia
argument needed for the present source is proved above.
\end{thebibliography}
\endgroup
\addtocontents{toc}{\protect\endgroup}
% END IMMUTABLE VERSION 45 BODY
% APPEND VERSION 46 HERE, BEFORE THE FINAL END-DOCUMENT COMMAND.


% BEGIN IMMUTABLE VERSION 46 BODY
\clearpage
\begin{center}
{\Large\bfseries Version 46: Two neutral source channels and a constraint-compatible longer interval}\\[0.5em]
{\large Exact zero-mode coupling, a graded active-coordinate change, and original histories through order $a^2$}
\end{center}
\makeatletter
\addtocontents{toc}{\protect\begingroup}
\addtocontents{toc}{\protect\makeatletter}
\addtocontents{toc}{\protect\def\protect\l@section{\protect\bfseries\protect\@dottedtocline{1}{0em}{2.5em}}}
\makeatother
\addcontentsline{toc}{part}{Version 46: Neutral forcing and constraint-compatible polynomial continuation}

\section{The zero-mode question and the unchanged selected family}
\label{v46:sec:context}

All statements in this body concern the original cutoff $h=1/3$, the
mixed-mode algebraic root of Version 39, and the same invariant initial
family selected in Version 44. The amplitude is $a>0$, and the physical
clock is still $t=a^3s$. No new first jet, initial-rate selection,
constraint projection, or field equation is introduced.

Version 45 leaves two questions distinct. Its five weak zero modes are
semisimple, but forcing from the canonical and active variables can still
produce a length-two zero Jordan block in the full layer matrix. Also,
an arbitrary $O(a)$ perturbation of such a block permits square-root
spectral scales; its example explains why the estimates used there stop
at $s\asymp a^{-1/2}$. The present body first evaluates the source
coupling on the flat nonconstant constraint tangent. It has rank exactly
two and is geometrically visible. We therefore do not try to remove
linear growth by declaring that coupling zero.

The useful upstream change is instead to use the exact nonconstant
constraints together with the exact active consistency equation. In
those coordinates the canonical/active derivative block is $O(a^2)$,
not merely $O(a)$. This structured estimate gives original histories
through $t=s_*a^2$ on the existing selected family. The argument is
finite-dimensional and retains every original mean constraint along
these histories. The constant $s_*>0$ is not numerically enclosed.
The distinction between constraints, multiplier determination, and a
modified evolution is standard in constrained discretization; see
\cite{v46:DBGP} for background, not for the estimates proved here.

The inputs reused without new coefficient regeneration are the V39 root
inclusion and full first-Poisson rank, the V41 signed stationary-source
certificate, the V44 locally invariant graph and its selected initializer,
and the V45 weak kernel and bounded nonfast weak semigroup. The new
finite checks are specified below. In particular no new high-order
stationary coefficient or nonlinear spacetime trajectory is reported.

\section{Exactly two neutral forcing channels survive the flat constraints}
\label{v46:sec:coupling}

Use $D$, $\mathsf K$, $B=-D^{-1}\mathsf K$, and the exact $29\times5$
kernel synthesis $N$ of \cref{v45:eq:N}. Its first three columns
synthesize constant shifts. Put
\begin{equation}
 H_0=N^TDN,\qquad
 \Pi_0=N H_0^{-1}N^TD.
 \label{v46:eq:Pi0}
\end{equation}
Since $D$ is nondegenerate on $\ker B$, this is the commuting spectral
projector onto $\ker B$, not an orthogonal projection in an arbitrarily
chosen positive norm.

Let $\mathcal L_{\rm con}$ be the original linear constraint map on
canonical fields $\xi=(U,p)$. In the normalized spatial pairing its
components are
\begin{equation}
 (\mathcal L_{\rm con}\xi)_N
  =-\delta_i\delta_jU_{ij}+\Delta_\delta\operatorname{tr}U,
 \qquad
 (\mathcal L_{\rm con}\xi)_{\beta,j}=-2\delta_i p_{ij}.
 \label{v46:eq:flat-constraint}
\end{equation}
All four component means vanish. Define
\[
 \mathcal Z_0=\{(\xi,\eta_A):\mathcal L_{\rm con}\xi=0\}.
\]
The full linear source map retained in V42 is
\begin{equation}
 \mathcal R(\xi,\eta_A)
 =-D^{-1}\!\left\{
    \mathcal C_\ell\xi+\mathcal K_W I_A\eta_A
       -L_*^TA_*^{-1}R_A\mathsf B_1\xi\right\}.
 \label{v46:eq:R}
\end{equation}
The active stationary cross block $L_*$ in this formula is not the
linear constraint map $\mathcal L_{\rm con}$.

\begin{theorem}[The source-weighted neutral coupling has rank two]
\label[theorem]{v46:thm:rank-two}
At the retained exact root,
\begin{equation}
 \boxed{\operatorname{rank}\bigl(\Pi_0\mathcal R|_{\mathcal Z_0}\bigr)=2.}
 \label{v46:eq:rank-two}
\end{equation}
Its range is exactly the $D$-orthogonal complement of the constant
shifts inside $\ker B$. The restriction of $D$ to that two-dimensional
range has inertia $(1,1,0)$, and the original weak geometric observation
is injective there. In particular the coupling is not removed by
imposing the flat nonconstant constraints.
\end{theorem}
\begin{proof}
Multiplication by \eqref{v46:eq:Pi0} gives the exact identity
\begin{equation}
 \Pi_0\mathcal R q=-N H_0^{-1}N^T
 \left\{\mathcal C_\ell\xi+\mathcal K_W I_A\eta_A
       -L_*^TA_*^{-1}R_A\mathsf B_1\xi\right\}.
 \label{v46:eq:coupling-formula}
\end{equation}
The inherited flat identity says that $\mathsf B_1$ factors through
$\mathcal L_{\rm con}$. For a constant shift $c$, the original quadratic
translation Hamiltonian has
$\{P_{c,2},H_2\}=0$ on all canonical fields. Hence
$Dd_{2,c}(Y)[\xi]=0$. Its first Poisson bracket is
\[
 \mathbb K_1(\xi)[c,\ell]
       =\langle\delta_c\ell,\mathcal L_{\rm con}\xi\rangle_h.
\]
Thus the first three entries of the last covector in
\eqref{v46:eq:coupling-formula} vanish on $\mathcal Z_0$; the active
term vanishes on those entries because
$\mathbb K_1(Y)[c,\mu]=\langle\delta_c\mu,
\mathcal L_{\rm con}Y\rangle_h=0$. Its rank is at most two.

For the reverse inequality it is already enough to set $\xi=0$ and
vary only the lapse. A basis-free check uses the inherited full Poisson
kernel, which consists of the four constant multipliers. The map
$k\mapsto\mathbb K_1(Y)[I_Wk,\cdot]$ on $\ker\mathsf K$ has
kernel precisely the three constant shifts. It therefore has rank two;
it annihilates weak and constant-lapse arguments and is detected on the
active complement. The explicit two-lapse certificate below separately
verifies this rank with the actual root data.

Let $I_{\rm c}:\mathbb R^3\to\mathbb R^{29}$ synthesize constant
weak shifts. The same certificate proves
\begin{equation}
                         I_{\rm c}^TD I_{\rm c}\succ26000 I_3
 \label{v46:eq:Dcc}
\end{equation}
in the retained spatial-sum coordinates. Equation
\eqref{v46:eq:coupling-formula} shows that the range is $D$-orthogonal
to $I_{\rm c}\mathbb R^3$. The latter is positive and nondegenerate.
Its complement in a kernel of inertia $(4,1,0)$ has dimension two and
inertia $(1,1,0)$, so the rank just proved makes the range equal to that
complement. It intersects the constant shifts only at zero. The
inherited literal weak shift-curvature observation has exactly those
constants as its kernel, proving injectivity. The signed inertia here
is not a physical energy or a dynamical instability assertion.
\end{proof}

\begin{proposition}[Two-column original-source enclosure]
\label[proposition]{v46:prop:minor}
Let $E=(e_2-\mathbf1_N/27,e_{25}-\mathbf1_N/27)$ denote two
mean-zero lapse point tests, zero in the shift slots. With $N_2$ the
last two columns of $N$, the exact matrix
\[
              T=N_2^T\mathsf W^T\widehat K E
\]
satisfies, on the entire retained $2^{-112}$ root cube,
\begin{equation}
                  \boxed{7244000<|\det T|<7244200.}
 \label{v46:eq:minor}
\end{equation}
Together with \eqref{v46:eq:Dcc}, this gives a root-specific check of
both nonzero neutral channels and their nonconstant observation.
\end{proposition}
\begin{proof}
Use the exact source midpoint $K_0$ from the six inherited Poisson
coefficient matrices and the dyadic kernel proposal $N_{2,0}$ already
enclosed in V45. If $e_K$ bounds $\|\widehat K-K_0\|_\infty$ and
$e_N$ bounds $\|N-N_0\|_\infty$, set
\begin{align*}
 T_0&=N_{2,0}^T\mathsf W^TK_0E,\\
 e_T&=29e_N\|\mathsf W^T\|_\infty(\|K_0\|_\infty+e_K)
                          \|E\|_\infty
       +\|N_{2,0}^T\|_\infty\|\mathsf W^T\|_\infty
                                      e_K\|E\|_\infty.
\end{align*}
This bounds $\|T-T_0\|_\infty$. Entrywise expansion of a two-by-two
determinant gives
\[
 |\det T-\det T_0|
       \le e_T\sum_{i,j}|(T_0)_{ij}|+2e_T^2.
\]
The exact rational checks give $e_T<2\,10^{-6}$ and determinant
error less than $1/100$. The midpoint determinant is approximately
$-7244078.072879051$, but only its exact rational value and the displayed
error are used to prove \eqref{v46:eq:minor}. The physical square-root
normalization cancels for these lapse tests:
$\mathsf W^T\widehat K E=\mathsf W^T\mathbb K_1 E$.

For \eqref{v46:eq:Dcc}, use the symmetrized rational Schur midpoint
$D_0$ and its inherited proven norm error $e_D$. Direct rational
Gershgorin bounds give
\[
 \min_{i<3}\left((D_0)_{ii}
            -\sum_{\substack{j<3\\j\ne i}}|(D_0)_{ij}|\right)-e_D
                                      >26000.
\]
The floating display is about $26264.4276675$. This proves positivity
for the exact symmetric block, without inferring inertia from a
floating eigensolve. No fresh reconstruction of the $17334$
stationary-source entries is represented by these consequences of their
inherited enclosures.
\end{proof}

\begin{corollary}[The constrained linear sector has two genuine zero Jordan chains]
\label[corollary]{v46:cor:Jordan}
On the flat constraint sector
$\{(\xi,\eta_A,w):\mathcal L_{\rm con}\xi=0\}$, the full layer
matrix has exactly two length-two zero Jordan blocks. All other zero
blocks have length one. Its nonfast restriction has dimension $318$:
there are $299$ length-one zero blocks, two length-two zero blocks,
nine decaying directions and six oscillatory directions. Its semigroup
is bounded above by $C(1+s)$ and is not bounded uniformly in $s$ in
any equivalent fixed positive norm.
\end{corollary}
\begin{proof}
At every nonzero spatial Fourier mode the scalar linear constraint
has rank one on symmetric $U$ and the vector constraint has rank
three on symmetric $p$. There are $26$ real nonconstant modes, so
$\operatorname{rank}\mathcal L_{\rm con}=104$ and its canonical
kernel has dimension $220$. The top block therefore has dimension
$220+78=298$. Split the weak space into its semisimple zero and
invertible nonzero parts. The bounded change
$w_{\ne0}\mapsto w_{\ne0}+B_{\ne0}^{-1}(I-\Pi_0)\mathcal Rq$
removes its constant forcing and leaves at zero
\[
              (q,w_0)'=(0,\Pi_0\mathcal Rq).
\]
This matrix squares to zero and has rank two by
\cref{v46:thm:rank-two}. Its zero algebraic dimension is $298+5=303$.
The asserted block count follows. The nine fast directions are in
the nonzero part and their removal leaves $318$ dimensions. A
nontrivial length-two chain gives a vector with norm growing at least
$cs-C$; changing to an equivalent norm cannot eliminate that growth.
These counts concern the flat nonconstant constraint tangent, not the
additional nonlinear mean-constraint tangent at nonzero amplitude.
\end{proof}

\section{An exact nonconstant-constraint chart, with the mean rows retained}
\label{v46:sec:constraint-chart}

Return to the regular normalized original field $\mathcal G(a,u)$ of
\cref{v44:lem:regular}. Write the base records as
$X_a^0=a\bar X(a)$ and $\lambda_a^0=e+a\bar\chi(a)$, and write
$u=(u_X,u_A,u_W)$ in the original fixed active/weak multiplier
synthesis. Let $\mathsf S:\mathbb R^{104}\to\mathbb R^{108}$
synthesize all componentwise mean-zero multiplier tests. Define
\begin{equation}
 \mathscr P(a,u)=\frac1a\mathsf S^T
 P\bigl(X_a^0+au_X,\lambda_a^0+a(I_Au_A+I_Wu_W)\bigr).
 \label{v46:eq:divided-constraints}
\end{equation}
At $a=0$ use its analytic extension. The pairing conversion from
spatial sum to mean is common and changes none of its zeros.

\begin{lemma}[Elimination of the nonconstant rows is analytic and invariant]
\label[lemma]{v46:lem:constraint-chart}
There is an analytic fixed-size chart for $\mathscr P=0$ near
$(a,u)=(0,0)$. It is parameterized by the $220$ components of
$\Pi_{\ker\mathcal L_{\rm con}}u_X$ and the $107$ normalized
multiplier components. The original normalized field is tangent to
this chart for small positive $a$. On it,
\begin{equation}
              \boxed{\mathsf B_A(a,u)=a^2\,b_A(a,u)}
 \label{v46:eq:active-division}
\end{equation}
with an analytic divided function in these chart coordinates.
No equation for a constant constraint is discarded from a claimed
physical history.
\end{lemma}
\begin{proof}
The original joint Taylor expansions at $(X,\lambda)=(0,e)$ are
\[
 P(X,e+\chi)=\mathcal L_{\rm con}X+P_{\ge2}(X,\chi),\qquad
 \mathsf B(X,e+\chi)=\mathsf B_1X+\mathsf B_{\ge2}(X,\chi).
\]
Their remainders have joint order at least two. In particular there
is no linear multiplier-only drift: the flat constraint Hamiltonian
vectors obey $\mathcal L_{\rm con}G=0$. Division in
\eqref{v46:eq:divided-constraints} is analytic and
\[
 \mathscr P(0,u)=\mathsf S^T\mathcal L_{\rm con}u_X.
\]
Its derivative on any fixed complement of the canonical kernel is an
isomorphism of dimension $104$. The analytic implicit-function theorem
eliminates these canonical components without a singular multiplier
inverse. This is a coordinate description of the constraint set, not
a projection applied to the dynamics.

For tangency, the original homogeneity identities give
$M\lambda=0$ and $\langle\lambda,\mathsf B\rangle_h=0$ throughout
the regular source chart. The normalized signed solve therefore
satisfies $\mathsf B+M\lambda_t=0$ on all rows, hence $P$ is constant
along that field on every constraint leaf. Its nonconstant zero leaf
is invariant. The selected initial data also have all four constant
rows zero, which remain zero by the same identity.

Finally $\mathsf B_1=\mathsf J_0\mathcal L_{\rm con}$, where
$\mathsf J_0$ has only a lapse output and takes the phase divergence
of the momentum-constraint output. This follows directly from
\eqref{v46:eq:flat-constraint} and the original quadratic field:
$\mathcal L_{\rm con}F_1=\mathsf J_0\mathcal L_{\rm con}$.
Since $\mathcal L_{\rm con}X$ has zero means, the equations
$\mathsf S^TP=0$ express it as the negative nonconstant projection
of $P_{\ge2}$. Substitution gives
\[
 \mathsf B_A
 =R_A\bigl(\mathsf B_{\ge2}
       -\mathsf J_0\operatorname{proj}_{\rm nc}P_{\ge2}\bigr).
\]
With $X=a(\bar X+u_X)$ and $\chi=a(\bar\chi+u_\lambda)$,
each term is divisible by $a^2$, proving
\eqref{v46:eq:active-division} with all chart derivatives.
\end{proof}

The new verifier independently forms the original flat maps on all
$324$ canonical coordinates and verifies
$\mathcal L_{\rm con}G=0$ and
$\mathcal L_{\rm con}F_1=\mathsf J_0\mathcal L_{\rm con}$ as
complete rational matrix identities. A nonzero modular $104$-minor,
together with the exact four constant annihilators, verifies the rank.
No constant constraint is inverted by a zero Fourier eigenvalue.

\section{The exact active equation removes the generic square-root feedback}
\label{v46:sec:shear}

Let $A(a,u)$, $L(a,u)$ and $C_W(a,u)$ be the analytic blocks of the
complete graded signed Hessian, so
\[
 M_{AA}=a^2A,\qquad M_{AW}=a^3L,\qquad M_{WW}=a^4C_W.
\]
The active block is invertible on the fixed normalized chart. Put
\begin{equation}
 \mathsf Y(a,u)=A(a,u)^{-1}L(a,u),\qquad
 \mathsf Y_*=A_*^{-1}L_*.
 \label{v46:eq:Y}
\end{equation}
Because $\mathcal G_\lambda=a^2\lambda_t$, the active original
consistency equation is exactly
\begin{equation}
 \boxed{A\mathcal G_A+aL\mathcal G_W=-\mathsf B_A.}
 \label{v46:eq:active-exact}
\end{equation}
Introduce the fixed-in-state analytical coordinate change
\begin{equation}
              p=u_A+a\mathsf Y_*u_W.
 \label{v46:eq:shear}
\end{equation}
It neither changes a record nor resets its multiplier. The original
variables are recovered by $u_A=p-a\mathsf Y_*u_W$.

\begin{proposition}[One additional amplitude factor in the upper derivative block]
\label[proposition]{v46:prop:upper}
On the nonconstant constraint chart, the top coordinates
$q=(\Pi_{\ker\mathcal L_{\rm con}}u_X,p)$ have derivative
\begin{equation}
 \begin{split}
 (\Pi_{\ker\mathcal L_{\rm con}}u_X)'&=
                        a^3\Pi_{\ker\mathcal L_{\rm con}}f(a,u),\\
 p'&=-a^2A^{-1}b_A+a(\mathsf Y_*-\mathsf Y)\mathcal G_W,
 \end{split}
 \label{v46:eq:top-field}
\end{equation}
where $f=F/a$ is analytic. For each fixed $K<\infty$, on
$\|u\|\le Ka^2$ the derivative of this top field with respect to
any uniformly regular constraint-chart coordinates has norm at most
$C_Ka^2$.
\end{proposition}
\begin{proof}
The canonical identity is $\mathcal G_X=a^2F=a^3f$; analyticity of
$f$ follows because the flat base canonical field vanishes. Subtracting
$a\mathsf Y\mathcal G_W$ using
\eqref{v46:eq:active-exact} and differentiating
\eqref{v46:eq:shear} gives the second exact identity in
\eqref{v46:eq:top-field}. Its first term is $a^2$ times a smooth
function by \cref{v46:lem:constraint-chart}.

Analyticity gives $\mathsf Y-\mathsf Y_*=O(a+\|u\|)$ and bounded
first derivatives. Also $\mathcal G_W(a,0)=O(a^4)$ for the fixed
$w=0$ base initializer, and $D_u\mathcal G_W=O(1)$; the weaker
bound $O(a^3+\|u\|)$ is sufficient. On $\|u\|\le Ka^2$, the
derivative of the remaining product is bounded by
\[
 Ca\bigl(\|\mathcal G_W\|+a+\|u\|\bigr)\le C_Ka^2.
\]
Coordinate maps with uniformly bounded derivatives preserve this
estimate. In the original active coordinates the term
$-a\mathsf Y_*\mathcal G_W$ has an $O(a)$ derivative. It was
this removable derivative term, together with the off-constraint
linear drift, that the unstructured V45 estimate had to pay.
\end{proof}

The two neutral Jordan chains from \cref{v46:cor:Jordan} remain at
$a=0$, because \eqref{v46:eq:shear} is then the identity. What has
changed is the estimate of the amplitude-dependent return into the
upper variables on the exact constraint leaf. The generic $O(a)$
Jordan perturbation in V45 does not have this $O(a^2)$ upper-block
property. We do not infer a spectral statement for every finite-
amplitude nonlinear trajectory from that distinction.

\section{Restriction of the same invariant graph to the exact constraint leaf}
\label{v46:sec:restricted-graph}

Write $\Pi_{\rm c}=I-\Pi_+$ for the nonfast weak projector. It has
range of dimension $20$, with bounded semigroup by V45. Use the
following coordinates on the intersection of the V44 graph and the
nonconstant constraint leaf:
\begin{equation}
 q=(\kappa,p),\quad
 \kappa=\Pi_{\ker\mathcal L_{\rm con}}u_X\in\mathbb R^{220},
 \quad p\in\mathbb R^{78},\quad
 w=\Pi_{\rm c}u_W\in\mathbb R^{20}.
 \label{v46:eq:coordinates}
\end{equation}
They are analytical coordinates, not additional physical readers.

\begin{lemma}[A regular $318$-coordinate description of the selected leaf]
\label[lemma]{v46:lem:intersection}
There is a $C^6$ local map $u=\Psi(a,q,w)$ describing that intersection,
with bounded fixed derivatives on a smaller neighborhood. It satisfies
$\Psi(a,0,0)=O(a^4)$. The induced original field
$(q',w')=(f_q(a,q,w),f_w(a,q,w))$ obeys, on
$\|(q,w)\|\le Ka^2$,
\begin{align}
 \|D f_q(a,q,w)\|&\le C_Ka^2,\label{v46:eq:upper-small}\\
 D f_w(a,q,w)&=(R_{\rm c},B_{\rm c})+O_K(a),
       \qquad\|e^{sB_{\rm c}}\|\le M_0.\label{v46:eq:lower-bounded}
\end{align}
For the unchanged selected initial records,
\begin{equation}
 (q_0,w_0)=O(a^2),\qquad
 \frac{q'(0)}{a^3}=b_q+O(a),\qquad
 \frac{w'(0)}{a^3}=O(a),\qquad R_{\rm c}b_q=0.
 \label{v46:eq:initial-data}
\end{equation}
Here $R_{\rm c}$ is the restriction of the inherited forcing to the
flat canonical constraint tangent followed by $\Pi_{\rm c}$.
\end{lemma}
\begin{proof}
Consider together the $104$ equations $\mathscr P=0$, the nine graph
equations $P_{\rm f}u-h(a,P_{\rm c}^{\rm full}u)=0$, and the
coordinates \eqref{v46:eq:coordinates}. At $(a,u)=(0,0)$ the first
block differentiates to $\mathsf S^T\mathcal L_{\rm con}u_X$.
The graph derivative is $P_{\rm f}u$, because $Dh(0,0)=0$.
The top canonical complement is thereby determined; the fast weak
component is then determined uniquely by the full spectral graph
$u_{W,+}=-B_+^{-1}\Pi_+\mathcal R(u_X,u_A)$. The remaining
coordinates are exactly \eqref{v46:eq:coordinates}. This block
linear map is an isomorphism. The finite-dimensional implicit-function
theorem supplies $\Psi$ with the stated regularity.

At $u=0$ the constraint and coordinate residuals are zero, while the
graph residual is $-h(a,0)=O(a^4)$ by V44. The uniformly bounded
implicit inverse gives $\Psi(a,0,0)=O(a^4)$. Both defining sets are
locally invariant for the original field, so their intersection is as
well. Its constant constraint rows are not required for this coordinate
chart; the selected records satisfy them and their own histories
preserve them.

At $a=0$ its induced linear field has zero top block and weak block
$R_{\rm c}q+B_{\rm c}w$. This follows either from the displayed
spectral graph or by projecting the full triangular layer system.
The $C^6$ bounds and $\Psi(a,0,0)=O(a^4)$ give
\eqref{v46:eq:lower-bounded}. Proposition~\ref{v46:prop:upper}
and the bounded coordinate derivatives give
\eqref{v46:eq:upper-small}. The old selected initializer has
$u(0)=O(a^2)$ and $\mathcal G(a,u(0))/a^3=b+O(a)$ with weak
leading component zero. Applying the coordinate map gives
\eqref{v46:eq:initial-data}; the additional term in $p$ is multiplied
by $a$ and does not change its leading velocity. Finally
$\mathbb Lb=0$ gives $R_{\rm c}b_q=0$.
\end{proof}

\section{Original-action histories through a physical interval of order \texorpdfstring{$a^2$}{a squared}}
\label{v46:sec:time}

\begin{theorem}[Constraint-compatible polynomial continuation]
\label[theorem]{v46:thm:time}
There are $s_*>0$, $a_*>0$, and finite $C$ such that the same V44
selected original histories exist for
\begin{equation}
             \boxed{0<a<a_*,\qquad 0\le s\le s_*/a,
                              \qquad 0\le t\le s_*a^2.}
 \label{v46:eq:time}
\end{equation}
They preserve all original constraints and remain on the same locally
invariant graph in the original stationary and signed-regular chart.
For their displacement $z(s)=(u(s)-u(0))/a^3$,
\begin{align}
 \|u(s)\|&\le Ca^2,\label{v46:eq:u}\\
 \|z'(s)-b\|+\|z''(s)\|&\le Ca(1+s),\label{v46:eq:velocity}\\
 \|z(s)-sb\|&\le Ca\,s(1+s).\label{v46:eq:position}
\end{align}
The constants and local chart radii are not numerically enclosed or
asserted uniform under spatial refinement.
\end{theorem}
\begin{proof}
Use the regular coordinates of \cref{v46:lem:intersection}, and set
$v_q=q'/a^3$, $v_w=w'/a^3$, $r_q=v_q-b_q$. Since the field is
autonomous at fixed amplitude, differentiating the actual ODE gives
\begin{align*}
 v_q'&=D f_q\,(v_q,v_w),\\
 v_w'&=R_{\rm c}r_q+B_{\rm c}v_w+
                E_a(s)(v_q,v_w),\qquad\|E_a(s)\|\le C_Ka.
\end{align*}
The term $R_{\rm c}b_q$ is zero, not estimated or deleted.

Choose a fixed velocity bootstrap $\|v_q\|+\|v_w\|\le B$ with
$B>2(\|b_q\|+1)$, and choose a fixed coordinate radius constant
$K$ larger than the initial $O(a^2)$ constant plus $2B$. For
$s\le s_*/a$, $s_*\le1$, direct integration gives
\[
 \|(q,w)(s)\|\le C_0a^2+a^3Bs\le(C_0+Bs_*)a^2<Ka^2.
\]
Thus the derivative bounds remain available throughout the bootstrap.
The upper equation and the initial data give
\begin{equation}
                 \|r_q(s)\|\le C_0a+C_1a^2s.
 \label{v46:eq:rq}
\end{equation}
Use the bounded \emph{weak} semigroup, rather than the linearly growing
full one, in the lower equation. It gives
\begin{align*}
 \|v_w(s)\|
 &\le M_0C_0a+
 M_0\int_0^s\bigl(\|R_{\rm c}\|\|r_q(r)\|+C_KaB\bigr)\,\dd r\\
 &\le C_2\bigl(a+as+a^2s^2\bigr).
\end{align*}
On $as\le s_*\le1$ the last term is at most $s_*as$, hence
\[
                 \|v_w(s)\|\le C_3a(1+s).
\]
Choose $s_*$ small enough, and then $a_*$ small enough, to make the
velocity bootstrap strict on its whole proposed interval. Equation
\eqref{v46:eq:rq} also gives $\|r_q\|\le C_4a$ there.

The coordinate reconstruction has bounded derivatives and
$\Psi(a,0,0)=O(a^4)$, so $u=O(a^2)$. The trajectory therefore
cannot exit the fixed original chart. Finite-dimensional continuation
establishes the whole interval. The selected graph and exact constraints
are invariant, and positivity and signed regularity persist there.

At $a=0$, $D_{(q,w)}\Psi(0,0,0)(b_q,0)=b$. On the bootstrap tube
this derivative differs from its base value by $O(a)$. Reconstructing
the velocity proves $\|z'-b\|\le Ca(1+s)$. Finally
$z''=D_u\mathcal G(a,u)z'$, $\mathbb Lb=0$, and
$D_u\mathcal G(a,u)=\mathbb L+O(a)$ give the acceleration bound
in \eqref{v46:eq:velocity}. Integrating the velocity difference gives
\eqref{v46:eq:position}. No Taylor expansion valid only at fixed $s$
has been substituted at a growing endpoint.
\end{proof}

The ratio of this interval to V45's $a^{5/2}$ interval tends to infinity.
This is a further polynomial extension for the same initial selection,
not a change of the physical clock or a new cancellation order. The
nonzero neutral forcing has been retained in \eqref{v46:eq:rq} and
in the integral for $v_w$.

\section{Physical observations and the strengthened preparation-accuracy cost}
\label{v46:sec:observations}

\begin{corollary}[Rate and two-curvature bounds]
\label[corollary]{v46:cor:observations}
On \eqref{v46:eq:time},
\begin{equation}
 \boxed{\|\lambda_{t,W}\|+\|\lambda_{t,A}-av_A\|
                                      \le Ca^2(1+s).}
 \label{v46:eq:rates}
\end{equation}
Each original metric Riemann tensor, literal electric link curvature,
and literal spatial link curvature is $O(a)$. For each observation
$\mathscr O$ separately,
\begin{equation}
 \boxed{\|\mathscr O_a(a^3s)-\mathscr O_a(0)\|
                                      \le Ca^2(1+s).}
 \label{v46:eq:curvature}
\end{equation}
In particular these changes and the weak rate are $o(a)$ on
$t\le a^2\varepsilon(a)$ whenever
$0\le\varepsilon(a)\le s_*$ and $\varepsilon(a)\to0$. A concrete
longer window is
\begin{equation}
 \boxed{T_a=a^{9/4},\qquad
 \lambda_{t,W}=O(a^{5/4}),\qquad
 \mathscr O_a(t)-\mathscr O_a(0)=O(a^{5/4}).}
 \label{v46:eq:example}
\end{equation}
On that window the normalized velocity error and the relative scaled
position error $\|z-sb\|/(1+s)$ are $O(a^{1/4})$.
\end{corollary}
\begin{proof}
Physical multiplier velocity equals $a z_\lambda'$. The leading weak
component of $b$ is zero, so \eqref{v46:eq:velocity} gives
\eqref{v46:eq:rates}, including the $O(a^2)$ initially selected weak
rate. The original state displacement is $a^4z$, bounded by
$C(a^4s+a^5s(1+s))$ and in particular by $Ca^3$ on the full interval.
The canonical field $F$ is uniformly smooth on the retained stationary
chart. Each curvature observation is the same smooth finite function
of $(X,\lambda,F,\lambda_t)$ used in V41--V45 and contains no
$\lambda_{tt}$ slot. Its initial size is $O(a)$. Comparing its arguments
with their own initial values proves \eqref{v46:eq:curvature}.
At $s\le\varepsilon(a)/a$ this bound divided by $a$ is at most
$C(a+\varepsilon(a))$. For the concrete example set
$\varepsilon(a)=a^{1/4}$. Equation \eqref{v46:eq:position} gives
the stated relative, rather than absolute, scaled-position estimate.
\end{proof}

At the full endpoint $s_*a^2$, curvature changes are proved $O(a)$,
not necessarily $o(a)$. On $a^{9/4}$ the absolute scaled position error
need not vanish; only the displayed relative error is asserted. The
corresponding original state error is small. No equality of the metric
and link curvatures, or uniform physical-time derivative of either, is
inferred.

\begin{proposition}[Normal repulsion still constrains preparation accuracy]
\label[proposition]{v46:prop:accuracy}
For a perturbed member of the same exact rate-access family, remaining
in the V44 fixed normalized neighborhood through $t=s_*a^2$ necessarily
requires its fast coefficient error to satisfy
\begin{equation}
              \|e(a)\|\le Ca^{-4}
                  \exp\!\left(-\frac{s_*}{122a}\right).
 \label{v46:eq:accuracy}
\end{equation}
This is not a sufficient accuracy condition or a lower bound on a
curvature component.
\end{proposition}
\begin{proof}
Reuse V44's exact lower initial normal displacement
$c_0a^4\|e(a)\|$ and normal growth at least $e^{s/122}$ while the
history remains in the neighborhood. The normal distance there is
bounded. Evaluate at $s=s_*/a$ and rearrange. No new repulsion
coefficient or nonlinear instability of the exactly selected family
is claimed.
\end{proof}

\section{What is closed, what is inherited, and the next actual source}
\label{v46:sec:status}

\paragraph{Proved here.}
The previously uncomputed neutral forcing has rank exactly two on the
flat nonconstant constraint tangent, is detected already by two lapse
tests, and has an injective original weak geometric observation. The
full restricted linear system really has two length-two zero Jordan
chains; a uniformly bounded full semigroup is not available. The exact
nonconstant constraints and a graded active-coordinate change instead
give an $O(a^2)$ upper derivative block. The same selected invariant
family therefore has unique original histories through $s_*a^2$, with
$O(a)$ curvatures and the stronger smaller-window estimates above.

\paragraph{Inherited and executed checks.}
The executable replays the V45 root, signed, weak-rank and spectral
certificates with their identified stationary-source enclosures. The
new checks bound the two-column neutral-source minor and the positive
constant-shift Schur block. They also construct the complete original
flat constraint, multiplier-vector, and quadratic-field matrices. The
identities $\mathcal L_{\rm con}G=0$ and
$\mathcal L_{\rm con}F_1=\mathsf J_0\mathcal L_{\rm con}$ are
checked on all entries over the rationals, and the rank $104$ is
checked with both a nonzero minor and the exact upper bound. The
optional original weak-bank regeneration is separately reported.
No fresh $17334$-entry stationary-source or $79224$-entry root-bank
regeneration is claimed. Small block examples check only the algebra
and exponent bookkeeping; the nonlinear longer-time result is the
written proof on the original constrained chart, not an ODE simulation.

\paragraph{Remaining question.}
On $s\asymp a^{-1}$ the two retained neutral forcing channels can
produce order-one normalized weak velocity changes. Their actual
amplitude-scale reduced forcing, including the selected initializer's
higher coefficients and the nonlinear motion on the graph, must be
evaluated before extending this argument to still longer intervals.
There is no generic square-root-return assumption left to remove, but
there is also no proved damping or bounded nonlinear center dynamics.
The graph is still normally repelling. No common positive physical
lifespan, spatial-refinement-uniform estimate, primitive selection, or
matter-coupled Einstein--Standard-Model continuum is asserted. The
independent unit-initial-multiplier stress question from V32 is unchanged.

The present coordinate changes do not supply a different physical
metric or a new action. Every observation is evaluated after restoring
the original canonical and multiplier coordinates, and every claimed
history satisfies all original constraints throughout its own interval.

\begingroup
\renewcommand{\refname}{Additional reference for Version 46}
\begin{thebibliography}{9}
\small\raggedright
\bibitem{v46:DBGP}
C.~Di Bartolo, R.~Gambini and J.~Pullin,
\emph{Consistent and mimetic discretizations in general relativity},
J. Math. Phys. \textbf{46} (2005), 032501.
\url{https://arxiv.org/abs/gr-qc/0404052}.
Primary-source context for simultaneous constraint and multiplier
consistency; not a source of the neutral-channel calculation, graded
coordinate estimate, or lifespan proved in this body.
\end{thebibliography}
\endgroup
\addtocontents{toc}{\protect\endgroup}
% END IMMUTABLE VERSION 46 BODY
% APPEND VERSION 47 HERE, BEFORE THE FINAL END-DOCUMENT COMMAND.


% BEGIN IMMUTABLE VERSION 47 BODY
\clearpage
\begin{center}
{\Large\bfseries Version 47: The inverse-amplitude time scale}\\[.5em]
{\large A seven-dimensional neutral equation and its signed return tensor}
\end{center}
\makeatletter
\addtocontents{toc}{\protect\begingroup}
\addtocontents{toc}{\protect\makeatletter}
\addtocontents{toc}{\protect\def\protect\l@section{\protect\bfseries\protect\@dottedtocline{1}{0em}{2.5em}}}
\makeatother
\addcontentsline{toc}{part}{Version 47: The inverse-amplitude time scale and the neutral return}

\section{The next question is the motion on the selected constraint leaf}
\label{v47:sec:context}

Version 46 proves original-action continuation through $s_*a^2$ and identifies
 two nonzero neutral forcing channels.  It does not identify the limiting
 motion at $t=a^2\tau$.  This body addresses that missing motion, rather than
 trying to eliminate the already certified Jordan chains.  Throughout, the
 cutoff is $h=1/3$, the leading fields are the exact V39 mixed-mode root, and
 the initial family and locally invariant graph are those selected in V44.
 No new initial rate is chosen.

The inputs retained from V44--V46 are the $C^6$ local invariant graph, the
 exact nonconstant-constraint chart, the graded active equation, the bounded
 nonfast weak semigroup, and the rank-two neutral coupling.  The matrix
 certificates for the algebraic root and signed forms are inherited and
 replayed; their complete source banks are not newly regenerated here.
 The additions are: the actual leading return through the active equations;
 a signed two-dimensional matrix enclosure; continuation for every fixed
 bounded interval in $\tau$; and a closed seven-dimensional equation for the
 limiting weak velocity.  The remaining higher coefficients of that equation
 are defined by derivatives of the original reduced field, not assigned
 values by a model or set to zero.

We use the V46 coordinates $(q,w)\in\mathbb R^{298}\times\mathbb R^{20}$,
 with reconstructed normalized state $u=\Psi(a,q,w)$ and field
 $(q',w')=(f_q,f_w)$.  A prime initially denotes $d/ds$, where $t=a^3s$.
 Write $R=R_{\rm c}$ and $B=B_{\rm c}$ on this nonfast space.  The exact
 spectral projector $\Pi_0$ selects its five-dimensional, semisimple kernel
 $\mathcal Z$.  The restriction $B_\perp$ to $(I-\Pi_0)\mathbb R^{20}$ is
 invertible; its semigroup is bounded.  Nine decaying and six oscillatory
 directions are retained there, not discarded as gauge variables.

\section{The leading return is determined by the original active equation}
\label{v47:sec:return}

Let $A\succ0$ and $D$ be the original leading active block and weak Schur
 form.  Let $K_{AW}$ and $K_{WA}=-K_{AW}^T$ be the active--weak blocks of
 the original first Poisson bracket in the same physical test coordinates.
 All common spatial-pairing factors must be used consistently.  Define
 \[
 R_0=\Pi_0R,\qquad \mathcal V=\operatorname{Ran}R_0\subset\mathcal Z.
 \]
 Version 46 proves $\dim\mathcal V=2$, that $\mathcal V$ is the
 $D$-orthogonal complement of the three constant shifts in $\mathcal Z$,
 and that $D|_{\mathcal V}$ has inertia $(1,1,0)$.

\begin{lemma}[One-step return from a neutral velocity]
\label[lemma]{v47:lem:return}
The limit
\[
 E=\lim_{a\to0}a^{-2}D_w f_q(a,0,0)|_{\mathcal Z}
\]
 exists.  In the split $q=(\kappa,p)$ of V46 it is exactly
\begin{equation}
 \boxed{Ey=\bigl(0,-A^{-1}K_{AW}y\bigr),\qquad y\in\mathcal Z.}
 \label{v47:eq:E}
\end{equation}
Consequently the neutral return $J=R_0E$ is
\begin{equation}
 \boxed{J=-\Pi_0D^{-1}K_{AW}^TA^{-1}K_{AW}|_{\mathcal Z}.}
 \label{v47:eq:J}
\end{equation}
\end{lemma}
\begin{proof}
The exact V46 top equations are
\[
 \kappa'=a^3\Pi_{\ker\mathcal L_{\rm con}}f,
 \qquad
 p'=-a^2A(a,u)^{-1}b_A(a,u)
        +a(\mathsf Y_*-\mathsf Y(a,u))\mathcal G_W(a,u).
\]
The first derivative contributes zero to the displayed limit.  At the
 leading selected record $b_A(0,0)=0$, by the full quadratic drift
 cancellation.  The original identity
 $D_\lambda\mathsf B[\mu]=\mathbb K\mu+D_X(M\mu)[F]$ implies
 $D_\lambda\mathsf B=a\mathbb K_1(Y_*)+O(a^2)$ at that record.
 Thus $D_{u_W}b_A(0,0)=K_{AW}$ on an actual weak variation.

The coordinate eliminations do not introduce another leading term.  In the
 divided constraints, a weak multiplier variation has derivative $MI_W$,
 whose active and weak orders are $a^3,a^4$.  The shear induces an active
 variation $-a\mathsf Y_*u_W$, also contributing only $O(a^3)$ to these
 constraint rows.  The eliminated canonical derivative is therefore
 $O(a^3)$.  At $a=0$ a variation $y\in\mathcal Z$ has no fast graph
 component and no active or canonical component.  The $O(a)$ graph and
 shear corrections affect only higher terms in $D b_A$.

Finally $\mathsf Y-\mathsf Y_*=O(a)$ at the origin and
 $D_w\mathcal G_W(0,0)y=By=0$.  Differentiating the remaining product
 and dividing by $a^2$ therefore gives zero on $\mathcal Z$;
 $\mathcal G_W(a,0)=O(a^3)$ suffices for its other term.  This proves
 \eqref{v47:eq:E}.  On a pure active variation the retained source is
 $R\eta_A=-\Pi_{\rm c}D^{-1}K_{WA}\eta_A$; its off-constraint canonical
 term is absent because that variation is zero.  Composition, followed
 by $\Pi_0$, and $K_{WA}=-K_{AW}^T$ gives \eqref{v47:eq:J}.
\end{proof}

Let $N=(N_c,N_2)$ be the exact five-column kernel synthesis retained from
 V45, with $N_c$ the three constant shifts.  Put $H=N^TDN$, partitioned
 according to $3+2$, and define
\begin{equation}
 N_v=N_2-N_cH_{cc}^{-1}H_{c2},\qquad
 D_v=H_{22}-H_{2c}H_{cc}^{-1}H_{c2},\qquad
 T=K_{AW}N_2.
 \label{v47:eq:two-basis}
\end{equation}
Then $N_v$ synthesizes $\mathcal V$ and $K_{AW}N_c=0$.

\begin{theorem}[The actual return has one positive and one negative eigenvalue]
\label[theorem]{v47:thm:signature}
The return annihilates constant shifts, has image $\mathcal V$, and in its
 two-dimensional basis is
\begin{equation}
 \boxed{J_v=-D_v^{-1}G_v,\qquad G_v=T^TA^{-1}T\succ0.}
 \label{v47:eq:Jv}
\end{equation}
It has two simple real nonzero eigenvalues of opposite signs.  Its three
 zero eigenvalues on $\mathcal Z$ are semisimple.
\end{theorem}
\begin{proof}
For $\Pi_0=NH^{-1}N^TD$, one has $\Pi_0D^{-1}=NH^{-1}N^T$.
 The constant columns of $K_{AW}N$ vanish, and the other two have rank
 two by the actual V46 coupling certificate.  Since $A\succ0$, $G_v$
 is positive definite.  Elimination of the positive constant block of $H$
 gives \eqref{v47:eq:Jv}.  The inertia of $D_v$ is $(1,1,0)$, so
 $\det J_v=\det G_v/\det D_v<0$.  A real two-by-two matrix with negative
 determinant has two distinct real roots of opposite signs.  The constant
 shifts are its exact kernel complement and $J_v$ is invertible, proving
 semisimplicity at zero.  All assertions concern this return tensor, not
 the spectrum of the complete slow evolution derived below.
\end{proof}

\section{An exact enclosure for the two-dimensional return}
\label{v47:sec:certificate}

\begin{proposition}[Source-specific spectral brackets]
\label[proposition]{v47:prop:brackets}
On the retained root cube of radius $2^{-112}$, the two nonzero eigenvalues
 of \eqref{v47:eq:Jv} satisfy
\begin{equation}
 \boxed{-0.543<\lambda_-<-0.542,\qquad
              13.68<\lambda_+<13.69.}
 \label{v47:eq:spectrum}
\end{equation}
\end{proposition}
\begin{proof}
The verifier first replays V46's root, signed form, kernel and source
 certificates.  It then forms \eqref{v47:eq:two-basis}--\eqref{v47:eq:Jv}
 with exact rational midpoints and norm errors.  In the raw active column
 list $\mathcal I_A$, physical grading is retained by
\[
 T=\operatorname{diag}(1\text{ on lapse},\,1/\sqrt3\text{ on shift})
       (\widehat K\mathsf W N_2)_{\mathcal I_A,:}.
\]
The certified bounds on $\sqrt3$, the source cube and the kernel synthesis
 are propagated in this multiplication.  Constant or weak additions to
 active representatives change neither $A$ nor this $T$ on $\ker\mathsf K$.

The approximate matrices, shown only to orient the certificate, are
\[
 D_v\simeq\begin{pmatrix}-45182.811534&2984.845166\\
                          2984.845166&624946.178317\end{pmatrix},\qquad
 G_v\simeq\begin{pmatrix}669919.403569&-749368.859853\\
                         -749368.859853&1151315.805701\end{pmatrix},
\]
\[
 J_v\simeq\begin{pmatrix}14.901376&-16.701697\\
                         1.127922&-1.762494\end{pmatrix}.
\]
The exact stored rational midpoint $J_{v,0}$ has the uniform enclosure
\begin{equation}
 \|J_v-J_{v,0}\|_\infty<1.264\,10^{-6}.
 \label{v47:eq:Jerror}
\end{equation}
For clarity, a product error is propagated before each inverse is used.
 If $X=X_0+\Delta X$, $\|\Delta X\|\le e_X$ and
 $\|X_0^{-1}\|e_X<1$, the inverse error is bounded by
 $\|X_0^{-1}\|^2e_X/(1-\|X_0^{-1}\|e_X)$.
 A dyadic proposed active inverse is accepted instead through its exact
 Neumann residual against the complete active enclosure.

For $p(\lambda)=\lambda^2-(\operatorname{tr}J_v)\lambda+\det J_v$,
 an error $e_J$ in the two-by-two matrix gives
\[
 |p(\lambda)-p_0(\lambda)|
 \le 2e_J|\lambda|+
 e_J\sum_{i,j}|(J_{v,0})_{ij}|+2e_J^2.
\]
Exact rational substitution certifies the following signs:
\begin{center}\small
\begin{tabular}{r r r c}
\toprule
$\lambda$ & Approximate $p_0(\lambda)$ & Error upper bound & Exact sign\\
\midrule
$-0.543$ & $0.00388934$ & $0.0000450$ & $+$\\
$-0.542$ & $-0.01033453$ & $0.0000450$ & $-$\\
$13.68$ & $-0.02288380$ & $0.0000782$ & $-$\\
$13.69$ & $0.11942736$ & $0.0000782$ & $+$\\
\bottomrule
\end{tabular}
\end{center}
Together with the preceding real-spectrum theorem, these strict signs
 prove \eqref{v47:eq:spectrum}.  The complete fractions and inequalities
 are in the executable certificate; no displayed decimal is used as an
 input proof of a sign.  The new computation contracts inherited source
 enclosures; it does not claim new regeneration of their original entries.
\end{proof}

The positive eigenvalue is a signed return coefficient in comparison time
 $\tau=t/a^2$, not a growth rate of the complete slow system.  That system
 includes a first-derivative matrix which has not yet been evaluated.
 A conclusion about actual excitation, damping or instability cannot be
 obtained by taking the square root of $\lambda_+$ alone.

\section{Uniform derivative structure and arbitrary bounded slow intervals}
\label{v47:sec:uniform}

The V46 proof used a small fixed endpoint in $\tau$.  Its exact upper
 equation permits a stronger continuation statement without any further
 spectral or preparation assumption.

\begin{lemma}[Derivative bounds on a larger shrinking tube]
\label[lemma]{v47:lem:jets}
For some fixed $c_0>0$ and all small positive $a$, on
 $\|(q,w)\|\le c_0a$ one has
\begin{equation}
 \|Df_q\|\le Ca^2,\qquad \|D^2f_q\|\le Ca,\qquad
 Df_w=(R,B)+O(a+\|(q,w)\|).
 \label{v47:eq:large-tube}
\end{equation}
There are fixed finite matrices
\begin{align}
 U&=\tfrac12\partial_a^2D_qf_q(0,0,0),&
 E_*&=\tfrac12\partial_a^2D_wf_q(0,0,0),\notag\\
 C&=\partial_aD_qf_w(0,0,0),&
 H_1&=\partial_aD_wf_w(0,0,0),
 \label{v47:eq:jets}
\end{align}
 such that, on each $O(a^2)$ state tube,
\begin{align}
 a^{-2}D_qf_q&=U+O(a),&a^{-2}D_wf_q&=E_*+O(a),\notag\\
 a^{-1}(D_qf_w-R)&=C+O(a),&
 a^{-1}(D_wf_w-B)&=H_1+O(a).
 \label{v47:eq:jet-limits}
\end{align}
The restriction of $E_*$ to $\mathcal Z$ is the already evaluated $E$.
\end{lemma}
\begin{proof}
Use the exact top field of V46.  The divided active term has its explicit
 $a^2$ factor and bounded derivatives.  The other term is
 $a(\mathsf Y_*-\mathsf Y)\mathcal G_W$, with
 $\mathsf Y-\mathsf Y_*=O(a+\|u\|)$,
 $\mathcal G_W=O(a^3+\|u\|)$, and bounded fixed derivatives.
 Its first derivative is $O(a^2)$ on $\|u\|=O(a)$; its second is
 $O(a)$.  The uniformly regular $C^6$ reconstruction preserves these
 orders.  The canonical prefactor is $a^3$.  Taylor expansion of the
 lower field gives the last bound in \eqref{v47:eq:large-tube}.
 At the origin, the $O(a^2)$ upper derivative has vanishing coefficients
 of orders zero and one, so its second amplitude derivative gives
 \eqref{v47:eq:jets}.  Moving the state by $O(a^2)$ changes that derivative
 by $O(a^3)$, using $D^2f_q=O(a)$; it changes the lower derivative by
 $O(a^2)$.  This proves the uniform limits.  Only finite derivatives of
 the retained $C^6$ chart are used.
\end{proof}

For the same selected initial family, let
\begin{equation}
 r_* =\lim_{a\downarrow0}\frac{a^{-3}q'(0)-b_q}{a}.
 \label{v47:eq:rinitial}
\end{equation}
This limit exists with remainder $O(a)$: V43--V44 retain the same analytic
 initial coefficients through order $a^2$, and their selected $w(a)$ is
 bounded with $w(a)=w^\sharp+O(a)$.  The changed fields start at $a^3$;
 the divided canonical velocity therefore has its usual first amplitude
 coefficient.  The exact selected multiplier rate gives the same assertion
 for the sheared active coordinates.  This defines $r_*$; it does not
 tabulate its value or presume that it vanishes.

\begin{theorem}[The same original histories on every fixed slow interval]
\label[theorem]{v47:thm:allT}
For every fixed $T<\infty$ there is $a_T>0$ such that the unchanged V44
 selected family has its unique original normalized histories on
\begin{equation}
                  \boxed{0\le t\le Ta^2\quad(0<a<a_T).}
 \label{v47:eq:allT}
\end{equation}
All original constraints, the selected invariant graph, positive lapse and
 metric, and the signed-regular chart persist.  Uniformly on that interval,
\[
 \|u\|\le C_Ta^2,\qquad
 \left\|a^{-1}(a^{-3}q'-b_q)\right\|+
 \|a^{-3}w'\|\le C_T.
\]
No constant here is asserted uniform as $T\to\infty$ or as $h\to0$.
\end{theorem}
\begin{proof}
Put $\tau=as=t/a^2$, $v_q=q'/a^3$, $v_w=w'/a^3$, and
 $r_a=(v_q-b_q)/a$.  Differentiating the actual autonomous field gives
\begin{align*}
 \dot r_a&=U_a(b_q+ar_a)+E_av_w,\\
 \dot v_w&=a^{-1}Bv_w+Rr_a+C_a(b_q+ar_a)+H_av_w,
\end{align*}
 where the four coefficients are the left sides of
 \eqref{v47:eq:jet-limits}; a dot denotes $d/d\tau$.  On the larger
 tube their norms are bounded independently of $a$, by
 \eqref{v47:eq:large-tube}.  Variation of constants for the second
 equation uses $\|e^{B\tau/a}\|\le M_0$.  Together with the first equation
 and the bounded initial values $r_a(0)$, $v_w(0)=O(a)$, Gronwall gives
 $\|r_a\|+\|v_w\|\le C_T$ while the trajectory stays in the tube.
 In original chart coordinates,
\[
 \frac{d(q,w)}{d\tau}=a^2(b_q+ar_a,v_w).
\]
 Integration yields $\|(q,w)\|\le C_T'a^2$.  For sufficiently small
 $a$, this is strictly inside the larger tube $c_0a$.  Continuation
 therefore proves the entire interval.  The exact graph and constraints
 are invariant, and $u=O_T(a^2)$ stays in the original regular chart.
 No repeated reinitialization or increase of the allowed graph error occurs.
\end{proof}

\section{A closed seven-dimensional equation for the neutral velocity}
\label{v47:sec:slow}

Define the finite constants
\begin{equation}
 \gamma=Ub_q,\qquad \beta=\Pi_0Cb_q,\qquad
 H_{\rm n}=\Pi_0H_1|_{\mathcal Z},\qquad
 d=R_0\gamma\in\mathcal V,\qquad
 \varpi=R_0r_*+\beta\in\mathcal Z.
 \label{v47:eq:slow-coeff}
\end{equation}
Unlike $J$, $H_{\rm n}$, $d$ and $\varpi$ have not been numerically
 evaluated in this body.  They retain the original higher initial
 coefficients and the derivatives through the constrained invariant graph.

\begin{theorem}[Neutral slow-time limit with all nonzero weak modes retained]
\label[theorem]{v47:thm:limit}
Let $(y,j)\in\mathcal Z\times\mathcal V$ solve
\begin{equation}
 \boxed{\dot y=H_{\rm n}y+j+\varpi,\qquad
        \dot j=Jy+d,\qquad y(0)=0,\quad j(0)=0.}
 \label{v47:eq:seven}
\end{equation}
This affine system has dimension seven and exists for every finite $\tau$.
 For every fixed $T$, the actual original histories satisfy
\begin{equation}
 \sup_{0\le\tau\le T}\|a^{-3}w'(\tau/a)-y(\tau)\|\le C_Ta.
 \label{v47:eq:weak-limit}
\end{equation}
More precisely, $r_a$ converges with error $O_T(a)$ to $r$ solving
 $\dot r=\gamma+Ey$, $r(0)=r_*$, and
\[
 j=R_0(r-r_*),\qquad
 \|(I-\Pi_0)a^{-3}w'(\tau/a)\|\le C_Ta.
\]
No decay assumption is made for the six oscillatory directions.
\end{theorem}
\begin{proof}
The preceding theorem places each trajectory in an $O_T(a^2)$ tube, so
 all four coefficients in its velocity equations converge to
 $(U,E_*,C,H_1)$ with error $O_T(a)$.  They also have bounded first
 $\tau$ derivatives: for example
 $\partial_\tau(a^{-2}Df_q)=a^{-2}D^2f_q\,(q,w)^{\boldsymbol\cdot}$
 is $O_T(a)$, and the lower coefficients have the same bound.

Write $v_w=y_a+z_a$, with $y_a=\Pi_0v_w$.  The equations for
 $r_a$ and $y_a$ have uniformly bounded right sides.  On the remaining
 weak space,
\[
 \dot z_a=a^{-1}B_\perp z_a+g_a(\tau)+L_a(\tau)z_a,
\]
 where $g_a$ and $\dot g_a$ are bounded on $[0,T]$, and $L_a$ is
 bounded.  Here $g_a$ contains only $r_a,y_a$ and the bounded
 coefficient functions, not a differentiated $z_a$.
 For any $C^1$ forcing $g$ in this space,
\begin{align*}
 \int_0^\tau e^{B_\perp(\tau-v)/a}g(v)\,dv
 ={}&-aB_\perp^{-1}g(\tau)
 +ae^{B_\perp\tau/a}B_\perp^{-1}g(0)\\
 &+a\int_0^\tau e^{B_\perp(\tau-v)/a}B_\perp^{-1}\dot g(v)\,dv.
\end{align*}
 The bounded semigroup and inverse make this $O_T(a)$.  Since
 $z_a(0)=O(a)$, a further integral Gronwall estimate gives
 $z_a=O_T(a)$.  This argument retains both decaying and oscillatory
 components; it is not a dissipative elimination of the latter.

Projecting the actual equations onto $\mathcal Z$ now gives, with
 uniform $O_T(a)$ errors,
\[
 \dot r_a=\gamma+Ey_a+O_T(a),\qquad
 \dot y_a=R_0r_a+\beta+H_{\rm n}y_a+O_T(a).
\]
 Their initial errors are $O(a)$.  Finite-dimensional Gronwall proves
 convergence to the displayed limiting equations.  Applying $R_0$ to
 the first equation gives \eqref{v47:eq:seven}, because $R_0E=J$.
 The range of $R_0$ is exactly $\mathcal V$, so this is seven rather
 than $303$ unknowns.  Constant coefficients ensure global existence
 of this comparison ODE, not a cutoff-uniform bound on its growth.
\end{proof}

This integration-by-parts removal of bounded rapid oscillations is a
 finite-dimensional singular-perturbation argument.  Averaging with
 oscillatory fast dynamics is established methodology
 \cite{v47:Tao}; no external result supplies the original return tensor
 or its eigenvalue signs.  The full slow generator, in coordinates
 $\mathcal Z=\operatorname{Ran}N_c\oplus\mathcal V$, is
\begin{equation}
 \begin{pmatrix}H_{\rm n}&\iota_{\mathcal V}\\
                 J_v\pi_{\mathcal V}&0_2\end{pmatrix},
 \label{v47:eq:generator}
\end{equation}
 not $J_v$.  Its spectrum has not been determined here.  The initial
 forcing is $(\varpi,d)$; a nonzero spectrum alone would not prove
 excitation by that forcing.

\section{Original fields, curvature observations and the remaining source bank}
\label{v47:sec:observations}

\begin{corollary}[Leading original motion on $t=a^2\tau$]
\label[corollary]{v47:cor:physical}
For every fixed $T$, uniformly on $0\le\tau\le T$,
\begin{align}
 a^{-1}\lambda_t(a^2\tau)&=v_*+I_Wy(\tau)+O_T(a),\notag\\
 a^{-3}\{X(a^2\tau)-X(0)\}&=\tau V+O_T(a),\notag\\
 a^{-3}\{\lambda(a^2\tau)-\lambda(0)\}
 &=\tau v_*+I_W\int_0^\tau y(v)\,dv+O_T(a).
 \label{v47:eq:physical}
\end{align}
For each retained smooth finite curvature observation $\mathscr O$, define
 $\mathcal O_Wy$ to be its flat derivative in the physical
 $\lambda_t$ slot, evaluated on $I_Wy$.  Then
\begin{equation}
 \boxed{\sup_{0\le\tau\le T}
 \left\|\frac{\mathscr O(a^2\tau)-\mathscr O(0)}a
               -\mathcal O_Wy(\tau)\right\|\le C_Ta.}
 \label{v47:eq:obs-limit}
\end{equation}
Each curvature is still $O_T(a)$.  An observation with no $\lambda_t$
 slot has zero leading change; for the spatial link curvature its change
 is $O_T(a^3)$.
\end{corollary}
\begin{proof}
The regular reconstruction has its base derivative sending $(b_q,0)$
 to the previously retained $b=(V,v_*,0)$, and sending $(0,y)$ to a
 pure weak multiplier variation when $y\in\mathcal Z$.  The latter
 has no fast graph component, since $By=0$.  Its derivative along the
 trajectory differs from this base derivative by $O_T(a)$.
 The preceding velocity limit therefore gives the first equation and,
 by integration over $\tau$, the two displacement statements.
 The original state arguments change by $O_T(a^3)$, and their smooth
 canonical field values have the same change.  The physical multiplier
 rate changes by $aI_Wy+O_T(a^2)$.  Taylor expansion of each retained
 finite observation at the flat state proves \eqref{v47:eq:obs-limit}.
 There is no $\lambda_{tt}$ slot in these observations.  This proof
 does not establish convergence of their time derivatives; the small
 retained oscillatory velocities can have non-small differentiated
 errors on the $\tau$ scale.
\end{proof}

\begin{proposition}[A finite test for disappearance of the entire slow weak response]
\label[proposition]{v47:prop:forcing}
The solution of \eqref{v47:eq:seven} is identically zero in its $y$
 component on an interval of positive length if and only if
\begin{equation}
                         \boxed{\varpi=0,\qquad d=0.}
 \label{v47:eq:zero-forcing}
\end{equation}
In particular, on any fixed interval $[0,T]$ with $T>0$, uniform
 $o(a)$ weak multiplier rate for this selected family is equivalent to
 \eqref{v47:eq:zero-forcing}.  This is not a claim that those native
 equalities hold.  Vanishing of a particular curvature observation is
 weaker and requires its own observation of $y$.
\end{proposition}
\begin{proof}
If $y=0$, the first equation gives $j=-\varpi$.  At $\tau=0$ this
 forces $\varpi=0$, hence $j=0$ and the second gives $d=0$.
 Conversely these values give the zero solution by uniqueness.
 The rate equivalence follows from \eqref{v47:eq:physical} and the
 injective weak synthesis, with the original leading transverse rate
 having zero weak component.  Analyticity of the constant-coefficient
 comparison ODE excludes a nonzero solution vanishing on an initial
 interval.  Its first derivatives are
 $\dot y(0)=\varpi$, $\ddot y(0)=H_{\rm n}\varpi+d$.
\end{proof}

Once the neutral basis is fixed, the remaining numerical source bank is
 $25$ entries of $H_{\rm n}$, five of $\varpi$, and two of $d$.
 Thus $32$ scalar entries, in addition to the now enclosed return,
 determine the complete seven-dimensional affine response.  For the narrower
question of vanishing leading weak velocity, only the seven forcing entries
in $(\varpi,d)$ need be tested: $H_{\rm n}$ does not enter
\eqref{v47:eq:zero-forcing}.  These entries
 require derivatives through the actual constrained selected initializer;
 they are not recoverable by discarding its uncomputed higher coefficients.
 No claim of nonzero excitation or a stable slow generator follows from
 the already evaluated return alone.

\section{Status, reproducibility and the fixed-cutoff boundary}
\label{v47:sec:status}

\paragraph{Proved in this body.}
The neutral-to-active-to-neutral return is the signed inverse of a positive
 rank-two Gram.  Its two nonzero eigenvalues have the explicit brackets
 \eqref{v47:eq:spectrum}.  The same selected original histories continue
 through $Ta^2$ for every fixed finite $T$, with an amplitude threshold
 allowed to depend on $T$.  Their normalized weak velocity converges to
 the seven-dimensional affine equation \eqref{v47:eq:seven}; the other
 fifteen nonfast weak velocity modes are $O_T(a)$, including the
 oscillatory ones.  The corresponding leading original field and
 curvature changes are identified without evolving a replacement action.

\paragraph{What is inherited or unevaluated.}
The V44 invariant graph, exact initialization, signed regularity, V45
 spectral decomposition and V46 constrained derivative identities are
 retained analytical inputs.  The new verifier replays the V46 executable
 certificates, then evaluates the return and its characteristic polynomial
 enclosures in exact arithmetic.  It does not regenerate the old stationary,
 weak-Poisson or root coefficient banks.  The $32$ remaining slow-source
 entries, the graph itself and the higher initial coefficients are not
 numerically tabulated.  Small exactly specified block tests validate the
 algebra and the singular-oscillation identity, not the original nonlinear
 spacetime flow.  The slow-time convergence is the written proof above.

\paragraph{Next source-specific task.}
Compute $H_{\rm n},\varpi,d$ in the declared basis and determine both the
 full generator and the actual forcing.  This replaces another generic
 lifespan bootstrap by a small, definite native response problem.  In
 particular, the positive return eigenvalue is not sufficient to assert
 instability, and the existence of an invariant graph is not attraction.
 The normal repulsion of V44 remains.  The present bounds do not permit
 $T$ to depend on $a$ without a new uniform estimate; they do not give
 a fixed positive physical lifespan.  Spatial refinement, primitive
 selection and the matter-coupled Einstein--Standard-Model continuum
 remain unresolved.  The independent V32 stress-compensation branch is
 unchanged.

\begingroup
\renewcommand{\refname}{Additional reference for Version 47}
\begin{thebibliography}{9}
\small\raggedright
\bibitem{v47:Tao}
M.~Tao, \emph{Simply improved averaging for coupled oscillators and weakly
 nonlinear waves}, arXiv:1806.07947 (2018).
\url{https://arxiv.org/abs/1806.07947}.
Primary-source methodological context for oscillatory averaging; the finite
 source, signed return and estimates above are derived here.
\end{thebibliography}
\endgroup
\addtocontents{toc}{\protect\endgroup}
% END IMMUTABLE VERSION 47 BODY
% APPEND VERSION 48 HERE, BEFORE THE FINAL END-DOCUMENT COMMAND.


% BEGIN IMMUTABLE VERSION 48 BODY
\clearpage
\begin{center}
{\Large\bfseries Version 48: The signed slow pencil and its harmonic return}\\[.5em]
{\large An exact covector reduction, evaluated harmonic coupling, and a growing slow pair}
\end{center}
\makeatletter
\addtocontents{toc}{\protect\begingroup}
\addtocontents{toc}{\protect\makeatletter}
\addtocontents{toc}{\protect\def\protect\l@section{\protect\bfseries\protect\@dottedtocline{1}{0em}{2.5em}}}
\makeatother
\addcontentsline{toc}{part}{Version 48: The signed slow pencil and its harmonic return}

\section{Companion audit and the actual remaining slow coefficients}
\label{v48:sec:audit}

Version 47 leaves a seven-dimensional forced equation on the original time
scale $t=a^2\tau$, with a known rank-two return and an apparently arbitrary
$5$-by-$5$ first-derivative matrix. The new companions suggest a more useful
upstream calculation than evaluating all of those entries: differentiate the
original covector consistency equation before taking its neutral projection.
This retains a signed skew structure. Three constant shifts then supply an
additional positive stiffness rather than disappearing from the slow equation.
The calculation below concerns the same Native V39 algebraic mixed-mode root,
$h=1/3$, and the same exactly constrained invariant initial family as V44--V47.
No new initial rate, action term, or time convention is selected.

\paragraph{Exact-action Einstein bridge.}
The supplied cumulative V36 \cite{v48:EB36}, especially
\nolinkurl{v34:derivative} and \nolinkurl{v36:Poisson-orders}, retains the exact
covector identity and derives the harmonic mixed Poisson coefficient from the
complete cubic-mean map. These identities concern the original scalar
envelopes and can be applied here with their hypotheses checked. Its spectral
counts cannot be copied: its V29 root has a different mixed-field population
and signed inertia $(96,11)$, whereas our root has inertia $(97,10)$ and nine,
not eleven, expanding fast directions. In particular its smaller canonical
spectral scale is not already a theorem about our seven-dimensional slow
response. We transfer the coefficient identities, not that root's spectrum.

\paragraph{Standard-Model source selection.}
The supplied V43 \cite{v48:SM43}, \nolinkurl{sel43:thm:physical} and its final
ledger, constructs an exactly circle-covariant count instrument by a specified
causal stopping policy with maximal accepted-event horizon seven. It retains
the actual independent renewal clock, proves minimality of that horizon for
the stated oriented setting, and distinguishes the count record from the
noncovariant fine word. Its mean accepted count is $197/40$. These are genuine
source-policy results, but the policy is not a damping mechanism for the
present deterministic Einstein flow and is not identified as its historical
matter generator. No stopping, forgetting, or feedback operation is inserted
in the histories analyzed below.

\paragraph{Perturbative continuum.}
The supplied V51 \cite{v48:PE51}, \nolinkurl{thm:v51-BV} and
\nolinkurl{eq:v51-next}, completes a weight-four polynomial Noether/BV section
in its specified four-equal-free-real-scalar sector. Its metric-only
construction retains the scalar descendants, including nonclosed one-form
terms. The next displayed identity still involves the actual new metric-sextic
and scalar-quintic sources. This is not a modification of our raw action, a
classical lifespan estimate, or a proof that the higher initializer forcing
below vanishes.

\paragraph{Yang--Mills.}
The supplied V79 \cite{v48:YM79}, \nolinkurl{r79:local} and
\nolinkurl{r79:unpaid}, retains the complete $3$-by-$3$ orientation contact of
the first covariance coefficient, its finite-range midpoint phases, and a
uniform second-moment remainder. The full compact remainder with a
volume-independent coupling threshold remains unpaid. The useful accounting
lesson here is to retain the entire harmonic cross block; neither the
Yang--Mills probability law nor a mass-gap assertion is imported.

The review covers these terminal bodies and the coefficient dependencies used
below, not an independent validation of every theorem in the four cumulative
companions. The new results are a signed-skew identity for our actual reduced
slow coefficient, a six-entry original harmonic-source enclosure, a forced
two-coordinate reduction, and its complete qualitative spectrum. Quadratic
pencils and gyroscopic reductions are established methods \cite{v48:TM};
the original-source identities, signs, and numerical inequalities here are
proved at the specified root.

\section{The neutral first derivative retains a signed skew form}
\label{v48:sec:covector}

Use the fixed normalized $107$-dimensional multiplier complement. To avoid
confusing the drift with a Poisson matrix, write the original objects as
\[
 P=H_\lambda,\quad M=H_{\lambda\lambda},\quad C_X=H_{\lambda X},
 \quad F=\mathbb JH_X,\quad K=C_X\mathbb JC_X^T,
 \quad \mathsf B=C_XF,\quad r=-M^{-1}\mathsf B.
\]
The inverse is the original signed inverse on the regular normalized chart.
Along the actual vector field $(F,r)$, a dot denotes its total derivative.

\begin{lemma}[The moving signed metric must be retained]
\label[lemma]{v48:lem:covector}
On that chart the exact identity is
\begin{equation}
 \boxed{M D_\lambda r=-K-\dot M.}
 \label{v48:eq:covector}
\end{equation}
\end{lemma}
\begin{proof}
Differentiate $Mr+\mathsf B=0$ with respect to the multiplier in a fixed
direction. The derivative of $\mathsf B=C_X\mathbb JH_X$ is
$K+D_XM[F]$. Symmetry of the third multiplier derivative gives
$D_\lambda M[\,\cdot\,]r=D_\lambda M[r]$. Their sum is precisely $\dot M$.
This proves the companion identity here without dropping its moving-metric
term. It is not, at finite amplitude, the formula $D_\lambda r=-M^{-1}K$.
\end{proof}

Let all leading matrices now be those of V41--V47 at the Native V39 root:
\begin{equation}
 \Gamma_*=
 \begin{pmatrix}A&L\\L^T&C_W\end{pmatrix},\qquad
 D=C_W-L^TA^{-1}L,\qquad Y_*=A^{-1}L,
 \qquad V_* = \begin{pmatrix}-Y_*\\I_{29}\end{pmatrix}.
 \label{v48:eq:leading}
\end{equation}
The exact grading is $M=D_a\Gamma(a,u)D_a$, where
$D_a=\operatorname{diag}(aI_{78},a^2I_{29})$. Put
$E_a=\operatorname{diag}(aI_{78},I_{29})$. At the base initial record expand
\begin{equation}
 K(a)=a\mathcal K(a),\qquad
 P_1=\left.\partial_a(E_a\mathcal K(a)E_a)\right|_{a=0}.
 \label{v48:eq:balanced}
\end{equation}
Every matrix in the balanced expansion is skew. Let $N_c$ synthesize the three
constant shifts and let $N_2$ be the two remaining kernel columns of V45. Set
\begin{equation}
 \begin{split}
 D_c&=N_c^TDN_c\succ0,\\
 N_v&=N_2-N_cD_c^{-1}N_c^TDN_2,\qquad Z=(N_c,N_v),\\
 D_v&=N_v^TDN_v,\qquad D_Z=Z^TDZ=\operatorname{diag}(D_c,D_v).
 \end{split}
 \label{v48:eq:neutral-basis}
\end{equation}
The retained kernel and signature results give
$\operatorname{Inertia}(D_v)=(1,1,0)$. The columns of $N_v$ span the actual
range $\mathcal V$ of the rank-two neutral forcing. Coordinates in this fixed
basis are used for $H_{\rm n}$ below.

\begin{theorem}[The actual neutral derivative is gyroscopic]
\label[theorem]{v48:thm:gyro}
For the reduced field and initializer of V47,
\begin{equation}
 \boxed{D_ZH_{\rm n}=-\Omega,\qquad
 \Omega=Z^TV_*^TP_1V_*Z,
 \qquad \Omega^T=-\Omega.}
 \label{v48:eq:gyro}
\end{equation}
Moreover the constant--constant block of $\Omega$ is zero. Hence, in the
retained real basis, there are $Q\in\mathbb R^{3\times2}$ and one real number
$\chi$ such that
\begin{equation}
 \boxed{
 \Omega=\begin{pmatrix}0_3&Q\\-Q^T&S\end{pmatrix},
 \qquad S=\begin{pmatrix}0&\chi\\-\chi&0\end{pmatrix}.}
 \label{v48:eq:blocks}
\end{equation}
These are identities of the actual amplitude derivative after the constraint,
shear, and invariant-graph reductions, not assumptions on a comparison system.
\end{theorem}
\begin{proof}
In the V44 normalized chart the multiplier field is
$\mathcal G_\lambda=a^2r$, and a multiplier derivative in $u$ supplies one
additional factor $a$. At fixed V46 sheared active coordinate, a weak variation
is $E_aV_*z$. Applying \eqref{v48:eq:covector} and the exact grading gives for
its weak output
\begin{equation}
 -R_W\Gamma^{-1}(E_a\mathcal K E_a)V_*z
 -a^3R_W\Gamma^{-1}\dot\Gamma V_*z.
 \label{v48:eq:balanced-derivative}
\end{equation}
At the base initializer $F=O(a)$ and $r=av_*+O(a^2)$. The normalized fields
therefore have bounded physical velocities, and $\dot\Gamma=O(1)$. The second
term of \eqref{v48:eq:balanced-derivative} has no coefficient of order $a$.
This order statement is at these prepared records; no global omission of
$\dot M$ is being made.

At $a=0$ the balanced Poisson matrix has only the weak block
$K_{1,WW}$. If $z\in\ker B=\ker K_{1,WW}$, it annihilates $V_*z$. The
first derivative of $\Gamma^{-1}$ therefore multiplies zero. Also
$R_W\Gamma_*^{-1}=D^{-1}V_*^T$. Projecting the remaining order-$a$ term by
$\Pi_0=Z D_Z^{-1}Z^TD$ gives exactly \eqref{v48:eq:gyro}.

We check the reductions in this differentiation. The nonconstant constraint
chart solves $P/a=0$ with a fixed canonical inverse. Its derivative in a
sheared weak direction is $ME_aV_*$, which is $O(a^3)$ or smaller by the
original block orders; thus the eliminated normalized canonical variation is
$O(a^3)$ and contributes nothing at order $a$. The invariant graph has first
slope $O(a)$ and fast tangent entirely in weak coordinates. Its possible
order-$a$ contribution to the weak field is $B$ applied to a fast vector;
$\Pi_0B=0$ removes it. The graph base point differs by higher order and does
not change this coefficient. The leading shear is precisely the fixed $Y_*$
in \eqref{v48:eq:leading}. Thus the calculation gives the $H_{\rm n}$ defined
from $f_w$ in V47, rather than an unreduced partial derivative.

Finally $K_1$ annihilates constant multipliers. The exact quadratic translation
generators $P_{c,2}=\langle p,\delta_cU\rangle_h$ commute, so the actual mutual
constant-shift bracket is $O(a^3)$ on exactly prepared records. Consequently
the constant--constant part of the order-$a$ balanced matrix, including the
shear terms, is zero. Skew symmetry gives \eqref{v48:eq:blocks}.
\end{proof}

\section{The harmonic cross block is fixed by the complete cubic means}
\label{v48:sec:harmonic}

Let $\mathscr R[c,\nu]$ denote the following complete mixed coefficient in
the original physical pairing:
\begin{equation}
 \mathscr R[c,\nu]
 =DP_{c,3}(Y_*)[G\nu]-DP_{\nu,2}(Y_*)[\delta_cY_*]
       -\langle\delta_c\nu,P_2(Y_*)\rangle_h.
 \label{v48:eq:Rsource}
\end{equation}
Write $\mathscr R_W$ for its three constant-shift rows restricted to the
original weak multiplier bank. The last term in \eqref{v48:eq:Rsource} is
part of the preparation; it cannot be omitted.

\begin{proposition}[A higher-preparation-free cross coefficient]
\label[proposition]{v48:prop:cross}
The block in \eqref{v48:eq:blocks} is
\begin{equation}
 \boxed{
 Q=\mathscr R_WN_v-(Y_*N_c)^TK_{AW}N_v
   =\mathscr R_WN_2-(Y_*N_c)^TT,
 \qquad T=K_{AW}N_2.}
 \label{v48:eq:Q}
\end{equation}
In particular $Q$ depends on the retained leading field and existing cubic
source bank, not on the unknown higher initializer or the new scalar $\chi$.
\end{proposition}
\begin{proof}
For a constant translation, the first constraint coefficient is zero and the
flat Hamiltonian vector of $P_{c,2}$ is $\delta_cY$. Expanding its Poisson
bracket with another constraint through order $a^2$ gives
\[
 DP_{c,2}(Y_2)[G\nu]+DP_{c,3}(Y_*)[G\nu]
                    -DP_{\nu,2}(Y_*)[\delta_cY_*].
\]
Exact preparation gives $\mathcal L_{\rm con}Y_2=-P_2(Y_*)$, turning the first
term into the last term of \eqref{v48:eq:Rsource}. This is the universal
coefficient identity in the supplied bridge; it uses the original phase
translation and complete envelopes. It shows that no unlisted $Y_2$ is needed.
For two constants, the quadratic generators commute identically on all fields,
so their order-$a^2$ coefficient is zero: $\mathscr R_WN_c=0$.

The balanced coefficient in \eqref{v48:eq:balanced} has blocks
\[
 P_1=\begin{pmatrix}0&K_{AW}\\-K_{AW}^T&K_{2,WW}\end{pmatrix}.
\]
Insert $V_*$ and the two parts of $Z$ into \eqref{v48:eq:gyro}.
Since $K_{AW}N_c=0$, the constant-to-$N_v$ cross is precisely
$N_c^TK_{2,WW}N_v-(Y_*N_c)^TK_{AW}N_v$.
Its first term is $\mathscr R_WN_v$. Both maps annihilate $N_c$, proving the
second equality of \eqref{v48:eq:Q}.
\end{proof}

Retain $G_v=T^TA^{-1}T\succ0$ from V47 and define
\begin{equation}
 \boxed{\overline G=G_v+Q^TD_c^{-1}Q\succ0,
 \qquad \overline J=-D_v^{-1}\overline G.}
 \label{v48:eq:effective}
\end{equation}
The added matrix is positive semidefinite because the actual constant-shift
metric $D_c$ is positive. It is not a new potential or an action counterterm.

\begin{theorem}[Exact-root harmonic and effective-return enclosures]
\label[theorem]{v48:thm:certificate}
On the retained cube $\|q-q_0\|_\infty\le2^{-112}$ containing the exact V39
root, the six entries of $Q$ obey the strict componentwise bounds
\begin{equation}
 \begin{pmatrix}3354&-1572\\10577&-8897\\24934&-27721\end{pmatrix}
 <Q<
 \begin{pmatrix}3355&-1571\\10578&-8896\\24935&-27720\end{pmatrix}.
 \label{v48:eq:Qboxes}
\end{equation}
The effective return has two simple real eigenvalues satisfying
\begin{equation}
 \boxed{-0.544<\overline\lambda_-<-0.543,
 \qquad13.96<\overline\lambda_+<13.97.}
 \label{v48:eq:Jbarspectrum}
\end{equation}
All coordinates are the original declared physical grading with a common
spatial-sum pairing, as in V47. The bounds are not cutoff-uniform norms.
\end{theorem}
\begin{proof}
The embedded parent certificate retains the exact root, $A,L,C_W$, the
kernel synthesis, and their entrywise errors. The existing quadratic bank
for the complete cubic-mean control gives $\widehat{\mathscr R}$ in the
rationalized multiplier coordinates. Since the weak bank is a pure shift,
\eqref{v48:eq:Q} is evaluated as
\[
 Q=\frac1{\sqrt3}\widehat{\mathscr R}\,\mathsf WN_2
                     -(Y_*N_c)^TT.
\]
The rational interval for $\sqrt3$, the physical active--weak block conversion,
and every original preparation term are retained. The 21 coefficient vectors
also give $\widehat{\mathscr R}\mathsf WN_c=0$ in exact rational arithmetic.

The machine-readable certificate stores exact rational centres and uniform
operator error radii. Its new enclosures satisfy
\begin{equation}
 \|Q-Q_0\|_\infty<\frac1{20000},\qquad
 \|\overline G-\overline G_0\|_\infty<0.031,\qquad
 \|\overline J-\overline J_0\|_\infty<0.000001282.
 \label{v48:eq:errors}
\end{equation}
For interpretation only, the first centre is approximately
\[
 Q_0\simeq
 \begin{pmatrix}
 3354.6861614&-1571.2704676\\
 10577.7113119&-8896.1288098\\
 24934.1681764&-27720.4106472
 \end{pmatrix}.
\]
It is the full stored rational centre, not these displayed decimals, that is
used in \eqref{v48:eq:errors}.

For reproducibility, product errors are propagated by
$\|AB-\widetilde A\widetilde B\|\le
\|\widetilde A\|e_B+e_A\|\widetilde B\|+e_Ae_B$.
If $R$ is an inverse proposal, $q=\|I-R\widetilde A\|+\|R\|e_A<1$
gives $\|A^{-1}\|\le\|R\|/(1-q)$ and
$\|A^{-1}-R\|\le q\|R\|/(1-q)$. These inequalities are applied to the
actual harmonic and two-dimensional metric solves, never to a substituted
positive weak Hessian. They give \eqref{v48:eq:errors} with all parent radii
included. The positivity of $G_v$, $D_c$ and the inertia of $D_v$ are inherited
certified inputs, so $\det\overline J<0$.

For $p(x)=\det(xI-\overline J)$, exact rational trace/determinant enclosures
give respectively the signs $+,-,-,+$ at
$-0.544,-0.543,13.96,13.97$. There is therefore one root in each stated
interval; the quadratic has no other roots. This proves
\eqref{v48:eq:Jbarspectrum} at the exact root and throughout the retained
cube. No floating eigenvalue determines a sign or a rank.
\end{proof}

\section{An exact two-coordinate equation retaining the harmonic history}
\label{v48:sec:forced}

Write $y=Z(h,z)$ in the five-dimensional neutral space, with
$h\in\mathbb R^3$ and $z\in\mathbb R^2$. Write $j\in\mathbb R^2$ for the
coordinates in $\mathcal V$ of V47's integrated return. Split the retained
forcing as $\varpi=(\varpi_c,\varpi_v)$ and $d=d_v$. Then V47's seven
first-order equations are exactly
\begin{equation}
 \begin{split}
 \dot h&=-D_c^{-1}Qz+\varpi_c,\\
 \dot z&=D_v^{-1}Q^Th-D_v^{-1}Sz+j+\varpi_v,\\
 \dot j&=-D_v^{-1}G_vz+d_v,
 \qquad h(0)=z(0)=j(0)=0.
 \end{split}
 \label{v48:eq:seven}
\end{equation}
Dots in this and subsequent sections mean $d/d\tau$, with the original
physical relation $t=a^2\tau$.

\begin{theorem}[Forced gyroscopic compression of the actual slow equation]
\label[theorem]{v48:thm:forced}
Put $F_v=D_vd_v+Q^T\varpi_c$. Equation \eqref{v48:eq:seven} is equivalent to
\begin{equation}
 \boxed{D_v\ddot z+S\dot z+\overline Gz=F_v,
 \qquad z(0)=0,\quad\dot z(0)=\varpi_v,}
 \label{v48:eq:two}
\end{equation}
with the original three harmonic components recovered by
\begin{equation}
 \begin{split}
 h(\tau)&=\tau\varpi_c-D_c^{-1}Q\int_0^\tau z(v)\,\dd v,\\
 j(\tau)&=\dot z+D_v^{-1}Sz-D_v^{-1}Q^Th-\varpi_v.
 \end{split}
 \label{v48:eq:reconstruction}
\end{equation}
The constant-force signed energy
\begin{equation}
 \frac12\dot z^TD_v\dot z+\frac12z^T\overline Gz-F_v^Tz
 \label{v48:eq:energy}
\end{equation}
is conserved. It is not positive definite.
\end{theorem}
\begin{proof}
Differentiate the second row of \eqref{v48:eq:seven} and substitute the
first and third rows. Moving the harmonic return to the left yields exactly
$Q^TD_c^{-1}Qz$ in addition to $G_vz$, and gives \eqref{v48:eq:two}.
Conversely define $h,j$ by \eqref{v48:eq:reconstruction} for a solution of
\eqref{v48:eq:two}. Differentiation recovers the first and third rows, while
the second is its definition. The initial derivative in \eqref{v48:eq:two}
gives $j(0)=0$. Thus no harmonic datum is lost. Finally multiply the equation
by $\dot z^T$ and use $\dot z^TS\dot z=0$ to prove conservation.
\end{proof}

Eliminating $h$ from the acceleration equation does not set the constant shifts
to zero. Their full integral and forcing are in \eqref{v48:eq:reconstruction}.
Likewise $\overline G$ is a return coefficient of this limiting system, not a
change to the finite action. The signed kinetic form $D_v$ does not by itself
identify negative-energy physical particles.

\section{The complete slow spectrum is determined without the last skew scalar}
\label{v48:sec:spectrum}

For the two-coordinate homogeneous equation define the original slow pencil
\begin{equation}
 \mathcal P(\zeta)=\zeta^2D_v+\zeta S+\overline G.
 \label{v48:eq:pencil}
\end{equation}
The variable $\zeta$ is spectral and is not an additional physical clock.

\begin{theorem}[A real pair cannot be removed by the retained gyroscopic term]
\label[theorem]{v48:thm:spectrum}
For every real value of the actual but unevaluated coefficient $\chi$, the
four roots of \eqref{v48:eq:pencil} are simple and are exactly
\begin{equation}
 \boxed{+\sigma,\ -\sigma,\ +\ii\omega,\ -\ii\omega,
 \qquad\sigma>0,\ \omega>0.}
 \label{v48:eq:roots}
\end{equation}
The full seven-dimensional homogeneous slow generator has these four simple
eigenvalues and three semisimple zero eigenvalues. At the retained root,
\begin{equation}
 \boxed{\sigma^2\ge\overline\lambda_+>13.96,
 \qquad\sigma>\frac{467}{125}=3.736.}
 \label{v48:eq:growth}
\end{equation}
These are spectral properties of the limiting slow equation, not a claim that
its selected forcing excites the positive eigenmode.
\end{theorem}
\begin{proof}
Let $\mu=-\det D_v>0$, $c_0=\det\overline G>0$, and
\[
 b_0=\operatorname{tr}(\operatorname{adj}(D_v)\overline G)+\chi^2.
\]
Direct expansion, retaining both off-diagonal entries, gives
\begin{equation}
 \boxed{\det\mathcal P(\zeta)=-\mu\zeta^4+b_0\zeta^2+c_0.}
 \label{v48:eq:quartic}
\end{equation}
The two roots in $x=\zeta^2$ have opposite signs, because their product is
$-c_0/\mu$. Explicitly,
\begin{equation}
 \sigma^2=\frac{b_0+\sqrt{b_0^2+4\mu c_0}}{2\mu},\qquad
 \omega^2=\frac{\sqrt{b_0^2+4\mu c_0}-b_0}{2\mu}.
 \label{v48:eq:rates}
\end{equation}
Both are strictly positive; the discriminant is positive and neither root is
zero. Thus the four spectral roots are simple.

For the full generator, set $m=j+D_v^{-1}Q^Th$. In the homogeneous system,
$(z,m)$ obey a four-dimensional autonomous equation while
$\dot h=-D_c^{-1}Qz$. Its characteristic polynomial is
$\zeta^3\det\mathcal P(\zeta)/\det D_v$. The four-dimensional block is
invertible. Therefore its coupling into $h$ can be removed by a fixed
triangular similarity, leaving three semisimple zeros. Equivalently its
kernel has $z=m=0$ and arbitrary $h$; the reconstructed $j$ is
$-D_v^{-1}Q^Th$.

Take a real nonzero nullvector $v$ of $\mathcal P(\sigma)$. Multiplication by
$v^T$ eliminates the skew term and gives
\[
 \sigma^2=-\frac{v^T\overline Gv}{v^TD_vv},\qquad v^TD_vv<0.
\]
The minimum of this quotient on the negative cone of $D_v$ is the positive
eigenvalue of $-D_v^{-1}\overline G$. To see this without a positive
replacement of $D_v$, diagonalize the real symmetric matrix
$\overline G^{-1/2}D_v\overline G^{-1/2}$. It has one negative eigenvalue
$d_-$ and one positive eigenvalue. On its negative cone, the quotient is at
least $-1/d_-$, exactly that positive generalized eigenvalue. Apply
\eqref{v48:eq:Jbarspectrum}; finally $(467/125)^2<13.96$.
\end{proof}

The statement is stronger than merely observing a positive eigenvalue of
V47's uncorrected return: it includes every constant-shift coupling and the
unknown skew derivative. No choice of $\chi$ can make all nonzero slow roots
imaginary. It gives no numerical upper bound for $\sigma$ without an enclosure
for $\chi$, and no comparison with a different companion root's frozen
canonical spectrum is made.

\section{The actual excitation and the four-row geometric test}
\label{v48:sec:forcing}

A positive homogeneous eigenvalue does not determine the solution from the
specific zero initial displacement and native forcing. Here that distinction
has an exact scalar test.

\begin{proposition}[The source coefficient of the growing slow mode]
\label[proposition]{v48:prop:residue}
Let $r,l$ be nonzero real right and left nullvectors of $\mathcal P(\sigma)$.
Then
$l^T(2\sigma D_v+S)r\ne0$. The coefficient of $e^{\sigma\tau}$ in the
solution $z$ of \eqref{v48:eq:two} is
\begin{equation}
 \boxed{
 r\,\frac{l^T(D_v\varpi_v+F_v/\sigma)}
              {l^T(2\sigma D_v+S)r}.}
 \label{v48:eq:residue}
\end{equation}
It vanishes if and only if
\begin{equation}
 \boxed{l^T(\sigma D_v\varpi_v+F_v)=0.}
 \label{v48:eq:no-excitation}
\end{equation}
\end{proposition}
\begin{proof}
For the constant-coefficient equation the Laplace transform, initially in a
right half-plane, is
\[
 \widehat z(\zeta)=\mathcal P(\zeta)^{-1}
                    (D_v\varpi_v+F_v/\zeta).
\]
The determinant has a simple zero at $\sigma$, so $\mathcal P(\sigma)$ has
one-dimensional right and left kernels. Its inverse has residue
$r l^T/[l^T\mathcal P'(\sigma)r]$; multiplication by the pencil and comparison
of the constant coefficient proves this formula and the nonzero denominator.
The residue yields \eqref{v48:eq:residue}. The other spectral roots and the
constant forcing yield different exponential or polynomial terms and cannot
cancel this one. This proves the equivalence.
\end{proof}

\begin{corollary}[Four scalar forcing conditions for the nonconstant geometric profile]
\label[corollary]{v48:cor:geometric}
For every $T>0$,
\begin{equation}
 \boxed{z(\tau)=0\text{ for all }0\le\tau\le T
 \quad\Longleftrightarrow\quad
 \varpi_v=0,\quad F_v=0.}
 \label{v48:eq:four}
\end{equation}
Let $\mathcal O_W$ be the retained complete flat weak geometric observation,
whose kernel consists exactly of the constant shifts. Then its leading slow
profile vanishes on this interval if and only if the same four scalar
equalities hold. In contrast, the complete weak rate profile $y$ vanishes
if and only if $\varpi_c=\varpi_v=d_v=0$, the seven conditions of V47.
\end{corollary}
\begin{proof}
If $z=0$, its initial derivative and acceleration equation imply respectively
$\varpi_v=0$ and $F_v=0$. The converse is uniqueness for
\eqref{v48:eq:two}. Under these equalities $h=\tau\varpi_c$ and
$j=\tau d_v$, with $d_v=-D_v^{-1}Q^T\varpi_c$. Thus a nonzero harmonic rate
may survive while this geometric observation is zero. By V46,
$\mathcal V=\Ran N_v$ has no nonzero constant vector and $\mathcal O_W$ is
injective on it. Since $\mathcal O_WN_c=0$, its profile is zero precisely
when $z=0$. For the complete rate, also require $h=0$; this forces
$\varpi_c=0$, and $F_v=0$ then forces $d_v=0$.
\end{proof}

The retained observation in this corollary is the complete weak shift-curvature
bank, not an assertion that every individual curvature component detects the
mode. In particular a nonzero residue in \eqref{v48:eq:residue} is visible in
that bank because $r\ne0$ lies in the nonconstant two-dimensional sector.
V47's proved $O_T(a)$ slow approximation transfers these profiles to the
original rates and the corresponding geometric observation on each fixed
bounded interval $t=a^2\tau$. No conclusion at $\tau\to\infty$ depending on
$a$ follows without an additional uniform error estimate.

\section{Verification, proved scope, and the smaller remaining source bank}
\label{v48:sec:status}

\paragraph{Proved in this body.}
The original covector consistency and prepared translation identities force
the actual $H_{\rm n}$ to have \eqref{v48:eq:blocks}. The harmonic cross is
evaluated with a uniform exact-root enclosure, and its positive return gives
\eqref{v48:eq:Jbarspectrum}. The complete seven-dimensional homogeneous slow
spectrum consists of three semisimple zeros, one real pair and one imaginary
pair; the positive exponent exceeds $3.736$ in the retained slow-time units.
The forced two-coordinate equation, the exact growing-source test, and the
four-row geometric criterion follow without assigning higher coefficients.

\paragraph{Inherited inputs and executed arithmetic.}
The standalone verifier embeds the unchanged V47 verifier with a checked hash.
Its default run replays the root and retained kernel/Schur enclosures; it does
not regenerate the older $79{,}224$-entry root bank, $17{,}334$ stationary-source
bank, or $5{,}046$ weak-Poisson bank. It evaluates the new harmonic block and
effective return by exact rational arithmetic, checks the four characteristic
signs, all retained constant annihilators, and the original commuting phase
derivatives. General finite symbolic tests check the balanced neutral
identity, the seven-state characteristic factorization, and the complete
forced reduction; those are algebraic regressions, not native trajectories.

The additional executed option \texttt{--regenerate-cubic} independently
rebuilds all $324$ entries of the complete cubic mixed-source map from the
original coefficient functional at the declared rational review point
\[
 (\xi_1,\xi_2,\xi_3,\xi_4)
       =(-5/16,313/1024,181/1024,1/5).
\]
It checks exact equality with the evaluated retained polynomial bank. This
point is not the algebraic root; the root-cube theorem uses the retained bank
with its certified root interval. The replay includes the two original
polarizations required by the physical grading. The verifier contains $23$
acceptance checks by default and $24$ with this independent source replay.
These are executable rational and symbolic checks with written analytical
proofs, not proof-assistant formalization or nonlinear spacetime integrations.

\paragraph{Unpaid source and dynamics.}
After the six harmonic entries have been evaluated, the old $32$-entry request
reduces to
\begin{equation}
 \boxed{\chi\quad(1),\qquad \varpi_c\quad(3),\qquad
         \varpi_v\quad(2),\qquad d_v\quad(2).}
 \label{v48:eq:remaining}
\end{equation}
None of these eight actual coefficients is numerically evaluated here. Four
forcing combinations suffice for the geometric-zero question, and one
additional scalar projection decides excitation of the positive mode once
the pencil is known. Their vanishing is not assumed. There is no new
lifespan theorem beyond V47's intervals $Ta^2$ for fixed finite $T$, no proof
of nonlinear instability of the chosen history, no attracting preparation,
and no spatial-refinement or coupled Einstein--SM continuum conclusion. The
separate V32 stress-compensation question remains unchanged. The useful next
calculation is the actual residual forcing, not another free-matrix stability
hypothesis or another rank certificate.

Every V47 byte before its final document command is preserved. New labels use
the prefix \texttt{v48:}. The supplied source hashes and the independently
executed cubic replay are recorded separately from inherited verification.

\begingroup
\renewcommand{\refname}{Additional references for Version 48}
\begin{thebibliography}{9}
\small\raggedright
\bibitem{v48:EB36}
\emph{Renewal Exact-Action Einstein Bridge Research Notes}, supplied V36,
\nolinkurl{Renewal_Exact_Action_Einstein_Bridge_Research_Notes_V36(2).tex}.
The transferred interfaces are \nolinkurl{v34:derivative} and
\nolinkurl{v36:Poisson-orders}; the root-dependent spectral statements remain
at their own source root.
\bibitem{v48:SM43}
\emph{Renewal SM Source Selection Research Notes}, supplied V43,
\nolinkurl{Renewal_SM_Source_Selection_Research_Notes_V43(2).tex},
\nolinkurl{sel43:thm:physical} and the terminal status ledger.
\bibitem{v48:PE51}
\emph{Renewal Perturbative ESM Continuum Research Notes}, supplied V51,
\nolinkurl{Renewal_Perturbative_ESM_Continuum_Research_Notes_V51(1).tex},
\nolinkurl{thm:v51-BV}, \nolinkurl{eq:v51-next}, and
\nolinkurl{sec:v51-ledgers}.
\bibitem{v48:YM79}
\emph{Renewal Yang--Mills Mass-Gap Research Notes}, supplied V79,
\nolinkurl{Renewal_Yang_Mills_Mass_Gap_Research_Notes_V79(1).tex},
\nolinkurl{r79:local} and \nolinkurl{r79:unpaid}.
\bibitem{v48:TM}
F.~Tisseur and K.~Meerbergen, \emph{The quadratic eigenvalue problem},
SIAM Review \textbf{43} (2001), 235--286,
\url{https://doi.org/10.1137/S0036144500381988}.
Primary-source context for polynomial and gyroscopic pencils; all particular
source identities, finite enclosures, and sign conclusions used here are
derived above.
\end{thebibliography}
\endgroup
\addtocontents{toc}{\protect\endgroup}
% END IMMUTABLE VERSION 48 BODY
% APPEND VERSION 49 HERE, BEFORE THE FINAL END-DOCUMENT COMMAND.



% BEGIN IMMUTABLE VERSION 49 BODY
\clearpage
\begin{center}
{\Large\bfseries Version 49: Exact removal of the slow source}\\[.5em]
{\large Five evaluated forcing coefficients, seven initial-rate controls,\newline
and a selected family with vanishing leading neutral motion}
\end{center}
\makeatletter
\addtocontents{toc}{\protect\begingroup}
\addtocontents{toc}{\protect\makeatletter}
\addtocontents{toc}{\protect\def\protect\l@section{\protect\bfseries\protect\@dottedtocline{1}{0em}{2.5em}}}
\makeatother
\addcontentsline{toc}{part}{Version 49: Exact removal of the slow source}

\section{The remaining forcing before another spectral calculation}
\label{v49:sec:context}

Version 48 fixes the qualitative slow spectrum but leaves eight scalar
coefficients unevaluated.  A positive exponent does not decide whether the
selected source excites it.  We first calculate the part of the forcing
which is already fixed by the quadratic and cubic envelopes.  It is nonzero
for the old choice of first rate.  We then go upstream of that choice:
five directions left unused by the first weak-source cancellation remove
five forcing coefficients, and two second-order active-rate directions
remove the remaining two.  All seven controls are initial selections in
the original constraint and consistency equations.

The cutoff is always $h=1/3$.  The algebraic leading fields $Y_*$ and first
multiplier tilt $\ell_*$ are exactly those of V39.  The signed source,
root isolation, weak kernel, and invariant-graph results of V39--V48 are
inherited.  In particular, this body neither regenerates the complete old
root bank nor rediscovers its signed inertia.  The new calculations use
that root's existing isolating cube of radius $2^{-112}$.  The companions
audited in V48 keep the same roles and boundaries; their other roots,
probability laws, and perturbative coefficients are not imported as
estimates of this deterministic initial-value problem.

The notation $\mathsf B$ denotes canonical constraint drift, $M$ denotes
the symmetric multiplier-rate Hessian, and $K=\mathbb K_1(Y_*)$ denotes
the \emph{first Poisson bracket}.  They are not interchangeable.  A rate
coefficient $v$ below is active, has zero lapse mean, and is fixed at the
initial cut; no selected rate is reimposed at later times.

\section{Five slow-source coefficients depend only on the first records and rate}
\label{v49:sec:forcing}

Let $I_A$ synthesize the 78-dimensional physical active space and $I_W$
the 29-dimensional weak space in the retained test coordinates.  We use
spatial-sum pairings for all covectors and matrices in a contraction;
the factor $1/27$ is restored only for reported physical squared norms.
Write
\[
 A\succ0,\quad L,\quad D=C_W-L^TA^{-1}L,\quad
 \mathsf Y=A^{-1}L,\quad K_{AW}=I_A^TKI_W.
\]
Let $N_c,N_v,D_c,D_v,T=K_{AW}N_v$ and $Q$ be the exact objects of V48.
Here $Q$ is the harmonic cross matrix, not the physical $\sqrt3$ grading.
The constant synthesis $I_WN_c$ is retained.  In particular,
\[
 D_c=N_c^TDN_c,\qquad D_v=N_v^TDN_v,\qquad
 Q=\mathscr R_WN_v-(\mathsf YN_c)^TT.
\]

At a flat-constraint-compatible leading field define the complete prepared
quadratic covector and cubic constant covector by
\begin{equation}
 \begin{aligned}
 B_2(Y,\ell)&=\widehat B_2(Y)+\mathbb K_1(Y)\ell,\\
 C_c(Y,\ell)&=t_c(Y)+\mathscr R_Y[c,\ell],\qquad c=e_1,e_2,e_3.
 \end{aligned}
 \label{v49:eq:envelopes}
\end{equation}
These are the original envelopes, with the preparation term in
$\mathscr R$ included.  At $(Y_*,\ell_*)$ all these covectors vanish.
Put $V=F_1Y_*+G\ell_*$; the exact flat relation is
$\mathcal L_{\rm con}V=0$.  For a physical active rate $v=I_Av_A$ put
\begin{equation}
 \begin{aligned}
 a_A(v)&=I_A^T\{D_YB_2(Y_*,\ell_*)[V]+K v\},\\
 c_3(v)&=D_Y C(Y_*,\ell_*)[V]+\mathscr R_{Y_*}[\,\cdot\,,v].
 \end{aligned}
 \label{v49:eq:two-covectors}
\end{equation}
Both are affine in $v$.  The symbol $a_A$ is a covector, not field amplitude.
Assume the already identified weak cancellation
\begin{equation}
 t_W+K_{WA}v_A=0.
 \label{v49:eq:first-weak}
\end{equation}
It ensures the same leading straight solution in the fast layer.

\begin{theorem}[Five higher-preparation-independent forcing coefficients]
\label[theorem]{v49:thm:forcing}
For each fixed finite active rate satisfying \eqref{v49:eq:first-weak},
use its exact original initializer and subsequent fast-graph selection.
The V47 slow coefficients obey
\begin{equation}
 \boxed{
 \begin{aligned}
 D_v d_v&=-T^TA^{-1}a_A(v),\\
 D_c\varpi_c&=(\mathsf YN_c)^Ta_A(v)-c_3(v).
 \end{aligned}}
 \label{v49:eq:five}
\end{equation}
Consequently the geometric force of V48 is
\begin{equation}
 \boxed{F_v=-T^TA^{-1}a_A(v)
 +Q^TD_c^{-1}\{(\mathsf YN_c)^Ta_A(v)-c_3(v)\}.}
 \label{v49:eq:F}
\end{equation}
No higher initial field coefficient enters these five scalar quantities.
The other two components $\varpi_v$ are not determined by this theorem.
\end{theorem}
\begin{proof}
On the exact nonconstant-constraint chart, the divided active drift has
leading value $B_2(Y,\ell)|_A$.  This follows by using the exact constraints
to replace the original linear drift, not by setting its second canonical
coefficient to zero.  Differentiate the exact sheared active equation of
V46 in its leading direction $b_q=(V,v_A)$.  The value of the divided
drift is zero at the selected point.  The derivative of $A^{-1}$ thus
multiplies zero.  Its remaining derivative is exactly $a_A(v)$.
The canonical component has the prefactor $a^3$, and the shear product
has no contribution to this limit because the leading weak field and
its derivative in direction $b_q$ are zero by
\eqref{v49:eq:first-weak}.  Hence V47's coefficient satisfies
\[
 Ub_q=(0,-A^{-1}a_A(v)).
\]
The original neutral coupling on an active variation $\eta_A$ is
$\Pi_0D^{-1}K_{AW}^T\eta_A$; in its two nonconstant coordinates this is
$D_v^{-1}T^T\eta_A$.  This proves the first identity of
\eqref{v49:eq:five}.

For the second identity it is useful to differentiate the full physical
consistency equation rather than a projected slow model.  Its leading
active differentiated source is $a^2a_A(v)$.  On a constant weak test the
leading differentiated drift is $a^3c_3(v)$.  Indeed the complete constant
cubic envelope on an exactly prepared record is $C(Y,\ell)$, independent
of its second canonical coefficient.  Its derivative along the original
leading tangent is $(V,v)$.  The required quadratic mean tangency also
holds: it is the constant part of the already vanishing quadratic drift.
Terms in $\dot M\lambda_t$ on such a weak test start at $a^4$, because
$\lambda_t=av+O(a^2)$ is active to leading order.  They therefore do not
enter its $a^3$ coefficient; they are not identically discarded.

Eliminate the active acceleration with the original graded blocks.  The
constant projection of the remaining weak source is
$c_3(v)-(\mathsf YN_c)^Ta_A(v)$.  Multiplying the weak acceleration by
$a$ and projecting onto the semisimple neutral kernel gives
\[
 D_c\lim_{a\to0} a(\lambda_{tt,W})_c
 =(\mathsf YN_c)^Ta_A(v)-c_3(v).
\]
This limit is precisely $D_c\varpi_c$: in the projected velocity equation,
$\Pi_0B=0$, the initial weak velocity is $O(a)$, and
$\dot y(0)=\varpi$.  A fast second weak rate does not change this
projection.  This proves the second identity.  Substitution in V48's
$F_v=D_vd_v+Q^T\varpi_c$ proves \eqref{v49:eq:F}.
\end{proof}

The differentiations in this proof are through the exact constraint chart.
They use its leading tangent and the complete prepared constant envelope.
Differentiating an unprepared cubic constant row would change $c_3$ and
would not establish \eqref{v49:eq:five}.

\section{The old initial-rate selection has nonzero geometric forcing}
\label{v49:sec:evaluation}

\begin{proposition}[An exact-root enclosure of the actual old forcing]
\label[proposition]{v49:prop:old}
For V40's original transverse rate $v_*$, and hence for the V44--V48
selected family built from that rate, the following componentwise intervals
hold:
\begin{equation}
 \begin{aligned}
 d_v&\in(2750,2760)\times(339,342),\\
 \varpi_c&\in(-539,-537)\times(1079,1082)\times(510,514),\\
 F_v&\in(-102000000,-100000000)\times(197000000,199000000).
 \end{aligned}
 \label{v49:eq:old-intervals}
\end{equation}
In particular $F_v\ne0$.  No choice of the uncomputed skew scalar or of
$\varpi_v$ can make this old family's leading geometric profile identically
zero on a nontrivial slow interval.
\end{proposition}
\begin{proof}
The verifier reconstructs the physical $Y_*,\ell_*,v_*$ on the retained
root cube, using the exact homogeneous $P_2,\Phi_3$ polarization formulas
and the inherited first-bracket bank.  It recomputes the full 108-component
covector in \eqref{v49:eq:two-covectors} before active restriction, and all
three components of $c_3$.  All phase shifts, off-diagonal Frobenius
occurrences, and physical $\sqrt3$ factors are retained.

The new exact entry-radius bounds for these full covectors are respectively
less than $0.000706$ and $0.002067$.  Contracting them with the inherited
root-cube enclosures for the signed blocks and neutral synthesis gives
approximate centres
\[
 d_v\simeq(2754.33410767,340.18218938),
\]
\[
 \varpi_c\simeq(-537.98852513,1080.13641959,511.69051067),
\]
\[
 F_v\simeq(-101054001.9142544,197868841.7101108).
\]
Their certified entry radii are less than $0.000199$, $0.0000402$, and
$8.878$, respectively.  Rational endpoint comparisons imply
\eqref{v49:eq:old-intervals}.  The decimals are explanatory, not proof
inputs.  By V48's exact two-coordinate equation, $z\equiv0$ would require
both $\varpi_v=0$ and $F_v=0$, which is impossible here.
\end{proof}

For the retained complete weak geometric observation, this implies a
nonzero leading profile on every fixed interval $[0,T]$, $T>0$.  Thus
V47's convergence theorem gives a positive limit for the supremum of the
amplitude-normalized change of that complete observation.  The value of
that supremum is not enclosed here.  This conclusion does not identify
the coefficient of its growing exponential: nonzero forcing can still be
orthogonal to the growing spectral residue.  Individual curvature
components need not all detect it.

\section{Five unused first-rate directions cancel the five evaluated coefficients}
\label{v49:sec:first-control}

The first cancellation \eqref{v49:eq:first-weak} imposes 26 independent
conditions on the 78 active directions.  It leaves 52 dimensions available.
We test their actual source map rather than assume that they can control
the slow forcing.

For reproducibility define $I_A$ by orthogonally projecting the original
V41 active raw columns off the weak space and the mean lapse.  This is an
exact rational projection; these columns are independent but not declared
orthonormal.  Let $E_W$ synthesize the original 26 nonconstant weak tests,
and let $Z_T$ be V40's 26 selected transverse columns.  Define
\begin{equation}
 H_T=E_W^TKZ_T,\qquad
 P_A=I_A-Z_TH_T^{-1}E_W^TKI_A.
 \label{v49:eq:Pk}
\end{equation}
Then $E_W^TKP_A=0$ exactly.  The constant weak rows vanish independently.
Let $E_5$ be columns $7,18,25,9,20$ of $P_A$, in that order, with zero-based
indexing of the retained 78-column active list.

Introduce the affine forcing vector
\[
 b_5(v)=\binom{D_vd_v(v)}{D_c\varpi_c(v)}
\]
and its linear part on physical active perturbations,
\begin{equation}
 \mathcal L_5\delta v=
 \binom{-T^TA^{-1}I_A^TK\delta v}
 {(\mathsf YN_c)^TI_A^TK\delta v-\mathscr R_{Y_*}[\,\cdot\,,\delta v]}.
 \label{v49:eq:map5}
\end{equation}

\begin{theorem}[A certified first-rate control of all five evaluated sources]
\label[theorem]{v49:thm:first-rate}
Throughout the retained root cube the matrix $H_5=\mathcal L_5E_5$ is
invertible.  Define at the exact root
\begin{equation}
 \boxed{\bar v=v_*-E_5H_5^{-1}b_5(v_*).}
 \label{v49:eq:vbar}
\end{equation}
Then $\bar v$ is active, has zero component means, and satisfies
\begin{equation}
 \boxed{t_W+K_{WA}\bar v_A=0,\qquad
        d_v(\bar v)=0,\qquad \varpi_c(\bar v)=0.}
 \label{v49:eq:first-zero}
\end{equation}
Its physical norm obeys
\begin{equation}
                         \boxed{3758<\|\bar v\|_h<3761.}
 \label{v49:eq:vbar-norm}
\end{equation}
The leading fields and first multiplier tilt remain $Y_*,\ell_*$.
\end{theorem}
\begin{proof}
The verifier supplies a fixed dyadic inverse proposal for $H_5$.  Three
rational Newton refinements, with explicit rounding and final residual
verification, give
\[
 \sup_{\text{cube}}\|I-P_5H_5\|_\infty<\frac1{6000},
 \qquad \|H_5^{-1}\|_\infty<\frac1{3000}.
\]
In the first strict bound the actual certified residual is below
$0.000144428$.  The inequality proves invertibility at every point of the
cube, not merely at its centre.  The exact solution definition in
\eqref{v49:eq:vbar} and the affine identity
$b_5(v+\delta v)=b_5(v)+\mathcal L_5\delta v$ imply the two zero sources.
The projector identity \eqref{v49:eq:Pk} preserves the previous weak
cancellation.  All synthesizing columns are active and mean-zero.

For orientation, the five coefficients in $\bar v-v_*=E_5\gamma$ are
approximately
\[
 \gamma\simeq(128.761400,-1101.675134,2499.757056,
                         -779.126166,761.973091).
\]
The exact-root coefficient error is less than $0.155$ per entry.  The
resulting 108 physical rate entries have error below $0.551$ relative to
the stored exact rational centres.  Bounding their squared norm above and
below componentwise, then dividing by $27$, gives the strict rational
inequalities $3758^2<\|\bar v\|_h^2<3761^2$.
No minimal-norm assertion is made.
\end{proof}

The first rate now has lapse as well as transverse-shift components.  The
V40 initializer extends to this case without a new hypothesis: only its
membership in the active space was used in the block orders of $aMv$.
That term has order $a^3$ in active rows and order $a^4$ in weak rows.
The three cubic constant equations and their verified root Jacobian are
therefore unchanged.  For every fixed finite $\bar v$, the same
block-triangular implicit-function proof gives exact $P=0$ records with
$\mathsf B+aM\bar v=0$.  Condition \eqref{v49:eq:first-zero} gives
$\mathbb Lb=0$ for $b=(V,\bar v_A,0)$.  The signed matrix and all its
leading spectral gaps depend on $Y_*,\ell_*$, not on this first-rate
choice, so the subsequent fast-graph selection also applies.  Its local
radii may change and are not numerically bounded here.

\section{Second-order active rates access the remaining two slow sources}
\label{v49:sec:second-control}

We now hold $\bar v$ fixed.  For a finite active coefficient $\rho_A$ and
weak coefficient $w$, seek exact initial records with actual rate
\begin{equation}
 \lambda_t(0)=a\bar v+a^2(I_A\rho_A+I_Ww).
 \label{v49:eq:second-rate}
\end{equation}
The fast part of $w$ will be selected by the same nine-dimensional
invariant-graph incidence.  It is not a free neutral slow forcing.

\begin{lemma}[Second active-rate access preserves the first two initial jets]
\label[lemma]{v49:lem:active-access}
For every fixed finite $\rho_A$, the exact constraint and consistency
initializer can be extended to \eqref{v49:eq:second-rate}.  Relative to
the family with $\rho_A=0$, its physical initial fields and multiplier
values differ by $O_{\rho_A}(a^3)$.  The fast weak coefficient can then
be selected once so that these initial records lie on a common locally
invariant fast graph.  The resulting selected families are $C^1$ in $a$,
with the same coefficients through order $a^2$.
\end{lemma}
\begin{proof}
Use the original analytic constraint initializer with
$\lambda=e+a\ell(q)+a^2m$.  The additional active-rate term is
$a^2MI_A\rho_A$: it is $O(a^4)$ on active rows and $O(a^5)$ on weak
rows.  Dividing the 107 consistency equations by $a^3$, the nonconstant
rows change by $O(a)$ and the three constant rows by $O(a^2)$.
At $a=0$ the latter do not depend on $m$, while the former have the
verified invertible $K|_{\mathcal E^\perp}$ derivative in $m$.
The joint root derivative stays block triangular and invertible.  Its
implicit solution has $\delta m=O(a)$ and $\delta q=O(a^2)$: the order-$a$
change of $m$ can affect the constant equations only at order $a^2$.
The original fields and multipliers therefore change by $O(a^3)$.
The final constant-lapse equation still follows from the original radial
identity and lapse normalization.

The extra weak rate has the higher active/weak consistency orders
$a^5,a^6$, as in V43.  In the fixed normalized chart, adding $\rho_A$
changes the first displacement forcing by $J_A\rho_A$ and changes
neither the leading layer matrix nor its first linear amplitude coefficient.
In V44's divided normal-velocity equation this is a finite additive term.
Its derivative in the nine fast weak controls remains the same isomorphism
$J_W$.  The same graph-incidence implicit-function argument selects those
controls, depending on $\rho_A$, and puts the exact records on the fixed
local graph of the original field.  The graph's auxiliary extension can
be kept fixed during this local parameter construction.  The selected
weak controls affect neither of the first two physical initial jets.
The graph and its selection have the inherited differentiability; no
analyticity of the final graph-selected family is asserted.
\end{proof}

\begin{theorem}[The remaining slow forcing shifts by the original neutral coupling]
\label[theorem]{v49:thm:shift}
Compare the selected families of \cref{v49:lem:active-access}, using their
own initial cuts and the same leading spectral bases.  The V47 slow
coefficients transform as
\begin{equation}
 \boxed{H_{\rm n}\mapsto H_{\rm n},\quad d\mapsto d,\quad
        \varpi\mapsto\varpi+R_0(0,\rho_A).}
 \label{v49:eq:shift}
\end{equation}
In particular $\varpi_c$ is unchanged and
\begin{equation}
                 \boxed{\varpi_v\mapsto
                       \varpi_v+D_v^{-1}T^T\rho_A.}
 \label{v49:eq:shift2}
\end{equation}
\end{theorem}
\begin{proof}
The change of base point in the normalized original field is $O(a^2)$.
The exact constraint chart, active shear and common graph reconstruction
therefore have unchanged first amplitude derivatives.  Thus $U,C,H_1$
of V47 do not change.  Equivalently, the upper field has its explicit
$a^2$ factor and its limiting second amplitude derivative is evaluated
at unchanged leading records; the lower field's first derivative changes
only by $O(a^2)$.  Hence $H_{\rm n}$, $\gamma=Ub_q$, $d=R_0\gamma$ and
$\beta=\Pi_0Cb_q$ are unchanged.

The actual initial canonical velocity has unchanged coefficients through
order $a^2$, so the canonical part of $r_*$ does not change.  The active
part of the sheared initial velocity, divided by $a^3$, changes by exactly
$a\rho_A+O(a^2)$.  The extra weak rate contributes to that shear only
at order $a^2$.  Hence
$r_*\mapsto r_*+(0,\rho_A)$, which proves \eqref{v49:eq:shift}.
The image of $R_0$ is the two-dimensional space synthesized by $N_v$,
and its pure active formula is $D_v^{-1}T^T\rho_A$.
This gives \eqref{v49:eq:shift2} and shows that the harmonic components
are unchanged.  No derivative of a numerically chosen gauge or of a
varying spectral rank is taken.
\end{proof}

\section{An exact initial selection makes every leading neutral source vanish}
\label{v49:sec:seven-zero}

Let $L_2$ be the two physical lapse columns
$e_2-\mathbf1/27$ and $e_{25}-\mathbf1/27$, extended by zero shifts.
Let $J_2$ be their \emph{active coordinate} matrix:
\[
 J_2=(I_A^TI_A)^{-1}I_A^TL_2.
\]
This inverse Gram is essential because the retained active synthesis is
not orthonormal.  Define
\begin{equation}
                   C_2=D_v^{-1}T^TJ_2.
 \label{v49:eq:C2}
\end{equation}

\begin{proposition}[Certified two-rate control]
\label[proposition]{v49:prop:two-control}
The matrix $C_2$ is invertible on the exact-root cube, with an exact
dyadic inverse proposal satisfying
\begin{equation}
 \sup\|I-P_2C_2\|_\infty<10^{-6},\qquad
                         \|C_2^{-1}\|_\infty<217.
 \label{v49:eq:two-bounds}
\end{equation}
Its rational centre is approximately
\[
 \begin{pmatrix}
 0.01272107&-0.05507105\\
 0.00475040&-0.00040433
 \end{pmatrix}.
\]
\end{proposition}
\begin{proof}
Form \eqref{v49:eq:C2} with the same exact-root signed and kernel
enclosures used above and the exact active Gram inverse.  The residual
of the fixed dyadic inverse is less than $2.095\,10^{-7}$, and its
Neumann inverse bound is less than $216.307$.  These strict rational
inequalities imply \eqref{v49:eq:two-bounds}.  They are consistent with
V46's independent two-lapse rank certificate, which is an inherited
input rather than another new rank theorem here.
\end{proof}

\begin{theorem}[A nonempty exact family with all seven slow forcing components zero]
\label[theorem]{v49:thm:zero-family}
Use $\bar v$ from \cref{v49:thm:first-rate}.  Let
$\varpi_v^{\rm base}$ be the actual two-component forcing of its selected
family at second active rate zero.  Define
\begin{equation}
            \boxed{\rho_A^\sharp=-J_2C_2^{-1}\varpi_v^{\rm base}.}
 \label{v49:eq:rho}
\end{equation}
Then there is a nonempty selected family of exact original initial records,
with the same $Y_*,\ell_*$, whose unique original continuation has
\begin{equation}
 \begin{aligned}
 P&=0,\\
 \mathsf B+M\{a\bar v+a^2(I_A\rho_A^\sharp+I_Ww(a))\}&=0,\\
 \lambda_t(0)&=a\bar v+a^2(I_A\rho_A^\sharp+I_Ww(a)),
 \end{aligned}
 \label{v49:eq:exact-family}
\end{equation}
and whose complete seven-dimensional slow forcing satisfies
\begin{equation}
                         \boxed{\varpi=0,\qquad d=0.}
 \label{v49:eq:allzero}
\end{equation}
Here $w(a)$ is the fixed initial fast-graph selection, not a feedback rule.
The numerical value of $\rho_A^\sharp$ is not tabulated.
\end{theorem}
\begin{proof}
The first selection gives $d_v=0$ and $\varpi_c=0$ by
\eqref{v49:eq:first-zero}.  The coefficient
$\varpi_v^{\rm base}$ exists as a finite derivative of the actual selected
initializer by V47.  \Cref{v49:prop:two-control} makes \eqref{v49:eq:rho}
a well-defined finite constant.  The access lemma constructs the exact
initial records and their graph selection for this constant.  The forcing
shift theorem then gives
\[
 \varpi_v^{\rm new}=\varpi_v^{\rm base}
                 -C_2C_2^{-1}\varpi_v^{\rm base}=0,
\]
without changing $d_v$ or $\varpi_c$.  These are all the components of
$\varpi,d$.  The original signed-Hessian regularity at $Y_*$ persists
for sufficiently small amplitude and makes the stated tangent the unique
normalized original rate.  Positive lapse and metric, and the original
mean normalizations at initialization, follow from the unchanged leading
chart.  No upper numerical bound on the finite second control or a uniform
radius for unbounded second controls is inferred.
\end{proof}

This theorem is not a computation of the unknown quartic initializer.
It proves that its two remaining finite forcing coefficients are reachable
by exact initial-rate access with a verified onto derivative.  They need
not be known numerically for the existence statement.  An implemented
preparation would still have to evaluate them and the graph selection.
The same distinction was already necessary for V39's analytic higher
initializer and V44's invariant graph.

\section{Original histories with vanishing leading slow rate and curvature change}
\label{v49:sec:dynamics}

\begin{corollary}[Second-order observation control on every fixed slow interval]
\label[corollary]{v49:cor:dynamics}
For the family in \cref{v49:thm:zero-family} and every fixed $T<\infty$,
all sufficiently small positive amplitudes have their original histories on
\begin{equation}
                              0\le t\le Ta^2.
 \label{v49:eq:interval}
\end{equation}
Every original constraint, the selected fast invariant graph, positive
lapse and metric, and signed regularity persist.  Uniformly there,
\begin{equation}
 \boxed{\lambda_{t,W}=O_T(a^2),\qquad
        \lambda_{t,A}-a\bar v_A=O_T(a^2).}
 \label{v49:eq:rates}
\end{equation}
For each separately retained metric Riemann, electric-link and spatial-link
curvature observation,
\begin{equation}
 \boxed{\|\mathscr O_a(t)\|=O_T(a),\qquad
        \|\mathscr O_a(t)-\mathscr O_a(0)\|=O_T(a^2).}
 \label{v49:eq:observations}
\end{equation}
The canonical motion is still nontrivial:
\[
 X(a^2\tau)-X(0)=a^3\tau V+O_T(a^4).
\]
\end{corollary}
\begin{proof}
The leading signed matrix, fast spectral gap, bounded nonfast weak
semigroup, exact constraint chart and graded active equation are unchanged.
The exact rate selection retains $\mathbb Lb=0$.  V47's growing-tube
continuation and velocity argument therefore applies with $v_*$ replaced
by the fixed finite $\bar v$ and the adjusted initial coefficient $r_*$.
It supplies the same interval \eqref{v49:eq:interval} and $O_T(a)$
velocity convergence to its actual slow equation.  For this equation
\eqref{v49:eq:allzero} and zero initial $y,j$ imply $y=j=0$ by uniqueness.
This conclusion is independent of the uncomputed skew coefficient $\chi$.
Undoing the original velocity scales gives \eqref{v49:eq:rates}.
The fixed-cutoff observation formulas of V47 are smooth in the original
fields, canonical rates and multiplier rate.  Their leading weak change
is $\mathcal O_Wy=0$, so their unscaled changes are $O_T(a^2)$.
Their initial values are $O(a)$.  The canonical displacement formula
follows from the same convergence theorem and its unchanged leading $V$.
The interval constant $T$ is fixed before $a\to0$; no substitution of a
growing $T$ into these estimates is made.
\end{proof}

The first-rate selection changes the order-$a$ initial multiplier velocity.
Thus the new family's initial curvature coefficient need not equal the
old family's.  All changes in \eqref{v49:eq:observations} are measured
from the \emph{new} initial cut.  No past trajectory is repaired, no
curvature component is subtracted, and no time rescaling is executed.
The nine normal repelling directions remain, as does the growing pair in
the homogeneous slow equation.  Their leading forcing has been selected
to vanish; neither spectrum nor the original dynamics has been made
contractive.

\section{Verification ledger and remaining mathematical work}
\label{v49:sec:status}

\paragraph{Proved in this body.}
Five slow forcing coefficients are expressed by the complete quadratic and
cubic envelopes and enclosed at the exact root.  For the old first-rate
choice their geometric forcing is nonzero.  A five-column actual active
control, inside the already required weak-cancellation kernel, is certified
onto and supplies an explicitly defined new first rate.  A second two-column
active control supplies exact access to the remaining two slow sources.
The initial incidence and invariant-graph arguments give an exactly
constrained selected family with all seven leading slow sources zero.
Its original histories have \eqref{v49:eq:rates}--\eqref{v49:eq:observations}
on every fixed $Ta^2$ interval.

\paragraph{Inherited and replayed.}
The algebraic root, original coefficient definitions, first-tilt equation,
nonlinear stationary chart, exact mean initializer, signed leading Hessian,
neutral basis, and fast invariant graph retain their earlier scopes.
The verifier replays V48's certificate before using these enclosures.
The old 79,224-entry root bank, 17,334-entry stationary-source bank and
5,046-entry weak-bracket bank are not represented as freshly regenerated.
The new source calculation recomputes the 108 differentiated drift entries
and three complete cubic-mean derivatives from the original envelopes,
using the retained bracket bank for its linear contraction.

The residual-verification methodology is standard \cite{v49:Rump}; the particular source contractions and control certificates are evaluated here.
All new finite matrix products and residuals use rational arithmetic.
For uniform entry radii $r_X,r_Y$ the matrix-product enclosure is
\[
 r_{XY}\le
 \max_i\sum_k|X_{0,ik}|\,r_Y
 +\max_j\sum_k|Y_{0,kj}|\,r_X+n r_Xr_Y.
\]
Proposed inverses are certified by $\|I-PM\|<1$ on the whole enclosure.
The selected five coefficients are enclosed from their equation residual,
not by interpreting a small floating residual as an exact zero.  The exact
zero identities follow from the displayed inverse-defined controls.

\paragraph{Unpaid.}
The two components $\varpi_v^{\rm base}$, their second active-rate value,
and the nonlinear fast graph are analytically defined but not numerically
tabulated.  An implemented source has not been shown to prepare these
records.  No improved power of the physical lifespan, attracting basin,
spatial-refinement estimate, differentiated-curvature bound, or complete
matter-coupled Einstein--SM continuum follows from this iteration.
The next dynamical issue is the next nonzero slow source for the selected
family, or an invariant slow selection of the original flow.  It is not
another rank calculation or a presumption that the remaining forcing of
the old family was zero.  The separate V32 unit-initial-multiplier stress
problem is unchanged.

Every V48 byte before its final document command is preserved.  The new
verification package distinguishes source recomputation, inherited replay,
exact linear certificates, and analytical existence statements.  Its
finite algebraic checks are not a proof-assistant formalization or a
nonlinear spacetime simulation.  New labels use the prefix \texttt{v49:}.

\begingroup
\renewcommand{\refname}{Additional source for Version 49}
\begin{thebibliography}{9}
\small\raggedright
\bibitem{v49:EB36}
\emph{Renewal Exact-Action Einstein Bridge Research Notes}, supplied
cumulative V36,
\nolinkurl{Renewal_Exact_Action_Einstein_Bridge_Research_Notes_V36(2).tex}.
The original prepared quadratic drift and cubic constant envelope, including
its higher-preparation independence, are used at the scopes established in
its V20--V21 and V34--V36 bodies.  Root-specific statements from that
companion are not substituted for this lineage's V39 root certificate.
\bibitem{v49:Rump}
S.~M.~Rump, \emph{Verification methods: Rigorous results using floating-point
arithmetic}, Acta Numerica \textbf{19} (2010), 287--449,
\url{https://doi.org/10.1017/S096249291000005X}.
Cited for established verification methodology, not for the source-specific
control or initialization theorems above.
\end{thebibliography}
\endgroup
\addtocontents{toc}{\protect\endgroup}
% END IMMUTABLE VERSION 49 BODY
% APPEND VERSION 50 HERE, BEFORE THE FINAL END-DOCUMENT COMMAND.



% BEGIN IMMUTABLE VERSION 50 BODY
\clearpage
\begin{center}
{\Large\bfseries Version 50: The next rate scale requires weak motion}\\[.5em]
{\large A certified persistence obstruction and a positive constrained\\
transport law at fixed spatial cutoff}
\end{center}
\makeatletter
\addtocontents{toc}{\protect\begingroup}
\addtocontents{toc}{\protect\makeatletter}
\addtocontents{toc}{\protect\def\protect\l@section{\protect\bfseries\protect\@dottedtocline{1}{0em}{2.5em}}}
\makeatother
\addcontentsline{toc}{part}{Version 50: The next rate scale requires weak motion}

\section{Persistence before another initial cancellation}
\label{v50:sec:context}

Version 49 constructs exact selected records with zero leading neutral source
on every fixed interval $t\le Ta^2$.  This is not an invariant statement about
the first-rate coefficients on the longer scale $t=a\theta$.  Before selecting
another higher coefficient, we test that persistence using the original
constraint and multiplier-consistency equations.  The resulting obstruction
is independent of every higher initial coefficient and of the choice of
bounded first active rate at the same leading fields.

Throughout this body $h=1/3$, and $Y_*,\ell_*$ are exactly the V39 algebraic
mixed-mode root and its first multiplier tilt.  They are not decimal
substitutes.  The parent V49 verifier has been rerun, including its original
108-component differentiated drift and three cubic-mean derivatives.  The
new numerical calculation contracts those source enclosures with the
original first Poisson bracket.  The root isolation, positive leading active
metric, signed stationary chart, weak-kernel rank and neutral coupling are
inherited at their previous scopes.  No new root, preparation graph or
physical lifespan is claimed.

The additions are a certified nonzero stationary-source mismatch, a
no-persistence theorem for a negligible nonconstant weak rate, and a closed
76-dimensional differential equation for the necessary nonconstant rate
limit when such a rate is allowed.  The latter is a conditional limit of
bounded-rate original histories, not a proof that those histories exist on
the longer interval.  The proof uses ordinary finite-dimensional
observability and residual verification; neither general method is claimed
as new \cite{v50:Hautus,v50:Rump}.

\section{A new original-source certificate: the active equilibrium is incompatible}
\label{v50:sec:certificate}

Let $I_A$ be the retained 78-column active synthesis.  Let $I_g$ consist of
the 26 nonconstant columns of the original weak synthesis $I_W$, and let
$I_c$ be its three constant-shift columns.  The active and nonconstant weak
spaces are orthogonal in the physical pairing, but these column bases need
not be orthonormal.  Every matrix below uses the same spatial-sum covector
convention.  Put
\begin{equation}
 K_A=I_A^TKI_A,\qquad C=I_A^TKI_g,\qquad
 K_g=I_g^TKI_g,\qquad K=\mathbb K_1(Y_*).
 \label{v50:eq:blocks}
\end{equation}
Thus $K_A^T=-K_A$, $K_g^T=-K_g$, and the opposite cross block is $-C^T$.
The symbol $A\succ0$ denotes the leading \emph{symmetric} active
multiplier-Hessian block from V41; it is not $K_A$.

Use the complete prepared quadratic covector from V49:
\[
 B_2(Y,\ell)=\widehat B_2(Y)+\mathbb K_1(Y)\ell,
 \qquad V=F_1Y_*+G\ell_*.
\]
Its derivative with the leading multiplier held fixed is
\begin{equation}
 \mathfrak t=D_YB_2(Y_*,\ell_*)[V],\qquad
 t_A=I_A^T\mathfrak t,\qquad t_g=I_g^T\mathfrak t.
 \label{v50:eq:t}
\end{equation}
In particular the term $\mathbb K_1(V)\ell_*$ is included in
$\mathfrak t$.  The parent source file contains
$\mathfrak t+Kv_*$ and $v_*$, so the new verifier subtracts the latter
contraction with its complete uncertainty.  No unknown quartic coefficient
is assigned a value.

\begin{theorem}[An exact-root active-source separation]
\label[theorem]{v50:thm:separation}
On the retained root cube of radius $2^{-112}$, $K_A$ is invertible and
\begin{equation}
 \boxed{\|K_A^{-1}\|_\infty<6.}
 \label{v50:eq:inverse}
\end{equation}
Define the active equilibrium and its nonconstant weak mismatch by
\begin{equation}
 v_{\rm e}=-K_A^{-1}t_A,\qquad
 \rho=t_g-C^Tv_{\rm e}=t_g+C^TK_A^{-1}t_A.
 \label{v50:eq:rho}
\end{equation}
For component $7$ in the zero-based nonconstant weak column order,
\begin{equation}
 \boxed{81476000<\rho_7<81478000.}
 \label{v50:eq:rho-interval}
\end{equation}
Consequently
\begin{equation}
 \operatorname{rank}\begin{pmatrix}K_A&t_A\\-C^T&t_g\end{pmatrix}=79,
 \qquad
 \operatorname{rank}\begin{pmatrix}K_A\\-C^T\end{pmatrix}=78.
 \label{v50:eq:rank}
\end{equation}
These are the actual first-source matrices at this root, not the full
symmetric multiplier Hessian.
\end{theorem}
\begin{proof}
The weak test in \eqref{v50:eq:rho-interval} is the original phase gradient
of the scalar point difference $e_7-e_{26}$, with no lapse component.
The verifier assembles $K_A,C,t_A,t_g$ from the inherited polynomial and
source enclosures with the physical $\sqrt3$ factors retained.  In
particular an active rate is converted to its coordinates using
$(I_A^TI_A)^{-1}I_A^T$, not $I_A^T$ alone.

A fixed rational inverse proposal $P_A$ satisfies the strict enclosure
\[
 \|I-P_AK_{A,0}\|_\infty
 +\|P_A\|_\infty\,78r_{K_A}<2^{-85},
\]
where $r_{K_A}$ is the uniform entry radius.  The actual left side is less
than $1.628\,10^{-27}$.  The Neumann bound gives an inverse norm below
$5.679$, proving \eqref{v50:eq:inverse}.  The exact rational centre of
$\rho_7$ is approximately $81476858.41736083$, with uniform entry radius
less than $537$.  Rational endpoint comparisons prove
\eqref{v50:eq:rho-interval}; the printed decimal is not used for a sign.

The first 78 rows give the rank of the original column matrix.  A solution
of its augmented stationary-source equation would have to equal
$-K_A^{-1}t_A$ in the active rows, and would then have weak residual
$\rho\ne0$.  This proves \eqref{v50:eq:rank}.  Equivalently, with $e$ the
coordinate-7 weak row, the inverse-defined covector
\[
             (e^TC^TK_A^{-1},\ e^T)
\]
annihilates $\binom{K_A}{-C^T}$ and pairs with $\binom{t_A}{t_g}$ to
$\rho_7$.  This is an exact annihilator identity, not a floating nullspace.
\end{proof}

All original constant-shift columns satisfy $KI_c=0$.  The verifier checks
this coefficientwise in the six matrices of the retained first-source
bank.  Allowing bounded harmonic shift rates therefore cannot change
\eqref{v50:eq:rho}.  The large numerical values above use unrescaled
spatial-sum tests; no physical coupling or cutoff-uniform size is inferred.

\section{No active transport can preserve the negligible-weak-rate condition}
\label{v50:sec:affine}

The following finite equation will be forced by the original action if
its nonconstant weak rate remains negligible on $t=a\theta$:
\begin{equation}
 A\dot v+t_A+K_Av=0,\qquad t_g-C^Tv=0.
 \label{v50:eq:flat-law}
\end{equation}
A dot denotes $d/d\theta$ only in this body.

\begin{lemma}[Exact persistence criterion and an all-initial-rate obstruction]
\label[lemma]{v50:lem:observability}
Set $\Omega=-A^{-1}K_A$.  There is a solution of
\eqref{v50:eq:flat-law} on a nonempty open time interval if and only if
$\rho=0$.  More precisely, when $\rho=0$ the admissible initial values are
\[
 v(0)=v_{\rm e}+h,\qquad
 C^T\Omega^kh=0\quad(0\le k<78).
\]
At the actual mixed-mode root no such solution exists, for any finite
initial active rate.
\end{lemma}
\begin{proof}
The differential equation has the entire solution
$v(\theta)=v_{\rm e}+e^{\Omega\theta}h$.  Its weak output is
$\rho-C^Te^{\Omega\theta}h$.  If this vanishes on an interval,
analyticity makes it vanish identically.  Differentiation gives
$C^T\Omega^kh=0$ for every $k\ge1$.  Since $\Omega$ is invertible,
its characteristic polynomial has a nonzero constant coefficient.
Cayley--Hamilton therefore expresses $I$ as a linear combination of
$\Omega,\ldots,\Omega^{78}$, giving $C^Th=0$.  The undifferentiated
identity implies $\rho=0$.  Conversely, the 78 displayed homogeneous
conditions and Cayley--Hamilton give zero output at every time, and
$h=0$ supplies at least one solution whenever $\rho=0$.
Apply \cref{v50:thm:separation} for the actual source.
\end{proof}

There is a positive finite-window version which does not bound the initial
rate in advance.  Put $O=-C^T$ and, for $T>0$, define
\begin{equation}
 \begin{split}
 \mathcal W_T&=\int_0^T e^{\Omega^T\theta}O^TOe^{\Omega\theta}\,d\theta,\\
 b_T&=\int_0^T e^{\Omega^T\theta}O^T\rho\,d\theta,\\
 \epsilon_T&=T\|\rho\|^2-b_T^T\mathcal W_T^\dagger b_T.
 \end{split}
 \label{v50:eq:distance}
\end{equation}
Here the Moore--Penrose inverse is on a finite positive Gram matrix; it
is not an inverse of the original indefinite action.

\begin{proposition}[A strictly positive output-distance floor]
\label[proposition]{v50:prop:distance}
For the actual source and every $T>0$,
\begin{equation}
 \boxed{\epsilon_T
  =\inf_h\int_0^T\|\rho+Oe^{\Omega\theta}h\|^2d\theta>0.}
 \label{v50:eq:distance-positive}
\end{equation}
No numerical lower enclosure of $\epsilon_T$ is supplied here.
\end{proposition}
\begin{proof}
If $h\in\ker\mathcal W_T$, then $Oe^{\Omega\theta}h=0$ almost
everywhere, so $h^Tb_T=0$.  Thus $b_T\in\Ran\mathcal W_T$, and
completion of the square gives \eqref{v50:eq:distance-positive} with
``$\ge0$''.  The range of the map $h\mapsto Oe^{\Omega\theta}h$ in
$L^2(0,T)$ is finite-dimensional and hence closed.  Zero distance would
therefore give an exact vanishing output, first almost everywhere and
then everywhere by analyticity.  The preceding lemma excludes it.
\end{proof}

\section{What bounded original histories must satisfy on the next time scale}
\label{v50:sec:native}

We state the entrance assumptions explicitly.  They are hypotheses for a
\emph{necessary limit}, not new existence claims.  Consider original,
exactly constrained histories on $0\le t\le aT$ in the inherited analytic
stationary chart, with mean lapse one and
\begin{equation}
 \begin{aligned}
 X_a(0)&=aY_*+O(a^2),&
 \lambda_a(0)&=e+a\ell_*+O(a^2),\\
 \sup_{t\le aT}\|\lambda_{t,a}(t)\|&\le M_Ta.&
 \end{aligned}
 \label{v50:eq:entrance}
\end{equation}
All norms here may be any fixed equivalent finite-cutoff norms.  Decompose
\begin{equation}
 a^{-1}\lambda_{t,a}(a\theta)
       =I_Av_a(\theta)+I_gw_a(\theta)+I_cc_a(\theta).
 \label{v50:eq:rates}
\end{equation}
The harmonic coefficient $c_a$ is retained and bounded; it is not set to
zero.  Suppose first that $v_a(0)\to v_0$.

\begin{theorem}[A necessary differential--algebraic rate limit]
\label[theorem]{v50:thm:limit}
Under \eqref{v50:eq:entrance}, the active rates converge uniformly and the
nonconstant weak rates converge weakly in $L^2$ and weak-star in $L^\infty$,
respectively, to the unique solution of
\begin{equation}
 \boxed{
 A\dot v+t_A+K_Av+Cw=0,\qquad
 t_g-C^Tv+K_gw=0,
 \qquad v(0)=v_0.}
 \label{v50:eq:dae}
\end{equation}
The primitives of $w_a$ converge uniformly to the primitive of $w$.
An initial compatibility condition specified below is necessary.
No strong convergence of $w_a$, convergence of its time derivative, or
leading equation for the three harmonic rates is asserted.
\end{theorem}
\begin{proof}
Bounded rates and the fixed-cutoff analytic field imply, uniformly on the
interval,
\[
 X_a(a\theta)-X_a(0)=a^2\theta V+O_T(a^3),\qquad
 \lambda_a(a\theta)-\lambda_a(0)=O_T(a^2).
\]
Indeed $F(0,e)=0$, so a short-time Gronwall bound first gives $X_a=O(a)$;
Taylor expansion then gives $F=aV+O_T(a^2)$.

It is necessary to eliminate the linear constraint drift before comparing
cubic coefficients.  The inherited flat identity is
$\mathsf B_1=\mathsf J_0\mathcal L_{\rm con}$.  Define the analytic
covector
\[
                 \widetilde{\mathsf B}=\mathsf B-\mathsf J_0P.
\]
It agrees with $\mathsf B$ on every exact history, has no linear term,
and its quadratic term is exactly the prepared $B_2(Y,\ell)$ used in
\eqref{v50:eq:t}.  Since $B_2(Y_*,\ell_*)=0$, Taylor expansion between
the initial and current records yields
\begin{equation}
 \frac{\mathsf B(a\theta)-\mathsf B(0)}{a^3}
 =\theta\mathfrak t
 +KI_A\int_0^\theta v_a(s)\,ds
 +KI_g\int_0^\theta w_a(s)\,ds+O_T(a).
 \label{v50:eq:primitive-source}
\end{equation}
The harmonic integral is annihilated by $KI_c=0$.  Unknown bounded second
initial coefficients affect only the stated remainder.  Without the exact
constraint subtraction the linear drift on the $O(a^3)$ canonical
remainder would have been incorrectly discarded.

On these same histories the original graded Hessian satisfies
\[
 M_{AA}=a^2(A+O_T(a)),\quad
 M_{AW}=O_T(a^3),\quad M_{WW}=O_T(a^4).
\]
The exact consistency equation therefore gives
$\mathsf B_A/a^3=-Av_a+O_T(a)$ and
$\mathsf B_g/a^3=O_T(a)$.  Restricting
\eqref{v50:eq:primitive-source} gives the two uniform integral relations
\begin{align}
 A(v_a(\theta)-v_a(0))
  +\theta t_A+K_A\int_0^\theta v_a
  +C\int_0^\theta w_a&=O_T(a),\label{v50:eq:active-primitive}\\
 \theta t_g-C^T\int_0^\theta v_a
                     +K_g\int_0^\theta w_a&=O_T(a).
 \label{v50:eq:weak-primitive}
\end{align}

The two primitives are uniformly Lipschitz.  Extract uniformly convergent
subsequences and a weak-star convergent subsequence of $w_a$.
Equation \eqref{v50:eq:active-primitive} then gives uniform convergence of
$v_a$ itself.  Differentiating the limiting integral equations in
distributions yields \eqref{v50:eq:dae}.  The algebraic reduction in the
next theorem makes the limiting solution unique and smooth, so every
subsequence has the same limits.  This proves the full convergence claims.
The harmonic coefficients did not enter the limiting equations; no
conclusion about their convergence follows from this argument.
\end{proof}

\section{One differentiation closes a positive 76-dimensional transport}
\label{v50:sec:reduction}

The inherited exact weak rank is $24$.  Choose any full-rank
$Z\in\R^{26\times2}$ with $\Ran Z=\ker K_g$ and set
\begin{equation}
 T=CZ,\quad G=T^TA^{-1}T,\quad R_g=K_g^\dagger,\quad
 \mathcal P=A^{-1}-A^{-1}TG^{-1}T^TA^{-1}.
 \label{v50:eq:projector}
\end{equation}
The rank-two coupling theorem makes $T$ injective and $G\succ0$.
Replacing the inherited $D$-orthogonal kernel representatives by these
mean-zero representatives adds only constant shifts, whose first Poisson
columns vanish.  No physical mode is removed in this change of coordinates.
The finite Moore--Penrose inverse obeys $R_g^T=-R_g$.

\begin{theorem}[Exact elimination of all nonconstant weak rates]
\label[theorem]{v50:thm:reduction}
Initial data for \eqref{v50:eq:dae} must lie on the affine space
\begin{equation}
                   \Sigma=\{v:T^Tv=Z^Tt_g\},\qquad\dim\Sigma=76.
 \label{v50:eq:Sigma}
\end{equation}
For $v\in\Sigma$ define
\begin{align}
 w_0(v)&=-R_g(t_g-C^Tv),\notag\\
 d_A(v)&=t_A+K_Av+Cw_0(v),\notag\\
 w(v)&=w_0(v)-ZG^{-1}T^TA^{-1}d_A(v).
 \label{v50:eq:w}
\end{align}
Then the original differential--algebraic system is equivalent to
\begin{equation}
                         \boxed{\dot v=-\mathcal P d_A(v),\qquad
                                       w=w(v),\qquad v(0)\in\Sigma.}
 \label{v50:eq:ODE}
\end{equation}
This affine ODE has a unique solution for every finite $\theta$.
For any two solutions its difference satisfies the exact positive identity
\begin{equation}
             \boxed{\frac{d}{d\theta}
                    \frac12(v_1-v_2)^TA(v_1-v_2)=0.}
 \label{v50:eq:energy}
\end{equation}
This is a property of the limiting constrained transport, not a
stability theorem for the full singularly perturbed original histories.
\end{theorem}
\begin{proof}
Pairing the algebraic equation with $Z$ gives \eqref{v50:eq:Sigma}.
Since a skew matrix has range equal to the orthogonal complement of its
kernel, every solution of that equation is $w=w_0(v)+Z\eta$.
The active differential equation is
$A\dot v=-d_A(v)-T\eta$.  Differentiating the two compatibility rows
once gives
\[
 0=T^T\dot v=-T^TA^{-1}d_A-G\eta.
\]
The positive $G$ has an inverse, so this determines $\eta$ and proves
\eqref{v50:eq:w}--\eqref{v50:eq:ODE}.  Conversely, these formulas satisfy
every algebraic and differential row, and $T^T\mathcal P=0$ preserves
$\Sigma$.  The resulting vector field is affine and global.

In expanded form
\[
 d_A(v)=t_A-CR_gt_g+K_{\rm eff}v,\qquad
 K_{\rm eff}=K_A+CR_gC^T,
 \qquad K_{\rm eff}^T=-K_{\rm eff}.
\]
For a difference $\delta v$ of two solutions,
$T^T\delta v=0$ and
$A\delta\dot v=-K_{\rm eff}\delta v-T\delta\eta$.
Pairing with $\delta v$ proves \eqref{v50:eq:energy}.
After subtracting any fixed point of the affine space $\Sigma$, the
remaining constant forcing gives a norm bound of the form $C(1+\theta)$.
The reconstructed $w(v)$ obeys the same type of bound.
\end{proof}

There is no conflict with V48's growing pair.  That pair belongs to the
faster seven-dimensional response on $t=a^2\tau$.  Equation
\eqref{v50:eq:ODE} is what a uniformly bounded family must satisfy on
$t=a\theta$ after its leading algebraic compatibility is enforced.  The
proof has not selected away the earlier growing mode or shown that a
native history tracks this latter equation.  Establishing such a lift
requires a separate multiscale preparation and propagation theorem.

\section{A source-specific obstruction to extending weak-rate flatness}
\label{v50:sec:obstruction}

\begin{theorem}[No bounded-rate weakly flat family at the next time scale]
\label[theorem]{v50:thm:no-flat}
For any fixed $T>0$, no family of original histories can satisfy
\eqref{v50:eq:entrance} and
\begin{equation}
                 \sup_{0\le t\le aT}
                    a^{-1}\|\lambda_{t,g}(t)\|\longrightarrow0.
 \label{v50:eq:weak-flat}
\end{equation}
Here $\lambda_{t,g}$ is the mean-zero curl-free shift-rate projection.
Harmonic rates of order $a$ are allowed.  The statement covers every
bounded initial active-rate choice and every bounded higher preparation
with the fixed $Y_*,\ell_*$, not only the V49 selection.
\end{theorem}
\begin{proof}
Bounded initial active rates have a convergent subsequence.  The necessary
limit theorem and \eqref{v50:eq:weak-flat} give \eqref{v50:eq:dae} with
$w=0$.  This is precisely \eqref{v50:eq:flat-law}, which is excluded by
the certified $\rho\ne0$ and \cref{v50:lem:observability}.
\end{proof}

The obstruction is also visible in an integrated positive rate norm.
Let $M_\Omega$ bound $\|e^{\Omega\theta}\|$ for $\theta\ge0$; it is
finite because $\Omega^TA+A\Omega=0$ and $A\succ0$.  Put
\[
 C_T=\|K_g\|+T\|C^T\|M_\Omega\|A^{-1}C\|>0.
\]

\begin{proposition}[A nonzero weak-rate floor for every bounded-rate lift]
\label[proposition]{v50:prop:floor}
Every solution of \eqref{v50:eq:dae} obeys
\begin{equation}
                     \int_0^T\|w(\theta)\|^2d\theta
                     \ge\frac{\epsilon_T}{C_T^2}>0.
 \label{v50:eq:floor}
\end{equation}
Consequently any original family satisfying \eqref{v50:eq:entrance}
has a positive lower limit for
\begin{equation}
             a^{-3}\int_0^{aT}\|\lambda_{t,g}(t)\|_h^2\,dt.
 \label{v50:eq:physical-floor}
\end{equation}
The lower constant is not numerically enclosed in this body.
\end{proposition}
\begin{proof}
Let $v^{\rm f}$ solve the free active equation with the same initial value.
Then
\[
 v-v^{\rm f}=-\int_0^\theta
 e^{\Omega(\theta-s)}A^{-1}Cw(s)\,ds.
\]
The algebraic equation gives
\[
 t_g-C^Tv^{\rm f}
 =-K_gw-C^T\int_0^\theta
 e^{\Omega(\theta-s)}A^{-1}Cw(s)\,ds.
\]
Young's convolution inequality bounds its $L^2$ norm by $C_T\|w\|_{L^2}$.
The output-distance proposition bounds that same norm below by
$\sqrt{\epsilon_T}$, proving \eqref{v50:eq:floor}.  For original
histories, take a subsequence realizing the lower limit and then a
convergent initial active subsequence.  Weak $L^2$ lower semicontinuity
and \cref{v50:thm:limit} apply.  The fixed positive Gram
$I_g^TI_g/27$ converts coordinate rate norms to the physical pairing.
Finally $t=a\theta$ and $\lambda_{t,g}=aI_gw_a$ give the factor $a^3$.
\end{proof}

This excludes persistence of the particular rate-flat strategy.  It does
not exclude a bounded original history with a nonzero order-$a$ weak rate,
a change of leading first fields, a different observation which does not
require this rate to vanish, or a different unbounded-control scaling.
An order-$a$ rate still vanishes without amplitude normalization.  No
nonexistence of classical Einstein fields, metric singularity, or
failure of a renormalized quantum continuum follows from this result.

\section{The V49 selection starts flat but the necessary transport leaves it}
\label{v50:sec:selected}

For V49's certified first active rate $\bar v_A$, let
\begin{equation}
                 \gamma_A=-A^{-1}a_A(\bar v),\qquad
                 b_g=-C^T\gamma_A.
 \label{v50:eq:selected}
\end{equation}
The exact V49 identities give
\[
 t_g-C^T\bar v_A=0,\qquad
 T^T\gamma_A=0.
\]
The second follows from its canceled two-component source $d_v$; it is
not inferred from a small interval enclosure.

\begin{proposition}[An evaluated nonzero first departure of the limiting weak rate]
\label[proposition]{v50:prop:selected}
The unique constrained transport with $v(0)=\bar v_A$ satisfies
\begin{equation}
               w(0)=0,\qquad \dot v(0)=\gamma_A,\qquad
               K_g\dot w(0)=-b_g\ne0.
 \label{v50:eq:departure}
\end{equation}
For coordinate $20$ in the same zero-based weak order, the exact-root
certificate proves
\begin{equation}
                 \boxed{-720000000<(b_g)_{20}<-340000000.}
 \label{v50:eq:departure-interval}
\end{equation}
Thus the conditional a-scale rate limit of the selected family is not
identically weakly flat, even though it starts with zero weak rate.
\end{proposition}
\begin{proof}
At $\bar v_A$, $w_0=0$ and $d_A=a_A(\bar v)$.  The neutral projection
term in \eqref{v50:eq:w} vanishes by $T^TA^{-1}a_A(\bar v)=0$.
This proves the initial values and active derivative.  Differentiate the
algebraic equation to obtain $K_g\dot w=C^T\dot v$, yielding
\eqref{v50:eq:departure}.

The new verifier forms $a_A(\bar v)$ by adding the complete $K$ contraction
of the certified V49 rate change to the parent active source.  Propagation
through the original positive active inverse gives a centre for
$(b_g)_{20}$ of approximately $-531704506.4762203$ and a conservative
entry radius below $185137262$.  Exact rational endpoint comparisons
prove \eqref{v50:eq:departure-interval}.  The relatively wide interval
retains the parent's rate uncertainty and is sufficient to exclude zero;
no uncomputed second active-rate coefficient enters this contraction.
\end{proof}

The derivative in \eqref{v50:eq:departure} belongs to the limiting
constrained equation.  The convergence theorem did not prove convergence
of derivatives of the original weak rates at $\theta=0$, so it must not
be read as a new initial-acceleration asymptotic for the native history.
Likewise the weak rate is not the full metric curvature: active rate
changes can contribute to the latter on this scale.  No curvature lower
bound is claimed here.

\section{Verification ledger and the noncircular next step}
\label{v50:sec:status}

\paragraph{Proved in this body.}
The original active Poisson block is uniformly invertible on the root
cube.  Its equilibrium source fails an explicit original weak row,
with a strict rational separation.  This excludes every active-only
persistent rate solution at the fixed leading fields, not just one
selected rate.  Any bounded-rate original lift to $t=a\theta$ must
instead satisfy the closed constrained transport above.  That limit is
well posed on a 76-dimensional affine space and has conserved positive
homogeneous difference energy.  It forces a nonzero integrated
nonconstant weak rate.  For the V49 initial selection, an additional
original-source contraction proves a nonzero first departure of that
limiting weak rate.

\paragraph{Inherited and executed.}
The default runner replays the recovered V49 package, including its
108 differentiated drift entries and three cubic-mean derivatives.
The new stage assembles the active and cross Poisson blocks from the
retained polynomial bank, certifies a new inverse on the full root
cube, and propagates the source and rate enclosures to the two reported
nonzero tests.  It checks constant-shift annihilation coefficientwise.
The old 79,224-entry root bank, 17,334-entry stationary-source bank and
5,046-entry weak-bracket bank are not claimed as fully regenerated.
The accompanying small rational examples check the projector identities,
all retained descriptor rows, positive difference energy, the dual
annihilator, Cayley--Hamilton closure, the finite-output Gram formula,
and physical rate scaling.  They are not native spacetime simulations.

\paragraph{Unpaid and next direction.}
Existence and uniform rate bounds for the actual source on $t\le aT$ have
not been proved.  The seven-dimensional growing mode on $t=a^2\tau$ is
still present; controlling its preparation and propagation is necessary
before the positive limiting equation can be lifted to native histories.
Neither its three harmonic rates nor a curvature lower bound is supplied
by this leading reduction.  The finite-window constants in
\eqref{v50:eq:floor} and the nonlinear invariant graph are not numerically
enclosed.  There is no common positive physical lifespan, cofinal spatial
estimate, primitive selection theorem or matter-coupled Einstein--SM
closure.

The next useful task is a well-prepared lift of the \emph{nonzero} weak
rate transport, or a change of leading fields which passes the explicit
source mismatch test.  Canceling more initial weak coefficients while
requiring $o(a)$ nonconstant weak rate throughout $t=a\theta$ cannot work
within the bounded-rate, fixed-leading-field regime proved above.  The
V49 theorem on fixed $Ta^2$ intervals is untouched, as is V32's separate
unit-initial-multiplier compensation problem.

Every Version 49 byte before its final document command is preserved.
The numerical source statements are exact-arithmetic inclusion
certificates; the limit and persistence results are the written proofs.
No proof-assistant formalization is asserted.  New labels have prefix
\texttt{v50:}.

\begingroup
\renewcommand{\refname}{Additional references for Version 50}
\begin{thebibliography}{9}
\small\raggedright
\bibitem{v50:Hautus}
M.~L.~J.~Hautus, \emph{Controllability and observability conditions of linear
autonomous systems}, Indagationes Mathematicae (Proceedings)
\textbf{72}(5) (1969), 443--448.
Cited for the established finite-dimensional observability framework;
the output-distance and native-source consequences used here are proved
explicitly above.
\bibitem{v50:Rump}
S.~M.~Rump, \emph{Verification methods: Rigorous results using floating-point
arithmetic}, Acta Numerica \textbf{19} (2010), 287--449,
\url{https://doi.org/10.1017/S096249291000005X}.
Cited for residual verification methodology, not as a source of the new
matrix values or physical conclusions.
\end{thebibliography}
\endgroup
\addtocontents{toc}{\protect\endgroup}
% END IMMUTABLE VERSION 50 BODY
% APPEND VERSION 51 HERE, BEFORE THE FINAL END-DOCUMENT COMMAND.



% BEGIN IMMUTABLE VERSION 51 BODY
\clearpage
\begin{center}
{\Large\bfseries Version 51: The harmonic rows obstruct the proposed lift}\\[.5em]
{\large Full Poisson reduction, a certified cubic source mismatch,\\
and a no-go for the bounded-rate $t=a\theta$ branch}
\end{center}
\makeatletter
\addtocontents{toc}{\protect\begingroup}
\addtocontents{toc}{\protect\makeatletter}
\addtocontents{toc}{\protect\def\protect\l@section{\protect\bfseries\protect\@dottedtocline{1}{0em}{2.5em}}}
\makeatother
\addcontentsline{toc}{part}{Version 51: Harmonic completion of the next rate limit}

\section{The missing original rows before a singular-perturbation lift}
\label{v51:sec:context}

Version 50 derives a necessary nonconstant rate transport on $t=a\theta$,
conditional on original histories with uniformly $O(a)$ multiplier rates.
It explicitly leaves the three harmonic rates without a leading equation.
Before attempting a lift of that transport, we return to the three original
constant-shift consistency rows. They occur one amplitude order later than
the nonconstant rows but remain necessary under the very same bounded-rate
hypotheses. The calculation below retains them rather than assigning the
harmonic rates by a gauge convention.

There are two new conclusions. First, maximal rank of the complete first
Poisson matrix implies a nondegenerate positive-energy reduction: V50's
76-dimensional affine transport has a unique equilibrium and bounded
oscillations about it, with no zero frequency. Second, the complete cubic
constant-shift source is incompatible with that transport, for every
initial rate, at the actual mixed-mode root. This excludes the proposed
bounded-rate lift at these leading records. It does not invalidate the
shorter exactly constructed histories or exclude histories with a different
rate scaling or different leading fields.

Throughout, $h=1/3$ and $(Y_*,\ell_*)$ is the exact V39 algebraic root.
The V50 verifier has been rerun, including its parent regeneration of all
108 differentiated quadratic-drift entries and the three complete cubic-mean
derivatives. The old root, stationary-source and weak-bracket banks are
inherited with their stored enclosures. The new calculation uses one full
104-dimensional original Poisson inverse and three new source contractions.
The method of structure-preserving algebraic elimination is established;
we use no external existence theorem for the desired native lift
\cite{v51:vdSM}. The new source values and harmonic-limit calculation are
proved here.

\section{Full Poisson rank removes the apparent secular freedom}
\label{v51:sec:reduction}

Retain all notation of \cref{v50:eq:blocks,v50:eq:projector}. Put
\begin{equation}
 \mathbb K=\begin{pmatrix}K_A&C\\-C^T&K_g\end{pmatrix},\qquad
 \mathbb t=\binom{t_A}{t_g},\qquad
 \mathbb I=(I_A,I_g).
 \label{v51:eq:full}
\end{equation}
The 104 columns of $\mathbb I$ span the componentwise mean-zero multiplier
space. The original 108-dimensional first Poisson matrix has rank 104 and
kernel exactly the four component constants at the root. Therefore
$\mathbb K=\mathbb I^TK\mathbb I$ is nonsingular. This statement uses the
\emph{complete} first Poisson matrix, not the symmetric multiplier Hessian.

Let $Z$ synthesize $\ker K_g$, $T=CZ$, $R_g=K_g^\dagger$, and
\[
 \mathcal S=\ker T^T,\qquad \dim\mathcal S=76,\qquad
 K_{\rm eff}=K_A+CR_gC^T.
\]
Choose any fixed basis matrix $E$ for $\mathcal S$, and write
\begin{equation}
 A_s=E^TAE\succ0,\qquad K_s=E^TK_{\rm eff}E.
 \label{v51:eq:reduced}
\end{equation}
The choice of $E$ affects coordinate norms but not the following conclusions.

\begin{lemma}[Full Poisson nondegeneracy is equivalent to reduced nondegeneracy]
\label[lemma]{v51:lem:nondegenerate}
Assume $T$ has full column rank. Then $\mathbb K$ is nonsingular if and
only if $K_s$ is nonsingular. In particular $K_s$ is nonsingular for the
actual root.
\end{lemma}
\begin{proof}
If $x\in\mathcal S$ and $E^TK_{\rm eff}x=0$, then
$K_{\rm eff}x\in\Ran T$. Choose $\eta$ with
$K_{\rm eff}x+T\eta=0$ and set
$y=R_gC^Tx+Z\eta$. Since $Z^TC^Tx=T^Tx=0$, one has
$K_gy=C^Tx$, and hence $\mathbb K(x,y)=0$. Nonsingularity of
$\mathbb K$ forces $x=y=0$.
Conversely, $\mathbb K(x,y)=0$ gives $T^Tx=0$ and
$y=R_gC^Tx+Z\eta$. The active equation gives
$E^TK_{\rm eff}x=0$. If $K_s$ is nonsingular, then $x=0$ and
$T\eta=0$, so $\eta=0$ and $y=0$.
\end{proof}

\begin{theorem}[Unique rate equilibrium and 38 nonzero oscillatory pairs]
\label[theorem]{v51:thm:oscillatory}
The V50 constrained transport has the unique equilibrium
\begin{equation}
             \boxed{\binom{v_\infty}{w_\infty}=-\mathbb K^{-1}\mathbb t.}
 \label{v51:eq:equilibrium}
\end{equation}
For every compatible initial value, its complete solution has the form
\begin{equation}
 \begin{split}
 v(\theta)&=v_\infty+E e^{\Omega_s\theta}q_0,\\
 w(\theta)&=w_\infty+W_s e^{\Omega_s\theta}q_0,\\
 \Omega_s&=-A_s^{-1}K_s,
 \end{split}
 \label{v51:eq:all-solutions}
\end{equation}
where $W_s$ is the linear reconstruction from \cref{v50:eq:w} on
$E\mathbb R^{76}$. The generator $\Omega_s$ is invertible and has 38
pairs $\pm i\omega_j$, $\omega_j>0$, counted with multiplicity; all
its eigenvalues are semisimple. The exact positive energy is
\begin{equation}
 \frac12(v-v_\infty)^TA(v-v_\infty)
              =\frac12q_0^TA_sq_0.
 \label{v51:eq:energy}
\end{equation}
Both rates are bounded for all real $\theta$, with a bound depending on
the initial value. Their Cesaro means converge to $(v_\infty,w_\infty)$
with error $O_{q_0}(T^{-1})$.
\end{theorem}
\begin{proof}
The stationary equations are exactly $\mathbb K(v,w)=-\mathbb t$.
Their weak projection implies compatibility, so the unique solution is an
equilibrium of the reduced ODE. Subtract it and use V50's elimination:
$A_s\dot q+K_sq=0$. The matrix
$A_s^{-1/2}K_sA_s^{-1/2}$ is real skew and nonsingular. Orthogonal
block diagonalization gives the stated spectrum and the energy identity.
The weak reconstruction is affine, giving \eqref{v51:eq:all-solutions}.
Finally
\[
 \int_0^T e^{\Omega_s\theta}\,d\theta
       =\Omega_s^{-1}(e^{\Omega_sT}-I)
\]
is uniformly bounded in the $A_s$ norm. This proves the mean estimate.
No simplicity of the individual frequencies is claimed.
\end{proof}

The equilibrium here is a \emph{rate equilibrium of the necessary limiting
system}, not a stationary spacetime or an attracting state of the original
flow. The bounded oscillatory result strengthens V50's at-most-linear
estimate, but by itself supplies no native trajectory. The harmonic rows
below decide whether that limiting transport is admissible at all.

\section{A certified original cubic mismatch at the unique rate equilibrium}
\label{v51:sec:certificate}

Let $I_c$ synthesize the three original constant shifts, and set
\begin{equation}
 L_c=I_A^TM^{[3]}I_c=L N_c,\qquad
 \mathscr R_A=\mathscr R_{Y_*}[\,\cdot\,,I_A],\quad
 \mathscr R_g=\mathscr R_{Y_*}[\,\cdot\,,I_g].
 \label{v51:eq:constant-blocks}
\end{equation}
Here the first equality denotes the retained coefficient
$M_{Ac}=a^3(L_c+O(a))$; the second expression uses the original weak
coordinates of V49. In particular $L_c$ is not a first-Poisson block.
Use the complete prepared cubic mean functional from
\cref{v39:eq:mean-functional,v49:eq:envelopes}, and define
\begin{equation}
 s_c=D_Y C(Y_*,\ell_*)[V],\qquad
 \rho_c=s_c+\mathscr R_Av_\infty+\mathscr R_gw_\infty.
 \label{v51:eq:harmonic-mismatch}
\end{equation}
The last preparation term in $\mathscr R$ is retained. The exact identity
\begin{equation}
                 \mathscr R_{Y_*}[\,\cdot\,,I_c]=0
 \label{v51:eq:Rc0}
\end{equation}
follows from the commuting quadratic translation generators as in
\cref{v48:prop:cross}. The verifier also checks every coefficient column
of this identity in the inherited 21-row polynomial bank.

\begin{theorem}[Full inverse and three nonzero harmonic source components]
\label[theorem]{v51:thm:source}
On the retained root cube of radius $2^{-112}$,
\begin{equation}
             \boxed{\|\mathbb K^{-1}\|_\infty<\frac23.}
 \label{v51:eq:full-inverse-bound}
\end{equation}
The nonconstant weak equilibrium has physical norm
\begin{equation}
                \boxed{2290<\|I_gw_\infty\|_h<2297.}
 \label{v51:eq:weak-norm}
\end{equation}
In the original constant-shift order, the cubic mismatch satisfies
\begin{equation}
 \boxed{
 \begin{aligned}
 -10665000&<(\rho_c)_1<-10662000,\\
 -27897000&<(\rho_c)_2<-27894000,\\
  12175000&<(\rho_c)_3< 12179000.
 \end{aligned}}
 \label{v51:eq:harmonic-intervals}
\end{equation}
In particular $\rho_c\ne0$. These use the common raw spatial-sum
covector normalization, not normalized continuum source densities.
\end{theorem}
\begin{proof}
The worker reconstructs the physical first Poisson matrix with its
$\sqrt3$ grading and forms $\mathbb K=\mathbb I^TK\mathbb I$.
The active synthesis is the same exact rational orthogonal projection
used in V49; it is not asserted orthonormal. All four component means
of $\mathbb I$ vanish exactly.

Let $P$ be the stored rational inverse proposal and write the uniform
entry enclosure as $\mathbb K\in\mathbb K_0+[-r_K,r_K]$. Exact
rational arithmetic verifies
\begin{equation}
 q=\|I-P\mathbb K_0\|_\infty+104\|P\|_\infty r_K<2^{-84},\qquad
 \frac{\|P\|_\infty}{1-q}<\frac23.
 \label{v51:eq:Neumann}
\end{equation}
The actual residual bound is below $3.244\,10^{-26}$ and the inverse
bound below $0.634$. The proposal was generated from floating arithmetic,
but only \eqref{v51:eq:Neumann} certifies it. Uniform-entry rational
interval multiplication then encloses \eqref{v51:eq:equilibrium}.
The 104 coordinate errors are below $0.017587$. The physical weak norm
is bounded by rational squaring and the exact positive $1/27$ pairing;
its resulting interval is contained in $(2290,2297)$.

The parent source regeneration supplies
$\mathfrak t+Kv_*$ and $s_c+\mathscr Rv_*$ with their full entry errors.
Subtracting the enclosed $Kv_*$ and $\mathscr Rv_*$ recovers the fixed-rate
sources $\mathfrak t,s_c$; no unknown higher initial coefficient is set
to zero. The three new contractions in \eqref{v51:eq:harmonic-mismatch}
have rational centres approximately
\[
 (-10663319.14095504,\ -27895554.93859927,\ 12176993.02278857)
\]
and a common entry radius below $1016$. Exact endpoint comparisons
prove \eqref{v51:eq:harmonic-intervals}. Every reported value encloses
the algebraic root, not just its printed approximation.
\end{proof}

The constant-rate extension of \eqref{v51:eq:equilibrium} is not unique
if harmonic coordinates are appended. Nevertheless both $KI_c=0$ and
\eqref{v51:eq:Rc0} hold, so no such harmonic addition changes $\rho_c$.
The next section shows why bounded, time-dependent harmonic rates cannot
change this leading compatibility test either.

\section{The original harmonic consistency equation on the same time scale}
\label{v51:sec:native}

Consider exactly the entrance hypotheses \eqref{v50:eq:entrance}.
In particular all multiplier rates, including harmonic rates, are uniformly
$O(a)$ on $0\le t\le aT$. Write their decomposition as in
\eqref{v50:eq:rates}. Set
\[
 V_a(\theta)=\int_0^\theta v_a(u)\,du,\qquad
 W_a(\theta)=\int_0^\theta w_a(u)\,du.
\]
No derivative bound on $v_a,w_a,c_a$ is assumed.

\begin{lemma}[Uniform prepared cubic expansion of the constant drift]
\label[lemma]{v51:lem:cubic-primitive}
For those original histories,
\begin{equation}
 \frac{\mathsf B_c(a\theta)-\mathsf B_c(0)}{a^4}
       =\theta s_c+\mathscr R_AV_a(\theta)
                      +\mathscr R_gW_a(\theta)+O_T(a)
 \label{v51:eq:cubic-primitive}
\end{equation}
uniformly for $0\le\theta\le T$. The coefficients are exactly those
of \eqref{v51:eq:constant-blocks}--\eqref{v51:eq:harmonic-mismatch}.
\end{lemma}
\begin{proof}
Use the invariant analytic nonconstant-constraint chart of
\cref{v46:lem:constraint-chart}. Its free canonical coordinates have
leading field in $\ker\mathcal L_{\rm con}$. On that exact chart the
constant-shift drift has no linear or prepared quadratic term; its
homogeneous cubic term is the complete function $C(Y,\ell)$ in
\eqref{v49:eq:envelopes}. This is the prepared coefficient identity,
not the unprepared $\{P_{c,2},H_3\}+\{P_{c,3},H_2\}$ alone.

For completeness, the potentially lower mixed term contains the linear
constraint acting on the second canonical coefficient. Substituting the
exact constraint row replaces that term by
$-\langle\delta_c\ell,P_2(Y)\rangle$, which is precisely the last
term of \eqref{v39:eq:mean-functional}. The commutation of quadratic
translations with the flat quadratic Hamiltonian removes the pure
quadratic drift. Analyticity on the constraint chart leaves an analytic
remainder of joint degree at least four. This coefficient statement is
valid for all nearby chart coordinates, not only at a chosen initializer.

The bounded-rate hypothesis and V50's short-time estimate give
\[
 X_a(a\theta)-X_a(0)=a^2\theta V+O_T(a^3),\quad
 \lambda_a(a\theta)-\lambda_a(0)
   =a^2\{I_AV_a+I_gW_a+I_c\textstyle\int_0^\theta c_a\}.
\]
The leading tangent $V$ is flat-constraint compatible. Taylor expansion
of the cubic chart coefficient between these two records therefore gives
\eqref{v51:eq:cubic-primitive}; the constant-shift integral is annihilated
by \eqref{v51:eq:Rc0}. Initial $O(a^2)$ coefficients change the cubic
derivative only by $O(a)$ after scaling. The difference of the analytic
fourth-order remainder is $O_T(a^5)$, since the record displacement is
$O_T(a^2)$. Uniform constants use only the stated bounded leading and
second initial coefficients and the fixed chart. No hidden derivative
bound on the rates is used.
\end{proof}

\begin{theorem}[A necessary harmonic completion of V50's transport]
\label[theorem]{v51:thm:harmonic-limit}
Under \eqref{v50:eq:entrance}, the limiting V50 rates must also satisfy
\begin{equation}
                 \boxed{L_c^T\dot v+s_c+
                         \mathscr R_Av+\mathscr R_gw=0.}
 \label{v51:eq:harmonic-law}
\end{equation}
This is an additional necessary equation, not a determination of the
three leading harmonic rates.
\end{theorem}
\begin{proof}
The literal constant row of $M\lambda_t=-\mathsf B$ has
$M_{cA}=a^3(L_c^T+O_T(a))$ and $M_{cW}=O_T(a^4)$.
Since every rate is $O(a)$,
\[
                  \mathsf B_c(a\theta)/a^4=-L_c^Tv_a(\theta)+O_T(a).
\]
Subtract the initial value and use \cref{v51:lem:cubic-primitive}:
\begin{equation}
 L_c^T\{v_a(\theta)-v_a(0)\}
 +\theta s_c+\mathscr R_AV_a(\theta)+\mathscr R_gW_a(\theta)
                       =O_T(a).
 \label{v51:eq:harmonic-primitive}
\end{equation}
V50 proves uniform convergence of the active rates and both primitives
along a convergent initial active subsequence. Passing to the limit in
\eqref{v51:eq:harmonic-primitive} and differentiating that limiting
integral identity proves \eqref{v51:eq:harmonic-law}. It is smooth
because the reduced transport is a finite affine ODE. No convergence
of original rate derivatives is inferred.
\end{proof}

The harmonic rate term in the literal constant row is one order smaller
under the entrance hypothesis. Omitting it here is a derived order
estimate, not an assertion of harmonic gauge symmetry. A different rate
scaling can change this conclusion and must retain that term explicitly.

\section{No initial active rate satisfies the completed transport}
\label{v51:sec:obstruction}

The new harmonic equation can be evaluated on every solution of the
nonconstant transport, without solving another time-dependent inverse.
Define the fixed output matrix
\begin{equation}
 \mathcal O_c=L_c^TE\Omega_s+\mathscr R_AE+\mathscr R_gW_s.
 \label{v51:eq:output}
\end{equation}
Inserting \eqref{v51:eq:all-solutions} in its left side gives
\begin{equation}
               \rho_c+\mathcal O_c e^{\Omega_s\theta}q_0.
 \label{v51:eq:output-trajectory}
\end{equation}

\begin{proposition}[Exact persistence criterion for the harmonic rows]
\label[proposition]{v51:prop:persistence}
The V50 transport and \eqref{v51:eq:harmonic-law} have a common solution
on a nonempty time interval if and only if $\rho_c=0$.
For the actual source there is no such solution, for any finite compatible
initial rate.
\end{proposition}
\begin{proof}
If \eqref{v51:eq:output-trajectory} vanishes on an interval, its entire
analytic continuation vanishes on the real line. Differentiate $k\ge1$
times to obtain $\mathcal O_c\Omega_s^kq_0=0$. Since $\Omega_s$ is
invertible, Cayley--Hamilton expresses $I$ through its first 76 positive
powers. Thus $\mathcal O_cq_0=0$, and the original equation gives
$\rho_c=0$. Equivalently, its Cesaro mean is $\rho_c$ by
\cref{v51:thm:oscillatory}. Conversely $q_0=0$ supplies a common
solution when $\rho_c=0$. Apply \eqref{v51:eq:harmonic-intervals}.
\end{proof}

\begin{theorem}[No uniformly amplitude-sized rate lift at the fixed leading root]
\label[theorem]{v51:thm:no-lift}
For any fixed $T>0$, no family of original exactly constrained histories
satisfies all the entrance hypotheses \eqref{v50:eq:entrance} on
$0\le t\le aT$ at the same $Y_*,\ell_*$. In particular the obstruction
now includes nonzero nonconstant weak rates and arbitrary bounded
harmonic rate coefficients, not only V50's weakly flat subcase.
\end{theorem}
\begin{proof}
If such a family existed, its bounded initial active rates would have a
convergent subsequence. V50 gives its unique smooth nonconstant transport
limit. \Cref{v51:thm:harmonic-limit} gives the three additional harmonic
equations for that same limit. This contradicts
\cref{v51:prop:persistence}. Every original mean row was retained; no
choice of unobserved harmonic rate within the assumed $O(a)$ bound can
supply the missing leading term.
\end{proof}

Thus, if original histories with the stated leading initial expansions
exist throughout $[0,aT]$ for a cofinal amplitude sequence within this
stationary branch, their full multiplier rates cannot have a bounded
ratio to $a$ along any subsequence. In that conditional sense
\[
                 a^{-1}\sup_{t\le aT}\|\lambda_t(t)\|\longrightarrow\infty.
\]
This is not a numerical lower power for the rate, does not identify which
component grows, and does not prove that a history reaches $aT$.
The alternatives include leaving the branch or a different initial or
rate scaling.

\section{A finite output-distance certificate and its limits}
\label{v51:sec:distance}

For $T>0$, put
\begin{equation}
 \begin{split}
 \mathcal W_T^c&=\int_0^T e^{\Omega_s^T\theta}
                         \mathcal O_c^T\mathcal O_c e^{\Omega_s\theta}\,d\theta,\\
 b_T^c&=\int_0^T e^{\Omega_s^T\theta}\mathcal O_c^T\rho_c\,d\theta,\\
 \epsilon_T^c&=T\|\rho_c\|^2-(b_T^c)^T(\mathcal W_T^c)^\dagger b_T^c.
 \end{split}
 \label{v51:eq:distance}
\end{equation}

\begin{proposition}[Positive harmonic output floor]
\label[proposition]{v51:prop:distance}
For the actual original-source coefficients and every $T>0$,
\begin{equation}
 \boxed{\epsilon_T^c
 =\inf_{q_0}\int_0^T
       \|\rho_c+\mathcal O_c e^{\Omega_s\theta}q_0\|^2d\theta>0.}
 \label{v51:eq:distance-positive}
\end{equation}
No numerical enclosure of $\epsilon_T^c$ is claimed here.
\end{proposition}
\begin{proof}
The map from $q_0$ to its output function has finite-dimensional closed
range in $L^2(0,T)$. Its Gram kernel annihilates $b_T^c$. Completion
of the square gives the pseudoinverse formula. If the minimum were
zero, the continuous output would vanish everywhere on the interval,
contradicting \cref{v51:prop:persistence}.
\end{proof}

This is a distance for the necessary rate equation. It is not a lower
bound for any curvature component, nor a prescription to add a residual
to the original action. The proof of \cref{v51:thm:no-lift} requires
only its strict positivity, not a computed value.

The positive oscillatory theorem also explains the scope of the
obstruction. The nonconstant transport, considered alone, is globally
bounded and has a nonzero mean weak rate. Its failure is not secular
blow-up of that transport. It is failure to satisfy three other original
rows which first become leading in the bounded-rate $a$-scale limit.
A well-prepared lift theorem cannot use the nonconstant energy identity
as a substitute for those rows.

\section{Verification, preserved results, and the next upstream question}
\label{v51:sec:status}

\paragraph{Proved.}
The original full mean-zero Poisson matrix is quantitatively invertible,
and its constrained 76-dimensional transport has no zero frequency,
a unique rate equilibrium and a conserved positive shifted energy.
All three cubic harmonic residuals at that equilibrium exclude zero.
Every bounded-rate original $a$-scale limit must satisfy the harmonic
completion, which is incompatible with the transport for every initial
rate. Thus the fixed-leading-root $O(a)$ total-rate lift is excluded,
not only the earlier weak-flat strategy.

\paragraph{Inherited and checked.}
A clean parent V50 run replays its dependencies and recomputes the
108 differentiated drift entries and three complete cubic derivatives.
The new worker forms and certifies a 104-dimensional original Poisson
inverse, its equilibrium, and the harmonic source contractions. It checks
all cubic constant-to-constant columns exactly in the retained polynomial
bank. The 79,224-entry root bank, 17,334-entry stationary-source bank and
5,046-entry weak-Poisson bank are inherited and replayed, not wholly
regenerated. The finite algebra regressions test nondegenerate reduction,
shifted energy, the harmonic primitive identity, observability and the
importance of retaining the harmonic rows. They are not native evolutions.

\paragraph{Unpaid and next direction.}
The shorter V49 original histories on each fixed $Ta^2$ interval remain
valid: that interval corresponds to $\theta\le Ta\to0$, so it does not
contradict a theorem for any fixed positive $\theta$ interval. The positive
reduced transport is still a mathematical equation, but it is no longer
a viable complete $O(a)$-rate limit at this leading root.

A noncircular continuation must change an actual entrance ingredient.
One possibility is to allow a faster harmonic rate and derive its complete
balance with the original metric observations; another is a different
leading field/tilt passing the new harmonic source test. The current proof
neither selects the correct new rate scale nor proves a geometric
coordinate comparison for it. In particular larger harmonic rates need
not automatically imply larger curvature. No fixed positive physical
lifespan, cofinal spatial result, native attraction, or coupled
Einstein--SM theorem follows. V32's independent compensation problem
is untouched.

Every V50 byte before its final document command is preserved. New labels
use \texttt{v51:}. Numerical bounds are rational inclusion certificates;
the limiting-harmonic argument is a written proof, not a spacetime
simulation or proof-assistant formalization.

\begingroup
\renewcommand{\refname}{Additional reference for Version 51}
\begin{thebibliography}{9}
\small\raggedright
\bibitem{v51:vdSM}
A.~van der Schaft and B.~Maschke, \emph{Generalized port-Hamiltonian DAE
systems}, Systems \& Control Letters \textbf{121} (2018), 31--37.
\url{https://arxiv.org/abs/1808.01845}.
General context for retaining algebraic constraints in conservative
systems; no source value or native existence result is imported.
\end{thebibliography}
\endgroup
\addtocontents{toc}{\protect\endgroup}
% END IMMUTABLE VERSION 51 BODY
% APPEND VERSION 52 HERE, BEFORE THE FINAL END-DOCUMENT COMMAND.



% BEGIN IMMUTABLE VERSION 52 BODY
\clearpage
\begin{center}
{\Large\bfseries Version 52: Retain the faster harmonic motion}\\[.5em]
{\large An original mixed-rate limit, one forced neutral direction,\\
and the geometric meaning of its constant-shift rate}
\end{center}
\makeatletter
\addtocontents{toc}{\protect\begingroup}
\addtocontents{toc}{\protect\makeatletter}
\addtocontents{toc}{\protect\def\protect\l@section{\protect\bfseries\protect\@dottedtocline{1}{0em}{2.5em}}}
\makeatother
\addcontentsline{toc}{part}{Version 52: Mixed harmonic rates and their original geometric readout}

\section{Change the rate entrance, not the original equations}
\label{v52:sec:context}

Version 51 excludes the total-$O(a)$ rate entrance on $t=a\theta$ at the
retained mixed-mode root. It does not exclude an $O(1)$ harmonic shift rate.
We examine that different entrance without suppressing its return into the
active or nonconstant weak equations. The leading harmonic multiplier value
can now vary by $O(a)$ on this interval; it must not be held at its initial
value while taking the limit.

All statements concern $h=1/3$ and the same exact V39 root $(Y_*,\ell_*)$.
The original constraints, physical clock, phase derivative, signed Hessian
and curvature readers are unchanged. The existing $t=O(a^2)$ selected
histories are not asserted to satisfy the new entrance on $t=O(a)$.
The new items are the evaluated cubic harmonic Poisson block, a necessary
mixed-rate limit, its complete positive reduction and forced zero mode, and
an original-rate geometric estimate. An exact initial population is proved
separately; a uniform original lift on $[0,aT]$ is not.

We use the V51 notation
\[
 \mathbb K=\begin{pmatrix}K_A&C\\-C^T&K_g\end{pmatrix},\qquad
 \mathbb t=\binom{t_A}{t_g},\qquad
 \mathscr R=(\mathscr R_A,\mathscr R_g),
\]
and its actual cubic derivative $s_c$. All matrices in the numerical
certificate use the common spatial-sum covector convention. The multiplier
syntheses are physical: $I_A$ has 78 active columns, $I_g$ has 26 nonconstant
curl-free columns, and $I_c$ consists of the three unit constant shifts.
The leading signed blocks are
\begin{equation}
 L_c=L I_c^W,\quad C_c=(I_c^W)^TC_WI_c^W,\quad
 D_c=C_c-L_c^TA^{-1}L_c\succ0,
 \label{v52:eq:mass-blocks}
\end{equation}
where $I_c^W$ inserts the first three coordinates in the original weak bank.
Positivity of $A,D_c$, rank $\mathbb K=104$, rank $K_g=24$, and the rank-two
coupling $C\ker K_g$ are inherited proved inputs. The supplied Einstein
bridge's cubic harmonic bracket formula is used at its coefficient scope
\cite{v52:EB36}; its spectral counts at its different root are not imported.

\section{The next harmonic Poisson coefficient is nonzero}
\label{v52:sec:cubic}

For $c,d\in\R^3$, let $\delta_c=c^i\delta_i$ and retain the complete original
cubic constant-shift envelope $P_{c,3}$. Define
\begin{equation}
 S[c,d]=DP_{c,3}(Y_*)[\delta_dY_*]
              -DP_{d,3}(Y_*)[\delta_cY_*].
 \label{v52:eq:S}
\end{equation}
This $3$-by-$3$ matrix is skew. The coefficient identity from the supplied
bridge gives, on exactly constrained records with bounded normalized first
coefficients,
\begin{equation}
 K_{RR}^{\rm orig}=a\mathbb K+O(a^2),\quad
 K_{cR}^{\rm orig}=a^2\mathscr R+O(a^3),\quad
 K_{cc}^{\rm orig}=a^3S+O(a^4).
 \label{v52:eq:Korders}
\end{equation}
These are Poisson blocks, not Hessian blocks. They are independent of bounded
higher initial preparation. In particular $\mathscr R I_c=0$ at the preceding
order does not imply $S=0$.

\begin{proposition}[An original cubic harmonic enclosure]
\label[proposition]{v52:prop:S}
On the retained root cube of radius $2^{-112}$, all nine derivatives in
\eqref{v52:eq:S} have been regenerated by exact rational interval arithmetic.
The resulting matrix has rational centre approximately
\begin{equation}
 S_0\simeq
 \begin{pmatrix}
 0&226.54000566&-424.10677731\\
 -226.54000566&0&-1014.29431412\\
 424.10677731&1014.29431412&0
 \end{pmatrix},
 \label{v52:eq:Scentre}
\end{equation}
and common componentwise error less than $2\cdot10^{-24}$. Each displayed
nonzero upper-triangular entry excludes zero.
\end{proposition}
\begin{proof}
The recovered source compiler supplies $P_{c,3}$ with its normal auxiliary,
shifted electric, and ordered magnetic terms. For each $d=e_j$, evaluate the
vector of three derivatives by the exact cubic polarization
\[
 DP_{c,3}(Y)[\delta_jY]
 =\tfrac12\{P_{c,3}(Y+\delta_jY)-P_{c,3}(Y-\delta_jY)\}
                         -P_{c,3}(\delta_jY).
\]
The physical momentum and phase grading are retained, with rational enclosing
endpoints for $\sqrt3$. Subtract the transpose. The arithmetic circuit uses
rational centres and outward-rounded radii, and includes the whole root cube.
The resulting radius is less than $1.063\cdot10^{-24}$. Rational endpoint
comparisons establish the stated bounds. The decimal matrix is descriptive;
it is not the defining source data or the proof of a zero.
\end{proof}

For clarity, the preparation independence in \eqref{v52:eq:Korders} follows
from the same original coefficient identities used in V48. One has
$P_{c,2}(Y)=\langle p,\delta_cU\rangle_h$ and
$\mathbb J\nabla P_{c,2}=\delta_cY$. The quadratic translation generators
commute on every field. Their polarization removes every second-field term
from the mutual cubic bracket, leaving \eqref{v52:eq:S}. In the mixed bracket,
the second constraint equation replaces its second-field contribution by
$-\langle\delta_c\nu,P_2(Y)\rangle_h$, yielding the full $\mathscr R$ already
retained in V51. This explanation does not reclassify those inherited
coefficient identities as new results.

\section{A necessary mixed-rate limit including every harmonic row}
\label{v52:sec:limit}

Consider original, exactly constrained histories on $0\le t\le aT$ in the
inherited stationary chart, with mean lapse one,
\begin{equation}
 X_a(0)=aY_*+O(a^2),\qquad
 \lambda_a(0)=e+a\ell_*+O(a^2),\qquad
 \lambda_t(a\theta)=aI_Av_a(\theta)+aI_gw_a(\theta)+I_cc_a(\theta),
 \label{v52:eq:entrance}
\end{equation}
and a uniform bound on $\|v_a\|+\|w_a\|+\|c_a\|$ on $[0,T]$.
This is a new hypothesis for a necessary limit, not a proved interval of
existence. Write $V_a,W_a,H_a$ for their time primitives, respectively.
Then
\begin{equation}
 \lambda(a\theta)-\lambda(0)=a^2(I_AV_a+I_gW_a)+aI_cH_a,
 \quad
 X(a\theta)-X(0)=a^2\theta V+O_T(a^3),
 \label{v52:eq:increments}
\end{equation}
where $V=F_1Y_*+G\ell_*$. The second assertion uses $GI_c=0$: a constant
shift changes the canonical field only at its next degree. It does not use
continuous translation symmetry of the finite action.

\begin{lemma}[Prepared mixed-rate source increments]
\label[lemma]{v52:lem:increments}
Under \eqref{v52:eq:entrance}, uniformly on $[0,T]$,
\begin{align}
 a^{-3}\Delta\mathsf B_A
 &=\theta t_A+K_AV_a+CW_a-\mathscr R_A^TH_a+O_T(a),\notag\\
 a^{-3}\Delta\mathsf B_g
 &=\theta t_g-C^TV_a+K_gW_a-\mathscr R_g^TH_a+O_T(a),\notag\\
 a^{-4}\Delta\mathsf B_c
 &=\theta s_c+\mathscr R_AV_a+\mathscr R_gW_a+SH_a+O_T(a).
 \label{v52:eq:source-increments}
\end{align}
Here $\Delta$ denotes the value at $a\theta$ minus the initial value.
\end{lemma}
\begin{proof}
Work on the exact nonconstant-constraint chart of V46. For the nonconstant
rows use $\widetilde{\mathsf B}=\mathsf B-\mathsf J_0P$, which agrees with
$\mathsf B$ on exact records and has no linear term. Its quadratic coefficient
is the prepared $B_2=\widehat B_2+\mathbb K_1\ell$. On the flat constraint
tangent, $\mathbb K_1(Y)I_c=0$, including its derivative along $V$. Therefore
that coefficient is unchanged by the bounded leading harmonic value $H_a$.
The other record increments in \eqref{v52:eq:increments} give precisely V50's
$\theta\mathfrak t+\mathbb K(V_a,W_a)$ terms.

The derivative in a harmonic multiplier direction is one order higher.
The exact identity is $D_\lambda\mathsf B=K+D_XM[F]$ before the constraint
chart is used. Here $M_{Rc}=O(a^3)$ and $F=O(a)$, so its second term is
$O(a^3)$, whereas $K_{Rc}=-a^2\mathscr R^T+O(a^3)$. The derivative of the
constraint elimination contributes only the higher order, since
$D_\lambda X=O(a^3)$ in a constant-shift direction. Integrating through the
$O(a)$ harmonic change gives the first two harmonic terms in
\eqref{v52:eq:source-increments}.

On the same chart the constant drift begins with the complete prepared cubic
coefficient $C(Y,\ell)$ from V51. Its derivative in a constant first tilt
vanishes, but the next derivative is $a^3S+O(a^4)$ by
\eqref{v52:eq:Korders}. The term $D_XM_{cc}[F]$ is $O(a^4)$, and the canonical
correction from the constraint chart again starts one order higher. Thus the
quartic chart coefficient is affine in the bounded harmonic first tilt with
coefficient $S$. This yields $SH_a$; changes of the other normalized variables
in its coefficient are $O(a)$. The original cubic coefficient gives
$\theta s_c+\mathscr R_AV_a+\mathscr R_gW_a$ exactly as in V51. Higher bounded
initial coefficients alter the displayed increments only by their remainders.

The leading graded Hessian is independent of a bounded first tilt by the
original weak tangent-Hessian cancellation: its first spatial connection is
$z_1(Y)$, and tilt changes in the second spatial connection are flat tangent
vectors annihilated on weak tests. This inherited identity is used uniformly
for bounded $H_a$, not just at its initial value. All Taylor bounds therefore
hold on the compact normalized coefficient set generated by the assumed
bound. No derivative of $v_a,w_a,c_a$ is needed.
\end{proof}

Define the physical rate mass and skew matrix in the order $(v,w,c)$:
\begin{equation}
 \mathbb M=
 \begin{pmatrix}A&0&L_c\\0&0&0\\L_c^T&0&C_c\end{pmatrix},\qquad
 \mathbb J_c=
 \begin{pmatrix}
 K_A&C&-\mathscr R_A^T\\
 -C^T&K_g&-\mathscr R_g^T\\
 \mathscr R_A&\mathscr R_g&S
 \end{pmatrix},\qquad
 f=\begin{pmatrix}t_A\\t_g\\s_c\end{pmatrix}.
 \label{v52:eq:pencil}
\end{equation}
In particular $\mathbb M\succeq0$, $\operatorname{rank}\mathbb M=81$, and
$\mathbb J_c^T=-\mathbb J_c$. Positivity follows from the Schur complement
$D_c\succ0$, not positivity of the complete original signed Hessian.

\begin{theorem}[Harmonic motion closes the necessary mixed-rate equations]
\label[theorem]{v52:thm:limit}
If the initial $(v_a(0),c_a(0))$ converge, the rates in
\eqref{v52:eq:entrance} have the following unique limits: $v_a,c_a$ converge
uniformly; $w_a$ converges weak-star in $L^\infty$ and weakly in $L^2$, and
its primitive converges uniformly. They solve
\begin{equation}
 \boxed{\mathbb M\dot q+\mathbb J_cq+f=0,\qquad q=(v,w,c),}
 \label{v52:eq:dae}
\end{equation}
with the inherited initial values for $v,c$. Initial compatibility is the
two-row condition in \eqref{v52:eq:compatibility} below. No convergence of
original rate derivatives, or native existence on $[0,aT]$, is asserted.
\end{theorem}
\begin{proof}
The leading consistency rows now read
\[
 a^{-3}\mathsf B_A=-Av_a-L_cc_a+O_T(a),\quad
 a^{-3}\mathsf B_g=O_T(a),\quad
 a^{-4}\mathsf B_c=-L_c^Tv_a-C_cc_a+O_T(a).
\]
The $C_cc_a$ term is leading and cannot be discarded. Subtract the initial
rows and insert \cref{v52:lem:increments}. This gives the integral form of
\eqref{v52:eq:dae} with a uniform $O_T(a)$ residual. The three rate primitives
are uniformly Lipschitz. Since the mass on $(v,c)$ is positive definite, its
integral identity gives uniform convergence of $v_a,c_a$ along extracted
subsequences. Pass to the limit, then differentiate in distributions.
The next theorem makes the limiting equation a unique smooth affine ODE on
its compatible data. Thus all extracted limits agree. Pointwise convergence
of $w_a$ at the initial endpoint need not hold.
\end{proof}

\section{All 107 rows reduce to a positive 79-dimensional transport}
\label{v52:sec:reduction}

Reorder the differential variables as $x=(v,c)\in\R^{81}$, and put
\begin{equation}
 M_x=\begin{pmatrix}A&L_c\\L_c^T&C_c\end{pmatrix}\succ0,\quad
 J_x=\begin{pmatrix}K_A&-\mathscr R_A^T\\\mathscr R_A&S\end{pmatrix},\quad
 E=\binom{C}{\mathscr R_g},\quad f_x=\binom{t_A}{s_c}.
 \label{v52:eq:reorder}
\end{equation}
Let $Z$ span $\ker K_g$ and define
\[
 T_x=EZ,\quad G_x=T_x^TM_x^{-1}T_x\succ0,\quad R_g=K_g^\dagger,
 \quad
 P_x=M_x^{-1}-M_x^{-1}T_xG_x^{-1}T_x^TM_x^{-1}.
\]
The upper part $CZ$ is injective by the inherited rank-two theorem, so
$T_x$ has full column rank. The admissible affine space is
\begin{equation}
                 \Sigma_x=\{x:T_x^Tx=Z^Tt_g\},\qquad \dim\Sigma_x=79.
 \label{v52:eq:compatibility}
\end{equation}

\begin{theorem}[Exact mixed-rate elimination and positive difference energy]
\label[theorem]{v52:thm:positive}
For $x\in\Sigma_x$, define
\begin{align}
 w_0(x)&=-R_g(t_g-E^Tx),\notag\\
 d_x(x)&=f_x+J_xx+Ew_0(x),\notag\\
 w(x)&=w_0(x)-ZG_x^{-1}T_x^TM_x^{-1}d_x(x).
 \label{v52:eq:reconstruction}
\end{align}
Then every row of \eqref{v52:eq:dae} is equivalent to
\begin{equation}
                  \boxed{\dot x=-P_xd_x(x),\qquad w=w(x).}
 \label{v52:eq:positive-ODE}
\end{equation}
This ODE is global on every finite comparison interval. For two solutions,
\begin{equation}
 \boxed{\frac{d}{d\theta}\frac12(x_1-x_2)^TM_x(x_1-x_2)=0.}
 \label{v52:eq:energy}
\end{equation}
\end{theorem}
\begin{proof}
The algebraic row is $t_g-E^Tx+K_gw=0$. Its compatibility condition is
\eqref{v52:eq:compatibility}, and all solutions are $w=w_0+Z\eta$.
Differentiate compatibility once in $M_x\dot x=-d_x-T_x\eta$.
The positive matrix $G_x$ fixes $\eta$ and gives
\eqref{v52:eq:reconstruction}. Conversely $T_x^TP_x=0$ preserves the affine
space and verifies all the equations. Since $R_g^T=-R_g$, the effective
matrix $J_x+ER_gE^T$ is skew. Pairing the difference equation with a tangent
difference $\delta x$, for which $T_x^T\delta x=0$, proves
\eqref{v52:eq:energy}. The reduced vector field is affine and hence global.
These are properties of the necessary mixed-rate equation, not a stability
or lift theorem for the original singular dynamics. The general descriptor
method is standard \cite{v52:vdSM}; the actual coefficients and their
harmonic coupling here are not imported from that reference.
\end{proof}

Unlike V51's 76-dimensional reduction, this odd-dimensional positive
transport must have a zero frequency. Its actual multiplicity and forcing
are decided next, rather than inferred from dimension alone.

\section{One zero mode is forced, with a certified harmonic slope}
\label{v52:sec:secular}

The complete nonconstant Poisson matrix is invertible. Define its harmonic
skew Schur complement and the original harmonic source residual by
\begin{equation}
 \Xi=S+\mathscr R\mathbb K^{-1}\mathscr R^T,\qquad
 \rho_c=s_c-\mathscr R\mathbb K^{-1}\mathbb t.
 \label{v52:eq:Xi}
\end{equation}
The second vector is exactly V51's source, not a new forcing choice.

\begin{proposition}[A rank-two harmonic Schur and a nonzero source projection]
\label[proposition]{v52:prop:Xi}
At the exact root, $\Xi$ is skew and has rank two. Its rational centre is
approximately
\begin{equation}
 \Xi_0\simeq
 \begin{pmatrix}
 0&180847.54568592&-89866.27756054\\
 -180847.54568592&0&387732.35818103\\
 89866.27756054&-387732.35818103&0
 \end{pmatrix},
 \label{v52:eq:Xi-centre}
\end{equation}
with common entry error below $10^{-15}$. Let $d_0$ be the fixed exact
positive rational centre of its $(1,2)$ entry, and define
\begin{equation}
 k=d_0^{-1}(\Xi_{23},-\Xi_{13},\Xi_{12})^T,\quad
 n=\mathbb K^{-1}\mathscr R^Tk,\quad r=(n,k).
 \label{v52:eq:null}
\end{equation}
Then $\ker\mathbb J_c=\R r$ and
\begin{equation}
 -24551000<\phi:=k^T\rho_c<-24542000,\qquad
 8810000<m:=r^T\mathbb M r<8812000.
 \label{v52:eq:phi-m}
\end{equation}
The normalization makes the centre of $k$ approximately
$(2.14397357,0.49691732,1)$; its third actual component is not declared
exactly one.
\end{proposition}
\begin{proof}
Assemble $\Xi$ from the new nine cubic derivatives and the retained V51
original Poisson inverse and cubic-mean map. The new common entry error is
less than $6.86\cdot10^{-17}$. Its $(1,2)$ entry excludes zero, so a
three-dimensional skew matrix has exactly rank two. Block elimination gives
$\operatorname{rank}\mathbb J_c=104+2=106$ and the kernel formula
\eqref{v52:eq:null}. The propagated radius of $k$ is below
$3.794\cdot10^{-22}$. Contract it with the full V51 source enclosure, whose
entry radius is below $1016$. The new projection centre is approximately
$-24546665.7594595$, with radius below $3699$. This proves its strict
separation.

For the mass use the manifestly positive formula
\[
 m=(n_A+A^{-1}L_ck)^TA(n_A+A^{-1}L_ck)+k^TD_ck.
\]
Its centre is approximately $8810892.63113434$, with radius below $10^{-12}$.
Exact rational comparisons give \eqref{v52:eq:phi-m}. No floating matrix
inverse or eigenvalue is used as the certificate.
\end{proof}

\begin{theorem}[The full reduced spectrum and its universal secular drift]
\label[theorem]{v52:thm:secular}
The homogeneous 79-dimensional transport has one semisimple zero and 39
nonzero purely imaginary pairs, all semisimple. It has no stationary solution
for the actual forcing. Every complete mixed-rate solution has
\begin{equation}
 \boxed{q(\theta)=-\frac{\phi}{m}\,r\theta+q_{\rm bounded}(\theta),}
 \label{v52:eq:secular}
\end{equation}
where $q_{\rm bounded}$ is bounded for all real comparison times. In
particular, the harmonic slope $s_c^{\rm sec}=-(\phi/m)k$ is independent of
initial data and satisfies
\begin{equation}
 \boxed{
 5.972<(s_c^{\rm sec})_1<5.974,\quad
 1.383<(s_c^{\rm sec})_2<1.386,\quad
 2.785<(s_c^{\rm sec})_3<2.787.}
 \label{v52:eq:slope}
\end{equation}
\end{theorem}
\begin{proof}
The nullspace correspondence in V51 applies to the full $(x,w)$ block:
the kernel of the skew form restricted to $\ker T_x^T$ is isomorphic to
$\ker\mathbb J_c$. Hence it is one-dimensional. Its reduced mass is positive,
so the homogeneous generator is skew-adjoint in a positive norm. Zero is
semisimple and the other 78 dimensions have nonzero semisimple imaginary
spectrum. There are 39 pairs with multiplicities; distinct frequencies are
not asserted.

Pair \eqref{v52:eq:dae} with $r$. Since $\mathbb J_cr=0$,
\[
             \frac{d}{d\theta}r^T\mathbb Mq=-r^Tf=-k^T\rho_c=-\phi.
\]
Positive orthogonal projection onto the one-dimensional reduced kernel gives
the linear coefficient $-\phi/m$. On its orthogonal complement the generator
is invertible and skew-adjoint; constant forcing has a bounded particular
solution there and homogeneous solutions are bounded. Reconstructing $w$
from the affine map in \eqref{v52:eq:reconstruction} gives
\eqref{v52:eq:secular} for all components. A stationary solution would require
$\phi=0$, which is excluded.

The scalar $-\phi/m$ has centre approximately $2.7859453959$ and radius
below $0.000420$. The vector centre is approximately
$(5.97299330,1.38438452,2.78594540)$, with common radius below $0.000901$.
Strict rational endpoint comparisons prove \eqref{v52:eq:slope}.
\end{proof}

This secular statement concerns the necessary limiting equation for finite
$\theta$. It does not permit taking $\theta\to\infty$ jointly with $a\to0$
in the original histories. It does show that a stationary faster harmonic
rate cannot satisfy every limiting row. If $c\equiv0$ on an interval, the
system would reduce to the incompatible V51 completion; thus every actual
mixed-rate solution has a nonzero harmonic response somewhere on each
nonempty interval. The long-time linear coefficient identifies that response
more precisely without proving its original-action lift.

\section{Faster harmonic initial rates are compatible with small curvature}
\label{v52:sec:geometry}

\begin{theorem}[Exact initial realization at the new rate scale]
\label[theorem]{v52:thm:initial}
For any fixed finite $(v_0,w_0,c_0)$ there are sufficiently small-amplitude,
exactly constrained initial records, with the same leading $Y_*,\ell_*$ and
normalized initial means, on which the unique normalized original rate is
\begin{equation}
                  \lambda_t(0)=aI_Av_0+aI_gw_0+I_cc_0.
 \label{v52:eq:initial-rate}
\end{equation}
Their lapse and spatial metric are positive and their original metric and
link curvatures are $O(a)$. This does not assert a uniform interval $[0,aT]$.
\end{theorem}
\begin{proof}
Use the original constraint initializer with
$\lambda=e+a\ell(q)+a^2m$, allowing the same three root variables $q$ and
104 normalized second multiplier coefficients $m$ as in V39. In
$\mathsf B+M(aI_Av_0+aI_gw_0+I_cc_0)$, the added harmonic term is $O(a^3)$
on active rows and $O(a^4)$ on all weak rows. After division by $a^3$, the
three leading constant equations are still exactly the verified cubic means
$\mathbf C(q)=0$. The 104 nonconstant equations have derivative
$\mathbb K|_{\mathcal E^\perp}$ in $m$. Their leading right side is shifted
by the finite $A v_0+L_cc_0$ term, not assumed zero. The same triangular
implicit-function derivative as V39 is invertible. It supplies exact initial
constraints and 107 consistency rows, and radial homogeneity supplies the
remaining row. Leading signed regularity is unchanged, so the original
normalized solve selects \eqref{v52:eq:initial-rate}. Positivity follows by
small amplitude. The curvature estimate follows from the next lemma.
Initial compatibility of a future bounded mixed-rate limit imposes
\eqref{v52:eq:compatibility} on $(v_0,c_0)$; this initial-existence theorem
does not dispense with it.
\end{proof}

\begin{lemma}[The harmonic rate does not create an order-one flat curvature]
\label[lemma]{v52:lem:curvature}
On the entrance \eqref{v52:eq:entrance}, or at the initial records just
constructed,
\begin{equation}
 \|X\|+\|X_t\|+\|X_{tt}\|+\|\lambda-e\|=O_T(a),\qquad
 \|\Riem(g_h)\|+\|E_h^{\rm link}\|+\|F_h^{\rm sp,link}\|=O_T(a).
 \label{v52:eq:curvature-bound}
\end{equation}
The constants may depend on the fixed mixed-rate bound. No multiplier
acceleration bound is used.
\end{lemma}
\begin{proof}
The canonical field is analytic and its derivative in a constant shift is
$D_\lambda F[I_cc]=a\delta_cY_*+O_T(a^2)$. Its flat value is zero.
Consequently $X_{tt}=D_XF[F]+D_\lambda F[\lambda_t]=O_T(a)$ even with
$c_a=O(1)$. Fixed-cutoff spatial derivatives have the same bounds.
The coordinate metric curvature is affine in the multiplier rates after
substituting the canonical equation. At the flat spatial metric, unit lapse
and zero nonconstant fields, an arbitrary time-dependent constant shift gives
the flat metric $-dt^2+|dx+\beta(t)dt|^2$. Thus its harmonic rate coefficient
vanishes at the flat record, and is $O(a)$ nearby. Every nonconstant rate is
$O(a)$ by assumption. This proves the metric bound without incorrectly
bounding each coordinate connection separately.

For the original spatial stationary connection, the constant-shift
multiplier derivative is $-ZK_{\rm tan}^{-1}J_c=O(a^2)$, by the inherited
weak second-source identity. Its field derivative is bounded and $X_t=O(a)$.
The spatial and temporal connection values remain $O(a)$; their literal
ordered curvature expressions therefore give the two link estimates. No
equality between the metric and link curvature coefficients is inferred.
\end{proof}

The geometric bound does not make the constant-rate response invisible at
order $a$. There is an explicit fixed-cutoff discrepancy. Use the Riemann
convention
\[
 R_{0i0j}=\tfrac12(\partial_0\partial_i g_{0j}
 +\partial_j\partial_0 g_{i0}-\partial_i\partial_jg_{00}
 -\partial_0^2g_{ij})+\hbox{the corresponding Christoffel products}.
\]
At the same record, compare rate choices differing only by $I_cc$, with
canonical acceleration calculated from the original Hamiltonian in both.

\begin{proposition}[The leading metric reader of a faster harmonic rate]
\label[proposition]{v52:prop:reader}
For bounded $c$ and the leading spatial metric $I+aU_*$,
\begin{equation}
 \boxed{\Delta_c R_{0i0j}
   =\frac a2\left(\partial_c-\delta_c\right)(U_*)_{ij}+O(a^2),\qquad
        \Delta_cE_h^{\rm link}=O(a^2).}
 \label{v52:eq:reader}
\end{equation}
With the opposite curvature sign convention both metric terms change sign.
The metric rate map in \eqref{v52:eq:reader} is injective on the three
harmonic directions at this root. In the normalized spatial $L^2$ Frobenius
norm its leading coefficient obeys
\begin{equation}
 \left\|\tfrac12(\partial_c-\delta_c)U_*\right\|_2
 \ge\frac{2\pi-3\sqrt3}{2}
       \left(\sum_i\alpha_i^2c_i^2\right)^{1/2}
 \ge\frac{2\pi-3\sqrt3}{2}|c|.
 \label{v52:eq:reader-lower}
\end{equation}
\end{proposition}
\begin{proof}
Changing the rate changes $\gamma_{tt}$ by $a\delta_cU_*+O(a^2)$,
because $P_{c,2}$ generates the original phase translation. In the second
metric derivatives the spatial derivative of $g_{0i,t}=\gamma_{ik}c^k+
\text{unchanged terms}$ contributes
$a c^k(\partial_iU_{jk}+\partial_jU_{ik})/2$.
The Christoffel product containing the order-one $\Gamma_{00}^k=c^k+O(a)$
subtracts
$a c^k(\partial_iU_{jk}+\partial_jU_{ik}-\partial_kU_{ij})/2$.
Combining with $-\Delta\gamma_{tt}/2$ proves \eqref{v52:eq:reader}.
All omitted rate-dependent terms contain two small field factors. The link
statement is the constant derivative of the spatial stationary connection
used in the preceding lemma; the temporal connection value does not depend
on multiplier rates.

The retained trigonometric interpolation has ordinary derivative $2\pi i k$
and original phase derivative $3\sqrt3\,i k$ on each axial frequency. The
three axial tensors have squared Frobenius norm two, their sine profiles
have squared spatial norm $1/2$, and their amplitudes are $(1,9/8,1)$.
They are orthogonal to one another and to the mixed frequencies. Keeping
just their nonnegative contribution proves \eqref{v52:eq:reader-lower}.
The remaining modes cannot cancel that norm contribution.
\end{proof}

This is a rate-slot comparison, not a lower bound for the complete curvature
along a history: active rates and other fields can contribute too. A genuine
geometric translation removes the ordinary translation part, not the
finite phase derivative. At this cutoff the discrepancy is nonzero. For a
fixed smooth Fourier bank its symbol is $O(h^2)$ under refinement, but no
cofinal control of the growing harmonic rate or the original histories is
proved here.

\section{Verification, scope, and the next lift problem}
\label{v52:sec:status}

\paragraph{Proved.}
The nine original cubic harmonic derivatives supply a nonzero skew block.
Retaining it and the harmonic inertia gives a complete necessary mixed-rate
system. Its 79-dimensional constrained reduction has positive homogeneous
difference energy, one semisimple zero and 39 oscillatory pairs. The actual
source has a nonzero projection on that zero mode, giving a universally
forced harmonic secular slope with explicit enclosures. Fixed faster
harmonic initial rates are exactly realizable, and their original curvature
remains $O(a)$. The fixed-cutoff rate-slot geometric discrepancy is identified
rather than calling those rates a gauge freedom.

\paragraph{Inherited and executed.}
A clean V51 replay recomputes its parent 108 differentiated drift entries and
three cubic-mean derivatives, and checks its original full Poisson inverse
and harmonic source intervals. The new worker separately regenerates all
nine $DP_{c,3}[\delta_dY_*]$ values on the root cube. A second worker forms
$\Xi$, its null vector, the source projection, positive mass and secular
slope in exact rational interval arithmetic. The old 79,224-entry root bank,
17,334-entry stationary-source bank, and 5,046-entry weak-bracket bank are
inherited and replayed, not wholly regenerated. Small-system algebra checks
verify the complete descriptor reconstruction, positive energy, kernel and
secular formulas, mixed scaling, and the curvature jet cancellation. They
are not original spacetime simulations or proof-assistant formalizations.

\paragraph{Unpaid and next direction.}
There is no original-action existence or strong rate-convergence theorem on
$[0,aT]$ at the new scaling. A compatible initial pair $(v_0,c_0)$ gives a
unique limiting transport, and a finite initial rate can be realized, but
neither fact controls the intervening fast modes or selects a well-prepared
native history. The next problem is that mixed-rate lift with its original
metric and link observations. The secular coefficient makes a stationary
harmonic-rate ansatz insufficient; the $O(a)$ geometric estimate shows why
V51's coordinate-rate obstruction is not a curvature obstruction.

The frozen canonical spectral pencil in the supplied Einstein companion is
a different, initial-matrix question at a different root. Its translation
energy term and real-pair theorem have not been replaced by this moving,
source-resolved necessary rate equation. No stability of the original flow,
normal attraction, fixed positive physical lifespan, spatially uniform
Einstein limit, or primitive Einstein--SM selection follows here. V32's
independent compensation question is unchanged.

Every V51 byte before its final document command is preserved. New labels
use \texttt{v52:}. The source intervals are arithmetic certificates; the
initialization, limiting-system, and geometric statements are written proofs.

\begingroup
\renewcommand{\refname}{Additional references for Version 52}
\begin{thebibliography}{9}
\small\raggedright
\bibitem{v52:EB36}
\emph{Renewal Exact-Action Einstein Bridge Research Notes}, supplied cumulative
Versions 1--36, file
\nolinkurl{Renewal_Exact_Action_Einstein_Bridge_Research_Notes_V36(2).tex}.
The used original coefficient identities are \nolinkurl{v20:graded-source}
and \nolinkurl{v36:Poisson-orders}, including its full cubic constant bracket.
Their different-root spectral results are not imported.
\bibitem{v52:vdSM}
A.~van der Schaft and B.~Maschke, \emph{Generalized port-Hamiltonian DAE
systems}, Systems \& Control Letters \textbf{121} (2018), 31--37.
\url{https://arxiv.org/abs/1808.01845}.
Primary-source context for conservative systems with algebraic constraints;
no original-source coefficient or native existence theorem is taken from it.
\end{thebibliography}
\endgroup
\addtocontents{toc}{\protect\endgroup}
% END IMMUTABLE VERSION 52 BODY
% APPEND VERSION 53 HERE, BEFORE THE FINAL END-DOCUMENT COMMAND.


% BEGIN IMMUTABLE VERSION 53 BODY
\clearpage
\begin{center}
{\Large\bfseries Version 53: Recover the modes hidden by the positive rate limit}\\[.5em]
{\large The full graded coefficient pencil, three spectral scales,\\
and an exact well-prepared coefficient-level lift}
\end{center}
\makeatletter
\addtocontents{toc}{\protect\begingroup}
\addtocontents{toc}{\protect\makeatletter}
\addtocontents{toc}{\protect\def\protect\l@section{\protect\bfseries\protect\@dottedtocline{1}{0em}{2.5em}}}
\makeatother
\addcontentsline{toc}{part}{Version 53: Graded rate pencil and a well-prepared coefficient-level lift}

\section{The positive limit does not control the discarded inertia}
\label{v53:sec:scope}

Version 52 supplies a necessary mixed-rate descriptor equation, an exact
initial-rate population, and a positive 79-dimensional reduction. It does
not prove that finite-amplitude original histories approach that reduction.
Before a nonlinear lift, the vanishing but signed inertial terms must be
retained: positivity of a limiting mass does not bound modes which disappear
when that mass becomes singular. This body resolves that linear-coefficient
question at the \emph{same} V39 algebraic root and cutoff $h=1/3$.

The new object is explicitly a \emph{frozen graded coefficient pencil}.
Its coefficients are the actual leading stationary-source and Poisson
coefficients already evaluated in Versions 41--52, but its equation is not
identified with the full nonlinear original flow. This distinction is
maintained in every lift and stability conclusion. No initial microscopic
law, action term, clock, or physical update is replaced by the pencil.
The construction gives an exact comparison and specifies which modes and
higher coefficients an original-flow proof still has to control.

We inherit the following proved inputs, rather than recalculate their primitive
source banks: $\operatorname{Inertia}\Gamma_*=(97,10,0)$; positivity of $A$
and $D_c$; the five-dimensional semisimple kernel of the weak generator;
its nine-dimensional growing, nine-dimensional decaying, and six-dimensional
positive oscillatory subspaces; $\operatorname{rank}K_g=24$; the two-dimensional
kernel metric of inertia $(1,1,0)$; the mixed-rate matrix
$\mathbb J_c$ of rank 106; and the positive 79-dimensional descriptor
reduction with one forced zero mode. The fresh V52 replay used here also
recomputes its nine cubic harmonic derivatives and the parent 111 drift/mean
derivatives. The new source calculation is an exact interval contraction of
these retained tensors, not a new 107-dimensional original-action evolution.

General quadratic-pencil and descriptor methods are established
\cite{v53:TM,v53:MMW}. The three-scale calculation, actual-root kernel
contraction, and prepared comparison below are proved explicitly. The signed
finite-amplitude mass is outside a blanket positive dissipative-pencil
argument; no stability theorem is imported from the positive descriptor.

\section{The exact graded mass and its singular limit}
\label{v53:sec:mass}

Use the physical synthesis order $(v,w,c)$ of Version 52, with dimensions
$(78,26,3)$ and $\theta=t/a$. Split the inherited weak blocks into
nonconstant and constant slots:
\[
 L=(L_c,L_g),\qquad
 C_W=\begin{pmatrix}C_c&C_{cg}\\C_{cg}^T&C_{gg}\end{pmatrix}.
\]
The complete leading Hessian in the order $(v,w,c)$ is denoted by
$\Gamma^{\rm ord}_*$. Define
\begin{equation}
 \begin{split}
 M(a)&=\begin{pmatrix}
 A&aL_g&L_c\\aL_g^T&a^2C_{gg}&aC_{cg}^T\\
 L_c^T&aC_{cg}&C_c
 \end{pmatrix},\\
 J&=\begin{pmatrix}
 K_A&C&-\mathscr R_A^T\\-C^T&K_g&-\mathscr R_g^T\\
 \mathscr R_A&\mathscr R_g&S
 \end{pmatrix},\qquad f=\begin{pmatrix}t_A\\t_g\\s_c\end{pmatrix}.
 \end{split}
 \label{v53:eq:pencil}
\end{equation}
Thus $M(a)=D_a\Gamma^{\rm ord}_*D_a$, where
$D_a=\operatorname{diag}(I_{78},aI_{26},I_3)$, and $J=\mathbb J_c$.
The coefficient-level comparison equation is
\begin{equation}
             M(a)\dot q+Jq+f=0,\qquad a>0.
 \label{v53:eq:regularized}
\end{equation}
It retains every graded inertial entry. For example, differentiating a rate
$aI_Av+aI_gw+I_cc$ in physical time and dividing the active, nonconstant-weak,
and harmonic rows by $a^2,a^2,a^3$, respectively, produces exactly the mass
in \eqref{v53:eq:pencil} at its retained leading coefficient. This calculation
does not discard the additional coefficients in the actual differentiated
source; \eqref{v53:eq:regularized} declares them absent only in the comparison.

\begin{proposition}[A signed regular pencil with a positive singular limit]
\label[proposition]{v53:prop:mass}
For $a>0$, $M(a)$ is nonsingular with inertia $(97,10,0)$.
At $a=0$, it is the Version-52 positive mass of rank 81.
The limiting pencil $zM(0)+J$ is regular, has 79 finite eigenvalues counted
algebraically, and has 28 eigenvalues at infinity. Its infinite part consists
of 24 blocks of size one and two blocks of size two. Its finite eigenvalues
are the one zero and 39 imaginary pairs of Version 52.
\end{proposition}
\begin{proof}
The first assertions follow from congruence and the Schur complement $D_c$.
Reorder as $x=(v,c)\in\R^{81}$ and $w\in\R^{26}$. Put
\begin{equation}
 M_x=\begin{pmatrix}A&L_c\\L_c^T&C_c\end{pmatrix}\succ0,
 \quad B_x=\binom{L_g}{C_{cg}},\quad
 E=\binom{C}{\mathscr R_g},\quad
 J_x=\begin{pmatrix}K_A&-\mathscr R_A^T\\\mathscr R_A&S\end{pmatrix}.
 \label{v53:eq:blocks}
\end{equation}
Choose a complement to $\ker K_g$, on which $K_g$ is nonsingular.
Eliminating these 24 algebraic variables leaves the 81 differential
coordinates and two multipliers $\eta$, with pencil
\[
 \begin{pmatrix}zM_x+J_x+EK_g^\dagger E^T&T_x\\-T_x^T&0\end{pmatrix},
 \qquad T_x=EZ,\quad Z\text{ spans }\ker K_g.
\]
The matrix $T_x$ has rank two. Split $x$ into its 79-dimensional tangent
$\ker T_x^T$ and a two-dimensional complement. Expanding the determinant
along the two constraint rows and the two multiplier columns leaves a
nonzero constant times the tangent pencil $zM_s+J_s$, where $M_s\succ0$.
Its determinant has degree 79. This proves regularity and the finite count.
The algebraic solution uses one differentiation of the two compatibility
conditions; no second derivative of the forcing is required in solving this
block system. Equivalently its nilpotent part has square zero. There are
$n-\operatorname{rank}M(0)=26$ nilpotent blocks, of total dimension
$107-79=28$. Thus two have size two and 24 have size one.
The finite spectrum is the already proved positive descriptor spectrum.
\end{proof}

In particular, dropping 26 inertial coordinates hides 28 dynamical modes:
two differential directions also become constrained. This count by itself
says neither that those modes decay nor that arbitrary initial data converge.

\section{Exact elimination reveals three spectral scales}
\label{v53:sec:scales}

Set
\begin{equation}
 H=M_x^{-1}B_x,\qquad \Delta=C_{gg}-B_x^TM_x^{-1}B_x.
 \label{v53:eq:delta}
\end{equation}
The congruence $x=\widehat x-aHw$ gives the exact transformed pencil
\begin{equation}
 \begin{pmatrix}
 zM_x+J_x&E-aJ_xH\\
 -E^T-aH^TJ_x&a^2z\Delta+K_g+aS_g+a^2H^TJ_xH
 \end{pmatrix},\qquad S_g=E^TH-H^TE.
 \label{v53:eq:sheared-pencil}
\end{equation}
All terms are retained. In particular the order-$a$ skew term $S_g$ is not
zero merely because it is absent from the descriptor at $a=0$.

For any fixed kernel synthesis $Z$ define
\begin{equation}
 D_0=Z^T\Delta Z,\qquad S_0=Z^TS_gZ,\qquad
 G_0=Z^TE^TM_x^{-1}EZ.
 \label{v53:eq:kernel-pencil}
\end{equation}
The inherited rank-two coupling gives $G_0\succ0$.

\begin{lemma}[The original kernel metric and stiffness reappear]
\label[lemma]{v53:lem:identification}
Take $Z$ to be the nonconstant part of the inherited two nonconstant weak
kernel columns. Then
\begin{equation}
 \boxed{D_0=D_v,\qquad
 G_0=G_v+Q^TD_c^{-1}Q=\overline G.}
 \label{v53:eq:identification}
\end{equation}
Hence $D_0$ has inertia $(1,1,0)$ and $G_0\succ0$.
The nonzero spectrum of $-\Delta^{-1}K_g$ is exactly the nonzero spectrum
of the inherited complete weak generator $-D^{-1}\mathsf K$.
Its zero eigenvalue has multiplicity two and is semisimple.
\end{lemma}
\begin{proof}
After eliminating $A$, eliminating the constant weak block gives
$\Delta=D_{gg}-D_{gc}D_c^{-1}D_{cg}$. The $D$-orthogonal weak synthesis is
$N_v=(-D_c^{-1}D_{cg}Z,Z)^T$. Thus $N_v^TDN_v=Z^T\Delta Z=D_v$.
Since $K_{AW}$ annihilates constant shifts and $\mathscr R_WI_c=0$,
\[
 EZ=\binom{T}{\mathscr R_WN_v},\qquad
 Q=\mathscr R_WN_v-L_c^TA^{-1}T.
\]
The block inverse of $M_x$ now gives the identity for $G_0$.
For a nonzero generalized weak eigenvalue, its constant row solves the
constant component as $-D_c^{-1}D_{cg}w$. The remaining equation is exactly
$(z\Delta+K_g)w=0$. Finally $\Delta$ is nonsingular by inertia congruence,
and its restriction to $\ker K_g$ is nondegenerate. Therefore
$\operatorname{Ran}(-\Delta^{-1}K_g)\cap\ker K_g=\{0\}$, proving
semisimplicity at zero. The two fast systems have identical nonzero
spectra, with algebraic multiplicities, by the same Schur determinant.
\end{proof}

\begin{theorem}[Three-scale spectral splitting of the full coefficient pencil]
\label[theorem]{v53:thm:three-scales}
For all sufficiently small $a>0$, the 107 eigenvalues of
$zM(a)+J$ split as follows.
\begin{enumerate}[label=\textup{(\roman*)},leftmargin=2.4em]
\item Twenty-four eigenvalues obey $a^2z\to\mu\ne0$, with the multiplicities
of $-\Delta^{-1}K_g$: nine lie in each open half-plane, and six are purely
imaginary.
\item Four eigenvalues obey $az\to\nu\ne0$, where
\begin{equation}
                 \det(\nu^2D_0+\nu S_0+G_0)=0.
 \label{v53:eq:middle-pencil}
\end{equation}
These are simple roots $+\sigma_0,-\sigma_0,+i\omega_0,-i\omega_0$, with
$\sigma_0,\omega_0>0$.
\item Seventy-nine eigenvalues remain bounded and approach the complete
Version-52 descriptor spectrum. Exactly one is zero. The other 78 are
purely imaginary and semisimple.
\end{enumerate}
In total there are ten eigenvalues in each open half-plane, counted
algebraically, 43 nonzero imaginary pairs, and one semisimple zero.
This is an asymptotic theorem; no numerical common upper bound on $a$ is
asserted.
\end{theorem}
\begin{proof}
For $z=\mu/a^2$ with $\mu$ away from zero, eliminate the upper block in
\eqref{v53:eq:sheared-pencil}. Its inverse is
$a^2\mu^{-1}M_x^{-1}+O(a^4)$. The lower Schur matrix tends uniformly on
compact nonzero $\mu$-sets to $\mu\Delta+K_g$. Its 24 nonzero roots persist
with their algebraic counts, by the argument principle on disjoint contours.
The inherited nonzero spectrum has nine growing, nine decaying, and six
positive-energy oscillatory dimensions.

For $z=\nu/a$, the upper inverse is $a\nu^{-1}M_x^{-1}+O(a^2)$.
Its lower Schur matrix is
\[
 K_g+a\{\nu\Delta+S_g+\nu^{-1}E^TM_x^{-1}E\}+O(a^2).
\]
Split $w$ into $\ker K_g$ and its nonsingular complement. Eliminating the
latter costs only $O(a^2)$ in the kernel equation. Divide by $a$ and multiply
by $\nu$. The limit is \eqref{v53:eq:middle-pencil}.
Write $S_0=\left(\begin{smallmatrix}0&\chi_0\\-\chi_0&0\end{smallmatrix}\right)$,
$d=\det D_0<0$ and $g=\det G_0>0$. Its determinant is
\begin{equation}
 d\nu^4+
 \{\operatorname{tr}(\operatorname{adj}D_0\,G_0)+\chi_0^2\}\nu^2+g.
 \label{v53:eq:quartic}
\end{equation}
The two roots in $\nu^2$ are simple and have opposite signs. Therefore
there are exactly the four displayed simple nonzero roots; analytic implicit
functions give $z=\nu/a+O(1)$. Real and skew/symmetric pencil symmetries
keep the real pair real and the isolated imaginary pair imaginary.

On fixed $z$-contours the pencil tends to the regular descriptor. Its degree-79
finite determinant therefore supplies 79 bounded roots, counting algebraic
multiplicity. These and the preceding 28 roots exhaust degree 107.
There are at least ten eigenvalues in the right half-plane. For the matrix
$L_a=-M(a)^{-1}J$, one has $L_a^TM(a)+M(a)L_a=0$. Its entire right-half-plane
spectral space is totally isotropic for $M(a)$: restriction of this identity
is a Sylvester equation with spectral sums having positive real parts.
An isotropic subspace has dimension at most ten because $M(a)$ has ten
negative directions. The count is thus exactly ten, paired nondegenerately
with the left-half-plane space; their joint inertia is $(10,10,0)$.
Their orthogonal complement has positive inertia $(87,0,0)$. There the
matrix is skew-adjoint in a positive norm, so all remaining roots are
imaginary and semisimple. The kernel has dimension one because
$\operatorname{rank}J=106$. This proves the claims for the oscillatory and
bounded groups as well, without assuming simple fast eigenvalues.
\end{proof}

Since $\theta=t/a$, these three groups correspond to the physical comparison
scales $t/a^3$, $t/a^2$, and $t/a$. The first growing sector has nine
coordinates, but a mixed-rate lift must also account for the additional
intermediate growing direction. No claim about a nonlinear native invariant
codimension follows solely from this linear count.

\section{An evaluated intermediate pencil at the actual root}
\label{v53:sec:certificate}

\begin{proposition}[Exact interval contraction of the retained source tensors]
\label[proposition]{v53:prop:certificate}
For the frozen pencil \eqref{v53:eq:pencil}, the new two-dimensional
coefficients have rational centres approximately
\begin{align}
 D_0&\simeq\begin{pmatrix}-45182.81153411&2984.84516621\\
 2984.84516621&624946.17831657\end{pmatrix},\notag\\[-1.8ex]
 G_0&\simeq\begin{pmatrix}683539.61284689&-763295.68459330\\
 -763295.68459330&1165762.85383392\end{pmatrix},\notag\\[-1.8ex]
 S_0&\simeq\begin{pmatrix}0&-65118.95947116\\65118.95947116&0\end{pmatrix}.
 \label{v53:eq:small-matrices}
\end{align}
Uniform entrywise radii are respectively less than
$0.00000511$, $0.002699$, and $0.000279$. The intermediate coefficients obey
\begin{equation}
 \boxed{3.756<\sigma_0<3.757,\qquad 0.733<\omega_0<0.734.}
 \label{v53:eq:middle-brackets}
\end{equation}
These bracket the scaled limits, not eigenvalues of the full nonlinear
original Jacobian at a finite amplitude.
\end{proposition}
\begin{proof}
The verifier reconstructs $A,L,C_W,K,\mathscr R$ from the replayed parent
source banks and applies the exact block inverses above. A rational-centre,
common-radius matrix product encloses all simultaneous root uncertainty.
The resulting radii are below $5.107\cdot10^{-6}$,
$0.002699$, and $0.000279$. It separately verifies $-\det D_0>0$,
$(G_0)_{11}>0$, and $\det G_0>0$ through rational interval comparisons.
The centres of $S_0$ are exactly skew. Evaluating
\eqref{v53:eq:quartic} at the rational real endpoints $3.756,3.757$
gives strictly positive and negative intervals. Evaluation at
$\nu=i\omega$ for $\omega=0.733,0.734$ gives the same two signs.
There is exactly one positive root in each case by the quadratic sign
argument, proving both enclosures. Floating values only propose endpoints;
no floating root, inverse or rank is used to certify them.
\end{proof}

\begin{proposition}[The omitted first skew coefficient matters at intermediate scale]
\label[proposition]{v53:prop:completion}
Replace the comparison by any specified analytic skew completion
$J(a)=J+aJ_1+O(a^2)$ and by a symmetric graded mass
$D_a(\Gamma_*^{\rm ord}+O(a))D_a$ with the same leading source.
The three spectral counts and the positive bounded spectral reduction
persist for sufficiently small $a$. The intermediate pencil has the same
$D_0,G_0$, but its skew coefficient becomes
\begin{equation}
 S_0^{\rm comp}=S_0+Z^T(J_1)_{gg}Z.
 \label{v53:eq:gyro-completion}
\end{equation}
For every such skew coefficient its positive exponent satisfies
$\sigma^2\ge\overline\lambda_+>13.96$, hence $\sigma>3.736$.
The numerical bracket \eqref{v53:eq:middle-brackets} is not asserted for an
uncomputed completion.
\end{proposition}
\begin{proof}
The same two Schur expansions show that only the order-$a$ nonconstant
weak skew block enters the intermediate kernel equation; all other new terms
are higher order there. On the fast and finite contours the leading limits
are unchanged. Rank 106 of an odd real skew $J(a)$ persists because a
nonzero 106-minor persists and odd skew rank is at most 106.
The preceding inertia proof then applies.
For a real nullvector of $\sigma^2D_0+\sigma S_0^{\rm comp}+G_0$, pairing
with that vector removes the skew term and gives
$\sigma^2=-v^TG_0v/(v^TD_0v)$ with negative denominator.
The minimum on this negative cone is the positive eigenvalue of
$-D_0^{-1}G_0=-D_v^{-1}\overline G$, whose inherited bound exceeds 13.96.
\end{proof}

The comparison gyro $\chi_0$ is \emph{not} an evaluation of V48's full slow
skew coefficient $\chi$. An $O(a)$ term invisible in Version 52's descriptor
can contribute at leading intermediate order. Moreover, the original rate
Jacobian contains differentiated metric and canonical couplings and has not
been shown to fall in the skew-completion class. The proposition cannot be
used to bypass those original-source calculations.

\section{An exact positive slow affine family for the full coefficient ODE}
\label{v53:sec:lift}

There is also a positive result: every solution of the Version-52 descriptor
has a strong, uniformly controlled lift to \emph{all 107 rows} of
\eqref{v53:eq:regularized}, with appropriately selected initial coefficients.
No physical source operation is defined by this selection.

Let $r$ be the fixed nonzero kernel vector of $J$ from Version 52 and put
\begin{equation}
 \phi=r^Tf\ne0,\qquad m(a)=r^TM(a)r,\qquad
 k(a)=-\phi/m(a),\qquad
 p(a)=-J^\dagger\{f+k(a)M(a)r\}.
 \label{v53:eq:particular}
\end{equation}
Here $m(0)>0$ is the certified parent mass, so these functions are analytic
for sufficiently small $a$. The right side in braces lies in $r^\perp=
\operatorname{Ran}J$. Consequently
\begin{equation}
                  q_{\rm sec,a}(\theta)=p(a)+k(a)r\theta
 \label{v53:eq:secular-solution}
\end{equation}
solves the full equation exactly, including at $a=0$ as a descriptor solution.
The linear coefficient retains the inherited secular slope in the limit.

\begin{theorem}[A 79-dimensional affine prepared lift with uniform error]
\label[theorem]{v53:thm:lift}
There are an analytic 79-dimensional spectral subspace $\mathcal E_s(a)$
and analytic isomorphisms $U_a:\mathcal E_s(0)\to\mathcal E_s(a)$,
$U_0=I$, such that $M(a)$ is uniformly positive on these spaces in their
pulled-back fixed coordinates. For each complete descriptor solution $q$
with initial value $q^0$, put $z^0=q^0-p(0)\in\mathcal E_s(0)$ and define
\begin{equation}
 q_a(\theta)=p(a)+k(a)r\theta+
 e^{-M(a)^{-1}J\theta}U_az^0.
 \label{v53:eq:prepared-lift}
\end{equation}
It is an exact solution of \eqref{v53:eq:regularized}. For all real $\theta$
and sufficiently small $a$,
\begin{equation}
 \boxed{\|q_a(\theta)-q(\theta)\|
 \le Ca(1+|\theta|)(1+\|z^0\|).}
 \label{v53:eq:lift-error}
\end{equation}
The same estimate holds for the first comparison-time derivatives.
In particular all rate components converge strongly on fixed finite
$\theta$-intervals. The affine initial set $p(a)+\mathcal E_s(a)$ is exactly
invariant and has codimension 28 in the coefficient ODE.
\end{theorem}
\begin{proof}
Choose a fixed complex contour enclosing the entire finite descriptor
spectrum and no other limiting point. Its boundary avoids every eigenvalue.
For small $a$, the matrix $zM(a)+J$ is invertible there, uniformly in $z$.
The operators
\begin{equation}
 P_s(a)=\frac1{2\pi i}\oint(zM(a)+J)^{-1}M(a)\,dz,\qquad
 A_s^{\rm amb}(a)=\frac1{2\pi i}\oint z(zM(a)+J)^{-1}M(a)\,dz
 \label{v53:eq:projector}
\end{equation}
are analytic in $a$, including $a=0$. The resolvent identity proves the
projection identities. At $a=0$, the algebraic polynomial part of the
resolvent has zero contour integral. The range is exactly the 79-dimensional
homogeneous descriptor consistency space, by the elimination in
\cref{v53:prop:mass}. Thus $\operatorname{rank}P_s(a)=79$ for small $a$.
On this range the second integral equals the restricted generator
$-M(a)^{-1}J$ for $a>0$ and the finite descriptor generator at zero.
Take $U_a=P_s(a)|_{\mathcal E_s(0)}$; it is an isomorphism near zero.

On the homogeneous consistency space, $q^TM(0)q=x^TM_xx>0$ unless $q=0$.
Indeed a state with $x=0$ has $w\in\ker K_g$, and the once-differentiated
compatibility forces $T_x^TM_x^{-1}T_x\eta=0$, hence $w=0$.
Compactness of the unit sphere and continuity give uniform positivity
of $U_a^TM(a)U_a$. The pulled-back generator $\widehat L_s(a)$ is analytic
and skew-adjoint in this positive metric. Its propagator is therefore
uniformly bounded for every real time and all small $a$.
The integral formula also gives
$\|\widehat L_s(a)-\widehat L_s(0)\|\le Ca$.
Duhamel's identity now bounds the difference of the two propagators by
$Ca|\theta|$, with no unstable exponential. The particular solution
\eqref{v53:eq:secular-solution} is analytic in $a$, so its difference from
$a=0$ is at most $Ca(1+|\theta|)$. This proves \eqref{v53:eq:lift-error};
multiplication by the bounded slow generators proves the derivative bound.
The nullvector $r$ belongs to $\mathcal E_s(a)$, so the secular motion
preserves the same affine set. All statements follow.
\end{proof}

This theorem is stronger than convergence of an algebraic projection:
the selected solutions satisfy every finite-$a$ differential row and have
strong nonconstant-rate convergence. It is narrower than the required
native theorem: the selected linear coefficients have not been realized as
an invariant family of the full constrained nonlinear action.
It does not identify V44's codimension-nine graph with this codimension-28
coefficient subspace. Removing only growing directions and removing all
rapid oscillatory/decaying transients are different preparations.

\begin{corollary}[The prepared coefficient rates have exact original initial access]
\label[corollary]{v53:cor:initial-access}
For each fixed compatible descriptor initial value $q^0$, write the initial
value in \eqref{v53:eq:prepared-lift} as
$q_a(0)=(v_a^0,w_a^0,c_a^0)$. There are original exactly constrained initial
records, analytic in sufficiently small amplitude, with the same leading
$Y_*,\ell_*$, on which the unique normalized original rate is
\begin{equation}
 \boxed{\lambda_t(0)=aI_Av_a^0+aI_gw_a^0+I_cc_a^0.}
 \label{v53:eq:original-initial-access}
\end{equation}
Their original metric and link curvatures are $O(a)$. No original subsequent
history is identified with \eqref{v53:eq:prepared-lift}.
\end{corollary}
\begin{proof}
The new initial triple is analytic in $a$ and has a finite limit. In the
Version-52 exact initializer, retain that triple as an additional parameter
near its limit. The joint derivative of the three cubic root equations and
the 104 nonconstant divided consistency equations is the same invertible
triangular matrix used in \cref{v52:thm:initial}. The parametric analytic
implicit-function theorem therefore supplies records uniformly on a
neighborhood of this finite parameter value. Substitute the analytic
triple $q_a(0)$. Exact $P=0$ and exact consistency give
\eqref{v53:eq:original-initial-access} by the already proved signed
regularity. Positivity and the $O(a)$ initial curvature bound follow from
\cref{v52:lem:curvature}. This initial matching does not control the
unmatched original differentiated-source terms after the initial cut.
\end{proof}

\section{Why finite-order matching alone does not justify a native lift}
\label{v53:sec:errors}

\begin{proposition}[Arbitrarily high algebraic initial accuracy can fail]
\label[proposition]{v53:prop:instability}
For every fixed $p>0$ and $T>0$, there exist exact solutions of the frozen
coefficient equation whose initial values differ by $a^p$ from a prepared
solution, but whose difference at $\theta=T$ diverges as $a\to0$.
They can be chosen along the simple growing branch with
$z_+(a)=\sigma_0/a+O(1)$, or along the inherited simple fastest positive
branch with $a^2z_f(a)\to\sigma_*>0.3382576$.
\end{proposition}
\begin{proof}
Choose a real unit eigenvector $e_+(a)$ of the simple positive branch and
add $a^pe^{z_+(a)\theta}e_+(a)$ to the exact prepared solution.
This is an exact homogeneous perturbation. Its norm is $a^p$ initially
and $a^pe^{z_+(a)T}$ at $T$, which diverges. The same proof applies to the
simple fastest branch. Along either real eigenvector,
$e^TM(a)e=0$, since $L_a^TM(a)+M(a)L_a=0$. Conservation of the signed
quadratic form therefore cannot repair the missing positive error bound.
\end{proof}

This is a sensitivity theorem for the restored coefficient equation,
not an original-flow instability theorem or a lower bound for every physical
observation. It explains why a polynomially small residual or initial mismatch
cannot simply be inserted into an unweighted full inverse estimate.

There is a precise missing transfer identity. For the actual original
consistency equation, differentiation gives
\begin{equation}
 M\lambda_{tt}+K\lambda_t+D_X\mathsf B[F]
 +2D_XM[F]\lambda_t+D_\lambda M[\lambda_t]\lambda_t=0.
 \label{v53:eq:true-differentiation}
\end{equation}
The factor two follows from $D_\lambda\mathsf B=K+D_XM[F]$ and from
$\dot M\lambda_t$. In mixed-rate variables, the omitted terms and the
moving canonical state must be divided by the same row scales as
\eqref{v53:eq:pencil}; an unscaled small remainder is insufficient.
A native proof must establish an exact constrained normal form, control
its projected remainder in the growing, oscillatory and slow sectors, and
show that original initial-data controls reach its prepared set. The
linear spectral projector is an analysis tool, not an allowed replacement
update or a proof of that nonlinear transversality.

\section{Scope and verification record}
\label{v53:sec:status}

\paragraph{Proved.}
The actual leading source defines a full signed regular coefficient pencil
whose positive descriptor limit hides 28 modes. Their scales are 24 modes
at $a^{-2}$, four at $a^{-1}$, and 79 bounded modes in comparison time.
There are ten growing directions, not just the nine fastest ones. The
intermediate frozen pencil is evaluated by new exact interval contractions.
Every descriptor solution has an exact 79-dimensional affine prepared lift
to all 107 frozen differential rows, with strong $O_T(a)$ convergence and
retained secular forcing. Its initial rates have exact original constrained
initial access, without an asserted subsequent shadowing theorem. Algebraic accuracy alone does not control the
excluded positive modes.

\paragraph{Inherited and executed.}
The V52 verifier was rerun in the present workspace and passed. It replays
its parent roots and enclosures and recomputes nine original cubic harmonic
derivatives and 111 parent drift/mean derivatives. The new source worker
performs ten exact checks of the two-dimensional kernel contraction,
its signs, and the four rational spectral endpoints. Large root, stationary,
and weak-Poisson coefficient banks are inherited, not wholly regenerated.
The new finite-algebra regressions separately check the graded congruence,
Schur formulas, infinite-block count, spectral scales on declared examples,
prepared affine solutions, and failure of naive initial matching.
Exploratory floating spectra of the 107-dimensional source pencil are
explicitly non-certifying; roundoff-created real parts are not counted as
additional modes. All spectral claims above rest on the written proofs
and the named exact certificates.

\paragraph{Unpaid and next direction.}
No original nonlinear history on $[0,aT]$, no complete native remainder
estimate in \eqref{v53:eq:true-differentiation}, and no original invariant family corresponding to
the full coefficient-level slow subspace is claimed. The frozen gyro
is not the uncomputed complete gyro. The next task is the projected
original remainder and initial-access calculation at the two distinct
fast scales, using this spectral decomposition rather than positivity of
the limiting descriptor alone. The previously proved shorter selected
histories, V32's independent stress problem, and the lack of a fixed
positive physical lifespan or cofinal Einstein--SM limit are unchanged.

Every V52 byte before its final document command is preserved.
The comparison and its prepared lift do not modify the microscopic law.
They are a resolved intermediate test for the native lift, not that lift.

\begingroup
\renewcommand{\refname}{Additional references for Version 53}
\begin{thebibliography}{9}
\small\raggedright
\bibitem{v53:TM}
F.~Tisseur and K.~Meerbergen, \emph{The quadratic eigenvalue problem},
SIAM Review \textbf{43} (2001), 235--286.
\url{https://doi.org/10.1137/S0036144500381988}.
Primary-source context for structured quadratic pencils, not a source of
the present numerical coefficients or native-flow bounds.
\bibitem{v53:MMW}
C.~Mehl, V.~Mehrmann and M.~Wojtylak,
\emph{Linear algebra properties of dissipative Hamiltonian descriptor systems},
arXiv:1801.02214 (2018).
\url{https://arxiv.org/abs/1801.02214}.
The positive descriptor context is distinguished from the signed
finite-amplitude pencil treated here.
\end{thebibliography}
\endgroup
\addtocontents{toc}{\protect\endgroup}
% END IMMUTABLE VERSION 53 BODY
% APPEND VERSION 54 HERE, BEFORE THE FINAL END-DOCUMENT COMMAND.


% BEGIN IMMUTABLE VERSION 54 BODY
\clearpage
\begin{center}
{\Large\bfseries Version 54: Time-dependent source work before a nonlinear lift}\\[.5em]
{\large The two differentiated weak-source coordinates, a two-scale Green estimate,\\
and exact prepared response to nonconstant covectors}
\end{center}
\makeatletter
\addtocontents{toc}{\protect\begingroup}
\addtocontents{toc}{\protect\makeatletter}
\addtocontents{toc}{\protect\def\protect\l@section{\protect\bfseries\protect\@dottedtocline{1}{0em}{2.5em}}}
\makeatother
\addcontentsline{toc}{part}{Version 54: Time-dependent source work and prepared response}

\section{The next missing estimate is a covector response, not a positive generator bound}
\label{v54:sec:scope}

Version 53 constructs an exact prepared lift for the constant source in its
full graded coefficient pencil.  Its final equation identifies the actual
nonlinear remainder but does not bound its response.  Repeating that constant-source
construction would not settle a moving source.  In particular, the inverse of
the small mass cannot simply multiply a small residual: it contains two different
fast scales and an index-two limiting constraint.

This body retains exactly $M(a),J$ of \cref{v53:eq:pencil}, at the same algebraic
mixed-mode root and $h=1/3$.  We now study
\begin{equation}
                  M(a)\dot q+Jq=g(\theta),\qquad 0\le\theta\le T,
\label{v54:eq:forced}
\end{equation}
where a dot denotes $d/d\theta$ and $T<\infty$ is fixed.  The old source is
$g=-f$.  An additional $g+f$ is an explicitly declared \emph{comparison
covector}; it is not an added term in the original action.  The results apply
to real sources and are proved in their complexifications when contours are used.
All norms below are fixed coordinate norms unless another norm is specified.
No constant is asserted uniform in the spatial cutoff or as $T\to\infty$.

The new conclusions are a rank-two differentiated-source map, its numerical
coefficient enclosure at the actual root, and a prepared source-response theorem
which pays both discarded fast scales.  Three continuous source derivatives
suffice for uniform $O(a)$ state convergence; four suffice for an analogous
first-derivative assertion with the stated additional preparation.  The proof
uses stable initial conditions, oscillatory initial conditions, and terminal
conditions on growing corrections.  These are conditions defining a comparison
solution, not observations, filtering, or repeated updates of a native history.

\paragraph{Audit and inherited inputs.}
The retained inputs are the ranks, definite oscillatory sectors, simple
intermediate roots, analytic finite spectral projection, and exact original
initial-rate access proved in Version 53.  The complete original differentiated
consistency equation is also inherited.  A targeted review of the preceding
native body and the supplied Einstein V36 and perturbative V51 sources did not
locate the nonconstant-covector response proved here.  This is not an independent
revalidation of those cumulative manuscripts.  Polynomial parts of index-two
transfer functions and structured-pencil methods are established subjects
\cite{v54:BGMehr,v54:Tisseur}; the original source contraction and specialized
two-scale estimates below are derived explicitly.

\section{The descriptor retains two differentiated weak-source coordinates}
\label{v54:sec:descriptor}

Use the reordered variables of \cref{v53:eq:blocks}:
$x=(v,c)\in\R^{81}$ and $w\in\R^{26}$.  Split $g=(g_x,g_g)$ in this order.
At $a=0$ the complete descriptor is
\begin{equation}
 M_x\dot x+J_xx+Ew=g_x,\qquad -E^Tx+K_gw=g_g.
\label{v54:eq:descriptor}
\end{equation}
The three harmonic shifts remain inside $x$; they are not discarded.
All formulas in this reordered section are transported back to $(v,w,c)$
by the same fixed permutation when physical row scalings are used.
A classical descriptor solution means $x\in C^1$ and $w\in C^0$;
assertions about $\dot w$ use the additional source regularity stated below.
Let $Z:\R^2\to\ker K_g$ be the retained full-rank kernel synthesis.  Set
\begin{equation}
\begin{gathered}
 R_g=K_g^\dagger,\quad T_0=EZ,\quad G_0=T_0^TM_x^{-1}T_0\succ0,\\
 \mathcal F=M_x^{-1}T_0G_0^{-1}Z^T,\\
 \mathcal P=M_x^{-1}-M_x^{-1}T_0G_0^{-1}T_0^TM_x^{-1},\\
 K_e=J_x+ER_gE^T,\qquad
 W_0=R_gE^T-ZG_0^{-1}T_0^TM_x^{-1}K_e.
\end{gathered}
\label{v54:eq:source-matrices}
\end{equation}
Here $\mathcal F$ is a source map, not the canonical vector field of earlier versions.

\begin{theorem}[Exact time-dependent constraint reconstruction]
\label[theorem]{v54:thm:descriptor}
For a continuously differentiable source, every classical solution of
\eqref{v54:eq:descriptor} has $x=y-\mathcal Fg_g$, with $T_0^Ty=0$, and satisfies
\begin{align}
 \dot y&=-\mathcal P K_e y+
 \mathcal P\{g_x-ER_gg_g+K_e\mathcal Fg_g\},\label{v54:eq:tangent}\\
 w&=W_0y+R_gg_g-W_0\mathcal Fg_g
 +ZG_0^{-1}T_0^TM_x^{-1}(g_x-ER_gg_g)
 +ZG_0^{-1}Z^T\dot g_g.\label{v54:eq:weak-reader}
\end{align}
Conversely these formulas solve every original descriptor row.  Arbitrary
initial $y(0)\in\ker T_0^T$ gives one such solution.  The coefficient of the
source derivative in the complete reconstructed state is
\begin{equation}
 \boxed{\quad
 \mathcal N_1=
 \begin{pmatrix}0_{81}&0\\0&ZG_0^{-1}Z^T\end{pmatrix},\qquad
 \mathcal N_1\succeq0,\quad\operatorname{rank}\mathcal N_1=2.
 \quad}
\label{v54:eq:N1}
\end{equation}
It depends on exactly the two covector coordinates $Z^T\dot g_g$.
The displayed map is unchanged under $Z\mapsto ZB$ for any invertible
$2\times2$ matrix $B$.
\end{theorem}
\begin{proof}
Pair the algebraic row with $Z$ to obtain
$T_0^Tx=-Z^Tg_g$.  Since $T_0^T\mathcal F=Z^T$, this is exactly
$T_0^Ty=0$.  The algebraic solution is
$w=R_g(g_g+E^Tx)+Z\eta$.  Substitution into the differential row and one
differentiation of the compatibility condition give
\[
 G_0\eta=T_0^TM_x^{-1}(g_x-K_ex-ER_gg_g)+Z^T\dot g_g.
\]
Solving this positive two-dimensional system and substituting
$x=y-\mathcal Fg_g$ proves \eqref{v54:eq:weak-reader}.  The differential row becomes
$\dot x=-\mathcal P K_ex+\mathcal P(g_x-ER_gg_g)-\mathcal F\dot g_g$,
which is \eqref{v54:eq:tangent}.  These substitutions are reversible,
and $T_0^T\mathcal P=0$ preserves the tangent condition.
Finally $Z$ is injective and $G_0\succ0$.  Congruence proves positivity
and rank two.  Replacing $Z$ by $ZB$ replaces $G_0$ by $B^TG_0B$ and leaves
$ZG_0^{-1}Z^T$ unchanged.
\end{proof}

There are two derivative-sensitive \emph{coordinates}, not two time
\emph{orders}: only the first time derivative appears in
\eqref{v54:eq:weak-reader}.  The source map's positivity is not positivity
of the restored signed mass $M(a)$.

For later use write the reconstruction as
\begin{equation}
 q_0=\mathcal Q y+\mathcal N_0g+\mathcal N_1\dot g,
 \qquad \mathcal Qy=(y,W_0y).
\label{v54:eq:feedthrough}
\end{equation}
The constant source matrix $\mathcal N_0$ is given explicitly by
\eqref{v54:eq:weak-reader} and $x=y-\mathcal Fg_g$; in block form,
\[
\mathcal N_0=
\begin{pmatrix}
0&-\mathcal F\\
ZG_0^{-1}T_0^TM_x^{-1}&R_g-W_0\mathcal F-ZG_0^{-1}T_0^TM_x^{-1}ER_g
\end{pmatrix}.
\]
Thus descriptor initial consistency for a moving source depends on
$g(0)$ and $Z^T\dot g_g(0)$, not just its instantaneous value.

\section{The actual coefficient and the failure of a source-amplitude test}
\label{v54:sec:certificate}

By \cref{v53:lem:identification}, the matrix in
\eqref{v54:eq:source-matrices} is exactly $G_0=\overline G$, not a new
adjustable comparison metric.

\begin{proposition}[Certified derivative-source coefficient]
\label[proposition]{v54:prop:inverse}
Throughout the inherited $2^{-112}$ root cube,
\begin{equation}
\boxed{
124179<\lambda_{\min}(G_0)<124180,\qquad
1725123<\lambda_{\max}(G_0)<1725124.
}
\label{v54:eq:G-spectrum}
\end{equation}
The exact rational proposal $P_G$ stored in the package has approximate value
\begin{equation}
 P_G\simeq10^{-6}
 \begin{pmatrix}5.44177388&3.56305959\\3.56305959&3.19075874\end{pmatrix},
 \qquad
 \boxed{\ \|G_0^{-1}-P_G\|_\infty<6\,10^{-13},\quad
 \|G_0^{-1}\|_\infty<\frac1{110000}.\ }
\label{v54:eq:G-inverse}
\end{equation}
In the stated kernel coordinates, the derivative contribution
$\eta_{\rm der}=G_0^{-1}Z^T\dot g_g$ consequently obeys
$\|\eta_{\rm der}\|_\infty<\|Z^T\dot g_g\|_\infty/110000$.
The physical weak vector is $Z\eta_{\rm der}$, so a bound on that vector
must also retain the norm of its synthesis $Z$.
\end{proposition}
\begin{proof}
The new verifier rounds the inherited exact stiffness centre to denominator
$2^{32}$ and proves that a common entry radius $1/300$ still contains the
entire original enclosure.  It uses exact rational determinant bounds for
$G_0-\lambda I$ at $124179,124180,1725123,1725124$, together with the
corresponding first principal entries.  Sylvester's criterion and the
intermediate determinant signs prove \eqref{v54:eq:G-spectrum}.
For the inverse, let $G_c$ be that centre, $r=1/300$, and let $P_G$ be the
rounded rational inverse proposal with denominator $2^{90}$.  The checks give
\[
 q:=\|I-P_GG_c\|_\infty+2r\|P_G\|_\infty<10^{-7}.
\]
The Neumann formula yields
$\|G_0^{-1}\|_\infty\le\|P_G\|_\infty/(1-q)$ and
$\|G_0^{-1}-P_G\|_\infty\le\|P_G\|_\infty q/(1-q)$.
Exact comparisons give the stated bounds.  This is a new small contraction
of the inherited source bank, not regeneration of its primitive entries.
\end{proof}

\begin{example}[Small covector amplitude does not control a differentiated rate]
\label[example]{v54:ex:resonance}
Consider the positive-mass comparison
\[
 \dot x+w=0,\qquad a^2\dot w-x=a\sin(\theta/a).
\]
Its descriptor has $x_0=-a\sin(\theta/a)$ and $w_0=\cos(\theta/a)$,
although the source has uniform norm $a$.  The exact regularized solution
with $(x(0),w(0))=(0,1)$ is
\begin{align*}
 x(\theta)&=\tfrac\theta2\cos(\theta/a)-\tfrac{3a}{2}\sin(\theta/a),\\
 w(\theta)&=\cos(\theta/a)+\tfrac\theta{2a}\sin(\theta/a).
\end{align*}
Both original rows hold by differentiation.  On every fixed positive
interval the supremum of $|w|$ is unbounded as $a\to0$, despite exact
agreement with the descriptor initial values and their initial derivatives.
This example contains no growing homogeneous eigenvalue.  It is a declared
comparison illustrating source resonance, not the actual Einstein source.
\end{example}

This distinguishes source amplitude from source variation.  Even after
unstable homogeneous errors have been selected away, an unweighted
$C^0$ residual estimate cannot justify the desired lift.  The three-derivative
hypothesis below is sufficient, not claimed minimal.

\section{Source moments and Green estimates at both fast scales}
\label{v54:sec:green}

For $a>0$ set $L_a=-M(a)^{-1}J$ and let $P_s(a)$ be the analytic finite
projection from \cref{v53:eq:projector}.  Let $P_f=I-P_s$ and
$B_s=P_sM(a)^{-1}$.  The letter $f$ on a projection means the full 28-dimensional
fast space; it does not denote the old constant source.  The fast generator
$L_f=L_a|_{\Ran P_f}$ is invertible.

\begin{lemma}[Bounded covector moments without a full mass-inverse loss]
\label[lemma]{v54:lem:moments}
The input map $B_s(a)$ is analytic and bounded through $a=0$.  For $j\ge0$,
put
\begin{equation}
 F_j(a)=-L_f^{-(j+1)}P_fM(a)^{-1}.
\label{v54:eq:moments}
\end{equation}
Every fixed $F_j$ is analytic through zero, and
\begin{equation}
 \boxed{F_0(0)=\mathcal N_0,\qquad F_1(0)=\mathcal N_1,\qquad
 F_j(0)=0\ (j\ge2).}
\label{v54:eq:moment-limits}
\end{equation}
In particular $\|F_0\|+\|F_1\|\le C$ and $\|F_2\|+\|F_3\|\le Ca$.
\end{lemma}
\begin{proof}
Choose a fixed positively oriented circle enclosing all finite descriptor
poles and meeting none, as in Version 53.  Write
$R_a(z)=(zM(a)+J)^{-1}$.  On this contour the inverse is uniformly analytic
in $a$, including zero.  Therefore
$B_s=(2\pi i)^{-1}\oint R_a(z)\,dz$ has the asserted extension.
For $z$ in its interior, the fast part of the source resolvent is
\[
 R_{f,a}(z)=\frac1{2\pi i}\oint\frac{R_a(\zeta)}{\zeta-z}\,d\zeta.
\]
The residue of every interior slow pole cancels in this expression; fast
poles lie outside.  At positive $a$ it equals
$(z-L_f)^{-1}P_fM(a)^{-1}$.  Its Taylor coefficients at zero are
precisely \eqref{v54:eq:moments}.  They are analytic in $a$.

The direct elimination in \cref{v54:thm:descriptor} gives a strictly proper
finite transfer and the polynomial part $\mathcal N_0+z\mathcal N_1$.
For example $M(0)\mathcal N_1=0$ and
$M(0)\mathcal N_0+J\mathcal N_1=0$ follow directly from the displayed
matrices.  Its finite state is $\mathcal Qy$, with equation
\eqref{v54:eq:tangent}.  Thus the same contour at $a=0$ returns exactly
that polynomial part.  This proves \eqref{v54:eq:moment-limits}; Taylor's
theorem in the parameter gives the final bounds.
\end{proof}

Split $P_f=P_++P_-+P_o$ into the ten growing, ten decaying and eight
rapidly oscillatory directions from Version 53.  All these projectors
commute with $L_a$; their separate coordinate norms are \emph{not} presumed
uniform.  In the next statement they always occur with the source moment.

\begin{lemma}[Two-scale covector Green bounds]
\label[lemma]{v54:lem:green}
There are constants $C,\kappa>0$ and a sufficiently small $a_0>0$ such
that, for $0<a<a_0$, $r\ge0$ and $m=3,4$,
\begin{align}
 \|e^{-rL_a}P_+L_f^{-m}P_fM(a)^{-1}\|
 +\|e^{rL_a}P_-L_f^{-m}P_fM(a)^{-1}\|
 &\le C\{a^{2m-2}e^{-\kappa r/a^2}
                    +a^{m-2}e^{-\kappa r/a}\},\label{v54:eq:hyper-green}\\
 \|e^{rL_a}P_oL_f^{-m}P_fM(a)^{-1}\|
 &\le C(a^{2m-2}+a^{m-2}).\label{v54:eq:osc-green}
\end{align}
The oscillatory estimate also holds for negative $r$.
\end{lemma}
\begin{proof}
We provide the scaling argument, because applying only
$\|M(a)^{-1}\|=O(a^{-2})$ would lose the conclusion.
Use Version 53's exact mass shear \eqref{v53:eq:sheared-pencil}; the shear
and its inverse remain bounded.  On a fixed contour $\mu$ around a nonzero
fastest cluster, set $z=\mu/a^2$.  The upper inverse is $O(a^2)$ and the
lower Schur complement tends to $\mu\Delta+K_g$, invertibly on that contour.
Hence $\|R_a(\mu/a^2)\|\le C$.  The contour formula for
$e^{rL_a}P L_f^{-m}P_fM(a)^{-1}$ has integrand
$e^{rz}z^{-m}R_a(z)$.  Its powers contribute $a^{2m}$ and its contour
measure contributes $a^{-2}$.  Contours within the left half-plane give
the first decaying term in \eqref{v54:eq:hyper-green}; contours within
the right half-plane give the corresponding backward estimate.  Nontrivial
Jordan blocks in these separated hyperbolic clusters cause no loss in this
contour argument.

For the intermediate clusters set $z=\nu/a$.  The upper inverse is
$O(a)$.  On the 24-dimensional complement of $\ker K_g$, the lower block
is uniformly invertible.  Its remaining two-dimensional Schur complement is
\[
 a\{\nu D_0+S_0+\nu^{-1}G_0\}+O(a^2).
\]
Away from the four nonzero roots, its inverse is $O(a^{-1})$.  The complete
resolvent is therefore $O(a^{-1})$ on a fixed $\nu$ contour.  Its measure
contributes another $a^{-1}$, whereas $z^{-m}$ contributes $a^m$.
This gives $a^{m-2}$.  The four roots are simple.  Each residue, including
the source-to-state normalization, consequently has this bound.  The real
pair yields the second exponential in \eqref{v54:eq:hyper-green}; the
imaginary pair gives its part of \eqref{v54:eq:osc-green}.

For completeness the six fastest imaginary modes need not have distinct
frequencies.  Their state projection, computed from
$R_a(z)M(a)$ on the $\mu$ contour, is bounded and has a limiting range
consisting of the corresponding pure weak vectors.  Restricted to that
range, $a^{-2}M(a)$ converges to the definite oscillatory weak metric used
in Version 53.  It remains uniformly definite in bounded transported
coordinates.  The rescaled generator is skew-adjoint in this metric, so
its propagator is uniformly bounded for all real times; its inverse powers
are bounded because the frequencies stay away from zero.  The same source
contour gives an $O(a^{-2})$ input norm and hence the factor $a^{2m-2}$.
This proves the remaining oscillatory estimate without discarding any mode
or assuming decay of an oscillation.
\end{proof}

\section{A prepared lift for arbitrary smooth time-dependent covectors}
\label{v54:sec:lift}

Let $q_0$ be the descriptor solution of \cref{v54:thm:descriptor}, with
chosen tangent initial data.  Its homogeneous finite component is denoted
$q_s^0=P_s(0)q_0(0)$.  Let
$U_a=P_s(a)|_{\Ran P_s(0)}$ be the analytic isomorphism used in Version 53.

\begin{theorem}[Exact prepared response, with both fast scales paid]
\label[theorem]{v54:thm:lift}
Fix $T<\infty$ and $g\in C^3([0,T];\R^{107})$.  For all sufficiently
small $a>0$ there is a unique solution of \eqref{v54:eq:forced} with the
following linear preparation conditions.  Put
\[
                 q_{f,\mathrm{app}}(\theta)
                  =F_0g+F_1\dot g+F_2\ddot g.
\]
The conditions are
\begin{align*}
 P_s(a)q_a(0)&=U_aq_s^0,\\
 (P_-+P_o)(q_a(0)-q_{f,\mathrm{app}}(0))&=0,\\
 P_+(q_a(T)-q_{f,\mathrm{app}}(T))&=0.
\end{align*}
They consist of $79+10+8+10=107$ scalar conditions.  The solution obeys
\begin{equation}
 \boxed{
 \sup_{0\le\theta\le T}\|q_a(\theta)-q_0(\theta)\|
 \le C_Ta\bigl(\|q_s^0\|+\|g\|_{C^3}\bigr).
 }
\label{v54:eq:lift}
\end{equation}
The map from $(q_s^0,g)$ to this solution is linear and uniformly bounded
from $\R^{79}\times C^3$ to $C^0$, after the stated identifications.
With $g\in C^4$, using one additional term $F_3g^{(3)}$ in the preparation
gives the same $O(a)$ estimate in $C^1$ against the descriptor solution.
\end{theorem}
\begin{proof}
First solve the finite component exactly:
\[
 q_{s,a}(\theta)=e^{\theta L_s(a)}U_aq_s^0+
       \int_0^\theta e^{(\theta-r)L_s(a)}B_s(a)g(r)\,dr.
\]
The analytic source map and the uniformly positive finite metric from
Version 53 imply that this converges, together with its first derivative,
at order $a$ on $[0,T]$.  This includes the forced finite zero mode; it
is not subtracted.  By \cref{v54:lem:moments},
$q_{f,\mathrm{app}}=\mathcal N_0g+\mathcal N_1\dot g+O(a\|g\|_{C^2})$.

The moment identities are $L_fF_0=-P_fM^{-1}$ and
$L_fF_j=F_{j-1}$ for $j\ge1$.  Hence the exact fast defect is
\[
 q_{f,\mathrm{app}}'-L_fq_{f,\mathrm{app}}-P_fM^{-1}g
                       =F_2g^{(3)}.
\]
It is canceled by the explicit correction
\begin{align}
 e_a(\theta)={}&-\int_0^\theta e^{(\theta-r)L_a}(P_-+P_o)F_2g^{(3)}(r)\,dr
 \\&+\int_\theta^T e^{(\theta-r)L_a}P_+F_2g^{(3)}(r)\,dr.
\label{v54:eq:correction}
\end{align}
Differentiation proves $e_a'-L_ae_a=-F_2g^{(3)}$, with exactly the two
fast boundary conditions.  The Green bounds with $m=3$ give
\[
 \|e_a\|_{C^0}\le
 C\{a^6+a^2+T(a^4+a)\}\,\|g^{(3)}\|_\infty
 \le C_Ta\|g\|_{C^3}.
\]
Thus $q_a=q_{s,a}+q_{f,\mathrm{app}}+e_a$ solves all 107 original
\emph{coefficient} rows and satisfies \eqref{v54:eq:lift}.  A homogeneous
difference satisfying the same conditions is zero on each invariant sector,
which proves uniqueness.  Linearity and the uniform bound follow directly.

For $C^4$ sources use $\sum_{j=0}^3F_jg^{(j)}$.  The defect is $F_3g^{(4)}$.
The same construction with $m=4$ bounds its correction by $C_Ta^2\|g\|_{C^4}$.
To differentiate that correction, multiply its kernels by $L_a$; this reduces
$m=4$ to $m=3$, and the endpoint terms are $F_3g^{(4)}=O(a)\|g\|_{C^4}$.
Thus its derivative is $O_T(a)$.  The additional approximate term and its
derivative are also $O(a)$ by the moment limits.  This proves the assertion.
\end{proof}

\begin{corollary}[A usable norm for a nonconstant comparison defect]
\label[corollary]{v54:cor:source-error}
Let $g_a=g_0+r_a$ with uniformly bounded $C^3$ norms, and give the prepared
solutions corresponding finite initial data.  Their error relative to the
$g_0$ descriptor is bounded by
\begin{equation}
 \boxed{C_T\{a(\|q_s^0\|+\|g_0\|_{C^3})+
             \|r_a\|_{C^3}+\|q_{s,a}^0-U_aq_s^0\|\}.}
\label{v54:eq:source-error}
\end{equation}
In particular $r_a=O(a)$ in this covector norm preserves an $O(a)$ rate
error.  No $a^{-2}$ multiplier from the full mass inverse and no
$e^{cT/a}$ or $e^{cT/a^2}$ appears in this prepared estimate.
\end{corollary}
\begin{proof}
Apply linearity and the uniform operator bound of the theorem to the
source and finite-initial-data differences, then apply
\eqref{v54:eq:lift} to $g_0$.
\end{proof}

The growing boundary condition uses the whole specified source segment.
It determines one initial preparation; it is not a causal decoder from a
single unknown sample, a future measurement, or feedback executed on a
history.  The preparation depends on $T$ and is not asserted invariant
under extending that horizon.  Keeping the growing initial error instead
would restore the exponential amplification exhibited in Version 53.
Rapid oscillatory modes have also been prepared, not assumed to decay.

\section{Exact original initial access and the remaining nonlinear source}
\label{v54:sec:native}

\begin{corollary}[The new initial rates are accessible, but not their future paths]
\label[corollary]{v54:cor:initial}
For a fixed source and descriptor initial value in
\cref{v54:thm:lift}, the bounded triple $q_a(0)=(v_a^0,w_a^0,c_a^0)$ has
exact original initial records with leading fields $Y_*,\ell_*$, for all
sufficiently small $a>0$.  Their original constraints and consistency hold,
and their unique original initial rate is
\[
             \lambda_t(0)=aI_Av_a^0+aI_gw_a^0+I_cc_a^0.
\]
Their initial metric and literal link curvatures are $O(a)$.
This assertion gives no subsequent original shadowing of $q_a$.
\end{corollary}
\begin{proof}
The prepared triples converge to the finite descriptor initial triple by
\eqref{v54:eq:lift}.  The parametric initializer in
\cref{v53:cor:initial-access}, inherited from Version 52, is analytic on a
neighborhood of any such finite triple, with the same invertible constraint
and consistency derivative.  Use that neighborhood uniformly for the
convergent triples.  Exact initialization and the signed Hessian give the
rate, and the inherited initial-curvature statement applies.  The dependence
of this new triple on $a$ need not be analytic because of the time integrals
in \eqref{v54:eq:correction}; no stronger amplitude regularity is asserted.
\end{proof}

For a putative original trajectory define the fixed rate synthesis and row
scaling
\[
 S_a=(aI_A,aI_g,I_c),\qquad
 R_a^{\rm row}=
 \begin{pmatrix}a^{-2}I_A^T\\a^{-2}I_g^T\\a^{-3}I_c^T\end{pmatrix},
 \qquad \lambda_t=S_aq,\quad t=a\theta.
\]
Use $M^{\rm orig},K^{\rm orig},\mathsf B,F$ for the actual moving original
objects.  The following quantities retain the raw covector coordinates:
\begin{align}
 \widetilde M_a&=a^{-1}R_a^{\rm row}M^{\rm orig}S_a,
 &\widetilde J_a&=R_a^{\rm row}K^{\rm orig}S_a,\notag\\
 \widetilde f_a&=R_a^{\rm row}\bigl\{
 D_X\mathsf B[F]+2D_XM^{\rm orig}[F]S_aq
 +D_\lambda M^{\rm orig}[S_aq]S_aq\bigr\}.
\label{v54:eq:original-divided}
\end{align}
The differentiated original consistency equation is exactly
$\widetilde M_a\dot q+\widetilde J_aq+\widetilde f_a=0$.
Hence the precise comparison defect in \eqref{v54:eq:forced} is
\begin{equation}
 \boxed{
 g=-f+\mathscr R_a,\qquad
 \mathscr R_a=(M(a)-\widetilde M_a)\dot q
             +(J-\widetilde J_a)q+f-\widetilde f_a.
 }
\label{v54:eq:native-defect}
\end{equation}
This is an identity, not an estimate.  It retains the mass defect multiplying
$\dot q$, the moving canonical state, the complete higher Poisson terms,
and both metric-derivative contributions.  In particular it does not
replace the actual native intermediate gyro by the frozen $S_0$.

The result identifies a concrete new source requirement:
$\mathscr R_a$ must be estimated after the displayed row scaling, and the
two contractions $Z^T(\mathscr R_a)_g$ have a differentiated effect already
at descriptor order.  A $C^3$ covector bound is one sufficient comparison
input; it has \emph{not} been proved for the actual nonlinear remainder.
Evaluating \eqref{v54:eq:native-defect} only along a proposed smooth curve
would also not prove it along a nearby exact trajectory.  The mass-defect
term requires control of that difference or a closed nonlinear normal form.

Likewise exact initial-rate access does not prove that the original
initializer reaches the boundary-prepared set of the nonlinear problem.
The ten growing terminal conditions, the eight oscillatory preparation
conditions, the stable transients, and the original constraints must be
handled together.  No terminal condition is imposed as an operational
update.  The present theorem gives the source-response estimate to be used
after such a construction, not a replacement for it.

\section{Status and reproducible checks}
\label{v54:sec:status}

\paragraph{Proved.}
The actual positive descriptor has a rank-two differentiated weak-source
map with an explicitly certified coefficient.  The restored signed
coefficient pencil admits bounded covector source moments and a Green
estimate paying both its $a^{-2}$ and $a^{-1}$ fast scales.  Every $C^3$
source and consistent descriptor initial value has an exact full-coefficient
prepared solution converging at order $a$ on a fixed time interval; an
additional derivative and source moment give a first-derivative estimate.
These initial rates are realized by exact original constrained initial
records.  None of these statements supplies a nonlinear original history
on $[0,aT]$.

\paragraph{Inherited and executed.}
The Version-53 standalone verifier was rerun successfully.  Its parent
replay regenerates nine cubic harmonic derivatives and 111 differentiated
source entries; its ten source contractions and forty algebraic checks pass.
The new verifier makes ten exact checks of the differentiated-source
coefficient and inverse enclosure, and thirty-five exact finite-algebra
checks of the descriptor reconstruction, polynomial transfer identities,
source moments, two-scale examples, resonance, and boundary correction.
The source matrix's large bank is inherited, not regenerated by the new
$2\times2$ contraction.  Example equations are explicitly distinguished
from the native source.  The Green bounds and original initial access are
written analytical proofs, not a nonlinear spacetime integration or
proof-assistant formalization.

\paragraph{Unpaid and next direction.}
The missing native work is now the actual projected and differentiated
remainder in \eqref{v54:eq:native-defect}, together with nonlinear initial
transversality to the prepared boundary set.  Small unscaled residuals,
positive descriptor energy, or another constant-source lift do not supply
those facts.  The known shorter original histories, lack of a fixed
positive physical lifespan or spatially cofinal Einstein--SM limit, and
V32's independent stress-compensation branch remain unchanged.  No primitive
selection law, attraction to the prepared set, or uniformly causal source
reader is asserted.

Every Version-53 byte before its final document command is preserved.
New labels use the prefix \texttt{v54:}.

\begingroup
\renewcommand{\refname}{Additional references for Version 54}
\begin{thebibliography}{9}
\small\raggedright
\bibitem{v54:BGMehr}
C.~A.~Beattie, S.~Gugercin and V.~Mehrmann,
\emph{Structure-preserving interpolatory model reduction for port-Hamiltonian
differential-algebraic systems}, arXiv:1910.05674 (2019).
\url{https://arxiv.org/abs/1910.05674}.
Primary-source context for polynomial transfer parts in index-two systems,
not a theorem about the present signed finite-amplitude family.
\bibitem{v54:Tisseur}
F.~Tisseur and K.~Meerbergen,
\emph{The quadratic eigenvalue problem}, SIAM Review \textbf{43} (2001),
235--286. \url{https://doi.org/10.1137/S0036144500381988}.
Structured-pencil context; the source maps and response bounds are proved above.
\end{thebibliography}
\endgroup
\addtocontents{toc}{\protect\endgroup}
% END IMMUTABLE VERSION 54 BODY
% APPEND VERSION 55 HERE, BEFORE THE FINAL END-DOCUMENT COMMAND.


% BEGIN IMMUTABLE VERSION 55 BODY
\clearpage
\begin{center}
{\Large\bfseries Version 55: Remove the moving mass before estimating the nonlinear source}\\[.5em]
{\large An exact graded congruence, the retained metric connection,\\
and an original-equivalent remainder with no rate derivative}
\end{center}
\makeatletter
\addtocontents{toc}{\protect\begingroup}
\addtocontents{toc}{\protect\makeatletter}
\addtocontents{toc}{\protect\def\protect\l@section{\protect\bfseries\protect\@dottedtocline{1}{0em}{2.5em}}}
\makeatother
\addcontentsline{toc}{part}{Version 55: Exact moving-mass reduction of the original equations}

\section{Upstream target: the derivative in the nonlinear mass defect}
\label{v55:sec:scope}

Version 54 ends with an exact original comparison remainder containing
$(M(a)-\widetilde M_a)\dot q$.  Its prepared linear response theorem does not
close a nonlinear iteration with this term: estimating that remainder in
$C^k$ ordinarily requires one additional derivative of the unknown rate.
The present body removes that particular derivative \emph{algebraically},
before estimating any forcing.  It does not assert that the other source
estimates or nonlinear prepared boundary conditions have thereby been proved.

The calculation is at the same V39 algebraic root and $h=1/3$.
All 107 normalized multiplier coordinates and both fast scales are retained.
The coordinate change is a state-dependent identification of rate variables,
not an alteration of the action, physical clock, or executed history.
The full signed mass is made equal to the \emph{same signed frozen mass};
it is not made positive.

\paragraph{Inputs and scope of the source audit.}
We use the original graded stationary chart and tilt-independent leading
Hessian of \cref{v41:thm:true-Hessian,v44:lem:regular}, the complete mixed scaling
of \cref{v53:eq:pencil,v54:eq:original-divided}, and the exact original
constraint and differentiated consistency identities.  The equation labels
used below specify these inputs independently if a title differs between
source renderings.  In particular the retained facts are: the leading
81-dimensional mixed differential block is positive; its 26-dimensional
Schur complement is nonsingular of inertia $(16,10,0)$; and the divided
coefficient functions are analytic on the inherited regular chart.  These
facts are inherited, not new numerical certificates.  A targeted text review
of the preceding native body and supplied Einstein V36 and perturbative V51
bodies did not locate the graded moving-mass construction below.  The review
is not an independent validation of all companion proofs.  Matrix functions
and isometries for indefinite scalar products are established mathematics
\cite{v55:HMT}; the required local formula and its estimates are proved here.

We first work with arbitrary nearby graded coefficients.  Application to an
original path is exact wherever the stated stationary chart is defined.
Small-$a$ estimates are local on bounded normalized coefficient sets, not a
new proof that an original trajectory stays in such a set for $t\le aT$.

\section{A bounded congruence respecting the two mass weights}
\label{v55:sec:congruence}

Reorder the rates as $q=(x,w)$, with $x=(v,c)\in\R^{81}$ and
$w\in\R^{26}$.  Write the frozen mass of Version 53 as
\begin{equation}
 M_a=\begin{pmatrix}A_0&aB_0\\aB_0^T&a^2C_0\end{pmatrix},\qquad
 H_0=A_0^{-1}B_0,\qquad
 \Delta_0=C_0-B_0^TA_0^{-1}B_0.
 \label{v55:eq:base-mass}
\end{equation}
Thus $A_0=M_x\succ0$, $B_0=B_x$, and
$\operatorname{Inertia}\Delta_0=(16,10,0)$.  The subscripts distinguish
these blocks from earlier active-only matrices.  A moving mass with the
\emph{same weights} is
\begin{equation}
 \widetilde M_a=
 \begin{pmatrix}\widetilde A&a\widetilde B\\
 a\widetilde B^T&a^2\widetilde C\end{pmatrix},\qquad
 \widetilde H=\widetilde A^{-1}\widetilde B,\quad
 \widetilde\Delta=\widetilde C-\widetilde B^T\widetilde A^{-1}\widetilde B.
 \label{v55:eq:moving-mass}
\end{equation}
The normalized moving coefficients are sufficiently close to their frozen
values.  They may depend on time or on the original state.

\begin{lemma}[Local isometry without a positive replacement]
\label[lemma]{v55:lem:square}
Let $G=G^T$ be nonsingular and let $\widetilde G=\widetilde G^T$ satisfy
$\|E_G\|\le r<1$, where $E_G=G^{-1}(\widetilde G-G)$ in a fixed
submultiplicative norm.  The convergent real matrix series
\begin{equation}
 R_G=(I+E_G)^{-1/2}
 =\sum_{n=0}^{\infty}(-1)^n
          \frac{\binom{2n}{n}}{4^n}E_G^n
 \label{v55:eq:binomial}
\end{equation}
is invertible, analytic in $\widetilde G$, and satisfies
\begin{equation}
 R_G^T\widetilde G R_G=G,\qquad
 \|R_G-I\|\le(1-r)^{-1/2}-1.
 \label{v55:eq:isometry}
\end{equation}
Its first derivative obeys
$\|DR_G[E]\|\le\tfrac12(1-r)^{-3/2}\|E\|$ when the derivative argument
is an increment of $E_G$.  Truncation after degree $m$ has error at most
$r^{m+1}/(1-r)$.
\end{lemma}
\begin{proof}
Symmetry gives $E_G^TG=GE_G$.  Every real polynomial in $E_G$, and then
the absolutely convergent series, obeys $R_G^TG=GR_G$.
The scalar power-series identity implies $R_G^2(I+E_G)=I$; hence $R_G$
is invertible and
$R_G^T\widetilde GR_G=GR_G(I+E_G)R_G=G$.
The absolute coefficient series is $(1-r)^{-1/2}$, proving the norm
bound.  Differentiating a matrix power gives the sum of its $n$ possible
insertions of the increment.  The absolute scalar derivative bounds this
sum, even when the matrices do not commute.  Finally
$\binom{2n}{n}/4^n\le1$ gives the geometric tail bound.  No square root
of the indefinite matrix $G$ itself is taken.
\end{proof}

Apply the lemma separately to $G=A_0$ and $G=\Delta_0$, with moving
matrices $\widetilde A,\widetilde\Delta$, and call the results $R_x,R_w$.
Define
\begin{equation}
 \boxed{
 \mathcal U_a=
 \begin{pmatrix}
 R_x&a(R_xH_0-\widetilde H R_w)\\0&R_w
 \end{pmatrix}.}
 \label{v55:eq:U}
\end{equation}

\begin{theorem}[Exact triangular normalization of the graded mass]
\label[theorem]{v55:thm:mass}
The matrix \eqref{v55:eq:U} is invertible and satisfies
\begin{equation}
                  \boxed{\mathcal U_a^T\widetilde M_a\mathcal U_a=M_a.}
 \label{v55:eq:mass-fixed}
\end{equation}
It and its inverse are uniformly bounded for $0<a\le a_0$ in a fixed
sufficiently small normalized coefficient neighborhood.  If
\begin{equation}
 \widetilde A=A_0+aA_1(a,\eta),\quad
 \widetilde B=B_0+aB_1(a,\eta),\quad
 \widetilde C=C_0+aC_1(a,\eta),
 \label{v55:eq:regular-coeff}
\end{equation}
with analytic bounded coefficients on a compact $\eta$-set, then
\begin{equation}
 \boxed{
 \mathcal U_a-I=
 \begin{pmatrix}O(a)&O(a^2)\\0&O(a)\end{pmatrix}.}
 \label{v55:eq:orders}
\end{equation}
These orders hold with any fixed number of $\eta$ derivatives on smaller
compact sets.  In particular neither the map nor its derivative estimates
contain the $a^{-2}$ norm of the full mass inverse.  The inertia remains
$(97,10,0)$ for $a>0$.
\end{theorem}
\begin{proof}
For $Q(H)=\left(\begin{smallmatrix}I&-aH\\0&I\end{smallmatrix}\right)$,
direct block multiplication gives
\[
 Q(\widetilde H)^T\widetilde M_aQ(\widetilde H)
       =\operatorname{diag}(\widetilde A,a^2\widetilde\Delta),\qquad
 Q(H_0)^TM_aQ(H_0)=\operatorname{diag}(A_0,a^2\Delta_0).
\]
The exact factorization
$\mathcal U_a=Q(\widetilde H)\operatorname{diag}(R_x,R_w)Q(H_0)^{-1}$
and the lemma prove \eqref{v55:eq:mass-fixed}.  All block inverses are
uniformly bounded near the nonsingular frozen blocks, proving boundedness
and inertia.  Under \eqref{v55:eq:regular-coeff},
$R_x-I,R_w-I,\widetilde H-H_0=O(a)$.  The constant terms of
$R_xH_0-\widetilde H R_w$ cancel, so its prefactor $a$ gives $O(a^2)$.
Analyticity gives the same orders for parameter derivatives.
\end{proof}

A generic square root of the full relatively normalized 107-dimensional
mass need not have these unweighted bounds: conjugating a generic $O(a)$
block by $\operatorname{diag}(I,aI)$ can create an $O(1)$ weak--differential
shear.  The block-triangular choice above avoids that artificial cost.
It fixes the original signed mass rather than diagonalizing or suppressing
its negative directions.

\section{First coefficients and the connection generated by a moving frame}
\label{v55:sec:connection}

The first coefficients of the map can be read without differentiating a
107-dimensional matrix square root.  In the following formulas $A_1,B_1,C_1$
mean the values of \eqref{v55:eq:regular-coeff} at $a=0$:
\begin{align}
 H_1&=A_0^{-1}(B_1-A_1H_0),\\
 \Delta_1&=C_1-B_1^TH_0-H_0^TB_1+H_0^TA_1H_0,\\
 X_1&=-\tfrac12A_0^{-1}A_1,\qquad
 Y_1=-\tfrac12\Delta_0^{-1}\Delta_1.
 \label{v55:eq:first-blocks}
\end{align}
Here and below all remainders retain the moving coefficient dependence.

\begin{proposition}[Explicit first frame coefficient]
\label[proposition]{v55:prop:first}
One has
\begin{equation}
 \mathcal U_a=
 \begin{pmatrix}
 I+aX_1+O(a^2)&a^2(X_1H_0-H_1-H_0Y_1)+O(a^3)\\
 0&I+aY_1+O(a^2)
 \end{pmatrix},
 \label{v55:eq:first-frame}
\end{equation}
These are identities in the actual normalized
coefficient jets, not numerical values for uncomputed higher native jets.
\end{proposition}
\begin{proof}
Differentiate $\widetilde A\widetilde H=\widetilde B$ and the Schur
complement to obtain \eqref{v55:eq:first-blocks}.  The first binomial
coefficient is $-1/2$.  Expanding \eqref{v55:eq:U} gives the formula; no
commutation of a matrix with its derivative is used.
\end{proof}

Let the normalized coefficients follow a differentiable path and put
\begin{equation}
 \Omega=\mathcal U_a^T\widetilde M_a\dot{\mathcal U}_a,
 \qquad \Omega_{\rm sk}=\tfrac12(\Omega-\Omega^T).
 \label{v55:eq:Omega}
\end{equation}
Dots denote $d/d\theta$ at fixed $a$.

\begin{lemma}[The metric-derivative term is retained exactly]
\label[lemma]{v55:lem:connection}
The normalization implies
\begin{equation}
 \Omega+\Omega^T=-\mathcal U_a^T\dot{\widetilde M}_a\mathcal U_a.
 \label{v55:eq:connection-id}
\end{equation}
If \eqref{v55:eq:regular-coeff} holds and the normalized path has bounded
velocity, then
\begin{equation}
 \Omega,\Omega_{\rm sk}=
 \begin{pmatrix}O(a)&O(a^2)\\O(a^2)&O(a^3)\end{pmatrix}.
 \label{v55:eq:connection-orders}
\end{equation}
\end{lemma}
\begin{proof}
Differentiate \eqref{v55:eq:mass-fixed}.  Its right side is independent of
$\theta$, which gives \eqref{v55:eq:connection-id}.  The differentiated
orders in \eqref{v55:eq:orders}, the blocks of
\eqref{v55:eq:moving-mass}, and direct multiplication give
\eqref{v55:eq:connection-orders}.
\end{proof}

For any skew $\widetilde J$ and column $\widetilde f$, the moving equation
$\widetilde M_a\dot q+\widetilde Jq+\widetilde f=0$, under
$q=\mathcal U_ap$, is therefore exactly
\begin{equation}
 \boxed{
 M_a\dot p+\widehat J_ap+\widehat f_a=0,
 \quad
 \widehat J_a=\mathcal U_a^T\widetilde J\mathcal U_a+\Omega_{\rm sk},
 \quad
 \widehat f_a=\mathcal U_a^T\widetilde f
       -\tfrac12\mathcal U_a^T\dot{\widetilde M}_a\mathcal U_ap.}
 \label{v55:eq:normal-eq}
\end{equation}
In particular $\widehat J_a^T=-\widehat J_a$.  Omitting the last term in
$\widehat f_a$ would change the equation.  Equivalently one can retain
$\Omega$ in the coefficient and use $\mathcal U_a^T\widetilde f$ as the
source.  The two forms are exactly equal by \eqref{v55:eq:connection-id}.

\section{An exact original-equivalent system with no derivative in its remainder}
\label{v55:sec:original}

Use the fixed reordered physical rate synthesis
\begin{equation}
 S_a=(aI_A,I_c,aI_g),\qquad R_a=a^{-3}S_a^T,
 \qquad \lambda_t=S_aq,\quad t=a\theta.
 \label{v55:eq:scales}
\end{equation}
Its range is the fixed mean-lapse-zero multiplier complement.  All matrices
of the preceding sections are in this reordered rate basis.  Define the
\emph{original} tensors at the moving record by
\begin{align}
 \widetilde M_a&=a^{-4}S_a^TM^{\rm orig}S_a,&
 \widetilde J_a&=a^{-3}S_a^TK^{\rm orig}S_a,\\
 \widetilde f_a&=a^{-3}S_a^T\bigl\{
 D_X\mathsf B[F]+2D_XM^{\rm orig}[F]S_aq
 +D_\lambda M^{\rm orig}[S_aq]S_aq\bigr\}.
 \label{v55:eq:original-objects}
\end{align}
These are precisely the Version-54 divided objects, after permutation.
The full factor two is retained.

For a candidate original state $Y=(X,\lambda)$ and a rate coordinate $p$,
set $q=\mathcal U_a(Y)p$ and define the directional derivative
\begin{equation}
 \mathfrak D_{a,p}Z(Y)=
 D_XZ(Y)[aF(Y)]+D_\lambda Z(Y)[aS_a\mathcal U_a(Y)p].
 \label{v55:eq:direction}
\end{equation}
This uses the canonical and multiplier \emph{first} equations.  It contains
$p$, but contains no $\dot p$.  In \eqref{v55:eq:Omega} and
\eqref{v55:eq:normal-eq}, replace dots on moving coefficients by
$\mathfrak D_{a,p}$, and evaluate \eqref{v55:eq:original-objects} at
$q=\mathcal U_ap$.

\begin{theorem}[Exact rate normalization of the nonlinear original problem]
\label[theorem]{v55:thm:original}
On the inherited stationary chart where the nearby graded congruence is
defined, the following system is exact for each $a>0$:
\begin{equation}
 \boxed{
 X_\theta=aF(X,\lambda),\qquad
 \lambda_\theta=aS_a\mathcal U_a(X,\lambda)p,\qquad
 M_a p_\theta+\widehat J_a(X,\lambda,p)p+
                    \widehat f_a(X,\lambda,p)=0.}
 \label{v55:eq:extended}
\end{equation}
Both coefficients in the last row are smooth functions of the displayed
state and rate variables, with no rate derivative in either coefficient.
Every original history produces a solution of this system.  Conversely,
a solution whose initial data satisfy
\begin{equation}
 P=0,\qquad
 \mathsf B+M^{\rm orig}S_a\mathcal U_ap=0,
 \qquad \overline{\lambda^N}=1
 \label{v55:eq:initial}
\end{equation}
is an original constrained history for as long as it remains in that chart.
No extra histories are introduced by differentiating consistency.
\end{theorem}
\begin{proof}
The exact derivative of
$\mathcal C=\mathsf B+M^{\rm orig}\lambda_t$ along the first two rows
has the form
\[
 \mathcal C_t=M^{\rm orig}\lambda_{tt}+K^{\rm orig}\lambda_t
 +D_X\mathsf B[F]+2D_XM^{\rm orig}[F]\lambda_t
 +D_\lambda M^{\rm orig}[\lambda_t]\lambda_t.
\]
Here $D_\lambda\mathsf B=K^{\rm orig}+D_XM^{\rm orig}[F]$,
which explains the factor two.  Multiplication by $R_a$ and
$\lambda_{tt}=a^{-1}S_a\dot q$ gives
$\widetilde M_a\dot q+\widetilde J_aq+\widetilde f_a=0$.
The chain rule gives
$\dot q=\mathcal U_a\dot p+(\mathfrak D_{a,p}\mathcal U_a)p$.
The exact congruence and connection identity prove the last row of
\eqref{v55:eq:extended}.  They also prove its asserted derivative dependence.

Conversely that row implies $S_a^T\mathcal C_t=0$, so
$S_a^T\mathcal C=0$ remains true from \eqref{v55:eq:initial}.
The synthesis range is the mean-lapse-zero hyperplane.  Hence $\mathcal C$
is a scalar multiple of its fixed normal.  Original homogeneity gives
$M^{\rm orig}\lambda=0$ and $\langle\lambda,\mathsf B\rangle=0$,
so $\langle\lambda,\mathcal C\rangle=0$ identically.  The fixed normal
pairs nontrivially with $\lambda$ because its mean lapse is one, preserved
by the second equation.  Therefore $\mathcal C=0$.  Finally
$P_t=D_XP[F]+M^{\rm orig}\lambda_t=\mathcal C=0$, preserving $P=0$.
The retained normalized signed regularity identifies this rate with the
unique original rate wherever that regularity holds.
\end{proof}

\begin{corollary}[The mass-defect derivative disappears from the exact comparison]
\label[corollary]{v55:cor:remainder}
Let $J,f$ be the frozen Version-53 coefficients.  The last equation in
\eqref{v55:eq:extended} can be written
\begin{equation}
 \boxed{M_a\dot p+Jp=-f+\mathscr R_a^\sharp(X,\lambda,p),}
 \label{v55:eq:fixed-comparison}
\end{equation}
where
\begin{equation}
 \boxed{
 \mathscr R_a^\sharp=(J-\widehat J_a)p+f-\widehat f_a.}
 \label{v55:eq:Rsharp}
\end{equation}
Unlike \eqref{v54:eq:native-defect}, this identity has no term containing
$\dot p$ on its right side.  All derivatives of the mass and frame are
retained through \eqref{v55:eq:direction}; no such derivative is set to zero.
\end{corollary}
\begin{proof}
Subtract the constant $Jp+f$ from \eqref{v55:eq:normal-eq}.  Theorem
\ref{v55:thm:original} establishes the stated state dependence.
\end{proof}

The inherited normalized stationary chart makes the moving blocks analytic
and nonsingular near the retained leading record.  On bounded sets with
$X=aY_*+O(a^2)$ and bounded first multiplier coefficients, the leading
weak-Hessian tilt cancellation in Version 52 gives
\eqref{v55:eq:regular-coeff}.  This statement fixes a local domain on which
the construction and its $O(a)$ bounds apply.  It does not establish the
existence of an original $[0,aT]$ history in that domain.  At fixed $a>0$
the exact equivalence above holds wherever the normalized blocks stay in
the specified neighborhood, even without an asymptotic rate estimate.


\begin{corollary}[Exact initial access in the normalized rate coordinates]
\label[corollary]{v55:cor:access}
Fix any finite normalized rate $p^0$.  For all sufficiently small $a>0$,
there are original initial records with the retained leading $Y_*,\ell_*$
that satisfy \eqref{v55:eq:initial} with normalized rate exactly $p^0$.
If $q_a^0$ is their physical mixed-rate coordinate, then
$q_a^0=p^0+O(a)$.  This is initial access only; it imposes no prepared
terminal conditions on a later nonlinear path.
\end{corollary}
\begin{proof}
The inherited parametric initializer assigns an exact initial record
$\mathcal I_a(q)$ to every $q$ in a neighborhood of the fixed finite
$p^0$, with original rate $S_aq$.  Its leading fields are independent
of $q$ and its higher coefficients are analytic with bounded derivatives
on a smaller neighborhood.  The graded mass expansion therefore gives
$\mathcal U_a(\mathcal I_a(q))=I+O(a)$ and
$D_q\mathcal U_a(\mathcal I_a(q))=O(a)$ there.  Solve
$q=\mathcal U_a(\mathcal I_a(q))p^0$ by contraction on a ball of radius
$C a$ around $p^0$.  The map takes the ball into itself for sufficiently
large fixed $C$, and its Lipschitz constant is $O(a\|p^0\|)<1$ for small
$a$.  The selected original record is already exactly constrained and
consistent.  By definition its normalized rate is $p^0$.
\end{proof}

The normalization also retains the complete signed work identity:
\begin{equation}
 \frac{d}{d\theta}\frac12p^TM_ap
 =-p^T\widehat f_a
 =-q^T\widetilde f_a+\frac12q^T\dot{\widetilde M}_aq.
 \label{v55:eq:work}
\end{equation}
This follows either from \eqref{v55:eq:normal-eq} and skew symmetry or
by differentiating $p^TM_ap=q^T\widetilde M_aq$.
It is conservation-of-work bookkeeping for an indefinite form, not a
positive energy estimate.  The physical rate is always reconstructed as
$S_a\mathcal U_ap$.  In fixed coordinate norms the nonconstant physical
rate correction is $O(a^2\|p\|)$, while the harmonic correction is
$O(a\|p\|)$; these follow directly from \eqref{v55:eq:orders} and
\eqref{v55:eq:scales}.

\section{The construction does not erase the native intermediate gyro}
\label{v55:sec:gyro}

A successful normalization must not silently replace the native intermediate
Poisson coefficient by the frozen one.  The following invariance identifies
what survives.

\begin{proposition}[Kernel-compressed first skew coefficient is unchanged]
\label[proposition]{v55:prop:gyro}
Suppose on a bounded normalized path
$\widetilde J_a=J+aJ_1+O(a^2)$ and
\eqref{v55:eq:regular-coeff} holds.  Let $Z$ be the retained synthesis of
$\ker K_g$, where $K_g=J_{gg}$.  Then
\begin{equation}
 \boxed{
 Z^T\left[\partial_a\widehat J_a\big|_{a=0}\right]_{gg}Z
                   =Z^T(J_1)_{gg}Z.}
 \label{v55:eq:gyro-invariant}
\end{equation}
In particular the intermediate pencil still retains
$S_0+Z^T(J_1)_{gg}Z$, as in Version 53; the congruence supplies no value
for the uncomputed native skew coefficient.
\end{proposition}
\begin{proof}
By \eqref{v55:eq:first-frame}, the first derivative of the frame is block
diagonal with blocks $X_1,Y_1$; its cross block begins at $a^2$.
Thus the order-$a$ weak block of
$\mathcal U_a^T\widetilde J_a\mathcal U_a$ is
$(J_1)_{gg}+Y_1^TK_g+K_gY_1$.
Both extra terms vanish when compressed by $Z$.  The weak block of the
connection is $O(a^3)$ by \eqref{v55:eq:connection-orders}; it contributes
nothing at order $a$.  This proves the identity.
\end{proof}

The leading descriptor and its rank-two differentiated-source map
$\mathcal N_1$ also remain unchanged, since $\mathcal U_0=I$.
The change removes a derivative of the \emph{unknown rate from the nonlinear
remainder}; it does not remove differentiation of the two weak source
coordinates in the descriptor.  These are distinct operations.
Likewise the transformed source in \eqref{v55:eq:Rsharp} need not have a
favorable sign.  The full fixed mass still has ten negative directions, and
none of the existing growing or oscillatory sectors is thereby discarded.

\section{A quantitative remainder bound and a certificate for the frame}
\label{v55:sec:bounds}

\begin{proposition}[The new remainder requires no extra rate derivative]
\label[proposition]{v55:prop:regularity}
On a compact normalized coefficient set with bounded $p$ and bounded
normalized state velocity, assume \eqref{v55:eq:regular-coeff}.  Then
\begin{equation}
 \|\mathscr R_a^\sharp\|\le C_K\bigl(
 a+\|\widetilde J_a-J\|+\|\widetilde f_a-f\|\bigr).
 \label{v55:eq:C0-bound}
\end{equation}
More generally, for each fixed $a>0$ and integer $k\ge0$, the right side
of \eqref{v55:eq:Rsharp} is a smooth composition on $C^k$ state/rate paths:
its $k$th time derivative involves rate derivatives only through order $k$,
not $k+1$.  On sets where the corresponding normalized coefficient
partial derivatives are uniformly bounded, the composition and local
Lipschitz constants are uniform as well.
\end{proposition}
\begin{proof}
The frame and its inverse are $I+O(a)$, while the connection and
$\mathfrak D_{a,p}\widetilde M_a$ are $O(a)$ in unweighted norm.
Use \eqref{v55:eq:normal-eq} and boundedness of the frozen matrices and $p$
to obtain \eqref{v55:eq:C0-bound}.  For the derivative assertion write the
coefficient derivatives as the explicit directional expression
\eqref{v55:eq:direction}.  It is a smooth function of the state and $p$.
Repeated finite-dimensional chain and product rules on this function
therefore introduce at most $k$ derivatives of $p$.  The same rules on a
compact convex neighborhood give the local Lipschitz estimate.  Uniform
higher partial-derivative bounds are stated explicitly, not inferred from
an unscaled analytic formula containing negative powers of $a$.
\end{proof}

This removes one avoidable derivative requirement from the nonlinear source
map.  It does \emph{not} remove Version 54's $C^3\to C^0$ source-response
regularity requirement.  In particular the proposition is not a Banach
fixed-point theorem for the native problem.  A closed function space,
control of the actual transformed forcing, and access to the nonlinear
prepared boundary set remain to be proved.

The frame itself admits a terminating coefficient-level enclosure.
Suppose certified bounds give
\[
 \|A_0^{-1}(\widetilde A-A_0)\|\le r_x<1,\qquad
 \|\Delta_0^{-1}(\widetilde\Delta-\Delta_0)\|\le r_w<1.
\]
For the degree-$m$ binomial polynomials $R_{x,m},R_{w,m}$, put
$e_x=r_x^{m+1}/(1-r_x)$ and $e_w=r_w^{m+1}/(1-r_w)$.
The block formula gives the rigorous synthesis bound
\begin{equation}
 \|\mathcal U_a-\mathcal U_{a,m}\|
 \le C_{\rm blk}\{e_x+e_w+
 a(e_x\|H_0\|+\|\widetilde H\|e_w)\}.
 \label{v55:eq:truncation}
\end{equation}
For a fixed compatible block norm $C_{\rm blk}$ is explicit (one may take
$2$ for the Euclidean block estimate).  The mass congruence residual is
bounded by
\begin{equation}
 \|\mathcal U_{a,m}^T\widetilde M_a\mathcal U_{a,m}-M_a\|
 \le\|\widetilde M_a\|\,
       (2\|\mathcal U_a\|e_U+e_U^2),
 \label{v55:eq:mass-certificate}
\end{equation}
where $e_U$ is the right side of \eqref{v55:eq:truncation}.
This bounds the full signed identity without an $a^{-2}$ factor.
Obtaining the input block enclosures for a proposed original state remains
an original stationary-coefficient calculation; it has not been executed
numerically here for an unknown nonlinear trajectory.

\section{Status, nonduplication, and the next native calculation}
\label{v55:sec:status}

\paragraph{Proved.}
The moving original mass has an exact, bounded, near-identity triangular
normalization to the retained frozen signed mass.  The weak-to-differential
cross correction is $O(a^2)$, rather than a singularly amplified generic
frame change.  Retaining the complete frame connection gives an exact
nonlinear system whose comparison remainder contains no derivative of the
unknown rate.  With original initial constraints and consistency, this
system is equivalent to the original action evolution; it does not create
additional differentiated-equation histories.  The native first
kernel-compressed skew correction is preserved.  Exact truncation bounds
supply a finite verification method for the new frame.

\paragraph{Inherited and executed.}
The Version-54 verifier was replayed with its embedded inputs checked.  The
new finite checks test the indefinite local square-root series, noncommuting
first coefficients, exact block congruence, metric-connection sign, the
factor-two original derivative identity, kernel-gyro invariance, and
constraint propagation in a separate explicitly identified homogeneous
Hamiltonian example.  The examples are algebraic tests, not original
spacetime integrations.  No new higher original matrix coefficient is
claimed evaluated.  No blanket revalidation of earlier nonlinear theorems
is inferred from a matrix-regression run.

\paragraph{Unpaid and next direction.}
The remaining native task is now the regularity and prepared response of
$\mathscr R_a^\sharp(X,\lambda,p)$ in \eqref{v55:eq:Rsharp}, together with
initial-access transversality for the coupled equations
\eqref{v55:eq:extended}.  The removable mass-defect derivative no longer
belongs in that task.  The genuine native skew coefficient, the two
source-sensitive weak coordinates, the faster growing modes, and the
horizon-dependent prepared conditions remain.  No original $[0,aT]$ lift,
fixed positive physical lifespan, spatially uniform estimate, primitive
selection, or Einstein--SM continuum theorem is proved in this body.
Earlier shorter original histories and the independent Version-32 stress
question remain at their stated scopes.

Every Version-54 byte before the final document command is preserved.
New labels use the prefix \texttt{v55:}.  The analytical construction is
not a proof-assistant formalization or a numerical nonlinear trajectory.

\begingroup
\renewcommand{\refname}{Additional reference for Version 55}
\begin{thebibliography}{9}
\small\raggedright
\bibitem{v55:HMT}
N.~J.~Higham, C.~Mehl and F.~Tisseur,
\emph{The canonical generalized polar decomposition},
SIAM Journal on Matrix Analysis and Applications \textbf{31} (2010),
2163--2180. \url{https://doi.org/10.1137/090765018}.
Primary-source context for matrix functions and indefinite isometries;
no original nonlinear evolution estimate is imported from that reference.
\end{thebibliography}
\endgroup
\addtocontents{toc}{\protect\endgroup}
% END IMMUTABLE VERSION 55 BODY
% APPEND VERSION 56 HERE, BEFORE THE FINAL END-DOCUMENT COMMAND.



% BEGIN IMMUTABLE VERSION 56 BODY
\clearpage
\begin{center}
{\Large\bfseries Version 56: The sharp source regularity before nonlinear closure}\\[.5em]
{\large An integrable second derivative, first-jet preparation,\\
and resonance in the retained intermediate oscillatory sector}
\end{center}
\makeatletter
\addtocontents{toc}{\protect\begingroup}
\addtocontents{toc}{\protect\makeatletter}
\addtocontents{toc}{\protect\def\protect\l@section{\protect\bfseries\protect\@dottedtocline{1}{0em}{2.5em}}}
\makeatother
\addcontentsline{toc}{part}{Version 56: Integrable second source derivatives and native-pencil resonance}

\section{Audit: removing a mass derivative did not close the response space}
\label{v56:sec:audit}

Version 55 constructs the exact graded mass normalization and the
original-equivalent remainder $\mathscr R_a^\sharp(X,\lambda,p)$ without a
rate derivative.  Its final proposition explicitly does not close a
nonlinear fixed-point space: Version 54 still used a $C^3$ covector norm
for a $C^0$ prepared response.  The first question here is whether that
three-derivative entrance is intrinsic or an artifact of preparing one
more source moment than necessary.

We retain the \emph{same frozen coefficient pencil} $M(a)\partial_\theta+J$
of \cref{v53:eq:pencil}, the same algebraic mixed-mode root, and $h=1/3$.
A dot denotes differentiation in $\theta=t/a$ and $0\le\theta\le T<\infty$.
The analysis is of a forced comparison equation
\begin{equation}
                 M(a)\dot q+Jq=g(\theta).
 \label{v56:eq:equation}
\end{equation}
It adds no force to the original action.  Native implications concern the
response norm available for a \emph{separately verified} original remainder,
not an assertion that an arbitrary comparison source is produced by an
original history.

The new theorem replaces $C^3$ by $W^{2,1}$, prepares only the first source
jet, and gives an explicit modulus of convergence.  A simple-pole
calculation in the \emph{actual retained intermediate pencil} then proves
that uniformly small $C^{1,\alpha}$ sources, $\alpha<1$, are insufficient.
This is not another use of the two-dimensional diagnostic oscillator from
Version 54.  The negative theorem uses the 107-dimensional retained
coefficient equation and its own preparation.

\paragraph{Inherited inputs and scope.}
We use the three-scale spectral splitting of \cref{v53:thm:three-scales},
the positive finite spectral subspace, the complete source-resolvent
construction of \cref{v54:lem:moments}, and the signed normalization of
\cref{v55:thm:mass}.  Their root and primitive coefficient banks are not
new computations.  The final source and status sections of Versions 53--55
were reviewed for overlap; their stated Green estimates use the third and
fourth inverse powers, rather than the second-power response proved here.
The polynomial part of an index-two descriptor is established context
\cite{v54:BGMehr}; the sharper estimates and original-pencil lower bound
below are derived in full.  No broader novelty claim for integration by
parts or oscillatory cancellation is made.

\section{One fewer source moment and the relevant Green bounds}
\label{v56:sec:green}

Let $L_a=-M(a)^{-1}J$, and retain the commuting projections $P_s,P_+,P_-,P_o$
onto respectively the 79 finite modes, ten growing modes, ten decaying
modes, and eight rapid oscillatory modes.  Write $P_f=I-P_s$.  For fixed
positive $a$, let $L_f$ be the invertible restriction of $L_a$ to its fast
space and put
\begin{equation}
 F_0=-L_f^{-1}P_fM(a)^{-1},\qquad
 F_1=-L_f^{-2}P_fM(a)^{-1}.
 \label{v56:eq:moments}
\end{equation}
Version 54 proves that these maps are bounded and analytic through $a=0$,
with $F_0(0)=\mathcal N_0$ and $F_1(0)=\mathcal N_1$.  We require only
\begin{equation}
                 q_{f,\mathrm{app}}=F_0g+F_1\dot g.
 \label{v56:eq:approximation}
\end{equation}
The source derivative here has well-defined endpoint values whenever
$g\in W^{2,1}(0,T)$.

\begin{lemma}[Second-power covector bounds]
\label[lemma]{v56:lem:green}
For sufficiently small $a>0$, there are constants $C,\kappa>0$ such that
\begin{align}
 \|e^{-rL_a}P_+F_1\|+\|e^{rL_a}P_-F_1\|
 &\le C\left(a^2e^{-\kappa r/a^2}+e^{-\kappa r/a}\right),\quad r\ge0,
 \label{v56:eq:hyper}\\
 \|e^{rL_a}P_{o,2}F_1\|&\le Ca^2,\quad r\in\R.
 \label{v56:eq:fast-osc}
\end{align}
Here $P_{o,2}$ contains the six fastest imaginary modes.  The remaining
intermediate imaginary projections $P_{o,\pm}$ are simple and satisfy
\begin{equation}
 L_aP_{o,\pm}=\pm i\omega_aP_{o,\pm},\qquad
 a\omega_a\longrightarrow\omega_0\in(0.733,0.734),\qquad
 \|P_{o,\pm}F_1\|\le C.
 \label{v56:eq:middle-osc}
\end{equation}
No uniform bound for the projections alone is assumed.
\end{lemma}
\begin{proof}
Use the exact sheared pencil in \cref{v53:eq:sheared-pencil}.  Its source
resolvent is $O(1)$ on $z=\mu/a^2$ contours around the nonzero fastest
clusters.  In the contour for $e^{rL_a}P L_f^{-2}P_fM^{-1}$, the factor
$z^{-2}$ contributes $a^4$ and the contour measure contributes $a^{-2}$.
This gives $a^2$.  Hyperbolic contours give the corresponding forward or
backward decay; the definite oscillatory metric gives a bounded propagator
for the six fastest imaginary modes, including repeated frequencies.
On intermediate contours $z=\nu/a$, the complete source resolvent is
$O(a^{-1})$.  The inverse square contributes $a^2$ and the contour measure
$a^{-1}$, giving a bounded product.  The two real poles yield the second
term in \eqref{v56:eq:hyper}.  The two imaginary poles are simple, stay
purely imaginary by the inherited definite-signature argument, and have
$az\to\pm i\omega_0$.  Their residues give
\eqref{v56:eq:middle-osc}.  These are the same contour and signature
arguments as Version 54 with the inverse power decreased to two; all
required power factors are displayed above.
\end{proof}

For $b\in L^1(0,T;\C^{107})$, define its correction by
\begin{align}
 (\mathcal E_ab)(\theta)={}&
 -\int_0^\theta e^{(\theta-r)L_a}(P_-+P_o)F_1b(r)\,dr
 \notag\\
 &+\int_\theta^T e^{(\theta-r)L_a}P_+F_1b(r)\,dr.
 \label{v56:eq:correction}
\end{align}
Both integrals are ordinary finite-dimensional Lebesgue integrals.

\begin{lemma}[A uniform integrable-source bound and one oscillatory derivative]
\label[lemma]{v56:lem:integral}
For fixed $T$, the correction satisfies
\begin{equation}
 \|\mathcal E_ab\|_{C^0}\le C_T\|b\|_{L^1}.
 \label{v56:eq:L1}
\end{equation}
If $b\in W^{1,1}(0,T)$, it also satisfies
\begin{equation}
 \|\mathcal E_ab\|_{C^0}
 \le C_Ta\left(\|b\|_{L^\infty}+\|\dot b\|_{L^1}\right).
 \label{v56:eq:W11}
\end{equation}
The statements apply componentwise to real sources.
\end{lemma}
\begin{proof}
All kernels in \cref{v56:lem:green} are uniformly bounded, which proves
\eqref{v56:eq:L1}.  For bounded $b$, integration of the hyperbolic kernels
costs at most $C(a^4+a)\|b\|_\infty$.  The fastest oscillatory term costs
$CTa^2\|b\|_\infty$.  For each intermediate pole, integrate
$e^{\pm i\omega_a(\theta-r)}B_ab(r)$ by parts, with
$B_a=P_{o,\pm}F_1$ uniformly bounded.  Its norm is at most
\[
 \frac C{\omega_a}
       \left(2\|b\|_\infty+\|\dot b\|_{L^1}\right).
\]
Since $\omega_a^{-1}\le Ca$, this proves \eqref{v56:eq:W11}.
The boundary terms are retained; oscillation is not replaced by decay.
\end{proof}

\section{Exact preparation for a source with only an integrable second derivative}
\label{v56:sec:main}

Use the norm
\begin{equation}
 \|g\|_{\mathcal X_T}:=\|g\|_{C^1}+\|\ddot g\|_{L^1}
 \quad\text{on }W^{2,1}(0,T;\R^{107}).
 \label{v56:eq:norm}
\end{equation}
Each such function has a unique $C^1$ representative.  For $b\in L^1$, set
\begin{equation}
 \mathfrak m_a(b):=
 \inf_{v\in W^{1,1}(0,T)}
 \left\{\|b-v\|_{L^1}
          +a\bigl(\|v\|_\infty+\|\dot v\|_{L^1}\bigr)\right\}.
 \label{v56:eq:modulus}
\end{equation}
This is an explicit approximation modulus, not a norm defined by the
unknown solution operator.  One has $\mathfrak m_a(b)\le\|b\|_1$ and
$\mathfrak m_a(b)\to0$ for every fixed $b\in L^1$.  The latter follows by
first approximating $b$ in $L^1$ by a smooth function and then taking $a\to0$.

\begin{theorem}[First-source-jet prepared lift on $W^{2,1}$]
\label[theorem]{v56:thm:lift}
Fix $T<\infty$, $g\in W^{2,1}(0,T;\R^{107})$, and a descriptor finite
initial component $q_s^0$.  Let $q_0$ be the complete descriptor solution
of \cref{v54:thm:descriptor}; its weak part is continuous, although it
need not have a continuous time derivative.  For all sufficiently small
$a>0$, \eqref{v56:eq:equation} has exactly one prepared solution satisfying
\begin{align}
 P_s(a)q_a(0)&=U_aq_s^0,\\
 (P_-+P_o)\bigl(q_a(0)-q_{f,\mathrm{app}}(0)\bigr)&=0,\\
 P_+\bigl(q_a(T)-q_{f,\mathrm{app}}(T)\bigr)&=0.
 \label{v56:eq:boundary}
\end{align}
Here $U_a$ is the inherited finite-space identification, and
$q_{f,\mathrm{app}}$ is \eqref{v56:eq:approximation}.  The solution is
$C^1$ for each fixed $a>0$ and obeys
\begin{equation}
 \boxed{\ \|q_a\|_{C^0}
       \le C_T\bigl(\|q_s^0\|+\|g\|_{\mathcal X_T}\bigr).\ }
 \label{v56:eq:bounded}
\end{equation}
Its descriptor error satisfies
\begin{equation}
 \boxed{\ \|q_a-q_0\|_{C^0}
 \le C_T\left[
       a\bigl(\|q_s^0\|+\|g\|_{C^1}\bigr)
       +\mathfrak m_a(\ddot g)\right].\ }
 \label{v56:eq:main-error}
\end{equation}
In particular, $q_a\to q_0$ uniformly for every fixed $W^{2,1}$ source.
Only $g$ and $\dot g$, not $\ddot g$, occur in the preparation traces.
\end{theorem}
\begin{proof}
Solve the finite spectral equation by the inherited bounded propagator
and source map.  Its difference from the finite descriptor solution is
bounded by $C_Ta(\|q_s^0\|+\|g\|_{C^0})$; the forced zero mode remains
inside this comparison.  The fast moment identities are
$L_fF_0=-P_fM^{-1}$ and $L_fF_1=F_0$.  Therefore, in the sense of
$L^1$ weak derivatives,
\[
 q_{f,\mathrm{app}}'-L_aq_{f,\mathrm{app}}-P_fM^{-1}g
                       =F_1\ddot g.
\]
Adding $\mathcal E_a\ddot g$ cancels this defect exactly and imposes the
two fast conditions in \eqref{v56:eq:boundary}.  The sum is absolutely
continuous and solves the coefficient equation almost everywhere.  Its
right-hand side is continuous, so the sum is in fact $C^1$.  A homogeneous
difference satisfying the three preparation conditions is zero in every
invariant sector, proving uniqueness.

The moment limits give
\[
 \|q_{f,\mathrm{app}}-
        (\mathcal N_0g+\mathcal N_1\dot g)\|_{C^0}
                                  \le Ca\|g\|_{C^1}.
\]
For any $v\in W^{1,1}$, apply \eqref{v56:eq:L1} to $\ddot g-v$ and
\eqref{v56:eq:W11} to $v$.  Taking the infimum proves
$\|\mathcal E_a\ddot g\|_{C^0}\le C_T\mathfrak m_a(\ddot g)$ and then
\eqref{v56:eq:main-error}.  The uniform bound follows from
$\mathfrak m_a(b)\le\|b\|_1$ and the descriptor formula.
\end{proof}

The $79+18+10$ preparation conditions still retain all 107 coordinates.
The growing terminal condition depends on the specified source segment
and on $T$.  It is a comparison preparation, not a causal reader, future
measurement, or update of a native history.  This preparation is different
from Version 54's additional $F_2\ddot g$ preparation.  The two coincide
for constant and affine sources, but are not silently identified for a
general moving source.

\section{Quantitative moduli, jumps of the second derivative, and source families}
\label{v56:sec:rates}

For continuous $b$ on $[0,T]$, let
\[
 \omega_b(\delta)=
 \sup_{|s-t|\le\delta}\|b(s)-b(t)\|.
\]
Extend $b$ constantly beyond the endpoints and take
$v(t)=a^{-1}\int_0^a b(t+s)\,ds$.  Then
$\|b-v\|_1\le T\omega_b(a)$,
$\|v'\|_1\le T\omega_b(a)/a$, and $\|v\|_\infty\le\|b\|_\infty$.
Consequently
\begin{equation}
 \mathfrak m_a(b)\le a\|b\|_\infty+2T\omega_b(a).
 \label{v56:eq:continuous-modulus}
\end{equation}

\begin{corollary}[H\"older and bounded-variation source entrances]
\label[corollary]{v56:cor:rates}
For $g\in C^2$, the error in \eqref{v56:eq:main-error} is bounded by
\begin{equation}
 C_T\left[a(\|q_s^0\|+\|g\|_{C^2})
                       +\omega_{\ddot g}(a)\right].
 \label{v56:eq:C2}
\end{equation}
If $\ddot g$ is H\"older continuous with exponent $0<\alpha\le1$, this is
$O_T(a^\alpha)$ on bounded $C^{2,\alpha}$ sets.  If instead
$\ddot g\in BV(0,T)$, then
\begin{equation}
 \boxed{\ \|q_a-q_0\|_{C^0}
 \le C_Ta\left(\|q_s^0\|+\|g\|_{C^1}
       +\|\ddot g\|_\infty+\operatorname{Var}(\ddot g)\right).\ }
 \label{v56:eq:BV}
\end{equation}
Jumps of the second derivative are permitted.  There is no point mass in
$g$ or $\dot g$, both of which retain their continuous representatives.
\end{corollary}
\begin{proof}
Use \eqref{v56:eq:continuous-modulus} and
$\omega_b(a)\le [b]_{C^{0,\alpha}}a^\alpha$.
For a bounded-variation representative, constant extension and the same
moving average give
$\|b-v\|_1\le a\operatorname{Var}(b)$ and
$\|v'\|_1\le\operatorname{Var}(b)$.
Thus $\mathfrak m_a(b)\le a(\|b\|_\infty+2\operatorname{Var}(b))$.
Substitute in the theorem.  Endpoint values of $b$ itself are not used
in the preparation, so their representatives do not alter the solution.
\end{proof}

\begin{corollary}[Uniform compact-source convergence]
\label[corollary]{v56:cor:compact}
Suppose the finite initial components and $\|g_a\|_{C^1}$ stay bounded,
and $\{\ddot g_a\}$ is relatively compact in $L^1(0,T)$.  Then the
prepared solution and the descriptor for the \emph{same} source $g_a$
satisfy $\|q_a-q_{0,a}\|_{C^0}\to0$.  If also $g_a\to g_0$ in $C^1$
and the finite initial components converge, then $q_a\to q_0[g_0]$
uniformly.  A perturbation $r_a=O(a)$ in $\mathcal X_T$ changes the
prepared response by $O_T(a)$.
\end{corollary}
\begin{proof}
The functional $\mathfrak m_a$ is $1$-Lipschitz in $L^1$, by its defining
infimum.  A finite $L^1$ net of the compact closure and its pointwise
convergence therefore give uniform convergence of
$\mathfrak m_a(\ddot g_a)$ to zero.  The descriptor is continuous from
$C^1$ source data to $C^0$ states by its explicit formula.  The last
assertion follows by linearity and \eqref{v56:eq:bounded}.
\end{proof}

As a separate diagnostic, consider $\dot x+w=0$,
$a^2\dot w-x=\tfrac12(\theta-\theta_0)_+^2$, with zero data before
$\theta_0$.  For $u=\theta-\theta_0\ge0$, the exact prepared solution is
\[
 x=-\tfrac12u^2+a^2(1-\cos(u/a)),\qquad
 w=u-a\sin(u/a).
\]
It differs by $O(a)$ from the descriptor $(-u^2/2,u)$ even though the
source second derivative jumps at $\theta_0$.  This checks the stated
regularity interpretation; it is not an original gravitational trajectory.

\section{The actual intermediate residue makes the lower regularity failure sharp}
\label{v56:sec:resonance}

The following result uses the original retained coefficient pencil, not
an auxiliary oscillator.  In the sheared variables of
\cref{v53:eq:sheared-pencil}, let
\[
 \mathcal P(\nu)=\nu^2D_0+\nu S_0+G_0,
 \qquad \nu_+=i\omega_0.
\]
The root $\nu_+$ is simple.  The kernel synthesis $Z$ and the three
matrices are precisely those of \cref{v53:eq:kernel-pencil}.

\begin{lemma}[Nonzero intermediate source residue]
\label[lemma]{v56:lem:residue}
Let $z_+(a)=i\omega_a$ be the positive intermediate imaginary pole, and
$\mathcal R_a(z)=(zM(a)+J)^{-1}$.  Then
\begin{equation}
 a^2\operatorname*{Res}_{z=z_+(a)}\mathcal R_a(z)
       =\mathcal C+O(a),\qquad
 \mathcal C=\begin{pmatrix}
 0&0\\0&
 \displaystyle Z\frac{\nu_+\operatorname{adj}\mathcal P(\nu_+)}
 {\partial_\nu\det\mathcal P(\nu_+)}Z^T
 \end{pmatrix}\ne0.
 \label{v56:eq:residue}
\end{equation}
The blocks are in the original reordered $(x,w)$ coordinates.  There
exists a fixed real covector $h=(0,h_g)$ in the two-dimensional source
bank $h_g\in\Ran Z$ with $\mathcal Ch\ne0$.  One of the two columns
$(0,Ze_1),(0,Ze_2)$ suffices.
\end{lemma}
\begin{proof}
For $z=\nu/a$, the upper inverse of the exact sheared pencil is
$a\nu^{-1}M_x^{-1}+O(a^2)$.  Eliminate it and then the 24-dimensional
nonkernel weak block.  On a fixed small contour around $\nu_+$, the
remaining two-dimensional block is
\[
 a\{\nu D_0+S_0+\nu^{-1}G_0\}+O(a^2).
\]
All eliminated inverses are analytic and uniformly bounded at their
stated scaling.  Hence, uniformly on that contour,
\[
 a\mathcal R_a(\nu/a)=
 \begin{pmatrix}0&0\\0&Z\nu\mathcal P(\nu)^{-1}Z^T\end{pmatrix}+O(a).
\]
The original mass shear is $I+O(a)$ and does not change the leading
coefficient.  Under $z=\nu/a$, a residue gains another factor $a^{-1}$,
which proves \eqref{v56:eq:residue} by contour integration.  A simple
zero of the determinant has a nonzero rank-one adjugate, so the residue
is nonzero.  Injectivity of $Z$ and nonsingularity of $Z^TZ$ show that it
cannot annihilate both displayed real source columns.  No numerical
pole residue or upper bound on the permissible $a$ is inferred.
\end{proof}

\begin{theorem}[Same-pencil resonance excludes a uniform $C^{1,\alpha}$ entrance]
\label[theorem]{v56:thm:sharp}
Choose the fixed real $h$ in \cref{v56:lem:residue}, and define
\begin{equation}
                 g_a(\theta)=a^2h\cos(\omega_a\theta).
 \label{v56:eq:resonant-source}
\end{equation}
Give its prepared solution the finite initial component zero and the
first-jet preparation \eqref{v56:eq:boundary}.  Then
\begin{equation}
 \|g_a\|_{C^{1,\alpha}}=O(a^{1-\alpha})\longrightarrow0
 \quad(0\le\alpha<1),\qquad
 \sup_a\|g_a\|_{C^2}<\infty,
 \label{v56:eq:source-norms}
\end{equation}
where $\alpha=0$ denotes the ordinary $C^1$ norm.  The descriptor for
$g_a$ satisfies $\|q_{0,a}\|_{C^0}=O_T(a)$, but
\begin{equation}
 \boxed{\ 
 \lim_{a\downarrow0}\|q_a-q_{0,a}\|_{L^2(0,T)}^2
       =\frac{T^3}{6}\|\mathcal Ch\|^2>0.
 \ }
 \label{v56:eq:resonant-lower}
\end{equation}
In particular, no uniformly bounded linear source-to-prepared-state
estimate from $C^{1,\alpha}$ to $C^0$ is possible for this pencil and
preparation when $\alpha<1$.  Bounded $C^2$ norms alone do not imply
uniform descriptor convergence for amplitude-dependent source families.
\end{theorem}
\begin{proof}
First use the complex source $g_a^c=a^2h e^{i\omega_a\theta}$.
The resonant source residue is $B_a=P_{o,+}M(a)^{-1}$, with
$a^2B_a=\mathcal C+O(a)$.  Its exact modal response is
\[
 e^{i\omega_a\theta}\left[P_{o,+}q_{f,\mathrm{app}}(0)
                                  +a^2B_ah\,\theta\right].
\]
The initial term equals $-2a^2z_+(a)^{-1}B_ah=O(a)$, by the two moment
identities.  The conjugate intermediate pole is separated by
$2\omega_a\asymp a^{-1}$ and contributes $O_T(a)$ to this complex
forcing.  The intermediate hyperbolic forward/backward Green integrals
and their prepared traces likewise contribute $O_T(a)$.  The fastest
source residues have size $O(a^{-2})$, but their spectral distance from
$i\omega_a$ is of order $a^{-2}$, so their oscillatory particular and
prepared terms are $O_T(a^2)$.  Their hyperbolic terms have the same
bound by forward/backward integration.  The finite source map is
bounded and its forcing is $O(a^2)$.  These estimates follow also by
using the unweighted resolvent on the separate contours in the proof of
\cref{v56:lem:green}; no norm of a bare intermediate state projector
is used.  Thus
\[
 q_a^c(\theta)=\theta e^{i\omega_a\theta}\mathcal Ch+O_T(a)
\]
uniformly.  Take real parts.  The squared norm of
$\operatorname{Re}(e^{i\omega_a\theta}v)$ has mean
$\|v\|^2/2$; its other terms have frequencies $2\omega_a\to\infty$.
Integration against $\theta^2$ gives
$T^3\|v\|^2/6+o(1)$.  The descriptor is $O(a)$ because its finite source
part is $O_T(a^2)$ and its differentiated feedthrough is $O(a)$.
This proves \eqref{v56:eq:resonant-lower}.

The norm estimates follow from $a\omega_a\to\omega_0$ and
$|\sin u-\sin v|\le\min(2,|u-v|)$:
$[\dot g_a]_{C^{0,\alpha}}\le C a^2\omega_a^{1+\alpha}$.
The second derivative is bounded but oscillates at frequency
$\omega_a$.  Its family is not relatively compact in $L^1$; otherwise
\cref{v56:cor:compact} would contradict \eqref{v56:eq:resonant-lower}.
Finally $\|q_a-q_{0,a}\|_{C^0}\ge T^{-1/2}\|q_a-q_{0,a}\|_{L^2}$.
\end{proof}

The same proof applies to any fixed linear observation $O$ for which
$O\mathcal Ch\ne0$, replacing the right side by
$T^3\|O\mathcal Ch\|^2/6$.  Nonannihilation must be checked for the
particular observation; it is not inferred for every curvature component.
The near-identity V55 frame preserves the existence of this normalized
rate response.  None of these assertions says that the actual nonlinear
remainder generates the deliberately chosen resonant covector.

\section{What changes for the normalized original remainder}
\label{v56:sec:native}

The uniform source perturbation requirement supplied by the new theorem is
\begin{equation}
 \boxed{\quad
 \|\mathscr R_a^\sharp\|_{C^1}
       +\|\partial_\theta^2\mathscr R_a^\sharp\|_{L^1}=O(a),
 \quad}
 \label{v56:eq:native-entrance}
\end{equation}
when the corresponding prepared boundary conditions and finite initial
comparison have independently been realized.  This replaces the former
sufficient $C^3$ covector bound, not the original action or its row scaling.
For a fixed smooth base source, it gives $O(a)$ state error by
\cref{v56:cor:compact}.  More generally the compact-source theorem permits
convergence without a uniform third derivative.

V55 shows that $\mathscr R_a^\sharp$ is a smooth function of the normalized
state and rate.  If a path $z$ lies in a bounded $W^{2,1}$ set and a source
function $R_a(z)$ has uniformly bounded first two state derivatives, then
\[
 \partial_\theta^2R_a(z)
   =D^2R_a(z)[\dot z,\dot z]+DR_a(z)\ddot z
\]
is integrable.  On compact state sets, its $W^{2,1}$ composition estimate
uses only two derivatives of $z$, since $W^{2,1}$ embeds into $C^1$ in
one time dimension.  If those source partial derivatives themselves are
$O(a)$, the resulting norm is $O(a)$ on bounded path sets.  These are
explicit conditional composition estimates; the required uniform native
partial bounds and bounded path set have not been established here.

There is still no same-space nonlinear contraction: the proved response
map is $W^{2,1}\to C^0$, not $W^{2,1}\to W^{2,1}$ uniformly.  The actual
resonance theorem rules out replacing this gap by a generic small-$C^1$
argument.  Furthermore the first-jet preparation depends on the whole
source segment through its growing terminal condition; no nonlinear
initial transversality or causal selection is inferred.

For a fixed source covered by the theorem, its initial prepared rate
converges to a finite descriptor initial rate.  The existing parametric
initializer and V55's exact normalized-rate access consequently realize
that rate on exact original constrained and consistent initial records.
This is an application of the inherited initial theorem, not a new claim
about the subsequent original trajectory.  The original nonlinear gyro,
all ten growing directions, and all eight rapid oscillatory directions
remain in the unpaid evolution problem.

\section{Status and verification record}
\label{v56:sec:status}

\paragraph{Proved.}
The full retained coefficient equation has a uniformly bounded prepared
response for sources with an integrable second derivative.  Its initial
and terminal preparation uses only the first source jet.  Every fixed
$W^{2,1}$ source converges to its complete descriptor; an explicit
approximation modulus quantifies the result.  H\"older second derivatives
give $O(a^\alpha)$ error, and bounded-variation second derivatives give
$O(a)$ error even with jumps.  In the same 107-dimensional pencil, the
actual intermediate imaginary pole has a nonzero order-$a^{-2}$ covector
residue.  It produces a resonant family tending to zero in every
$C^{1,\alpha}$, $\alpha<1$, whose prepared rate error has a strictly
positive $L^2$ limit.  Thus a generic one-derivative source entrance is
not available.

\paragraph{Inherited and executed.}
The V55 verifier is replayed with its embedded hashes and inherited source
checks.  New exact finite-algebra checks verify moment-defect identities,
the two-derivative preparation, the scaled intermediate resolvent and its
simple nonzero residue in a separately declared constant-coefficient
model, and exact polynomial and piecewise-polynomial source responses.
Additional numerical quadrature tests of the modulus and resonance
identities are explicitly diagnostics, not proof certificates.  The
same-root resonance theorem is a written argument from the inherited
original intermediate contour and simple-pole certificates.  No new
107-dimensional native coefficient or original spacetime trajectory is
claimed to have been evaluated.

\paragraph{Unpaid and next direction.}
The remaining nonlinear calculation is the original transformed remainder
in the weaker norm \eqref{v56:eq:native-entrance}, or a structured
principal treatment of its resonant feedback, together with prepared
initial-access transversality.  The new estimate removes an unnecessary
third source derivative; it does not close a nonlinear existence argument
by itself.  No original $[0,aT]$ lift, larger lifespan, spatially uniform
continuum theorem, primitive preparation mechanism, or Einstein--SM
closure is proved.  The earlier original histories and the independent
Version-32 stress-compensation problem remain at their stated scopes.

Every Version-55 byte before its final document command is preserved.
New labels use the prefix \texttt{v56:}.  Exact finite checks are not a
proof-assistant formalization of the cumulative programme.

\addtocontents{toc}{\protect\endgroup}
% END IMMUTABLE VERSION 56 BODY
% APPEND VERSION 57 HERE, BEFORE THE FINAL END-DOCUMENT COMMAND.



% BEGIN IMMUTABLE VERSION 57 BODY
\clearpage
\begin{center}
{\Large\bfseries Version 57: The metric feedback closes in continuous rate space}\\[.5em]
{\large Exact original-source factorization, weighted Green bounds,\\
and a same-space nonlinear prepared-response theorem}
\end{center}
\makeatletter
\addtocontents{toc}{\protect\begingroup}
\addtocontents{toc}{\protect\makeatletter}
\addtocontents{toc}{\protect\def\protect\l@section{\protect\bfseries\protect\@dottedtocline{1}{0em}{2.5em}}}
\makeatother
\addcontentsline{toc}{part}{Version 57: Exact metric feedback and a continuous-rate contraction}

\section{The next obstruction is structural, not another derivative count}
\label{v57:sec:scope}

Version 56 proves a prepared $W^{2,1}$ covector response but leaves a
same-space nonlinear iteration open.  Its resonance theorem concerns
arbitrary covectors; it does not say that every term generated by the
original action has that worst source scaling.  We return to the exact
Version-55 rate equation and separate its rate dependence before applying
a response estimate.  This gives a positive distinction: all terms caused
by differentiating the multiplier Hessian and its normalizing frame have
a stronger weak-row weight.  Their complete, genuinely quadratic rate
feedback admits a uniform contraction in $C^0$ on a fixed sufficiently
short comparison interval.

This is not a proof of the full original nonlinear lift.  The state-dependent
Poisson matrix and the rate-independent differentiated drift are retained
separately.  The new contraction treats the metric-feedback part about the
retained constant pencil; it does not freeze the other two terms and then
call the resulting path an original spacetime.  Everything is at the same
V39 root and $h=1/3$.

The inherited inputs are the exact objects and chain rule in
\cref{v55:thm:original}, the triangular frame in \cref{v55:thm:mass}, and
the three-scale resolvent and definite oscillatory subspaces of Versions
53--56.  The source moments, general $W^{2,1}$ estimate, and resonance
counterexample are not reproved as new results.  A targeted search of the
cumulative body found no factorization of its complete quadratic rate
term in the form below.  Dependence of descriptor perturbation estimates
on source structure is established context \cite{v57:Gear}; the identities,
weights and nonlinear theorem used here are proved explicitly.

\section{The original rate dependence has an exact quadratic factorization}
\label{v57:sec:factor}

Use the reordered coordinates $q=(x,w)\in\R^{81}\oplus\R^{26}$, with
all three harmonic shifts in $x$.  Keep the physical synthesis
$S_a=(aI_A,I_c,aI_g)$ and the original Hamiltonian objects
$P=H_\lambda$, $M^{\rm orig}=H_{\lambda\lambda}$,
$F=\mathbb JH_X$, and $\mathsf B=P_XF$.
Put $Y=(X,\lambda)$ and retain
\begin{equation}
 \widetilde M=a^{-4}S_a^TM^{\rm orig}S_a,\qquad
 \widetilde J=a^{-3}S_a^TK^{\rm orig}S_a,\qquad
 \widetilde f_0=a^{-3}S_a^TD_X\mathsf B[F].
 \label{v57:eq:objects}
\end{equation}
The first two matrices are respectively symmetric and skew.  Write
$U=\mathcal U_a(Y)$ for the exact triangular frame and $M_a$ for its
constant signed mass, so $U^T\widetilde M U=M_a$ and $q=Up$.
The letter $p$ in this body is a rate coordinate, not canonical momentum.

Differentiate the frame separately in its two original state directions:
\begin{equation}
 C_X=U^{-1}D_XU[aF],\qquad
 C_\lambda[p]=U^{-1}D_\lambda U[aS_aUp].
 \label{v57:eq:connections}
\end{equation}
The second expression is linear in $p$.  Neither expression contains a
rate derivative.  Define the rate-independent coefficients
\begin{equation}
 J_U=U^T\widetilde J U,\qquad f_U=U^T\widetilde f_0.
 \label{v57:eq:JU}
\end{equation}

\begin{theorem}[Exact original-source rate split]
\label[theorem]{v57:thm:factor}
On the Version-55 stationary chart the original-equivalent rate row is
exactly
\begin{equation}
 \boxed{\ M_a p_\theta+J_U(Y)p+f_U(Y)=\mathcal Q_a(Y,p),\ }
 \label{v57:eq:split}
\end{equation}
where $J_U^T=-J_U$ and
\begin{equation}
 \boxed{\ \mathcal Q_a(Y,p)
 =(2C_X^TM_a+M_aC_X)p+C_\lambda[p]^TM_ap.\ }
 \label{v57:eq:quadratic}
\end{equation}
At every fixed state this is a polynomial of degree at most two in the
rate.  Its quadratic Hessian is the explicit bilinear map
\begin{equation}
 D_p^2\mathcal Q_a[u,v]
 =C_\lambda[u]^TM_av+C_\lambda[v]^TM_au.
 \label{v57:eq:rate-hessian}
\end{equation}
All metric-derivative and moving-frame terms of the original equation
are included.  There are no higher powers of the rate in this part.
\end{theorem}
\begin{proof}
Use $\theta=t/a$ at fixed $a$.  The twice differentiated original
consistency equation, in the divided coordinates, reads
\[
 \widetilde M q_\theta+\widetilde Jq+\widetilde f_0
 +2D_X\widetilde M[aF]q
 +D_\lambda\widetilde M[aS_aq]q=0.
\]
The factors follow directly from \eqref{v57:eq:objects}; for example
$D_X\widetilde M[aF]q=a^{-3}S_a^TD_XM^{\rm orig}[F]S_aq$.
This retains the original factor two.

Differentiate $U^T\widetilde M U=M_a$ in either of the two directions in
\eqref{v57:eq:connections}.  If the corresponding connection is $C$, then
\[
 U^T(D\widetilde M)U=-C^TM_a-M_aC,
 \qquad U^T\widetilde M(DU)=M_aC.
\]
Substitute $q=Up$.  For the $X$ direction the sum of connection and metric
terms is
$M_aC_X-2(C_X^TM_a+M_aC_X)=-2C_X^TM_a-M_aC_X$.
For the $\lambda$ direction it is
$M_aC_\lambda-(C_\lambda^TM_a+M_aC_\lambda)=-C_\lambda^TM_a$.
Move these terms to the right.  This proves \eqref{v57:eq:split}--
\eqref{v57:eq:quadratic}.  Linearity of $C_\lambda[p]$ gives
\eqref{v57:eq:rate-hessian}.  The argument is pointwise on the chart;
it does not first assume a solution or assign values to higher initializer
coefficients.
\end{proof}

Together with $X_\theta=aF$ and $\lambda_\theta=aS_aUp$, this is the
same original-equivalent system as Version 55.  Its preservation of the
undifferentiated constraints from consistent initial data is inherited
from \cref{v55:thm:original}.  No new evolution term has been inserted.
In particular, the signed work identity is still
\[
 \frac d{d\theta}\frac12p^TM_ap=-p^Tf_U+p^T\mathcal Q_a.
\]
The right side has no sign asserted here.

\section{The entire metric feedback has a favorable weak-row weight}
\label{v57:sec:weights}

Use a bounded normalized chart with
$X=aY_*+a^2\xi$ and $\lambda=e+a\ell$.  The regular coefficient
hypotheses of Version 55 hold on its compact subsets.  In these coordinates
$F/a$ is bounded, and the two directions in \eqref{v57:eq:connections}
have normalized velocities $F/a$ and $S_aUp$, respectively.
Hence differentiating the established frame orders gives
\begin{equation}
 C_X=\begin{pmatrix}O(a)&O(a^2)\\0&O(a)\end{pmatrix},\qquad
 C_\lambda[p]=\begin{pmatrix}O(a\|p\|)&O(a^2\|p\|)\\0&O(a\|p\|)\end{pmatrix}.
 \label{v57:eq:Corders}
\end{equation}
These bounds also hold with any fixed number of normalized-state
partial derivatives on smaller compact sets.  They are local coefficient
bounds, not a proof of a native lifespan in that chart.

Define the source injection
\begin{equation}
 \boxed{\ \mathcal D_a=\operatorname{diag}(aI_{81},a^2I_{26}).\ }
 \label{v57:eq:injection}
\end{equation}
It is a weight on covectors, distinct from the mass and from the physical
rate synthesis $S_a$.

\begin{proposition}[Automatic source-weighted rate bounds]
\label[proposition]{v57:prop:weighted}
There is a polynomial $\mathcal F_a(Y,p)$ of degree at most two in $p$ such
that
\begin{equation}
 \boxed{\ \mathcal Q_a(Y,p)=\mathcal D_a\mathcal F_a(Y,p).\ }
 \label{v57:eq:F}
\end{equation}
Its coefficients and any fixed number of normalized-state derivatives
extend boundedly through $a=0$ on the stated compact sets.  In particular,
for $\|p\|,\|\widetilde p\|\le R$,
\begin{align}
 \|\mathcal F_a(Y,p)\|&\le C(R+R^2),\\
 \|\mathcal F_a(Y,p)-\mathcal F_a(Y,\widetilde p)\|
 &\le C(1+R)\|p-\widetilde p\|.
 \label{v57:eq:FLip}
\end{align}
No time derivative of $p$ is used in these inequalities.
\end{proposition}
\begin{proof}
Write $M_a=\left(\begin{smallmatrix}A&aB\\aB^T&a^2C\end{smallmatrix}\right)$
and $C_*=\left(\begin{smallmatrix}aX&a^2Y\\0&aZ\end{smallmatrix}\right)$.
Direct multiplication shows that both $C_*^TM_a$ and $M_aC_*$ have
blocks of orders
\[
 \begin{pmatrix}O(a)&O(a^2)\\O(a^2)&O(a^3)\end{pmatrix}.
\]
Left multiplication by $\mathcal D_a^{-1}$ removes no more powers than
these blocks contain.  Apply this to $C_X$ and to the matrix linear in
$p$ in \eqref{v57:eq:Corders}.  The coefficients are analytic in the
normalized variables and their powers of $a$ are nonnegative after
weighting.  The polynomial norm and difference estimates follow, as do
the derivative statements.
\end{proof}

An unstructured estimate $\mathcal Q_a=O(a)$ loses the additional weak
factor.  That factor is precisely what matters at the intermediate
resonance, whose raw weak source residue is of order $a^{-2}$.  The
proposition does not put the state-dependent Poisson difference or the
rate-independent $f_U$ in the same weighted class.

\section{A continuous-source Green operator for the retained signed pencil}
\label{v57:sec:Green}

Use the same frozen $M_a,J$ and commuting spectral projections as Version
56.  For a continuous $r:[0,T]\to\R^{107}$, let $\mathcal G_ar$ solve
\begin{equation}
 M_a e_\theta+Je=\mathcal D_ar,
 \qquad P_s e(0)=0,\quad(P_-+P_o)e(0)=0,\quad P_+e(T)=0.
 \label{v57:eq:Green-bc}
\end{equation}
The preparation is homogeneous for the correction.  It is not the
first-source-jet preparation imposed on the complete nonlinear source
in Version 56.

\begin{lemma}[Source-weighted kernels without source differentiation]
\label[lemma]{v57:lem:kernels}
For $0<T\le T_0<\infty$ and all sufficiently small $a>0$,
\begin{align}
 \|e^{sL_a}P_sM_a^{-1}\mathcal D_a\|&\le C_{T_0}a,
                   &&0\le s\le T_0,\\
 \|e^{-sL_a}P_+M_a^{-1}\mathcal D_a\|
 +\|e^{sL_a}P_-M_a^{-1}\mathcal D_a\|
 &\le C(e^{-\kappa s/a^2}+e^{-\kappa s/a}),&&s\ge0,\\
 \|e^{sL_a}P_oM_a^{-1}\mathcal D_a\|&\le C,&&s\in\R.
 \label{v57:eq:kernels}
\end{align}
If $F_0=-L_f^{-1}P_fM_a^{-1}$, its projected weighted propagators obey
the same bounds with additional factors $a^2$ on the fastest groups
and $a$ on the intermediate groups.
\end{lemma}
\begin{proof}
We display the source powers, which differ from Version 56.  In the
exact mass shear of Version 53, multiplication by $\mathcal D_a$ gives
an upper source $ar_x$ and a lower source
$a^2(r_w-H_0^Tr_x)$.  On $z=\mu/a^2$ contours around a nonzero fastest
cluster, the upper inverse is $O(a^2)$ and the remaining weak Schur
inverse is $O(1)$.  The full sheared source solution is $O(a^2)$; undoing
the bounded shear does not change that bound.  The contour measure
contributes $a^{-2}$, giving a bounded source propagator.

On $z=\nu/a$ intermediate contours, the upper inverse is $O(a)$.
The effective lower source is $O(a^2)$ even after its upper-block return.
Eliminating the rank-24 weak block leaves an $O(a^{-1})$ two-dimensional
Schur inverse, so the full source resolvent is $O(a)$.  Contour measure
$a^{-1}$ again gives a bounded source propagator.  The inherited
hyperbolic gaps give the first two decays.  The definite metrics on the
oscillatory clusters give the bounded real-time estimates, including
repeated fastest imaginary frequencies; no decay of oscillations is
used.  On the fixed finite contour the source resolvent is $O(a)$,
giving the first estimate.  Inserting $-z^{-1}$ in the fast integrals
for $F_0$ supplies respectively $a^2$ and $a$.  These arguments use the
same uniformly nonsingular eliminated blocks and definite-signature
bases as the retained V53--V56 proofs, not uniform bounds for bare
unweighted intermediate projections.
\end{proof}

\begin{theorem}[Uniform continuous-source response on the weighted bank]
\label[theorem]{v57:thm:Green}
Equation \eqref{v57:eq:Green-bc} has a unique $C^1$ solution for continuous
$r$, and
\begin{equation}
 \boxed{\ \|\mathcal G_ar\|_{C^0}\le C_{T_0}(T+a)\|r\|_{C^0}.\ }
 \label{v57:eq:C0Green}
\end{equation}
For $r\in W^{1,1}$ it has the stronger estimate
\begin{equation}
 \boxed{\ \|\mathcal G_ar\|_{C^0}
 \le C_{T_0}a(\|r\|_\infty+\|\dot r\|_{L^1}).\ }
 \label{v57:eq:smoothGreen}
\end{equation}
All 107 equations and all 28 fast modes are retained.
\end{theorem}
\begin{proof}
The exact solution is
\[
 \mathcal G_ar(\theta)=
 \int_0^\theta e^{(\theta-s)L_a}(P_s+P_-+P_o)M_a^{-1}\mathcal D_ar(s)\,ds
 -\int_\theta^T e^{(\theta-s)L_a}P_+M_a^{-1}\mathcal D_ar(s)\,ds.
\]
The kernels give respectively $CaT$, $C(a^2+a)$ and $CT$.  For $a\le1$
these prove \eqref{v57:eq:C0Green}.  The formula differentiates to the
original equation.  A homogeneous solution with its stated initial and
terminal projections zero is zero, proving uniqueness.

For \eqref{v57:eq:smoothGreen}, integrate the fast terms once by parts,
using $L_fF_0=-P_fM_a^{-1}$.  The result is $F_0\mathcal D_ar(\theta)$,
minus its forward and backward propagated endpoint values, plus Green
integrals with $F_0\mathcal D_a\dot r$.  Every projected propagator in
this expression is $O(a)$ by the last statement of the lemma.  The
endpoint terms are retained.  The finite part costs $CaT\|r\|_\infty$.
This proves the estimate without a second source derivative.
\end{proof}

The terminal condition still uses the whole specified interval; this is a
prepared boundary-value response, not a causal Green function.  It imposes
79 finite initial, 18 decaying/oscillatory initial and ten growing terminal
conditions.  No physical measurement is associated with those projections.

\section{A same-space nonlinear contraction, including a prepared error of order a}
\label{v57:sec:nonlinear}

Consider a reference solution $p_{b,a}$ of the retained linear equation,
$M_ap_{b,a}'+Jp_{b,a}=g_{b,a}$, with continuous $g_{b,a}$ and a uniform
bound on $p_{b,a}$.  Give the correction $p-p_{b,a}$ exactly the homogeneous
conditions in \eqref{v57:eq:Green-bc}.  The nonlinear comparison is
\begin{equation}
 M_ap'+Jp=g_{b,a}+\mathcal D_a\mathcal F_a(\theta,p).
 \label{v57:eq:semilinear}
\end{equation}
For metric feedback, $\mathcal F_a$ is the polynomial in
\cref{v57:prop:weighted} evaluated on a prescribed bounded normalized
coefficient path.  A prescribed path is not assumed to satisfy the
original kinematic equations.

\begin{theorem}[Closed continuous-rate prepared feedback problem]
\label[theorem]{v57:thm:nonlinear}
Suppose $\mathcal F_a$ is continuous in $\theta,p$ and, on a fixed tube
$\|p-p_{b,a}\|\le r$, has uniform bound $B_r$ and rate Lipschitz constant
$L_r$.  If
\begin{equation}
 C_{T_0}(T+a)B_r\le r,\qquad
 \kappa:=C_{T_0}(T+a)L_r<1,
 \label{v57:eq:contraction}
\end{equation}
then \eqref{v57:eq:semilinear} has exactly one solution in that tube with
the stated correction preparation.  It is $C^1$ for every $a>0$.
The conditions hold for all sufficiently small $a$ on a fixed sufficiently
short positive interval $0\le\theta\le T_*$, independent of $a$.

If in addition the reference feedback obeys
\begin{equation}
 \|\mathcal F_a(\cdot,p_{b,a})\|_\infty
 +\|\partial_\theta\mathcal F_a(\cdot,p_{b,a})\|_{L^1}\le B_* ,
 \label{v57:eq:reference-smooth}
\end{equation}
uniformly, then
\begin{equation}
 \boxed{\ \|p_a-p_{b,a}\|_{C^0}
       \le\frac{C_{T_0}aB_*}{1-\kappa}.\ }
 \label{v57:eq:nonlinear-Oa}
\end{equation}
No derivative bound on the unknown nonlinear rate $p_a$ is assumed.
\end{theorem}
\begin{proof}
On the complete closed ball in $C^0$, apply Banach's theorem to
$p\mapsto p_{b,a}+\mathcal G_a\mathcal F_a(\cdot,p)$.
The first inequality makes the ball invariant and the second is its
contraction bound, both by \eqref{v57:eq:C0Green}.  Choose $T_*$ and
then $a_*$ so that these inequalities hold with, for example,
$\kappa\le1/2$.  The exact Green formula and continuity of the right
side give $C^1$ regularity and every differential row at the fixed point.
Subtract the map evaluated at $p_{b,a}$ and use
\eqref{v57:eq:smoothGreen} only on that reference feedback.  This gives
\[
 \|p_a-p_{b,a}\|_\infty
 \le\kappa\|p_a-p_{b,a}\|_\infty+C_{T_0}aB_*.
\]
Rearrangement proves \eqref{v57:eq:nonlinear-Oa}.  Derivatives of the
unknown iterate have never appeared in the contraction norm or its
Lipschitz estimate.
\end{proof}

This includes the complete actual metric-derivative polynomial at a fixed
regular original state, and prescribed coefficient paths with uniformly
integrable first derivatives.  For a uniformly $C^1$ prepared reference,
\eqref{v57:eq:reference-smooth} follows by the ordinary polynomial chain
rule.  The nonlinear solution itself may have rapid derivatives; they
are not substituted into that reference-only estimate.  If the reference
has an $O(a)$ descriptor limit, the nonlinear comparison has the same
limit at $O(a)$.

The theorem is not an original-action existence theorem.  In particular,
freezing $J_U(Y)$ at the reference matrix $J$ while keeping the metric
polynomial defines the stated comparison subproblem.  It does not keep the
other original state-dependent terms by implication.

\section{The weighted gain cannot be treated as an arbitrary small feedback}
\label{v57:sec:sharp}

The factor $T+a$ in \eqref{v57:eq:C0Green} is substantive.  For arbitrary
continuous amplitudes $r_a$ it cannot be replaced by a factor tending to
zero with $a$ on a fixed interval.

\begin{proposition}[Same-pencil sharpness of the continuous-source gain]
\label[proposition]{v57:prop:sharp}
Let $h_g\in\Ran Z$ be one of the real source directions selected in
\cref{v56:lem:residue}, let $h=(0,h_g)$, and put
$r_a(\theta)=(0,h_g\cos(\omega_a\theta))$.  For the homogeneous correction
preparation in \eqref{v57:eq:Green-bc},
\begin{equation}
 \boxed{\ \lim_{a\downarrow0}\|\mathcal G_ar_a\|_{L^2(0,T)}^2
        =\frac{T^3}{6}\|\mathcal Ch\|^2>0.\ }
 \label{v57:eq:sharp}
\end{equation}
Here $\mathcal C$ is the same nonzero intermediate source residue as in
Version 56.  This is a forcing experiment for the retained coefficient
pencil, not a claim about the nonlinear original remainder.
\end{proposition}
\begin{proof}
The source is exactly $\mathcal D_ar_a=a^2h\cos(\omega_a\theta)$.
Version 56 gives its resonant term
$\theta\operatorname{Re}(e^{i\omega_a\theta}\mathcal Ch)$ with a
uniform $O_T(a)$ remainder under first-source-jet preparation.  Changing
that preparation to homogeneous fast correction data subtracts its
forward/backward propagated $F_0g+F_1\dot g$ endpoint values.  On the
intermediate modes these products are $O(a)$: the source residue is
$O(a^{-2})$, the source is $O(a^2)$, and each inverse frequency contributes
$a$.  The fastest products are $O(a^2)$ and their prepared hyperbolic or
definite oscillatory propagation is bounded.  The finite initial data
are zero in both cases.  The changed solution therefore has the same
leading resonant term.  Squaring and averaging the oscillation gives
\eqref{v57:eq:sharp} exactly as in the retained proof.
\end{proof}

There is no conflict with \cref{v57:thm:nonlinear}.  Small comparison time
makes its Lipschitz gain less than one, while only the reference feedback,
not every possible oscillatory argument, is required to satisfy
\eqref{v57:eq:reference-smooth}.  The theorem can consequently give an
$O(a)$ nonlinear solution without assuming an $O(a)$ operator norm on
all continuous sources.  Nor does it contradict Version 56's failure of
a generic $C^1$ source bound: the present unweighted source is restricted
to $\mathcal D_ar$, a much smaller amplitude-dependent bank.

\section{The exact remaining nonlinear problem}
\label{v57:sec:remaining}

The original system is now separated into three explicit objects:
\begin{equation}
 \boxed{
 M_ap'+Jp=-f+\mathcal D_a\mathcal F_a(Y,p)
              +(J-J_U(Y))p+f-f_U(Y).}
 \label{v57:eq:remaining}
\end{equation}
The first feedback on the right is completely derived above.  It contains
all rate-quadratic metric terms and all rate-linear metric connections;
it has a closed continuous-rate response on the retained pencil.

The last two terms are not included in that conclusion.  The weak rows of
$J_U-J$ can be $O(a)$ rather than $O(a^2)$.  In particular the invariant
kernel-compressed first skew coefficient of \cref{v55:prop:gyro} remains
part of the principal intermediate pencil.  It must not be treated as a
weighted small forcing merely because it tends to zero without row
weights.  The rate-independent $f_U$ must also be coupled to the original
state equations, rather than evaluated only on a convenient prescribed
curve.  Making both terms small along a comparison path is not a proof
of the corresponding statement along the sought history.

The next source calculation is therefore more focused: absorb or control
the actual state-dependent Poisson coefficient, and estimate the coupled
state-only drift in a response-compatible space.  The new theorem removes
the need to demand two derivatives of the unknown rate merely to handle
its Hessian-generated self-interactions.  It does not furnish nonlinear
transversality to the source-dependent prepared set or identify that set
with a native invariant family.

\section{Status, preservation, and verification}
\label{v57:sec:status}

The new exact original-source assertion is the factorization
\eqref{v57:eq:quadratic}, including the factor two and the sign of the
quadratic connection.  It proves the stronger weak-row weight without
computing or setting to zero any higher initializer coefficient.  The new
analytical assertion is the uniform $C^0\to C^0$ prepared response on
that weighted bank and its same-space nonlinear contraction.  An $O(a)$
prepared error follows when the feedback is regular only on the reference
path.  The sharpness experiment retains the same original coefficient
pencil and explains why small interval or small Lipschitz gain is still
needed.

No new numerical 107-dimensional native coefficient is claimed.  The
verifier checks exact differentiated-Hamiltonian identities, the complete
frame algebra, weighted block divisibility, and polynomial nonlinear
feedback identities.  Lower-dimensional examples are labelled algebraic
checks, not original trajectories.  Separate numerical oscillator tests
are diagnostics only.  The retained V56 verifier is replayed, with its
own source-regeneration and inheritance scopes unchanged.

Every Version-56 byte before its final document command is preserved.
New labels use the prefix \texttt{v57:}.  No nonlinear original $[0,aT]$
lift, larger original lifespan, spatial-refinement theorem, microscopic
preparation mechanism, or coupled Einstein--SM continuum is inferred.
The shorter original histories and the independent Version-32
stress-compensation problem remain at their stated scopes.

\begingroup
\renewcommand{\refname}{Additional source for Version 57}
\begin{thebibliography}{9}
\small\raggedright
\bibitem{v57:Gear}
C.~W.~Gear, \emph{Differential algebraic equations, indices, and integral
algebraic equations}, SIAM Journal on Numerical Analysis \textbf{27}
(1990), 1527--1534. \url{https://doi.org/10.1137/0727089}.
The reference supplies context for source differentiation in descriptor
perturbation estimates; the same-root weighted response and original
metric-source factorization are established in this body.
\end{thebibliography}
\endgroup
\addtocontents{toc}{\protect\endgroup}
% END IMMUTABLE VERSION 57 BODY
% APPEND VERSION 58 HERE, BEFORE THE FINAL END-DOCUMENT COMMAND.



% BEGIN IMMUTABLE VERSION 58 BODY
\clearpage
\begin{center}
{\Large\bfseries Version 58: Retain the moving Poisson term in the principal equation}\\[.5em]
{\large A companion audit, two-scale moving spectral bundles,\\
and a source-weighted prepared response}
\end{center}
\makeatletter
\addtocontents{toc}{\protect\begingroup}
\addtocontents{toc}{\protect\makeatletter}
\addtocontents{toc}{\protect\def\protect\l@section{\protect\bfseries\protect\@dottedtocline{1}{0em}{2.5em}}}
\makeatother
\addcontentsline{toc}{part}{Version 58: Moving principal coefficients and prepared response}

\section{Companion audit and the nonduplicating target}
\label{v58:sec:audit}

Version 57 isolates a favorable metric-feedback source, but its comparison
still freezes the Poisson matrix.  The term left outside that theorem can
have weak rows of order $a$, whereas its continuous-source estimate expects
weak rows of order $a^2$.  It must therefore be retained in the principal
operator rather than renamed a small weighted source.

The four new companions were read at their terminal developments and at the
dependencies used below.  This is not a new audit of every theorem in their
cumulative prefixes.  Their relevant contents and boundaries are as follows.
\begin{center}\small
\begin{tabularx}{\textwidth}{@{}>{\raggedright\arraybackslash}p{.21\textwidth}X@{}}
\toprule
Source & Transfer and boundary\\\midrule
Exact-Action Bridge V45 & Its V34 already proves a nonautonomous mixed
Green estimate with the moving Hessian retained.  V45 proves exact weak
nullity, two constrained Jordan chains, and a positive rank-four
intermediate pencil at its own V29 root.  That root has different signed
spectral data from the native V39 root used here.  Neither its two
oscillatory intermediate pairs nor its $ca^2$ histories are transferred.\\[3pt]
SM Source Selection V52 & The exact one-query optima on its specified
full source are $(3-2\sqrt2)/3$ without a quantum reference and
$(2-\sqrt3)/3$ with an invariant reference.  Its preparation-support and
noise tests separate positive executable preparation from signed
reconstruction.  They do not supply the present horizon-dependent initial
selection or a classical damping mechanism.\\[3pt]
Perturbative ESM V64 & One common first-order functional produces the
single insertion $\frac12\operatorname{STr}(\mathbb G H_1)$, its complete
metric derivatives and the linked graph-cut multiplicities.  The actual
source-complete critical coefficient and paired order-two master/reference
identity remain open there.  This is not a classical evolution estimate.\\[3pt]
YM V100\_f17 & Exact electric resonances prohibit eliminating the complete
one-link off-diagonal bank.  A local averaging step retains its resonant
part and bounds the remainder.  Its strong-electric window and comparison
gap are not a four-dimensional physical gap.  The useful direction is to
retain resonant principal data, not assume a whole-band inverse.\\
\bottomrule
\end{tabularx}
\end{center}
These statements are used at the scopes of the supplied sources
\cite{v58:EB45,v58:SM52,v58:PE64,v58:YM100}.  In particular the general
forward/backward Green construction is not claimed as new: the Einstein
companion already supplies it on a different one-scale split.  What must
be established here is a uniform split at \emph{both} fast scales of the
native V39 pencil, and the conversion back to its physical rate and source
weights.  The source polynomial of \cref{v57:thm:factor} is inherited.

The outcome is a prescribed-coefficient-path theorem.  It accommodates a
nonconstant first skew correction of arbitrary fixed bounded size.  Its
kernel compression remains in the intermediate pencil.  It does not
compute that compression for the actual higher initializer, or close the
state-dependent, rate-independent original drift.  No new original
$[0,aT]$ lifespan is asserted.  Spectral transport and nonself-adjoint
adiabatic methods are established context \cite{v58:Joye}; the estimates
needed here are proved below, without substituting an adiabatic-following
assertion for the mixed boundary problem.

\section{A precisely stated class of moving same-root pencils}
\label{v58:sec:class}

Keep the original reordered coefficients of \cref{v53:eq:blocks} and write
\begin{equation}
 M_a=\begin{pmatrix}M_x&aB_x\\aB_x^T&a^2C_{gg}\end{pmatrix},\qquad
 J=\begin{pmatrix}J_x&E\\-E^T&K_g\end{pmatrix},\qquad
 \mathcal D_a=\operatorname{diag}(aI_{81},a^2I_{26}).
 \label{v58:eq:base}
\end{equation}
Here $M_x\succ0$, $\operatorname{rank}K_g=24$, and the full symmetric
mass has inertia $(97,10,0)$ for $a>0$.  Let $Z$ synthesize $\ker K_g$.
The inherited matrices
\begin{equation}
 H=M_x^{-1}B_x,\quad \Delta=C_{gg}-B_x^TM_x^{-1}B_x,\quad
 D_0=Z^T\Delta Z,\quad G_0=Z^TE^TM_x^{-1}EZ
 \label{v58:eq:kernel-data}
\end{equation}
have respectively $\operatorname{Inertia}(D_0)=(1,1,0)$ and $G_0\succ0$.
The nonzero part of $-\Delta^{-1}K_g$ consists of nine directions in
each half-plane and a six-dimensional positive oscillatory space.  These
are the same-root inputs of Versions 41--53, not the spectral data of the
new Einstein companion.

Fix $T_0,R<\infty$.  The prescribed principal paths in this body are
\begin{equation}
 J_a(\theta)=J+aK(\theta)+a^2B_a(\theta),\qquad
 K^T=-K,\quad B_a^T=-B_a,
 \label{v58:eq:paths}
\end{equation}
on $0\le\theta\le T\le T_0$, with
\begin{equation}
 \|K\|_{W^{1,\infty}}+\sup_{0<a\le a_0}\|B_a\|_{W^{1,\infty}}
 \le R.
 \label{v58:eq:path-norm}
\end{equation}
All norms are finite-dimensional coordinate norms.  Coefficients may depend
on amplitude through $B_a$; no analyticity of the time path is required.
All assertions below hold for sufficiently small $a$, with a threshold
allowed to depend on $R,T_0$ and the fixed root data.  A source bound uniform
in spatial refinement, in $R$, or on unbounded intervals is not asserted.

The equation under consideration is
\begin{equation}
             M_a e_\theta+J_a(\theta)e=\mathcal D_ar(\theta).
 \label{v58:eq:equation}
\end{equation}
It is a nonautonomous coefficient equation, not by itself the original
nonlinear history.  The time remains $\theta=t/a$.

\begin{lemma}[Balancing converts the complete source to order $a$]
\label[lemma]{v58:lem:balance}
Put $D_a=\operatorname{diag}(I_{81},aI_{26})$, $u=D_ae$, and
\[
 \Gamma=\begin{pmatrix}M_x&B_x\\B_x^T&C_{gg}\end{pmatrix}.
\]
Then \eqref{v58:eq:equation} is exactly
\begin{equation}
 u_\theta=\mathscr A_a(\theta)u+a\Gamma^{-1}r,\qquad
 \mathscr A_a=-\Gamma^{-1}D_a^{-T}J_aD_a^{-1},\qquad
 \mathscr A_a^T\Gamma+\Gamma\mathscr A_a=0.
 \label{v58:eq:balanced}
\end{equation}
The balanced matrix has the expansion
\begin{equation}
 \mathscr A_a=a^{-2}\mathscr A_{-2}
              +a^{-1}\mathscr A_{-1}(\theta)+\mathscr A_0(a,\theta),
 \label{v58:eq:levels}
\end{equation}
where
\[
 \mathscr A_{-2}=-\Gamma^{-1}\begin{pmatrix}0&0\\0&K_g\end{pmatrix},\qquad
 \mathscr A_{-1}=-\Gamma^{-1}
             \begin{pmatrix}0&E\\-E^T&K_{gg}(\theta)\end{pmatrix},
\]
and $\mathscr A_0$ is uniformly bounded in $W^{1,\infty}$.
\end{lemma}
\begin{proof}
$M_a=D_a^T\Gamma D_a$ and $D_a^{-T}\mathcal D_a=aI$.
Multiplication of \eqref{v58:eq:equation} by $D_a^{-T}$ proves the first
formula.  Its skew identity follows from $J_a^T=-J_a$.  Block expansion
proves \eqref{v58:eq:levels}, including the order-$a^{-1}$ term
$K_{gg}$.  No inverse norm of $M_a$ is used.
\end{proof}

\section{The first skew correction survives in a uniform intermediate split}
\label{v58:sec:split}

\begin{proposition}[Two-scale spectral structure for every bounded skew path]
\label[proposition]{v58:prop:spectrum}
For \eqref{v58:eq:paths}--\eqref{v58:eq:path-norm}, the instantaneous
balanced matrix has the following uniform small-$a$ splitting.
Twenty-four directions have eigenvalues of order $a^{-2}$, with dimensions
$9,9,6$ in the right half-plane, left half-plane, and imaginary axis.
Four directions have eigenvalues of order $a^{-1}$: one positive real,
one negative real, and one imaginary pair.  The other 79 directions have
bounded eigenvalues.  The union of all imaginary/zero groups, of dimension
87, is uniformly positive for $\Gamma$ in balanced coordinates.

The leading four intermediate roots are the roots of
\begin{equation}
 \det\{\nu^2D_0+\nu\Sigma(\theta)+G_0\}=0,\qquad
 \Sigma(\theta)=S_0+Z^TK_{gg}(\theta)Z,
 \quad S_0=Z^T(E^TH-H^TE)Z.
 \label{v58:eq:middle}
\end{equation}
They are simple and uniformly separated from zero and from one another
on the bounded coefficient class.  In particular the first skew
compression is retained, not discarded.  The complete instantaneous
matrix has ten growing and ten decaying directions, 43 imaginary pairs,
and one semisimple zero.  Internal coincidences within the 79-dimensional
finite group need not be resolved.
\end{proposition}
\begin{proof}
The kernel of $\mathscr A_{-2}$ is
$\mathcal K=\{(x,Zy):x\in\R^{81},y\in\R^2\}$.
The restriction of $\Gamma$ to it is
\[
 \Gamma_{\mathcal K}=
 \begin{pmatrix}M_x&B_xZ\\Z^TB_x^T&Z^TC_{gg}Z\end{pmatrix},
\]
with inertia $(82,1,0)$, by its Schur complement $D_0$.
It is nondegenerate.  For a $\Gamma$-skew matrix, its range is the
$\Gamma$-orthogonal complement of its kernel; nondegeneracy therefore
excludes a Jordan chain at zero.  The other 24 eigenvalues and their
positive oscillatory restriction are inherited.  Uniform resolvent
bounds on fixed contours around these nonzero scaled clusters give the
24-dimensional spectral block of $a^2\mathscr A_a$.  Its complementary
83-dimensional block and its identification with $\mathcal K$ are
$O(a)$ close in $W^{1,\infty}$.

Restrict the next skew form to $\mathcal K$.  Its matrix is
$\left(\begin{smallmatrix}0&EZ\\-Z^TE^T&Z^TK_{gg}Z\end{smallmatrix}\right)$.
The shear $x=\widehat x-HZy$ diagonalizes its mass as $(M_x,D_0)$ and
adds $S_0$ to the lower skew block.  Eliminating $\widehat x$ at a nonzero
scaled eigenvalue $\nu$ gives exactly \eqref{v58:eq:middle}.
Let $\Sigma=\left(\begin{smallmatrix}0&\chi\\-\chi&0\end{smallmatrix}\right)$.
Its determinant is
\begin{equation}
 (\det D_0)\nu^4+
 [\operatorname{tr}(\operatorname{adj}D_0G_0)+\chi^2]\nu^2
 +\det G_0.
 \label{v58:eq:quartic}
\end{equation}
The leading coefficient is negative and the constant positive.  The two
roots in $\nu^2$ are simple, nonzero, and of opposite signs.  Bounded
$\chi$ confines them to a compact separated set.  Their simple real and
imaginary pairs persist for small $a$, including their reality, by the
real and signed-skew symmetries.

The zero kernel at this second reduction is
$\{(x,0):(EZ)^Tx=0\}$, of dimension 79 and positive for $M_x$.
It is nondegenerate and hence semisimple.  Its orthogonal complement in
$\mathcal K$ has signature $(3,1)$.  The real intermediate pair accounts
for $(1,1)$; its oscillatory plane is therefore positive.  The inherited
six fastest oscillatory directions are positive as well.  Separated
spectral subspaces for a signed-skew matrix are orthogonal except when
their spectral values sum to zero.  The combined 87-dimensional
nonhyperbolic subspace is thus positive, with a uniform margin obtained
from its compact limiting subspaces.  On it the exact matrix is similar
to a skew-symmetric matrix, so its eigenvalues are imaginary and
semisimple.  The finite eigenvalues converge to the regular descriptor
spectrum because $zM_a+J_a\to zM_0+J$ on fixed contours.  Finally,
$\operatorname{rank}J=106$ and a fixed nonzero minor persist uniformly;
odd-dimensional skew symmetry gives the reverse rank bound.  Hence the
exact zero is unique.  This proves all dimensions without assuming a
positive full mass or counting rounded eigenvalues.
\end{proof}

\begin{lemma}[Bounded moving coordinates, including their derivatives]
\label[lemma]{v58:lem:bundle}
There are invertible, absolutely continuous coordinates $u=V_a(\theta)z$
with $V_a,V_a^{-1}$ and $V_a'$ uniformly bounded, which split the
instantaneous matrix into its seven groups (three fastest, three
intermediate, and one finite).  In these coordinates
\begin{equation}
 z'=\Lambda_a(\theta)z+C_a(\theta)z+aB_a^\sharp(\theta)r,
 \quad C_a=-V_a^{-1}V_a',\quad
 B_a^\sharp=V_a^{-1}\Gamma^{-1},
 \label{v58:eq:bundle}
\end{equation}
with block-diagonal $\Lambda_a$ and bounded $C_a,B_a^\sharp$.
The finite 79-dimensional columns satisfy the stronger estimate
\begin{equation}
       \|D_a^{-1}V_{a,s}\|\le C.
 \label{v58:eq:finite-synthesis}
\end{equation}
With $W^{2,\infty}$ principal paths, the corresponding second frame
derivatives are uniformly bounded.  In that case the inverse of the
full 28-dimensional fast diagonal block, even after adding its bounded
frame connection, and its first time derivative are both $O(a)$.
\end{lemma}
\begin{proof}
Use the two rescaled resolvents in the preceding proof, not a resolvent
bound across all diverging eigenvalues at once.  The first spectral
identification is $I+O(a)$ with derivative $O(a)$.  On the 83-dimensional
block divide the restricted matrix by its single remaining $a^{-1}$
factor.  Its four nonzero eigenvalues are uniformly separated; their
projections and first time derivatives are bounded by contour inversion.
The 79-dimensional limiting kernel is independent of $K_{gg}$ and lies
in $w=0$.  Its exact bundle differs by $O(a)$ with the same regularity.
This proves \eqref{v58:eq:finite-synthesis}.

For any complete family of the resulting projections $P_j$, the bounded
connection $\mathcal A=\sum_jP_j'P_j$ obeys
$[\mathcal A,P_j]=P_j'$.  Its matrix ODE transports fixed endpoint
ranges to the moving ranges.  Starting with uniformly conditioned bases
produces the stated $V_a$ and bounded inverse.  The bounds are uniform
on $[0,T_0]$; local coordinate charts can equivalently be concatenated.
Repeated contour differentiation proves the $W^{2,\infty}$ assertion.
The fast diagonal inverse has orders $a^2$ and $a$ on its two groups.
Their inverse derivatives have those same orders.  A bounded frame
connection is absorbed by a Neumann series; differentiating that series
preserves the $O(a)$ bound.  No inverse of a finite or zero frequency is
taken.  Substitution in $u'=\mathscr A_au+a\Gamma^{-1}r$ proves
\eqref{v58:eq:bundle}, so the full frame connection is retained.
\end{proof}

\section{Uniform prepared response along the moving coefficient path}
\label{v58:sec:green}

Let $P_+(a,\theta),P_-(a,\theta),P_c(a,\theta)$ be the instantaneous
balanced projections onto growing, decaying, and nonhyperbolic spaces.
Their ranks are $10,10,87$.  Impose
\begin{equation}
 P_-(a,0)u(0)=P_c(a,0)u(0)=0,\qquad
 P_+(a,T)u(T)=0,\qquad u=D_ae.
 \label{v58:eq:BC}
\end{equation}
These are mixed conditions on the complete correction, not a physical
operation or a deletion of its oscillatory modes.  They reduce to
Version 57's conditions when the principal matrix is constant.

\begin{theorem}[Moving-principal continuous-source response]
\label[theorem]{v58:thm:green}
For every principal path \eqref{v58:eq:paths}--\eqref{v58:eq:path-norm}
and continuous $r$, the full equation \eqref{v58:eq:equation} with
\eqref{v58:eq:BC} has a unique solution for sufficiently small $a>0$.
For $0<T\le T_0$,
\begin{equation}
 \boxed{\ \|e\|_{C^0}\le C_{R,T_0}(T+a)\|r\|_{C^0}.\ }
 \label{v58:eq:green-bound}
\end{equation}
In balanced coordinates its hyperbolic components are $O(a^2\|r\|)$;
the nonhyperbolic components are $O(aT\|r\|)$, with constants as above.
No derivative of the source or unknown rate occurs in this estimate.
\end{theorem}
\begin{proof}
Use \cref{v58:lem:bundle}.  On the positive nonhyperbolic block the
restriction of $V_a^T\Gamma V_a$ is uniformly positive with bounded
derivative.  Its principal matrix is skew for that metric.  Including
the diagonal frame connection, its propagator is bounded by $C e^{C|t-s|}$.
This estimate charges metric and frame motion, not the frequencies
$a^{-2}$ and $a^{-1}$.

On the fastest stable block choose a positive Lyapunov metric for the
\emph{rescaled} stable generator.  The inherited uniform gap gives
uniform metric, inverse and derivative bounds through its integral
formula or its Sylvester equation.  On the intermediate stable line the
same conclusion follows directly from the separated negative real root.
After restoring their time factors, both are bounded by
$Ce^{-\kappa(t-s)/a}$ forward, for small $a$.  The bounded connection
between them is absorbed by reducing $a_0$.  Unstable blocks have the
same estimate backward.  All diagonal frame terms are included.

Let $h$ collect the twenty hyperbolic variables and $c$ the other 87.
For a proposed $h$, solve the central Volterra equation exactly, including
its center--center connection.  With zero initial value it obeys
\[
 \|c[h]\|_\infty\le C_{R,T_0}T(\|h\|_\infty+a\|r\|_\infty).
\]
The stable and unstable integral equations, with their prescribed opposite
endpoints, have kernel integrals at most $Ca$.  Their mutual bounded
couplings give
\[
 \|h\|_\infty\le Ca(\|h\|_\infty+\|c[h]\|_\infty+a\|r\|_\infty).
\]
For sufficiently small $a$, substitution is a contraction on $C^0$.
It gives
\[
 \|h\|_\infty\le C_{R,T_0}a^2\|r\|_\infty,\qquad
 \|c\|_\infty\le C_{R,T_0}aT\|r\|_\infty.
\]
The same estimates for differences prove uniqueness.  Differentiating
these integral equations verifies every differential row and the mixed
conditions.  Finally $\|D_a^{-1}\|=a^{-1}$ for $0<a\le1$ and $V_a$ is
bounded, which proves \eqref{v58:eq:green-bound}.  No assertion that a
solution follows the instantaneous spectral spaces is used: all their
mutual connection terms were solved in the integral equations.
\end{proof}

\begin{proposition}[An order-$a$ response on a regular reference source]
\label[proposition]{v58:prop:smooth}
If the coefficient bounds are in $W^{2,\infty}$ and $r\in W^{1,\infty}$,
the same prepared solution satisfies
\begin{equation}
 \boxed{\ \|e\|_{C^0}\le C_{R,T_0}a\|r\|_{W^{1,\infty}}.\ }
 \label{v58:eq:smooth-bound}
\end{equation}
This is a sufficient regularity statement, not a claim of optimal
coefficient or source differentiability.
\end{proposition}
\begin{proof}
In the moving frame split the finite 79 coordinates $s$ from the fast
28 coordinates $f$.  Include the bounded diagonal connection in their
principal equations and call the fast matrix $A_f$.  Its inverse and
inverse derivative are $O(a)$ by \cref{v58:lem:bundle}.  For the finite
Volterra equation,
\[
 \|s\|_\infty+\|s'\|_\infty
 \le C_{R,T_0}(a\|r\|_\infty+\|f\|_\infty).
\]
The fast equation is
$f'=A_ff+C_{fs}s+aB_fr$.  Put
$k=-A_f^{-1}(C_{fs}s+aB_fr)$.  It obeys
\[
 \|k\|_\infty+\|k'\|_\infty
 \le C_{R,T_0}a(\|s\|_\infty+\|s'\|_\infty
                      +a\|r\|_{W^{1,\infty}}).
\]
The fast correction $f-k$ has source $-k'$ and endpoint data given by
$-k$ on its stable/oscillatory initial and unstable terminal subspaces.
The preceding Green argument, now with only eight positive oscillatory
central coordinates, bounds its solution by
$C_{R,T_0}(\|k\|_\infty+\|k'\|_\infty)$.
Thus
\[
 \|f\|_\infty\le C_{R,T_0}
                (a^2\|r\|_{W^{1,\infty}}+a\|f\|_\infty).
\]
Absorb the last term.  Then $f=O(a^2)$ and $s=O(a)$ in the indicated
source norm.  The finite synthesis \eqref{v58:eq:finite-synthesis} costs
no inverse power of $a$; the fast synthesis costs at most $a^{-1}$.
Returning to $e$ proves \eqref{v58:eq:smooth-bound}.  The derivatives
used here belong to the prescribed coefficients and reference source,
not to a nonlinear unknown iterate.
\end{proof}

\section{Metric feedback with a nonconstant principal path}
\label{v58:sec:nonlinear}

For the same prescribed $J_a(\theta)$, let $p_{b,a}$ be a prepared
reference for $M_ap_b'+J_ap_b=g_{b,a}$.  Use the whole weighted metric
polynomial $\mathcal F_a$ of \cref{v57:prop:weighted} along a prescribed
regular state/coefficient path, or any continuous polynomial with the
same bounds.  It is not replaced by a selected quadratic summand.

\begin{theorem}[A moving-principal same-space metric contraction]
\label[theorem]{v58:thm:nonlinear}
On a tube $\|p-p_b\|_\infty\le r_0$, suppose the feedback has bound
$B_{r_0}$ and rate Lipschitz constant $L_{r_0}$.  If
\begin{equation}
 C_{R,T_0}(T+a)B_{r_0}\le r_0,\qquad
 \kappa=C_{R,T_0}(T+a)L_{r_0}<1,
 \label{v58:eq:contract}
\end{equation}
then
\begin{equation}
 M_ap'+J_a(\theta)p
       =g_{b,a}+\mathcal D_a\mathcal F_a(\theta,p)
 \label{v58:eq:semilinear}
\end{equation}
has a unique solution in that tube with \eqref{v58:eq:BC} imposed on
$p-p_b$.  With $W^{2,\infty}$ coefficient bounds and
$\|\mathcal F_a(\cdot,p_b)\|_{W^{1,\infty}}\le B_*$, it satisfies
\begin{equation}
 \boxed{\ \|p-p_b\|_\infty\le
             \frac{C_{R,T_0}aB_*}{1-\kappa}.\ }
 \label{v58:eq:semilinear-error}
\end{equation}
\end{theorem}
\begin{proof}
Let $\mathcal G_{a,K,B}$ be the linear Green operator just proved.
The map $p\mapsto p_b+\mathcal G_{a,K,B}\mathcal F_a(\cdot,p)$ is a
contraction of the stated complete $C^0$ ball by \eqref{v58:eq:contract}.
Its fixed point is a solution of every row, and the linear uniqueness
proves the nonlinear uniqueness in the tube.  Apply
\cref{v58:prop:smooth} only to $\mathcal F_a(\cdot,p_b)$ and apply the
continuous-source estimate to the feedback difference.  Absorption gives
\eqref{v58:eq:semilinear-error}.  No uniform derivative of the unknown
nonlinear rate is assumed.
\end{proof}

This extends the metric-feedback comparison of Version 57 without
freezing the first moving skew coefficient.  For each fixed coefficient
bound it supplies a fixed sufficiently short comparison interval.
It remains a prescribed-path statement: assigning its coefficients from
an unknown state makes an additional, genuinely coupled problem.

\section{Why pathwise bounds do not automatically give a coupled contraction}
\label{v58:sec:sensitivity}

The companion's warning to retain resonances also applies to coefficient
\emph{sensitivity}.  A uniform bound along each path is not a uniform
Lipschitz bound with respect to those paths.

\begin{proposition}[An exact coefficient-sensitivity diagnostic]
\label[proposition]{v58:prop:sensitivity}
Consider the declared two-dimensional positive comparison, written as a
complex scalar,
\begin{equation}
 z'+\frac{i(1+\eta)}a z=e^{-i\theta/a},\qquad z(0)=0.
 \label{v58:eq:test}
\end{equation}
It has mass $a^2I_2$, skew coefficient $a(1+\eta)
\left(\begin{smallmatrix}0&-1\\1&0\end{smallmatrix}\right)$,
and weighted covector of size $a^2$.  Compare $\eta=0$ with $\eta=a$.
Their first skew-coefficient paths differ by $a$ in every $C^k$ norm,
and their actual skew matrices differ by $a^2$.  Nevertheless at $\theta=1$,
\begin{equation}
 |z_a(1)-z_0(1)|\ge1-\cos1>\frac{11}{24}.
 \label{v58:eq:sensitivity}
\end{equation}
\end{proposition}
\begin{proof}
The exact solutions are
\[
 z_0(\theta)=\theta e^{-i\theta/a},\qquad
 z_a(\theta)=e^{-i\theta/a}\frac{1-e^{-i\theta}}i.
\]
The imaginary part of their phase-removed difference at one has magnitude
$1-\cos1$.  The alternating Taylor bound
$\cos1<1-1/2+1/24$ gives the strict inequality.  Substitution verifies
both original equations.  The source amplitude in the equation after
division by $a^2$ is uniformly one, but its time derivatives are not
uniformly bounded.  The example therefore does not contradict the
regular-reference estimate.  It is a coefficient comparison, not a
native Einstein trajectory or an assertion of its actual forcing.
\end{proof}

The example excludes a general coefficient-path contraction based only
on an unscaled small matrix difference and a bounded continuous source.
It neither excludes a structured native contraction nor asserts that
the original nonlinear feedback excites this chosen resonance.  The
new estimates do not permit an unevaluated first skew compression to be
set to zero; it changes \eqref{v58:eq:middle} at leading order.

\section{Original-source application test, preservation, and remaining work}
\label{v58:sec:status}

The exact transformed original equation is still
\[
 M_ap'+J_U(Y)p+f_U(Y)=\mathcal D_a\mathcal F_a(Y,p),
 \qquad Y_\theta=(aF,\ aS_a\mathcal U_ap).
\]
For an independently supplied bounded normalized coefficient path, the
new principal theorem applies once the actual skew coefficient is shown
to have the decomposition \eqref{v58:eq:paths} with the stated uniform
time bounds.  Smoothness at each fixed $a$ is not that uniform assertion.
One must differentiate through the same constrained state path, retain
the full kernel compression, and retain the original state-only $f_U$.
A path chosen for a comparison need not be generated by the displayed
state equations.  Applying the new theorem along such a path does not
produce an original history by relabeling it.

The coefficient sensitivity in \cref{v58:prop:sensitivity} identifies
why the remaining coupled step is not immediate.  The source polynomial
has a uniform rate-Lipschitz bound, but changing the unknown state also
changes the rapidly oscillating principal propagator and its endpoint
subspaces.  The actual state response, state-only source and the
transversality of consistent initial records to this moving prepared
set must still be estimated together.  The frame and endpoint selections
are analytical constructions depending on the full path; no causal
selection or future physical intervention is claimed.

\paragraph{Proved.}
The balanced same-root pencil retains a uniform two-scale split for every
bounded first-order skew path.  Its intermediate gyro is kept exactly.
The complete nonautonomous prepared equation has a uniform $C^0$ response
on the physical weighted source bank, and an $O(a)$ bound for regular
reference sources.  The complete metric-feedback comparison then closes
on the same $C^0$ space with the principal path moving.

\paragraph{Inherited and not recomputed.}
The native root, signed mass, ranks, definite oscillatory subspaces,
and descriptor positivity are the stated Versions 39--57 inputs.
The general mixed Green argument was already present in the Einstein
companion's V34; the added work here is the two-level bundle reduction,
87-dimensional positive center, physical source balancing, and smooth
reference conversion.  No new numerical native Poisson coefficient,
root, frequency, or full coefficient bank is claimed.

\paragraph{Still open.}
There is no new original $[0,aT]$ lifespan, cofinal spatial estimate,
microscopic preparation mechanism, or coupled Einstein--SM continuum
result.  The shorter selected original histories and the independent
Version-32 stress question are unchanged.  The next nonduplicating task
is the actual constrained coefficient-path and state-only source test,
not another frozen spectral count.

\paragraph{Verification.}
The accompanying checker tests exact balanced expansions, signed
restrictions, rank-four intermediate factorization, the gyro quartic,
the positive-complement signature count, the integral-inequality algebra,
and the coefficient-sensitivity solution.  Separately labelled small
model computations test a moving intermediate gyro and the prepared
boundary solve.  Such models are diagnostics, not substitutes for the
native tensors.  The parent replay and its scope are recorded separately.
Every byte of Version 57 before its final document command is preserved;
new labels have prefix \texttt{v58:}.  No proof-assistant formalization
or numerical original spacetime evolution is asserted.

\begingroup
\renewcommand{\refname}{Sources used in Version 58}
\begin{thebibliography}{9}
\small\raggedright
\bibitem{v58:EB45}
\emph{Renewal Exact-Action Einstein Bridge Research Notes V45}, supplied
file \nolinkurl{Renewal_Exact_Action_Einstein_Bridge_Research_Notes_V45(3).tex}.
Used at V34, \texttt{v34:green}, and V45, \texttt{v45:weak-rank},
\texttt{v45:frequency-formula}, \texttt{v45:status}; different root retained.
\bibitem{v58:SM52}
\emph{Renewal SM Source Selection Research Notes V52}, supplied file
\nolinkurl{Renewal_SM_Source_Selection_Research_Notes_V52(3).tex},
\texttt{sel52:thm:global}, \texttt{sel52:thm:prepface},
\texttt{sel52:prop:bank}, \texttt{sel52:status}.
\bibitem{v58:PE64}
\emph{Renewal Perturbative ESM Continuum Research Notes V64}, supplied file
\nolinkurl{Renewal_Perturbative_ESM_Continuum_Research_Notes_V64(3).tex},
\texttt{thm:v64-insertion}, \texttt{prop:v64-marks},
\texttt{sec:v64-ledgers}.
\bibitem{v58:YM100}
\emph{Renewal Yang--Mills Mass-Gap Research Notes V100\_f17}, supplied file
\nolinkurl{Renewal_Yang_Mills_Mass_Gap_Research_Notes_V100_f17(2).tex},
\texttt{r100:no-inverse}, \texttt{r100:average-lemma},
\texttt{r100:local-normal-theorem}.
\bibitem{v58:Joye}
A.~Joye, \emph{General Adiabatic Evolution with a Gap Condition},
Communications in Mathematical Physics \textbf{275} (2007), 139--162,
\url{https://doi.org/10.1007/s00220-007-0299-y};
\url{https://arxiv.org/abs/math-ph/0608059}.
Context for transporting separated nonself-adjoint spectral groups;
its analytic adiabatic theorem is not substituted for the finite
$W^{1,\infty}$ mixed-boundary proof above.
\end{thebibliography}
\endgroup
\addtocontents{toc}{\protect\endgroup}
% END IMMUTABLE VERSION 58 BODY
% APPEND VERSION 59 HERE, BEFORE THE FINAL END-DOCUMENT COMMAND.


% BEGIN IMMUTABLE VERSION 59 BODY
\clearpage
\begin{center}
{\Large\bfseries Version 59: State-integrated principal feedback}\\[.5em]
{\large A one-derivative Green estimate, fixed reference endpoint tests,\\
and a coupled continuous-rate preparation theorem}
\end{center}
\makeatletter
\addtocontents{toc}{\protect\begingroup}
\addtocontents{toc}{\protect\makeatletter}
\addtocontents{toc}{\protect\def\protect\l@section{\protect\bfseries\protect\@dottedtocline{1}{0em}{2.5em}}}
\makeatother
\addcontentsline{toc}{part}{Version 59: Integrated-state coefficient feedback and its preparation}

\section{The remaining coefficient dependence is not an arbitrary time path}
\label{v59:sec:scope}

Version 58 proves a Green estimate for a prescribed moving principal matrix,
but explicitly leaves its dependence on the unknown state open.  Its resonant
coefficient-sensitivity example does not exclude a contraction for coefficients
generated by an integrated state.  This body exploits precisely that
distinction.  It first sharpens the finite--fast frame estimate and reduces
the regular-reference assumption from two bounded coefficient derivatives to
one.  It then proves a different source estimate for an order-$a$ covector
which is a time primitive.  Finally these two estimates close a genuinely
coupled state/principal/metric comparison on a continuous-rate ball.

The coupled theorem has explicit \emph{pointwise coefficient hypotheses};
they are not asserted without verification for the full original Poisson and
state-only source maps.  The original kinematic map does satisfy its regular
state-equation hypotheses, as proved below.  The metric polynomial is inherited
from Version 57.  An evaluation of the remaining divided native coefficient
jets and simultaneous original initial consistency is still needed before
the coupled comparison can be called an original $[0,aT]$ history.
There is no new spatial-refinement or physical-lifespan conclusion.

All same-root spectral data, the two-scale separation, the positive
87-dimensional nonhyperbolic metric, and the continuous weighted-source
estimate are inherited from Version 58.  The new proofs use them, not a
recomputation of the root or its Hessian.  A targeted review of the supplied
Einstein Bridge V45 and the preceding native bodies found the mixed Green
construction, but not the integrated-state comparison proved here.
The nonautonomous spectral-transport warning in \cite{v58:Joye} remains in
force: instantaneous subspaces are not invariant trajectories.  Every frame
connection and every endpoint condition is retained below.

\section{A refined finite--fast frame uses only one coefficient derivative}
\label{v59:sec:frame}

Keep
\[
 D_a=\operatorname{diag}(I_{81},aI_{26}),\qquad
 \mathcal D_a=\operatorname{diag}(aI_{81},a^2I_{26}),\qquad
 M_a=D_a^T\Gamma D_a .
\]
The fixed $\Gamma$ has inertia $(97,10,0)$.  Principal paths have the form
\begin{equation}
 J_a(\theta)=J+aK_a(\theta)+a^2B_a(\theta),\qquad
 K_a^T=-K_a,\quad B_a^T=-B_a,
 \quad \|K_a\|_{W^{1,\infty}}+\|B_a\|_{W^{1,\infty}}\le R .
 \label{v59:eq:principal}
\end{equation}
Allowing bounded amplitude dependence in $K_a$ changes none of the
two rescaled contour arguments of Version 58.  Constants below may depend on
$R,T_0$ and the fixed root, never on $0<a\le a_0$ or $0<T\le T_0$.
Primes denote $d/d\theta$.

Let $P_s$ be the balanced spectral projector on the 79 finite modes, and
$P_f=I-P_s$ on all 28 fast modes.  Write
\[
 \mathscr A_a=-\Gamma^{-1}D_a^{-T}J_aD_a^{-1}.
\]
The seven-group moving frame of Version 58 can be chosen with finite
coordinates $s$ and fast coordinates $z_f$.

\begin{lemma}[Small finite--fast connection and weak-source projection]
\label[lemma]{v59:lem:frame}
One may choose a uniformly bounded frame and inverse, with bounded first
derivatives, so that
\begin{align}
 s'&=A_s s+C_{sf}z_f+b_s, \notag\\
 z_f'&=\Lambda_f z_f+C_{ff}z_f+C_{fs}s+b_f ,
 \label{v59:eq:frame}
\end{align}
where
\begin{equation}
 \|A_s\|+\|C_{ff}\|\le C,\qquad
 \|C_{sf}\|+\|C_{fs}\|\le Ca,\qquad
 \|\Lambda_f^{-1}\|+\|(\Lambda_f^{-1})'\|\le Ca .
 \label{v59:eq:frame-bounds}
\end{equation}
No derivative of $C_{ff},C_{sf},C_{fs}$ is required.  The finite synthesis
satisfies $\|D_a^{-1}V_s\|\le C$.  With
$\iota_w h=(0,h)$, one also has
\begin{equation}
 \|P_s\Gamma^{-1}\iota_w\|_{W^{1,\infty}}\le Ca .
 \label{v59:eq:weak-projection}
\end{equation}
\end{lemma}
\begin{proof}
The two contour reductions in \cref{v58:lem:bundle} give
$P_s=P_s^0+O_{W^{1,\infty}}(a)$, where the limiting range is
\[
 \mathcal S_0=\{(x,0):(EZ)^Tx=0\}.
\]
Both this range and its $\Gamma$-orthogonal projector $P_s^0$ are independent
of $K_{a,gg}$.  Nondegeneracy follows from $M_x\succ0$ on this range.
For $s_0=(x,0)\in\mathcal S_0$,
$s_0^T\Gamma\Gamma^{-1}\iota_w h=s_0^T\iota_w h=0$.
Thus $P_s^0\Gamma^{-1}\iota_w=0$, proving
\eqref{v59:eq:weak-projection} including its first derivative.

Transport all seven projectors with the complete Kato connection used in
Version 58.  Off-diagonal frame entries between the \emph{union} $P_s$ and
its complement are controlled by $P_s'=O(a)$.  Indeed, for two distinct
transported groups, differentiation of $P_iP_j=0$ gives
$P_iP_j'=-P_i'P_j$.  Summing this identity over the fast groups proves
both finite--fast bounds in \eqref{v59:eq:frame-bounds}; internal fast
connections need only be bounded.  The instantaneous finite block is
bounded.  The instantaneous fast diagonal blocks have inverse orders
$a^2$ and $a$, with the same orders for their inverse derivatives by the
rescaled contour calculus and first coefficient derivatives.  Put only these
instantaneous blocks in $\Lambda_f$ and leave the whole bounded frame
connection in $C_{ff}$.  This proves the remaining assertions without ever
differentiating a frame connection.  The finite-synthesis bound is the
retained \cref{v58:eq:finite-synthesis}.
\end{proof}

For later use the fast subsystem, including $C_{ff}$, has a mixed Green
bound
\begin{equation}
 \|z_f\|_\infty\le C\bigl(
 |\zeta_{\rm end}|+\|b_f\|_{L^1}
 \bigr).
 \label{v59:eq:fast-L1}
\end{equation}
Here the eight oscillatory and ten decaying data are at the initial
endpoint and the ten growing data at the final endpoint.
To prove \eqref{v59:eq:fast-L1}, use the bounded positive propagator on the
eight oscillatory coordinates and the forward/backward decaying kernels
on the hyperbolic coordinates.  An $L^1$ forcing is bounded by its $L^1$
norm on each kernel.  Mutual bounded couplings are solved as in the
Version-58 integral equations: the hyperbolic integrations contribute
$Ca$, and the oscillatory Volterra equation contributes a fixed $C_{T_0}$.
Small $a$ absorbs the hyperbolic return.  The same proof treats endpoint
data.  This estimate does not assert decay of the oscillatory modes.

\section{One derivative suffices for a regular weighted source or a primitive source}
\label{v59:sec:sources}

Initially impose the instantaneous mixed preparation of Version 58 on the
balanced correction $u=D_ae$:
\begin{equation}
 (P_-+P_c)(0)u(0)=0,\qquad P_+(T)u(T)=0 .
 \label{v59:eq:instant-bc}
\end{equation}
Here $P_c$ includes all 87 nonhyperbolic directions.
For an absolutely continuous vector $h$, put
\[
 \|h\|_{\mathsf A_T}=\|h\|_\infty+\|h'\|_{L^1(0,T)}.
\]

\begin{theorem}[A Lipschitz-principal, one-source-derivative estimate]
\label[theorem]{v59:thm:source}
For principal paths \eqref{v59:eq:principal} the prepared solutions satisfy
\begin{align}
 M_ae'+J_ae=\mathcal D_ar,\quad r\in W^{1,1}
 &\quad\Longrightarrow\quad
 \|e\|_\infty\le Ca\|r\|_{\mathsf A_T}, \label{v59:eq:weighted-A}\\
 M_ae'+J_ae=ah,\quad h\in W^{1,1}
 &\quad\Longrightarrow\quad
 \|e\|_\infty\le C\|h\|_{\mathsf A_T}. \label{v59:eq:primitive-A}
\end{align}
For a continuous weighted amplitude one retains
$\|e\|_\infty\le C(T+a)\|r\|_\infty$.
The coefficient assumptions are only $W^{1,\infty}$ in all these estimates.
\end{theorem}
\begin{proof}
For $\mathcal D_ar$ the balanced forcing is $a\Gamma^{-1}r$.
In \eqref{v59:eq:frame}, its finite and fast parts are
$O(a\|r\|)$, with
$\|b_f'\|_{L^1}\le Ca(\|r'\|_{L^1}+T\|r\|_\infty)$.
Let $k=-\Lambda_f^{-1}b_f$.  Then
\[
 \|k\|_\infty+\|k'\|_{L^1}\le Ca^2\|r\|_{\mathsf A_T}.
\]
The fast equation for $z_f-k$ has forcing
$-k'+C_{ff}k+C_{fs}s$ and endpoint values from $-k$.
Crucially, no derivative of $C_{fs}s$ is taken.
By \eqref{v59:eq:fast-L1},
\[
 \|z_f\|_\infty\le
 C a^2\|r\|_{\mathsf A_T}+CaT\|s\|_\infty .
\]
The finite Volterra equation gives
$\|s\|_\infty\le CaT(\|r\|_\infty+\|z_f\|_\infty)$.
Absorbing their small return proves $z_f=O(a^2\|r\|_{\mathsf A_T})$
and $s=O(a\|r\|_{\mathsf A_T})$.
Only the fast synthesis pays $a^{-1}$, proving \eqref{v59:eq:weighted-A}.

For the source $ah$, the balanced forcing is
$\Gamma^{-1}(ah_x,h_w)$.  Its finite component is $O(a\|h\|)$ by
\eqref{v59:eq:weak-projection}; its fast component and first derivative
have the unweighted bounds $O(\|h\|)$ and
$O(\|h'\|+\|h\|)$.  Thus the same $k$ now obeys
\[
 \|k\|_\infty+\|k'\|_{L^1}\le Ca\|h\|_{\mathsf A_T}.
\]
The two estimates become
\[
 \|z_f\|_\infty\le Ca\|h\|_{\mathsf A_T}+CaT\|s\|_\infty,\qquad
 \|s\|_\infty\le CaT(\|h\|_\infty+\|z_f\|_\infty).
\]
They prove \eqref{v59:eq:primitive-A} after physical synthesis.
The continuous weighted estimate is \cref{v58:thm:green}.
\end{proof}

The first assertion sharpens, rather than contradicts,
\cref{v58:prop:smooth}: the small finite--fast connection avoids its use
of a second coefficient derivative.  The second assertion concerns a
different source normalization.  It does not supply a uniformly bounded
response for arbitrary unweighted continuous covectors.

\begin{corollary}[A primitive has a short-time gain]
\label[corollary]{v59:cor:primitive}
If $h(0)=0$ and $\|h'\|_\infty\le L\|e_1-e_2\|_\infty$, then the response
to $ah$ satisfies
\begin{equation}
 \|\operatorname{Response}_{J_a}(ah)\|_\infty
 \le 2CLT\|e_1-e_2\|_\infty .
 \label{v59:eq:primitive-gain}
\end{equation}
A uniformly Lipschitz multiplier of $h$ changes only the constant.
\end{corollary}
\begin{proof}
Both $\|h\|_\infty$ and $\|h'\|_{L^1}$ are bounded by
$LT\|e_1-e_2\|_\infty$.  Apply \eqref{v59:eq:primitive-A}.
For a multiplier $B$, use $(Bh)'=B'h+Bh'$.
\end{proof}

\section{Preparation tests can be anchored at a reference path}
\label{v59:sec:endpoints}

Fix a reference coefficient path $J_{b,a}$ in the same class, and put
$P_{b,+}=P_+[J_{b,a}]$, $P_{b,-}+P_{b,c}=I-P_{b,+}$.
Instead of changing the preparation tests with each unknown path, impose
\begin{equation}
 (I-P_{b,+}(0))D_ae(0)=0,\qquad
 P_{b,+}(T)D_ae(T)=0 .
 \label{v59:eq:anchor-bc}
\end{equation}
These are fixed linear endpoint covectors for the \emph{correction}.
The actual principal equation remains $J_a(\theta)$.

\begin{lemma}[Fixed reference endpoints preserve the source estimates]
\label[lemma]{v59:lem:anchor}
Suppose $J_a(0)=J_{b,a}(0)$ and
\[
 \|K_a-K_{b,a}\|_\infty
   +a\|B_a-B_{b,a}\|_\infty\le aR_1 .
\]
Then \eqref{v59:eq:anchor-bc} defines a unique prepared solution for small $a$,
and all three estimates of \cref{v59:thm:source} remain valid, with constants
also depending on $R_1$.  The instantaneous terminal projectors obey
\begin{equation}
 \|P_+[J_a](T)-P_{b,+}(T)\|
 \le C\bigl(\|K_a-K_{b,a}\|_\infty
           +a\|B_a-B_{b,a}\|_\infty\bigr).
 \label{v59:eq:projector-Lip}
\end{equation}
\end{lemma}
\begin{proof}
The two rescaled contour inversions used in Version 58 give bounded
first derivatives of the total unstable projector with respect to
$K_a$ and $aB_a$.  They do not resolve spacings inside either unstable
group.  This proves \eqref{v59:eq:projector-Lip}.

In balanced coordinates a homogeneous solution with zero initial stable
and center data and prescribed instantaneous terminal unstable data $h$
has norm at most $C|h|$, by the hyperbolic/center integral equations of
Version 58.  Let $\mathcal Hh$ be that solution.  At the terminal endpoint
$P_+(T)\mathcal Hh(T)=h$.  After identifying the close unstable ranges,
the map $h\mapsto P_{b,+}(T)\mathcal Hh(T)$ is the identity plus an operator
of norm at most $C\|P_+(T)-P_{b,+}(T)\|$.  It is invertible for small $a$.

If $u_0$ is the instantaneous-prepared inhomogeneous solution, the required
terminal correction has size at most
$C\|P_+(T)-P_{b,+}(T)\|\|u_0\|_\infty$.
For a continuous weighted source,
$\|u_0\|_\infty\le Ca(T+a)\|r\|_\infty$.
For a regular weighted source it is at most $Ca\|r\|_{\mathsf A_T}$,
and for an $ah$ source at most $Ca\|h\|_{\mathsf A_T}$.
The projector error is $O(a)$; the additional homogeneous correction
therefore stays within each claimed physical bound after multiplication
by $D_a^{-1}$.  The initial subspaces are exactly equal.
The same argument proves uniqueness.
\end{proof}

Fixing these tests is not freezing the principal field, nor is it a
measurement or projection applied to a history.  It is a reference-dependent
mixed boundary prescription.  It depends on the chosen horizon.  It is
different from imposing each iterate's own instantaneous endpoint spaces.

\section{Original kinematics gives a uniformly regular integrated-state map}
\label{v59:sec:kinematics}

Use the normalized original state
\[
 Y=(X,\lambda),\qquad
 X=aY_*+a^2\xi,\qquad \lambda=e+a\ell,\qquad y=(\xi,\ell).
\]
The mean-lapse normalization is retained in the multiplier coordinate space.
On a compact subchart from Versions 55--57, write the same triangular
mass frame as $U_a(y)$.  For an \emph{input} continuous rate $p$,
the first original equations become
\begin{equation}
 y'=\mathfrak f_a(y,p):=
 \left(
 \frac{F(aY_*+a^2\xi,e+a\ell)}a,\quad S_aU_a(y)p
 \right),\qquad S_a=(aI_A,I_c,aI_g).
 \label{v59:eq:kin}
\end{equation}

\begin{proposition}[Kinematic smoothing and a Volterra comparison]
\label[proposition]{v59:prop:kin}
The right-hand side of \eqref{v59:eq:kin} extends smoothly through $a=0$
on the stated compact subchart, with uniform fixed-order state and rate
derivative bounds.  For rates in a fixed $C^0$ ball and a common initial
state, it has solutions on a fixed sufficiently short interval and
\begin{align}
 \|y[p_1]-y[p_2]\|_\infty
   &\le LTe^{LT}\|p_1-p_2\|_\infty,\notag\\
 \|(y[p_1]-y[p_2])'\|_\infty
   &\le C_{T_0}\|p_1-p_2\|_\infty .
 \label{v59:eq:kin-Lip}
\end{align}
No rate derivative is used.  For any uniformly $C^2$ coefficient $A_a(y)$,
the difference $A_a(y[p_1])-A_a(y[p_2])$ vanishes initially and has
$\mathsf A_T$ norm at most $C_{T_0}T\|p_1-p_2\|_\infty$.
\end{proposition}
\begin{proof}
The inherited stationary Hamiltonian gives $F(0,e)=0$ and analytic
dependence on its normalized chart.  Hadamard's integral identity divides
$F(aY_*+a^2\xi,e+a\ell)$ by $a$ with no pole.  The established frame
orders and bounded $S_a$ give all other bounds.  Picard existence on a
compact interior tube and Gronwall applied to the state difference give
\eqref{v59:eq:kin-Lip}.  For the coefficient derivative subtract
\[
 DA_a(y_1)\mathfrak f_a(y_1,p_1)
     -DA_a(y_2)\mathfrak f_a(y_2,p_2).
\]
The $C^2$ bound and \eqref{v59:eq:kin-Lip} bound this by
$C\|p_1-p_2\|_\infty$.  Integration and the common initial value prove
the last assertion.
\end{proof}

This result concerns the \emph{kinematic input map}.  An arbitrary input
rate need not preserve $P=0$ or its consistency equation.  It is not called
an original history until those extra equations and their initial data
are satisfied.  A rate-dependent initial selection would add its actual
initial-state difference to \eqref{v59:eq:kin-Lip}; it cannot be omitted.

For a uniformly smooth divided coefficient with a pointwise expansion
\begin{equation}
 J_a(y)=J+aK(y)+a^2B_a(y),\qquad K^T=-K,\quad B_a^T=-B_a,
 \label{v59:eq:pointwise-J}
\end{equation}
the proposition immediately gives the required $W^{1,\infty}$ path bounds
from bounded continuous rates.  No $p'$ or second coefficient derivative
in time is needed.  A proof of \eqref{v59:eq:pointwise-J} on the relevant
trial tube is a separate coefficient identity; smoothness at each $a>0$
does not imply its uniform divided extension.

\section{A coupled state-dependent principal and metric theorem}
\label{v59:sec:coupled}

We state the class precisely rather than assign unverified divided jets
to the original $J_U,f_U$.  Let $\mathfrak f_a(y,p)$ be any uniformly $C^2$
state equation on a fixed finite-dimensional tube, including the original
kinematic map of \cref{v59:prop:kin}.  Let \eqref{v59:eq:pointwise-J} hold
with uniform $C^2$ bounds.  Let $\mathcal F_a(y,p)$ be uniformly $C^2$ and
bounded on rate balls; the complete metric polynomial of Version 57 is
an example at the stated regular states.

Choose a uniformly $W^{1,\infty}$ reference $p_b$, let
$y_b'= \mathfrak f_a(y_b,p_b)$ with $y_b(0)=y_0$, and define the declared
reference covector
\[
 g_b=M_ap_b'+J_a(y_b)p_b .
\]
This definition does not identify $g_b$ with the original state-only source.

\begin{theorem}[A closed state/principal/metric preparation in continuous rates]
\label[theorem]{v59:thm:coupled}
There are $R_0<\infty$, $T_*>0$, and $a_*>0$, depending on the displayed
coefficient and reference bounds, such that for $0<T\le T_*$ and
$0<a<a_*$ the system
\begin{equation}
 \boxed{
 y'=\mathfrak f_a(y,p),\quad y(0)=y_0,\qquad
 M_ap'+J_a(y)p=g_b+\mathcal D_a\mathcal F_a(y,p)
 }
 \label{v59:eq:coupled}
\end{equation}
has a unique solution in the tube
$\|p-p_b\|_\infty\le aR_0$, with
\eqref{v59:eq:anchor-bc} imposed on $e=p-p_b$.
It satisfies
\begin{equation}
 \|p-p_b\|_\infty\le Ca,\qquad
 \|y-y_b\|_\infty\le CaT.
 \label{v59:eq:coupled-error}
\end{equation}
No derivative bound on the unknown nonlinear rate is an input.
The correction $p-p_b$ and state $y$ are $C^1$ for each fixed $a>0$;
$p$ is $W^{1,\infty}$ and the rate equation holds almost everywhere.
If $p_b$ is $C^1$, then so is $p$.
\end{theorem}
\begin{proof}
For an input $e$ in the complete $C^0$ ball of radius $aR_0$, solve
the state equation with rate $p_b+e$ and call its state $y_e$.
Put $J_e=J_a(y_e)$, $J_b=J_a(y_b)$ and
\[
 H_e=a^{-1}(J_e-J_b)p_b,\qquad
 R_e=\mathcal F_a(y_e,p_b+e).
\]
The map sends $e$ to the exact solution $u$ of
\begin{equation}
 M_au'+J_eu=\mathcal D_aR_e-aH_e
 \label{v59:eq:iteration}
\end{equation}
with the fixed reference endpoint tests.  The kinematic estimates prove
\[
 \|K(y_e)-K(y_b)\|_\infty+
 a\|B_a(y_e)-B_a(y_b)\|_\infty\le CaT R_0 .
\]
Thus \cref{v59:lem:anchor} applies uniformly to all iterates.
They also give $H_e(0)=0$ and
\begin{equation}
 \|H_{e_1}-H_{e_2}\|_{\mathsf A_T}
 \le CT\|e_1-e_2\|_\infty .
 \label{v59:eq:H-Lip}
\end{equation}
This uses the time derivative of $p_b$, not of either unknown input.
The reference amplitude
$R_b=\mathcal F_a(y_b,p_b)$ has uniformly bounded $\mathsf A_T$ norm.
Also
$\|R_e-R_b\|_\infty\le C\|e\|_\infty$.
The two source estimates and \eqref{v59:eq:H-Lip} give
\[
 \|u\|_\infty\le C_0a+C_1(T+a)\|e\|_\infty.
\]
Obtain these constants on the fixed larger rate set
$\|p\|\le\sup_a\|p_b\|_\infty+1$ before choosing $R_0$.
Choose $R_0\ge2C_0$, then impose $aR_0\le1$ and take $T_*,a_*$
small enough to map the ball into itself.

For contraction compare two outputs $u_1,u_2$.  Their endpoint tests are
identical, and subtraction yields
\[
 M_a(u_1-u_2)'+J_{e_1}(u_1-u_2)
 =\mathcal D_a(R_{e_1}-R_{e_2})
  -a(H_{e_1}-H_{e_2})-(J_{e_1}-J_{e_2})u_2.
\]
The first two terms have responses bounded by
$C(T+a)\|e_1-e_2\|_\infty$.  For the last term use the already proved
\emph{output} bound $\|u_2\|_\infty\le aR_0$:
\[
 \left\|\mathcal D_a^{-1}
   (J_{e_1}-J_{e_2})u_2\right\|_\infty
 \le CR_0T\|e_1-e_2\|_\infty .
\]
Here $J_{e_1}-J_{e_2}=a\,O(T\|e_1-e_2\|_\infty)$, so the product is
order $a^2$; it is therefore in the weighted source class without
differentiating $u_2$.  Its response costs
$C(T+a)R_0T\|e_1-e_2\|_\infty$.
The resulting contraction constant is at most
\[
 C_2(T+a)(1+R_0T),
\]
which is less than one after reducing $T_*,a_*$.  Banach's theorem proves
existence and uniqueness in the tube.  At the fixed point,
\eqref{v59:eq:iteration} is exactly \eqref{v59:eq:coupled}.  The invertible
positive-$a$ mass and continuous right-hand side in
\eqref{v59:eq:iteration} make $e$ continuously differentiable, and the
state has the same regularity.  Since $p_b$ was only required to be
$W^{1,\infty}$, the full rate is asserted in that space unless $p_b$ is
$C^1$; its original displayed equation then holds almost everywhere.
The state error is \eqref{v59:eq:kin-Lip}.
\end{proof}

There is no contradiction with \cref{v58:prop:sensitivity}.  In that example
the reference response has fast oscillations, not a uniform
$W^{1,\infty}$ bound.  Here the coefficient difference is a primitive of a
rate difference, the leading reference is regular, and the coefficient
difference multiplies a correction of size $O(a)$.  These three facts,
together with fixed reference endpoint covectors, are all used in the proof.

\begin{corollary}[A state-only term with one explicit amplitude factor]
\label[corollary]{v59:cor:state-source}
Suppose a state-only covector has the pointwise decomposition
\begin{equation}
 f_a(y)=f_*+a h_a(y),\qquad \|h_a\|_{C^2}\le C,
 \label{v59:eq:fjet}
\end{equation}
and a uniformly regular reference satisfies
$M_ap_b'+J_a(y_b)p_b+f_a(y_b)=0$.
Then the theorem holds for
\[
 M_ap'+J_a(y)p+f_a(y)=\mathcal D_a\mathcal F_a(y,p)
\]
with the same kinematic equation and prepared error.  A term
$\mathcal D_a k_a(y)$ with bounded $C^2$ coefficients may be included in
$f_a$ as well.
\end{corollary}
\begin{proof}
In \eqref{v59:eq:iteration} add
$-a[h_a(y_e)-h_a(y_b)]$.  The bracket has zero initial value and
$\mathsf A_T$ Lipschitz bound $CT\|e_1-e_2\|_\infty$ by
\cref{v59:prop:kin}; use \eqref{v59:eq:primitive-A}.
The extra weighted term is handled with $R_e$.
\end{proof}

\section{What must still be checked for the original action}
\label{v59:sec:native}

The full original-equivalent system is
\[
 y'=\mathfrak f_a(y,p),\qquad
 M_ap'+J_U(a,y)p+f_U(a,y)=\mathcal D_a\mathcal F_a(y,p).
\]
The kinematic regularity and complete metric polynomial required by the
coupled theorem are proved original-source inputs.  The following two
applications have \emph{not} been established in this body.

First, the actual divided Poisson map must satisfy
\eqref{v59:eq:pointwise-J}, and the actual state-only covector must satisfy a
usable form such as \eqref{v59:eq:fjet}, on the very trial states used by
the iteration.  These are pointwise, uniformly differentiated coefficient
statements, not conclusions from smoothness for each positive $a$.
If $J_U$ extends $C^4$ to the chosen amplitude chart and
$J_U(0,y)=J$ throughout its trial set, Taylor's integral formula gives
\[
 K(y)=\partial_aJ_U(0,y),\qquad
 B_a(y)=\int_0^1(1-s)\partial_a^2J_U(sa,y)\,ds .
\]
An analogous identity gives $h_a$ from $f_U(0,y)=f_*$.
The fixed leading value must hold on the relevant set, not just at one
initial point.  No value of these unevaluated native jets is supplied here.
When a constrained coordinate chart is used instead, its vector field and
its fixed-initial-state comparison have to be used consistently.

Second, the fixed-state comparison above leaves the unstable initial rate
coordinates to be selected by the terminal preparation.  A fixed original
state already determines its normalized rate through
$\mathsf B+M^{\rm orig}\lambda_t=0$.  Therefore the comparison initial
condition is \emph{not automatically an original consistent initializer}.
The original initializer changes the state when its initial rate is varied.
Its actual state response must be retained in the transversality or coupled
selection calculation.  The common-initial-state zero in
\eqref{v59:eq:H-Lip} cannot simply be asserted for that different map.

Subject to those additional identifications, the exact equivalence of
Version 55 recovers $P=0$ and the original consistency equation from their
initial values.  Without them the new theorem is an exact, nonempty
state-coupled comparison class, not a claimed original $[0,aT]$ lift.
Its time is still $\theta=t/a$; no later-time implication is inferred.
The pointwise coefficient tests and the original initial-state dependence
are the remaining task, rather than a generic coefficient-path Lipschitz
estimate invalidated by resonance.

\section{Verification, preservation, and scope}
\label{v59:sec:status}

\paragraph{Proved.}
The finite--fast connection is $O(a)$ under one bounded principal derivative.
A regular weighted source has an order-$a$ prepared response with this
weaker coefficient regularity.  An order-$a$ primitive covector has a
uniform one-derivative response and a short-time Volterra gain.
Fixed reference endpoint tests retain these bounds in an order-$a$
coefficient tube.  The original kinematics maps bounded continuous input
rates into uniformly Lipschitz state and coefficient paths.
Those estimates close an explicit coupled state/principal/metric
comparison in a complete $C^0$ rate ball, including a state-only source
when its displayed pointwise amplitude factor is independently verified.

\paragraph{Inherited.}
All native numerical tensors, root certificates, signatures, two-scale
spectral groups, and the positive nonhyperbolic metric are retained from the
prior versions.  The general mixed Green construction is not reproved as new.
The new proof avoids differentiating the frame connection and pays the
coefficient sensitivity through the integrated state and a reference-sized
tube.  No new native numerical coefficient or native trajectory is claimed.

\paragraph{Not proved.}
The full native divided coefficient identities on the required trial set,
the original initializer's simultaneous endpoint transversality, an original
$[0,aT]$ existence theorem, a fixed positive physical lifespan, cofinal
spatial control, and microscopic selection remain open.  The previously
constructed shorter histories and the independent Version-32 stress
question are unchanged.

\paragraph{Executed checks.}
The companion script verifies the small connection identities in explicit
moving-projector models, the weak-source projection cancellation, exact
one-derivative oscillator responses, Volterra difference estimates,
source-weight arithmetic in the contraction, and the fixed-endpoint
perturbation algebra.  Separately labeled numerical checks solve an
endogenous-coefficient comparison, not the native action.  The V58
verifier is replayed at its stated scope, without regeneration of the older
root, stationary-source, or weak-Poisson banks.  These diagnostics do not
replace the written functional estimates or constitute a proof-assistant
formalization.  Every V58 byte preceding its final document command is
preserved; new labels have prefix \texttt{v59:}.

\addtocontents{toc}{\protect\endgroup}
% END IMMUTABLE VERSION 59 BODY
% APPEND VERSION 60 HERE, BEFORE THE FINAL END-DOCUMENT COMMAND.

% BEGIN IMMUTABLE VERSION 60 BODY
\clearpage
\begin{center}
{\Large\bfseries Version 60: Initial consistency in the prepared state--rate problem}\\[.5em]
{\large Delayed weak-rate susceptibility, physical unstable directions,\\
and a coupled selection with the original rate-dependent initializer}
\end{center}
\makeatletter
\addtocontents{toc}{\protect\begingroup}
\addtocontents{toc}{\protect\makeatletter}
\addtocontents{toc}{\protect\def\protect\l@section{\protect\bfseries\protect\@dottedtocline{1}{0em}{2.5em}}}
\makeatother
\addcontentsline{toc}{part}{Version 60: Original initial consistency in the prepared state--rate comparison}

\section{The initial state cannot be held fixed while its consistent rate changes}
\label{v60:sec:scope}

Version 59 closes a continuous-rate comparison with a common initial state.
Its last section correctly leaves a different problem: an original initial
record determines its rate through the consistency equation. Changing that
rate changes the record. Repeating the common-initial-state contraction
would not establish original initial consistency. This body estimates the
\emph{actual parametric initializer} and then includes it inside the
prepared fixed-point map.

There are two complementary estimates. A nonconstant weak initial rate
first changes the physical initialized fields at order $a^3$, whereas an
active or harmonic rate can change them at order $a^2$. The ten rate
directions left free by the initial endpoint tests have active/harmonic
components smaller than their weak components by a factor $a$. Their
combined normalized initial-state response is therefore $O(a)$, not an
unspecified order-one map. This is the additional small factor required
when the two state paths no longer start at the same point.

All statements concern the same V39 mixed-mode root and $h=1/3$.
The exact initializer, its mean and Poisson inverses, the leading signed
Hessian, the two-scale spectral splitting, and the Version-59 prepared
source estimates are inherited. We use the original order statements in
\cref{v39:lem:analytic-prep,v52:thm:initial,v55:cor:access}; their source
coefficients are not numerically regenerated. In particular the distinction
between active and nonconstant weak rates is the original one, not a new
choice of a favorable norm.

The resulting coupled theorem still has Version 59's explicit hypotheses
on the divided principal and state-only coefficients and its regular
reference. Those hypotheses are not established for the complete original
source here. What is removed is the separate common-initial-state
substitution: the selected comparison starts on exactly constrained,
exactly consistent original records with its \emph{own} initial rate.
Geometric boundary transversality has a substantial literature
\cite{v60:TinKopellJones}; neither an external exchange lemma nor an
unverified native transversality assumption is used in place of the
amplitude calculation below.

\section{The original initializer delays nonconstant weak-rate dependence}
\label{v60:sec:initializer}

Write a physical mixed-rate coordinate as $r=(x,w)$, where
$x=(v,c)\in\R^{81}$ and $w\in\R^{26}$, with original rate
$S_ar=aI_Av+I_cc+aI_gw$. Fix a finite reference value of $r$ and a
sufficiently small, relatively compact rate neighborhood. All bounds below
are local and uniform on that neighborhood, not uniform over unbounded
initial rates.

Let $\mathcal I_a(r)=(X_a(r),\lambda_a(r))$ be the analytic initial
selection used in Version 52, with all other free preparation coordinates
held fixed. It obeys
\begin{equation}
 P(\mathcal I_a(r))=0,\qquad
 \mathsf B(\mathcal I_a(r))+M^{\rm orig}(\mathcal I_a(r))S_ar=0.
 \label{v60:eq:initial-exact}
\end{equation}
Here and below $M^{\rm orig}$ is the original symmetric multiplier
Hessian, not the first Poisson matrix.

\begin{proposition}[A graded derivative of the actual initial selection]
\label[proposition]{v60:prop:initializer}
The inherited analytic selection can be taken on the fixed local rate
neighborhood so that
\begin{align}
 \|D_xX_a\|+\|D_x\lambda_a\|&\le Ca^2,\notag\\
 \|D_wX_a\|+\|D_w\lambda_a\|&\le Ca^3.
 \label{v60:eq:physical-response}
\end{align}
The same orders hold after a fixed finite number of additional derivatives
in the rate parameters on a smaller neighborhood.
For the normalized initial state
\[
 \iota_a(r)=\left(a^{-2}(X_a(r)-aY_*),\ a^{-1}(\lambda_a(r)-e)\right),
\]
one consequently has
\begin{equation}
 \|D_x\iota_a\|\le C,\qquad \|D_w\iota_a\|\le Ca.
 \label{v60:eq:normalized-response}
\end{equation}
\end{proposition}
\begin{proof}
Use the three root variables $z$ and the 104 normalized second multiplier
values $m$ from \cref{v39:lem:analytic-prep}. Before consistency is imposed,
the exact constraint initializer has the form
\begin{align*}
 X(a,z,m)&=aY(z)+a^2X_2(z)+a^3X_3(a,z,m),\\
 \lambda(a,z,m)&=e+a\ell(z)+a^2m.
\end{align*}
The independence of $X_2$ from $m$ is the original re-preparation order in
that lemma: changing $a^2m$ changes the canonical constraint initializer
only from its third coefficient onward. No higher coefficient is set to
zero.

Divide the three constant-shift consistency rows and the 104
nonconstant rows by $a^3$. Denote their analytic extensions by
$\mathcal A_c(a,z,m;x,w)$ and $\mathcal A_\perp(a,z,m;x,w)$.
The original block orders give
\begin{align*}
 \mathcal A_c(0,z,m;x,w)&=\mathbf C(z),\\
 \mathcal A_\perp(0,z,m;x,w)
 &=B_3^\perp(z)+\mathcal K(z)m+\mathcal T_x(z)x,
\end{align*}
where $\mathcal K=\mathbb K_1|_{\mathcal E^\perp}$ is invertible and
$\mathcal T_xx$ includes the original $Av+L_cc$ term. In addition,
\begin{equation}
 D_w\mathcal A_c=O(a^2),\qquad
 D_w\mathcal A_\perp=a\mathcal T_g(z)+O(a^2),\qquad
 D_m\mathcal A_c=O(a).
 \label{v60:eq:divided-orders}
\end{equation}
Indeed $M_{Ag}=O(a^3)$ and $M_{Wg}=O(a^4)$, and the physical weak
rate is $aw$: its contributions before division are $O(a^4)$ on
active rows and $O(a^5)$ on all weak rows, including constants. This
also identifies $\mathcal T_g$ as the nonconstant restriction of
$\lim_{a\to0}a^{-3}M^{\rm orig}I_g$. It is not a new chosen control.

The triangular implicit-function derivative from Version 39 is invertible.
At $a=0$ it yields $z(0,r)=z_*$ and a finite $m_0(x)$ independent of
$w$. Differentiate the divided equations with respect to $w$. The lower
block first gives
$D_wm=O(a)+O(D_wz)$; the upper block, using
\eqref{v60:eq:divided-orders}, gives
$D_wz=O(a^2)+O(aD_wm)$. Hence
\[
 D_wz=O(a^2),\qquad D_wm=O(a).
\]
For $x$, the same two block equations instead give
$D_xz=O(a)$ and $D_xm=O(1)$. Equivalently,
$z=z_*+az_1(x)+a^2z_2(a,x,w)$ and
$m=m_0(x)+am_1(a,x,w)$, with analytic bounded coefficients.
Substitution into the displayed constraint initializer proves
\eqref{v60:eq:physical-response}. Analyticity on a compactly contained
parameter set permits the additional fixed-order derivatives. Division
by the two state weights gives \eqref{v60:eq:normalized-response}.
\end{proof}

For example the leading weak susceptibility of the second multiplier
control is explicitly
\begin{equation}
 D_wm_1(0,x,w)=-\mathcal K(z_*)^{-1}\mathcal T_g(z_*).
 \label{v60:eq:m1}
\end{equation}
This formula does not require computing the unknown absolute higher
preparation baseline. No numerical value for this matrix is claimed.

The physical mixed rate is not yet the mass-normalized coordinate of
Version 55. Let $\mathfrak I_a(p)$ denote the normalized initial state
obtained by solving
\begin{equation}
 r=U_a(\iota_a(r))p,\qquad \mathfrak I_a(p)=\iota_a(r).
 \label{v60:eq:frame-initial}
\end{equation}
The contraction proving \cref{v55:cor:access} applies parametrically.

\begin{corollary}[The delayed response survives the triangular mass frame]
\label[corollary]{v60:cor:frame-initial}
The exact normalized-rate initializer satisfies
\begin{equation}
 \|D_{p_x}\mathfrak I_a\|\le C,\qquad
 \|D_{p_w}\mathfrak I_a\|\le Ca,
 \label{v60:eq:I-response}
\end{equation}
with the same fixed-order derivative convention. Its physical records
satisfy \eqref{v60:eq:initial-exact} with rate $S_aU_ap$.
\end{corollary}
\begin{proof}
The frame has diagonal blocks $I+O(a)$, identically zero lower-left block,
and upper-right block $O(a^2)$, with bounded divided state derivatives.
By \eqref{v60:eq:normalized-response}, differentiating the frame through
$\iota_a$ in a weak input direction costs an additional $a$.
Implicit differentiation of \eqref{v60:eq:frame-initial} therefore gives
$D_{p_w}r_x=O(a^2)$ and $D_{p_w}r_w=I+O(a)$, whereas
$D_{p_x}r$ stays bounded. The derivative matrix used in this implicit
solve is $I+O(a)$, uniformly invertible. Compose with
\eqref{v60:eq:normalized-response}. The exact initial identities follow
from the definition, not from the derivative estimates.
\end{proof}

\section{The free growing-rate directions have small upper components}
\label{v60:sec:unstable}

Retain Version 59's principal class and a fixed reference path. Let
$P_{b,+}(0)$ be its \emph{balanced} projector onto the ten growing
directions, and let $D_a=\operatorname{diag}(I_{81},aI_{26})$.
The physical rate subspace left free by the initial tests is
\begin{equation}
 \mathcal E^+_{b,a}=D_a^{-1}\operatorname{Ran}P_{b,+}(0).
 \label{v60:eq:physical-plus}
\end{equation}
It is this ten-dimensional subspace, not the whole 28-dimensional fast
space, which is used below.

\begin{lemma}[Upper-to-weak ratio on the physical growing subspace]
\label[lemma]{v60:lem:plus}
Uniformly in the bounded principal class and for small $a>0$,
\begin{equation}
 d=(d_x,d_w)\in\mathcal E^+_{b,a}
 \quad\Longrightarrow\quad |d_x|\le Ca|d_w|.
 \label{v60:eq:plus-ratio}
\end{equation}
In particular projection of this subspace onto its weak coordinates is
injective. One may choose a uniformly conditioned ten-column synthesis
$\bigl(aX_+(a),W_+(a)\bigr)^T$ in physical coordinates.
\end{lemma}
\begin{proof}
Use the two separated spectral reductions of Versions 53 and 58. In
balanced coordinates $(u_x,u_w)=(d_x,ad_w)$ the nine fastest growing
directions have a nonsingular limiting weak projection. A nonzero
fastest eigenvector with $u_w=0$ would be annihilated by the leading
$a^{-2}$ generator and cannot belong to its nonzero spectral group.
More precisely its weak images form the nine-dimensional growing
subspace of $-\Delta_0^{-1}K_g$. Since the kernel form on
$Z=\ker K_g$ is nondegenerate, the sum of its nonzero spectral groups
is $\Delta_0$-orthogonal to $Z$ and intersects $Z$ trivially.

The one intermediate growing direction has a limiting nonzero weak
component in $Z$, determined by the simple positive-real root of the
retained two-dimensional quadratic pencil. Its upper balanced component
is a bounded linear function of that weak component, obtained from the
invertible upper Schur equation. Thus the combined ten-dimensional
limiting growing space has injective weak projection: the fastest weak
images and this kernel direction are independent. The two contour
reductions give convergence of the combined subspace, including at
multiplicities inside the fastest group; individual eigenvectors there
are not needed.

On the compact closure of a bounded coefficient set, the intermediate
root remains simple and bounded away from zero, and the preceding
injectivity persists with a uniform positive minimum singular value.
Consequently $|u_x|\le C|u_w|$ on the actual growing subspace for
small $a$. Returning to physical coordinates proves
\eqref{v60:eq:plus-ratio}. A local orthonormal weak-image basis and
its graph supply the stated synthesis. No analogous graph over 26 weak
coordinates is claimed for all 28 fast directions.
\end{proof}

\begin{theorem}[Small original initial-state response on the free rate fibre]
\label[theorem]{v60:thm:initial-fibre}
Fix a bounded reference initial rate $p_b^0$. On a small fixed ball in
its affine fibre $p_b^0+\mathcal E^+_{b,a}$,
\begin{equation}
 \boxed{
 \|\mathfrak I_a(p_b^0+d_1)-\mathfrak I_a(p_b^0+d_2)\|
 \le Ca\|d_1-d_2\|.}
 \label{v60:eq:fibre-Lip}
\end{equation}
The corresponding physical records differ by at most
$Ca^3\|d_1-d_2\|$ in both $X$ and $\lambda$. On a rate-correction
ball $\|d\|\le aR$, their physical initial changes from the reference
are therefore $O(a^4)$.
\end{theorem}
\begin{proof}
Integrate \eqref{v60:eq:I-response} along the line segment in the
fixed affine fibre, and use \eqref{v60:eq:plus-ratio}. For the physical
records use \eqref{v60:eq:physical-response}; the triangular
initial frame preserves the same estimates as in
\cref{v60:cor:frame-initial}. There is no assertion that the normalized
initial map is $O(a)$-Lipschitz on arbitrary active/harmonic variations.
\end{proof}

\section{The state estimate acquires an amplitude term, not an order-one gap}
\label{v60:sec:state}

For an input correction $e\in C^0([0,T])$ with
$e(0)\in\mathcal E^+_{b,a}$, solve the original kinematic input equation
\begin{equation}
 y_e'=\mathfrak f_a(y_e,p_b+e),\qquad
 y_e(0)=\mathfrak I_a(p_b(0)+e(0)).
 \label{v60:eq:moving-initial-state}
\end{equation}
Write $y_b=y_0$. This is still an input-state map: until its rate equation
is the original one, it is not called an original history.

\begin{proposition}[Integrated differences with an exactly consistent moving initial cut]
\label[proposition]{v60:prop:state}
On a fixed compact kinematic tube, for bounded reference and input rates,
\begin{align}
 \|y_{e_1}-y_{e_2}\|_\infty
 &\le C(a+T)\|e_1-e_2\|_\infty,\notag\\
 \|(y_{e_1}-y_{e_2})'\|_\infty
 &\le C\|e_1-e_2\|_\infty.
 \label{v60:eq:state-Lip}
\end{align}
For any uniformly $C^2$ coefficient $B_a(y)$,
\begin{equation}
 \|B_a(y_{e_1})-B_a(y_{e_2})\|_{\mathsf A_T}
 \le C(a+T)\|e_1-e_2\|_\infty,
 \qquad \|h\|_{\mathsf A_T}=\|h\|_\infty+\|h'\|_{L^1}.
 \label{v60:eq:coefficient-Lip}
\end{equation}
The coefficient difference need not vanish at time zero; its actual
initial value is $O(a\|e_1-e_2\|_\infty)$.
\end{proposition}
\begin{proof}
The initial difference is \eqref{v60:eq:fibre-Lip}. The uniformly
Lipschitz original kinematics in \cref{v59:prop:kin} and Gronwall give
\eqref{v60:eq:state-Lip}. Subtract the two chain rules
$DB_a(y_e)\mathfrak f_a(y_e,p_b+e)$. Their difference is bounded by
$C\|e_1-e_2\|_\infty$ without differentiating either input rate.
Integrating this derivative bound and retaining the nonzero initial
coefficient difference proves \eqref{v60:eq:coefficient-Lip}.
\end{proof}

The initial principal paths now also differ. The following minor
extension of the Version-59 endpoint lemma is required.

\begin{lemma}[Fixed tests allow small projector changes at both ends]
\label[lemma]{v60:lem:endpoints}
Let $J_a,J_{b,a}$ be in the Version-59 class and suppose the balanced
growing projectors differ by at most $C_0a$ at each endpoint. Impose the
reference tests
\begin{equation}
 (I-P_{b,+}(0))D_au(0)=0,\qquad P_{b,+}(T)D_au(T)=0.
 \label{v60:eq:anchor-bc}
\end{equation}
All three source estimates of \cref{v59:thm:source} remain valid. Exact
equality of the initial projectors is not required.
\end{lemma}
\begin{proof}
The instantaneous mixed problem has a homogeneous endpoint-solution map
bounded uniformly in balanced coordinates. Express both the initial
non-growing data and terminal growing data of this map in the nearby
reference test spaces. The resulting 107-dimensional boundary map is
its invertible instantaneous counterpart plus an operator of norm
$O(a)$; its inverse follows from a Neumann series.

For an instantaneous prepared inhomogeneous solution $u_0$ in balanced
coordinates, the required endpoint correction has size at most
$Ca\|u_0\|_\infty$. The inherited estimates give respectively
$\|u_0\|_\infty\le Ca(T+a)\|r\|_\infty$ for a continuous weighted
source, $Ca\|r\|_{\mathsf A_T}$ for a regular weighted source, and
$Ca\|h\|_{\mathsf A_T}$ for an $ah$ source. Undoing balance costs
at most $a^{-1}$ on the added homogeneous correction. It therefore
preserves each stated physical estimate. The same boundary-map inversion
proves uniqueness. This argument retains the finite coordinates and all
18 decaying/oscillatory initial and ten growing terminal conditions.
\end{proof}

\section{A coupled prepared theorem with the original initializer inside the map}
\label{v60:sec:coupled}

Assume exactly the coefficient class of \cref{v59:thm:coupled} on the trial
states: uniformly $C^2$ functions
\begin{equation}
 J_a(y)=J+aK(y)+a^2B_a(y),\qquad K^T=-K,\quad B_a^T=-B_a,
 \label{v60:eq:principal}
\end{equation}
and the complete bounded rate-Lipschitz weighted feedback
$\mathcal D_a\mathcal F_a(y,p)$. Use the original kinematic input map.
Choose a uniformly $W^{1,\infty}$ reference $p_b$, take
$y_b(0)=\mathfrak I_a(p_b(0))$, and let
\[
 g_b=M_ap_b'+J_a(y_b)p_b.
\]
As in Version 59, this definition alone does not identify $g_b$ with
the original state-only covector.

\begin{theorem}[Simultaneous prepared rate and exact original initial consistency]
\label[theorem]{v60:thm:coupled}
There are fixed $T_*,a_*,R_*>0$ such that, for $0<T\le T_*$ and
$0<a<a_*$, the system
\begin{equation}
 \boxed{
 \begin{aligned}
 y'&=\mathfrak f_a(y,p),&y(0)&=\mathfrak I_a(p(0)),\\
 M_ap'+J_a(y)p&=g_b+\mathcal D_a\mathcal F_a(y,p)
 \end{aligned}}
 \label{v60:eq:coupled}
\end{equation}
with \eqref{v60:eq:anchor-bc} on $e=p-p_b$ has a unique solution in
$\|e\|_\infty\le aR_*$. It satisfies
\begin{equation}
 \|p-p_b\|_\infty\le Ca,\quad
 \|y-y_b\|_\infty\le Ca(a+T),\quad
 \|y(0)-y_b(0)\|\le Ca^2.
 \label{v60:eq:errors}
\end{equation}
Its original initial record is exactly constrained and consistent with
its own initial normalized rate. The physical initial records differ
from the reference by $O(a^4)$. No rate derivative of an unknown iterate
is an input.
\end{theorem}
\begin{proof}
Work on the complete closed ball
\[
 \mathcal B_a=\{e\in C^0:\ \|e\|_\infty\le aR_*,\quad
 e(0)\in\mathcal E^+_{b,a}\}.
\]
For an input $e$ solve \eqref{v60:eq:moving-initial-state}. Put
$J_e=J_a(y_e)$,
$H_e=a^{-1}(J_e-J_b)p_b$, and
$R_e=\mathcal F_a(y_e,p_b+e)$. Send $e$ to the exact prepared solution
$u$ of
\begin{equation}
 M_au'+J_eu=\mathcal D_aR_e-aH_e
 \label{v60:eq:iteration}
\end{equation}
with the fixed reference tests. Its initial value automatically belongs
to $\mathcal E^+_{b,a}$.

The initial coefficient difference is $O(a\|e(0)\|)$ in the divided
$K$-coordinate, and the full path difference is
$O((a+T)\|e\|_\infty)$. On $\mathcal B_a$ both endpoint projector
differences are $O(a)$, so \cref{v60:lem:endpoints} applies.
By \cref{v60:prop:state} and the uniform reference derivative,
\[
 \|H_{e_1}-H_{e_2}\|_{\mathsf A_T}
 \le C(a+T)\|e_1-e_2\|_\infty.
\]
Unlike Version 59, $H_e(0)$ is not asserted to vanish. Its actual
$O(a\|e(0)\|)$ term is included. The reference amplitude
$R_b=\mathcal F_a(y_b,p_b)$ has bounded $\mathsf A_T$ norm, while
$\|R_e-R_b\|_\infty\le C\|e\|_\infty$. The regular weighted,
continuous weighted, and $ah$ estimates therefore give
\[
 \|u\|_\infty\le C_0a+C_1(a+T)\|e\|_\infty.
\]
Choose coefficient bounds on the fixed larger rate ball before choosing
$R_*\ge2C_0$; then reduce $a_*,T_*$ to obtain self-mapping.

For two outputs, subtract \eqref{v60:eq:iteration}. All boundary
covectors are identical. In addition to the two Lipschitz source
differences, there is $(J_{e_2}-J_{e_1})u_2$. The already established
$\|u_2\|\le aR_*$ gives
\[
 \|\mathcal D_a^{-1}(J_{e_1}-J_{e_2})u_2\|_\infty
 \le CR_*(a+T)\|e_1-e_2\|_\infty.
\]
No derivative of $u_2$ is taken. The response estimates bound the total
contraction constant by
\begin{equation}
 \kappa_a\le C(a+T)\bigl(1+R_*(a+T)\bigr)<1
 \label{v60:eq:contraction}
\end{equation}
after reducing the two parameters. Banach's theorem supplies a fixed
point in the same complete $C^0$ space. At that fixed point the state
uses its own initial rate, and \eqref{v60:eq:iteration} is exactly
\eqref{v60:eq:coupled}. The estimates follow from
\cref{v60:thm:initial-fibre,v60:prop:state}. Exact original initial
consistency follows from \eqref{v60:eq:frame-initial}, not from an
asymptotic residual. Regularity for positive $a$ is the same as in
\cref{v59:thm:coupled}.
\end{proof}

\begin{corollary}[Original recovery after the remaining coefficient identification]
\label[corollary]{v60:cor:original}
The theorem also applies to a state-only covector
$f_a(y)=f_*+ah_a(y)+\mathcal D_ak_a(y)$ with the uniform derivative
bounds and source-compatible regular reference stated in
\cref{v59:cor:state-source}. If these coefficients are independently
identified with the full original $J_U,f_U$ on the trial tube, the resulting
solution is an original constrained history on $0\le t\le aT$.
No additional initial endpoint-transversality hypothesis is then needed.
\end{corollary}
\begin{proof}
Add $-a(h_a(y_e)-h_a(y_b))$ to \eqref{v60:eq:iteration} and include
$k_a$ in the weighted source. Equation \eqref{v60:eq:coefficient-Lip}
gives the same contraction. After the stated original identification,
\cref{v55:thm:original} recovers all constraints and the undifferentiated
consistency equation from the exact initial values already supplied by
the theorem. The corollary is explicitly conditional on that coefficient
and reference identification; it is not an assertion that it has been
performed in this body.
\end{proof}

\section{Why ordinary initial access alone would not have sufficed}
\label{v60:sec:diagnostic}

A simple exact comparison isolates the role of the delayed initial map.
This is not a native Hamiltonian model. Let $a,T>0$ and consider
\[
 e'=a^{-1}(e-y),\qquad y'=0,\qquad e(T)=0,\qquad y(0)=\beta e(0).
\]
For an input initial correction $\eta$, the output initial correction is
\begin{equation}
 \eta\longmapsto \beta(1-e^{-T/a})\eta.
 \label{v60:eq:toy}
\end{equation}
With $\beta=1$, its contraction factor approaches one, and with $\beta>1$
it exceeds one for small $a$. Merely proving that each initial rate has a
corresponding initial state supplies no uniform prepared contraction.
With the actual order represented by $\beta=a$, the factor is at most
$a$. The comparison illustrates the role of
\eqref{v60:eq:fibre-Lip}; it is not used to prove that estimate for the
original initializer.

The new result fixes neither the leading state-only source nor the
principal amplitude jets. In particular the identities
$J_U(0,y)=J$ and a suitable leading value of $f_U(0,y)$ must hold on
the states used by the iteration, not just at the initial root. A
constrained trial chart also needs its actual induced kinematics. The
reference used for a putative original lift must satisfy the corresponding
source condition, rather than being assigned a convenient $g_b$.
The endpoint selection is horizon-dependent; no causal source rule or
microscopic attraction is implied. The ten growing directions remain.

\clearpage
\section{Verification, preservation, and precise remaining scope}
\label{v60:sec:status}

\paragraph{Proved.}
The original initializer has delayed weak-rate dependence, and the mass
frame preserves it. The physical growing rate fibre has an additional
amplitude factor in its upper components. Their composition gives a small
original initial-state response. Incorporating that response, both moving
endpoint projectors, and the integrated kinematics closes the prepared
comparison with exact original initial consistency, under the stated
coefficient/reference class. The fixed initial state used in Version 59
has not been silently reused.

\paragraph{Inherited.}
The mixed-mode root, original initial implicit-function inverses,
stationary Hessian grades, spectral splitting, and one-derivative Green
bounds are inherited at their literal scopes. The derivative statements
here are analytic consequences of those original equations. No new
numerical native coefficient or native trajectory is claimed.

\paragraph{Not proved.}
The divided original $J_U,f_U$ identities on the necessary trial tube and
a source-compatible uniformly regular reference remain unverified. Thus
an original $[0,aT]$ lift is not yet asserted, despite removal of the
separate initializer/endpoint selection problem. No fixed positive
physical lifespan, spatial-refinement uniformity, microscopic preparation
mechanism, or complete Einstein--Standard-Model limit follows. The older
shorter histories and the independent Version-32 stress problem are
unchanged.

\paragraph{Executed checks.}
The accompanying script checks the graded implicit-function orders,
triangular initial frame, physical unstable projection in declared finite
models, two-endpoint boundary perturbation algebra, contraction weights,
and a genuinely rate-dependent initial-state comparison. The latter
numerical diagnostics are not original spacetime evolutions. Parent
replay is stated explicitly in the provenance; historical native tensor
banks are not described as regenerated. Every parent byte before its
final document command is retained. These checks support the written
arguments but are not proof-assistant formalization.

\begingroup
\renewcommand{\refname}{Source used in Version 60}
\begin{thebibliography}{9}
\small\raggedright
\bibitem{v60:TinKopellJones}
S.-K. Tin, N. Kopell, and C. K. R. T. Jones,
\emph{Invariant Manifolds and Singularly Perturbed Boundary Value Problems},
SIAM Journal on Numerical Analysis \textbf{31} (1994), 1558--1576,
\url{https://doi.org/10.1137/0731081}.
Context for the distinction between initial-manifold access and boundary
transversality. Its exchange lemma is not substituted for the graded
initializer proof above.
\end{thebibliography}
\endgroup
\addtocontents{toc}{\protect\endgroup}
% END IMMUTABLE VERSION 60 BODY
% APPEND VERSION 61 HERE, BEFORE THE FINAL END-DOCUMENT COMMAND.


% BEGIN IMMUTABLE VERSION 61 BODY
\clearpage
\begin{center}
{\Large\bfseries Version 61: The original principal coefficient on reachable trial states}\\[.5em]
{\large An exact harmonic Ward defect, fourth-order constraint tracking,\\
and removal of the principal-coefficient hypothesis}
\end{center}
\makeatletter
\addtocontents{toc}{\protect\begingroup}
\addtocontents{toc}{\protect\makeatletter}
\addtocontents{toc}{\protect\def\protect\l@section{\protect\bfseries\protect\@dottedtocline{1}{0em}{2.5em}}}
\makeatother
\addcontentsline{toc}{part}{Version 61: Original principal coefficients on reachable trial states}

\section{Audit: a pointwise coefficient condition is stronger than a path condition}
\label{v61:sec:audit}

Version 60 includes the original rate-dependent initializer in the prepared
comparison, but retains the pointwise principal hypothesis
$J_U(a,y)=J+aK(y)+a^2B_a(y)$.  Its final section correctly requires that
hypothesis to be checked on the actual trial states.  We perform that check
here rather than add another comparison theorem.  The raw statement fails
on a full ambient normalized-state neighborhood.  A constant-shift Poisson
column detects the order-$a^2$ constraint defect there, and dividing by the
mixed rate weights makes its contribution order one.

The failure does not block the reachable trial paths.  Their initial
constraints are exact, their first canonical field and first nonconstant
multiplier tilt are fixed, and their quadratic prepared drift cancels.
These three facts imply an $O(a^4)$ constraint defect on every bounded
kinematic input path over a fixed short comparison interval, even though
that input is not yet required to solve the rate equation.  Adding a
\emph{tracked} normalized constraint residual to the trial-state variables
then gives the precise pointwise principal class required by Version 60.
No constraint projection, correction to the action, or modification of a
physical history is used.

All statements are at the same V39 mixed-mode root and $h=1/3$.  The
analytic stationary chart, original initial selection, first Poisson
normalization, signed mass frame and prepared Green estimates are inherited.
The specific coefficient inputs are the original linear constraint map,
the quadratic translation generator, the full mixed cubic coefficient
\cref{v48:eq:Rsource}, the harmonic cubic coefficient \cref{v52:eq:S}, and
the quadratic-drift cancellation at $(Y_*,\ell_*)$.  The companion's
\texttt{v36:Poisson-orders} proves the corresponding constrained-curve
expansion; we use its coefficient proof, not its different-root spectral
counts \cite{v61:EB45}.  The off-constraint identity and the reachable-path
tracking below are the new steps.  This is a distinction between an
extension used for analysis and the original constrained evolution, not a
claim of a new mimetic discretization \cite{v61:DBGP}.

\section{The original constant-shift column measures the constraint defect}
\label{v61:sec:ward}

Use the original mean pairing $\langle\cdot,\cdot\rangle_h$ on the
$3^3$ grid.  Let $\mathcal R$ be the 104-dimensional space of multipliers
whose four component means vanish, let $\mathcal H$ be the three constant
shifts, and let $e$ be the unit constant lapse.  Adjoints below use the
retained physical pairings; no orthonormality of a previously printed
synthesis is presumed.  For $c\in\mathcal H$, put
$\delta_c=c^i\delta_i$, where $\delta_i=3(S_i-S_i^{-1})$.

The flat canonical constraint map $L$ is
\begin{equation}
 (LX)_N=-\delta_i\delta_jU_{ij}+\Delta_\delta\tr U,
 \qquad (LX)_{\beta,j}=-2\delta_i p_{ij}.
 \label{v61:eq:L}
\end{equation}
Its generators satisfy
\begin{equation}
 P_{c,1}=0,\qquad P_{c,2}(U,p)=\langle p,\delta_cU\rangle_h,
 \qquad \mathbb J\nabla P_{c,2}=\delta_c(U,p).
 \label{v61:eq:translation}
\end{equation}
In particular, for \emph{arbitrary} symmetric canonical fields $Z$,
\begin{equation}
 \boxed{\mathbb K_1(Z)[c,\nu]
             =\langle\delta_c\nu,LZ\rangle_h.}
 \label{v61:eq:Ward1}
\end{equation}
This identity is an inherited original coefficient identity; the new use is
that $Z$ need not be an admissible second prepared field.  One direct proof
uses $G c=0$, $\mathbb K_1(Z)[c,\nu]=DP_{c,2}(Z)[G\nu]$, skewness of
$\delta_c$, and commutation of $L$ with the phase derivatives.  It does not
use a phase Leibniz rule.

Consider the ambient chart
\begin{equation}
 X=aY_*+a^2\xi,\qquad \lambda=e+a\ell,
 \label{v61:eq:ambient}
\end{equation}
with bounded $\xi,\ell$ and the inherited mean-lapse normalization.  Set
\begin{equation}
 r_2(\xi)=L\xi+P_2(Y_*).
 \label{v61:eq:r2}
\end{equation}
Since $LY_*=0$ and $M^{\rm orig}=O(\|X\|^2)$,
$P=a^2r_2+O(a^3)$, uniformly with any fixed finite number of normalized
state derivatives on a smaller compact chart.

\begin{proposition}[Exact leading off-constraint harmonic defect]
\label[proposition]{v61:prop:defect}
The original Poisson blocks on \eqref{v61:eq:ambient} satisfy
\begin{align}
 K^{\rm orig}_{\mathcal R\mathcal R}
    &=a\mathbb K_1(Y_*)|_{\mathcal R}+O(a^2),\\
 K^{\rm orig}[c,\nu]
    &=a^2\bigl(\mathscr R[c,\nu]
          +\langle\delta_c\nu,r_2(\xi)\rangle_h\bigr)+O(a^3),\\
 K^{\rm orig}[c,d]&=a^3\mathscr S[c,d]+O(a^4).
 \label{v61:eq:off-orders}
\end{align}
Here $\mathscr R$ and $\mathscr S$ are exactly the prepared coefficients
in \cref{v48:eq:Rsource,v52:eq:S}.  The displayed leading coefficients,
other than the explicit $r_2$ term, are independent of $\xi$ and of every
bounded first multiplier tilt $\ell$.
\end{proposition}
\begin{proof}
The mixed bracket through order $a^2$ is
\[
 DP_{c,2}(\xi)[G\nu]+DP_{c,3}(Y_*)[G\nu]
                   -DP_{\nu,2}(Y_*)[\delta_cY_*].
\]
The first term is $\langle\delta_c\nu,L\xi\rangle_h$ by
\eqref{v61:eq:Ward1}. Add and subtract
$\langle\delta_c\nu,P_2(Y_*)\rangle_h$ to obtain the second line
of \eqref{v61:eq:off-orders}, with its sign fixed by the complete
prepared coefficient.  The harmonic stationary first source vanishes
on every canonical field.  Consequently no order-$a$ multiplier tilt
contributes to $D_XP_c$ through order $a^2$, precisely as in the
companion coefficient proof.  The other gradients are
$D_XP_\nu=L_\nu+aDP_{\nu,2}(Y_*)+O(a^2)$.
For two constant shifts the quadratic generators commute on every field.
Their polarized identity removes all second-field contributions at order
$a^3$; the remaining term is $\mathscr S$.  Harmonic first-tilt terms
start one order later.  The nonconstant block follows from the original
first-bracket expansion.  Analytic Taylor division gives the uniform
remainder and derivative statements.  No new numerical value of the
prepared blocks is inferred here.
\end{proof}

\begin{example}[An exact nonzero native witness]
\label[example]{v61:ex:witness}
Let $c_0=(1,-1/2,-1/2)$ and $s_0=(0,1,-1)$ on the first circle.
Choose the second-field increment
$Z=(c_0(x_1)I_3,0)$, the lapse test $\nu=(s_0(x_1),0)$ and the unit
constant shift $c=e_1$. Then
\[
 (LZ)_N=-54c_0,\qquad \delta_1s_0=6c_0,
 \qquad \langle c_0,c_0\rangle_h=\tfrac12.
\]
Therefore
\begin{equation}
 \boxed{\mathbb K_1(Z)[c,\nu]=-162,
        \qquad\mathbb K_1(Z)[\nu,c]=162.}
 \label{v61:eq:witness}
\end{equation}
The second number is also regenerated directly from the original shifted
$P_2$ compiler by quadratic polarization; its spatial-sum value is $4374$.
Thus two ambient trial states with $\xi$ differing by $\epsilon Z$ have
an order-one divided mixed-block difference $162\epsilon+O(a)$.
A uniform identity $J_U(0,y)=J$ cannot hold on a full ambient neighborhood
containing those states.  A frame $U=I+O(a)$ cannot repair this: if
$U^T\widetilde JU=J+O(a)$ held uniformly, its bounded inverse would imply
$\widetilde J=J+O(a)$, contradicting the displayed difference.
\end{example}

The witness is deliberately off the constraint surface. It does not
contradict the prepared-curve expansions or the comparison theorems of
Versions 58--60. It excludes a direct verification of their stronger
pointwise hypothesis on an unrestricted ambient state ball.

\section{An exact principal decomposition, without changing an equation}
\label{v61:sec:decomposition}

For a constraint covector $r$, define an antisymmetric bilinear form on
multipliers by
\begin{equation}
 \mathcal W(r)[u,v]
  =\langle\delta_{c(u)}v,r\rangle_h
       -\langle\delta_{c(v)}u,r\rangle_h,
 \label{v61:eq:W}
\end{equation}
where $c(u)$ is the mean shift of $u$.  It vanishes on two nonconstant
arguments, on two harmonic arguments, and on the constant lapse. It is
linear in $r$, depends only on its nonconstant part, and has rank at most
six.  These statements do not say that it is a Poisson bracket.

Let $S_a=(aI_A,I_c,aI_g)$ in the fixed reordered rate coordinates and set
$\mathfrak W(r)=S_1^T\mathcal W(r)S_1$.  Since the only nonzero blocks
couple a constant shift to a nonconstant multiplier,
\begin{equation}
 a^{-3}S_a^T\mathcal W(P)S_a=a^{-2}\mathfrak W(P).
 \label{v61:eq:weightW}
\end{equation}
The actual divided principal matrix is
$\widetilde J=a^{-3}S_a^TK^{\rm orig}S_a$. Define the auxiliary difference
\begin{equation}
 \widetilde J^{\rm reg}
   =\widetilde J-a^{-2}\mathfrak W(P),\qquad
 J_U^{\rm reg}=U^T\widetilde J^{\rm reg}U.
 \label{v61:eq:regular}
\end{equation}
This is a decomposition of coefficients, not the equation to be solved.
In particular, the $P$ term is not dropped in any trial evolution below.

\begin{theorem}[Uniform regular part of the original principal coefficient]
\label[theorem]{v61:thm:principal-part}
On the bounded ambient chart there are uniformly smooth skew matrices
$K_{\rm reg}(y)$ and $B_{\rm reg}(a,y)$ such that
\begin{equation}
 \boxed{
 J_U=J+aK_{\rm reg}(y)+a^2B_{\rm reg}(a,y)
           +a^{-2}U^T\mathfrak W(P)U.}
 \label{v61:eq:principal-exact}
\end{equation}
This is an exact identity for positive $a$.  The regular part extends
through $a=0$ with the same frozen $J$ used in Versions 52--60. In
particular,
\begin{equation}
 \|J_U-J\|\le C\left(a+a^{-2}\|P_\perp\|\right).
 \label{v61:eq:residual-budget}
\end{equation}
\end{theorem}
\begin{proof}
Subtract \eqref{v61:eq:weightW} in the mixed block of
\eqref{v61:eq:off-orders}. Because $P/a^2=r_2+O(a)$, the explicit defect
cancels, leaving the original prepared $\mathscr R$. The other leading
blocks are unchanged. Thus $\widetilde J^{\rm reg}(0,y)=J$ with no
$\xi$ or $\ell$ dependence. The divisions extend analytically by the
block Taylor expansions. The inherited frame satisfies $U=I+O(a)$
with uniform normalized derivatives, so its congruence preserves that
leading value. Taylor's integral formula gives
\[
 K_{\rm reg}(y)=\partial_a J_U^{\rm reg}(0,y),\qquad
 B_{\rm reg}(a,y)=\int_0^1(1-s)
             \partial_a^2J_U^{\rm reg}(sa,y)\,ds.
\]
Skew symmetry is preserved by every step. The finite phase derivatives
bound $\|\mathfrak W(r)\|$ by $C_h\|r_\perp\|$, and the frame is
bounded, proving \eqref{v61:eq:residual-budget} at the fixed cutoff.
No small-curvature or exact-constraint assumption on the ambient candidate
is used in this estimate.
\end{proof}

\section{Reachable kinematic trial paths have a fourth-order constraint defect}
\label{v61:sec:tracking}

A bounded mixed-rate input does not explore arbitrary first nonconstant
multiplier tilts. Use the finer coordinates
\begin{equation}
 X=aY_*+a^2\xi,\qquad
 \lambda=e+a\ell_*+aI_c\gamma+a^2I_R\eta,
 \qquad I_R=(I_A,I_g).
 \label{v61:eq:refined}
\end{equation}
Here $I_R$ parametrizes the same 104-dimensional nonconstant space;
coefficient extraction uses its fixed coordinate left inverse, denoted
$\operatorname{coord}_R$.  No orthonormality is implicit.  Likewise
$\operatorname{coord}_c$ extracts the three constant-shift coefficients.
The means of $\ell_*$ vanish, and the lapse mean of $\lambda$ is one.
All constants below are on a compact normalized chart and a bounded rate
neighborhood.

Let $\mathsf J_0$ denote the original flat drift map, so
$\mathsf B_1=\mathsf J_0L$.  Its only nonzero output is the lapse row,
the phase divergence of the shift covector.  In particular
$\mathsf J_0^2=0$.  The exact prepared quadratic coefficient is
$\widehat B_2(Y)=\mathsf B_2(Y)-\mathsf J_0P_2(Y)$, and the verified
root obeys
\begin{equation}
 \widehat B_2(Y_*)+\mathbb K_1(Y_*)\ell_*=0,
 \qquad \mathbb K_1(Y_*)I_c=0.
 \label{v61:eq:root-cancel}
\end{equation}

\begin{lemma}[A regular constraint-defect equation from the original jets]
\label[lemma]{v61:lem:defect-ode}
On \eqref{v61:eq:refined}, the functions
\begin{equation}
 E_a=\frac{\mathsf B-\mathsf J_0P}{a^3},
 \qquad N_a=\frac{M^{\rm orig}S_a}{a^3}
 \label{v61:eq:EandN}
\end{equation}
extend smoothly through $a=0$, with uniform fixed-order normalized-state
derivatives.  The kinematic input equation is uniformly regular:
\begin{equation}
 \xi'=F/a,\qquad
 \eta'=\operatorname{coord}_R(S_1Up),\qquad
 \gamma'=\operatorname{coord}_c(S_1Up).
 \label{v61:eq:kin-refined}
\end{equation}
For any continuous bounded input $p$, its actual constraint covector obeys
\begin{equation}
 \boxed{P'=a\mathsf J_0P+a^4R_a(y,p),
        \qquad R_a=E_a+N_aUp.}
 \label{v61:eq:P-ode}
\end{equation}
A prime in this section is $d/d\theta$, with $t=a\theta$.
\end{lemma}
\begin{proof}
On the ambient chart, direct joint Taylor expansion gives
\[
 \mathsf B-\mathsf J_0P
   =a^2\{\widehat B_2(Y_*)+\mathbb K_1(Y_*)\ell\}+O(a^3).
\]
The linear terms cancel identically, including their second-field
contribution.  Substitute
$\ell=\ell_*+I_c\gamma+aI_R\eta$ and use
\eqref{v61:eq:root-cancel}; every coefficient through degree two then
vanishes for all bounded $\xi,\eta,\gamma$. Taylor division proves the
first extension in \eqref{v61:eq:EandN}. The original Hessian grades
$M_{AA}=O(a^2)$, $M_{AW}=O(a^3)$ and $M_{WW}=O(a^4)$ give
$M^{\rm orig}S_a=O(a^3)$, with the same divisible analytic coefficients.
The radial mean-lapse row has no worse order by $M^{\rm orig}\lambda=0$.
This proves the second extension.

The leading field velocity is $a(F_1Y_*+G\ell_*)$ because $GI_c=0$.
All remaining terms in $F/a$ are regular. The other two equations follow
by extracting the actual nonconstant and harmonic coordinates from
$\lambda'=aS_aUp$. Finally $P'=a(\mathsf B+M^{\rm orig}S_aUp)$ by
the chain rule and the original first field equations. Substitution of
\eqref{v61:eq:EandN} yields \eqref{v61:eq:P-ode}. This argument uses
neither the rate equation nor a derivative of the input $p$.
\end{proof}

\begin{theorem}[Automatic fourth-order residual on initialized trial paths]
\label[theorem]{v61:thm:residual}
Initialize \eqref{v61:eq:kin-refined} using the original exact initializer
for $p(0)$.  On a fixed sufficiently short interval on which the bounded
kinematic paths stay in the chart,
\begin{equation}
 \boxed{\|P(\theta)\|\le C_Ta^4\theta,
        \qquad\|P'\|_\infty\le C_Ta^4.}
 \label{v61:eq:P4}
\end{equation}
These statements hold for every bounded continuous input in the specified
rate ball, not only for a solution of a comparison rate equation.
\end{theorem}
\begin{proof}
$P(0)=0$ is exact. Variation of constants in \eqref{v61:eq:P-ode} gives
$P(\theta)=a^4\int_0^\theta e^{a(\theta-v)\mathsf J_0}R_a(v)\,dv$.
On fixed bounded intervals the exponential is bounded, and $R_a$ is
bounded by the preceding lemma. This proves both estimates. The fixed
kinematic interval follows from the uniform vector field and compact
initial-state bounds, independently of the fast rate equation. The bound
is not a proof that these unconstrained trial paths are original histories.
\end{proof}

\section{An augmented trial state satisfies the required principal class exactly}
\label{v61:sec:augmented}

Introduce a \emph{tracked} covector $\pi\in\mathbb R^{108}$ and solve
\begin{equation}
 \boxed{\pi'=a\mathsf J_0\pi+R_a(y,p),\qquad\pi(0)=0}
 \label{v61:eq:pi}
\end{equation}
together with the unmodified kinematic equations. This adds redundant
coordinates to the analysis; it is not a constraint penalty or a new
physical variable.

\begin{proposition}[Exact tracking, not projection]
\label[proposition]{v61:prop:tracking}
Every initialized augmented path satisfies
\begin{equation}
                         \boxed{P(X,\lambda)=a^4\pi}
 \label{v61:eq:track-exact}
\end{equation}
identically.  The augmented state equation has bounded state and rate
partials on its compact normalized chart, including at $a=0$.
\end{proposition}
\begin{proof}
Subtract $a^4$ times \eqref{v61:eq:pi} from \eqref{v61:eq:P-ode}.
The difference $P-a^4\pi$ solves $D'=a\mathsf J_0D$ with $D(0)=0$,
and hence is exactly zero.  All right-hand sides have already been proved
regular. No modification of the canonical velocity is involved.
\end{proof}

Pull the regular coefficients of \cref{v61:thm:principal-part} back by
\eqref{v61:eq:refined}. Their Taylor expansions are again regular because
$\ell=\ell_*+I_c\gamma+aI_R\eta$ has bounded derivatives.  Denote the
resulting first two coefficients by $K_*(y)$ and $B_*(a,y)$.
For $\widehat y=(y,\pi)$ define
\begin{equation}
 \boxed{
 \mathcal J_a(\widehat y)
 =J+aK_*(y)+a^2\{B_*(a,y)+U^T\mathfrak W(\pi)U\}.}
 \label{v61:eq:Ja}
\end{equation}

\begin{theorem}[The original principal hypothesis on reachable trials]
\label[theorem]{v61:thm:principal}
The function \eqref{v61:eq:Ja} has exactly the skew, uniformly
differentiated principal form required in Versions 59--60 on a fixed
augmented state tube. On every initialized augmented kinematic input path,
\begin{equation}
                      \boxed{\mathcal J_a(\widehat y)=J_U(a,y)}
 \label{v61:eq:J-identification}
\end{equation}
exactly.  Hence the principal path has the required uniform
$W^{1,\infty}$ bounds for bounded continuous rates, without assuming any
rate derivative.  Its constraint-defect contribution is $O_T(a^2)$.
\end{theorem}
\begin{proof}
All coefficient functions in \eqref{v61:eq:Ja} are uniformly smooth;
$\mathfrak W$ is linear.  Skewness is preserved by congruence.
Insert \eqref{v61:eq:track-exact} into the exact identity
\eqref{v61:eq:principal-exact}. This gives
\eqref{v61:eq:J-identification}, including the residual term rather than
deleting it. Composition with the regular augmented input equation
provides one bounded time derivative of the divided coefficient functions.
The $a^2$ order follows from bounded $\pi$.
\end{proof}

The refined initial coordinates have the same delayed susceptibility as
Version 60. In fact its physical derivative bounds give
$D_x\eta=O(1)$ and $D_w\eta=O(a)$ after division by $a^2$;
$\gamma(0)=0$ for the fixed initial means, and $\pi(0)=0$.
Together with the existing estimates for $\xi$, this proves the next
consequence without freezing the initial state.

\begin{corollary}[The prepared initial-state gain is retained]
\label[corollary]{v61:cor:Lip}
For two inputs in Version 60's anchored initial growing-rate fibre,
\begin{align}
 \|\widehat y_1-\widehat y_2\|_\infty
    &\le C(a+T)\|p_1-p_2\|_\infty,\\
 \|\widehat y_1'-\widehat y_2'\|_\infty
    &\le C\|p_1-p_2\|_\infty.
 \label{v61:eq:Lip}
\end{align}
The same $C(a+T)$ bound holds for the supremum norm plus the $L^1$ norm
of the time derivative of the difference of any uniformly $C^2$
coefficient composed with these paths.  Thus the original principal map
can be used in the Version-60 iteration without its previously unverified
principal-coefficient assumption.
\end{corollary}
\begin{proof}
The initial augmented-state difference is $O(a\|p_1(0)-p_2(0)\|)$
by the refined initializer orders and the physical growing-fibre estimate.
The regular augmented vector field gives \eqref{v61:eq:Lip} by Gronwall.
For a coefficient $A$, subtract
$DA(\widehat y_1)\widehat y_1'-DA(\widehat y_2)\widehat y_2'$;
bounded second derivatives give the time-derivative estimate, whose
integral has the factor $T$.  No input rate is differentiated.
\end{proof}

\section{What is closed, and what still prevents an original long-time theorem}
\label{v61:sec:status}

\paragraph{Proved.}
The raw principal fixed-leading-value hypothesis is false on a full
ambient normalized state ball, with the exact native witness
\eqref{v61:eq:witness}.  Its leading defect is entirely the harmonic
translation of the second-order constraint residual.  On the actually
initialized mixed-rate input paths the residual is $O(a^4)$ and admits a
regular exact tracking equation.  That equation supplies a uniformly
smooth principal representation with the original frozen $J$, keeps the
complete residual term, and retains the delayed initial-state gain.
The original Poisson and harmonic cubic blocks are not replaced by
matrices chosen for favorable spectra.

\paragraph{Inherited.}
The analytic original coefficient envelopes, mixed-mode root and its
quadratic cancellation, signed mass frame, parametric initializer, and
prepared response bounds retain their earlier scopes.  The constrained
harmonic block formulas were already proved in the companion and in
Versions 48 and 52.  They are not reclassified as new here. The new native
arithmetic consists of a nonzero off-constraint witness and direct
quadratic-compiler checks of its underlying Ward identity.

\paragraph{Not proved.}
The state-only divided source $f_U$ has not been evaluated on this
augmented trial chart, and a source-compatible uniformly regular
reference is not supplied. The leading residual identity for $J_U$ does
not establish a corresponding identity for $f_U$. The complete metric
feedback remains the inherited source-weighted polynomial, but that fact
does not control the omitted state-only covector. Consequently no original
$[0,aT]$ lift, new lifespan, spatial-refinement uniformity, microscopic
selection mechanism or coupled Einstein--Standard-Model limit is asserted.
The older shorter original histories and Version 32's stress problem are
unchanged.

\paragraph{Next source calculation.}
The pointwise test should now be made on
$(\xi,\eta,\gamma,\pi)$ with the identity $P=a^4\pi$ retained.
It is the leading state-only source and its source-compatible reference,
not another moving-principal Green estimate or initializer-transversality
condition.  A further analytic representation must preserve this exact
tracking relation rather than set the trial constraint defect to zero.

\paragraph{Verification.}
The accompanying executable reconstructs three entire harmonic columns
of the original quadratic Poisson coefficient on a declared populated
rational canonical field (324 exact row equalities), regenerates the
witness by the original shifted $P_2$ polarization, and checks the phase
normalizations, skew residual form, regular Taylor division and exact
residual-tracking algebra. The populated field is a coefficient test, not
the V39 root or an original trajectory. Additional finite examples are
explicitly labeled comparison algebra. No higher native coefficient bank,
nonlinear original history, or proof-assistant formalization is claimed.
The parent replay and preservation are recorded separately.

\begingroup
\renewcommand{\refname}{Sources used in Version 61}
\begin{thebibliography}{9}
\small\raggedright
\bibitem{v61:EB45}
\emph{Renewal Exact-Action Einstein Bridge Research Notes}, supplied
cumulative Version 45, file
\nolinkurl{Renewal_Exact_Action_Einstein_Bridge_Research_Notes_V45(3).tex}.
Used at \texttt{v5:linear-multiplier}, \texttt{v21:translation-incidence},
and \texttt{v36:Poisson-orders}; its different-root spectral data are not
transferred. The original $P_2$ compiler is the retained Version-35
standalone implementation, with its file hash recorded in the package.
\bibitem{v61:DBGP}
C.~Di Bartolo, R.~Gambini and J.~Pullin,
\emph{Consistent and mimetic discretizations in general relativity},
J. Math. Phys. \textbf{46} (2005), 032501,
\url{https://arxiv.org/abs/gr-qc/0404052}.
Context for distinguishing constraint-consistent discrete evolution from
other extensions; it supplies none of the source coefficients proved here.
\end{thebibliography}
\endgroup
\addtocontents{toc}{\protect\endgroup}
% END IMMUTABLE VERSION 61 BODY
% APPEND VERSION 62 HERE, BEFORE THE FINAL END-DOCUMENT COMMAND.

\end{document}
